% Two editions from one source. The Kindle edition is selected by defining
% \KindleEdition before the class is read, which the Makefile does on the
% command line:  pdflatex "\def\KindleEdition{}% Two editions from one source. The Kindle edition is selected by defining
% \KindleEdition before the class is read, which the Makefile does on the
% command line:  pdflatex "\def\KindleEdition{}% Two editions from one source. The Kindle edition is selected by defining
% \KindleEdition before the class is read, which the Makefile does on the
% command line:  pdflatex "\def\KindleEdition{}% Two editions from one source. The Kindle edition is selected by defining
% \KindleEdition before the class is read, which the Makefile does on the
% command line:  pdflatex "\def\KindleEdition{}\input{main.tex}"
\ifdefined\KindleEdition
  \documentclass[10pt,oneside]{book}
\else
  \documentclass[11pt,a4paper,twoside,openright]{book}
\fi
\input{preamble}
\renewcommand*{\thisvolume}{derivatives}

\title{\Huge Derivatives \& Partial Derivatives\\[0.4em]
  \large A Training Sheet with Exercises and Worked Solutions}
\author{}
\date{}

\begin{document}
\frontmatter
\maketitle

\chapter*{About this sheet}
\addcontentsline{toc}{chapter}{About this sheet}

This sheet drills the computation of derivatives in one and
several variables, with the applications most used in economics: marginal quantities, elasticities, tangents and linearisations, rates of
substitution, and the second-order conditions behind optimisation.

The structure is the same in every chapter: the \textbf{Core Rules} section
states the rules compactly and says when each one applies; \textbf{Practice
Problems} drills them; \textbf{Solutions} works every problem through in
full, showing the intermediate step rather than only the answer.

Chapters 1--8 cover the one-variable theory: the differentiation rules and
the standard functions, higher derivatives and convexity, Taylor
approximation, implicit differentiation, the classical applications, the
indeterminate forms with L'H\^opital's rule, and a chapter that assembles all
of it into the complete curve discussion. Chapters 9--14 are the
multivariable half, opening with the geometry --- surfaces, level curves and
the sub- and superlevel sets --- before partial derivatives and the gradient
(with the topological vocabulary of domains), the Hessian and definiteness,
the multivariable chain rule and total differential, directional derivatives
and vector-valued functions, and finally the Implicit Function Theorem with
Jacobians and Euler's theorem. Most chapters close with a block of
\emph{additional problems} for a second pass through the material.

Work the problems before reading the solutions. A derivative you have
recognised is worth more than ten you have read.

\tableofcontents

\mainmatter
\include{chapters/ch01}
\include{chapters/ch02}
\include{chapters/ch03}
\include{chapters/ch04}
\include{chapters/ch05}
\include{chapters/ch06}
\include{chapters/ch07}
\include{chapters/ch08}
\include{chapters/ch09}
\include{chapters/ch10}
\include{chapters/ch11}
\include{chapters/ch12}
\include{chapters/ch13}
\include{chapters/ch14}

\end{document}
"
\ifdefined\KindleEdition
  \documentclass[10pt,oneside]{book}
\else
  \documentclass[11pt,a4paper,twoside,openright]{book}
\fi
% Shared preamble for the Derivatives training sheet.
\usepackage[T1]{fontenc}
\usepackage[utf8]{inputenc}
\usepackage{lmodern}
\ifdefined\KindleEdition
  % 6-inch e-ink page (Kindle Paperwhite): the page *is* the screen, so the
  % margins are only what the eye needs, not what a printer needs.
  \usepackage[paperwidth=90mm,paperheight=122mm,
              top=4mm,bottom=6mm,left=4.5mm,right=4.5mm]{geometry}
\else
  \usepackage[a4paper,margin=2.4cm]{geometry}
\fi
\usepackage{amsmath,amssymb,amsthm}
\usepackage{microtype}
\usepackage{enumitem}
\usepackage{xcolor}
\usepackage[most]{tcolorbox}
\usepackage{fancyhdr}
\usepackage[hidelinks]{hyperref}

\definecolor{accent}{HTML}{1F4E79}
\definecolor{soft}{HTML}{F4F7FA}

% --- function graphs ------------------------------------------------------
\usepackage{pgfplots}
\pgfplotsset{compat=1.18}
\usetikzlibrary{arrows.meta}

% One consistent look for every function graph in these books.
\pgfplotsset{
  fnaxis/.style={
    width=5.9cm, height=4.6cm,
    axis lines=middle,
    axis line style={-{Stealth[length=4pt]}, gray!65!black, line width=0.3pt},
    tick style={gray!65!black},
    tick label style={font=\tiny, gray!50!black},
    label style={font=\scriptsize},
    xlabel={$x$}, ylabel={$y$},
    every axis x label/.style={at={(ticklabel* cs:1.0)}, anchor=west},
    every axis y label/.style={at={(ticklabel* cs:1.0)}, anchor=south east,
      xshift=-1pt},
    every axis plot/.append style={accent, very thick, smooth, samples=201},
    clip=true, enlargelimits=false,
    title style={font=\scriptsize, yshift=4pt},
  },
  fnaxis/.append style={/pgfplots/legend style={font=\tiny, draw=none,
    fill=none, at={(0.98,0.03)}, anchor=south east}},
}

% A row of small graphs, evenly spaced with \hfill between them, plus an
% optional caption line: \begin{gallery} <pic>\hfill<pic> \gcap{...} \end{gallery}
\newenvironment{gallery}
  {\par\medskip\noindent\ignorespaces}
  {\par\medskip\noindent}
\newcommand{\gcap}[1]{\par\smallskip{\centering\footnotesize #1\par}}

% --- plane regions, implicit curves, surfaces -----------------------------
\usepgfplotslibrary{fillbetween}
\definecolor{regionfill}{HTML}{9DB4D8}
\pgfplotsset{colormap={house}{color=(soft) color=(accent!45) color=(accent)}}

\pgfplotsset{
  % Squared, gridded axis for plane regions and implicitly given curves.
  geoaxis/.style={
    fnaxis,
    width=5.9cm, height=5.9cm,
    grid=both,
    grid style={gray!22, line width=0.2pt},
    minor tick num=1,
  },
  % Level-curve maps: many thin curves, no single curve emphasised.
  contouraxis/.style={
    geoaxis,
    every axis plot/.append style={accent!75, line width=0.5pt, smooth,
      samples=201},
  },
  % Surfaces: \addplot3[surf, domain=..., y domain=...] {expr};
  surfaxis/.style={
    width=6.0cm, height=5.0cm,
    view={130}{22},
    colormap name=house,
    ticks=none,
    xlabel={$x$}, ylabel={$y$}, zlabel={$f$},
    label style={font=\tiny, gray!50!black},
    title style={font=\scriptsize},
    every axis plot/.append style={samples=25, shader=faceted interp,
      draw=accent!35, line width=0.1pt},
  },
}

% Shorthand for a shaded region between two named paths.
% \region{<path A>}{<path B>}{<soft clip options>}
\newcommand{\region}[3]{%
  \addplot[regionfill, opacity=0.55, forget plot, sharp plot]
    fill between[of=#1 and #2, #3];}

% --- rule boxes -----------------------------------------------------------
\newtcolorbox{rulebox}[1]{%
  breakable, enhanced, colback=soft, colframe=accent,
  boxrule=0.4pt, left=7pt, right=7pt, top=6pt, bottom=6pt,
  fonttitle=\bfseries, coltitle=white, colbacktitle=accent,
  title={#1}}

% --- worked examples ------------------------------------------------------
% One fully narrated solution per chapter, carrying the chapter's central
% technique and the mistakes that technique invites:
%   \begin{worked}{Task} ... \wstep{Step 1.} ... \wtrap{...} ... \end{worked}
\newtcolorbox{worked}[1]{%
  breakable, enhanced, colback=white, colframe=accent!65,
  boxrule=0.4pt, left=7pt, right=7pt, top=6pt, bottom=6pt,
  fonttitle=\bfseries, coltitle=accent, colbacktitle=soft,
  title={#1}}

% Run-in step heading inside a worked example.
\newcommand{\wstep}[1]{\par\smallskip\noindent\textbf{#1}\enspace\ignorespaces}

% A pitfall, set off from the running solution so it can be found again.
\newcommand{\wtrap}[1]{%
  \par\smallskip
  {\leftskip=1.3em\noindent
   \llap{\textcolor{accent}{\rule[0.22ex]{0.9em}{0.6pt}}\hspace{0.4em}}%
   \textbf{Trap.}\enspace\itshape #1\par}%
  \smallskip}

% --- problem / solution lists --------------------------------------------
\newlist{problems}{enumerate}{1}
\setlist[problems]{label=\textbf{\arabic*.}, leftmargin=2.1em,
  itemsep=0.55ex, parsep=0pt, topsep=0.8ex}
\newlist{solutions}{enumerate}{1}
\setlist[solutions]{label=\textbf{\arabic*.}, leftmargin=2.1em,
  itemsep=1.0ex, parsep=0pt, topsep=0.8ex}

% --- running heads --------------------------------------------------------
\pagestyle{fancy}
\fancyhf{}
\fancyhead[LE,RO]{\thepage}
\fancyhead[RE]{\small\itshape\nouppercase{\leftmark}}
\fancyhead[LO]{\small\itshape\nouppercase{\rightmark}}
\setlength{\headheight}{14pt}
\renewcommand{\headrulewidth}{0.3pt}
\renewcommand{\chaptermark}[1]{\markboth{\chaptername\ \thechapter.\ #1}{}}
\renewcommand{\sectionmark}[1]{\markright{\thesection\ #1}}

% --- shorthands used throughout -------------------------------------------
\newcommand{\dd}[2]{\frac{\mathrm{d}#1}{\mathrm{d}#2}}
\newcommand{\pd}[2]{\frac{\partial #1}{\partial #2}}
\newcommand{\pdd}[3]{\frac{\partial^{2} #1}{\partial #2\,\partial #3}}
\newcommand{\R}{\mathbb{R}}
\newcommand{\grad}{\nabla}
\DeclareMathOperator{\sgn}{sgn}
\DeclareMathOperator{\Hess}{H}

% German gloss, used on first mention of a technical term
\newcommand{\de}[1]{\emph{#1}}

% --- cross-references to chapters ------------------------------------------
% Chapter files are shared by the standalone volumes and the compendium, so a
% hard-coded "Chapter 6" is wrong in one of them. Write \chref{derivatives}{6}
% instead: it prints "Chapter 6" inside the same volume, "Volume 5, Chapter 6"
% from another volume, and the running number ("Chapter 34") in the
% compendium. \chrefs{vol}{n} is the abbreviated form (Ch. / Vol., Ch.),
% \chnum{vol}{n} the bare number in the current build, and \volref{vol} names
% the volume ("Volume 5", or "Part V" in the compendium). Each volume's
% main.tex sets \thisvolume; the compendium's sets \incompendiumtrue.
\newif\ifincompendium
\newcommand*{\thisvolume}{}
\makeatletter
\@namedef{chref@vol@algebra}{1}\@namedef{chref@off@algebra}{0}
\@namedef{chref@vol@functions}{2}\@namedef{chref@off@functions}{9}
\@namedef{chref@vol@matrices}{3}\@namedef{chref@off@matrices}{14}
\@namedef{chref@vol@foundations}{4}\@namedef{chref@off@foundations}{24}
\@namedef{chref@vol@derivatives}{5}\@namedef{chref@off@derivatives}{28}
\@namedef{chref@vol@integrals}{6}\@namedef{chref@off@integrals}{42}
\@namedef{chref@vol@optimization}{7}\@namedef{chref@off@optimization}{50}
\@namedef{chref@vol@linprog}{8}\@namedef{chref@off@linprog}{60}
\@namedef{chref@vol@mixed}{9}\@namedef{chref@off@mixed}{68}
\newcommand*{\chref@check}[1]{%
  \@ifundefined{chref@vol@#1}{\PackageError{chref}{Unknown volume `#1'}{}}{}}
\newcommand*{\chnum}[2]{\chref@check{#1}%
  \ifincompendium\the\numexpr\csname chref@off@#1\endcsname+#2\relax
  \else#2\fi}
% \chref@same{vol}{same volume}{other volume}
\newcommand*{\chref@same}[3]{\def\chref@tmp{#1}%
  \ifx\chref@tmp\thisvolume\expandafter\@firstoftwo
  \else\expandafter\@secondoftwo\fi{#2}{#3}}
\newcommand*{\chref}[2]{\chref@check{#1}%
  \ifincompendium Chapter~\chnum{#1}{#2}%
  \else\chref@same{#1}{Chapter~#2}%
    {Volume~\csname chref@vol@#1\endcsname, Chapter~#2}\fi}
\newcommand*{\chrefs}[2]{\chref@check{#1}%
  \ifincompendium Ch.~\chnum{#1}{#2}%
  \else\chref@same{#1}{Ch.~#2}%
    {Vol.~\csname chref@vol@#1\endcsname, Ch.~#2}\fi}
\newcommand*{\volref}[1]{\chref@check{#1}%
  \ifincompendium Part~\@Roman{\csname chref@vol@#1\endcsname}%
  \else Volume~\csname chref@vol@#1\endcsname\fi}
\makeatother

% --- Kindle edition: compact layout for a 6-inch screen -------------------
% Everything below applies only when \KindleEdition is defined on the
% command line; the A4 edition is untouched.
\ifdefined\KindleEdition
  \usepackage{etoolbox}
  \usepackage{adjustbox}
  \usepackage{titlesec}

  % The 8cm measure leaves few break points; let TeX loosen a tight line
  % rather than push it into the margin.
  \setlength{\emergencystretch}{1.5em}

  % No running heads: on a 90mm page the header costs more than it gives.
  % A centred page number in the foot is kept so the contents list means
  % something.
  \pagestyle{plain}
  \fancypagestyle{plain}{\fancyhf{}\fancyfoot[C]{\small\thepage}%
    \renewcommand{\headrulewidth}{0pt}}

  % Compact headings: the book class leaves a third of a page above a
  % chapter title, which is unaffordable here.
  \titleformat{\chapter}[display]
    {\normalfont\Large\bfseries\color{accent}}{\normalsize\chaptertitlename\ \thechapter}{0.4em}{\Large}
  \titlespacing*{\chapter}{0pt}{0pt}{1.2em}
  \titlespacing*{\section}{0pt}{1.4ex plus .3ex}{0.7ex}
  \titlespacing*{\subsection}{0pt}{1.2ex plus .2ex}{0.5ex}

  % Boxes: trim the padding, which was set for a 16cm measure.
  \tcbset{left=3.5pt,right=3.5pt,top=3pt,bottom=3pt,
          fonttitle=\bfseries\small,
          halign title=flush left}

  % Problem and solution lists: reclaim the wide left margin.
  \setlist[problems]{label=\textbf{\arabic*.}, leftmargin=1.5em,
    itemsep=0.45ex, parsep=0pt, topsep=0.6ex}
  \setlist[solutions]{label=\textbf{\arabic*.}, leftmargin=1.5em,
    itemsep=0.8ex, parsep=0pt, topsep=0.6ex}

  % Figures and wide tables shrink to the measure instead of running off it.
  % adjustbox only scales down: anything already narrow is left alone.
  \BeforeBeginEnvironment{tikzpicture}{\begin{adjustbox}{max width=\linewidth}}
  \AfterEndEnvironment{tikzpicture}{\end{adjustbox}}
  \BeforeBeginEnvironment{tabular}{\begin{adjustbox}{max width=\linewidth}}
  \AfterEndEnvironment{tabular}{\end{adjustbox}}

  % Display maths was written for a 16cm measure. On 8cm the long chained
  % equations would run off the screen, so any display wider than the text
  % block is scaled down to fit; narrower ones are left at full size. The
  % whole corpus uses \[ ... \] and no other display environment, so this
  % one redefinition covers every case.
  \newsavebox{\kndbox}
  \renewcommand{\[}{\par\addvspace{0.7ex}%
    \begin{lrbox}{\kndbox}$\displaystyle}
  \renewcommand{\]}{$\end{lrbox}%
    {\centering
      \ifdim\wd\kndbox>\linewidth
        \resizebox{\linewidth}{!}{\usebox{\kndbox}}%
      \else
        \usebox{\kndbox}%
      \fi
     \par}%
    \addvspace{0.7ex}}

  % A narrow measure makes perfect justification impossible; prefer a loose
  % line to a line that runs into the margin.
  \sloppy
  \raggedbottom
\fi

\renewcommand*{\thisvolume}{derivatives}

\title{\Huge Derivatives \& Partial Derivatives\\[0.4em]
  \large A Training Sheet with Exercises and Worked Solutions}
\author{}
\date{}

\begin{document}
\frontmatter
\maketitle

\chapter*{About this sheet}
\addcontentsline{toc}{chapter}{About this sheet}

This sheet drills the computation of derivatives in one and
several variables, with the applications most used in economics: marginal quantities, elasticities, tangents and linearisations, rates of
substitution, and the second-order conditions behind optimisation.

The structure is the same in every chapter: the \textbf{Core Rules} section
states the rules compactly and says when each one applies; \textbf{Practice
Problems} drills them; \textbf{Solutions} works every problem through in
full, showing the intermediate step rather than only the answer.

Chapters 1--8 cover the one-variable theory: the differentiation rules and
the standard functions, higher derivatives and convexity, Taylor
approximation, implicit differentiation, the classical applications, the
indeterminate forms with L'H\^opital's rule, and a chapter that assembles all
of it into the complete curve discussion. Chapters 9--14 are the
multivariable half, opening with the geometry --- surfaces, level curves and
the sub- and superlevel sets --- before partial derivatives and the gradient
(with the topological vocabulary of domains), the Hessian and definiteness,
the multivariable chain rule and total differential, directional derivatives
and vector-valued functions, and finally the Implicit Function Theorem with
Jacobians and Euler's theorem. Most chapters close with a block of
\emph{additional problems} for a second pass through the material.

Work the problems before reading the solutions. A derivative you have
recognised is worth more than ten you have read.

\tableofcontents

\mainmatter
\chapter{Basic Differentiation Rules}

The derivative measures the instantaneous rate of change of a function. In
economics this is the language of \emph{marginal} quantities: marginal cost,
marginal product, marginal revenue. Before specialising to particular
functional forms, master the structural rules below; almost every derivative
computed later is built from a handful of them.

\section{Core Rules}

\begin{rulebox}{Definition (limit form)}
For a function $f:\R\to\R$ the derivative at $x$ is
\[
  f'(x) \;=\; \lim_{h\to 0}\frac{f(x+h)-f(x)}{h},
\]
provided this limit exists. Equivalent notations include $\dd{f}{x}$,
$\dfrac{\mathrm{d}}{\mathrm{d}x}f(x)$, and $Df(x)$.
\end{rulebox}

\begin{rulebox}{Linearity}
For constants $c,\lambda\in\R$ and differentiable $u,v$:
\begin{itemize}[leftmargin=1.6em,itemsep=0.3ex]
  \item \textbf{Constant rule:} $\dfrac{\mathrm{d}}{\mathrm{d}x}(c)=0$
  \item \textbf{Scalar multiple:} $(\lambda u)'=\lambda u'$
  \item \textbf{Sum rule:} $(u+v)'=u'+v'$
\end{itemize}
\end{rulebox}

\begin{rulebox}{Power rule}
For any real exponent $n$ (where the expression is defined):
\[
  \frac{\mathrm{d}}{\mathrm{d}x}x^{n} \;=\; n\,x^{n-1} .
\]
Special cases worth memorising: $(\sqrt{x})'=\dfrac{1}{2\sqrt{x}}$,
$\Bigl(\dfrac{1}{x}\Bigr)'=-\dfrac{1}{x^{2}}$, and $(x)'=1$.
\end{rulebox}

\begin{rulebox}{Product rule}
If $u,v$ are differentiable,
\[
  (uv)' \;=\; u'v+uv' .
\]
\end{rulebox}

\begin{rulebox}{Quotient rule}
If $u,v$ are differentiable and $v\neq 0$,
\[
  \left(\frac{u}{v}\right)' \;=\; \frac{u'v-uv'}{v^{2}} .
\]
\end{rulebox}

\begin{rulebox}{Chain rule}
If $f$ is differentiable at $g(x)$ and $g$ is differentiable at $x$,
\[
  (f\circ g)'(x) \;=\; f'\bigl(g(x)\bigr)\cdot g'(x) .
\]
In Leibniz notation, with $y=f(u)$ and $u=g(x)$, this reads
$\dd{y}{x}=\dd{y}{u}\cdot\dd{u}{x}$.
\end{rulebox}

\begin{rulebox}{Higher-order derivatives (preview)}
The second derivative is $f''(x)=(f'(x))'$, written also
$\dfrac{\mathrm{d}^{2}f}{\mathrm{d}x^{2}}$. In economics $f'>0$ and $f''<0$
characterise an increasing, concave production or utility function. Higher
orders $f^{(n)}$ appear in Taylor expansions and in concavity tests.
\end{rulebox}

\section{Practice Problems}

Differentiate each function with respect to $x$ (or the indicated variable).
Simplify where natural.

\begin{problems}
  \item $f(x)=x^{7}$
  \item $f(x)=(3x+5)^{6}$
  \item $f(x)=\dfrac{x}{x+1}$
  \item $f(x)=(x^{2}+1)(x^{3}-2x)$
  \item $f(x)=x^{5}-4x^{3}+9x-2$
  \item $f(x)=\sqrt{2x^{2}+1}$
  \item $C(q)=100+4q+0.05q^{2}$ \quad (find $C'(q)$, the marginal cost)
  \item $f(x)=\dfrac{2x-3}{x^{2}+1}$
  \item $f(x)=(3x-1)(x^{2}+4x+5)$
  \item $f(x)=\sqrt{x}$
  \item $f(x)=(x^{2}-4x+1)^{4}$
  \item $f(x)=\dfrac{x^{2}+1}{x^{2}-1}$
  \item $f(x)=x\sqrt{x+1}$
  \item $f(x)=6x^{4}+2x^{2}-x$
  \item $f(x)=\bigl((x+1)^{2}+1\bigr)^{3}$ \quad (double chain)
  \item $f(x)=\dfrac{\sqrt{x}}{x+1}$
  \item $f(x)=(x^{2}+2)(\sqrt{x}-3)$
  \item $f(x)=x^{2/3}$
  \item $R(q)=20q-0.1q^{2}$ \quad (find $R'(q)$, the marginal revenue)
  \item $f(x)=\dfrac{1}{(x^{2}+3)^{5}}$
  \item $f(x)=(x^{3}-x)(2x^{2}+7)$
  \item $f(x)=\dfrac{x^{3}-2x}{x+4}$
  \item $f(x)=\dfrac{1}{x^{4}}$
  \item $f(x)=\bigl(\sqrt{x}+2x\bigr)^{4}$
  \item $f(x)=5x^{4}-\tfrac{1}{2}x^{2}+7$
  \item $\Pi(L)=L\cdot(10-L)$ \quad (a profit-like product)
  \item $f(x)=\sqrt[3]{(x^{2}+x+1)^{2}}$
  \item Average cost: $\displaystyle AC(q)=\frac{100+2q^{2}}{q}$
  \item $f(x)=3\sqrt[3]{x}+\dfrac{2}{x^{2}}$
  \item $f(x)=(5x^{2}-1)\Bigl(\dfrac{1}{x}+2\Bigr)$
  \item $f(x)=\dfrac{1-x^{2}}{1+x^{2}}$
  \item $f(x)=\pi x^{3}-\sqrt{2}\,x+e$
  \item $f(x)=\bigl(\sqrt{x}+1\bigr)^{2}$ \quad \emph{(expand first)}
  \item $f(x)=x^{1/2}+x^{-1/2}$
  \item $f(x)=\dfrac{x^{3}-5x+1}{x}$
  \item $f(x)=\tfrac{1}{3}x^{6}+\tfrac{2}{5}x^{5}-4$
  \item $f(x)=(x^{2}-9)^{5}$
  \item $f(x)=\dfrac{x+2}{x-2}$
  \item $f(x)=(4x+3)(x^{2}-x)$
  \item $f(x)=2x^{8}-3x^{5}+x$
  \item $f(x)=\sqrt{4-x^{2}}$
  \item Profit: $\Pi(q)=(20-0.1q)q-(100+4q+0.05q^{2})$ \quad (find
        $\Pi'(q)$)
  \item $f(x)=\dfrac{1}{2x+7}$
  \item $f(x)=(x^{2}+1)^{3}(x-1)$
  \item $f(x)=\dfrac{x^{2}}{1+x}$
  \item $f(x)=7\sqrt{x}-\dfrac{3}{x}$
  \item $f(x)=\bigl(1+x^{3}\bigr)^{10}$
  \item $f(x)=(x+1)^{2}(x-1)^{3}$
  \item Marginal product: $Q(L)=5L^{1/2}$ \quad (find $Q'(L)$)
  \item $f(x)=\dfrac{3x+1}{(x+2)^{2}}$
  \item $f(x)=x^{3}-5x^{2}+7x-4$
  \item $f(x)=\sqrt{x^{2}+9}$
  \item $f(x)=(2x^{3}+1)(x^{2}+x)$
  \item $f(x)=\dfrac{x^{2}-1}{x^{2}+1}$
  \item $f(x)=\Bigl(\dfrac{1}{x}+x\Bigr)^{4}$
  \item $f(x)=-3x^{5}+2x^{-3}$
  \item $f(x)=\dfrac{1}{\sqrt{x^{2}+1}}$
  \item $f(x)=(x-1)(x-2)(x-3)$
  \item Cost: $C(q)=50+\sqrt{q^{2}+4}$ \quad (find $C'(q)$)
  \item $f(x)=\dfrac{x+1}{x^{2}+2}$
  \item $f(x)=(5-2x)^{7}$
  \item $f(x)=x^{2}\sqrt{1+x}$
  \item $f(x)=\dfrac{1}{x^{3}+1}$
  \item $f(x)=(x^{4}+3x)^{6}$
  \item $f(x)=\tfrac{2}{3}x^{3/2}-4x^{-1/2}$
  \item Revenue per unit time: $R(t)=(50+2t)^{3}$ \quad (find $R'(t)$)
  \item $f(x)=\dfrac{\sqrt{x}+1}{\sqrt{x}-1}$
  \item $f(x)=(x^{2}+1)(x^{2}-1)(x+3)$
  \item $f(x)=\sqrt{(x+1)(x+2)}$
  \item $f(x)=\dfrac{x^{2}+2x}{(x+1)^{2}}$
  \item $f(x)=\bigl(3x^{2}-2\sqrt{x}\bigr)^{5}$
\end{problems}

\section{Solutions}

\begin{solutions}
  \item $f'(x)=7x^{6}$ (power rule).

  \item Outer $u^{6}$, inner $u=3x+5$. $f'(x)=6(3x+5)^{5}\cdot 3
        =18(3x+5)^{5}$.

  \item Quotient rule with $u=x$, $v=x+1$, $u'=1$, $v'=1$:
        \[
          f'(x)=\frac{1\cdot(x+1)-x\cdot 1}{(x+1)^{2}}=\frac{1}{(x+1)^{2}} .
        \]

  \item Let $u=x^{2}+1$, $v=x^{3}-2x$. Then $u'=2x$, $v'=3x^{2}-2$.
        \[
          f'(x)=2x(x^{3}-2x)+(x^{2}+1)(3x^{2}-2)=5x^{4}-3x^{2}-2 .
        \]

  \item $f'(x)=5x^{4}-12x^{2}+9$.

  \item $f(x)=(2x^{2}+1)^{1/2}$, inner $u=2x^{2}+1$, $u'=4x$.
        \[
          f'(x)=\tfrac12(2x^{2}+1)^{-1/2}\cdot 4x=\frac{2x}{\sqrt{2x^{2}+1}} .
        \]

  \item $C'(q)=4+0.1q$. Marginal cost is linear in output.

  \item $u=2x-3$, $v=x^{2}+1$, $u'=2$, $v'=2x$.
        \[
          f'(x)=\frac{2(x^{2}+1)-(2x-3)(2x)}{(x^{2}+1)^{2}}
               =\frac{-2x^{2}+6x+2}{(x^{2}+1)^{2}} .
        \]

  \item $u=3x-1$, $v=x^{2}+4x+5$, $u'=3$, $v'=2x+4$.
        \[
          f'(x)=3(x^{2}+4x+5)+(3x-1)(2x+4)=9x^{2}+22x+11 .
        \]

  \item Write $\sqrt{x}=x^{1/2}$, so $f'(x)=\tfrac12x^{-1/2}
        =\dfrac{1}{2\sqrt{x}}$.

  \item Inner $u=x^{2}-4x+1$, $u'=2x-4$. $f'(x)=4(x^{2}-4x+1)^{3}(2x-4)
        =8(x-2)(x^{2}-4x+1)^{3}$.

  \item $u=x^{2}+1$, $v=x^{2}-1$, $u'=2x$, $v'=2x$.
        \[
          f'(x)=\frac{2x(x^{2}-1)-(x^{2}+1)(2x)}{(x^{2}-1)^{2}}
               =\frac{-4x}{(x^{2}-1)^{2}} .
        \]

  \item $u=x$, $v=(x+1)^{1/2}$, $u'=1$, $v'=\tfrac12(x+1)^{-1/2}$.
        \[
          f'(x)=\sqrt{x+1}+\frac{x}{2\sqrt{x+1}}
               =\frac{2(x+1)+x}{2\sqrt{x+1}}=\frac{3x+2}{2\sqrt{x+1}} .
        \]

  \item $f'(x)=24x^{3}+4x-1$.

  \item Two layers. Let $g(x)=(x+1)^{2}+1$, $g'(x)=2(x+1)$. Then
        \[
          f'(x)=3\,g(x)^{2}\cdot g'(x)=6(x+1)\bigl((x+1)^{2}+1\bigr)^{2} .
        \]

  \item $u=x^{1/2}$, $v=x+1$, $u'=\tfrac12x^{-1/2}$, $v'=1$.
        \[
          f'(x)=\frac{\frac{1}{2\sqrt{x}}(x+1)-\sqrt{x}}{(x+1)^{2}}
               =\frac{(x+1)-2x}{2\sqrt{x}(x+1)^{2}}
               =\frac{1-x}{2\sqrt{x}(x+1)^{2}} .
        \]

  \item $u=x^{2}+2$, $v=x^{1/2}-3$, $u'=2x$, $v'=\tfrac12x^{-1/2}$.
        \[
          f'(x)=2x(\sqrt{x}-3)+(x^{2}+2)\cdot\frac{1}{2\sqrt{x}}
               =2x\sqrt{x}-6x+\frac{x^{2}+2}{2\sqrt{x}} .
        \]

  \item $f'(x)=\tfrac23x^{-1/3}=\dfrac{2}{3\sqrt[3]{x}}$.

  \item $R'(q)=20-0.2q$. Marginal revenue falls with $q$.

  \item $f(x)=(x^{2}+3)^{-5}$, inner $u'=2x$.
        \[
          f'(x)=-5(x^{2}+3)^{-6}\cdot 2x=-\frac{10x}{(x^{2}+3)^{6}} .
        \]

  \item $u=x^{3}-x$, $v=2x^{2}+7$, $u'=3x^{2}-1$, $v'=4x$.
        \[
          f'(x)=(3x^{2}-1)(2x^{2}+7)+(x^{3}-x)(4x)=10x^{4}+15x^{2}-7 .
        \]

  \item $u=x^{3}-2x$, $v=x+4$, $u'=3x^{2}-2$, $v'=1$.
        \[
          f'(x)=\frac{(3x^{2}-2)(x+4)-(x^{3}-2x)}{(x+4)^{2}}
               =\frac{2x^{3}+12x^{2}-8}{(x+4)^{2}} .
        \]

  \item Write $f(x)=x^{-4}$, so $f'(x)=-4x^{-5}=-\dfrac{4}{x^{5}}$.

  \item Inner $u=\sqrt{x}+2x$, $u'=\dfrac{1}{2\sqrt{x}}+2$.
        \[
          f'(x)=4\bigl(\sqrt{x}+2x\bigr)^{3}\left(\frac{1}{2\sqrt{x}}+2\right) .
        \]

  \item $f'(x)=20x^{3}-x$.

  \item $\Pi'(L)=1\cdot(10-L)+L\cdot(-1)=10-2L$.

  \item $f(x)=(x^{2}+x+1)^{2/3}$, inner $u'=2x+1$.
        \[
          f'(x)=\tfrac23(x^{2}+x+1)^{-1/3}(2x+1)
               =\frac{2(2x+1)}{3\sqrt[3]{x^{2}+x+1}} .
        \]

  \item Easier: $AC(q)=100q^{-1}+2q$, so $AC'(q)=-100q^{-2}+2
        =2-\dfrac{100}{q^{2}}$.

  \item $f(x)=3x^{1/3}+2x^{-2}$, hence $f'(x)=x^{-2/3}-4x^{-3}
        =\dfrac{1}{x^{2/3}}-\dfrac{4}{x^{3}}$.

  \item $u=5x^{2}-1$, $v=x^{-1}+2$, $u'=10x$, $v'=-x^{-2}$.
        \[
          f'(x)=10x(x^{-1}+2)+(5x^{2}-1)(-x^{-2})
               =10+20x-5+x^{-2}=5+20x+\frac{1}{x^{2}} .
        \]

  \item $u=1-x^{2}$, $v=1+x^{2}$, $u'=-2x$, $v'=2x$.
        \[
          f'(x)=\frac{-2x(1+x^{2})-(1-x^{2})(2x)}{(1+x^{2})^{2}}
               =\frac{-4x}{(1+x^{2})^{2}} .
        \]

  \item $f'(x)=3\pi x^{2}-\sqrt{2}$. ($e$ is a constant.)

  \item Expand: $f(x)=x+2\sqrt{x}+1$, so $f'(x)=1+\dfrac{1}{\sqrt{x}}$.

  \item $f'(x)=\tfrac12x^{-1/2}-\tfrac12x^{-3/2}
        =\dfrac{1}{2\sqrt{x}}-\dfrac{1}{2x^{3/2}}$.

  \item Simplify first: $f(x)=x^{2}-5+x^{-1}$, so $f'(x)=2x-x^{-2}
        =2x-\dfrac{1}{x^{2}}$.

  \item $f'(x)=2x^{5}+2x^{4}$.

  \item Chain rule with $u=x^{2}-9$, $u'=2x$:
        \[
          f'(x)=5(x^{2}-9)^{4}\cdot 2x=10x(x^{2}-9)^{4} .
        \]

  \item $u=x+2$, $v=x-2$, $u'=1$, $v'=1$.
        \[
          f'(x)=\frac{1\cdot(x-2)-(x+2)\cdot 1}{(x-2)^{2}}
               =\frac{-4}{(x-2)^{2}} .
        \]

  \item Product rule with $u=4x+3$, $v=x^{2}-x$, $u'=4$, $v'=2x-1$.
        \[
          f'(x)=4(x^{2}-x)+(4x+3)(2x-1)=12x^{2}-2x-3 .
        \]

  \item $f'(x)=16x^{7}-15x^{4}+1$.

  \item $f(x)=(4-x^{2})^{1/2}$, inner $u'=-2x$.
        \[
          f'(x)=\tfrac12(4-x^{2})^{-1/2}\cdot(-2x)=-\frac{x}{\sqrt{4-x^{2}}} .
        \]

  \item Expand: $\Pi(q)=20q-0.1q^{2}-100-4q-0.05q^{2}
        =-100+16q-0.15q^{2}$.
        \[
          \Pi'(q)=16-0.3q .
        \]
        Marginal profit is zero at $q^{*}=16/0.3\approx 53.33$.

  \item Write $f(x)=(2x+7)^{-1}$. Chain rule: $f'(x)=-(2x+7)^{-2}\cdot 2
        =-\dfrac{2}{(2x+7)^{2}}$.

  \item Product rule with $u=(x^{2}+1)^{3}$, $v=x-1$. $u'=3(x^{2}+1)^{2}
        \cdot 2x=6x(x^{2}+1)^{2}$, $v'=1$.
        \[
          f'(x)=6x(x^{2}+1)^{2}(x-1)+(x^{2}+1)^{3}
               =(x^{2}+1)^{2}\bigl[6x(x-1)+(x^{2}+1)\bigr]
               =(x^{2}+1)^{2}(7x^{2}-6x+1) .
        \]

  \item Quotient rule with $u=x^{2}$, $v=1+x$, $u'=2x$, $v'=1$.
        \[
          f'(x)=\frac{2x(1+x)-x^{2}}{(1+x)^{2}}=\frac{x^{2}+2x}{(1+x)^{2}}
               =\frac{x(x+2)}{(1+x)^{2}} .
        \]

  \item $f(x)=7x^{1/2}-3x^{-1}$, so $f'(x)=\tfrac72x^{-1/2}+3x^{-2}
        =\dfrac{7}{2\sqrt{x}}+\dfrac{3}{x^{2}}$.

  \item Inner $u=1+x^{3}$, $u'=3x^{2}$.
        \[
          f'(x)=10(1+x^{3})^{9}\cdot 3x^{2}=30x^{2}(1+x^{3})^{9} .
        \]

  \item Product rule with $u=(x+1)^{2}$, $v=(x-1)^{3}$, $u'=2(x+1)$,
        $v'=3(x-1)^{2}$.
        \[
          f'(x)=2(x+1)(x-1)^{3}+3(x+1)^{2}(x-1)^{2}
               =(x+1)(x-1)^{2}\bigl[2(x-1)+3(x+1)\bigr]
               =(x+1)(x-1)^{2}(5x+1) .
        \]

  \item $Q(L)=5L^{1/2}$, so $Q'(L)=\tfrac52L^{-1/2}=\dfrac{5}{2\sqrt{L}}$.
        Marginal product is positive but diminishing.

  \item Quotient rule with $u=3x+1$, $v=(x+2)^{2}$, $u'=3$,
        $v'=2(x+2)$.
        \[
          f'(x)=\frac{3(x+2)^{2}-(3x+1)\cdot 2(x+2)}{(x+2)^{4}}
               =\frac{3(x+2)-2(3x+1)}{(x+2)^{3}}=\frac{-3x+4}{(x+2)^{3}} .
        \]

  \item $f'(x)=3x^{2}-10x+7$.

  \item $f(x)=(x^{2}+9)^{1/2}$, $u'=2x$.
        \[
          f'(x)=\tfrac12(x^{2}+9)^{-1/2}\cdot 2x=\frac{x}{\sqrt{x^{2}+9}} .
        \]

  \item Product rule with $u=2x^{3}+1$, $v=x^{2}+x$, $u'=6x^{2}$,
        $v'=2x+1$.
        \[
          f'(x)=6x^{2}(x^{2}+x)+(2x^{3}+1)(2x+1)=10x^{4}+8x^{3}+2x+1 .
        \]

  \item Quotient rule with $u=x^{2}-1$, $v=x^{2}+1$, $u'=2x$, $v'=2x$.
        \[
          f'(x)=\frac{2x(x^{2}+1)-(x^{2}-1)(2x)}{(x^{2}+1)^{2}}
               =\frac{4x}{(x^{2}+1)^{2}} .
        \]

  \item Inner $u=x^{-1}+x$, $u'=-x^{-2}+1$.
        \[
          f'(x)=4\Bigl(\frac{1}{x}+x\Bigr)^{3}\left(1-\frac{1}{x^{2}}\right) .
        \]

  \item $f'(x)=-15x^{4}-6x^{-4}=-15x^{4}-\dfrac{6}{x^{4}}$.

  \item $f(x)=(x^{2}+1)^{-1/2}$, inner $u'=2x$.
        \[
          f'(x)=-\tfrac12(x^{2}+1)^{-3/2}\cdot 2x=-\frac{x}{(x^{2}+1)^{3/2}} .
        \]

  \item Expand or use the product rule iteratively. Quickest:
        $(x-1)(x-2)(x-3)=x^{3}-6x^{2}+11x-6$, so $f'(x)=3x^{2}-12x+11$.

  \item $C(q)=50+(q^{2}+4)^{1/2}$, $u'=2q$.
        \[
          C'(q)=\tfrac12(q^{2}+4)^{-1/2}\cdot 2q=\frac{q}{\sqrt{q^{2}+4}} .
        \]

  \item Quotient rule with $u=x+1$, $v=x^{2}+2$, $u'=1$, $v'=2x$.
        \[
          f'(x)=\frac{1\cdot(x^{2}+2)-(x+1)(2x)}{(x^{2}+2)^{2}}
               =\frac{-x^{2}-2x+2}{(x^{2}+2)^{2}} .
        \]

  \item Inner $u=5-2x$, $u'=-2$.
        \[
          f'(x)=7(5-2x)^{6}\cdot(-2)=-14(5-2x)^{6} .
        \]

  \item Product rule with $u=x^{2}$, $v=(1+x)^{1/2}$, $u'=2x$,
        $v'=\tfrac12(1+x)^{-1/2}$.
        \[
          f'(x)=2x\sqrt{1+x}+\frac{x^{2}}{2\sqrt{1+x}}
               =\frac{4x(1+x)+x^{2}}{2\sqrt{1+x}}=\frac{5x^{2}+4x}{2\sqrt{1+x}} .
        \]

  \item $f(x)=(x^{3}+1)^{-1}$, inner $u'=3x^{2}$.
        \[
          f'(x)=-(x^{3}+1)^{-2}\cdot 3x^{2}=-\frac{3x^{2}}{(x^{3}+1)^{2}} .
        \]

  \item Inner $u=x^{4}+3x$, $u'=4x^{3}+3$.
        \[
          f'(x)=6(x^{4}+3x)^{5}(4x^{3}+3) .
        \]

  \item $f'(x)=\tfrac23\cdot\tfrac32x^{1/2}-4\cdot(-\tfrac12)x^{-3/2}
        =x^{1/2}+2x^{-3/2}=\sqrt{x}+\dfrac{2}{x^{3/2}}$.

  \item Inner $u=50+2t$, $u'=2$.
        \[
          R'(t)=3(50+2t)^{2}\cdot 2=6(50+2t)^{2} .
        \]

  \item Quotient rule with $u=\sqrt{x}+1$, $v=\sqrt{x}-1$,
        $u'=v'=\dfrac{1}{2\sqrt{x}}$.
        \[
          f'(x)=\frac{\frac{1}{2\sqrt{x}}(\sqrt{x}-1)
               -(\sqrt{x}+1)\cdot\frac{1}{2\sqrt{x}}}{(\sqrt{x}-1)^{2}}
               =\frac{-1/\sqrt{x}}{(\sqrt{x}-1)^{2}}
               =-\frac{1}{\sqrt{x}(\sqrt{x}-1)^{2}} .
        \]

  \item Group: note $(x^{2}+1)(x^{2}-1)=x^{4}-1$, so
        $f(x)=(x^{4}-1)(x+3)=x^{5}+3x^{4}-x-3$. Hence
        $f'(x)=5x^{4}+12x^{3}-1$.

  \item $f(x)=\bigl((x+1)(x+2)\bigr)^{1/2}=(x^{2}+3x+2)^{1/2}$, inner
        $u'=2x+3$.
        \[
          f'(x)=\tfrac12(x^{2}+3x+2)^{-1/2}(2x+3)
               =\frac{2x+3}{2\sqrt{(x+1)(x+2)}} .
        \]

  \item Simplify first: $x^{2}+2x=x(x+2)$ and $(x+1)^{2}=x^{2}+2x+1$.
        Quotient rule with $u=x^{2}+2x$, $v=(x+1)^{2}$, $u'=2x+2=2(x+1)$,
        $v'=2(x+1)$.
        \[
          f'(x)=\frac{2(x+1)(x+1)^{2}-(x^{2}+2x)\cdot 2(x+1)}{(x+1)^{4}}
               =\frac{2(x+1)^{2}-2(x^{2}+2x)}{(x+1)^{3}}=\frac{2}{(x+1)^{3}} .
        \]

  \item Inner $u=3x^{2}-2x^{1/2}$, $u'=6x-x^{-1/2}$.
        \[
          f'(x)=5\bigl(3x^{2}-2\sqrt{x}\bigr)^{4}\left(6x-\frac{1}{\sqrt{x}}\right) .
        \]
\end{solutions}

\section{Additional Problems}

\begin{problems}
  \item $f(x)=4x^{5}-7x^{3}+2x-9$
  \item $f(x)=(2x+3)^{8}$
  \item $f(x)=\dfrac{x-1}{x+1}$
  \item $f(x)=(x^{3}-2)(x^{2}+5)$
  \item $f(x)=\sqrt{3x^{2}+4}$
  \item $f(x)=\dfrac{x^{2}}{x^{3}+1}$
  \item Cost: $C(q)=200+6q+0.02q^{3}$ \quad (find $C'(q)$)
  \item $f(x)=(x^{2}+x)^{5}$
  \item $f(x)=\dfrac{1}{(2x-5)^{3}}$
  \item $f(x)=x\sqrt{x^{2}+1}$
  \item $f(x)=\bigl((x^{2}+1)^{2}-3\bigr)^{4}$ \quad (double chain)
  \item Marginal product: $Q(L)=8L^{2/3}$ \quad (find $Q'(L)$)
  \item $f(x)=\dfrac{x^{2}+3x}{(x-1)^{2}}$
  \item $f(x)=\sqrt[4]{x^{3}}$
  \item Profit: $\Pi(q)=(30-0.2q)\,q-(50+2q+0.1q^{2})$ \quad (find
        $\Pi'(q)$)
\end{problems}

\section{Additional Solutions}

\begin{solutions}
  \item $f'(x)=20x^{4}-21x^{2}+2$.

  \item Chain rule: $f'(x)=8(2x+3)^{7}\cdot 2=16(2x+3)^{7}$.

  \item Quotient rule with $u=x-1$, $v=x+1$, $u'=1$, $v'=1$:
        \[
          f'(x)=\frac{(x+1)-(x-1)}{(x+1)^{2}}=\frac{2}{(x+1)^{2}} .
        \]

  \item Product rule: $u=x^{3}-2$, $v=x^{2}+5$, $u'=3x^{2}$, $v'=2x$.
        \[
          f'(x)=3x^{2}(x^{2}+5)+(x^{3}-2)(2x)=5x^{4}+15x^{2}-4x .
        \]

  \item $f(x)=(3x^{2}+4)^{1/2}$, inner $u'=6x$.
        \[
          f'(x)=\tfrac12(3x^{2}+4)^{-1/2}\cdot 6x=\frac{3x}{\sqrt{3x^{2}+4}} .
        \]

  \item Quotient rule with $u=x^{2}$, $v=x^{3}+1$, $u'=2x$, $v'=3x^{2}$.
        \[
          f'(x)=\frac{2x(x^{3}+1)-x^{2}\cdot 3x^{2}}{(x^{3}+1)^{2}}
               =\frac{-x^{4}+2x}{(x^{3}+1)^{2}}
               =\frac{x(2-x^{3})}{(x^{3}+1)^{2}} .
        \]

  \item $C'(q)=6+0.06q^{2}$. Marginal cost grows quadratically with output.

  \item Chain rule: inner $u=x^{2}+x$, $u'=2x+1$.
        $f'(x)=5(x^{2}+x)^{4}(2x+1)$.

  \item $f(x)=(2x-5)^{-3}$. Chain rule: $f'(x)=-3(2x-5)^{-4}\cdot 2
        =-\dfrac{6}{(2x-5)^{4}}$.

  \item Product rule with $u=x$, $v=(x^{2}+1)^{1/2}$, $u'=1$,
        $v'=x/(x^{2}+1)^{1/2}$.
        \[
          f'(x)=\sqrt{x^{2}+1}+\frac{x^{2}}{\sqrt{x^{2}+1}}
               =\frac{2x^{2}+1}{\sqrt{x^{2}+1}} .
        \]

  \item Two layers. Let $g(x)=(x^{2}+1)^{2}-3$,
        $g'(x)=2(x^{2}+1)\cdot 2x=4x(x^{2}+1)$.
        \[
          f'(x)=4\,g(x)^{3}\cdot g'(x)
               =16x(x^{2}+1)\bigl((x^{2}+1)^{2}-3\bigr)^{3} .
        \]

  \item $Q'(L)=8\cdot\tfrac23L^{-1/3}=\dfrac{16}{3L^{1/3}}$. Marginal
        product is positive but diminishing.

  \item Quotient rule with $u=x^{2}+3x$, $v=(x-1)^{2}$, $u'=2x+3$,
        $v'=2(x-1)$.
        \[
          f'(x)=\frac{(2x+3)(x-1)^{2}-(x^{2}+3x)\cdot 2(x-1)}{(x-1)^{4}}
               =\frac{(2x+3)(x-1)-2(x^{2}+3x)}{(x-1)^{3}}
               =\frac{-5x-3}{(x-1)^{3}} .
        \]

  \item $f(x)=x^{3/4}$, so $f'(x)=\tfrac34x^{-1/4}=\dfrac{3}{4\sqrt[4]{x}}$.

  \item Expand: $\Pi(q)=30q-0.2q^{2}-50-2q-0.1q^{2}
        =-50+28q-0.3q^{2}$.
        \[
          \Pi'(q)=28-0.6q .
        \]
        Marginal profit is zero at $q^{*}=28/0.6\approx 46.67$.
\end{solutions}

\chapter{Derivatives of Standard Functions}

The basic rules of Section 1 reduce every derivative to a finite list of
\emph{building blocks}: exponentials, logarithms, trigonometric, inverse
trigonometric and hyperbolic functions. Memorise the table below. In
economics, exponentials govern continuous growth and discounting, while
logarithms encode percentage changes and elasticities.

\section{Core Rules}

\begin{rulebox}{Exponential functions}
\begin{itemize}[leftmargin=1.6em,itemsep=0.3ex]
  \item $\dfrac{\mathrm{d}}{\mathrm{d}x}\mathrm{e}^{x}=\mathrm{e}^{x}$
  \item $\dfrac{\mathrm{d}}{\mathrm{d}x}\mathrm{e}^{ax}=a\,\mathrm{e}^{ax}$
        \quad ($a$ constant; chain rule)
  \item $\dfrac{\mathrm{d}}{\mathrm{d}x}a^{x}=a^{x}\ln a$ \quad
        ($a>0$, $a\neq 1$)
\end{itemize}
\end{rulebox}

\begin{rulebox}{Logarithmic functions}
\begin{itemize}[leftmargin=1.6em,itemsep=0.3ex]
  \item $\dfrac{\mathrm{d}}{\mathrm{d}x}\ln x=\dfrac{1}{x}$ \quad ($x>0$)
  \item $\dfrac{\mathrm{d}}{\mathrm{d}x}\ln|x|=\dfrac{1}{x}$ \quad
        ($x\neq 0$)
  \item $\dfrac{\mathrm{d}}{\mathrm{d}x}\log_{a}x=\dfrac{1}{x\ln a}$ \quad
        ($a>0$, $a\neq 1$)
\end{itemize}
\end{rulebox}

\begin{rulebox}{Trigonometric functions}
Angles in radians.
\begin{itemize}[leftmargin=1.6em,itemsep=0.3ex]
  \item $(\sin x)'=\cos x$
  \item $(\cos x)'=-\sin x$
  \item $(\tan x)'=\sec^{2}x=\dfrac{1}{\cos^{2}x}$
  \item $(\cot x)'=-\csc^{2}x=-\dfrac{1}{\sin^{2}x}$
  \item $(\sec x)'=\sec x\tan x$
  \item $(\csc x)'=-\csc x\cot x$
\end{itemize}
\end{rulebox}

\begin{rulebox}{Inverse trigonometric functions}
\begin{itemize}[leftmargin=1.6em,itemsep=0.3ex]
  \item $(\arcsin x)'=\dfrac{1}{\sqrt{1-x^{2}}}$, \quad $|x|<1$
  \item $(\arccos x)'=-\dfrac{1}{\sqrt{1-x^{2}}}$, \quad $|x|<1$
  \item $(\arctan x)'=\dfrac{1}{1+x^{2}}$
\end{itemize}
\end{rulebox}

\begin{rulebox}{Hyperbolic functions}
Recall $\sinh x=\tfrac12(\mathrm{e}^{x}-\mathrm{e}^{-x})$,
$\cosh x=\tfrac12(\mathrm{e}^{x}+\mathrm{e}^{-x})$.
\begin{itemize}[leftmargin=1.6em,itemsep=0.3ex]
  \item $(\sinh x)'=\cosh x$
  \item $(\cosh x)'=\sinh x$ \quad (note: \emph{no} sign change)
  \item $(\tanh x)'=\operatorname{sech}^{2}x=1-\tanh^{2}x$
\end{itemize}
\end{rulebox}

\begin{rulebox}{Summary table}
\[
\begin{array}{lll}
\hline
f(x) & f'(x) & \text{Domain note} \\
\hline
\mathrm{e}^{x} & \mathrm{e}^{x} & \R \\
a^{x} & a^{x}\ln a & a>0,\ a\neq 1 \\
\ln x & 1/x & x>0 \\
\log_{a}x & 1/(x\ln a) & x>0 \\
\sin x & \cos x & \R \\
\cos x & -\sin x & \R \\
\tan x & \sec^{2}x & x\neq \tfrac{\pi}{2}+k\pi \\
\cot x & -\csc^{2}x & x\neq k\pi \\
\sec x & \sec x\tan x & x\neq \tfrac{\pi}{2}+k\pi \\
\csc x & -\csc x\cot x & x\neq k\pi \\
\arcsin x & 1/\sqrt{1-x^{2}} & |x|<1 \\
\arccos x & -1/\sqrt{1-x^{2}} & |x|<1 \\
\arctan x & 1/(1+x^{2}) & \R \\
\sinh x & \cosh x & \R \\
\cosh x & \sinh x & \R \\
\tanh x & \operatorname{sech}^{2}x & \R \\
\hline
\end{array}
\]
Combine each entry with the chain rule whenever the argument is a
non-trivial function of $x$.
\end{rulebox}

\section{Practice Problems}

Differentiate each of the following with respect to $x$ (or the indicated
variable).

\begin{problems}
  \item $f(x)=\mathrm{e}^{3x}$
  \item $f(x)=\ln(5x)$
  \item $f(x)=\sin(3x)$
  \item $f(x)=\arcsin(2x)$
  \item $f(x)=\sinh(2x)$
  \item $f(x)=\mathrm{e}^{-x^{2}}$
  \item $f(x)=\cos(x^{2})$
  \item $f(x)=\ln(x^{2}+1)$
  \item $f(x)=x\,\mathrm{e}^{x}$
  \item $f(x)=\arctan(x^{2})$
  \item $f(x)=\tan(2x+1)$
  \item Continuous compounding: $V(t)=V_{0}\,\mathrm{e}^{rt}$ \quad
        (constants $V_{0},r$)
  \item $f(x)=x\ln x$
  \item $f(x)=\cosh(x^{2})$
  \item $f(x)=\mathrm{e}^{\sqrt{x}}$
  \item $f(x)=\arccos(\sqrt{x})$
  \item $f(x)=x\sin x$
  \item $f(x)=\ln(\sqrt{x})$
  \item $f(x)=\sin^{2}x$ \quad (i.e.\ $(\sin x)^{2}$)
  \item $f(x)=2^{x}$
  \item $f(x)=\mathrm{e}^{x^{2}+1}$
  \item $f(x)=(\ln x)^{2}$
  \item $f(x)=\dfrac{\sin x}{x}$
  \item $f(x)=x\arctan x$
  \item $f(x)=\tanh(3x)$
  \item $f(x)=\log_{10}(x)$
  \item $f(x)=\sin(\cos x)$
  \item $f(x)=\mathrm{e}^{\sin x}$
  \item $f(x)=\arctan\Bigl(\dfrac{1}{x}\Bigr)$, \quad $x>0$
  \item $f(x)=\ln\Bigl(\dfrac{x}{x+1}\Bigr)$
  \item $f(x)=\sec(3x)$
  \item Log-utility: $U(c)=\ln(c)$ \quad (find marginal utility)
  \item $f(x)=\ln(\cos x)$
  \item $f(x)=\tan^{3}x$
  \item Logistic: $P(t)=\dfrac{1}{1+\mathrm{e}^{-t}}$
  \item $f(x)=\dfrac{\mathrm{e}^{x}}{1+\mathrm{e}^{x}}$
  \item $f(x)=\mathrm{e}^{x}\sin(2x)$
  \item Cobb-Douglas log: $f(K)=\ln(K^{\alpha})$ \quad ($\alpha$ constant)
  \item $f(x)=x^{x}$ \quad (use logarithmic differentiation)
  \item $f(x)=\mathrm{e}^{2x}\cos(3x)$
  \item $f(x)=\ln\bigl(x+\sqrt{x^{2}+1}\bigr)$
  \item $f(x)=\sin(x)\cos(x)$
  \item Discount factor: $D(t)=\mathrm{e}^{-rt}$ \quad (find $D'(t)$ and
        interpret)
  \item $f(x)=\arctan(\mathrm{e}^{x})$
  \item $f(x)=\tanh(x^{2}+1)$
  \item $f(x)=\dfrac{\ln x}{x}$
  \item $f(x)=\mathrm{e}^{x}\ln x$
  \item $f(x)=\sqrt{\sin x}$
  \item Exponential growth: $N(t)=1000\cdot\mathrm{e}^{0.04t}$ \quad (find
        $N'(t)$ and the growth rate $N'/N$)
  \item $f(x)=5^{x^{2}}$
  \item $f(x)=\arcsin(x^{2})$
  \item $f(x)=\cot(x)+\csc(x)$
  \item $f(x)=\sinh(x)\cosh(x)$
  \item Log of a product: $f(x)=\ln\bigl(x\sqrt{x^{2}+1}\bigr)$
  \item $f(x)=\cos^{4}(x)$
  \item $f(x)=\mathrm{e}^{x^{3}}$
  \item Iso-elastic utility: $\displaystyle U(c)=\frac{c^{1-\sigma}-1}{1-\sigma}$,
        $\sigma\neq 1$ \quad (find $U'(c)$)
  \item $f(x)=\arctan(\sqrt{x})$
  \item $f(x)=\tan(x)\ln(x)$
  \item $f(x)=\dfrac{1}{\ln x}$
  \item $f(x)=\mathrm{e}^{ax}\sin(bx)$ \quad (constants $a,b$)
  \item $f(x)=\log_{2}(x^{2}+1)$
  \item $f(x)=\cos(2x)\mathrm{e}^{-x}$
  \item $f(x)=\operatorname{arccot}(x)\cdot\arctan(x)$ \quad [use
        $\operatorname{arccot}(x)=\tfrac{\pi}{2}-\arctan(x)$]
  \item $f(x)=\ln(\ln(x))$, \quad $x>1$
  \item $f(x)=x^{2}\mathrm{e}^{-x}$
  \item Per-capita growth: $g(t)=\ln(N(t))$ where $N(t)=N_{0}\mathrm{e}^{rt}$
  \item $f(x)=\sin(\mathrm{e}^{x})$
  \item $f(x)=(1+\mathrm{e}^{x})^{3}$
  \item $f(x)=x^{\sin x}$ \quad (use logarithmic differentiation)
  \item $f(x)=\arctan(\tan x)$ \quad (restrict $x\in(-\tfrac{\pi}{2},
        \tfrac{\pi}{2})$)
\end{problems}

\section{Solutions}

\begin{solutions}
  \item Inner $u=3x$, $u'=3$. $f'(x)=\mathrm{e}^{3x}\cdot 3=3\mathrm{e}^{3x}$.

  \item Either expand $\ln(5x)=\ln 5+\ln x$, or chain:
        $f'(x)=\dfrac{1}{5x}\cdot 5=\dfrac{1}{x}$.

  \item Inner $u=3x$, $u'=3$. $f'(x)=\cos(3x)\cdot 3=3\cos(3x)$.

  \item Inner $u=2x$, $u'=2$.
        \[
          f'(x)=\frac{1}{\sqrt{1-(2x)^{2}}}\cdot 2=\frac{2}{\sqrt{1-4x^{2}}} .
        \]

  \item Inner $u=2x$, $u'=2$. $f'(x)=\cosh(2x)\cdot 2=2\cosh(2x)$.

  \item Inner $u=-x^{2}$, $u'=-2x$.
        $f'(x)=\mathrm{e}^{-x^{2}}\cdot(-2x)=-2x\,\mathrm{e}^{-x^{2}}$.

  \item Inner $u=x^{2}$, $u'=2x$.
        $f'(x)=-\sin(x^{2})\cdot 2x=-2x\sin(x^{2})$.

  \item Inner $u=x^{2}+1$, $u'=2x$. $f'(x)=\dfrac{2x}{x^{2}+1}$.

  \item Product rule: $f'(x)=1\cdot\mathrm{e}^{x}+x\cdot\mathrm{e}^{x}
        =(1+x)\,\mathrm{e}^{x}$.

  \item Inner $u=x^{2}$, $u'=2x$. $f'(x)=\dfrac{2x}{1+x^{4}}$.

  \item Inner $u=2x+1$, $u'=2$. $f'(x)=2\sec^{2}(2x+1)$.

  \item $V'(t)=V_{0}\cdot r\,\mathrm{e}^{rt}=r\,V(t)$. The growth rate is
        constant: $V'/V=r$.

  \item Product rule: $f'(x)=1\cdot\ln x+x\cdot\dfrac{1}{x}=\ln x+1$.

  \item Inner $u=x^{2}$, $u'=2x$. $f'(x)=\sinh(x^{2})\cdot 2x
        =2x\sinh(x^{2})$.

  \item Inner $u=\sqrt{x}$, $u'=\dfrac{1}{2\sqrt{x}}$.
        \[
          f'(x)=\mathrm{e}^{\sqrt{x}}\cdot\frac{1}{2\sqrt{x}}
               =\frac{\mathrm{e}^{\sqrt{x}}}{2\sqrt{x}} .
        \]

  \item Inner $u=\sqrt{x}=x^{1/2}$, $u'=\dfrac{1}{2\sqrt{x}}$.
        \[
          f'(x)=-\frac{1}{\sqrt{1-x}}\cdot\frac{1}{2\sqrt{x}}
               =-\frac{1}{2\sqrt{x-x^{2}}} .
        \]

  \item Product rule: $f'(x)=\sin x+x\cos x$.

  \item Simplify: $\ln(\sqrt{x})=\tfrac12\ln x$, so $f'(x)=\dfrac{1}{2x}$.

  \item Chain: $f'(x)=2\sin x\cdot\cos x=\sin(2x)$.

  \item $f'(x)=2^{x}\ln 2$ (general exponential rule).

  \item Inner $u=x^{2}+1$, $u'=2x$. $f'(x)=2x\,\mathrm{e}^{x^{2}+1}$.

  \item Chain on outer square: $f'(x)=2\ln x\cdot\dfrac{1}{x}
        =\dfrac{2\ln x}{x}$.

  \item Quotient rule:
        \[
          f'(x)=\frac{\cos x\cdot x-\sin x\cdot 1}{x^{2}}
               =\frac{x\cos x-\sin x}{x^{2}} .
        \]

  \item Product rule: $f'(x)=\arctan x+x\cdot\dfrac{1}{1+x^{2}}
        =\arctan x+\dfrac{x}{1+x^{2}}$.

  \item Inner $u=3x$, $u'=3$. $f'(x)=\operatorname{sech}^{2}(3x)\cdot 3
        =3\operatorname{sech}^{2}(3x)$.

  \item $f'(x)=\dfrac{1}{x\ln 10}$.

  \item Outer $\sin$, inner $\cos x$: $f'(x)=\cos(\cos x)\cdot(-\sin x)
        =-\sin x\cos(\cos x)$.

  \item Chain: $f'(x)=\mathrm{e}^{\sin x}\cdot\cos x=\cos x\,\mathrm{e}^{\sin x}$.

  \item Inner $u=1/x$, $u'=-1/x^{2}$.
        \[
          f'(x)=\frac{1}{1+(1/x)^{2}}\cdot\left(-\frac{1}{x^{2}}\right)
               =-\frac{1}{x^{2}+1} .
        \]
        (Aside: this confirms $\arctan x+\arctan(1/x)=\tfrac{\pi}{2}$ for
        $x>0$.)

  \item Expand $\ln\dfrac{x}{x+1}=\ln x-\ln(x+1)$, so
        $f'(x)=\dfrac{1}{x}-\dfrac{1}{x+1}=\dfrac{1}{x(x+1)}$.

  \item Inner $u=3x$, $u'=3$. $f'(x)=\sec(3x)\tan(3x)\cdot 3
        =3\sec(3x)\tan(3x)$.

  \item $U'(c)=\dfrac{1}{c}$. Marginal utility of consumption falls as $c$
        grows, diminishing marginal utility.

  \item Outer $\ln$, inner $\cos x$. $f'(x)=\dfrac{1}{\cos x}\cdot(-\sin x)
        =-\tan x$.

  \item Outer $u^{3}$, inner $\tan x$: $f'(x)=3\tan^{2}x\cdot\sec^{2}x$.

  \item Quotient rule with $u=1$, $v=1+\mathrm{e}^{-t}$. $u'=0$,
        $v'=-\mathrm{e}^{-t}$.
        \[
          P'(t)=\frac{\mathrm{e}^{-t}}{(1+\mathrm{e}^{-t})^{2}} .
        \]
        Equivalently, $P'(t)=P(t)\bigl(1-P(t)\bigr)$.

  \item Quotient rule. $u=\mathrm{e}^{x}$, $v=1+\mathrm{e}^{x}$,
        $u'=\mathrm{e}^{x}$, $v'=\mathrm{e}^{x}$.
        \[
          f'(x)=\frac{\mathrm{e}^{x}(1+\mathrm{e}^{x})-\mathrm{e}^{x}\cdot\mathrm{e}^{x}}
               {(1+\mathrm{e}^{x})^{2}}=\frac{\mathrm{e}^{x}}{(1+\mathrm{e}^{x})^{2}} .
        \]

  \item Product rule: $u=\mathrm{e}^{x}$, $v=\sin(2x)$, $u'=\mathrm{e}^{x}$,
        $v'=2\cos(2x)$.
        \[
          f'(x)=\mathrm{e}^{x}\bigl(\sin(2x)+2\cos(2x)\bigr) .
        \]

  \item Use the log power rule first: $f(K)=\alpha\ln K$, hence
        $f'(K)=\dfrac{\alpha}{K}$. (This is the elasticity-style structure
        that makes Cobb-Douglas convenient.)

  \item Logarithmic differentiation. Let $y=x^{x}$. Then $\ln y=x\ln x$.
        Differentiate:
        \[
          \frac{y'}{y}=\ln x+1,\qquad y'=x^{x}(\ln x+1) .
        \]

  \item Product rule with $u=\mathrm{e}^{2x}$, $v=\cos(3x)$,
        $u'=2\mathrm{e}^{2x}$, $v'=-3\sin(3x)$.
        \[
          f'(x)=2\mathrm{e}^{2x}\cos(3x)-3\mathrm{e}^{2x}\sin(3x)
               =\mathrm{e}^{2x}\bigl(2\cos(3x)-3\sin(3x)\bigr) .
        \]

  \item Inner $u=x+\sqrt{x^{2}+1}$. Note $u'=1+\dfrac{x}{\sqrt{x^{2}+1}}
        =\dfrac{\sqrt{x^{2}+1}+x}{\sqrt{x^{2}+1}}=\dfrac{u}{\sqrt{x^{2}+1}}$.
        Hence
        \[
          f'(x)=\frac{1}{u}\cdot\frac{u}{\sqrt{x^{2}+1}}=\frac{1}{\sqrt{x^{2}+1}} .
        \]
        (This is $\operatorname{arsinh}(x)$, so the result matches the
        derivative of inverse hyperbolic sine.)

  \item Product rule: $f'(x)=\cos x\cdot\cos x+\sin x\cdot(-\sin x)
        =\cos^{2}x-\sin^{2}x=\cos(2x)$.

  \item $D'(t)=-r\,\mathrm{e}^{-rt}=-r\,D(t)$. Discount factor decays at
        rate $r$; instantaneous proportional decay $D'/D=-r$.

  \item Inner $u=\mathrm{e}^{x}$, $u'=\mathrm{e}^{x}$.
        \[
          f'(x)=\frac{1}{1+\mathrm{e}^{2x}}\cdot\mathrm{e}^{x}
               =\frac{\mathrm{e}^{x}}{1+\mathrm{e}^{2x}} .
        \]

  \item Inner $u=x^{2}+1$, $u'=2x$. $f'(x)=\operatorname{sech}^{2}(x^{2}+1)
        \cdot 2x=2x\operatorname{sech}^{2}(x^{2}+1)$.

  \item Quotient rule with $u=\ln x$, $v=x$, $u'=1/x$, $v'=1$.
        \[
          f'(x)=\frac{(1/x)\cdot x-\ln x}{x^{2}}=\frac{1-\ln x}{x^{2}} .
        \]

  \item Product rule: $f'(x)=\mathrm{e}^{x}\ln x+\mathrm{e}^{x}\cdot\dfrac{1}{x}
        =\mathrm{e}^{x}\Bigl(\ln x+\dfrac{1}{x}\Bigr)$.

  \item $f(x)=(\sin x)^{1/2}$, inner $u'=\cos x$.
        \[
          f'(x)=\tfrac12(\sin x)^{-1/2}\cos x=\frac{\cos x}{2\sqrt{\sin x}} .
        \]

  \item $N'(t)=1000\cdot 0.04\,\mathrm{e}^{0.04t}=0.04\,N(t)$, so growth
        rate $N'/N=0.04=4\%$ per unit time.

  \item Rewrite $5^{x^{2}}=\mathrm{e}^{x^{2}\ln 5}$, so
        $f'(x)=\mathrm{e}^{x^{2}\ln 5}\cdot 2x\ln 5=2x\ln 5\cdot 5^{x^{2}}$.

  \item Inner $u=x^{2}$, $u'=2x$.
        \[
          f'(x)=\frac{2x}{\sqrt{1-x^{4}}} .
        \]

  \item $f'(x)=-\csc^{2}x-\csc x\cot x=-\csc x(\csc x+\cot x)$.

  \item Product rule: $f'(x)=\cosh^{2}x+\sinh^{2}x=\cosh(2x)$. (Use the
        hyperbolic double-angle identity.)

  \item Simplify first: $f(x)=\ln x+\tfrac12\ln(x^{2}+1)$, so
        \[
          f'(x)=\frac{1}{x}+\frac12\cdot\frac{2x}{x^{2}+1}
               =\frac{1}{x}+\frac{x}{x^{2}+1} .
        \]

  \item Outer $u^{4}$, inner $\cos x$, $u'=-\sin x$: $f'(x)=4\cos^{3}x\cdot
        (-\sin x)=-4\cos^{3}x\sin x$.

  \item Inner $u=x^{3}$, $u'=3x^{2}$. $f'(x)=3x^{2}\mathrm{e}^{x^{3}}$.

  \item $U'(c)=\dfrac{(1-\sigma)c^{-\sigma}}{1-\sigma}=c^{-\sigma}$.
        Marginal utility is positive and decreasing for $\sigma>0$.

  \item Inner $u=\sqrt{x}$, $u'=\dfrac{1}{2\sqrt{x}}$.
        \[
          f'(x)=\frac{1}{1+x}\cdot\frac{1}{2\sqrt{x}}
               =\frac{1}{2\sqrt{x}(1+x)} .
        \]

  \item Product rule: $f'(x)=\sec^{2}x\cdot\ln x+\tan x\cdot\dfrac{1}{x}$.

  \item $f(x)=(\ln x)^{-1}$, so $f'(x)=-(\ln x)^{-2}\cdot\dfrac{1}{x}
        =-\dfrac{1}{x(\ln x)^{2}}$.

  \item Product rule with $u=\mathrm{e}^{ax}$, $v=\sin(bx)$,
        $u'=a\mathrm{e}^{ax}$, $v'=b\cos(bx)$.
        \[
          f'(x)=\mathrm{e}^{ax}\bigl(a\sin(bx)+b\cos(bx)\bigr) .
        \]

  \item Use $\log_{2}(u)=\dfrac{\ln u}{\ln 2}$ with $u=x^{2}+1$:
        \[
          f'(x)=\frac{1}{\ln 2}\cdot\frac{2x}{x^{2}+1}
               =\frac{2x}{(x^{2}+1)\ln 2} .
        \]

  \item Product rule with $u=\cos(2x)$, $v=\mathrm{e}^{-x}$,
        $u'=-2\sin(2x)$, $v'=-\mathrm{e}^{-x}$.
        \[
          f'(x)=-2\sin(2x)\mathrm{e}^{-x}-\cos(2x)\mathrm{e}^{-x}
               =-\mathrm{e}^{-x}\bigl(2\sin(2x)+\cos(2x)\bigr) .
        \]

  \item Using $\operatorname{arccot}(x)=\tfrac{\pi}{2}-\arctan(x)$:
        \[
          f(x)=\Bigl(\frac{\pi}{2}-\arctan x\Bigr)\arctan x
              =\frac{\pi}{2}\arctan x-(\arctan x)^{2} .
        \]
        \[
          f'(x)=\frac{\pi}{2}\cdot\frac{1}{1+x^{2}}
               -2\arctan x\cdot\frac{1}{1+x^{2}}
               =\frac{\frac{\pi}{2}-2\arctan x}{1+x^{2}} .
        \]

  \item Outer $\ln$, inner $\ln x$. $f'(x)=\dfrac{1}{\ln x}\cdot\dfrac{1}{x}
        =\dfrac{1}{x\ln x}$.

  \item Product rule: $f'(x)=2x\,\mathrm{e}^{-x}+x^{2}(-\mathrm{e}^{-x})
        =x\,\mathrm{e}^{-x}(2-x)$.

  \item $g(t)=\ln(N_{0}\mathrm{e}^{rt})=\ln N_{0}+rt$, so $g'(t)=r$. Log
        of an exponential is linear: continuous growth rate is the slope
        of $\ln N$.

  \item Inner $u=\mathrm{e}^{x}$, $u'=\mathrm{e}^{x}$.
        $f'(x)=\cos(\mathrm{e}^{x})\cdot\mathrm{e}^{x}
        =\mathrm{e}^{x}\cos(\mathrm{e}^{x})$.

  \item Inner $u=1+\mathrm{e}^{x}$, $u'=\mathrm{e}^{x}$.
        $f'(x)=3(1+\mathrm{e}^{x})^{2}\cdot\mathrm{e}^{x}
        =3\mathrm{e}^{x}(1+\mathrm{e}^{x})^{2}$.

  \item Logarithmic differentiation. $\ln y=\sin x\cdot\ln x$.
        Differentiate:
        \[
          \frac{y'}{y}=\cos x\cdot\ln x+\sin x\cdot\frac{1}{x},
        \]
        so $y'=x^{\sin x}\Bigl(\cos x\ln x+\dfrac{\sin x}{x}\Bigr)$.

  \item For $x\in(-\tfrac{\pi}{2},\tfrac{\pi}{2})$ we have
        $\arctan(\tan x)=x$, so $f'(x)=1$. (Equivalently, chain rule:
        $\dfrac{1}{1+\tan^{2}x}\cdot\sec^{2}x=\dfrac{\sec^{2}x}{\sec^{2}x}=1$.)
\end{solutions}

\section{Additional Problems}

\begin{problems}
  \item $f(x)=\mathrm{e}^{5x}$
  \item $f(x)=\ln(7x+3)$
  \item $f(x)=\cos(4x)$
  \item $f(x)=\arctan(3x)$
  \item $f(x)=\mathrm{e}^{x^{3}}$
  \item $f(x)=\sin^{3}x$
  \item $f(x)=\ln(x^{3}+2)$
  \item $f(x)=x^{2}\,\mathrm{e}^{-x}$
  \item Continuous discount: $V(t)=1000\,\mathrm{e}^{-0.05t}$ \quad (find
        $V'(t)$)
  \item $f(x)=\dfrac{\ln x}{x^{2}}$
  \item $f(x)=\mathrm{e}^{x}\cos x$
  \item $f(x)=\arcsin(3x)$
  \item $f(x)=\tan(x^{2})$
  \item $f(x)=\ln(\sin x)$
  \item Log-utility variant: $U(x)=\ln(1+x)$ \quad (find $U'(x)$)
  \item $f(x)=3^{x^{2}+1}$
  \item $f(x)=(\ln x)^{3}$
  \item $f(x)=\mathrm{e}^{x^{2}}\sin x$
  \item $f(x)=\dfrac{\mathrm{e}^{x}-1}{\mathrm{e}^{x}+1}$
  \item $f(x)=(\cos x)^{x}$ \quad (use logarithmic differentiation,
        $\cos x>0$)
  \item $f(x)=\sinh(x)-\cosh(x)$
  \item $f(x)=\ln\Bigl(\dfrac{x^{2}}{x+1}\Bigr)$
  \item Per-capita capital: $k(t)=\ln(K(t)/L(t))$ with
        $K(t)=K_{0}\mathrm{e}^{at}$, $L(t)=L_{0}\mathrm{e}^{bt}$ \quad
        (find $k'(t)$)
  \item $f(x)=\arctan(\ln x)$
  \item $f(x)=\mathrm{e}^{\cos x}$
\end{problems}

\section{Additional Solutions}

\begin{solutions}
  \item Chain: $f'(x)=5\mathrm{e}^{5x}$.

  \item Chain: $f'(x)=\dfrac{7}{7x+3}$.

  \item Chain: $f'(x)=-4\sin(4x)$.

  \item Chain: $f'(x)=\dfrac{3}{1+9x^{2}}$.

  \item Chain: inner $u=x^{3}$, $u'=3x^{2}$. $f'(x)=3x^{2}\mathrm{e}^{x^{3}}$.

  \item Chain: $f'(x)=3\sin^{2}(x)\cos(x)$.

  \item Chain: $f'(x)=\dfrac{3x^{2}}{x^{3}+2}$.

  \item Product rule: $f'(x)=2x\,\mathrm{e}^{-x}-x^{2}\mathrm{e}^{-x}
        =x(2-x)\,\mathrm{e}^{-x}$.

  \item $V'(t)=1000\cdot(-0.05)\,\mathrm{e}^{-0.05t}=-50\,\mathrm{e}^{-0.05t}
        =-0.05\,V(t)$. Decays at rate $5\%$.

  \item Quotient rule with $u=\ln x$, $v=x^{2}$, $u'=1/x$, $v'=2x$.
        \[
          f'(x)=\frac{(1/x)x^{2}-\ln x\cdot 2x}{x^{4}}
               =\frac{x-2x\ln x}{x^{4}}=\frac{1-2\ln x}{x^{3}} .
        \]

  \item Product rule: $f'(x)=\mathrm{e}^{x}\cos x-\mathrm{e}^{x}\sin x
        =\mathrm{e}^{x}(\cos x-\sin x)$.

  \item Chain: $f'(x)=\dfrac{3}{\sqrt{1-9x^{2}}}$.

  \item Chain: inner $u=x^{2}$, $u'=2x$. $f'(x)=2x\sec^{2}(x^{2})$.

  \item Chain: $f'(x)=\dfrac{\cos x}{\sin x}=\cot x$.

  \item $U'(x)=\dfrac{1}{1+x}$. Marginal utility positive but diminishing
        as $x$ grows.

  \item Write $3^{x^{2}+1}=\mathrm{e}^{(x^{2}+1)\ln 3}$.
        $f'(x)=\mathrm{e}^{(x^{2}+1)\ln 3}\cdot 2x\ln 3
        =2x\ln 3\cdot 3^{x^{2}+1}$.

  \item Chain: $f'(x)=3(\ln x)^{2}\cdot\dfrac{1}{x}=\dfrac{3(\ln x)^{2}}{x}$.

  \item Product rule with $u=\mathrm{e}^{x^{2}}$, $v=\sin x$,
        $u'=2x\,\mathrm{e}^{x^{2}}$, $v'=\cos x$.
        \[
          f'(x)=2x\,\mathrm{e}^{x^{2}}\sin x+\mathrm{e}^{x^{2}}\cos x
               =\mathrm{e}^{x^{2}}(2x\sin x+\cos x) .
        \]

  \item Quotient rule with $u=\mathrm{e}^{x}-1$, $v=\mathrm{e}^{x}+1$,
        $u'=\mathrm{e}^{x}$, $v'=\mathrm{e}^{x}$.
        \[
          f'(x)=\frac{\mathrm{e}^{x}(\mathrm{e}^{x}+1)
                -(\mathrm{e}^{x}-1)\mathrm{e}^{x}}{(\mathrm{e}^{x}+1)^{2}}
               =\frac{2\mathrm{e}^{x}}{(\mathrm{e}^{x}+1)^{2}} .
        \]

  \item Logarithmic differentiation. $\ln y=x\ln(\cos x)$. Differentiate:
        \[
          \frac{y'}{y}=\ln(\cos x)+x\cdot\frac{-\sin x}{\cos x}
                      =\ln(\cos x)-x\tan x .
        \]
        So $y'=(\cos x)^{x}\bigl[\ln(\cos x)-x\tan x\bigr]$.

  \item $f'(x)=\cosh x-\sinh x$. (Equivalently, $f(x)=\mathrm{e}^{-x}$ and
        $f'(x)=-\mathrm{e}^{-x}=\cosh x-\sinh x$.)

  \item Simplify: $f(x)=2\ln x-\ln(x+1)$, so
        $f'(x)=\dfrac{2}{x}-\dfrac{1}{x+1}=\dfrac{x+2}{x(x+1)}$.

  \item $k(t)=\ln K_{0}+at-\ln L_{0}-bt=(a-b)t+\text{const}$, so
        $k'(t)=a-b$ (constant; difference of growth rates).

  \item Chain: outer $\arctan$, inner $\ln x$.
        $f'(x)=\dfrac{1}{1+(\ln x)^{2}}\cdot\dfrac{1}{x}
        =\dfrac{1}{x(1+(\ln x)^{2})}$.

  \item Chain: $f'(x)=\mathrm{e}^{\cos x}\cdot(-\sin x)
        =-\sin x\,\mathrm{e}^{\cos x}$.
\end{solutions}

\chapter{Higher-Order Derivatives, Convexity, Inflection}

Once $f'$ exists and is itself differentiable, we can differentiate again to
obtain $f''$, and iterate to higher orders. The second derivative records
how the slope changes: it controls convexity/concavity and powers the
second-order condition for optimisation. Higher derivatives of standard
functions follow simple closed-form patterns, and the Leibniz rule extends
the product rule to all orders.

\section{Core Rules}

\begin{rulebox}{Definition}
For $f:I\subseteq\R\to\R$ smooth enough, write
\[
  f''(x)=\frac{\mathrm{d}^{2}f}{\mathrm{d}x^{2}},\qquad
  f^{(n)}(x)=\frac{\mathrm{d}^{n}f}{\mathrm{d}x^{n}},\qquad
  f^{(0)}:=f .
\]
\end{rulebox}

\begin{rulebox}{Polynomials}
\[
  \frac{\mathrm{d}^{n}}{\mathrm{d}x^{n}}x^{k}
    =k(k-1)\cdots(k-n+1)\,x^{k-n} .
\]
In particular $(x^{n})^{(n)}=n!$ and $(x^{n})^{(n+1)}=0$.
\end{rulebox}

\begin{rulebox}{Exponentials}
\[
  \frac{\mathrm{d}^{n}}{\mathrm{d}x^{n}}\mathrm{e}^{ax}=a^{n}\mathrm{e}^{ax} .
\]
\end{rulebox}

\begin{rulebox}{Logarithm}
\[
  \frac{\mathrm{d}^{n}}{\mathrm{d}x^{n}}\ln(x)
    =\frac{(-1)^{n-1}(n-1)!}{x^{n}}\qquad\text{for } n\geq 1 .
\]
\end{rulebox}

\begin{rulebox}{Trig}
\[
  \frac{\mathrm{d}^{n}}{\mathrm{d}x^{n}}\sin(x)=\sin\Bigl(x+\frac{n\pi}{2}\Bigr),
  \qquad
  \frac{\mathrm{d}^{n}}{\mathrm{d}x^{n}}\cos(x)=\cos\Bigl(x+\frac{n\pi}{2}\Bigr) .
\]
\end{rulebox}

\begin{rulebox}{Convexity / concavity}
$f''\geq 0$ on $I$ implies $f$ is convex on $I$; $f''\leq 0$ implies
concave. Strict inequalities give strict (con)cavity.
\end{rulebox}

\begin{rulebox}{Inflection point}
$x_{0}$ is an inflection point if $f''$ \emph{changes sign} at $x_{0}$. The
condition $f''(x_{0})=0$ alone is insufficient (e.g.\ $x^{4}$ at $x_{0}=0$).
\end{rulebox}

\begin{rulebox}{Second-derivative test}
If $f'(x_{0})=0$ and $f''(x_{0})>0$, $x_{0}$ is a local minimum;
$f''(x_{0})<0$ gives a local maximum.
\end{rulebox}

\begin{rulebox}{Leibniz rule (higher product)}
\[
  (fg)^{(n)}=\sum_{k=0}^{n}\binom{n}{k}f^{(k)}g^{(n-k)} .
\]
\end{rulebox}

\begin{rulebox}{Taylor polynomial}
The Taylor polynomial of order $n$ at $a$:
\[
  T_{n}(x)=\sum_{k=0}^{n}\frac{f^{(k)}(a)}{k!}(x-a)^{k},
\]
with remainder
\[
  R_{n}(x)=\frac{f^{(n+1)}(\xi)}{(n+1)!}(x-a)^{n+1}
\]
for some $\xi$ between $a$ and $x$.
\end{rulebox}

\section{Practice Problems}

\begin{problems}
  \item Compute $f''(x)$ for $f(x)=2x^{4}-5x^{3}+x-7$.
  \item Determine the intervals on which $f(x)=x^{3}-6x^{2}+5$ is convex and
        on which it is concave.
  \item Find a closed form for $f^{(n)}(x)$ when $f(x)=\mathrm{e}^{-x}$.
  \item Find all inflection points of $f(x)=x^{6}$.
  \item Compute $f''(x)$ for $f(x)=\mathrm{e}^{-2x}$.
  \item For a cost function $C(Q)=Q^{3}-12Q^{2}+50Q+100$, determine the
        output level where marginal cost transitions from decreasing to
        increasing.
  \item Find a closed form for $f^{(n)}(x)$ when $f(x)=\ln(1+x)$.
  \item Compute $f''(x)$ for $f(x)=x^{2}\mathrm{e}^{x}$.
  \item Find all inflection points of $f(x)=x^{3}-9x^{2}+24x-3$.
  \item Determine the convexity/concavity intervals of $f(x)=x\mathrm{e}^{-x}$.
  \item Compute $f''(x)$ for $f(x)=\ln(2x+5)$.
  \item Find a closed form for $f^{(n)}(x)$ when $f(x)=\dfrac{1}{x+a}$
        (constant $a$).
  \item Show that $f(x)=x^{4}$ has $f''(0)=0$ but $x=0$ is \emph{not} an
        inflection point.
  \item Compute $f''(x)$ for $f(x)=3x^{2/3}-x^{-2}$.
  \item Determine the convexity intervals of $f(x)=\ln(1+\mathrm{e}^{x})$
        (softplus).
  \item Find all inflection points of $f(x)=x^{4}-12x^{2}+5$.
  \item Compute $f''(x)$ for $f(x)=x^{2}\ln(x)$.
  \item Find a closed form for $f^{(n)}(x)$ when $f(x)=\cos(x)$.
  \item For utility $U(x)=\ln(x)$, $x>0$, verify that $U$ is strictly
        concave (the standard ``risk-aversion'' condition).
  \item Compute $f''(x)$ for $f(x)=\tan(x)$.
  \item Find all inflection points of $f(x)=\dfrac{x^{2}}{x^{2}+3}$.
  \item Determine convexity/concavity intervals of $f(x)=\mathrm{e}^{x}-x$.
  \item Compute $f''(x)$ for $f(x)=\sin(x)\cos(x)$.
  \item Find a closed form for $f^{(n)}(x)$ when $f(x)=\dfrac{1}{(1-x)^{2}}$.
  \item Find all inflection points of $f(x)=x^{2}\mathrm{e}^{-x}$.
  \item Compute $f''(x)$ for $f(x)=\dfrac{x}{x-3}$.
  \item Determine convexity/concavity of $f(x)=x^{4}-8x^{3}+1$ and the
        corresponding inflection points.
  \item Find a closed form for $f^{(n)}(x)$ when $f(x)=\mathrm{e}^{ax}$ for
        a constant $a$.
  \item Compute $f''(x)$ for $f(x)=(x^{2}+4)^{3/2}$.
  \item For $f(x)=x^{5}$, determine whether $x=0$ is an inflection point.
  \item Compute $f''(x)$ for $f(x)=\sqrt{1+x^{2}}$.
  \item Determine the inflection points of $f(x)=\arctan(x)$.
  \item Find $f^{(n)}(x)$ for $f(x)=\sin(2x)$.
  \item Compute $f''(x)$ for $f(x)=\dfrac{1}{1+x^{2}}$.
  \item Determine convexity/concavity of $f(x)=x\ln(x)$, $x>0$.
  \item For the logistic function $f(x)=\dfrac{1}{1+\mathrm{e}^{-x}}$, find
        the inflection point.
  \item Compute $f''(x)$ for $f(x)=x^{3}\ln(x)$.
  \item For a profit function $\pi(Q)=-Q^{3}+15Q^{2}-27Q-10$, where is
        $\pi$ concave?
  \item Find $f^{(n)}(0)$ for $f(x)=\ln(1+x)$ at $x=0$ (i.e.\ the
        coefficients in its Taylor series at zero).
  \item Compute $f''(x)$ for $f(x)=\mathrm{e}^{x^{2}}$.
  \item Determine all inflection points of $f(x)=x^{5}-5x^{4}$.
  \item For $f(x)=x|x|$ (i.e.\ $x^{2}$ for $x\geq 0$, $-x^{2}$ for $x<0$),
        determine whether $x=0$ is an inflection point and whether
        $f''(0)$ exists.
  \item Compute $f''(x)$ for $f(x)=\arctan(x)$.
  \item Apply the Leibniz rule to compute $(x^{2}\mathrm{e}^{x})^{(n)}$ in
        closed form.
  \item Find the inflection points of $f(x)=(x-1)^{4}$.
  \item For $f(x)=\sin(x)$, find the Taylor polynomial $T_{5}(x)$ at $a=0$.
  \item Compute $f''(x)$ for $f(x)=\ln(\cos(x))$ on
        $\bigl(-\tfrac{\pi}{2},\tfrac{\pi}{2}\bigr)$.
  \item Find a closed form for $f^{(n)}(x)$ when $f(x)=x\mathrm{e}^{x}$,
        using the Leibniz rule.
  \item Determine the convexity intervals of $f(x)=\dfrac{1}{1+x^{2}}$ (the
        standard Cauchy-type bell curve).
  \item For a Cobb--Douglas production function $f(L)=AL^{\alpha}$ with
        $0<\alpha<1$ and $A>0$, verify that $f$ is strictly concave in
        $L>0$ (the law of diminishing returns).
\end{problems}

\section{Solutions}

\begin{solutions}
  \item $f'(x)=8x^{3}-15x^{2}+1$, $f''(x)=24x^{2}-30x$.

  \item $f'(x)=3x^{2}-12x$, $f''(x)=6x-12=6(x-2)$. So $f''(x)\geq 0$ iff
        $x\geq 2$: convex on $[2,\infty)$, concave on $(-\infty,2]$.

  \item Each differentiation introduces a factor of $-1$:
        $f^{(n)}(x)=(-1)^{n}\mathrm{e}^{-x}$.

  \item $f'(x)=6x^{5}$, $f''(x)=30x^{4}\geq 0$ for all $x$, with equality
        only at $x=0$. Since $f''$ does not change sign, $x=0$ is
        \emph{not} an inflection point. $f$ has no inflection points.

  \item $f'(x)=-2\mathrm{e}^{-2x}$, $f''(x)=4\mathrm{e}^{-2x}$.

  \item $MC(Q)=C'(Q)=3Q^{2}-24Q+50$. $MC'(Q)=6Q-24=0\Rightarrow Q=4$.
        $MC''(Q)=6>0$ so $Q=4$ minimises $MC$. At $Q=4$, $C''(Q)$ changes
        sign there: $C''(Q)=6Q-24$, negative for $Q<4$, positive for
        $Q>4$, so $Q=4$ is the inflection point of $C$ and the transition
        output for marginal cost.

  \item With $u=1+x$: $f'(x)=u^{-1}$, $f''(x)=-u^{-2}$,
        $f'''(x)=2u^{-3}$, in general
        \[
          f^{(n)}(x)=\frac{(-1)^{n-1}(n-1)!}{(1+x)^{n}}\qquad (n\geq 1) .
        \]

  \item $f'(x)=2x\mathrm{e}^{x}+x^{2}\mathrm{e}^{x}=(x^{2}+2x)\mathrm{e}^{x}$,
        then
        \[
          f''(x)=(2x+2)\mathrm{e}^{x}+(x^{2}+2x)\mathrm{e}^{x}
                =(x^{2}+4x+2)\mathrm{e}^{x} .
        \]

  \item $f'(x)=3x^{2}-18x+24$, $f''(x)=6x-18=6(x-3)$. Sign change at
        $x=3$. $f(3)=27-81+72-3=15$. Inflection at $(3,15)$.

  \item $f'(x)=\mathrm{e}^{-x}-x\mathrm{e}^{-x}=(1-x)\mathrm{e}^{-x}$.
        $f''(x)=-\mathrm{e}^{-x}-(1-x)\mathrm{e}^{-x}=(x-2)\mathrm{e}^{-x}$.
        Since $\mathrm{e}^{-x}>0$, the sign of $f''$ is that of $x-2$.
        Concave on $(-\infty,2]$, convex on $[2,\infty)$. Inflection at
        $x=2$.

  \item $f'(x)=\dfrac{2}{2x+5}$, $f''(x)=-\dfrac{4}{(2x+5)^{2}}$.

  \item $f(x)=(x+a)^{-1}$, so
        $f^{(n)}(x)=(-1)^{n}n!\,(x+a)^{-(n+1)}=\dfrac{(-1)^{n}n!}{(x+a)^{n+1}}$.

  \item $f'(x)=4x^{3}$, $f''(x)=12x^{2}$. Then $f''(0)=0$. But
        $f''(x)=12x^{2}\geq 0$ everywhere, no sign change. Hence $x=0$
        is not an inflection point. Indeed $f$ is convex on all of $\R$.

  \item $f'(x)=2x^{-1/3}+2x^{-3}$,
        $f''(x)=-\tfrac{2}{3}x^{-4/3}-6x^{-4}$.

  \item $f'(x)=\dfrac{\mathrm{e}^{x}}{1+\mathrm{e}^{x}}$. By the quotient
        rule
        \[
          f''(x)=\frac{\mathrm{e}^{x}(1+\mathrm{e}^{x})-\mathrm{e}^{x}\cdot\mathrm{e}^{x}}
                       {(1+\mathrm{e}^{x})^{2}}
                =\frac{\mathrm{e}^{x}}{(1+\mathrm{e}^{x})^{2}}>0
          \text{ for all } x .
        \]
        $f$ is strictly convex on all of $\R$.

  \item $f'(x)=4x^{3}-24x$, $f''(x)=12x^{2}-24=12(x^{2}-2)$. Sign changes
        at $x=\pm\sqrt{2}$. $f(\pm\sqrt{2})=4-24+5=-15$. Inflections at
        $(\pm\sqrt{2},-15)$.

  \item $f'(x)=2x\ln(x)+x$. Then $f''(x)=2\ln(x)+2+1=2\ln(x)+3$.

  \item $f^{(n)}(x)=\cos\bigl(x+\tfrac{n\pi}{2}\bigr)$.

  \item $U'(x)=1/x$, $U''(x)=-1/x^{2}<0$ for all $x>0$. So $U$ is strictly
        concave on $(0,\infty)$.

  \item $f'(x)=\sec^{2}(x)$. Then
        $f''(x)=2\sec(x)\cdot\sec(x)\tan(x)=2\sec^{2}(x)\tan(x)$.

  \item Write $f(x)=1-\dfrac{3}{x^{2}+3}$. Then
        $f'(x)=\dfrac{6x}{(x^{2}+3)^{2}}$ and
        \[
          f''(x)=\frac{6(x^{2}+3)^{2}-6x\cdot 2(x^{2}+3)\cdot 2x}
                       {(x^{2}+3)^{4}}
                =\frac{6(x^{2}+3)-24x^{2}}{(x^{2}+3)^{3}}
                =\frac{18-18x^{2}}{(x^{2}+3)^{3}} .
        \]
        Sign changes at $x=\pm 1$. $f(\pm 1)=1/4$. Inflections at
        $(\pm 1,\tfrac14)$.

  \item $f'(x)=\mathrm{e}^{x}-1$, $f''(x)=\mathrm{e}^{x}>0$ for all $x$.
        Strictly convex on $\R$. No concave region.

  \item $f(x)=\tfrac12\sin(2x)$, $f'(x)=\cos(2x)$, $f''(x)=-2\sin(2x)$.

  \item Note $\dfrac{1}{(1-x)^{2}}=\dfrac{\mathrm{d}}{\mathrm{d}x}
        \dfrac{1}{1-x}$. Since
        $\dfrac{\mathrm{d}^{k}}{\mathrm{d}x^{k}}\dfrac{1}{1-x}
        =\dfrac{k!}{(1-x)^{k+1}}$, we get
        \[
          f^{(n)}(x)=\frac{(n+1)!}{(1-x)^{n+2}} .
        \]

  \item $f'(x)=2x\mathrm{e}^{-x}-x^{2}\mathrm{e}^{-x}=(2x-x^{2})\mathrm{e}^{-x}$.
        $f''(x)=(2-2x)\mathrm{e}^{-x}-(2x-x^{2})\mathrm{e}^{-x}
               =(x^{2}-4x+2)\mathrm{e}^{-x}$. Roots: $x=2\pm\sqrt{2}$.
        Inflections at $x=2\pm\sqrt{2}$; $y$-values
        $(2\pm\sqrt{2})^{2}\mathrm{e}^{-(2\pm\sqrt{2})}
        =(6\pm4\sqrt{2})\mathrm{e}^{-(2\pm\sqrt{2})}$.

  \item Write $f(x)=1+\dfrac{3}{x-3}$. Then $f'(x)=-\dfrac{3}{(x-3)^{2}}$,
        $f''(x)=\dfrac{6}{(x-3)^{3}}$.

  \item $f'(x)=4x^{3}-24x^{2}$, $f''(x)=12x^{2}-48x=12x(x-4)$. Sign
        changes at $x=0$ and $x=4$. Convex on $(-\infty,0]\cup[4,\infty)$,
        concave on $[0,4]$. Inflections at $(0,1)$ and
        $(4,4^{4}-8\cdot 64+1)=(4,256-512+1)=(4,-255)$.

  \item Each derivative pulls out a factor $a$: $f^{(n)}(x)=a^{n}\mathrm{e}^{ax}$.

  \item $f'(x)=\tfrac32(x^{2}+4)^{1/2}\cdot 2x=3x\sqrt{x^{2}+4}$. Then
        \[
          f''(x)=3\sqrt{x^{2}+4}+3x\cdot\frac{x}{\sqrt{x^{2}+4}}
                =\frac{3(x^{2}+4)+3x^{2}}{\sqrt{x^{2}+4}}
                =\frac{6x^{2}+12}{\sqrt{x^{2}+4}}
                =\frac{6(x^{2}+2)}{\sqrt{x^{2}+4}} .
        \]

  \item $f'(x)=5x^{4}$, $f''(x)=20x^{3}$. $f''(x)<0$ for $x<0$,
        $f''(x)>0$ for $x>0$: sign change at $x=0$. Hence $x=0$ is an
        inflection point (with $f(0)=0$).

  \item $f'(x)=\dfrac{x}{\sqrt{1+x^{2}}}$. Then
        \[
          f''(x)=\frac{\sqrt{1+x^{2}}-x\cdot\frac{x}{\sqrt{1+x^{2}}}}
                       {1+x^{2}}
                =\frac{(1+x^{2})-x^{2}}{(1+x^{2})^{3/2}}
                =\frac{1}{(1+x^{2})^{3/2}} .
        \]

  \item $f'(x)=\dfrac{1}{1+x^{2}}$. Then
        $f''(x)=-\dfrac{2x}{(1+x^{2})^{2}}$. Sign change at $x=0$, with
        $f''(x)>0$ for $x<0$ and $f''(x)<0$ for $x>0$. Inflection at
        $(0,0)$.

  \item Write $\sin(2x)$. Each differentiation multiplies the
        argument-shift trick by $2$:
        \[
          f^{(n)}(x)=2^{n}\sin\Bigl(2x+\frac{n\pi}{2}\Bigr) .
        \]

  \item $f(x)=(1+x^{2})^{-1}$. $f'(x)=-2x(1+x^{2})^{-2}$. Then
        \[
          f''(x)=-2(1+x^{2})^{-2}+(-2x)\cdot(-2)(1+x^{2})^{-3}\cdot 2x
                =-\frac{2}{(1+x^{2})^{2}}+\frac{8x^{2}}{(1+x^{2})^{3}}
                =\frac{-2(1+x^{2})+8x^{2}}{(1+x^{2})^{3}}
                =\frac{6x^{2}-2}{(1+x^{2})^{3}} .
        \]

  \item $f'(x)=\ln(x)+1$, $f''(x)=1/x>0$ for all $x>0$. Strictly convex
        on $(0,\infty)$.

  \item $f(x)=\dfrac{1}{1+\mathrm{e}^{-x}}=\dfrac{\mathrm{e}^{x}}{1+\mathrm{e}^{x}}$.
        Then $f'(x)=\dfrac{\mathrm{e}^{x}}{(1+\mathrm{e}^{x})^{2}}$.
        Differentiate again using the quotient rule (with $u=\mathrm{e}^{x}$,
        $v=(1+\mathrm{e}^{x})^{2}$, $v'=2(1+\mathrm{e}^{x})\mathrm{e}^{x}$):
        \[
          f''(x)=\frac{\mathrm{e}^{x}(1+\mathrm{e}^{x})^{2}
                       -\mathrm{e}^{x}\cdot 2(1+\mathrm{e}^{x})\mathrm{e}^{x}}
                       {(1+\mathrm{e}^{x})^{4}}
                =\frac{\mathrm{e}^{x}(1+\mathrm{e}^{x})-2\mathrm{e}^{2x}}
                       {(1+\mathrm{e}^{x})^{3}}
                =\frac{\mathrm{e}^{x}(1-\mathrm{e}^{x})}{(1+\mathrm{e}^{x})^{3}} .
        \]
        Sign of $f''$ is the sign of $1-\mathrm{e}^{x}$: positive for
        $x<0$, negative for $x>0$. Inflection at $x=0$, where
        $f(0)=1/2$. Inflection point $(0,\tfrac12)$.

  \item $f'(x)=3x^{2}\ln(x)+x^{2}=x^{2}(3\ln(x)+1)$. Then
        \[
          f''(x)=2x(3\ln(x)+1)+x^{2}\cdot\frac{3}{x}
                =2x(3\ln(x)+1)+3x
                =6x\ln(x)+5x=x(6\ln(x)+5) .
        \]

  \item $\pi'(Q)=-3Q^{2}+30Q-27$, $\pi''(Q)=-6Q+30=6(5-Q)$. $\pi$ is
        concave on $\{Q:\pi''(Q)\leq 0\}=[5,\infty)$ and convex on
        $(-\infty,5]$. (Economic note: the profit function changes from
        convex to concave at the inflection $Q=5$.)

  \item $f^{(n)}(x)=\dfrac{(-1)^{n-1}(n-1)!}{(1+x)^{n}}$ for $n\geq 1$. At
        $x=0$: $f^{(n)}(0)=(-1)^{n-1}(n-1)!$ for $n\geq 1$, and $f(0)=0$.
        So the Taylor series is
        $\displaystyle\sum_{n\geq 1}\frac{(-1)^{n-1}}{n}x^{n}
        =x-\frac{x^{2}}{2}+\frac{x^{3}}{3}-\cdots$.

  \item $f'(x)=2x\mathrm{e}^{x^{2}}$. Then
        $f''(x)=2\mathrm{e}^{x^{2}}+2x\cdot 2x\mathrm{e}^{x^{2}}
               =(2+4x^{2})\mathrm{e}^{x^{2}}$.

  \item $f'(x)=5x^{4}-20x^{3}$, $f''(x)=20x^{3}-60x^{2}=20x^{2}(x-3)$.
        Now $20x^{2}\geq 0$, so the sign of $f''$ is determined by $x-3$
        except at $x=0$ where $f''=0$. Check signs: for $x<0$, $x^{2}>0$
        and $x-3<0$, so $f''<0$. For $0<x<3$, $x^{2}>0$ and $x-3<0$, so
        $f''<0$. No sign change at $x=0$. For $x>3$, $f''>0$. Sign change
        only at $x=3$. $f(3)=243-405=-162$. Inflection at $(3,-162)$.
        ($x=0$ is \emph{not} an inflection.)

  \item For $x>0$: $f(x)=x^{2}$, $f''(x)=2$. For $x<0$: $f(x)=-x^{2}$,
        $f''(x)=-2$. So $f''$ has different one-sided limits at $0$ ($2$
        from the right, $-2$ from the left). Hence $f''(0)$ does not
        exist. But $f'(x)=2|x|$ exists and is continuous everywhere, and
        $f''$ changes sign at $0$ (in the sense of one-sided values). The
        graph changes from concave (left) to convex (right) at $x=0$, so
        $x=0$ is an inflection point in the geometric sense even though
        $f''(0)$ itself is undefined.

  \item $f'(x)=\dfrac{1}{1+x^{2}}$. Then
        $f''(x)=-\dfrac{2x}{(1+x^{2})^{2}}$.

  \item Take $u=x^{2}$ and $v=\mathrm{e}^{x}$. Then $u'=2x$, $u''=2$,
        $u^{(k)}=0$ for $k\geq 3$, and $v^{(k)}=\mathrm{e}^{x}$. By
        Leibniz:
        \[
          (x^{2}\mathrm{e}^{x})^{(n)}
            =\binom{n}{0}x^{2}\mathrm{e}^{x}
             +\binom{n}{1}\cdot 2x\cdot\mathrm{e}^{x}
             +\binom{n}{2}\cdot 2\cdot\mathrm{e}^{x}
            =\mathrm{e}^{x}\bigl[x^{2}+2nx+n(n-1)\bigr] .
        \]

  \item $f'(x)=4(x-1)^{3}$, $f''(x)=12(x-1)^{2}\geq 0$ everywhere, with
        equality only at $x=1$. No sign change. Hence $x=1$ is \emph{not}
        an inflection point; $f$ is convex on all of $\R$.

  \item Using $f^{(k)}(0)=\sin(k\pi/2)$: values are $0,1,0,-1,0,1$ for
        $k=0,1,2,3,4,5$. Hence
        \[
          T_{5}(x)=x-\frac{x^{3}}{3!}+\frac{x^{5}}{5!}
                  =x-\frac{x^{3}}{6}+\frac{x^{5}}{120} .
        \]

  \item $f'(x)=-\tan(x)$. Then $f''(x)=-\sec^{2}(x)<0$ on
        $\bigl(-\tfrac{\pi}{2},\tfrac{\pi}{2}\bigr)$. So $f$ is strictly
        concave on that interval.

  \item Take $u=x$, $v=\mathrm{e}^{x}$. Then $u'=1$, $u^{(k)}=0$ for
        $k\geq 2$, and $v^{(k)}=\mathrm{e}^{x}$. By Leibniz:
        \[
          (x\mathrm{e}^{x})^{(n)}
            =\binom{n}{0}x\mathrm{e}^{x}+\binom{n}{1}\cdot 1\cdot\mathrm{e}^{x}
            =(x+n)\mathrm{e}^{x} .
        \]

  \item From problem 34: $f''(x)=\dfrac{6x^{2}-2}{(1+x^{2})^{3}}$. Sign of
        $f''$ equals sign of $6x^{2}-2$, i.e.\ positive iff
        $|x|>1/\sqrt{3}$. So $f$ is convex on
        $(-\infty,-1/\sqrt{3}]\cup[1/\sqrt{3},\infty)$ and concave on
        $[-1/\sqrt{3},1/\sqrt{3}]$. Inflection points at $x=\pm 1/\sqrt{3}$.

  \item $f'(L)=A\alpha L^{\alpha-1}$, $f''(L)=A\alpha(\alpha-1)L^{\alpha-2}$.
        With $A>0$, $0<\alpha<1$: $\alpha(\alpha-1)<0$ and
        $L^{\alpha-2}>0$, so $f''(L)<0$ for all $L>0$. Strictly concave on
        $(0,\infty)$, diminishing marginal product of labour.
\end{solutions}

\section{Additional Problems}

\begin{problems}
  \item Compute $f''(x)$ for $f(x)=3x^{5}-2x^{4}+x^{2}-7$.
  \item Determine the convexity/concavity intervals of
        $f(x)=x^{3}-3x^{2}-9x+5$.
  \item Find a closed form for $f^{(n)}(x)$ when $f(x)=\mathrm{e}^{3x}$.
  \item Find all inflection points of $f(x)=x^{4}-6x^{2}+8$.
  \item Compute $f''(x)$ for $f(x)=x\mathrm{e}^{-2x}$.
  \item For a cost function $C(Q)=Q^{3}-15Q^{2}+72Q+50$, determine the
        output level where marginal cost transitions from decreasing to
        increasing.
  \item Find a closed form for $f^{(n)}(x)$ when $f(x)=\ln(2-x)$, on
        $x<2$.
  \item Compute $f''(x)$ for $f(x)=\sqrt{x^{2}+4}$.
  \item Find all inflection points of $f(x)=x^{3}-12x^{2}+45x-7$.
  \item Determine the convexity intervals of $f(x)=\ln(x^{2}+1)$.
  \item Compute $f''(x)$ for $f(x)=\sin(x)\,\mathrm{e}^{x}$.
  \item Find a closed form for $f^{(n)}(x)$ when
        $f(x)=\dfrac{1}{(1-2x)}$.
  \item For utility $U(x)=\sqrt{x}$, $x>0$, verify that $U$ is strictly
        concave (diminishing marginal utility).
  \item Compute $f''(x)$ for $f(x)=x^{3}\mathrm{e}^{-x}$.
  \item Apply the Leibniz rule to compute $(x^{3}\mathrm{e}^{x})^{(n)}$ in
        closed form.
\end{problems}

\section{Additional Solutions}

\begin{solutions}
  \item $f'(x)=15x^{4}-8x^{3}+2x$, $f''(x)=60x^{3}-24x^{2}+2$.

  \item $f'(x)=3x^{2}-6x-9$, $f''(x)=6x-6=6(x-1)$. So $f''\geq 0$ iff
        $x\geq 1$: convex on $[1,\infty)$, concave on $(-\infty,1]$.
        Inflection at $x=1$, $f(1)=-6$.

  \item Each differentiation introduces a factor of $3$:
        $f^{(n)}(x)=3^{n}\mathrm{e}^{3x}$.

  \item $f'(x)=4x^{3}-12x$, $f''(x)=12x^{2}-12=12(x^{2}-1)$. Sign changes
        at $x=\pm 1$. $f(\pm 1)=1-6+8=3$. Inflections at $(\pm 1,3)$.

  \item Product rule:
        $f'(x)=\mathrm{e}^{-2x}-2x\mathrm{e}^{-2x}=(1-2x)\mathrm{e}^{-2x}$.
        Then
        \[
          f''(x)=-2\mathrm{e}^{-2x}-2(1-2x)\mathrm{e}^{-2x}
                =(4x-4)\mathrm{e}^{-2x}=4(x-1)\mathrm{e}^{-2x} .
        \]

  \item $MC(Q)=C'(Q)=3Q^{2}-30Q+72$. $MC'(Q)=6Q-30=0\Rightarrow Q=5$.
        $MC''(Q)=6>0$, so $Q=5$ minimises $MC$. Equivalently,
        $C''(Q)=6Q-30$ changes sign at $Q=5$: the cost function
        transitions from concave to convex at the inflection $Q=5$.

  \item With $u=2-x$: $f'(x)=-1/(2-x)$, $f''(x)=-1/(2-x)^{2}$,
        $f'''(x)=-2/(2-x)^{3}$. By induction:
        \[
          f^{(n)}(x)=-\frac{(n-1)!}{(2-x)^{n}}\qquad (n\geq 1) .
        \]

  \item $f(x)=(x^{2}+4)^{1/2}$. $f'(x)=\dfrac{x}{\sqrt{x^{2}+4}}$. Then
        \[
          f''(x)=\frac{\sqrt{x^{2}+4}-x\cdot\frac{x}{\sqrt{x^{2}+4}}}
                       {x^{2}+4}
                =\frac{(x^{2}+4)-x^{2}}{(x^{2}+4)^{3/2}}
                =\frac{4}{(x^{2}+4)^{3/2}} .
        \]

  \item $f'(x)=3x^{2}-24x+45$, $f''(x)=6x-24=6(x-4)$. Sign change at
        $x=4$. $f(4)=64-192+180-7=45$. Inflection at $(4,45)$.

  \item $f'(x)=\dfrac{2x}{x^{2}+1}$. By the quotient rule
        \[
          f''(x)=\frac{2(x^{2}+1)-2x\cdot 2x}{(x^{2}+1)^{2}}
                =\frac{2-2x^{2}}{(x^{2}+1)^{2}}
                =\frac{2(1-x^{2})}{(x^{2}+1)^{2}} .
        \]
        Sign of $f''$ matches sign of $1-x^{2}$. Convex on $[-1,1]$,
        concave on $(-\infty,-1]\cup[1,\infty)$. Inflections at
        $x=\pm 1$.

  \item $f'(x)=\cos(x)\mathrm{e}^{x}+\sin(x)\mathrm{e}^{x}
        =(\cos x+\sin x)\mathrm{e}^{x}$. Then
        \[
          f''(x)=(-\sin x+\cos x)\mathrm{e}^{x}
                 +(\cos x+\sin x)\mathrm{e}^{x}
                =2\cos(x)\,\mathrm{e}^{x} .
        \]

  \item $f(x)=(1-2x)^{-1}$. By induction:
        $f^{(n)}(x)=\dfrac{2^{n}\,n!}{(1-2x)^{n+1}}$.

  \item $U'(x)=\tfrac12 x^{-1/2}$, $U''(x)=-\tfrac14 x^{-3/2}<0$ for all
        $x>0$. So $U$ is strictly concave on $(0,\infty)$.

  \item $f'(x)=3x^{2}\mathrm{e}^{-x}-x^{3}\mathrm{e}^{-x}
        =(3x^{2}-x^{3})\mathrm{e}^{-x}$. Then
        \[
          f''(x)=(6x-3x^{2})\mathrm{e}^{-x}-(3x^{2}-x^{3})\mathrm{e}^{-x}
                =(x^{3}-6x^{2}+6x)\mathrm{e}^{-x}
                =x(x^{2}-6x+6)\mathrm{e}^{-x} .
        \]

  \item Take $u=x^{3}$ and $v=\mathrm{e}^{x}$. Then $u'=3x^{2}$,
        $u''=6x$, $u'''=6$, $u^{(k)}=0$ for $k\geq 4$, and
        $v^{(k)}=\mathrm{e}^{x}$. By Leibniz:
        \[
          (x^{3}\mathrm{e}^{x})^{(n)}
            =\binom{n}{0}x^{3}\mathrm{e}^{x}
             +\binom{n}{1}\cdot 3x^{2}\cdot\mathrm{e}^{x}
             +\binom{n}{2}\cdot 6x\cdot\mathrm{e}^{x}
             +\binom{n}{3}\cdot 6\cdot\mathrm{e}^{x}
            =\mathrm{e}^{x}\bigl[x^{3}+3nx^{2}+3n(n-1)x+n(n-1)(n-2)\bigr] .
        \]
\end{solutions}

\chapter{Taylor Approximation}

The previous chapter closed with a formula for the Taylor polynomial;
this chapter turns it into a working tool. Most functions that appear in
economics (utility, production, growth) are not polynomials, yet
polynomials are the easiest objects to evaluate and to reason about.
Taylor's theorem says that near a chosen base point every sufficiently
smooth function \emph{is} a polynomial, up to a controllable error. This
single idea underlies the linearisation of models around a steady state,
the log-linearisation trick $\ln(1+g)\approx g$ for growth rates, the
second-order conditions of optimisation, and quick numerical estimates
such as $\sqrt{4.1}$ or $\mathrm{e}^{0.1}$ without a calculator.

\section{Core Rules}

\begin{rulebox}{Taylor polynomial of order $n$}
Let $f$ be $n$ times differentiable at $a$. The \emph{Taylor polynomial}
of order $n$ of $f$ at $a$ is
\[
  T_{n}(x)=\sum_{k=0}^{n}\frac{f^{(k)}(a)}{k!}(x-a)^{k}
          =f(a)+f'(a)(x-a)+\frac{f''(a)}{2!}(x-a)^{2}+\cdots
           +\frac{f^{(n)}(a)}{n!}(x-a)^{n}.
\]
It is the \emph{unique} polynomial of degree $\leq n$ whose value and
first $n$ derivatives at $a$ all agree with those of $f$:
$T_{n}^{(k)}(a)=f^{(k)}(a)$ for $k=0,1,\ldots,n$.
\end{rulebox}

\begin{rulebox}{Linear and quadratic approximation ($n=1,2$)}
\[
  T_{1}(x)=f(a)+f'(a)(x-a),
  \qquad
  T_{2}(x)=f(a)+f'(a)(x-a)+\tfrac12 f''(a)(x-a)^{2}.
\]
$T_{1}$ is the tangent line at $a$ (the \emph{linearisation} of $f$);
$T_{2}$ is the best-fitting parabola, matching value, slope and
curvature at $a$. Each extra order typically buys one further correct
decimal place near $a$.
\end{rulebox}

\begin{rulebox}{Lagrange remainder and error bound}
If $f$ is $(n+1)$ times differentiable on an interval containing $a$ and
$x$, then
\[
  f(x)=T_{n}(x)+R_{n}(x),
  \qquad
  R_{n}(x)=\frac{f^{(n+1)}(\xi)}{(n+1)!}(x-a)^{n+1}
\]
for some $\xi$ strictly between $a$ and $x$. Since $\xi$ is unknown,
bound it: if $|f^{(n+1)}(t)|\leq M$ for all $t$ between $a$ and $x$,
then
\[
  \bigl|f(x)-T_{n}(x)\bigr|\leq\frac{M}{(n+1)!}\,|x-a|^{n+1}.
\]
The error shrinks like $|x-a|^{n+1}$: halving the distance to $a$
divides the error of $T_{n}$ by about $2^{n+1}$.
\end{rulebox}

\begin{rulebox}{Maclaurin case ($a=0$): standard expansions}
The Taylor polynomial at $a=0$ is called the \emph{Maclaurin}
polynomial. The classical expansions are
\[
\begin{aligned}
  \mathrm{e}^{x}&=1+x+\frac{x^{2}}{2!}+\frac{x^{3}}{3!}+\cdots
                  +\frac{x^{n}}{n!}+\cdots
                  &&\text{(all }x\text{)},\\
  \ln(1+x)&=x-\frac{x^{2}}{2}+\frac{x^{3}}{3}-\frac{x^{4}}{4}+\cdots
                  &&(-1<x\leq 1),\\
  (1+x)^{k}&=1+kx+\frac{k(k-1)}{2!}x^{2}
                  +\frac{k(k-1)(k-2)}{3!}x^{3}+\cdots
                  &&(|x|<1),\\
  \sin(x)&=x-\frac{x^{3}}{3!}+\frac{x^{5}}{5!}-\cdots
                  &&\text{(all }x\text{)},\\
  \cos(x)&=1-\frac{x^{2}}{2!}+\frac{x^{4}}{4!}-\cdots
                  &&\text{(all }x\text{)}.
\end{aligned}
\]
The first-order case $\ln(1+x)\approx x$ is the workhorse of
log-linearisation in economics: a small growth rate $g$ and the log
change $\ln(1+g)$ are interchangeable to first order.
\end{rulebox}

\begin{rulebox}{Quadratic approximation in economics}
Near a stationary point $a$ (where $f'(a)=0$) the linear term vanishes
and
\[
  f(x)\approx f(a)+\tfrac12 f''(a)(x-a)^{2}:
\]
locally, $f$ \emph{is} a parabola sitting on top of $f(a)$. If
$f''(a)<0$ the parabola opens downward and $a$ is a local maximum; if
$f''(a)>0$ it opens upward and $a$ is a local minimum. This is exactly
where the second-order condition comes from. In several variables the
role of $f''(a)$ is played by the Hessian, and the sign condition
becomes definiteness.
\end{rulebox}

\begin{rulebox}{Caution: validity is local}
A Taylor polynomial is a \emph{local} approximation: the error grows
like $|x-a|^{n+1}$, so accuracy is excellent near $a$ and can be
useless far away. For $\ln(1+x)$ and $(1+x)^{k}$ the expansions are
only valid on $(-1,1]$ resp.\ $(-1,1)$: outside, adding more terms
makes things \emph{worse}. Even inside: $T_{1}(x)=x$ for $\ln(1+x)$
gives $\ln(2)\approx 1$ against the true $0.693$, a $44\%$ error,
versus $0.5\%$ at $x=0.1$. Rule of thumb: expand around a base point
close to where you evaluate, and use the remainder bound to check.
\end{rulebox}

\section{Worked Example}

\begin{worked}{Expand $\ln(1+x)$ to third order about $0$, bound the error at $x=0.5$, and read the economics off the second-order term}
\wstep{Step 1. Collect the derivatives at the expansion point.}
\[
  f(x)=\ln(1+x),\quad
  f'=\frac{1}{1+x},\quad
  f''=\frac{-1}{(1+x)^{2}},\quad
  f'''=\frac{2}{(1+x)^{3}},\quad
  f^{(4)}=\frac{-6}{(1+x)^{4}} .
\]
At $a=0$: $f(0)=0$, $f'(0)=1$, $f''(0)=-1$, $f'''(0)=2$, and
$f^{(4)}(0)=-6$.
\wtrap{Evaluate \emph{after} differentiating, never before. Substituting
$x=0$ into $f'$ and then differentiating the resulting constant gives zeros
all the way down --- a surprisingly frequent way to produce the polynomial
$P_{3}(x)=x$.}

\wstep{Step 2. Assemble the Taylor polynomial.}
\[
  P_{n}(x)=\sum_{k=0}^{n}\frac{f^{(k)}(a)}{k!}(x-a)^{k}
  \quad\Longrightarrow\quad
  P_{3}(x)=0+x-\frac{1}{2!}x^{2}+\frac{2}{3!}x^{3}
  =\boxed{x-\frac{x^{2}}{2}+\frac{x^{3}}{3}}
\]
The factorials are what turn the raw derivatives $1,-1,2$ into the
coefficients $1,-\tfrac12,\tfrac13$.
\wtrap{Do not omit $k!$, and do not omit $(x-a)^{k}$ when $a\neq0$. Expanding
about $a=2$ and writing powers of $x$ rather than of $x-2$ is the standard
error, and it silently produces a polynomial that is wrong everywhere.}

\wstep{Step 3. Evaluate at $x=0.5$ and compare.}
\[
  P_{3}(0.5)=0.5-0.125+0.041\overline{6}=0.4166\overline{6},
  \qquad
  \ln(1.5)=0.405465 ,
\]
so the actual error is $0.011202$.

\wstep{Step 4. Bound the error \emph{before} knowing the true value.}
The Lagrange remainder is
\[
  R_{3}(x)=\frac{f^{(4)}(\xi)}{4!}x^{4}
  \qquad\text{for some }\xi\text{ between }0\text{ and }x .
\]
$\xi$ is unknown, so take the worst case. On $[0,0.5]$,
\[
  \bigl|f^{(4)}(\xi)\bigr|=\frac{6}{(1+\xi)^{4}}\le\frac{6}{1^{4}}=6
  \qquad(\text{largest at }\xi=0),
\]
\[
  |R_{3}(0.5)|\le\frac{6}{24}(0.5)^{4}=0.25\cdot0.0625=0.015625 .
\]
And indeed $0.011201\le0.015625$. \checkmark
\wtrap{Maximise $|f^{(n+1)}|$ over the whole interval between $a$ and $x$ ---
do not substitute $\xi=x$ or $\xi=a$ and call it a bound. Here $1/(1+\xi)^{4}$
is decreasing, so the worst case is at the \emph{left} end; for an
increasing derivative it would be at the right. Check which end before
committing.}

\wstep{Step 5. Note how fast the accuracy degrades with distance.}
At $x=0.1$ the same bound gives $\tfrac14(0.1)^{4}=0.000025$; at $x=0.5$ it
is $625$ times larger. The error carries the factor $x^{4}$, so halving the
distance to the expansion point cuts the error sixteenfold. At $x=1$ the
series still converges, but at $x=1.1$ it diverges outright --- the
expansion of $\ln(1+x)$ is valid only for $|x|<1$.
\wtrap{A Taylor polynomial is a \textbf{local} statement. It is excellent
near $a$, useless far from it, and there is no warning in the formula --- 
$P_{3}(5)=5-12.5+41.7=34.2$, while $\ln6=1.79$. Always ask how far from $a$
you are being asked to work.}

\wstep{Step 6. The first two terms are the economics.}
Keeping $n=1$ gives the \textbf{linear approximation} $\ln(1+x)\approx x$:
the familiar ``a log difference is a percentage change'', accurate to about
half a percentage point for $|x|\le0.1$. Keeping $n=2$ gives
\[
  f(x)\approx f(a)+f'(a)(x-a)+\tfrac12 f''(a)(x-a)^{2},
\]
whose quadratic term carries the \textbf{curvature}. In expected-utility
theory that term is exactly what produces the risk premium: a concave $u$
has $u''<0$, so the quadratic term is negative, and its size relative to
$u'$ is the Arrow--Pratt coefficient $-u''/u'$. Risk aversion \emph{is} the
second-order term of a Taylor expansion.

\wstep{Step 7. The expansions worth knowing by heart.}
\[
  \mathrm{e}^{x}=1+x+\tfrac{x^{2}}{2}+\tfrac{x^{3}}{6}+\cdots,
  \qquad
  \ln(1+x)=x-\tfrac{x^{2}}{2}+\tfrac{x^{3}}{3}-\cdots,
\]
\[
  (1+x)^{\alpha}=1+\alpha x+\tfrac{\alpha(\alpha-1)}{2}x^{2}+\cdots,
  \qquad
  \frac{1}{1-x}=1+x+x^{2}+\cdots
\]
The last is the geometric series of Volume~4, which is why a perpetuity
formula and a Taylor expansion are the same computation twice.
\end{worked}

\section{Practice Problems}

\begin{problems}
  \item Find $T_{1}$ of $f(x)=\ln(x)$ at $a=1$.
  \item Find $T_{1}$ of $f(x)=\sqrt{x}$ at $a=4$.
  \item Find $T_{1}$ of $f(x)=\dfrac{1}{x}$ at $a=1$.
  \item Find $T_{1}$ of $f(x)=x^{3}-2x$ at $a=2$.
  \item Find $T_{1}$ of $f(x)=\mathrm{e}^{2x}$ at $a=\ln 2$.
  \item Find $T_{2}$ of $f(x)=\mathrm{e}^{x}$ at $a=0$ and use it to
        approximate $\mathrm{e}^{0.1}$. Compare with the true value
        $\mathrm{e}^{0.1}=1.105171$.
  \item Find $T_{2}$ of $f(x)=\ln(x)$ at $a=1$ and use it to
        approximate $\ln(1.05)$. Compare with the true value
        $\ln(1.05)=0.048790$.
  \item Find $T_{2}$ of $f(x)=\sqrt{x}$ at $a=4$ and use it to
        approximate $\sqrt{4.1}$. Compare with the true value
        $\sqrt{4.1}=2.024846$.
  \item Find $T_{2}$ of $f(x)=\dfrac{1}{1-x}$ at $a=0$.
  \item Find $T_{2}$ of $f(x)=x\ln(x)$ at $a=1$.
  \item Find $T_{2}$ of $f(x)=\mathrm{e}^{-x^{2}}$ at $a=0$.
  \item Find $T_{3}$ of $f(x)=\ln(1+x)$ at $a=0$ and use it to
        approximate $\ln(0.9)$. Compare with the true value
        $\ln(0.9)=-0.105361$.
  \item Use the binomial expansion to find $T_{2}$ of
        $f(x)=(1+x)^{1/2}$ at $a=0$, and use it to approximate
        $\sqrt{1.1}$. Compare with the true value
        $\sqrt{1.1}=1.048809$.
  \item Find the Maclaurin polynomial $T_{4}$ of $f(x)=\cos(x)$ and use
        it to approximate $\cos(0.2)$. Compare with the true value
        $\cos(0.2)=0.980067$.
  \item A firm produces $Q(L)=10\sqrt{L}$ units with $L$ hours of
        labour. Find $T_{2}$ of $Q$ at $a=100$ and use it to
        approximate output at $L=104$. Compare with the true value
        $Q(104)=101.980390$.
  \item (Log-linearisation.) Output grows from $Y_{0}$ to
        $Y_{1}=Y_{0}(1+g)$. Show, using $T_{1}$ of $\ln(1+x)$ at $0$,
        that $\ln Y_{1}-\ln Y_{0}\approx g$, and quantify the
        approximation for $g=0.03$ given $\ln(1.03)=0.029559$.
  \item Use the Lagrange remainder to bound the error of the $T_{1}$
        approximation of $\sqrt{4.1}$ at $a=4$ (Problem 2), and compare
        the bound with the actual error.
  \item Use the Lagrange remainder to bound the error of the $T_{2}$
        approximation of $\mathrm{e}^{0.1}$ at $a=0$ (Problem 6), and
        compare the bound with the actual error.
  \item Find the smallest $n$ such that the Maclaurin polynomial
        $T_{n}$ of $\mathrm{e}^{x}$ approximates $\mathrm{e}^{0.5}$
        with error at most $10^{-4}$. (Use $\mathrm{e}^{0.5}<2$ in the
        remainder bound.)
  \item Prove, using $T_{1}$ with Lagrange remainder: if $f'(a)=0$ and
        $f''$ is continuous with $f''(a)<0$, then $a$ is a strict local
        maximum of $f$.
\end{problems}

\section{Solutions}

\begin{solutions}
  \item $f(1)=\ln(1)=0$; $f'(x)=1/x$, so $f'(1)=1$. Hence
        \[
          T_{1}(x)=0+1\cdot(x-1)=x-1 .
        \]

  \item $f(4)=2$; $f'(x)=\dfrac{1}{2\sqrt{x}}$, so
        $f'(4)=\dfrac{1}{2\cdot 2}=\dfrac14$. Hence
        \[
          T_{1}(x)=2+\tfrac14(x-4) .
        \]

  \item $f(1)=1$; $f'(x)=-1/x^{2}$, so $f'(1)=-1$. Hence
        \[
          T_{1}(x)=1-(x-1)=2-x .
        \]

  \item $f(2)=8-4=4$; $f'(x)=3x^{2}-2$, so $f'(2)=12-2=10$. Hence
        \[
          T_{1}(x)=4+10(x-2) .
        \]

  \item $f(\ln 2)=\mathrm{e}^{2\ln 2}=(\mathrm{e}^{\ln 2})^{2}=4$;
        $f'(x)=2\mathrm{e}^{2x}$, so $f'(\ln 2)=2\cdot 4=8$. Hence
        \[
          T_{1}(x)=4+8(x-\ln 2) .
        \]

  \item $f^{(k)}(x)=\mathrm{e}^{x}$ for every $k$, so
        $f(0)=f'(0)=f''(0)=1$ and
        \[
          T_{2}(x)=1+x+\frac{x^{2}}{2} .
        \]
        At $x=0.1$: $T_{2}(0.1)=1+0.1+\tfrac{0.01}{2}=1.105$. True
        value $1.105171$; the error is about $1.7\cdot 10^{-4}$.

  \item $f(1)=0$; $f'(x)=1/x$, $f'(1)=1$; $f''(x)=-1/x^{2}$,
        $f''(1)=-1$. Hence
        \[
          T_{2}(x)=(x-1)-\tfrac12(x-1)^{2} .
        \]
        At $x=1.05$:
        $T_{2}(1.05)=0.05-\tfrac12(0.0025)=0.05-0.00125=0.04875$. True
        value $0.048790$; the error is about $4.0\cdot 10^{-5}$.

  \item From Problem 2, $f(4)=2$ and $f'(4)=\tfrac14$. Further
        $f''(x)=-\tfrac14 x^{-3/2}$, so
        $f''(4)=-\tfrac14\cdot\tfrac18=-\tfrac{1}{32}$ and the
        quadratic coefficient is
        $\tfrac12 f''(4)=-\tfrac{1}{64}$:
        \[
          T_{2}(x)=2+\tfrac14(x-4)-\tfrac{1}{64}(x-4)^{2} .
        \]
        At $x=4.1$:
        \[
          T_{2}(4.1)=2+0.025-\tfrac{0.01}{64}
                    =2.025-0.00015625=2.02484375 .
        \]
        True value $2.024846$; the error is about $1.9\cdot 10^{-6}$
        (the linear estimate $2.025$ alone is off by
        $1.5\cdot 10^{-4}$).

  \item $f(0)=1$; $f'(x)=\dfrac{1}{(1-x)^{2}}$, $f'(0)=1$;
        $f''(x)=\dfrac{2}{(1-x)^{3}}$, $f''(0)=2$, so
        $\tfrac12 f''(0)=1$. Hence
        \[
          T_{2}(x)=1+x+x^{2} ,
        \]
        the start of the geometric series $1+x+x^{2}+\cdots$.

  \item $f(1)=1\cdot\ln(1)=0$; $f'(x)=\ln(x)+1$, $f'(1)=1$;
        $f''(x)=1/x$, $f''(1)=1$. Hence
        \[
          T_{2}(x)=(x-1)+\tfrac12(x-1)^{2} .
        \]

  \item $f(0)=1$; $f'(x)=-2x\mathrm{e}^{-x^{2}}$, $f'(0)=0$;
        $f''(x)=(-2+4x^{2})\mathrm{e}^{-x^{2}}$, $f''(0)=-2$. Hence
        \[
          T_{2}(x)=1-x^{2} .
        \]
        (Equivalently: substitute $u=-x^{2}$ into
        $\mathrm{e}^{u}\approx 1+u$.)

  \item $f(0)=0$; $f'(x)=\dfrac{1}{1+x}$, $f'(0)=1$;
        $f''(x)=-\dfrac{1}{(1+x)^{2}}$, $f''(0)=-1$;
        $f'''(x)=\dfrac{2}{(1+x)^{3}}$, $f'''(0)=2$. The coefficients
        are $1$, $-\tfrac{1}{2!}=-\tfrac12$ and
        $\tfrac{2}{3!}=\tfrac13$:
        \[
          T_{3}(x)=x-\frac{x^{2}}{2}+\frac{x^{3}}{3} .
        \]
        At $x=-0.1$:
        \[
          T_{3}(-0.1)=-0.1-\frac{0.01}{2}+\frac{-0.001}{3}
                     =-0.1-0.005-0.000333=-0.105333 .
        \]
        True value $-0.105361$; the error is about $2.7\cdot 10^{-5}$.

  \item With $k=\tfrac12$ in the binomial expansion: the linear
        coefficient is $k=\tfrac12$ and the quadratic coefficient is
        \[
          \frac{k(k-1)}{2!}
            =\frac{\tfrac12\cdot\bigl(-\tfrac12\bigr)}{2}
            =-\frac18 ,
        \]
        so
        \[
          T_{2}(x)=1+\frac{x}{2}-\frac{x^{2}}{8} .
        \]
        At $x=0.1$: $T_{2}(0.1)=1+0.05-0.00125=1.04875$. True value
        $1.048809$; the error is about $5.9\cdot 10^{-5}$.

  \item Derivatives of $\cos$ at $0$ cycle through
        $1,\,0,\,-1,\,0,\,1$ for $k=0,\ldots,4$, so
        \[
          T_{4}(x)=1-\frac{x^{2}}{2!}+\frac{x^{4}}{4!}
                  =1-\frac{x^{2}}{2}+\frac{x^{4}}{24} .
        \]
        At $x=0.2$: $(0.2)^{2}=0.04$ and $(0.2)^{4}=0.0016$, so
        \[
          T_{4}(0.2)=1-0.02+\frac{0.0016}{24}
                    =1-0.02+0.000067=0.980067 ,
        \]
        agreeing with the true value $0.980067$ to all six decimals
        shown (the next term is of size $\tfrac{(0.2)^{6}}{720}\approx
        9\cdot 10^{-8}$).

  \item $Q(100)=10\cdot 10=100$; $Q'(L)=5L^{-1/2}$,
        $Q'(100)=\tfrac{5}{10}=0.5$; $Q''(L)=-\tfrac52 L^{-3/2}$,
        $Q''(100)=-\tfrac{2.5}{1000}=-0.0025$, so
        $\tfrac12 Q''(100)=-0.00125$. Hence
        \[
          T_{2}(L)=100+0.5\,(L-100)-0.00125\,(L-100)^{2} .
        \]
        At $L=104$ (so $L-100=4$):
        \[
          T_{2}(104)=100+0.5\cdot 4-0.00125\cdot 16=102-0.02=101.98 .
        \]
        True value $101.980390$; the error is about
        $3.9\cdot 10^{-4}$, four extra labour hours raise output by
        just under $2$ units, the shortfall against the linear estimate
        $102$ reflecting diminishing returns ($Q''<0$).

  \item Taking logs, $\ln Y_{1}=\ln Y_{0}+\ln(1+g)$. For
        $f(x)=\ln(1+x)$ at $a=0$: $f(0)=0$ and $f'(0)=1$, so
        $T_{1}(x)=x$ and
        \[
          \ln Y_{1}-\ln Y_{0}=\ln(1+g)\approx g .
        \]
        For $g=0.03$ the log change is $\ln(1.03)=0.029559$, while the
        growth rate is $0.03$: the approximation is off by about
        $4.4\cdot 10^{-4}$, i.e.\ by roughly $1.5\%$ of the quantity
        itself. (The second-order term $-\tfrac12 g^{2}=-0.00045$
        accounts for almost all of the gap.)

  \item Here $n=1$, $a=4$, $x=4.1$ and $f''(t)=-\tfrac14 t^{-3/2}$.
        For $t\in[4,4.1]$, $t^{-3/2}$ is decreasing, so
        \[
          |f''(t)|\leq\tfrac14\cdot 4^{-3/2}=\tfrac14\cdot\tfrac18
                 =\tfrac{1}{32}=:M .
        \]
        The remainder bound gives
        \[
          |R_{1}(4.1)|\leq\frac{M}{2!}\,(0.1)^{2}
            =\frac{1}{32}\cdot\frac{0.01}{2}
            =\frac{0.01}{64}\approx 1.56\cdot 10^{-4} .
        \]
        Actual error: $T_{1}(4.1)=2.025$ versus $\sqrt{4.1}=2.024846$,
        i.e.\ $1.54\cdot 10^{-4}$, just inside the bound.

  \item Here $n=2$, $a=0$, $x=0.1$ and
        $f'''(t)=\mathrm{e}^{t}$. For $t\in(0,0.1)$,
        $\mathrm{e}^{t}<\mathrm{e}^{0.1}<1.2=:M$, so
        \[
          |R_{2}(0.1)|\leq\frac{M}{3!}\,(0.1)^{3}
            =\frac{1.2\cdot 0.001}{6}=2\cdot 10^{-4} .
        \]
        Actual error: $T_{2}(0.1)=1.105$ versus
        $\mathrm{e}^{0.1}=1.105171$, i.e.\ $1.7\cdot 10^{-4}$,
        consistent with the bound.

  \item With $f(x)=\mathrm{e}^{x}$, $a=0$, $x=0.5$:
        $f^{(n+1)}(\xi)=\mathrm{e}^{\xi}\leq\mathrm{e}^{0.5}<2$ for
        $\xi\in(0,0.5)$, so
        \[
          |R_{n}(0.5)|\leq\frac{2\,(0.5)^{n+1}}{(n+1)!} .
        \]
        Tabulate:
        \[
          n=3:\ \frac{2\cdot 0.0625}{24}\approx 5.2\cdot 10^{-3},
          \qquad
          n=4:\ \frac{2\cdot 0.03125}{120}\approx 5.2\cdot 10^{-4},
        \]
        \[
          n=5:\ \frac{2\cdot 0.015625}{720}\approx 4.3\cdot 10^{-5} .
        \]
        The first value below $10^{-4}$ occurs at $n=5$. Check with the
        true value $\mathrm{e}^{0.5}=1.648721$: $T_{4}(0.5)=1.648438$
        (error $2.8\cdot 10^{-4}$, indeed too large) while
        $T_{5}(0.5)=1.648698$ (error $2.3\cdot 10^{-5}$). So $n=5$.

  \item By Taylor's theorem with $n=1$: for each $x$ near $a$ there is
        $\xi$ between $a$ and $x$ with
        \[
          f(x)=f(a)+f'(a)(x-a)+\tfrac12 f''(\xi)(x-a)^{2}
              =f(a)+\tfrac12 f''(\xi)(x-a)^{2},
        \]
        using $f'(a)=0$. Since $f''$ is continuous and $f''(a)<0$,
        there is a neighbourhood $U$ of $a$ on which $f''<0$. For
        $x\in U$, $x\neq a$, the point $\xi$ also lies in $U$, so
        $f''(\xi)<0$ while $(x-a)^{2}>0$, giving
        \[
          f(x)-f(a)=\tfrac12 f''(\xi)(x-a)^{2}<0 .
        \]
        Hence $f(x)<f(a)$ for all $x\in U\setminus\{a\}$: $a$ is a
        strict local maximum. (This is the second-order condition in
        its cleanest form: near a stationary point the function
        \emph{is} its osculating parabola, up to higher-order terms.)
\end{solutions}

\section{Additional Problems}

\begin{problems}
  \item Find $T_{1}$ of $f(x)=(1+x)^{5}$ at $a=0$.
  \item Find $T_{1}$ of $f(x)=\sin(x)$ at $a=0$ and use it to
        approximate $\sin(0.1)$. Compare with the true value
        $\sin(0.1)=0.099833$.
  \item Find $T_{2}$ of $f(x)=x^{4}$ at $a=1$.
  \item Find $T_{2}$ of $f(x)=\dfrac{1}{x}$ at $a=2$.
  \item Find $T_{3}$ of $f(x)=\mathrm{e}^{x}$ at $a=0$ and use it to
        approximate $\mathrm{e}^{0.2}$. Compare with the true value
        $\mathrm{e}^{0.2}=1.221403$.
  \item Find $T_{2}$ of $f(x)=\ln(1+2x)$ at $a=0$.
  \item (Cobb--Douglas.) For $f(K)=K^{0.3}$, find $T_{2}$ at $a=1$ and
        use it to approximate $1.1^{0.3}$. Compare with the true value
        $1.1^{0.3}=1.029006$.
  \item Use the Lagrange remainder to bound the error of approximating
        $\ln(1.05)$ by $T_{1}(x)=x-1$ of $\ln(x)$ at $a=1$, and compare
        the bound with the actual error (true value
        $\ln(1.05)=0.048790$).
  \item Find the smallest $n$ such that the Maclaurin polynomial
        $T_{n}$ of $\ln(1+x)$ approximates $\ln(1.1)$ with error at
        most $10^{-6}$ (true value $\ln(1.1)=0.095310$).
  \item Suppose $|f'''(t)|\leq M$ for all $t$ between $a$ and $x$.
        Derive a bound for $|f(x)-T_{2}(x)|$ in terms of $M$ and
        $|x-a|$.
\end{problems}

\section{Additional Solutions}

\begin{solutions}
  \item $f(0)=1$; $f'(x)=5(1+x)^{4}$, $f'(0)=5$. Hence
        \[
          T_{1}(x)=1+5x .
        \]

  \item $f(0)=0$; $f'(x)=\cos(x)$, $f'(0)=1$. Hence $T_{1}(x)=x$ (the
        small-angle approximation) and $\sin(0.1)\approx 0.1$. True
        value $0.099833$; the error is about $1.7\cdot 10^{-4}$
        (matching the size of the first omitted term,
        $\tfrac{x^{3}}{3!}=\tfrac{0.001}{6}$).

  \item $f(1)=1$; $f'(x)=4x^{3}$, $f'(1)=4$; $f''(x)=12x^{2}$,
        $f''(1)=12$, so $\tfrac12 f''(1)=6$. Hence
        \[
          T_{2}(x)=1+4(x-1)+6(x-1)^{2} .
        \]

  \item $f(2)=\tfrac12$; $f'(x)=-1/x^{2}$, $f'(2)=-\tfrac14$;
        $f''(x)=2/x^{3}$, $f''(2)=\tfrac28=\tfrac14$, so
        $\tfrac12 f''(2)=\tfrac18$. Hence
        \[
          T_{2}(x)=\frac12-\frac14(x-2)+\frac18(x-2)^{2} .
        \]

  \item All derivatives equal $1$ at $0$, so
        \[
          T_{3}(x)=1+x+\frac{x^{2}}{2}+\frac{x^{3}}{6} .
        \]
        At $x=0.2$:
        \[
          T_{3}(0.2)=1+0.2+\frac{0.04}{2}+\frac{0.008}{6}
                    =1+0.2+0.02+0.001333=1.221333 .
        \]
        True value $1.221403$; the error is about $6.9\cdot 10^{-5}$.

  \item $f(0)=0$; $f'(x)=\dfrac{2}{1+2x}$, $f'(0)=2$;
        $f''(x)=-\dfrac{4}{(1+2x)^{2}}$, $f''(0)=-4$, so
        $\tfrac12 f''(0)=-2$. Hence
        \[
          T_{2}(x)=2x-2x^{2} .
        \]
        (Equivalently: substitute $u=2x$ into
        $\ln(1+u)\approx u-\tfrac12 u^{2}$.)

  \item $f(1)=1$; $f'(K)=0.3\,K^{-0.7}$, $f'(1)=0.3$;
        $f''(K)=0.3\cdot(-0.7)\,K^{-1.7}=-0.21\,K^{-1.7}$,
        $f''(1)=-0.21$, so $\tfrac12 f''(1)=-0.105$. Hence
        \[
          T_{2}(K)=1+0.3\,(K-1)-0.105\,(K-1)^{2} .
        \]
        At $K=1.1$:
        \[
          T_{2}(1.1)=1+0.03-0.105\cdot 0.01=1.03-0.00105=1.02895 .
        \]
        True value $1.029006$; the error is about $5.6\cdot 10^{-5}$.

  \item Here $n=1$, $a=1$, $x=1.05$ and $f''(t)=-1/t^{2}$. On
        $[1,1.05]$, $|f''(t)|=1/t^{2}\leq 1=:M$. The remainder bound
        gives
        \[
          |R_{1}(1.05)|\leq\frac{M}{2!}\,(0.05)^{2}
            =\frac{0.0025}{2}=1.25\cdot 10^{-3} .
        \]
        Actual error: $T_{1}(1.05)=0.05$ versus $0.048790$, i.e.\
        $1.21\cdot 10^{-3}$, close to the bound, because
        $1/t^{2}\approx 1$ throughout the interval.

  \item For $f(x)=\ln(1+x)$,
        $f^{(n+1)}(t)=\dfrac{(-1)^{n}\,n!}{(1+t)^{n+1}}$, so on
        $[0,0.1]$ we have $|f^{(n+1)}(t)|\leq n!$ and
        \[
          |R_{n}(0.1)|\leq\frac{n!}{(n+1)!}\,(0.1)^{n+1}
            =\frac{(0.1)^{n+1}}{n+1} .
        \]
        Tabulate:
        \[
          n=4:\ \frac{(0.1)^{5}}{5}=2\cdot 10^{-6},
          \qquad
          n=5:\ \frac{(0.1)^{6}}{6}\approx 1.7\cdot 10^{-7} .
        \]
        The first value below $10^{-6}$ occurs at $n=5$. Check against
        the true value $0.095310$: $T_{4}(0.1)=0.1-0.005+0.000333
        -0.000025=0.095308$ (error $1.9\cdot 10^{-6}$, indeed too
        large), while $T_{5}(0.1)=0.095308+0.000002=0.095310$ (error
        about $1.5\cdot 10^{-7}$). So $n=5$.

  \item Apply the Lagrange form with $n=2$: there is $\xi$ between $a$
        and $x$ with
        \[
          f(x)-T_{2}(x)=\frac{f'''(\xi)}{3!}(x-a)^{3},
        \]
        hence
        \[
          |f(x)-T_{2}(x)|\leq\frac{|f'''(\xi)|}{6}\,|x-a|^{3}
            \leq\frac{M}{6}\,|x-a|^{3} .
        \]
\end{solutions}

\chapter{Implicit Differentiation}

Most curves in the plane are not graphs. The circle $x^{2}+y^{2}=25$ fails
the vertical line test, the equation $\mathrm{e}^{y}+xy=1$ cannot be solved
for $y$ in closed form at all, and an indifference curve is handed to us as
a level set rather than as a formula. Yet each of them has a tangent at
almost every point, and economics asks for its slope constantly: that slope
\emph{is} the marginal rate of substitution.

Implicit differentiation gets it without solving for $y$. The move rests on
one idea, and the whole chapter is that idea applied repeatedly: pretend the
equation has already been solved near the point of interest, so that $y$ is
a function $y(x)$; substituting it makes the equation an identity in $x$;
identities may be differentiated. Because every $y$ in sight is secretly the
composite $y(x)$, the chain rule leaves a factor $y'$ behind wherever a $y$
was, and the resulting linear equation is solved for $y'$.

Two questions have to be kept apart, and the chapter answers both. \emph{Is}
there such a function $y(x)$ near the point --- and if so, what is its
slope? The first question is the Implicit Function Theorem, and the honest
answer is sometimes no: at a point where the curve crosses itself or stands
vertically, no function $y(x)$ describes it, and the slope formula returns a
division by zero to say so.

\section{Core Rules}

\begin{rulebox}{Why a factor $y'$ appears}
Suppose that near a point of the curve the equation really does pin down one
$y$ for each $x$, and call that function $y(x)$. Substituting it turns the
equation into an \textbf{identity in $x$}, true for every $x$ near $x_{0}$:
\[
  x^{2}+\bigl(y(x)\bigr)^{2}=25 .
\]
An identity may be differentiated on both sides. Every $y$ in the expression
is the composite $y(x)$, so it is differentiated by the \textbf{chain rule}
and leaves a factor $y'$ behind:
\[
  \frac{\mathrm{d}}{\mathrm{d}x}\bigl[y^{2}\bigr]
  =\frac{\mathrm{d}}{\mathrm{d}x}\bigl[(y(x))^{2}\bigr]
  =2\,y(x)\,y'(x)=2yy' .
\]
Hence $2x+2yy'=0$ and $y'=-x/y$. Nothing was solved for $y$, and the answer
is \emph{allowed} to contain $y$: it is the slope at a point of the curve,
and the two points with the same $x$ --- here $(3,4)$ and $(3,-4)$ --- have
slopes $-3/4$ and $+3/4$.
\end{rulebox}

\begin{rulebox}{Setup}
Suppose $F(x,y)=0$ defines $y=y(x)$ implicitly. Differentiate every term
with respect to $x$. Whenever an expression involves $y$, the chain rule
produces an extra factor $\dd{y}{x}$:
\[
  \frac{\mathrm{d}}{\mathrm{d}x}\bigl[g(y)\bigr]=g'(y)\cdot\dd{y}{x} .
\]
So $\frac{\mathrm{d}}{\mathrm{d}x}[\sin y]=\cos(y)\,y'$,
$\frac{\mathrm{d}}{\mathrm{d}x}[\mathrm{e}^{y}]=\mathrm{e}^{y}y'$, and
$\frac{\mathrm{d}}{\mathrm{d}x}[\ln y]=y'/y$.
\end{rulebox}

\begin{rulebox}{Mixed terms}
For a product such as $x^{3}y^{2}$ apply the product rule, treating $y$ as
a function of $x$:
\[
  \frac{\mathrm{d}}{\mathrm{d}x}\bigl[x^{3}y^{2}\bigr]=3x^{2}y^{2}+x^{3}\cdot 2yy' .
\]
\end{rulebox}

\begin{rulebox}{Solving for $y'$}
Collect all terms containing $y'$ on one side and factor:
\[
  A(x,y)\,y'+B(x,y)=0\;\Longrightarrow\;y'=-\frac{B(x,y)}{A(x,y)} .
\]
The equation is always \emph{linear} in $y'$, however complicated it looks,
because $y'$ only ever enters through one chain-rule step.
\end{rulebox}

\begin{rulebox}{The curve is a level set --- and that gives the formula at once}
Write the equation as $F(x,y)=0$. Then the curve is the \textbf{level set of
the surface $z=F(x,y)$ at height $0$}: the shoreline where the surface cuts
the horizontal plane $z=0$. Walking along the curve, $F$ never changes.
Writing that down with the chain rule for $F\bigl(x,y(x)\bigr)=0$,
\[
  \pd{F}{x}+\pd{F}{y}\cdot\dd{y}{x}=0
  \qquad\Longrightarrow\qquad
  \boxed{\;\dd{y}{x}=-\frac{F_{x}(x,y)}{F_{y}(x,y)}\;}
  \qquad\text{provided } F_{y}(x,y)\neq0 .
\]
This is not a new method, it is the same computation done once and for all:
for $F=x^{2}+y^{2}-25$ it returns $-2x/2y=-x/y$, as above. In economics the
curve is an indifference curve $U(x,y)=\bar u$ and the formula reads
$\dd{y}{x}=-U_{x}/U_{y}$: the marginal rate of substitution is the ratio of
marginal utilities.
\end{rulebox}

\begin{gallery}
\begin{tikzpicture}
  \begin{axis}[surfaxis, title={$z=F(x,y)=x^{3}-3x^{2}+y^{2}-2y+1$}]
    \addplot3[surf, domain=-1.2:3.2, y domain=-1.6:3.6]{x^3-3*x^2+y^2-2*y+1};
  \end{axis}
\end{tikzpicture}\hfill
\begin{tikzpicture}
  \begin{axis}[geoaxis, xmin=-1.6, xmax=3.6, ymin=-1.9, ymax=3.9,
      xtick={-1,1,2,3}, ytick={-1,1,2,3}, title={the level set $F=0$}]
    \addplot[domain=-1.2:3, samples=241]{1+abs(x)*sqrt(3-x)};
    \addplot[domain=-1.2:3, samples=241]{1-abs(x)*sqrt(3-x)};
    \addplot[only marks, mark=*, mark size=1.3pt, accent, forget plot]
      coordinates {(0,1) (3,1)};
  \end{axis}
\end{tikzpicture}
\gcap{The surface and its zero level set $k$. At $(0,1)$ the curve crosses
itself; at $(3,1)$ it stands vertical. Everywhere else a local $y(x)$
exists.}
\end{gallery}

\begin{rulebox}{When there is no $y(x)$: the Implicit Function Theorem}
\textbf{Theorem.} If $F$ is $C^{1}$ near $(x_{0},y_{0})$, if
$F(x_{0},y_{0})=0$, and if
\[
  F_{y}(x_{0},y_{0})\neq0 ,
\]
then there are neighbourhoods on which the equation $F(x,y)=0$ defines
exactly one differentiable function $y=y(x)$ with $y(x_{0})=y_{0}$, and its
derivative is $-F_{x}/F_{y}$.

The condition is exactly what the formula needs, and its failure is
visible in the picture. Where $F_{y}=0$:
\begin{itemize}[leftmargin=1.6em,itemsep=0.3ex]
  \item if $F_{x}\neq0$, the curve has a \textbf{vertical tangent}. No
        function $y(x)$ can describe it there --- but the roles may be
        swapped: $x=x(y)$ exists, with
        $\dd{x}{y}=-F_{y}/F_{x}$, and that derivative is $0$.
  \item if $F_{x}=0$ as well, the point is \textbf{singular}: typically the
        curve crosses itself, and \emph{two} branches with two different
        slopes pass through. No single function describes either the curve
        or its slope, and $-F_{x}/F_{y}$ degenerates to $0/0$. Separate the
        branches first, then differentiate each one.
\end{itemize}
Reading it off a picture: look for vertical tangents and self-intersections.
Those are the only places the theorem can fail. A failed condition is a
warning, not a proof that no $y(x)$ exists: $y^{3}-x=0$ has $F_{y}=0$ at the
origin, yet $y=x^{1/3}$ solves it for every $x$ (with a vertical tangent, so
not differentiably). Chapter~14 collects such cases.
\end{rulebox}

\begin{rulebox}{Logarithmic differentiation}
If $y=f(x)>0$ is a complicated product, quotient or power, take $\ln$
first:
\[
  \ln y=\ln f(x)
  \;\Longrightarrow\;
  \frac{y'}{y}=\frac{\mathrm{d}}{\mathrm{d}x}\ln f(x)
  \;\Longrightarrow\;
  y'=y\cdot\frac{\mathrm{d}}{\mathrm{d}x}\ln f(x) .
\]
Particularly useful for $y=u(x)^{v(x)}$: $\ln y=v(x)\ln u(x)$. This is
implicit differentiation applied to $\ln y-\ln f(x)=0$, which is why the
factor $y'/y$ appears.
\end{rulebox}

\begin{rulebox}{Higher-order implicit}
To find $y''$, differentiate $y'$ a second time with respect to $x$, again
remembering that any $y$ that appears contributes a factor $y'$. Then
substitute the expression already found for $y'$ to eliminate it.
\end{rulebox}

\begin{rulebox}{Tangent line via implicit slope}
The tangent at $(x_{0},y_{0})$ on $F(x,y)=0$ is
$y-y_{0}=y'(x_{0})\cdot(x-x_{0})$, where $y'(x_{0})$ is computed
implicitly. Both coordinates are substituted into $-F_{x}/F_{y}$ before the
line is written down.
\end{rulebox}

\section{Worked Example}

\begin{worked}{The folium $x^{3}+y^{3}=6xy$: find $y'$, the tangent at $(3,3)$, and every point where the tangent is vertical}
\wstep{Step 1. Confirm the point is on the curve.}
$27+27=54$ and $6\cdot3\cdot3=54$. \checkmark\ Do this first: an implicit
slope formula evaluated at a point that is not on the curve returns a
perfectly plausible number that means nothing.

\wstep{Step 2. Differentiate every term with respect to $x$, treating $y$ as $y(x)$.}
Left side: $\frac{\mathrm{d}}{\mathrm{d}x}\bigl[x^{3}+y^{3}\bigr]
=3x^{2}+3y^{2}y'$ --- the chain rule leaves a factor $y'$ behind wherever a
$y$ was.

Right side: $6xy$ is a \textbf{product} of two functions of $x$, so the
product rule applies:
\[
  \frac{\mathrm{d}}{\mathrm{d}x}\bigl[6xy\bigr]=6y+6xy' .
\]
\wtrap{The right side is where this problem is lost. Writing $6y$ alone
(treating $y$ as a constant) or $6xy'$ alone (forgetting the product rule)
are both common. Every occurrence of $y$ needs the chain rule, \emph{and}
every product needs the product rule --- they compound rather than replace
each other.}

\wstep{Step 3. Collect the $y'$ terms; the equation is always linear in $y'$.}
\[
  3x^{2}+3y^{2}y'=6y+6xy'
  \;\Longrightarrow\;
  3y^{2}y'-6xy'=6y-3x^{2}
  \;\Longrightarrow\;
  y'\bigl(3y^{2}-6x\bigr)=6y-3x^{2} .
\]
\[
  \boxed{\;y'=\frac{6y-3x^{2}}{3y^{2}-6x}=\frac{2y-x^{2}}{y^{2}-2x}\;}
\]
The answer contains $y$, and it must: two points of the curve can share an
$x$ and have different slopes, so $x$ alone cannot determine $y'$.

\wstep{Step 4. Evaluate at the point.}
\[
  y'(3,3)=\frac{2\cdot3-9}{9-6}=\frac{-3}{3}=-1 ,
\]
so the tangent is
\[
  y-3=-1(x-3)\quad\Longrightarrow\quad y=6-x .
\]
By the symmetry of the equation in $x$ and $y$, the curve is symmetric about
the line $y=x$, and $(3,3)$ lies on that line --- so a slope of $-1$,
perpendicular to it, is exactly what the geometry predicts. \checkmark

\wstep{Step 5. Cross-check with the level-set formula.}
Writing $F(x,y)=x^{3}+y^{3}-6xy$,
\[
  F_{x}=3x^{2}-6y,\qquad F_{y}=3y^{2}-6x,
  \qquad
  y'=-\frac{F_{x}}{F_{y}}=-\frac{3x^{2}-6y}{3y^{2}-6x}=\frac{6y-3x^{2}}{3y^{2}-6x},
\]
the same expression. The two routes --- differentiate term by term, or
apply $-F_{x}/F_{y}$ --- are the same computation; use whichever you find
faster, but know that they must agree.

\wstep{Step 6. Vertical tangents: where the denominator vanishes.}
$y'$ blows up when $F_{y}=3y^{2}-6x=0$, i.e.\ $x=y^{2}/2$. Substituting into
the curve:
\[
  \frac{y^{6}}{8}+y^{3}=6\cdot\frac{y^{2}}{2}\cdot y=3y^{3}
  \;\Longrightarrow\;
  \frac{y^{6}}{8}=2y^{3}
  \;\Longrightarrow\;
  y^{3}\bigl(y^{3}-16\bigr)=0 .
\]
So $y=0$ (giving $x=0$) or $y=16^{1/3}=2\sqrt[3]{2}\approx2.520$, giving
$x=y^{2}/2\approx3.175$. At the latter the tangent is genuinely vertical.
\wtrap{At the origin the numerator $2y-x^{2}$ also vanishes, so $y'$ is
$0/0$ there --- undefined for a different reason. The folium
\emph{crosses itself} at $(0,0)$, with two distinct tangents, so no single
function $y(x)$ describes the curve near it. That is precisely the case the
Implicit Function Theorem excludes, and a $0/0$ in the slope formula is how
it announces itself.}

\wstep{Step 7. The economic version of the same calculation.}
An indifference curve is a level set $U(x,y)=\bar u$, and the identical
manipulation gives
\[
  \frac{\mathrm{d}y}{\mathrm{d}x}=-\frac{U_{x}}{U_{y}}=-\mathrm{MRS} .
\]
For $U=x^{1/2}y^{1/2}=4$: $U_{x}=\tfrac12x^{-1/2}y^{1/2}$,
$U_{y}=\tfrac12x^{1/2}y^{-1/2}$, so $y'=-y/x$. At $(4,4)$ the slope is
$-1$; at $(8,2)$ it is $-\tfrac14$ --- the indifference curve flattens as
$x$ becomes abundant, which is diminishing marginal rate of substitution
stated as a derivative.
\end{worked}

\section{Practice Problems}

\begin{problems}
  \item Compute $\dd{y}{x}$ implicitly for $x^{2}+4y^{2}=16$.
  \item Find the equation of the tangent line to $x^{2}+4y^{2}=20$ at
        $(2,2)$.
  \item Use logarithmic differentiation to find $y'$ for $y=x^{\sin x}$
        ($x>0$, $\sin x>0$).
  \item Compute $\dd{^{2}y}{x^{2}}$ implicitly for $x^{2}+4y^{2}=4$.
  \item Compute $\dd{y}{x}$ implicitly for $x^{3}-y^{3}=3xy$.
  \item Compute $\dd{y}{x}$ implicitly for $x^{2}y^{3}=8$.
  \item Find the equation of the tangent line to $x^{2}y+y^{2}x=6$ at
        $(1,2)$.
  \item Use logarithmic differentiation to find $y'$ for $y=(\ln x)^{x}$
        ($x>1$).
  \item Compute $\dd{^{2}y}{x^{2}}$ implicitly for $x^{2}y=1$.
  \item Compute $\dd{y}{x}$ implicitly for $\cos(y)+x=y$.
  \item Find the equation of the tangent line to $y^{2}-x^{2}=3$ at
        $(1,2)$.
  \item Use logarithmic differentiation to find $y'$ for
        \[
          y=\frac{(2x+1)^{5}\sqrt{x^{2}+3}}{(x-4)^{3}}\qquad (x>4).
        \]
  \item Compute $\dd{y}{x}$ implicitly for $\mathrm{e}^{x+y}=xy$.
  \item Compute $\dd{^{2}y}{x^{2}}$ implicitly for $y^{3}=x^{2}$.
  \item Compute $\dd{y}{x}$ implicitly for $x^{2}+2xy+3y^{2}=12$.
  \item Find the equation of the tangent line to $\mathrm{e}^{y}+xy=\mathrm{e}+1$
        at $(1,1)$.
  \item Use logarithmic differentiation to find $y'$ for $y=(x^{2}+1)^{x}$
        ($x>0$).
  \item Compute $\dd{y}{x}$ implicitly for $\sqrt{x+y}=x-y$.
  \item Compute $\dd{^{2}y}{x^{2}}$ implicitly for $xy+y^{2}=1$.
  \item Find the equation of the tangent line to $x^{2}+y^{2}=4y$ at
        $(\sqrt{3},1)$.
  \item Compute $\dd{y}{x}$ implicitly for $\tan(xy)=x$.
  \item Use logarithmic differentiation to find $y'$ for $y=x^{1/x}$
        ($x>0$).
  \item Compute $\dd{y}{x}$ implicitly for $\ln(y)+xy=2$.
  \item Compute $\dd{y}{x}$ implicitly for $y^{3}+y=x^{2}$.
  \item Find the tangent line to the astroid $x^{2/3}+y^{2/3}=2$ at
        $(1,1)$.
  \item Compute $\dd{y}{x}$ implicitly for $\sin(x+y)=y$.
  \item Use logarithmic differentiation to find $y'$ for
        $y=\dfrac{\mathrm{e}^{2x}(x+1)^{3}}{\sqrt{x^{2}+5}}$ (assume $x>-1$).
  \item Compute $\dd{^{2}y}{x^{2}}$ implicitly for $x^{2}+y^{2}=25$.
  \item For an indifference curve given by $U(x,y)=x^{1/2}y^{1/2}=\bar U$
        (constant), compute $\dd{y}{x}$, the marginal rate of
        substitution.
  \item Find all points on the curve $x^{2}+y^{2}=2x+4y$ where the tangent
        line is horizontal.
  \item Compute $\dd{y}{x}$ implicitly for $\arctan(y/x)=x$.
  \item Use logarithmic differentiation to find $y'$ for $y=x^{x}$
        ($x>0$).
  \item Find the equation of the tangent line to $x^{3}+y^{3}=9xy$ at
        $(2,4)$ (point on the folium of Descartes).
  \item Compute $\dd{^{2}y}{x^{2}}$ implicitly for $y^{2}=4x$.
  \item Compute $\dd{y}{x}$ implicitly for $x\mathrm{e}^{y}+y=x$.
  \item Find all points on the ellipse $\dfrac{x^{2}}{9}+\dfrac{y^{2}}{4}=1$
        where the tangent line has slope $-1$.
  \item Use logarithmic differentiation to find $y'$ for
        $y=(\sin x)^{\cos x}$ on $(0,\pi/2)$.
  \item Compute $\dd{y}{x}$ implicitly for $x^{2}y-xy^{2}+x^{2}+y^{2}=0$.
  \item For the production isoquant $K^{1/3}L^{2/3}=\bar Q$ (constant),
        compute $\dd{K}{L}$, the marginal rate of technical
        substitution.
  \item Compute $\dd{^{2}y}{x^{2}}$ implicitly for $\sin(y)=x$.
  \item Find the tangent line to $xy=4$ at $x=2$.
  \item Use logarithmic differentiation to find $y'$ for
        $y=\dfrac{\sqrt{x+1}\sqrt[3]{x-2}}{(x+5)^{4}}$ (assume $x>2$).
  \item Compute $\dd{y}{x}$ implicitly for $\mathrm{e}^{y}\ln(x)+x^{2}=y$.
  \item Show that on the circle $x^{2}+y^{2}=r^{2}$ the tangent at
        $(x_{0},y_{0})$ is perpendicular to the radius. (I.e.\ derive the
        tangent slope and compare with the radius slope.)
  \item Use logarithmic differentiation to find $y'$ for $y=7^{x^{2}}$
        ($y>0$).
  \item Compute $\dd{^{2}y}{x^{2}}$ implicitly for $x^{2}+xy+y^{2}=3$.
  \item For the budget surface $px+qy=m$ (a line, $p,q,m>0$ constants),
        differentiate implicitly to recover $\dd{y}{x}=-p/q$.
  \item Find the tangent line to $x^{2}-xy+y^{2}=7$ at $(3,2)$.
  \item Compute $\dd{y}{x}$ implicitly for $\cos(xy)=x+y$.
  \item Compute $\dd{^{2}y}{x^{2}}$ implicitly for $x^{3}+y^{3}=6xy$ at the
        point $(3,3)$ (folium of Descartes). Express the answer as a
        number.
\end{problems}

\section{Solutions}

\begin{solutions}
  \item $2x+8yy'=0\Rightarrow y'=-\dfrac{x}{4y}$.

  \item Check $(2,2)$: $4+16=20$. From $2x+8yy'=0$: $y'=-x/(4y)$. At
        $(2,2)$: $y'=-2/8=-1/4$. Tangent: $y-2=-\tfrac14(x-2)$, i.e.\
        $x+4y=10$.

  \item $\ln y=\sin(x)\ln(x)\Rightarrow\dfrac{y'}{y}=\cos(x)\ln(x)+\dfrac{\sin x}{x}$.
        So
        \[
          y'=x^{\sin x}\left[\cos(x)\ln(x)+\frac{\sin x}{x}\right] .
        \]

  \item $2x+8yy'=0\Rightarrow y'=-\dfrac{x}{4y}$. Differentiate again:
        \[
          y''=-\frac{4y\cdot 1-x\cdot 4y'}{(4y)^{2}}=-\frac{y-xy'}{4y^{2}} .
        \]
        Substitute $y'=-x/(4y)$:
        $y-xy'=y+\dfrac{x^{2}}{4y}=\dfrac{4y^{2}+x^{2}}{4y}=\dfrac{4}{4y}
        =\dfrac{1}{y}$ (using $x^{2}+4y^{2}=4$). Hence
        $y''=-\dfrac{1}{4y^{3}}$.

  \item $3x^{2}-3y^{2}y'=3y+3xy'\Rightarrow -3(y^{2}+x)y'=3(y-x^{2})
        \Rightarrow y'=\dfrac{x^{2}-y}{y^{2}+x}$.

  \item $2xy^{3}+3x^{2}y^{2}y'=0\Rightarrow
        y'=-\dfrac{2xy^{3}}{3x^{2}y^{2}}=-\dfrac{2y}{3x}$.

  \item Differentiate $x^{2}y+y^{2}x=6$: $2xy+x^{2}y'+2yxy'+y^{2}=0
        \Rightarrow y'(x^{2}+2xy)=-(2xy+y^{2})
        \Rightarrow y'=-\dfrac{y(2x+y)}{x(x+2y)}$. At $(1,2)$:
        $y'=-\dfrac{2\cdot 4}{1\cdot 5}=-\dfrac{8}{5}$. Tangent:
        $y-2=-\tfrac85(x-1)$, i.e.\ $8x+5y=18$.

  \item $\ln y=x\ln(\ln x)\Rightarrow\dfrac{y'}{y}=\ln(\ln x)+x\cdot
        \dfrac{1/x}{\ln x}=\ln(\ln x)+\dfrac{1}{\ln x}$. So
        $y'=(\ln x)^{x}\left[\ln(\ln x)+\dfrac{1}{\ln x}\right]$.

  \item Explicit: $y=x^{-2}$, $y'=-2x^{-3}$, $y''=6x^{-4}$. Implicit
        check: $2xy+x^{2}y'=0\Rightarrow y'=-2y/x$. Differentiate:
        \[
          y''=-\frac{2y'x-2y}{x^{2}}=-\frac{2(-2y/x)x-2y}{x^{2}}
             =-\frac{-4y-2y}{x^{2}}=\frac{6y}{x^{2}}=\frac{6}{x^{4}} .
        \]

  \item $-\sin(y)y'+1=y'\Rightarrow y'(1+\sin y)=1\Rightarrow
        y'=\dfrac{1}{1+\sin y}$.

  \item $2yy'-2x=0\Rightarrow y'=x/y$. At $(1,2)$: $y'=1/2$. Tangent:
        $y-2=\tfrac12(x-1)$, i.e.\ $x-2y+3=0$.

  \item $\ln y=5\ln(2x+1)+\tfrac12\ln(x^{2}+3)-3\ln(x-4)$. Then
        \[
          \frac{y'}{y}=\frac{10}{2x+1}+\frac{x}{x^{2}+3}-\frac{3}{x-4} .
        \]
        Hence
        \[
          y'=\frac{(2x+1)^{5}\sqrt{x^{2}+3}}{(x-4)^{3}}
             \left[\frac{10}{2x+1}+\frac{x}{x^{2}+3}-\frac{3}{x-4}\right] .
        \]

  \item $\mathrm{e}^{x+y}(1+y')=y+xy'\Rightarrow
        y'(\mathrm{e}^{x+y}-x)=y-\mathrm{e}^{x+y}\Rightarrow
        y'=\dfrac{y-\mathrm{e}^{x+y}}{\mathrm{e}^{x+y}-x}$.

  \item $3y^{2}y'=2x\Rightarrow y'=\dfrac{2x}{3y^{2}}$. Differentiate
        again:
        \[
          y''=\frac{2\cdot 3y^{2}-2x\cdot 6yy'}{9y^{4}}
             =\frac{6y^{2}-12xy\cdot y'}{9y^{4}}
             =\frac{2}{3y^{2}}-\frac{4xy'}{3y^{3}} .
        \]
        Substitute $y'=2x/(3y^{2})$:
        \[
          y''=\frac{2}{3y^{2}}-\frac{4x}{3y^{3}}\cdot\frac{2x}{3y^{2}}
             =\frac{2}{3y^{2}}-\frac{8x^{2}}{9y^{5}}
             =\frac{6y^{3}-8x^{2}}{9y^{5}}
             =\frac{6x^{2}-8x^{2}}{9y^{5}}=-\frac{2x^{2}}{9y^{5}},
        \]
        using $y^{3}=x^{2}$. Equivalently, $y''=-\dfrac{2y^{3}}{9y^{5}}
        =-\dfrac{2}{9y^{2}}$.

  \item $2x+2y+2xy'+6yy'=0\Rightarrow y'(2x+6y)=-(2x+2y)\Rightarrow
        y'=-\dfrac{x+y}{x+3y}$.

  \item Differentiate $\mathrm{e}^{y}+xy=\mathrm{e}+1$:
        $\mathrm{e}^{y}y'+y+xy'=0\Rightarrow y'(\mathrm{e}^{y}+x)=-y
        \Rightarrow y'=-\dfrac{y}{\mathrm{e}^{y}+x}$. At $(1,1)$:
        $y'=-\dfrac{1}{\mathrm{e}+1}$. Tangent:
        $y-1=-\dfrac{1}{\mathrm{e}+1}(x-1)$.

  \item $\ln y=x\ln(x^{2}+1)\Rightarrow\dfrac{y'}{y}=\ln(x^{2}+1)
        +x\cdot\dfrac{2x}{x^{2}+1}=\ln(x^{2}+1)+\dfrac{2x^{2}}{x^{2}+1}$.
        So
        \[
          y'=(x^{2}+1)^{x}\left[\ln(x^{2}+1)+\frac{2x^{2}}{x^{2}+1}\right] .
        \]

  \item $\dfrac{1+y'}{2\sqrt{x+y}}=1-y'$. Multiply by $2\sqrt{x+y}$:
        $1+y'=2\sqrt{x+y}(1-y')\Rightarrow y'(1+2\sqrt{x+y})
        =2\sqrt{x+y}-1\Rightarrow y'=\dfrac{2\sqrt{x+y}-1}{1+2\sqrt{x+y}}$.

  \item $y+xy'+2yy'=0\Rightarrow y'(x+2y)=-y\Rightarrow
        y'=-\dfrac{y}{x+2y}$. Differentiate again using the quotient
        rule:
        \[
          y''=-\frac{y'(x+2y)-y(1+2y')}{(x+2y)^{2}}
             =-\frac{y'(x+2y)-y-2yy'}{(x+2y)^{2}}
             =-\frac{xy'-y}{(x+2y)^{2}} .
        \]
        Substitute $y'=-y/(x+2y)$:
        \[
          xy'-y=-\frac{xy}{x+2y}-y=-\frac{xy+y(x+2y)}{x+2y}
               =-\frac{2xy+2y^{2}}{x+2y}=-\frac{2y(x+y)}{x+2y} .
        \]
        Hence
        \[
          y''=\frac{2y(x+y)}{(x+2y)^{3}} .
        \]

  \item $2x+2yy'=4y'\Rightarrow y'(2y-4)=-2x\Rightarrow
        y'=\dfrac{x}{2-y}$. At $(\sqrt{3},1)$: $y'=\dfrac{\sqrt{3}}{1}
        =\sqrt{3}$. Tangent: $y-1=\sqrt{3}(x-\sqrt{3})$, i.e.\
        $y=\sqrt{3}x-2$.

  \item $\sec^{2}(xy)\cdot(y+xy')=1\Rightarrow y+xy'=\cos^{2}(xy)
        \Rightarrow y'=\dfrac{\cos^{2}(xy)-y}{x}$.

  \item $\ln y=\dfrac{\ln x}{x}\Rightarrow\dfrac{y'}{y}
        =\dfrac{(1/x)\cdot x-\ln x\cdot 1}{x^{2}}=\dfrac{1-\ln x}{x^{2}}$.
        So $y'=x^{1/x}\cdot\dfrac{1-\ln x}{x^{2}}$.

  \item $\dfrac{y'}{y}+y+xy'=0\Rightarrow y'\left(\dfrac1y+x\right)=-y
        \Rightarrow y'=-\dfrac{y^{2}}{1+xy}$.

  \item $3y^{2}y'+y'=2x\Rightarrow y'=\dfrac{2x}{3y^{2}+1}$.

  \item Differentiate $x^{2/3}+y^{2/3}=2$: $\tfrac23x^{-1/3}
        +\tfrac23y^{-1/3}y'=0\Rightarrow y'=-\dfrac{y^{1/3}}{x^{1/3}}
        =-\sqrt[3]{y/x}$. At $(1,1)$: $y'=-1$. Tangent: $y-1=-(x-1)$,
        i.e.\ $y=-x+2$.

  \item $\cos(x+y)(1+y')=y'\Rightarrow\cos(x+y)+\cos(x+y)y'=y'
        \Rightarrow\cos(x+y)=y'[1-\cos(x+y)]\Rightarrow
        y'=\dfrac{\cos(x+y)}{1-\cos(x+y)}$.

  \item $\ln y=2x+3\ln(x+1)-\tfrac12\ln(x^{2}+5)$. Differentiate:
        \[
          \frac{y'}{y}=2+\frac{3}{x+1}-\frac{x}{x^{2}+5} .
        \]
        Hence
        \[
          y'=\frac{\mathrm{e}^{2x}(x+1)^{3}}{\sqrt{x^{2}+5}}
             \left[2+\frac{3}{x+1}-\frac{x}{x^{2}+5}\right] .
        \]

  \item $2x+2yy'=0\Rightarrow y'=-x/y$. Differentiate again:
        \[
          y''=-\frac{y\cdot 1-x\cdot y'}{y^{2}}
             =-\frac{y-x(-x/y)}{y^{2}}
             =-\frac{y+x^{2}/y}{y^{2}}
             =-\frac{y^{2}+x^{2}}{y^{3}}=-\frac{25}{y^{3}},
        \]
        using $x^{2}+y^{2}=25$.

  \item Differentiate $x^{1/2}y^{1/2}=\bar U$:
        $\tfrac12x^{-1/2}y^{1/2}+x^{1/2}\cdot\tfrac12y^{-1/2}y'=0$.
        Multiply by $2x^{1/2}y^{1/2}$: $y+xy'=0\Rightarrow y'=-y/x$. So
        $\dd{y}{x}=-\dfrac{y}{x}$ (negative slope along the indifference
        curve, standard MRS).

  \item $2x+2yy'=2+4y'\Rightarrow y'(2y-4)=2-2x\Rightarrow
        y'=\dfrac{1-x}{y-2}$. Horizontal tangent: $y'=0$ requires $x=1$
        (and $y\neq 2$). Plug into the curve: $1+y^{2}=2+4y\Rightarrow
        y^{2}-4y-1=0\Rightarrow y=2\pm\sqrt{5}$. So horizontal tangents at
        $(1,2+\sqrt{5})$ and $(1,2-\sqrt{5})$.

  \item Differentiate $\arctan(y/x)=x$:
        $\dfrac{1}{1+(y/x)^{2}}\cdot\dfrac{y'x-y}{x^{2}}=1$, i.e.\
        $\dfrac{y'x-y}{x^{2}+y^{2}}=1$. Hence $y'x-y=x^{2}+y^{2}$, so
        $y'=\dfrac{x^{2}+y^{2}+y}{x}$.

  \item $\ln y=x\ln x\Rightarrow\dfrac{y'}{y}=\ln x+1\Rightarrow
        y'=x^{x}(\ln x+1)$.

  \item Check $(2,4)$ on $x^{3}+y^{3}=9xy$: $8+64=72$, $9\cdot 8=72$.
        Good. Differentiate: $3x^{2}+3y^{2}y'=9y+9xy'\Rightarrow
        y'(3y^{2}-9x)=9y-3x^{2}\Rightarrow
        y'=\dfrac{9y-3x^{2}}{3y^{2}-9x}=\dfrac{3y-x^{2}}{y^{2}-3x}$.
        At $(2,4)$: $y'=\dfrac{12-4}{16-6}=\dfrac{8}{10}=\dfrac{4}{5}$.
        Tangent: $y-4=\tfrac45(x-2)$, i.e.\ $4x-5y+12=0$.

  \item $2yy'=4\Rightarrow y'=2/y$. Differentiate again:
        $y''=-\dfrac{2y'}{y^{2}}=-\dfrac{2(2/y)}{y^{2}}=-\dfrac{4}{y^{3}}$.

  \item $\mathrm{e}^{y}+x\mathrm{e}^{y}y'+y'=1\Rightarrow
        y'(x\mathrm{e}^{y}+1)=1-\mathrm{e}^{y}\Rightarrow
        y'=\dfrac{1-\mathrm{e}^{y}}{1+x\mathrm{e}^{y}}$.

  \item $\dfrac{2x}{9}+\dfrac{2y}{4}y'=0\Rightarrow y'=-\dfrac{4x}{9y}$.
        Set $-\dfrac{4x}{9y}=-1\Rightarrow 4x=9y\Rightarrow y=\dfrac{4x}{9}$.
        Plug into the ellipse: $\dfrac{x^{2}}{9}+\dfrac{(4x/9)^{2}}{4}=1
        \Rightarrow\dfrac{x^{2}}{9}+\dfrac{16x^{2}}{4\cdot 81}=1
        \Rightarrow\dfrac{x^{2}}{9}+\dfrac{4x^{2}}{81}=1
        \Rightarrow\dfrac{9x^{2}+4x^{2}}{81}=1\Rightarrow
        13x^{2}=81\Rightarrow x=\pm 9/\sqrt{13}$. Corresponding
        $y=\pm 4/\sqrt{13}$ (same sign as $x$). Points:
        $\left(\dfrac{9}{\sqrt{13}},\dfrac{4}{\sqrt{13}}\right)$ and
        $\left(-\dfrac{9}{\sqrt{13}},-\dfrac{4}{\sqrt{13}}\right)$.

  \item $\ln y=\cos(x)\ln(\sin x)\Rightarrow\dfrac{y'}{y}
        =-\sin(x)\ln(\sin x)+\cos(x)\cdot\dfrac{\cos x}{\sin x}
        =-\sin(x)\ln(\sin x)+\dfrac{\cos^{2}x}{\sin x}$. So
        \[
          y'=(\sin x)^{\cos x}\left[-\sin(x)\ln(\sin x)
             +\frac{\cos^{2}x}{\sin x}\right] .
        \]

  \item $2xy+x^{2}y'-y^{2}-x\cdot 2yy'+2x+2yy'=0$. Collect $y'$:
        $y'(x^{2}-2xy+2y)=y^{2}-2xy-2x$. So
        $y'=\dfrac{y^{2}-2xy-2x}{x^{2}-2xy+2y}$.

  \item $\tfrac13K^{-2/3}L^{2/3}\dd{K}{L}\cdot L^{2/3}\ldots$
        carefully: differentiate $K^{1/3}L^{2/3}=\bar Q$ with respect to
        $L$, treating $K=K(L)$:
        \[
          \tfrac13K^{-2/3}K'(L)\cdot L^{2/3}+K^{1/3}\cdot\tfrac23L^{-1/3}=0 .
        \]
        Multiply by $3K^{2/3}L^{1/3}$: $LK'(L)+2K=0\Rightarrow
        K'(L)=-\dfrac{2K}{L}$. So $\dd{K}{L}=-\dfrac{2K}{L}$ (the MRTS).

  \item $\cos(y)y'=1\Rightarrow y'=\dfrac{1}{\cos y}=\sec(y)$.
        Differentiate again: $y''=\sec(y)\tan(y)\cdot y'
        =\sec^{2}(y)\tan(y)$.

  \item At $x=2$, $y=2$. Differentiate $xy=4$: $y+xy'=0\Rightarrow
        y'=-y/x$. At $(2,2)$: $y'=-1$. Tangent: $y-2=-(x-2)$, i.e.\
        $y=-x+4$.

  \item $\ln y=\tfrac12\ln(x+1)+\tfrac13\ln(x-2)-4\ln(x+5)$.
        Differentiate:
        \[
          \frac{y'}{y}=\frac{1}{2(x+1)}+\frac{1}{3(x-2)}-\frac{4}{x+5} .
        \]
        Hence
        \[
          y'=\frac{\sqrt{x+1}\sqrt[3]{x-2}}{(x+5)^{4}}
             \left[\frac{1}{2(x+1)}+\frac{1}{3(x-2)}-\frac{4}{x+5}\right] .
        \]

  \item $\mathrm{e}^{y}y'\ln(x)+\dfrac{\mathrm{e}^{y}}{x}+2x=y'\Rightarrow
        y'[\mathrm{e}^{y}\ln(x)-1]=-\dfrac{\mathrm{e}^{y}}{x}-2x\Rightarrow
        y'=\dfrac{\mathrm{e}^{y}/x+2x}{1-\mathrm{e}^{y}\ln(x)}$.

  \item Differentiate $x^{2}+y^{2}=r^{2}$: $2x+2yy'=0\Rightarrow
        y'=-x/y$. The radius from origin to $(x_{0},y_{0})$ has slope
        $y_{0}/x_{0}$. Product of slopes:
        $(-x_{0}/y_{0})\cdot(y_{0}/x_{0})=-1$, confirming
        perpendicularity.

  \item $\ln y=x^{2}\ln(7)\Rightarrow\dfrac{y'}{y}=2x\ln(7)\Rightarrow
        y'=2x\ln(7)\cdot 7^{x^{2}}$.

  \item Differentiate $x^{2}+xy+y^{2}=3$: $2x+y+xy'+2yy'=0\Rightarrow
        y'(x+2y)=-(2x+y)\Rightarrow y'=-\dfrac{2x+y}{x+2y}$.
        Differentiate again using the quotient rule:
        \[
          y''=-\frac{(2+y')(x+2y)-(2x+y)(1+2y')}{(x+2y)^{2}} .
        \]
        Expand the numerator: $(2+y')(x+2y)=2x+4y+xy'+2yy'$, and
        $(2x+y)(1+2y')=2x+y+4xy'+2yy'$. Subtract:
        $N:=(2+y')(x+2y)-(2x+y)(1+2y')=(2x+4y+xy'+2yy')
        -(2x+y+4xy'+2yy')=3y-3xy'$. So
        $y''=-\dfrac{3y-3xy'}{(x+2y)^{2}}=\dfrac{3(xy'-y)}{(x+2y)^{2}}$.
        Substitute $y'=-\dfrac{2x+y}{x+2y}$:
        \[
          xy'-y=-\frac{x(2x+y)}{x+2y}-y
               =-\frac{x(2x+y)+y(x+2y)}{x+2y}
        \]
        \[
          \phantom{xy'-y}=-\frac{2x^{2}+xy+xy+2y^{2}}{x+2y}
               =-\frac{2(x^{2}+xy+y^{2})}{x+2y},
        \]
        using $x^{2}+xy+y^{2}=3$. Hence
        \[
          y''=\frac{3}{(x+2y)^{2}}\cdot\left(-\frac{6}{x+2y}\right)
             =-\frac{18}{(x+2y)^{3}} .
        \]

  \item Differentiate $px+qy=m$ w.r.t.\ $x$: $p+qy'=0\Rightarrow
        y'=-p/q$, as required (the budget line has constant slope).

  \item Differentiate $x^{2}-xy+y^{2}=7$: $2x-y-xy'+2yy'=0\Rightarrow
        y'(2y-x)=y-2x\Rightarrow y'=\dfrac{y-2x}{2y-x}$. At $(3,2)$:
        $y'=\dfrac{2-6}{4-3}=-4$. Tangent: $y-2=-4(x-3)$, i.e.\
        $y=-4x+14$.

  \item Differentiate $\cos(xy)=x+y$: $-\sin(xy)\cdot(y+xy')=1+y'
        \Rightarrow -y\sin(xy)-x\sin(xy)y'=1+y'\Rightarrow
        -y'(x\sin(xy)+1)=1+y\sin(xy)\Rightarrow
        y'=-\dfrac{1+y\sin(xy)}{1+x\sin(xy)}$.

  \item First verify $(3,3)$: $27+27=54$ and $6\cdot 9=54$. Good.
        Differentiate $x^{3}+y^{3}=6xy$: $3x^{2}+3y^{2}y'=6y+6xy'
        \Rightarrow y'(3y^{2}-6x)=6y-3x^{2}\Rightarrow
        y'=\dfrac{6y-3x^{2}}{3y^{2}-6x}=\dfrac{2y-x^{2}}{y^{2}-2x}$.
        At $(3,3)$: $y'=\dfrac{6-9}{9-6}=-1$. For $y''$, differentiate
        $y'=\dfrac{2y-x^{2}}{y^{2}-2x}$:
        \[
          y''=\frac{(2y'-2x)(y^{2}-2x)-(2y-x^{2})(2yy'-2)}{(y^{2}-2x)^{2}} .
        \]
        At $(3,3)$ with $y'=-1$: numerator pieces:
        \[
          (2y'-2x)(y^{2}-2x)=(-2-6)(9-6)=(-8)(3)=-24,
        \]
        \[
          (2y-x^{2})(2yy'-2)=(6-9)(6\cdot(-1)-2)=(-3)(-8)=24 .
        \]
        So numerator $=-24-24=-48$ and denominator $=(9-6)^{2}=9$.
        Hence $y''(3)=-\dfrac{48}{9}=-\dfrac{16}{3}$.
\end{solutions}

\section{Additional Problems}

\begin{problems}
  \item Compute $\dd{y}{x}$ implicitly for $x^{2}+9y^{2}=36$.
  \item Find the equation of the tangent line to $x^{2}+y^{2}=25$ at
        $(3,4)$.
  \item Use logarithmic differentiation to find $y'$ for $y=x^{\cos x}$
        ($x>0$).
  \item Compute $\dd{y}{x}$ implicitly for $x^{2}y+xy^{2}=6$.
  \item Compute $\dd{^{2}y}{x^{2}}$ implicitly for $x^{2}-y^{2}=4$.
  \item Find the equation of the tangent line to $x^{3}+xy+y^{3}=11$ at
        $(1,2)$.
  \item Use logarithmic differentiation to find $y'$ for
        \[
          y=\frac{(x+1)^{2}\sqrt{x-3}}{(2x+5)^{4}}\qquad (x>3).
        \]
  \item Compute $\dd{y}{x}$ implicitly for $\mathrm{e}^{xy}=x+y$.
  \item For an indifference curve $U(x,y)=x^{1/3}y^{2/3}=\bar U$
        (constant), compute $\dd{y}{x}$, the marginal rate of
        substitution.
  \item Compute $\dd{y}{x}$ implicitly for $\sin(x+y)=xy$.
  \item Find the equation of the tangent line to $x^{2}-2xy+3y^{2}=9$ at
        $(1,2)$.
  \item Compute $\dd{^{2}y}{x^{2}}$ implicitly for $x^{2}+y^{2}=9$.
  \item Find all points on the curve $x^{2}+y^{2}=4x-2y$ where the tangent
        line is vertical.
  \item For the production isoquant $K^{2/5}L^{3/5}=\bar Q$ (constant),
        compute $\dd{K}{L}$, the marginal rate of technical
        substitution.
  \item Use logarithmic differentiation to find $y'$ for
        $y=(x^{2}+1)^{\sin x}$ ($x\in\R$).

  \item The curve $k$ below is the set of points satisfying
        \[
          x^{3}-3x^{2}+y^{2}-2y=-1 .
        \]
        Verify that $(0,1)$, $(3,1)$, $(-1,3)$ and $(-1,-1)$ all lie on
        $k$.

\begin{gallery}
\begin{tikzpicture}
  \begin{axis}[geoaxis, xmin=-1.6, xmax=3.6, ymin=-1.9, ymax=3.9,
      xtick={-1,1,2,3}, ytick={-1,1,2,3}, title={the curve $k$}]
    \addplot[domain=-1.2:3, samples=241]{1+abs(x)*sqrt(3-x)};
    \addplot[domain=-1.2:3, samples=241]{1-abs(x)*sqrt(3-x)};
  \end{axis}
\end{tikzpicture}\hfill
\begin{tikzpicture}
  \begin{axis}[geoaxis, xmin=-1.4, xmax=1.4, ymin=-0.9, ymax=0.9,
      xtick={-1,1}, ytick={-0.5,0.5}, title={the curve $\ell$}]
    \addplot[domain=-1:1, samples=241]{x*sqrt(1-x^2)};
    \addplot[domain=-1:1, samples=241]{-x*sqrt(1-x^2)};
  \end{axis}
\end{tikzpicture}
\end{gallery}

  \item Name every point of $k$ at which the hypotheses of the implicit
        function theorem fail for writing $y$ as a function of $x$, and say
        which graphical feature identifies each one.
  \item Confirm the previous answer by computation: solve $F_{y}=0$
        together with $F=0$.
  \item Compute the slope of the tangent to $k$ at $(-1,3)$.
  \item Compute the slope of the tangent to $k$ at $(-1,-1)$, and explain
        why it is the negative of the previous answer.
  \item Find every point of $k$ at which the tangent is horizontal.
  \item Near $(3,1)$ no function $y(x)$ describes $k$. Show that
        $x=x(y)$ does exist there, and compute $\dd{x}{y}$ at that point.
  \item The curve $\ell$ above is $y^{2}=x^{2}\bigl(1-x^{2}\bigr)$. Write it
        as $F(x,y)=0$ and give $F_{x}$ and $F_{y}$.
  \item Name every point of $\ell$ at which the implicit function theorem
        fails, and the graphical feature of each.
  \item Compute the slope of the tangent to $\ell$ at
        $\bigl(\tfrac12,\tfrac{\sqrt3}{4}\bigr)$.
  \item Compute the slope of the tangent to $\ell$ at
        $\bigl(\tfrac12,-\tfrac{\sqrt3}{4}\bigr)$.
  \item Find the four points of $\ell$ at which the tangent is horizontal.
  \item At the origin two branches of $\ell$ cross. Separate them by solving
        for $y$, and compute the slope of each branch there.
  \item Explain in one or two sentences why $-F_{x}/F_{y}$ cannot deliver
        those two slopes, and what has to be done instead.
\end{problems}

\section{Additional Solutions}

\begin{solutions}
  \item $2x+18yy'=0\Rightarrow y'=-\dfrac{x}{9y}$.

  \item Check $(3,4)$: $9+16=25$. From $2x+2yy'=0$: $y'=-x/y$. At
        $(3,4)$: $y'=-3/4$. Tangent: $y-4=-\tfrac34(x-3)$, i.e.\
        $3x+4y=25$.

  \item $\ln y=\cos(x)\ln(x)\Rightarrow\dfrac{y'}{y}
        =-\sin(x)\ln(x)+\dfrac{\cos x}{x}$. So
        \[
          y'=x^{\cos x}\left[-\sin(x)\ln(x)+\frac{\cos x}{x}\right] .
        \]

  \item Differentiate $x^{2}y+xy^{2}=6$: $2xy+x^{2}y'+y^{2}+2xyy'=0
        \Rightarrow y'(x^{2}+2xy)=-(2xy+y^{2})
        \Rightarrow y'=-\dfrac{y(2x+y)}{x(x+2y)}$.

  \item $2x-2yy'=0\Rightarrow y'=x/y$. Differentiate again:
        \[
          y''=\frac{y-xy'}{y^{2}}=\frac{y-x(x/y)}{y^{2}}
             =\frac{y^{2}-x^{2}}{y^{3}}=\frac{-4}{y^{3}},
        \]
        using $x^{2}-y^{2}=4\Rightarrow y^{2}-x^{2}=-4$.

  \item Check $(1,2)$: $1+2+8=11$. Differentiate:
        $3x^{2}+y+xy'+3y^{2}y'=0\Rightarrow y'(x+3y^{2})=-(3x^{2}+y)
        \Rightarrow y'=-\dfrac{3x^{2}+y}{x+3y^{2}}$. At $(1,2)$:
        $y'=-\dfrac{3+2}{1+12}=-\dfrac{5}{13}$. Tangent:
        $y-2=-\tfrac{5}{13}(x-1)$, i.e.\ $5x+13y=31$.

  \item $\ln y=2\ln(x+1)+\tfrac12\ln(x-3)-4\ln(2x+5)$. Then
        \[
          \frac{y'}{y}=\frac{2}{x+1}+\frac{1}{2(x-3)}-\frac{8}{2x+5} .
        \]
        Hence
        \[
          y'=\frac{(x+1)^{2}\sqrt{x-3}}{(2x+5)^{4}}
             \left[\frac{2}{x+1}+\frac{1}{2(x-3)}-\frac{8}{2x+5}\right] .
        \]

  \item Differentiate $\mathrm{e}^{xy}=x+y$:
        $\mathrm{e}^{xy}(y+xy')=1+y'\Rightarrow
        y'(x\mathrm{e}^{xy}-1)=1-y\mathrm{e}^{xy}\Rightarrow
        y'=\dfrac{1-y\mathrm{e}^{xy}}{x\mathrm{e}^{xy}-1}$.

  \item Differentiate $x^{1/3}y^{2/3}=\bar U$:
        $\tfrac13x^{-2/3}y^{2/3}+x^{1/3}\cdot\tfrac23y^{-1/3}y'=0$.
        Multiply by $3x^{2/3}y^{1/3}$: $y+2xy'=0\Rightarrow y'=-y/(2x)$.
        So $\dd{y}{x}=-\dfrac{y}{2x}$ (the MRS).

  \item Differentiate $\sin(x+y)=xy$: $\cos(x+y)(1+y')=y+xy'
        \Rightarrow y'(\cos(x+y)-x)=y-\cos(x+y)\Rightarrow
        y'=\dfrac{y-\cos(x+y)}{\cos(x+y)-x}$.

  \item Check $(1,2)$: $1-4+12=9$. Good. Differentiate:
        $2x-2y-2xy'+6yy'=0\Rightarrow y'(6y-2x)=2y-2x\Rightarrow
        y'=\dfrac{y-x}{3y-x}$. At $(1,2)$: $y'=\dfrac{1}{5}$. Tangent:
        $y-2=\tfrac15(x-1)$, i.e.\ $x-5y+9=0$.

  \item $2x+2yy'=0\Rightarrow y'=-x/y$. Differentiate again:
        \[
          y''=-\frac{y-xy'}{y^{2}}=-\frac{y-x(-x/y)}{y^{2}}
             =-\frac{y+x^{2}/y}{y^{2}}
             =-\frac{y^{2}+x^{2}}{y^{3}}=-\frac{9}{y^{3}},
        \]
        using $x^{2}+y^{2}=9$.

  \item $2x+2yy'=4-2y'\Rightarrow y'(2y+2)=4-2x\Rightarrow
        y'=\dfrac{2-x}{y+1}$. Vertical tangent: $y'$ undefined when
        $y+1=0\Rightarrow y=-1$ (and $x\neq 2$). Plug into the curve:
        $x^{2}+1=4x+2\Rightarrow x^{2}-4x-1=0\Rightarrow
        x=2\pm\sqrt{5}$. Vertical tangents at $(2+\sqrt{5},-1)$ and
        $(2-\sqrt{5},-1)$.

  \item Differentiate $K^{2/5}L^{3/5}=\bar Q$ w.r.t.\ $L$, treating
        $K=K(L)$:
        \[
          \tfrac25K^{-3/5}K'(L)\cdot L^{3/5}+K^{2/5}\cdot\tfrac35L^{-2/5}=0 .
        \]
        Multiply by $5K^{3/5}L^{2/5}$: $2LK'(L)+3K=0\Rightarrow
        K'(L)=-\dfrac{3K}{2L}$. So $\dd{K}{L}=-\dfrac{3K}{2L}$ (the
        MRTS).

  \item $\ln y=\sin(x)\ln(x^{2}+1)\Rightarrow\dfrac{y'}{y}
        =\cos(x)\ln(x^{2}+1)+\sin(x)\cdot\dfrac{2x}{x^{2}+1}$. So
        \[
          y'=(x^{2}+1)^{\sin x}\left[\cos(x)\ln(x^{2}+1)
             +\frac{2x\sin x}{x^{2}+1}\right] .
        \]

  \item With $F(x,y)=x^{3}-3x^{2}+y^{2}-2y+1$:
        $F(0,1)=0-0+1-2+1=0$; \;$F(3,1)=27-27+1-2+1=0$;
        \;$F(-1,3)=-1-3+9-6+1=0$; \;$F(-1,-1)=-1-3+1+2+1=0$.
        All four are on $k$.
  \item Graphically: at $(0,1)$ the curve \textbf{crosses itself}, and at
        $(3,1)$ it has a \textbf{vertical tangent}. At every other point of
        $k$ a local $y(x)$ exists.
  \item $F_{y}=2y-2=0\iff y=1$. Substituting $y=1$ into $F=0$ gives
        $x^{3}-3x^{2}+1-2+1=x^{2}(x-3)=0$, so $x=0$ (a double root) or
        $x=3$: the points $(0,1)$ and $(3,1)$. At $(3,1)$,
        $F_{x}=3x^{2}-6x=9\neq0$, so this is the vertical-tangent case; at
        $(0,1)$, $F_{x}=0$ as well, the singular case.
  \item $F_{x}=3x^{2}-6x$, $F_{y}=2y-2$, so
        $y'=-\dfrac{3x^{2}-6x}{2y-2}$. At $(-1,3)$:
        $y'=-\dfrac{3+6}{6-2}=-\dfrac94$.
  \item At $(-1,-1)$: $y'=-\dfrac{3+6}{-2-2}=\dfrac94$. The equation
        contains $y$ only through $(y-1)^{2}$, so $k$ is symmetric about the
        line $y=1$; reflecting a curve in a horizontal line negates every
        slope, and $(-1,3)$ and $(-1,-1)$ are mirror images.
  \item Horizontal tangent needs $F_{x}=0$ and $F_{y}\neq0$:
        $3x^{2}-6x=3x(x-2)=0$ gives $x=0$ or $x=2$. The value $x=0$ forces
        $y=1$, where $F_{y}=0$ too --- that is the singular point, not a
        horizontal tangent. For $x=2$: $(y-1)^{2}=x^{2}(3-x)=4$, so $y=3$ or
        $y=-1$. The tangent is horizontal at $(2,3)$ and $(2,-1)$.
  \item At $(3,1)$ we have $F_{x}=9\neq0$, so the theorem applies with the
        roles of the variables exchanged and gives a function $x=x(y)$.
        Differentiating $F\bigl(x(y),y\bigr)=0$ with respect to $y$:
        $F_{x}\,x'(y)+F_{y}=0$, hence
        $\dd{x}{y}=-\dfrac{F_{y}}{F_{x}}=-\dfrac{0}{9}=0$ --- the curve is
        momentarily vertical, which is exactly a horizontal tangent for
        $x(y)$.
  \item $F(x,y)=y^{2}-x^{2}+x^{4}$, with $F_{x}=-2x+4x^{3}=2x(2x^{2}-1)$
        and $F_{y}=2y$.
  \item $F_{y}=2y=0\iff y=0$, and then $x^{2}(1-x^{2})=0$ gives
        $x=0,\pm1$. At $(\pm1,0)$: $F_{x}=\pm2\neq0$, \textbf{vertical
        tangents} --- the two tips of the figure eight. At $(0,0)$:
        $F_{x}=0$ too, the \textbf{self-intersection} in the middle.
  \item $y'=-\dfrac{2x(2x^{2}-1)}{2y}$. At
        $\bigl(\tfrac12,\tfrac{\sqrt3}{4}\bigr)$ the numerator is
        $2\cdot\tfrac12\bigl(\tfrac12-1\bigr)=-\tfrac12$ and the denominator
        is $\tfrac{\sqrt3}{2}$, so
        $y'=\dfrac{1/2}{\sqrt3/2}=\dfrac{1}{\sqrt3}=\dfrac{\sqrt3}{3}$.
  \item Only the sign of $y$ changes, so
        $y'=-\dfrac{\sqrt3}{3}$.
  \item $F_{x}=2x(2x^{2}-1)=0$ with $F_{y}\neq0$ gives $x=\pm
        \tfrac{1}{\sqrt2}$ (the root $x=0$ is the singular point). Then
        $y^{2}=x^{2}(1-x^{2})=\tfrac12\cdot\tfrac12=\tfrac14$, so
        $y=\pm\tfrac12$: the four points
        $\bigl(\pm\tfrac{1}{\sqrt2},\pm\tfrac12\bigr)$.
  \item Solving gives $y=x\sqrt{1-x^{2}}$ and $y=-x\sqrt{1-x^{2}}$. For the
        first,
        $y'=\sqrt{1-x^{2}}-\dfrac{x^{2}}{\sqrt{1-x^{2}}}$, which at $x=0$
        equals $1$; the second branch gives $-1$. The two branches cross the
        origin at $+45^{\circ}$ and $-45^{\circ}$.
  \item At the origin both partials vanish, so $-F_{x}/F_{y}$ is $0/0$: it
        carries no information, and a single number could not describe two
        different slopes anyway. The branches have to be separated first ---
        here by solving for $y$ --- and each differentiated on its own.
\end{solutions}

\chapter{Applications: Tangent, Linearization, Marginals, Elasticity, Optimization}

The derivative is the bridge between calculus and economics. Geometrically it
gives the slope of the tangent (and hence of the normal); analytically it
produces the linearisation that turns nonlinear functions into
easy-to-evaluate affine ones; in economics it is the marginal quantity
(cost, revenue, utility) and, after rescaling, the elasticity. The first-
and second-order conditions $f'=0$, $f''<0$ (or $>0$) drive one-variable
optimisation.

\section{Core Rules}

\begin{rulebox}{Tangent line}
At $(a,f(a))$: $y = f(a) + f'(a)(x-a)$.
\end{rulebox}

\begin{rulebox}{Normal line}
At $(a,f(a))$: perpendicular to the tangent, slope $-1/f'(a)$ when
$f'(a)\neq 0$:
\[
  y = f(a) - \frac{1}{f'(a)}(x-a).
\]
\end{rulebox}

\begin{rulebox}{Linear approximation (linearisation)}
For $x$ near $a$,
\[
  f(x) \approx L(x) := f(a) + f'(a)(x-a), \qquad \Delta f \approx f'(a)\,\Delta x .
\]
\end{rulebox}

\begin{rulebox}{Marginal quantities}
If $C(Q)$, $R(Q)$, $\pi(Q)=R(Q)-C(Q)$ are total cost, revenue and profit,
and $U(x)$ is utility, then
\[
  MC(Q)=C'(Q), \quad MR(Q)=R'(Q), \quad M\pi(Q)=\pi'(Q), \quad MU(x)=U'(x).
\]
Profit-maximising output satisfies $MR(Q^*)=MC(Q^*)$ (FOC) and
$\pi''(Q^*)<0$ (SOC).
\end{rulebox}

\begin{rulebox}{Elasticity}
The (point) elasticity of $f$ with respect to $x$ is
\[
  \varepsilon = \frac{x}{f(x)}f'(x) = \frac{\mathrm{d}f/f}{\mathrm{d}x/x}
    = \frac{\mathrm{d}\ln f}{\mathrm{d}\ln x} .
\]
For demand $q=q(p)$: $\displaystyle \varepsilon_d = \frac{p}{q}q'(p)$
(often negative). Demand is \emph{elastic} when $|\varepsilon_d|>1$,
\emph{inelastic} when $|\varepsilon_d|<1$, \emph{unit elastic} when
$|\varepsilon_d|=1$.
\end{rulebox}

\begin{rulebox}{1D unconstrained optimisation}
\begin{enumerate}[leftmargin=1.6em,itemsep=0.3ex]
  \item Solve $f'(x)=0$ on the relevant domain to find critical points.
  \item Classify with the second-derivative test: $f''(x^*)<0 \Rightarrow$
        local max, $f''(x^*)>0 \Rightarrow$ local min.
  \item Compare with boundary values and endpoint behaviour to identify
        global extrema.
\end{enumerate}
\end{rulebox}

\section{Worked Example}

\begin{worked}{Demand is $q=100-2p$. Find the price elasticity at $p=20$ and at the revenue-maximising price, and approximate $\sqrt{101}$ by linearisation}
\wstep{Step 1. Elasticity is a derivative times a ratio --- not a derivative.}
\[
  \varepsilon=\frac{\mathrm{d}q}{\mathrm{d}p}\cdot\frac{p}{q}
   =\frac{\%\ \text{change in }q}{\%\ \text{change in }p} .
\]
Here $\mathrm{d}q/\mathrm{d}p=-2$ everywhere, but the elasticity is not
constant, because $p/q$ moves along the curve.
\wtrap{$\varepsilon\neq\mathrm{d}q/\mathrm{d}p$. The slope has units
(quantity per euro) and changes if you re-denominate the price; the
elasticity is a pure number and does not. That invariance is the main
reason economists use it.}

\wstep{Step 2. Evaluate at $p=20$.}
$q=100-40=60$, so
\[
  \varepsilon=-2\cdot\frac{20}{60}=-\tfrac23\approx-0.67 .
\]
$|\varepsilon|<1$: demand is \textbf{inelastic} here. A $1\%$ price rise
costs only $0.67\%$ of quantity, so revenue \emph{rises} --- the firm is
pricing below the revenue maximum.
\wtrap{The sign is negative for any downward-sloping demand, and dropping it
is a real error even though the classification uses $|\varepsilon|$. Report
$-0.67$ and then say ``inelastic''; do not report $0.67$.}

\wstep{Step 3. Maximise revenue directly.}
\[
  R(p)=pq=p(100-2p)=100p-2p^{2},
  \qquad
  R'(p)=100-4p=0
  \quad\Longrightarrow\quad
  p^{*}=25 .
\]
$R''=-4<0$, so this is a maximum. Then $q^{*}=50$ and $R^{*}=1250$.

\wstep{Step 4. Check the elasticity there --- and find the general rule.}
\[
  \varepsilon(25)=-2\cdot\frac{25}{50}=-1 .
\]
Unit elastic, exactly at the revenue peak. That is no coincidence:
\[
  R=pq
  \;\Longrightarrow\;
  \frac{\mathrm{d}R}{\mathrm{d}p}=q+p\frac{\mathrm{d}q}{\mathrm{d}p}
  =q\Bigl(1+\varepsilon\Bigr),
\]
which vanishes exactly when $\varepsilon=-1$. Revenue rises where demand is
inelastic, falls where it is elastic, and peaks in between.

\wstep{Step 5. Marginal revenue, and the same fact in quantity terms.}
Inverting demand, $p=50-q/2$, so $R(q)=50q-q^{2}/2$ and
\[
  MR(q)=50-q .
\]
$MR=0$ at $q=50$ --- the same point. Note $MR$ falls twice as fast as the
demand curve $p=50-q/2$: the standard ``$MR$ has twice the slope'' result,
here visible as $-1$ against $-\tfrac12$.
\wtrap{Marginal revenue is $\mathrm{d}R/\mathrm{d}q$, not $R/q$ (that is
\emph{average} revenue, which equals the price). Confusing marginal with
average is the most consequential error in applied differentiation, and it
recurs for cost: $MC=C'(q)$, while $AC=C(q)/q$.}

\wstep{Step 6. Linearisation: the tangent line as a calculator.}
For $f(x)=\sqrt{x}$ near $a=100$, where the value is known exactly:
\[
  f(x)\approx f(a)+f'(a)(x-a),
  \qquad
  f'(x)=\frac{1}{2\sqrt x},\quad f'(100)=\frac{1}{20} .
\]
\[
  \sqrt{101}\approx10+\frac{1}{20}(1)=10.05 .
\]
The true value is $10.0499$, so the error is $10^{-4}$.
\wtrap{Expand about the nearest point whose value you actually know. Using
$a=101$ makes the formula exact and useless; using $a=64$ would be legal but
far away, and the error grows with $(x-a)^{2}$. Choose $a$ first, and choose
it well.}

\wstep{Step 7. The same tangent, three names.}
\[
  \underbrace{y=f(a)+f'(a)(x-a)}_{\text{tangent line}}
  \;=\;
  \underbrace{f(a)+f'(a)\,\Delta x}_{\text{linearisation}}
  \;=\;
  \underbrace{f(a)+\text{marginal}\times\text{one more unit}}_{\text{economics}} .
\]
``Marginal cost is the cost of one more unit'' is a linear approximation:
strictly $C'(q)$ is a rate, and $C(q+1)-C(q)\approx C'(q)$ with an error of
order $C''/2$. The approximation is good precisely when the curvature is
small.
\end{worked}

\section{Practice Problems}

\begin{problems}
  \item Find the equation of the tangent to $f(x)=x^3-2x+5$ at $x=1$.
  \item Use the linearisation of $f(x)=\sqrt{x}$ at $a=25$ to estimate
        $\sqrt{25.4}$.
  \item For $C(Q)=0.2Q^2+5Q+80$, find $MC(Q)$ and $MC(15)$.
  \item For demand $q(p)=200-5p$, compute $\varepsilon_d(p)$ and evaluate at
        $p=10$.
  \item Maximise profit $\pi(Q)=R(Q)-C(Q)$ where $R(Q)=60Q-Q^2$ and
        $C(Q)=Q^2+10Q+25$. Verify the SOC.
  \item Find the tangent to $f(x)=1/x$ at $x=2$.
  \item Use the linearisation of $f(x)=x^{1/4}$ at $a=16$ to estimate
        $16.3^{1/4}$.
  \item Inverse demand $p(Q)=80-0.4Q$. Find $R(Q)$, $MR(Q)$, and $MR(50)$.
  \item For the Cobb--Douglas production $Y=AK^\alpha L^\beta$ with $L$
        fixed, compute the output elasticity $\varepsilon_{Y,K}$.
  \item Minimise the average cost $AC(Q)=C(Q)/Q$ for
        $C(Q)=2Q^2+8Q+200$, over $Q>0$.
  \item Find the tangent to $f(x)=e^{2x}$ at $x=0$.
  \item Estimate $\ln(0.98)$ using the linearisation of $f(x)=\ln(x)$ at
        $a=1$.
  \item Using $C$ from problem 3 and $R$ from problem 8, find the
        profit-maximising output $Q^*$ and verify SOC.
  \item For demand $q(p)=80e^{-0.05p}$, compute $\varepsilon_d(p)$ and at
        $p=20$.
  \item A monopolist faces $p(Q)=150-Q$ and $C(Q)=Q^2+30Q+100$. Find $Q^*$,
        $p^*$, and $\pi^*$.
  \item Find the normal line to $f(x)=\sqrt{x}$ at $x=9$.
  \item Estimate $e^{-0.2}$ using the linearisation of $f(x)=e^x$ at $a=0$.
  \item A monopolist faces $p(Q)=200-3Q$ and $C(Q)=2Q^2+10Q+50$. Find
        $Q^*$, $p^*$, and $\pi^*$.
  \item For $q(p)=200-5p$, find the price at which demand is unit elastic.
  \item A firm sells at the fixed price $p=24$ per unit; cost is
        $C(Q)=Q^2+8Q+10$. Find the profit-maximising output and verify SOC.
  \item Find the normal line to $f(x)=\ln(x)$ at $x=e$.
  \item A firm's cost is $C(Q)=0.5Q^2+8Q+200$. Use linear approximation at
        $Q=20$ to estimate $C(21.5)$.
  \item Marginal cost is $MC(Q)=0.06Q^2-2.4Q+30$. Find the production level
        that minimises $MC$ and verify the SOC.
  \item Demand $q(p)=\dfrac{300}{p^2+4}$. Compute $\varepsilon_d(p)$ and
        evaluate at $p=2$.
  \item Maximise utility $U(x)=4\ln(x)-0.5x$ for $x>0$.
  \item Find the tangent to $f(x)=\cos(x)$ at $x=\pi/3$.
  \item Use linearisation of $f(x)=(1+x)^{10}$ at $a=0$ to estimate
        $(1.02)^{10}$.
  \item For $C(Q)=Q^3-9Q^2+30Q+40$, compute $MC(Q)$ and locate the $Q$
        where $MC$ attains its minimum.
  \item For demand $q(p)=Ap^{-\alpha}$ with $A,\alpha>0$, show that the
        (point) elasticity is the constant $-\alpha$.
  \item A producer faces inverse demand $p(Q)=100-2Q$ and constant
        marginal cost $MC=20$ (so $C(Q)=20Q+F$ with fixed cost $F$). Find
        $Q^*$, $p^*$, and verify SOC.
  \item Find the normal line to $f(x)=x^2$ at $x=1$.
  \item Use linearisation of $f(x)=\tan(x)$ at $a=0$ to estimate
        $\tan(0.05)$.
  \item For the cubic cost $C(Q)=Q^3-6Q^2+15Q+20$, find marginal cost and
        average cost, and the $Q>0$ where $AC$ is minimised. Verify SOC.
  \item For demand $q(p)=100-2p^2$ (defined on $0\le p\le\sqrt{50}$),
        compute $\varepsilon_d(p)$ at $p=4$ and classify.
  \item Maximise revenue $R(Q)=Q\cdot p(Q)$ when $p(Q)=120-2Q$ and verify
        SOC.
  \item Find the tangent and the normal to $f(x)=x^2-3x+2$ at $x=2$.
  \item For a firm with $C(Q)=4Q^2+20Q+100$, what is the change in cost
        (approximately) when output increases from 10 to 10.5?
  \item For $C(Q)=Q^2+10Q+9$ and $R(Q)=30Q$, find the profit-maximising
        output and the break-even points (where $\pi=0$).
  \item For demand $q(p)=400-p^2$, find the price at which
        $|\varepsilon_d|=1$.
  \item Find positive $x$ that maximises $f(x)=x^2e^{-x}$ on
        $(0,\infty)$. Verify SOC.
  \item Find the tangent to $f(x)=\arctan(x)$ at $x=1$.
  \item Estimate $\sqrt[3]{27.9}$ using the linearisation of
        $f(x)=x^{1/3}$ at $a=27$.
  \item For utility $U(x)=\sqrt{x}$ and consumption $x=100$, compute the
        marginal utility $MU(x)$ and its elasticity $\varepsilon_{U,x}$.
        Comment briefly on signs.
  \item A monopolist faces demand $q(p)=120-3p$ and constant unit cost
        $c=10$. Reformulate via the inverse demand to find $Q^*$, $p^*$,
        $\pi^*$. Verify SOC.
  \item Minimise $f(x)=x+\dfrac{16}{x}$ on $(0,\infty)$. Verify SOC.
  \item Find both the tangent and the normal to $f(x)=\sin(x)$ at
        $x=\pi/2$.
  \item For revenue $R(Q)=p(Q)\cdot Q$ with $p(Q)=100-Q$, express $MR$ in
        terms of $p$ and the (price-)elasticity of demand: verify the
        identity $MR = p\left(1+\dfrac{1}{\varepsilon_d}\right)$ at
        $Q=20$.
  \item Use linearisation of $f(x)=\ln(1+x)$ at $a=0$ to estimate
        $\ln(1.05)$, and compute the linearisation error.
  \item A firm has revenue $R(Q)=20Q-Q^2$ and cost
        $C(Q)=2Q+\frac12Q^2+10$, with the constraint that $Q\in[0,8]$.
        Find the profit-maximising $Q$ on this closed interval, comparing
        the interior critical point with the endpoints.
  \item For demand $q(p)=50/p$, compute $\varepsilon_d$ for all $p>0$ and
        explain what it implies for total revenue $R(p)=p\cdot q(p)$.
  \item Maximise $\pi(Q)=12Q-Q^2-\dfrac{Q^3}{3}$ for $Q\ge0$. (Take the
        unique non-negative critical point and verify SOC.)
  \item For inverse demand $p(Q)=50-Q$ and cost $C(Q)=10Q+Q^2$, compute
        $\pi(Q)$, $Q^*$, $p^*$, $\pi^*$ and also $\varepsilon_d$ at $Q^*$.
  \item A firm produces according to $Y=\sqrt{L}$ (a strictly concave
        one-input production function). The wage is $w=2$ and the output
        price is $p=8$. Choose $L\ge0$ to maximise profit
        $\pi(L)=p\cdot Y(L)-wL$. Verify SOC.
\end{problems}

\section{Solutions}

\begin{solutions}
  \item $f'(x)=3x^2-2$. $f(1)=4$, $f'(1)=1$. Tangent:
        $y=4+1\cdot(x-1)=x+3$.

  \item $f(25)=5$, $f'(x)=1/(2\sqrt{x})$, $f'(25)=1/10=0.1$. So
        $\sqrt{25.4}\approx 5+0.1\cdot0.4=5.04$. (True: $5.0399\ldots$)

  \item $MC(Q)=C'(Q)=0.4Q+5$. $MC(15)=6+5=11$.

  \item $q'(p)=-5$.
        $\varepsilon_d(p)=\dfrac{p}{200-5p}\cdot(-5)=\dfrac{-5p}{200-5p}
        =\dfrac{-p}{40-p}$. At $p=10$: $\varepsilon_d=-10/30=-1/3$, so
        $|\varepsilon_d|<1$, inelastic.

  \item $\pi(Q)=60Q-Q^2-Q^2-10Q-25=50Q-2Q^2-25$. FOC:
        $\pi'(Q)=50-4Q=0\Rightarrow Q^*=12.5$. SOC: $\pi''(Q)=-4<0$, so a
        maximum. $\pi(12.5)=625-312.5-25=287.5$.

  \item $f'(x)=-1/x^2$. $f(2)=1/2$, $f'(2)=-1/4$. Tangent:
        $y=\frac12-\frac14(x-2)=-\frac{x}{4}+1$.

  \item $f'(x)=\frac14x^{-3/4}$. $f(16)=2$,
        $f'(16)=\frac14\cdot16^{-3/4}=\frac14\cdot\frac18=\frac1{32}$. So
        $16.3^{1/4}\approx2+\frac{0.3}{32}=2+0.009375=2.009375$.

  \item $R(Q)=p(Q)Q=80Q-0.4Q^2$. $MR(Q)=80-0.8Q$. $MR(50)=80-40=40$.

  \item $\dfrac{\partial Y}{\partial K}=\alpha AK^{\alpha-1}L^\beta$. So
        \[
          \varepsilon_{Y,K}=\frac{K}{Y}\cdot\alpha AK^{\alpha-1}L^\beta
            =\frac{\alpha AK^\alpha L^\beta}{AK^\alpha L^\beta}=\alpha .
        \]
        The Cobb--Douglas exponent equals the elasticity.

  \item $AC(Q)=2Q+8+\dfrac{200}{Q}$.
        $AC'(Q)=2-\dfrac{200}{Q^2}=0\Rightarrow Q^2=100\Rightarrow Q^*=10$
        (positive root). $AC''(Q)=\dfrac{400}{Q^3}>0$ at $Q=10$, so a
        minimum. $AC(10)=20+8+20=48$.

  \item $f'(x)=2e^{2x}$. $f(0)=1$, $f'(0)=2$. Tangent: $y=1+2x$.

  \item $f(1)=0$, $f'(1)=1$. $\ln(0.98)\approx0+1\cdot(-0.02)=-0.02$.
        (True: $-0.02020\ldots$)

  \item Set $MR(Q)=MC(Q)$: $80-0.8Q=0.4Q+5\Rightarrow75=1.2Q\Rightarrow
        Q^*=62.5$. SOC: $\pi''(Q)=R''(Q)-C''(Q)=-0.8-0.4=-1.2<0$, so this
        is a maximum.

  \item $q'(p)=80\cdot(-0.05)e^{-0.05p}=-4e^{-0.05p}$. So
        \[
          \varepsilon_d(p)=\frac{p}{80e^{-0.05p}}\cdot(-4e^{-0.05p})
            =-0.05p .
        \]
        At $p=20$: $\varepsilon_d=-1$, exactly unit elastic.

  \item $R(Q)=p(Q)Q=150Q-Q^2$.
        $\pi(Q)=150Q-Q^2-Q^2-30Q-100=120Q-2Q^2-100$. FOC:
        $\pi'(Q)=120-4Q=0\Rightarrow Q^*=30$. SOC: $\pi''(Q)=-4<0$.
        $p^*=150-30=120$. $\pi^*=120\cdot30-2\cdot900-100
        =3600-1800-100=1700$.

  \item $f'(x)=\dfrac{1}{2\sqrt{x}}$. $f(9)=3$, $f'(9)=1/6$. Normal slope
        $=-1/f'(9)=-6$. Normal: $y=3-6(x-9)=-6x+57$.

  \item $f(0)=1$, $f'(0)=1$. $e^{-0.2}\approx1+1\cdot(-0.2)=0.8$. (True:
        $0.8187\ldots$)

  \item $R(Q)=p(Q)Q=200Q-3Q^2$.
        $\pi(Q)=200Q-3Q^2-2Q^2-10Q-50=190Q-5Q^2-50$. FOC:
        $\pi'(Q)=190-10Q=0\Rightarrow Q^*=19$. SOC: $\pi''(Q)=-10<0$, max.
        $p^*=200-57=143$. $\pi^*=190\cdot19-5\cdot361-50
        =3610-1805-50=1755$.

  \item Using problem 4: $|\varepsilon_d|=1 \Leftrightarrow
        \dfrac{p}{40-p}=1\Rightarrow p=40-p\Rightarrow p=20$. (At $p=20$,
        demand $q=200-100=100$.)

  \item $\pi(Q)=24Q-Q^2-8Q-10=16Q-Q^2-10$. FOC:
        $\pi'(Q)=16-2Q=0\Rightarrow Q^*=8$. SOC: $\pi''(Q)=-2<0$.
        $\pi(8)=128-64-10=54$.

  \item $f'(x)=1/x$. $f(e)=1$, $f'(e)=1/e$. Normal slope $=-e$. Normal:
        $y=1-e(x-e)=1-ex+e^2$.

  \item $C'(Q)=Q+8$. $C(20)=200+160+200=560$, $C'(20)=28$. So
        $C(21.5)\approx560+28\cdot1.5=602$. (Exact:
        $0.5\cdot462.25+172+200=603.125$; error $\approx1.1$.)

  \item $MC'(Q)=0.12Q-2.4=0\Rightarrow Q=20$. $MC''(Q)=0.12>0$, so $Q=20$
        minimises $MC$. $MC(20)=0.06\cdot400-48+30=24-48+30=6$.

  \item $q(p)=300(p^2+4)^{-1}$,
        $q'(p)=-\dfrac{600p}{(p^2+4)^2}$. Then
        \[
          \varepsilon_d(p)=\frac{p}{q}\cdot q'(p)
            =\frac{p(p^2+4)}{300}\cdot\left(-\frac{600p}{(p^2+4)^2}\right)
            =\frac{-2p^2}{p^2+4} .
        \]
        At $p=2$: $\varepsilon_d=-8/8=-1$ (unit elastic).

  \item $U'(x)=4/x-0.5=0\Rightarrow x^*=8$. $U''(x)=-4/x^2<0$, max.
        $U(8)=4\ln(8)-4=12\ln(2)-4\approx4.318$.

  \item $f'(x)=-\sin(x)$. $f(\pi/3)=1/2$, $f'(\pi/3)=-\sqrt3/2$. Tangent:
        $y=\frac12-\frac{\sqrt3}{2}\left(x-\frac{\pi}{3}\right)$.

  \item $f(0)=1$, $f'(x)=10(1+x)^9$, $f'(0)=10$. So
        $(1.02)^{10}\approx1+10\cdot0.02=1.2$. (True: $1.21899\ldots$,
        error $\approx0.019$.)

  \item $MC(Q)=C'(Q)=3Q^2-18Q+30$. $MC'(Q)=6Q-18=0\Rightarrow Q=3$.
        $MC''(Q)=6>0$, min. $MC(3)=27-54+30=3$.

  \item $q'(p)=A\cdot(-\alpha)p^{-\alpha-1}=-\alpha\dfrac{q}{p}$. Hence
        $\varepsilon_d=\dfrac{p}{q}q'(p)=\dfrac{p}{q}\cdot
        \left(-\alpha\dfrac{q}{p}\right)=-\alpha$. Constant elasticity
        demand has elasticity equal to the negative exponent.

  \item $R(Q)=p(Q)Q=100Q-2Q^2$. $\pi(Q)=100Q-2Q^2-20Q-F=80Q-2Q^2-F$. FOC:
        $\pi'(Q)=80-4Q=0\Rightarrow Q^*=20$. SOC: $\pi''(Q)=-4<0$, max.
        $p^*=100-40=60$. Cross-check via $MR=MC$: $100-4Q=20\Rightarrow
        Q=20$. Consistent. (Profit $\pi^*=80\cdot20-2\cdot400-F
        =1600-800-F=800-F$.)

  \item $f'(x)=2x$. $f(1)=1$, $f'(1)=2$. Normal slope $=-1/2$. Normal:
        $y=1-\frac12(x-1)=-\frac{x}{2}+\frac32$.

  \item $f(0)=0$, $f'(x)=\sec^2(x)$, $f'(0)=1$. So
        $\tan(0.05)\approx0+1\cdot0.05=0.05$. (True: $0.050042\ldots$)

  \item $MC(Q)=C'(Q)=3Q^2-12Q+15$. $AC(Q)=Q^2-6Q+15+20/Q$.
        $AC'(Q)=2Q-6-20/Q^2$. Set $2Q-6-20/Q^2=0$, multiply by $Q^2$:
        $2Q^3-6Q^2-20=0\Rightarrow Q^3-3Q^2-10=0$. Try $Q=$ roots: at
        $Q=3$, $27-27-10=-10$; at $Q=4$, $64-48-10=6$; root between 3 and
        4. Newton step from $Q_0=3.7$: $Q_0^3-3Q_0^2-10
        =50.653-41.07-10=-0.417$, derivative $3Q^2-6Q=3\cdot13.69-22.2
        =18.87$, $Q_1=3.7+0.417/18.87\approx3.722$. So $Q^*\approx3.72$.
        SOC: $AC''(Q)=2+40/Q^3>0$ for $Q>0$, so it is a minimum. (Note: at
        the optimum $AC(Q^*)=MC(Q^*)$, a standard identity.)

  \item $q'(p)=-4p$. At $p=4$: $q(4)=100-32=68$, $q'(4)=-16$.
        $\varepsilon_d(4)=\dfrac{4}{68}\cdot(-16)=-\dfrac{64}{68}
        =-\dfrac{16}{17}\approx-0.94$. Since $|\varepsilon_d|<1$, demand
        is \emph{inelastic} at $p=4$.

  \item $R(Q)=120Q-2Q^2$. $R'(Q)=120-4Q=0\Rightarrow Q^*=30$.
        $R''(Q)=-4<0$, max. $R(30)=3600-1800=1800$.

  \item $f'(x)=2x-3$. $f(2)=4-6+2=0$, $f'(2)=1$. Tangent:
        $y=0+1\cdot(x-2)=x-2$. Normal slope $=-1$. Normal:
        $y=-(x-2)=-x+2$.

  \item $C'(Q)=8Q+20$. $\Delta C\approx C'(10)\cdot0.5
        =(80+20)\cdot0.5=100\cdot0.5=50$.

  \item $\pi(Q)=30Q-Q^2-10Q-9=20Q-Q^2-9$. FOC:
        $\pi'(Q)=20-2Q=0\Rightarrow Q^*=10$. SOC: $\pi''(Q)=-2<0$, max.
        $\pi(10)=200-100-9=91$. Break-even: $20Q-Q^2-9=0\Rightarrow
        Q^2-20Q+9=0\Rightarrow Q=\dfrac{20\pm\sqrt{400-36}}{2}
        =10\pm\sqrt{91}$. Both positive: $Q\approx10-9.54=0.46$ and
        $Q\approx19.54$.

  \item $q(p)=400-p^2$, $q'(p)=-2p$.
        $\varepsilon_d=\dfrac{p}{400-p^2}\cdot(-2p)=\dfrac{-2p^2}{400-p^2}$.
        Set $\dfrac{2p^2}{400-p^2}=1\Rightarrow2p^2=400-p^2
        \Rightarrow3p^2=400\Rightarrow p=\dfrac{20}{\sqrt3}\approx11.55$.

  \item $f'(x)=2xe^{-x}-x^2e^{-x}=x(2-x)e^{-x}$. On $(0,\infty)$:
        $f'>0$ for $0<x<2$, $f'<0$ for $x>2$. Critical point at $x=2$.
        $f''(x)=(2-2x)e^{-x}-(2x-x^2)e^{-x}=(x^2-4x+2)e^{-x}$. At $x=2$:
        $f''(2)=(4-8+2)e^{-2}=-2e^{-2}<0$, confirming local (and global on
        $(0,\infty)$) maximum. $f(2)=4e^{-2}\approx0.541$.

  \item $f'(x)=\dfrac{1}{1+x^2}$. $f(1)=\pi/4$, $f'(1)=1/2$. Tangent:
        $y=\dfrac{\pi}{4}+\dfrac12(x-1)$.

  \item $f(27)=3$, $f'(x)=\frac13x^{-2/3}$,
        $f'(27)=\frac13\cdot27^{-2/3}=\frac13\cdot\frac19=\frac1{27}$. So
        $\sqrt[3]{27.9}\approx3+\frac1{27}\cdot0.9=3+1/30\approx3.0333$.
        (True: $3.03296\ldots$)

  \item $U'(x)=\dfrac{1}{2\sqrt{x}}>0$ for $x>0$ (monotone preferences);
        at $x=100$: $MU=\dfrac{1}{20}=0.05>0$. Elasticity:
        $\varepsilon_{U,x}=\dfrac{x}{U(x)}U'(x)=\dfrac{x}{\sqrt{x}}
        \cdot\dfrac{1}{2\sqrt{x}}=\dfrac12$. Constant elasticity $1/2$,
        positive (utility rises with $x$). Marginal utility itself is
        positive but its elasticity is positive and less than 1.

  \item Invert demand: $q=120-3p\Rightarrow p(Q)=40-Q/3$.
        $R(Q)=Qp=40Q-Q^2/3$. $\pi(Q)=40Q-Q^2/3-10Q=30Q-Q^2/3$. FOC:
        $\pi'(Q)=30-2Q/3=0\Rightarrow Q^*=45$. SOC: $\pi''(Q)=-2/3<0$,
        max. $p^*=40-15=25$, $\pi^*=30\cdot45-45^2/3=1350-675=675$.

  \item $f'(x)=1-16/x^2=0\Rightarrow x^2=16\Rightarrow x=4$ on
        $(0,\infty)$. $f''(x)=32/x^3>0$ for $x>0$, so $x=4$ is a minimum.
        $f(4)=4+4=8$.

  \item $f'(x)=\cos(x)$. $f(\pi/2)=1$, $f'(\pi/2)=0$. The tangent is
        horizontal: $y=1$. The normal is the vertical line $x=\pi/2$
        (slope $-1/0$ undefined; normal is vertical).

  \item At $Q=20$: $p=80$. $R(Q)=100Q-Q^2$, $MR(Q)=100-2Q$, $MR(20)=60$.
        From $q=100-p$, $q'(p)=-1$,
        $\varepsilon_d(p)=\dfrac{p}{q}\cdot(-1)=-\dfrac{p}{100-p}$. At
        $p=80$: $\varepsilon_d=-80/20=-4$. Then
        $p\left(1+\dfrac{1}{\varepsilon_d}\right)=80\left(1-\dfrac14\right)
        =80\cdot\dfrac34=60=MR$. Verified.

  \item $f(0)=0$, $f'(x)=1/(1+x)$, $f'(0)=1$. So
        $\ln(1.05)\approx0+1\cdot0.05=0.05$. True value:
        $\ln(1.05)=0.048790\ldots$ Linearisation error
        $\approx0.05-0.04879=0.00121$.

  \item $\pi(Q)=R(Q)-C(Q)=(20Q-Q^2)-(2Q+\frac12Q^2+10)
        =18Q-\frac32Q^2-10$. FOC: $\pi'(Q)=18-3Q=0\Rightarrow Q=6$.
        $Q=6\in[0,8]$. SOC: $\pi''(Q)=-3<0$. $\pi(6)=108-54-10=44$. Check
        endpoints: $\pi(0)=-10$, $\pi(8)=144-96-10=38$. So $Q^*=6$
        maximises profit; $\pi^*=44$.

  \item $q(p)=50/p$, $q'(p)=-50/p^2$.
        $\varepsilon_d=\dfrac{p}{q}q'(p)=\dfrac{p}{50/p}\cdot(-50/p^2)
        =\dfrac{p^2}{50}\cdot(-50/p^2)=-1$ for all $p>0$. The elasticity
        is identically $-1$, unit-elastic demand. Total revenue
        $R(p)=p\cdot50/p=50$, constant, independent of $p$. This is the
        hallmark of unit-elastic demand.

  \item $\pi'(Q)=12-2Q-Q^2$. Solve $Q^2+2Q-12=0\Rightarrow
        Q=\dfrac{-2+\sqrt{4+48}}{2}=\dfrac{-2+\sqrt{52}}{2}
        =-1+\sqrt{13}$. So $Q^*=\sqrt{13}-1\approx2.606$ (the positive
        root). SOC: $\pi''(Q)=-2-2Q$. At $Q^*$:
        $\pi''(Q^*)=-2-2(\sqrt{13}-1)=-2\sqrt{13}<0$, confirming maximum.

  \item $R(Q)=(50-Q)Q=50Q-Q^2$. $\pi(Q)=50Q-Q^2-10Q-Q^2=40Q-2Q^2$. FOC:
        $\pi'(Q)=40-4Q=0\Rightarrow Q^*=10$. SOC: $\pi''(Q)=-4<0$, max.
        $p^*=50-10=40$, $\pi^*=40\cdot10-2\cdot100=400-200=200$.
        Elasticity: $q(p)=50-p$, $q'(p)=-1$,
        $\varepsilon_d(p)=\dfrac{p}{50-p}\cdot(-1)$. At $p^*=40$:
        $\varepsilon_d=-40/10=-4$ (elastic, as expected at a monopolist's
        optimum where $|\varepsilon_d|>1$).

  \item $\pi(L)=8\sqrt{L}-2L$.
        $\pi'(L)=\dfrac{4}{\sqrt{L}}-2=0\Rightarrow\sqrt{L}=2\Rightarrow
        L^*=4$. SOC: $\pi''(L)=-\dfrac{2}{L^{3/2}}<0$ for $L>0$, so
        $L^*=4$ maximises profit. $Y(4)=2$,
        $\pi^*=8\cdot2-2\cdot4=16-8=8$.
\end{solutions}

\section{Additional Problems}

\begin{problems}
  \item Find the equation of the tangent to $f(x)=x^2+3x-1$ at $x=2$.
  \item Use the linearisation of $f(x)=\sqrt[3]{x}$ at $a=8$ to estimate
        $\sqrt[3]{8.6}$.
  \item For $C(Q)=0.05Q^2+8Q+150$, find $MC(Q)$ and $MC(20)$.
  \item For demand $q(p)=300-4p$, compute $\varepsilon_d(p)$ at $p=25$ and
        classify.
  \item Maximise profit $\pi(Q)=R(Q)-C(Q)$ where $R(Q)=90Q-1.5Q^2$ and
        $C(Q)=Q^2+20Q+30$.
  \item Find the tangent and normal to $f(x)=x^3-4x$ at $x=1$.
  \item Use the linearisation of $f(x)=\mathrm{e}^{x}$ at $a=0$ to estimate
        $\mathrm{e}^{0.1}$.
  \item Inverse demand $p(Q)=100-0.5Q$. Find $R(Q)$, $MR(Q)$, and $MR(40)$.
  \item For the Cobb--Douglas production $Y=AK^{\alpha}L^{1-\alpha}$ with
        $K$ fixed, compute the output elasticity $\varepsilon_{Y,L}$.
  \item Minimise the average cost $AC(Q)=C(Q)/Q$ for
        $C(Q)=3Q^2+12Q+108$, over $Q>0$.
  \item A monopolist faces $p(Q)=80-2Q$ and $C(Q)=Q^2+8Q+40$. Find $Q^*$,
        $p^*$, and $\pi^*$.
  \item For demand $q(p)=120\mathrm{e}^{-0.04p}$, compute
        $\varepsilon_d(p)$ and at $p=10$.
  \item For utility $U(x)=2\sqrt{x}-0.1x$ for $x>0$, find the maximiser
        and verify SOC.
  \item Estimate $\ln(1.04)$ using the linearisation of $f(x)=\ln(1+x)$
        at $a=0$.
  \item Minimise $f(x)=2x+\dfrac{50}{x}$ on $(0,\infty)$ and verify SOC.
\end{problems}

\section{Additional Solutions}

\begin{solutions}
  \item $f'(x)=2x+3$. $f(2)=9$, $f'(2)=7$. Tangent: $y=9+7(x-2)=7x-5$.

  \item $f(8)=2$, $f'(x)=\tfrac{1}{3}x^{-2/3}$,
        $f'(8)=\tfrac{1}{3}\cdot 8^{-2/3}=\tfrac{1}{3}\cdot\tfrac{1}{4}
        =\tfrac{1}{12}$. So
        $\sqrt[3]{8.6}\approx 2+\tfrac{0.6}{12}=2+0.05=2.05$. (True:
        $2.0494\ldots$)

  \item $MC(Q)=C'(Q)=0.1Q+8$. $MC(20)=2+8=10$.

  \item $q'(p)=-4$.
        $\varepsilon_d(p)=\dfrac{p}{300-4p}\cdot(-4)=\dfrac{-4p}{300-4p}$.
        At $p=25$: $\varepsilon_d=\dfrac{-100}{200}=-\dfrac{1}{2}$.
        $|\varepsilon_d|<1$: \emph{inelastic}.

  \item $\pi(Q)=90Q-1.5Q^2-Q^2-20Q-30=70Q-2.5Q^2-30$. FOC:
        $\pi'(Q)=70-5Q=0\Rightarrow Q^*=14$. SOC: $\pi''(Q)=-5<0$, max.
        $\pi(14)=980-490-30=460$.

  \item $f'(x)=3x^2-4$. $f(1)=-3$, $f'(1)=-1$. Tangent:
        $y=-3-(x-1)=-x-2$. Normal slope $=-1/(-1)=1$. Normal:
        $y=-3+1\cdot(x-1)=x-4$.

  \item $f(0)=1$, $f'(0)=1$. $\mathrm{e}^{0.1}\approx 1+0.1=1.1$. (True:
        $1.10517\ldots$)

  \item $R(Q)=p(Q)Q=100Q-0.5Q^2$. $MR(Q)=100-Q$. $MR(40)=60$.

  \item $\dfrac{\partial Y}{\partial L}=(1-\alpha)AK^{\alpha}L^{-\alpha}$.
        So
        \[
          \varepsilon_{Y,L}=\frac{L}{Y}\cdot(1-\alpha)AK^{\alpha}L^{-\alpha}
            =\frac{(1-\alpha)AK^{\alpha}L^{1-\alpha}}{AK^{\alpha}L^{1-\alpha}}
            =1-\alpha .
        \]
        Standard Cobb--Douglas: elasticity equals the exponent.

  \item $AC(Q)=3Q+12+\dfrac{108}{Q}$.
        $AC'(Q)=3-\dfrac{108}{Q^2}=0\Rightarrow Q^2=36\Rightarrow Q^*=6$.
        $AC''(Q)=\dfrac{216}{Q^3}>0$ at $Q=6$: minimum.
        $AC(6)=18+12+18=48$.

  \item $R(Q)=p(Q)Q=80Q-2Q^2$.
        $\pi(Q)=80Q-2Q^2-Q^2-8Q-40=72Q-3Q^2-40$. FOC:
        $\pi'(Q)=72-6Q=0\Rightarrow Q^*=12$. SOC: $\pi''(Q)=-6<0$.
        $p^*=80-24=56$. $\pi^*=72\cdot12-3\cdot144-40=864-432-40=392$.

  \item $q'(p)=120\cdot(-0.04)\mathrm{e}^{-0.04p}
        =-4.8\mathrm{e}^{-0.04p}$. So
        \[
          \varepsilon_d(p)=\frac{p}{120\mathrm{e}^{-0.04p}}
            \cdot(-4.8\mathrm{e}^{-0.04p})=-0.04p .
        \]
        At $p=10$: $\varepsilon_d=-0.4$. $|\varepsilon_d|<1$:
        \emph{inelastic}.

  \item $U'(x)=\tfrac{1}{\sqrt{x}}-0.1=0\Rightarrow\sqrt{x}=10
        \Rightarrow x^*=100$. $U''(x)=-\tfrac{1}{2}x^{-3/2}<0$ for $x>0$,
        max. $U(100)=20-10=10$.

  \item $f(0)=0$, $f'(0)=1$. $\ln(1.04)\approx 0+1\cdot0.04=0.04$.
        (True: $0.03922\ldots$)

  \item $f'(x)=2-\dfrac{50}{x^2}=0\Rightarrow x^2=25\Rightarrow x=5$ on
        $(0,\infty)$. $f''(x)=\dfrac{100}{x^3}>0$ for $x>0$, so $x=5$ is
        a minimum. $f(5)=10+10=20$.
\end{solutions}

\chapter{First-Order Partial Derivatives and the Gradient}

When a function depends on more than one input, we measure how it responds
to a change in \emph{one} variable at a time, holding the others fixed,
the ceteris paribus principle. This is exactly what a partial derivative
computes. From these one-variable slopes we assemble the gradient, which
packages all first-order information into a single vector and unlocks the
geometry of multivariate optimisation.

\section{Core Rules}

\begin{rulebox}{Topological vocabulary of domains}
\begin{itemize}[leftmargin=1.6em,itemsep=0.3ex]
  \item The \emph{open ball} (or \emph{neighbourhood}) of radius $r>0$ around $\mathbf{x}_0\in\R^n$ is
        $B_r(\mathbf{x}_0)=\{\mathbf{x}\in\R^n :
        \|\mathbf{x}-\mathbf{x}_0\|<r\}$.
  \item $\mathbf{x}_0\in M$ is an \emph{interior point} of $M$ if some open ball around
        $\mathbf{x}_0$ lies entirely in $M$.
  \item $M$ is \emph{open} if \emph{every} point of $M$ is
        an interior point.
  \item $\mathbf{x}_0$ is a \emph{boundary point} of
        $M$ if every ball around $\mathbf{x}_0$ contains points of $M$
        \emph{and} points of its complement.
  \item $M$ is \emph{closed} if it contains all
        its boundary points; equivalently, the complement
        $\R^n\setminus M$ is open.
  \item $M$ is \emph{bounded} if
        $M\subseteq B_r(\mathbf{0})$ for some radius $r$.
  \item $M\subseteq\R^n$ is \emph{compact} exactly when
        $M$ is closed \emph{and} bounded.
\end{itemize}
Economic examples: the open positive orthant $\{(x,y): x>0,\ y>0\}$
(the natural Cobb--Douglas domain) is open and unbounded; a budget set
$\{\mathbf{x}\in\R^n : \mathbf{p}\cdot\mathbf{x}\le m,\
\mathbf{x}\ge\mathbf{0}\}$ is closed and bounded, hence compact.
\end{rulebox}

\begin{rulebox}{Why it matters}
Differentiability (and hence every partial derivative below) is
defined at \emph{interior} points of the domain, where $f$ can be
probed in all directions. And by the \emph{Weierstrass theorem}, a
continuous function on a compact set attains its maximum and minimum:
this is why budget-constrained optimisation problems are
guaranteed to have solutions.
\end{rulebox}

\begin{rulebox}{Definition}
Let $f:\R^n\to\R$. The partial derivative of $f$ with respect to $x_i$ at
$\mathbf{x}=(x_1,\ldots,x_n)$ is
\[
  \pd{f}{x_i}(\mathbf{x}) \;=\; \lim_{h\to0}
    \frac{f(x_1,\ldots,x_i+h,\ldots,x_n)-f(x_1,\ldots,x_n)}{h} .
\]
\end{rulebox}

\begin{rulebox}{Practical computation}
To compute $\partial f/\partial x_i$, differentiate $f$ with respect to
$x_i$ and treat every other variable as a constant. All single-variable
rules (power, product, quotient, chain) apply unchanged.
\end{rulebox}

\begin{rulebox}{Notation}
The following all refer to the same object:
\[
  \pd{f}{x} = f_x = \partial_x f = D_x f .
\]
\end{rulebox}

\begin{rulebox}{Gradient}
The vector of all first-order partial derivatives is
\[
  \grad f(\mathbf{x}) \;=\;
  \begin{pmatrix} f_{x_1}(\mathbf{x}) \\ f_{x_2}(\mathbf{x}) \\ \vdots \\
    f_{x_n}(\mathbf{x}) \end{pmatrix} .
\]
\end{rulebox}

\begin{rulebox}{Geometric meaning}
\begin{itemize}[leftmargin=1.6em,itemsep=0.3ex]
  \item $\grad f(\mathbf{x}_0)$ points in the direction of \emph{steepest
        ascent} at $\mathbf{x}_0$.
  \item $\|\grad f(\mathbf{x}_0)\|$ is the maximal rate of increase at that
        point.
  \item $\grad f(\mathbf{x}_0)$ is \emph{perpendicular} to the level set
        $\{\mathbf{x} : f(\mathbf{x})=f(\mathbf{x}_0)\}$ at $\mathbf{x}_0$.
\end{itemize}
\end{rulebox}

\begin{rulebox}{Economic interpretation}
\begin{itemize}[leftmargin=1.6em,itemsep=0.3ex]
  \item For a production function $Q=F(K,L)$: $F_K=\mathrm{MP}_K$
        (marginal product of capital), $F_L=\mathrm{MP}_L$ (marginal
        product of labour).
  \item For a utility function $U(x,y)$: $U_x, U_y$ are the marginal
        utilities; the marginal rate of substitution is
        $\mathrm{MRS}_{xy}=U_x/U_y$.
  \item For production: the marginal rate of technical substitution is
        $\mathrm{MRTS}_{LK}=F_L/F_K$.
\end{itemize}
\end{rulebox}

\begin{rulebox}{Directional derivative}
Let $\mathbf{u}\in\R^n$ be a \emph{unit} vector ($\|\mathbf{u}\|=1$). The
directional derivative of $f$ at $\mathbf{x}_0$ in the direction
$\mathbf{u}$ is
\[
  D_{\mathbf{u}}f(\mathbf{x}_0) \;=\; \grad f(\mathbf{x}_0)\cdot\mathbf{u} .
\]
\end{rulebox}

\begin{rulebox}{Caution}
If you are given a non-unit direction $\mathbf{v}$, you \emph{must}
normalise first: $\mathbf{u}=\mathbf{v}/\|\mathbf{v}\|$. Forgetting this is
the single most common error in directional-derivative problems.
\end{rulebox}

\section{Practice Problems}

\begin{problems}
  \item Compute $f_x$ and $f_y$ for $f(x,y)=5x^4-2x^2y^3+y^5$.
  \item For $F(K,L)=12\,K^{1/4}L^{3/4}$, find $F_K$ and $F_L$ at
        $(K,L)=(16,81)$.
  \item Compute $f_x$ and $f_y$ for $f(x,y)=x^2e^{3y}$.
  \item Compute $\nabla f$ for $f(x,y)=4x^3y-xy^4+7$.
  \item Compute the directional derivative of $f(x,y)=x^2-4xy+y^3$ at
        $(2,1)$ in the direction $\mathbf{v}=(5,-12)$.
  \item For $F(K,L)=6K^{2/3}L^{1/3}$, compute $\mathrm{MP}_K$ and
        $\mathrm{MP}_L$ at $(K,L)=(27,8)$.
  \item Compute $f_x$ and $f_y$ for $f(x,y)=(3x-2y)^4$.
  \item Compute $f_x$ and $f_y$ for $f(x,y)=e^{xy^2}$.
  \item For $F(K,L)=25K^{0.6}L^{0.5}$, find $F_K$ and $F_L$.
  \item Compute $\nabla f$ for $f(x,y,z)=x^2+3yz-z^2e^x$.
  \item Compute the directional derivative of $f(x,y)=e^x\ln(1+y)$ at
        $(0,0)$ in the direction $\mathbf{v}=(1,2)$.
  \item For $U(x,y)=2\ln(x)+5\ln(y)$, derive $\mathrm{MRS}_{xy}=U_x/U_y$ as
        a function of $(x,y)$ and evaluate at $(3,5)$.
  \item Simplify and compute $f_x, f_y$ for
        $f(x,y)=\dfrac{x^2-y^2}{x+y}$.
  \item Compute $f_x$ and $f_y$ for $f(x,y)=\ln(x^2y+1)$.
  \item For $U(x_1,x_2,x_3)=x_1^{1/2}x_2^{1/3}x_3^{1/6}$, compute the
        three partials.
  \item Compute the directional derivative of $f(x,y,z)=xy-2yz+z^2$ at
        $(1,2,3)$ in the direction $\mathbf{v}=(1,2,2)$.
  \item For $F(K,L)=K^{0.25}L^{0.75}$, derive $\mathrm{MRTS}_{LK}=F_L/F_K$
        and evaluate at $(K,L)=(16,4)$.
  \item Expand and then compute $f_x, f_y$ for
        $f(x,y)=(x^2+y)(x-3y^2)$.
  \item Compute $f_x$ and $f_y$ for $f(x,y)=xe^{-y/x}$ with $x\neq0$.
  \item For $F(K,L)=(3K^{1/2}+2L^{1/2})^2$, find $F_K$ and $F_L$.
  \item Compute $\nabla f$ for $f(x,y)=\arctan(y/x)$ with $x>0$.
  \item For $f(x,y)=x^{1/2}y^{1/2}$ at $(K,L)=(4,9)$, find the unit
        direction of steepest ascent and the maximum rate of increase.
  \item For the CES function $F(K,L)=(K^{1/2}+L^{1/2})^2$, verify Euler's
        theorem $KF_K+LF_L=rF$ for the appropriate degree of homogeneity
        $r$.
  \item Compute $f_x$ and $f_y$ for $f(x,y)=\dfrac{2x+3y}{x^2+y^2}$.
  \item Compute $f_x$ and $f_y$ for $f(x,y)=\dfrac{e^x}{1+e^y}$.
  \item For $F(K,L)=A(\alpha K^\rho+(1-\alpha)L^\rho)^{\nu/\rho}$ (CES),
        find $F_L$.
  \item Compute the directional derivative of $f(x,y)=\ln(x^2+y^2)$ at
        $(3,4)$ in the direction $\mathbf{v}=(-1,1)$.
  \item Compute $f_x, f_y, f_z$ for $f(x,y,z)=x^2yz^3-4xy^2z+5z$.
  \item Compute $f_x$ and $f_y$ for $f(x,y)=\sqrt{4x^2+9y^2}$.
  \item For $F(K,L)=(K^{1/2}+L^{1/2})^4$, compute $F_K$ at
        $(K,L)=(4,9)$.
  \item Compute $\nabla f$ for $f(x,y,z)=\ln(x^2+y^2+z^2)$.
  \item Compute $f_x$ and $f_y$ for
        $f(x,y)=\ln\left(\sqrt{x^2+y^2}\right)$.
  \item Cost minimisation. A firm uses inputs $(x,y)$ with cost
        $C(x,y)=4x+9y$ and output $Q(x,y)=x^{1/2}y^{1/2}$. Compute
        $\nabla C$ and $\nabla Q$, and verify that at the cost-minimising
        input mix $\nabla C=\lambda\nabla Q$ for some $\lambda$.
  \item Compute $f_x$ and $f_y$ for $f(x,y)=(1+xy)^5$.
  \item Compute the directional derivative of $f(x,y)=x^3-2xy+y^2$ at
        $(1,2)$ in the direction $\mathbf{v}=(3,4)$.
  \item For $F(K,L)=10K^{0.3}L^{0.7}$ (Cobb--Douglas), compute $F_K$ and
        $F_L$ at $(K,L)=(100,100)$ and verify $KF_K+LF_L=F$.
  \item Compute $f_x$ and $f_y$ for $f(x,y)=\sin(xy)+\cos(x-y)$.
  \item Compute $\nabla f$ for $f(x,y)=x^2y-xy^2+3x-2y$ at the point
        $(1,-1)$.
  \item For $U(x,y)=x^\alpha y^{1-\alpha}$ (Cobb--Douglas,
        $0<\alpha<1$), derive $\mathrm{MRS}_{xy}=U_x/U_y$ and show it
        equals $\dfrac{\alpha}{1-\alpha}\cdot\dfrac{y}{x}$.
  \item Compute the directional derivative of $f(x,y)=e^{xy}$ at $(1,0)$
        in the direction $\mathbf{v}=(-3,4)$.
  \item Compute $f_x$ and $f_y$ for $f(x,y)=x^y$ (with $x>0$).
        \emph{Hint:} write $x^y=e^{y\ln x}$.
  \item For $F(K,L)=\min\{2K,3L\}$ (Leontief, fixed-proportions, on the
        smooth region away from the kink), compute $F_K$ and $F_L$ in
        each of the two regions $\{2K<3L\}$ and $\{2K>3L\}$.
  \item Compute $\nabla f$ for $f(x,y,z)=xyz+x^2-z^2$ at the point
        $(2,-1,3)$.
  \item Compute the unit vector of steepest \emph{descent} for
        $f(x,y)=x^2+4y^2$ at $(3,1)$, and the corresponding rate of
        decrease.
  \item For $f(x,y)=xe^y-ye^x$, compute $f_x, f_y$.
  \item Compute $f_x$ and $f_y$ for $f(x,y)=\dfrac{x}{x^2+y^2}$.
  \item For $U(x,y)=x+2\sqrt{y}$ (quasi-linear), derive
        $\mathrm{MRS}_{xy}$.
  \item Compute the directional derivative of $f(x,y,z)=x^2+y^2-z^2$ at
        $(2,-1,2)$ in the direction $\mathbf{v}=(2,2,1)$.
  \item Compute $f_x$ and $f_y$ for $f(x,y)=\tan(x+y)+\tan(x-y)$.
  \item For $F(K,L)=(\alpha K+\beta L)^r$ (linear-aggregate,
        $\alpha,\beta,r>0$), compute $F_K, F_L$, and verify
        $KF_K+LF_L=rF$.
  \item Compute $\nabla f$ for $f(x,y)=e^{-x}\sin(y)+e^x\cos(y)$.
  \item Compute the directional derivative of $f(x,y)=\sqrt{x^2+y^2}$ at
        $(3,4)$ in the direction $\mathbf{v}=(1,0)$.
  \item For $U(x,y)=\min\{2x,3y\}$, compute the marginal utilities on
        each smooth region $\{2x<3y\}$ and $\{2x>3y\}$.
  \item Compute $f_x, f_y$ for $f(x,y)=\dfrac{\ln(x+y)}{x-y}$ (assume
        $x>y>0$).
  \item For $F(K,L)=AK^\alpha L^\beta$ general Cobb--Douglas, derive
        $\mathrm{MRTS}_{LK}$ and show it depends only on the ratio $K/L$
        (not on scale).
  \item Compute the unit direction of steepest ascent for
        $f(x,y)=x^2-y^2$ at $(2,-1)$, and the maximum rate.
  \item Compute $f_x$ and $f_y$ for $f(x,y)=\cosh(x+y)+\sinh(x-y)$.
  \item For $f(x,y,z)=e^{x+y+z}$, compute $\nabla f$ at $(0,0,0)$ and the
        directional derivative in $\mathbf{v}=(1,1,1)$.
  \item Compute $f_x$ and $f_y$ for $f(x,y)=(\ln x)(\ln y)$ with
        $x,y>1$.
  \item For $F(K,L)=K+2K^{1/2}L^{1/2}+L$ (which equals
        $(K^{1/2}+L^{1/2})^2$), compute $F_K, F_L$ directly and confirm
        agreement with the CES form.
  \item Compute the directional derivative of $f(x,y)=x^2y+y^3$ at
        $(-1,2)$ in the direction $\mathbf{v}=(4,-3)$.
  \item For $U(x,y)=\alpha\ln x+\beta\ln y$ (log-Cobb--Douglas), derive
        $\mathrm{MRS}_{xy}$ and evaluate at $(x,y)=(\alpha,\beta)$.
  \item Compute $\nabla f$ for $f(x,y)=(x^2+y^2)^{-1/2}$ (Newtonian
        potential in 2D).
  \item Compute $f_x$ and $f_y$ for $f(x,y)=\arctan(xy)$.
  \item For $F(K,L)=(K^\rho+L^\rho)^{1/\rho}$ with $\rho<1$, $\rho\neq0$
        (CES with constant returns), compute $F_K$ and show
        $F_K=(K/F)^{\rho-1}$.
  \item Compute the directional derivative of $f(x,y,z)=xyz$ at
        $(2,3,-1)$ in the direction $\mathbf{v}=(2,-1,2)$.
\end{problems}

\section{Solutions}

\begin{solutions}
  \item $f_x=20x^3-4xy^3$, $f_y=-6x^2y^2+5y^4$.

  \item $F_K=12\cdot\frac14K^{-3/4}L^{3/4}=3(L/K)^{3/4}$;
        $F_L=12\cdot\frac34K^{1/4}L^{-1/4}=9(K/L)^{1/4}$. At $(16,81)$:
        $(L/K)^{3/4}=(81/16)^{3/4}=27/8$ so $F_K=81/8$;
        $(K/L)^{1/4}=(16/81)^{1/4}=2/3$ so $F_L=6$.

  \item $f_x=2xe^{3y}$, $f_y=3x^2e^{3y}$.

  \item $\nabla f=(12x^2y-y^4,\ 4x^3-4xy^3)^T$.

  \item $f_x=2x-4y$, $f_y=-4x+3y^2$. At $(2,1)$: $\nabla f=(0,-5)^T$.
        Normalise: $\|\mathbf{v}\|=\sqrt{25+144}=13$,
        $\mathbf{u}=(5/13,-12/13)^T$. Then
        \[
          D_{\mathbf{u}}f=0\cdot\frac{5}{13}+(-5)\cdot
            \left(-\frac{12}{13}\right)=\frac{60}{13} .
        \]

  \item $F_K=6\cdot\frac23K^{-1/3}L^{1/3}=4(L/K)^{1/3}$,
        $F_L=2(K/L)^{2/3}$. At $(27,8)$: $(L/K)^{1/3}=2/3$ so
        $\mathrm{MP}_K=8/3$; $(K/L)^{2/3}=(27/8)^{2/3}=9/4$ so
        $\mathrm{MP}_L=9/2$.

  \item Chain rule with $u=3x-2y$: $f_x=4(3x-2y)^3\cdot3=12(3x-2y)^3$,
        $f_y=-8(3x-2y)^3$.

  \item $f_x=y^2e^{xy^2}$, $f_y=2xye^{xy^2}$.

  \item $F_K=15K^{-0.4}L^{0.5}$, $F_L=12.5K^{0.6}L^{-0.5}$.

  \item $f_x=2x-z^2e^x$, $f_y=3z$, $f_z=3y-2ze^x$. So
        $\nabla f=(2x-z^2e^x,\ 3z,\ 3y-2ze^x)^T$.

  \item $f_x=e^x\ln(1+y)$, $f_y=e^x/(1+y)$. At $(0,0)$:
        $\nabla f=(0,1)^T$. Normalise: $\|\mathbf{v}\|=\sqrt5$,
        $\mathbf{u}=(1/\sqrt5,2/\sqrt5)^T$. So
        \[
          D_{\mathbf{u}}f=0\cdot\frac{1}{\sqrt5}+1\cdot\frac{2}{\sqrt5}
            =\frac{2\sqrt5}{5} .
        \]

  \item $U_x=2/x$, $U_y=5/y$. $\mathrm{MRS}_{xy}=U_x/U_y
        =(2/x)/(5/y)=2y/(5x)$. At $(3,5)$: $\mathrm{MRS}=10/15=2/3$.

  \item For $x+y\neq0$: $f(x,y)=x-y$, so $f_x=1$, $f_y=-1$.

  \item With $u=x^2y+1$: $f_x=2xy/(x^2y+1)$, $f_y=x^2/(x^2y+1)$.

  \item $U_{x_1}=\frac12x_1^{-1/2}x_2^{1/3}x_3^{1/6}$,
        $U_{x_2}=\frac13x_1^{1/2}x_2^{-2/3}x_3^{1/6}$,
        $U_{x_3}=\frac16x_1^{1/2}x_2^{1/3}x_3^{-5/6}$. Equivalently
        $U_{x_i}=(\alpha_i/x_i)U$ with $\alpha=(1/2,1/3,1/6)$.

  \item $f_x=y$, $f_y=x-2z$, $f_z=-2y+2z$. At $(1,2,3)$:
        $\nabla f=(2,-5,2)^T$. Normalise: $\|\mathbf{v}\|=\sqrt{1+4+4}=3$,
        $\mathbf{u}=(1/3,2/3,2/3)^T$. Then
        $D_{\mathbf{u}}f=(2-10+4)/3=-4/3$.

  \item $F_K=0.25K^{-0.75}L^{0.75}$, $F_L=0.75K^{0.25}L^{-0.25}$. So
        $\mathrm{MRTS}_{LK}=F_L/F_K=3K/L$. At $(16,4)$:
        $\mathrm{MRTS}=48/4=12$.

  \item Expand: $f=x^3-3x^2y^2+xy-3y^3$. Hence $f_x=3x^2-6xy^2+y$,
        $f_y=-6x^2y+x-9y^2$.

  \item With $u=-y/x$: $\partial u/\partial x=y/x^2$,
        $\partial u/\partial y=-1/x$. Product rule:
        \[
          f_x=e^{-y/x}+x\cdot e^{-y/x}\cdot\frac{y}{x^2}
            =e^{-y/x}\left(1+\frac{y}{x}\right), \qquad
          f_y=x\cdot e^{-y/x}\cdot\left(-\frac1x\right)=-e^{-y/x} .
        \]

  \item Let $g=3K^{1/2}+2L^{1/2}$, so $F=g^2$. Then
        \[
          F_K=2g\cdot\frac32K^{-1/2}=\frac{3g}{K^{1/2}}
            =\frac{3(3K^{1/2}+2L^{1/2})}{K^{1/2}}, \qquad
          F_L=2g\cdot L^{-1/2}=\frac{2(3K^{1/2}+2L^{1/2})}{L^{1/2}} .
        \]

  \item With $u=y/x$, $f=\arctan u$, $f'=1/(1+u^2)$:
        \[
          f_x=\frac{1}{1+y^2/x^2}\cdot\left(-\frac{y}{x^2}\right)
            =\frac{-y}{x^2+y^2}, \qquad
          f_y=\frac{x}{x^2+y^2} .
        \]
        So $\nabla f=(-y/(x^2+y^2),\ x/(x^2+y^2))^T$.

  \item $f_K=\frac12K^{-1/2}L^{1/2}$, $f_L=\frac12K^{1/2}L^{-1/2}$. At
        $(4,9)$: $\nabla f=(\frac12\cdot\frac32,\ \frac12\cdot\frac23)^T
        =(3/4,1/3)^T$. Then $\|\nabla f\|^2=9/16+1/9=97/144$, so the
        maximum rate is $\sqrt{97}/12\approx0.821$ in direction
        $\mathbf{u}^*=(9,4)^T/\sqrt{97}$.

  \item $F=(K^{1/2}+L^{1/2})^2$. Replacing $(K,L)\to(\lambda K,\lambda L)$
        gives $\lambda F$, so $r=1$. With $g=K^{1/2}+L^{1/2}$:
        $F_K=2g\cdot\frac12K^{-1/2}=g/K^{1/2}$, $F_L=g/L^{1/2}$. Hence
        \[
          KF_K+LF_L=K^{1/2}g+L^{1/2}g=g(K^{1/2}+L^{1/2})=g^2=F. \checkmark
        \]

  \item Quotient rule, $D=x^2+y^2$:
        \[
          f_x=\frac{2D-(2x+3y)(2x)}{D^2}=\frac{-2x^2-6xy+2y^2}{(x^2+y^2)^2},
          \qquad f_y=\frac{3x^2-4xy-3y^2}{(x^2+y^2)^2} .
        \]

  \item Treat denominator as constant for $f_x$: $f_x=e^x/(1+e^y)$. For
        $f_y$:
        \[
          f_y=e^x\cdot\frac{-e^y}{(1+e^y)^2}=-\frac{e^{x+y}}{(1+e^y)^2} .
        \]

  \item Let $S=\alpha K^\rho+(1-\alpha)L^\rho$, $F=AS^{\nu/\rho}$.
        Treating $K$ constant, $\partial S/\partial L=\rho(1-\alpha)L^{\rho-1}$:
        \[
          F_L=A\cdot\frac{\nu}{\rho}S^{\nu/\rho-1}\cdot\rho(1-\alpha)L^{\rho-1}
            =\nu(1-\alpha)AL^{\rho-1}S^{\nu/\rho-1} .
        \]

  \item $f=\ln(x^2+y^2)$, $f_x=2x/(x^2+y^2)$, $f_y=2y/(x^2+y^2)$. At
        $(3,4)$: $x^2+y^2=25$, $\nabla f=(6/25,8/25)^T$. Normalise:
        $\|\mathbf{v}\|=\sqrt2$, $\mathbf{u}=(-1/\sqrt2,1/\sqrt2)^T$. Then
        \[
          D_{\mathbf{u}}f=\frac{6}{25}\cdot\left(-\frac{1}{\sqrt2}\right)
            +\frac{8}{25}\cdot\frac{1}{\sqrt2}=\frac{2}{25\sqrt2}
            =\frac{\sqrt2}{25} .
        \]

  \item $f_x=2xyz^3-4y^2z$, $f_y=x^2z^3-8xyz$,
        $f_z=3x^2yz^2-4xy^2+5$.

  \item Chain rule, $u=4x^2+9y^2$:
        \[
          f_x=\frac{8x}{2\sqrt u}=\frac{4x}{\sqrt{4x^2+9y^2}}, \qquad
          f_y=\frac{9y}{\sqrt{4x^2+9y^2}} .
        \]

  \item With $g=K^{1/2}+L^{1/2}$, $F=g^4$,
        $F_K=4g^3\cdot\frac12K^{-1/2}=2g^3/K^{1/2}$. At $(4,9)$: $g=5$,
        $K^{1/2}=2$, so $F_K=2\cdot125/2=125$.

  \item Let $r^2=x^2+y^2+z^2$. Then $f=\ln(r^2)$ and $f_{x_i}=2x_i/r^2$.
        So $\nabla f=(2/(x^2+y^2+z^2))(x,y,z)^T$.

  \item Use $\ln\sqrt{u}=\frac12\ln u$: $f=\frac12\ln(x^2+y^2)$. So
        $f_x=x/(x^2+y^2)$, $f_y=y/(x^2+y^2)$.

  \item $\nabla C=(4,9)^T$. With $Q=x^{1/2}y^{1/2}$: $Q_x=\frac12(y/x)^{1/2}$,
        $Q_y=\frac12(x/y)^{1/2}$, so
        $\nabla Q=\frac12\left((y/x)^{1/2},(x/y)^{1/2}\right)^T$. The
        condition $\nabla C=\lambda\nabla Q$ requires
        $4/9=Q_x/Q_y=y/x$, i.e.\ $y=4x/9$. Along this locus
        $\nabla Q\propto(2/3,3/2)^T\propto(4,9)^T=\nabla C$. \checkmark

  \item Chain rule: $f_x=5(1+xy)^4\cdot y=5y(1+xy)^4$,
        $f_y=5x(1+xy)^4$.

  \item $f_x=3x^2-2y$, $f_y=-2x+2y$. At $(1,2)$:
        $\nabla f=(3-4,-2+4)^T=(-1,2)^T$. Normalise: $\|\mathbf{v}\|=5$,
        $\mathbf{u}=(3/5,4/5)^T$. Then
        \[
          D_{\mathbf{u}}f=-1\cdot\frac35+2\cdot\frac45
            =\frac{-3+8}{5}=1 .
        \]

  \item $F_K=10\cdot0.3K^{-0.7}L^{0.7}=3(L/K)^{0.7}$;
        $F_L=7(K/L)^{0.3}$. At $K=L=100$: $F_K=3$, $F_L=7$. And
        $F=10\cdot100^{0.3}\cdot100^{0.7}=10\cdot100=1000$. Then
        $KF_K+LF_L=100\cdot3+100\cdot7=1000=F$. \checkmark (homogeneity
        degree $r=0.3+0.7=1$).

  \item $f_x=y\cos(xy)-\sin(x-y)$, $f_y=x\cos(xy)+\sin(x-y)$.

  \item $f_x=2xy-y^2+3$, $f_y=x^2-2xy-2$. At $(1,-1)$: $f_x=-2-1+3=0$,
        $f_y=1+2-2=1$. So $\nabla f(1,-1)=(0,1)^T$.

  \item $U_x=\alpha x^{\alpha-1}y^{1-\alpha}$,
        $U_y=(1-\alpha)x^\alpha y^{-\alpha}$. So
        \[
          \mathrm{MRS}_{xy}=\frac{U_x}{U_y}
            =\frac{\alpha x^{\alpha-1}y^{1-\alpha}}{(1-\alpha)x^\alpha y^{-\alpha}}
            =\frac{\alpha}{1-\alpha}\cdot\frac{y}{x} . \checkmark
        \]

  \item $f_x=ye^{xy}$, $f_y=xe^{xy}$. At $(1,0)$: $\nabla f=(0,1)^T$.
        Normalise: $\|\mathbf{v}\|=5$, $\mathbf{u}=(-3/5,4/5)^T$. Then
        $D_{\mathbf{u}}f=0\cdot(-3/5)+1\cdot(4/5)=4/5$.

  \item Write $f=e^{y\ln x}$. Then $f_x=e^{y\ln x}\cdot y/x=yx^{y-1}$ and
        $f_y=e^{y\ln x}\cdot\ln x=x^y\ln x$.

  \item On $\{2K<3L\}$ the binding factor is capital: $F=2K$, so
        $F_K=2$, $F_L=0$. On $\{2K>3L\}$ labour binds: $F=3L$, so
        $F_K=0$, $F_L=3$. (Non-differentiable on $\{2K=3L\}$.)

  \item $f_x=yz+2x$, $f_y=xz$, $f_z=xy-2z$. At $(2,-1,3)$:
        $f_x=-3+4=1$, $f_y=6$, $f_z=-2-6=-8$. So
        $\nabla f=(1,6,-8)^T$.

  \item $\nabla f=(2x,8y)^T$. At $(3,1)$: $\nabla f=(6,8)^T$. Steepest
        descent is $-\nabla f/\|\nabla f\|$. $\|\nabla f\|=10$, so unit
        descent direction is $(-3/5,-4/5)^T$, with rate of decrease
        $-\|\nabla f\|=-10$ (i.e.\ $f$ decreases at rate 10).

  \item $f_x=e^y-ye^x$, $f_y=xe^y-e^x$.

  \item Quotient rule with $D=x^2+y^2$:
        \[
          f_x=\frac{D-x\cdot2x}{D^2}=\frac{y^2-x^2}{(x^2+y^2)^2}, \qquad
          f_y=\frac{-x\cdot2y}{D^2}=\frac{-2xy}{(x^2+y^2)^2} .
        \]

  \item $U_x=1$, $U_y=1/\sqrt{y}$. So $\mathrm{MRS}_{xy}=U_x/U_y
        =\sqrt{y}$ (depends only on $y$, quasi-linear utility).

  \item $f_x=2x$, $f_y=2y$, $f_z=-2z$. At $(2,-1,2)$:
        $\nabla f=(4,-2,-4)^T$. Normalise: $\|\mathbf{v}\|=\sqrt{4+4+1}=3$,
        $\mathbf{u}=(2/3,2/3,1/3)^T$. Then
        \[
          D_{\mathbf{u}}f=4\cdot\frac23+(-2)\cdot\frac23+(-4)\cdot\frac13
            =\frac{8-4-4}{3}=0 .
        \]
        (We are moving tangent to the level surface.)

  \item $f_x=\sec^2(x+y)+\sec^2(x-y)$, $f_y=\sec^2(x+y)-\sec^2(x-y)$.

  \item Let $S=\alpha K+\beta L$, $F=S^r$. Then $F_K=rS^{r-1}\alpha$,
        $F_L=rS^{r-1}\beta$. Hence
        \[
          KF_K+LF_L=rS^{r-1}(\alpha K+\beta L)=rS^{r-1}\cdot S=rS^r=rF.
            \checkmark
        \]

  \item $\nabla f=(-e^{-x}\sin y+e^x\cos y,\ e^{-x}\cos y-e^x\sin y)^T$.

  \item $f_x=x/\sqrt{x^2+y^2}$, $f_y=y/\sqrt{x^2+y^2}$. At $(3,4)$:
        $\sqrt{x^2+y^2}=5$, $\nabla f=(3/5,4/5)^T$. The direction
        $\mathbf{v}=(1,0)$ is already unit. So
        $D_{\mathbf{u}}f=3/5\cdot1+4/5\cdot0=3/5$.

  \item On $\{2x<3y\}$: $U=2x$, so $U_x=2$, $U_y=0$. On $\{2x>3y\}$:
        $U=3y$, so $U_x=0$, $U_y=3$. (Non-differentiable on the ray
        $2x=3y$.)

  \item Let $L=\ln(x+y)$, $D=x-y$. By the quotient rule:
        \[
          f_x=\frac{(1/(x+y))(x-y)-\ln(x+y)\cdot1}{(x-y)^2}
            =\frac{1/(x+y)\cdot(x-y)-\ln(x+y)}{(x-y)^2} .
        \]
        Simplifying
        $f_x=\dfrac{(x-y)/(x+y)-\ln(x+y)}{(x-y)^2}$. Similarly,
        \[
          f_y=\frac{(1/(x+y))(x-y)-\ln(x+y)\cdot(-1)}{(x-y)^2}
            =\frac{(x-y)/(x+y)+\ln(x+y)}{(x-y)^2} .
        \]

  \item $F_K=A\alpha K^{\alpha-1}L^\beta$, $F_L=A\beta K^\alpha L^{\beta-1}$.
        So
        \[
          \mathrm{MRTS}_{LK}=\frac{F_L}{F_K}=\frac{\beta}{\alpha}\cdot
            \frac{K}{L} .
        \]
        This depends only on the ratio $K/L$ and the exponent ratio
        $\beta/\alpha$: it is scale-invariant.

  \item $\nabla f=(2x,-2y)^T$. At $(2,-1)$: $\nabla f=(4,2)^T$. Direction
        of steepest ascent: $\mathbf{u}^*=(4,2)/\sqrt{20}=(2,1)/\sqrt5$.
        Maximum rate $=\|\nabla f\|=\sqrt{20}=2\sqrt5$.

  \item $f_x=\sinh(x+y)+\cosh(x-y)$, $f_y=\sinh(x+y)-\cosh(x-y)$.

  \item $\nabla f=(e^{x+y+z},e^{x+y+z},e^{x+y+z})^T$. At $(0,0,0)$:
        $\nabla f=(1,1,1)^T$. Normalise $\mathbf{v}=(1,1,1)$:
        $\|\mathbf{v}\|=\sqrt3$, $\mathbf{u}=(1,1,1)/\sqrt3$. Then
        $D_{\mathbf{u}}f=(1+1+1)/\sqrt3=\sqrt3$.

  \item $f_x=(\ln y)/x$, $f_y=(\ln x)/y$.

  \item $F_K=1+L^{1/2}/K^{1/2}$, $F_L=1+K^{1/2}/L^{1/2}$. From
        $F=(K^{1/2}+L^{1/2})^2$ with $g=K^{1/2}+L^{1/2}$:
        $F_K=g/K^{1/2}=(K^{1/2}+L^{1/2})/K^{1/2}=1+L^{1/2}/K^{1/2}$.
        \checkmark Same for $F_L$.

  \item $f_x=2xy$, $f_y=x^2+3y^2$. At $(-1,2)$: $\nabla f=(-4,1+12)^T
        =(-4,13)^T$. Normalise: $\|\mathbf{v}\|=5$,
        $\mathbf{u}=(4/5,-3/5)^T$. Then
        \[
          D_{\mathbf{u}}f=-4\cdot\frac45+13\cdot\left(-\frac35\right)
            =-\frac{16+39}{5}=-11 .
        \]

  \item $U_x=\alpha/x$, $U_y=\beta/y$.
        $\mathrm{MRS}_{xy}=\dfrac{\alpha/x}{\beta/y}=\dfrac{\alpha y}{\beta x}$.
        At $(\alpha,\beta)$: $\mathrm{MRS}=(\alpha\cdot\beta)/(\beta\cdot\alpha)=1$.

  \item $f=(x^2+y^2)^{-1/2}$. Chain rule:
        \[
          f_x=-\frac12(x^2+y^2)^{-3/2}\cdot2x=\frac{-x}{(x^2+y^2)^{3/2}},
          \qquad f_y=\frac{-y}{(x^2+y^2)^{3/2}} .
        \]
        So $\nabla f=-(x,y)^T/(x^2+y^2)^{3/2}$.

  \item Chain rule with $u=xy$:
        \[
          f_x=\frac{y}{1+x^2y^2}, \qquad f_y=\frac{x}{1+x^2y^2} .
        \]

  \item Let $S=K^\rho+L^\rho$, $F=S^{1/\rho}$. Then
        $F_K=\frac1\rho S^{1/\rho-1}\cdot\rho K^{\rho-1}
        =K^{\rho-1}S^{1/\rho-1}$. Since $S^{1/\rho}=F$, we have $S=F^\rho$
        and so $S^{1/\rho-1}=F^{1-\rho}\cdot F^{-1}\cdot F=F^{1-\rho}$.
        Thus $F_K=K^{\rho-1}F^{1-\rho}=(F/K)^{1-\rho}=(K/F)^{\rho-1}$.
        \checkmark

  \item $f_x=yz$, $f_y=xz$, $f_z=xy$. At $(2,3,-1)$:
        $\nabla f=(3\cdot(-1),2\cdot(-1),2\cdot3)^T=(-3,-2,6)^T$.
        Normalise $\mathbf{v}=(2,-1,2)$: $\|\mathbf{v}\|=3$,
        $\mathbf{u}=(2/3,-1/3,2/3)^T$. Then
        \[
          D_{\mathbf{u}}f=-3\cdot\frac23+(-2)\cdot\left(-\frac13\right)
            +6\cdot\frac23=-2+\frac23+4=2+\frac23=\frac83 .
        \]
\end{solutions}

\section{Additional Problems}

\begin{problems}
  \item Compute $f_x$ and $f_y$ for $f(x,y)=3x^5-2x^3y^2+4y^4$.
  \item For $F(K,L)=20\,K^{1/3}L^{2/3}$, find $F_K$ and $F_L$ at
        $(K,L)=(8,27)$.
  \item Compute $f_x$ and $f_y$ for $f(x,y)=e^{xy}$.
  \item Compute $\nabla f$ for $f(x,y)=2x^2y-3xy^3+5$.
  \item Compute the directional derivative of $f(x,y)=x^2+2xy-y^2$ at
        $(1,2)$ in the direction $\mathbf{v}=(3,4)$.
  \item For $F(K,L)=8K^{3/4}L^{1/4}$, compute $\mathrm{MP}_K$ and
        $\mathrm{MP}_L$ at $(K,L)=(16,16)$.
  \item Compute $f_x$ and $f_y$ for $f(x,y)=\ln(2x+3y)$.
  \item For $U(x,y)=3\ln(x)+4\ln(y)$, derive $\mathrm{MRS}_{xy}$ and
        evaluate at $(2,3)$.
  \item Compute $f_x$ and $f_y$ for $f(x,y)=(x+y^2)^4$.
  \item Compute $\nabla f$ for $f(x,y,z)=xy+yz+zx$.
  \item For a CES production $F(K,L)=5(K^{1/3}+L^{1/3})^3$, find $F_K$.
  \item Compute the directional derivative of $f(x,y)=x^2y$ at $(2,1)$
        in the direction $\mathbf{v}=(1,-1)$.
  \item For $U(x,y)=x^\alpha y^\beta$ (general Cobb--Douglas,
        $\alpha,\beta>0$), derive $\mathrm{MRS}_{xy}$ and show it depends
        only on the ratio $y/x$.
  \item Compute $f_x$ and $f_y$ for $f(x,y)=\dfrac{x-y}{x+y}$.
  \item For $f(K,L)=K^{1/2}L^{1/2}$ at $(K,L)=(16,9)$, find the unit
        direction of steepest ascent and the maximum rate of increase.
  \item Classify the open positive orthant
        $Q=\{(x,y)\in\R^2 : x>0,\ y>0\}$: is it open, closed, or
        neither? Is it bounded? Is it compact?
  \item Classify the budget set
        $B=\{(x,y)\in\R^2 : 2x+3y\le12,\ x\ge0,\ y\ge0\}$: open,
        closed, or neither? Bounded? Compact?
  \item Classify $M=\{(x,y)\in\R^2 : xy\ge1,\ x>0\}$: open, closed, or
        neither? Bounded? Compact? \emph{Hint:} check whether a
        sequence in $M$ can converge to a point with $x=0$.
  \item Classify the open disc
        $D=\{(x,y)\in\R^2 : x^2+y^2<4\}$: open, closed, or neither?
        Bounded? Compact?
  \item Classify (a) the half-open strip
        $S=\{(x,y)\in\R^2 : 0\le x<1\}$ and (b) the finite set
        $E=\{(0,0),\ (1,2),\ (3,1)\}$: open, closed, or neither?
        Bounded? Compact?
  \item Does the Weierstrass theorem guarantee a maximiser for the
        continuous function $u(x,y)=x^{1/2}y^{1/2}$ (a) on the budget
        set $B=\{(x,y): 2x+3y\le12,\ x\ge0,\ y\ge0\}$, (b) on the open
        positive orthant $Q=\{(x,y): x>0,\ y>0\}$? Justify each
        answer.
\end{problems}

\section{Additional Solutions}

\begin{solutions}
  \item $f_x=15x^4-6x^2y^2$, $f_y=-4x^3y+16y^3$.

  \item $F_K=20\cdot\frac13K^{-2/3}L^{2/3}=\frac{20}{3}(L/K)^{2/3}$;
        $F_L=20\cdot\frac23K^{1/3}L^{-1/3}=\frac{40}{3}(K/L)^{1/3}$. At
        $(8,27)$: $(L/K)^{2/3}=(27/8)^{2/3}=9/4$, so $F_K=15$;
        $(K/L)^{1/3}=(8/27)^{1/3}=2/3$, so
        $F_L=\frac{40}{3}\cdot\frac23=\frac{80}{9}$.

  \item $f_x=ye^{xy}$, $f_y=xe^{xy}$.

  \item $\nabla f=(4xy-3y^3,\ 2x^2-9xy^2)^T$.

  \item $f_x=2x+2y$, $f_y=2x-2y$. At $(1,2)$: $\nabla f=(6,-2)^T$.
        Normalise: $\|\mathbf{v}\|=5$, $\mathbf{u}=(3/5,4/5)^T$. Then
        \[
          D_{\mathbf{u}}f=6\cdot\frac35+(-2)\cdot\frac45
            =\frac{18-8}{5}=2 .
        \]

  \item $F_K=8\cdot\frac34K^{-1/4}L^{1/4}=6(L/K)^{1/4}$,
        $F_L=2(K/L)^{3/4}$. At $(16,16)$: $\mathrm{MP}_K=6$,
        $\mathrm{MP}_L=2$.

  \item $f_x=\dfrac{2}{2x+3y}$, $f_y=\dfrac{3}{2x+3y}$.

  \item $U_x=3/x$, $U_y=4/y$.
        $\mathrm{MRS}_{xy}=U_x/U_y=\dfrac{3y}{4x}$. At $(2,3)$:
        $\mathrm{MRS}=\dfrac98$.

  \item Chain: $f_x=4(x+y^2)^3$,
        $f_y=4(x+y^2)^3\cdot2y=8y(x+y^2)^3$.

  \item $f_x=y+z$, $f_y=x+z$, $f_z=x+y$. So
        $\nabla f=(y+z,\ x+z,\ x+y)^T$.

  \item Let $g=K^{1/3}+L^{1/3}$, so $F=5g^3$. Then
        $F_K=5\cdot3g^2\cdot\frac13K^{-2/3}=5g^2K^{-2/3}
        =\dfrac{5(K^{1/3}+L^{1/3})^2}{K^{2/3}}$.

  \item $f_x=2xy$, $f_y=x^2$. At $(2,1)$: $\nabla f=(4,4)^T$. Normalise:
        $\|\mathbf{v}\|=\sqrt2$, $\mathbf{u}=(1/\sqrt2,-1/\sqrt2)^T$.
        Then
        \[
          D_{\mathbf{u}}f=4\cdot\frac{1}{\sqrt2}
            +4\cdot\left(-\frac{1}{\sqrt2}\right)=0 .
        \]
        (Tangent to the level curve.)

  \item $U_x=\alpha x^{\alpha-1}y^\beta$,
        $U_y=\beta x^\alpha y^{\beta-1}$. So
        \[
          \mathrm{MRS}_{xy}=\frac{U_x}{U_y}
            =\frac{\alpha}{\beta}\cdot\frac{y}{x} .
        \]
        Depends only on $y/x$ (and the exponent ratio), not on the
        scale. \checkmark

  \item Quotient rule with $D=x+y$:
        \[
          f_x=\frac{D-(x-y)}{D^2}=\frac{2y}{(x+y)^2}, \qquad
          f_y=\frac{-D-(x-y)}{D^2}=\frac{-2x}{(x+y)^2} .
        \]

  \item $f_K=\frac12K^{-1/2}L^{1/2}$, $f_L=\frac12K^{1/2}L^{-1/2}$. At
        $(16,9)$: $f_K=\frac12\cdot\frac34=3/8$;
        $f_L=\frac12\cdot\frac43=2/3$. So $\nabla f=(3/8,2/3)^T$.
        Magnitude:
        $\|\nabla f\|=\sqrt{9/64+4/9}=\sqrt{81/576+256/576}
        =\sqrt{337/576}=\sqrt{337}/24\approx0.765$. Unit direction of
        steepest ascent:
        $\mathbf{u}^*=(3/8,2/3)^T/(\sqrt{337}/24)=(9,16)^T/\sqrt{337}$.

  \item \textbf{Open.} For any $(x,y)\in Q$ the ball of radius
        $r=\min\{x,y\}>0$ around $(x,y)$ stays inside $Q$, so every
        point is interior. \textbf{Not closed:} the boundary points
        $(x,0)$ and $(0,y)$ with $x,y\ge0$ (the coordinate half-axes,
        including the origin) do not belong to $Q$. \textbf{Not
        bounded:} $Q$ contains $(t,t)$ for every $t>0$, and
        $\|(t,t)\|=t\sqrt2\to\infty$. Hence $Q$ is open, unbounded and
        \textbf{not compact}.

  \item \textbf{Closed.} $B$ is the intersection of the three sets
        $\{2x+3y\le12\}$, $\{x\ge0\}$ and $\{y\ge0\}$, each defined by
        a non-strict inequality on a continuous function; each is
        closed, and an intersection of closed sets is closed (every
        boundary point of $B$ satisfies all three inequalities, so it
        lies in $B$). \textbf{Not open:} $(0,0)\in B$, but every ball
        around the origin contains points with negative coordinates,
        which are outside $B$; so $(0,0)$ is not interior.
        \textbf{Bounded:} from the constraints, $0\le x\le6$ and
        $0\le y\le4$, so $B\subseteq B_8(\mathbf{0})$. Closed and
        bounded in $\R^2$: $B$ is \textbf{compact}.

  \item First note the constraint $x>0$ together with $xy\ge1$ forces
        $y\ge1/x>0$. \textbf{Closed:} suppose
        $(x_n,y_n)\in M$ with $(x_n,y_n)\to(a,b)$. If $a=0$, then
        $y_n\to b$ finite would give $x_ny_n\to0<1$, contradicting
        $x_ny_n\ge1$; so $a>0$, and passing to the limit in
        $x_ny_n\ge1$ gives $ab\ge1$, hence $(a,b)\in M$. $M$ contains
        all its limit (boundary) points, so it is closed: even
        though it is written with a strict inequality, that inequality
        is never active. \textbf{Not open:} $(1,1)\in M$ lies on the
        hyperbola $xy=1$, and every ball around it contains points
        with $xy<1$. \textbf{Not bounded:} $(t,t)\in M$ for all
        $t\ge1$. Hence $M$ is closed, unbounded and \textbf{not
        compact}.

  \item \textbf{Open.} For $(x,y)\in D$ put
        $r=2-\sqrt{x^2+y^2}>0$; the ball $B_r(x,y)$ lies in $D$ by the
        triangle inequality. \textbf{Not closed:} the boundary circle
        $\{x^2+y^2=4\}$ is not contained in $D$ (e.g.\ $(2,0)\notin
        D$). \textbf{Bounded:} $D\subseteq B_3(\mathbf{0})$. Bounded
        but not closed: $D$ is \textbf{not compact}.

  \item (a) \textbf{Neither open nor closed.} The points $(0,y)$
        belong to $S$ but are boundary points (every ball around them
        contains points with $x<0$), so $S$ is not open; the boundary
        points $(1,y)$ do not belong to $S$, so $S$ is not closed.
        \textbf{Unbounded:} $(0,t)\in S$ for all $t$. Hence
        \textbf{not compact}. (b) \textbf{Closed:} around any point
        not in $E$ there is a ball avoiding all three points (a
        finite set always contains all its boundary points:
        indeed $E$ is its own boundary). \textbf{Not open:} no ball
        around $(0,0)$ lies inside $E$. \textbf{Bounded:}
        $E\subseteq B_4(\mathbf{0})$. Closed and bounded: $E$ is
        \textbf{compact}.

  \item (a) \textbf{Yes.} $B$ is compact (closed and bounded, see the
        budget-set problem above) and $u$ is continuous on $B$, so by
        the Weierstrass theorem $u$ attains both a maximum and a
        minimum on $B$; a maximiser exists. (b) \textbf{No.} $Q$ is
        not compact (it is neither closed nor bounded), so Weierstrass
        does not apply: and indeed no maximiser exists:
        $u(t,t)=t\to\infty$ as $t\to\infty$, so $u$ is unbounded above
        on $Q$.
\end{solutions}

\chapter{Curve Sketching: The Complete Discussion}

Everything needed to draw a curve has now been assembled in pieces:
domains and asymptotes in the functions volume, limits in the analysis
volume, the first derivative for monotonicity, the second for curvature.
What is missing is the habit of running all of it in one fixed order and
finishing with a picture. That is the \de{Kurvendiskussion}, and it is worth
drilling as a single routine for two reasons.

First, examinations ask for it directly. Second, and more usefully, the
routine is a self-checking mechanism. A sign table that says the function
falls on $(0,2)$ and a computed local maximum at $x=2$ cannot both be right;
the contradiction surfaces before the picture is drawn. A curve sketch is
the cheapest way there is to catch an algebra slip.

The chapter runs the protocol on rational, exponential, logarithmic and root
functions, and closes with the discussion an economist actually performs:
locating the inflection of a total cost curve, the minimum of average cost,
and the output that maximises profit --- all of them the same routine
applied to a curve with a name.

\section{Core Rules}

\begin{rulebox}{The protocol}
Always in this order, because each step feeds the next:
\begin{enumerate}[leftmargin=1.8em,itemsep=0.35ex]
  \item \textbf{Domain.} Denominators, even roots, logarithms. Note which
        points are excluded --- they are the candidates for poles.
  \item \textbf{Symmetry.} $f(-x)=f(x)$ even, $f(-x)=-f(x)$ odd. Either one
        halves the work.
  \item \textbf{Intercepts and roots.} $f(0)$, and the solutions of
        $f(x)=0$ with their multiplicities.
  \item \textbf{Behaviour at the edges.} Limits at $\pm\infty$ and at every
        excluded point, one side at a time. This produces the
        \textbf{asymptotes}: vertical at poles, horizontal or oblique at
        infinity.
  \item \textbf{First derivative.} Solve $f'(x)=0$ for the critical points,
        build a \textbf{sign table} for $f'$, and read off intervals of
        increase and decrease and the type of each extremum.
  \item \textbf{Second derivative.} Solve $f''(x)=0$, build a sign table
        for $f''$, and read off convex and concave intervals and the
        \textbf{inflection points}.
  \item \textbf{Value table} for any remaining doubt, then \textbf{draw}:
        asymptotes first as dashed guides, then the special points, then
        join them respecting monotonicity and curvature.
\end{enumerate}
\end{rulebox}

\begin{rulebox}{Sign tables}
A sign table is the bookkeeping device that prevents guesswork. List every
point where the expression can change sign --- its roots and its poles ---
in increasing order; they cut the domain into intervals. On each interval
the sign is constant, so \emph{one} test value settles it. For
$f'(x)=3x(x-2)$:
\[
  \begin{array}{c|ccccc}
    x & (-\infty,0) & 0 & (0,2) & 2 & (2,\infty)\\\hline
    f' & + & 0 & - & 0 & +\\
    f & \nearrow & \text{max} & \searrow & \text{min} & \nearrow
  \end{array}
\]
A sign change $+\to-$ is a local maximum, $-\to+$ a local minimum, and
\emph{no} change means no extremum, however much the derivative vanishes
(as at $x=0$ for $f=x^{3}$).
\end{rulebox}

\begin{rulebox}{First- and second-derivative tests}
Let $f'(x_{0})=0$.
\begin{itemize}[leftmargin=1.6em,itemsep=0.3ex]
  \item \textbf{First-derivative test:} the sign change of $f'$ across
        $x_{0}$ decides the type. Always applicable, and the only reliable
        test when $f''(x_{0})=0$.
  \item \textbf{Second-derivative test:} $f''(x_{0})>0$ gives a local
        minimum, $f''(x_{0})<0$ a local maximum. If $f''(x_{0})=0$ the test
        is \textbf{silent} --- fall back on the sign table. Compare
        $x^{4}$ (minimum), $-x^{4}$ (maximum) and $x^{3}$ (neither), all
        three with $f'=f''=0$ at the origin.
  \item \textbf{Inflection} at $x_{0}$ requires $f''$ to \emph{change
        sign} there. $f''(x_{0})=0$ alone is not enough: $f=x^{4}$ has
        $f''(0)=0$ and no inflection.
\end{itemize}
\end{rulebox}

\begin{rulebox}{Asymptotes, recalled}
For $f=p/q$ with $p,q$ polynomials, after cancelling common factors:
vertical asymptote at each zero of $q$; horizontal $y=0$ if
$\deg p<\deg q$; horizontal $y=$ ratio of leading coefficients if
$\deg p=\deg q$; and if $\deg p=\deg q+1$, polynomial division gives
\[
  f(x)=\underbrace{mx+c}_{\text{oblique asymptote}}+\frac{\text{remainder}}{q(x)} .
\]
For non-rational functions the same information comes from limits directly:
$\lim_{x\to\pm\infty}f(x)$, and for a possible oblique asymptote
$m=\lim_{x\to\infty}f(x)/x$ followed by $c=\lim_{x\to\infty}(f(x)-mx)$.
\end{rulebox}

\begin{gallery}
\hspace*{\fill}%
\begin{tikzpicture}
  \begin{axis}[fnaxis, width=8.4cm, height=7.4cm,
      xmin=-4.4, xmax=5.4, ymin=-6.4, ymax=10.4,
      xtick={-4,-2,2,4}, ytick={-4,4,8}]
    \addplot[gray!55, dashed, thin, samples=2, forget plot]
      coordinates {(1,-6.2)(1,10.2)};
    \addplot[gray!55, dashed, thin, samples=2, domain=-4.2:5.2, forget plot]{x+1};
    \addplot[domain=-4.2:0.86, samples=201]{x^2/(x-1)};
    \addplot[domain=1.14:5.2, samples=201]{x^2/(x-1)};
    \addplot[only marks, mark=*, mark size=1.3pt, black, forget plot]
      coordinates {(0,0) (2,4)};
    \node[font=\tiny, anchor=north east] at (axis cs:-0.1,-0.15) {$(0,0)$};
    \node[font=\tiny, anchor=west] at (axis cs:2.15,4) {$(2,4)$};
  \end{axis}
\end{tikzpicture}%
\hspace*{\fill}
\gcap{The finished discussion of $f(x)=x^{2}/(x-1)$: pole $x=1$, oblique
asymptote $y=x+1$, local maximum $(0,0)$ below the local minimum $(2,4)$.}
\end{gallery}

\begin{rulebox}{The economist's curve discussion}
Three facts fall straight out of the protocol applied to a total cost
function $C(q)$, and they are worth knowing as results:
\begin{itemize}[leftmargin=1.6em,itemsep=0.4ex]
  \item Marginal cost is $MC=C'$. The \textbf{inflection point of $C$} is
        exactly where $MC$ is smallest, since $MC'=C''$ changes sign there.
        The S-shaped cost curve turns from concave to convex at its most
        efficient scale.
  \item Average cost is $AC=C(q)/q$. Differentiating,
        \[
          AC'=\frac{C'(q)q-C(q)}{q^{2}}=\frac{MC-AC}{q},
        \]
        so $AC'=0$ exactly when $MC=AC$: \textbf{marginal cost cuts average
        cost at the minimum of average cost}. That is a theorem, not a
        coincidence of a particular example.
  \item Profit at price $p$ is $\pi=pq-C(q)$, so $\pi'=p-MC$ and the
        optimum has $p=MC$, confirmed as a maximum by $\pi''=-C''<0$,
        i.e.\ by cost being convex there.
\end{itemize}
\end{rulebox}

\section{Worked Example}

\begin{worked}{The complete discussion of $f(x)=\dfrac{x^{2}}{x-1}$}
\wstep{Step 1. Domain, and the reason to find it first.}
The denominator vanishes at $x=1$, so
$D_{f}=\R\setminus\{1\}$. Every later step --- sign tables, monotonicity,
convexity --- must be stated \emph{per interval}, and $x=1$ is the wall
between them.

\wstep{Step 2. Symmetry and roots.}
$f(-x)=\frac{x^{2}}{-x-1}$, which is neither $f(x)$ nor $-f(x)$: no
symmetry. The numerator vanishes at $x=0$ with multiplicity two, so the
curve \textbf{touches} the axis at the origin rather than crossing it, and
$f(0)=0$.

\wstep{Step 3. Behaviour at the pole, from both sides.}
Near $x=1$ the numerator is close to $1>0$, so the sign of $f$ is the sign
of $x-1$:
\[
  x\to1^{-}:\ \frac{1^{+}}{0^{-}}\to-\infty,
  \qquad
  x\to1^{+}:\ \frac{1^{+}}{0^{+}}\to+\infty .
\]
\wtrap{Take one-sided limits at every pole. Writing ``$f\to\infty$ at
$x=1$'' loses the sign information and makes the sketch impossible to draw
correctly.}

\wstep{Step 4. Behaviour at infinity: divide, do not guess.}
Numerator degree $2$ exceeds denominator degree $1$ by exactly one, so there
is an \textbf{oblique} asymptote. Polynomial division gives it:
\[
  f(x)=x+1+\frac{1}{x-1} .
\]
So $y=x+1$ is the asymptote, and the remainder $\frac{1}{x-1}$ says how the
curve sits relative to it: \emph{above} for $x>1$, \emph{below} for $x<1$.
\wtrap{Horizontal asymptote when the degrees are equal or the numerator's is
smaller; oblique when the numerator's degree is exactly one larger; neither
when it is larger by two or more. And the division is the whole method ---
guessing $y=x$ from the leading behaviour misses the $+1$.}

\wstep{Step 5. First derivative and critical points.}
Differentiating the divided form is easiest:
\[
  f'(x)=1-\frac{1}{(x-1)^{2}}=\frac{(x-1)^{2}-1}{(x-1)^{2}}=\frac{x(x-2)}{(x-1)^{2}} .
\]
Zeros at $x=0$ and $x=2$; the denominator is positive wherever it is
defined, so the sign of $f'$ is the sign of $x(x-2)$.
\begin{center}\footnotesize
\begin{tabular}{lcccc}
  \hline
  interval & $x<0$ & $0<x<1$ & $1<x<2$ & $x>2$\\
  \hline
  $f'$ & $+$ & $-$ & $-$ & $+$\\
  $f$ & rising & falling & falling & rising\\
  \hline
\end{tabular}
\end{center}

\wstep{Step 6. Classify, and notice something strange.}
$f''(x)=\dfrac{2}{(x-1)^{3}}$, so $f''(0)=-2<0$ and $f''(2)=2>0$:
\[
  \text{local maximum at }(0,0),
  \qquad
  \text{local minimum at }(2,4).
\]
The local maximum lies \textbf{below} the local minimum. That is impossible
for a continuous function --- and possible here only because the pole at
$x=1$ separates them. The two branches never see each other.
\wtrap{Local extrema are local. A local maximum of $0$ coexisting with a
local minimum of $4$ is not an arithmetic error, and neither point is a
global extremum: $f$ is unbounded in both directions.}

\wstep{Step 7. Convexity, and the absent inflection.}
$f''=\frac{2}{(x-1)^{3}}$ is never zero, so there is \textbf{no inflection
point}. It is negative for $x<1$ and positive for $x>1$:
concave on $(-\infty,1)$, convex on $(1,\infty)$. The change of curvature
happens at the pole, where the function does not exist --- so it is not an
inflection.
\wtrap{A change of convexity across a point of the domain is an inflection;
across a gap in the domain it is nothing at all. Check that the candidate
point is in the domain before naming it.}

\wstep{Step 8. Draw it.}
\begin{center}
\begin{tikzpicture}
\begin{axis}[fnaxis, width=8.4cm, height=6.0cm,
             xmin=-4.2, xmax=6.2, ymin=-9, ymax=13,
             xtick={-3,-1,1,2,4}, ytick={4,8},
             restrict y to domain=-40:40]
  \addplot[domain=-4:0.88, samples=200] {x^2/(x-1)};
  \addplot[domain=1.12:6, samples=200] {x^2/(x-1)};
  \addplot[accent!40, thick, dashed, samples=2, domain=-4:6, forget plot]
    {x+1};
  \addplot[gray!55, thin, dashed, samples=2, forget plot]
    coordinates {(1,-9) (1,13)};
  \addplot[only marks, mark=*, mark size=1.2pt, accent!70, forget plot]
    coordinates {(0,0) (2,4)};
\end{axis}
\end{tikzpicture}
\end{center}
Everything checks against the figure: the touch at the origin, the two
branches, the pole at $x=1$, the asymptote $y=x+1$ approached from below on
the left and from above on the right, and the maximum sitting lower than the
minimum.

\wstep{Step 9. The protocol, in the order that saves work.}
\begin{center}\footnotesize
\begin{tabular}{ll}
  \hline
  1. domain & 6. sign table for $f'$\\
  2. symmetry & 7. extrema, classified by $f''$\\
  3. roots and intercepts & 8. $f''$, convexity, inflections\\
  4. poles, one-sided limits & 9. sketch\\
  5. asymptotes (divide!) & 10. check the sketch against 1--8\\
  \hline
\end{tabular}
\end{center}
Step~10 is not optional. A sketch that contradicts the sign table means one
of them is wrong, and finding out which takes a minute --- far less than
losing the marks.
\end{worked}

\section{Practice Problems}

In problems 1--8 give the intervals of increase and decrease and all local
extrema, using a sign table.

\begin{problems}
  \item $f(x)=x^{3}-12x$
  \item $f(x)=x^{4}-8x^{2}$
  \item $f(x)=x\mathrm{e}^{-x}$
  \item $f(x)=\dfrac{x}{x^{2}+1}$
  \item $f(x)=x-\ln x$ on $x>0$
  \item $f(x)=x^{3}$ --- and explain why $f'(0)=0$ yields no extremum.
  \item $f(x)=\dfrac{x^{2}}{x-1}$
  \item $f(x)=\sqrt{x}-x$ on $x\ge0$
\end{problems}

In problems 9--14 give the convex and concave intervals and all inflection
points.

\begin{problems}[resume]
  \item $f(x)=x^{3}-3x^{2}$
  \item $f(x)=x^{4}-4x^{2}$
  \item $f(x)=\mathrm{e}^{-x^{2}}$
  \item $f(x)=x^{4}$ --- and explain why $f''(0)=0$ gives no inflection.
  \item $f(x)=\dfrac{\ln x}{x}$ on $x>0$
  \item $f(x)=x^{2}\mathrm{e}^{-x}$
\end{problems}

In problems 15--20 give every asymptote.

\begin{problems}[resume]
  \item $f(x)=\dfrac{x-1}{x+2}$
  \item $f(x)=\dfrac{1}{x^{2}-4}$
  \item $f(x)=\dfrac{x^{2}}{x-1}$
  \item $f(x)=\dfrac{x^{3}}{x^{2}+1}$
  \item $f(x)=\dfrac{\ln x}{x}$
  \item $f(x)=x+\mathrm{e}^{-x}$
\end{problems}

Problems 21--32 ask for a \emph{complete} discussion: domain, symmetry,
intercepts, asymptotes and edge behaviour, monotonicity and extrema,
curvature and inflections, and a sketch.

\begin{problems}[resume]
  \item $f(x)=x^{3}-3x^{2}$
  \item $f(x)=x^{4}-4x^{2}$
  \item $f(x)=\dfrac{x}{x^{2}+1}$
  \item $f(x)=\dfrac{x^{2}}{x-1}$
  \item $f(x)=\mathrm{e}^{-x^{2}}$
  \item $f(x)=x\mathrm{e}^{-x}$
  \item $f(x)=\dfrac{\ln x}{x}$
  \item $f(x)=\dfrac{x-1}{x+2}$
  \item $f(x)=x^{2}\mathrm{e}^{-x}$
  \item $f(x)=\dfrac{1}{x^{2}-4}$
  \item $f(x)=\sqrt{x}-x$
  \item $f(x)=x^{3}-3x+2$
\end{problems}

\begin{problems}[resume]
  \item A function satisfies $f'>0$ and $f''<0$ throughout an interval.
        Describe the shape of its graph there in one sentence.
  \item The graph of $f'$ is a parabola opening upwards with roots at
        $x=-1$ and $x=3$. What does $f$ do, and where are its extrema?
  \item The graph of $f'$ is a straight line through the origin with
        positive slope. What is $f$, up to a constant, and what is its
        shape?
  \item $f$ is differentiable with $f(1)=2$, $f'(1)=0$ and $f''(1)=-3$.
        Classify $x=1$ and give the tangent line there.
  \item Sketch a single curve with all of: domain $\R\setminus\{2\}$, a
        pole at $x=2$, horizontal asymptote $y=1$, a local minimum at
        $(0,0)$, and convex on $(2,\infty)$.
  \item Why can a function with a vertical asymptote at $x=a$ have a local
        maximum lying below a local minimum? Illustrate with problem 24.
\end{problems}

\begin{problems}[resume]
  \item For $C(q)=q^{3}-6q^{2}+15q$ on $q\ge0$: show $C$ is increasing
        throughout, find its inflection point, and show that marginal cost
        is minimal exactly there.
  \item For the same $C$: find the output minimising average cost and
        verify that $MC=AC$ at that output.
  \item With the same $C$ and price $p=15$, find the profit-maximising
        output and the maximal profit.
  \item For $R(q)=q(200-4q)$: discuss the curve on its economic domain and
        find the revenue-maximising quantity.
  \item For $C(q)=0.5q^{2}+2q+8$: find the minimum of $AC$ and confirm
        $MC=AC$ there.
  \item The adoption path is $y(t)=\dfrac{100}{1+\mathrm{e}^{-0.5t}}$.
        Find its asymptotes, show it is increasing, and locate its
        inflection point.
  \item Explain, using $AC'=(MC-AC)/q$, why average cost must be falling
        wherever marginal cost lies below it.
\end{problems}

\section{Solutions}

\begin{solutions}
  \item $f'=3x^{2}-12=3(x-2)(x+2)$: $+$ on $(-\infty,-2)$, $-$ on
        $(-2,2)$, $+$ on $(2,\infty)$. Local maximum $(-2,16)$, local
        minimum $(2,-16)$.
  \item $f'=4x^{3}-16x=4x(x-2)(x+2)$: $-$, $+$, $-$, $+$ on the four
        intervals cut by $-2,0,2$. Local minima $(\pm2,-16)$, local maximum
        $(0,0)$.
  \item $f'=(1-x)\mathrm{e}^{-x}$: positive for $x<1$, negative for $x>1$.
        Single local (and global) maximum $(1,1/\mathrm{e})$.
  \item $f'=\dfrac{1-x^{2}}{(x^{2}+1)^{2}}$: negative, positive, negative
        on $(-\infty,-1)$, $(-1,1)$, $(1,\infty)$. Local minimum
        $(-1,-\tfrac12)$, local maximum $(1,\tfrac12)$.
  \item $f'=1-\dfrac1x=\dfrac{x-1}{x}$: negative on $(0,1)$, positive on
        $(1,\infty)$. Local (global) minimum $(1,1)$.
  \item $f'=3x^{2}\ge0$ with equality only at $0$: the sign of $f'$ never
        changes, so $f$ increases throughout and $x=0$ is a
        \emph{saddle-like} stationary point, not an extremum. A vanishing
        derivative is necessary, never sufficient.
  \item $f'=\dfrac{x(x-2)}{(x-1)^{2}}$. The denominator is always positive,
        so the sign is that of $x(x-2)$: $+$, $-$, $-$, $+$ across
        $0,1,2$ (with $x=1$ excluded). Local maximum $(0,0)$, local minimum
        $(2,4)$; the function decreases on $(0,1)$ and on $(1,2)$.
  \item $f'=\dfrac{1}{2\sqrt{x}}-1$, zero when $\sqrt{x}=\tfrac12$, i.e.
        $x=\tfrac14$: positive before, negative after. Global maximum
        $\bigl(\tfrac14,\tfrac14\bigr)$.
  \item $f''=6x-6$: negative for $x<1$ (concave), positive for $x>1$
        (convex). Inflection $(1,-2)$.
  \item $f''=12x^{2}-8$: zero at $x=\pm\sqrt{2/3}$; convex outside,
        concave between. Inflections at
        $\bigl(\pm\sqrt{2/3},-\tfrac{20}{9}\bigr)$.
  \item $f'=-2x\mathrm{e}^{-x^{2}}$ and
        $f''=(4x^{2}-2)\mathrm{e}^{-x^{2}}$: zero at
        $x=\pm1/\sqrt2$, convex outside, concave between. Inflections at
        $\bigl(\pm1/\sqrt2,\mathrm{e}^{-1/2}\bigr)$.
  \item $f''=12x^{2}\ge0$, vanishing only at $x=0$ but \emph{not changing
        sign} there: the curve is convex on both sides, so there is no
        inflection. $f''=0$ is necessary, not sufficient.
  \item $f'=\dfrac{1-\ln x}{x^{2}}$, $f''=\dfrac{2\ln x-3}{x^{3}}$: zero at
        $\ln x=\tfrac32$, i.e. $x=\mathrm{e}^{3/2}$; concave before, convex
        after. Inflection at
        $\bigl(\mathrm{e}^{3/2},\tfrac32\mathrm{e}^{-3/2}\bigr)$.
  \item $f'=(2x-x^{2})\mathrm{e}^{-x}$,
        $f''=(x^{2}-4x+2)\mathrm{e}^{-x}$: zero at $x=2\pm\sqrt2$, convex
        outside those, concave between.
  \item Pole $x=-2$; equal degrees, so horizontal asymptote $y=1$.
  \item Poles $x=\pm2$; numerator degree smaller, so $y=0$.
  \item Pole $x=1$; degrees differ by one, and
        $f=x+1+\dfrac{1}{x-1}$, so the oblique asymptote is $y=x+1$.
  \item No real pole ($x^{2}+1>0$); degrees differ by one, and
        $x^{3}=x(x^{2}+1)-x$, so $f=x-\dfrac{x}{x^{2}+1}$ and the oblique
        asymptote is $y=x$.
  \item Vertical asymptote $x=0$ (as $x\to0^{+}$, $f\to-\infty$) and
        horizontal $y=0$ as $x\to\infty$, since $\ln x\ll x$.
  \item As $x\to+\infty$, $\mathrm{e}^{-x}\to0$, so $y=x$ is an oblique
        asymptote. As $x\to-\infty$ the exponential dominates and
        $f\to+\infty$ with no asymptote.
  \item Domain $\R$; no symmetry; $f=x^{2}(x-3)$ so roots $0$ (double,
        touching) and $3$; $f(0)=0$; ends $-\infty$ and $+\infty$. From
        problem 9's table: increasing on $(-\infty,0)$, decreasing on
        $(0,2)$, increasing on $(2,\infty)$, local maximum $(0,0)$, local
        minimum $(2,-4)$; concave then convex with inflection $(1,-2)$. See
        the figure.
  \item Domain $\R$; \emph{even}; $f=x^{2}(x^{2}-4)$ so roots $0$ (double)
        and $\pm2$; both ends $+\infty$; local maximum $(0,0)$, local
        minima $(\pm\sqrt2,-4)$; inflections
        $\bigl(\pm\sqrt{2/3},-\tfrac{20}{9}\bigr)$. A symmetric W.
  \item Domain $\R$; \emph{odd}; single root $0$; horizontal asymptote
        $y=0$ at both ends; local minimum $(-1,-\tfrac12)$, local maximum
        $(1,\tfrac12)$;
        $f''=\dfrac{2x(x^{2}-3)}{(x^{2}+1)^{3}}$ gives inflections at
        $(0,0)$ and $\bigl(\pm\sqrt3,\pm\tfrac{\sqrt3}{4}\bigr)$. See the
        figure.
  \item The worked example of the Core Rules: domain $\R\setminus\{1\}$,
        double root at $0$, pole $x=1$, oblique asymptote $y=x+1$, local
        maximum $(0,0)$, local minimum $(2,4)$, no inflection, concave left
        of the pole and convex right of it.
  \item Domain $\R$; \emph{even}; no roots ($f>0$ always); horizontal
        asymptote $y=0$ at both ends; global maximum $(0,1)$; inflections
        at $\bigl(\pm1/\sqrt2,\mathrm{e}^{-1/2}\bigr)$; convex outside
        them, concave between. The bell curve.
  \item Domain $\R$; no symmetry; single root $0$; as $x\to+\infty$,
        $f\to0^{+}$ (asymptote $y=0$), as $x\to-\infty$, $f\to-\infty$;
        global maximum $(1,1/\mathrm{e})$; inflection
        $(2,2/\mathrm{e}^{2})$, concave before, convex after. See the
        figure.
  \item Domain $(0,\infty)$; root $x=1$; vertical asymptote $x=0$ with
        $f\to-\infty$, horizontal asymptote $y=0$ as $x\to\infty$; global
        maximum $(\mathrm{e},1/\mathrm{e})$; inflection at
        $x=\mathrm{e}^{3/2}$. See the figure.
  \item Domain $\R\setminus\{-2\}$; root $x=1$; $f(0)=-\tfrac12$; pole
        $x=-2$ ($-\infty$ from the left, $+\infty$ from the right);
        horizontal asymptote $y=1$; $f'=\dfrac{3}{(x+2)^{2}}>0$ so strictly
        increasing on each branch with \emph{no} extrema;
        $f''=\dfrac{-6}{(x+2)^{3}}$, convex left of the pole, concave right
        of it.
  \item Domain $\R$; no symmetry; double root $0$; as $x\to+\infty$,
        $f\to0^{+}$; as $x\to-\infty$, $f\to+\infty$; local minimum
        $(0,0)$, local maximum $(2,4/\mathrm{e}^{2})$; inflections at
        $x=2\pm\sqrt2$.
  \item Domain $\R\setminus\{\pm2\}$; \emph{even}; no roots; poles
        $x=\pm2$; horizontal asymptote $y=0$;
        $f'=\dfrac{-2x}{(x^{2}-4)^{2}}$, so increasing for $x<0$ and
        decreasing for $x>0$, giving a local maximum
        $\bigl(0,-\tfrac14\bigr)$ --- a maximum with a negative value,
        since the middle branch lies entirely below the axis.
  \item Domain $[0,\infty)$; roots $0$ and $1$; $f\to-\infty$ as
        $x\to\infty$; global maximum $\bigl(\tfrac14,\tfrac14\bigr)$;
        $f''=-\tfrac14x^{-3/2}<0$, so concave throughout with no
        inflection; the tangent at $0$ is vertical.
  \item Domain $\R$; $f=(x-1)^{2}(x+2)$, so a double root at $1$ (touching)
        and a simple root at $-2$; $f(0)=2$; ends $\mp\infty$; local
        maximum $(-1,4)$, local minimum $(1,0)$; inflection $(0,2)$.
  \item It rises, but ever more slowly: increasing and concave, bending
        away from its own tangent downwards --- the shape of a production
        function with diminishing returns.
  \item $f'<0$ between the roots and $>0$ outside them, so $f$ increases on
        $(-\infty,-1)$, decreases on $(-1,3)$, increases on $(3,\infty)$: a
        local maximum at $x=-1$ and a local minimum at $x=3$.
  \item $f'(x)=kx$ with $k>0$ gives $f(x)=\tfrac{k}{2}x^{2}+c$: an upward
        parabola with vertex at $x=0$, convex everywhere.
  \item $f'(1)=0$ with $f''(1)<0$: a local \textbf{maximum} at $(1,2)$ by
        the second-derivative test. The tangent is horizontal, $y=2$.
  \item One valid answer: $f(x)=1-\dfrac{4}{(x-2)^{2}}$ has the pole and the
        asymptote $y=1$ but a maximum, not a minimum, so adjust ---
        $f(x)=\dfrac{x^{2}}{(x-2)^{2}}$ works: $f(0)=0$ is a local minimum,
        the pole is at $x=2$, and $f\to1$ at both ends. Any sketch with a
        left branch dipping to $(0,0)$ and rising to $+\infty$ at $x=2^{-}$,
        and a right branch descending convexly from $+\infty$ towards
        $y=1$, is correct.
  \item Because a local extremum is only a statement about a
        \emph{neighbourhood}. The pole cuts the domain into two intervals
        with no path between them, so the values on one side need not
        compare with those on the other. In problem 24 the maximum $(0,0)$
        lies on the left branch, the minimum $(2,4)$ on the right, and
        between them the function runs off to $\mp\infty$.
  \item $C'=3q^{2}-12q+15=3(q^{2}-4q+5)=3\bigl((q-2)^{2}+1\bigr)\ge3>0$, so
        $C$ increases for every $q$. $C''=6q-12$ vanishes at $q=2$ and
        changes sign there: inflection at $(2,14)$. Since $MC=C'$ and
        $MC'=C''$, marginal cost has its minimum exactly at $q=2$, where
        $MC=3$.
  \item $AC=q^{2}-6q+15=(q-3)^{2}+6$, minimal at $q=3$ with $AC=6$. There
        $MC=3\cdot9-12\cdot3+15=6=AC$, as the identity
        $AC'=(MC-AC)/q$ requires.
  \item $\pi=15q-(q^{3}-6q^{2}+15q)=-q^{3}+6q^{2}$. Then
        $\pi'=-3q^{2}+12q=-3q(q-4)$, zero at $q=4$ (and $q=0$), and
        $\pi''=-6q+12=-12<0$ at $q=4$: a maximum, with
        $\pi(4)=-64+96=32$.
  \item $R=200q-4q^{2}=-4(q-25)^{2}+2500$ on the economic domain
        $[0,50]$ (where the price $200-4q$ stays non-negative): a downward
        parabola, increasing on $[0,25)$, decreasing on $(25,50]$, concave
        throughout, with maximum revenue $2500$ at $q=25$ and zeros at both
        ends.
  \item $AC=0.5q+2+\dfrac{8}{q}$, so $AC'=0.5-\dfrac{8}{q^{2}}=0$ gives
        $q^{2}=16$, $q=4$, with $AC(4)=2+2+2=6$. And $MC=q+2=6$ at $q=4$:
        equal, as required.
  \item As $t\to\infty$, $y\to100$; as $t\to-\infty$, $y\to0$: two
        horizontal asymptotes. $y'=\dfrac{50\mathrm{e}^{-0.5t}}
        {(1+\mathrm{e}^{-0.5t})^{2}}>0$ always, so $y$ increases
        throughout. The inflection is where the numerator's growth turns,
        at $\mathrm{e}^{-0.5t}=1$, i.e. $t=0$, with $y=50$ --- exactly half
        the ceiling, which is the signature of the logistic curve.
  \item $q>0$, so the sign of $AC'$ is the sign of $MC-AC$. If
        $MC<AC$ then $AC'<0$ and average cost falls: producing one more
        unit costs less than the current average, which pulls the average
        down. Once $MC$ rises above $AC$ the average is pulled up, and the
        crossing point is therefore the minimum.
\end{solutions}

\begin{gallery}
\begin{tikzpicture}
  \begin{axis}[fnaxis, xmin=-1.6, xmax=3.9, ymin=-5.4, ymax=4.4,
      xtick={-1,1,2,3}, ytick={-4,-2,2}, title={$x^{3}-3x^{2}$}]
    \addplot[domain=-1.2:3.6, samples=201]{x^3-3*x^2};
    \addplot[only marks, mark=*, mark size=1.2pt, black, forget plot]
      coordinates {(0,0) (2,-4) (1,-2)};
  \end{axis}
\end{tikzpicture}\hfill
\begin{tikzpicture}
  \begin{axis}[fnaxis, xmin=-5.4, xmax=5.4, ymin=-0.75, ymax=0.75,
      xtick={-4,-2,2,4}, ytick={-0.5,0.5}, title={$x/(x^{2}+1)$}]
    \addplot[domain=-5.2:5.2, samples=201]{x/(x^2+1)};
    \addplot[only marks, mark=*, mark size=1.2pt, black, forget plot]
      coordinates {(1,0.5) (-1,-0.5) (0,0)};
  \end{axis}
\end{tikzpicture}
\end{gallery}

\begin{gallery}
\begin{tikzpicture}
  \begin{axis}[fnaxis, xmin=-1.4, xmax=5.4, ymin=-1.4, ymax=0.75,
      xtick={-1,1,2,4}, ytick={-1,0.5}, title={$x\mathrm{e}^{-x}$}]
    \addplot[domain=-1.2:5.2, samples=201]{x*exp(-x)};
    \addplot[only marks, mark=*, mark size=1.2pt, black, forget plot]
      coordinates {(1,0.3679) (2,0.2707)};
  \end{axis}
\end{tikzpicture}\hfill
\begin{tikzpicture}
  \begin{axis}[fnaxis, xmin=-0.5, xmax=9.4, ymin=-1.9, ymax=0.75,
      xtick={2,4,6,8}, ytick={-1,0.5}, title={$\ln x/x$}]
    \addplot[domain=0.14:9.2, samples=241]{ln(x)/x};
    \addplot[only marks, mark=*, mark size=1.2pt, black, forget plot]
      coordinates {(2.7183,0.3679) (4.4817,0.3347)};
  \end{axis}
\end{tikzpicture}
\gcap{Four of the completed discussions: problems 21, 23, 26 and 27, with
their extrema and inflection points marked.}
\end{gallery}

\section{Additional Problems}

\begin{problems}
  \item Monotonicity and extrema of $f(x)=2x^{3}-9x^{2}+12x$.
  \item Monotonicity and extrema of $f(x)=\dfrac{x^{2}}{x^{2}+3}$.
  \item Monotonicity and extrema of $f(x)=x\ln x$ on $x>0$.
  \item Convexity and inflections of $f(x)=x^{3}-6x^{2}+9x$.
  \item Convexity and inflections of $f(x)=\dfrac{1}{1+x^{2}}$.
  \item Convexity of $f(x)=\mathrm{e}^{x}+\mathrm{e}^{-x}$, and its
        minimum.
  \item All asymptotes of $f(x)=\dfrac{2x^{2}-3}{x+1}$.
  \item All asymptotes of $f(x)=\dfrac{x^{2}-1}{x^{2}+1}$.
  \item All asymptotes of $f(x)=x\mathrm{e}^{1/x}$.
  \item Complete discussion of $f(x)=x^{3}-6x^{2}+9x$.
  \item Complete discussion of $f(x)=\dfrac{x^{2}-4}{x^{2}}$.
  \item Complete discussion of $f(x)=(x-2)\mathrm{e}^{-x}$.
  \item Complete discussion of $f(x)=x^{2}\ln x$ on $x>0$.
  \item Complete discussion of $f(x)=\dfrac{4x}{x^{2}+4}$.
  \item Complete discussion of $f(x)=\sqrt{4-x^{2}}$.
  \item A cubic has a local maximum at $(0,4)$ and a local minimum at
        $(2,0)$. Determine it.
  \item A function has $f'(x)>0$ everywhere and $f''(x)<0$ everywhere, with
        horizontal asymptote $y=5$. Sketch it and give an example formula.
  \item Where does $f(x)=x^{4}-4x^{3}$ have inflections, and why does the
        second-derivative test fail at one of its critical points?
  \item Show that $f(x)=x+1/x$ has a local minimum at $x=1$ and a local
        maximum at $x=-1$, with the maximum value below the minimum value,
        and explain.
  \item For which $a$ does $f(x)=x^{3}+ax$ have two local extrema?
  \item For $C(q)=q^{3}-9q^{2}+30q$: find the inflection, the minimum of
        average cost, and verify $MC=AC$ there.
  \item For $C(q)=2q^{2}+50$ and price $p=40$: find the profit-maximising
        output, the profit, and the break-even outputs.
  \item A monopolist faces $p(q)=100-2q$ with $C(q)=q^{2}+20q$. Discuss the
        profit curve and find the optimum.
  \item Show that if $AC$ has a minimum at $q^{*}$ then the $AC$ curve is
        tangent to no ray from the origin except the one through
        $(q^{*},AC(q^{*}))$, and relate this to $MC=AC$.
  \item The production function $Q(L)=12L^{2}-L^{3}$ describes output from
        labour $L\in[0,12]$. Find where average product is maximal and show
        that marginal product equals average product there.
\end{problems}

\section{Additional Solutions}

\begin{solutions}
  \item $f'=6x^{2}-18x+12=6(x-1)(x-2)$: $+,-,+$. Local maximum $(1,5)$,
        local minimum $(2,4)$.
  \item $f'=\dfrac{6x}{(x^{2}+3)^{2}}$: negative for $x<0$, positive for
        $x>0$. Global minimum $(0,0)$; $f\to1$ at both ends, so no
        maximum.
  \item $f'=\ln x+1$, zero at $x=\mathrm{e}^{-1}$: negative before,
        positive after. Global minimum
        $\bigl(\mathrm{e}^{-1},-\mathrm{e}^{-1}\bigr)$.
  \item $f''=6x-12$: concave for $x<2$, convex for $x>2$; inflection
        $(2,2)$.
  \item $f'=\dfrac{-2x}{(1+x^{2})^{2}}$,
        $f''=\dfrac{6x^{2}-2}{(1+x^{2})^{3}}$: zero at
        $x=\pm1/\sqrt3$, convex outside, concave between. Inflections at
        $\bigl(\pm1/\sqrt3,\tfrac34\bigr)$.
  \item $f''=\mathrm{e}^{x}+\mathrm{e}^{-x}>0$ everywhere: convex, no
        inflection. $f'=\mathrm{e}^{x}-\mathrm{e}^{-x}=0$ at $x=0$, giving
        the global minimum $(0,2)$.
  \item Pole $x=-1$; degrees differ by one, and
        $2x^{2}-3=(x+1)(2x-2)-1$, so $f=2x-2-\dfrac{1}{x+1}$: oblique
        asymptote $y=2x-2$.
  \item No poles; equal degrees, so horizontal asymptote $y=1$ at both
        ends.
  \item As $x\to\pm\infty$, $\mathrm{e}^{1/x}\to1$ and more precisely
        $x\mathrm{e}^{1/x}=x+1+\tfrac{1}{2x}+\cdots$, so $y=x+1$ is an
        oblique asymptote. As $x\to0^{+}$, $f\to+\infty$ (vertical
        asymptote $x=0$); as $x\to0^{-}$, $f\to0$.
  \item $f=x(x-3)^{2}$: roots $0$ and $3$ (double); ends $\mp\infty$;
        $f'=3x^{2}-12x+9=3(x-1)(x-3)$ gives local maximum $(1,4)$ and local
        minimum $(3,0)$; $f''=6x-12$ gives inflection $(2,2)$; no
        symmetry, no asymptotes.
  \item $f=1-\dfrac{4}{x^{2}}$: domain $\R\setminus\{0\}$, \emph{even},
        roots $\pm2$, pole $x=0$ with $f\to-\infty$ from both sides,
        horizontal asymptote $y=1$; $f'=\dfrac{8}{x^{3}}$, so decreasing
        for $x<0$ and increasing for $x>0$ with no extrema;
        $f''=-\dfrac{24}{x^{4}}<0$: concave on both branches.
  \item Domain $\R$; root $x=2$; $f\to0^{+}$ as $x\to+\infty$
        (asymptote $y=0$) and $f\to-\infty$ as $x\to-\infty$;
        $f'=(3-x)\mathrm{e}^{-x}$ gives a maximum
        $(3,\mathrm{e}^{-3})$; $f''=(x-4)\mathrm{e}^{-x}$ gives an
        inflection $(4,2\mathrm{e}^{-4})$.
  \item Domain $(0,\infty)$; root $x=1$; $f\to0$ as $x\to0^{+}$ and
        $f\to\infty$ as $x\to\infty$; $f'=2x\ln x+x=x(2\ln x+1)$, zero at
        $x=\mathrm{e}^{-1/2}$, a minimum with value
        $-\tfrac{1}{2\mathrm{e}}$; $f''=2\ln x+3$, zero at
        $x=\mathrm{e}^{-3/2}$, an inflection.
  \item Domain $\R$; \emph{odd}; root $0$; horizontal asymptote $y=0$;
        $f'=\dfrac{4(4-x^{2})}{(x^{2}+4)^{2}}$ gives local minimum
        $(-2,-1)$ and local maximum $(2,1)$; inflections at $0$ and
        $\pm2\sqrt3$.
  \item Domain $[-2,2]$; \emph{even}; roots $\pm2$; range $[0,2]$; it is
        the upper semicircle of radius $2$, so concave throughout, maximum
        $(0,2)$, and the tangents at $x=\pm2$ are vertical. No asymptotes.
  \item Write $f=ax^{3}+bx^{2}+cx+d$. From the maximum at $(0,4)$:
        $d=4$ and $f'(0)=c=0$. From $f'(2)=12a+4b=0$ and $f(2)=8a+4b+4=0$:
        $b=-3a$ and $8a-12a+4=0$, so $a=1$, $b=-3$. Hence
        $f(x)=x^{3}-3x^{2}+4$.
  \item Increasing and concave with a ceiling: it rises steeply at first
        and flattens onto $y=5$ without reaching it. Example:
        $f(x)=5-\mathrm{e}^{-x}$.
  \item $f'=4x^{3}-12x^{2}=4x^{2}(x-3)$, critical at $x=0$ and $x=3$;
        $f''=12x^{2}-24x=12x(x-2)$, zero at $x=0$ and $x=2$ and changing
        sign at both, so inflections at $(0,0)$ and $(2,-16)$. At $x=0$ the
        second-derivative test fails because $f''(0)=0$; the sign table for
        $f'$ shows $4x^{2}(x-3)<0$ on both sides of $0$, so $x=0$ is no
        extremum --- it is the inflection.
  \item $f'=1-\dfrac{1}{x^{2}}$, zero at $x=\pm1$;
        $f''=\dfrac{2}{x^{3}}$, positive at $x=1$ (minimum, value $2$) and
        negative at $x=-1$ (maximum, value $-2$). The maximum value $-2$
        lies below the minimum value $2$ because the pole at $x=0$
        separates the two branches, so neither extremum is global.
  \item $f'=3x^{2}+a$ has two real roots exactly when $a<0$; then
        $x=\pm\sqrt{-a/3}$ and the sign of $f'$ changes at both, giving a
        maximum and a minimum. For $a\ge0$ the function is monotone
        increasing.
  \item $C''=6q-18=0$ at $q=3$: inflection, and $MC=3q^{2}-18q+30$ is
        minimal there with $MC=3$. $AC=q^{2}-9q+30=(q-4.5)^{2}+9.75$,
        minimal at $q=4.5$ with $AC=9.75$; and
        $MC(4.5)=60.75-81+30=9.75=AC$.
  \item $\pi=40q-2q^{2}-50$, so $\pi'=40-4q=0$ at $q=10$ with
        $\pi''=-4<0$: maximum profit $\pi(10)=400-200-50=150$. Break-even
        needs $\pi=0$: $2q^{2}-40q+50=0$, i.e. $q^{2}-20q+25=0$, so
        $q=10\pm\sqrt{75}$, about $1.34$ and $18.66$.
  \item $\pi=(100-2q)q-(q^{2}+20q)=80q-3q^{2}$, a downward parabola with
        zeros at $q=0$ and $q=80/3$, concave throughout. $\pi'=80-6q=0$ at
        $q=\tfrac{40}{3}\approx13.33$, with
        $\pi=\tfrac{1600}{3}\approx533.3$ and price
        $p=100-\tfrac{80}{3}\approx73.3$.
  \item The slope of the ray from the origin through $(q,C(q))$ is exactly
        $C(q)/q=AC(q)$. Minimising $AC$ therefore means finding the
        \emph{flattest} such ray, and the flattest ray that still meets the
        cost curve must touch it --- be tangent to it. At a tangency the
        ray's slope equals the curve's slope, i.e. $AC=MC$. That is the
        geometric content of $AC'=(MC-AC)/q$.
  \item $AP=Q/L=12L-L^{2}$, maximal where $12-2L=0$, i.e. $L=6$, with
        $AP=36$. And $MP=Q'=24L-3L^{2}$, so $MP(6)=144-108=36=AP$. The
        same tangency argument as for average cost applies, with the ray
        from the origin now touching the production curve.
\end{solutions}

\chapter{Functions of Several Variables: Surfaces and Level Curves}

A function of one variable is understood by looking at its graph. For a
function of two variables there are two pictures, and both are needed. The
\textbf{graph} is a surface in space: mountains, valleys, passes. The
\textbf{contour map} is what a cartographer draws instead --- the level
curves of that landscape, projected flat onto the plane. Economics almost
never draws the surface. It draws the contour map and gives the curves
names: isoquants, indifference curves, iso-profit lines.

The reason to spend a chapter on pictures before computing anything is that
most of multivariable calculus is easier to see than to calculate. Whether a
critical point is a maximum, a minimum or a saddle is visible at a glance in
a contour map. So is whether a function is convex, whether its upper contour
sets are convex, and where it is steep. And the geometry is not decoration:
the Lagrange method is the statement that two families of level curves touch
without crossing, and Kuhn--Tucker is a statement about which boundary of a
shaded region binds.

This chapter therefore does two things. It sets up domains, graphs, level
curves and the sub- and superlevel sets, with the vocabulary in both the
German and the English forms; and it drills reading properties off a picture
in both directions --- picture to property, and formula to picture.

\section{Core Rules}

\begin{rulebox}{Functions of two variables: domain, graph, range}
A function $f\colon D\to\R$ with $D\subseteq\R^{2}$ assigns a number
$f(x,y)$ to each point of a \textbf{region} $D$ in the plane.
\begin{itemize}[leftmargin=1.6em,itemsep=0.3ex]
  \item The \textbf{domain} $D$ is now a subset of the plane, not of the
        line, and the same three obstructions decide it: no zero
        denominator, no negative radicand under an even root, no
        non-positive argument in a logarithm. So
        $f(x,y)=\sqrt{9-x^{2}-y^{2}}$ has $D=\{(x,y):x^{2}+y^{2}\le9\}$, a
        closed disc, and $f(x,y)=\ln(x+y)$ has the open half-plane
        $x+y>0$.
  \item The \textbf{graph} is the surface
        $\{(x,y,f(x,y)) : (x,y)\in D\}$ in $\R^{3}$: above each point of
        the domain, a single height.
  \item The \textbf{range} is the set of heights attained, exactly as
        before.
\end{itemize}
Economic domains are usually the non-negative quadrant $K,L\ge0$, and often
its interior.
\end{rulebox}

\begin{rulebox}{Level curves and the contour map}
For $c\in\R$ the \textbf{level set} \de{(Niveaumenge, Niveaulinie)} is
\[
  N_{c}=\{(x,y)\in D : f(x,y)=c\} ,
\]
the set of points at height $c$. Drawing $N_{c}$ for several equally spaced
values of $c$ gives the \textbf{contour map}. To compute one, set
$f(x,y)=c$ and solve for $y$ --- the level set is a curve in the plane, and
often a familiar one.

Reading a contour map:
\begin{itemize}[leftmargin=1.6em,itemsep=0.3ex]
  \item \textbf{Contours cannot cross}, since a point has one height ---
        with one exception: at a \textbf{saddle} the level curve through
        the saddle point crosses itself.
  \item \textbf{Spacing is steepness.} Contours crowded together mean the
        surface rises fast; widely spaced contours mean gentle ground.
        Steepest ascent runs perpendicular to the contours.
  \item \textbf{Closed nested loops} shrinking to a point mark a
        \textbf{peak} or a \textbf{basin}; which one is decided by whether
        the labels increase or decrease inwards.
  \item A contour that \textbf{pinches and crosses} marks a saddle: two
        directions of ascent and two of descent alternate around it.
\end{itemize}
\end{rulebox}

\begin{rulebox}{Sub- and superlevel sets}
Two \emph{solid} sets accompany each level curve:
\[
  W(c)=\{(x,y): f(x,y)\le c\},
  \qquad
  B(c)=\{(x,y): f(x,y)\ge c\} .
\]
$W(c)$ is the \textbf{sublevel set} \de{(Subniveaumenge, untere
Konturmenge)}, $B(c)$ the \textbf{superlevel set} \de{(Superniveaumenge,
obere Konturmenge)}. Written through a point $P$ rather than a level,
$W(P)=W(f(P))$ and $B(P)=B(f(P))$: everything at most as high as $P$, and
everything at least as high. The same sets appear in the optimization
volume as $L_{c}$ and $U_{c}$ --- lower and upper contour set --- and in
consumer theory $B(P)$ is the \emph{better set} of the bundle $P$.

The link to curvature, drilled later and worth seeing now:
\begin{itemize}[leftmargin=1.6em,itemsep=0.3ex]
  \item $f$ convex $\Longrightarrow$ every $W(c)$ is convex;
  \item $f$ concave $\Longrightarrow$ every $B(c)$ is convex;
  \item every $B(c)$ convex is the \emph{definition} of
        \textbf{quasiconcave}, which is weaker: $f(x,y)=xy$ on the positive
        quadrant is quasiconcave but not concave.
\end{itemize}
\end{rulebox}

\begin{gallery}
\begin{tikzpicture}
  \begin{axis}[geoaxis, xmin=-3.2, xmax=3.2, ymin=-3.2, ymax=3.2,
      xtick={-2,2}, ytick={-2,2}, title={$W(4)$ for $f=x^{2}+y^{2}$}]
    \addplot[draw=none, name path=up, domain=-2:2, samples=201]{sqrt(4-x^2)};
    \addplot[draw=none, name path=lo, domain=-2:2, samples=201]{-sqrt(4-x^2)};
    \region{up}{lo}{}
    \addplot[domain=0:360, samples=181]({2*cos(x)},{2*sin(x)});
  \end{axis}
\end{tikzpicture}\hfill
\begin{tikzpicture}
  \begin{axis}[geoaxis, xmin=-0.4, xmax=8.4, ymin=-0.4, ymax=8.4,
      xtick={2,4,6,8}, ytick={2,4,6,8}, title={$B(4)$ for $f=xy$, $x,y>0$}]
    \addplot[draw=none, name path=hyp, domain=0.5:8, samples=201]{4/x};
    \addplot[draw=none, name path=top, domain=0.5:8, samples=2]{8.4};
    \region{hyp}{top}{}
    \addplot[domain=0.5:8, samples=201]{4/x};
  \end{axis}
\end{tikzpicture}
\gcap{A sublevel set of a convex function and a superlevel set of a
quasiconcave one. Both are convex sets; the second is the \emph{better set}
of a Cobb--Douglas bundle.}
\end{gallery}

\begin{rulebox}{The standard shapes}
Six surfaces cover most of what appears in exercises, and each has a
signature contour map.
\begin{itemize}[leftmargin=1.6em,itemsep=0.3ex]
  \item \textbf{Plane} $f=ax+by+c$: contours are \emph{parallel straight
        lines}, equally spaced. A budget line is one of them.
  \item \textbf{Paraboloid} $f=x^{2}+y^{2}$: a bowl; contours are
        concentric circles of radius $\sqrt{c}$, so they get
        \emph{closer together} outwards --- the bowl steepens.
  \item \textbf{Cone} $f=\sqrt{x^{2}+y^{2}}$: contours are concentric
        circles of radius $c$, \emph{equally spaced} --- constant slope.
  \item \textbf{Saddle} $f=x^{2}-y^{2}$: contours are hyperbolas; the level
        $c=0$ degenerates into the crossing pair $y=\pm x$.
  \item \textbf{Rotated saddle} $f=xy$: hyperbolas again, but with the
        \emph{axes} as asymptotes. Restricted to the positive quadrant this
        is the Cobb--Douglas picture.
  \item \textbf{Parabolic cylinder} $f=x^{2}$: the surface is a trough,
        constant in $y$; contours are \emph{pairs of vertical lines}
        $x=\pm\sqrt{c}$.
\end{itemize}
Circles alone do not identify the surface --- bowl, cone and dome all have
circular contours. The \emph{spacing} and the direction in which the labels
grow settle it.
\end{rulebox}

\begin{rulebox}{Reading curvature off the picture}
\begin{itemize}[leftmargin=1.6em,itemsep=0.3ex]
  \item \textbf{Convex} (bowl-shaped): every chord drawn between two points
        of the surface lies \emph{above} it; no part of the surface bulges
        over a chord. All sublevel sets are convex.
  \item \textbf{Concave} (dome-shaped): every chord lies \emph{below} the
        surface. All superlevel sets are convex.
  \item \textbf{Neither}: a saddle. Some chords lie above, some below.
  \item \textbf{Quasiconcave without being concave}: the superlevel sets
        are still convex --- nested ``better sets'' bulging towards the
        origin --- while the surface itself is not a dome. The Cobb--Douglas
        hill is the standard example.
  \item A \textbf{linear} function is both convex and concave, and both
        quasiconvex and quasiconcave.
\end{itemize}
\end{rulebox}

\begin{rulebox}{Economic readings}
\begin{itemize}[leftmargin=1.6em,itemsep=0.3ex]
  \item An \textbf{isoquant} is a level curve of a production function:
        all input combinations yielding the same output. An
        \textbf{indifference curve} is a level curve of utility.
  \item The slope of a level curve is the \textbf{marginal rate of
        substitution}, computed implicitly as $-f_{x}/f_{y}$
        (Chapter~5). Convex-to-the-origin level curves mean a diminishing
        rate of substitution.
  \item A \textbf{budget set} $\{p_{x}x+p_{y}y\le M,\ x,y\ge0\}$ is the
        sublevel set of a linear function intersected with the quadrant ---
        a triangle, and convex.
  \item The consumer's problem is to reach the highest indifference curve
        that still touches the budget triangle: the contour map plus the
        shaded set is the whole picture.
\end{itemize}
\end{rulebox}

\section{Worked Example}

\begin{worked}{Read the topography of $f(x,y)=x^{2}+4y^{2}$ and of $g(x,y)=x^{2}-y^{2}$: level curves, sub- and superlevel sets, and the surfaces}
\wstep{Step 1. A level curve is what you get by \emph{fixing the height}.}
Set $f(x,y)=c$ and solve for the curve in the plane:
\[
  x^{2}+4y^{2}=c .
\]
For $c>0$ divide by $c$ to reach the standard form of an ellipse:
\[
  \frac{x^{2}}{c}+\frac{y^{2}}{c/4}=1,
  \qquad\text{semi-axes } \sqrt c \text{ and } \tfrac{\sqrt c}{2} .
\]
So the contour map is a family of concentric ellipses, twice as wide as they
are tall. For $c=0$ the level ``curve'' collapses to the single point
$(0,0)$; for $c<0$ it is empty, since $f\ge0$ always.
\wtrap{Check which values of $c$ occur before drawing. Level sets can be
empty, a single point, or several disconnected pieces --- $x^{2}+4y^{2}=-1$
has no solutions, and a level set of $x^{2}-y^{2}$ has two branches.}

\wstep{Step 2. Spacing carries the steepness.}
The contours for $c=1,2,3,4$ have semi-axes $1,\sqrt2,\sqrt3,2$ in $x$. The
gaps shrink as you move out --- $1,\ 0.41,\ 0.32,\ 0.27$ --- so the surface
gets \emph{steeper} away from the origin. Contours that crowd together mean
a steep slope; contours far apart mean a plateau. That is the one skill a
contour map is for.

\wstep{Step 3. Sub- and superlevel sets.}
\[
  \{f\le c\}=\text{the filled ellipse: convex},
  \qquad
  \{f\ge c\}=\text{everything outside it: not convex.}
\]
$f$ is a convex function, so its \textbf{sub}level sets are convex. For a
concave function such as a utility $U$, it is the \textbf{super}level sets
$\{U\ge\bar u\}$ --- the ``at least as good as'' sets --- that are convex,
and that is precisely the convexity of preferences.
\wtrap{Convex function and convex set are different words. Match them
carefully: convex function $\Rightarrow$ sublevel sets convex; concave
function $\Rightarrow$ superlevel sets convex. The direction is chosen by
the inequality, not by the word ``convex''.}

\wstep{Step 4. The surface itself.}
$z=x^{2}+4y^{2}$ is an \textbf{elliptic paraboloid}: a bowl, with its
minimum at the origin, rising in every direction, steeper along $y$ than
along $x$ (the coefficient $4$ against $1$).

\wstep{Step 5. Now $g(x,y)=x^{2}-y^{2}$, where the sign flips.}
\[
  g=c:\quad x^{2}-y^{2}=c ,
\]
a family of \textbf{hyperbolas} --- opening left--right for $c>0$, up--down
for $c<0$. The level set $c=0$ degenerates into the \emph{pair of lines}
$y=\pm x$, which is where the two families change over.

\wstep{Step 6. Two pictures, side by side.}
\begin{center}
\begin{tikzpicture}
\begin{axis}[contouraxis, width=5.2cm, height=5.2cm, xmin=-2.6, xmax=2.6,
             ymin=-2.6, ymax=2.6, title={$x^{2}+4y^{2}=c$}]
  \foreach \c in {0.5,1,2,3,4} {
    \addplot[domain=0:360, samples=161, variable=\t]
      ({sqrt(\c)*cos(\t)}, {sqrt(\c)/2*sin(\t)});
  }
\end{axis}
\end{tikzpicture}
\hfill
\begin{tikzpicture}
\begin{axis}[contouraxis, width=5.2cm, height=5.2cm, xmin=-2.6, xmax=2.6,
             ymin=-2.6, ymax=2.6, title={$x^{2}-y^{2}=c$}]
  \foreach \c in {0.5,2} {
    \addplot[domain=-1.6:1.6, samples=121] ({sqrt(\c)*cosh(x)}, {sqrt(\c)*sinh(x)});
    \addplot[domain=-1.6:1.6, samples=121] ({-sqrt(\c)*cosh(x)}, {sqrt(\c)*sinh(x)});
    \addplot[domain=-1.6:1.6, samples=121] ({sqrt(\c)*sinh(x)}, {sqrt(\c)*cosh(x)});
    \addplot[domain=-1.6:1.6, samples=121] ({sqrt(\c)*sinh(x)}, {-sqrt(\c)*cosh(x)});
  }
  \addplot[gray!60, thin, samples=2, domain=-2.6:2.6, forget plot] {x};
  \addplot[gray!60, thin, samples=2, domain=-2.6:2.6, forget plot] {-x};
\end{axis}
\end{tikzpicture}
\end{center}
\begin{center}
\begin{tikzpicture}
\begin{axis}[surfaxis, title={$z=x^{2}+4y^{2}$}]
  \addplot3[surf, domain=-2:2, y domain=-1:1] {x^2+4*y^2};
\end{axis}
\end{tikzpicture}
\hfill
\begin{tikzpicture}
\begin{axis}[surfaxis, title={$z=x^{2}-y^{2}$}]
  \addplot3[surf, domain=-2:2, y domain=-2:2] {x^2-y^2};
\end{axis}
\end{tikzpicture}
\end{center}

\wstep{Step 7. Read the origin off each map.}
For $f$ the contours \emph{close around} the origin, so it is a minimum: go
any direction and you climb. For $g$ the contours \emph{cross} at the
origin, and the crossing is the signature of a \textbf{saddle}: along the
$x$-axis $g=x^{2}$ rises, along the $y$-axis $g=-y^{2}$ falls. The point is
a minimum in one direction and a maximum in another, so it is neither.
\wtrap{Closed loops shrinking to a point mean an extremum; crossing curves
mean a saddle. Checking only one direction --- for instance only along the
axes of $g$ where the answers disagree --- is exactly how a saddle gets
misreported as an optimum in Volume~7.}

\wstep{Step 8. The economic vocabulary.}
Level curves of a utility function are \textbf{indifference curves}; of a
production function, \textbf{isoquants}; of a cost function,
\textbf{isocost lines}. Their slope is the marginal rate of substitution,
computed in Chapter~5 by implicit differentiation, and the standard
``convex to the origin'' shape is exactly the convexity of the superlevel
sets in Step~3.
\end{worked}

\section{Practice Problems}

State the domain of each function as a region in the plane, in problems
1--8, and describe it in words.

\begin{problems}
  \item $f(x,y)=\sqrt{16-x^{2}-y^{2}}$
  \item $f(x,y)=\ln(x-y)$
  \item $f(x,y)=\dfrac{1}{x^{2}+y^{2}-1}$
  \item $f(x,y)=\dfrac{xy}{x-y}$
  \item $f(x,y)=\sqrt{x}+\sqrt{y}$
  \item $f(x,y)=\ln\bigl(xy\bigr)$
  \item $f(x,y)=\sqrt{x^{2}+y^{2}-4}$
  \item $f(K,L)=K^{0.4}L^{0.6}$, and state the economic domain too.
\end{problems}

In problems 9--20, give the level set $N_{c}$ for the stated level, solved
for $y$ where possible, and name the curve.

\begin{problems}[resume]
  \item $f(x,y)=x+2y$, \;$c=4$.
  \item $f(x,y)=x^{2}+y^{2}$, \;$c=9$.
  \item $f(x,y)=x^{2}+y^{2}$, \;$c=0$.
  \item $f(x,y)=x^{2}+4y^{2}$, \;$c=4$.
  \item $f(x,y)=xy$, \;$c=6$.
  \item $f(x,y)=x^{2}-y^{2}$, \;$c=0$.
  \item $f(x,y)=x^{2}-y^{2}$, \;$c=4$.
  \item $f(x,y)=y-x^{2}$, \;$c=2$.
  \item $f(x,y)=\sqrt{x^{2}+y^{2}}$, \;$c=3$.
  \item $f(x,y)=\mathrm{e}^{x+y}$, \;$c=1$.
  \item $f(K,L)=\sqrt{KL}$, \;$c=10$.
  \item $f(x,y)=2x+3y$, \;$c=12$; and say what this curve is called when
        $f$ is a cost.
\end{problems}

Problems 21--28 refer to the contour map below, which shows the level
curves of $f(x,y)=x^{3}-3x+y^{2}$ at the levels $c=-1,0,2,4$, together with
six marked points.

\begin{gallery}
\hspace*{\fill}%
\begin{tikzpicture}
  \begin{axis}[contouraxis, width=8.2cm, height=8.2cm,
      xmin=-2.6, xmax=2.6, ymin=-3.2, ymax=3.2,
      xtick={-2,-1,1,2}, ytick={-2,-1,1,2}]
    % c = 4
    \addplot[domain=-2.2:2.196, samples=201]{ sqrt(4-x^3+3*x)};
    \addplot[domain=-2.2:2.196, samples=201]{-sqrt(4-x^3+3*x)};
    % c = 2 (through the saddle at (-1,0))
    \addplot[domain=-2.2:2, samples=201]{ sqrt(2-x^3+3*x)};
    \addplot[domain=-2.2:2, samples=201]{-sqrt(2-x^3+3*x)};
    % c = 0, two pieces
    \addplot[domain=-2.2:-1.7321, samples=101]{ sqrt(-x^3+3*x)};
    \addplot[domain=-2.2:-1.7321, samples=101]{-sqrt(-x^3+3*x)};
    \addplot[domain=0:1.7321, samples=101]{ sqrt(-x^3+3*x)};
    \addplot[domain=0:1.7321, samples=101]{-sqrt(-x^3+3*x)};
    % c = -1, two pieces
    \addplot[domain=-2.2:-1.8794, samples=61]{ sqrt(-1-x^3+3*x)};
    \addplot[domain=-2.2:-1.8794, samples=61]{-sqrt(-1-x^3+3*x)};
    \addplot[domain=0.3473:1.5321, samples=101]{ sqrt(-1-x^3+3*x)};
    \addplot[domain=0.3473:1.5321, samples=101]{-sqrt(-1-x^3+3*x)};
    % marked points
    \addplot[only marks, mark=*, mark size=1.4pt, black, forget plot]
      coordinates {(0,0) (1,0) (-1,0) (2,1) (0,2) (-2,1)};
    \node[font=\tiny, anchor=north east] at (axis cs:-0.06,-0.12) {$P_{1}$};
    \node[font=\tiny, anchor=north west] at (axis cs:1.08,-0.12) {$P_{2}$};
    \node[font=\tiny, anchor=south east] at (axis cs:-1.12,0.12) {$P_{3}$};
    \node[font=\tiny, anchor=south west] at (axis cs:2.05,1.05) {$P_{4}$};
    \node[font=\tiny, anchor=south west] at (axis cs:0.05,2.05) {$P_{5}$};
    \node[font=\tiny, anchor=south east] at (axis cs:-2.05,1.05) {$P_{6}$};
  \end{axis}
\end{tikzpicture}%
\hspace*{\fill}
\gcap{Level curves of $f(x,y)=x^{3}-3x+y^{2}$ at $c=-1,0,2,4$.}
\end{gallery}

\begin{problems}[resume]
  \item Give $f(P_{i})$ for all six marked points
        $P_{1}=(0,0)$, $P_{2}=(1,0)$, $P_{3}=(-1,0)$, $P_{4}=(2,1)$,
        $P_{5}=(0,2)$, $P_{6}=(-2,1)$.
  \item Which of the six points lie on one and the same level curve?
  \item One of the drawn contours crosses itself. At which point, and what
        does that tell you about the surface there?
  \item One of the drawn contours has shrunk to a single point. Which
        level, which point, and what feature of the surface is it?
  \item From the picture alone, name the two critical points of $f$ and
        classify each as a local maximum, a local minimum, or a saddle.
  \item Confirm the previous answer by computation: solve
        $f_{x}=f_{y}=0$.
  \item Is the sublevel set $W(-1)$ connected? Answer from the picture and
        explain what the two pieces mean.
  \item Moving away from $P_{2}$, in which direction does $f$ increase
        fastest? Answer qualitatively from the contours, then check with
        $\grad f$ at a nearby point, say $(1.5,0.5)$.
\end{problems}

\begin{problems}[resume]
  \item Match each surface (a)--(f) below to its contour map (1)--(6).
  \item For each of the six functions in the previous problem, say whether
        it is convex, concave, both, or neither on all of $\R^{2}$.
\end{problems}

\begin{gallery}
\begin{tikzpicture}
  \begin{axis}[surfaxis, title={(a)}]
    \addplot3[surf, domain=-2:2, y domain=-2:2]{x+y};
  \end{axis}
\end{tikzpicture}\hfill
\begin{tikzpicture}
  \begin{axis}[surfaxis, title={(b)}, z post scale=1.2]
    \addplot3[surf, domain=-2:2, y domain=-2:2]{x^2+y^2};
  \end{axis}
\end{tikzpicture}\hfill
\begin{tikzpicture}
  \begin{axis}[surfaxis, title={(c)}, z post scale=1.2]
    \addplot3[surf, domain=-2:2, y domain=-2:2]{sqrt(x^2+y^2)};
  \end{axis}
\end{tikzpicture}
\end{gallery}

\begin{gallery}
\begin{tikzpicture}
  \begin{axis}[surfaxis, title={(d)}, z post scale=1.2]
    \addplot3[surf, domain=-2:2, y domain=-2:2]{x^2-y^2};
  \end{axis}
\end{tikzpicture}\hfill
\begin{tikzpicture}
  \begin{axis}[surfaxis, title={(e)}, z post scale=1.2, view={105}{32}]
    \addplot3[surf, domain=-2:2, y domain=-2:2]{x*y};
  \end{axis}
\end{tikzpicture}\hfill
\begin{tikzpicture}
  \begin{axis}[surfaxis, title={(f)}, z post scale=1.2]
    \addplot3[surf, domain=-2:2, y domain=-2:2]{x^2};
  \end{axis}
\end{tikzpicture}
\end{gallery}

\begin{gallery}
\begin{tikzpicture}
  \begin{axis}[contouraxis, xmin=-2.4, xmax=2.4, ymin=-2.4, ymax=2.4,
      xtick={-2,2}, ytick={-2,2}, title={(1)}]
    \foreach \c in {-2,-1,0,1,2} { \addplot[domain=-2.3:2.3]{\c-x}; }
  \end{axis}
\end{tikzpicture}\hfill
\begin{tikzpicture}
  \begin{axis}[contouraxis, xmin=-2.4, xmax=2.4, ymin=-2.4, ymax=2.4,
      xtick={-2,2}, ytick={-2,2}, title={(2)}]
    \foreach \r in {0.25,0.5,0.75,1,1.25,1.5,1.75,2} {
      \addplot[domain=0:360, samples=121]({\r*cos(x)},{\r*sin(x)}); }
  \end{axis}
\end{tikzpicture}\hfill
\begin{tikzpicture}
  \begin{axis}[contouraxis, xmin=-2.4, xmax=2.4, ymin=-2.4, ymax=2.4,
      xtick={-2,2}, ytick={-2,2}, title={(3)}]
    \foreach \r in {0.707,1,1.225,1.414,1.581,1.732,1.871,2} {
      \addplot[domain=0:360, samples=121]({\r*cos(x)},{\r*sin(x)}); }
  \end{axis}
\end{tikzpicture}
\end{gallery}

\begin{gallery}
\begin{tikzpicture}
  \begin{axis}[contouraxis, xmin=-2.4, xmax=2.4, ymin=-2.4, ymax=2.4,
      xtick={-2,2}, ytick={-2,2}, title={(4)}]
    \addplot[domain=-2.3:2.3]{x};
    \addplot[domain=-2.3:2.3]{-x};
    \foreach \s in {1,-1} {
      \addplot[domain=-1.4:1.4, samples=81]({\s*cosh(x)},{sinh(x)});
      \addplot[domain=-1.4:1.4, samples=81]({sinh(x)},{\s*cosh(x)}); }
  \end{axis}
\end{tikzpicture}\hfill
\begin{tikzpicture}
  \begin{axis}[contouraxis, xmin=-2.4, xmax=2.4, ymin=-2.4, ymax=2.4,
      xtick={-2,2}, ytick={-2,2}, title={(5)}]
    \foreach \c in {1,2} {
      \addplot[domain=0.42:2.3, samples=101]{\c/x};
      \addplot[domain=-2.3:-0.42, samples=101]{\c/x};
      \addplot[domain=0.42:2.3, samples=101]{-\c/x};
      \addplot[domain=-2.3:-0.42, samples=101]{-\c/x}; }
  \end{axis}
\end{tikzpicture}\hfill
\begin{tikzpicture}
  \begin{axis}[contouraxis, xmin=-2.4, xmax=2.4, ymin=-2.4, ymax=2.4,
      xtick={-2,2}, ytick={-2,2}, title={(6)}]
    \foreach \v in {-2,-1.414,-1,1,1.414,2} {
      \addplot[samples=2] coordinates {(\v,-2.3) (\v,2.3)}; }
  \end{axis}
\end{tikzpicture}
\end{gallery}

In problems 31--38, describe the stated sub- or superlevel set as a region,
and say whether it is convex.

\begin{problems}[resume]
  \item $W(4)$ for $f(x,y)=x^{2}+y^{2}$.
  \item $B(4)$ for $f(x,y)=x^{2}+y^{2}$.
  \item $W(6)$ for $f(x,y)=2x+3y$.
  \item $B(4)$ for $f(x,y)=xy$ restricted to $x,y>0$.
  \item $W(0)$ for $f(x,y)=x^{2}-y^{2}$.
  \item $B(1)$ for $f(x,y)=\sqrt{x^{2}+y^{2}}$.
  \item $W(P)$ for $f(x,y)=x^{2}+y^{2}$ with $P=(3,4)$.
  \item $B(P)$ for $f(K,L)=\sqrt{KL}$ with $P=(4,9)$, on $K,L>0$.
\end{problems}

\begin{problems}[resume]
  \item Sketch the level sets $W(2)$, $B(2)$, $W(-1)$ and $B(-1)$ of the
        one-variable function $g(x)=x^{3}-3x$ as subsets of the $x$-axis,
        and say which of the four are convex sets.
  \item Explain why for a function of one variable a level set
        $W(c)$ is convex exactly when it is an interval.
  \item A firm's isoquants are the curves $KL=\bar Q$. Sketch three of them
        for $\bar Q=1,4,9$ on $K,L>0$, and say whether the isoquants get
        closer together or further apart as output rises along the ray
        $K=L$.
  \item For the budget set $\{2x+5y\le100,\ x,y\ge0\}$: give the two
        intercepts, sketch the set, and say why it is convex.
  \item The indifference curves of $U(x,y)=xy$ and the budget line
        $2x+5y=100$ are drawn on the same axes. Describe, in words and
        without computing, the geometric condition that locates the optimal
        bundle.
  \item Two contour maps show circles as level curves. In the first the
        circles are equally spaced; in the second they crowd together
        outwards. Which is the cone and which the paraboloid, and why?
  \item A contour map has no closed loops and no crossings anywhere, only
        parallel straight lines. What can be said about the function?
\end{problems}

\section{Solutions}

\begin{solutions}
  \item $16-x^{2}-y^{2}\ge0\iff x^{2}+y^{2}\le16$: the closed disc of
        radius $4$ about the origin.
  \item $x-y>0$: the open half-plane strictly below the line $y=x$.
  \item $x^{2}+y^{2}\neq1$: the whole plane except the unit circle.
  \item $x\neq y$: the plane without the line $y=x$.
  \item $x\ge0$ and $y\ge0$: the closed non-negative quadrant.
  \item $xy>0$: the open first and third quadrants (the axes excluded).
  \item $x^{2}+y^{2}\ge4$: everything outside or on the circle of radius
        $2$ --- the complement of an open disc.
  \item The exponents are not integers, so $K\ge0$ and $L\ge0$: the closed
        non-negative quadrant. Economically the same set, usually taken
        open ($K,L>0$) so that marginal products exist.
  \item $y=\dfrac{4-x}{2}=2-\tfrac12x$: a straight line of slope
        $-\tfrac12$. All level sets are parallel lines.
  \item $x^{2}+y^{2}=9$, i.e. $y=\pm\sqrt{9-x^{2}}$: the circle of radius
        $3$.
  \item $x^{2}+y^{2}=0$ forces $x=y=0$: the level set is the single point
        $(0,0)$, the bottom of the bowl.
  \item $\dfrac{x^{2}}{4}+y^{2}=1$, so $y=\pm\sqrt{1-x^{2}/4}$: an ellipse
        with semi-axes $2$ and $1$.
  \item $y=6/x$: a hyperbola in the first and third quadrants, with both
        axes as asymptotes.
  \item $x^{2}=y^{2}\iff y=\pm x$: the \emph{two crossing lines} through
        the origin. This degenerate level is the signature of a saddle.
  \item $x^{2}-y^{2}=4$, so $y=\pm\sqrt{x^{2}-4}$, defined for
        $|x|\ge2$: a hyperbola opening left and right.
  \item $y=x^{2}+2$: a parabola, the standard parabola shifted up by $2$.
        All level sets are vertical translates of it.
  \item $\sqrt{x^{2}+y^{2}}=3\iff x^{2}+y^{2}=9$: the circle of radius
        $3$. Note that for the cone the radius equals $c$, not
        $\sqrt{c}$.
  \item $\mathrm{e}^{x+y}=1\iff x+y=0$, so $y=-x$: a straight line. Every
        level set of $\mathrm{e}^{x+y}$ is a line $x+y=\ln c$, defined only
        for $c>0$.
  \item $\sqrt{KL}=10\iff KL=100\iff L=100/K$: a hyperbola --- the
        isoquant for output $10$.
  \item $y=\dfrac{12-2x}{3}=4-\tfrac23x$: a straight line. For a cost
        function this is an \textbf{isocost line}; the analogous level
        curve of a budget constraint is the budget line.
  \item $f(x,y)=x^{3}-3x+y^{2}$ gives
        $f(P_{1})=0$, \;$f(P_{2})=1-3+0=-2$, \;$f(P_{3})=-1+3+0=2$,
        \;$f(P_{4})=8-6+1=3$, \;$f(P_{5})=0-0+4=4$,
        \;$f(P_{6})=-8+6+1=-1$.
  \item None of them share a level: the six values
        $0,-2,2,3,4,-1$ are distinct. $P_{3}$ lies on the drawn contour
        $c=2$, $P_{5}$ on $c=4$ and $P_{6}$ on $c=-1$; $P_{1}$, $P_{2}$ and
        $P_{4}$ lie between drawn contours.
  \item The contour $c=2$ crosses itself at $P_{3}=(-1,0)$. A
        self-crossing level curve marks a \textbf{saddle point}: from
        $(-1,0)$ the surface rises in two opposite directions and falls in
        the two others.
  \item The level $c=-2$ is a single point, $(1,0)$. Nested contours
        shrinking to a point mark a \textbf{local extremum}; since the
        levels decrease inwards, it is a local \emph{minimum}, with value
        $-2$.
  \item From the picture: a local minimum at $(1,0)$ (contours closing in
        on it) and a saddle at $(-1,0)$ (the crossing contour). There is no
        local maximum --- no contours close in on a point with the labels
        increasing inwards.
  \item $f_{x}=3x^{2}-3=0\Rightarrow x=\pm1$; $f_{y}=2y=0\Rightarrow y=0$.
        Critical points $(1,0)$ and $(-1,0)$. With $f_{xx}=6x$,
        $f_{yy}=2$, $f_{xy}=0$: at $(1,0)$ the determinant is $12>0$ with
        $f_{xx}>0$, a local minimum; at $(-1,0)$ it is $-12<0$, a saddle.
  \item No: $W(-1)$ has two pieces --- a small oval around the minimum
        $(1,0)$ and an unbounded piece on the far left. A sublevel set of a
        non-convex function may fall apart into components, which is
        exactly what convexity of $f$ would forbid.
  \item Away from $P_{2}=(1,0)$ the contours are closed ovals around it, so
        $f$ increases in \emph{every} direction, fastest where the ovals
        are most tightly packed --- to the right. Checking,
        $\grad f=(3x^{2}-3,\,2y)$, so at $(1.5,0.5)$,
        $\grad f=(3.75,\,1)$: pointing right and slightly up, and
        perpendicular to the contour there.
  \item (a)$\to$(1) plane, parallel lines. \;(b)$\to$(3) paraboloid:
        circles of radius $\sqrt{c}$, crowding outwards.
        \;(c)$\to$(2) cone: circles of radius $c$, equally spaced.
        \;(d)$\to$(4) saddle $x^{2}-y^{2}$: hyperbolas with the diagonals
        $y=\pm x$ as the degenerate level. \;(e)$\to$(5) $xy$: hyperbolas
        with the \emph{axes} as asymptotes. \;(f)$\to$(6) parabolic
        cylinder: pairs of vertical lines.
  \item (a) linear, hence \textbf{both} convex and concave. (b) convex (a
        bowl). (c) convex (a cone opening upwards; not differentiable at
        the origin, but still convex). (d) neither --- a saddle.
        (e) neither on $\R^{2}$; on the positive quadrant it is
        quasiconcave. (f) convex: $x^{2}$ is convex and adding no
        $y$-dependence preserves it.
  \item The closed disc $x^{2}+y^{2}\le4$ of radius $2$. Convex.
  \item The complement of the open disc, $x^{2}+y^{2}\ge4$: everything
        outside the circle. \textbf{Not} convex --- the segment joining
        $(-2,0)$ to $(2,0)$ passes through the origin, which is not in the
        set.
  \item The half-plane $2x+3y\le6$, bounded by the line through $(3,0)$ and
        $(0,2)$. Convex: every half-plane is.
  \item $\{(x,y): x,y>0,\ xy\ge4\}$, the region on the far side of the
        hyperbola $y=4/x$. Convex --- this is precisely the statement that
        $xy$ is quasiconcave on the quadrant.
  \item $x^{2}-y^{2}\le0\iff|y|\ge|x|$: the two opposite ``vertical'' sectors
        between the lines $y=\pm x$. Not convex: $(0,1)$ and $(0,-1)$ lie
        in it, but the midpoint $(0,0)$ is on the boundary and points just
        off the segment leave the set --- more simply, the set is not even
        connected in the way a convex set must be, since joining
        $(0,1)$ to $(1,2)$ stays inside but joining $(0,1)$ to $(0,-1)$
        passes through the pinch at the origin only.
  \item $\sqrt{x^{2}+y^{2}}\ge1$: everything outside or on the unit circle.
        Not convex, for the same reason as problem 32.
  \item $f(P)=9+16=25$, so $W(P)=\{x^{2}+y^{2}\le25\}$, the closed disc of
        radius $5$ through $P$. Convex.
  \item $f(P)=\sqrt{36}=6$, so
        $B(P)=\{(K,L):K,L>0,\ \sqrt{KL}\ge6\}=\{KL\ge36\}$: the region
        beyond the isoquant $L=36/K$. Convex.
  \item $g(x)=x^{3}-3x$ has $g(2)=2$, $g(-1)=2$, $g(1)=-2$, $g(-2)=-2$.
        $W(2)=\{x\le2\}$ minus nothing --- solving $x^{3}-3x\le2$ gives
        $(x+1)^{2}(x-2)\le0$, i.e. $x\le2$: the ray $(-\infty,2]$,
        \textbf{convex}. $B(2)=[2,\infty)\cup\{-1\}$: the ray from $2$
        \emph{plus the isolated point} $x=-1$, \textbf{not convex}.
        $W(-1)$: solving $x^{3}-3x\le-1$ gives
        $(-\infty,-1.879]\cup[0.347,1.532]$, two intervals, \textbf{not
        convex}. $B(-1)=[-1.879,0.347]\cup[1.532,\infty)$, also two
        pieces, \textbf{not convex}.
  \item In $\R^{1}$ a set is convex exactly when it contains the whole
        segment between any two of its points, and a subset of the line
        with that property is precisely an interval (possibly unbounded or
        empty). So ``convex'' and ``interval'' are the same statement one
        dimension down --- which is why the one-variable picture is the
        right place to first meet the definition.
  \item $L=\bar Q/K$: three hyperbola branches, $L=1/K$, $L=4/K$ and
        $L=9/K$. Along the ray $K=L$ they cut it at
        $K=1,2,3$: the isoquants for equally spaced \emph{square roots} of
        output, so for equally spaced output levels they move
        \emph{further apart}. Rising output requires ever larger input
        increases --- diminishing returns along the ray.
  \item Intercepts $(50,0)$ and $(0,20)$. The set is the triangle with
        vertices $(0,0)$, $(50,0)$, $(0,20)$. It is convex because it is an
        intersection of three half-planes ($2x+5y\le100$, $x\ge0$,
        $y\ge0$), and any intersection of convex sets is convex.
  \item The optimum is where the budget line is \textbf{tangent} to an
        indifference curve: the highest curve of the family that still
        touches the triangle touches it at exactly one point, and there the
        two curves have the same slope. Crossing would mean a higher curve
        is still reachable; missing the set entirely means the bundle is
        unaffordable.
  \item Equally spaced circles are the \textbf{cone}
        $\sqrt{x^{2}+y^{2}}$, whose level $c$ is a circle of radius $c$ ---
        constant slope, so equal height steps need equal horizontal steps.
        Circles crowding outwards are the \textbf{paraboloid}
        $x^{2}+y^{2}$, level $c$ having radius $\sqrt{c}$: the surface gets
        steeper, so equal height steps need ever smaller horizontal steps.
  \item Parallel straight contours, equally spaced or not, mean the
        function depends on the two variables only through one linear
        combination: $f(x,y)=h(ax+by)$ for some one-variable $h$. If the
        spacing is also constant, $h$ is linear and the surface is a plane.
        There are no critical points of extremum type anywhere.
\end{solutions}

\section{Additional Problems}

\begin{problems}
  \item Domain of $f(x,y)=\dfrac{\sqrt{y-x^{2}}}{x}$, as a region.
  \item Domain of $f(x,y)=\ln(4-x^{2}-y^{2})$, as a region.
  \item Domain and range of $f(x,y)=\mathrm{e}^{-(x^{2}+y^{2})}$.
  \item Level set of $f(x,y)=x^{2}+y^{2}-4x$ at $c=0$; name and locate the
        curve.
  \item Level set of $f(x,y)=x^{2}+y^{2}-4x$ at $c=-4$.
  \item Level set of $f(x,y)=y\,\mathrm{e}^{x}$ at $c=1$.
  \item Level set of $f(x,y)=\dfrac{y}{x}$ at $c=2$, and why the origin is
        excluded.
  \item Level sets of $f(x,y)=\max(x,y)$ at $c=1$: sketch and describe.
  \item For $f(x,y)=x^{2}+4y^{2}$, sketch the contours $c=4,16,36$ and say
        whether they crowd inwards or outwards.
  \item For $f(x,y)=x-y^{2}$, give the contours at $c=-1,0,1$ and describe
        the family.
  \item Give a function whose level curves are concentric circles that
        crowd \emph{inwards} towards the origin.
  \item Give a function of two variables whose level sets are all straight
        lines through the origin.
  \item $W(1)$ and $B(1)$ for $f(x,y)=|x|+|y|$: describe both and say which
        is convex.
  \item $W(0)$ for $f(x,y)=xy$ on all of $\R^{2}$: describe it and decide
        convexity.
  \item $B(2)$ for $f(x,y)=x+y$: describe it and decide convexity.
  \item $W(4)$ for $f(x,y)=x^{2}$: describe it and decide convexity.
  \item For $f(x,y)=x^{2}+y^{2}$ and $P=(1,2)$, give $W(P)$ and $B(P)$ and
        the level curve through $P$.
  \item A function has convex sublevel sets \emph{and} convex superlevel
        sets at every level. What kind of function can it be? Give an
        example.
  \item Sketch the surface $f(x,y)=4-x^{2}-y^{2}$ in words, give its
        contours, and classify it as convex or concave.
  \item Two functions have exactly the same family of level curves but
        different surfaces. Give such a pair and explain how the contour
        maps can nevertheless be told apart.
  \item Classify $f(x,y)=x^{2}+y^{2}$, $g(x,y)=-x^{2}-y^{2}$,
        $h(x,y)=x^{2}-y^{2}$ and $k(x,y)=3x-2y$ as convex, concave, both or
        neither.
  \item Which of the four functions in the previous problem are
        quasiconcave?
  \item The isoquants of a production function are straight lines with
        slope $-1$. What does that say about substitutability of the two
        inputs, and give such a production function.
  \item An indifference map consists of circles centred at $(4,4)$, with
        utility \emph{increasing} inwards. Interpret the point $(4,4)$
        economically and say why such preferences are unusual.
  \item On the contour map of problem 21, a walker starts at $P_{6}=(-2,1)$
        and always walks perpendicular to the contours, downhill. Describe
        where they end up and why the answer is not the minimum at
        $(1,0)$.
\end{problems}

\section{Additional Solutions}

\begin{solutions}
  \item $y-x^{2}\ge0$ and $x\neq0$: the region on or above the parabola
        $y=x^{2}$, with the $y$-axis removed.
  \item $4-x^{2}-y^{2}>0\iff x^{2}+y^{2}<4$: the \emph{open} disc of
        radius $2$ (boundary excluded, because the logarithm needs a
        strictly positive argument).
  \item $D=\R^{2}$. The exponent $-(x^{2}+y^{2})$ runs over $(-\infty,0]$,
        so the range is $(0,1]$ --- a bell-shaped hill with peak $1$ at the
        origin.
  \item $x^{2}-4x+y^{2}=0\iff(x-2)^{2}+y^{2}=4$: the circle of radius $2$
        centred at $(2,0)$, passing through the origin.
  \item $(x-2)^{2}+y^{2}=0$: the single point $(2,0)$ --- the minimum of
        the shifted bowl, with value $-4$.
  \item $y\,\mathrm{e}^{x}=1\iff y=\mathrm{e}^{-x}$: a decaying
        exponential curve.
  \item $y/x=2\iff y=2x$ with $x\neq0$: the line of slope $2$ through the
        origin, \emph{minus} the origin itself, since $f$ is undefined
        where $x=0$.
  \item $\max(x,y)=1$ is the union of the ray $\{x=1,\ y\le1\}$ and the ray
        $\{y=1,\ x\le1\}$: an L-shaped corner opening down-left, with its
        vertex at $(1,1)$.
  \item Ellipses $\dfrac{x^{2}}{c}+\dfrac{4y^{2}}{c}=1$ with semi-axes
        $\sqrt{c}$ and $\sqrt{c}/2$: for $c=4,16,36$ the semi-axes are
        $2,4,6$ and $1,2,3$. Equally spaced, in fact --- but for
        \emph{equally spaced} levels the contours would crowd outwards, as
        for any quadratic bowl.
  \item $x=y^{2}+c$: parabolas opening to the right, the whole family being
        horizontal translates of $x=y^{2}$ by $c$. For $c=-1,0,1$ the
        vertices sit at $(-1,0)$, $(0,0)$, $(1,0)$.
  \item Any function whose height grows ever more slowly with the radius,
        for instance $f(x,y)=\sqrt[4]{x^{2}+y^{2}}$ or
        $f(x,y)=\ln(x^{2}+y^{2})$: the level $c$ has radius growing faster
        than linearly in $c$, so equal height steps take ever larger
        radial steps and the contours thin out --- equivalently they crowd
        towards the origin.
  \item $f(x,y)=y/x$ (defined for $x\neq0$): the level $c$ is the punctured
        line $y=cx$. Any function of the ratio $y/x$ alone has this
        property; such functions are exactly the homogeneous functions of
        degree $0$.
  \item $W(1)=\{|x|+|y|\le1\}$ is the closed diamond (square rotated
        $45^{\circ}$) with vertices $(\pm1,0)$, $(0,\pm1)$ --- convex.
        $B(1)$ is everything outside it --- not convex.
  \item $xy\le0$ is the union of the second and fourth quadrants. Not
        convex: $(-1,1)$ and $(1,-1)$ both lie in it, but their midpoint
        $(0,0)$ does lie in it too --- take instead $(-1,2)$ and $(2,-1)$,
        whose midpoint $(0.5,0.5)$ has $xy>0$ and is outside.
  \item $x+y\ge2$: a closed half-plane, convex.
  \item $x^{2}\le4\iff-2\le x\le2$: the vertical strip between the lines
        $x=-2$ and $x=2$, convex (an intersection of two half-planes).
  \item $f(P)=1+4=5$: the level curve is the circle of radius $\sqrt5$,
        $W(P)$ the closed disc of radius $\sqrt5$, and $B(P)$ everything
        outside or on that circle.
  \item If every sublevel \emph{and} every superlevel set is convex, the
        function is both quasiconvex and quasiconcave. The level sets must
        then be ``flat'': any function of the form $h(ax+by)$ with $h$
        monotone qualifies, the simplest being a linear function
        $f(x,y)=ax+by$.
  \item A dome: the paraboloid turned upside down, with its peak of height
        $4$ above the origin and falling away in every direction. Contours
        are circles $x^{2}+y^{2}=4-c$ of radius $\sqrt{4-c}$, shrinking as
        $c$ rises to $4$. The function is \textbf{concave}.
  \item $f(x,y)=x^{2}+y^{2}$ and $g(x,y)=\sqrt{x^{2}+y^{2}}$ both have
        concentric circles as level sets. The maps differ in the
        \emph{labels}: for equally spaced levels $f$'s contours crowd
        outwards while $g$'s stay equally spaced. A contour map without its
        level labels does not determine the function.
  \item $f$ convex; $g$ concave; $h$ neither (a saddle); $k$ linear, hence
        both.
  \item $f$: quasiconvex, not quasiconcave (its superlevel sets are
        complements of discs). $g$: concave, hence quasiconcave. $h$:
        neither. $k$: linear, hence both quasiconvex and quasiconcave.
  \item Straight isoquants of slope $-1$ mean the marginal rate of
        technical substitution is constant at $1$: the inputs are
        \textbf{perfect substitutes}, one unit of either always replacing
        one unit of the other. The production function is
        $Q(K,L)=K+L$ (or any increasing function of $K+L$).
  \item $(4,4)$ is a \textbf{bliss point}: a most-preferred bundle, beyond
        which more of either good makes the consumer worse off. Standard
        preferences are monotone --- more is better --- so utility should
        increase without bound in the direction of both goods; circular
        indifference curves with utility rising inwards describe satiation
        instead.
  \item Walking perpendicular to the contours downhill from $P_{6}=(-2,1)$
        means following $-\grad f$. At $P_{6}$,
        $\grad f=(3\cdot4-3,\,2)=(9,2)$, so the walker moves left and
        slightly down --- into the region $x<-2$, where the contours of
        \emph{lower} level continue without bound. The function
        $x^{3}-3x+y^{2}$ tends to $-\infty$ as $x\to-\infty$, so the
        descent never ends: the walker never reaches the local minimum at
        $(1,0)$, because gradient descent finds a \emph{local} basin only
        if it starts in one.
\end{solutions}

\chapter{First-Order Partial Derivatives and the Gradient}

A function of one variable has a graph in the plane and a derivative that is
the slope of a tangent line. A function of two variables has a graph that is
a \emph{surface} in space, and there is no single slope any more: standing
at a point of a hillside, the ground rises steeply in one direction, falls
in another, and is level in a third. The first task of multivariable
calculus is to make that observation quantitative.

The device is to look one direction at a time. Freeze $y$ at a value and let
only $x$ move: the surface is cut by a vertical plane, the cut is an
ordinary curve, and the slope of \emph{that} curve is the partial derivative
$\partial f/\partial x$. This is the \emph{ceteris paribus} principle of
economics, and it is why partial derivatives are the natural language for
marginal products and marginal utilities.

Collecting the partial derivatives into one vector gives the gradient. It
turns out to carry all first-order information at once: it points in the
direction of steepest ascent, its length is the steepest slope available,
and it stands perpendicular to the level curve through the point. Those
three facts are proved below rather than asserted --- the third one in
particular is what the Lagrange method rests on several chapters later.

\section{Core Rules}

\begin{rulebox}{Topological vocabulary of domains}
\begin{itemize}[leftmargin=1.6em,itemsep=0.3ex]
  \item The \emph{open ball} (or \emph{neighbourhood}) of radius $r>0$ around $\mathbf{x}_0\in\R^n$ is
        $B_r(\mathbf{x}_0)=\{\mathbf{x}\in\R^n :
        \|\mathbf{x}-\mathbf{x}_0\|<r\}$.
  \item $\mathbf{x}_0\in M$ is an \emph{interior point} of $M$ if some open ball around
        $\mathbf{x}_0$ lies entirely in $M$.
  \item $M$ is \emph{open} if \emph{every} point of $M$ is
        an interior point.
  \item $\mathbf{x}_0$ is a \emph{boundary point} of
        $M$ if every ball around $\mathbf{x}_0$ contains points of $M$
        \emph{and} points of its complement.
  \item $M$ is \emph{closed} if it contains all
        its boundary points; equivalently, the complement
        $\R^n\setminus M$ is open.
  \item $M$ is \emph{bounded} if
        $M\subseteq B_r(\mathbf{0})$ for some radius $r$.
  \item $M\subseteq\R^n$ is \emph{compact} exactly when
        $M$ is closed \emph{and} bounded.
\end{itemize}
Economic examples: the open positive orthant $\{(x,y): x>0,\ y>0\}$
(the natural Cobb--Douglas domain) is open and unbounded; a budget set
$\{\mathbf{x}\in\R^n : \mathbf{p}\cdot\mathbf{x}\le m,\
\mathbf{x}\ge\mathbf{0}\}$ is closed and bounded, hence compact.
\end{rulebox}

\begin{rulebox}{Why it matters}
Differentiability (and hence every partial derivative below) is
defined at \emph{interior} points of the domain, where $f$ can be
probed in all directions. And by the \emph{Weierstrass theorem}, a
continuous function on a compact set attains its maximum and minimum:
this is why budget-constrained optimisation problems are
guaranteed to have solutions.
\end{rulebox}

\begin{rulebox}{Definition}
Let $f:\R^n\to\R$. The partial derivative of $f$ with respect to $x_i$ at
$\mathbf{x}=(x_1,\ldots,x_n)$ is
\[
  \pd{f}{x_i}(\mathbf{x}) \;=\; \lim_{h\to0}
    \frac{f(x_1,\ldots,x_i+h,\ldots,x_n)-f(x_1,\ldots,x_n)}{h} .
\]
Only the $i$-th slot changes; all other coordinates sit still. It is the
ordinary one-variable derivative of the function
$t\mapsto f(x_1,\ldots,t,\ldots,x_n)$, evaluated at $t=x_i$.
\end{rulebox}

\begin{rulebox}{What it measures: a slice of the surface}
Fix $y=y_{0}$. The vertical plane $y=y_{0}$ cuts the surface $z=f(x,y)$
along a curve, and along that curve only $x$ varies, so the curve is the
graph of the one-variable function $g(x)=f(x,y_{0})$. Then
\[
  \pd{f}{x}(x_{0},y_{0})=g'(x_{0}),
\]
the ordinary slope of the cut curve at $x_{0}$. Symmetrically,
$\partial f/\partial y$ is the slope of the cut made by the plane
$x=x_{0}$. A partial derivative is therefore never a slope \emph{of the
surface}; it is the slope of the surface \emph{along one chosen direction},
and different directions give different numbers. What the surface does in
all the remaining directions is the subject of the gradient below and of the
directional derivative.
\end{rulebox}

\begin{rulebox}{Practical computation}
To compute $\partial f/\partial x_i$, differentiate $f$ with respect to
$x_i$ and treat every other variable as a constant. All single-variable
rules (power, product, quotient, chain) apply unchanged. Worked example, for
$f(x,y)=x^{3}y^{2}+\mathrm{e}^{xy}-\ln y$:
\[
  \pd{f}{x}=3x^{2}y^{2}+y\,\mathrm{e}^{xy},
  \qquad
  \pd{f}{y}=2x^{3}y+x\,\mathrm{e}^{xy}-\frac{1}{y} .
\]
In the first computation $y$ is a number: $y^{2}$ is a constant factor, and
the chain rule on $\mathrm{e}^{xy}$ produces the inner derivative $y$. In
the second, $x$ is the constant and $\ln y$ finally contributes.
\end{rulebox}

\begin{rulebox}{Notation}
The following all refer to the same object:
\[
  \pd{f}{x} = f_x = \partial_x f = D_x f .
\]
\end{rulebox}

\begin{rulebox}{Gradient}
The vector of all first-order partial derivatives is
\[
  \grad f(\mathbf{x}) \;=\;
  \begin{pmatrix} f_{x_1}(\mathbf{x}) \\ f_{x_2}(\mathbf{x}) \\ \vdots \\
    f_{x_n}(\mathbf{x}) \end{pmatrix} .
\]
\end{rulebox}

\begin{rulebox}{Gradient and derivative: column and row}
The gradient is a \textbf{column} vector. The \textbf{derivative} of
$f\colon\R^{n}\to\R$ at $\mathbf{x}$, written $Df(\mathbf{x})$, is the same
list of partials arranged as a $1\times n$ \textbf{row}:
\[
  Df(\mathbf{x})=\bigl(f_{x_1}(\mathbf{x})\ \ f_{x_2}(\mathbf{x})\ \ \cdots\ \
  f_{x_n}(\mathbf{x})\bigr)=\grad f(\mathbf{x})^{T} .
\]
Same numbers, different shape. The row is the form that multiplies a column
of changes $\mathbf{h}$:
$Df(\mathbf{x})\,\mathbf{h}=\grad f(\mathbf{x})\cdot\mathbf{h}
=\sum_{i}f_{x_i}h_{i}$, which is the total differential of Chapter~12 and
the linear part of the tangent plane below. In Chapter~13 these rows are
stacked into the Jacobian matrix.
\end{rulebox}

\begin{rulebox}{Functions of $n$ variables written with sums}
Many economic functions are sums over goods, periods or inputs. To
differentiate with respect to one fixed variable $x_{k}$, ask which terms of
the sum contain $x_{k}$; every other term is a constant and drops out.
\begin{itemize}[leftmargin=1.6em,itemsep=0.3ex]
  \item \textbf{One function of each variable.} For
        $f=\sum_{i=1}^{n}g(x_{i})$ only the term $i=k$ survives:
        \[
          \pd{f}{x_{k}}=g'(x_{k}) .
        \]
        So $\sum_{i}x_{i}^{3}$ has $\partial f/\partial x_{k}=3x_{k}^{2}$.
  \item \textbf{The index inside the term.} $\sum_{i=1}^{n}x_{i}^{\,i}
        =x_{1}+x_{2}^{2}+\cdots+x_{n}^{n}$: the term containing $x_{k}$ is
        $x_{k}^{k}$, so $\partial f/\partial x_{k}=k\,x_{k}^{k-1}$.
        Likewise $\sum_{j=1}^{n}j/x_{j}^{4}=\sum_{j}j\,x_{j}^{-4}$ gives
        $\partial f/\partial x_{k}=-4k\,x_{k}^{-5}=-4k/x_{k}^{5}$.
  \item \textbf{The sum inside a function.} For $f=\bigl(\sum_{i}x_{i}\bigr)^{2}$
        use the chain rule; the inner derivative is
        $\partial\bigl(\sum_{i}x_{i}\bigr)/\partial x_{k}=1$, so
        $\partial f/\partial x_{k}=2\sum_{i=1}^{n}x_{i}$, the same for
        every $k$.
  \item \textbf{Products.} $f=\prod_{i}x_{i}$ gives
        $\partial f/\partial x_{k}=\prod_{i\neq k}x_{i}$: the product with
        the factor $x_{k}$ left out.
\end{itemize}
The summation index is a dummy and the differentiation variable is fixed:
call them by different letters ($i$ and $k$). Writing
$\partial/\partial x_{i}$ of a sum over $i$ is the standard source of
confusion.
\end{rulebox}

\begin{rulebox}{Exponentials to another base}
For $a>0$, $a^{x}=\mathrm{e}^{x\ln a}$, so
\[
  \frac{\mathrm{d}}{\mathrm{d}x}\,a^{x}=a^{x}\ln a,
  \qquad
  \pd{}{x}\,a^{g(x,y)}=a^{g(x,y)}\ln a\cdot g_{x}(x,y) .
\]
For example $\partial(2^{xy})/\partial x=y\,2^{xy}\ln2$. Keep the three
shapes apart: $x^{a}$ (variable base, power rule), $a^{x}$ (variable
exponent, factor $\ln a$), and $x^{y}$ (both, written as
$\mathrm{e}^{y\ln x}$).
\end{rulebox}

\begin{gallery}
\begin{tikzpicture}
  \begin{axis}[surfaxis, title={$f(K,L)=\sqrt{KL}$, sliced at $L=2$},
      xlabel={$K$}, ylabel={$L$}, zlabel={}, view={135}{26},
      zmin=0, zmax=3.2, z post scale=1.7]
    \addplot3[surf, domain=0.15:3, y domain=0.15:3, opacity=0.85]{sqrt(x*y)};
    \addplot3[samples=80, samples y=0, domain=0.15:3, line width=1.1pt,
      draw=black]({x},{2},{sqrt(2*x)});
  \end{axis}
\end{tikzpicture}\hfill
\begin{tikzpicture}
  \begin{axis}[contouraxis, xmin=-0.3, xmax=3.4, ymin=-0.3, ymax=3.4,
      xtick={1,2,3}, ytick={1,2,3},
      title={isoquants of $\sqrt{KL}$ with $\grad f$},
      xlabel={$K$}, ylabel={$L$}]
    \addplot[accent!40, line width=0.5pt, domain=0.22:3.3]{0.64/x};
    \addplot[accent!40, line width=0.5pt, domain=0.48:3.3]{1.44/x};
    \addplot[accent!40, line width=0.5pt, domain=0.85:3.3]{2.56/x};
    \addplot[accent!40, line width=0.5pt, domain=1.34:3.3]{4/x};
    \draw[-{Stealth[length=6pt,width=4pt]}, accent, line width=1pt]
      (axis cs:1,1) -- (axis cs:1.62,1.62);
    \draw[-{Stealth[length=6pt,width=4pt]}, accent, line width=1pt]
      (axis cs:2,0.5) -- (axis cs:2.16,1.14);
    \draw[-{Stealth[length=6pt,width=4pt]}, accent, line width=1pt]
      (axis cs:0.5,2) -- (axis cs:1.14,2.16);
    \addplot[only marks, mark=*, mark size=1pt, black, forget plot]
      coordinates {(1,1) (2,0.5) (0.5,2)};
  \end{axis}
\end{tikzpicture}
\gcap{Left: the plane $L=2$ cuts the surface; the slope of the cut curve is
$\partial f/\partial K$. Right: the gradient at three points, each one
perpendicular to the isoquant through it.}
\end{gallery}

\begin{rulebox}{Geometric meaning --- and why}
All three facts follow from one identity: the rate of change of $f$ in the
direction of a unit vector $\mathbf{u}$ is $\grad f\cdot\mathbf{u}$ (the
directional derivative below).
\begin{itemize}[leftmargin=1.6em,itemsep=0.4ex]
  \item \textbf{Steepest ascent.} Since
        $\grad f\cdot\mathbf{u}=\|\grad f\|\,\|\mathbf{u}\|\cos\theta
        =\|\grad f\|\cos\theta$, the rate is largest when $\cos\theta=1$,
        i.e. when $\mathbf{u}$ points \emph{along} $\grad f$. So
        $\grad f(\mathbf{x}_0)$ is the direction of steepest ascent, and
        $-\grad f$ of steepest descent.
  \item \textbf{Its length is that steepest rate.} Putting $\theta=0$ above
        leaves exactly $\|\grad f(\mathbf{x}_0)\|$.
  \item \textbf{Perpendicular to the level set.} Let $\mathbf{c}(t)$ be any
        curve running inside the level set through $\mathbf{x}_{0}$, so that
        $f(\mathbf{c}(t))$ is \emph{constant}. Differentiating with the
        chain rule gives
        \[
          0=\frac{\mathrm{d}}{\mathrm{d}t}f(\mathbf{c}(t))
           =\grad f(\mathbf{c}(t))\cdot\mathbf{c}'(t) .
        \]
        At $t_{0}$ this says $\grad f(\mathbf{x}_{0})$ is orthogonal to
        $\mathbf{c}'(t_{0})$, the tangent to the level set. Since the curve
        was arbitrary, the gradient is perpendicular to the whole level set.
\end{itemize}
The third fact is the one that matters most later: an optimum on a
constraint is a point where the level set of the objective and the
constraint set are \emph{tangent}, which is the same as their gradients
being parallel --- and that is the Lagrange condition.
\end{rulebox}

\begin{rulebox}{Economic interpretation}
\begin{itemize}[leftmargin=1.6em,itemsep=0.3ex]
  \item For a production function $Q=F(K,L)$: $F_K=\mathrm{MP}_K$
        (marginal product of capital), $F_L=\mathrm{MP}_L$ (marginal
        product of labour).
  \item For a utility function $U(x,y)$: $U_x, U_y$ are the marginal
        utilities; the marginal rate of substitution is
        $\mathrm{MRS}_{xy}=U_x/U_y$.
  \item For production: the marginal rate of technical substitution is
        $\mathrm{MRTS}_{LK}=F_L/F_K$.
\end{itemize}
\end{rulebox}

\begin{rulebox}{Directional derivative}
Let $\mathbf{u}\in\R^n$ be a \emph{unit} vector ($\|\mathbf{u}\|=1$). The
directional derivative of $f$ at $\mathbf{x}_0$ in the direction
$\mathbf{u}$ is
\[
  D_{\mathbf{u}}f(\mathbf{x}_0) \;=\; \grad f(\mathbf{x}_0)\cdot\mathbf{u} .
\]
\end{rulebox}

\begin{rulebox}{Caution}
If you are given a non-unit direction $\mathbf{v}$, you \emph{must}
normalise first: $\mathbf{u}=\mathbf{v}/\|\mathbf{v}\|$. Forgetting this is
the single most common error in directional-derivative problems.
\end{rulebox}

\begin{rulebox}{Differentiability in several variables}
Let $D\subseteq\R^{n}$ be open and $\mathbf{x}_{0}\in D$. The function $f$
is \textbf{differentiable} \de{(differenzierbar)} at $\mathbf{x}_{0}$ if
there is a vector $\mathbf{a}$ with
\[
  f(\mathbf{x}_{0}+\mathbf{h})=f(\mathbf{x}_{0})+\mathbf{a}\cdot\mathbf{h}
  +r(\mathbf{h}),
  \qquad
  \frac{r(\mathbf{h})}{\|\mathbf{h}\|}\to0\ \text{ as }\ \mathbf{h}\to\mathbf{0}.
\]
Near $\mathbf{x}_{0}$, $f$ is a linear function plus an error that vanishes
faster than the distance. For $n=2$ the graph of the linear part is a plane,
so differentiability says that the surface has a \textbf{tangent plane} at
the point: zoom in far enough and the surface looks flat.

If $f$ is differentiable at $\mathbf{x}_{0}$, then
\begin{itemize}[leftmargin=1.6em,itemsep=0.3ex]
  \item $f$ is continuous at $\mathbf{x}_{0}$;
  \item every partial derivative exists there, and
        $\mathbf{a}=\grad f(\mathbf{x}_{0})$ (take
        $\mathbf{h}=t\mathbf{e}_{k}$ and divide by $t$);
  \item every directional derivative exists and equals
        $\grad f(\mathbf{x}_{0})\cdot\mathbf{u}$.
\end{itemize}
\textbf{Sufficient criterion.} If all partials exist on an open
neighbourhood of $\mathbf{x}_{0}$ and are continuous there ($f$ is
$C^{1}$, \emph{continuously differentiable}), then $f$ is differentiable at
$\mathbf{x}_{0}$. Polynomials, exponentials, logarithms of positive
expressions, roots of positive expressions, quotients with non-zero
denominator, and all their sums, products and compositions are $C^{1}$ on
the open set where their formulas apply.

\textbf{Partials alone are not enough.} Take
$f(x,y)=\dfrac{xy}{x^{2}+y^{2}}$ with $f(0,0)=0$. On both axes $f=0$, so
$f_{x}(0,0)=f_{y}(0,0)=0$. But on the diagonal $y=x$,
$f=x^{2}/(2x^{2})=\tfrac12$ for every $x\neq0$. So $f$ is not even
continuous at the origin, and therefore not differentiable: the partials
see only two directions out of infinitely many.
\end{rulebox}

\begin{rulebox}{Where differentiability fails, and finding the set}
Compute the partials by the rules. The formulas are valid on the open set
where the rules applied; the points where a formula breaks down (a zero
denominator, a zero radicand, the kink of an absolute value or a minimum)
are the suspects, and each suspect is settled by the difference quotient.
\begin{itemize}[leftmargin=1.6em,itemsep=0.3ex]
  \item \textbf{Cone apex.} $f=\sqrt{(x-a)^{2}+(y-b)^{2}}$ has
        $f_{x}=(x-a)/f$, undefined at $(a,b)$. Directly,
        $f(a+h,b)-f(a,b)=\sqrt{h^{2}}=|h|$, and $|h|/h=\pm1$ has no limit:
        $f_{x}(a,b)$ does not exist, nor does $f_{y}(a,b)$. The same
        happens for the elliptic cone $\sqrt{4(x-a)^{2}+y^{2}}$ at $(a,0)$,
        with quotient $2|h|/h$. For the root of a sum of squares of
        \emph{linear} expressions, such as these cones, $f$ is
        differentiable exactly where the radicand is non-zero, and the
        graph has a sharp tip (or crease) where it vanishes. With higher
        powers under the root this can change: $\sqrt{x^{4}+y^{4}}$ is
        differentiable even at the origin, where its radicand is $0$.
  \item \textbf{$|x|+|y|$.} $f_{x}=1$ for $x>0$ and $-1$ for $x<0$; at
        $x=0$ the quotient is $|h|/h$ again. So $f_{x}$ fails on the whole
        $y$-axis and $f_{y}$ on the whole $x$-axis, and $f$ is
        differentiable exactly off the axes. The graph is an upside-down
        square pyramid, creased along both axes.
  \item \textbf{$\min\{ax,by\}$} with $a,b>0$. Where $ax<by$, $f=ax$;
        where $ax>by$, $f=by$. On the kink line $ax=by$ (in the quadrant
        $x,y>0$ of economics, a ray from the origin), raising $x$ leaves
        $f$ at $by$ (slope $0$) while lowering $x$ lowers $f$ at rate $a$:
        the one-sided slopes differ, so $f_{x}$ does not exist there, and
        nor does $f_{y}$. The graph is a roof whose ridge lies above that
        line.
\end{itemize}
A breaking formula is a warning, not a verdict: the partial may still
exist at that point, and only the difference quotient decides. And partials
that exist at a point still do not guarantee differentiability; continuity
of the partials near the point does.
\end{rulebox}

\begin{rulebox}{Proving differentiability from the definition}
When the $C^{1}$ criterion is not available, or a proof is asked for,
work with the remainder directly. At a point $\mathbf{x}^{*}$:
\begin{enumerate}[leftmargin=1.6em,itemsep=0.3ex]
  \item \textbf{Candidate.} Compute $\grad f(\mathbf{x}^{*})$; it is the
        only possible linear part. If a partial does not exist, stop: $f$
        is not differentiable.
  \item \textbf{Remainder.}
        $r(\mathbf{h})=f(\mathbf{x}^{*}+\mathbf{h})-f(\mathbf{x}^{*})
        -\grad f(\mathbf{x}^{*})\cdot\mathbf{h}$, expanded and simplified.
        For a polynomial, every term left over contains at least two
        factors $h_{i}$.
  \item \textbf{Estimate.} Show $|r(\mathbf{h})|/\|\mathbf{h}\|\to0$ with
        the standard inequalities
        \[
          |h_{i}|\le\|\mathbf{h}\|,\qquad
          \frac{h_{i}^{2}}{\|\mathbf{h}\|}\le\|\mathbf{h}\|,\qquad
          \frac{|h_{1}h_{2}|}{\|\mathbf{h}\|}\le\tfrac12\|\mathbf{h}\| ,
        \]
        and, for a product of three factors $h_{i}$, a bound
        $\|\mathbf{h}\|^{2}$.
  \item \textbf{Transcendental factors.} Replace them by their linear part
        plus a remainder: $\mathrm{e}^{t}=1+t+R(t)$ with $|R(t)|\le2t^{2}$
        for $|t|\le1$, and $\ln(1+t)=t+R(t)$ with $|R(t)|\le2t^{2}$ for
        $|t|\le\tfrac12$ (both from the Lagrange remainder of Chapter~4).
        The quadratic bound on $R$ is what makes $R/\|\mathbf{h}\|\to0$.
\end{enumerate}
\textbf{Roots at the origin.} For $f=(x^{2}+y^{2})^{p}$ with $p>0$,
$f(\mathbf{h})=\|\mathbf{h}\|^{2p}$. The difference quotient on an axis is
$|h|^{2p}/h$: it tends to $0$ if $2p>1$, and has no limit if $2p\le1$. So
$f$ is differentiable at the origin exactly when $p>\tfrac12$; the cone
$p=\tfrac12$ is the borderline case, and $p=\tfrac13$ is worse than a cone.
\end{rulebox}

\begin{rulebox}{Tangent plane and linearisation}
If $f$ is differentiable at $P=(x_{0},y_{0})$, the \textbf{tangent plane}
\de{(Tangentialebene)} to its graph at $(x_{0},y_{0},f(P))$ is
\[
  z=f(P)+f_{x}(P)\,(x-x_{0})+f_{y}(P)\,(y-y_{0}) .
\]
In $n$ variables, with $P=\mathbf{p}=(p_{1},\ldots,p_{n})$, the
\textbf{tangent hyperplane} is
\[
  z=f(\mathbf{p})+\grad f(\mathbf{p})\cdot(\mathbf{x}-\mathbf{p})
   =f(\mathbf{p})+\sum_{i=1}^{n}f_{x_i}(\mathbf{p})\,(x_{i}-p_{i}) .
\]
\begin{itemize}[leftmargin=1.6em,itemsep=0.3ex]
  \item The right-hand side is the \textbf{linearisation} $L(\mathbf{x})$
        of $f$ at $\mathbf{p}$, the multivariable version of the tangent
        line $f(a)+f'(a)(x-a)$ of Chapter~6. With
        $\mathrm{d}x_{i}=x_{i}-p_{i}$ it reads $L=f(\mathbf{p})+\mathrm{d}f$:
        the old value plus the total differential.
  \item The partials are evaluated \emph{at the point} $\mathbf{p}$; they
        are numbers, which is what makes the result a plane.
  \item For a point $Q$ near $P$, $f(Q)\approx L(Q)$. The difference
        $f(Q)-L(Q)$ is the remainder $r(\mathbf{h})$ of the definition, so
        comparing the two numbers measures how good the approximation is.
        For $f$ with continuous second partials it is at most of the order
        $\|Q-P\|^{2}$: halve the step and the error falls to about a
        quarter, or even less.
  \item Equivalently, $f_{x}(P)(x-x_{0})+f_{y}(P)(y-y_{0})-(z-f(P))=0$: the
        vector $(f_{x}(P),\,f_{y}(P),\,-1)$ is normal to the plane.
\end{itemize}
Where $f$ is not differentiable there is no tangent plane: a cone has none
at its apex.
\end{rulebox}

\section{Worked Examples}

\begin{worked}{For $f(x,y)=x^{3}y^{2}+\mathrm{e}^{2x}\ln y$, compute both partials and the gradient at $(1,1)$; then read the marginal products of $Q=10K^{0.3}L^{0.7}$}
\wstep{Step 1. Differentiating in $x$ means freezing $y$ --- literally.}
Treat $y$ as if it were the number $7$. Then $x^{3}y^{2}$ is
``$y^{2}$ times $x^{3}$'', a constant times a power, and $\ln y$ is a
constant multiplying $\mathrm{e}^{2x}$:
\[
  f_{x}=3x^{2}y^{2}+2\mathrm{e}^{2x}\ln y .
\]
\wtrap{Freezing means the frozen variable never produces a derivative
factor. $\frac{\partial}{\partial x}\bigl[\ln y\bigr]=0$, not $1/y$. The
symbol $\partial$ instead of $\mathrm{d}$ is the reminder that something is
being held fixed.}

\wstep{Step 2. Now freeze $x$ instead.}
$x^{3}$ is a constant multiplying $y^{2}$, and $\mathrm{e}^{2x}$ is a
constant multiplying $\ln y$:
\[
  f_{y}=2x^{3}y+\frac{\mathrm{e}^{2x}}{y} .
\]

\wstep{Step 3. Evaluate, and assemble the gradient.}
At $(1,1)$, with $\ln1=0$ and $\mathrm{e}^{2}\approx7.389$:
\[
  f_{x}(1,1)=3\cdot1\cdot1+2\mathrm{e}^{2}\cdot0=3,
  \qquad
  f_{y}(1,1)=2\cdot1\cdot1+\frac{\mathrm{e}^{2}}{1}=2+\mathrm{e}^{2}\approx9.389 .
\]
\[
  \grad f(1,1)=\begin{pmatrix}3\\2+\mathrm{e}^{2}\end{pmatrix}\approx\begin{pmatrix}3\\9.389\end{pmatrix}
\]
\wtrap{The gradient is a \textbf{vector}, and it is a vector in the
\emph{input} space --- here a direction in the $xy$-plane, not on the
surface. Reporting a single number, or a three-component object including
the height, are the two standard confusions.}

\wstep{Step 4. What each partial measures.}
$f_{x}$ is the slope of the curve obtained by slicing the surface with the
plane $y=1$ --- an ordinary one-variable derivative of the function
$x\mapsto f(x,1)$. Indeed $f(x,1)=x^{3}$, whose derivative at $x=1$ is $3$.
\checkmark\ Two partials, two slices, two ordinary derivatives.

\wstep{Step 5. What the gradient does that the partials alone do not.}
\begin{itemize}[leftmargin=1.6em,itemsep=0.2ex]
  \item It points in the direction of \textbf{steepest ascent}. Here that
        is roughly $(0.30,\,0.95)$ after normalising: mostly ``north''.
  \item Its length $\|\grad f\|=\sqrt{9+(2+\mathrm{e}^{2})^{2}}\approx9.86$
        is the steepest rate of climb available at that point.
  \item It is \textbf{perpendicular} to the level curve through the point.
\end{itemize}
The last is why the gradient governs constrained optimisation: at an
optimum the objective's gradient must be parallel to the constraint's, which
is the Lagrange condition.

\wstep{Step 6. Directional derivative: normalise, then take an inner product.}
In the direction of $\mathbf v=(3,4)$:
\[
  \mathbf u=\frac{\mathbf v}{\|\mathbf v\|}=\Bigl(\tfrac35,\tfrac45\Bigr),
  \qquad
  D_{\mathbf u}f=\grad f\cdot\mathbf u
  =\tfrac35\cdot3+\tfrac45(2+\mathrm{e}^{2})\approx9.31 .
\]
\wtrap{Normalise first. Using $(3,4)$ unnormalised gives $46.6$, which is
five times too large and has no meaning --- ``rate of change per unit
distance'' requires a unit vector. And $9.31<9.86$, as it must be: no
direction beats the gradient.}

\wstep{Step 7. The economic reading, on a production function.}
For $Q=10K^{0.3}L^{0.7}$:
\[
  MP_{K}=\pd{Q}{K}=3K^{-0.7}L^{0.7}=3\Bigl(\frac{L}{K}\Bigr)^{0.7},
  \qquad
  MP_{L}=\pd{Q}{L}=7\Bigl(\frac{K}{L}\Bigr)^{0.3} .
\]
At $K=100$, $L=50$: $MP_{K}=3(0.5)^{0.7}\approx1.85$ and
$MP_{L}=7(2)^{0.3}\approx8.62$. So one more unit of labour adds about $8.6$
units of output, one more of capital about $1.85$.

Two structural facts fall out. The \textbf{marginal rate of technical
substitution} is
\[
  \mathrm{MRTS}=\frac{MP_{L}}{MP_{K}}=\frac{0.7}{0.3}\cdot\frac{K}{L}\approx4.67 ,
\]
and each marginal product depends only on the \emph{ratio} $K/L$ --- the
signature of constant returns to scale, which is why Cobb--Douglas with
exponents summing to $1$ behaves the way it does.
\wtrap{$MP_{K}$ falls as $K$ rises (the exponent $-0.7$ is negative):
diminishing marginal product. But that is diminishing \emph{marginal},
not diminishing total --- output still increases. Confusing the two is the
most common misreading of a partial derivative in economics.}
\end{worked}

\begin{worked}{Find the tangent plane of $f(x,y)=x\ln(x+2y)$ at $P=(1,0)$, and test it at the nearby point $Q=(1.1,\,0.05)$}
\wstep{Step 1. Make sure there is a plane to find.}
The domain is the open half-plane $x+2y>0$, and $P$ lies in it since
$1+0>0$. On that half-plane $f$ is a polynomial times the logarithm of a
positive linear expression, so its partials exist and are continuous: $f$
is $C^{1}$, hence differentiable at $P$, and the tangent plane exists.

\wstep{Step 2. The height at $P$.}
$f(1,0)=1\cdot\ln1=0$. The plane touches the surface at $(1,0,0)$.

\wstep{Step 3. The two partials, then their values at $P$.}
In $x$ the function is a product of $x$ and $\ln(x+2y)$:
\[
  f_{x}=\ln(x+2y)+x\cdot\frac{1}{x+2y} .
\]
In $y$ the factor $x$ is a constant, and the chain rule supplies the inner
derivative $2$:
\[
  f_{y}=\frac{2x}{x+2y} .
\]
At $P$: $f_{x}(1,0)=\ln1+\tfrac11=1$ and $f_{y}(1,0)=\tfrac21=2$.
\wtrap{$x$ appears twice in $f$, once as a factor and once inside the
logarithm, so $f_{x}$ needs the product rule. Writing $f_{x}=x/(x+2y)$
alone loses the term $\ln(x+2y)$; here it happens to vanish at $P$, which
hides the error until the next point.}

\wstep{Step 4. Assemble the plane.}
\[
  z=f(P)+f_{x}(P)(x-1)+f_{y}(P)(y-0)=0+1\cdot(x-1)+2y=x+2y-1 .
\]
\wtrap{The partials go in as \emph{numbers} evaluated at $P$. Leaving them
as functions of $x$ and $y$ gives a curved surface, not a plane.}

\wstep{Step 5. Compare at $Q$.}
\[
  f(Q)=1.1\cdot\ln(1.1+0.1)=1.1\ln1.2\approx0.200554,
  \qquad
  L(Q)=1.1+0.1-1=0.200000 .
\]
The plane misses by about $0.000554$, while $f$ itself changed by about
$0.2$: the error is under $0.3\%$ of the change. The step has length
$\|Q-P\|=\sqrt{0.1^{2}+0.05^{2}}\approx0.112$. The error is far below that
length, and even below its square $0.0125$, as the rule box predicts: the
remainder vanishes faster than the distance.

\wstep{Step 6. The same number as a total differential.}
With $\mathrm{d}x=0.1$ and $\mathrm{d}y=0.05$,
\[
  \mathrm{d}f=f_{x}(P)\,\mathrm{d}x+f_{y}(P)\,\mathrm{d}y=0.1+0.1=0.2,
\]
and $L(Q)=f(P)+\mathrm{d}f=0.2$. \checkmark\ The tangent plane and the
total differential are one idea written two ways.
\end{worked}

\begin{worked}{Prove from the definition that $f(x,y)=x^{2}y$ is differentiable at every point; then decide at the origin for $(x^{2}+y^{2})^{1/3}$ and $(x^{2}+y^{2})^{2/3}$}
\wstep{Step 1. The candidate for the linear part.}
Fix a point $(a,b)$. If $f$ is differentiable there, the linear part must
be $\grad f(a,b)\cdot\mathbf{h}$, so compute the partials:
$f_{x}=2xy$, $f_{y}=x^{2}$, hence $\grad f(a,b)=(2ab,\,a^{2})$.

\wstep{Step 2. Form the remainder and let everything cancel that can.}
With $\mathbf{h}=(h_{1},h_{2})$,
\[
  f(a+h_{1},b+h_{2})=(a^{2}+2ah_{1}+h_{1}^{2})(b+h_{2})
  =a^{2}b+a^{2}h_{2}+2abh_{1}+2ah_{1}h_{2}+bh_{1}^{2}+h_{1}^{2}h_{2} .
\]
Subtract $f(a,b)=a^{2}b$ and the linear part $2abh_{1}+a^{2}h_{2}$:
\[
  r(\mathbf{h})=2a\,h_{1}h_{2}+b\,h_{1}^{2}+h_{1}^{2}h_{2} .
\]
Every surviving term contains at least \emph{two} factors $h_{i}$. That is
the sign that the candidate was right.

\wstep{Step 3. Bound the quotient by something that tends to $0$.}
Use $|h_{i}|\le\|\mathbf{h}\|$ and $|h_{1}h_{2}|\le\tfrac12\|\mathbf{h}\|^{2}$
(from $(|h_{1}|-|h_{2}|)^{2}\ge0$):
\[
  \frac{|r(\mathbf{h})|}{\|\mathbf{h}\|}
  \le\frac{2|a|\cdot\tfrac12\|\mathbf{h}\|^{2}+|b|\,\|\mathbf{h}\|^{2}
          +\|\mathbf{h}\|^{3}}{\|\mathbf{h}\|}
  =(|a|+|b|)\,\|\mathbf{h}\|+\|\mathbf{h}\|^{2}\;\longrightarrow\;0 .
\]
So $f$ is differentiable at $(a,b)$, for every $(a,b)$, with derivative
$(2ab,\,a^{2})$.
\wtrap{Checking $r/\|\mathbf{h}\|\to0$ along the axes, or along every
line, is not a proof: the bound must hold for \emph{all} small $\mathbf{h}$
at once. The bound may depend on the fixed point $(a,b)$, but only through
constants.}

\wstep{Step 4. A root at the origin: $g=(x^{2}+y^{2})^{1/3}$.}
First the partials, from the difference quotient:
\[
  \frac{g(h,0)-g(0,0)}{h}=\frac{(h^{2})^{1/3}}{h}
  =\frac{|h|^{2/3}}{h}=\pm|h|^{-1/3},
\]
which is unbounded as $h\to0$. So $g_{x}(0,0)$ does not exist (nor, by
symmetry, $g_{y}(0,0)$), and $g$ is \textbf{not differentiable} at the
origin, although it is continuous there. Its graph has a tip even sharper
than a cone's.

\wstep{Step 5. A higher power: $k=(x^{2}+y^{2})^{2/3}$.}
Now $\dfrac{k(h,0)-k(0,0)}{h}=\dfrac{|h|^{4/3}}{h}=\pm|h|^{1/3}\to0$, so both
partials are $0$ and the candidate linear part is $0$. The remainder is
$r(\mathbf{h})=k(\mathbf{h})=\|\mathbf{h}\|^{4/3}$, and
\[
  \frac{|r(\mathbf{h})|}{\|\mathbf{h}\|}=\|\mathbf{h}\|^{1/3}\to0 .
\]
So $k$ \textbf{is} differentiable at the origin, with the tangent plane
$z=0$. In general $(x^{2}+y^{2})^{p}=\|\mathbf{h}\|^{2p}$ at the origin, and
the verdict depends only on whether $2p>1$.
\end{worked}

\section{Practice Problems}

\begin{problems}
  \item Compute $f_x$ and $f_y$ for $f(x,y)=5x^4-2x^2y^3+y^5$.
  \item For $F(K,L)=12\,K^{1/4}L^{3/4}$, find $F_K$ and $F_L$ at
        $(K,L)=(16,81)$.
  \item Compute $f_x$ and $f_y$ for $f(x,y)=x^2e^{3y}$.
  \item Compute $\nabla f$ for $f(x,y)=4x^3y-xy^4+7$.
  \item Compute the directional derivative of $f(x,y)=x^2-4xy+y^3$ at
        $(2,1)$ in the direction $\mathbf{v}=(5,-12)$.
  \item For $F(K,L)=6K^{2/3}L^{1/3}$, compute $\mathrm{MP}_K$ and
        $\mathrm{MP}_L$ at $(K,L)=(27,8)$.
  \item Compute $f_x$ and $f_y$ for $f(x,y)=(3x-2y)^4$.
  \item Compute $f_x$ and $f_y$ for $f(x,y)=e^{xy^2}$.
  \item For $F(K,L)=25K^{0.6}L^{0.5}$, find $F_K$ and $F_L$.
  \item Compute $\nabla f$ for $f(x,y,z)=x^2+3yz-z^2e^x$.
  \item Compute the directional derivative of $f(x,y)=e^x\ln(1+y)$ at
        $(0,0)$ in the direction $\mathbf{v}=(1,2)$.
  \item For $U(x,y)=2\ln(x)+5\ln(y)$, derive $\mathrm{MRS}_{xy}=U_x/U_y$ as
        a function of $(x,y)$ and evaluate at $(3,5)$.
  \item Simplify and compute $f_x, f_y$ for
        $f(x,y)=\dfrac{x^2-y^2}{x+y}$.
  \item Compute $f_x$ and $f_y$ for $f(x,y)=\ln(x^2y+1)$.
  \item For $U(x_1,x_2,x_3)=x_1^{1/2}x_2^{1/3}x_3^{1/6}$, compute the
        three partials.
  \item Compute the directional derivative of $f(x,y,z)=xy-2yz+z^2$ at
        $(1,2,3)$ in the direction $\mathbf{v}=(1,2,2)$.
  \item For $F(K,L)=K^{0.25}L^{0.75}$, derive $\mathrm{MRTS}_{LK}=F_L/F_K$
        and evaluate at $(K,L)=(16,4)$.
  \item Expand and then compute $f_x, f_y$ for
        $f(x,y)=(x^2+y)(x-3y^2)$.
  \item Compute $f_x$ and $f_y$ for $f(x,y)=xe^{-y/x}$ with $x\neq0$.
  \item For $F(K,L)=(3K^{1/2}+2L^{1/2})^2$, find $F_K$ and $F_L$.
  \item Compute $\nabla f$ for $f(x,y)=\arctan(y/x)$ with $x>0$.
  \item For $f(x,y)=x^{1/2}y^{1/2}$ at $(K,L)=(4,9)$, find the unit
        direction of steepest ascent and the maximum rate of increase.
  \item For the CES function $F(K,L)=(K^{1/2}+L^{1/2})^2$, verify Euler's
        theorem $KF_K+LF_L=rF$ for the appropriate degree of homogeneity
        $r$.
  \item Compute $f_x$ and $f_y$ for $f(x,y)=\dfrac{2x+3y}{x^2+y^2}$.
  \item Compute $f_x$ and $f_y$ for $f(x,y)=\dfrac{e^x}{1+e^y}$.
  \item For $F(K,L)=A(\alpha K^\rho+(1-\alpha)L^\rho)^{\nu/\rho}$ (CES),
        find $F_L$.
  \item Compute the directional derivative of $f(x,y)=\ln(x^2+y^2)$ at
        $(3,4)$ in the direction $\mathbf{v}=(-1,1)$.
  \item Compute $f_x, f_y, f_z$ for $f(x,y,z)=x^2yz^3-4xy^2z+5z$.
  \item Compute $f_x$ and $f_y$ for $f(x,y)=\sqrt{4x^2+9y^2}$.
  \item For $F(K,L)=(K^{1/2}+L^{1/2})^4$, compute $F_K$ at
        $(K,L)=(4,9)$.
  \item Compute $\nabla f$ for $f(x,y,z)=\ln(x^2+y^2+z^2)$.
  \item Compute $f_x$ and $f_y$ for
        $f(x,y)=\ln\left(\sqrt{x^2+y^2}\right)$.
  \item Cost minimisation. A firm uses inputs $(x,y)$ with cost
        $C(x,y)=4x+9y$ and output $Q(x,y)=x^{1/2}y^{1/2}$. Compute
        $\nabla C$ and $\nabla Q$, and verify that at the cost-minimising
        input mix $\nabla C=\lambda\nabla Q$ for some $\lambda$.
  \item Compute $f_x$ and $f_y$ for $f(x,y)=(1+xy)^5$.
  \item Compute the directional derivative of $f(x,y)=x^3-2xy+y^2$ at
        $(1,2)$ in the direction $\mathbf{v}=(3,4)$.
  \item For $F(K,L)=10K^{0.3}L^{0.7}$ (Cobb--Douglas), compute $F_K$ and
        $F_L$ at $(K,L)=(100,100)$ and verify $KF_K+LF_L=F$.
  \item Compute $f_x$ and $f_y$ for $f(x,y)=\sin(xy)+\cos(x-y)$.
  \item Compute $\nabla f$ for $f(x,y)=x^2y-xy^2+3x-2y$ at the point
        $(1,-1)$.
  \item For $U(x,y)=x^\alpha y^{1-\alpha}$ (Cobb--Douglas,
        $0<\alpha<1$), derive $\mathrm{MRS}_{xy}=U_x/U_y$ and show it
        equals $\dfrac{\alpha}{1-\alpha}\cdot\dfrac{y}{x}$.
  \item Compute the directional derivative of $f(x,y)=e^{xy}$ at $(1,0)$
        in the direction $\mathbf{v}=(-3,4)$.
  \item Compute $f_x$ and $f_y$ for $f(x,y)=x^y$ (with $x>0$).
        \emph{Hint:} write $x^y=e^{y\ln x}$.
  \item For $F(K,L)=\min\{2K,3L\}$ (Leontief, fixed-proportions, on the
        smooth region away from the kink), compute $F_K$ and $F_L$ in
        each of the two regions $\{2K<3L\}$ and $\{2K>3L\}$.
  \item Compute $\nabla f$ for $f(x,y,z)=xyz+x^2-z^2$ at the point
        $(2,-1,3)$.
  \item Compute the unit vector of steepest \emph{descent} for
        $f(x,y)=x^2+4y^2$ at $(3,1)$, and the corresponding rate of
        decrease.
  \item For $f(x,y)=xe^y-ye^x$, compute $f_x, f_y$.
  \item Compute $f_x$ and $f_y$ for $f(x,y)=\dfrac{x}{x^2+y^2}$.
  \item For $U(x,y)=x+2\sqrt{y}$ (quasi-linear), derive
        $\mathrm{MRS}_{xy}$.
  \item Compute the directional derivative of $f(x,y,z)=x^2+y^2-z^2$ at
        $(2,-1,2)$ in the direction $\mathbf{v}=(2,2,1)$.
  \item Compute $f_x$ and $f_y$ for $f(x,y)=\tan(x+y)+\tan(x-y)$.
  \item For $F(K,L)=(\alpha K+\beta L)^r$ (linear-aggregate,
        $\alpha,\beta,r>0$), compute $F_K, F_L$, and verify
        $KF_K+LF_L=rF$.
  \item Compute $\nabla f$ for $f(x,y)=e^{-x}\sin(y)+e^x\cos(y)$.
  \item Compute the directional derivative of $f(x,y)=\sqrt{x^2+y^2}$ at
        $(3,4)$ in the direction $\mathbf{v}=(1,0)$.
  \item For $U(x,y)=\min\{2x,3y\}$, compute the marginal utilities on
        each smooth region $\{2x<3y\}$ and $\{2x>3y\}$.
  \item Compute $f_x, f_y$ for $f(x,y)=\dfrac{\ln(x+y)}{x-y}$ (assume
        $x>y>0$).
  \item For $F(K,L)=AK^\alpha L^\beta$ general Cobb--Douglas, derive
        $\mathrm{MRTS}_{LK}$ and show it depends only on the ratio $K/L$
        (not on scale).
  \item Compute the unit direction of steepest ascent for
        $f(x,y)=x^2-y^2$ at $(2,-1)$, and the maximum rate.
  \item Compute $f_x$ and $f_y$ for $f(x,y)=\cosh(x+y)+\sinh(x-y)$.
  \item For $f(x,y,z)=e^{x+y+z}$, compute $\nabla f$ at $(0,0,0)$ and the
        directional derivative in $\mathbf{v}=(1,1,1)$.
  \item Compute $f_x$ and $f_y$ for $f(x,y)=(\ln x)(\ln y)$ with
        $x,y>1$.
  \item For $F(K,L)=K+2K^{1/2}L^{1/2}+L$ (which equals
        $(K^{1/2}+L^{1/2})^2$), compute $F_K, F_L$ directly and confirm
        agreement with the CES form.
  \item Compute the directional derivative of $f(x,y)=x^2y+y^3$ at
        $(-1,2)$ in the direction $\mathbf{v}=(4,-3)$.
  \item For $U(x,y)=\alpha\ln x+\beta\ln y$ (log-Cobb--Douglas), derive
        $\mathrm{MRS}_{xy}$ and evaluate at $(x,y)=(\alpha,\beta)$.
  \item Compute $\nabla f$ for $f(x,y)=(x^2+y^2)^{-1/2}$ (Newtonian
        potential in 2D).
  \item Compute $f_x$ and $f_y$ for $f(x,y)=\arctan(xy)$.
  \item For $F(K,L)=(K^\rho+L^\rho)^{1/\rho}$ with $\rho<1$, $\rho\neq0$
        (CES with constant returns), compute $F_K$ and show
        $F_K=(K/F)^{\rho-1}$.
  \item Compute the directional derivative of $f(x,y,z)=xyz$ at
        $(2,3,-1)$ in the direction $\mathbf{v}=(2,-1,2)$.
\end{problems}

\section{Solutions}

\begin{solutions}
  \item $f_x=20x^3-4xy^3$, $f_y=-6x^2y^2+5y^4$.

  \item $F_K=12\cdot\frac14K^{-3/4}L^{3/4}=3(L/K)^{3/4}$;
        $F_L=12\cdot\frac34K^{1/4}L^{-1/4}=9(K/L)^{1/4}$. At $(16,81)$:
        $(L/K)^{3/4}=(81/16)^{3/4}=27/8$ so $F_K=81/8$;
        $(K/L)^{1/4}=(16/81)^{1/4}=2/3$ so $F_L=6$.

  \item $f_x=2xe^{3y}$, $f_y=3x^2e^{3y}$.

  \item $\nabla f=(12x^2y-y^4,\ 4x^3-4xy^3)^T$.

  \item $f_x=2x-4y$, $f_y=-4x+3y^2$. At $(2,1)$: $\nabla f=(0,-5)^T$.
        Normalise: $\|\mathbf{v}\|=\sqrt{25+144}=13$,
        $\mathbf{u}=(5/13,-12/13)^T$. Then
        \[
          D_{\mathbf{u}}f=0\cdot\frac{5}{13}+(-5)\cdot
            \left(-\frac{12}{13}\right)=\frac{60}{13} .
        \]

  \item $F_K=6\cdot\frac23K^{-1/3}L^{1/3}=4(L/K)^{1/3}$,
        $F_L=2(K/L)^{2/3}$. At $(27,8)$: $(L/K)^{1/3}=2/3$ so
        $\mathrm{MP}_K=8/3$; $(K/L)^{2/3}=(27/8)^{2/3}=9/4$ so
        $\mathrm{MP}_L=9/2$.

  \item Chain rule with $u=3x-2y$: $f_x=4(3x-2y)^3\cdot3=12(3x-2y)^3$,
        $f_y=-8(3x-2y)^3$.

  \item $f_x=y^2e^{xy^2}$, $f_y=2xye^{xy^2}$.

  \item $F_K=15K^{-0.4}L^{0.5}$, $F_L=12.5K^{0.6}L^{-0.5}$.

  \item $f_x=2x-z^2e^x$, $f_y=3z$, $f_z=3y-2ze^x$. So
        $\nabla f=(2x-z^2e^x,\ 3z,\ 3y-2ze^x)^T$.

  \item $f_x=e^x\ln(1+y)$, $f_y=e^x/(1+y)$. At $(0,0)$:
        $\nabla f=(0,1)^T$. Normalise: $\|\mathbf{v}\|=\sqrt5$,
        $\mathbf{u}=(1/\sqrt5,2/\sqrt5)^T$. So
        \[
          D_{\mathbf{u}}f=0\cdot\frac{1}{\sqrt5}+1\cdot\frac{2}{\sqrt5}
            =\frac{2\sqrt5}{5} .
        \]

  \item $U_x=2/x$, $U_y=5/y$. $\mathrm{MRS}_{xy}=U_x/U_y
        =(2/x)/(5/y)=2y/(5x)$. At $(3,5)$: $\mathrm{MRS}=10/15=2/3$.

  \item For $x+y\neq0$: $f(x,y)=x-y$, so $f_x=1$, $f_y=-1$.

  \item With $u=x^2y+1$: $f_x=2xy/(x^2y+1)$, $f_y=x^2/(x^2y+1)$.

  \item $U_{x_1}=\frac12x_1^{-1/2}x_2^{1/3}x_3^{1/6}$,
        $U_{x_2}=\frac13x_1^{1/2}x_2^{-2/3}x_3^{1/6}$,
        $U_{x_3}=\frac16x_1^{1/2}x_2^{1/3}x_3^{-5/6}$. Equivalently
        $U_{x_i}=(\alpha_i/x_i)U$ with $\alpha=(1/2,1/3,1/6)$.

  \item $f_x=y$, $f_y=x-2z$, $f_z=-2y+2z$. At $(1,2,3)$:
        $\nabla f=(2,-5,2)^T$. Normalise: $\|\mathbf{v}\|=\sqrt{1+4+4}=3$,
        $\mathbf{u}=(1/3,2/3,2/3)^T$. Then
        $D_{\mathbf{u}}f=(2-10+4)/3=-4/3$.

  \item $F_K=0.25K^{-0.75}L^{0.75}$, $F_L=0.75K^{0.25}L^{-0.25}$. So
        $\mathrm{MRTS}_{LK}=F_L/F_K=3K/L$. At $(16,4)$:
        $\mathrm{MRTS}=48/4=12$.

  \item Expand: $f=x^3-3x^2y^2+xy-3y^3$. Hence $f_x=3x^2-6xy^2+y$,
        $f_y=-6x^2y+x-9y^2$.

  \item With $u=-y/x$: $\partial u/\partial x=y/x^2$,
        $\partial u/\partial y=-1/x$. Product rule:
        \[
          f_x=e^{-y/x}+x\cdot e^{-y/x}\cdot\frac{y}{x^2}
            =e^{-y/x}\left(1+\frac{y}{x}\right), \qquad
          f_y=x\cdot e^{-y/x}\cdot\left(-\frac1x\right)=-e^{-y/x} .
        \]

  \item Let $g=3K^{1/2}+2L^{1/2}$, so $F=g^2$. Then
        \[
          F_K=2g\cdot\frac32K^{-1/2}=\frac{3g}{K^{1/2}}
            =\frac{3(3K^{1/2}+2L^{1/2})}{K^{1/2}}, \qquad
          F_L=2g\cdot L^{-1/2}=\frac{2(3K^{1/2}+2L^{1/2})}{L^{1/2}} .
        \]

  \item With $u=y/x$, $f=\arctan u$, $f'=1/(1+u^2)$:
        \[
          f_x=\frac{1}{1+y^2/x^2}\cdot\left(-\frac{y}{x^2}\right)
            =\frac{-y}{x^2+y^2}, \qquad
          f_y=\frac{x}{x^2+y^2} .
        \]
        So $\nabla f=(-y/(x^2+y^2),\ x/(x^2+y^2))^T$.

  \item $f_K=\frac12K^{-1/2}L^{1/2}$, $f_L=\frac12K^{1/2}L^{-1/2}$. At
        $(4,9)$: $\nabla f=(\frac12\cdot\frac32,\ \frac12\cdot\frac23)^T
        =(3/4,1/3)^T$. Then $\|\nabla f\|^2=9/16+1/9=97/144$, so the
        maximum rate is $\sqrt{97}/12\approx0.821$ in direction
        $\mathbf{u}^*=(9,4)^T/\sqrt{97}$.

  \item $F=(K^{1/2}+L^{1/2})^2$. Replacing $(K,L)\to(\lambda K,\lambda L)$
        gives $\lambda F$, so $r=1$. With $g=K^{1/2}+L^{1/2}$:
        $F_K=2g\cdot\frac12K^{-1/2}=g/K^{1/2}$, $F_L=g/L^{1/2}$. Hence
        \[
          KF_K+LF_L=K^{1/2}g+L^{1/2}g=g(K^{1/2}+L^{1/2})=g^2=F. \checkmark
        \]

  \item Quotient rule, $D=x^2+y^2$:
        \[
          f_x=\frac{2D-(2x+3y)(2x)}{D^2}=\frac{-2x^2-6xy+2y^2}{(x^2+y^2)^2},
          \qquad f_y=\frac{3x^2-4xy-3y^2}{(x^2+y^2)^2} .
        \]

  \item Treat denominator as constant for $f_x$: $f_x=e^x/(1+e^y)$. For
        $f_y$:
        \[
          f_y=e^x\cdot\frac{-e^y}{(1+e^y)^2}=-\frac{e^{x+y}}{(1+e^y)^2} .
        \]

  \item Let $S=\alpha K^\rho+(1-\alpha)L^\rho$, $F=AS^{\nu/\rho}$.
        Treating $K$ constant, $\partial S/\partial L=\rho(1-\alpha)L^{\rho-1}$:
        \[
          F_L=A\cdot\frac{\nu}{\rho}S^{\nu/\rho-1}\cdot\rho(1-\alpha)L^{\rho-1}
            =\nu(1-\alpha)AL^{\rho-1}S^{\nu/\rho-1} .
        \]

  \item $f=\ln(x^2+y^2)$, $f_x=2x/(x^2+y^2)$, $f_y=2y/(x^2+y^2)$. At
        $(3,4)$: $x^2+y^2=25$, $\nabla f=(6/25,8/25)^T$. Normalise:
        $\|\mathbf{v}\|=\sqrt2$, $\mathbf{u}=(-1/\sqrt2,1/\sqrt2)^T$. Then
        \[
          D_{\mathbf{u}}f=\frac{6}{25}\cdot\left(-\frac{1}{\sqrt2}\right)
            +\frac{8}{25}\cdot\frac{1}{\sqrt2}=\frac{2}{25\sqrt2}
            =\frac{\sqrt2}{25} .
        \]

  \item $f_x=2xyz^3-4y^2z$, $f_y=x^2z^3-8xyz$,
        $f_z=3x^2yz^2-4xy^2+5$.

  \item Chain rule, $u=4x^2+9y^2$:
        \[
          f_x=\frac{8x}{2\sqrt u}=\frac{4x}{\sqrt{4x^2+9y^2}}, \qquad
          f_y=\frac{9y}{\sqrt{4x^2+9y^2}} .
        \]

  \item With $g=K^{1/2}+L^{1/2}$, $F=g^4$,
        $F_K=4g^3\cdot\frac12K^{-1/2}=2g^3/K^{1/2}$. At $(4,9)$: $g=5$,
        $K^{1/2}=2$, so $F_K=2\cdot125/2=125$.

  \item Let $r^2=x^2+y^2+z^2$. Then $f=\ln(r^2)$ and $f_{x_i}=2x_i/r^2$.
        So $\nabla f=(2/(x^2+y^2+z^2))(x,y,z)^T$.

  \item Use $\ln\sqrt{u}=\frac12\ln u$: $f=\frac12\ln(x^2+y^2)$. So
        $f_x=x/(x^2+y^2)$, $f_y=y/(x^2+y^2)$.

  \item $\nabla C=(4,9)^T$. With $Q=x^{1/2}y^{1/2}$: $Q_x=\frac12(y/x)^{1/2}$
        and $Q_y=\frac12(x/y)^{1/2}$, so
        \[
          \nabla Q=\tfrac12\left((y/x)^{1/2},\,(x/y)^{1/2}\right)^{T} .
        \]
        The condition $\nabla C=\lambda\nabla Q$ requires
        $4/9=Q_x/Q_y=y/x$, i.e.\ $y=4x/9$. Along this locus
        $\nabla Q\propto(2/3,3/2)^T\propto(4,9)^T=\nabla C$. \checkmark

  \item Chain rule: $f_x=5(1+xy)^4\cdot y=5y(1+xy)^4$,
        $f_y=5x(1+xy)^4$.

  \item $f_x=3x^2-2y$, $f_y=-2x+2y$. At $(1,2)$:
        $\nabla f=(3-4,-2+4)^T=(-1,2)^T$. Normalise: $\|\mathbf{v}\|=5$,
        $\mathbf{u}=(3/5,4/5)^T$. Then
        \[
          D_{\mathbf{u}}f=-1\cdot\frac35+2\cdot\frac45
            =\frac{-3+8}{5}=1 .
        \]

  \item $F_K=10\cdot0.3K^{-0.7}L^{0.7}=3(L/K)^{0.7}$;
        $F_L=7(K/L)^{0.3}$. At $K=L=100$: $F_K=3$, $F_L=7$. And
        $F=10\cdot100^{0.3}\cdot100^{0.7}=10\cdot100=1000$. Then
        $KF_K+LF_L=100\cdot3+100\cdot7=1000=F$. \checkmark (homogeneity
        degree $r=0.3+0.7=1$).

  \item $f_x=y\cos(xy)-\sin(x-y)$, $f_y=x\cos(xy)+\sin(x-y)$.

  \item $f_x=2xy-y^2+3$, $f_y=x^2-2xy-2$. At $(1,-1)$: $f_x=-2-1+3=0$,
        $f_y=1+2-2=1$. So $\nabla f(1,-1)=(0,1)^T$.

  \item $U_x=\alpha x^{\alpha-1}y^{1-\alpha}$,
        $U_y=(1-\alpha)x^\alpha y^{-\alpha}$. So
        \[
          \mathrm{MRS}_{xy}=\frac{U_x}{U_y}
            =\frac{\alpha x^{\alpha-1}y^{1-\alpha}}{(1-\alpha)x^\alpha y^{-\alpha}}
            =\frac{\alpha}{1-\alpha}\cdot\frac{y}{x} . \checkmark
        \]

  \item $f_x=ye^{xy}$, $f_y=xe^{xy}$. At $(1,0)$: $\nabla f=(0,1)^T$.
        Normalise: $\|\mathbf{v}\|=5$, $\mathbf{u}=(-3/5,4/5)^T$. Then
        $D_{\mathbf{u}}f=0\cdot(-3/5)+1\cdot(4/5)=4/5$.

  \item Write $f=e^{y\ln x}$. Then $f_x=e^{y\ln x}\cdot y/x=yx^{y-1}$ and
        $f_y=e^{y\ln x}\cdot\ln x=x^y\ln x$.

  \item On $\{2K<3L\}$ the binding factor is capital: $F=2K$, so
        $F_K=2$, $F_L=0$. On $\{2K>3L\}$ labour binds: $F=3L$, so
        $F_K=0$, $F_L=3$. (Non-differentiable on $\{2K=3L\}$.)

  \item $f_x=yz+2x$, $f_y=xz$, $f_z=xy-2z$. At $(2,-1,3)$:
        $f_x=-3+4=1$, $f_y=6$, $f_z=-2-6=-8$. So
        $\nabla f=(1,6,-8)^T$.

  \item $\nabla f=(2x,8y)^T$. At $(3,1)$: $\nabla f=(6,8)^T$. Steepest
        descent is $-\nabla f/\|\nabla f\|$. $\|\nabla f\|=10$, so unit
        descent direction is $(-3/5,-4/5)^T$, with rate of decrease
        $-\|\nabla f\|=-10$ (i.e.\ $f$ decreases at rate 10).

  \item $f_x=e^y-ye^x$, $f_y=xe^y-e^x$.

  \item Quotient rule with $D=x^2+y^2$:
        \[
          f_x=\frac{D-x\cdot2x}{D^2}=\frac{y^2-x^2}{(x^2+y^2)^2}, \qquad
          f_y=\frac{-x\cdot2y}{D^2}=\frac{-2xy}{(x^2+y^2)^2} .
        \]

  \item $U_x=1$, $U_y=1/\sqrt{y}$. So $\mathrm{MRS}_{xy}=U_x/U_y
        =\sqrt{y}$ (depends only on $y$, quasi-linear utility).

  \item $f_x=2x$, $f_y=2y$, $f_z=-2z$. At $(2,-1,2)$:
        $\nabla f=(4,-2,-4)^T$. Normalise: $\|\mathbf{v}\|=\sqrt{4+4+1}=3$,
        $\mathbf{u}=(2/3,2/3,1/3)^T$. Then
        \[
          D_{\mathbf{u}}f=4\cdot\frac23+(-2)\cdot\frac23+(-4)\cdot\frac13
            =\frac{8-4-4}{3}=0 .
        \]
        (We are moving tangent to the level surface.)

  \item $f_x=\sec^2(x+y)+\sec^2(x-y)$, $f_y=\sec^2(x+y)-\sec^2(x-y)$.

  \item Let $S=\alpha K+\beta L$, $F=S^r$. Then $F_K=rS^{r-1}\alpha$,
        $F_L=rS^{r-1}\beta$. Hence
        \[
          KF_K+LF_L=rS^{r-1}(\alpha K+\beta L)=rS^{r-1}\cdot S=rS^r=rF.
            \checkmark
        \]

  \item $\nabla f=(-e^{-x}\sin y+e^x\cos y,\ e^{-x}\cos y-e^x\sin y)^T$.

  \item $f_x=x/\sqrt{x^2+y^2}$, $f_y=y/\sqrt{x^2+y^2}$. At $(3,4)$:
        $\sqrt{x^2+y^2}=5$, $\nabla f=(3/5,4/5)^T$. The direction
        $\mathbf{v}=(1,0)$ is already unit. So
        $D_{\mathbf{u}}f=3/5\cdot1+4/5\cdot0=3/5$.

  \item On $\{2x<3y\}$: $U=2x$, so $U_x=2$, $U_y=0$. On $\{2x>3y\}$:
        $U=3y$, so $U_x=0$, $U_y=3$. (Non-differentiable on the ray
        $2x=3y$.)

  \item Let $L=\ln(x+y)$, $D=x-y$. By the quotient rule:
        \[
          f_x=\frac{(1/(x+y))(x-y)-\ln(x+y)\cdot1}{(x-y)^2}
            =\frac{1/(x+y)\cdot(x-y)-\ln(x+y)}{(x-y)^2} .
        \]
        Simplifying
        $f_x=\dfrac{(x-y)/(x+y)-\ln(x+y)}{(x-y)^2}$. Similarly,
        \[
          f_y=\frac{(1/(x+y))(x-y)-\ln(x+y)\cdot(-1)}{(x-y)^2}
            =\frac{(x-y)/(x+y)+\ln(x+y)}{(x-y)^2} .
        \]

  \item $F_K=A\alpha K^{\alpha-1}L^\beta$, $F_L=A\beta K^\alpha L^{\beta-1}$.
        So
        \[
          \mathrm{MRTS}_{LK}=\frac{F_L}{F_K}=\frac{\beta}{\alpha}\cdot
            \frac{K}{L} .
        \]
        This depends only on the ratio $K/L$ and the exponent ratio
        $\beta/\alpha$: it is scale-invariant.

  \item $\nabla f=(2x,-2y)^T$. At $(2,-1)$: $\nabla f=(4,2)^T$. Direction
        of steepest ascent: $\mathbf{u}^*=(4,2)/\sqrt{20}=(2,1)/\sqrt5$.
        Maximum rate $=\|\nabla f\|=\sqrt{20}=2\sqrt5$.

  \item $f_x=\sinh(x+y)+\cosh(x-y)$, $f_y=\sinh(x+y)-\cosh(x-y)$.

  \item $\nabla f=(e^{x+y+z},e^{x+y+z},e^{x+y+z})^T$. At $(0,0,0)$:
        $\nabla f=(1,1,1)^T$. Normalise $\mathbf{v}=(1,1,1)$:
        $\|\mathbf{v}\|=\sqrt3$, $\mathbf{u}=(1,1,1)/\sqrt3$. Then
        $D_{\mathbf{u}}f=(1+1+1)/\sqrt3=\sqrt3$.

  \item $f_x=(\ln y)/x$, $f_y=(\ln x)/y$.

  \item $F_K=1+L^{1/2}/K^{1/2}$, $F_L=1+K^{1/2}/L^{1/2}$. From
        $F=(K^{1/2}+L^{1/2})^2$ with $g=K^{1/2}+L^{1/2}$:
        $F_K=g/K^{1/2}=(K^{1/2}+L^{1/2})/K^{1/2}=1+L^{1/2}/K^{1/2}$.
        \checkmark Same for $F_L$.

  \item $f_x=2xy$, $f_y=x^2+3y^2$. At $(-1,2)$: $\nabla f=(-4,1+12)^T
        =(-4,13)^T$. Normalise: $\|\mathbf{v}\|=5$,
        $\mathbf{u}=(4/5,-3/5)^T$. Then
        \[
          D_{\mathbf{u}}f=-4\cdot\frac45+13\cdot\left(-\frac35\right)
            =-\frac{16+39}{5}=-11 .
        \]

  \item $U_x=\alpha/x$, $U_y=\beta/y$.
        $\mathrm{MRS}_{xy}=\dfrac{\alpha/x}{\beta/y}=\dfrac{\alpha y}{\beta x}$.
        At $(\alpha,\beta)$: $\mathrm{MRS}=(\alpha\cdot\beta)/(\beta\cdot\alpha)=1$.

  \item $f=(x^2+y^2)^{-1/2}$. Chain rule:
        \[
          f_x=-\frac12(x^2+y^2)^{-3/2}\cdot2x=\frac{-x}{(x^2+y^2)^{3/2}},
          \qquad f_y=\frac{-y}{(x^2+y^2)^{3/2}} .
        \]
        So $\nabla f=-(x,y)^T/(x^2+y^2)^{3/2}$.

  \item Chain rule with $u=xy$:
        \[
          f_x=\frac{y}{1+x^2y^2}, \qquad f_y=\frac{x}{1+x^2y^2} .
        \]

  \item Let $S=K^\rho+L^\rho$, $F=S^{1/\rho}$. Then
        $F_K=\frac1\rho S^{1/\rho-1}\cdot\rho K^{\rho-1}
        =K^{\rho-1}S^{1/\rho-1}$. Since $S^{1/\rho}=F$, we have $S=F^\rho$
        and so $S^{1/\rho-1}=F^{1-\rho}\cdot F^{-1}\cdot F=F^{1-\rho}$.
        Thus $F_K=K^{\rho-1}F^{1-\rho}=(F/K)^{1-\rho}=(K/F)^{\rho-1}$.
        \checkmark

  \item $f_x=yz$, $f_y=xz$, $f_z=xy$. At $(2,3,-1)$:
        $\nabla f=(3\cdot(-1),2\cdot(-1),2\cdot3)^T=(-3,-2,6)^T$.
        Normalise $\mathbf{v}=(2,-1,2)$: $\|\mathbf{v}\|=3$,
        $\mathbf{u}=(2/3,-1/3,2/3)^T$. Then
        \[
          D_{\mathbf{u}}f=-3\cdot\frac23+(-2)\cdot\left(-\frac13\right)
            +6\cdot\frac23=-2+\frac23+4=2+\frac23=\frac83 .
        \]
\end{solutions}

\section{Additional Problems}

\begin{problems}
  \item Compute $f_x$ and $f_y$ for $f(x,y)=3x^5-2x^3y^2+4y^4$.
  \item For $F(K,L)=20\,K^{1/3}L^{2/3}$, find $F_K$ and $F_L$ at
        $(K,L)=(8,27)$.
  \item Compute $f_x$ and $f_y$ for $f(x,y)=e^{xy}$.
  \item Compute $\nabla f$ for $f(x,y)=2x^2y-3xy^3+5$.
  \item Compute the directional derivative of $f(x,y)=x^2+2xy-y^2$ at
        $(1,2)$ in the direction $\mathbf{v}=(3,4)$.
  \item For $F(K,L)=8K^{3/4}L^{1/4}$, compute $\mathrm{MP}_K$ and
        $\mathrm{MP}_L$ at $(K,L)=(16,16)$.
  \item Compute $f_x$ and $f_y$ for $f(x,y)=\ln(2x+3y)$.
  \item For $U(x,y)=3\ln(x)+4\ln(y)$, derive $\mathrm{MRS}_{xy}$ and
        evaluate at $(2,3)$.
  \item Compute $f_x$ and $f_y$ for $f(x,y)=(x+y^2)^4$.
  \item Compute $\nabla f$ for $f(x,y,z)=xy+yz+zx$.
  \item For a CES production $F(K,L)=5(K^{1/3}+L^{1/3})^3$, find $F_K$.
  \item Compute the directional derivative of $f(x,y)=x^2y$ at $(2,1)$
        in the direction $\mathbf{v}=(1,-1)$.
  \item For $U(x,y)=x^\alpha y^\beta$ (general Cobb--Douglas,
        $\alpha,\beta>0$), derive $\mathrm{MRS}_{xy}$ and show it depends
        only on the ratio $y/x$.
  \item Compute $f_x$ and $f_y$ for $f(x,y)=\dfrac{x-y}{x+y}$.
  \item For $f(K,L)=K^{1/2}L^{1/2}$ at $(K,L)=(16,9)$, find the unit
        direction of steepest ascent and the maximum rate of increase.
  \item Classify the open positive orthant
        $Q=\{(x,y)\in\R^2 : x>0,\ y>0\}$: is it open, closed, or
        neither? Is it bounded? Is it compact?
  \item Classify the budget set
        $B=\{(x,y)\in\R^2 : 2x+3y\le12,\ x\ge0,\ y\ge0\}$: open,
        closed, or neither? Bounded? Compact?
  \item Classify $M=\{(x,y)\in\R^2 : xy\ge1,\ x>0\}$: open, closed, or
        neither? Bounded? Compact? \emph{Hint:} check whether a
        sequence in $M$ can converge to a point with $x=0$.
  \item Classify the open disc
        $D=\{(x,y)\in\R^2 : x^2+y^2<4\}$: open, closed, or neither?
        Bounded? Compact?
  \item Classify (a) the half-open strip
        $S=\{(x,y)\in\R^2 : 0\le x<1\}$ and (b) the finite set
        $E=\{(0,0),\ (1,2),\ (3,1)\}$: open, closed, or neither?
        Bounded? Compact?
  \item Does the Weierstrass theorem guarantee a maximiser for the
        continuous function $u(x,y)=x^{1/2}y^{1/2}$ (a) on the budget
        set $B=\{(x,y): 2x+3y\le12,\ x\ge0,\ y\ge0\}$, (b) on the open
        positive orthant $Q=\{(x,y): x>0,\ y>0\}$? Justify each
        answer.
\end{problems}

In problems 22--29, for each function: state the domain; compute the
partials and the set on which they exist; decide, with a justification,
where $f$ is differentiable; and describe the graph near the points where
differentiability fails.

\begin{problems}[resume]
  \item $f(x,y)=\sqrt{(x-1)^{2}+(y+2)^{2}}$.
  \item $f(x,y)=\sqrt{4x^{2}+(y-3)^{2}}$.
  \item $f(x,y)=|x-2|+|y|$.
  \item $f(x,y)=\min\{2x,\,5y\}$ on $x,y>0$.
  \item $f(x,y)=\dfrac{3xy}{2x^{2}+y^{2}}$ for $(x,y)\neq(0,0)$, with
        $f(0,0)=0$. Also decide whether $f$ is continuous at the origin.
  \item $f(x,y)=\sqrt{x^{4}+y^{4}}$. The formula for $f_{x}$ breaks down
        at the origin; decide what really happens there.
  \item $f(x,y)=\sqrt{9-x^{2}-y^{2}}$.
  \item $f(x,y)=|x-y|$.
\end{problems}

In problems 30--37, find the tangent plane of $f$ at $P$ (the tangent
hyperplane, for three variables), then evaluate both $f$ and the plane at
the nearby point $Q$ and compare. Give both values to six decimals.

\begin{problems}[resume]
  \item $f(x,y)=\ln(2x+y^{2})$, \;$P=(0,1)$, \;$Q=(0.05,\,0.98)$.
  \item $f(x,y)=\mathrm{e}^{x-2y}$, \;$P=(2,1)$, \;$Q=(2.1,\,1.02)$.
  \item $f(x,y)=xy^{2}\,\mathrm{e}^{x-1}$, \;$P=(1,2)$,
        \;$Q=(0.98,\,2.03)$.
  \item $f(x,y)=\dfrac{x+y}{x-y}$, \;$P=(3,1)$, \;$Q=(3.1,\,0.9)$.
  \item $F(K,L)=K^{1/4}L^{3/4}$, \;$P=(16,1)$, \;$Q=(16.5,\,1.1)$; read
        the plane as an output estimate.
  \item $f(x,y,z)=\ln(x+2y+3z)$, \;$P=(2,-1,\tfrac13)$,
        \;$Q=(2.02,\,-0.99,\,0.35)$.
  \item $f(x,y,z)=x\,\mathrm{e}^{yz}$, \;$P=(2,0,5)$,
        \;$Q=(2.05,\,0.01,\,4.9)$.
  \item $f(x,y,z)=\dfrac{xy}{z}$, \;$P=(2,3,6)$, \;$Q=(2.1,\,2.9,\,6.2)$.
\end{problems}

In problems 38--46, compute the partial derivatives, and give the gradient
(a column) and the derivative (a row) at the stated point.

\begin{problems}[resume]
  \item $f(x_{1},\ldots,x_{n})=\sum_{i=1}^{n}\ln(1+x_{i}^{2})$, at
        $(1,1,\ldots,1)$.
  \item $f(x_{1},\ldots,x_{n})=\Bigl(\sum_{i=1}^{n}x_{i}\Bigr)^{3}$; then
        for $n=4$ at $(1,2,3,4)$.
  \item $f(x_{1},\ldots,x_{n})=\sum_{i=1}^{n}x_{i}^{\,i+1}$, at
        $(1,1,\ldots,1)$; then for $n=3$ at $(2,1,-1)$.
  \item $f(x_{1},\ldots,x_{n})=\sum_{j=1}^{n}\dfrac{j^{2}}{x_{j}^{3}}$ with
        all $x_{j}>0$, at $(1,1,\ldots,1)$.
  \item $f(x_{1},\ldots,x_{n})=\sum_{i=1}^{n}x_{i}^{2}
        -\dfrac1n\Bigl(\sum_{i=1}^{n}x_{i}\Bigr)^{2}$; show that the
        gradient vanishes exactly when all $x_{i}$ are equal, then evaluate
        for $n=3$ at $(1,2,6)$.
  \item $f(x,y,z)=y^{2}\ln(x^{2}+3yz)$, at $P=(1,1,1)$.
  \item $f(x,y,z)=x^{3}\,\mathrm{e}^{y^{2}-2xz}$, at $P=(1,0,1)$.
  \item $f(x,y)=2^{x}3^{y}+5^{xy}$, at $(1,1)$.
  \item $f(x_{1},\ldots,x_{n})=\sum_{i=1}^{n}3^{\,i x_{i}}$, at
        $(0,0,\ldots,0)$.
\end{problems}

In problems 47--52, decide differentiability directly from the definition,
by forming the remainder and estimating it, as in the third worked example.

\begin{problems}[resume]
  \item $f(x,y)=3xy$, at every point $(a,b)$.
  \item $f(x,y)=x^{2}+y^{2}$, first at the origin, then at every point.
  \item $f(x,y)=xy^{2}$, at every point $(a,b)$.
  \item $f(x,y)=y\,\mathrm{e}^{3x}$, at every point $(a,b)$.
  \item At the origin: (a) $f(x,y)=(x^{2}+4y^{2})^{1/3}$;
        (b) $g(x,y)=(x^{2}+y^{2})^{3/4}$.
  \item $f(x,y)=|xy|$: at the origin, and at the point $(2,0)$.
\end{problems}

In problems 53--54, $f(0,0)=0$. Show with the difference quotient that both
partials exist at the origin and are $0$; compute the limit along every
straight line through the origin and along the stated parabola; and decide
continuity and differentiability at the origin.

\begin{problems}[resume]
  \item $f(x,y)=\dfrac{x^{2}y}{x^{4}+y^{2}}$ for $(x,y)\neq(0,0)$; parabola
        $y=x^{2}$.
  \item $f(x,y)=\dfrac{2xy^{2}}{x^{2}+y^{4}}$ for $(x,y)\neq(0,0)$;
        parabola $x=y^{2}$.
\end{problems}

Problems 55--56 refer to the three surfaces below, each drawn over a square
around its special point.

\begin{gallery}
\begin{tikzpicture}
  \begin{axis}[surfaxis, title={(a)}, z post scale=1.1]
    \addplot3[surf, domain=-1:3, y domain=-2:2]{(x-1)^2+2*y^2};
  \end{axis}
\end{tikzpicture}\hfill
\begin{tikzpicture}
  \begin{axis}[surfaxis, title={(b)}, z post scale=1.1]
    \addplot3[surf, domain=-1:3, y domain=-2:2]{sqrt((x-1)^2+4*y^2)};
  \end{axis}
\end{tikzpicture}\hfill
\begin{tikzpicture}
  \begin{axis}[surfaxis, title={(c)}, z post scale=1.1]
    \addplot3[surf, domain=-2:2, y domain=-2:2]{2-abs(x)-abs(y)};
  \end{axis}
\end{tikzpicture}
\end{gallery}

\begin{problems}[resume]
  \item For each surface (a)--(c), say from the picture alone where the
        function is differentiable and where it is not, and why.
  \item Surface (b) has its tip above the point $(1,0)$, and its contours
        are ellipses twice as long in $x$ as in $y$. Which of the formulas
        $(x-1)^{2}+4y^{2}$, $\sqrt{(x-1)^{2}+4y^{2}}$ and
        $\sqrt{4(x-1)^{2}+y^{2}}$ can it be the graph of? Confirm the
        verdict of the previous problem for that formula with the
        difference quotient at $(1,0)$.
\end{problems}

\section{Additional Solutions}

\begin{solutions}
  \item $f_x=15x^4-6x^2y^2$, $f_y=-4x^3y+16y^3$.

  \item $F_K=20\cdot\frac13K^{-2/3}L^{2/3}=\frac{20}{3}(L/K)^{2/3}$;
        $F_L=20\cdot\frac23K^{1/3}L^{-1/3}=\frac{40}{3}(K/L)^{1/3}$. At
        $(8,27)$: $(L/K)^{2/3}=(27/8)^{2/3}=9/4$, so $F_K=15$;
        $(K/L)^{1/3}=(8/27)^{1/3}=2/3$, so
        $F_L=\frac{40}{3}\cdot\frac23=\frac{80}{9}$.

  \item $f_x=ye^{xy}$, $f_y=xe^{xy}$.

  \item $\nabla f=(4xy-3y^3,\ 2x^2-9xy^2)^T$.

  \item $f_x=2x+2y$, $f_y=2x-2y$. At $(1,2)$: $\nabla f=(6,-2)^T$.
        Normalise: $\|\mathbf{v}\|=5$, $\mathbf{u}=(3/5,4/5)^T$. Then
        \[
          D_{\mathbf{u}}f=6\cdot\frac35+(-2)\cdot\frac45
            =\frac{18-8}{5}=2 .
        \]

  \item $F_K=8\cdot\frac34K^{-1/4}L^{1/4}=6(L/K)^{1/4}$,
        $F_L=2(K/L)^{3/4}$. At $(16,16)$: $\mathrm{MP}_K=6$,
        $\mathrm{MP}_L=2$.

  \item $f_x=\dfrac{2}{2x+3y}$, $f_y=\dfrac{3}{2x+3y}$.

  \item $U_x=3/x$, $U_y=4/y$.
        $\mathrm{MRS}_{xy}=U_x/U_y=\dfrac{3y}{4x}$. At $(2,3)$:
        $\mathrm{MRS}=\dfrac98$.

  \item Chain: $f_x=4(x+y^2)^3$,
        $f_y=4(x+y^2)^3\cdot2y=8y(x+y^2)^3$.

  \item $f_x=y+z$, $f_y=x+z$, $f_z=x+y$. So
        $\nabla f=(y+z,\ x+z,\ x+y)^T$.

  \item Let $g=K^{1/3}+L^{1/3}$, so $F=5g^3$. Then
        $F_K=5\cdot3g^2\cdot\frac13K^{-2/3}=5g^2K^{-2/3}
        =\dfrac{5(K^{1/3}+L^{1/3})^2}{K^{2/3}}$.

  \item $f_x=2xy$, $f_y=x^2$. At $(2,1)$: $\nabla f=(4,4)^T$. Normalise:
        $\|\mathbf{v}\|=\sqrt2$, $\mathbf{u}=(1/\sqrt2,-1/\sqrt2)^T$.
        Then
        \[
          D_{\mathbf{u}}f=4\cdot\frac{1}{\sqrt2}
            +4\cdot\left(-\frac{1}{\sqrt2}\right)=0 .
        \]
        (Tangent to the level curve.)

  \item $U_x=\alpha x^{\alpha-1}y^\beta$,
        $U_y=\beta x^\alpha y^{\beta-1}$. So
        \[
          \mathrm{MRS}_{xy}=\frac{U_x}{U_y}
            =\frac{\alpha}{\beta}\cdot\frac{y}{x} .
        \]
        Depends only on $y/x$ (and the exponent ratio), not on the
        scale. \checkmark

  \item Quotient rule with $D=x+y$:
        \[
          f_x=\frac{D-(x-y)}{D^2}=\frac{2y}{(x+y)^2}, \qquad
          f_y=\frac{-D-(x-y)}{D^2}=\frac{-2x}{(x+y)^2} .
        \]

  \item $f_K=\frac12K^{-1/2}L^{1/2}$, $f_L=\frac12K^{1/2}L^{-1/2}$. At
        $(16,9)$: $f_K=\frac12\cdot\frac34=3/8$;
        $f_L=\frac12\cdot\frac43=2/3$. So $\nabla f=(3/8,2/3)^T$.
        Magnitude:
        $\|\nabla f\|=\sqrt{9/64+4/9}=\sqrt{81/576+256/576}
        =\sqrt{337/576}=\sqrt{337}/24\approx0.765$. Unit direction of
        steepest ascent:
        $\mathbf{u}^*=(3/8,2/3)^T/(\sqrt{337}/24)=(9,16)^T/\sqrt{337}$.

  \item \textbf{Open.} For any $(x,y)\in Q$ the ball of radius
        $r=\min\{x,y\}>0$ around $(x,y)$ stays inside $Q$, so every
        point is interior. \textbf{Not closed:} the boundary points
        $(x,0)$ and $(0,y)$ with $x,y\ge0$ (the coordinate half-axes,
        including the origin) do not belong to $Q$. \textbf{Not
        bounded:} $Q$ contains $(t,t)$ for every $t>0$, and
        $\|(t,t)\|=t\sqrt2\to\infty$. Hence $Q$ is open, unbounded and
        \textbf{not compact}.

  \item \textbf{Closed.} $B$ is the intersection of the three sets
        $\{2x+3y\le12\}$, $\{x\ge0\}$ and $\{y\ge0\}$, each defined by
        a non-strict inequality on a continuous function; each is
        closed, and an intersection of closed sets is closed (every
        boundary point of $B$ satisfies all three inequalities, so it
        lies in $B$). \textbf{Not open:} $(0,0)\in B$, but every ball
        around the origin contains points with negative coordinates,
        which are outside $B$; so $(0,0)$ is not interior.
        \textbf{Bounded:} from the constraints, $0\le x\le6$ and
        $0\le y\le4$, so $B\subseteq B_8(\mathbf{0})$. Closed and
        bounded in $\R^2$: $B$ is \textbf{compact}.

  \item First note the constraint $x>0$ together with $xy\ge1$ forces
        $y\ge1/x>0$. \textbf{Closed:} suppose
        $(x_n,y_n)\in M$ with $(x_n,y_n)\to(a,b)$. If $a=0$, then
        $y_n\to b$ finite would give $x_ny_n\to0<1$, contradicting
        $x_ny_n\ge1$; so $a>0$, and passing to the limit in
        $x_ny_n\ge1$ gives $ab\ge1$, hence $(a,b)\in M$. $M$ contains
        all its limit (boundary) points, so it is closed: even
        though it is written with a strict inequality, that inequality
        is never active. \textbf{Not open:} $(1,1)\in M$ lies on the
        hyperbola $xy=1$, and every ball around it contains points
        with $xy<1$. \textbf{Not bounded:} $(t,t)\in M$ for all
        $t\ge1$. Hence $M$ is closed, unbounded and \textbf{not
        compact}.

  \item \textbf{Open.} For $(x,y)\in D$ put
        $r=2-\sqrt{x^2+y^2}>0$; the ball $B_r(x,y)$ lies in $D$ by the
        triangle inequality. \textbf{Not closed:} the boundary circle
        $\{x^2+y^2=4\}$ is not contained in $D$ (e.g.\ $(2,0)\notin
        D$). \textbf{Bounded:} $D\subseteq B_3(\mathbf{0})$. Bounded
        but not closed: $D$ is \textbf{not compact}.

  \item (a) \textbf{Neither open nor closed.} The points $(0,y)$
        belong to $S$ but are boundary points (every ball around them
        contains points with $x<0$), so $S$ is not open; the boundary
        points $(1,y)$ do not belong to $S$, so $S$ is not closed.
        \textbf{Unbounded:} $(0,t)\in S$ for all $t$. Hence
        \textbf{not compact}. (b) \textbf{Closed:} around any point
        not in $E$ there is a ball avoiding all three points (a
        finite set always contains all its boundary points:
        indeed $E$ is its own boundary). \textbf{Not open:} no ball
        around $(0,0)$ lies inside $E$. \textbf{Bounded:}
        $E\subseteq B_4(\mathbf{0})$. Closed and bounded: $E$ is
        \textbf{compact}.

  \item (a) \textbf{Yes.} $B$ is compact (closed and bounded, see the
        budget-set problem above) and $u$ is continuous on $B$, so by
        the Weierstrass theorem $u$ attains both a maximum and a
        minimum on $B$; a maximiser exists. (b) \textbf{No.} $Q$ is
        not compact (it is neither closed nor bounded), so Weierstrass
        does not apply: and indeed no maximiser exists:
        $u(t,t)=t\to\infty$ as $t\to\infty$, so $u$ is unbounded above
        on $Q$.

  \item \textbf{Domain:} the radicand is a sum of squares, never
        negative, so $D=\R^{2}$. \textbf{Partials:} where the radicand is
        positive, i.e.\ for $(x,y)\neq(1,-2)$, the chain rule gives
        \[
          f_{x}=\frac{x-1}{\sqrt{(x-1)^{2}+(y+2)^{2}}},\qquad
          f_{y}=\frac{y+2}{\sqrt{(x-1)^{2}+(y+2)^{2}}} .
        \]
        At $(1,-2)$ use the quotient:
        $\bigl(f(1+h,-2)-f(1,-2)\bigr)/h=\sqrt{h^{2}}/h=|h|/h$, which is $1$
        for $h>0$ and $-1$ for $h<0$, so it has no limit. $f_{x}(1,-2)$
        does not exist, and in the same way neither does $f_{y}(1,-2)$. The
        partials exist exactly on $\R^{2}\setminus\{(1,-2)\}$.
        \textbf{Differentiability:} on that open set the partials are
        continuous ($f$ is $C^{1}$), so $f$ is differentiable there. At
        $(1,-2)$ it is not, since differentiability would force the
        partials to exist. Note that $f$ \emph{is} continuous at $(1,-2)$:
        continuity is not enough. \textbf{Graph:} a circular cone with its
        vertex at $(1,-2,0)$, opening upwards with slope $1$ in every
        direction; the contours are circles of radius $c$ about $(1,-2)$.
        At the sharp tip there is no tangent plane.

  \item \textbf{Domain:} $\R^{2}$, the radicand being a sum of squares.
        \textbf{Partials:} for $(x,y)\neq(0,3)$,
        \[
          f_{x}=\frac{4x}{\sqrt{4x^{2}+(y-3)^{2}}},\qquad
          f_{y}=\frac{y-3}{\sqrt{4x^{2}+(y-3)^{2}}} .
        \]
        At $(0,3)$: $f(h,3)-f(0,3)=\sqrt{4h^{2}}=2|h|$, so the quotient
        $2|h|/h=\pm2$ has no limit; and $f(0,3+h)-f(0,3)=|h|$ gives
        $\pm1$. Neither partial exists at $(0,3)$; both exist everywhere
        else. \textbf{Differentiability:} $C^{1}$, hence differentiable, on
        $\R^{2}\setminus\{(0,3)\}$; not differentiable at $(0,3)$.
        \textbf{Graph:} an elliptic cone with apex $(0,3,0)$. The contours
        $4x^{2}+(y-3)^{2}=c^{2}$ are ellipses about $(0,3)$ with semi-axes
        $c/2$ in the $x$-direction and $c$ in the $y$-direction: slope $2$
        across $x$, slope $1$ along $y$, and a sharp point at the apex.

  \item \textbf{Domain:} $\R^{2}$. \textbf{Partials:} for $x\neq2$,
        $f_{x}=1$ if $x>2$ and $f_{x}=-1$ if $x<2$; at a point $(2,y_{0})$,
        $\bigl(f(2+h,y_{0})-f(2,y_{0})\bigr)/h=|h|/h=\pm1$, so $f_{x}$ does
        not exist anywhere on the line $x=2$. Likewise $f_{y}=\pm1$ (the
        sign of $y$) for $y\neq0$, and $f_{y}$ does not exist on the line
        $y=0$. Both partials exist exactly off the two lines $x=2$ and
        $y=0$. \textbf{Differentiability:} off the two lines $f$ is, near
        each point, one of the four linear functions $\pm(x-2)\pm y$, so it
        is differentiable there; on the lines it is not, since a partial
        is missing. \textbf{Graph:} an upside-down square pyramid with its
        vertex at $(2,0,0)$; the contours $|x-2|+|y|=c$ are diamonds with
        corners $(2\pm c,0)$ and $(2,\pm c)$, and the surface is creased
        above the lines $x=2$ and $y=0$.

  \item \textbf{Domain:} the open quadrant $x,y>0$, as given.
        \textbf{Partials:} where $2x<5y$ (above the ray
        $y=\tfrac25x$), $f=2x$, so $f_{x}=2$, $f_{y}=0$. Where $2x>5y$,
        $f=5y$, so $f_{x}=0$, $f_{y}=5$. On the ray, at $(x_{0},y_{0})$ with
        $2x_{0}=5y_{0}$:
        \[
          f(x_{0}+h,y_{0})-f(x_{0},y_{0})=\min\{2x_{0}+2h,\,5y_{0}\}-5y_{0}
          =\min\{2h,\,0\},
        \]
        so the quotient is $0$ for $h>0$ and $2$ for $h<0$: no limit, no
        $f_{x}$. Similarly $f(x_{0},y_{0}+k)-f(x_{0},y_{0})=\min\{0,5k\}$
        gives one-sided slopes $0$ and $5$: no $f_{y}$. The partials exist
        exactly off the ray $y=\tfrac25x$. \textbf{Differentiability:}
        differentiable off the ray (locally linear), not on it.
        \textbf{Graph:} two planes, $z=2x$ and $z=5y$, meeting in a ridge
        above the ray: a roof. The contours are L-shaped corners, the level
        $c$ having its corner at $(c/2,\,c/5)$ on the ray, which is the
        Leontief isoquant picture.

  \item \textbf{Domain:} the denominator $2x^{2}+y^{2}$ vanishes only at
        the origin, where $f$ is defined separately, so $D=\R^{2}$.
        \textbf{Partials off the origin}, quotient rule with
        $N=2x^{2}+y^{2}$:
        \[
          f_{x}=\frac{3y\,N-3xy\cdot4x}{N^{2}}
               =\frac{3y\,(y^{2}-2x^{2})}{(2x^{2}+y^{2})^{2}},\qquad
          f_{y}=\frac{3x\,N-3xy\cdot2y}{N^{2}}
               =\frac{3x\,(2x^{2}-y^{2})}{(2x^{2}+y^{2})^{2}} .
        \]
        \textbf{At the origin}, from the definition: $f(h,0)=0$ and
        $f(0,h)=0$ for all $h$, so both quotients are $0$ and
        $f_{x}(0,0)=f_{y}(0,0)=0$. The partials exist \emph{everywhere}.
        \textbf{Continuity:} on the line $y=mx$, $x\neq0$,
        \[
          f(x,mx)=\frac{3mx^{2}}{(2+m^{2})x^{2}}=\frac{3m}{2+m^{2}},
        \]
        a constant that depends on $m$; on $y=x$ it is $1\neq f(0,0)$. So
        $f$ is not continuous at the origin, and hence not differentiable
        there, even though both partials exist. Away from the origin it is
        $C^{1}$ and differentiable. \textbf{Graph:} each ray from the origin
        carries its own constant height $3m/(2+m^{2})$, ranging between
        $-\tfrac{3\sqrt2}{4}$ and $\tfrac{3\sqrt2}{4}$ (attained at
        $m=\mp\sqrt2$); all these heights meet above the origin, so the
        surface is a twisted sheet torn at that point.

  \item \textbf{Domain:} $\R^{2}$. \textbf{Partials} off the origin:
        $f_{x}=\dfrac{2x^{3}}{\sqrt{x^{4}+y^{4}}}$,
        $f_{y}=\dfrac{2y^{3}}{\sqrt{x^{4}+y^{4}}}$. At the origin the
        formula fails, so use the quotient:
        $f(h,0)=\sqrt{h^{4}}=h^{2}$, and $h^{2}/h=h\to0$. So
        $f_{x}(0,0)=0$, and likewise $f_{y}(0,0)=0$: the partials exist
        everywhere, despite the broken formula. \textbf{Differentiability
        at the origin:} with $\mathbf{a}=\mathbf{0}$ the remainder is
        $r(\mathbf{h})=\sqrt{h_{1}^{4}+h_{2}^{4}}$. Since
        $h_{1}^{4}+h_{2}^{4}\le(h_{1}^{2}+h_{2}^{2})^{2}$ (the missing
        cross term $2h_{1}^{2}h_{2}^{2}$ is non-negative),
        \[
          0\le\frac{r(\mathbf{h})}{\|\mathbf{h}\|}
          \le\frac{h_{1}^{2}+h_{2}^{2}}{\|\mathbf{h}\|}=\|\mathbf{h}\|\to0 .
        \]
        So $f$ is differentiable at the origin with tangent plane $z=0$,
        and, being $C^{1}$ elsewhere, differentiable on all of $\R^{2}$.
        \textbf{Graph:} a bowl with a flat bottom at the origin; its
        contours $x^{4}+y^{4}=c^{2}$ are rounded squares. Nothing fails.

  \item \textbf{Domain:} $9-x^{2}-y^{2}\ge0$, the closed disc of radius $3$.
        \textbf{Partials:} for $x^{2}+y^{2}<9$,
        \[
          f_{x}=\frac{-x}{\sqrt{9-x^{2}-y^{2}}},\qquad
          f_{y}=\frac{-y}{\sqrt{9-x^{2}-y^{2}}},
        \]
        continuous on the open disc. On the boundary circle the radicand
        is $0$, the formulas divide by $0$, and, more basically, the points
        are not interior: every neighbourhood of a boundary point contains
        points outside the disc, where $f$ is not defined, so the two-sided
        difference quotients needed for the partials do not exist. \textbf{Differentiability:}
        $f$ is differentiable exactly on the open disc $x^{2}+y^{2}<9$.
        \textbf{Graph:} the upper hemisphere of radius $3$. Its tangent
        plane is horizontal at the top $(0,0,3)$, tilts more and more
        steeply outwards, and becomes vertical at the rim, where
        $\|\grad f\|\to\infty$.

  \item \textbf{Domain:} $\R^{2}$. \textbf{Partials:} where $x>y$,
        $f=x-y$, so $f_{x}=1$, $f_{y}=-1$; where $x<y$, $f=y-x$, so
        $f_{x}=-1$, $f_{y}=1$. At a point $(x_{0},x_{0})$ of the diagonal,
        $f(x_{0}+h,x_{0})-f(x_{0},x_{0})=|h|$, and $|h|/h=\pm1$ has no limit;
        the same holds in $y$. The partials exist exactly off the line
        $y=x$. \textbf{Differentiability:} differentiable off $y=x$, not on
        it, although $f$ is continuous everywhere. \textbf{Graph:} a
        V-shaped valley whose crease, at height $0$, runs along the line
        $y=x$; the contours are pairs of parallel lines $y=x\pm c$, the
        tilted-trough picture of Chapter~9 with a sharp floor.

  \item $f(P)=\ln(0+1)=0$. Chain rule:
        $f_{x}=\dfrac{2}{2x+y^{2}}$, $f_{y}=\dfrac{2y}{2x+y^{2}}$, so at
        $P$ (denominator $1$): $f_{x}=2$, $f_{y}=2$. Tangent plane:
        \[
          z=0+2(x-0)+2(y-1)=2x+2y-2 .
        \]
        At $Q$: $f(Q)=\ln(0.1+0.9604)=\ln1.0604\approx0.058646$, and
        $L(Q)=0.1+1.96-2=0.060000$. The plane overestimates by about
        $0.001354$.

  \item $f(P)=\mathrm{e}^{0}=1$. $f_{x}=\mathrm{e}^{x-2y}$,
        $f_{y}=-2\mathrm{e}^{x-2y}$, so at $P$: $f_{x}=1$, $f_{y}=-2$.
        Tangent plane:
        \[
          z=1+(x-2)-2(y-1)=x-2y+1 .
        \]
        At $Q$: $f(Q)=\mathrm{e}^{2.1-2.04}=\mathrm{e}^{0.06}\approx1.061837$,
        and $L(Q)=2.1-2.04+1=1.060000$. The plane underestimates by about
        $0.001837$: the exponential is convex, so it lies above all its
        tangent planes.

  \item $f(P)=1\cdot4\cdot\mathrm{e}^{0}=4$. Product rule in $x$, where
        both $x$ and $\mathrm{e}^{x-1}$ depend on $x$:
        \[
          f_{x}=y^{2}\mathrm{e}^{x-1}+xy^{2}\mathrm{e}^{x-1}
               =y^{2}\mathrm{e}^{x-1}(1+x),\qquad
          f_{y}=2xy\,\mathrm{e}^{x-1} .
        \]
        At $P$: $f_{x}=4\cdot1\cdot2=8$, $f_{y}=2\cdot1\cdot2\cdot1=4$.
        Tangent plane:
        \[
          z=4+8(x-1)+4(y-2) .
        \]
        At $Q$: $f(Q)=0.98\cdot2.03^{2}\cdot\mathrm{e}^{-0.02}
        =0.98\cdot4.1209\cdot0.980199\approx3.958515$, and
        $L(Q)=4-0.16+0.12=3.960000$. The plane is off by about
        $0.001485$.

  \item $f(P)=\dfrac{4}{2}=2$. Quotient rule with denominator $(x-y)^{2}$:
        \[
          f_{x}=\frac{(x-y)-(x+y)}{(x-y)^{2}}=\frac{-2y}{(x-y)^{2}},\qquad
          f_{y}=\frac{(x-y)+(x+y)}{(x-y)^{2}}=\frac{2x}{(x-y)^{2}} .
        \]
        At $P$: $f_{x}=-\tfrac24=-\tfrac12$, $f_{y}=\tfrac64=\tfrac32$.
        Tangent plane:
        \[
          z=2-\tfrac12(x-3)+\tfrac32(y-1) .
        \]
        At $Q$: $f(Q)=\dfrac{4.0}{2.2}\approx1.818182$, and
        $L(Q)=2-0.05-0.15=1.800000$. The error, about $0.018182$, is
        several times larger than in the neighbouring problems. The step
        keeps $x+y=4$ but moves the denominator $x-y$ from $2$ to $2.2$,
        and along that direction $f=4/(x-y)$ is strongly curved. The plane
        is a local approximation, and its quality depends on the curvature
        as much as on the length of the step.

  \item $F(P)=16^{1/4}\cdot1=2$.
        $F_{K}=\tfrac14K^{-3/4}L^{3/4}$, $F_{L}=\tfrac34K^{1/4}L^{-1/4}$; at
        $P$: $F_{K}=\tfrac14\cdot\tfrac18=\tfrac1{32}$ and
        $F_{L}=\tfrac34\cdot2=\tfrac32$. Tangent plane:
        \[
          z=2+\tfrac1{32}(K-16)+\tfrac32(L-1) .
        \]
        Economically: output $2$ plus marginal product times input change,
        for each input. At $Q$: $L(Q)=2+\tfrac{0.5}{32}+\tfrac32\cdot0.1
        =2+0.015625+0.15=2.165625$, while
        $F(Q)=16.5^{1/4}\cdot1.1^{3/4}\approx2.164789$. The estimate is
        high by about $0.000836$: with exponents summing to $1$ this
        Cobb--Douglas function is concave, so its graph lies below every
        tangent plane.

  \item $f(P)=\ln(2-2+1)=\ln1=0$. With $s=x+2y+3z$:
        $f_{x}=\tfrac1s$, $f_{y}=\tfrac2s$, $f_{z}=\tfrac3s$, and $s=1$ at
        $P$, so $\grad f(P)=(1,2,3)^{T}$. Tangent hyperplane:
        \[
          w=0+(x-2)+2(y+1)+3\bigl(z-\tfrac13\bigr)=x+2y+3z-1 .
        \]
        At $Q$: $s=2.02-1.98+1.05=1.09$, so $f(Q)=\ln1.09\approx0.086178$,
        and $L(Q)=1.09-1=0.090000$. The hyperplane overestimates by about
        $0.003822$, since $\ln$ is concave.

  \item $f(P)=2\mathrm{e}^{0}=2$. $f_{x}=\mathrm{e}^{yz}$,
        $f_{y}=xz\,\mathrm{e}^{yz}$, $f_{z}=xy\,\mathrm{e}^{yz}$; at $P$:
        $f_{x}=1$, $f_{y}=2\cdot5=10$, $f_{z}=2\cdot0=0$. Tangent
        hyperplane:
        \[
          w=2+(x-2)+10(y-0)+0\cdot(z-5)=x+10y .
        \]
        The change in $z$ does not enter at all, because $f_{z}(P)=0$. At
        $Q$: $f(Q)=2.05\,\mathrm{e}^{0.049}\approx2.152952$, and
        $L(Q)=2.05+0.1=2.150000$; error about $0.002952$.

  \item $f(P)=\tfrac66=1$. $f_{x}=\tfrac yz$, $f_{y}=\tfrac xz$,
        $f_{z}=-\tfrac{xy}{z^{2}}$; at $P$: $f_{x}=\tfrac12$,
        $f_{y}=\tfrac13$, $f_{z}=-\tfrac6{36}=-\tfrac16$. Tangent
        hyperplane:
        \[
          w=1+\tfrac12(x-2)+\tfrac13(y-3)-\tfrac16(z-6) .
        \]
        At $Q$: $L(Q)=1+0.05-0.033333-0.033333=0.983333$, and
        $f(Q)=\dfrac{2.1\cdot2.9}{6.2}=\dfrac{6.09}{6.2}\approx0.982258$;
        error about $0.001075$.

  \item Only the term $i=k$ contains $x_{k}$, so
        \[
          \pd{f}{x_{k}}=\frac{2x_{k}}{1+x_{k}^{2}},\qquad k=1,\ldots,n .
        \]
        At $(1,\ldots,1)$ every partial is $\tfrac22=1$:
        $\grad f=(1,1,\ldots,1)^{T}$ and $Df=(1\ \ 1\ \ \cdots\ \ 1)$.

  \item Chain rule with inner function $s=\sum_{i}x_{i}$, whose partial in
        $x_{k}$ is $1$:
        \[
          \pd{f}{x_{k}}=3\Bigl(\sum_{i=1}^{n}x_{i}\Bigr)^{2}\cdot1 ,
        \]
        the same for every $k$. For $n=4$ at $(1,2,3,4)$: $s=10$, so each
        partial is $300$, $\grad f=(300,300,300,300)^{T}$ and
        $Df=(300\ \ 300\ \ 300\ \ 300)$.

  \item The term containing $x_{k}$ is $x_{k}^{k+1}$, so
        $\pd{f}{x_{k}}=(k+1)\,x_{k}^{k}$. At $(1,\ldots,1)$:
        $\grad f=(2,3,\ldots,n+1)^{T}$, $Df=(2\ \ 3\ \ \cdots\ \ n+1)$. For
        $n=3$, $f=x_{1}^{2}+x_{2}^{3}+x_{3}^{4}$ with partials
        $2x_{1},\,3x_{2}^{2},\,4x_{3}^{3}$; at $(2,1,-1)$:
        $\grad f=(4,3,-4)^{T}$ and $Df=(4\ \ 3\ \ {-4})$.

  \item Write $f=\sum_{j}j^{2}x_{j}^{-3}$. The term with $j=k$ gives
        \[
          \pd{f}{x_{k}}=k^{2}\cdot(-3)x_{k}^{-4}=-\frac{3k^{2}}{x_{k}^{4}} .
        \]
        At $(1,\ldots,1)$: $\grad f=-3\,(1,4,9,\ldots,n^{2})^{T}$ and
        $Df=(-3\ \ {-12}\ \ {-27}\ \ \cdots\ \ {-3n^{2}})$.

  \item The first sum contributes $2x_{k}$; the second term, by the chain
        rule as in problem 39, contributes
        $-\tfrac1n\cdot2\sum_{i}x_{i}$. With the mean
        $\bar x=\tfrac1n\sum_{i}x_{i}$:
        \[
          \pd{f}{x_{k}}=2x_{k}-\frac2n\sum_{i=1}^{n}x_{i}=2\,(x_{k}-\bar x) .
        \]
        The gradient is zero exactly when $x_{k}=\bar x$ for every $k$,
        i.e.\ when all $x_{i}$ are equal. For $n=3$ at $(1,2,6)$:
        $\bar x=3$, so $\grad f=(-4,-2,6)^{T}$ and $Df=(-4\ \ {-2}\ \ 6)$.
        (The components always add up to $0$.)

  \item With $s=x^{2}+3yz$: $f_{x}=y^{2}\cdot\dfrac{2x}{s}$, since
        $y^{2}$ is a constant factor; $f_{y}$ needs the product rule,
        because $y$ appears in both factors:
        \[
          f_{y}=2y\ln s+y^{2}\cdot\frac{3z}{s};
        \]
        and $f_{z}=y^{2}\cdot\dfrac{3y}{s}=\dfrac{3y^{3}}{s}$. At
        $P=(1,1,1)$, $s=4$:
        \[
          f_{x}=\tfrac24=\tfrac12,\qquad
          f_{y}=2\ln4+\tfrac34\approx3.522589,\qquad
          f_{z}=\tfrac34 .
        \]
        So $\grad f(P)=\bigl(\tfrac12,\ 2\ln4+\tfrac34,\ \tfrac34\bigr)^{T}$
        and $Df(P)=\bigl(\tfrac12\ \ \ 2\ln4+\tfrac34\ \ \ \tfrac34\bigr)$.

  \item Write $E=\mathrm{e}^{y^{2}-2xz}$. In $x$ both factors move, so
        product and chain rule:
        \[
          f_{x}=3x^{2}E+x^{3}E\cdot(-2z)=x^{2}(3-2xz)\,E .
        \]
        In $y$ and $z$ only the exponential moves:
        $f_{y}=x^{3}E\cdot2y=2x^{3}y\,E$ and
        $f_{z}=x^{3}E\cdot(-2x)=-2x^{4}E$. At $P=(1,0,1)$,
        $E=\mathrm{e}^{-2}$:
        $f_{x}=1\cdot(3-2)\mathrm{e}^{-2}=\mathrm{e}^{-2}$, $f_{y}=0$,
        $f_{z}=-2\mathrm{e}^{-2}$. So
        $\grad f(P)=\mathrm{e}^{-2}(1,0,-2)^{T}\approx(0.135335,\,0,\,-0.270671)^{T}$
        and $Df(P)=\mathrm{e}^{-2}(1\ \ 0\ \ {-2})$.

  \item Use $\partial a^{g}/\partial x=a^{g}\ln a\cdot g_{x}$. In the first
        term $3^{y}$ is a constant factor; in the second the exponent $xy$
        has $x$-derivative $y$:
        \[
          f_{x}=2^{x}3^{y}\ln2+y\,5^{xy}\ln5,\qquad
          f_{y}=2^{x}3^{y}\ln3+x\,5^{xy}\ln5 .
        \]
        At $(1,1)$, $2^{x}3^{y}=6$ and $5^{xy}=5$:
        \[
          \grad f(1,1)=\begin{pmatrix}6\ln2+5\ln5\\6\ln3+5\ln5\end{pmatrix}
          \approx\begin{pmatrix}12.206073\\14.638863\end{pmatrix},
        \]
        and $Df(1,1)$ is the same pair written as a row.

  \item Only the term $i=k$ contains $x_{k}$, and its exponent $kx_{k}$ has
        derivative $k$:
        \[
          \pd{f}{x_{k}}=3^{\,kx_{k}}\ln3\cdot k=k\ln3\cdot3^{\,kx_{k}} .
        \]
        At the origin every $3^{0}=1$, so
        $\grad f=\ln3\,(1,2,\ldots,n)^{T}$ and
        $Df=\ln3\,(1\ \ 2\ \ \cdots\ \ n)$.

  \item Candidate: $f_{x}=3y$, $f_{y}=3x$, so $\grad f(a,b)=(3b,3a)$.
        Remainder:
        \[
          r(\mathbf{h})=3(a+h_{1})(b+h_{2})-3ab-3bh_{1}-3ah_{2}=3h_{1}h_{2} .
        \]
        Estimate:
        $\dfrac{|r(\mathbf{h})|}{\|\mathbf{h}\|}\le\dfrac{3\cdot\frac12
        \|\mathbf{h}\|^{2}}{\|\mathbf{h}\|}=\tfrac32\|\mathbf{h}\|\to0$.
        So $f$ is differentiable at every point, with derivative
        $(3b,\,3a)$.
  \item At the origin $\grad f=(0,0)$ and
        $r(\mathbf{h})=h_{1}^{2}+h_{2}^{2}=\|\mathbf{h}\|^{2}$, so
        $|r|/\|\mathbf{h}\|=\|\mathbf{h}\|\to0$: differentiable, with tangent
        plane $z=0$. At $(a,b)$: $\grad f=(2a,2b)$ and
        \[
          r(\mathbf{h})=(a+h_{1})^{2}+(b+h_{2})^{2}-a^{2}-b^{2}-2ah_{1}-2bh_{2}
          =h_{1}^{2}+h_{2}^{2},
        \]
        the same remainder, so the same estimate works everywhere.
  \item Candidate: $f_{x}=y^{2}$, $f_{y}=2xy$, so $\grad f(a,b)=(b^{2},2ab)$.
        Expand $(a+h_{1})(b^{2}+2bh_{2}+h_{2}^{2})
        =ab^{2}+2abh_{2}+ah_{2}^{2}+b^{2}h_{1}+2bh_{1}h_{2}+h_{1}h_{2}^{2}$
        and subtract $ab^{2}+b^{2}h_{1}+2abh_{2}$:
        \[
          r(\mathbf{h})=a\,h_{2}^{2}+2b\,h_{1}h_{2}+h_{1}h_{2}^{2} .
        \]
        Estimate:
        $\dfrac{|r|}{\|\mathbf{h}\|}\le|a|\,\|\mathbf{h}\|
        +|b|\,\|\mathbf{h}\|+\|\mathbf{h}\|^{2}\to0$. Differentiable
        everywhere, derivative $(b^{2},\,2ab)$.
  \item Candidate: $f_{x}=3y\mathrm{e}^{3x}$, $f_{y}=\mathrm{e}^{3x}$, so
        $\grad f(a,b)=\mathrm{e}^{3a}(3b,\,1)$. Write
        $\mathrm{e}^{3(a+h_{1})}=\mathrm{e}^{3a}\bigl(1+3h_{1}+R(3h_{1})\bigr)$:
        \[
          f(a+h_{1},b+h_{2})=\mathrm{e}^{3a}(b+h_{2})\bigl(1+3h_{1}+R\bigr)
          =\mathrm{e}^{3a}\bigl(b+3bh_{1}+h_{2}+3h_{1}h_{2}+(b+h_{2})R\bigr).
        \]
        Subtracting $b\mathrm{e}^{3a}$ and $\mathrm{e}^{3a}(3bh_{1}+h_{2})$
        leaves
        \[
          r(\mathbf{h})=\mathrm{e}^{3a}\bigl(3h_{1}h_{2}+(b+h_{2})\,R(3h_{1})\bigr).
        \]
        For $\|\mathbf{h}\|\le\tfrac13$ we have $|3h_{1}|\le1$, so
        $|R(3h_{1})|\le2\cdot9h_{1}^{2}\le18\|\mathbf{h}\|^{2}$, and
        \[
          \frac{|r|}{\|\mathbf{h}\|}\le\mathrm{e}^{3a}\Bigl(\tfrac32
          \|\mathbf{h}\|+18\bigl(|b|+\|\mathbf{h}\|\bigr)\|\mathbf{h}\|\Bigr)
          \to0 .
        \]
        Differentiable at every point.
  \item (a) On the $x$-axis $f(h,0)=(h^{2})^{1/3}=|h|^{2/3}$, so the
        difference quotient $|h|^{2/3}/h=\pm|h|^{-1/3}$ is unbounded:
        $f_{x}(0,0)$ does not exist (likewise $f_{y}$, with quotient
        $\pm4^{1/3}|h|^{-1/3}$). Not differentiable.
        (b) $g(\mathbf{h})=\|\mathbf{h}\|^{3/2}$. On an axis the quotient is
        $|h|^{3/2}/h=\pm|h|^{1/2}\to0$, so both partials are $0$. Then
        $r(\mathbf{h})=\|\mathbf{h}\|^{3/2}$ and
        $|r|/\|\mathbf{h}\|=\|\mathbf{h}\|^{1/2}\to0$: differentiable, with
        tangent plane $z=0$. The exponents decide: $2p=\tfrac23<1$ fails,
        $2p=\tfrac32>1$ works.
  \item \emph{Origin.} $f=0$ on both axes, so both partials are $0$. The
        remainder is $r(\mathbf{h})=|h_{1}h_{2}|\le\tfrac12\|\mathbf{h}\|^{2}$,
        so $|r|/\|\mathbf{h}\|\le\tfrac12\|\mathbf{h}\|\to0$: differentiable,
        with tangent plane $z=0$. \emph{At $(2,0)$.}
        $f(2,h)-f(2,0)=|2h|=2|h|$, and the quotient $2|h|/h=\pm2$ has no
        limit: $f_{y}(2,0)$ does not exist, so $f$ is not differentiable
        there. The graph is creased along both axes, except at the origin,
        where the two creases flatten out.
  \item \emph{Partials.} $f(h,0)=0$ and $f(0,h)=0$ for $h\neq0$, so both
        difference quotients are $0/h=0$: $f_{x}(0,0)=f_{y}(0,0)=0$.
        \emph{Lines.} On the $y$-axis $f=0$. On $y=mx$, $x\neq0$:
        \[
          f(x,mx)=\frac{mx^{3}}{x^{4}+m^{2}x^{2}}=\frac{mx}{x^{2}+m^{2}}
          \;\longrightarrow\;0\quad(x\to0),
        \]
        for $m\neq0$, and $f=0$ for $m=0$. Every straight line gives $0$.
        \emph{Parabola.} On $y=x^{2}$: $f=\dfrac{x^{4}}{x^{4}+x^{4}}=\tfrac12$
        for every $x\neq0$. So $f$ has no limit at the origin and is
        \textbf{not continuous} there, hence \textbf{not differentiable},
        although both partials exist and every line agrees.
  \item \emph{Partials.} $f=0$ on both axes, so $f_{x}(0,0)=f_{y}(0,0)=0$.
        \emph{Lines.} On $y=mx$:
        $f=\dfrac{2m^{2}x^{3}}{x^{2}+m^{4}x^{4}}=\dfrac{2m^{2}x}{1+m^{4}x^{2}}
        \to0$; on $x=0$, $f=0$. \emph{Parabola.} On $x=y^{2}$:
        $f=\dfrac{2y^{4}}{y^{4}+y^{4}}=1$ for $y\neq0$. Not continuous at the
        origin, hence not differentiable.
  \item (a) A smooth bowl, with no edges, creases or tips anywhere: it has a
        tangent plane at every point (horizontal at the bottom), so the
        function is differentiable everywhere. (b) A cone with a sharp tip:
        differentiable everywhere except at the tip, where every vertical
        cut through the tip is a V and no plane fits. (c) A pyramid whose
        four flat faces meet along two creases running over the axes:
        differentiable on the faces, not on the two creases (including
        the top point).
  \item The bowl $(x-1)^{2}+4y^{2}$ has no tip, so it is out. The level
        $c>0$ of $\sqrt{(x-1)^{2}+4y^{2}}$ is
        $\dfrac{(x-1)^{2}}{c^{2}}+\dfrac{y^{2}}{(c/2)^{2}}=1$, semi-axes $c$
        in $x$ and $c/2$ in $y$: twice as long in $x$, as in the picture.
        The third formula has semi-axes $c/2$ in $x$ and $c$ in $y$, the
        wrong way round. So (b) is $\sqrt{(x-1)^{2}+4y^{2}}$. At $(1,0)$:
        $f(1+h,0)-f(1,0)=\sqrt{h^{2}}=|h|$, and $|h|/h=\pm1$ has no limit;
        likewise $f(1,h)-f(1,0)=2|h|$. Neither partial exists, so $f$ is
        not differentiable at the tip, as the picture said.
\end{solutions}

\chapter{Higher-Order Partials, Young's Theorem, the Hessian, and Definiteness}

First-order partials describe slopes; second-order partials describe
\emph{curvature}. In one variable we used $f''$ to distinguish maxima from
minima. In several variables we need a whole \emph{matrix} of second
derivatives (the Hessian) and curvature is captured by its
definiteness, which we test with the leading principal minors (Sylvester's
criterion).

\section{Core Rules}

\begin{rulebox}{Direct second-order partials}
For $f(x,y)$,
\[
  f_{xx} \;=\; \pd{}{x}\!\left(\pd{f}{x}\right) \;=\; \frac{\partial^{2} f}{\partial x^{2}},
  \qquad
  f_{yy} \;=\; \frac{\partial^{2} f}{\partial y^{2}} .
\]
\end{rulebox}

\begin{rulebox}{Cross (mixed) partials}
\[
  f_{xy} \;=\; \pd{}{y}\!\left(\pd{f}{x}\right) \;=\; \frac{\partial^{2} f}{\partial y\,\partial x},
  \qquad
  f_{yx} \;=\; \pd{}{x}\!\left(\pd{f}{y}\right) .
\]
\end{rulebox}

\begin{rulebox}{Young's Theorem (Schwarz / Clairaut)}
If $f$ is $C^{2}$ on an open set (i.e.\ all second partials exist and
are continuous) then the order of differentiation is irrelevant:
\[
  f_{xy} \;=\; f_{yx} .
\]
This is the symmetry of the Hessian. All standard polynomial, exponential,
logarithmic, Cobb--Douglas, and CES forms are $C^{2}$ on their natural
domains.
\end{rulebox}

\begin{rulebox}{Hessian matrix}
For $f:\R^{n}\to\R$,
\[
  \Hess_f(\mathbf{x}) \;=\; \bigl[f_{x_ix_j}(\mathbf{x})\bigr]_{i,j=1}^{n}
  \;=\;
  \begin{pmatrix}
    f_{x_1x_1} & f_{x_1x_2} & \cdots & f_{x_1x_n} \\
    f_{x_2x_1} & f_{x_2x_2} & \cdots & f_{x_2x_n} \\
    \vdots     & \vdots     & \ddots & \vdots     \\
    f_{x_nx_1} & f_{x_nx_2} & \cdots & f_{x_nx_n}
  \end{pmatrix} .
\]
For $C^{2}$ functions $\Hess_f$ is symmetric.
\end{rulebox}

\begin{rulebox}{Definiteness}
A symmetric matrix $H$ is
\begin{itemize}[leftmargin=1.6em,itemsep=0.3ex]
  \item \emph{positive definite (PD)} if $\mathbf{v}^{T}H\mathbf{v}>0$ for
        every $\mathbf{v}\neq\mathbf{0}$ (curvature ``up'' in every
        direction $\Rightarrow$ local minimum;
  \item \emph{negative definite (ND)} if $\mathbf{v}^{T}H\mathbf{v}<0$ for
        every $\mathbf{v}\neq\mathbf{0}$) curvature ``down'' in every
        direction $\Rightarrow$ local maximum;
  \item \emph{indefinite} if $\mathbf{v}^{T}H\mathbf{v}$ takes both
        positive and negative values $\Rightarrow$ saddle point.
\end{itemize}
\end{rulebox}

\begin{rulebox}{Sylvester's criterion ($2\times2$ case)}
Let $H=\begin{pmatrix}a&b\\b&c\end{pmatrix}$ with leading principal minors
$D_1=a$ and $D_2=\det H=ac-b^{2}$. Then
\[
  \begin{aligned}
    D_1>0 \text{ and } D_2>0 &\;\Rightarrow\; H \text{ is PD}, \\
    D_1<0 \text{ and } D_2>0 &\;\Rightarrow\; H \text{ is ND}, \\
    D_2<0 &\;\Rightarrow\; H \text{ is indefinite (saddle)}, \\
    D_2=0 &\;\Rightarrow\; \text{test is inconclusive (semi-definite case)}.
  \end{aligned}
\]
Note: ND requires $D_1<0$ \emph{and} $D_2>0$; the alternating-sign
condition is what one generalises to higher $n$ (signs $-,+,-,+,\dots$ for
ND).
\end{rulebox}

\begin{rulebox}{Classification of critical points}
A point $\mathbf{x}_0$ is a \emph{critical point} of $f$ if
$\grad f(\mathbf{x}_0)=\mathbf{0}$. Then evaluate $\Hess_f(\mathbf{x}_0)$:
\begin{itemize}[leftmargin=1.6em,itemsep=0.3ex]
  \item $H$ PD at $\mathbf{x}_0$ $\Rightarrow$ strict local minimum;
  \item $H$ ND at $\mathbf{x}_0$ $\Rightarrow$ strict local maximum;
  \item $H$ indefinite at $\mathbf{x}_0$ $\Rightarrow$ saddle point;
  \item $D_2=0$ $\Rightarrow$ inconclusive, higher-order analysis
        required.
\end{itemize}
\end{rulebox}

\section{Worked Example}

\begin{worked}{Find and classify every critical point of $f(x,y)=x^{3}-3xy+y^{3}$}
\wstep{Step 1. First-order conditions: both partials zero, simultaneously.}
\[
  f_{x}=3x^{2}-3y=0
  \quad\Longrightarrow\quad y=x^{2},
\]
\[
  f_{y}=-3x+3y^{2}=0
  \quad\Longrightarrow\quad x=y^{2} .
\]
\wtrap{A critical point needs \emph{every} partial to vanish at the
\emph{same} point. Solving each equation separately and pairing the answers
arbitrarily is not the same thing --- the two equations form a system.}

\wstep{Step 2. Solve the system by substitution.}
Substituting $y=x^{2}$ into $x=y^{2}$:
\[
  x=x^{4}
  \;\Longrightarrow\;
  x^{4}-x=0
  \;\Longrightarrow\;
  x\bigl(x^{3}-1\bigr)=0
  \;\Longrightarrow\;
  x=0\ \text{or}\ x=1 .
\]
The corresponding $y=x^{2}$ gives the two critical points
\[
  (0,0)\qquad\text{and}\qquad(1,1).
\]
\wtrap{Do not divide by $x$ to get $x^{3}=1$. That silently discards the
root $x=0$ --- and $(0,0)$ turns out to be the more interesting of the two
points. Factor; never divide by something that may vanish.}

\wstep{Step 3. Build the Hessian once, symbolically.}
\[
  f_{xx}=6x,\qquad f_{yy}=6y,\qquad f_{xy}=f_{yx}=-3 ,
\]
\[
  \Hess f(x,y)=\begin{pmatrix}6x&-3\\-3&6y\end{pmatrix}.
\]
The two mixed partials agree, as Young's theorem guarantees for a function
with continuous second derivatives --- differentiating in $x$ then $y$ gives
the same result as $y$ then $x$. Computing both is a free check, and here
$\frac{\partial}{\partial y}(3x^{2}-3y)=-3$ and
$\frac{\partial}{\partial x}(-3x+3y^{2})=-3$. \checkmark

\wstep{Step 4. Classify $(0,0)$.}
\[
  \Hess f(0,0)=\begin{pmatrix}0&-3\\-3&0\end{pmatrix},
  \qquad
  \det=0-9=-9<0 .
\]
A negative determinant means the two eigenvalues have opposite signs, so the
form is indefinite:
\[
  \boxed{(0,0)\ \text{is a saddle point.}}
\]
No further test is needed, and the value of $f_{xx}$ is irrelevant here.
\wtrap{$\det\Hess<0$ settles the matter immediately --- saddle, whatever
$f_{xx}$ does. It is only when $\det\Hess>0$ that $f_{xx}$ has to be
consulted. Checking $f_{xx}$ first wastes effort and invites the wrong
conclusion.}

\wstep{Step 5. Classify $(1,1)$.}
\[
  \Hess f(1,1)=\begin{pmatrix}6&-3\\-3&6\end{pmatrix},
  \qquad
  \det=36-9=27>0,
  \qquad
  f_{xx}=6>0 .
\]
Positive determinant with a positive $f_{xx}$: positive definite, so
\[
  \boxed{(1,1)\ \text{is a strict local minimum},\qquad f(1,1)=1-3+1=-1 .}
\]

\wstep{Step 6. Verify the saddle directly, because it is cheap.}
Along the line $y=x$: $f(x,x)=2x^{3}-3x^{2}$, whose value near $0$ is
$\approx-3x^{2}<0$. Along $y=-x$: $f(x,-x)=x^{3}+3x^{2}-x^{3}=3x^{2}>0$. So
$f$ takes both signs arbitrarily close to $(0,0)$, where $f=0$ --- a saddle,
confirmed without any matrix.

\wstep{Step 7. The decision rule, complete.}
\begin{center}\footnotesize
\begin{tabular}{lll}
  \hline
  $\det\Hess$ & $f_{xx}$ & conclusion\\
  \hline
  $>0$ & $>0$ & strict local minimum\\
  $>0$ & $<0$ & strict local maximum\\
  $<0$ & --- & saddle point\\
  $=0$ & --- & \textbf{test fails}; investigate directly\\
  \hline
\end{tabular}
\end{center}
\wtrap{$\det\Hess=0$ means the test is silent, not that the point is a
saddle. $f=x^{4}+y^{4}$ (minimum), $f=-x^{4}-y^{4}$ (maximum) and
$f=x^{4}-y^{4}$ (saddle) all have the zero Hessian at the origin. When the
determinant vanishes you must argue directly, as in Step~6.}

\wstep{Step 8. Local is not global.}
Neither point is a global extremum: $f(x,0)=x^{3}$ runs to $\pm\infty$. On
an unbounded domain, second-order conditions can never certify a global
optimum by themselves --- that needs convexity of $f$ on the whole domain,
or a compact feasible set and Weierstrass. Volume~7 makes this the
difference between finding a candidate and having an answer.
\end{worked}

\section{Practice Problems}

\begin{problems}
  \item Compute all second partials of $f(x,y) = 4x^{3} - 5x^{2}y + 7y^{4}$
        and verify Young's theorem.
  \item Construct the Hessian of
        $f(x,y) = 5x^{2} + 6xy + 2y^{2} - 3x + 4y$.
  \item Classify the definiteness of
        $\Hess = \begin{pmatrix}6&2\\2&3\end{pmatrix}$ using Sylvester's
        criterion.
  \item Find and classify the critical points of
        $f(x,y) = x^{2} + 4y^{2} - 6x + 8y + 11$.
  \item Compute all second partials of $f(x,y) = (x+2y)^{3}$ and verify
        $f_{xy} = f_{yx}$.
  \item Construct the Hessian of $f(x,y) = x^{3} + 2xy^{2} - 4y$.
  \item Classify $\Hess = \begin{pmatrix}-4&2\\2&-5\end{pmatrix}$.
  \item Find and classify the critical points of
        $f(x,y) = -x^{2} - 3y^{2} + 8x - 6y + 5$.
  \item Compute all second partials of $f(x,y) = x^{2}\ln(y)$ ($y>0$) and
        verify Young's theorem.
  \item Construct the Hessian of $f(x,y) = \mathrm{e}^{2x-y}$.
  \item Classify $\Hess = \begin{pmatrix}2&4\\4&2\end{pmatrix}$.
  \item Find and classify the critical points of
        $f(x,y) = x^{2} - y^{2} + 6x - 4y$.
  \item Compute all second partials of $f(x,y) = \mathrm{e}^{x^{2}+y}$ and
        verify Young's theorem.
  \item Construct the Hessian of the Cobb--Douglas function
        $f(K,L) = K^{1/3}L^{2/3}$, $K,L>0$.
  \item Classify $\Hess = \begin{pmatrix}7&-3\\-3&1\end{pmatrix}$.
  \item Find and classify the critical points of
        $f(x,y) = x^{3} + y^{3} - 12x - 27y$.
  \item Compute all second partials of $f(x,y) = 1/(x+y)$ and verify
        Young's theorem.
  \item Construct the Hessian of
        $f(x,y,z) = x^{2} + 2y^{2} + 3z^{2} - xy + 2yz$.
  \item Classify $\Hess = \begin{pmatrix}-3&1\\1&-3\end{pmatrix}$.
  \item Find and classify the critical points of
        $f(x,y) = x^{2} + 4xy + y^{2} + 2y$.
  \item Compute all second partials of $f(x,y) = x^{1/2}y^{1/2}$
        ($x,y>0$) and verify $f_{xy} = f_{yx}$.
  \item Construct the Hessian of $f(x,y) = \ln(x^{2}+y^{2})$ at a generic
        point $(x,y)\neq(0,0)$.
  \item Find and classify the critical points of
        $f(x,y) = x^{4} + y^{4} - 4xy$.
  \item Compute all second partials of $f(x,y) = x\mathrm{e}^{y} + y\mathrm{e}^{x}$
        and verify Young's theorem.
  \item Construct the Hessian of $f(x,y) = x^{3} - 3xy^{2}$ (a harmonic
        function).
  \item Classify $\Hess = \begin{pmatrix}5&0\\0&-2\end{pmatrix}$ and state
        what this implies if $H$ is the Hessian at a critical point.
  \item Find and classify the critical points of
        $f(x,y) = xy - x^{2} - y^{2} + 6x + 3y$.
  \item Compute all second partials of $f(x,y) = \ln(xy)$ ($x,y>0$).
  \item Construct the Hessian of $f(x,y) = (x^{2}+y^{2})^{2}$ at a generic
        point.
  \item Classify $\Hess = \begin{pmatrix}-1&2\\2&-1\end{pmatrix}$.
  \item Find and classify the critical points of
        $f(x,y) = x^{3} - 3xy + y^{3}$.
  \item Compute all second partials of $f(x,y) = \sin(x)\cos(y)$ and
        verify Young's theorem.
  \item Construct the Hessian of $f(x,y,z) = \mathrm{e}^{xyz}$ at the point
        $(1,1,1)$ (evaluate numerically).
  \item Classify $\Hess = \begin{pmatrix}4&-2\\-2&4\end{pmatrix}$.
  \item Find and classify the critical points of the profit function
        $\pi(L,K) = 10L^{1/2}K^{1/2} - 2L - 3K$ (assume $L,K>0$).
  \item Compute all second partials of $f(x,y) = x^{2}y^{3}$ and verify
        Young's theorem.
  \item Construct the Hessian of
        $f(x,y) = \mathrm{e}^{x}\sin(y) - \mathrm{e}^{x}\cos(y)$.
  \item Classify $\Hess = \begin{pmatrix}9&6\\6&4\end{pmatrix}$. (Note the
        boundary case.)
  \item Find and classify the critical points of
        $f(x,y) = x^{2} + xy + y^{2} - 3x$.
  \item Compute all second partials of $f(x,y) = \arctan(y/x)$ ($x>0$) and
        verify Young's theorem.
  \item Construct the Hessian of the CES function
        $f(K,L) = (K^{1/2}+L^{1/2})^{2}$ at a generic point.
  \item Find and classify the critical points of
        $f(x,y) = -x^{4} - y^{4} + 8xy$.
  \item Compute all second partials of $f(x,y) = x^{2}\mathrm{e}^{-y}$ and
        verify Young's theorem.
  \item Construct the Hessian of $f(x,y) = \cos(xy)$ at the origin
        $(0,0)$.
  \item Classify $\Hess = \begin{pmatrix}0&3\\3&0\end{pmatrix}$.
  \item Find and classify the critical points of
        $f(x,y) = x^{3} + 3xy^{2} - 15x - 12y$.
  \item Construct the Hessian of $f(x,y) = \ln(1+x^{2}+y^{2})$ at the
        origin.
  \item Find and classify the critical points of
        $f(x,y) = \mathrm{e}^{-(x^{2}+y^{2})}$.
\end{problems}

\section{Solutions}

\begin{solutions}
  \item $f_x = 12x^{2} - 10xy$, $f_y = -5x^{2} + 28y^{3}$. Then
        $f_{xx} = 24x - 10y$, $f_{yy} = 84y^{2}$, and
        \[
          f_{xy} = \pd{}{y}\bigl(12x^{2}-10xy\bigr) = -10x, \qquad
          f_{yx} = \pd{}{x}\bigl(-5x^{2}+28y^{3}\bigr) = -10x . \; \checkmark
        \]

  \item $f_x = 10x + 6y - 3$, $f_y = 6x + 4y + 4$. Hessian
        \[
          \Hess = \begin{pmatrix}10&6\\6&4\end{pmatrix} .
        \]

  \item $D_1 = 6 > 0$, $D_2 = 18 - 4 = 14 > 0$. \textbf{PD.}

  \item $\grad f = (2x-6,\,8y+8) = (0,0) \Rightarrow x=3,\,y=-1$. Hessian
        $\Hess = \begin{pmatrix}2&0\\0&8\end{pmatrix}$: $D_1=2>0$,
        $D_2=16>0$, \textbf{PD}, so $(3,-1)$ is a strict \textbf{local
        minimum} with $f(3,-1) = 9+4-18-8+11 = -2$.

  \item Chain rule: $f_x = 3(x+2y)^{2}$, $f_y = 6(x+2y)^{2}$. So
        $f_{xx} = 6(x+2y)$, $f_{yy} = 24(x+2y)$, and
        \[
          f_{xy} = 12(x+2y) = f_{yx} . \; \checkmark
        \]

  \item $f_x = 3x^{2} + 2y^{2}$, $f_y = 4xy - 4$. Hessian
        \[
          \Hess = \begin{pmatrix}6x&4y\\4y&4x\end{pmatrix} .
        \]

  \item $D_1 = -4 < 0$, $D_2 = 20 - 4 = 16 > 0$. \textbf{ND.}

  \item $\grad f = (-2x+8,\,-6y-6) = (0,0) \Rightarrow x=4,\,y=-1$. Hessian
        $\Hess = \begin{pmatrix}-2&0\\0&-6\end{pmatrix}$: $D_1=-2<0$,
        $D_2=12>0$, \textbf{ND}, so $(4,-1)$ is a strict \textbf{local
        maximum} with $f(4,-1) = -16-3+32+6+5 = 24$.

  \item $f_x = 2x\ln(y)$, $f_y = x^{2}/y$. So $f_{xx} = 2\ln(y)$,
        $f_{yy} = -x^{2}/y^{2}$, and
        \[
          f_{xy} = 2x/y = f_{yx} . \; \checkmark
        \]

  \item $f_x = 2\mathrm{e}^{2x-y}$, $f_y = -\mathrm{e}^{2x-y}$. So
        $f_{xx} = 4\mathrm{e}^{2x-y}$, $f_{yy} = \mathrm{e}^{2x-y}$,
        $f_{xy} = -2\mathrm{e}^{2x-y}$. Hessian
        \[
          \Hess = \mathrm{e}^{2x-y}\begin{pmatrix}4&-2\\-2&1\end{pmatrix} .
        \]

  \item $D_1 = 2 > 0$, $D_2 = 4 - 16 = -12 < 0$. \textbf{Indefinite
        (saddle).}

  \item $\grad f = (2x+6,\,-2y-4) = (0,0) \Rightarrow x=-3,\,y=-2$. Hessian
        $\Hess = \begin{pmatrix}2&0\\0&-2\end{pmatrix}$: $D_2 = -4 < 0$,
        \textbf{indefinite}, so $(-3,-2)$ is a \textbf{saddle point}.

  \item Let $E = \mathrm{e}^{x^{2}+y}$. Then $f_x = 2xE$, $f_y = E$. So
        $f_{xx} = 2E + 4x^{2}E = (2+4x^{2})E$, $f_{yy} = E$, and
        \[
          f_{xy} = \pd{}{y}(2xE) = 2xE, \qquad
          f_{yx} = \pd{}{x}(E) = 2xE . \; \checkmark
        \]

  \item $f_K = \tfrac13K^{-2/3}L^{2/3}$, $f_L = \tfrac23K^{1/3}L^{-1/3}$.
        So
        \[
          f_{KK} = -\tfrac29K^{-5/3}L^{2/3}, \quad
          f_{LL} = -\tfrac29K^{1/3}L^{-4/3}, \quad
          f_{KL} = \tfrac29K^{-2/3}L^{-1/3} = f_{LK} . \; \checkmark
        \]
        Hessian
        \[
          \Hess = \tfrac29
          \begin{pmatrix}
            -K^{-5/3}L^{2/3} & K^{-2/3}L^{-1/3} \\
            K^{-2/3}L^{-1/3} & -K^{1/3}L^{-4/3}
          \end{pmatrix} .
        \]

  \item $D_1 = 7 > 0$, $D_2 = 7 - 9 = -2 < 0$. \textbf{Indefinite.}

  \item $\grad f = (3x^{2}-12,\,3y^{2}-27) = (0,0) \Rightarrow x^{2}=4,\,
        y^{2}=9$, giving $(\pm2,\pm3)$. Hessian
        $\Hess = \begin{pmatrix}6x&0\\0&6y\end{pmatrix}$, so $D_1=6x$,
        $D_2=36xy$.
        \begin{itemize}
          \item $(2,3)$: $D_1=12>0$, $D_2=216>0$, \textbf{PD}: local
                minimum, $f = 8+27-24-81 = -70$.
          \item $(2,-3)$: $D_2=-216<0$, \textbf{saddle}.
          \item $(-2,3)$: $D_2=-216<0$, \textbf{saddle}.
          \item $(-2,-3)$: $D_1=-12<0$, $D_2=216>0$, \textbf{ND}: local
                maximum, $f = -8-27+24+81 = 70$.
        \end{itemize}

  \item Let $S=x+y$, $f=S^{-1}$, $f_x = -S^{-2} = f_y$. So
        $f_{xx} = 2S^{-3} = f_{yy}$, and
        \[
          f_{xy} = 2S^{-3} = f_{yx} = \frac{2}{(x+y)^{3}} . \; \checkmark
        \]

  \item $f_x = 2x-y$, $f_y = 4y-x+2z$, $f_z = 6z+2y$. Hessian
        \[
          \Hess = \begin{pmatrix}2&-1&0\\-1&4&2\\0&2&6\end{pmatrix} .
        \]

  \item $D_1 = -3 < 0$, $D_2 = 9 - 1 = 8 > 0$. \textbf{ND.}

  \item $\grad f = (2x+4y,\,4x+2y+2) = (0,0)$. From the first, $x=-2y$;
        substituting, $-8y+2y+2=0 \Rightarrow y=1/3$, $x=-2/3$. Hessian
        $\Hess = \begin{pmatrix}2&4\\4&2\end{pmatrix}$: $D_1=2>0$,
        $D_2=4-16=-12<0$, \textbf{indefinite: saddle point} at
        $(-2/3,1/3)$.

  \item $f_x = \tfrac12x^{-1/2}y^{1/2}$, $f_y = \tfrac12x^{1/2}y^{-1/2}$.
        So
        \[
          f_{xx} = -\tfrac14x^{-3/2}y^{1/2}, \quad
          f_{yy} = -\tfrac14x^{1/2}y^{-3/2}, \quad
          f_{xy} = \tfrac14x^{-1/2}y^{-1/2} = f_{yx} . \; \checkmark
        \]

  \item Let $S=x^{2}+y^{2}$, $f=\ln S$, $f_x = 2x/S$, $f_y = 2y/S$. By the
        quotient rule:
        \[
          f_{xx} = \frac{2S-2x\cdot2x}{S^{2}} = \frac{2(y^{2}-x^{2})}{(x^{2}+y^{2})^{2}}, \quad
          f_{yy} = \frac{2(x^{2}-y^{2})}{(x^{2}+y^{2})^{2}}, \quad
          f_{xy} = \frac{-4xy}{(x^{2}+y^{2})^{2}} .
        \]
        Hessian
        \[
          \Hess = \frac{2}{(x^{2}+y^{2})^{2}}
          \begin{pmatrix}y^{2}-x^{2} & -2xy \\ -2xy & x^{2}-y^{2}\end{pmatrix} .
        \]
        Note $f_{xx}+f_{yy}=0$: $\ln(x^{2}+y^{2})$ is harmonic.

  \item $\grad f = (4x^{3}-4y,\,4y^{3}-4x) = (0,0) \Rightarrow y=x^{3}$
        and $x=y^{3}=x^{9}$. So $x(x^{8}-1)=0$, giving $x\in\{-1,0,1\}$.
        Critical points: $(0,0),(1,1),(-1,-1)$. Hessian
        $\Hess = \begin{pmatrix}12x^{2}&-4\\-4&12y^{2}\end{pmatrix}$.
        \begin{itemize}
          \item $(0,0)$: $\Hess = \begin{pmatrix}0&-4\\-4&0\end{pmatrix}$,
                $D_2=-16<0$, \textbf{saddle}.
          \item $(1,1)$: $\Hess = \begin{pmatrix}12&-4\\-4&12\end{pmatrix}$,
                $D_1=12>0$, $D_2=144-16=128>0$, \textbf{PD}: local
                minimum, $f=1+1-4=-2$.
          \item $(-1,-1)$: same $H$, \textbf{PD}: local minimum, $f=-2$.
        \end{itemize}

  \item $f_x = \mathrm{e}^{y}+y\mathrm{e}^{x}$, $f_y = x\mathrm{e}^{y}+\mathrm{e}^{x}$. So
        $f_{xx} = y\mathrm{e}^{x}$, $f_{yy} = x\mathrm{e}^{y}$, and
        \[
          f_{xy} = \mathrm{e}^{y}+\mathrm{e}^{x} = f_{yx} . \; \checkmark
        \]

  \item $f_x = 3x^{2}-3y^{2}$, $f_y = -6xy$. So $f_{xx}=6x$,
        $f_{yy}=-6x$, $f_{xy}=-6y=f_{yx}$. Hessian
        \[
          \Hess = \begin{pmatrix}6x&-6y\\-6y&-6x\end{pmatrix} .
        \]
        Note $f_{xx}+f_{yy}=0$: $f$ is harmonic.

  \item $D_1 = 5 > 0$, $D_2 = -10 < 0$. \textbf{Indefinite.} At a critical
        point this would mean a \textbf{saddle}.

  \item $\grad f = (y-2x+6,\,x-2y+3) = (0,0)$. So $2x-y=6$ and
        $x-2y=-3$. From the first, $y=2x-6$; substituting:
        $x-2(2x-6)=-3 \Rightarrow -3x+12=-3 \Rightarrow x=5$, $y=4$.
        Hessian $\Hess = \begin{pmatrix}-2&1\\1&-2\end{pmatrix}$:
        $D_1=-2<0$, $D_2=4-1=3>0$, \textbf{ND}, so $(5,4)$ is a
        \textbf{local maximum} with $f(5,4) = 20-25-16+30+12 = 21$.

  \item $\ln(xy) = \ln x + \ln y$. So $f_x = 1/x$, $f_y = 1/y$,
        $f_{xx} = -1/x^{2}$, $f_{yy} = -1/y^{2}$, $f_{xy} = 0 = f_{yx}$.
        \; $\checkmark$

  \item $f_x = 2(x^{2}+y^{2})\cdot2x = 4x(x^{2}+y^{2})$,
        $f_y = 4y(x^{2}+y^{2})$. Then
        \[
          f_{xx} = 4(x^{2}+y^{2})+4x\cdot2x = 4(3x^{2}+y^{2}), \quad
          f_{yy} = 4(x^{2}+3y^{2}), \quad
          f_{xy} = 4x\cdot2y = 8xy = f_{yx} .
        \]
        Hessian
        \[
          \Hess = 4\begin{pmatrix}3x^{2}+y^{2}&2xy\\2xy&x^{2}+3y^{2}\end{pmatrix} .
        \]

  \item $D_1 = -1 < 0$, $D_2 = 1 - 4 = -3 < 0$. Since $D_2<0$,
        \textbf{indefinite.}

  \item $\grad f = (3x^{2}-3y,\,-3x+3y^{2}) = (0,0) \Rightarrow y=x^{2}$
        and $x=y^{2}=x^{4}$, so $x(x^{3}-1)=0$, giving $x\in\{0,1\}$.
        Critical points: $(0,0)$ and $(1,1)$. Hessian
        $\Hess = \begin{pmatrix}6x&-3\\-3&6y\end{pmatrix}$.
        \begin{itemize}
          \item $(0,0)$: $\Hess = \begin{pmatrix}0&-3\\-3&0\end{pmatrix}$,
                $D_2=-9<0$, \textbf{saddle}.
          \item $(1,1)$: $\Hess = \begin{pmatrix}6&-3\\-3&6\end{pmatrix}$,
                $D_1=6>0$, $D_2=36-9=27>0$, \textbf{PD}: local minimum,
                $f=1-3+1=-1$.
        \end{itemize}

  \item $f_x = \cos(x)\cos(y)$, $f_y = -\sin(x)\sin(y)$. So
        $f_{xx} = -\sin(x)\cos(y)$, $f_{yy} = -\sin(x)\cos(y)$, and
        \[
          f_{xy} = -\cos(x)\sin(y) = f_{yx} . \; \checkmark
        \]

  \item Let $u=xyz$, $f=\mathrm{e}^{u}$. Then $f_x = yz\mathrm{e}^{u}$,
        $f_y = xz\mathrm{e}^{u}$, $f_z = xy\mathrm{e}^{u}$. Direct
        partials: $f_{xx} = (yz)^{2}\mathrm{e}^{u}$,
        $f_{yy} = (xz)^{2}\mathrm{e}^{u}$, $f_{zz} = (xy)^{2}\mathrm{e}^{u}$.
        Cross-partials (using the product rule on
        $f_x = yz\,\mathrm{e}^{xyz}$):
        \[
          f_{xy} = z\mathrm{e}^{xyz} + yz\cdot xz\mathrm{e}^{xyz}
                 = z\mathrm{e}^{xyz}(1+xyz) .
        \]
        Similarly $f_{xz} = y\mathrm{e}^{xyz}(1+xyz)$ and
        $f_{yz} = x\mathrm{e}^{xyz}(1+xyz)$. (Symmetry $f_{xy}=f_{yx}$
        etc.\ confirms Young's theorem.) At $(1,1,1)$: $u=1$,
        $\mathrm{e}^{u}=\mathrm{e}$, $1+xyz=2$, and $yz=xz=xy=1$, so
        \[
          \Hess(1,1,1) = \mathrm{e}\begin{pmatrix}1&2&2\\2&1&2\\2&2&1\end{pmatrix} .
        \]

  \item $D_1 = 4 > 0$, $D_2 = 16 - 4 = 12 > 0$. \textbf{PD.}

  \item $\pi_L = 5L^{-1/2}K^{1/2}-2$, $\pi_K = 5L^{1/2}K^{-1/2}-3$.
        Setting both to zero gives $5K^{1/2}/L^{1/2}=2$ and
        $5L^{1/2}/K^{1/2}=3$. Multiplying these two equations: $25=6$,
        which is a contradiction. So $\pi$ has \textbf{no interior
        critical point}. Why? The production function
        $Q=10L^{1/2}K^{1/2}$ is homogeneous of degree 1 (constant returns
        to scale), so along any ray $(\lambda L,\lambda K)$ revenue and
        cost both scale linearly: $\pi(\lambda L,\lambda K)=\lambda\pi(L,K)$.
        Hence profit is either identically zero, or unbounded above
        (depending on sign of $\pi$ along some direction). With
        $\max_{L/K}(10\sqrt{LK}-2L-3K)/K$: setting $t=L/K$ gives
        $\pi/K = 10\sqrt{t}-2t-3$, maximised at $t=25/4$ giving
        $\pi/K = 25-25/2-3 = 19/2 > 0$. So $\pi\to+\infty$ along the ray
        $L=(25/4)K$. The Hessian test does not apply because there is no
        critical point to classify. (Lesson: an inconsistent
        first-order system is itself a meaningful answer about the
        geometry of the problem.)

  \item $f_x = 2xy^{3}$, $f_y = 3x^{2}y^{2}$. So $f_{xx} = 2y^{3}$,
        $f_{yy} = 6x^{2}y$, and
        \[
          f_{xy} = 6xy^{2} = f_{yx} . \; \checkmark
        \]

  \item Let $g=\sin y-\cos y$, so $f=\mathrm{e}^{x}g$. Then
        $f_x = \mathrm{e}^{x}g$, $f_y = \mathrm{e}^{x}(\cos y+\sin y)$. So
        $f_{xx} = \mathrm{e}^{x}g = \mathrm{e}^{x}(\sin y-\cos y)$,
        $f_{yy} = \mathrm{e}^{x}(-\sin y+\cos y) = -\mathrm{e}^{x}g$, and
        \[
          f_{xy} = \mathrm{e}^{x}(\cos y+\sin y) = f_{yx} . \; \checkmark
        \]
        Hessian
        \[
          \Hess = \mathrm{e}^{x}
          \begin{pmatrix}
            \sin y-\cos y & \cos y+\sin y \\
            \cos y+\sin y & \cos y-\sin y
          \end{pmatrix} .
        \]

  \item $D_1 = 9 > 0$, $D_2 = 36 - 36 = 0$. \textbf{Inconclusive
        (positive semi-definite boundary case)}: Sylvester gives no
        verdict; we cannot rule out PSD with a zero eigenvalue. (In fact
        $H = (3,2)^{T}(3,2)$, so it is PSD with one zero eigenvalue.)

  \item $\grad f = (2x+y-3,\,x+2y) = (0,0)$. From the second: $x=-2y$;
        substituting: $-4y+y-3=0 \Rightarrow y=-1$, $x=2$. Critical point
        $(2,-1)$. Hessian $\Hess = \begin{pmatrix}2&1\\1&2\end{pmatrix}$:
        $D_1=2>0$, $D_2=4-1=3>0$, \textbf{PD}, so $(2,-1)$ is a strict
        \textbf{local minimum} with $f(2,-1) = 4-2+1-6 = -3$.

  \item With $u=y/x$, $f=\arctan u$: $f_x = -y/(x^{2}+y^{2})$,
        $f_y = x/(x^{2}+y^{2})$. Let $S=x^{2}+y^{2}$. Then
        \[
          f_{xx} = \frac{2xy}{S^{2}}, \qquad
          f_{yy} = \frac{-2xy}{S^{2}}, \qquad
          f_{xy} = \frac{-S+2y^{2}}{S^{2}} = \frac{y^{2}-x^{2}}{(x^{2}+y^{2})^{2}}
                  = f_{yx} . \; \checkmark
        \]
        (And $f_{xx}+f_{yy}=0$: $\arctan(y/x)$ is harmonic.)

  \item Let $g=K^{1/2}+L^{1/2}$. From sec06 sol.\ 23, $f_K = g/K^{1/2}$,
        $f_L = g/L^{1/2}$. Then
        \[
          f_{KK} = \frac{\tfrac12K^{-1/2}\cdot K^{1/2}-g\cdot\tfrac12K^{-1/2}}{K}
                  = \frac{\tfrac12-g/(2K^{1/2})}{K}
                  = \frac{1}{2K}-\frac{g}{2K^{3/2}} .
        \]
        More cleanly: $f_K = 1+(L/K)^{1/2}$, so
        $f_{KK} = -\tfrac12L^{1/2}K^{-3/2}$, and by symmetry
        $f_{LL} = -\tfrac12K^{1/2}L^{-3/2}$. Cross-partial:
        $f_{KL} = \tfrac12K^{-1/2}L^{-1/2} = f_{LK}$. \; $\checkmark$
        Hessian
        \[
          \Hess = \tfrac12
          \begin{pmatrix}
            -L^{1/2}K^{-3/2} & K^{-1/2}L^{-1/2} \\
            K^{-1/2}L^{-1/2} & -K^{1/2}L^{-3/2}
          \end{pmatrix} .
        \]

  \item $\grad f = (-4x^{3}+8y,\,-4y^{3}+8x) = (0,0) \Rightarrow x^{3}=2y$
        and $y^{3}=2x$. Substituting: $y=x^{3}/2$, then
        $(x^{3}/2)^{3}=2x \Rightarrow x^{9}/8=2x \Rightarrow x(x^{8}-16)=0$,
        so $x\in\{0,\pm16^{1/8}\}=\{0,\pm\sqrt2\}$. Points: $(0,0)$,
        $(\sqrt2,\sqrt2)$ (since $y=x^{3}/2=2\sqrt2/2=\sqrt2$), and
        $(-\sqrt2,-\sqrt2)$. Hessian
        $\Hess = \begin{pmatrix}-12x^{2}&8\\8&-12y^{2}\end{pmatrix}$.
        \begin{itemize}
          \item $(0,0)$: $\Hess = \begin{pmatrix}0&8\\8&0\end{pmatrix}$,
                $D_2=-64<0$, \textbf{saddle}.
          \item $(\sqrt2,\sqrt2)$: $x^{2}=y^{2}=2$,
                $\Hess = \begin{pmatrix}-24&8\\8&-24\end{pmatrix}$,
                $D_1=-24<0$, $D_2=576-64=512>0$, \textbf{ND}: local
                maximum, $f=-4-4+16=8$.
          \item $(-\sqrt2,-\sqrt2)$: same $H$, $f=-4-4+16=8$. \textbf{Local
                maximum.}
        \end{itemize}

  \item $f_x = 2x\mathrm{e}^{-y}$, $f_y = -x^{2}\mathrm{e}^{-y}$. So
        $f_{xx} = 2\mathrm{e}^{-y}$, $f_{yy} = x^{2}\mathrm{e}^{-y}$, and
        \[
          f_{xy} = -2x\mathrm{e}^{-y} = f_{yx} . \; \checkmark
        \]

  \item $f_x = -y\sin(xy)$, $f_y = -x\sin(xy)$. Then
        $f_{xx} = -y^{2}\cos(xy)$, $f_{yy} = -x^{2}\cos(xy)$, and
        (product rule)
        \[
          f_{xy} = -\sin(xy) - y\cdot x\cos(xy) = -\sin(xy) - xy\cos(xy) .
        \]
        At $(0,0)$: $\sin 0=0$, $\cos 0=1$. So $f_{xx}(0,0)=0$,
        $f_{yy}(0,0)=0$, $f_{xy}(0,0)=0$. Hessian
        \[
          \Hess(0,0) = \begin{pmatrix}0&0\\0&0\end{pmatrix} .
        \]

  \item $D_1 = 0$. Sylvester's stated criterion is inconclusive because
        $D_1$ is not strictly positive or negative. Directly:
        $D_2 = 0-9 = -9 < 0$, which by the third clause of Sylvester
        (sufficient on its own) gives \textbf{indefinite.}

  \item $\grad f = (3x^{2}+3y^{2}-15,\,6xy-12) = (0,0)$. From the second:
        $xy=2$. From the first: $x^{2}+y^{2}=5$. Combined with $xy=2$:
        $(x+y)^{2}=5+4=9$ so $x+y=\pm3$; $(x-y)^{2}=5-4=1$ so
        $x-y=\pm1$. Solutions $(x,y)\in\{(2,1),(1,2),(-2,-1),(-1,-2)\}$.
        Hessian $\Hess = \begin{pmatrix}6x&6y\\6y&6x\end{pmatrix}$, so
        $D_1=6x$, $D_2=36(x^{2}-y^{2})$.
        \begin{itemize}
          \item $(2,1)$: $D_1=12>0$, $D_2=36(4-1)=108>0$, \textbf{PD}:
                local minimum, $f=8+6-30-12=-28$.
          \item $(1,2)$: $D_1=6>0$, $D_2=36(1-4)=-108<0$, \textbf{saddle}.
          \item $(-2,-1)$: $D_1=-12<0$, $D_2=108>0$, \textbf{ND}: local
                maximum, $f=-8-6+30+12=28$.
          \item $(-1,-2)$: $D_1=-6<0$, $D_2=-108<0$, $D_2<0$ gives
                \textbf{saddle}.
        \end{itemize}

  \item Let $S=1+x^{2}+y^{2}$, $f=\ln S$, $f_x = 2x/S$, $f_y = 2y/S$.
        Then (using quotient rule):
        \[
          f_{xx} = \frac{2S-2x\cdot2x}{S^{2}} = \frac{2(1+y^{2}-x^{2})}{S^{2}}, \quad
          f_{yy} = \frac{2(1+x^{2}-y^{2})}{S^{2}}, \quad
          f_{xy} = \frac{-4xy}{S^{2}} = f_{yx} . \; \checkmark
        \]
        At the origin $(0,0)$: $S=1$, so
        \[
          \Hess(0,0) = \begin{pmatrix}2&0\\0&2\end{pmatrix} .
        \]

  \item $\grad f = \mathrm{e}^{-(x^{2}+y^{2})}(-2x,-2y) = (0,0)
        \Rightarrow x=y=0$, since $\mathrm{e}^{-r^{2}}>0$. Critical point
        $(0,0)$. Compute the Hessian: $f_x = -2x\mathrm{e}^{-(x^{2}+y^{2})}$,
        and at $(0,0)$, $f_{xx} = -2\mathrm{e}^{0}+(-2x)(-2x)\mathrm{e}^{0}\big|_{0,0}=-2$.
        Similarly $f_{yy}(0,0)=-2$, and
        $f_{xy}(0,0) = (-2x)(-2y)\mathrm{e}^{0}\big|_{0,0}=0$. Hessian
        $\Hess(0,0) = \begin{pmatrix}-2&0\\0&-2\end{pmatrix}$: $D_1=-2<0$,
        $D_2=4>0$, \textbf{ND: strict local maximum} at $(0,0)$ with
        $f(0,0)=1$. (Indeed, this is the global maximum, since $f\to0$
        at infinity and $f>0$ everywhere.)
\end{solutions}

\section{Additional Problems}

\begin{problems}
  \item Compute all second partials of $f(x,y) = 3x^{3} - 4x^{2}y + 5y^{3}$
        and verify Young's theorem.
  \item Construct the Hessian of
        $f(x,y) = 3x^{2} - 4xy + 2y^{2} + x - 5y$.
  \item Classify the definiteness of
        $\Hess = \begin{pmatrix}5&-1\\-1&4\end{pmatrix}$ using Sylvester's
        criterion.
  \item Find and classify the critical points of
        $f(x,y) = 2x^{2} + 3y^{2} - 4x + 6y + 1$.
  \item Compute all second partials of $f(x,y) = (2x+y)^{4}$ and verify
        $f_{xy} = f_{yx}$.
  \item Construct the Hessian of $f(x,y) = \mathrm{e}^{x+2y}$.
  \item Classify $\Hess = \begin{pmatrix}-3&1\\1&-2\end{pmatrix}$.
  \item Find and classify the critical points of
        $f(x,y) = -2x^{2} - 4y^{2} + 8x + 8y - 3$.
  \item Compute all second partials of $f(x,y) = xy\ln y$ ($y>0$) and
        verify Young's theorem.
  \item Construct the Hessian of $f(x,y) = x^{3}y + xy^{3}$.
  \item Classify $\Hess = \begin{pmatrix}2&3\\3&4\end{pmatrix}$.
  \item Find and classify the critical points of
        $f(x,y) = x^{2} - 2xy + 3y^{2} + 2x - 6y$.
  \item Compute all second partials of the Cobb--Douglas
        $f(K,L) = K^{1/2}L^{1/2}$ ($K,L>0$) and verify Young's theorem.
  \item Find and classify the critical points of
        $f(x,y) = x^{3} + y^{3} - 3x - 3y$.
  \item Classify $\Hess = \begin{pmatrix}-5&0\\0&-3\end{pmatrix}$ and state
        what it implies if $H$ is the Hessian at a critical point.
\end{problems}

\section{Additional Solutions}

\begin{solutions}
  \item $f_x = 9x^{2} - 8xy$, $f_y = -4x^{2} + 15y^{2}$. Then
        $f_{xx} = 18x - 8y$, $f_{yy} = 30y$, and
        \[
          f_{xy} = -8x = f_{yx} . \; \checkmark
        \]

  \item $f_x = 6x - 4y + 1$, $f_y = -4x + 4y - 5$. Hessian
        \[
          \Hess = \begin{pmatrix}6&-4\\-4&4\end{pmatrix} .
        \]

  \item $D_1 = 5 > 0$, $D_2 = 20 - 1 = 19 > 0$. \textbf{PD.}

  \item $\grad f = (4x-4,\,6y+6) = (0,0) \Rightarrow x=1,\,y=-1$. Hessian
        $\Hess = \begin{pmatrix}4&0\\0&6\end{pmatrix}$: $D_1=4>0$,
        $D_2=24>0$, \textbf{PD}, so $(1,-1)$ is a strict \textbf{local
        minimum} with $f(1,-1) = 2+3-4-6+1 = -4$.

  \item Chain rule: $f_x = 4(2x+y)^{3}\cdot2 = 8(2x+y)^{3}$,
        $f_y = 4(2x+y)^{3}$. So $f_{xx} = 48(2x+y)^{2}$,
        $f_{yy} = 12(2x+y)^{2}$, and
        \[
          f_{xy} = 24(2x+y)^{2} = f_{yx} . \; \checkmark
        \]

  \item $f_x = \mathrm{e}^{x+2y}$, $f_y = 2\mathrm{e}^{x+2y}$. So
        $f_{xx} = \mathrm{e}^{x+2y}$, $f_{yy} = 4\mathrm{e}^{x+2y}$,
        $f_{xy} = 2\mathrm{e}^{x+2y}$. Hessian
        \[
          \Hess = \mathrm{e}^{x+2y}\begin{pmatrix}1&2\\2&4\end{pmatrix} .
        \]
        (Note $\det\Hess = 0$: this is rank-1 PSD, consistent with $f$
        being a function of the single combination $x+2y$.)

  \item $D_1 = -3 < 0$, $D_2 = 6 - 1 = 5 > 0$. \textbf{ND.}

  \item $\grad f = (-4x+8,\,-8y+8) = (0,0) \Rightarrow x=2,\,y=1$. Hessian
        $\Hess = \begin{pmatrix}-4&0\\0&-8\end{pmatrix}$: $D_1=-4<0$,
        $D_2=32>0$, \textbf{ND}, so $(2,1)$ is a strict \textbf{local
        maximum} with $f(2,1) = -8-4+16+8-3 = 9$.

  \item $f_x = y\ln y$, $f_y = x\ln y + x = x(\ln y + 1)$. So
        $f_{xx} = 0$, $f_{yy} = x/y$, and
        \[
          f_{xy} = \ln y + 1 = f_{yx} . \; \checkmark
        \]

  \item $f_x = 3x^{2}y + y^{3}$, $f_y = x^{3} + 3xy^{2}$. So
        $f_{xx} = 6xy$, $f_{yy} = 6xy$,
        $f_{xy} = 3x^{2} + 3y^{2} = f_{yx}$. Hessian
        \[
          \Hess = \begin{pmatrix}
            6xy & 3x^{2}+3y^{2} \\
            3x^{2}+3y^{2} & 6xy
          \end{pmatrix} .
        \]

  \item $D_1 = 2 > 0$, $D_2 = 8 - 9 = -1 < 0$. \textbf{Indefinite
        (saddle).}

  \item $\grad f = (2x-2y+2,\,-2x+6y-6) = (0,0)$. From the first:
        $x = y-1$. Substitute: $-2(y-1)+6y-6 = 4y-4 = 0 \Rightarrow y=1$,
        $x=0$. Hessian $\Hess = \begin{pmatrix}2&-2\\-2&6\end{pmatrix}$:
        $D_1=2>0$, $D_2=12-4=8>0$, \textbf{PD}, so $(0,1)$ is a strict
        \textbf{local minimum} with $f(0,1) = 0-0+3+0-6 = -3$.

  \item $f_K = \tfrac12K^{-1/2}L^{1/2}$, $f_L = \tfrac12K^{1/2}L^{-1/2}$.
        So
        \[
          f_{KK} = -\tfrac14K^{-3/2}L^{1/2}, \quad
          f_{LL} = -\tfrac14K^{1/2}L^{-3/2}, \quad
          f_{KL} = \tfrac14K^{-1/2}L^{-1/2} = f_{LK} . \; \checkmark
        \]

  \item $\grad f = (3x^{2}-3,\,3y^{2}-3) = (0,0) \Rightarrow x^{2}=1,\,
        y^{2}=1$, giving $(\pm1,\pm1)$. Hessian
        $\Hess = \begin{pmatrix}6x&0\\0&6y\end{pmatrix}$, so $D_1=6x$,
        $D_2=36xy$.
        \begin{itemize}
          \item $(1,1)$: $D_1=6>0$, $D_2=36>0$, \textbf{PD}: local
                minimum, $f = 1+1-3-3 = -4$.
          \item $(1,-1)$: $D_2=-36<0$, \textbf{saddle}.
          \item $(-1,1)$: $D_2=-36<0$, \textbf{saddle}.
          \item $(-1,-1)$: $D_1=-6<0$, $D_2=36>0$, \textbf{ND}: local
                maximum, $f = -1-1+3+3 = 4$.
        \end{itemize}

  \item $D_1 = -5 < 0$, $D_2 = 15 > 0$. \textbf{ND}: at a critical point
        this would mean a \textbf{strict local maximum.}
\end{solutions}

\chapter{Multivariable Chain Rule and Total Differential}

When several inputs move at once, the total differential bundles the partial
derivatives into one linear approximation, and the multivariable chain rule
pipes that approximation through any composed function. Together they
generate the marginal rate of substitution along an indifference curve, the
marginal rate of technical substitution along an isoquant, and (as a
preview) the envelope theorem along an optimal value function.

\section{Core Rules}

\begin{rulebox}{Total Differential}
If $z=f(x,y)$ is differentiable at $(x,y)$, the \textbf{total differential}
is
\[
  \mathrm{d}z \;=\; f_{x}(x,y)\,\mathrm{d}x \;+\; f_{y}(x,y)\,\mathrm{d}y .
\]

\textbf{Interpretation.} $\mathrm{d}z$ is the best linear estimate of the
actual change $\Delta z$ when $x$ and $y$ shift by $\mathrm{d}x$ and
$\mathrm{d}y$. In economics, $\mathrm{d}U \approx MU_{x}\,\mathrm{d}x +
MU_{y}\,\mathrm{d}y$ tells how utility responds to a small basket change;
$\mathrm{d}Q \approx MP_{K}\,\mathrm{d}K + MP_{L}\,\mathrm{d}L$ does the same
for output.

For $n$ inputs, $z=f(x_{1},\ldots,x_{n})$:
\[
  \mathrm{d}z \;=\; \sum_{i=1}^{n} \pd{f}{x_{i}}\,\mathrm{d}x_{i} .
\]
\end{rulebox}

\begin{rulebox}{Chain Rule, Single Driver $t$}
If $z=f(x,y)$ with $x=x(t)$ and $y=y(t)$,
\[
  \dd{z}{t} \;=\; f_{x}\,\dd{x}{t} \;+\; f_{y}\,\dd{y}{t} .
\]
\end{rulebox}

\begin{rulebox}{Chain Rule, Two Drivers $s,t$}
If $z=f(x,y)$ with $x=x(s,t)$, $y=y(s,t)$,
\[
  \pd{z}{s} \;=\; f_{x}\,\pd{x}{s} + f_{y}\,\pd{y}{s},\qquad
  \pd{z}{t} \;=\; f_{x}\,\pd{x}{t} + f_{y}\,\pd{y}{t} .
\]
\end{rulebox}

\begin{rulebox}{Level Sets, MRS, MRTS}
On a level curve $f(x,y)=c$ the total differential is zero:
\[
  f_{x}\,\mathrm{d}x + f_{y}\,\mathrm{d}y = 0
  \;\Longrightarrow\;
  \dd{y}{x}\bigg|_{f=c} = -\frac{f_{x}}{f_{y}} .
\]
For a utility $U(x,y)$, the (signed) slope $-U_{x}/U_{y}$ is minus the
\textbf{MRS}; for a production function $Q(K,L)$, $-Q_{K}/Q_{L}$ is minus
the \textbf{MRTS}. We will state explicitly which sign convention we use in
each problem.
\end{rulebox}

\section{Worked Example}

\begin{worked}{With $z=x^{2}y$, $x=t^{2}$, $y=\ln t$, find $\mathrm{d}z/\mathrm{d}t$ two ways; then use the total differential on $Q=K^{0.3}L^{0.7}$}
\wstep{Step 1. Draw the dependency before writing a formula.}
$t$ drives both $x$ and $y$; $x$ and $y$ together drive $z$. So there are
\textbf{two paths} from $t$ to $z$, and the chain rule adds one term per
path:
\[
  \dd{z}{t}=\pd{z}{x}\dd{x}{t}+\pd{z}{y}\dd{y}{t} .
\]
\wtrap{One term per path. Losing a term is the standard error, and it is
invisible in the answer --- the result still looks like a derivative.
Sketching the tree (one node $t$, two arrows to $x$ and $y$, two arrows into
$z$) takes five seconds and prevents it.}

\wstep{Step 2. Note which symbol goes where.}
$\partial z/\partial x$ is \emph{partial}, because $z$ depends on two
things. $\mathrm{d}x/\mathrm{d}t$ is \emph{total}, because $x$ depends on
$t$ alone. And $\mathrm{d}z/\mathrm{d}t$ on the left is total, because once
the substitutions are made $z$ is a function of $t$ only. The notation is
carrying real information here.

\wstep{Step 3. Compute the four pieces.}
\[
  \pd{z}{x}=2xy,\qquad
  \pd{z}{y}=x^{2},\qquad
  \dd{x}{t}=2t,\qquad
  \dd{y}{t}=\frac1t .
\]

\wstep{Step 4. Assemble, then substitute the paths.}
\[
  \dd{z}{t}=2xy\cdot2t+x^{2}\cdot\frac1t
  =2t^{2}\ln t\cdot2t+\frac{t^{4}}{t}
  =\boxed{4t^{3}\ln t+t^{3}}
\]
\wtrap{Substitute $x=t^{2}$ and $y=\ln t$ only \emph{after} assembling the
formula. Substituting first turns the exercise into a one-variable problem
and defeats its purpose --- and in the many cases where $x(t)$ is not known
explicitly, it is not even possible.}

\wstep{Step 5. Check by direct substitution, which here is available.}
$z=(t^{2})^{2}\ln t=t^{4}\ln t$, so
\[
  \dd{z}{t}=4t^{3}\ln t+t^{4}\cdot\frac1t=4t^{3}\ln t+t^{3} . \checkmark
\]
Two routes, one answer. When both are available, use the direct one as the
check and the chain rule as the method.

\wstep{Step 6. The total differential: the chain rule with the $t$ removed.}
\[
  \mathrm{d}z=\pd{z}{x}\mathrm{d}x+\pd{z}{y}\mathrm{d}y
\]
is the first-order approximation to the change in $z$ caused by small
changes in $x$ and $y$ --- the two-variable linearisation. For
$Q=K^{0.3}L^{0.7}$:
\[
  \mathrm{d}Q=0.3K^{-0.7}L^{0.7}\,\mathrm{d}K+0.7K^{0.3}L^{-0.3}\,\mathrm{d}L .
\]

\wstep{Step 7. Divide by $Q$ and watch it turn into percentages.}
\[
  \frac{\mathrm{d}Q}{Q}
  =0.3\,\frac{\mathrm{d}K}{K}+0.7\,\frac{\mathrm{d}L}{L} .
\]
Every power cancels. This is the growth-accounting equation: the percentage
change in output is the \emph{exponent-weighted average} of the percentage
changes in inputs, and the exponents are the elasticities. If capital rises
$2\%$ while labour falls $1\%$:
\[
  \frac{\mathrm{d}Q}{Q}=0.3(0.02)+0.7(-0.01)=0.006-0.007=-0.001=-0.1\% .
\]
\wtrap{This identity holds only for Cobb--Douglas, where the elasticities
are the constant exponents. For a general $F(K,L)$ the weights are
$\frac{K F_{K}}{F}$ and $\frac{L F_{L}}{F}$, which vary from point to
point. Applying the constant-weight version to a CES function is a real
error.}

\wstep{Step 8. Two drivers instead of one.}
If $z=f(x,y)$ with $x=x(s,t)$ and $y=y(s,t)$, the same path-counting gives
one equation per driver:
\[
  \pd{z}{t}=\pd{z}{x}\pd{x}{t}+\pd{z}{y}\pd{y}{t},
  \qquad
  \pd{z}{s}=\pd{z}{x}\pd{x}{s}+\pd{z}{y}\pd{y}{s} .
\]
All derivatives are now partial, because two independent variables remain.
In matrix form these two lines are $\grad_{(s,t)}z=J^{\!\top}\grad_{(x,y)}z$
with $J$ the Jacobian of \chref{derivatives}{13}: the same rule, written once.
\end{worked}

\section{Practice Problems}

\begin{problems}
  \item Compute $\mathrm{d}z$ for $z=4x^{2}-5xy+2y^{2}$.
  \item Find $\dd{z}{t}$ for $z=x^{2}-y^{2}$ with $x=2t+1$, $y=t^{2}$.
  \item For $U(x,y)=x^{2}+4y=c$, find $\dd{y}{x}$ and state the unsigned
        MRS.
  \item Compute $\pd{z}{s}$ and $\pd{z}{t}$ for $z=x^{2}-3y$ with
        $x=s+2t$, $y=2s-t$.
  \item Compute $\mathrm{d}z$ for $z=x^{4}y^{-2}$.
  \item Find $\dd{z}{t}$ for $z=\mathrm{e}^{xy}$ with $x=\sin t$,
        $y=\cos t$.
  \item For the isoquant $Q(K,L)=K^{0.4}L^{0.6}=Q_{0}$, find $\dd{L}{K}$
        and the unsigned MRTS.
  \item Compute $\mathrm{d}z$ for $z=\ln(3x+2y)$.
  \item Compute $\pd{z}{s}$ and $\pd{z}{t}$ for $z=\mathrm{e}^{x}\cos y$
        with $x=st$, $y=s+t$.
  \item Find $\dd{z}{t}$ for $z=x^{2}y$ with $x=\ln t$, $y=t^{3}$.
  \item For $U(x,y)=x^{2/3}y^{1/3}=c$, find $\dd{y}{x}$ and the unsigned
        MRS.
  \item Compute $\mathrm{d}z$ for $z=\mathrm{e}^{x^{2}+y}$.
  \item Find $\dd{z}{t}$ for $z=\sqrt{x+y^{2}}$ with $x=t^{2}$, $y=t$.
  \item Compute $\pd{z}{s}$ and $\pd{z}{t}$ for $z=\ln(x^{2}+y^{2})$ with
        $x=s/t$, $y=t/s$.
  \item For $U(x,y)=\sqrt{x}+\sqrt{y}=c$, find $\dd{y}{x}$ and the unsigned
        MRS.
  \item Compute $\mathrm{d}z$ for $z=\dfrac{x+y}{x-y}$.
  \item Find $\dd{z}{t}$ for $z=\arctan(y/x)$ with $x=\mathrm{e}^{t}$,
        $y=\mathrm{e}^{2t}$.
  \item For $Q(K,L)=(aK^{1/2}+bL^{1/2})^{2}=Q_{0}$ (with $a,b>0$), find
        $\dd{L}{K}$ and the unsigned MRTS.
  \item Compute $\pd{z}{s}$ and $\pd{z}{t}$ for $z=x^{1/2}y^{1/2}$ with
        $x=s^{2}+t^{2}$, $y=st$.
  \item Compute $\mathrm{d}z$ for $z=x\sin y+y\cos x$.
  \item For $U(x,y)=\ln(x+1)+\ln(y+1)=c$, find $\dd{y}{x}$ and the
        unsigned MRS.
  \item Find $\dd{z}{t}$ for $z=x\ln y$ with $x=t^{2}+1$, $y=\mathrm{e}^{t}$.
  \item Compute $\mathrm{d}z$ for $z=x^{3}\mathrm{e}^{2y}$.
  \item Find $\dd{z}{t}$ for $z=x^{2}+y^{2}$ with $x=\cos t$, $y=\sin t$,
        and verify the result.
  \item For $U(x,y)=x+y^{2}=c$, find $\dd{y}{x}$ and state where the
        result is well-defined.
  \item Compute $\pd{z}{s}$ and $\pd{z}{t}$ for $z=xy^{2}$ with $x=s+t$,
        $y=s-t$.
  \item Compute $\mathrm{d}z$ for $z=\arctan(xy)$.
  \item Find $\dd{z}{t}$ for $z=\ln(x^{2}+y^{2})$ with $x=t$, $y=t^{2}$.
  \item For the isoquant $Q(K,L)=AK^{\alpha}L^{\beta}=Q_{0}$ (general
        Cobb--Douglas with $\alpha,\beta>0$), derive the MRTS in terms of
        $\alpha,\beta,K,L$.
  \item Compute $\mathrm{d}z$ for $z=\dfrac{x^{2}}{y}+\dfrac{y^{2}}{x}$.
  \item For $U(x,y)=\mathrm{e}^{x}+\mathrm{e}^{y}=c$, find $\dd{y}{x}$ and
        the unsigned MRS.
  \item Find $\dd{z}{t}$ for $z=\mathrm{e}^{x}\sin y$ with $x=2t$,
        $y=t^{2}$.
  \item Compute $\pd{z}{s}$ and $\pd{z}{t}$ for $z=x^{2}y-y^{3}$ with
        $x=s/t$, $y=st$.
  \item Compute $\mathrm{d}z$ for $z=(x+y)\mathrm{e}^{x-y}$.
  \item For the indifference curve $U(x,y)=\alpha\ln x+\beta\ln y=c$ (with
        $\alpha,\beta>0$), find $\dd{y}{x}$ and the MRS.
  \item Find $\dd{z}{t}$ for $z=\sqrt{xy}$ with $x=t^{3}$, $y=t$.
  \item Compute $\mathrm{d}z$ for $z=\ln(x^{2}+y^{2}+1)$.
  \item Compute $\pd{z}{s}$ and $\pd{z}{t}$ for $z=\mathrm{e}^{xy}$ with
        $x=s+t$, $y=s-t$.
  \item For $Q(K,L)=K+2\sqrt{KL}+L=Q_{0}$ (a perfect-square form), find
        the MRTS.
  \item Find $\dd{z}{t}$ for $z=\cos(xy)$ with $x=t$, $y=1/t$ (for $t>0$).
        Verify by direct substitution.
  \item For the indifference curve $U(x,y)=xy+x+y=c$, find $\dd{y}{x}$
        and the unsigned MRS.
  \item Use the total differential of $V(P,Y)=Y/P$ (real income) to
        approximate $\mathrm{d}V$ when $\mathrm{d}P=0.02P$ and
        $\mathrm{d}Y=0.03Y$. Interpret.
  \item Compute $\pd{z}{s}$ and $\pd{z}{t}$ for $z=x^{2}+2xy$ with
        $x=s\cos t$, $y=s\sin t$.
  \item For $U(x_{1},x_{2},x_{3})=x_{1}+2x_{2}+3x_{3}=c$, find
        $\pd{x_{3}}{x_{1}}$ and $\pd{x_{3}}{x_{2}}$ along the level set,
        and interpret as marginal rates of substitution.
\end{problems}

\section{Solutions}

\begin{solutions}
  \item $f_{x}=8x-5y$, $f_{y}=-5x+4y$.
        \[
          \mathrm{d}z \;=\; (8x-5y)\,\mathrm{d}x+(-5x+4y)\,\mathrm{d}y .
        \]

  \item $f_{x}=2x$, $f_{y}=-2y$, $\dd{x}{t}=2$, $\dd{y}{t}=2t$.
        \[
          \dd{z}{t}=2x\cdot 2+(-2y)\cdot 2t=4(2t+1)-4t^{3}=8t+4-4t^{3} .
        \]
        \textbf{Check.} $z=(2t+1)^{2}-t^{4}$, $\dd{z}{t}=8t+4-4t^{3}$.
        \checkmark

  \item $U_{x}=2x$, $U_{y}=4$. Signed slope $\dd{y}{x}=-2x/4=-x/2$.
        \textbf{MRS} $=x/2$.

  \item $f_{x}=2x$, $f_{y}=-3$; $x_{s}=1$, $x_{t}=2$, $y_{s}=2$,
        $y_{t}=-1$.
        \[
          \pd{z}{s}=2x(1)+(-3)(2)=2(s+2t)-6,\qquad
          \pd{z}{t}=2x(2)+(-3)(-1)=4(s+2t)+3 .
        \]

  \item $f_{x}=4x^{3}y^{-2}$, $f_{y}=-2x^{4}y^{-3}$.
        \[
          \mathrm{d}z \;=\; \frac{4x^{3}}{y^{2}}\,\mathrm{d}x
            -\frac{2x^{4}}{y^{3}}\,\mathrm{d}y .
        \]

  \item $f_{x}=y\mathrm{e}^{xy}$, $f_{y}=x\mathrm{e}^{xy}$;
        $\dd{x}{t}=\cos t$, $\dd{y}{t}=-\sin t$.
        \[
          \dd{z}{t}=\mathrm{e}^{xy}(y\cos t-x\sin t)
            =\mathrm{e}^{\sin t\cos t}(\cos^{2}t-\sin^{2}t)
            =\mathrm{e}^{\sin t\cos t}\cos(2t) .
        \]

  \item $Q_{K}=0.4K^{-0.6}L^{0.6}$, $Q_{L}=0.6K^{0.4}L^{-0.4}$.
        \[
          \frac{Q_{K}}{Q_{L}}=\frac{0.4}{0.6}\cdot\frac{L}{K}
            =\frac{2}{3}\,\frac{L}{K} .
        \]
        $\dd{L}{K}=-\tfrac{2}{3}L/K$. \textbf{MRTS} $=\tfrac{2}{3}L/K$.

  \item $f_{x}=3/(3x+2y)$, $f_{y}=2/(3x+2y)$.
        \[
          \mathrm{d}z \;=\; \frac{3\,\mathrm{d}x+2\,\mathrm{d}y}{3x+2y} .
        \]

  \item $f_{x}=\mathrm{e}^{x}\cos y$, $f_{y}=-\mathrm{e}^{x}\sin y$;
        $x_{s}=t$, $x_{t}=s$, $y_{s}=y_{t}=1$.
        \[
          \pd{z}{s}=\mathrm{e}^{st}\bigl(t\cos(s+t)-\sin(s+t)\bigr),\qquad
          \pd{z}{t}=\mathrm{e}^{st}\bigl(s\cos(s+t)-\sin(s+t)\bigr) .
        \]

  \item $f_{x}=2xy$, $f_{y}=x^{2}$; $\dd{x}{t}=1/t$, $\dd{y}{t}=3t^{2}$.
        \[
          \dd{z}{t}=2xy\cdot\frac{1}{t}+x^{2}\cdot 3t^{2}
            =\frac{2(\ln t)t^{3}}{t}+3t^{2}(\ln t)^{2}
            =2t^{2}\ln t+3t^{2}(\ln t)^{2} .
        \]

  \item $U_{x}=\tfrac{2}{3}x^{-1/3}y^{1/3}$,
        $U_{y}=\tfrac{1}{3}x^{2/3}y^{-2/3}$.
        \[
          \frac{U_{x}}{U_{y}}=2\,\frac{y}{x} .
        \]
        $\dd{y}{x}=-2y/x$. \textbf{MRS} $=2y/x$.

  \item $f_{x}=2x\mathrm{e}^{x^{2}+y}$, $f_{y}=\mathrm{e}^{x^{2}+y}$.
        \[
          \mathrm{d}z \;=\; \mathrm{e}^{x^{2}+y}(2x\,\mathrm{d}x+\mathrm{d}y) .
        \]

  \item $f_{x}=1/(2\sqrt{x+y^{2}})$, $f_{y}=y/\sqrt{x+y^{2}}$;
        $\dd{x}{t}=2t$, $\dd{y}{t}=1$.
        \[
          \dd{z}{t}=\frac{2t}{2\sqrt{x+y^{2}}}+\frac{y}{\sqrt{x+y^{2}}}
            =\frac{t+y}{\sqrt{x+y^{2}}}=\frac{2t}{\sqrt{2t^{2}}}
            =\frac{\sqrt{2}\,t}{|t|} .
        \]
        For $t>0$: $\dd{z}{t}=\sqrt{2}$. (Indeed $z=\sqrt{2}\,|t|$.)

  \item $f_{x}=2x/(x^{2}+y^{2})$, $f_{y}=2y/(x^{2}+y^{2})$;
        $x_{s}=1/t$, $x_{t}=-s/t^{2}$, $y_{s}=-t/s^{2}$, $y_{t}=1/s$.
        \[
          \pd{z}{s}=\frac{2}{x^{2}+y^{2}}\Bigl(\frac{x}{t}-\frac{yt}{s^{2}}\Bigr)
            =\frac{2}{x^{2}+y^{2}}\Bigl(\frac{s}{t^{2}}-\frac{t^{2}}{s^{3}}\Bigr) .
        \]
        \[
          \pd{z}{t}=\frac{2}{x^{2}+y^{2}}\Bigl(\frac{t}{s^{2}}-\frac{s^{2}}{t^{3}}\Bigr) .
        \]

  \item $U_{x}=1/(2\sqrt{x})$, $U_{y}=1/(2\sqrt{y})$.
        \[
          \frac{U_{x}}{U_{y}}=\sqrt{\frac{y}{x}} .
        \]
        $\dd{y}{x}=-\sqrt{y/x}$. \textbf{MRS} $=\sqrt{y/x}$.

  \item Quotient rule.
        \[
          f_{x}=\frac{(x-y)-(x+y)}{(x-y)^{2}}=\frac{-2y}{(x-y)^{2}},\qquad
          f_{y}=\frac{(x-y)+(x+y)}{(x-y)^{2}}=\frac{2x}{(x-y)^{2}} .
        \]
        \[
          \mathrm{d}z \;=\; \frac{-2y\,\mathrm{d}x+2x\,\mathrm{d}y}{(x-y)^{2}} .
        \]

  \item Let $u=y/x$. $f_{x}=-y/(x^{2}+y^{2})$, $f_{y}=x/(x^{2}+y^{2})$;
        $\dd{x}{t}=\mathrm{e}^{t}$, $\dd{y}{t}=2\mathrm{e}^{2t}$.
        \[
          \dd{z}{t}=\frac{-y\mathrm{e}^{t}+2x\mathrm{e}^{2t}}{x^{2}+y^{2}} .
        \]
        Substitute $x=\mathrm{e}^{t}$, $y=\mathrm{e}^{2t}$,
        $x^{2}+y^{2}=\mathrm{e}^{2t}(1+\mathrm{e}^{2t})$:
        \[
          \dd{z}{t}=\frac{-\mathrm{e}^{3t}+2\mathrm{e}^{3t}}
            {\mathrm{e}^{2t}(1+\mathrm{e}^{2t})}
            =\frac{\mathrm{e}^{t}}{1+\mathrm{e}^{2t}} .
        \]

  \item Let $A=aK^{1/2}+bL^{1/2}$, $Q=A^{2}$. $Q_{K}=2A\cdot(a/2)K^{-1/2}
        =aAK^{-1/2}$, $Q_{L}=bAL^{-1/2}$.
        \[
          \frac{Q_{K}}{Q_{L}}=\frac{a}{b}\sqrt{\frac{L}{K}} .
        \]
        $\dd{L}{K}=-(a/b)\sqrt{L/K}$. \textbf{MRTS} $=(a/b)\sqrt{L/K}$.

  \item $f_{x}=\tfrac12\sqrt{y/x}$, $f_{y}=\tfrac12\sqrt{x/y}$;
        $x_{s}=2s$, $x_{t}=2t$, $y_{s}=t$, $y_{t}=s$.
        \[
          \pd{z}{s}=s\sqrt{\frac{y}{x}}+\frac{t}{2}\sqrt{\frac{x}{y}},\qquad
          \pd{z}{t}=t\sqrt{\frac{y}{x}}+\frac{s}{2}\sqrt{\frac{x}{y}} .
        \]

  \item $f_{x}=\sin y-y\sin x$, $f_{y}=x\cos y+\cos x$.
        \[
          \mathrm{d}z \;=\; (\sin y-y\sin x)\,\mathrm{d}x
            +(x\cos y+\cos x)\,\mathrm{d}y .
        \]

  \item $U_{x}=1/(x+1)$, $U_{y}=1/(y+1)$.
        \[
          \frac{U_{x}}{U_{y}}=\frac{y+1}{x+1} .
        \]
        $\dd{y}{x}=-(y+1)/(x+1)$. \textbf{MRS} $=(y+1)/(x+1)$.

  \item $f_{x}=\ln y$, $f_{y}=x/y$; $\dd{x}{t}=2t$, $\dd{y}{t}=\mathrm{e}^{t}$.
        \[
          \dd{z}{t}=(\ln y)(2t)+\frac{x}{y}\mathrm{e}^{t}
            =t(2t)+\frac{(t^{2}+1)\mathrm{e}^{t}}{\mathrm{e}^{t}}
            =3t^{2}+1 .
        \]
        \textbf{Check.} $z=(t^{2}+1)t=t^{3}+t$, $\dd{z}{t}=3t^{2}+1$.
        \checkmark

  \item $f_{x}=3x^{2}\mathrm{e}^{2y}$, $f_{y}=2x^{3}\mathrm{e}^{2y}$.
        \[
          \mathrm{d}z \;=\; x^{2}\mathrm{e}^{2y}(3\,\mathrm{d}x+2x\,\mathrm{d}y) .
        \]

  \item $f_{x}=2x$, $f_{y}=2y$; $\dd{x}{t}=-\sin t$, $\dd{y}{t}=\cos t$.
        \[
          \dd{z}{t}=2x(-\sin t)+2y(\cos t)=-2\cos t\sin t+2\sin t\cos t=0 .
        \]
        \textbf{Check.} $z=\cos^{2}t+\sin^{2}t=1$, so $\dd{z}{t}=0$.
        \checkmark

  \item $U_{x}=1$, $U_{y}=2y$. Signed slope: $\dd{y}{x}=-1/(2y)$,
        well-defined for $y\neq 0$. The MRS in the usual sense is
        $1/(2y)$.

  \item $f_{x}=y^{2}$, $f_{y}=2xy$; $x_{s}=1$, $x_{t}=1$, $y_{s}=1$,
        $y_{t}=-1$.
        \[
          \pd{z}{s}=y^{2}(1)+2xy(1)=y^{2}+2xy=(s-t)^{2}+2(s+t)(s-t) .
        \]
        Simplify: $(s-t)^{2}+2(s^{2}-t^{2})=s^{2}-2st+t^{2}+2s^{2}-2t^{2}
          =3s^{2}-2st-t^{2}$.
        \[
          \pd{z}{t}=y^{2}(1)+2xy(-1)=y^{2}-2xy=(s-t)^{2}-2(s+t)(s-t)
            =-s^{2}-2st+3t^{2} .
        \]

  \item $f_{x}=y/(1+x^{2}y^{2})$, $f_{y}=x/(1+x^{2}y^{2})$.
        \[
          \mathrm{d}z \;=\; \frac{y\,\mathrm{d}x+x\,\mathrm{d}y}{1+x^{2}y^{2}} .
        \]

  \item $f_{x}=2x/(x^{2}+y^{2})$, $f_{y}=2y/(x^{2}+y^{2})$;
        $\dd{x}{t}=1$, $\dd{y}{t}=2t$.
        \[
          \dd{z}{t}=\frac{2x+2y(2t)}{x^{2}+y^{2}}
            =\frac{2t+4t^{3}}{t^{2}+t^{4}}
            =\frac{2t(1+2t^{2})}{t^{2}(1+t^{2})}
            =\frac{2(1+2t^{2})}{t(1+t^{2})} .
        \]

  \item $Q_{K}=A\alpha K^{\alpha-1}L^{\beta}$,
        $Q_{L}=A\beta K^{\alpha}L^{\beta-1}$.
        \[
          \frac{Q_{K}}{Q_{L}}=\frac{\alpha}{\beta}\cdot\frac{L}{K} .
        \]
        $\dd{L}{K}=-(\alpha/\beta)(L/K)$.
        \textbf{MRTS} $=(\alpha/\beta)(L/K)$.

  \item $f_{x}=2x/y-y^{2}/x^{2}$, $f_{y}=-x^{2}/y^{2}+2y/x$.
        \[
          \mathrm{d}z \;=\; \Bigl(\frac{2x}{y}-\frac{y^{2}}{x^{2}}\Bigr)\mathrm{d}x
            +\Bigl(\frac{2y}{x}-\frac{x^{2}}{y^{2}}\Bigr)\mathrm{d}y .
        \]

  \item $U_{x}=\mathrm{e}^{x}$, $U_{y}=\mathrm{e}^{y}$.
        \[
          \frac{U_{x}}{U_{y}}=\mathrm{e}^{x-y} .
        \]
        $\dd{y}{x}=-\mathrm{e}^{x-y}$. \textbf{MRS} $=\mathrm{e}^{x-y}$.

  \item $f_{x}=\mathrm{e}^{x}\sin y$, $f_{y}=\mathrm{e}^{x}\cos y$;
        $\dd{x}{t}=2$, $\dd{y}{t}=2t$.
        \[
          \dd{z}{t}=2\mathrm{e}^{x}\sin y+2t\mathrm{e}^{x}\cos y
            =2\mathrm{e}^{2t}(\sin t^{2}+t\cos t^{2}) .
        \]

  \item $f_{x}=2xy$, $f_{y}=x^{2}-3y^{2}$; $x_{s}=1/t$, $x_{t}=-s/t^{2}$,
        $y_{s}=t$, $y_{t}=s$. Note $xy=(s/t)(st)=s^{2}$.
        \[
          \pd{z}{s}=2xy\cdot\frac{1}{t}+(x^{2}-3y^{2})\cdot t
            =\frac{2s^{2}}{t}+t\Bigl(\frac{s^{2}}{t^{2}}-3s^{2}t^{2}\Bigr)
            =\frac{2s^{2}}{t}+\frac{s^{2}}{t}-3s^{2}t^{3}
            =\frac{3s^{2}}{t}-3s^{2}t^{3} .
        \]
        \[
          \pd{z}{t}=2xy\cdot\Bigl(-\frac{s}{t^{2}}\Bigr)+(x^{2}-3y^{2})\cdot s
            =-\frac{2s^{3}}{t^{2}}+s\Bigl(\frac{s^{2}}{t^{2}}-3s^{2}t^{2}\Bigr)
            =-\frac{s^{3}}{t^{2}}-3s^{3}t^{2} .
        \]
        \textbf{Check.} $z=(s/t)^{2}(st)-(st)^{3}=s^{3}/t-s^{3}t^{3}$. Then
        $\pd{z}{s}=3s^{2}/t-3s^{2}t^{3}$ and
        $\pd{z}{t}=-s^{3}/t^{2}-3s^{3}t^{2}$. \checkmark

  \item Product rule: $f_{x}=\mathrm{e}^{x-y}+(x+y)\mathrm{e}^{x-y}
        =(1+x+y)\mathrm{e}^{x-y}$, and
        $f_{y}=\mathrm{e}^{x-y}-(x+y)\mathrm{e}^{x-y}=(1-x-y)\mathrm{e}^{x-y}$.
        \[
          \mathrm{d}z \;=\; \mathrm{e}^{x-y}\bigl((1+x+y)\,\mathrm{d}x
            +(1-x-y)\,\mathrm{d}y\bigr) .
        \]

  \item $U_{x}=\alpha/x$, $U_{y}=\beta/y$.
        \[
          \frac{U_{x}}{U_{y}}=\frac{\alpha}{\beta}\cdot\frac{y}{x} .
        \]
        $\dd{y}{x}=-(\alpha/\beta)(y/x)$.
        \textbf{MRS} $=(\alpha/\beta)(y/x)$. (Standard log-Cobb-Douglas.)

  \item $f_{x}=\tfrac12\sqrt{y/x}$, $f_{y}=\tfrac12\sqrt{x/y}$;
        $\dd{x}{t}=3t^{2}$, $\dd{y}{t}=1$.
        \[
          \dd{z}{t}=\frac{3t^{2}}{2}\sqrt{\frac{y}{x}}
            +\frac{1}{2}\sqrt{\frac{x}{y}} .
        \]
        Substitute $y/x=1/t^{2}$, so $\sqrt{y/x}=1/t$ (for $t>0$), and
        $\sqrt{x/y}=t$:
        \[
          \dd{z}{t}=\frac{3t^{2}}{2}\cdot\frac{1}{t}+\frac{t}{2}
            =\frac{3t}{2}+\frac{t}{2}=2t .
        \]
        \textbf{Check.} $z=\sqrt{t^{3}\cdot t}=\sqrt{t^{4}}=t^{2}$,
        $\dd{z}{t}=2t$. \checkmark

  \item $f_{x}=2x/(x^{2}+y^{2}+1)$, $f_{y}=2y/(x^{2}+y^{2}+1)$.
        \[
          \mathrm{d}z \;=\; \frac{2x\,\mathrm{d}x+2y\,\mathrm{d}y}{x^{2}+y^{2}+1} .
        \]

  \item $f_{x}=y\mathrm{e}^{xy}$, $f_{y}=x\mathrm{e}^{xy}$; $x_{s}=1$,
        $x_{t}=1$, $y_{s}=1$, $y_{t}=-1$.
        \[
          \pd{z}{s}=y\mathrm{e}^{xy}+x\mathrm{e}^{xy}=(x+y)\mathrm{e}^{xy}
            =2s\,\mathrm{e}^{s^{2}-t^{2}} .
        \]
        \[
          \pd{z}{t}=y\mathrm{e}^{xy}-x\mathrm{e}^{xy}=(y-x)\mathrm{e}^{xy}
            =-2t\,\mathrm{e}^{s^{2}-t^{2}} .
        \]

  \item Note $Q=(\sqrt{K}+\sqrt{L})^{2}$. So $Q_{K}=2(\sqrt{K}+\sqrt{L})
        \cdot\dfrac{1}{2\sqrt{K}}=(\sqrt{K}+\sqrt{L})/\sqrt{K}$, and
        $Q_{L}=(\sqrt{K}+\sqrt{L})/\sqrt{L}$.
        \[
          \frac{Q_{K}}{Q_{L}}=\frac{\sqrt{L}}{\sqrt{K}}=\sqrt{\frac{L}{K}} .
        \]
        $\dd{L}{K}=-\sqrt{L/K}$. \textbf{MRTS} $=\sqrt{L/K}$.

  \item $f_{x}=-y\sin(xy)$, $f_{y}=-x\sin(xy)$; $\dd{x}{t}=1$,
        $\dd{y}{t}=-1/t^{2}$.
        \[
          \dd{z}{t}=-y\sin(xy)+\bigl(-x\sin(xy)\bigr)\Bigl(-\frac{1}{t^{2}}\Bigr)
            =-\sin(xy)\Bigl(y-\frac{x}{t^{2}}\Bigr) .
        \]
        With $x=t$, $y=1/t$: $xy=1$, and $y-x/t^{2}=1/t-1/t=0$. Hence
        $\dd{z}{t}=0$. \textbf{Check.} $z=\cos(xy)=\cos 1$ is constant.
        \checkmark

  \item $U_{x}=y+1$, $U_{y}=x+1$.
        \[
          \frac{U_{x}}{U_{y}}=\frac{y+1}{x+1} .
        \]
        $\dd{y}{x}=-(y+1)/(x+1)$. \textbf{MRS} $=(y+1)/(x+1)$.

  \item $V_{P}=-Y/P^{2}$, $V_{Y}=1/P$.
        \[
          \mathrm{d}V \;=\; -\frac{Y}{P^{2}}\,\mathrm{d}P+\frac{1}{P}\,\mathrm{d}Y
            =\frac{V}{1}\Bigl(-\frac{\mathrm{d}P}{P}+\frac{\mathrm{d}Y}{Y}\Bigr)\cdot 1
        \]
        since $V=Y/P$, we have $-Y\,\mathrm{d}P/P^{2}=-V\,\mathrm{d}P/P$
        and $\mathrm{d}Y/P=V\,\mathrm{d}Y/Y$. So
        \[
          \frac{\mathrm{d}V}{V} \;=\; \frac{\mathrm{d}Y}{Y}-\frac{\mathrm{d}P}{P}
            \;=\; 0.03-0.02 \;=\; 0.01 .
        \]
        Real income rises by about $1\%$. Interpretation: nominal income
        growth less inflation gives real income growth (log-linearised).

  \item $f_{x}=2x+2y$, $f_{y}=2x$; $x_{s}=\cos t$, $x_{t}=-s\sin t$,
        $y_{s}=\sin t$, $y_{t}=s\cos t$.
        \[
          \pd{z}{s}=(2x+2y)\cos t+2x\sin t .
        \]
        Substitute $x=s\cos t$, $y=s\sin t$: $2x+2y=2s(\cos t+\sin t)$,
        so
        \[
          \pd{z}{s}=2s(\cos t+\sin t)\cos t+2s\cos t\sin t
            =2s\cos^{2}t+4s\sin t\cos t
            =2s(\cos^{2}t+2\sin t\cos t) .
        \]
        \[
          \pd{z}{t}=(2x+2y)(-s\sin t)+2x(s\cos t)
            =-2s^{2}(\cos t+\sin t)\sin t+2s^{2}\cos^{2}t .
        \]
        Expanding: $-2s^{2}(\cos t\sin t+\sin^{2}t)+2s^{2}\cos^{2}t
          =2s^{2}(\cos^{2}t-\sin^{2}t-\sin t\cos t)
          =2s^{2}\bigl(\cos 2t-\tfrac12\sin 2t\bigr)$.

  \item On $U=c$, $\mathrm{d}U=\mathrm{d}x_{1}+2\,\mathrm{d}x_{2}
        +3\,\mathrm{d}x_{3}=0$. Solve for $\mathrm{d}x_{3}$:
        \[
          \mathrm{d}x_{3}=-\tfrac13\,\mathrm{d}x_{1}-\tfrac23\,\mathrm{d}x_{2}
          \;\Longrightarrow\;
          \pd{x_{3}}{x_{1}}\bigg|_{U=c}=-\tfrac13,\qquad
          \pd{x_{3}}{x_{2}}\bigg|_{U=c}=-\tfrac23 .
        \]
        Unsigned MRSes: $1/3$ (between $x_{3}$ and $x_{1}$) and $2/3$
        (between $x_{3}$ and $x_{2}$).
\end{solutions}

\section{Additional Problems}

\begin{problems}
  \item Compute $\mathrm{d}z$ for $z=3x^{2}+2xy-5y^{2}$.
  \item Find $\dd{z}{t}$ for $z=x^{2}+2y$ with $x=3t+1$, $y=t^{2}-1$.
  \item For $U(x,y)=x^{2}+y^{2}=c$, find $\dd{y}{x}$ and the unsigned MRS.
  \item Compute $\pd{z}{s}$ and $\pd{z}{t}$ for $z=x^{2}+y^{2}$ with
        $x=s+t$, $y=s-t$.
  \item Compute $\mathrm{d}z$ for $z=x^{3}\sqrt{y}$.
  \item Find $\dd{z}{t}$ for $z=\ln(x^{2}+y^{2})$ with $x=\mathrm{e}^{t}$,
        $y=\mathrm{e}^{-t}$.
  \item For the isoquant $Q(K,L)=K^{0.3}L^{0.7}=Q_{0}$, find $\dd{L}{K}$
        and the unsigned MRTS.
  \item Compute $\mathrm{d}z$ for $z=\mathrm{e}^{xy}$.
  \item Compute $\pd{z}{s}$ and $\pd{z}{t}$ for $z=xy^{2}$ with
        $x=s^{2}+t$, $y=s-t^{2}$.
  \item For $U(x,y)=\sqrt{xy}=c$, find $\dd{y}{x}$ and the unsigned MRS.
  \item Find $\dd{z}{t}$ for $z=\mathrm{e}^{x+y}$ with $x=\sin t$,
        $y=\cos t$.
  \item Compute $\mathrm{d}z$ for $z=\dfrac{x}{x+y}$.
  \item For $U(x,y)=x+\ln y=c$, find $\dd{y}{x}$ and the unsigned MRS.
  \item Find $\dd{z}{t}$ for $z=x^{2}y-y^{3}$ with $x=t^{2}$, $y=2t$.
  \item Use the total differential of $W(P,Q)=PQ$ (total revenue) to
        approximate $\mathrm{d}W$ when $\mathrm{d}P=0.01P$ and
        $\mathrm{d}Q=-0.005Q$. Interpret.
\end{problems}

In problems 16--18, approximate the change $\Delta f$ by the total
differential $\mathrm{d}f$ at the point $P$, then compute the exact change
and compare.

\begin{problems}[resume]
  \item $Q(x,y,z)=4\sqrt{x^{2}+yz}$ at $P=(3,2,8)$, with
        $\mathrm{d}x=0.1$, $\mathrm{d}y=-0.2$, $\mathrm{d}z=0.3$.
  \item $Y(x,y,z)=2\sqrt{xy+z^{2}}$ at $P=(2,8,3)$, with
        $\mathrm{d}x=0.05$, $\mathrm{d}y=0.1$, $\mathrm{d}z=-0.05$.
  \item $f(x,y)=\ln(x^{2}y)-xy$ at $P=(1,2)$, with
        $\mathrm{d}x=\mathrm{d}y=0.1$; explain why the approximation is
        poorer here.
\end{problems}

\section{Additional Solutions}

\begin{solutions}
  \item $f_{x}=6x+2y$, $f_{y}=2x-10y$.
        \[
          \mathrm{d}z \;=\; (6x+2y)\,\mathrm{d}x+(2x-10y)\,\mathrm{d}y .
        \]

  \item $f_{x}=2x$, $f_{y}=2$; $\dd{x}{t}=3$, $\dd{y}{t}=2t$.
        \[
          \dd{z}{t} \;=\; 2x\cdot 3+2\cdot 2t \;=\; 6(3t+1)+4t
            \;=\; 22t+6 .
        \]
        \textbf{Check.} $z=(3t+1)^{2}+2(t^{2}-1)=9t^{2}+6t+1+2t^{2}-2
        =11t^{2}+6t-1$, $\dd{z}{t}=22t+6$. \checkmark

  \item $U_{x}=2x$, $U_{y}=2y$. Signed slope: $\dd{y}{x}=-x/y$.
        \textbf{MRS} $=x/y$ (unsigned).

  \item $f_{x}=2x$, $f_{y}=2y$; $x_{s}=1$, $x_{t}=1$, $y_{s}=1$,
        $y_{t}=-1$.
        \[
          \pd{z}{s}=2x+2y=2(s+t)+2(s-t)=4s .
        \]
        \[
          \pd{z}{t}=2x-2y=2(s+t)-2(s-t)=4t .
        \]
        \textbf{Check.} $z=(s+t)^{2}+(s-t)^{2}=2(s^{2}+t^{2})$,
        $\pd{z}{s}=4s$, $\pd{z}{t}=4t$. \checkmark

  \item $f_{x}=3x^{2}\sqrt{y}$, $f_{y}=\dfrac{x^{3}}{2\sqrt{y}}$.
        \[
          \mathrm{d}z \;=\; 3x^{2}\sqrt{y}\,\mathrm{d}x
            +\frac{x^{3}}{2\sqrt{y}}\,\mathrm{d}y .
        \]

  \item $f_{x}=2x/(x^{2}+y^{2})$, $f_{y}=2y/(x^{2}+y^{2})$;
        $\dd{x}{t}=\mathrm{e}^{t}$, $\dd{y}{t}=-\mathrm{e}^{-t}$. With
        $x=\mathrm{e}^{t}$, $y=\mathrm{e}^{-t}$,
        $x^{2}+y^{2}=\mathrm{e}^{2t}+\mathrm{e}^{-2t}$.
        \[
          \dd{z}{t} \;=\;
            \frac{2\mathrm{e}^{t}\cdot\mathrm{e}^{t}
              +2\mathrm{e}^{-t}\cdot(-\mathrm{e}^{-t})}
              {\mathrm{e}^{2t}+\mathrm{e}^{-2t}}
          \;=\; \frac{2\mathrm{e}^{2t}-2\mathrm{e}^{-2t}}
              {\mathrm{e}^{2t}+\mathrm{e}^{-2t}}
          \;=\; \frac{2(\mathrm{e}^{2t}-\mathrm{e}^{-2t})}
              {\mathrm{e}^{2t}+\mathrm{e}^{-2t}}
          =2\tanh(2t) .
        \]

  \item $Q_{K}=0.3K^{-0.7}L^{0.7}$, $Q_{L}=0.7K^{0.3}L^{-0.3}$.
        \[
          \frac{Q_{K}}{Q_{L}}=\frac{0.3}{0.7}\cdot\frac{L}{K}
            =\frac{3}{7}\cdot\frac{L}{K} .
        \]
        $\dd{L}{K}=-(3/7)(L/K)$. \textbf{MRTS} $=(3/7)(L/K)$.

  \item $f_{x}=y\mathrm{e}^{xy}$, $f_{y}=x\mathrm{e}^{xy}$.
        \[
          \mathrm{d}z \;=\; \mathrm{e}^{xy}(y\,\mathrm{d}x+x\,\mathrm{d}y) .
        \]

  \item $f_{x}=y^{2}$, $f_{y}=2xy$; $x_{s}=2s$, $x_{t}=1$, $y_{s}=1$,
        $y_{t}=-2t$.
        \[
          \pd{z}{s}=y^{2}\cdot 2s+2xy\cdot 1=2sy^{2}+2xy .
        \]
        \[
          \pd{z}{t}=y^{2}\cdot 1+2xy\cdot(-2t)=y^{2}-4txy .
        \]

  \item $U_{x}=\tfrac{1}{2}\sqrt{y/x}$, $U_{y}=\tfrac{1}{2}\sqrt{x/y}$.
        \[
          \frac{U_{x}}{U_{y}}=\sqrt{\frac{y/x}{x/y}}=\frac{y}{x} .
        \]
        $\dd{y}{x}=-y/x$. \textbf{MRS} $=y/x$.

  \item $f_{x}=f_{y}=\mathrm{e}^{x+y}$; $\dd{x}{t}=\cos t$,
        $\dd{y}{t}=-\sin t$.
        \[
          \dd{z}{t}=\mathrm{e}^{x+y}(\cos t-\sin t)
            =\mathrm{e}^{\sin t+\cos t}(\cos t-\sin t) .
        \]

  \item Quotient rule with $u=x$, $v=x+y$:
        \[
          f_{x}=\frac{(x+y)-x}{(x+y)^{2}}=\frac{y}{(x+y)^{2}},\qquad
          f_{y}=\frac{-x}{(x+y)^{2}} .
        \]
        \[
          \mathrm{d}z=\frac{y\,\mathrm{d}x-x\,\mathrm{d}y}{(x+y)^{2}} .
        \]

  \item $U_{x}=1$, $U_{y}=1/y$. $U_{x}/U_{y}=y$. Signed slope
        $\dd{y}{x}=-y$. \textbf{MRS} $=y$.

  \item $f_{x}=2xy$, $f_{y}=x^{2}-3y^{2}$; $\dd{x}{t}=2t$,
        $\dd{y}{t}=2$.
        \[
          \dd{z}{t}=2xy\cdot 2t+(x^{2}-3y^{2})\cdot 2
            =4t\cdot t^{2}\cdot 2t+2(t^{4}-3\cdot 4t^{2})
            =8t^{4}+2t^{4}-24t^{2}=10t^{4}-24t^{2} .
        \]
        \textbf{Check.} $z=t^{4}\cdot 2t-8t^{3}=2t^{5}-8t^{3}$,
        $\dd{z}{t}=10t^{4}-24t^{2}$. \checkmark

  \item $W_{P}=Q$, $W_{Q}=P$.
        \[
          \mathrm{d}W=Q\,\mathrm{d}P+P\,\mathrm{d}Q
            =W\!\left(\frac{\mathrm{d}P}{P}+\frac{\mathrm{d}Q}{Q}\right) .
        \]
        So $\mathrm{d}W/W=0.01+(-0.005)=0.005$, i.e.\ revenue rises by
        about $0.5\%$. (Log-linear: revenue percent change is the sum of
        price and quantity percent changes.)

  \item At $P$: $x^{2}+yz=9+16=25$, so $Q(P)=4\cdot5=20$. Partials, by
        the chain rule with $w=x^{2}+yz$:
        \[
          Q_{x}=\frac{4\cdot2x}{2\sqrt w}=\frac{4x}{\sqrt w}=\frac{12}{5},\quad
          Q_{y}=\frac{2z}{\sqrt w}=\frac{16}{5},\quad
          Q_{z}=\frac{2y}{\sqrt w}=\frac45 .
        \]
        \[
          \mathrm{d}Q=2.4(0.1)+3.2(-0.2)+0.8(0.3)=0.24-0.64+0.24=-0.16 .
        \]
        Exact: $Q(3.1,1.8,8.3)=4\sqrt{9.61+14.94}=4\sqrt{24.55}
        \approx19.819183$, so $\Delta Q\approx-0.180817$. The differential
        gets the sign and the size about right; the error $0.021$ comes
        from the curvature of the root.
  \item At $P$: $xy+z^{2}=16+9=25$, $Y(P)=10$. With $w=xy+z^{2}$:
        \[
          Y_{x}=\frac{y}{\sqrt w}=\frac85,\quad Y_{y}=\frac{x}{\sqrt w}=\frac25,
          \quad Y_{z}=\frac{2z}{\sqrt w}=\frac65 .
        \]
        \[
          \mathrm{d}Y=1.6(0.05)+0.4(0.1)+1.2(-0.05)=0.08+0.04-0.06=0.06 .
        \]
        Exact: $Y(2.05,8.1,2.95)=2\sqrt{16.605+8.7025}=2\sqrt{25.3075}
        \approx10.061312$, so $\Delta Y\approx0.061312$: the error is about
        $0.0013$, two percent of the change.
  \item $f(P)=\ln2-2\approx-1.306853$. $f=2\ln x+\ln y-xy$, so
        $f_{x}=\dfrac2x-y=0$ and $f_{y}=\dfrac1y-x=-\tfrac12$ at $P$.
        \[
          \mathrm{d}f=0\cdot0.1-\tfrac12\cdot0.1=-0.05 .
        \]
        Exact: $f(1.1,2.1)=\ln(1.21\cdot2.1)-2.31=\ln2.541-2.31\approx
        -1.377442$, so $\Delta f\approx-0.070589$. The differential misses
        by $0.02$, about $30\%$ of the change. At $P$ the slope in the
        $x$-direction is $0$, so the linear term sees nothing of the
        $x$-step, while the curvature, with $f_{xx}=-2/x^{2}=-2$,
        $f_{xy}=-1$ and $f_{yy}=-1/y^{2}=-\tfrac14$, contributes about
        $\tfrac12\bigl(-2(0.01)+2(-1)(0.01)-\tfrac14(0.01)\bigr)\approx-0.021$:
        indeed $-0.05-0.021\approx-0.071$.
\end{solutions}

\chapter{Directional Derivatives and Vector-Valued Functions}

The partial derivatives $f_{x}$ and $f_{y}$ measure the response of $f$
when one input moves and the others stand still. Real adjustments are
rarely so tidy: a firm re-balances capital and labour together, a consumer
slides along a budget line, a hiker walks north-east rather than due
north. The \textbf{directional derivative} measures the rate of change of
$f$ per unit distance travelled in any prescribed direction, and a single
dot product with the gradient computes it. The gradient thereby earns its
geometric role: it points where $f$ climbs fastest, its length is that
fastest rate, and it stands perpendicular to the level curves, to the
indifference curves and isoquants of the previous chapters.

The second half of the chapter lets the function itself be vector-valued.
A map $\mathbf{f}:\R^{n}\to\R^{m}$ (a production plan sending inputs to
(output, cost), a coordinate change, a demand system) is differentiated
by collecting the partial derivatives of all components into the
\textbf{Jacobian matrix}, the correct ``derivative as linear
approximation'' for vector-valued functions. In this language the chain
rule becomes matrix multiplication. Determinants, invertibility and the
Implicit Function Theorem built on top of this matrix are the subject of
the next chapter.

\section{Core Rules}

\begin{rulebox}{Directional Derivative}
Let $f:\R^{n}\to\R$ be differentiable at $\mathbf{x}_{0}$ and let
$\mathbf{u}\in\R^{n}$ be a \textbf{unit} vector, $|\mathbf{u}|=1$. The
directional derivative of $f$ at $\mathbf{x}_{0}$ in the direction
$\mathbf{u}$ is
\[
  D_{\mathbf{u}}f(\mathbf{x}_{0})
  \;=\; \lim_{h\to 0}\frac{f(\mathbf{x}_{0}+h\mathbf{u})
        -f(\mathbf{x}_{0})}{h}
  \;=\; \grad f(\mathbf{x}_{0})\cdot\mathbf{u} .
\]
Two variables: $D_{\mathbf{u}}f=f_{x}u_{1}+f_{y}u_{2}$. The partials are
the special cases $\mathbf{u}=(1,0)$ and $\mathbf{u}=(0,1)$. The unit
length is what makes $D_{\mathbf{u}}f$ a rate \emph{per unit distance}.
\end{rulebox}

\begin{rulebox}{Normalising the Direction}
If a direction is given by a nonzero vector $\mathbf{v}$ that is not of
unit length, normalise first:
\[
  \mathbf{u} \;=\; \frac{\mathbf{v}}{|\mathbf{v}|}
  \;=\; \frac{\mathbf{v}}{\sqrt{v_{1}^{2}+\cdots+v_{n}^{2}}} .
\]
``The direction from $P$ to $Q$'' means $\mathbf{v}=Q-P$, then normalise.
Forgetting this step is a frequent error in directional-derivative
problems.
\end{rulebox}

\begin{rulebox}{Steepest Ascent and Descent}
Since $D_{\mathbf{u}}f=\grad f\cdot\mathbf{u}=|\grad f|\cos\theta$, where
$\theta$ is the angle between $\mathbf{u}$ and $\grad f$, at any point
with $\grad f(\mathbf{x}_{0})\neq\mathbf{0}$:
\begin{itemize}[leftmargin=1.6em,itemsep=0.3ex]
  \item $D_{\mathbf{u}}f$ is \emph{maximal} for
        $\mathbf{u}=\grad f/|\grad f|$, with maximum value $|\grad f|$
        (steepest ascent);
  \item $D_{\mathbf{u}}f$ is \emph{minimal} for
        $\mathbf{u}=-\grad f/|\grad f|$, with value $-|\grad f|$
        (steepest descent);
  \item $D_{\mathbf{u}}f=0$ exactly when $\mathbf{u}\perp\grad f$.
\end{itemize}
For output $Q(K,L)$, the gradient $(Q_{K},Q_{L})$ is the input adjustment
that raises output fastest per unit of adjustment, and $|\grad Q|$ is that
maximal marginal gain.

\textbf{A fixed amount of adjustment.} If the change vector
$(\mathrm{d}K,\mathrm{d}L)$ must have a given length $I$, output rises
fastest when it points along the gradient:
\[
  \mathrm{d}K:\mathrm{d}L=Q_{K}:Q_{L},\qquad
  (\mathrm{d}K,\mathrm{d}L)=\frac{I}{|\grad Q|}\,(Q_{K},Q_{L}),\qquad
  \mathrm{d}Q\approx I\,|\grad Q| .
\]
A \emph{budget} $\mathrm{d}K+\mathrm{d}L=I$ with $\mathrm{d}K,\mathrm{d}L\ge0$ is
a different constraint: the linear estimate $Q_{K}\mathrm{d}K+Q_{L}\mathrm{d}L$ is then largest at a
corner, with everything spent on the input with the larger marginal
product.
\end{rulebox}

\begin{rulebox}{Gradient and Level Curves}
Along the level set $f=c$ the function does not change, so every tangent
direction $\mathbf{u}$ of the level set satisfies
$D_{\mathbf{u}}f=\grad f\cdot\mathbf{u}=0$:
\[
  \grad f \;\perp\; \text{level curve (surface) through the point.}
\]
In two variables (with $f_{y}\neq0$) the tangent direction is proportional
to $(f_{y},-f_{x})$, consistent with the implicit slope
$\dd{y}{x}=-f_{x}/f_{y}$ of the previous chapters. Economically: the
marginal-utility vector $(U_{x},U_{y})$ is normal to the indifference
curve, and $(Q_{K},Q_{L})$ is normal to the isoquant.
\end{rulebox}

\begin{rulebox}{Tangent lines: to a level curve, and to the graph}
\textbf{Tangent to the level curve} $f(x,y)=c$ through
$\mathbf{p}=(x_{0},y_{0})$, where $\grad f(\mathbf{p})\neq\mathbf{0}$. The
gradient is a normal vector, so the tangent line is
\[
  \grad f(\mathbf{p})\cdot(\mathbf{x}-\mathbf{p})=0,
  \qquad\text{i.e.}\qquad
  f_{x}(\mathbf{p})(x-x_{0})+f_{y}(\mathbf{p})(y-y_{0})=0
\]
(normal form), or $\mathbf{x}=\mathbf{p}+t\mathbf{h}$ with
$\mathbf{h}\perp\grad f(\mathbf{p})$, for instance
$\mathbf{h}=(f_{y},-f_{x})$ (parametric form). Both describe the same
line in the plane.

\textbf{Tangent to the graph in a direction.} The vertical plane through
$P=(x_{0},y_{0})$ along the unit vector $\mathbf{u}$ cuts the surface
$z=f(x,y)$ in the curve $t\mapsto f(P+t\mathbf{u})$, whose slope at $t=0$
is $D_{\mathbf{u}}f(P)$. That is the geometric meaning of the directional
derivative: the slope of the surface in the direction $\mathbf{u}$. The
tangent line to the surface in that direction is a line in $\R^{3}$:
\[
  (x,y,z)=\bigl(x_{0},\,y_{0},\,f(P)\bigr)
          +t\,\bigl(u_{1},\,u_{2},\,D_{\mathbf{u}}f(P)\bigr),\qquad t\in\R .
\]
\begin{itemize}[leftmargin=1.6em,itemsep=0.3ex]
  \item $x$-direction ($\mathbf{u}=\mathbf{e}_{1}$):
        $\bigl(x_{0}+t,\;y_{0},\;f(P)+t\,f_{x}(P)\bigr)$;
        $y$-direction ($\mathbf{u}=\mathbf{e}_{2}$):
        $\bigl(x_{0},\;y_{0}+t,\;f(P)+t\,f_{y}(P)\bigr)$.
  \item For a direction vector $\mathbf{h}$ that is not normalised,
        $\bigl(P+t\mathbf{h},\;f(P)+t\,\grad f(P)\cdot\mathbf{h}\bigr)$ is
        the \emph{same} line, run through at a different speed. Only the
        slope, rise per unit of horizontal distance, needs the unit vector.
  \item All these lines lie in the tangent plane of \chref{derivatives}{10}.
\end{itemize}
\end{rulebox}

\begin{rulebox}{Second Directional Derivative}
For a $C^{2}$ function $f$, differentiating once more in the same unit
direction $\mathbf{u}$ gives
\[
  D^{2}_{\mathbf{u}}f
  \;=\; \mathbf{u}^{T}\Hess_{f}\,\mathbf{u}
  \;=\; \sum_{i,j} f_{x_{i}x_{j}}\,u_{i}u_{j} ;
  \qquad\text{two variables:}\quad
  D^{2}_{\mathbf{u}}f
  \;=\; f_{xx}u_{1}^{2}+2f_{xy}u_{1}u_{2}+f_{yy}u_{2}^{2} .
\]
It is the second derivative at $t=0$ of the profile
$t\mapsto f(\mathbf{x}_{0}+t\mathbf{u})$: the curvature of $f$ along the
line through $\mathbf{x}_{0}$ in direction $\mathbf{u}$. Definiteness of
the Hessian (\chref{derivatives}{11}) is exactly the statement that
$D^{2}_{\mathbf{u}}f$ has one sign for \emph{every} direction
$\mathbf{u}$.
\end{rulebox}

\begin{rulebox}{Curvature in a direction: computing it, and what it proves}
\textbf{The profile method.} Without the Hessian: set
$\varphi(t)=f(\mathbf{x}_{0}+t\mathbf{u})$, a function of one variable,
and differentiate twice. Then $\varphi'(0)=D_{\mathbf{u}}f(\mathbf{x}_{0})$
and $\varphi''(0)=D^{2}_{\mathbf{u}}f(\mathbf{x}_{0})$. Doing it both ways
is a built-in check.

\textbf{Non-unit directions.} For any $\mathbf{h}\neq\mathbf{0}$, with
$\mathbf{u}=\mathbf{h}/|\mathbf{h}|$,
\[
  \mathbf{h}^{T}H\mathbf{h}=|\mathbf{h}|^{2}\,D^{2}_{\mathbf{u}}f,
  \qquad
  \grad f\cdot\mathbf{h}=|\mathbf{h}|\,D_{\mathbf{u}}f .
\]
For the \emph{sign} of the curvature, normalising is unnecessary; for the
value of $D_{\mathbf{u}}f$ or $D^{2}_{\mathbf{u}}f$, divide by $|\mathbf{h}|$
or $|\mathbf{h}|^{2}$.

\textbf{Directional verdicts.} At $\mathbf{x}_{0}$, $f$ is
\emph{convex in the direction $\mathbf{u}$} if $D^{2}_{\mathbf{u}}f>0$,
\emph{concave in that direction} if $D^{2}_{\mathbf{u}}f<0$; the profile
$\varphi$ bends up or down there. A value $0$ gives no verdict.

\textbf{Refuting global convexity.} A convex $f$ has a PSD Hessian at
\emph{every} point, so $D^{2}_{\mathbf{u}}f\ge0$ everywhere in every
direction; a concave $f$ has $D^{2}_{\mathbf{u}}f\le0$ everywhere. Hence one
point and direction with $D^{2}_{\mathbf{u}}f<0$ shows that $f$ is not
convex, and one with $D^{2}_{\mathbf{v}}f>0$ (for example at a strict local
minimum with PD Hessian) shows that $f$ is not concave. Both together prove
that $f$ is \textbf{neither}.

\textbf{Where this goes next.} The box ``Curvature at a point and in a
direction'' in \chref{optimization}{3} turns these directional verdicts into
statements about convexity near a point and on a whole domain, including
the case of a semidefinite, singular Hessian where no local conclusion is
possible.
\end{rulebox}

\begin{rulebox}{Vector-Valued Functions and the Jacobian Matrix}
For $\mathbf{f}:\R^{n}\to\R^{m}$ with components
$\mathbf{f}=(f_{1},\ldots,f_{m})$, the \textbf{Jacobian matrix} at
$\mathbf{x}$ is the $m\times n$ matrix
\[
  J_{\mathbf{f}}(\mathbf{x}) \;=\;
  \begin{pmatrix}
    \partial f_{1}/\partial x_{1} & \cdots & \partial f_{1}/\partial x_{n}\\
    \vdots & \ddots & \vdots\\
    \partial f_{m}/\partial x_{1} & \cdots & \partial f_{m}/\partial x_{n}
  \end{pmatrix}:
  \quad\text{row $i$ is the gradient of $f_{i}$.}
\]
It is the best linear approximation of $\mathbf{f}$ near $\mathbf{x}_{0}$:
\[
  \mathbf{f}(\mathbf{x}_{0}+\mathbf{h})
  \;\approx\; \mathbf{f}(\mathbf{x}_{0})
  +J_{\mathbf{f}}(\mathbf{x}_{0})\,\mathbf{h} ,
\]
the matrix version of $f(x+h)\approx f(x)+f'(x)h$. Special cases: $m=1$
gives the gradient as a $1\times n$ row; $n=1$ gives the velocity
(tangent) column vector $\mathbf{r}'(t)$ of a curve.
\end{rulebox}

\begin{rulebox}{Differentiability of vector-valued maps}
$\mathbf{f}\colon D\to\R^{m}$, $D\subseteq\R^{n}$ open, is
\textbf{differentiable} at $\mathbf{x}_{0}$ if there is an $m\times n$
matrix $J$ with
\[
  \mathbf{f}(\mathbf{x}_{0}+\mathbf{h})=\mathbf{f}(\mathbf{x}_{0})
  +J\mathbf{h}+\mathbf{r}(\mathbf{h}),
  \qquad
  \frac{\|\mathbf{r}(\mathbf{h})\|}{\|\mathbf{h}\|}\to0\ \text{ as }
  \mathbf{h}\to\mathbf{0} .
\]
The ``$\approx$'' of the Jacobian box is made exact by the remainder.
\begin{itemize}[leftmargin=1.6em,itemsep=0.3ex]
  \item The matrix $J$ is then unique and equals the Jacobian
        $J_{\mathbf{f}}(\mathbf{x}_{0})$ of partial derivatives. It is
        called the \textbf{derivative} (or total derivative,
        \de{Ableitung, Jacobi-Matrix}) of $\mathbf{f}$ at $\mathbf{x}_{0}$.
  \item $\mathbf{f}$ is differentiable exactly when every component
        $f_{i}$ is: the remainder vector tends to $0$ faster than
        $\|\mathbf{h}\|$ exactly when each of its entries does.
  \item Existence of all partials is \emph{necessary} but not sufficient
        (the function $xy/(x^{2}+y^{2})$ of \chref{derivatives}{10} is a single component that fails).
        Continuous partials on an open set are sufficient: a $C^{1}$ map is
        differentiable.
  \item Geometry: for $n=m=1$ the linear part is the tangent line; for
        $n=2$, $m=1$ it is the tangent plane; for a curve ($n=1$) it is the
        tangent vector $\mathbf{r}'(t)$.
\end{itemize}
\end{rulebox}

\begin{rulebox}{The shape of a Jacobian, and which compositions exist}
$J_{\mathbf{f}}$ has one \textbf{row per output} and one \textbf{column per
input}: (number of outputs)$\times$(number of inputs).
\begin{itemize}[leftmargin=1.6em,itemsep=0.3ex]
  \item A curve $\R\to\R^{m}$ gives an $m\times1$ column; a scalar function
        $\R^{n}\to\R$ a $1\times n$ row; $\R^{3}\to\R^{2}$ a $2\times3$
        (wide) matrix; $\R^{2}\to\R^{4}$ a $4\times2$ (tall) one.
  \item \textbf{Linear and affine maps} have constant Jacobians:
        $J(A\mathbf{x}+\mathbf{b})=A$ for an $m\times n$ matrix $A$;
        $J(c\,\mathbf{x}-d\,\mathbf{a})=c\,I_{n}$; the identity map has
        $J=I_{n}$; $J(\mathbf{a}^{T}\mathbf{x})=\mathbf{a}^{T}$.
  \item \textbf{Patterns in $n$ dimensions.} If component $k$ depends on
        $x_{k}$ alone, $f_{k}=g_{k}(x_{k})$, the Jacobian is diagonal,
        $\operatorname{diag}\bigl(g_{1}'(x_{1}),\ldots,g_{n}'(x_{n})\bigr)$.
\end{itemize}
\textbf{Which composition exists.} $\mathbf{g}\circ\mathbf{f}$ means
``first $\mathbf{f}$, then $\mathbf{g}$''; it is defined when the number of
outputs of $\mathbf{f}$ equals the number of inputs of $\mathbf{g}$. With
$\mathbf{f}\colon\R^{n}\to\R^{m}$ and $\mathbf{g}\colon\R^{m}\to\R^{p}$ the
chain rule multiplies $(p\times m)(m\times n)$; the inner sizes must
match, which is the same condition. Check it before computing anything:
often only one of $\mathbf{g}\circ\mathbf{f}$ and $\mathbf{f}\circ\mathbf{g}$
makes sense.
\end{rulebox}

\begin{rulebox}{Chain Rule in Matrix Form}
If $\mathbf{f}:\R^{n}\to\R^{m}$ and $\mathbf{g}:\R^{m}\to\R^{p}$ are
differentiable, then $\mathbf{h}=\mathbf{g}\circ\mathbf{f}$ is
differentiable and
\[
  J_{\mathbf{h}}(\mathbf{x})
  \;=\; J_{\mathbf{g}}\bigl(\mathbf{f}(\mathbf{x})\bigr)\,
        J_{\mathbf{f}}(\mathbf{x}),
  \qquad
  (p\times n) \;=\; (p\times m)(m\times n) .
\]
Evaluate the \emph{outer} Jacobian at the \emph{inner value}
$\mathbf{f}(\mathbf{x})$; the order of the factors is fixed by the matrix
shapes. Every chain rule of \chref{derivatives}{12} is the case $p=1$: for
$z=f(\mathbf{r}(t))$,
\[
  \dd{z}{t} \;=\; \grad f\bigl(\mathbf{r}(t)\bigr)\cdot\mathbf{r}'(t),
\]
a $1\times m$ row times an $m\times 1$ column.
\end{rulebox}

\section{Worked Example}

\begin{worked}{For $f(x,y)=x^{2}+3xy-y^{2}$ at $P=(1,2)$: the derivative in the direction $(3,4)$, the steepest ascent, a direction of no change; then the Jacobian of $F(x,y)=(x^{2}-y,\;xy)$}
\wstep{Step 1. The gradient first --- every question below is an inner product with it.}
\[
  f_{x}=2x+3y,\qquad f_{y}=3x-2y,
\]
\[
  \grad f(1,2)=\bigl(2+6,\;3-4\bigr)=(8,-1).
\]

\wstep{Step 2. Normalise the direction, which the formula requires.}
\[
  \|(3,4)\|=5,
  \qquad
  \mathbf u=\Bigl(\tfrac35,\tfrac45\Bigr).
\]
\[
  D_{\mathbf u}f(P)=\grad f\cdot\mathbf u=8\cdot\tfrac35+(-1)\cdot\tfrac45
  =\frac{24-4}{5}=\boxed{4}.
\]
\wtrap{Without normalising you would get $\grad f\cdot(3,4)=20$, five times
too large. The directional derivative is a rate \emph{per unit distance
travelled}, so doubling the length of the direction vector must not double
the answer. If the problem hands you a non-unit vector --- and it usually
does --- divide by its length first.}

\wstep{Step 3. Steepest ascent and descent.}
The steepest increase is along the gradient itself:
\[
  \text{direction }\frac{(8,-1)}{\sqrt{65}},
  \qquad
  \text{rate }\|\grad f\|=\sqrt{65}\approx8.06 .
\]
The steepest \emph{decrease} is the opposite direction, at rate
$-\sqrt{65}$. Every other direction lies strictly between, since
\[
  D_{\mathbf u}f=\grad f\cdot\mathbf u=\|\grad f\|\cos\theta ,
\]
and $\cos\theta\in[-1,1]$. Our $4$ sits comfortably inside
$[-8.06,\,8.06]$. \checkmark
\wtrap{The value $8.06$ is the \emph{rate}; the direction is the normalised
gradient $(0.992,-0.124)$, not $(8,-1)$. Reporting the unnormalised gradient
as ``the direction'' conflates a direction with a magnitude.}

\wstep{Step 4. A direction of no change --- and what it is tangent to.}
$D_{\mathbf u}f=0$ requires $\mathbf u\perp\grad f$. Rotating $(8,-1)$ by a
right angle gives $(1,8)$, so
\[
  \mathbf u=\frac{(1,8)}{\sqrt{65}}
  \quad\Longrightarrow\quad
  D_{\mathbf u}f=\frac{8\cdot1+(-1)\cdot8}{\sqrt{65}}=0 . \checkmark
\]
Moving in that direction keeps $f$ constant to first order, so $(1,8)$ is
tangent to the level curve of $f$ through $P$ --- which is the geometric
content of ``the gradient is perpendicular to the level curve''.

\wstep{Step 5. A vector-valued function needs a matrix of derivatives.}
$F(x,y)=(x^{2}-y,\ xy)$ has two components, each with two partials. Collect
them with \textbf{one row per component function}:
\[
  J_{F}(x,y)=
  \begin{pmatrix}
    \pd{F_{1}}{x} & \pd{F_{1}}{y}\\[0.6ex]
    \pd{F_{2}}{x} & \pd{F_{2}}{y}
  \end{pmatrix}
  =\begin{pmatrix}2x&-1\\ y&x\end{pmatrix}.
\]
\wtrap{Rows are components, columns are variables. Transposing this is the
standard error, and it matters as soon as the Jacobian is multiplied by
anything. A check: for a scalar function ($m=1$) the Jacobian is a single
\emph{row}, and the gradient is conventionally a column --- so
$J=\grad^{\!\top}$.}

\wstep{Step 6. Evaluate, and read the determinant.}
\[
  J_{F}(1,2)=\begin{pmatrix}2&-1\\2&1\end{pmatrix},
  \qquad
  \det J_{F}(1,2)=2+2=4\neq0 .
\]
A non-zero Jacobian determinant says the map is locally invertible near
$(1,2)$: the system
\[
  x^{2}-y=u,\qquad xy=v
\]
can be solved for $(x,y)$ in terms of $(u,v)$ near that point. That is the
Implicit Function Theorem of \chref{derivatives}{14} in its multivariate form, and
$\det J$ is the condition it turns on.

\wstep{Step 7. Why the Jacobian is the right generalisation.}
\begin{center}\footnotesize
\begin{tabular}{lll}
  \hline
  $f:\R\to\R$ & derivative $f'$ & a number\\
  $f:\R^{n}\to\R$ & gradient $\grad f$ & a vector\\
  $F:\R^{n}\to\R^{m}$ & Jacobian $J_{F}$ & an $m\times n$ matrix\\
  \hline
\end{tabular}
\end{center}
In each row the object is the best \emph{linear} approximation at a point,
and the chain rule is matrix multiplication:
$J_{G\circ F}=J_{G}\cdot J_{F}$, with the order fixed by the order of the
composition. The one-variable chain rule $g'(f(x))f'(x)$ is the
$1\times1$ case.
\end{worked}

\section{Practice Problems}

\begin{problems}
  \item Find $D_{\mathbf{u}}f(1,1)$ for $f(x,y)=x^{2}y$ in the direction
        $\mathbf{u}=(1/\sqrt{2},\,1/\sqrt{2})$.
  \item For $f(x,y)=\ln(x^{2}+y^{2})$ at $(1,1)$, find the directional
        derivative toward $\mathbf{v}=(2,1)$.
  \item For $f(x,y)=x^{3}-3xy^{2}$ at $(2,1)$, find the directional
        derivative in the direction of the vector from $(2,1)$ to
        $(5,5)$.
  \item For $f(x,y,z)=xyz$ at $(1,2,3)$, find $D_{\mathbf{u}}f$ where
        $\mathbf{u}=\tfrac{1}{3}(1,2,2)$.
  \item For $f(x,y)=\mathrm{e}^{x+y}$ at $(0,0)$, find $D_{\mathbf{u}}f$
        where $\mathbf{u}$ makes a $60^{\circ}$ angle with the positive
        $x$-axis.
  \item For $f(x,y)=\sqrt{x^{2}+y^{2}}$ at $(3,4)$, find the directional
        derivative in the direction $(1,0)$.
  \item A consumer with utility $U(x,y)=xy$ currently consumes
        $(x,y)=(2,3)$. Her basket shifts in the direction
        $\mathbf{v}=(3,-4)$, more $x$, less $y$. Find the rate of
        change of utility per unit distance moved and interpret the sign.
  \item For $f(x,y)=x^{2}+xy+y^{2}$ at $(1,2)$, find the direction of
        steepest ascent and the maximum rate of increase.
  \item For $f(x,y)=x^{2}+2y^{2}$ at $(1,1)$, find the direction in which
        $f$ decreases fastest and the corresponding rate.
  \item Find all unit vectors $\mathbf{u}$ along which
        $f(x,y)=x^{2}-y^{2}$ has zero rate of change at $(2,1)$.
  \item For $f(x,y,z)=x^{2}+y^{2}-z^{2}$ at $(1,2,2)$, determine the
        maximum rate of increase and the direction in which it is
        attained.
  \item A hiker stands at $(1,1)$ on the hill $z=4-x^{2}-y^{2}$. In which
        direction should she walk to descend most quickly, and how steep
        is that descent?
  \item Output is $Q(K,L)=K^{1/2}L^{1/2}$ with current input mix
        $(K,L)=(4,9)$. Find the rate of change of output in the balanced
        expansion direction $\mathbf{u}=(1,1)/\sqrt{2}$; then find the
        input direction that raises output fastest and the corresponding
        maximal rate.
  \item Show that at every point of a level curve $f(x,y)=c$ the gradient
        $\grad f$ is orthogonal to the tangent direction of the curve.
        Conclude that the marginal-utility vector $(U_{x},U_{y})$ is
        normal to the indifference curve.
  \item For $f(x,y)=x^{2}y$ at $(1,2)$ with $\mathbf{u}=(3/5,\,4/5)$,
        compute both $D_{\mathbf{u}}f$ and the second directional
        derivative
        $D^{2}_{\mathbf{u}}f=\mathbf{u}^{T}\Hess_{f}\,\mathbf{u}$.
  \item For $f(x,y)=x^{3}+y^{3}-3xy$: show that $(1,1)$ is a critical
        point, then compute $D^{2}_{\mathbf{u}}f(1,1)$ for
        $\mathbf{u}=(1,1)/\sqrt{2}$ and for
        $\mathbf{v}=(1,-1)/\sqrt{2}$. What do the signs suggest?
  \item For $f(x,y)=\mathrm{e}^{x}\cos y$ at $(0,0)$, find all unit
        vectors $\mathbf{u}$ with $D^{2}_{\mathbf{u}}f=0$.
  \item Compute the Jacobian of $\mathbf{f}(x,y)=(x+y,\;x-y)$ and its
        determinant. What is special about the Jacobian of a linear map?
  \item Compute the Jacobian of
        $\mathbf{f}(x,y)=(x^{2}y,\;x+y^{2},\;xy)$, state its shape, and
        evaluate it at $(1,2)$.
  \item Compute the Jacobian of
        $\mathbf{f}(x,y,z)=(x+y,\;y+z,\;x+z)$ and its determinant.
  \item For $\mathbf{f}(x,y)=(x^{2}-y,\;xy)$: compute $J_{\mathbf{f}}$,
        then use the linear approximation
        $\mathbf{f}(\mathbf{x}_{0}+\mathbf{h})\approx
        \mathbf{f}(\mathbf{x}_{0})+J_{\mathbf{f}}(\mathbf{x}_{0})\mathbf{h}$
        at $\mathbf{x}_{0}=(1,1)$ to estimate $\mathbf{f}(1.1,\,0.9)$,
        and compare with the exact value.
  \item Demand and supply are $D(p,q)=100-p+0.5q$ and
        $S(p,q)=-10+2p+q$, and the two equilibrium conditions are
        $G_{1}=D-S=0$ and $G_{2}=p+q-20=0$. Build the Jacobian matrix
        $\partial(G_{1},G_{2})/\partial(p,q)$ and compute its
        determinant.
  \item Let $\mathbf{g}(x,y)=(x+y,\;x-y)$ and
        $\mathbf{F}(u,v)=(u^{2},\;uv)$. Compute the Jacobian of
        $\mathbf{F}\circ\mathbf{g}$ by the matrix chain rule
        $J_{\mathbf{F}\circ\mathbf{g}}
        =J_{\mathbf{F}}(\mathbf{g}(x,y))\,J_{\mathbf{g}}(x,y)$, and
        verify by composing first.
  \item For $f(x,y)=xy^{2}$ and the curve
        $\mathbf{r}(t)=(\mathrm{e}^{t},\,t)$, compute $\dd{z}{t}$ for
        $z=f(\mathbf{r}(t))$ as the matrix product
        $J_{f}(\mathbf{r}(t))\,J_{\mathbf{r}}(t)$ (a $1\times2$ row
        times a $2\times1$ column) and verify directly.
  \item A firm converts inputs into (output, cost) via
        $\mathbf{g}(K,L)=\bigl(K^{1/2}L^{1/2},\;2K+L\bigr)$ and earns
        profit $\pi(Q,C)=10Q-C$. Use the matrix chain rule to find the
        marginal profits $(\pi_{K},\pi_{L})$ at $(K,L)=(1,1)$, and verify
        by substituting first.
\end{problems}

\section{Solutions}

\begin{solutions}
  \item $\grad f=(2xy,\,x^{2})$. At $(1,1)$: $\grad f=(2,1)$, and
        $\mathbf{u}$ is already a unit vector.
        \[
          D_{\mathbf{u}}f=(2)\frac{1}{\sqrt{2}}+(1)\frac{1}{\sqrt{2}}
            =\frac{3}{\sqrt{2}}=\frac{3\sqrt{2}}{2} .
        \]

  \item $\grad f=\Bigl(\dfrac{2x}{x^{2}+y^{2}},\,
        \dfrac{2y}{x^{2}+y^{2}}\Bigr)$. At $(1,1)$: $\grad f=(1,1)$.
        Normalise: $|\mathbf{v}|=\sqrt{4+1}=\sqrt{5}$, so
        $\mathbf{u}=(2,1)/\sqrt{5}$.
        \[
          D_{\mathbf{u}}f=\frac{1\cdot 2+1\cdot 1}{\sqrt{5}}
            =\frac{3}{\sqrt{5}}=\frac{3\sqrt{5}}{5} .
        \]

  \item Direction vector $\mathbf{v}=(5,5)-(2,1)=(3,4)$,
        $|\mathbf{v}|=5$, so $\mathbf{u}=(3/5,\,4/5)$.
        $\grad f=(3x^{2}-3y^{2},\,-6xy)$; at $(2,1)$:
        $\grad f=(12-3,\,-12)=(9,-12)$.
        \[
          D_{\mathbf{u}}f=9\cdot\frac{3}{5}+(-12)\cdot\frac{4}{5}
            =\frac{27}{5}-\frac{48}{5}=-\frac{21}{5} .
        \]

  \item $\grad f=(yz,\,xz,\,xy)$. At $(1,2,3)$: $\grad f=(6,3,2)$. The
        vector $\mathbf{u}=\tfrac13(1,2,2)$ is unit since
        $\sqrt{1+4+4}/3=1$.
        \[
          D_{\mathbf{u}}f=\frac{6\cdot 1+3\cdot 2+2\cdot 2}{3}
            =\frac{6+6+4}{3}=\frac{16}{3} .
        \]

  \item $\mathbf{u}=(\cos 60^{\circ},\,\sin 60^{\circ})
        =(1/2,\,\sqrt{3}/2)$, a unit vector.
        $\grad f=(\mathrm{e}^{x+y},\,\mathrm{e}^{x+y})$; at $(0,0)$:
        $\grad f=(1,1)$.
        \[
          D_{\mathbf{u}}f=\frac{1}{2}+\frac{\sqrt{3}}{2}
            =\frac{1+\sqrt{3}}{2} .
        \]

  \item $\grad f=\Bigl(\dfrac{x}{\sqrt{x^{2}+y^{2}}},\,
        \dfrac{y}{\sqrt{x^{2}+y^{2}}}\Bigr)$. At $(3,4)$:
        $\grad f=(3/5,\,4/5)$. With $\mathbf{u}=(1,0)$ (already unit):
        \[
          D_{\mathbf{u}}f=\frac{3}{5} .
        \]

  \item $\grad U=(y,\,x)$; at $(2,3)$: $\grad U=(3,2)$. Normalise:
        $|\mathbf{v}|=\sqrt{9+16}=5$, $\mathbf{u}=(3/5,\,-4/5)$.
        \[
          D_{\mathbf{u}}U=3\cdot\frac{3}{5}+2\cdot\Bigl(-\frac{4}{5}\Bigr)
            =\frac{9}{5}-\frac{8}{5}=\frac{1}{5} .
        \]
        The rate is positive: the marginal-utility gain from the extra
        $x$ slightly outweighs the loss from the reduced $y$, so utility
        rises, at $1/5$ per unit distance moved.

  \item $\grad f=(2x+y,\,x+2y)$. At $(1,2)$: $\grad f=(4,5)$,
        $|\grad f|=\sqrt{16+25}=\sqrt{41}$.
        \[
          \mathbf{u}=\frac{(4,5)}{\sqrt{41}},\qquad
          \text{maximum rate}=\sqrt{41} .
        \]

  \item $\grad f=(2x,\,4y)$. At $(1,1)$: $\grad f=(2,4)$,
        $|\grad f|=\sqrt{4+16}=2\sqrt{5}$. Steepest descent:
        \[
          \mathbf{u}=-\frac{(2,4)}{2\sqrt{5}}=-\frac{(1,2)}{\sqrt{5}},
          \qquad \text{rate}=-2\sqrt{5} .
        \]

  \item $\grad f=(2x,\,-2y)$. At $(2,1)$: $\grad f=(4,-2)$. We need
        $\mathbf{u}\perp(4,-2)$: $4u_{1}-2u_{2}=0$, i.e.
        $u_{2}=2u_{1}$. Normalising $(1,2)$:
        \[
          \mathbf{u}=\pm\frac{(1,2)}{\sqrt{5}} .
        \]
        \textbf{Check.} $(4,-2)\cdot(1,2)=4-4=0$. \checkmark

  \item $\grad f=(2x,\,2y,\,-2z)$. At $(1,2,2)$: $\grad f=(2,4,-4)$,
        $|\grad f|=\sqrt{4+16+16}=6$. Maximum rate $6$, attained in the
        direction
        \[
          \mathbf{u}=\frac{(2,4,-4)}{6}=\frac{(1,2,-2)}{3} .
        \]

  \item $\grad z=(-2x,\,-2y)$; at $(1,1)$: $\grad z=(-2,-2)$,
        $|\grad z|=2\sqrt{2}$. Steepest descent direction:
        \[
          \mathbf{u}=-\frac{(-2,-2)}{2\sqrt{2}}=\frac{(1,1)}{\sqrt{2}} ,
        \]
        i.e. north-east, directly away from the summit above the origin.
        The slope of that descent is $-|\grad z|=-2\sqrt{2}$.

  \item $Q_{K}=\tfrac12K^{-1/2}L^{1/2}$, $Q_{L}=\tfrac12K^{1/2}L^{-1/2}$.
        At $(4,9)$:
        \[
          Q_{K}=\frac{1}{2}\cdot\frac{1}{2}\cdot 3=\frac{3}{4},\qquad
          Q_{L}=\frac{1}{2}\cdot 2\cdot\frac{1}{3}=\frac{1}{3},
        \]
        so $\grad Q=(3/4,\,1/3)$. Balanced direction
        $\mathbf{u}=(1,1)/\sqrt{2}$:
        \[
          D_{\mathbf{u}}Q=\frac{3/4+1/3}{\sqrt{2}}
            =\frac{13/12}{\sqrt{2}}=\frac{13\sqrt{2}}{24} .
        \]
        Fastest increase along $\grad Q/|\grad Q|$ with
        \[
          |\grad Q|=\sqrt{\frac{9}{16}+\frac{1}{9}}
            =\sqrt{\frac{81+16}{144}}=\frac{\sqrt{97}}{12},\qquad
          \mathbf{u}^{*}=\frac{12}{\sqrt{97}}\Bigl(\frac34,\frac13\Bigr)
            =\frac{(9,4)}{\sqrt{97}} .
        \]
        The maximal rate is $\sqrt{97}/12$ (larger than $13\sqrt{2}/24$,
        as it must be).

  \item Let $\mathbf{r}(t)=(x(t),y(t))$ be a differentiable curve lying
        in the level set, $f(\mathbf{r}(t))=c$ for all $t$.
        Differentiate with the chain rule:
        \[
          0=\dd{}{t}\,c=\grad f(\mathbf{r}(t))\cdot\mathbf{r}'(t) .
        \]
        Hence the tangent vector $\mathbf{r}'(t)$ is orthogonal to
        $\grad f$ at every point of the curve. For $U(x,y)=c$ an
        indifference curve, this says exactly that
        $(U_{x},U_{y})$ is normal to it.

  \item First order: $\grad f=(2xy,\,x^{2})$; at $(1,2)$:
        $\grad f=(4,1)$.
        \[
          D_{\mathbf{u}}f=4\cdot\frac{3}{5}+1\cdot\frac{4}{5}
            =\frac{12}{5}+\frac{4}{5}=\frac{16}{5} .
        \]
        Second order: $f_{xx}=2y=4$, $f_{xy}=2x=2$, $f_{yy}=0$ at
        $(1,2)$.
        \[
          D^{2}_{\mathbf{u}}f
            =4\Bigl(\frac{3}{5}\Bigr)^{2}
             +2\cdot 2\cdot\frac{3}{5}\cdot\frac{4}{5}
             +0\cdot\Bigl(\frac{4}{5}\Bigr)^{2}
            =\frac{36}{25}+\frac{48}{25}=\frac{84}{25} .
        \]

  \item $\grad f=(3x^{2}-3y,\,3y^{2}-3x)$; at $(1,1)$:
        $\grad f=(0,0)$, so $(1,1)$ is a critical point. Second
        partials: $f_{xx}=6x=6$, $f_{yy}=6y=6$, $f_{xy}=-3$, so
        $\Hess_{f}(1,1)=\begin{pmatrix}6&-3\\-3&6\end{pmatrix}$. For
        $\mathbf{u}=(1,1)/\sqrt{2}$:
        \[
          D^{2}_{\mathbf{u}}f=6\cdot\frac12
            +2(-3)\cdot\frac12+6\cdot\frac12=3-3+3=3 .
        \]
        For $\mathbf{v}=(1,-1)/\sqrt{2}$:
        \[
          D^{2}_{\mathbf{v}}f=6\cdot\frac12
            +2(-3)\cdot\Bigl(-\frac12\Bigr)+6\cdot\frac12=3+3+3=9 .
        \]
        Both are positive, the surface curves upward in both
        directions. Indeed $D_{1}=6>0$ and
        $D_{2}=36-9=27>0$, so the Hessian is positive definite and
        $(1,1)$ is a strict local minimum (\chref{derivatives}{11}).

  \item $f_{xx}=\mathrm{e}^{x}\cos y$, $f_{xy}=-\mathrm{e}^{x}\sin y$,
        $f_{yy}=-\mathrm{e}^{x}\cos y$. At $(0,0)$: $f_{xx}=1$,
        $f_{xy}=0$, $f_{yy}=-1$. Writing
        $\mathbf{u}=(\cos\theta,\,\sin\theta)$:
        \[
          D^{2}_{\mathbf{u}}f=\cos^{2}\theta-\sin^{2}\theta=\cos 2\theta .
        \]
        This vanishes for $2\theta=\tfrac{\pi}{2}+k\pi$, i.e.
        $\theta=\tfrac{\pi}{4}+k\tfrac{\pi}{2}$:
        \[
          \mathbf{u}=\pm\frac{(1,1)}{\sqrt{2}}
          \quad\text{and}\quad
          \mathbf{u}=\pm\frac{(1,-1)}{\sqrt{2}} .
        \]

  \item
        \[
          J_{\mathbf{f}}=\begin{pmatrix}1&1\\1&-1\end{pmatrix},\qquad
          \det J_{\mathbf{f}}=1\cdot(-1)-1\cdot 1=-2 .
        \]
        The map is linear, $\mathbf{f}(\mathbf{x})=A\mathbf{x}$ with
        $A=\bigl(\begin{smallmatrix}1&1\\1&-1\end{smallmatrix}\bigr)$,
        and for a linear map $f_{i}=\sum_{j}A_{ij}x_{j}$ gives
        $\partial f_{i}/\partial x_{j}=A_{ij}$: the Jacobian is the
        constant matrix $A$ itself: a linear map is its own linear
        approximation.

  \item With $f_{1}=x^{2}y$, $f_{2}=x+y^{2}$, $f_{3}=xy$:
        \[
          J_{\mathbf{f}}=\begin{pmatrix}2xy&x^{2}\\1&2y\\y&x\end{pmatrix}
          \;\in\;\R^{3\times 2} .
        \]
        A $3\times 2$ matrix: one row per component, one column per
        variable; being non-square, it has no determinant. At $(1,2)$:
        \[
          J_{\mathbf{f}}(1,2)=\begin{pmatrix}4&1\\1&4\\2&1\end{pmatrix} .
        \]

  \item
        \[
          J_{\mathbf{f}}=\begin{pmatrix}1&1&0\\0&1&1\\1&0&1\end{pmatrix} .
        \]
        Expanding along the first row:
        \[
          \det J_{\mathbf{f}}
            =1\,(1\cdot 1-1\cdot 0)-1\,(0\cdot 1-1\cdot 1)+0
            =1+1=2 .
        \]

  \item $J_{\mathbf{f}}=\begin{pmatrix}2x&-1\\y&x\end{pmatrix}$. At
        $\mathbf{x}_{0}=(1,1)$:
        \[
          \mathbf{f}(1,1)=(0,\,1),\qquad
          J_{\mathbf{f}}(1,1)=\begin{pmatrix}2&-1\\1&1\end{pmatrix} .
        \]
        With $\mathbf{h}=(0.1,\,-0.1)$:
        \[
          J_{\mathbf{f}}(1,1)\,\mathbf{h}
            =\begin{pmatrix}2(0.1)+(-1)(-0.1)\\
              1(0.1)+1(-0.1)\end{pmatrix}
            =\begin{pmatrix}0.3\\0\end{pmatrix},
        \]
        so $\mathbf{f}(1.1,\,0.9)\approx(0,1)+(0.3,\,0)=(0.3,\,1)$.
        Exact value:
        \[
          \mathbf{f}(1.1,\,0.9)=(1.21-0.9,\;1.1\cdot 0.9)
            =(0.31,\;0.99) .
        \]
        The approximation is off by only $0.01$ in each component.

  \item Excess demand:
        \[
          G_{1}=D-S=(100-p+0.5q)-(-10+2p+q)=110-3p-0.5q ,
        \]
        and $G_{2}=p+q-20$. Hence
        \[
          \frac{\partial(G_{1},G_{2})}{\partial(p,q)}
            =\begin{pmatrix}-3&-0.5\\1&1\end{pmatrix},\qquad
          \det=-3\cdot 1-(-0.5)\cdot 1=-2.5 .
        \]
        The determinant is nonzero: exactly the condition the next
        chapter's Implicit Function Theorem needs to solve the system
        locally for $(p,q)$ when a parameter shifts.

  \item Inner and outer Jacobians:
        \[
          J_{\mathbf{g}}=\begin{pmatrix}1&1\\1&-1\end{pmatrix},\qquad
          J_{\mathbf{F}}(u,v)=\begin{pmatrix}2u&0\\v&u\end{pmatrix} .
        \]
        Evaluate the outer at the inner value
        $(u,v)=\mathbf{g}(x,y)=(x+y,\,x-y)$ and multiply:
        \[
          J_{\mathbf{F}\circ\mathbf{g}}
            =\begin{pmatrix}2(x+y)&0\\x-y&x+y\end{pmatrix}
             \begin{pmatrix}1&1\\1&-1\end{pmatrix}
            =\begin{pmatrix}2(x+y)&2(x+y)\\
              (x-y)+(x+y)&(x-y)-(x+y)\end{pmatrix}
        \]
        \[
          \phantom{J_{\mathbf{F}\circ\mathbf{g}}}
            =\begin{pmatrix}2(x+y)&2(x+y)\\2x&-2y\end{pmatrix} .
        \]
        \textbf{Check.} Composing first:
        $(\mathbf{F}\circ\mathbf{g})(x,y)
        =\bigl((x+y)^{2},\,(x+y)(x-y)\bigr)
        =\bigl((x+y)^{2},\,x^{2}-y^{2}\bigr)$, whose Jacobian is
        \[
          \begin{pmatrix}2(x+y)&2(x+y)\\2x&-2y\end{pmatrix} .
          \quad\checkmark
        \]

  \item $J_{f}=\bigl(y^{2}\;\;2xy\bigr)$ (a row) and
        $J_{\mathbf{r}}(t)=\begin{pmatrix}\mathrm{e}^{t}\\1\end{pmatrix}$
        (a column). At $\mathbf{r}(t)=(\mathrm{e}^{t},\,t)$ the row is
        $\bigl(t^{2}\;\;2t\mathrm{e}^{t}\bigr)$, so
        \[
          \dd{z}{t}
            =\bigl(t^{2}\;\;2t\mathrm{e}^{t}\bigr)
             \begin{pmatrix}\mathrm{e}^{t}\\1\end{pmatrix}
            =t^{2}\mathrm{e}^{t}+2t\mathrm{e}^{t}
            =\mathrm{e}^{t}(t^{2}+2t) .
        \]
        \textbf{Check.} $z=\mathrm{e}^{t}\,t^{2}$, so
        $\dd{z}{t}=\mathrm{e}^{t}t^{2}+2t\mathrm{e}^{t}$. \checkmark

  \item Inner Jacobian:
        \[
          J_{\mathbf{g}}(K,L)=\begin{pmatrix}
            \tfrac12K^{-1/2}L^{1/2}&\tfrac12K^{1/2}L^{-1/2}\\
            2&1\end{pmatrix},\qquad
          J_{\mathbf{g}}(1,1)=\begin{pmatrix}1/2&1/2\\2&1\end{pmatrix} .
        \]
        Outer Jacobian (a row, constant since $\pi$ is linear):
        $J_{\pi}=(10\;\;-1)$. Matrix chain rule:
        \[
          (\pi_{K},\,\pi_{L})
            =(10\;\;-1)\begin{pmatrix}1/2&1/2\\2&1\end{pmatrix}
            =\bigl(10\cdot\tfrac12-1\cdot 2,\;
              10\cdot\tfrac12-1\cdot 1\bigr)
            =(3,\;4) .
        \]
        \textbf{Check.} Substituting first,
        $\pi(K,L)=10K^{1/2}L^{1/2}-2K-L$, so
        $\pi_{K}=5K^{-1/2}L^{1/2}-2=5-2=3$ and
        $\pi_{L}=5K^{1/2}L^{-1/2}-1=5-1=4$ at $(1,1)$. \checkmark
        At the current input mix an extra unit of capital raises profit
        by about $3$, an extra unit of labour by about $4$.
\end{solutions}

\section{Additional Problems}

\begin{problems}
  \item For $f(x,y)=x^{2}-y^{2}+3xy$ at $(1,2)$, find the directional
        derivative toward $\mathbf{v}=(4,3)$.
  \item For $f(x,y,z)=x^{2}+2y^{2}+3z^{2}$ at $(1,1,1)$, find the
        directional derivative toward $\mathbf{v}=(1,-1,1)$.
  \item For $U(x,y)=\ln x+2\ln y$ at $(1,2)$, find the direction of
        steepest ascent and the maximum rate of increase.
  \item Find all unit vectors $\mathbf{u}$ with $D_{\mathbf{u}}f=0$ for
        $f(x,y)=xy$ at $(2,1)$.
  \item At a point $P$ the directional derivatives of $f$ toward $(1,0)$
        and toward $(0,1)$ are $3$ and $-4$. Find $\grad f(P)$, the
        maximum rate of increase at $P$, and $D_{\mathbf{u}}f$ toward
        $\mathbf{u}=(3/5,\,4/5)$.
  \item For $f(x,y)=\ln(x^{2}+y^{2})$ at $(1,0)$ with
        $\mathbf{u}=(0,1)$, compute $D_{\mathbf{u}}f$ and
        $D^{2}_{\mathbf{u}}f$, and interpret the two results.
  \item Compute the Jacobian of
        $\mathbf{f}(x,y)=(x^{2}-y^{2},\;3xy+y^{2})$ and evaluate it at
        $(1,1)$.
  \item For $f(x,y)=xy$ and $\mathbf{r}(t)=(t^{2},\,t^{3})$, compute
        $\dd{z}{t}$ for $z=f(\mathbf{r}(t))$ by the matrix chain rule
        and verify directly.
  \item Let $\mathbf{g}(s,t)=(st,\;s+t)$ and
        $\mathbf{F}(x,y)=(x+y,\;xy)$. Compute
        $J_{\mathbf{F}\circ\mathbf{g}}$ by the matrix chain rule and
        verify by composing first.
  \item Inputs produce two intermediate goods $x_{1}=2K+L$ and
        $x_{2}=K+3L$, and final output is $y=x_{1}x_{2}$. Use the matrix
        chain rule to find $\pd{y}{K}$ and $\pd{y}{L}$ at $(K,L)=(1,1)$,
        and verify by expanding $y$ as a polynomial in $K$ and $L$.
\end{problems}

Problems 11--14: tangent lines to the graph of $f$ in a given direction.

\begin{problems}[resume]
  \item For $f(x,y)=x^{2}y+y^{3}$ at $P=(1,2)$, give the slope of the graph
        and the tangent line to the graph in the $x$-direction, in the
        $y$-direction, and in the direction of $\mathbf{v}=(3,-4)$.
  \item For $f(x,y)=\ln(x^{2}+3y)$ at $P=(2,0)$, give the slope of the
        graph in the direction of steepest ascent and the tangent line to
        the graph in that direction.
  \item For $f(x,y)=(x+y)\,\mathrm{e}^{x-y}$ at $P=(1,1)$, give the slope in
        the direction from $P$ towards $Q=(3,2)$ and the tangent line in
        that direction, once with the unit vector and once with
        $\mathbf{h}=Q-P$. Show that both describe the same line.
  \item For $Q(K,L)=4K^{1/2}L^{1/4}$ at $(16,16)$, give the tangent line to
        the graph in the $L$-direction, read it as an output estimate for
        $L=17$, and compare with the exact value.
\end{problems}

In problems 15--17, for the point $P$: (a) find the level $c=f(P)$; (b)
solve the level curve $f=c$ explicitly; (c) compute $\grad f(P)$ and check
that it is perpendicular to the tangent of the level curve; (d) give the
tangent line to the level curve in normal and in parametric form; (e) give
$D_{\mathbf{u}}f(P)$ in the direction of the gradient.

\begin{problems}[resume]
  \item $f(x,y)=\sqrt{2x^{2}+3y+1}$, \;$P=(1,2)$.
  \item $f(x,y)=x^{1/3}y^{2/3}$ on $x,y>0$, \;$P=(1,8)$.
  \item $f(x,y)=x\,\mathrm{e}^{y}+y$, \;$P=(2,0)$; solve the level curve
        for $x$.
\end{problems}

Problems 18--20: differentiability of vector-valued maps.

\begin{problems}[resume]
  \item Explain why
        $\mathbf{F}(x,y)=\bigl(x^{2}y,\;x-3y^{2},\;xy\bigr)$ is
        differentiable at every point. Give $J_{\mathbf{F}}(1,1)$, and
        compute the remainder vector $\mathbf{r}(\mathbf{h})$ at $(1,1)$
        to show directly that $\|\mathbf{r}(\mathbf{h})\|/\|\mathbf{h}\|\to0$.
  \item For $\mathbf{F}(x,y)=\bigl(\ln(3x^{2}+y^{4}),\;xy^{2},\;|y-1|\bigr)$,
        give the domain, the set on which $\mathbf{F}$ is differentiable,
        and the Jacobian there.
  \item Answer in words: when is a map $\mathbf{f}\colon\R^{n}\to\R^{m}$
        differentiable at $\mathbf{x}_{0}$; what is its derivative called
        and what shape does it have; how is it related to the partial
        derivatives of the components; and what is the geometric meaning
        for $n=m=1$ and for $n=2$, $m=1$?
\end{problems}

Problems 21--26: Jacobians of all shapes.

\begin{problems}[resume]
  \item $\mathbf{r}(t)=\bigl(\ln(1+t^{2}),\;\mathrm{e}^{-2t},\;\cos3t,\;5\bigr)$:
        the Jacobian, its shape, and its value at $t=0$.
  \item $\mathbf{F}(x,y,z)=(xyz,\;x^{2}-z)$ at $(1,2,3)$, and
        $g(x_{1},x_{2},x_{3},x_{4})=x_{1}x_{2}+x_{3}\mathrm{e}^{x_{4}}$ at
        $(1,1,1,0)$: shapes and values of the Jacobians.
  \item $\mathbf{F}(s,t)=\bigl(s+t,\;st,\;s^{2},\;\mathrm{e}^{s-t}\bigr)$:
        the Jacobian, its shape, and its value at $(1,1)$.
  \item Jacobians in $n$ dimensions: (a) the identity map
        $\mathbf{F}(\mathbf{x})=\mathbf{x}$ on $\R^{n}$; (b) the map
        $\mathbf{F}\colon\R^{n}\to\R^{n}$ with components
        $F_{k}(\mathbf{x})=k\,\mathrm{e}^{x_{k}}$, that is
        $(\mathrm{e}^{x_{1}},2\mathrm{e}^{x_{2}},\ldots,n\,\mathrm{e}^{x_{n}})$;
        (c) the curve $\mathbf{r}\colon\R\to\R^{n}$ with components
        $r_{k}(t)=t^{k}/k$, that is $(t,t^{2}/2,\ldots,t^{n}/n)$; (d)
        $f(\mathbf{x})=\sum_{k=1}^{n}x_{k}^{k}/k$.
  \item Jacobians of linear and affine maps: (a) $\mathbf{F}(\mathbf{x})=
        A\mathbf{x}$ with $A=\begin{pmatrix}2&-1&0\\1&0&3\\0&4&-2
        \end{pmatrix}$; (b) $\mathbf{F}(\mathbf{x})=3\mathbf{x}-2\mathbf{a}$
        for a fixed $\mathbf{a}\in\R^{n}$; (c) $f(\mathbf{x})=
        \mathbf{a}^{T}\mathbf{x}$; (d) $\mathbf{G}(\mathbf{x})=B\mathbf{x}
        +\mathbf{c}$ with $B$ a $2\times3$ matrix.
  \item Let $\mathbf{F}\colon\R^{3}\to\R^{2}$, $\mathbf{G}\colon\R^{2}\to
        \R^{4}$, $H\colon\R^{4}\to\R$ and $\mathbf{K}\colon\R\to\R^{3}$.
        Which of $\mathbf{G}\circ\mathbf{F}$, $\mathbf{F}\circ\mathbf{G}$,
        $H\circ\mathbf{G}$, $\mathbf{G}\circ H$, $\mathbf{K}\circ H$,
        $H\circ\mathbf{K}$, $\mathbf{F}\circ\mathbf{K}$,
        $\mathbf{K}\circ\mathbf{F}$ and $H\circ\mathbf{G}\circ\mathbf{F}$
        are defined, and what is the shape of the Jacobian of each one that
        is?
\end{problems}

Problems 27--30: the matrix chain rule with rectangular Jacobians. Verify
each result by composing first.

\begin{problems}[resume]
  \item $\mathbf{f}(x,y,z)=(x+yz,\;xz)$ and
        $\mathbf{g}(u,v)=(u,\;v,\;uv,\;u-v)$: the Jacobian of
        $\mathbf{g}\circ\mathbf{f}$ at $(1,2,1)$.
  \item $g(x,y,z)=x+2y+z^{2}$ and $\mathbf{f}(s)=(s,\;s^{2},\;\mathrm{e}^{s})$:
        say why both $\mathbf{f}\circ g$ and $g\circ\mathbf{f}$ exist; give
        the Jacobian of $\mathbf{f}\circ g$ at $(1,0,0)$ and the derivative
        of $g\circ\mathbf{f}$.
  \item $\mathbf{f}(x,y)=(x,\;y,\;xy)$ and $\mathbf{g}(u,v,w)=(u+w,\;vw)$:
        the Jacobian of $\mathbf{g}\circ\mathbf{f}$ at $(1,2)$ and of
        $\mathbf{f}\circ\mathbf{g}$ at $(1,1,1)$.
  \item $\mathbf{r}(t)=(\cos t,\;\sin t,\;t)$ and $F(x,y,z)=x^{2}+y^{2}+z^{2}$:
        $\dd{}{t}F(\mathbf{r}(t))$ by the chain rule, simplified.
\end{problems}

In problems 31--33, compute $D_{\mathbf{u}}f(P)$ and $D^{2}_{\mathbf{u}}f(P)$
for the direction of $\mathbf{h}$, both with the gradient and Hessian and
with the profile $\varphi(t)=f(P+t\mathbf{h})$, and say whether $f$ is
convex or concave in that direction at $P$.

\begin{problems}[resume]
  \item $f(x,y,z)=x\ln y+yz^{2}$, \;$P=(2,1,1)$, \;$\mathbf{h}=(1,2,-2)$.
  \item $f(x,y,z)=x\,\mathrm{e}^{yz}$, \;$P=(1,0,2)$, \;$\mathbf{h}=(2,1,2)$.
  \item $f(x,y,z)=xyz+\ln(xz)$, \;$P=(1,1,1)$, \;$\mathbf{h}=(1,-1,1)$.
  \item Redo practice problem 15 ($f=x^{2}y$ at $(1,2)$) with the
        \emph{unnormalised} $\mathbf{h}=(3,4)$: compute
        $\grad f\cdot\mathbf{h}$, $\mathbf{h}^{T}H\mathbf{h}$, and the
        profile derivatives $\varphi'(0)$, $\varphi''(0)$. Explain the
        factors $5$ and $25$.
\end{problems}

In problems 35--37, give the directional verdicts asked for, then decide
whether $f$ is convex, concave or neither on its whole domain.

\begin{problems}[resume]
  \item $f(x,y)=x^{3}-3x+y^{2}$: at $Q=(-1,1)$ in the direction of
        $\mathbf{h}=(3,4)$, compute $D_{\mathbf{u}}f$ and
        $D^{2}_{\mathbf{u}}f$; then use the point $(1,0)$.
  \item $f(x,y)=x^{4}-4xy+y^{2}$: classify the critical points, then give
        $D^{2}_{\mathbf{u}}f(0,0)$ for $\mathbf{u}=(1,1)/\sqrt2$ and for
        $\mathbf{u}=(0,1)$.
  \item $f(x,y)=x\ln y$ on $y>0$, at $P=(1,1)$: $D_{\mathbf{u}}f$ and
        $D^{2}_{\mathbf{u}}f$ for $\mathbf{u}=(1,0)$, $(0,1)$ and
        $(1,1)/\sqrt2$.
\end{problems}

Problems 38--40: statements to check. Justify every answer.

\begin{problems}[resume]
  \item True or false? (a) Walking along the contour line through $P$,
        the rate of change of $f$ is $0$. (b) The direction of steepest
        descent is $-\grad f(P)$. (c) If $\grad f(P)=\mathbf{0}$, then $P$
        is a peak or a valley. (d) $|\grad f(P)|$ is the slope of the
        steepest path through $P$. (e) The gradient at $P$ is tangent to
        the level curve through $P$.
  \item A hiker stands on a hillside with neither map nor formula. How can
        she find the direction of the gradient, and how does the slope in
        other directions compare with the slope in that one?
  \item True or false? (a) For $f(x,y)$ of two variables: if
        $D_{\mathbf{u}}f(P)=0$ for two non-parallel unit vectors $\mathbf{u}$, then $\grad f(P)=\mathbf{0}$. (b) If
        $D_{\mathbf{u}}f(P)=3$, then $D_{-\mathbf{u}}f(P)=-3$. (c) A
        function can have $D_{\mathbf{u}}f(P)=5$ for some unit vector
        $\mathbf{u}$ while $|\grad f(P)|=4$.
\end{problems}

Problems 41--42: an input change vector of fixed length is to raise output
as fast as possible.

\begin{problems}[resume]
  \item $Y=10K^{0.4}L^{0.6}$ at $K=32768$, $L=7776$ (so that
        $K^{0.4}=64$, $L^{0.6}=216$). The change vector
        $(\mathrm{d}K,\mathrm{d}L)$ has length $100$. Find the best ratio
        $\mathrm{d}K:\mathrm{d}L$, the amounts, the estimated and the exact
        rise in output; then compare with spending a budget
        $\mathrm{d}K+\mathrm{d}L=100$.
  \item $Y=\sqrt{KL}$ at $K=400$, $L=900$, with a change vector of length
        $20$: the same questions.
\end{problems}

\section{Additional Solutions}

\begin{solutions}
  \item $\grad f=(2x+3y,\,-2y+3x)$; at $(1,2)$:
        $\grad f=(2+6,\,-4+3)=(8,-1)$. Normalise:
        $|\mathbf{v}|=\sqrt{16+9}=5$, $\mathbf{u}=(4/5,\,3/5)$.
        \[
          D_{\mathbf{u}}f=8\cdot\frac{4}{5}+(-1)\cdot\frac{3}{5}
            =\frac{32}{5}-\frac{3}{5}=\frac{29}{5} .
        \]

  \item $\grad f=(2x,\,4y,\,6z)$; at $(1,1,1)$: $\grad f=(2,4,6)$.
        Normalise: $|\mathbf{v}|=\sqrt{1+1+1}=\sqrt{3}$,
        $\mathbf{u}=(1,-1,1)/\sqrt{3}$.
        \[
          D_{\mathbf{u}}f=\frac{2-4+6}{\sqrt{3}}
            =\frac{4}{\sqrt{3}}=\frac{4\sqrt{3}}{3} .
        \]

  \item $\grad U=(1/x,\,2/y)$; at $(1,2)$: $\grad U=(1,1)$,
        $|\grad U|=\sqrt{2}$. Steepest ascent along
        \[
          \mathbf{u}=\frac{(1,1)}{\sqrt{2}},\qquad
          \text{maximum rate}=\sqrt{2} .
        \]

  \item $\grad f=(y,\,x)$; at $(2,1)$: $\grad f=(1,2)$. We need
        $\mathbf{u}\perp(1,2)$: $u_{1}+2u_{2}=0$, so
        $\mathbf{u}=\pm(2,-1)/\sqrt{5}$.
        \textbf{Check.} $(1,2)\cdot(2,-1)=2-2=0$. \checkmark

  \item The directional derivatives toward $(1,0)$ and $(0,1)$ are
        precisely $f_{x}(P)$ and $f_{y}(P)$, so $\grad f(P)=(3,-4)$.
        Maximum rate: $|\grad f(P)|=\sqrt{9+16}=5$, attained toward
        $(3,-4)/5$. Toward $\mathbf{u}=(3/5,\,4/5)$:
        \[
          D_{\mathbf{u}}f=3\cdot\frac{3}{5}+(-4)\cdot\frac{4}{5}
            =\frac{9}{5}-\frac{16}{5}=-\frac{7}{5} .
        \]

  \item $\grad f=\Bigl(\dfrac{2x}{x^{2}+y^{2}},\,
        \dfrac{2y}{x^{2}+y^{2}}\Bigr)$; at $(1,0)$: $\grad f=(2,0)$, so
        \[
          D_{\mathbf{u}}f=(2,0)\cdot(0,1)=0 :
        \]
        $\mathbf{u}$ is tangent to the level curve $x^{2}+y^{2}=1$
        through $(1,0)$. Second partials:
        \[
          f_{xx}=\frac{2y^{2}-2x^{2}}{(x^{2}+y^{2})^{2}},\qquad
          f_{yy}=\frac{2x^{2}-2y^{2}}{(x^{2}+y^{2})^{2}},\qquad
          f_{xy}=\frac{-4xy}{(x^{2}+y^{2})^{2}} ;
        \]
        at $(1,0)$: $f_{xx}=-2$, $f_{yy}=2$, $f_{xy}=0$. Hence with
        $\mathbf{u}=(0,1)$:
        \[
          D^{2}_{\mathbf{u}}f=f_{yy}=2>0 .
        \]
        Along $\mathbf{u}$ the function is flat to first order but
        curves upward: moving off $(1,0)$ vertically raises $f$ at
        second order. \textbf{Check.}
        $f(1,t)=\ln(1+t^{2})\approx t^{2}$ for small $t$: derivative
        $0$, second derivative $2$. \checkmark

  \item
        \[
          J_{\mathbf{f}}=\begin{pmatrix}2x&-2y\\3y&3x+2y\end{pmatrix},
          \qquad
          J_{\mathbf{f}}(1,1)=\begin{pmatrix}2&-2\\3&5\end{pmatrix} .
        \]

  \item $J_{f}=\bigl(y\;\;x\bigr)$, evaluated at
        $\mathbf{r}(t)=(t^{2},\,t^{3})$: $\bigl(t^{3}\;\;t^{2}\bigr)$;
        $J_{\mathbf{r}}(t)=\begin{pmatrix}2t\\3t^{2}\end{pmatrix}$.
        \[
          \dd{z}{t}=\bigl(t^{3}\;\;t^{2}\bigr)
            \begin{pmatrix}2t\\3t^{2}\end{pmatrix}
            =2t^{4}+3t^{4}=5t^{4} .
        \]
        \textbf{Check.} $z=t^{2}\cdot t^{3}=t^{5}$, so
        $\dd{z}{t}=5t^{4}$. \checkmark

  \item Inner and outer Jacobians:
        \[
          J_{\mathbf{g}}(s,t)=\begin{pmatrix}t&s\\1&1\end{pmatrix},\qquad
          J_{\mathbf{F}}(x,y)=\begin{pmatrix}1&1\\y&x\end{pmatrix} .
        \]
        Evaluate the outer at $(x,y)=\mathbf{g}(s,t)=(st,\,s+t)$ and
        multiply:
        \[
          J_{\mathbf{F}\circ\mathbf{g}}
            =\begin{pmatrix}1&1\\s+t&st\end{pmatrix}
             \begin{pmatrix}t&s\\1&1\end{pmatrix}
            =\begin{pmatrix}t+1&s+1\\
              (s+t)t+st&(s+t)s+st\end{pmatrix}
            =\begin{pmatrix}t+1&s+1\\
              t^{2}+2st&s^{2}+2st\end{pmatrix} .
        \]
        \textbf{Check.} Composing first:
        $(\mathbf{F}\circ\mathbf{g})(s,t)
        =\bigl(st+s+t,\;s^{2}t+st^{2}\bigr)$, with Jacobian
        \[
          \begin{pmatrix}t+1&s+1\\2st+t^{2}&s^{2}+2st\end{pmatrix} .
          \quad\checkmark
        \]

  \item Inner map $\mathbf{g}(K,L)=(2K+L,\;K+3L)$ with constant
        Jacobian
        $J_{\mathbf{g}}=\begin{pmatrix}2&1\\1&3\end{pmatrix}$; outer
        function $y=x_{1}x_{2}$ with row Jacobian
        $\grad y=(x_{2}\;\;x_{1})$. At $(K,L)=(1,1)$:
        $(x_{1},x_{2})=(3,4)$, so the row is $(4\;\;3)$ and
        \[
          \Bigl(\pd{y}{K},\,\pd{y}{L}\Bigr)
            =(4\;\;3)\begin{pmatrix}2&1\\1&3\end{pmatrix}
            =\bigl(4\cdot 2+3\cdot 1,\;4\cdot 1+3\cdot 3\bigr)
            =(11,\;13) .
        \]
        \textbf{Check.} Expanding,
        $y=(2K+L)(K+3L)=2K^{2}+7KL+3L^{2}$, so
        $\pd{y}{K}=4K+7L=11$ and $\pd{y}{L}=7K+6L=13$ at $(1,1)$.
        \checkmark

  \item $f(P)=2+8=10$ and $\grad f=(2xy,\;x^{2}+3y^{2})=(4,13)$ at $P$.
        \begin{itemize}[leftmargin=1.2em,itemsep=0.2ex]
          \item $x$-direction: slope $f_{x}(P)=4$, line
                $(1+t,\;2,\;10+4t)$.
          \item $y$-direction: slope $f_{y}(P)=13$, line
                $(1,\;2+t,\;10+13t)$.
          \item $|\mathbf{v}|=5$, $\mathbf{u}=(\tfrac35,-\tfrac45)$,
                $D_{\mathbf{u}}f=\tfrac{12}{5}-\tfrac{52}{5}=-8$: the
                surface falls with slope $8$ that way. Line
                $\bigl(1+\tfrac35t,\;2-\tfrac45t,\;10-8t\bigr)$.
        \end{itemize}
  \item $f(P)=\ln4$. With $w=x^{2}+3y$, $\grad f=\bigl(\tfrac{2x}{w},\,
        \tfrac{3}{w}\bigr)=(1,\tfrac34)$ at $P$, where $w=4$. Steepest
        ascent: $\mathbf{u}=\grad f/|\grad f|$ with $|\grad f|=\tfrac54$, so
        $\mathbf{u}=(\tfrac45,\tfrac35)$ and the slope is
        $D_{\mathbf{u}}f=|\grad f|=\tfrac54$. Tangent line:
        \[
          (x,y,z)=\bigl(2+\tfrac45t,\;\tfrac35t,\;\ln4+\tfrac54t\bigr) .
        \]
  \item $f(P)=2\mathrm{e}^{0}=2$. Product rule:
        $f_{x}=\mathrm{e}^{x-y}(x+y+1)=3$,
        $f_{y}=\mathrm{e}^{x-y}(1-x-y)=-1$ at $P$. $\mathbf{h}=Q-P=(2,1)$,
        $|\mathbf{h}|=\sqrt5$, $\mathbf{u}=(2,1)/\sqrt5$:
        \[
          D_{\mathbf{u}}f=\frac{3\cdot2-1\cdot1}{\sqrt5}=\sqrt5\approx2.236 .
        \]
        With the unit vector:
        $\bigl(1+\tfrac{2}{\sqrt5}t,\;1+\tfrac{1}{\sqrt5}t,\;2+\sqrt5\,t\bigr)$.
        With $\mathbf{h}$, $\grad f\cdot\mathbf{h}=5$:
        $(1+2s,\;1+s,\;2+5s)$. Putting $s=t/\sqrt5$ in the second gives the
        first, so it is the same line; only the parameter runs at a
        different speed.
  \item $Q(16,16)=4\cdot4\cdot2=32$,
        $Q_{L}=K^{1/2}L^{-3/4}=4\cdot\tfrac18=\tfrac12$. Tangent line in the
        $L$-direction: $(16,\;16+t,\;32+\tfrac12t)$. At $L=17$ ($t=1$) it
        gives $32.5$; exactly, $Q(16,17)=16\cdot17^{1/4}\approx32.4887$. The
        slope $\tfrac12$ is the marginal product of labour, and the line
        slightly overestimates because $Q$ is concave in $L$.
  \item (a) $c=\sqrt{2+6+1}=3$. (b) $2x^{2}+3y+1=9$, so
        $y=\dfrac{8-2x^{2}}{3}$, a downward parabola. Its slope at $x=1$ is
        $y'=-\tfrac{4x}{3}=-\tfrac43$, tangent direction $(3,-4)$.
        (c) $\grad f=\Bigl(\dfrac{2x}{f},\,\dfrac{3}{2f}\Bigr)=
        \bigl(\tfrac23,\tfrac12\bigr)$, and
        $\bigl(\tfrac23,\tfrac12\bigr)\cdot(3,-4)=2-2=0$. \checkmark\ (d)
        Normal form: $\tfrac23(x-1)+\tfrac12(y-2)=0$, i.e.\
        $4(x-1)+3(y-2)=0$, $y=2-\tfrac43(x-1)$. Parametric:
        $(x,y)=(1,2)+t(3,-4)$. (e) $D_{\mathbf{u}}f=|\grad f|=
        \sqrt{\tfrac49+\tfrac14}=\sqrt{\tfrac{25}{36}}=\tfrac56$.
  \item (a) $c=1\cdot8^{2/3}=4$. (b) $y^{2/3}=4x^{-1/3}$, so
        $y=8x^{-1/2}$; slope $y'=-4x^{-3/2}=-4$ at $x=1$, direction $(1,-4)$.
        (c) $f_{x}=\tfrac13\dfrac fx=\tfrac43$,
        $f_{y}=\tfrac23\dfrac fy=\tfrac23\cdot\tfrac48=\tfrac13$, and
        $\bigl(\tfrac43,\tfrac13\bigr)\cdot(1,-4)=0$. \checkmark\ (d)
        $\tfrac43(x-1)+\tfrac13(y-8)=0$, i.e.\ $y=12-4x$; parametric
        $(1,8)+t(1,-4)$. (e) $|\grad f|=\sqrt{\tfrac{16}{9}+\tfrac19}=
        \tfrac{\sqrt{17}}{3}\approx1.374$.
  \item (a) $c=2\cdot1+0=2$. (b) $x\mathrm{e}^{y}=2-y$, so
        $x=(2-y)\mathrm{e}^{-y}$. Then
        $\dd{x}{y}=-\mathrm{e}^{-y}-(2-y)\mathrm{e}^{-y}=(y-3)\mathrm{e}^{-y}
        =-3$ at $y=0$: tangent direction $(\mathrm{d}x,\mathrm{d}y)=(-3,1)$.
        (c) $\grad f=(\mathrm{e}^{y},\;x\mathrm{e}^{y}+1)=(1,3)$, and
        $(1,3)\cdot(-3,1)=0$. \checkmark\ (d) $(x-2)+3y=0$, i.e.\
        $x=2-3y$; parametric $(2,0)+t(-3,1)$. (e) $|\grad f|=\sqrt{10}
        \approx3.162$.
  \item Each component is a polynomial, so its partials exist and are
        continuous everywhere; a $C^{1}$ map is differentiable, and a map is
        differentiable when each component is.
        $J_{\mathbf{F}}=\begin{pmatrix}2xy&x^{2}\\1&-6y\\y&x\end{pmatrix}$,
        and $J_{\mathbf{F}}(1,1)=\begin{pmatrix}2&1\\1&-6\\1&1\end{pmatrix}$.
        Remainder at $(1,1)$, component by component, with
        $\mathbf{F}(1,1)=(1,-2,1)$:
        \[\begin{aligned}
          r_{1}&=(1+h_{1})^{2}(1+h_{2})-1-2h_{1}-h_{2}
               =2h_{1}h_{2}+h_{1}^{2}+h_{1}^{2}h_{2},\\
          r_{2}&=(1+h_{1})-3(1+h_{2})^{2}+2-h_{1}+6h_{2}=-3h_{2}^{2},\\
          r_{3}&=(1+h_{1})(1+h_{2})-1-h_{1}-h_{2}=h_{1}h_{2} .
        \end{aligned}\]
        Each is bounded by a constant times $\|\mathbf{h}\|^{2}$ (for
        $\|\mathbf{h}\|\le1$), so
        $\|\mathbf{r}\|\le C\|\mathbf{h}\|^{2}$ and
        $\|\mathbf{r}\|/\|\mathbf{h}\|\to0$.
  \item \emph{Domain.} $3x^{2}+y^{4}>0$ fails only at the origin: $D=\R^{2}
        \setminus\{(0,0)\}$. \emph{Differentiability.} The first two
        components are $C^{1}$ on $D$. The third, $|y-1|$, has
        $\partial/\partial y$ equal to $+1$ above and $-1$ below the line
        $y=1$, and on the line the quotient $|h|/h$ has no limit. So
        $\mathbf{F}$ is differentiable exactly on
        $\{(x,y)\neq(0,0):y\neq1\}$, with
        \[
          J_{\mathbf{F}}=\begin{pmatrix}
            \dfrac{6x}{3x^{2}+y^{4}}&\dfrac{4y^{3}}{3x^{2}+y^{4}}\\[1.2ex]
            y^{2}&2xy\\ 0&\operatorname{sign}(y-1)\end{pmatrix} .
        \]
  \item $\mathbf{f}$ is differentiable at $\mathbf{x}_{0}$ if
        $\mathbf{f}(\mathbf{x}_{0}+\mathbf{h})=\mathbf{f}(\mathbf{x}_{0})
        +J\mathbf{h}+\mathbf{r}(\mathbf{h})$ for a matrix $J$, with
        $\|\mathbf{r}(\mathbf{h})\|/\|\mathbf{h}\|\to0$: near
        $\mathbf{x}_{0}$ the map is a linear map plus an error that vanishes
        faster than the step. The matrix is the derivative (total
        derivative), the Jacobian matrix, of shape $m\times n$. Its entries
        are the partials $\partial f_{i}/\partial x_{j}$; their existence is
        necessary but not sufficient, while continuity of all of them is
        sufficient, and $\mathbf{f}$ is differentiable exactly when each
        component is. For $n=m=1$ the linear part is the tangent line of
        the graph, for $n=2$, $m=1$ the tangent plane.
  \item $J_{\mathbf{r}}(t)=\Bigl(\dfrac{2t}{1+t^{2}},\;-2\mathrm{e}^{-2t},
        \;-3\sin3t,\;0\Bigr)^{T}$, a $4\times1$ column (four outputs, one
        input); the constant component gives $0$. At $t=0$:
        $(0,-2,0,0)^{T}$.
  \item $J_{\mathbf{F}}=\begin{pmatrix}yz&xz&xy\\2x&0&-1\end{pmatrix}$,
        $2\times3$; at $(1,2,3)$: $\begin{pmatrix}6&3&2\\2&0&-1\end{pmatrix}$.
        $J_{g}=\bigl(x_{2}\ \ x_{1}\ \ \mathrm{e}^{x_{4}}\ \
        x_{3}\mathrm{e}^{x_{4}}\bigr)$, a $1\times4$ row; at $(1,1,1,0)$:
        $(1\ \ 1\ \ 1\ \ 1)$.
  \item $J_{\mathbf{F}}=\begin{pmatrix}1&1\\t&s\\2s&0\\
        \mathrm{e}^{s-t}&-\mathrm{e}^{s-t}\end{pmatrix}$, $4\times2$ (four
        outputs, two inputs); at $(1,1)$:
        $\begin{pmatrix}1&1\\1&1\\2&0\\1&-1\end{pmatrix}$.
  \item (a) $F_{k}=x_{k}$, so $\partial F_{k}/\partial x_{j}$ is $1$ for
        $j=k$ and $0$ otherwise: $J=I_{n}$. (b) $F_{k}=k\mathrm{e}^{x_{k}}$
        depends on $x_{k}$ alone: $J=\operatorname{diag}\bigl(\mathrm{e}^{x_{1}},
        2\mathrm{e}^{x_{2}},\ldots,n\,\mathrm{e}^{x_{n}}\bigr)$, $n\times n$.
        (c) $r_{k}=t^{k}/k$ has $r_{k}'=t^{k-1}$: the $n\times1$ column
        $(1,t,t^{2},\ldots,t^{n-1})^{T}$. (d) The term with $x_{k}$ is
        $x_{k}^{k}/k$, with derivative $x_{k}^{k-1}$: the $1\times n$ row
        $(1\ \ x_{2}\ \ x_{3}^{2}\ \ \cdots\ \ x_{n}^{n-1})$.
  \item (a) Component $i$ is $\sum_{j}a_{ij}x_{j}$, with partials $a_{ij}$:
        $J=A$ at every point. (b) Component $k$ is $3x_{k}-2a_{k}$:
        $J=3I_{n}$, and the constant shift drops out. (c) $\partial f/
        \partial x_{k}=a_{k}$: $J=\mathbf{a}^{T}$, a $1\times n$ row. (d)
        $J=B$, $2\times3$, whatever $\mathbf{c}$ is.
  \item Match outputs of the inner map with inputs of the outer map.
        \begin{itemize}[leftmargin=1.2em,itemsep=0.2ex]
          \item $\mathbf{G}\circ\mathbf{F}$: $2=2$, defined,
                $\R^{3}\to\R^{4}$, $4\times3$. $\mathbf{F}\circ\mathbf{G}$:
                $4\neq3$, undefined.
          \item $H\circ\mathbf{G}$: $4=4$, defined, $\R^{2}\to\R$,
                $1\times2$. $\mathbf{G}\circ H$: $1\neq2$, undefined.
          \item $\mathbf{K}\circ H$: $1=1$, defined, $\R^{4}\to\R^{3}$,
                $3\times4$. $H\circ\mathbf{K}$: $3\neq4$, undefined.
          \item $\mathbf{F}\circ\mathbf{K}$: $3=3$, defined, $\R\to\R^{2}$,
                $2\times1$. $\mathbf{K}\circ\mathbf{F}$: $2\neq1$,
                undefined.
          \item $H\circ\mathbf{G}\circ\mathbf{F}$: $\R^{3}\to\R^{2}\to\R^{4}
                \to\R$, defined, $1\times3=(1\times4)(4\times2)(2\times3)$.
        \end{itemize}
  \item $\mathbf{f}(1,2,1)=(1+2,\;1)=(3,1)$.
        $J_{\mathbf{f}}=\begin{pmatrix}1&z&y\\z&0&x\end{pmatrix}
        =\begin{pmatrix}1&1&2\\1&0&1\end{pmatrix}$ and
        $J_{\mathbf{g}}=\begin{pmatrix}1&0\\0&1\\v&u\\1&-1\end{pmatrix}
        =\begin{pmatrix}1&0\\0&1\\1&3\\1&-1\end{pmatrix}$ at $(u,v)=(3,1)$.
        \[
          J_{\mathbf{g}\circ\mathbf{f}}=
          \begin{pmatrix}1&0\\0&1\\1&3\\1&-1\end{pmatrix}
          \begin{pmatrix}1&1&2\\1&0&1\end{pmatrix}
          =\begin{pmatrix}1&1&2\\1&0&1\\4&1&5\\0&1&1\end{pmatrix},
        \]
        a $(4\times2)(2\times3)=4\times3$ product. \textbf{Check.}
        $\mathbf{g}\circ\mathbf{f}=\bigl(x+yz,\;xz,\;x^{2}z+xyz^{2},\;
        x+yz-xz\bigr)$; the third row is
        $(2xz+yz^{2},\;xz^{2},\;x^{2}+2xyz)=(4,1,5)$ and the fourth
        $(1-z,\;z,\;y-x)=(0,1,1)$ at $(1,2,1)$. \checkmark
  \item $g\colon\R^{3}\to\R$ and $\mathbf{f}\colon\R\to\R^{3}$: in
        $\mathbf{f}\circ g$ the one output of $g$ feeds the one input of
        $\mathbf{f}$, and in $g\circ\mathbf{f}$ the three outputs of
        $\mathbf{f}$ feed the three inputs of $g$. Both exist.
        \emph{$\mathbf{f}\circ g\colon\R^{3}\to\R^{3}$.} The chain rule
        gives a column times a row,
        $J=\mathbf{f}'(g(\mathbf{x}))\,\grad g(\mathbf{x})^{T}$, a
        rank-one $3\times3$ matrix. At $(1,0,0)$: $g=1$,
        $\mathbf{f}'(1)=(1,\;2,\;\mathrm{e})^{T}$, $\grad g=(1,2,2z)=(1,2,0)$:
        \[
          J=\begin{pmatrix}1\\2\\\mathrm{e}\end{pmatrix}(1\ \ 2\ \ 0)
           =\begin{pmatrix}1&2&0\\2&4&0\\\mathrm{e}&2\mathrm{e}&0\end{pmatrix} .
        \]
        Check: $\mathbf{f}(g)=(g,\;g^{2},\;\mathrm{e}^{g})$ has rows
        $\grad g$, $2g\grad g$, $\mathrm{e}^{g}\grad g$. \checkmark\
        \emph{$g\circ\mathbf{f}\colon\R\to\R$.}
        $\grad g(\mathbf{f}(s))\cdot\mathbf{f}'(s)=(1,\;2,\;2\mathrm{e}^{s})
        \cdot(1,\;2s,\;\mathrm{e}^{s})=1+4s+2\mathrm{e}^{2s}$. Check:
        $g(\mathbf{f}(s))=s+2s^{2}+\mathrm{e}^{2s}$. \checkmark
  \item \emph{$\mathbf{g}\circ\mathbf{f}\colon\R^{2}\to\R^{2}$} at $(1,2)$:
        $\mathbf{f}(1,2)=(1,2,2)$,
        $J_{\mathbf{f}}=\begin{pmatrix}1&0\\0&1\\y&x\end{pmatrix}
        =\begin{pmatrix}1&0\\0&1\\2&1\end{pmatrix}$,
        $J_{\mathbf{g}}=\begin{pmatrix}1&0&1\\0&w&v\end{pmatrix}
        =\begin{pmatrix}1&0&1\\0&2&2\end{pmatrix}$, so
        \[
          J_{\mathbf{g}\circ\mathbf{f}}=\begin{pmatrix}1&0&1\\0&2&2\end{pmatrix}
          \begin{pmatrix}1&0\\0&1\\2&1\end{pmatrix}
          =\begin{pmatrix}3&1\\4&4\end{pmatrix} .
        \]
        Check: $\mathbf{g}(\mathbf{f}(x,y))=(x+xy,\;xy^{2})$, Jacobian
        $\begin{pmatrix}1+y&x\\y^{2}&2xy\end{pmatrix}$, which is
        $\begin{pmatrix}3&1\\4&4\end{pmatrix}$ at $(1,2)$. \checkmark\
        \emph{$\mathbf{f}\circ\mathbf{g}\colon\R^{3}\to\R^{3}$} at
        $(1,1,1)$: $\mathbf{g}(1,1,1)=(2,1)$, so
        $J_{\mathbf{f}}(2,1)=\begin{pmatrix}1&0\\0&1\\1&2\end{pmatrix}$ and
        $J_{\mathbf{g}}(1,1,1)=\begin{pmatrix}1&0&1\\0&1&1\end{pmatrix}$:
        \[
          J_{\mathbf{f}\circ\mathbf{g}}=\begin{pmatrix}1&0\\0&1\\1&2\end{pmatrix}
          \begin{pmatrix}1&0&1\\0&1&1\end{pmatrix}
          =\begin{pmatrix}1&0&1\\0&1&1\\1&2&3\end{pmatrix} .
        \]
        Check: $\mathbf{f}(\mathbf{g})=\bigl(u+w,\;vw,\;(u+w)vw\bigr)$; the
        last row is $\bigl(vw,\;(u+w)w,\;vw+(u+w)v\bigr)=(1,2,3)$.
        \checkmark
  \item $\grad F=(2x,2y,2z)$ and $\mathbf{r}'(t)=(-\sin t,\cos t,1)$:
        \[
          \dd{}{t}F(\mathbf{r}(t))=2\cos t(-\sin t)+2\sin t\cos t+2t\cdot1
          =2t .
        \]
        Check: $F(\mathbf{r}(t))=\cos^{2}t+\sin^{2}t+t^{2}=1+t^{2}$.
        \checkmark
  \item $|\mathbf{h}|=3$. $\grad f=(\ln y,\;x/y+z^{2},\;2yz)=(0,3,2)$, so
        $\grad f\cdot\mathbf{h}=0+6-4=2$ and $D_{\mathbf{u}}f=\tfrac23$.
        $H=\begin{pmatrix}0&1/y&0\\1/y&-x/y^{2}&2z\\0&2z&2y\end{pmatrix}
        =\begin{pmatrix}0&1&0\\1&-2&2\\0&2&2\end{pmatrix}$ and
        \[
          \mathbf{h}^{T}H\mathbf{h}=2\cdot1\cdot(1)(2)-2\cdot4
          +2\cdot2\cdot(2)(-2)+2\cdot4=4-8-16+8=-12,
        \]
        so $D^{2}_{\mathbf{u}}f=-\tfrac{12}{9}=-\tfrac43$. \emph{Profile.}
        $\varphi(t)=(2+t)\ln(1+2t)+(1+2t)(1-2t)^{2}$, and
        $(1+2t)(1-2t)^{2}=1-2t-4t^{2}+8t^{3}$. The log part has derivative
        $\ln(1+2t)+\dfrac{2(2+t)}{1+2t}$, equal to $4$ at $0$, and second
        derivative $\dfrac{2}{1+2t}+\dfrac{2(1+2t)-4(2+t)}{(1+2t)^{2}}$,
        equal to $2-6=-4$. So $\varphi'(0)=4-2=2$ and
        $\varphi''(0)=-4-8=-12$. \checkmark\ Concave in this direction.
  \item $|\mathbf{h}|=3$. $\grad f=\mathrm{e}^{yz}(1,\;xz,\;xy)=(1,2,0)$;
        $D_{\mathbf{u}}f=\tfrac{2+2+0}{3}=\tfrac43$. Second partials at $P$
        ($\mathrm{e}^{yz}=1$): $f_{xx}=0$, $f_{xy}=z\mathrm{e}^{yz}=2$,
        $f_{xz}=y\mathrm{e}^{yz}=0$, $f_{yy}=xz^{2}\mathrm{e}^{yz}=4$,
        $f_{yz}=x(1+yz)\mathrm{e}^{yz}=1$, $f_{zz}=xy^{2}\mathrm{e}^{yz}=0$.
        \[
          \mathbf{h}^{T}H\mathbf{h}=2\cdot2\cdot(2)(1)+4\cdot1+2\cdot1\cdot(1)(2)
          =8+4+4=16,\qquad D^{2}_{\mathbf{u}}f=\tfrac{16}{9} .
        \]
        \emph{Profile.} $\varphi(t)=(1+2t)\,\mathrm{e}^{t(2+2t)}
        =(1+2t)\,\mathrm{e}^{2t+2t^{2}}$. With $E=\mathrm{e}^{2t+2t^{2}}$,
        $E'=(2+4t)E$:
        $\varphi'=\bigl(2+(1+2t)(2+4t)\bigr)E$, so $\varphi'(0)=4$, and
        $\varphi''=\bigl(4+8t+4+8t\bigr)E+\bigl(2+(1+2t)(2+4t)\bigr)(2+4t)E$,
        so $\varphi''(0)=8+4\cdot2=16$. \checkmark\ Convex in this direction.
  \item $|\mathbf{h}|=\sqrt3$. $\grad f=(yz+\tfrac1x,\;xz,\;xy+\tfrac1z)
        =(2,1,2)$, $\grad f\cdot\mathbf{h}=2-1+2=3$,
        $D_{\mathbf{u}}f=3/\sqrt3=\sqrt3$.
        $H=\begin{pmatrix}-1/x^{2}&z&y\\z&0&x\\y&x&-1/z^{2}\end{pmatrix}
        =\begin{pmatrix}-1&1&1\\1&0&1\\1&1&-1\end{pmatrix}$, and
        \[
          \mathbf{h}^{T}H\mathbf{h}=-1+0-1+2\bigl[(1)(-1)+(1)(1)+(-1)(1)\bigr]
          =-2-2=-4,\qquad D^{2}_{\mathbf{u}}f=-\tfrac43 .
        \]
        \emph{Profile.} $\varphi(t)=(1+t)(1-t)(1+t)+\ln\bigl((1+t)^{2}\bigr)
        =1+t-t^{2}-t^{3}+2\ln(1+t)$. $\varphi'(t)=1-2t-3t^{2}+\tfrac{2}{1+t}$,
        $\varphi'(0)=3$; $\varphi''(t)=-2-6t-\tfrac{2}{(1+t)^{2}}$,
        $\varphi''(0)=-4$. \checkmark\ Concave in this direction.
  \item At $(1,2)$: $\grad f=(2xy,x^{2})=(4,1)$ and
        $H=\begin{pmatrix}2y&2x\\2x&0\end{pmatrix}=\begin{pmatrix}4&2\\2&0
        \end{pmatrix}$. With $\mathbf{h}=(3,4)$:
        $\grad f\cdot\mathbf{h}=12+4=16$ and
        $\mathbf{h}^{T}H\mathbf{h}=4\cdot9+2\cdot2\cdot12+0=84$. Profile:
        $\varphi(t)=(1+3t)^{2}(2+4t)=2+16t+42t^{2}+36t^{3}$, so
        $\varphi'(0)=16$ and $\varphi''(0)=84$. \checkmark\ Practice
        problem 15 found $D_{\mathbf{u}}f=\tfrac{16}{5}$ and
        $D^{2}_{\mathbf{u}}f=\tfrac{84}{25}$ with $\mathbf{u}=\mathbf{h}/5$:
        moving along $\mathbf{h}$ covers $5$ units of distance per unit of
        $t$, so first derivatives scale by $|\mathbf{h}|=5$ and second
        derivatives by $|\mathbf{h}|^{2}=25$. The sign, and so the
        convexity verdict, is the same either way.
  \item $\grad f=(3x^{2}-3,\;2y)$, $H=\begin{pmatrix}6x&0\\0&2\end{pmatrix}$.
        At $Q=(-1,1)$: $\grad f=(0,2)$, $|\mathbf{h}|=5$,
        $\mathbf{u}=(\tfrac35,\tfrac45)$, so $D_{\mathbf{u}}f=\tfrac85$;
        $H=\begin{pmatrix}-6&0\\0&2\end{pmatrix}$ and
        $D^{2}_{\mathbf{u}}f=\dfrac{-6\cdot9+2\cdot16}{25}=-\dfrac{22}{25}<0$:
        rising, but concave in this direction at $Q$. At $(1,0)$ the
        gradient is $\mathbf{0}$ and $H=\begin{pmatrix}6&0\\0&2\end{pmatrix}$
        is PD: a strict local minimum, with positive curvature in every
        direction. Negative curvature somewhere rules out convexity,
        positive curvature somewhere rules out concavity: $f$ is
        \textbf{neither}.
  \item $f_{x}=4x^{3}-4y$, $f_{y}=-4x+2y$. From $f_{y}=0$: $y=2x$; then
        $4x^{3}=8x$, so $x=0$ or $x^{2}=2$. Critical points $(0,0)$ and
        $\pm(\sqrt2,2\sqrt2)$. $H=\begin{pmatrix}12x^{2}&-4\\-4&2\end{pmatrix}$.
        At $\pm(\sqrt2,2\sqrt2)$: $\begin{pmatrix}24&-4\\-4&2\end{pmatrix}$,
        $D_{1}=24$, $D_{2}=48-16=32$: PD, strict local minima, with value
        $4-16+8=-4$. At $(0,0)$: $\begin{pmatrix}0&-4\\-4&2\end{pmatrix}$,
        $D_{2}=-16$: a saddle. There, for $\mathbf{u}=(1,1)/\sqrt2$,
        $D^{2}_{\mathbf{u}}f=\tfrac{0-8+2}{2}=-3<0$ (concave that way), and
        for $\mathbf{u}=(0,1)$, $D^{2}_{\mathbf{u}}f=2>0$ (convex that way).
        Globally: negative curvature at the origin and positive definite
        Hessians at the minima, so $f$ is neither convex nor concave.
  \item $\grad f=(\ln y,\;x/y)=(0,1)$ and
        $H=\begin{pmatrix}0&1/y\\1/y&-x/y^{2}\end{pmatrix}=
        \begin{pmatrix}0&1\\1&-1\end{pmatrix}$ at $P$.
        \begin{itemize}[leftmargin=1.2em,itemsep=0.2ex]
          \item $\mathbf{u}=(1,0)$: $D_{\mathbf{u}}f=0$,
                $D^{2}_{\mathbf{u}}f=0$. In fact $f(x,1)=0$: flat, no
                verdict beyond ``linear''.
          \item $\mathbf{u}=(0,1)$: $D_{\mathbf{u}}f=1$,
                $D^{2}_{\mathbf{u}}f=-1$: concave in this direction.
          \item $\mathbf{u}=(1,1)/\sqrt2$: $D_{\mathbf{u}}f=\tfrac{1}{\sqrt2}$,
                $D^{2}_{\mathbf{u}}f=\tfrac12(0+2-1)=\tfrac12$: convex in
                this direction.
        \end{itemize}
        Both signs at one point already decide it: $f$ is neither convex
        nor concave. Indeed $\det H=-1/y^{2}<0$ everywhere.
  \item (a) True: along the contour $f$ is constant; equivalently the
        tangent $\mathbf{u}$ is perpendicular to $\grad f$, so
        $D_{\mathbf{u}}f=0$. (b) True when $\grad f(P)\neq\mathbf{0}$: the
        unit vector $-\grad f/|\grad f|$ gives the most negative rate,
        $-|\grad f|$. (c) False: $f=xy$ has $\grad f(0,0)=\mathbf{0}$ at a
        saddle. (d) True: the largest directional derivative is
        $|\grad f(P)|$. (e) False: it is perpendicular (normal) to the level
        curve.
  \item She walks a short way in the direction in which she neither climbs
        nor descends: she is following the contour line. The gradient
        points at right angles to that path, on the uphill side. Its
        direction is the steepest one; in a direction at angle $\theta$ to
        it the slope is $|\grad f|\cos\theta$: smaller, zero along the
        contour, and negative once she turns downhill.
  \item (a) True: $\grad f\cdot\mathbf{u}_{1}=\grad f\cdot\mathbf{u}_{2}=0$
        with $\mathbf{u}_{1},\mathbf{u}_{2}$ independent forces $\grad f$ to
        be perpendicular to the whole plane, hence $\mathbf{0}$. With three
        variables two directions are not enough: $f(x,y,z)=z$ has
        $D_{\mathbf{e}_{1}}f=D_{\mathbf{e}_{2}}f=0$ but $\grad f=(0,0,1)$.
        (b) True:
        $D_{-\mathbf{u}}f=\grad f\cdot(-\mathbf{u})=-D_{\mathbf{u}}f$.
        (c) False: by Cauchy--Schwarz,
        $D_{\mathbf{u}}f=\grad f\cdot\mathbf{u}\le|\grad f|\,|\mathbf{u}|=4$.
  \item $Y=10\cdot64\cdot216=138240$. In output-ratio form
        $Y_{K}=0.4\,Y/K=\tfrac{27}{16}=1.6875$ and
        $Y_{L}=0.6\,Y/L=\tfrac{32}{3}\approx10.667$. Best ratio
        $\mathrm{d}K:\mathrm{d}L=Y_{K}:Y_{L}=81:512$.
        $|\grad Y|=\sqrt{1.6875^{2}+10.667^{2}}\approx10.7993$, so
        \[
          (\mathrm{d}K,\mathrm{d}L)=\frac{100}{10.7993}(1.6875,\;10.667)
          \approx(15.63,\;98.77),
          \qquad \mathrm{d}Y\approx100\cdot10.7993\approx1079.9 .
        \]
        Exactly, $Y(32783.63,\;7874.77)-Y\approx1077.5$. A budget
        $\mathrm{d}K+\mathrm{d}L=100$ is maximised in the linear estimate by
        spending everything on labour, the input with the larger marginal
        product: $\mathrm{d}Y\approx1066.7$ (exactly $1063.9$). That is
        less, but it also uses less: the steepest-ascent vector has
        $\mathrm{d}K+\mathrm{d}L\approx114.4$.
  \item $Y=\sqrt{400\cdot900}=600$, $Y_{K}=\dfrac{Y}{2K}=0.75$,
        $Y_{L}=\dfrac{Y}{2L}=\tfrac13$. Ratio $\mathrm{d}K:\mathrm{d}L=
        0.75:\tfrac13=9:4$. $|\grad Y|=\sqrt{0.5625+0.1111}\approx0.82074$,
        so $(\mathrm{d}K,\mathrm{d}L)\approx\tfrac{20}{0.82074}(0.75,\;0.3333)
        \approx(18.28,\;8.12)$ and $\mathrm{d}Y\approx20\cdot0.82074\approx
        16.41$; exactly $16.32$. A budget $\mathrm{d}K+\mathrm{d}L=20$ would
        go entirely into capital: $\mathrm{d}Y\approx15$ (exactly $14.82$).
\end{solutions}

\chapter{Inverse and Implicit Function Theorems, Jacobians, Homogeneity, Euler}

When a relation between variables is given implicitly, the Implicit
Function Theorem (IFT) extracts derivatives without solving for one
variable in terms of the other, exactly the situation behind
indifference curves and isoquants. The Jacobian generalises the derivative
to vector-valued maps; its determinant is the local area / volume scaling
factor, and the Inverse Function Theorem turns that determinant into a test
for whether the map can be undone near a point --- and, when it can, gives
the derivative of the inverse without ever computing the inverse.
Homogeneous functions and Euler's theorem then give a clean
structural identity that ties marginal products to total output,
underpinning constant-returns-to-scale models.

\section{Core Rules}

\begin{rulebox}{IFT, One Equation in Two Unknowns}
Suppose $F:\R^{2}\to\R$ is $C^{1}$ and $F(x_{0},y_{0})=0$ with
$F_{y}(x_{0},y_{0})\neq 0$. Then near $(x_{0},y_{0})$ there is a $C^{1}$
function $y=g(x)$ satisfying $F(x,g(x))=0$, and
\[
  \dd{y}{x} \;=\; -\frac{F_{x}}{F_{y}} .
\]
\end{rulebox}

\begin{rulebox}{What the condition $F_{y}\neq0$ rules out}
The theorem is local and its hypothesis is sharp. Where $F_{y}=0$ the
conclusion genuinely fails, and the failure is visible in the picture of the
curve $F=0$ (Chapter~5):
\begin{itemize}[leftmargin=1.6em,itemsep=0.3ex]
  \item $F_{y}=0$, $F_{x}\neq0$: a \textbf{vertical tangent}. No $y(x)$
        exists near the point, but $x=x(y)$ does, with
        $\dd{x}{y}=-F_{y}/F_{x}=0$. The theorem is symmetric in the two
        variables: at least one of the two representations survives as long
        as $\grad F\neq\mathbf{0}$.
  \item $F_{x}=F_{y}=0$: a \textbf{singular point}, typically a
        self-intersection where two branches with different slopes meet.
        Neither representation exists, and $-F_{x}/F_{y}$ degenerates to
        $0/0$.
\end{itemize}
So the working rule is: the level set of $F$ is a well-behaved curve, with a
well-defined tangent, exactly where $\grad F\neq\mathbf{0}$. This is the
same gradient condition that reappears as the \emph{constraint
qualification} in the Lagrange and Kuhn--Tucker chapters, where it rules out
precisely these kinks and crossings in the constraint set.
\end{rulebox}

\begin{rulebox}{IFT, Multivariate}
If $F(x_{1},\ldots,x_{n},y)=0$ with $F_{y}\neq 0$, then
$y=g(x_{1},\ldots,x_{n})$ locally and
\[
  \pd{y}{x_{i}} \;=\; -\frac{F_{x_{i}}}{F_{y}} .
\]
\end{rulebox}

\begin{rulebox}{MRS and MRTS via IFT}
For $U(x,y)=c$ set $F=U(x,y)-c$. Then $F_{x}=U_{x}$, $F_{y}=U_{y}$, so
\[
  \dd{y}{x}\bigg|_{U=c} = -\frac{U_{x}}{U_{y}}\quad\text{(signed)},\qquad
  \mathrm{MRS} = \frac{U_{x}}{U_{y}}\quad\text{(unsigned)} .
\]
Likewise on $Q(K,L)=Q_{0}$: $\dd{L}{K}=-Q_{K}/Q_{L}$,
$\mathrm{MRTS}=Q_{K}/Q_{L}$.
\end{rulebox}

\begin{rulebox}{Jacobian Matrix}
For $\mathbf{f}:\R^{n}\to\R^{m}$,
$\mathbf{f}(x_{1},\ldots,x_{n})=(f_{1},\ldots,f_{m})$, the \textbf{Jacobian}
is
\[
  J_{\mathbf{f}} \;=\;
  \begin{pmatrix}
    \partial f_{1}/\partial x_{1} & \cdots & \partial f_{1}/\partial x_{n}\\
    \vdots & \ddots & \vdots\\
    \partial f_{m}/\partial x_{1} & \cdots & \partial f_{m}/\partial x_{n}
  \end{pmatrix} \;\in\; \R^{m\times n} .
\]
For $m=n$, $\det J_{\mathbf{f}}$ is the local volume scaling factor. Whether
it vanishes is exactly what decides local invertibility, which is the next
three boxes.
\end{rulebox}

\begin{rulebox}{Inverse Function Theorem, one variable}
Let $f$ be continuously differentiable near $x_{0}$ and set $y_{0}=f(x_{0})$.
If
\[
  f'(x_{0})\neq0 ,
\]
then $f$ is invertible on some open interval around $x_{0}$, the inverse is
differentiable at $y_{0}$, and
\[
  \boxed{\;\bigl(f^{-1}\bigr)'(y_{0})=\frac{1}{f'(x_{0})}
        =\frac{1}{f'\bigl(f^{-1}(y_{0})\bigr)}\;}
\]
The formula is the chain rule read backwards. Differentiating the identity
$f\bigl(f^{-1}(y)\bigr)=y$ gives
$f'\bigl(f^{-1}(y)\bigr)\cdot\bigl(f^{-1}\bigr)'(y)=1$, and dividing
finishes it --- which is legitimate precisely because $f'\neq0$.

Note which point goes where: the derivative of the \emph{inverse} is
evaluated at $y_{0}$, the derivative of $f$ at $x_{0}$. Geometrically the
graph of $f^{-1}$ is the graph of $f$ mirrored in $y=x$, and mirroring turns
a slope into its reciprocal --- a vertical tangent for $f$ becoming a
horizontal one for $f^{-1}$, which is why $f'=0$ has to be excluded.

The derivatives of $\arcsin$, $\arccos$ and $\arctan$ in Chapter~30 are this
rule applied to $\sin$, $\cos$ and $\tan$; nothing else is needed to get
them.
\end{rulebox}

\begin{rulebox}{Inverse Function Theorem, several variables}
Let $\mathbf{f}:\R^{n}\to\R^{n}$ be continuously differentiable near
$\mathbf{x}_{0}$ and set $\mathbf{y}_{0}=\mathbf{f}(\mathbf{x}_{0})$. If
\[
  \det J_{\mathbf{f}}(\mathbf{x}_{0})\neq0 ,
\]
then $\mathbf{f}$ maps some open neighbourhood of $\mathbf{x}_{0}$
one-to-one onto an open neighbourhood of $\mathbf{y}_{0}$, the local inverse
is continuously differentiable, and its Jacobian is the \textbf{matrix
inverse} of the original:
\[
  \boxed{\;J_{\mathbf{f}^{-1}}(\mathbf{y}_{0})
        =\bigl[J_{\mathbf{f}}(\mathbf{x}_{0})\bigr]^{-1}\;}
\]
For $n=1$ the matrix is $1\times1$ and this is the previous box.
The condition $\det J\neq0$ is the exact analogue of $f'\neq0$: it says the
linear approximation at the point is invertible, and the theorem promises
that the map itself inherits that property nearby.

Note that the two named theorems of this chapter are the same statement
twice. Applying the Implicit Function Theorem to
$\mathbf{F}(\mathbf{x},\mathbf{y})=\mathbf{f}(\mathbf{x})-\mathbf{y}$
--- solving for $\mathbf{x}$ in terms of $\mathbf{y}$ --- has the condition
$\det J_{\mathbf{f}}\neq0$ and delivers exactly the inverse.
\end{rulebox}

\begin{rulebox}{Local, and only local}
Both theorems are \textbf{local}, and both directions of the obvious
over-reading are false.

\emph{A non-vanishing derivative everywhere does not give a global
inverse.} The map
$\mathbf{f}(x,y)=(\mathrm{e}^{x}\cos y,\ \mathrm{e}^{x}\sin y)$ has
$\det J=\mathrm{e}^{2x}>0$ at every point, yet
$\mathbf{f}(x,y)=\mathbf{f}(x,y+2\pi)$: it is invertible near each point and
nowhere invertible globally.

\emph{A vanishing derivative does not rule out an inverse.} $f(x)=x^{3}$ has
$f'(0)=0$, and is nevertheless a bijection of $\R$. What fails is not
invertibility but \emph{differentiability of the inverse}: $f^{-1}(y)=
y^{1/3}$ has a vertical tangent at $0$, exactly as the reciprocal formula
predicts.

So the condition is sufficient for a local inverse, necessary for a
\emph{differentiable} one, and silent about global behaviour.
\end{rulebox}

\begin{rulebox}{Homogeneous Functions}
$f:\R^{n}\to\R$ is \textbf{homogeneous of degree $k$} if for all $t>0$ and
all $\mathbf{x}$ in the domain,
\[
  f(t\mathbf{x}) \;=\; t^{k}f(\mathbf{x}) .
\]
Constant returns to scale (CRS) means $k=1$.
\end{rulebox}

\begin{rulebox}{Euler's Theorem}
If $f$ is differentiable and homogeneous of degree $k$ on an open cone,
then
\[
  \sum_{i=1}^{n} x_{i}\pd{f}{x_{i}} \;=\; k\,f(\mathbf{x}) .
\]
Two-variable case: $xf_{x}+yf_{y}=kf$. \textbf{Verification rule.} Always
compute the LHS $xf_{x}+yf_{y}$ and the RHS $kf$ \emph{separately}, then
check that they agree.
\end{rulebox}

\section{Worked Examples}

\begin{worked}{$F(x,y)=x^{2}y+y^{3}-10=0$ at $(1,2)$: does it define $y(x)$, and what is $y'$? Then verify Euler's theorem on $f(K,L)=K^{0.3}L^{0.7}$}
\wstep{Step 1. Check the point lies on the curve.}
$1^{2}\cdot2+2^{3}-10=2+8-10=0$. \checkmark

\wstep{Step 2. Test the hypothesis of the Implicit Function Theorem.}
The theorem needs $F_{y}\neq0$ at the point:
\[
  F_{y}=x^{2}+3y^{2},
  \qquad
  F_{y}(1,2)=1+12=13\neq0 .
\]
So there is a neighbourhood of $x=1$ on which the equation defines a unique
differentiable function $y(x)$ with $y(1)=2$.
\wtrap{The condition is on $F_{y}$ --- the derivative with respect to the
variable you are \textbf{solving for}. If you want $x$ as a function of $y$
instead, the condition is $F_{x}\neq0$. Testing the wrong partial can
license a conclusion that is false: at a point with $F_{y}=0$ the curve is
vertical and no function $y(x)$ exists there.}

\wstep{Step 3. Apply the formula.}
\[
  F_{x}=2xy,
  \qquad
  \dd{y}{x}=-\frac{F_{x}}{F_{y}}=-\frac{2xy}{x^{2}+3y^{2}},
  \qquad
  y'(1)=-\frac{4}{13}\approx-0.308 .
\]
\wtrap{The minus sign is part of the formula, and it is dropped constantly.
It comes from the chain rule applied to the identity
$F\bigl(x,y(x)\bigr)=0$: $F_{x}+F_{y}y'=0$, so $y'=-F_{x}/F_{y}$. Derive it
once and the sign stops being something to remember.}

\wstep{Step 4. Check by differentiating implicitly, from scratch.}
\[
  \frac{\mathrm{d}}{\mathrm{d}x}\bigl[x^{2}y+y^{3}-10\bigr]=0
  \;\Longrightarrow\;
  2xy+x^{2}y'+3y^{2}y'=0
  \;\Longrightarrow\;
  y'=\frac{-2xy}{x^{2}+3y^{2}} . \checkmark
\]
At $(1,2)$: $4+y'+12y'=0$, so $y'=-4/13$. The theorem's formula is not a new
technique --- it is Chapter~5 done once in general.

\wstep{Step 5. A numerical confirmation.}
Take $x=1.01$ and solve $1.0201y+y^{3}=10$ numerically: $y\approx1.996908$.
The linear prediction is $2-0.30769\times0.01=1.996923$. \checkmark\ The two
agree to four decimals over a step of $0.01$; the residual $1.5\times10^{-5}$
is the second-order term, and it should shrink by a factor of $100$ if the
step is halved twice --- which is the signature of a first-order
approximation.

\wstep{Step 6. Where this is used: comparative statics without solving.}
An equilibrium condition $F(p,\alpha)=0$ relates a price $p$ to a parameter
$\alpha$, usually with no closed-form solution for $p$. The theorem still
gives the response:
\[
  \dd{p}{\alpha}=-\frac{F_{\alpha}}{F_{p}} ,
\]
and its \emph{sign} --- which is often all that is asked --- follows from
the signs of two partials. This is the engine of every comparative-statics
result in economics.

\wstep{Step 7. Homogeneity: check the degree by scaling.}
$f$ is homogeneous of degree $r$ if $f(tK,tL)=t^{r}f(K,L)$ for all $t>0$.
For $f=K^{0.3}L^{0.7}$:
\[
  f(tK,tL)=(tK)^{0.3}(tL)^{0.7}=t^{0.3+0.7}K^{0.3}L^{0.7}=t\,f(K,L),
\]
so $r=1$: constant returns to scale. For Cobb--Douglas the degree is simply
the \emph{sum of the exponents}.
\wtrap{Scale \emph{every} argument by the same $t$, and require the identity
for all $t>0$. $f(K,L)=K^{0.3}L^{0.7}+5$ is \textbf{not} homogeneous --- the
constant does not scale --- and Euler's theorem fails for it. Homogeneity is
a strong property, not a formality.}

\wstep{Step 8. Euler's theorem, verified.}
For $f$ homogeneous of degree $r$,
\[
  K f_{K}+L f_{L}=r\,f .
\]
Here $f_{K}=0.3K^{-0.7}L^{0.7}$ and $f_{L}=0.7K^{0.3}L^{-0.3}$, so
\[
  Kf_{K}+Lf_{L}=0.3K^{0.3}L^{0.7}+0.7K^{0.3}L^{0.7}
  =(0.3+0.7)f=1\cdot f . \checkmark
\]

\wstep{Step 9. Read the economics --- this is the product-exhaustion theorem.}
Under constant returns to scale ($r=1$), paying each factor its marginal
product exhausts output exactly:
\[
  \underbrace{K\cdot MP_{K}}_{\text{capital's share}}
  +\underbrace{L\cdot MP_{L}}_{\text{labour's share}}
  =\underbrace{Q}_{\text{total output}} .
\]
Nothing is left over and nothing is missing. With increasing returns
($r>1$) the factor payments would exceed output, which is why perfect
competition and increasing returns cannot coexist --- a theorem of this
chapter, not a modelling assumption.
\end{worked}

\begin{worked}{The inverse theorem both ways: $f(x)=x^{5}+2x+1$ at $y=4$, and $\mathbf{f}(x,y)=(x^{2}-y^{2},\,2xy)$ at $(1,1)$}
\wstep{Step 1. One variable: check the hypothesis before anything else.}
\[
  f(x)=x^{5}+2x+1,
  \qquad
  f'(x)=5x^{4}+2\ \ge 2>0\quad\text{for every }x .
\]
The derivative never vanishes, so $f$ is strictly increasing and invertible
on all of $\R$, and its inverse is differentiable everywhere.

\wstep{Step 2. Locate the point --- in the \emph{domain}.}
We are asked about $y=4$, so we need the $x$ with $f(x)=4$:
\[
  x^{5}+2x+1=4
  \quad\Longrightarrow\quad
  x=1\ \ (1+2+1=4). \checkmark
\]
\wtrap{The formula needs $x_{0}=f^{-1}(y_{0})$, not $y_{0}$. Substituting
$4$ into $f'$ gives $5\cdot4^{4}+2=1282$ and the wrong answer $1/1282$.
Find the pre-image first; in an examination it is always a value you can
guess and verify, because no one can invert a quintic either.}

\wstep{Step 3. Apply the formula.}
\[
  \bigl(f^{-1}\bigr)'(4)=\frac{1}{f'(1)}=\frac{1}{5+2}=\boxed{\tfrac17}
\]
Note what was \emph{not} done: $f^{-1}$ was never computed. It cannot be ---
a general quintic has no closed-form root --- and that is exactly the point
of the theorem. It hands you a derivative of a function you cannot write
down.

\wstep{Step 4. Sanity-check it numerically.}
$f(1.01)=1.01^{5}+2.02+1\approx4.0710$, so a step of $0.0710$ in $y$ answers
to a step of $0.01$ in $x$:
\[
  \frac{\Delta x}{\Delta y}\approx\frac{0.01}{0.0710}\approx0.141
  \approx\tfrac17=0.1429 . \checkmark
\]

\wstep{Step 5. Several variables: build the Jacobian.}
For $\mathbf{f}(x,y)=(x^{2}-y^{2},\ 2xy)$,
\[
  J_{\mathbf{f}}=\begin{pmatrix}2x&-2y\\ 2y&2x\end{pmatrix},
  \qquad
  \det J_{\mathbf{f}}=4x^{2}+4y^{2}=4(x^{2}+y^{2}).
\]
This vanishes only at the origin, so the theorem applies at every other
point.

\wstep{Step 6. Evaluate at $(1,1)$ and invert the matrix.}
$\mathbf{f}(1,1)=(1-1,\,2)=(0,2)$, and
\[
  J_{\mathbf{f}}(1,1)=\begin{pmatrix}2&-2\\2&2\end{pmatrix},
  \qquad
  \det=4+4=8\neq0 .
\]
So a local inverse exists near $(0,2)$, and its Jacobian is the matrix
inverse:
\[
  J_{\mathbf{f}^{-1}}(0,2)
  =\begin{pmatrix}2&-2\\2&2\end{pmatrix}^{-1}
  =\frac18\begin{pmatrix}2&2\\-2&2\end{pmatrix}
  =\boxed{\begin{pmatrix}1/4&1/4\\-1/4&1/4\end{pmatrix}}
\]
Check by multiplying: $\begin{pmatrix}2&-2\\2&2\end{pmatrix}
\begin{pmatrix}1/4&1/4\\-1/4&1/4\end{pmatrix}
=\begin{pmatrix}1&0\\0&1\end{pmatrix}$. \checkmark
\wtrap{Invert the \emph{matrix}; do not take reciprocals entry by entry.
$\bigl[J\bigr]^{-1}$ is a matrix inverse, and the $2\times2$ rule swaps the
diagonal and negates the off-diagonal before dividing by the determinant.
Entrywise reciprocals would give $\begin{pmatrix}1/2&-1/2\\1/2&1/2
\end{pmatrix}$, which is not the inverse of anything relevant.}

\wstep{Step 7. Now show the inverse is only local.}
$\det J\neq0$ holds at every point except the origin, which invites the
conclusion that $\mathbf{f}$ is invertible on $\R^{2}\setminus\{0\}$. It is
not:
\[
  \mathbf{f}(-1,-1)=\bigl((-1)^{2}-(-1)^{2},\ 2(-1)(-1)\bigr)=(0,2)
  =\mathbf{f}(1,1).
\]
Two different points share an image, so no global inverse exists. In general
$\mathbf{f}(-x,-y)=\mathbf{f}(x,y)$: the map is two-to-one everywhere.
\wtrap{``$\det J\neq0$ everywhere'' gives a local inverse \emph{at each
point} and says nothing about fitting them together into one function. This
is the chapter's sharpest trap, and the map above is the standard
counterexample --- it is squaring, in disguise, and squaring is two-to-one.}

\wstep{Step 8. What the two halves have in common.}
\begin{center}\footnotesize
\begin{tabular}{lll}
  \hline
   & one variable & several variables\\
  \hline
  hypothesis & $f'(x_{0})\neq0$ & $\det J(\mathbf{x}_{0})\neq0$\\
  conclusion & local inverse, differentiable & the same\\
  derivative & $1/f'(x_{0})$ & $\bigl[J(\mathbf{x}_{0})\bigr]^{-1}$\\
  evaluated at & $y_{0}=f(x_{0})$ & $\mathbf{y}_{0}=\mathbf{f}(\mathbf{x}_{0})$\\
  says nothing about & global behaviour & global behaviour\\
  \hline
\end{tabular}
\end{center}
The right-hand column is the left-hand one with ``non-zero number''
replaced by ``invertible matrix'' --- which is what $\det\neq0$ means, and
why Volume~3 had to come first.
\end{worked}

\section{Practice Problems}

\begin{problems}
  \item Use IFT to find $\dd{y}{x}$ for $x^{3}+y^{3}=6xy$.
  \item Verify homogeneity and find the degree of
        $f(x,y)=x^{4}+5x^{2}y^{2}+y^{4}$.
  \item Compute the Jacobian and its determinant for
        $\mathbf{f}(x,y)=(3x-2y,\;x+4y)$.
  \item For $U(x,y)=x+2\ln y$, find the signed slope and unsigned MRS
        using IFT.
  \item Verify Euler's theorem for $f(x,y)=x^{2}+5xy-2y^{2}$ (degree $2$).
  \item Use IFT to find $\dd{y}{x}$ for $x^{2}y+xy^{3}=10$.
  \item Verify homogeneity and find the degree of
        $f(x,y)=\dfrac{x^{3}+y^{3}}{x-y}$.
  \item Compute the Jacobian and its determinant for
        $\mathbf{f}(x,y)=(x^{2}+y^{2},\;x^{2}-y^{2})$.
  \item For $U(x,y)=(x+1)(y+2)$, find the signed slope and unsigned MRS.
  \item Verify Euler's theorem for $Q(K,L)=4K^{1/4}L^{3/4}$ (degree $1$).
  \item Use IFT to find $\dd{y}{x}$ for $\mathrm{e}^{x+y}=x+y+1$.
  \item Verify homogeneity and find the degree of $f(x,y)=\sqrt{xy}$.
  \item Compute the Jacobian and its determinant for
        $\mathbf{f}(x,y)=(x\mathrm{e}^{y},\;y\mathrm{e}^{x})$.
  \item For $U(x,y)=x^{a}y^{b}$ ($a,b>0$), find the signed slope and MRS
        via IFT.
  \item Verify Euler's theorem for $Q(K,L)=K^{1/2}L^{1/2}+K+L$
        (degree $1$).
  \item Use IFT to find $\dd{y}{x}$ for $\ln(x^{2}+y)=x-y$.
  \item Verify homogeneity for
        $f(x,y,z)=\dfrac{xy+yz+zx}{x+y+z}$ and find the degree.
  \item Compute the Jacobian and its determinant for the cylindrical map
        \[
          \mathbf{f}(\rho,\theta,z)=(\rho\cos\theta,\;\rho\sin\theta,\;z).
        \]
  \item For the Leontief $U(x,y)=\min\{2x,3y\}$, discuss MRS on each
        smooth piece and at the kink.
  \item Verify Euler's theorem for the harmonic mean
        $f(x,y)=(x^{-1}+y^{-1})^{-1}$ (degree $1$).
  \item Use IFT to find $\dd{y}{x}$ for $x\cos y+y\cos x=1$.
  \item Verify homogeneity and find the degree of
        $f(x,y)=x\ln(y/x)+y\ln(x/y)$.
  \item Compute the Jacobian of
        $\mathbf{f}(x,y)=(x+y,\;xy,\;x^{2}+y^{2})$ and identify where the
        rank drops.
  \item For $Q(K,L)=5K+3L$, find the MRTS.
  \item Verify Euler's theorem for
        $Q(K,L)=\bigl(\alpha K^{\rho}+(1-\alpha)L^{\rho}\bigr)^{1/\rho}$
        (CES, degree $1$).
  \item Use IFT to find $\dd{y}{x}$ for $x^{y}=y^{x}$ with $x,y>0$.
  \item Verify whether $f(x,y)=x^{2}+xy+y$ is homogeneous; if so, find the
        degree.
  \item Compute the Jacobian and determinant of
        $\mathbf{f}(x,y)=(\mathrm{e}^{x}\cos y,\;\mathrm{e}^{x}\sin y)$.
        Interpret.
  \item For $Q(K,L)=K^{2/3}L^{1/3}$, find the MRTS via IFT.
  \item Verify Euler's theorem for $f(x,y,z)=x^{3}+y^{3}+z^{3}+3xyz$
        (degree $3$).
  \item Use IFT to find $\dd{y}{x}$ for $\sqrt{x}+\sqrt{y}+\sqrt{xy}=5$.
  \item Verify homogeneity and find the degree of
        $f(x,y)=\dfrac{x^{2}+y^{2}}{xy}$.
  % VERIFY: right-hand entries of the spherical map obscured by a printer's
  % page-turn mark in the scan; reconstructed from the standard spherical
  % parametrization, and confirmed by the worked solution below.
  \item Compute the Jacobian and determinant of the spherical map
        \[
          \mathbf{f}(\rho,\theta,\phi)=(\rho\sin\phi\cos\theta,\;
          \rho\sin\phi\sin\theta,\;\rho\cos\phi).
        \]
  \item For $U(x,y)=ax+by$ ($a,b>0$, perfect substitutes), find the MRS.
  \item Verify Euler's theorem for $f(x,y)=\dfrac{x^{3}+y^{3}}{x+y}$
        (degree $2$).
  \item Use IFT to find $\pd{z}{x}$ and $\pd{z}{y}$ for
        $F(x,y,z)=x^{2}+y^{2}+z^{2}-R^{2}=0$.
  \item Verify homogeneity of $f(x,y)=x\mathrm{e}^{y/x}$ and find the
        degree.
  \item Compute the Jacobian and determinant of
        $\mathbf{f}(u,v)=(u^{2}-v^{2},\;2uv)$.
  \item For $U(x,y)=\alpha\ln x+\beta\ln y$, find MRS via IFT.
  \item Verify Euler's theorem for $f(x,y)=\sqrt{x^{2}+y^{2}}$
        (degree $1$).
  \item Use IFT to find $\dd{y}{x}$ for $\sin(xy)+\cos(x+y)=1/2$.
  \item Decide whether $f(x,y)=x^{2}+y^{3}$ is homogeneous, with
        justification.
  \item Compute the Jacobian determinant of the polar map
        $\mathbf{f}(r,\theta)=(r\cos\theta,\;r\sin\theta)$ and connect to
        the area element.
  \item For $Q(K,L)=(K^{\rho}+L^{\rho})^{1/\rho}$ ($\rho\neq 0$), find the
        MRTS.
  \item Verify Euler's theorem for the Leontief-like
        $f(x,y)=\min\{x,y\}$ in the interior of each smooth region.
  \item Use IFT to find $\pd{y}{x_{1}}$ and $\pd{y}{x_{2}}$ at any point
        where $F(x_{1},x_{2},y)=x_{1}^{2}+x_{2}y+y^{3}-7=0$, $F_{y}\neq 0$.
  \item Verify homogeneity and find the degree of
        $f(x,y,z)=\dfrac{x^{2}y+y^{2}z+z^{2}x}{xyz}$.
  \item Compute the Jacobian determinant of
        $\mathbf{f}(x,y)=(x+\sin y,\;y+\sin x)$ and find where it
        vanishes.
  \item For $U(x,y)=(x^{\rho}+y^{\rho})^{1/\rho}$ (CES utility), find the
        MRS.
  \item Verify Euler's theorem (3-variable form) for
        $f(x,y,z)=xy+yz+zx$ (degree $2$).
  \item Use IFT to verify that on the unit circle $x^{2}+y^{2}=1$ the
        slope at $(\sqrt{3}/2,1/2)$ equals $-\sqrt{3}$.
  \item Show that $f(x,y)=(x^{1/2}+y^{1/2})^{2}$ is homogeneous of
        degree $1$, then verify Euler's theorem (LHS and RHS separately).
  \item Compute the Jacobian determinant of
        $\mathbf{f}(x,y)=(x/y,\;y/x)$.
  \item For the translog $\ln Q = \alpha_{0}+\alpha_{K}\ln K+\alpha_{L}\ln L$,
        find the MRTS in $(K,L)$ space.
  \item Verify Euler's theorem for the production function
        $Q(K,L,M)=K^{1/3}L^{1/3}M^{1/3}$ (three-factor CRS Cobb--Douglas).
  \item Compute the Jacobian and Jacobian determinant of
        $\mathbf{f}(x,y,z)=(x+y+z,\;xy+yz+zx,\;xyz)$, and explain the link
        to elementary symmetric polynomials.
  \item Show that $f(x)=x^{5}+2x+1$ is invertible on $\R$, and compute
        $\bigl(f^{-1}\bigr)'(4)$ without inverting $f$.
  \item For $f(x)=x^{3}+3x-2$, compute $\bigl(f^{-1}\bigr)'(2)$.
  \item For $f(x)=\mathrm{e}^{x}+x$, compute $\bigl(f^{-1}\bigr)'(1)$.
  \item For $f(x)=x+\ln x$ on $x>0$, compute $\bigl(f^{-1}\bigr)'(1)$.
  \item Derive $\bigl(f^{-1}\bigr)'(y)=1/f'\bigl(f^{-1}(y)\bigr)$ by
        differentiating the identity $f\bigl(f^{-1}(y)\bigr)=y$.
  \item $f(x)=x^{3}$ is a bijection of $\R$, yet $f'(0)=0$. Which
        conclusion of the Inverse Function Theorem fails at $0$, and which
        survives?
  \item Where does the Inverse Function Theorem apply to $f(x)=x^{2}$, and
        what is the largest interval containing $x=2$ on which $f$ is
        invertible?
  \item Give an interval containing $\pi/3$ on which $\sin$ is invertible,
        and compute $\bigl(\sin^{-1}\bigr)'\bigl(\tfrac{\sqrt3}{2}\bigr)$
        two ways: by the theorem, and from the known derivative of
        $\arcsin$.
  \item For $\mathbf{f}(x,y)=(x+y,\;x-y)$: compute $J$ and $\det J$, invert
        the map explicitly, and check that $J_{\mathbf{f}^{-1}}=J^{-1}$.
  \item For $\mathbf{f}(x,y)=(x^{2}-y^{2},\;2xy)$, compute $\det J$ at
        $(1,2)$ and say whether a local inverse exists there.
  \item Show that the map of the previous problem is \emph{not} globally
        injective, despite $\det J\neq0$ off the origin.
  \item For $\mathbf{f}(x,y)=(\mathrm{e}^{x}\cos y,\;\mathrm{e}^{x}\sin y)$,
        compute $\det J$ and reconcile ``non-zero everywhere'' with the fact
        that $\mathbf{f}$ is not injective.
  \item For the polar map
        $\mathbf{f}(r,\theta)=(r\cos\theta,\;r\sin\theta)$, give $\det J$
        and the set of points where the theorem gives nothing.
  \item For $\mathbf{f}(x,y)=(x+y^{3},\;y+x^{3})$, verify the hypothesis at
        the origin and give $J_{\mathbf{f}^{-1}}$ there.
  \item For $\mathbf{f}(x,y,z)=(x,\;x+y,\;x+y+z)$: compute $\det J$, write
        the inverse map explicitly, and confirm
        $J_{\mathbf{f}^{-1}}=J^{-1}$.
  \item A demand system $\mathbf{q}=\mathbf{f}(\mathbf{p})$ has
        $\det J_{\mathbf{f}}(\mathbf{p}_{0})\neq0$. What does the Inverse
        Function Theorem permit an economist to say, and what does it not?
\end{problems}

\section{Solutions}

\begin{solutions}
  \item $F=x^{3}+y^{3}-6xy$. $F_{x}=3x^{2}-6y$, $F_{y}=3y^{2}-6x$.
        \[
          \dd{y}{x} = -\frac{3x^{2}-6y}{3y^{2}-6x} = \frac{2y-x^{2}}{y^{2}-2x} .
        \]

  \item $f(tx,ty)=(tx)^{4}+5(tx)^{2}(ty)^{2}+(ty)^{4}
        =t^{4}(x^{4}+5x^{2}y^{2}+y^{4})=t^{4}f(x,y)$. \textbf{Degree}
        $k=4$.

  \item
        \[
          J=\begin{pmatrix} 3 & -2\\ 1 & 4\end{pmatrix},\qquad
          \det J=3\cdot 4-(-2)\cdot 1=14 .
        \]
        Globally invertible; areas scale by $14$.

  \item $U_{x}=1$, $U_{y}=2/y$. By IFT $\dd{y}{x}=-U_{x}/U_{y}=-y/2$.
        \textbf{MRS} $=y/2$.

  \item $f_{x}=2x+5y$, $f_{y}=5x-4y$.
        \[
          \mathrm{LHS}=x(2x+5y)+y(5x-4y)=2x^{2}+10xy-4y^{2} .
        \]
        \[
          \mathrm{RHS}=2f=2x^{2}+10xy-4y^{2} .
        \]
        $\mathrm{LHS}=\mathrm{RHS}$. \checkmark

  \item $F=x^{2}y+xy^{3}-10$. $F_{x}=2xy+y^{3}$, $F_{y}=x^{2}+3xy^{2}$.
        \[
          \dd{y}{x} = -\frac{2xy+y^{3}}{x^{2}+3xy^{2}}
            = -\frac{y(2x+y^{2})}{x(x+3y^{2})} .
        \]

  \item
        \[
          f(tx,ty)=\frac{(tx)^{3}+(ty)^{3}}{tx-ty}
            =\frac{t^{3}(x^{3}+y^{3})}{t(x-y)}=t^{2}f(x,y) .
        \]
        \textbf{Degree} $k=2$.

  \item
        \[
          J=\begin{pmatrix} 2x & 2y\\ 2x & -2y\end{pmatrix},\qquad
          \det J=(2x)(-2y)-(2y)(2x)=-8xy .
        \]
        Locally invertible whenever $x\neq 0$ and $y\neq 0$.

  \item $U_{x}=y+2$, $U_{y}=x+1$. $\dd{y}{x}=-(y+2)/(x+1)$.
        \textbf{MRS} $=(y+2)/(x+1)$.

  \item $Q_{K}=K^{-3/4}L^{3/4}$, $Q_{L}=3K^{1/4}L^{-1/4}$.
        \[
          \mathrm{LHS}=K\cdot K^{-3/4}L^{3/4}+L\cdot 3K^{1/4}L^{-1/4}
            =K^{1/4}L^{3/4}+3K^{1/4}L^{3/4}=4K^{1/4}L^{3/4} .
        \]
        \[
          \mathrm{RHS}=1\cdot Q=4K^{1/4}L^{3/4} .
        \]
        $\mathrm{LHS}=\mathrm{RHS}$. \checkmark

  \item $F=\mathrm{e}^{x+y}-x-y-1$. $F_{x}=\mathrm{e}^{x+y}-1=F_{y}$.
        \[
          \dd{y}{x} = -\frac{F_{x}}{F_{y}} = -1 .
        \]
        Valid away from $\mathrm{e}^{x+y}=1$, i.e., $x+y\neq 0$. (The
        implicit relation forces $x+y$ to be locally constant on any
        branch.)

  \item $f(tx,ty)=\sqrt{(tx)(ty)}=t\sqrt{xy}=t f(x,y)$ for $t>0$.
        \textbf{Degree} $k=1$.

  \item
        \[
          J=\begin{pmatrix} \mathrm{e}^{y} & x\mathrm{e}^{y}\\
                             y\mathrm{e}^{x} & \mathrm{e}^{x}\end{pmatrix},
          \qquad \det J=\mathrm{e}^{x+y}(1-xy) .
        \]
        Singular along $xy=1$.

  \item $U_{x}=ax^{a-1}y^{b}$, $U_{y}=bx^{a}y^{b-1}$.
        $U_{x}/U_{y}=(a/b)(y/x)$. By IFT $\dd{y}{x}=-(a/b)(y/x)$.
        \textbf{MRS} $=(a/b)(y/x)$.

  \item $Q_{K}=\tfrac12K^{-1/2}L^{1/2}+1$, $Q_{L}=\tfrac12K^{1/2}L^{-1/2}+1$.
        \[
          \mathrm{LHS}=K\bigl(\tfrac12K^{-1/2}L^{1/2}+1\bigr)
            +L\bigl(\tfrac12K^{1/2}L^{-1/2}+1\bigr)=K^{1/2}L^{1/2}+K+L .
        \]
        \[
          \mathrm{RHS}=1\cdot Q=K^{1/2}L^{1/2}+K+L .
        \]
        $\mathrm{LHS}=\mathrm{RHS}$. \checkmark

  \item $F=\ln(x^{2}+y)-x+y$. $F_{x}=\dfrac{2x}{x^{2}+y}-1$,
        $F_{y}=\dfrac{1}{x^{2}+y}+1$.
        \[
          \dd{y}{x} = -\frac{2x/(x^{2}+y)-1}{1/(x^{2}+y)+1}
            = -\frac{2x-(x^{2}+y)}{1+x^{2}+y}
            = \frac{x^{2}+y-2x}{1+x^{2}+y} .
        \]

  \item
        \[
          f(tx,ty,tz)=\frac{t^{2}(xy+yz+zx)}{t(x+y+z)}=t f(x,y,z) .
        \]
        \textbf{Degree} $k=1$.

  \item
        \[
          J=\begin{pmatrix}
            \cos\theta & -\rho\sin\theta & 0\\
            \sin\theta & \rho\cos\theta & 0\\
            0 & 0 & 1
          \end{pmatrix} .
        \]
        Expanding along the third row:
        $\det J = 1\cdot(\rho\cos^{2}\theta+\rho\sin^{2}\theta)=\rho$.
        Volume element $\mathrm{d}V=\rho\,\mathrm{d}\rho\,\mathrm{d}\theta\,\mathrm{d}z$.

  \item Leontief, non-differentiable on $\{2x=3y\}$. On $\{2x<3y\}$:
        $U=2x$, so $U_{x}=2$, $U_{y}=0$; signed slope $\dd{y}{x}$ is
        undefined (level curve is vertical), MRS is infinite: $x$ is
        the binding factor. On $\{2x>3y\}$: $U=3y$, $U_{x}=0$, $U_{y}=3$;
        signed slope $=0$, MRS $=0$: $y$ is binding. At the kink the
        two-sided derivative does not exist.

  \item $f=(x^{-1}+y^{-1})^{-1}=xy/(x+y)$. $f_{x}=y^{2}/(x+y)^{2}$,
        $f_{y}=x^{2}/(x+y)^{2}$.
        \[
          \mathrm{LHS}=x\cdot\frac{y^{2}}{(x+y)^{2}}
            +y\cdot\frac{x^{2}}{(x+y)^{2}}
            =\frac{xy(y+x)}{(x+y)^{2}}=\frac{xy}{x+y} .
        \]
        \[
          \mathrm{RHS}=1\cdot f=\frac{xy}{x+y} .
        \]
        $\mathrm{LHS}=\mathrm{RHS}$. \checkmark

  \item $F=x\cos y+y\cos x-1$. $F_{x}=\cos y-y\sin x$,
        $F_{y}=-x\sin y+\cos x$.
        \[
          \dd{y}{x} = -\frac{\cos y-y\sin x}{\cos x-x\sin y}
            = \frac{y\sin x-\cos y}{\cos x-x\sin y} .
        \]

  \item
        \[
          f(tx,ty)=(tx)\ln(ty/tx)+(ty)\ln(tx/ty)
            =tx\ln(y/x)+ty\ln(x/y)=t f(x,y)
        \]
        (the $t$ inside the log cancels). \textbf{Degree} $k=1$.

  \item
        \[
          J=\begin{pmatrix} 1 & 1\\ y & x\\ 2x & 2y\end{pmatrix}
            \in\R^{3\times 2} .
        \]
        Not square, so there is no determinant. The top $2\times 2$
        minor is $x-y$, which vanishes precisely on $x=y$; the rank
        drops to $1$ along that diagonal (the other minors also vanish
        there).

  \item Linear: $Q_{K}=5$, $Q_{L}=3$. $\dd{L}{K}=-5/3$. \textbf{MRTS}
        $=5/3$ (perfect substitutes with constant exchange rate).

  \item Let $S=\alpha K^{\rho}+(1-\alpha)L^{\rho}$, so $Q=S^{1/\rho}$.
        \[
          Q_{K}=\frac{1}{\rho}S^{1/\rho-1}\cdot\alpha\rho K^{\rho-1}
            =\alpha S^{1/\rho-1}K^{\rho-1},\qquad
          Q_{L}=(1-\alpha)S^{1/\rho-1}L^{\rho-1} .
        \]
        \[
          \mathrm{LHS}=K\alpha S^{1/\rho-1}K^{\rho-1}
            +L(1-\alpha)S^{1/\rho-1}L^{\rho-1}
            =S^{1/\rho-1}\bigl(\alpha K^{\rho}+(1-\alpha)L^{\rho}\bigr)
            =S^{1/\rho} .
        \]
        \[
          \mathrm{RHS}=1\cdot Q=S^{1/\rho} .
        \]
        $\mathrm{LHS}=\mathrm{RHS}$. \checkmark

  \item Take logs: $y\ln x-x\ln y=0$. Set $F=y\ln x-x\ln y$.
        \[
          F_{x}=\frac{y}{x}-\ln y,\qquad F_{y}=\ln x-\frac{x}{y} .
        \]
        \[
          \dd{y}{x} = -\frac{y/x-\ln y}{\ln x-x/y}
            = \frac{\ln y-y/x}{\ln x-x/y}
            = \frac{y(x\ln y-y)}{x(y\ln x-x)} .
        \]

  \item Test: $f(tx,ty)=(tx)^{2}+(tx)(ty)+(ty)=t^{2}x^{2}+t^{2}xy+ty$.
        This cannot be written as $t^{k}(x^{2}+xy+y)$ for a single $k$:
        the first two terms scale by $t^{2}$, the last by $t$.
        \textbf{Not homogeneous.}

  \item $f_{1}=\mathrm{e}^{x}\cos y$, $f_{2}=\mathrm{e}^{x}\sin y$.
        \[
          J=\begin{pmatrix} \mathrm{e}^{x}\cos y & -\mathrm{e}^{x}\sin y\\
                             \mathrm{e}^{x}\sin y & \mathrm{e}^{x}\cos y
            \end{pmatrix},\qquad
          \det J=\mathrm{e}^{2x}(\cos^{2}y+\sin^{2}y)=\mathrm{e}^{2x} .
        \]
        Always strictly positive: the map is locally invertible
        everywhere. (It is the complex exponential $w=\mathrm{e}^{x+iy}$;
        conformal, orientation-preserving.)

  \item $Q_{K}=\tfrac23K^{-1/3}L^{1/3}$, $Q_{L}=\tfrac13K^{2/3}L^{-2/3}$.
        \[
          \frac{Q_{K}}{Q_{L}}=2\,\frac{L}{K} .
        \]
        $\dd{L}{K}=-2L/K$. \textbf{MRTS} $=2L/K$.

  \item $f_{x}=3x^{2}+3yz$, $f_{y}=3y^{2}+3xz$, $f_{z}=3z^{2}+3xy$.
        \[
          \mathrm{LHS}=x(3x^{2}+3yz)+y(3y^{2}+3xz)+z(3z^{2}+3xy)
            =3(x^{3}+y^{3}+z^{3})+9xyz=3f .
        \]
        \[
          \mathrm{RHS}=3f .
        \]
        $\mathrm{LHS}=\mathrm{RHS}$. \checkmark

  \item $F=\sqrt{x}+\sqrt{y}+\sqrt{xy}-5$. Compute:
        \[
          F_{x}=\frac{1}{2\sqrt{x}}+\frac{y}{2\sqrt{xy}}
            =\frac{1}{2\sqrt{x}}+\frac{1}{2}\sqrt{\frac{y}{x}}
            =\frac{1+\sqrt{y}}{2\sqrt{x}} .
        \]
        By symmetry, $F_{y}=\dfrac{1+\sqrt{x}}{2\sqrt{y}}$.
        \[
          \dd{y}{x} = -\frac{F_{x}}{F_{y}}
            = -\frac{\sqrt{y}(1+\sqrt{y})}{\sqrt{x}(1+\sqrt{x})} .
        \]

  \item
        \[
          f(tx,ty)=\frac{t^{2}x^{2}+t^{2}y^{2}}{(tx)(ty)}
            =\frac{t^{2}(x^{2}+y^{2})}{t^{2}xy}=f(x,y) .
        \]
        \textbf{Degree} $k=0$.

  \item
        \[
          J=\begin{pmatrix}
            \sin\phi\cos\theta & -\rho\sin\phi\sin\theta & \rho\cos\phi\cos\theta\\
            \sin\phi\sin\theta & \rho\sin\phi\cos\theta & \rho\cos\phi\sin\theta\\
            \cos\phi & 0 & -\rho\sin\phi
          \end{pmatrix} .
        \]
        Expand along the third row:
        \[
          \det J = \cos\phi\Bigl[(-\rho\sin\phi\sin\theta)(\rho\cos\phi\sin\theta)
            -(\rho\cos\phi\cos\theta)(\rho\sin\phi\cos\theta)\Bigr]
        \]
        \[
          \qquad{}+(-\rho\sin\phi)\Bigl[(\sin\phi\cos\theta)(\rho\sin\phi\cos\theta)
            -(-\rho\sin\phi\sin\theta)(\sin\phi\sin\theta)\Bigr] .
        \]
        First bracket:
        \[
          -\rho^{2}\sin\phi\cos\phi(\sin^{2}\theta+\cos^{2}\theta)
            =-\rho^{2}\sin\phi\cos\phi .
        \]
        Multiplied by $\cos\phi$ this is
        $-\rho^{2}\sin\phi\cos^{2}\phi$. Second bracket:
        $\rho\sin^{2}\phi(\cos^{2}\theta+\sin^{2}\theta)=\rho\sin^{2}\phi$.
        Multiplied by $-\rho\sin\phi$: $-\rho^{2}\sin^{3}\phi$. Sum:
        $-\rho^{2}\sin\phi(\cos^{2}\phi+\sin^{2}\phi)=-\rho^{2}\sin\phi$.
        Thus $|\det J|=\rho^{2}\sin\phi$, the spherical volume element.

  \item $U_{x}=a$, $U_{y}=b$. $\dd{y}{x}=-a/b$. \textbf{MRS} $=a/b$
        (constant, perfect substitutes).

  \item Simplify: $f=\dfrac{(x+y)(x^{2}-xy+y^{2})}{x+y}=x^{2}-xy+y^{2}$.
        Then $f_{x}=2x-y$, $f_{y}=-x+2y$.
        \[
          \mathrm{LHS}=x(2x-y)+y(-x+2y)=2x^{2}-2xy+2y^{2} .
        \]
        \[
          \mathrm{RHS}=2f=2(x^{2}-xy+y^{2})=2x^{2}-2xy+2y^{2} .
        \]
        $\mathrm{LHS}=\mathrm{RHS}$. \checkmark

  \item $F=x^{2}+y^{2}+z^{2}-R^{2}$. $F_{x}=2x$, $F_{y}=2y$, $F_{z}=2z$.
        Provided $z\neq 0$:
        \[
          \pd{z}{x}=-\frac{F_{x}}{F_{z}}=-\frac{x}{z},\qquad
          \pd{z}{y}=-\frac{F_{y}}{F_{z}}=-\frac{y}{z} .
        \]
        (Geometrically, the outward normal on the sphere has direction
        $(x,y,z)/R$.)

  \item $f(tx,ty)=(tx)\mathrm{e}^{(ty)/(tx)}=tx\,\mathrm{e}^{y/x}
        =t\,f(x,y)$. \textbf{Degree} $k=1$.

  \item
        \[
          J=\begin{pmatrix} 2u & -2v\\ 2v & 2u\end{pmatrix},\qquad
          \det J=4u^{2}+4v^{2}=4(u^{2}+v^{2}) .
        \]
        Vanishes only at the origin. (This is the complex squaring
        $w=(u+iv)^{2}$.)

  \item $U_{x}=\alpha/x$, $U_{y}=\beta/y$. $U_{x}/U_{y}=(\alpha/\beta)(y/x)$.
        $\dd{y}{x}=-(\alpha/\beta)(y/x)$. \textbf{MRS}
        $=(\alpha/\beta)(y/x)$.

  \item $f=(x^{2}+y^{2})^{1/2}$. $f_{x}=x/\sqrt{x^{2}+y^{2}}$,
        $f_{y}=y/\sqrt{x^{2}+y^{2}}$.
        \[
          \mathrm{LHS}=\frac{x^{2}}{\sqrt{x^{2}+y^{2}}}
            +\frac{y^{2}}{\sqrt{x^{2}+y^{2}}}
            =\frac{x^{2}+y^{2}}{\sqrt{x^{2}+y^{2}}}=\sqrt{x^{2}+y^{2}} .
        \]
        \[
          \mathrm{RHS}=1\cdot f=\sqrt{x^{2}+y^{2}} .
        \]
        $\mathrm{LHS}=\mathrm{RHS}$. \checkmark

  \item $F=\sin(xy)+\cos(x+y)-1/2$. $F_{x}=y\cos(xy)-\sin(x+y)$,
        $F_{y}=x\cos(xy)-\sin(x+y)$.
        \[
          \dd{y}{x} = -\frac{y\cos(xy)-\sin(x+y)}{x\cos(xy)-\sin(x+y)}
            = \frac{\sin(x+y)-y\cos(xy)}{x\cos(xy)-\sin(x+y)} .
        \]

  \item $f(tx,ty)=t^{2}x^{2}+t^{3}y^{3}$. The two terms scale with
        different powers of $t$, so they cannot be factored as
        $t^{k}(x^{2}+y^{3})$. \textbf{Not homogeneous.}

  \item
        \[
          J=\begin{pmatrix} \cos\theta & -r\sin\theta\\
                             \sin\theta & r\cos\theta\end{pmatrix},\qquad
          \det J=r\cos^{2}\theta+r\sin^{2}\theta=r .
        \]
        Hence $|\det J|=r$, and the area element is
        $\mathrm{d}A=r\,\mathrm{d}r\,\mathrm{d}\theta$, the standard
        polar substitution.

  \item Let $S=K^{\rho}+L^{\rho}$. Then $Q=S^{1/\rho}$,
        $Q_{K}=S^{1/\rho-1}K^{\rho-1}$, $Q_{L}=S^{1/\rho-1}L^{\rho-1}$.
        \[
          \frac{Q_{K}}{Q_{L}}=\Bigl(\frac{K}{L}\Bigr)^{\rho-1}
            =\Bigl(\frac{L}{K}\Bigr)^{1-\rho} .
        \]
        $\dd{L}{K}=-(L/K)^{1-\rho}$. \textbf{MRTS} $=(L/K)^{1-\rho}$.

  \item On the smooth piece $\{x<y\}$: $f=x$, $f_{x}=1$, $f_{y}=0$.
        \[
          \mathrm{LHS}=x\cdot 1+y\cdot 0=x=f,\qquad \mathrm{RHS}=1\cdot f=x .
        \]
        $\mathrm{LHS}=\mathrm{RHS}$. \checkmark On $\{x>y\}$ by symmetry
        $\mathrm{LHS}=y=f=\mathrm{RHS}$. \checkmark Euler's identity
        holds wherever the function is differentiable. (At the ridge
        $x=y$ the function is not differentiable.)

  \item By IFT, $\pd{y}{x_{i}}=-F_{x_{i}}/F_{y}$ with
        $F_{y}=x_{2}+3y^{2}$:
        \[
          \pd{y}{x_{1}}=-\frac{2x_{1}}{x_{2}+3y^{2}},\qquad
          \pd{y}{x_{2}}=-\frac{y}{x_{2}+3y^{2}} .
        \]

  \item Numerator has degree $3$, denominator degree $3$. Concretely
        \[
          f(tx,ty,tz)=\frac{t^{3}(x^{2}y+y^{2}z+z^{2}x)}{t^{3}xyz}=f(x,y,z) .
        \]
        \textbf{Degree} $k=0$.

  \item $J=\begin{pmatrix} 1 & \cos y\\ \cos x & 1\end{pmatrix}$,
        $\det J=1-\cos x\cos y$. Vanishes when $\cos x\cos y=1$, i.e.,
        $\cos x=\cos y=1$ or $\cos x=\cos y=-1$, equivalently
        $(x,y)\in 2\pi\mathbb{Z}\times 2\pi\mathbb{Z}$ or
        $(2\mathbb{Z}+1)\pi\times(2\mathbb{Z}+1)\pi$.

  \item Let $S=x^{\rho}+y^{\rho}$, $U=S^{1/\rho}$.
        $U_{x}=S^{1/\rho-1}x^{\rho-1}$, $U_{y}=S^{1/\rho-1}y^{\rho-1}$.
        \[
          \frac{U_{x}}{U_{y}}=\Bigl(\frac{x}{y}\Bigr)^{\rho-1}
            =\Bigl(\frac{y}{x}\Bigr)^{1-\rho} .
        \]
        $\dd{y}{x}=-(y/x)^{1-\rho}$. \textbf{MRS} $=(y/x)^{1-\rho}$.

  \item $f_{x}=y+z$, $f_{y}=x+z$, $f_{z}=x+y$.
        \[
          \mathrm{LHS}=x(y+z)+y(x+z)+z(x+y)=2(xy+yz+zx)=2f .
        \]
        \[
          \mathrm{RHS}=2f .
        \]
        $\mathrm{LHS}=\mathrm{RHS}$. \checkmark

  \item $F=x^{2}+y^{2}-1$. $F_{x}=2x$, $F_{y}=2y$. $\dd{y}{x}=-x/y$.
        At $(\sqrt{3}/2,1/2)$: $\dd{y}{x}=-(\sqrt{3}/2)/(1/2)=-\sqrt{3}$.
        \checkmark

  \item $f(tx,ty)=\bigl((tx)^{1/2}+(ty)^{1/2}\bigr)^{2}
        =(t^{1/2})^{2}(\sqrt{x}+\sqrt{y})^{2}=t f(x,y)$. \textbf{Degree}
        $k=1$. For Euler's theorem, write $A=\sqrt{x}+\sqrt{y}$,
        $f=A^{2}$. Then $f_{x}=2A\cdot\dfrac{1}{2\sqrt{x}}=A/\sqrt{x}$,
        $f_{y}=A/\sqrt{y}$.
        \[
          \mathrm{LHS}=x\cdot\frac{A}{\sqrt{x}}+y\cdot\frac{A}{\sqrt{y}}
            =A(\sqrt{x}+\sqrt{y})=A^{2}=f .
        \]
        \[
          \mathrm{RHS}=1\cdot f=A^{2} .
        \]
        $\mathrm{LHS}=\mathrm{RHS}$. \checkmark

  \item
        \[
          J=\begin{pmatrix} 1/y & -x/y^{2}\\ -y/x^{2} & 1/x\end{pmatrix},
          \qquad \det J=\frac{1}{xy}-\frac{xy}{x^{2}y^{2}}
            =\frac{1}{xy}-\frac{1}{xy}=0 .
        \]
        The Jacobian is identically singular: because
        $(x/y)\cdot(y/x)=1$, the image lies on the curve $f_{1}f_{2}=1$
        (one-dimensional), so the map is not locally invertible
        anywhere.

  \item From $\ln Q=\alpha_{0}+\alpha_{K}\ln K+\alpha_{L}\ln L$, we get
        $Q=AK^{\alpha_{K}}L^{\alpha_{L}}$ with $A=\mathrm{e}^{\alpha_{0}}$
        (Cobb--Douglas). Then $Q_{K}=A\alpha_{K}K^{\alpha_{K}-1}L^{\alpha_{L}}$,
        $Q_{L}=A\alpha_{L}K^{\alpha_{K}}L^{\alpha_{L}-1}$.
        \[
          \frac{Q_{K}}{Q_{L}}=\frac{\alpha_{K}}{\alpha_{L}}\cdot\frac{L}{K} .
        \]
        $\dd{L}{K}=-(\alpha_{K}/\alpha_{L})(L/K)$. \textbf{MRTS}
        $=(\alpha_{K}/\alpha_{L})(L/K)$.

  \item $Q=K^{1/3}L^{1/3}M^{1/3}$ is homogeneous of degree
        $1/3+1/3+1/3=1$. $Q_{K}=\tfrac13K^{-2/3}L^{1/3}M^{1/3}=Q/(3K)$,
        and similarly $Q_{L}=Q/(3L)$, $Q_{M}=Q/(3M)$.
        \[
          \mathrm{LHS}=K\cdot\frac{Q}{3K}+L\cdot\frac{Q}{3L}
            +M\cdot\frac{Q}{3M}=\frac{Q}{3}+\frac{Q}{3}+\frac{Q}{3}=Q .
        \]
        \[
          \mathrm{RHS}=1\cdot Q=Q .
        \]
        $\mathrm{LHS}=\mathrm{RHS}$. \checkmark

  \item Write $f_{1}=x+y+z$, $f_{2}=xy+yz+zx$, $f_{3}=xyz$ (the
        elementary symmetric polynomials $e_{1},e_{2},e_{3}$). Then
        \[
          J=\begin{pmatrix}
            1 & 1 & 1\\
            y+z & x+z & x+y\\
            yz & xz & xy
          \end{pmatrix} .
        \]
        For the determinant, subtract column $1$ from columns $2$ and
        $3$ to get
        \[
          \det J = \det\begin{pmatrix}
            1 & 0 & 0\\
            y+z & x-y & x-z\\
            yz & z(x-y) & y(x-z)
          \end{pmatrix}
          = 1\cdot\det\begin{pmatrix}
            x-y & x-z\\ z(x-y) & y(x-z)
          \end{pmatrix}
        \]
        \[
          = (x-y)(x-z)\det\begin{pmatrix} 1 & 1\\ z & y\end{pmatrix}
          = (x-y)(x-z)(y-z) .
        \]
        Thus $\det J=-(x-y)(y-z)(z-x)$ (the Vandermonde discriminant, up
        to sign), which vanishes exactly when two of $x,y,z$ coincide:
        the symmetric-polynomial map fails to be locally invertible
        at points with repeated coordinates, consistent with the fact
        that an unordered triple $\{x,y,z\}$ has a unique sorted
        representative.
  \item $f'(x)=5x^{4}+2\ge2>0$ for all $x$, so $f$ is strictly increasing,
        hence injective, and the theorem applies everywhere. Solve
        $x^{5}+2x+1=4$: $x=1$ works ($1+2+1=4$), and monotonicity makes it
        the only root. So
        $\bigl(f^{-1}\bigr)'(4)=1/f'(1)=1/7$.
  \item $f(1)=1+3-2=2$, and $f'(x)=3x^{2}+3$, so $f'(1)=6$ and
        $\bigl(f^{-1}\bigr)'(2)=\tfrac16$. ($f'>0$ everywhere, so the
        inverse is global.)
  \item $f(0)=1+0=1$, and $f'(x)=\mathrm{e}^{x}+1$, so $f'(0)=2$ and
        $\bigl(f^{-1}\bigr)'(1)=\tfrac12$.
  \item $f(1)=1+\ln1=1$, and $f'(x)=1+1/x$, so $f'(1)=2$ and
        $\bigl(f^{-1}\bigr)'(1)=\tfrac12$.
  \item Differentiate both sides with respect to $y$. The left side is a
        composite, so the chain rule gives
        $f'\bigl(f^{-1}(y)\bigr)\cdot\bigl(f^{-1}\bigr)'(y)$; the right side
        gives $1$. Hence
        $\bigl(f^{-1}\bigr)'(y)=1/f'\bigl(f^{-1}(y)\bigr)$, the division
        being legitimate exactly when $f'\neq0$ there.
  \item Invertibility survives: $x^{3}$ is strictly increasing and onto, so
        $f^{-1}(y)=y^{1/3}$ exists globally. What fails is
        \emph{differentiability of the inverse at $0$}:
        $\bigl(y^{1/3}\bigr)'=\tfrac13y^{-2/3}\to\infty$ as $y\to0$. The
        hypothesis $f'\neq0$ is sufficient for a differentiable local
        inverse, not necessary for an inverse.
  \item $f'(x)=2x\neq0$ for $x\neq0$, so the theorem applies at every point
        except the origin. Around $x=2$ the largest interval of
        injectivity is $(0,\infty)$: on it $f$ is strictly increasing, while
        any interval containing $0$ in its interior contains a pair
        $\pm a$ with equal images.
  \item On $(-\pi/2,\pi/2)$ the sine is strictly increasing, and
        $\pi/3$ lies inside. By the theorem, with
        $\sin(\pi/3)=\sqrt3/2$ and $\cos(\pi/3)=\tfrac12$,
        \[
          \bigl(\sin^{-1}\bigr)'\Bigl(\tfrac{\sqrt3}{2}\Bigr)
          =\frac{1}{\cos(\pi/3)}=2 .
        \]
        Directly: $\arcsin'(y)=1/\sqrt{1-y^{2}}
        =1/\sqrt{1-\tfrac34}=1/\tfrac12=2$. \checkmark
  \item $J=\begin{pmatrix}1&1\\1&-1\end{pmatrix}$, $\det J=-2\neq0$
        everywhere. Solving $u=x+y$, $v=x-y$ gives
        $x=\tfrac{u+v}{2}$, $y=\tfrac{u-v}{2}$, so
        $J_{\mathbf{f}^{-1}}
        =\begin{pmatrix}1/2&1/2\\1/2&-1/2\end{pmatrix}$, which is indeed
        $J^{-1}=-\tfrac12\begin{pmatrix}-1&-1\\-1&1\end{pmatrix}$.
        \checkmark\ The map is linear with $\det\neq0$, so here local and
        global invertibility coincide.
  \item $J=\begin{pmatrix}2x&-2y\\2y&2x\end{pmatrix}$ and
        $\det J=4(x^{2}+y^{2})$, so at $(1,2)$ it is $4\cdot5=20\neq0$: a
        differentiable local inverse exists near
        $\mathbf{f}(1,2)=(-3,4)$.
  \item $\mathbf{f}(-x,-y)=\bigl((-x)^{2}-(-y)^{2},\,2(-x)(-y)\bigr)
        =(x^{2}-y^{2},\,2xy)=\mathbf{f}(x,y)$, so every point and its
        negative share an image --- for instance
        $\mathbf{f}(1,1)=\mathbf{f}(-1,-1)=(0,2)$. The map is two-to-one,
        and no global inverse exists. Local invertibility at every point
        does not assemble into global invertibility.
  \item $J=\mathrm{e}^{x}\begin{pmatrix}\cos y&-\sin y\\
        \sin y&\cos y\end{pmatrix}$, so
        $\det J=\mathrm{e}^{2x}(\cos^{2}y+\sin^{2}y)=\mathrm{e}^{2x}>0$
        everywhere. Yet $\mathbf{f}(x,y+2\pi)=\mathbf{f}(x,y)$. There is no
        contradiction: the theorem promises an inverse on \emph{some
        neighbourhood} of each point, and those neighbourhoods (width less
        than $2\pi$ in $y$) never join up into one.
  \item $J=\begin{pmatrix}\cos\theta&-r\sin\theta\\
        \sin\theta&r\cos\theta\end{pmatrix}$, so $\det J=r$. The theorem
        says nothing on the line $r=0$ --- the origin, where every $\theta$
        gives the same point and the map is infinitely many-to-one.
  \item $J=\begin{pmatrix}1&3y^{2}\\3x^{2}&1\end{pmatrix}$, which at
        $(0,0)$ is the identity, with $\det=1\neq0$. Hence a local inverse
        exists near $\mathbf{f}(0,0)=(0,0)$ and
        $J_{\mathbf{f}^{-1}}(0,0)=I^{-1}=I$.
  \item $J=\begin{pmatrix}1&0&0\\1&1&0\\1&1&1\end{pmatrix}$, lower
        triangular, so $\det J=1\neq0$. From $u=x$, $v=x+y$, $w=x+y+z$:
        $x=u$, $y=v-u$, $z=w-v$, giving
        $J_{\mathbf{f}^{-1}}=\begin{pmatrix}1&0&0\\-1&1&0\\
        0&-1&1\end{pmatrix}$. Multiplying the two matrices returns $I$.
        \checkmark
  \item It permits: near $\mathbf{p}_{0}$ the system can in principle be
        solved for prices in terms of quantities, the inverse demand
        $\mathbf{p}(\mathbf{q})$ is differentiable, and its derivative
        matrix is $\bigl[J_{\mathbf{f}}\bigr]^{-1}$ --- so the price
        responses can be computed from the quantity responses without ever
        solving the system. It does not permit: any claim that the same
        quantities cannot arise from a different price vector elsewhere.
        Global identification is a separate question, and $\det J\neq0$ is
        no evidence for it.
\end{solutions}

\section{Additional Problems}

\begin{problems}
  \item Use IFT to find $\dd{y}{x}$ for $x^{2}+2xy+y^{2}=9$.
  \item Verify homogeneity and find the degree of $f(x,y)=x^{3}y+xy^{3}$.
  \item Compute the Jacobian and its determinant for
        $\mathbf{f}(x,y)=(2x+y,\;x-y)$.
  \item For $U(x,y)=x^{2}+y$, find the signed slope and unsigned MRS
        using IFT.
  \item Verify Euler's theorem for $f(x,y)=x^{3}+2x^{2}y-y^{3}$
        (degree $3$).
  \item Use IFT to find $\dd{y}{x}$ for $x^{3}+y^{3}+xy=5$.
  \item Verify homogeneity and find the degree of
        $f(x,y)=\dfrac{x^{4}+y^{4}}{x^{2}+y^{2}}$.
  \item Compute the Jacobian and its determinant for
        $\mathbf{f}(x,y)=(x^{2}y,\;xy^{2})$.
  \item For $U(x,y)=x^{1/3}y^{2/3}$, find the signed slope and MRS via
        IFT.
  \item Verify Euler's theorem for $Q(K,L)=K^{1/2}L^{1/2}$ (degree $1$).
  \item Use IFT to find $\dd{y}{x}$ for $\ln(xy)+x+y=4$.
  \item Decide whether $f(x,y)=x+y^{2}$ is homogeneous, with
        justification.
  \item Compute the Jacobian determinant of
        $\mathbf{f}(x,y,z)=(x+y,\;y+z,\;x+z)$.
  \item For $Q(K,L)=aK+bL$ ($a,b>0$), find the MRTS.
  \item Verify Euler's theorem for $f(x,y,z)=xy+yz+xz$ (degree $2$,
        three-variable).
  \item For $f(x)=x^{7}+x^{3}+x$, compute $\bigl(f^{-1}\bigr)'(3)$.
  \item For $f(x)=2x+\sin x$, compute $\bigl(f^{-1}\bigr)'(0)$.
  \item For $f(x)=\ln x+x^{2}$ on $x>0$, compute $\bigl(f^{-1}\bigr)'(1)$.
  \item For $f(x)=\tan x$ on $(-\pi/2,\pi/2)$, compute
        $\bigl(f^{-1}\bigr)'(1)$ and check it against $\arctan'$.
  \item For $\mathbf{f}(x,y)=(2x+3y,\;x+2y)$, give $\det J$ and
        $J_{\mathbf{f}^{-1}}$, and say why the inverse is global here.
  \item For $\mathbf{f}(x,y)=(x^{3},\;y^{3})$: where does $\det J$ vanish,
        and is the map invertible there anyway?
  \item For $\mathbf{f}(x,y)=(x+\mathrm{e}^{y},\;y+\mathrm{e}^{x})$,
        evaluate $\det J$ at the origin and say what the theorem concludes.
  \item Let $\mathbf{f}(\mathbf{x})=A\mathbf{x}$ with $A$ an $n\times n$
        matrix. Show that $\det J$ is constant, and that local and global
        invertibility coincide.
  \item For $\mathbf{f}(x,y)=(x^{2}+y^{2},\;x^{2}-y^{2})$ on $x,y>0$,
        compute $\det J$ and decide local invertibility.
  \item State in one sentence each what $f'(x_{0})\neq0$ and
        $\det J(\mathbf{x}_{0})\neq0$ do and do not guarantee.
\end{problems}

\section{Additional Solutions}

\begin{solutions}
  \item $F=x^{2}+2xy+y^{2}-9=(x+y)^{2}-9$. $F_{x}=2(x+y)$,
        $F_{y}=2(x+y)$.
        \[
          \dd{y}{x}=-\frac{F_{x}}{F_{y}}=-1 .
        \]
        (The curve is two lines $x+y=\pm 3$, each with slope $-1$.)

  \item $f(tx,ty)=(tx)^{3}(ty)+(tx)(ty)^{3}=t^{4}(x^{3}y+xy^{3})
        =t^{4}f(x,y)$. \textbf{Degree $k=4$.}

  \item
        \[
          J=\begin{pmatrix} 2 & 1\\ 1 & -1\end{pmatrix},\qquad
          \det J=-2-1=-3 .
        \]
        Globally invertible; areas scale by $|{-3}|=3$ (with orientation
        reversal).

  \item $U_{x}=2x$, $U_{y}=1$. By IFT $\dd{y}{x}=-2x$. \textbf{MRS}
        $=2x$ (unsigned).

  \item $f_{x}=3x^{2}+4xy$, $f_{y}=2x^{2}-3y^{2}$.
        \[
          \mathrm{LHS}=x(3x^{2}+4xy)+y(2x^{2}-3y^{2})
          =3x^{3}+4x^{2}y+2x^{2}y-3y^{3}=3x^{3}+6x^{2}y-3y^{3} .
        \]
        \[
          \mathrm{RHS}=3f=3(x^{3}+2x^{2}y-y^{3})
          =3x^{3}+6x^{2}y-3y^{3} .
        \]
        $\mathrm{LHS}=\mathrm{RHS}$. \checkmark

  \item $F=x^{3}+y^{3}+xy-5$. $F_{x}=3x^{2}+y$, $F_{y}=3y^{2}+x$.
        \[
          \dd{y}{x}=-\frac{3x^{2}+y}{3y^{2}+x} .
        \]

  \item $f(tx,ty)=\dfrac{t^{4}(x^{4}+y^{4})}{t^{2}(x^{2}+y^{2})}
        =t^{2}f(x,y)$. \textbf{Degree $k=2$.}

  \item
        \[
          J=\begin{pmatrix} 2xy & x^{2}\\ y^{2} & 2xy\end{pmatrix},\qquad
          \det J=4x^{2}y^{2}-x^{2}y^{2}=3x^{2}y^{2} .
        \]
        Singular when $x=0$ or $y=0$.

  \item $U_{x}=\tfrac13 x^{-2/3}y^{2/3}$,
        $U_{y}=\tfrac23 x^{1/3}y^{-1/3}$.
        $U_{x}/U_{y}=\tfrac12\cdot y/x$. By IFT $\dd{y}{x}=-y/(2x)$.
        \textbf{MRS} $=y/(2x)$.

  \item $Q_{K}=\tfrac12 K^{-1/2}L^{1/2}$,
        $Q_{L}=\tfrac12 K^{1/2}L^{-1/2}$.
        \[
          \mathrm{LHS}=K\cdot\tfrac12 K^{-1/2}L^{1/2}
          +L\cdot\tfrac12 K^{1/2}L^{-1/2}
          =\tfrac12 K^{1/2}L^{1/2}+\tfrac12 K^{1/2}L^{1/2}
          =K^{1/2}L^{1/2} .
        \]
        \[
          \mathrm{RHS}=1\cdot Q=K^{1/2}L^{1/2} .
        \]
        $\mathrm{LHS}=\mathrm{RHS}$. \checkmark

  \item $F=\ln(xy)+x+y-4=\ln x+\ln y+x+y-4$. $F_{x}=1/x+1$,
        $F_{y}=1/y+1$.
        \[
          \dd{y}{x}=-\frac{1/x+1}{1/y+1}=-\frac{y(1+x)}{x(1+y)} .
        \]

  \item $f(tx,ty)=tx+t^{2}y^{2}$. The two terms scale with different
        powers of $t$, so they cannot share a common $t^{k}$ factor.
        \textbf{Not homogeneous.}

  \item
        \[
          J=\begin{pmatrix}
            1 & 1 & 0\\
            0 & 1 & 1\\
            1 & 0 & 1
          \end{pmatrix} .
        \]
        Expanding along the first row:
        $\det J=1\cdot(1\cdot 1-1\cdot 0)-1\cdot(0\cdot 1-1\cdot 1)+0
        =1-(-1)+0=2$.

  \item Linear: $Q_{K}=a$, $Q_{L}=b$. $\dd{L}{K}=-a/b$. \textbf{MRTS}
        $=a/b$ (constant, perfect substitutes).

  \item $f_{x}=y+z$, $f_{y}=x+z$, $f_{z}=x+y$.
        \[
          \mathrm{LHS}=x(y+z)+y(x+z)+z(x+y)=2(xy+yz+xz)=2f .
        \]
        \[
          \mathrm{RHS}=2f .
        \]
        $\mathrm{LHS}=\mathrm{RHS}$. \checkmark
  \item $f(1)=1+1+1=3$ and $f'(x)=7x^{6}+3x^{2}+1$, so $f'(1)=11$ and
        $\bigl(f^{-1}\bigr)'(3)=\tfrac{1}{11}$. ($f'>0$ everywhere, so the
        pre-image is unique.)
  \item $f(0)=0$ and $f'(x)=2+\cos x\ge1>0$, so $f'(0)=3$ and
        $\bigl(f^{-1}\bigr)'(0)=\tfrac13$.
  \item $f(1)=\ln1+1=1$ and $f'(x)=1/x+2x$, so $f'(1)=3$ and
        $\bigl(f^{-1}\bigr)'(1)=\tfrac13$.
  \item $\tan(\pi/4)=1$ and $f'(x)=1/\cos^{2}x$, so
        $f'(\pi/4)=1/(\sqrt2/2)^{2}=2$ and
        $\bigl(f^{-1}\bigr)'(1)=\tfrac12$. Directly:
        $\arctan'(1)=1/(1+1^{2})=\tfrac12$. \checkmark
  \item $J=\begin{pmatrix}2&3\\1&2\end{pmatrix}$ is constant with
        $\det J=4-3=1\neq0$, and
        $J_{\mathbf{f}^{-1}}=J^{-1}
        =\begin{pmatrix}2&-3\\-1&2\end{pmatrix}$. The map is linear, so the
        local inverse supplied by the theorem is the global one:
        $\mathbf{x}=A^{-1}\mathbf{y}$ is defined on all of $\R^{2}$.
  \item $J=\begin{pmatrix}3x^{2}&0\\0&3y^{2}\end{pmatrix}$ and
        $\det J=9x^{2}y^{2}$, which vanishes on both axes. The map is
        nevertheless a bijection of $\R^{2}$, with inverse
        $(u,v)\mapsto(u^{1/3},v^{1/3})$ --- but that inverse is not
        differentiable where $\det J=0$. Same lesson as $x^{3}$ in one
        variable.
  \item $J=\begin{pmatrix}1&\mathrm{e}^{y}\\ \mathrm{e}^{x}&1\end{pmatrix}$,
        so at $(0,0)$ it is
        $\begin{pmatrix}1&1\\1&1\end{pmatrix}$ with $\det J=0$. The theorem
        concludes \textbf{nothing}: its hypothesis fails, and that is not
        the same as showing no local inverse exists. A separate argument
        would be needed either way.
  \item Each component of $A\mathbf{x}$ is linear, so
        $\partial(A\mathbf{x})_{i}/\partial x_{j}=a_{ij}$ and $J=A$ at every
        point --- constant, with $\det J=\det A$. If $\det A\neq0$ then
        $\mathbf{x}=A^{-1}\mathbf{y}$ inverts the map on all of $\R^{n}$; if
        $\det A=0$ then $A$ has a non-trivial null space and the map is not
        injective anywhere. So for linear maps the local condition is
        already global.
  \item $J=\begin{pmatrix}2x&2y\\2x&-2y\end{pmatrix}$ and
        $\det J=-4xy-4xy=-8xy$, which is non-zero throughout $x,y>0$. So a
        differentiable local inverse exists at every point of the open
        quadrant.
  \item $f'(x_{0})\neq0$ guarantees an interval around $x_{0}$ on which $f$
        is invertible with a differentiable inverse of derivative
        $1/f'(x_{0})$; it does not guarantee that $f$ is injective anywhere
        beyond that interval. $\det J(\mathbf{x}_{0})\neq0$ guarantees the
        same locally in $n$ dimensions, with the inverse's Jacobian equal to
        $\bigl[J(\mathbf{x}_{0})\bigr]^{-1}$; it likewise says nothing about
        global injectivity.
\end{solutions}


\end{document}
"
\ifdefined\KindleEdition
  \documentclass[10pt,oneside]{book}
\else
  \documentclass[11pt,a4paper,twoside,openright]{book}
\fi
% Shared preamble for the Derivatives training sheet.
\usepackage[T1]{fontenc}
\usepackage[utf8]{inputenc}
\usepackage{lmodern}
\ifdefined\KindleEdition
  % 6-inch e-ink page (Kindle Paperwhite): the page *is* the screen, so the
  % margins are only what the eye needs, not what a printer needs.
  \usepackage[paperwidth=90mm,paperheight=122mm,
              top=4mm,bottom=6mm,left=4.5mm,right=4.5mm]{geometry}
\else
  \usepackage[a4paper,margin=2.4cm]{geometry}
\fi
\usepackage{amsmath,amssymb,amsthm}
\usepackage{microtype}
\usepackage{enumitem}
\usepackage{xcolor}
\usepackage[most]{tcolorbox}
\usepackage{fancyhdr}
\usepackage[hidelinks]{hyperref}

\definecolor{accent}{HTML}{1F4E79}
\definecolor{soft}{HTML}{F4F7FA}

% --- function graphs ------------------------------------------------------
\usepackage{pgfplots}
\pgfplotsset{compat=1.18}
\usetikzlibrary{arrows.meta}

% One consistent look for every function graph in these books.
\pgfplotsset{
  fnaxis/.style={
    width=5.9cm, height=4.6cm,
    axis lines=middle,
    axis line style={-{Stealth[length=4pt]}, gray!65!black, line width=0.3pt},
    tick style={gray!65!black},
    tick label style={font=\tiny, gray!50!black},
    label style={font=\scriptsize},
    xlabel={$x$}, ylabel={$y$},
    every axis x label/.style={at={(ticklabel* cs:1.0)}, anchor=west},
    every axis y label/.style={at={(ticklabel* cs:1.0)}, anchor=south east,
      xshift=-1pt},
    every axis plot/.append style={accent, very thick, smooth, samples=201},
    clip=true, enlargelimits=false,
    title style={font=\scriptsize, yshift=4pt},
  },
  fnaxis/.append style={/pgfplots/legend style={font=\tiny, draw=none,
    fill=none, at={(0.98,0.03)}, anchor=south east}},
}

% A row of small graphs, evenly spaced with \hfill between them, plus an
% optional caption line: \begin{gallery} <pic>\hfill<pic> \gcap{...} \end{gallery}
\newenvironment{gallery}
  {\par\medskip\noindent\ignorespaces}
  {\par\medskip\noindent}
\newcommand{\gcap}[1]{\par\smallskip{\centering\footnotesize #1\par}}

% --- plane regions, implicit curves, surfaces -----------------------------
\usepgfplotslibrary{fillbetween}
\definecolor{regionfill}{HTML}{9DB4D8}
\pgfplotsset{colormap={house}{color=(soft) color=(accent!45) color=(accent)}}

\pgfplotsset{
  % Squared, gridded axis for plane regions and implicitly given curves.
  geoaxis/.style={
    fnaxis,
    width=5.9cm, height=5.9cm,
    grid=both,
    grid style={gray!22, line width=0.2pt},
    minor tick num=1,
  },
  % Level-curve maps: many thin curves, no single curve emphasised.
  contouraxis/.style={
    geoaxis,
    every axis plot/.append style={accent!75, line width=0.5pt, smooth,
      samples=201},
  },
  % Surfaces: \addplot3[surf, domain=..., y domain=...] {expr};
  surfaxis/.style={
    width=6.0cm, height=5.0cm,
    view={130}{22},
    colormap name=house,
    ticks=none,
    xlabel={$x$}, ylabel={$y$}, zlabel={$f$},
    label style={font=\tiny, gray!50!black},
    title style={font=\scriptsize},
    every axis plot/.append style={samples=25, shader=faceted interp,
      draw=accent!35, line width=0.1pt},
  },
}

% Shorthand for a shaded region between two named paths.
% \region{<path A>}{<path B>}{<soft clip options>}
\newcommand{\region}[3]{%
  \addplot[regionfill, opacity=0.55, forget plot, sharp plot]
    fill between[of=#1 and #2, #3];}

% --- rule boxes -----------------------------------------------------------
\newtcolorbox{rulebox}[1]{%
  breakable, enhanced, colback=soft, colframe=accent,
  boxrule=0.4pt, left=7pt, right=7pt, top=6pt, bottom=6pt,
  fonttitle=\bfseries, coltitle=white, colbacktitle=accent,
  title={#1}}

% --- worked examples ------------------------------------------------------
% One fully narrated solution per chapter, carrying the chapter's central
% technique and the mistakes that technique invites:
%   \begin{worked}{Task} ... \wstep{Step 1.} ... \wtrap{...} ... \end{worked}
\newtcolorbox{worked}[1]{%
  breakable, enhanced, colback=white, colframe=accent!65,
  boxrule=0.4pt, left=7pt, right=7pt, top=6pt, bottom=6pt,
  fonttitle=\bfseries, coltitle=accent, colbacktitle=soft,
  title={#1}}

% Run-in step heading inside a worked example.
\newcommand{\wstep}[1]{\par\smallskip\noindent\textbf{#1}\enspace\ignorespaces}

% A pitfall, set off from the running solution so it can be found again.
\newcommand{\wtrap}[1]{%
  \par\smallskip
  {\leftskip=1.3em\noindent
   \llap{\textcolor{accent}{\rule[0.22ex]{0.9em}{0.6pt}}\hspace{0.4em}}%
   \textbf{Trap.}\enspace\itshape #1\par}%
  \smallskip}

% --- problem / solution lists --------------------------------------------
\newlist{problems}{enumerate}{1}
\setlist[problems]{label=\textbf{\arabic*.}, leftmargin=2.1em,
  itemsep=0.55ex, parsep=0pt, topsep=0.8ex}
\newlist{solutions}{enumerate}{1}
\setlist[solutions]{label=\textbf{\arabic*.}, leftmargin=2.1em,
  itemsep=1.0ex, parsep=0pt, topsep=0.8ex}

% --- running heads --------------------------------------------------------
\pagestyle{fancy}
\fancyhf{}
\fancyhead[LE,RO]{\thepage}
\fancyhead[RE]{\small\itshape\nouppercase{\leftmark}}
\fancyhead[LO]{\small\itshape\nouppercase{\rightmark}}
\setlength{\headheight}{14pt}
\renewcommand{\headrulewidth}{0.3pt}
\renewcommand{\chaptermark}[1]{\markboth{\chaptername\ \thechapter.\ #1}{}}
\renewcommand{\sectionmark}[1]{\markright{\thesection\ #1}}

% --- shorthands used throughout -------------------------------------------
\newcommand{\dd}[2]{\frac{\mathrm{d}#1}{\mathrm{d}#2}}
\newcommand{\pd}[2]{\frac{\partial #1}{\partial #2}}
\newcommand{\pdd}[3]{\frac{\partial^{2} #1}{\partial #2\,\partial #3}}
\newcommand{\R}{\mathbb{R}}
\newcommand{\grad}{\nabla}
\DeclareMathOperator{\sgn}{sgn}
\DeclareMathOperator{\Hess}{H}

% German gloss, used on first mention of a technical term
\newcommand{\de}[1]{\emph{#1}}

% --- cross-references to chapters ------------------------------------------
% Chapter files are shared by the standalone volumes and the compendium, so a
% hard-coded "Chapter 6" is wrong in one of them. Write \chref{derivatives}{6}
% instead: it prints "Chapter 6" inside the same volume, "Volume 5, Chapter 6"
% from another volume, and the running number ("Chapter 34") in the
% compendium. \chrefs{vol}{n} is the abbreviated form (Ch. / Vol., Ch.),
% \chnum{vol}{n} the bare number in the current build, and \volref{vol} names
% the volume ("Volume 5", or "Part V" in the compendium). Each volume's
% main.tex sets \thisvolume; the compendium's sets \incompendiumtrue.
\newif\ifincompendium
\newcommand*{\thisvolume}{}
\makeatletter
\@namedef{chref@vol@algebra}{1}\@namedef{chref@off@algebra}{0}
\@namedef{chref@vol@functions}{2}\@namedef{chref@off@functions}{9}
\@namedef{chref@vol@matrices}{3}\@namedef{chref@off@matrices}{14}
\@namedef{chref@vol@foundations}{4}\@namedef{chref@off@foundations}{24}
\@namedef{chref@vol@derivatives}{5}\@namedef{chref@off@derivatives}{28}
\@namedef{chref@vol@integrals}{6}\@namedef{chref@off@integrals}{42}
\@namedef{chref@vol@optimization}{7}\@namedef{chref@off@optimization}{50}
\@namedef{chref@vol@linprog}{8}\@namedef{chref@off@linprog}{60}
\@namedef{chref@vol@mixed}{9}\@namedef{chref@off@mixed}{68}
\newcommand*{\chref@check}[1]{%
  \@ifundefined{chref@vol@#1}{\PackageError{chref}{Unknown volume `#1'}{}}{}}
\newcommand*{\chnum}[2]{\chref@check{#1}%
  \ifincompendium\the\numexpr\csname chref@off@#1\endcsname+#2\relax
  \else#2\fi}
% \chref@same{vol}{same volume}{other volume}
\newcommand*{\chref@same}[3]{\def\chref@tmp{#1}%
  \ifx\chref@tmp\thisvolume\expandafter\@firstoftwo
  \else\expandafter\@secondoftwo\fi{#2}{#3}}
\newcommand*{\chref}[2]{\chref@check{#1}%
  \ifincompendium Chapter~\chnum{#1}{#2}%
  \else\chref@same{#1}{Chapter~#2}%
    {Volume~\csname chref@vol@#1\endcsname, Chapter~#2}\fi}
\newcommand*{\chrefs}[2]{\chref@check{#1}%
  \ifincompendium Ch.~\chnum{#1}{#2}%
  \else\chref@same{#1}{Ch.~#2}%
    {Vol.~\csname chref@vol@#1\endcsname, Ch.~#2}\fi}
\newcommand*{\volref}[1]{\chref@check{#1}%
  \ifincompendium Part~\@Roman{\csname chref@vol@#1\endcsname}%
  \else Volume~\csname chref@vol@#1\endcsname\fi}
\makeatother

% --- Kindle edition: compact layout for a 6-inch screen -------------------
% Everything below applies only when \KindleEdition is defined on the
% command line; the A4 edition is untouched.
\ifdefined\KindleEdition
  \usepackage{etoolbox}
  \usepackage{adjustbox}
  \usepackage{titlesec}

  % The 8cm measure leaves few break points; let TeX loosen a tight line
  % rather than push it into the margin.
  \setlength{\emergencystretch}{1.5em}

  % No running heads: on a 90mm page the header costs more than it gives.
  % A centred page number in the foot is kept so the contents list means
  % something.
  \pagestyle{plain}
  \fancypagestyle{plain}{\fancyhf{}\fancyfoot[C]{\small\thepage}%
    \renewcommand{\headrulewidth}{0pt}}

  % Compact headings: the book class leaves a third of a page above a
  % chapter title, which is unaffordable here.
  \titleformat{\chapter}[display]
    {\normalfont\Large\bfseries\color{accent}}{\normalsize\chaptertitlename\ \thechapter}{0.4em}{\Large}
  \titlespacing*{\chapter}{0pt}{0pt}{1.2em}
  \titlespacing*{\section}{0pt}{1.4ex plus .3ex}{0.7ex}
  \titlespacing*{\subsection}{0pt}{1.2ex plus .2ex}{0.5ex}

  % Boxes: trim the padding, which was set for a 16cm measure.
  \tcbset{left=3.5pt,right=3.5pt,top=3pt,bottom=3pt,
          fonttitle=\bfseries\small,
          halign title=flush left}

  % Problem and solution lists: reclaim the wide left margin.
  \setlist[problems]{label=\textbf{\arabic*.}, leftmargin=1.5em,
    itemsep=0.45ex, parsep=0pt, topsep=0.6ex}
  \setlist[solutions]{label=\textbf{\arabic*.}, leftmargin=1.5em,
    itemsep=0.8ex, parsep=0pt, topsep=0.6ex}

  % Figures and wide tables shrink to the measure instead of running off it.
  % adjustbox only scales down: anything already narrow is left alone.
  \BeforeBeginEnvironment{tikzpicture}{\begin{adjustbox}{max width=\linewidth}}
  \AfterEndEnvironment{tikzpicture}{\end{adjustbox}}
  \BeforeBeginEnvironment{tabular}{\begin{adjustbox}{max width=\linewidth}}
  \AfterEndEnvironment{tabular}{\end{adjustbox}}

  % Display maths was written for a 16cm measure. On 8cm the long chained
  % equations would run off the screen, so any display wider than the text
  % block is scaled down to fit; narrower ones are left at full size. The
  % whole corpus uses \[ ... \] and no other display environment, so this
  % one redefinition covers every case.
  \newsavebox{\kndbox}
  \renewcommand{\[}{\par\addvspace{0.7ex}%
    \begin{lrbox}{\kndbox}$\displaystyle}
  \renewcommand{\]}{$\end{lrbox}%
    {\centering
      \ifdim\wd\kndbox>\linewidth
        \resizebox{\linewidth}{!}{\usebox{\kndbox}}%
      \else
        \usebox{\kndbox}%
      \fi
     \par}%
    \addvspace{0.7ex}}

  % A narrow measure makes perfect justification impossible; prefer a loose
  % line to a line that runs into the margin.
  \sloppy
  \raggedbottom
\fi

\renewcommand*{\thisvolume}{derivatives}

\title{\Huge Derivatives \& Partial Derivatives\\[0.4em]
  \large A Training Sheet with Exercises and Worked Solutions}
\author{}
\date{}

\begin{document}
\frontmatter
\maketitle

\chapter*{About this sheet}
\addcontentsline{toc}{chapter}{About this sheet}

This sheet drills the computation of derivatives in one and
several variables, with the applications most used in economics: marginal quantities, elasticities, tangents and linearisations, rates of
substitution, and the second-order conditions behind optimisation.

The structure is the same in every chapter: the \textbf{Core Rules} section
states the rules compactly and says when each one applies; \textbf{Practice
Problems} drills them; \textbf{Solutions} works every problem through in
full, showing the intermediate step rather than only the answer.

Chapters 1--8 cover the one-variable theory: the differentiation rules and
the standard functions, higher derivatives and convexity, Taylor
approximation, implicit differentiation, the classical applications, the
indeterminate forms with L'H\^opital's rule, and a chapter that assembles all
of it into the complete curve discussion. Chapters 9--14 are the
multivariable half, opening with the geometry --- surfaces, level curves and
the sub- and superlevel sets --- before partial derivatives and the gradient
(with the topological vocabulary of domains), the Hessian and definiteness,
the multivariable chain rule and total differential, directional derivatives
and vector-valued functions, and finally the Implicit Function Theorem with
Jacobians and Euler's theorem. Most chapters close with a block of
\emph{additional problems} for a second pass through the material.

Work the problems before reading the solutions. A derivative you have
recognised is worth more than ten you have read.

\tableofcontents

\mainmatter
\chapter{Basic Differentiation Rules}

The derivative measures the instantaneous rate of change of a function. In
economics this is the language of \emph{marginal} quantities: marginal cost,
marginal product, marginal revenue. Before specialising to particular
functional forms, master the structural rules below; almost every derivative
computed later is built from a handful of them.

\section{Core Rules}

\begin{rulebox}{Definition (limit form)}
For a function $f:\R\to\R$ the derivative at $x$ is
\[
  f'(x) \;=\; \lim_{h\to 0}\frac{f(x+h)-f(x)}{h},
\]
provided this limit exists. Equivalent notations include $\dd{f}{x}$,
$\dfrac{\mathrm{d}}{\mathrm{d}x}f(x)$, and $Df(x)$.
\end{rulebox}

\begin{rulebox}{Linearity}
For constants $c,\lambda\in\R$ and differentiable $u,v$:
\begin{itemize}[leftmargin=1.6em,itemsep=0.3ex]
  \item \textbf{Constant rule:} $\dfrac{\mathrm{d}}{\mathrm{d}x}(c)=0$
  \item \textbf{Scalar multiple:} $(\lambda u)'=\lambda u'$
  \item \textbf{Sum rule:} $(u+v)'=u'+v'$
\end{itemize}
\end{rulebox}

\begin{rulebox}{Power rule}
For any real exponent $n$ (where the expression is defined):
\[
  \frac{\mathrm{d}}{\mathrm{d}x}x^{n} \;=\; n\,x^{n-1} .
\]
Special cases worth memorising: $(\sqrt{x})'=\dfrac{1}{2\sqrt{x}}$,
$\Bigl(\dfrac{1}{x}\Bigr)'=-\dfrac{1}{x^{2}}$, and $(x)'=1$.
\end{rulebox}

\begin{rulebox}{Product rule}
If $u,v$ are differentiable,
\[
  (uv)' \;=\; u'v+uv' .
\]
\end{rulebox}

\begin{rulebox}{Quotient rule}
If $u,v$ are differentiable and $v\neq 0$,
\[
  \left(\frac{u}{v}\right)' \;=\; \frac{u'v-uv'}{v^{2}} .
\]
\end{rulebox}

\begin{rulebox}{Chain rule}
If $f$ is differentiable at $g(x)$ and $g$ is differentiable at $x$,
\[
  (f\circ g)'(x) \;=\; f'\bigl(g(x)\bigr)\cdot g'(x) .
\]
In Leibniz notation, with $y=f(u)$ and $u=g(x)$, this reads
$\dd{y}{x}=\dd{y}{u}\cdot\dd{u}{x}$.
\end{rulebox}

\begin{rulebox}{Higher-order derivatives (preview)}
The second derivative is $f''(x)=(f'(x))'$, written also
$\dfrac{\mathrm{d}^{2}f}{\mathrm{d}x^{2}}$. In economics $f'>0$ and $f''<0$
characterise an increasing, concave production or utility function. Higher
orders $f^{(n)}$ appear in Taylor expansions and in concavity tests.
\end{rulebox}

\section{Practice Problems}

Differentiate each function with respect to $x$ (or the indicated variable).
Simplify where natural.

\begin{problems}
  \item $f(x)=x^{7}$
  \item $f(x)=(3x+5)^{6}$
  \item $f(x)=\dfrac{x}{x+1}$
  \item $f(x)=(x^{2}+1)(x^{3}-2x)$
  \item $f(x)=x^{5}-4x^{3}+9x-2$
  \item $f(x)=\sqrt{2x^{2}+1}$
  \item $C(q)=100+4q+0.05q^{2}$ \quad (find $C'(q)$, the marginal cost)
  \item $f(x)=\dfrac{2x-3}{x^{2}+1}$
  \item $f(x)=(3x-1)(x^{2}+4x+5)$
  \item $f(x)=\sqrt{x}$
  \item $f(x)=(x^{2}-4x+1)^{4}$
  \item $f(x)=\dfrac{x^{2}+1}{x^{2}-1}$
  \item $f(x)=x\sqrt{x+1}$
  \item $f(x)=6x^{4}+2x^{2}-x$
  \item $f(x)=\bigl((x+1)^{2}+1\bigr)^{3}$ \quad (double chain)
  \item $f(x)=\dfrac{\sqrt{x}}{x+1}$
  \item $f(x)=(x^{2}+2)(\sqrt{x}-3)$
  \item $f(x)=x^{2/3}$
  \item $R(q)=20q-0.1q^{2}$ \quad (find $R'(q)$, the marginal revenue)
  \item $f(x)=\dfrac{1}{(x^{2}+3)^{5}}$
  \item $f(x)=(x^{3}-x)(2x^{2}+7)$
  \item $f(x)=\dfrac{x^{3}-2x}{x+4}$
  \item $f(x)=\dfrac{1}{x^{4}}$
  \item $f(x)=\bigl(\sqrt{x}+2x\bigr)^{4}$
  \item $f(x)=5x^{4}-\tfrac{1}{2}x^{2}+7$
  \item $\Pi(L)=L\cdot(10-L)$ \quad (a profit-like product)
  \item $f(x)=\sqrt[3]{(x^{2}+x+1)^{2}}$
  \item Average cost: $\displaystyle AC(q)=\frac{100+2q^{2}}{q}$
  \item $f(x)=3\sqrt[3]{x}+\dfrac{2}{x^{2}}$
  \item $f(x)=(5x^{2}-1)\Bigl(\dfrac{1}{x}+2\Bigr)$
  \item $f(x)=\dfrac{1-x^{2}}{1+x^{2}}$
  \item $f(x)=\pi x^{3}-\sqrt{2}\,x+e$
  \item $f(x)=\bigl(\sqrt{x}+1\bigr)^{2}$ \quad \emph{(expand first)}
  \item $f(x)=x^{1/2}+x^{-1/2}$
  \item $f(x)=\dfrac{x^{3}-5x+1}{x}$
  \item $f(x)=\tfrac{1}{3}x^{6}+\tfrac{2}{5}x^{5}-4$
  \item $f(x)=(x^{2}-9)^{5}$
  \item $f(x)=\dfrac{x+2}{x-2}$
  \item $f(x)=(4x+3)(x^{2}-x)$
  \item $f(x)=2x^{8}-3x^{5}+x$
  \item $f(x)=\sqrt{4-x^{2}}$
  \item Profit: $\Pi(q)=(20-0.1q)q-(100+4q+0.05q^{2})$ \quad (find
        $\Pi'(q)$)
  \item $f(x)=\dfrac{1}{2x+7}$
  \item $f(x)=(x^{2}+1)^{3}(x-1)$
  \item $f(x)=\dfrac{x^{2}}{1+x}$
  \item $f(x)=7\sqrt{x}-\dfrac{3}{x}$
  \item $f(x)=\bigl(1+x^{3}\bigr)^{10}$
  \item $f(x)=(x+1)^{2}(x-1)^{3}$
  \item Marginal product: $Q(L)=5L^{1/2}$ \quad (find $Q'(L)$)
  \item $f(x)=\dfrac{3x+1}{(x+2)^{2}}$
  \item $f(x)=x^{3}-5x^{2}+7x-4$
  \item $f(x)=\sqrt{x^{2}+9}$
  \item $f(x)=(2x^{3}+1)(x^{2}+x)$
  \item $f(x)=\dfrac{x^{2}-1}{x^{2}+1}$
  \item $f(x)=\Bigl(\dfrac{1}{x}+x\Bigr)^{4}$
  \item $f(x)=-3x^{5}+2x^{-3}$
  \item $f(x)=\dfrac{1}{\sqrt{x^{2}+1}}$
  \item $f(x)=(x-1)(x-2)(x-3)$
  \item Cost: $C(q)=50+\sqrt{q^{2}+4}$ \quad (find $C'(q)$)
  \item $f(x)=\dfrac{x+1}{x^{2}+2}$
  \item $f(x)=(5-2x)^{7}$
  \item $f(x)=x^{2}\sqrt{1+x}$
  \item $f(x)=\dfrac{1}{x^{3}+1}$
  \item $f(x)=(x^{4}+3x)^{6}$
  \item $f(x)=\tfrac{2}{3}x^{3/2}-4x^{-1/2}$
  \item Revenue per unit time: $R(t)=(50+2t)^{3}$ \quad (find $R'(t)$)
  \item $f(x)=\dfrac{\sqrt{x}+1}{\sqrt{x}-1}$
  \item $f(x)=(x^{2}+1)(x^{2}-1)(x+3)$
  \item $f(x)=\sqrt{(x+1)(x+2)}$
  \item $f(x)=\dfrac{x^{2}+2x}{(x+1)^{2}}$
  \item $f(x)=\bigl(3x^{2}-2\sqrt{x}\bigr)^{5}$
\end{problems}

\section{Solutions}

\begin{solutions}
  \item $f'(x)=7x^{6}$ (power rule).

  \item Outer $u^{6}$, inner $u=3x+5$. $f'(x)=6(3x+5)^{5}\cdot 3
        =18(3x+5)^{5}$.

  \item Quotient rule with $u=x$, $v=x+1$, $u'=1$, $v'=1$:
        \[
          f'(x)=\frac{1\cdot(x+1)-x\cdot 1}{(x+1)^{2}}=\frac{1}{(x+1)^{2}} .
        \]

  \item Let $u=x^{2}+1$, $v=x^{3}-2x$. Then $u'=2x$, $v'=3x^{2}-2$.
        \[
          f'(x)=2x(x^{3}-2x)+(x^{2}+1)(3x^{2}-2)=5x^{4}-3x^{2}-2 .
        \]

  \item $f'(x)=5x^{4}-12x^{2}+9$.

  \item $f(x)=(2x^{2}+1)^{1/2}$, inner $u=2x^{2}+1$, $u'=4x$.
        \[
          f'(x)=\tfrac12(2x^{2}+1)^{-1/2}\cdot 4x=\frac{2x}{\sqrt{2x^{2}+1}} .
        \]

  \item $C'(q)=4+0.1q$. Marginal cost is linear in output.

  \item $u=2x-3$, $v=x^{2}+1$, $u'=2$, $v'=2x$.
        \[
          f'(x)=\frac{2(x^{2}+1)-(2x-3)(2x)}{(x^{2}+1)^{2}}
               =\frac{-2x^{2}+6x+2}{(x^{2}+1)^{2}} .
        \]

  \item $u=3x-1$, $v=x^{2}+4x+5$, $u'=3$, $v'=2x+4$.
        \[
          f'(x)=3(x^{2}+4x+5)+(3x-1)(2x+4)=9x^{2}+22x+11 .
        \]

  \item Write $\sqrt{x}=x^{1/2}$, so $f'(x)=\tfrac12x^{-1/2}
        =\dfrac{1}{2\sqrt{x}}$.

  \item Inner $u=x^{2}-4x+1$, $u'=2x-4$. $f'(x)=4(x^{2}-4x+1)^{3}(2x-4)
        =8(x-2)(x^{2}-4x+1)^{3}$.

  \item $u=x^{2}+1$, $v=x^{2}-1$, $u'=2x$, $v'=2x$.
        \[
          f'(x)=\frac{2x(x^{2}-1)-(x^{2}+1)(2x)}{(x^{2}-1)^{2}}
               =\frac{-4x}{(x^{2}-1)^{2}} .
        \]

  \item $u=x$, $v=(x+1)^{1/2}$, $u'=1$, $v'=\tfrac12(x+1)^{-1/2}$.
        \[
          f'(x)=\sqrt{x+1}+\frac{x}{2\sqrt{x+1}}
               =\frac{2(x+1)+x}{2\sqrt{x+1}}=\frac{3x+2}{2\sqrt{x+1}} .
        \]

  \item $f'(x)=24x^{3}+4x-1$.

  \item Two layers. Let $g(x)=(x+1)^{2}+1$, $g'(x)=2(x+1)$. Then
        \[
          f'(x)=3\,g(x)^{2}\cdot g'(x)=6(x+1)\bigl((x+1)^{2}+1\bigr)^{2} .
        \]

  \item $u=x^{1/2}$, $v=x+1$, $u'=\tfrac12x^{-1/2}$, $v'=1$.
        \[
          f'(x)=\frac{\frac{1}{2\sqrt{x}}(x+1)-\sqrt{x}}{(x+1)^{2}}
               =\frac{(x+1)-2x}{2\sqrt{x}(x+1)^{2}}
               =\frac{1-x}{2\sqrt{x}(x+1)^{2}} .
        \]

  \item $u=x^{2}+2$, $v=x^{1/2}-3$, $u'=2x$, $v'=\tfrac12x^{-1/2}$.
        \[
          f'(x)=2x(\sqrt{x}-3)+(x^{2}+2)\cdot\frac{1}{2\sqrt{x}}
               =2x\sqrt{x}-6x+\frac{x^{2}+2}{2\sqrt{x}} .
        \]

  \item $f'(x)=\tfrac23x^{-1/3}=\dfrac{2}{3\sqrt[3]{x}}$.

  \item $R'(q)=20-0.2q$. Marginal revenue falls with $q$.

  \item $f(x)=(x^{2}+3)^{-5}$, inner $u'=2x$.
        \[
          f'(x)=-5(x^{2}+3)^{-6}\cdot 2x=-\frac{10x}{(x^{2}+3)^{6}} .
        \]

  \item $u=x^{3}-x$, $v=2x^{2}+7$, $u'=3x^{2}-1$, $v'=4x$.
        \[
          f'(x)=(3x^{2}-1)(2x^{2}+7)+(x^{3}-x)(4x)=10x^{4}+15x^{2}-7 .
        \]

  \item $u=x^{3}-2x$, $v=x+4$, $u'=3x^{2}-2$, $v'=1$.
        \[
          f'(x)=\frac{(3x^{2}-2)(x+4)-(x^{3}-2x)}{(x+4)^{2}}
               =\frac{2x^{3}+12x^{2}-8}{(x+4)^{2}} .
        \]

  \item Write $f(x)=x^{-4}$, so $f'(x)=-4x^{-5}=-\dfrac{4}{x^{5}}$.

  \item Inner $u=\sqrt{x}+2x$, $u'=\dfrac{1}{2\sqrt{x}}+2$.
        \[
          f'(x)=4\bigl(\sqrt{x}+2x\bigr)^{3}\left(\frac{1}{2\sqrt{x}}+2\right) .
        \]

  \item $f'(x)=20x^{3}-x$.

  \item $\Pi'(L)=1\cdot(10-L)+L\cdot(-1)=10-2L$.

  \item $f(x)=(x^{2}+x+1)^{2/3}$, inner $u'=2x+1$.
        \[
          f'(x)=\tfrac23(x^{2}+x+1)^{-1/3}(2x+1)
               =\frac{2(2x+1)}{3\sqrt[3]{x^{2}+x+1}} .
        \]

  \item Easier: $AC(q)=100q^{-1}+2q$, so $AC'(q)=-100q^{-2}+2
        =2-\dfrac{100}{q^{2}}$.

  \item $f(x)=3x^{1/3}+2x^{-2}$, hence $f'(x)=x^{-2/3}-4x^{-3}
        =\dfrac{1}{x^{2/3}}-\dfrac{4}{x^{3}}$.

  \item $u=5x^{2}-1$, $v=x^{-1}+2$, $u'=10x$, $v'=-x^{-2}$.
        \[
          f'(x)=10x(x^{-1}+2)+(5x^{2}-1)(-x^{-2})
               =10+20x-5+x^{-2}=5+20x+\frac{1}{x^{2}} .
        \]

  \item $u=1-x^{2}$, $v=1+x^{2}$, $u'=-2x$, $v'=2x$.
        \[
          f'(x)=\frac{-2x(1+x^{2})-(1-x^{2})(2x)}{(1+x^{2})^{2}}
               =\frac{-4x}{(1+x^{2})^{2}} .
        \]

  \item $f'(x)=3\pi x^{2}-\sqrt{2}$. ($e$ is a constant.)

  \item Expand: $f(x)=x+2\sqrt{x}+1$, so $f'(x)=1+\dfrac{1}{\sqrt{x}}$.

  \item $f'(x)=\tfrac12x^{-1/2}-\tfrac12x^{-3/2}
        =\dfrac{1}{2\sqrt{x}}-\dfrac{1}{2x^{3/2}}$.

  \item Simplify first: $f(x)=x^{2}-5+x^{-1}$, so $f'(x)=2x-x^{-2}
        =2x-\dfrac{1}{x^{2}}$.

  \item $f'(x)=2x^{5}+2x^{4}$.

  \item Chain rule with $u=x^{2}-9$, $u'=2x$:
        \[
          f'(x)=5(x^{2}-9)^{4}\cdot 2x=10x(x^{2}-9)^{4} .
        \]

  \item $u=x+2$, $v=x-2$, $u'=1$, $v'=1$.
        \[
          f'(x)=\frac{1\cdot(x-2)-(x+2)\cdot 1}{(x-2)^{2}}
               =\frac{-4}{(x-2)^{2}} .
        \]

  \item Product rule with $u=4x+3$, $v=x^{2}-x$, $u'=4$, $v'=2x-1$.
        \[
          f'(x)=4(x^{2}-x)+(4x+3)(2x-1)=12x^{2}-2x-3 .
        \]

  \item $f'(x)=16x^{7}-15x^{4}+1$.

  \item $f(x)=(4-x^{2})^{1/2}$, inner $u'=-2x$.
        \[
          f'(x)=\tfrac12(4-x^{2})^{-1/2}\cdot(-2x)=-\frac{x}{\sqrt{4-x^{2}}} .
        \]

  \item Expand: $\Pi(q)=20q-0.1q^{2}-100-4q-0.05q^{2}
        =-100+16q-0.15q^{2}$.
        \[
          \Pi'(q)=16-0.3q .
        \]
        Marginal profit is zero at $q^{*}=16/0.3\approx 53.33$.

  \item Write $f(x)=(2x+7)^{-1}$. Chain rule: $f'(x)=-(2x+7)^{-2}\cdot 2
        =-\dfrac{2}{(2x+7)^{2}}$.

  \item Product rule with $u=(x^{2}+1)^{3}$, $v=x-1$. $u'=3(x^{2}+1)^{2}
        \cdot 2x=6x(x^{2}+1)^{2}$, $v'=1$.
        \[
          f'(x)=6x(x^{2}+1)^{2}(x-1)+(x^{2}+1)^{3}
               =(x^{2}+1)^{2}\bigl[6x(x-1)+(x^{2}+1)\bigr]
               =(x^{2}+1)^{2}(7x^{2}-6x+1) .
        \]

  \item Quotient rule with $u=x^{2}$, $v=1+x$, $u'=2x$, $v'=1$.
        \[
          f'(x)=\frac{2x(1+x)-x^{2}}{(1+x)^{2}}=\frac{x^{2}+2x}{(1+x)^{2}}
               =\frac{x(x+2)}{(1+x)^{2}} .
        \]

  \item $f(x)=7x^{1/2}-3x^{-1}$, so $f'(x)=\tfrac72x^{-1/2}+3x^{-2}
        =\dfrac{7}{2\sqrt{x}}+\dfrac{3}{x^{2}}$.

  \item Inner $u=1+x^{3}$, $u'=3x^{2}$.
        \[
          f'(x)=10(1+x^{3})^{9}\cdot 3x^{2}=30x^{2}(1+x^{3})^{9} .
        \]

  \item Product rule with $u=(x+1)^{2}$, $v=(x-1)^{3}$, $u'=2(x+1)$,
        $v'=3(x-1)^{2}$.
        \[
          f'(x)=2(x+1)(x-1)^{3}+3(x+1)^{2}(x-1)^{2}
               =(x+1)(x-1)^{2}\bigl[2(x-1)+3(x+1)\bigr]
               =(x+1)(x-1)^{2}(5x+1) .
        \]

  \item $Q(L)=5L^{1/2}$, so $Q'(L)=\tfrac52L^{-1/2}=\dfrac{5}{2\sqrt{L}}$.
        Marginal product is positive but diminishing.

  \item Quotient rule with $u=3x+1$, $v=(x+2)^{2}$, $u'=3$,
        $v'=2(x+2)$.
        \[
          f'(x)=\frac{3(x+2)^{2}-(3x+1)\cdot 2(x+2)}{(x+2)^{4}}
               =\frac{3(x+2)-2(3x+1)}{(x+2)^{3}}=\frac{-3x+4}{(x+2)^{3}} .
        \]

  \item $f'(x)=3x^{2}-10x+7$.

  \item $f(x)=(x^{2}+9)^{1/2}$, $u'=2x$.
        \[
          f'(x)=\tfrac12(x^{2}+9)^{-1/2}\cdot 2x=\frac{x}{\sqrt{x^{2}+9}} .
        \]

  \item Product rule with $u=2x^{3}+1$, $v=x^{2}+x$, $u'=6x^{2}$,
        $v'=2x+1$.
        \[
          f'(x)=6x^{2}(x^{2}+x)+(2x^{3}+1)(2x+1)=10x^{4}+8x^{3}+2x+1 .
        \]

  \item Quotient rule with $u=x^{2}-1$, $v=x^{2}+1$, $u'=2x$, $v'=2x$.
        \[
          f'(x)=\frac{2x(x^{2}+1)-(x^{2}-1)(2x)}{(x^{2}+1)^{2}}
               =\frac{4x}{(x^{2}+1)^{2}} .
        \]

  \item Inner $u=x^{-1}+x$, $u'=-x^{-2}+1$.
        \[
          f'(x)=4\Bigl(\frac{1}{x}+x\Bigr)^{3}\left(1-\frac{1}{x^{2}}\right) .
        \]

  \item $f'(x)=-15x^{4}-6x^{-4}=-15x^{4}-\dfrac{6}{x^{4}}$.

  \item $f(x)=(x^{2}+1)^{-1/2}$, inner $u'=2x$.
        \[
          f'(x)=-\tfrac12(x^{2}+1)^{-3/2}\cdot 2x=-\frac{x}{(x^{2}+1)^{3/2}} .
        \]

  \item Expand or use the product rule iteratively. Quickest:
        $(x-1)(x-2)(x-3)=x^{3}-6x^{2}+11x-6$, so $f'(x)=3x^{2}-12x+11$.

  \item $C(q)=50+(q^{2}+4)^{1/2}$, $u'=2q$.
        \[
          C'(q)=\tfrac12(q^{2}+4)^{-1/2}\cdot 2q=\frac{q}{\sqrt{q^{2}+4}} .
        \]

  \item Quotient rule with $u=x+1$, $v=x^{2}+2$, $u'=1$, $v'=2x$.
        \[
          f'(x)=\frac{1\cdot(x^{2}+2)-(x+1)(2x)}{(x^{2}+2)^{2}}
               =\frac{-x^{2}-2x+2}{(x^{2}+2)^{2}} .
        \]

  \item Inner $u=5-2x$, $u'=-2$.
        \[
          f'(x)=7(5-2x)^{6}\cdot(-2)=-14(5-2x)^{6} .
        \]

  \item Product rule with $u=x^{2}$, $v=(1+x)^{1/2}$, $u'=2x$,
        $v'=\tfrac12(1+x)^{-1/2}$.
        \[
          f'(x)=2x\sqrt{1+x}+\frac{x^{2}}{2\sqrt{1+x}}
               =\frac{4x(1+x)+x^{2}}{2\sqrt{1+x}}=\frac{5x^{2}+4x}{2\sqrt{1+x}} .
        \]

  \item $f(x)=(x^{3}+1)^{-1}$, inner $u'=3x^{2}$.
        \[
          f'(x)=-(x^{3}+1)^{-2}\cdot 3x^{2}=-\frac{3x^{2}}{(x^{3}+1)^{2}} .
        \]

  \item Inner $u=x^{4}+3x$, $u'=4x^{3}+3$.
        \[
          f'(x)=6(x^{4}+3x)^{5}(4x^{3}+3) .
        \]

  \item $f'(x)=\tfrac23\cdot\tfrac32x^{1/2}-4\cdot(-\tfrac12)x^{-3/2}
        =x^{1/2}+2x^{-3/2}=\sqrt{x}+\dfrac{2}{x^{3/2}}$.

  \item Inner $u=50+2t$, $u'=2$.
        \[
          R'(t)=3(50+2t)^{2}\cdot 2=6(50+2t)^{2} .
        \]

  \item Quotient rule with $u=\sqrt{x}+1$, $v=\sqrt{x}-1$,
        $u'=v'=\dfrac{1}{2\sqrt{x}}$.
        \[
          f'(x)=\frac{\frac{1}{2\sqrt{x}}(\sqrt{x}-1)
               -(\sqrt{x}+1)\cdot\frac{1}{2\sqrt{x}}}{(\sqrt{x}-1)^{2}}
               =\frac{-1/\sqrt{x}}{(\sqrt{x}-1)^{2}}
               =-\frac{1}{\sqrt{x}(\sqrt{x}-1)^{2}} .
        \]

  \item Group: note $(x^{2}+1)(x^{2}-1)=x^{4}-1$, so
        $f(x)=(x^{4}-1)(x+3)=x^{5}+3x^{4}-x-3$. Hence
        $f'(x)=5x^{4}+12x^{3}-1$.

  \item $f(x)=\bigl((x+1)(x+2)\bigr)^{1/2}=(x^{2}+3x+2)^{1/2}$, inner
        $u'=2x+3$.
        \[
          f'(x)=\tfrac12(x^{2}+3x+2)^{-1/2}(2x+3)
               =\frac{2x+3}{2\sqrt{(x+1)(x+2)}} .
        \]

  \item Simplify first: $x^{2}+2x=x(x+2)$ and $(x+1)^{2}=x^{2}+2x+1$.
        Quotient rule with $u=x^{2}+2x$, $v=(x+1)^{2}$, $u'=2x+2=2(x+1)$,
        $v'=2(x+1)$.
        \[
          f'(x)=\frac{2(x+1)(x+1)^{2}-(x^{2}+2x)\cdot 2(x+1)}{(x+1)^{4}}
               =\frac{2(x+1)^{2}-2(x^{2}+2x)}{(x+1)^{3}}=\frac{2}{(x+1)^{3}} .
        \]

  \item Inner $u=3x^{2}-2x^{1/2}$, $u'=6x-x^{-1/2}$.
        \[
          f'(x)=5\bigl(3x^{2}-2\sqrt{x}\bigr)^{4}\left(6x-\frac{1}{\sqrt{x}}\right) .
        \]
\end{solutions}

\section{Additional Problems}

\begin{problems}
  \item $f(x)=4x^{5}-7x^{3}+2x-9$
  \item $f(x)=(2x+3)^{8}$
  \item $f(x)=\dfrac{x-1}{x+1}$
  \item $f(x)=(x^{3}-2)(x^{2}+5)$
  \item $f(x)=\sqrt{3x^{2}+4}$
  \item $f(x)=\dfrac{x^{2}}{x^{3}+1}$
  \item Cost: $C(q)=200+6q+0.02q^{3}$ \quad (find $C'(q)$)
  \item $f(x)=(x^{2}+x)^{5}$
  \item $f(x)=\dfrac{1}{(2x-5)^{3}}$
  \item $f(x)=x\sqrt{x^{2}+1}$
  \item $f(x)=\bigl((x^{2}+1)^{2}-3\bigr)^{4}$ \quad (double chain)
  \item Marginal product: $Q(L)=8L^{2/3}$ \quad (find $Q'(L)$)
  \item $f(x)=\dfrac{x^{2}+3x}{(x-1)^{2}}$
  \item $f(x)=\sqrt[4]{x^{3}}$
  \item Profit: $\Pi(q)=(30-0.2q)\,q-(50+2q+0.1q^{2})$ \quad (find
        $\Pi'(q)$)
\end{problems}

\section{Additional Solutions}

\begin{solutions}
  \item $f'(x)=20x^{4}-21x^{2}+2$.

  \item Chain rule: $f'(x)=8(2x+3)^{7}\cdot 2=16(2x+3)^{7}$.

  \item Quotient rule with $u=x-1$, $v=x+1$, $u'=1$, $v'=1$:
        \[
          f'(x)=\frac{(x+1)-(x-1)}{(x+1)^{2}}=\frac{2}{(x+1)^{2}} .
        \]

  \item Product rule: $u=x^{3}-2$, $v=x^{2}+5$, $u'=3x^{2}$, $v'=2x$.
        \[
          f'(x)=3x^{2}(x^{2}+5)+(x^{3}-2)(2x)=5x^{4}+15x^{2}-4x .
        \]

  \item $f(x)=(3x^{2}+4)^{1/2}$, inner $u'=6x$.
        \[
          f'(x)=\tfrac12(3x^{2}+4)^{-1/2}\cdot 6x=\frac{3x}{\sqrt{3x^{2}+4}} .
        \]

  \item Quotient rule with $u=x^{2}$, $v=x^{3}+1$, $u'=2x$, $v'=3x^{2}$.
        \[
          f'(x)=\frac{2x(x^{3}+1)-x^{2}\cdot 3x^{2}}{(x^{3}+1)^{2}}
               =\frac{-x^{4}+2x}{(x^{3}+1)^{2}}
               =\frac{x(2-x^{3})}{(x^{3}+1)^{2}} .
        \]

  \item $C'(q)=6+0.06q^{2}$. Marginal cost grows quadratically with output.

  \item Chain rule: inner $u=x^{2}+x$, $u'=2x+1$.
        $f'(x)=5(x^{2}+x)^{4}(2x+1)$.

  \item $f(x)=(2x-5)^{-3}$. Chain rule: $f'(x)=-3(2x-5)^{-4}\cdot 2
        =-\dfrac{6}{(2x-5)^{4}}$.

  \item Product rule with $u=x$, $v=(x^{2}+1)^{1/2}$, $u'=1$,
        $v'=x/(x^{2}+1)^{1/2}$.
        \[
          f'(x)=\sqrt{x^{2}+1}+\frac{x^{2}}{\sqrt{x^{2}+1}}
               =\frac{2x^{2}+1}{\sqrt{x^{2}+1}} .
        \]

  \item Two layers. Let $g(x)=(x^{2}+1)^{2}-3$,
        $g'(x)=2(x^{2}+1)\cdot 2x=4x(x^{2}+1)$.
        \[
          f'(x)=4\,g(x)^{3}\cdot g'(x)
               =16x(x^{2}+1)\bigl((x^{2}+1)^{2}-3\bigr)^{3} .
        \]

  \item $Q'(L)=8\cdot\tfrac23L^{-1/3}=\dfrac{16}{3L^{1/3}}$. Marginal
        product is positive but diminishing.

  \item Quotient rule with $u=x^{2}+3x$, $v=(x-1)^{2}$, $u'=2x+3$,
        $v'=2(x-1)$.
        \[
          f'(x)=\frac{(2x+3)(x-1)^{2}-(x^{2}+3x)\cdot 2(x-1)}{(x-1)^{4}}
               =\frac{(2x+3)(x-1)-2(x^{2}+3x)}{(x-1)^{3}}
               =\frac{-5x-3}{(x-1)^{3}} .
        \]

  \item $f(x)=x^{3/4}$, so $f'(x)=\tfrac34x^{-1/4}=\dfrac{3}{4\sqrt[4]{x}}$.

  \item Expand: $\Pi(q)=30q-0.2q^{2}-50-2q-0.1q^{2}
        =-50+28q-0.3q^{2}$.
        \[
          \Pi'(q)=28-0.6q .
        \]
        Marginal profit is zero at $q^{*}=28/0.6\approx 46.67$.
\end{solutions}

\chapter{Derivatives of Standard Functions}

The basic rules of Section 1 reduce every derivative to a finite list of
\emph{building blocks}: exponentials, logarithms, trigonometric, inverse
trigonometric and hyperbolic functions. Memorise the table below. In
economics, exponentials govern continuous growth and discounting, while
logarithms encode percentage changes and elasticities.

\section{Core Rules}

\begin{rulebox}{Exponential functions}
\begin{itemize}[leftmargin=1.6em,itemsep=0.3ex]
  \item $\dfrac{\mathrm{d}}{\mathrm{d}x}\mathrm{e}^{x}=\mathrm{e}^{x}$
  \item $\dfrac{\mathrm{d}}{\mathrm{d}x}\mathrm{e}^{ax}=a\,\mathrm{e}^{ax}$
        \quad ($a$ constant; chain rule)
  \item $\dfrac{\mathrm{d}}{\mathrm{d}x}a^{x}=a^{x}\ln a$ \quad
        ($a>0$, $a\neq 1$)
\end{itemize}
\end{rulebox}

\begin{rulebox}{Logarithmic functions}
\begin{itemize}[leftmargin=1.6em,itemsep=0.3ex]
  \item $\dfrac{\mathrm{d}}{\mathrm{d}x}\ln x=\dfrac{1}{x}$ \quad ($x>0$)
  \item $\dfrac{\mathrm{d}}{\mathrm{d}x}\ln|x|=\dfrac{1}{x}$ \quad
        ($x\neq 0$)
  \item $\dfrac{\mathrm{d}}{\mathrm{d}x}\log_{a}x=\dfrac{1}{x\ln a}$ \quad
        ($a>0$, $a\neq 1$)
\end{itemize}
\end{rulebox}

\begin{rulebox}{Trigonometric functions}
Angles in radians.
\begin{itemize}[leftmargin=1.6em,itemsep=0.3ex]
  \item $(\sin x)'=\cos x$
  \item $(\cos x)'=-\sin x$
  \item $(\tan x)'=\sec^{2}x=\dfrac{1}{\cos^{2}x}$
  \item $(\cot x)'=-\csc^{2}x=-\dfrac{1}{\sin^{2}x}$
  \item $(\sec x)'=\sec x\tan x$
  \item $(\csc x)'=-\csc x\cot x$
\end{itemize}
\end{rulebox}

\begin{rulebox}{Inverse trigonometric functions}
\begin{itemize}[leftmargin=1.6em,itemsep=0.3ex]
  \item $(\arcsin x)'=\dfrac{1}{\sqrt{1-x^{2}}}$, \quad $|x|<1$
  \item $(\arccos x)'=-\dfrac{1}{\sqrt{1-x^{2}}}$, \quad $|x|<1$
  \item $(\arctan x)'=\dfrac{1}{1+x^{2}}$
\end{itemize}
\end{rulebox}

\begin{rulebox}{Hyperbolic functions}
Recall $\sinh x=\tfrac12(\mathrm{e}^{x}-\mathrm{e}^{-x})$,
$\cosh x=\tfrac12(\mathrm{e}^{x}+\mathrm{e}^{-x})$.
\begin{itemize}[leftmargin=1.6em,itemsep=0.3ex]
  \item $(\sinh x)'=\cosh x$
  \item $(\cosh x)'=\sinh x$ \quad (note: \emph{no} sign change)
  \item $(\tanh x)'=\operatorname{sech}^{2}x=1-\tanh^{2}x$
\end{itemize}
\end{rulebox}

\begin{rulebox}{Summary table}
\[
\begin{array}{lll}
\hline
f(x) & f'(x) & \text{Domain note} \\
\hline
\mathrm{e}^{x} & \mathrm{e}^{x} & \R \\
a^{x} & a^{x}\ln a & a>0,\ a\neq 1 \\
\ln x & 1/x & x>0 \\
\log_{a}x & 1/(x\ln a) & x>0 \\
\sin x & \cos x & \R \\
\cos x & -\sin x & \R \\
\tan x & \sec^{2}x & x\neq \tfrac{\pi}{2}+k\pi \\
\cot x & -\csc^{2}x & x\neq k\pi \\
\sec x & \sec x\tan x & x\neq \tfrac{\pi}{2}+k\pi \\
\csc x & -\csc x\cot x & x\neq k\pi \\
\arcsin x & 1/\sqrt{1-x^{2}} & |x|<1 \\
\arccos x & -1/\sqrt{1-x^{2}} & |x|<1 \\
\arctan x & 1/(1+x^{2}) & \R \\
\sinh x & \cosh x & \R \\
\cosh x & \sinh x & \R \\
\tanh x & \operatorname{sech}^{2}x & \R \\
\hline
\end{array}
\]
Combine each entry with the chain rule whenever the argument is a
non-trivial function of $x$.
\end{rulebox}

\section{Practice Problems}

Differentiate each of the following with respect to $x$ (or the indicated
variable).

\begin{problems}
  \item $f(x)=\mathrm{e}^{3x}$
  \item $f(x)=\ln(5x)$
  \item $f(x)=\sin(3x)$
  \item $f(x)=\arcsin(2x)$
  \item $f(x)=\sinh(2x)$
  \item $f(x)=\mathrm{e}^{-x^{2}}$
  \item $f(x)=\cos(x^{2})$
  \item $f(x)=\ln(x^{2}+1)$
  \item $f(x)=x\,\mathrm{e}^{x}$
  \item $f(x)=\arctan(x^{2})$
  \item $f(x)=\tan(2x+1)$
  \item Continuous compounding: $V(t)=V_{0}\,\mathrm{e}^{rt}$ \quad
        (constants $V_{0},r$)
  \item $f(x)=x\ln x$
  \item $f(x)=\cosh(x^{2})$
  \item $f(x)=\mathrm{e}^{\sqrt{x}}$
  \item $f(x)=\arccos(\sqrt{x})$
  \item $f(x)=x\sin x$
  \item $f(x)=\ln(\sqrt{x})$
  \item $f(x)=\sin^{2}x$ \quad (i.e.\ $(\sin x)^{2}$)
  \item $f(x)=2^{x}$
  \item $f(x)=\mathrm{e}^{x^{2}+1}$
  \item $f(x)=(\ln x)^{2}$
  \item $f(x)=\dfrac{\sin x}{x}$
  \item $f(x)=x\arctan x$
  \item $f(x)=\tanh(3x)$
  \item $f(x)=\log_{10}(x)$
  \item $f(x)=\sin(\cos x)$
  \item $f(x)=\mathrm{e}^{\sin x}$
  \item $f(x)=\arctan\Bigl(\dfrac{1}{x}\Bigr)$, \quad $x>0$
  \item $f(x)=\ln\Bigl(\dfrac{x}{x+1}\Bigr)$
  \item $f(x)=\sec(3x)$
  \item Log-utility: $U(c)=\ln(c)$ \quad (find marginal utility)
  \item $f(x)=\ln(\cos x)$
  \item $f(x)=\tan^{3}x$
  \item Logistic: $P(t)=\dfrac{1}{1+\mathrm{e}^{-t}}$
  \item $f(x)=\dfrac{\mathrm{e}^{x}}{1+\mathrm{e}^{x}}$
  \item $f(x)=\mathrm{e}^{x}\sin(2x)$
  \item Cobb-Douglas log: $f(K)=\ln(K^{\alpha})$ \quad ($\alpha$ constant)
  \item $f(x)=x^{x}$ \quad (use logarithmic differentiation)
  \item $f(x)=\mathrm{e}^{2x}\cos(3x)$
  \item $f(x)=\ln\bigl(x+\sqrt{x^{2}+1}\bigr)$
  \item $f(x)=\sin(x)\cos(x)$
  \item Discount factor: $D(t)=\mathrm{e}^{-rt}$ \quad (find $D'(t)$ and
        interpret)
  \item $f(x)=\arctan(\mathrm{e}^{x})$
  \item $f(x)=\tanh(x^{2}+1)$
  \item $f(x)=\dfrac{\ln x}{x}$
  \item $f(x)=\mathrm{e}^{x}\ln x$
  \item $f(x)=\sqrt{\sin x}$
  \item Exponential growth: $N(t)=1000\cdot\mathrm{e}^{0.04t}$ \quad (find
        $N'(t)$ and the growth rate $N'/N$)
  \item $f(x)=5^{x^{2}}$
  \item $f(x)=\arcsin(x^{2})$
  \item $f(x)=\cot(x)+\csc(x)$
  \item $f(x)=\sinh(x)\cosh(x)$
  \item Log of a product: $f(x)=\ln\bigl(x\sqrt{x^{2}+1}\bigr)$
  \item $f(x)=\cos^{4}(x)$
  \item $f(x)=\mathrm{e}^{x^{3}}$
  \item Iso-elastic utility: $\displaystyle U(c)=\frac{c^{1-\sigma}-1}{1-\sigma}$,
        $\sigma\neq 1$ \quad (find $U'(c)$)
  \item $f(x)=\arctan(\sqrt{x})$
  \item $f(x)=\tan(x)\ln(x)$
  \item $f(x)=\dfrac{1}{\ln x}$
  \item $f(x)=\mathrm{e}^{ax}\sin(bx)$ \quad (constants $a,b$)
  \item $f(x)=\log_{2}(x^{2}+1)$
  \item $f(x)=\cos(2x)\mathrm{e}^{-x}$
  \item $f(x)=\operatorname{arccot}(x)\cdot\arctan(x)$ \quad [use
        $\operatorname{arccot}(x)=\tfrac{\pi}{2}-\arctan(x)$]
  \item $f(x)=\ln(\ln(x))$, \quad $x>1$
  \item $f(x)=x^{2}\mathrm{e}^{-x}$
  \item Per-capita growth: $g(t)=\ln(N(t))$ where $N(t)=N_{0}\mathrm{e}^{rt}$
  \item $f(x)=\sin(\mathrm{e}^{x})$
  \item $f(x)=(1+\mathrm{e}^{x})^{3}$
  \item $f(x)=x^{\sin x}$ \quad (use logarithmic differentiation)
  \item $f(x)=\arctan(\tan x)$ \quad (restrict $x\in(-\tfrac{\pi}{2},
        \tfrac{\pi}{2})$)
\end{problems}

\section{Solutions}

\begin{solutions}
  \item Inner $u=3x$, $u'=3$. $f'(x)=\mathrm{e}^{3x}\cdot 3=3\mathrm{e}^{3x}$.

  \item Either expand $\ln(5x)=\ln 5+\ln x$, or chain:
        $f'(x)=\dfrac{1}{5x}\cdot 5=\dfrac{1}{x}$.

  \item Inner $u=3x$, $u'=3$. $f'(x)=\cos(3x)\cdot 3=3\cos(3x)$.

  \item Inner $u=2x$, $u'=2$.
        \[
          f'(x)=\frac{1}{\sqrt{1-(2x)^{2}}}\cdot 2=\frac{2}{\sqrt{1-4x^{2}}} .
        \]

  \item Inner $u=2x$, $u'=2$. $f'(x)=\cosh(2x)\cdot 2=2\cosh(2x)$.

  \item Inner $u=-x^{2}$, $u'=-2x$.
        $f'(x)=\mathrm{e}^{-x^{2}}\cdot(-2x)=-2x\,\mathrm{e}^{-x^{2}}$.

  \item Inner $u=x^{2}$, $u'=2x$.
        $f'(x)=-\sin(x^{2})\cdot 2x=-2x\sin(x^{2})$.

  \item Inner $u=x^{2}+1$, $u'=2x$. $f'(x)=\dfrac{2x}{x^{2}+1}$.

  \item Product rule: $f'(x)=1\cdot\mathrm{e}^{x}+x\cdot\mathrm{e}^{x}
        =(1+x)\,\mathrm{e}^{x}$.

  \item Inner $u=x^{2}$, $u'=2x$. $f'(x)=\dfrac{2x}{1+x^{4}}$.

  \item Inner $u=2x+1$, $u'=2$. $f'(x)=2\sec^{2}(2x+1)$.

  \item $V'(t)=V_{0}\cdot r\,\mathrm{e}^{rt}=r\,V(t)$. The growth rate is
        constant: $V'/V=r$.

  \item Product rule: $f'(x)=1\cdot\ln x+x\cdot\dfrac{1}{x}=\ln x+1$.

  \item Inner $u=x^{2}$, $u'=2x$. $f'(x)=\sinh(x^{2})\cdot 2x
        =2x\sinh(x^{2})$.

  \item Inner $u=\sqrt{x}$, $u'=\dfrac{1}{2\sqrt{x}}$.
        \[
          f'(x)=\mathrm{e}^{\sqrt{x}}\cdot\frac{1}{2\sqrt{x}}
               =\frac{\mathrm{e}^{\sqrt{x}}}{2\sqrt{x}} .
        \]

  \item Inner $u=\sqrt{x}=x^{1/2}$, $u'=\dfrac{1}{2\sqrt{x}}$.
        \[
          f'(x)=-\frac{1}{\sqrt{1-x}}\cdot\frac{1}{2\sqrt{x}}
               =-\frac{1}{2\sqrt{x-x^{2}}} .
        \]

  \item Product rule: $f'(x)=\sin x+x\cos x$.

  \item Simplify: $\ln(\sqrt{x})=\tfrac12\ln x$, so $f'(x)=\dfrac{1}{2x}$.

  \item Chain: $f'(x)=2\sin x\cdot\cos x=\sin(2x)$.

  \item $f'(x)=2^{x}\ln 2$ (general exponential rule).

  \item Inner $u=x^{2}+1$, $u'=2x$. $f'(x)=2x\,\mathrm{e}^{x^{2}+1}$.

  \item Chain on outer square: $f'(x)=2\ln x\cdot\dfrac{1}{x}
        =\dfrac{2\ln x}{x}$.

  \item Quotient rule:
        \[
          f'(x)=\frac{\cos x\cdot x-\sin x\cdot 1}{x^{2}}
               =\frac{x\cos x-\sin x}{x^{2}} .
        \]

  \item Product rule: $f'(x)=\arctan x+x\cdot\dfrac{1}{1+x^{2}}
        =\arctan x+\dfrac{x}{1+x^{2}}$.

  \item Inner $u=3x$, $u'=3$. $f'(x)=\operatorname{sech}^{2}(3x)\cdot 3
        =3\operatorname{sech}^{2}(3x)$.

  \item $f'(x)=\dfrac{1}{x\ln 10}$.

  \item Outer $\sin$, inner $\cos x$: $f'(x)=\cos(\cos x)\cdot(-\sin x)
        =-\sin x\cos(\cos x)$.

  \item Chain: $f'(x)=\mathrm{e}^{\sin x}\cdot\cos x=\cos x\,\mathrm{e}^{\sin x}$.

  \item Inner $u=1/x$, $u'=-1/x^{2}$.
        \[
          f'(x)=\frac{1}{1+(1/x)^{2}}\cdot\left(-\frac{1}{x^{2}}\right)
               =-\frac{1}{x^{2}+1} .
        \]
        (Aside: this confirms $\arctan x+\arctan(1/x)=\tfrac{\pi}{2}$ for
        $x>0$.)

  \item Expand $\ln\dfrac{x}{x+1}=\ln x-\ln(x+1)$, so
        $f'(x)=\dfrac{1}{x}-\dfrac{1}{x+1}=\dfrac{1}{x(x+1)}$.

  \item Inner $u=3x$, $u'=3$. $f'(x)=\sec(3x)\tan(3x)\cdot 3
        =3\sec(3x)\tan(3x)$.

  \item $U'(c)=\dfrac{1}{c}$. Marginal utility of consumption falls as $c$
        grows, diminishing marginal utility.

  \item Outer $\ln$, inner $\cos x$. $f'(x)=\dfrac{1}{\cos x}\cdot(-\sin x)
        =-\tan x$.

  \item Outer $u^{3}$, inner $\tan x$: $f'(x)=3\tan^{2}x\cdot\sec^{2}x$.

  \item Quotient rule with $u=1$, $v=1+\mathrm{e}^{-t}$. $u'=0$,
        $v'=-\mathrm{e}^{-t}$.
        \[
          P'(t)=\frac{\mathrm{e}^{-t}}{(1+\mathrm{e}^{-t})^{2}} .
        \]
        Equivalently, $P'(t)=P(t)\bigl(1-P(t)\bigr)$.

  \item Quotient rule. $u=\mathrm{e}^{x}$, $v=1+\mathrm{e}^{x}$,
        $u'=\mathrm{e}^{x}$, $v'=\mathrm{e}^{x}$.
        \[
          f'(x)=\frac{\mathrm{e}^{x}(1+\mathrm{e}^{x})-\mathrm{e}^{x}\cdot\mathrm{e}^{x}}
               {(1+\mathrm{e}^{x})^{2}}=\frac{\mathrm{e}^{x}}{(1+\mathrm{e}^{x})^{2}} .
        \]

  \item Product rule: $u=\mathrm{e}^{x}$, $v=\sin(2x)$, $u'=\mathrm{e}^{x}$,
        $v'=2\cos(2x)$.
        \[
          f'(x)=\mathrm{e}^{x}\bigl(\sin(2x)+2\cos(2x)\bigr) .
        \]

  \item Use the log power rule first: $f(K)=\alpha\ln K$, hence
        $f'(K)=\dfrac{\alpha}{K}$. (This is the elasticity-style structure
        that makes Cobb-Douglas convenient.)

  \item Logarithmic differentiation. Let $y=x^{x}$. Then $\ln y=x\ln x$.
        Differentiate:
        \[
          \frac{y'}{y}=\ln x+1,\qquad y'=x^{x}(\ln x+1) .
        \]

  \item Product rule with $u=\mathrm{e}^{2x}$, $v=\cos(3x)$,
        $u'=2\mathrm{e}^{2x}$, $v'=-3\sin(3x)$.
        \[
          f'(x)=2\mathrm{e}^{2x}\cos(3x)-3\mathrm{e}^{2x}\sin(3x)
               =\mathrm{e}^{2x}\bigl(2\cos(3x)-3\sin(3x)\bigr) .
        \]

  \item Inner $u=x+\sqrt{x^{2}+1}$. Note $u'=1+\dfrac{x}{\sqrt{x^{2}+1}}
        =\dfrac{\sqrt{x^{2}+1}+x}{\sqrt{x^{2}+1}}=\dfrac{u}{\sqrt{x^{2}+1}}$.
        Hence
        \[
          f'(x)=\frac{1}{u}\cdot\frac{u}{\sqrt{x^{2}+1}}=\frac{1}{\sqrt{x^{2}+1}} .
        \]
        (This is $\operatorname{arsinh}(x)$, so the result matches the
        derivative of inverse hyperbolic sine.)

  \item Product rule: $f'(x)=\cos x\cdot\cos x+\sin x\cdot(-\sin x)
        =\cos^{2}x-\sin^{2}x=\cos(2x)$.

  \item $D'(t)=-r\,\mathrm{e}^{-rt}=-r\,D(t)$. Discount factor decays at
        rate $r$; instantaneous proportional decay $D'/D=-r$.

  \item Inner $u=\mathrm{e}^{x}$, $u'=\mathrm{e}^{x}$.
        \[
          f'(x)=\frac{1}{1+\mathrm{e}^{2x}}\cdot\mathrm{e}^{x}
               =\frac{\mathrm{e}^{x}}{1+\mathrm{e}^{2x}} .
        \]

  \item Inner $u=x^{2}+1$, $u'=2x$. $f'(x)=\operatorname{sech}^{2}(x^{2}+1)
        \cdot 2x=2x\operatorname{sech}^{2}(x^{2}+1)$.

  \item Quotient rule with $u=\ln x$, $v=x$, $u'=1/x$, $v'=1$.
        \[
          f'(x)=\frac{(1/x)\cdot x-\ln x}{x^{2}}=\frac{1-\ln x}{x^{2}} .
        \]

  \item Product rule: $f'(x)=\mathrm{e}^{x}\ln x+\mathrm{e}^{x}\cdot\dfrac{1}{x}
        =\mathrm{e}^{x}\Bigl(\ln x+\dfrac{1}{x}\Bigr)$.

  \item $f(x)=(\sin x)^{1/2}$, inner $u'=\cos x$.
        \[
          f'(x)=\tfrac12(\sin x)^{-1/2}\cos x=\frac{\cos x}{2\sqrt{\sin x}} .
        \]

  \item $N'(t)=1000\cdot 0.04\,\mathrm{e}^{0.04t}=0.04\,N(t)$, so growth
        rate $N'/N=0.04=4\%$ per unit time.

  \item Rewrite $5^{x^{2}}=\mathrm{e}^{x^{2}\ln 5}$, so
        $f'(x)=\mathrm{e}^{x^{2}\ln 5}\cdot 2x\ln 5=2x\ln 5\cdot 5^{x^{2}}$.

  \item Inner $u=x^{2}$, $u'=2x$.
        \[
          f'(x)=\frac{2x}{\sqrt{1-x^{4}}} .
        \]

  \item $f'(x)=-\csc^{2}x-\csc x\cot x=-\csc x(\csc x+\cot x)$.

  \item Product rule: $f'(x)=\cosh^{2}x+\sinh^{2}x=\cosh(2x)$. (Use the
        hyperbolic double-angle identity.)

  \item Simplify first: $f(x)=\ln x+\tfrac12\ln(x^{2}+1)$, so
        \[
          f'(x)=\frac{1}{x}+\frac12\cdot\frac{2x}{x^{2}+1}
               =\frac{1}{x}+\frac{x}{x^{2}+1} .
        \]

  \item Outer $u^{4}$, inner $\cos x$, $u'=-\sin x$: $f'(x)=4\cos^{3}x\cdot
        (-\sin x)=-4\cos^{3}x\sin x$.

  \item Inner $u=x^{3}$, $u'=3x^{2}$. $f'(x)=3x^{2}\mathrm{e}^{x^{3}}$.

  \item $U'(c)=\dfrac{(1-\sigma)c^{-\sigma}}{1-\sigma}=c^{-\sigma}$.
        Marginal utility is positive and decreasing for $\sigma>0$.

  \item Inner $u=\sqrt{x}$, $u'=\dfrac{1}{2\sqrt{x}}$.
        \[
          f'(x)=\frac{1}{1+x}\cdot\frac{1}{2\sqrt{x}}
               =\frac{1}{2\sqrt{x}(1+x)} .
        \]

  \item Product rule: $f'(x)=\sec^{2}x\cdot\ln x+\tan x\cdot\dfrac{1}{x}$.

  \item $f(x)=(\ln x)^{-1}$, so $f'(x)=-(\ln x)^{-2}\cdot\dfrac{1}{x}
        =-\dfrac{1}{x(\ln x)^{2}}$.

  \item Product rule with $u=\mathrm{e}^{ax}$, $v=\sin(bx)$,
        $u'=a\mathrm{e}^{ax}$, $v'=b\cos(bx)$.
        \[
          f'(x)=\mathrm{e}^{ax}\bigl(a\sin(bx)+b\cos(bx)\bigr) .
        \]

  \item Use $\log_{2}(u)=\dfrac{\ln u}{\ln 2}$ with $u=x^{2}+1$:
        \[
          f'(x)=\frac{1}{\ln 2}\cdot\frac{2x}{x^{2}+1}
               =\frac{2x}{(x^{2}+1)\ln 2} .
        \]

  \item Product rule with $u=\cos(2x)$, $v=\mathrm{e}^{-x}$,
        $u'=-2\sin(2x)$, $v'=-\mathrm{e}^{-x}$.
        \[
          f'(x)=-2\sin(2x)\mathrm{e}^{-x}-\cos(2x)\mathrm{e}^{-x}
               =-\mathrm{e}^{-x}\bigl(2\sin(2x)+\cos(2x)\bigr) .
        \]

  \item Using $\operatorname{arccot}(x)=\tfrac{\pi}{2}-\arctan(x)$:
        \[
          f(x)=\Bigl(\frac{\pi}{2}-\arctan x\Bigr)\arctan x
              =\frac{\pi}{2}\arctan x-(\arctan x)^{2} .
        \]
        \[
          f'(x)=\frac{\pi}{2}\cdot\frac{1}{1+x^{2}}
               -2\arctan x\cdot\frac{1}{1+x^{2}}
               =\frac{\frac{\pi}{2}-2\arctan x}{1+x^{2}} .
        \]

  \item Outer $\ln$, inner $\ln x$. $f'(x)=\dfrac{1}{\ln x}\cdot\dfrac{1}{x}
        =\dfrac{1}{x\ln x}$.

  \item Product rule: $f'(x)=2x\,\mathrm{e}^{-x}+x^{2}(-\mathrm{e}^{-x})
        =x\,\mathrm{e}^{-x}(2-x)$.

  \item $g(t)=\ln(N_{0}\mathrm{e}^{rt})=\ln N_{0}+rt$, so $g'(t)=r$. Log
        of an exponential is linear: continuous growth rate is the slope
        of $\ln N$.

  \item Inner $u=\mathrm{e}^{x}$, $u'=\mathrm{e}^{x}$.
        $f'(x)=\cos(\mathrm{e}^{x})\cdot\mathrm{e}^{x}
        =\mathrm{e}^{x}\cos(\mathrm{e}^{x})$.

  \item Inner $u=1+\mathrm{e}^{x}$, $u'=\mathrm{e}^{x}$.
        $f'(x)=3(1+\mathrm{e}^{x})^{2}\cdot\mathrm{e}^{x}
        =3\mathrm{e}^{x}(1+\mathrm{e}^{x})^{2}$.

  \item Logarithmic differentiation. $\ln y=\sin x\cdot\ln x$.
        Differentiate:
        \[
          \frac{y'}{y}=\cos x\cdot\ln x+\sin x\cdot\frac{1}{x},
        \]
        so $y'=x^{\sin x}\Bigl(\cos x\ln x+\dfrac{\sin x}{x}\Bigr)$.

  \item For $x\in(-\tfrac{\pi}{2},\tfrac{\pi}{2})$ we have
        $\arctan(\tan x)=x$, so $f'(x)=1$. (Equivalently, chain rule:
        $\dfrac{1}{1+\tan^{2}x}\cdot\sec^{2}x=\dfrac{\sec^{2}x}{\sec^{2}x}=1$.)
\end{solutions}

\section{Additional Problems}

\begin{problems}
  \item $f(x)=\mathrm{e}^{5x}$
  \item $f(x)=\ln(7x+3)$
  \item $f(x)=\cos(4x)$
  \item $f(x)=\arctan(3x)$
  \item $f(x)=\mathrm{e}^{x^{3}}$
  \item $f(x)=\sin^{3}x$
  \item $f(x)=\ln(x^{3}+2)$
  \item $f(x)=x^{2}\,\mathrm{e}^{-x}$
  \item Continuous discount: $V(t)=1000\,\mathrm{e}^{-0.05t}$ \quad (find
        $V'(t)$)
  \item $f(x)=\dfrac{\ln x}{x^{2}}$
  \item $f(x)=\mathrm{e}^{x}\cos x$
  \item $f(x)=\arcsin(3x)$
  \item $f(x)=\tan(x^{2})$
  \item $f(x)=\ln(\sin x)$
  \item Log-utility variant: $U(x)=\ln(1+x)$ \quad (find $U'(x)$)
  \item $f(x)=3^{x^{2}+1}$
  \item $f(x)=(\ln x)^{3}$
  \item $f(x)=\mathrm{e}^{x^{2}}\sin x$
  \item $f(x)=\dfrac{\mathrm{e}^{x}-1}{\mathrm{e}^{x}+1}$
  \item $f(x)=(\cos x)^{x}$ \quad (use logarithmic differentiation,
        $\cos x>0$)
  \item $f(x)=\sinh(x)-\cosh(x)$
  \item $f(x)=\ln\Bigl(\dfrac{x^{2}}{x+1}\Bigr)$
  \item Per-capita capital: $k(t)=\ln(K(t)/L(t))$ with
        $K(t)=K_{0}\mathrm{e}^{at}$, $L(t)=L_{0}\mathrm{e}^{bt}$ \quad
        (find $k'(t)$)
  \item $f(x)=\arctan(\ln x)$
  \item $f(x)=\mathrm{e}^{\cos x}$
\end{problems}

\section{Additional Solutions}

\begin{solutions}
  \item Chain: $f'(x)=5\mathrm{e}^{5x}$.

  \item Chain: $f'(x)=\dfrac{7}{7x+3}$.

  \item Chain: $f'(x)=-4\sin(4x)$.

  \item Chain: $f'(x)=\dfrac{3}{1+9x^{2}}$.

  \item Chain: inner $u=x^{3}$, $u'=3x^{2}$. $f'(x)=3x^{2}\mathrm{e}^{x^{3}}$.

  \item Chain: $f'(x)=3\sin^{2}(x)\cos(x)$.

  \item Chain: $f'(x)=\dfrac{3x^{2}}{x^{3}+2}$.

  \item Product rule: $f'(x)=2x\,\mathrm{e}^{-x}-x^{2}\mathrm{e}^{-x}
        =x(2-x)\,\mathrm{e}^{-x}$.

  \item $V'(t)=1000\cdot(-0.05)\,\mathrm{e}^{-0.05t}=-50\,\mathrm{e}^{-0.05t}
        =-0.05\,V(t)$. Decays at rate $5\%$.

  \item Quotient rule with $u=\ln x$, $v=x^{2}$, $u'=1/x$, $v'=2x$.
        \[
          f'(x)=\frac{(1/x)x^{2}-\ln x\cdot 2x}{x^{4}}
               =\frac{x-2x\ln x}{x^{4}}=\frac{1-2\ln x}{x^{3}} .
        \]

  \item Product rule: $f'(x)=\mathrm{e}^{x}\cos x-\mathrm{e}^{x}\sin x
        =\mathrm{e}^{x}(\cos x-\sin x)$.

  \item Chain: $f'(x)=\dfrac{3}{\sqrt{1-9x^{2}}}$.

  \item Chain: inner $u=x^{2}$, $u'=2x$. $f'(x)=2x\sec^{2}(x^{2})$.

  \item Chain: $f'(x)=\dfrac{\cos x}{\sin x}=\cot x$.

  \item $U'(x)=\dfrac{1}{1+x}$. Marginal utility positive but diminishing
        as $x$ grows.

  \item Write $3^{x^{2}+1}=\mathrm{e}^{(x^{2}+1)\ln 3}$.
        $f'(x)=\mathrm{e}^{(x^{2}+1)\ln 3}\cdot 2x\ln 3
        =2x\ln 3\cdot 3^{x^{2}+1}$.

  \item Chain: $f'(x)=3(\ln x)^{2}\cdot\dfrac{1}{x}=\dfrac{3(\ln x)^{2}}{x}$.

  \item Product rule with $u=\mathrm{e}^{x^{2}}$, $v=\sin x$,
        $u'=2x\,\mathrm{e}^{x^{2}}$, $v'=\cos x$.
        \[
          f'(x)=2x\,\mathrm{e}^{x^{2}}\sin x+\mathrm{e}^{x^{2}}\cos x
               =\mathrm{e}^{x^{2}}(2x\sin x+\cos x) .
        \]

  \item Quotient rule with $u=\mathrm{e}^{x}-1$, $v=\mathrm{e}^{x}+1$,
        $u'=\mathrm{e}^{x}$, $v'=\mathrm{e}^{x}$.
        \[
          f'(x)=\frac{\mathrm{e}^{x}(\mathrm{e}^{x}+1)
                -(\mathrm{e}^{x}-1)\mathrm{e}^{x}}{(\mathrm{e}^{x}+1)^{2}}
               =\frac{2\mathrm{e}^{x}}{(\mathrm{e}^{x}+1)^{2}} .
        \]

  \item Logarithmic differentiation. $\ln y=x\ln(\cos x)$. Differentiate:
        \[
          \frac{y'}{y}=\ln(\cos x)+x\cdot\frac{-\sin x}{\cos x}
                      =\ln(\cos x)-x\tan x .
        \]
        So $y'=(\cos x)^{x}\bigl[\ln(\cos x)-x\tan x\bigr]$.

  \item $f'(x)=\cosh x-\sinh x$. (Equivalently, $f(x)=\mathrm{e}^{-x}$ and
        $f'(x)=-\mathrm{e}^{-x}=\cosh x-\sinh x$.)

  \item Simplify: $f(x)=2\ln x-\ln(x+1)$, so
        $f'(x)=\dfrac{2}{x}-\dfrac{1}{x+1}=\dfrac{x+2}{x(x+1)}$.

  \item $k(t)=\ln K_{0}+at-\ln L_{0}-bt=(a-b)t+\text{const}$, so
        $k'(t)=a-b$ (constant; difference of growth rates).

  \item Chain: outer $\arctan$, inner $\ln x$.
        $f'(x)=\dfrac{1}{1+(\ln x)^{2}}\cdot\dfrac{1}{x}
        =\dfrac{1}{x(1+(\ln x)^{2})}$.

  \item Chain: $f'(x)=\mathrm{e}^{\cos x}\cdot(-\sin x)
        =-\sin x\,\mathrm{e}^{\cos x}$.
\end{solutions}

\chapter{Higher-Order Derivatives, Convexity, Inflection}

Once $f'$ exists and is itself differentiable, we can differentiate again to
obtain $f''$, and iterate to higher orders. The second derivative records
how the slope changes: it controls convexity/concavity and powers the
second-order condition for optimisation. Higher derivatives of standard
functions follow simple closed-form patterns, and the Leibniz rule extends
the product rule to all orders.

\section{Core Rules}

\begin{rulebox}{Definition}
For $f:I\subseteq\R\to\R$ smooth enough, write
\[
  f''(x)=\frac{\mathrm{d}^{2}f}{\mathrm{d}x^{2}},\qquad
  f^{(n)}(x)=\frac{\mathrm{d}^{n}f}{\mathrm{d}x^{n}},\qquad
  f^{(0)}:=f .
\]
\end{rulebox}

\begin{rulebox}{Polynomials}
\[
  \frac{\mathrm{d}^{n}}{\mathrm{d}x^{n}}x^{k}
    =k(k-1)\cdots(k-n+1)\,x^{k-n} .
\]
In particular $(x^{n})^{(n)}=n!$ and $(x^{n})^{(n+1)}=0$.
\end{rulebox}

\begin{rulebox}{Exponentials}
\[
  \frac{\mathrm{d}^{n}}{\mathrm{d}x^{n}}\mathrm{e}^{ax}=a^{n}\mathrm{e}^{ax} .
\]
\end{rulebox}

\begin{rulebox}{Logarithm}
\[
  \frac{\mathrm{d}^{n}}{\mathrm{d}x^{n}}\ln(x)
    =\frac{(-1)^{n-1}(n-1)!}{x^{n}}\qquad\text{for } n\geq 1 .
\]
\end{rulebox}

\begin{rulebox}{Trig}
\[
  \frac{\mathrm{d}^{n}}{\mathrm{d}x^{n}}\sin(x)=\sin\Bigl(x+\frac{n\pi}{2}\Bigr),
  \qquad
  \frac{\mathrm{d}^{n}}{\mathrm{d}x^{n}}\cos(x)=\cos\Bigl(x+\frac{n\pi}{2}\Bigr) .
\]
\end{rulebox}

\begin{rulebox}{Convexity / concavity}
$f''\geq 0$ on $I$ implies $f$ is convex on $I$; $f''\leq 0$ implies
concave. Strict inequalities give strict (con)cavity.
\end{rulebox}

\begin{rulebox}{Inflection point}
$x_{0}$ is an inflection point if $f''$ \emph{changes sign} at $x_{0}$. The
condition $f''(x_{0})=0$ alone is insufficient (e.g.\ $x^{4}$ at $x_{0}=0$).
\end{rulebox}

\begin{rulebox}{Second-derivative test}
If $f'(x_{0})=0$ and $f''(x_{0})>0$, $x_{0}$ is a local minimum;
$f''(x_{0})<0$ gives a local maximum.
\end{rulebox}

\begin{rulebox}{Leibniz rule (higher product)}
\[
  (fg)^{(n)}=\sum_{k=0}^{n}\binom{n}{k}f^{(k)}g^{(n-k)} .
\]
\end{rulebox}

\begin{rulebox}{Taylor polynomial}
The Taylor polynomial of order $n$ at $a$:
\[
  T_{n}(x)=\sum_{k=0}^{n}\frac{f^{(k)}(a)}{k!}(x-a)^{k},
\]
with remainder
\[
  R_{n}(x)=\frac{f^{(n+1)}(\xi)}{(n+1)!}(x-a)^{n+1}
\]
for some $\xi$ between $a$ and $x$.
\end{rulebox}

\section{Practice Problems}

\begin{problems}
  \item Compute $f''(x)$ for $f(x)=2x^{4}-5x^{3}+x-7$.
  \item Determine the intervals on which $f(x)=x^{3}-6x^{2}+5$ is convex and
        on which it is concave.
  \item Find a closed form for $f^{(n)}(x)$ when $f(x)=\mathrm{e}^{-x}$.
  \item Find all inflection points of $f(x)=x^{6}$.
  \item Compute $f''(x)$ for $f(x)=\mathrm{e}^{-2x}$.
  \item For a cost function $C(Q)=Q^{3}-12Q^{2}+50Q+100$, determine the
        output level where marginal cost transitions from decreasing to
        increasing.
  \item Find a closed form for $f^{(n)}(x)$ when $f(x)=\ln(1+x)$.
  \item Compute $f''(x)$ for $f(x)=x^{2}\mathrm{e}^{x}$.
  \item Find all inflection points of $f(x)=x^{3}-9x^{2}+24x-3$.
  \item Determine the convexity/concavity intervals of $f(x)=x\mathrm{e}^{-x}$.
  \item Compute $f''(x)$ for $f(x)=\ln(2x+5)$.
  \item Find a closed form for $f^{(n)}(x)$ when $f(x)=\dfrac{1}{x+a}$
        (constant $a$).
  \item Show that $f(x)=x^{4}$ has $f''(0)=0$ but $x=0$ is \emph{not} an
        inflection point.
  \item Compute $f''(x)$ for $f(x)=3x^{2/3}-x^{-2}$.
  \item Determine the convexity intervals of $f(x)=\ln(1+\mathrm{e}^{x})$
        (softplus).
  \item Find all inflection points of $f(x)=x^{4}-12x^{2}+5$.
  \item Compute $f''(x)$ for $f(x)=x^{2}\ln(x)$.
  \item Find a closed form for $f^{(n)}(x)$ when $f(x)=\cos(x)$.
  \item For utility $U(x)=\ln(x)$, $x>0$, verify that $U$ is strictly
        concave (the standard ``risk-aversion'' condition).
  \item Compute $f''(x)$ for $f(x)=\tan(x)$.
  \item Find all inflection points of $f(x)=\dfrac{x^{2}}{x^{2}+3}$.
  \item Determine convexity/concavity intervals of $f(x)=\mathrm{e}^{x}-x$.
  \item Compute $f''(x)$ for $f(x)=\sin(x)\cos(x)$.
  \item Find a closed form for $f^{(n)}(x)$ when $f(x)=\dfrac{1}{(1-x)^{2}}$.
  \item Find all inflection points of $f(x)=x^{2}\mathrm{e}^{-x}$.
  \item Compute $f''(x)$ for $f(x)=\dfrac{x}{x-3}$.
  \item Determine convexity/concavity of $f(x)=x^{4}-8x^{3}+1$ and the
        corresponding inflection points.
  \item Find a closed form for $f^{(n)}(x)$ when $f(x)=\mathrm{e}^{ax}$ for
        a constant $a$.
  \item Compute $f''(x)$ for $f(x)=(x^{2}+4)^{3/2}$.
  \item For $f(x)=x^{5}$, determine whether $x=0$ is an inflection point.
  \item Compute $f''(x)$ for $f(x)=\sqrt{1+x^{2}}$.
  \item Determine the inflection points of $f(x)=\arctan(x)$.
  \item Find $f^{(n)}(x)$ for $f(x)=\sin(2x)$.
  \item Compute $f''(x)$ for $f(x)=\dfrac{1}{1+x^{2}}$.
  \item Determine convexity/concavity of $f(x)=x\ln(x)$, $x>0$.
  \item For the logistic function $f(x)=\dfrac{1}{1+\mathrm{e}^{-x}}$, find
        the inflection point.
  \item Compute $f''(x)$ for $f(x)=x^{3}\ln(x)$.
  \item For a profit function $\pi(Q)=-Q^{3}+15Q^{2}-27Q-10$, where is
        $\pi$ concave?
  \item Find $f^{(n)}(0)$ for $f(x)=\ln(1+x)$ at $x=0$ (i.e.\ the
        coefficients in its Taylor series at zero).
  \item Compute $f''(x)$ for $f(x)=\mathrm{e}^{x^{2}}$.
  \item Determine all inflection points of $f(x)=x^{5}-5x^{4}$.
  \item For $f(x)=x|x|$ (i.e.\ $x^{2}$ for $x\geq 0$, $-x^{2}$ for $x<0$),
        determine whether $x=0$ is an inflection point and whether
        $f''(0)$ exists.
  \item Compute $f''(x)$ for $f(x)=\arctan(x)$.
  \item Apply the Leibniz rule to compute $(x^{2}\mathrm{e}^{x})^{(n)}$ in
        closed form.
  \item Find the inflection points of $f(x)=(x-1)^{4}$.
  \item For $f(x)=\sin(x)$, find the Taylor polynomial $T_{5}(x)$ at $a=0$.
  \item Compute $f''(x)$ for $f(x)=\ln(\cos(x))$ on
        $\bigl(-\tfrac{\pi}{2},\tfrac{\pi}{2}\bigr)$.
  \item Find a closed form for $f^{(n)}(x)$ when $f(x)=x\mathrm{e}^{x}$,
        using the Leibniz rule.
  \item Determine the convexity intervals of $f(x)=\dfrac{1}{1+x^{2}}$ (the
        standard Cauchy-type bell curve).
  \item For a Cobb--Douglas production function $f(L)=AL^{\alpha}$ with
        $0<\alpha<1$ and $A>0$, verify that $f$ is strictly concave in
        $L>0$ (the law of diminishing returns).
\end{problems}

\section{Solutions}

\begin{solutions}
  \item $f'(x)=8x^{3}-15x^{2}+1$, $f''(x)=24x^{2}-30x$.

  \item $f'(x)=3x^{2}-12x$, $f''(x)=6x-12=6(x-2)$. So $f''(x)\geq 0$ iff
        $x\geq 2$: convex on $[2,\infty)$, concave on $(-\infty,2]$.

  \item Each differentiation introduces a factor of $-1$:
        $f^{(n)}(x)=(-1)^{n}\mathrm{e}^{-x}$.

  \item $f'(x)=6x^{5}$, $f''(x)=30x^{4}\geq 0$ for all $x$, with equality
        only at $x=0$. Since $f''$ does not change sign, $x=0$ is
        \emph{not} an inflection point. $f$ has no inflection points.

  \item $f'(x)=-2\mathrm{e}^{-2x}$, $f''(x)=4\mathrm{e}^{-2x}$.

  \item $MC(Q)=C'(Q)=3Q^{2}-24Q+50$. $MC'(Q)=6Q-24=0\Rightarrow Q=4$.
        $MC''(Q)=6>0$ so $Q=4$ minimises $MC$. At $Q=4$, $C''(Q)$ changes
        sign there: $C''(Q)=6Q-24$, negative for $Q<4$, positive for
        $Q>4$, so $Q=4$ is the inflection point of $C$ and the transition
        output for marginal cost.

  \item With $u=1+x$: $f'(x)=u^{-1}$, $f''(x)=-u^{-2}$,
        $f'''(x)=2u^{-3}$, in general
        \[
          f^{(n)}(x)=\frac{(-1)^{n-1}(n-1)!}{(1+x)^{n}}\qquad (n\geq 1) .
        \]

  \item $f'(x)=2x\mathrm{e}^{x}+x^{2}\mathrm{e}^{x}=(x^{2}+2x)\mathrm{e}^{x}$,
        then
        \[
          f''(x)=(2x+2)\mathrm{e}^{x}+(x^{2}+2x)\mathrm{e}^{x}
                =(x^{2}+4x+2)\mathrm{e}^{x} .
        \]

  \item $f'(x)=3x^{2}-18x+24$, $f''(x)=6x-18=6(x-3)$. Sign change at
        $x=3$. $f(3)=27-81+72-3=15$. Inflection at $(3,15)$.

  \item $f'(x)=\mathrm{e}^{-x}-x\mathrm{e}^{-x}=(1-x)\mathrm{e}^{-x}$.
        $f''(x)=-\mathrm{e}^{-x}-(1-x)\mathrm{e}^{-x}=(x-2)\mathrm{e}^{-x}$.
        Since $\mathrm{e}^{-x}>0$, the sign of $f''$ is that of $x-2$.
        Concave on $(-\infty,2]$, convex on $[2,\infty)$. Inflection at
        $x=2$.

  \item $f'(x)=\dfrac{2}{2x+5}$, $f''(x)=-\dfrac{4}{(2x+5)^{2}}$.

  \item $f(x)=(x+a)^{-1}$, so
        $f^{(n)}(x)=(-1)^{n}n!\,(x+a)^{-(n+1)}=\dfrac{(-1)^{n}n!}{(x+a)^{n+1}}$.

  \item $f'(x)=4x^{3}$, $f''(x)=12x^{2}$. Then $f''(0)=0$. But
        $f''(x)=12x^{2}\geq 0$ everywhere, no sign change. Hence $x=0$
        is not an inflection point. Indeed $f$ is convex on all of $\R$.

  \item $f'(x)=2x^{-1/3}+2x^{-3}$,
        $f''(x)=-\tfrac{2}{3}x^{-4/3}-6x^{-4}$.

  \item $f'(x)=\dfrac{\mathrm{e}^{x}}{1+\mathrm{e}^{x}}$. By the quotient
        rule
        \[
          f''(x)=\frac{\mathrm{e}^{x}(1+\mathrm{e}^{x})-\mathrm{e}^{x}\cdot\mathrm{e}^{x}}
                       {(1+\mathrm{e}^{x})^{2}}
                =\frac{\mathrm{e}^{x}}{(1+\mathrm{e}^{x})^{2}}>0
          \text{ for all } x .
        \]
        $f$ is strictly convex on all of $\R$.

  \item $f'(x)=4x^{3}-24x$, $f''(x)=12x^{2}-24=12(x^{2}-2)$. Sign changes
        at $x=\pm\sqrt{2}$. $f(\pm\sqrt{2})=4-24+5=-15$. Inflections at
        $(\pm\sqrt{2},-15)$.

  \item $f'(x)=2x\ln(x)+x$. Then $f''(x)=2\ln(x)+2+1=2\ln(x)+3$.

  \item $f^{(n)}(x)=\cos\bigl(x+\tfrac{n\pi}{2}\bigr)$.

  \item $U'(x)=1/x$, $U''(x)=-1/x^{2}<0$ for all $x>0$. So $U$ is strictly
        concave on $(0,\infty)$.

  \item $f'(x)=\sec^{2}(x)$. Then
        $f''(x)=2\sec(x)\cdot\sec(x)\tan(x)=2\sec^{2}(x)\tan(x)$.

  \item Write $f(x)=1-\dfrac{3}{x^{2}+3}$. Then
        $f'(x)=\dfrac{6x}{(x^{2}+3)^{2}}$ and
        \[
          f''(x)=\frac{6(x^{2}+3)^{2}-6x\cdot 2(x^{2}+3)\cdot 2x}
                       {(x^{2}+3)^{4}}
                =\frac{6(x^{2}+3)-24x^{2}}{(x^{2}+3)^{3}}
                =\frac{18-18x^{2}}{(x^{2}+3)^{3}} .
        \]
        Sign changes at $x=\pm 1$. $f(\pm 1)=1/4$. Inflections at
        $(\pm 1,\tfrac14)$.

  \item $f'(x)=\mathrm{e}^{x}-1$, $f''(x)=\mathrm{e}^{x}>0$ for all $x$.
        Strictly convex on $\R$. No concave region.

  \item $f(x)=\tfrac12\sin(2x)$, $f'(x)=\cos(2x)$, $f''(x)=-2\sin(2x)$.

  \item Note $\dfrac{1}{(1-x)^{2}}=\dfrac{\mathrm{d}}{\mathrm{d}x}
        \dfrac{1}{1-x}$. Since
        $\dfrac{\mathrm{d}^{k}}{\mathrm{d}x^{k}}\dfrac{1}{1-x}
        =\dfrac{k!}{(1-x)^{k+1}}$, we get
        \[
          f^{(n)}(x)=\frac{(n+1)!}{(1-x)^{n+2}} .
        \]

  \item $f'(x)=2x\mathrm{e}^{-x}-x^{2}\mathrm{e}^{-x}=(2x-x^{2})\mathrm{e}^{-x}$.
        $f''(x)=(2-2x)\mathrm{e}^{-x}-(2x-x^{2})\mathrm{e}^{-x}
               =(x^{2}-4x+2)\mathrm{e}^{-x}$. Roots: $x=2\pm\sqrt{2}$.
        Inflections at $x=2\pm\sqrt{2}$; $y$-values
        $(2\pm\sqrt{2})^{2}\mathrm{e}^{-(2\pm\sqrt{2})}
        =(6\pm4\sqrt{2})\mathrm{e}^{-(2\pm\sqrt{2})}$.

  \item Write $f(x)=1+\dfrac{3}{x-3}$. Then $f'(x)=-\dfrac{3}{(x-3)^{2}}$,
        $f''(x)=\dfrac{6}{(x-3)^{3}}$.

  \item $f'(x)=4x^{3}-24x^{2}$, $f''(x)=12x^{2}-48x=12x(x-4)$. Sign
        changes at $x=0$ and $x=4$. Convex on $(-\infty,0]\cup[4,\infty)$,
        concave on $[0,4]$. Inflections at $(0,1)$ and
        $(4,4^{4}-8\cdot 64+1)=(4,256-512+1)=(4,-255)$.

  \item Each derivative pulls out a factor $a$: $f^{(n)}(x)=a^{n}\mathrm{e}^{ax}$.

  \item $f'(x)=\tfrac32(x^{2}+4)^{1/2}\cdot 2x=3x\sqrt{x^{2}+4}$. Then
        \[
          f''(x)=3\sqrt{x^{2}+4}+3x\cdot\frac{x}{\sqrt{x^{2}+4}}
                =\frac{3(x^{2}+4)+3x^{2}}{\sqrt{x^{2}+4}}
                =\frac{6x^{2}+12}{\sqrt{x^{2}+4}}
                =\frac{6(x^{2}+2)}{\sqrt{x^{2}+4}} .
        \]

  \item $f'(x)=5x^{4}$, $f''(x)=20x^{3}$. $f''(x)<0$ for $x<0$,
        $f''(x)>0$ for $x>0$: sign change at $x=0$. Hence $x=0$ is an
        inflection point (with $f(0)=0$).

  \item $f'(x)=\dfrac{x}{\sqrt{1+x^{2}}}$. Then
        \[
          f''(x)=\frac{\sqrt{1+x^{2}}-x\cdot\frac{x}{\sqrt{1+x^{2}}}}
                       {1+x^{2}}
                =\frac{(1+x^{2})-x^{2}}{(1+x^{2})^{3/2}}
                =\frac{1}{(1+x^{2})^{3/2}} .
        \]

  \item $f'(x)=\dfrac{1}{1+x^{2}}$. Then
        $f''(x)=-\dfrac{2x}{(1+x^{2})^{2}}$. Sign change at $x=0$, with
        $f''(x)>0$ for $x<0$ and $f''(x)<0$ for $x>0$. Inflection at
        $(0,0)$.

  \item Write $\sin(2x)$. Each differentiation multiplies the
        argument-shift trick by $2$:
        \[
          f^{(n)}(x)=2^{n}\sin\Bigl(2x+\frac{n\pi}{2}\Bigr) .
        \]

  \item $f(x)=(1+x^{2})^{-1}$. $f'(x)=-2x(1+x^{2})^{-2}$. Then
        \[
          f''(x)=-2(1+x^{2})^{-2}+(-2x)\cdot(-2)(1+x^{2})^{-3}\cdot 2x
                =-\frac{2}{(1+x^{2})^{2}}+\frac{8x^{2}}{(1+x^{2})^{3}}
                =\frac{-2(1+x^{2})+8x^{2}}{(1+x^{2})^{3}}
                =\frac{6x^{2}-2}{(1+x^{2})^{3}} .
        \]

  \item $f'(x)=\ln(x)+1$, $f''(x)=1/x>0$ for all $x>0$. Strictly convex
        on $(0,\infty)$.

  \item $f(x)=\dfrac{1}{1+\mathrm{e}^{-x}}=\dfrac{\mathrm{e}^{x}}{1+\mathrm{e}^{x}}$.
        Then $f'(x)=\dfrac{\mathrm{e}^{x}}{(1+\mathrm{e}^{x})^{2}}$.
        Differentiate again using the quotient rule (with $u=\mathrm{e}^{x}$,
        $v=(1+\mathrm{e}^{x})^{2}$, $v'=2(1+\mathrm{e}^{x})\mathrm{e}^{x}$):
        \[
          f''(x)=\frac{\mathrm{e}^{x}(1+\mathrm{e}^{x})^{2}
                       -\mathrm{e}^{x}\cdot 2(1+\mathrm{e}^{x})\mathrm{e}^{x}}
                       {(1+\mathrm{e}^{x})^{4}}
                =\frac{\mathrm{e}^{x}(1+\mathrm{e}^{x})-2\mathrm{e}^{2x}}
                       {(1+\mathrm{e}^{x})^{3}}
                =\frac{\mathrm{e}^{x}(1-\mathrm{e}^{x})}{(1+\mathrm{e}^{x})^{3}} .
        \]
        Sign of $f''$ is the sign of $1-\mathrm{e}^{x}$: positive for
        $x<0$, negative for $x>0$. Inflection at $x=0$, where
        $f(0)=1/2$. Inflection point $(0,\tfrac12)$.

  \item $f'(x)=3x^{2}\ln(x)+x^{2}=x^{2}(3\ln(x)+1)$. Then
        \[
          f''(x)=2x(3\ln(x)+1)+x^{2}\cdot\frac{3}{x}
                =2x(3\ln(x)+1)+3x
                =6x\ln(x)+5x=x(6\ln(x)+5) .
        \]

  \item $\pi'(Q)=-3Q^{2}+30Q-27$, $\pi''(Q)=-6Q+30=6(5-Q)$. $\pi$ is
        concave on $\{Q:\pi''(Q)\leq 0\}=[5,\infty)$ and convex on
        $(-\infty,5]$. (Economic note: the profit function changes from
        convex to concave at the inflection $Q=5$.)

  \item $f^{(n)}(x)=\dfrac{(-1)^{n-1}(n-1)!}{(1+x)^{n}}$ for $n\geq 1$. At
        $x=0$: $f^{(n)}(0)=(-1)^{n-1}(n-1)!$ for $n\geq 1$, and $f(0)=0$.
        So the Taylor series is
        $\displaystyle\sum_{n\geq 1}\frac{(-1)^{n-1}}{n}x^{n}
        =x-\frac{x^{2}}{2}+\frac{x^{3}}{3}-\cdots$.

  \item $f'(x)=2x\mathrm{e}^{x^{2}}$. Then
        $f''(x)=2\mathrm{e}^{x^{2}}+2x\cdot 2x\mathrm{e}^{x^{2}}
               =(2+4x^{2})\mathrm{e}^{x^{2}}$.

  \item $f'(x)=5x^{4}-20x^{3}$, $f''(x)=20x^{3}-60x^{2}=20x^{2}(x-3)$.
        Now $20x^{2}\geq 0$, so the sign of $f''$ is determined by $x-3$
        except at $x=0$ where $f''=0$. Check signs: for $x<0$, $x^{2}>0$
        and $x-3<0$, so $f''<0$. For $0<x<3$, $x^{2}>0$ and $x-3<0$, so
        $f''<0$. No sign change at $x=0$. For $x>3$, $f''>0$. Sign change
        only at $x=3$. $f(3)=243-405=-162$. Inflection at $(3,-162)$.
        ($x=0$ is \emph{not} an inflection.)

  \item For $x>0$: $f(x)=x^{2}$, $f''(x)=2$. For $x<0$: $f(x)=-x^{2}$,
        $f''(x)=-2$. So $f''$ has different one-sided limits at $0$ ($2$
        from the right, $-2$ from the left). Hence $f''(0)$ does not
        exist. But $f'(x)=2|x|$ exists and is continuous everywhere, and
        $f''$ changes sign at $0$ (in the sense of one-sided values). The
        graph changes from concave (left) to convex (right) at $x=0$, so
        $x=0$ is an inflection point in the geometric sense even though
        $f''(0)$ itself is undefined.

  \item $f'(x)=\dfrac{1}{1+x^{2}}$. Then
        $f''(x)=-\dfrac{2x}{(1+x^{2})^{2}}$.

  \item Take $u=x^{2}$ and $v=\mathrm{e}^{x}$. Then $u'=2x$, $u''=2$,
        $u^{(k)}=0$ for $k\geq 3$, and $v^{(k)}=\mathrm{e}^{x}$. By
        Leibniz:
        \[
          (x^{2}\mathrm{e}^{x})^{(n)}
            =\binom{n}{0}x^{2}\mathrm{e}^{x}
             +\binom{n}{1}\cdot 2x\cdot\mathrm{e}^{x}
             +\binom{n}{2}\cdot 2\cdot\mathrm{e}^{x}
            =\mathrm{e}^{x}\bigl[x^{2}+2nx+n(n-1)\bigr] .
        \]

  \item $f'(x)=4(x-1)^{3}$, $f''(x)=12(x-1)^{2}\geq 0$ everywhere, with
        equality only at $x=1$. No sign change. Hence $x=1$ is \emph{not}
        an inflection point; $f$ is convex on all of $\R$.

  \item Using $f^{(k)}(0)=\sin(k\pi/2)$: values are $0,1,0,-1,0,1$ for
        $k=0,1,2,3,4,5$. Hence
        \[
          T_{5}(x)=x-\frac{x^{3}}{3!}+\frac{x^{5}}{5!}
                  =x-\frac{x^{3}}{6}+\frac{x^{5}}{120} .
        \]

  \item $f'(x)=-\tan(x)$. Then $f''(x)=-\sec^{2}(x)<0$ on
        $\bigl(-\tfrac{\pi}{2},\tfrac{\pi}{2}\bigr)$. So $f$ is strictly
        concave on that interval.

  \item Take $u=x$, $v=\mathrm{e}^{x}$. Then $u'=1$, $u^{(k)}=0$ for
        $k\geq 2$, and $v^{(k)}=\mathrm{e}^{x}$. By Leibniz:
        \[
          (x\mathrm{e}^{x})^{(n)}
            =\binom{n}{0}x\mathrm{e}^{x}+\binom{n}{1}\cdot 1\cdot\mathrm{e}^{x}
            =(x+n)\mathrm{e}^{x} .
        \]

  \item From problem 34: $f''(x)=\dfrac{6x^{2}-2}{(1+x^{2})^{3}}$. Sign of
        $f''$ equals sign of $6x^{2}-2$, i.e.\ positive iff
        $|x|>1/\sqrt{3}$. So $f$ is convex on
        $(-\infty,-1/\sqrt{3}]\cup[1/\sqrt{3},\infty)$ and concave on
        $[-1/\sqrt{3},1/\sqrt{3}]$. Inflection points at $x=\pm 1/\sqrt{3}$.

  \item $f'(L)=A\alpha L^{\alpha-1}$, $f''(L)=A\alpha(\alpha-1)L^{\alpha-2}$.
        With $A>0$, $0<\alpha<1$: $\alpha(\alpha-1)<0$ and
        $L^{\alpha-2}>0$, so $f''(L)<0$ for all $L>0$. Strictly concave on
        $(0,\infty)$, diminishing marginal product of labour.
\end{solutions}

\section{Additional Problems}

\begin{problems}
  \item Compute $f''(x)$ for $f(x)=3x^{5}-2x^{4}+x^{2}-7$.
  \item Determine the convexity/concavity intervals of
        $f(x)=x^{3}-3x^{2}-9x+5$.
  \item Find a closed form for $f^{(n)}(x)$ when $f(x)=\mathrm{e}^{3x}$.
  \item Find all inflection points of $f(x)=x^{4}-6x^{2}+8$.
  \item Compute $f''(x)$ for $f(x)=x\mathrm{e}^{-2x}$.
  \item For a cost function $C(Q)=Q^{3}-15Q^{2}+72Q+50$, determine the
        output level where marginal cost transitions from decreasing to
        increasing.
  \item Find a closed form for $f^{(n)}(x)$ when $f(x)=\ln(2-x)$, on
        $x<2$.
  \item Compute $f''(x)$ for $f(x)=\sqrt{x^{2}+4}$.
  \item Find all inflection points of $f(x)=x^{3}-12x^{2}+45x-7$.
  \item Determine the convexity intervals of $f(x)=\ln(x^{2}+1)$.
  \item Compute $f''(x)$ for $f(x)=\sin(x)\,\mathrm{e}^{x}$.
  \item Find a closed form for $f^{(n)}(x)$ when
        $f(x)=\dfrac{1}{(1-2x)}$.
  \item For utility $U(x)=\sqrt{x}$, $x>0$, verify that $U$ is strictly
        concave (diminishing marginal utility).
  \item Compute $f''(x)$ for $f(x)=x^{3}\mathrm{e}^{-x}$.
  \item Apply the Leibniz rule to compute $(x^{3}\mathrm{e}^{x})^{(n)}$ in
        closed form.
\end{problems}

\section{Additional Solutions}

\begin{solutions}
  \item $f'(x)=15x^{4}-8x^{3}+2x$, $f''(x)=60x^{3}-24x^{2}+2$.

  \item $f'(x)=3x^{2}-6x-9$, $f''(x)=6x-6=6(x-1)$. So $f''\geq 0$ iff
        $x\geq 1$: convex on $[1,\infty)$, concave on $(-\infty,1]$.
        Inflection at $x=1$, $f(1)=-6$.

  \item Each differentiation introduces a factor of $3$:
        $f^{(n)}(x)=3^{n}\mathrm{e}^{3x}$.

  \item $f'(x)=4x^{3}-12x$, $f''(x)=12x^{2}-12=12(x^{2}-1)$. Sign changes
        at $x=\pm 1$. $f(\pm 1)=1-6+8=3$. Inflections at $(\pm 1,3)$.

  \item Product rule:
        $f'(x)=\mathrm{e}^{-2x}-2x\mathrm{e}^{-2x}=(1-2x)\mathrm{e}^{-2x}$.
        Then
        \[
          f''(x)=-2\mathrm{e}^{-2x}-2(1-2x)\mathrm{e}^{-2x}
                =(4x-4)\mathrm{e}^{-2x}=4(x-1)\mathrm{e}^{-2x} .
        \]

  \item $MC(Q)=C'(Q)=3Q^{2}-30Q+72$. $MC'(Q)=6Q-30=0\Rightarrow Q=5$.
        $MC''(Q)=6>0$, so $Q=5$ minimises $MC$. Equivalently,
        $C''(Q)=6Q-30$ changes sign at $Q=5$: the cost function
        transitions from concave to convex at the inflection $Q=5$.

  \item With $u=2-x$: $f'(x)=-1/(2-x)$, $f''(x)=-1/(2-x)^{2}$,
        $f'''(x)=-2/(2-x)^{3}$. By induction:
        \[
          f^{(n)}(x)=-\frac{(n-1)!}{(2-x)^{n}}\qquad (n\geq 1) .
        \]

  \item $f(x)=(x^{2}+4)^{1/2}$. $f'(x)=\dfrac{x}{\sqrt{x^{2}+4}}$. Then
        \[
          f''(x)=\frac{\sqrt{x^{2}+4}-x\cdot\frac{x}{\sqrt{x^{2}+4}}}
                       {x^{2}+4}
                =\frac{(x^{2}+4)-x^{2}}{(x^{2}+4)^{3/2}}
                =\frac{4}{(x^{2}+4)^{3/2}} .
        \]

  \item $f'(x)=3x^{2}-24x+45$, $f''(x)=6x-24=6(x-4)$. Sign change at
        $x=4$. $f(4)=64-192+180-7=45$. Inflection at $(4,45)$.

  \item $f'(x)=\dfrac{2x}{x^{2}+1}$. By the quotient rule
        \[
          f''(x)=\frac{2(x^{2}+1)-2x\cdot 2x}{(x^{2}+1)^{2}}
                =\frac{2-2x^{2}}{(x^{2}+1)^{2}}
                =\frac{2(1-x^{2})}{(x^{2}+1)^{2}} .
        \]
        Sign of $f''$ matches sign of $1-x^{2}$. Convex on $[-1,1]$,
        concave on $(-\infty,-1]\cup[1,\infty)$. Inflections at
        $x=\pm 1$.

  \item $f'(x)=\cos(x)\mathrm{e}^{x}+\sin(x)\mathrm{e}^{x}
        =(\cos x+\sin x)\mathrm{e}^{x}$. Then
        \[
          f''(x)=(-\sin x+\cos x)\mathrm{e}^{x}
                 +(\cos x+\sin x)\mathrm{e}^{x}
                =2\cos(x)\,\mathrm{e}^{x} .
        \]

  \item $f(x)=(1-2x)^{-1}$. By induction:
        $f^{(n)}(x)=\dfrac{2^{n}\,n!}{(1-2x)^{n+1}}$.

  \item $U'(x)=\tfrac12 x^{-1/2}$, $U''(x)=-\tfrac14 x^{-3/2}<0$ for all
        $x>0$. So $U$ is strictly concave on $(0,\infty)$.

  \item $f'(x)=3x^{2}\mathrm{e}^{-x}-x^{3}\mathrm{e}^{-x}
        =(3x^{2}-x^{3})\mathrm{e}^{-x}$. Then
        \[
          f''(x)=(6x-3x^{2})\mathrm{e}^{-x}-(3x^{2}-x^{3})\mathrm{e}^{-x}
                =(x^{3}-6x^{2}+6x)\mathrm{e}^{-x}
                =x(x^{2}-6x+6)\mathrm{e}^{-x} .
        \]

  \item Take $u=x^{3}$ and $v=\mathrm{e}^{x}$. Then $u'=3x^{2}$,
        $u''=6x$, $u'''=6$, $u^{(k)}=0$ for $k\geq 4$, and
        $v^{(k)}=\mathrm{e}^{x}$. By Leibniz:
        \[
          (x^{3}\mathrm{e}^{x})^{(n)}
            =\binom{n}{0}x^{3}\mathrm{e}^{x}
             +\binom{n}{1}\cdot 3x^{2}\cdot\mathrm{e}^{x}
             +\binom{n}{2}\cdot 6x\cdot\mathrm{e}^{x}
             +\binom{n}{3}\cdot 6\cdot\mathrm{e}^{x}
            =\mathrm{e}^{x}\bigl[x^{3}+3nx^{2}+3n(n-1)x+n(n-1)(n-2)\bigr] .
        \]
\end{solutions}

\chapter{Taylor Approximation}

The previous chapter closed with a formula for the Taylor polynomial;
this chapter turns it into a working tool. Most functions that appear in
economics (utility, production, growth) are not polynomials, yet
polynomials are the easiest objects to evaluate and to reason about.
Taylor's theorem says that near a chosen base point every sufficiently
smooth function \emph{is} a polynomial, up to a controllable error. This
single idea underlies the linearisation of models around a steady state,
the log-linearisation trick $\ln(1+g)\approx g$ for growth rates, the
second-order conditions of optimisation, and quick numerical estimates
such as $\sqrt{4.1}$ or $\mathrm{e}^{0.1}$ without a calculator.

\section{Core Rules}

\begin{rulebox}{Taylor polynomial of order $n$}
Let $f$ be $n$ times differentiable at $a$. The \emph{Taylor polynomial}
of order $n$ of $f$ at $a$ is
\[
  T_{n}(x)=\sum_{k=0}^{n}\frac{f^{(k)}(a)}{k!}(x-a)^{k}
          =f(a)+f'(a)(x-a)+\frac{f''(a)}{2!}(x-a)^{2}+\cdots
           +\frac{f^{(n)}(a)}{n!}(x-a)^{n}.
\]
It is the \emph{unique} polynomial of degree $\leq n$ whose value and
first $n$ derivatives at $a$ all agree with those of $f$:
$T_{n}^{(k)}(a)=f^{(k)}(a)$ for $k=0,1,\ldots,n$.
\end{rulebox}

\begin{rulebox}{Linear and quadratic approximation ($n=1,2$)}
\[
  T_{1}(x)=f(a)+f'(a)(x-a),
  \qquad
  T_{2}(x)=f(a)+f'(a)(x-a)+\tfrac12 f''(a)(x-a)^{2}.
\]
$T_{1}$ is the tangent line at $a$ (the \emph{linearisation} of $f$);
$T_{2}$ is the best-fitting parabola, matching value, slope and
curvature at $a$. Each extra order typically buys one further correct
decimal place near $a$.
\end{rulebox}

\begin{rulebox}{Lagrange remainder and error bound}
If $f$ is $(n+1)$ times differentiable on an interval containing $a$ and
$x$, then
\[
  f(x)=T_{n}(x)+R_{n}(x),
  \qquad
  R_{n}(x)=\frac{f^{(n+1)}(\xi)}{(n+1)!}(x-a)^{n+1}
\]
for some $\xi$ strictly between $a$ and $x$. Since $\xi$ is unknown,
bound it: if $|f^{(n+1)}(t)|\leq M$ for all $t$ between $a$ and $x$,
then
\[
  \bigl|f(x)-T_{n}(x)\bigr|\leq\frac{M}{(n+1)!}\,|x-a|^{n+1}.
\]
The error shrinks like $|x-a|^{n+1}$: halving the distance to $a$
divides the error of $T_{n}$ by about $2^{n+1}$.
\end{rulebox}

\begin{rulebox}{Maclaurin case ($a=0$): standard expansions}
The Taylor polynomial at $a=0$ is called the \emph{Maclaurin}
polynomial. The classical expansions are
\[
\begin{aligned}
  \mathrm{e}^{x}&=1+x+\frac{x^{2}}{2!}+\frac{x^{3}}{3!}+\cdots
                  +\frac{x^{n}}{n!}+\cdots
                  &&\text{(all }x\text{)},\\
  \ln(1+x)&=x-\frac{x^{2}}{2}+\frac{x^{3}}{3}-\frac{x^{4}}{4}+\cdots
                  &&(-1<x\leq 1),\\
  (1+x)^{k}&=1+kx+\frac{k(k-1)}{2!}x^{2}
                  +\frac{k(k-1)(k-2)}{3!}x^{3}+\cdots
                  &&(|x|<1),\\
  \sin(x)&=x-\frac{x^{3}}{3!}+\frac{x^{5}}{5!}-\cdots
                  &&\text{(all }x\text{)},\\
  \cos(x)&=1-\frac{x^{2}}{2!}+\frac{x^{4}}{4!}-\cdots
                  &&\text{(all }x\text{)}.
\end{aligned}
\]
The first-order case $\ln(1+x)\approx x$ is the workhorse of
log-linearisation in economics: a small growth rate $g$ and the log
change $\ln(1+g)$ are interchangeable to first order.
\end{rulebox}

\begin{rulebox}{Quadratic approximation in economics}
Near a stationary point $a$ (where $f'(a)=0$) the linear term vanishes
and
\[
  f(x)\approx f(a)+\tfrac12 f''(a)(x-a)^{2}:
\]
locally, $f$ \emph{is} a parabola sitting on top of $f(a)$. If
$f''(a)<0$ the parabola opens downward and $a$ is a local maximum; if
$f''(a)>0$ it opens upward and $a$ is a local minimum. This is exactly
where the second-order condition comes from. In several variables the
role of $f''(a)$ is played by the Hessian, and the sign condition
becomes definiteness.
\end{rulebox}

\begin{rulebox}{Caution: validity is local}
A Taylor polynomial is a \emph{local} approximation: the error grows
like $|x-a|^{n+1}$, so accuracy is excellent near $a$ and can be
useless far away. For $\ln(1+x)$ and $(1+x)^{k}$ the expansions are
only valid on $(-1,1]$ resp.\ $(-1,1)$: outside, adding more terms
makes things \emph{worse}. Even inside: $T_{1}(x)=x$ for $\ln(1+x)$
gives $\ln(2)\approx 1$ against the true $0.693$, a $44\%$ error,
versus $0.5\%$ at $x=0.1$. Rule of thumb: expand around a base point
close to where you evaluate, and use the remainder bound to check.
\end{rulebox}

\section{Worked Example}

\begin{worked}{Expand $\ln(1+x)$ to third order about $0$, bound the error at $x=0.5$, and read the economics off the second-order term}
\wstep{Step 1. Collect the derivatives at the expansion point.}
\[
  f(x)=\ln(1+x),\quad
  f'=\frac{1}{1+x},\quad
  f''=\frac{-1}{(1+x)^{2}},\quad
  f'''=\frac{2}{(1+x)^{3}},\quad
  f^{(4)}=\frac{-6}{(1+x)^{4}} .
\]
At $a=0$: $f(0)=0$, $f'(0)=1$, $f''(0)=-1$, $f'''(0)=2$, and
$f^{(4)}(0)=-6$.
\wtrap{Evaluate \emph{after} differentiating, never before. Substituting
$x=0$ into $f'$ and then differentiating the resulting constant gives zeros
all the way down --- a surprisingly frequent way to produce the polynomial
$P_{3}(x)=x$.}

\wstep{Step 2. Assemble the Taylor polynomial.}
\[
  P_{n}(x)=\sum_{k=0}^{n}\frac{f^{(k)}(a)}{k!}(x-a)^{k}
  \quad\Longrightarrow\quad
  P_{3}(x)=0+x-\frac{1}{2!}x^{2}+\frac{2}{3!}x^{3}
  =\boxed{x-\frac{x^{2}}{2}+\frac{x^{3}}{3}}
\]
The factorials are what turn the raw derivatives $1,-1,2$ into the
coefficients $1,-\tfrac12,\tfrac13$.
\wtrap{Do not omit $k!$, and do not omit $(x-a)^{k}$ when $a\neq0$. Expanding
about $a=2$ and writing powers of $x$ rather than of $x-2$ is the standard
error, and it silently produces a polynomial that is wrong everywhere.}

\wstep{Step 3. Evaluate at $x=0.5$ and compare.}
\[
  P_{3}(0.5)=0.5-0.125+0.041\overline{6}=0.4166\overline{6},
  \qquad
  \ln(1.5)=0.405465 ,
\]
so the actual error is $0.011202$.

\wstep{Step 4. Bound the error \emph{before} knowing the true value.}
The Lagrange remainder is
\[
  R_{3}(x)=\frac{f^{(4)}(\xi)}{4!}x^{4}
  \qquad\text{for some }\xi\text{ between }0\text{ and }x .
\]
$\xi$ is unknown, so take the worst case. On $[0,0.5]$,
\[
  \bigl|f^{(4)}(\xi)\bigr|=\frac{6}{(1+\xi)^{4}}\le\frac{6}{1^{4}}=6
  \qquad(\text{largest at }\xi=0),
\]
\[
  |R_{3}(0.5)|\le\frac{6}{24}(0.5)^{4}=0.25\cdot0.0625=0.015625 .
\]
And indeed $0.011201\le0.015625$. \checkmark
\wtrap{Maximise $|f^{(n+1)}|$ over the whole interval between $a$ and $x$ ---
do not substitute $\xi=x$ or $\xi=a$ and call it a bound. Here $1/(1+\xi)^{4}$
is decreasing, so the worst case is at the \emph{left} end; for an
increasing derivative it would be at the right. Check which end before
committing.}

\wstep{Step 5. Note how fast the accuracy degrades with distance.}
At $x=0.1$ the same bound gives $\tfrac14(0.1)^{4}=0.000025$; at $x=0.5$ it
is $625$ times larger. The error carries the factor $x^{4}$, so halving the
distance to the expansion point cuts the error sixteenfold. At $x=1$ the
series still converges, but at $x=1.1$ it diverges outright --- the
expansion of $\ln(1+x)$ is valid only for $|x|<1$.
\wtrap{A Taylor polynomial is a \textbf{local} statement. It is excellent
near $a$, useless far from it, and there is no warning in the formula --- 
$P_{3}(5)=5-12.5+41.7=34.2$, while $\ln6=1.79$. Always ask how far from $a$
you are being asked to work.}

\wstep{Step 6. The first two terms are the economics.}
Keeping $n=1$ gives the \textbf{linear approximation} $\ln(1+x)\approx x$:
the familiar ``a log difference is a percentage change'', accurate to about
half a percentage point for $|x|\le0.1$. Keeping $n=2$ gives
\[
  f(x)\approx f(a)+f'(a)(x-a)+\tfrac12 f''(a)(x-a)^{2},
\]
whose quadratic term carries the \textbf{curvature}. In expected-utility
theory that term is exactly what produces the risk premium: a concave $u$
has $u''<0$, so the quadratic term is negative, and its size relative to
$u'$ is the Arrow--Pratt coefficient $-u''/u'$. Risk aversion \emph{is} the
second-order term of a Taylor expansion.

\wstep{Step 7. The expansions worth knowing by heart.}
\[
  \mathrm{e}^{x}=1+x+\tfrac{x^{2}}{2}+\tfrac{x^{3}}{6}+\cdots,
  \qquad
  \ln(1+x)=x-\tfrac{x^{2}}{2}+\tfrac{x^{3}}{3}-\cdots,
\]
\[
  (1+x)^{\alpha}=1+\alpha x+\tfrac{\alpha(\alpha-1)}{2}x^{2}+\cdots,
  \qquad
  \frac{1}{1-x}=1+x+x^{2}+\cdots
\]
The last is the geometric series of Volume~4, which is why a perpetuity
formula and a Taylor expansion are the same computation twice.
\end{worked}

\section{Practice Problems}

\begin{problems}
  \item Find $T_{1}$ of $f(x)=\ln(x)$ at $a=1$.
  \item Find $T_{1}$ of $f(x)=\sqrt{x}$ at $a=4$.
  \item Find $T_{1}$ of $f(x)=\dfrac{1}{x}$ at $a=1$.
  \item Find $T_{1}$ of $f(x)=x^{3}-2x$ at $a=2$.
  \item Find $T_{1}$ of $f(x)=\mathrm{e}^{2x}$ at $a=\ln 2$.
  \item Find $T_{2}$ of $f(x)=\mathrm{e}^{x}$ at $a=0$ and use it to
        approximate $\mathrm{e}^{0.1}$. Compare with the true value
        $\mathrm{e}^{0.1}=1.105171$.
  \item Find $T_{2}$ of $f(x)=\ln(x)$ at $a=1$ and use it to
        approximate $\ln(1.05)$. Compare with the true value
        $\ln(1.05)=0.048790$.
  \item Find $T_{2}$ of $f(x)=\sqrt{x}$ at $a=4$ and use it to
        approximate $\sqrt{4.1}$. Compare with the true value
        $\sqrt{4.1}=2.024846$.
  \item Find $T_{2}$ of $f(x)=\dfrac{1}{1-x}$ at $a=0$.
  \item Find $T_{2}$ of $f(x)=x\ln(x)$ at $a=1$.
  \item Find $T_{2}$ of $f(x)=\mathrm{e}^{-x^{2}}$ at $a=0$.
  \item Find $T_{3}$ of $f(x)=\ln(1+x)$ at $a=0$ and use it to
        approximate $\ln(0.9)$. Compare with the true value
        $\ln(0.9)=-0.105361$.
  \item Use the binomial expansion to find $T_{2}$ of
        $f(x)=(1+x)^{1/2}$ at $a=0$, and use it to approximate
        $\sqrt{1.1}$. Compare with the true value
        $\sqrt{1.1}=1.048809$.
  \item Find the Maclaurin polynomial $T_{4}$ of $f(x)=\cos(x)$ and use
        it to approximate $\cos(0.2)$. Compare with the true value
        $\cos(0.2)=0.980067$.
  \item A firm produces $Q(L)=10\sqrt{L}$ units with $L$ hours of
        labour. Find $T_{2}$ of $Q$ at $a=100$ and use it to
        approximate output at $L=104$. Compare with the true value
        $Q(104)=101.980390$.
  \item (Log-linearisation.) Output grows from $Y_{0}$ to
        $Y_{1}=Y_{0}(1+g)$. Show, using $T_{1}$ of $\ln(1+x)$ at $0$,
        that $\ln Y_{1}-\ln Y_{0}\approx g$, and quantify the
        approximation for $g=0.03$ given $\ln(1.03)=0.029559$.
  \item Use the Lagrange remainder to bound the error of the $T_{1}$
        approximation of $\sqrt{4.1}$ at $a=4$ (Problem 2), and compare
        the bound with the actual error.
  \item Use the Lagrange remainder to bound the error of the $T_{2}$
        approximation of $\mathrm{e}^{0.1}$ at $a=0$ (Problem 6), and
        compare the bound with the actual error.
  \item Find the smallest $n$ such that the Maclaurin polynomial
        $T_{n}$ of $\mathrm{e}^{x}$ approximates $\mathrm{e}^{0.5}$
        with error at most $10^{-4}$. (Use $\mathrm{e}^{0.5}<2$ in the
        remainder bound.)
  \item Prove, using $T_{1}$ with Lagrange remainder: if $f'(a)=0$ and
        $f''$ is continuous with $f''(a)<0$, then $a$ is a strict local
        maximum of $f$.
\end{problems}

\section{Solutions}

\begin{solutions}
  \item $f(1)=\ln(1)=0$; $f'(x)=1/x$, so $f'(1)=1$. Hence
        \[
          T_{1}(x)=0+1\cdot(x-1)=x-1 .
        \]

  \item $f(4)=2$; $f'(x)=\dfrac{1}{2\sqrt{x}}$, so
        $f'(4)=\dfrac{1}{2\cdot 2}=\dfrac14$. Hence
        \[
          T_{1}(x)=2+\tfrac14(x-4) .
        \]

  \item $f(1)=1$; $f'(x)=-1/x^{2}$, so $f'(1)=-1$. Hence
        \[
          T_{1}(x)=1-(x-1)=2-x .
        \]

  \item $f(2)=8-4=4$; $f'(x)=3x^{2}-2$, so $f'(2)=12-2=10$. Hence
        \[
          T_{1}(x)=4+10(x-2) .
        \]

  \item $f(\ln 2)=\mathrm{e}^{2\ln 2}=(\mathrm{e}^{\ln 2})^{2}=4$;
        $f'(x)=2\mathrm{e}^{2x}$, so $f'(\ln 2)=2\cdot 4=8$. Hence
        \[
          T_{1}(x)=4+8(x-\ln 2) .
        \]

  \item $f^{(k)}(x)=\mathrm{e}^{x}$ for every $k$, so
        $f(0)=f'(0)=f''(0)=1$ and
        \[
          T_{2}(x)=1+x+\frac{x^{2}}{2} .
        \]
        At $x=0.1$: $T_{2}(0.1)=1+0.1+\tfrac{0.01}{2}=1.105$. True
        value $1.105171$; the error is about $1.7\cdot 10^{-4}$.

  \item $f(1)=0$; $f'(x)=1/x$, $f'(1)=1$; $f''(x)=-1/x^{2}$,
        $f''(1)=-1$. Hence
        \[
          T_{2}(x)=(x-1)-\tfrac12(x-1)^{2} .
        \]
        At $x=1.05$:
        $T_{2}(1.05)=0.05-\tfrac12(0.0025)=0.05-0.00125=0.04875$. True
        value $0.048790$; the error is about $4.0\cdot 10^{-5}$.

  \item From Problem 2, $f(4)=2$ and $f'(4)=\tfrac14$. Further
        $f''(x)=-\tfrac14 x^{-3/2}$, so
        $f''(4)=-\tfrac14\cdot\tfrac18=-\tfrac{1}{32}$ and the
        quadratic coefficient is
        $\tfrac12 f''(4)=-\tfrac{1}{64}$:
        \[
          T_{2}(x)=2+\tfrac14(x-4)-\tfrac{1}{64}(x-4)^{2} .
        \]
        At $x=4.1$:
        \[
          T_{2}(4.1)=2+0.025-\tfrac{0.01}{64}
                    =2.025-0.00015625=2.02484375 .
        \]
        True value $2.024846$; the error is about $1.9\cdot 10^{-6}$
        (the linear estimate $2.025$ alone is off by
        $1.5\cdot 10^{-4}$).

  \item $f(0)=1$; $f'(x)=\dfrac{1}{(1-x)^{2}}$, $f'(0)=1$;
        $f''(x)=\dfrac{2}{(1-x)^{3}}$, $f''(0)=2$, so
        $\tfrac12 f''(0)=1$. Hence
        \[
          T_{2}(x)=1+x+x^{2} ,
        \]
        the start of the geometric series $1+x+x^{2}+\cdots$.

  \item $f(1)=1\cdot\ln(1)=0$; $f'(x)=\ln(x)+1$, $f'(1)=1$;
        $f''(x)=1/x$, $f''(1)=1$. Hence
        \[
          T_{2}(x)=(x-1)+\tfrac12(x-1)^{2} .
        \]

  \item $f(0)=1$; $f'(x)=-2x\mathrm{e}^{-x^{2}}$, $f'(0)=0$;
        $f''(x)=(-2+4x^{2})\mathrm{e}^{-x^{2}}$, $f''(0)=-2$. Hence
        \[
          T_{2}(x)=1-x^{2} .
        \]
        (Equivalently: substitute $u=-x^{2}$ into
        $\mathrm{e}^{u}\approx 1+u$.)

  \item $f(0)=0$; $f'(x)=\dfrac{1}{1+x}$, $f'(0)=1$;
        $f''(x)=-\dfrac{1}{(1+x)^{2}}$, $f''(0)=-1$;
        $f'''(x)=\dfrac{2}{(1+x)^{3}}$, $f'''(0)=2$. The coefficients
        are $1$, $-\tfrac{1}{2!}=-\tfrac12$ and
        $\tfrac{2}{3!}=\tfrac13$:
        \[
          T_{3}(x)=x-\frac{x^{2}}{2}+\frac{x^{3}}{3} .
        \]
        At $x=-0.1$:
        \[
          T_{3}(-0.1)=-0.1-\frac{0.01}{2}+\frac{-0.001}{3}
                     =-0.1-0.005-0.000333=-0.105333 .
        \]
        True value $-0.105361$; the error is about $2.7\cdot 10^{-5}$.

  \item With $k=\tfrac12$ in the binomial expansion: the linear
        coefficient is $k=\tfrac12$ and the quadratic coefficient is
        \[
          \frac{k(k-1)}{2!}
            =\frac{\tfrac12\cdot\bigl(-\tfrac12\bigr)}{2}
            =-\frac18 ,
        \]
        so
        \[
          T_{2}(x)=1+\frac{x}{2}-\frac{x^{2}}{8} .
        \]
        At $x=0.1$: $T_{2}(0.1)=1+0.05-0.00125=1.04875$. True value
        $1.048809$; the error is about $5.9\cdot 10^{-5}$.

  \item Derivatives of $\cos$ at $0$ cycle through
        $1,\,0,\,-1,\,0,\,1$ for $k=0,\ldots,4$, so
        \[
          T_{4}(x)=1-\frac{x^{2}}{2!}+\frac{x^{4}}{4!}
                  =1-\frac{x^{2}}{2}+\frac{x^{4}}{24} .
        \]
        At $x=0.2$: $(0.2)^{2}=0.04$ and $(0.2)^{4}=0.0016$, so
        \[
          T_{4}(0.2)=1-0.02+\frac{0.0016}{24}
                    =1-0.02+0.000067=0.980067 ,
        \]
        agreeing with the true value $0.980067$ to all six decimals
        shown (the next term is of size $\tfrac{(0.2)^{6}}{720}\approx
        9\cdot 10^{-8}$).

  \item $Q(100)=10\cdot 10=100$; $Q'(L)=5L^{-1/2}$,
        $Q'(100)=\tfrac{5}{10}=0.5$; $Q''(L)=-\tfrac52 L^{-3/2}$,
        $Q''(100)=-\tfrac{2.5}{1000}=-0.0025$, so
        $\tfrac12 Q''(100)=-0.00125$. Hence
        \[
          T_{2}(L)=100+0.5\,(L-100)-0.00125\,(L-100)^{2} .
        \]
        At $L=104$ (so $L-100=4$):
        \[
          T_{2}(104)=100+0.5\cdot 4-0.00125\cdot 16=102-0.02=101.98 .
        \]
        True value $101.980390$; the error is about
        $3.9\cdot 10^{-4}$, four extra labour hours raise output by
        just under $2$ units, the shortfall against the linear estimate
        $102$ reflecting diminishing returns ($Q''<0$).

  \item Taking logs, $\ln Y_{1}=\ln Y_{0}+\ln(1+g)$. For
        $f(x)=\ln(1+x)$ at $a=0$: $f(0)=0$ and $f'(0)=1$, so
        $T_{1}(x)=x$ and
        \[
          \ln Y_{1}-\ln Y_{0}=\ln(1+g)\approx g .
        \]
        For $g=0.03$ the log change is $\ln(1.03)=0.029559$, while the
        growth rate is $0.03$: the approximation is off by about
        $4.4\cdot 10^{-4}$, i.e.\ by roughly $1.5\%$ of the quantity
        itself. (The second-order term $-\tfrac12 g^{2}=-0.00045$
        accounts for almost all of the gap.)

  \item Here $n=1$, $a=4$, $x=4.1$ and $f''(t)=-\tfrac14 t^{-3/2}$.
        For $t\in[4,4.1]$, $t^{-3/2}$ is decreasing, so
        \[
          |f''(t)|\leq\tfrac14\cdot 4^{-3/2}=\tfrac14\cdot\tfrac18
                 =\tfrac{1}{32}=:M .
        \]
        The remainder bound gives
        \[
          |R_{1}(4.1)|\leq\frac{M}{2!}\,(0.1)^{2}
            =\frac{1}{32}\cdot\frac{0.01}{2}
            =\frac{0.01}{64}\approx 1.56\cdot 10^{-4} .
        \]
        Actual error: $T_{1}(4.1)=2.025$ versus $\sqrt{4.1}=2.024846$,
        i.e.\ $1.54\cdot 10^{-4}$, just inside the bound.

  \item Here $n=2$, $a=0$, $x=0.1$ and
        $f'''(t)=\mathrm{e}^{t}$. For $t\in(0,0.1)$,
        $\mathrm{e}^{t}<\mathrm{e}^{0.1}<1.2=:M$, so
        \[
          |R_{2}(0.1)|\leq\frac{M}{3!}\,(0.1)^{3}
            =\frac{1.2\cdot 0.001}{6}=2\cdot 10^{-4} .
        \]
        Actual error: $T_{2}(0.1)=1.105$ versus
        $\mathrm{e}^{0.1}=1.105171$, i.e.\ $1.7\cdot 10^{-4}$,
        consistent with the bound.

  \item With $f(x)=\mathrm{e}^{x}$, $a=0$, $x=0.5$:
        $f^{(n+1)}(\xi)=\mathrm{e}^{\xi}\leq\mathrm{e}^{0.5}<2$ for
        $\xi\in(0,0.5)$, so
        \[
          |R_{n}(0.5)|\leq\frac{2\,(0.5)^{n+1}}{(n+1)!} .
        \]
        Tabulate:
        \[
          n=3:\ \frac{2\cdot 0.0625}{24}\approx 5.2\cdot 10^{-3},
          \qquad
          n=4:\ \frac{2\cdot 0.03125}{120}\approx 5.2\cdot 10^{-4},
        \]
        \[
          n=5:\ \frac{2\cdot 0.015625}{720}\approx 4.3\cdot 10^{-5} .
        \]
        The first value below $10^{-4}$ occurs at $n=5$. Check with the
        true value $\mathrm{e}^{0.5}=1.648721$: $T_{4}(0.5)=1.648438$
        (error $2.8\cdot 10^{-4}$, indeed too large) while
        $T_{5}(0.5)=1.648698$ (error $2.3\cdot 10^{-5}$). So $n=5$.

  \item By Taylor's theorem with $n=1$: for each $x$ near $a$ there is
        $\xi$ between $a$ and $x$ with
        \[
          f(x)=f(a)+f'(a)(x-a)+\tfrac12 f''(\xi)(x-a)^{2}
              =f(a)+\tfrac12 f''(\xi)(x-a)^{2},
        \]
        using $f'(a)=0$. Since $f''$ is continuous and $f''(a)<0$,
        there is a neighbourhood $U$ of $a$ on which $f''<0$. For
        $x\in U$, $x\neq a$, the point $\xi$ also lies in $U$, so
        $f''(\xi)<0$ while $(x-a)^{2}>0$, giving
        \[
          f(x)-f(a)=\tfrac12 f''(\xi)(x-a)^{2}<0 .
        \]
        Hence $f(x)<f(a)$ for all $x\in U\setminus\{a\}$: $a$ is a
        strict local maximum. (This is the second-order condition in
        its cleanest form: near a stationary point the function
        \emph{is} its osculating parabola, up to higher-order terms.)
\end{solutions}

\section{Additional Problems}

\begin{problems}
  \item Find $T_{1}$ of $f(x)=(1+x)^{5}$ at $a=0$.
  \item Find $T_{1}$ of $f(x)=\sin(x)$ at $a=0$ and use it to
        approximate $\sin(0.1)$. Compare with the true value
        $\sin(0.1)=0.099833$.
  \item Find $T_{2}$ of $f(x)=x^{4}$ at $a=1$.
  \item Find $T_{2}$ of $f(x)=\dfrac{1}{x}$ at $a=2$.
  \item Find $T_{3}$ of $f(x)=\mathrm{e}^{x}$ at $a=0$ and use it to
        approximate $\mathrm{e}^{0.2}$. Compare with the true value
        $\mathrm{e}^{0.2}=1.221403$.
  \item Find $T_{2}$ of $f(x)=\ln(1+2x)$ at $a=0$.
  \item (Cobb--Douglas.) For $f(K)=K^{0.3}$, find $T_{2}$ at $a=1$ and
        use it to approximate $1.1^{0.3}$. Compare with the true value
        $1.1^{0.3}=1.029006$.
  \item Use the Lagrange remainder to bound the error of approximating
        $\ln(1.05)$ by $T_{1}(x)=x-1$ of $\ln(x)$ at $a=1$, and compare
        the bound with the actual error (true value
        $\ln(1.05)=0.048790$).
  \item Find the smallest $n$ such that the Maclaurin polynomial
        $T_{n}$ of $\ln(1+x)$ approximates $\ln(1.1)$ with error at
        most $10^{-6}$ (true value $\ln(1.1)=0.095310$).
  \item Suppose $|f'''(t)|\leq M$ for all $t$ between $a$ and $x$.
        Derive a bound for $|f(x)-T_{2}(x)|$ in terms of $M$ and
        $|x-a|$.
\end{problems}

\section{Additional Solutions}

\begin{solutions}
  \item $f(0)=1$; $f'(x)=5(1+x)^{4}$, $f'(0)=5$. Hence
        \[
          T_{1}(x)=1+5x .
        \]

  \item $f(0)=0$; $f'(x)=\cos(x)$, $f'(0)=1$. Hence $T_{1}(x)=x$ (the
        small-angle approximation) and $\sin(0.1)\approx 0.1$. True
        value $0.099833$; the error is about $1.7\cdot 10^{-4}$
        (matching the size of the first omitted term,
        $\tfrac{x^{3}}{3!}=\tfrac{0.001}{6}$).

  \item $f(1)=1$; $f'(x)=4x^{3}$, $f'(1)=4$; $f''(x)=12x^{2}$,
        $f''(1)=12$, so $\tfrac12 f''(1)=6$. Hence
        \[
          T_{2}(x)=1+4(x-1)+6(x-1)^{2} .
        \]

  \item $f(2)=\tfrac12$; $f'(x)=-1/x^{2}$, $f'(2)=-\tfrac14$;
        $f''(x)=2/x^{3}$, $f''(2)=\tfrac28=\tfrac14$, so
        $\tfrac12 f''(2)=\tfrac18$. Hence
        \[
          T_{2}(x)=\frac12-\frac14(x-2)+\frac18(x-2)^{2} .
        \]

  \item All derivatives equal $1$ at $0$, so
        \[
          T_{3}(x)=1+x+\frac{x^{2}}{2}+\frac{x^{3}}{6} .
        \]
        At $x=0.2$:
        \[
          T_{3}(0.2)=1+0.2+\frac{0.04}{2}+\frac{0.008}{6}
                    =1+0.2+0.02+0.001333=1.221333 .
        \]
        True value $1.221403$; the error is about $6.9\cdot 10^{-5}$.

  \item $f(0)=0$; $f'(x)=\dfrac{2}{1+2x}$, $f'(0)=2$;
        $f''(x)=-\dfrac{4}{(1+2x)^{2}}$, $f''(0)=-4$, so
        $\tfrac12 f''(0)=-2$. Hence
        \[
          T_{2}(x)=2x-2x^{2} .
        \]
        (Equivalently: substitute $u=2x$ into
        $\ln(1+u)\approx u-\tfrac12 u^{2}$.)

  \item $f(1)=1$; $f'(K)=0.3\,K^{-0.7}$, $f'(1)=0.3$;
        $f''(K)=0.3\cdot(-0.7)\,K^{-1.7}=-0.21\,K^{-1.7}$,
        $f''(1)=-0.21$, so $\tfrac12 f''(1)=-0.105$. Hence
        \[
          T_{2}(K)=1+0.3\,(K-1)-0.105\,(K-1)^{2} .
        \]
        At $K=1.1$:
        \[
          T_{2}(1.1)=1+0.03-0.105\cdot 0.01=1.03-0.00105=1.02895 .
        \]
        True value $1.029006$; the error is about $5.6\cdot 10^{-5}$.

  \item Here $n=1$, $a=1$, $x=1.05$ and $f''(t)=-1/t^{2}$. On
        $[1,1.05]$, $|f''(t)|=1/t^{2}\leq 1=:M$. The remainder bound
        gives
        \[
          |R_{1}(1.05)|\leq\frac{M}{2!}\,(0.05)^{2}
            =\frac{0.0025}{2}=1.25\cdot 10^{-3} .
        \]
        Actual error: $T_{1}(1.05)=0.05$ versus $0.048790$, i.e.\
        $1.21\cdot 10^{-3}$, close to the bound, because
        $1/t^{2}\approx 1$ throughout the interval.

  \item For $f(x)=\ln(1+x)$,
        $f^{(n+1)}(t)=\dfrac{(-1)^{n}\,n!}{(1+t)^{n+1}}$, so on
        $[0,0.1]$ we have $|f^{(n+1)}(t)|\leq n!$ and
        \[
          |R_{n}(0.1)|\leq\frac{n!}{(n+1)!}\,(0.1)^{n+1}
            =\frac{(0.1)^{n+1}}{n+1} .
        \]
        Tabulate:
        \[
          n=4:\ \frac{(0.1)^{5}}{5}=2\cdot 10^{-6},
          \qquad
          n=5:\ \frac{(0.1)^{6}}{6}\approx 1.7\cdot 10^{-7} .
        \]
        The first value below $10^{-6}$ occurs at $n=5$. Check against
        the true value $0.095310$: $T_{4}(0.1)=0.1-0.005+0.000333
        -0.000025=0.095308$ (error $1.9\cdot 10^{-6}$, indeed too
        large), while $T_{5}(0.1)=0.095308+0.000002=0.095310$ (error
        about $1.5\cdot 10^{-7}$). So $n=5$.

  \item Apply the Lagrange form with $n=2$: there is $\xi$ between $a$
        and $x$ with
        \[
          f(x)-T_{2}(x)=\frac{f'''(\xi)}{3!}(x-a)^{3},
        \]
        hence
        \[
          |f(x)-T_{2}(x)|\leq\frac{|f'''(\xi)|}{6}\,|x-a|^{3}
            \leq\frac{M}{6}\,|x-a|^{3} .
        \]
\end{solutions}

\chapter{Implicit Differentiation}

Most curves in the plane are not graphs. The circle $x^{2}+y^{2}=25$ fails
the vertical line test, the equation $\mathrm{e}^{y}+xy=1$ cannot be solved
for $y$ in closed form at all, and an indifference curve is handed to us as
a level set rather than as a formula. Yet each of them has a tangent at
almost every point, and economics asks for its slope constantly: that slope
\emph{is} the marginal rate of substitution.

Implicit differentiation gets it without solving for $y$. The move rests on
one idea, and the whole chapter is that idea applied repeatedly: pretend the
equation has already been solved near the point of interest, so that $y$ is
a function $y(x)$; substituting it makes the equation an identity in $x$;
identities may be differentiated. Because every $y$ in sight is secretly the
composite $y(x)$, the chain rule leaves a factor $y'$ behind wherever a $y$
was, and the resulting linear equation is solved for $y'$.

Two questions have to be kept apart, and the chapter answers both. \emph{Is}
there such a function $y(x)$ near the point --- and if so, what is its
slope? The first question is the Implicit Function Theorem, and the honest
answer is sometimes no: at a point where the curve crosses itself or stands
vertically, no function $y(x)$ describes it, and the slope formula returns a
division by zero to say so.

\section{Core Rules}

\begin{rulebox}{Why a factor $y'$ appears}
Suppose that near a point of the curve the equation really does pin down one
$y$ for each $x$, and call that function $y(x)$. Substituting it turns the
equation into an \textbf{identity in $x$}, true for every $x$ near $x_{0}$:
\[
  x^{2}+\bigl(y(x)\bigr)^{2}=25 .
\]
An identity may be differentiated on both sides. Every $y$ in the expression
is the composite $y(x)$, so it is differentiated by the \textbf{chain rule}
and leaves a factor $y'$ behind:
\[
  \frac{\mathrm{d}}{\mathrm{d}x}\bigl[y^{2}\bigr]
  =\frac{\mathrm{d}}{\mathrm{d}x}\bigl[(y(x))^{2}\bigr]
  =2\,y(x)\,y'(x)=2yy' .
\]
Hence $2x+2yy'=0$ and $y'=-x/y$. Nothing was solved for $y$, and the answer
is \emph{allowed} to contain $y$: it is the slope at a point of the curve,
and the two points with the same $x$ --- here $(3,4)$ and $(3,-4)$ --- have
slopes $-3/4$ and $+3/4$.
\end{rulebox}

\begin{rulebox}{Setup}
Suppose $F(x,y)=0$ defines $y=y(x)$ implicitly. Differentiate every term
with respect to $x$. Whenever an expression involves $y$, the chain rule
produces an extra factor $\dd{y}{x}$:
\[
  \frac{\mathrm{d}}{\mathrm{d}x}\bigl[g(y)\bigr]=g'(y)\cdot\dd{y}{x} .
\]
So $\frac{\mathrm{d}}{\mathrm{d}x}[\sin y]=\cos(y)\,y'$,
$\frac{\mathrm{d}}{\mathrm{d}x}[\mathrm{e}^{y}]=\mathrm{e}^{y}y'$, and
$\frac{\mathrm{d}}{\mathrm{d}x}[\ln y]=y'/y$.
\end{rulebox}

\begin{rulebox}{Mixed terms}
For a product such as $x^{3}y^{2}$ apply the product rule, treating $y$ as
a function of $x$:
\[
  \frac{\mathrm{d}}{\mathrm{d}x}\bigl[x^{3}y^{2}\bigr]=3x^{2}y^{2}+x^{3}\cdot 2yy' .
\]
\end{rulebox}

\begin{rulebox}{Solving for $y'$}
Collect all terms containing $y'$ on one side and factor:
\[
  A(x,y)\,y'+B(x,y)=0\;\Longrightarrow\;y'=-\frac{B(x,y)}{A(x,y)} .
\]
The equation is always \emph{linear} in $y'$, however complicated it looks,
because $y'$ only ever enters through one chain-rule step.
\end{rulebox}

\begin{rulebox}{The curve is a level set --- and that gives the formula at once}
Write the equation as $F(x,y)=0$. Then the curve is the \textbf{level set of
the surface $z=F(x,y)$ at height $0$}: the shoreline where the surface cuts
the horizontal plane $z=0$. Walking along the curve, $F$ never changes.
Writing that down with the chain rule for $F\bigl(x,y(x)\bigr)=0$,
\[
  \pd{F}{x}+\pd{F}{y}\cdot\dd{y}{x}=0
  \qquad\Longrightarrow\qquad
  \boxed{\;\dd{y}{x}=-\frac{F_{x}(x,y)}{F_{y}(x,y)}\;}
  \qquad\text{provided } F_{y}(x,y)\neq0 .
\]
This is not a new method, it is the same computation done once and for all:
for $F=x^{2}+y^{2}-25$ it returns $-2x/2y=-x/y$, as above. In economics the
curve is an indifference curve $U(x,y)=\bar u$ and the formula reads
$\dd{y}{x}=-U_{x}/U_{y}$: the marginal rate of substitution is the ratio of
marginal utilities.
\end{rulebox}

\begin{gallery}
\begin{tikzpicture}
  \begin{axis}[surfaxis, title={$z=F(x,y)=x^{3}-3x^{2}+y^{2}-2y+1$}]
    \addplot3[surf, domain=-1.2:3.2, y domain=-1.6:3.6]{x^3-3*x^2+y^2-2*y+1};
  \end{axis}
\end{tikzpicture}\hfill
\begin{tikzpicture}
  \begin{axis}[geoaxis, xmin=-1.6, xmax=3.6, ymin=-1.9, ymax=3.9,
      xtick={-1,1,2,3}, ytick={-1,1,2,3}, title={the level set $F=0$}]
    \addplot[domain=-1.2:3, samples=241]{1+abs(x)*sqrt(3-x)};
    \addplot[domain=-1.2:3, samples=241]{1-abs(x)*sqrt(3-x)};
    \addplot[only marks, mark=*, mark size=1.3pt, accent, forget plot]
      coordinates {(0,1) (3,1)};
  \end{axis}
\end{tikzpicture}
\gcap{The surface and its zero level set $k$. At $(0,1)$ the curve crosses
itself; at $(3,1)$ it stands vertical. Everywhere else a local $y(x)$
exists.}
\end{gallery}

\begin{rulebox}{When there is no $y(x)$: the Implicit Function Theorem}
\textbf{Theorem.} If $F$ is $C^{1}$ near $(x_{0},y_{0})$, if
$F(x_{0},y_{0})=0$, and if
\[
  F_{y}(x_{0},y_{0})\neq0 ,
\]
then there are neighbourhoods on which the equation $F(x,y)=0$ defines
exactly one differentiable function $y=y(x)$ with $y(x_{0})=y_{0}$, and its
derivative is $-F_{x}/F_{y}$.

The condition is exactly what the formula needs, and its failure is
visible in the picture. Where $F_{y}=0$:
\begin{itemize}[leftmargin=1.6em,itemsep=0.3ex]
  \item if $F_{x}\neq0$, the curve has a \textbf{vertical tangent}. No
        function $y(x)$ can describe it there --- but the roles may be
        swapped: $x=x(y)$ exists, with
        $\dd{x}{y}=-F_{y}/F_{x}$, and that derivative is $0$.
  \item if $F_{x}=0$ as well, the point is \textbf{singular}: typically the
        curve crosses itself, and \emph{two} branches with two different
        slopes pass through. No single function describes either the curve
        or its slope, and $-F_{x}/F_{y}$ degenerates to $0/0$. Separate the
        branches first, then differentiate each one.
\end{itemize}
Reading it off a picture: look for vertical tangents and self-intersections.
Those are the only places the theorem can fail. A failed condition is a
warning, not a proof that no $y(x)$ exists: $y^{3}-x=0$ has $F_{y}=0$ at the
origin, yet $y=x^{1/3}$ solves it for every $x$ (with a vertical tangent, so
not differentiably). Chapter~14 collects such cases.
\end{rulebox}

\begin{rulebox}{Logarithmic differentiation}
If $y=f(x)>0$ is a complicated product, quotient or power, take $\ln$
first:
\[
  \ln y=\ln f(x)
  \;\Longrightarrow\;
  \frac{y'}{y}=\frac{\mathrm{d}}{\mathrm{d}x}\ln f(x)
  \;\Longrightarrow\;
  y'=y\cdot\frac{\mathrm{d}}{\mathrm{d}x}\ln f(x) .
\]
Particularly useful for $y=u(x)^{v(x)}$: $\ln y=v(x)\ln u(x)$. This is
implicit differentiation applied to $\ln y-\ln f(x)=0$, which is why the
factor $y'/y$ appears.
\end{rulebox}

\begin{rulebox}{Higher-order implicit}
To find $y''$, differentiate $y'$ a second time with respect to $x$, again
remembering that any $y$ that appears contributes a factor $y'$. Then
substitute the expression already found for $y'$ to eliminate it.
\end{rulebox}

\begin{rulebox}{Tangent line via implicit slope}
The tangent at $(x_{0},y_{0})$ on $F(x,y)=0$ is
$y-y_{0}=y'(x_{0})\cdot(x-x_{0})$, where $y'(x_{0})$ is computed
implicitly. Both coordinates are substituted into $-F_{x}/F_{y}$ before the
line is written down.
\end{rulebox}

\section{Worked Example}

\begin{worked}{The folium $x^{3}+y^{3}=6xy$: find $y'$, the tangent at $(3,3)$, and every point where the tangent is vertical}
\wstep{Step 1. Confirm the point is on the curve.}
$27+27=54$ and $6\cdot3\cdot3=54$. \checkmark\ Do this first: an implicit
slope formula evaluated at a point that is not on the curve returns a
perfectly plausible number that means nothing.

\wstep{Step 2. Differentiate every term with respect to $x$, treating $y$ as $y(x)$.}
Left side: $\frac{\mathrm{d}}{\mathrm{d}x}\bigl[x^{3}+y^{3}\bigr]
=3x^{2}+3y^{2}y'$ --- the chain rule leaves a factor $y'$ behind wherever a
$y$ was.

Right side: $6xy$ is a \textbf{product} of two functions of $x$, so the
product rule applies:
\[
  \frac{\mathrm{d}}{\mathrm{d}x}\bigl[6xy\bigr]=6y+6xy' .
\]
\wtrap{The right side is where this problem is lost. Writing $6y$ alone
(treating $y$ as a constant) or $6xy'$ alone (forgetting the product rule)
are both common. Every occurrence of $y$ needs the chain rule, \emph{and}
every product needs the product rule --- they compound rather than replace
each other.}

\wstep{Step 3. Collect the $y'$ terms; the equation is always linear in $y'$.}
\[
  3x^{2}+3y^{2}y'=6y+6xy'
  \;\Longrightarrow\;
  3y^{2}y'-6xy'=6y-3x^{2}
  \;\Longrightarrow\;
  y'\bigl(3y^{2}-6x\bigr)=6y-3x^{2} .
\]
\[
  \boxed{\;y'=\frac{6y-3x^{2}}{3y^{2}-6x}=\frac{2y-x^{2}}{y^{2}-2x}\;}
\]
The answer contains $y$, and it must: two points of the curve can share an
$x$ and have different slopes, so $x$ alone cannot determine $y'$.

\wstep{Step 4. Evaluate at the point.}
\[
  y'(3,3)=\frac{2\cdot3-9}{9-6}=\frac{-3}{3}=-1 ,
\]
so the tangent is
\[
  y-3=-1(x-3)\quad\Longrightarrow\quad y=6-x .
\]
By the symmetry of the equation in $x$ and $y$, the curve is symmetric about
the line $y=x$, and $(3,3)$ lies on that line --- so a slope of $-1$,
perpendicular to it, is exactly what the geometry predicts. \checkmark

\wstep{Step 5. Cross-check with the level-set formula.}
Writing $F(x,y)=x^{3}+y^{3}-6xy$,
\[
  F_{x}=3x^{2}-6y,\qquad F_{y}=3y^{2}-6x,
  \qquad
  y'=-\frac{F_{x}}{F_{y}}=-\frac{3x^{2}-6y}{3y^{2}-6x}=\frac{6y-3x^{2}}{3y^{2}-6x},
\]
the same expression. The two routes --- differentiate term by term, or
apply $-F_{x}/F_{y}$ --- are the same computation; use whichever you find
faster, but know that they must agree.

\wstep{Step 6. Vertical tangents: where the denominator vanishes.}
$y'$ blows up when $F_{y}=3y^{2}-6x=0$, i.e.\ $x=y^{2}/2$. Substituting into
the curve:
\[
  \frac{y^{6}}{8}+y^{3}=6\cdot\frac{y^{2}}{2}\cdot y=3y^{3}
  \;\Longrightarrow\;
  \frac{y^{6}}{8}=2y^{3}
  \;\Longrightarrow\;
  y^{3}\bigl(y^{3}-16\bigr)=0 .
\]
So $y=0$ (giving $x=0$) or $y=16^{1/3}=2\sqrt[3]{2}\approx2.520$, giving
$x=y^{2}/2\approx3.175$. At the latter the tangent is genuinely vertical.
\wtrap{At the origin the numerator $2y-x^{2}$ also vanishes, so $y'$ is
$0/0$ there --- undefined for a different reason. The folium
\emph{crosses itself} at $(0,0)$, with two distinct tangents, so no single
function $y(x)$ describes the curve near it. That is precisely the case the
Implicit Function Theorem excludes, and a $0/0$ in the slope formula is how
it announces itself.}

\wstep{Step 7. The economic version of the same calculation.}
An indifference curve is a level set $U(x,y)=\bar u$, and the identical
manipulation gives
\[
  \frac{\mathrm{d}y}{\mathrm{d}x}=-\frac{U_{x}}{U_{y}}=-\mathrm{MRS} .
\]
For $U=x^{1/2}y^{1/2}=4$: $U_{x}=\tfrac12x^{-1/2}y^{1/2}$,
$U_{y}=\tfrac12x^{1/2}y^{-1/2}$, so $y'=-y/x$. At $(4,4)$ the slope is
$-1$; at $(8,2)$ it is $-\tfrac14$ --- the indifference curve flattens as
$x$ becomes abundant, which is diminishing marginal rate of substitution
stated as a derivative.
\end{worked}

\section{Practice Problems}

\begin{problems}
  \item Compute $\dd{y}{x}$ implicitly for $x^{2}+4y^{2}=16$.
  \item Find the equation of the tangent line to $x^{2}+4y^{2}=20$ at
        $(2,2)$.
  \item Use logarithmic differentiation to find $y'$ for $y=x^{\sin x}$
        ($x>0$, $\sin x>0$).
  \item Compute $\dd{^{2}y}{x^{2}}$ implicitly for $x^{2}+4y^{2}=4$.
  \item Compute $\dd{y}{x}$ implicitly for $x^{3}-y^{3}=3xy$.
  \item Compute $\dd{y}{x}$ implicitly for $x^{2}y^{3}=8$.
  \item Find the equation of the tangent line to $x^{2}y+y^{2}x=6$ at
        $(1,2)$.
  \item Use logarithmic differentiation to find $y'$ for $y=(\ln x)^{x}$
        ($x>1$).
  \item Compute $\dd{^{2}y}{x^{2}}$ implicitly for $x^{2}y=1$.
  \item Compute $\dd{y}{x}$ implicitly for $\cos(y)+x=y$.
  \item Find the equation of the tangent line to $y^{2}-x^{2}=3$ at
        $(1,2)$.
  \item Use logarithmic differentiation to find $y'$ for
        \[
          y=\frac{(2x+1)^{5}\sqrt{x^{2}+3}}{(x-4)^{3}}\qquad (x>4).
        \]
  \item Compute $\dd{y}{x}$ implicitly for $\mathrm{e}^{x+y}=xy$.
  \item Compute $\dd{^{2}y}{x^{2}}$ implicitly for $y^{3}=x^{2}$.
  \item Compute $\dd{y}{x}$ implicitly for $x^{2}+2xy+3y^{2}=12$.
  \item Find the equation of the tangent line to $\mathrm{e}^{y}+xy=\mathrm{e}+1$
        at $(1,1)$.
  \item Use logarithmic differentiation to find $y'$ for $y=(x^{2}+1)^{x}$
        ($x>0$).
  \item Compute $\dd{y}{x}$ implicitly for $\sqrt{x+y}=x-y$.
  \item Compute $\dd{^{2}y}{x^{2}}$ implicitly for $xy+y^{2}=1$.
  \item Find the equation of the tangent line to $x^{2}+y^{2}=4y$ at
        $(\sqrt{3},1)$.
  \item Compute $\dd{y}{x}$ implicitly for $\tan(xy)=x$.
  \item Use logarithmic differentiation to find $y'$ for $y=x^{1/x}$
        ($x>0$).
  \item Compute $\dd{y}{x}$ implicitly for $\ln(y)+xy=2$.
  \item Compute $\dd{y}{x}$ implicitly for $y^{3}+y=x^{2}$.
  \item Find the tangent line to the astroid $x^{2/3}+y^{2/3}=2$ at
        $(1,1)$.
  \item Compute $\dd{y}{x}$ implicitly for $\sin(x+y)=y$.
  \item Use logarithmic differentiation to find $y'$ for
        $y=\dfrac{\mathrm{e}^{2x}(x+1)^{3}}{\sqrt{x^{2}+5}}$ (assume $x>-1$).
  \item Compute $\dd{^{2}y}{x^{2}}$ implicitly for $x^{2}+y^{2}=25$.
  \item For an indifference curve given by $U(x,y)=x^{1/2}y^{1/2}=\bar U$
        (constant), compute $\dd{y}{x}$, the marginal rate of
        substitution.
  \item Find all points on the curve $x^{2}+y^{2}=2x+4y$ where the tangent
        line is horizontal.
  \item Compute $\dd{y}{x}$ implicitly for $\arctan(y/x)=x$.
  \item Use logarithmic differentiation to find $y'$ for $y=x^{x}$
        ($x>0$).
  \item Find the equation of the tangent line to $x^{3}+y^{3}=9xy$ at
        $(2,4)$ (point on the folium of Descartes).
  \item Compute $\dd{^{2}y}{x^{2}}$ implicitly for $y^{2}=4x$.
  \item Compute $\dd{y}{x}$ implicitly for $x\mathrm{e}^{y}+y=x$.
  \item Find all points on the ellipse $\dfrac{x^{2}}{9}+\dfrac{y^{2}}{4}=1$
        where the tangent line has slope $-1$.
  \item Use logarithmic differentiation to find $y'$ for
        $y=(\sin x)^{\cos x}$ on $(0,\pi/2)$.
  \item Compute $\dd{y}{x}$ implicitly for $x^{2}y-xy^{2}+x^{2}+y^{2}=0$.
  \item For the production isoquant $K^{1/3}L^{2/3}=\bar Q$ (constant),
        compute $\dd{K}{L}$, the marginal rate of technical
        substitution.
  \item Compute $\dd{^{2}y}{x^{2}}$ implicitly for $\sin(y)=x$.
  \item Find the tangent line to $xy=4$ at $x=2$.
  \item Use logarithmic differentiation to find $y'$ for
        $y=\dfrac{\sqrt{x+1}\sqrt[3]{x-2}}{(x+5)^{4}}$ (assume $x>2$).
  \item Compute $\dd{y}{x}$ implicitly for $\mathrm{e}^{y}\ln(x)+x^{2}=y$.
  \item Show that on the circle $x^{2}+y^{2}=r^{2}$ the tangent at
        $(x_{0},y_{0})$ is perpendicular to the radius. (I.e.\ derive the
        tangent slope and compare with the radius slope.)
  \item Use logarithmic differentiation to find $y'$ for $y=7^{x^{2}}$
        ($y>0$).
  \item Compute $\dd{^{2}y}{x^{2}}$ implicitly for $x^{2}+xy+y^{2}=3$.
  \item For the budget surface $px+qy=m$ (a line, $p,q,m>0$ constants),
        differentiate implicitly to recover $\dd{y}{x}=-p/q$.
  \item Find the tangent line to $x^{2}-xy+y^{2}=7$ at $(3,2)$.
  \item Compute $\dd{y}{x}$ implicitly for $\cos(xy)=x+y$.
  \item Compute $\dd{^{2}y}{x^{2}}$ implicitly for $x^{3}+y^{3}=6xy$ at the
        point $(3,3)$ (folium of Descartes). Express the answer as a
        number.
\end{problems}

\section{Solutions}

\begin{solutions}
  \item $2x+8yy'=0\Rightarrow y'=-\dfrac{x}{4y}$.

  \item Check $(2,2)$: $4+16=20$. From $2x+8yy'=0$: $y'=-x/(4y)$. At
        $(2,2)$: $y'=-2/8=-1/4$. Tangent: $y-2=-\tfrac14(x-2)$, i.e.\
        $x+4y=10$.

  \item $\ln y=\sin(x)\ln(x)\Rightarrow\dfrac{y'}{y}=\cos(x)\ln(x)+\dfrac{\sin x}{x}$.
        So
        \[
          y'=x^{\sin x}\left[\cos(x)\ln(x)+\frac{\sin x}{x}\right] .
        \]

  \item $2x+8yy'=0\Rightarrow y'=-\dfrac{x}{4y}$. Differentiate again:
        \[
          y''=-\frac{4y\cdot 1-x\cdot 4y'}{(4y)^{2}}=-\frac{y-xy'}{4y^{2}} .
        \]
        Substitute $y'=-x/(4y)$:
        $y-xy'=y+\dfrac{x^{2}}{4y}=\dfrac{4y^{2}+x^{2}}{4y}=\dfrac{4}{4y}
        =\dfrac{1}{y}$ (using $x^{2}+4y^{2}=4$). Hence
        $y''=-\dfrac{1}{4y^{3}}$.

  \item $3x^{2}-3y^{2}y'=3y+3xy'\Rightarrow -3(y^{2}+x)y'=3(y-x^{2})
        \Rightarrow y'=\dfrac{x^{2}-y}{y^{2}+x}$.

  \item $2xy^{3}+3x^{2}y^{2}y'=0\Rightarrow
        y'=-\dfrac{2xy^{3}}{3x^{2}y^{2}}=-\dfrac{2y}{3x}$.

  \item Differentiate $x^{2}y+y^{2}x=6$: $2xy+x^{2}y'+2yxy'+y^{2}=0
        \Rightarrow y'(x^{2}+2xy)=-(2xy+y^{2})
        \Rightarrow y'=-\dfrac{y(2x+y)}{x(x+2y)}$. At $(1,2)$:
        $y'=-\dfrac{2\cdot 4}{1\cdot 5}=-\dfrac{8}{5}$. Tangent:
        $y-2=-\tfrac85(x-1)$, i.e.\ $8x+5y=18$.

  \item $\ln y=x\ln(\ln x)\Rightarrow\dfrac{y'}{y}=\ln(\ln x)+x\cdot
        \dfrac{1/x}{\ln x}=\ln(\ln x)+\dfrac{1}{\ln x}$. So
        $y'=(\ln x)^{x}\left[\ln(\ln x)+\dfrac{1}{\ln x}\right]$.

  \item Explicit: $y=x^{-2}$, $y'=-2x^{-3}$, $y''=6x^{-4}$. Implicit
        check: $2xy+x^{2}y'=0\Rightarrow y'=-2y/x$. Differentiate:
        \[
          y''=-\frac{2y'x-2y}{x^{2}}=-\frac{2(-2y/x)x-2y}{x^{2}}
             =-\frac{-4y-2y}{x^{2}}=\frac{6y}{x^{2}}=\frac{6}{x^{4}} .
        \]

  \item $-\sin(y)y'+1=y'\Rightarrow y'(1+\sin y)=1\Rightarrow
        y'=\dfrac{1}{1+\sin y}$.

  \item $2yy'-2x=0\Rightarrow y'=x/y$. At $(1,2)$: $y'=1/2$. Tangent:
        $y-2=\tfrac12(x-1)$, i.e.\ $x-2y+3=0$.

  \item $\ln y=5\ln(2x+1)+\tfrac12\ln(x^{2}+3)-3\ln(x-4)$. Then
        \[
          \frac{y'}{y}=\frac{10}{2x+1}+\frac{x}{x^{2}+3}-\frac{3}{x-4} .
        \]
        Hence
        \[
          y'=\frac{(2x+1)^{5}\sqrt{x^{2}+3}}{(x-4)^{3}}
             \left[\frac{10}{2x+1}+\frac{x}{x^{2}+3}-\frac{3}{x-4}\right] .
        \]

  \item $\mathrm{e}^{x+y}(1+y')=y+xy'\Rightarrow
        y'(\mathrm{e}^{x+y}-x)=y-\mathrm{e}^{x+y}\Rightarrow
        y'=\dfrac{y-\mathrm{e}^{x+y}}{\mathrm{e}^{x+y}-x}$.

  \item $3y^{2}y'=2x\Rightarrow y'=\dfrac{2x}{3y^{2}}$. Differentiate
        again:
        \[
          y''=\frac{2\cdot 3y^{2}-2x\cdot 6yy'}{9y^{4}}
             =\frac{6y^{2}-12xy\cdot y'}{9y^{4}}
             =\frac{2}{3y^{2}}-\frac{4xy'}{3y^{3}} .
        \]
        Substitute $y'=2x/(3y^{2})$:
        \[
          y''=\frac{2}{3y^{2}}-\frac{4x}{3y^{3}}\cdot\frac{2x}{3y^{2}}
             =\frac{2}{3y^{2}}-\frac{8x^{2}}{9y^{5}}
             =\frac{6y^{3}-8x^{2}}{9y^{5}}
             =\frac{6x^{2}-8x^{2}}{9y^{5}}=-\frac{2x^{2}}{9y^{5}},
        \]
        using $y^{3}=x^{2}$. Equivalently, $y''=-\dfrac{2y^{3}}{9y^{5}}
        =-\dfrac{2}{9y^{2}}$.

  \item $2x+2y+2xy'+6yy'=0\Rightarrow y'(2x+6y)=-(2x+2y)\Rightarrow
        y'=-\dfrac{x+y}{x+3y}$.

  \item Differentiate $\mathrm{e}^{y}+xy=\mathrm{e}+1$:
        $\mathrm{e}^{y}y'+y+xy'=0\Rightarrow y'(\mathrm{e}^{y}+x)=-y
        \Rightarrow y'=-\dfrac{y}{\mathrm{e}^{y}+x}$. At $(1,1)$:
        $y'=-\dfrac{1}{\mathrm{e}+1}$. Tangent:
        $y-1=-\dfrac{1}{\mathrm{e}+1}(x-1)$.

  \item $\ln y=x\ln(x^{2}+1)\Rightarrow\dfrac{y'}{y}=\ln(x^{2}+1)
        +x\cdot\dfrac{2x}{x^{2}+1}=\ln(x^{2}+1)+\dfrac{2x^{2}}{x^{2}+1}$.
        So
        \[
          y'=(x^{2}+1)^{x}\left[\ln(x^{2}+1)+\frac{2x^{2}}{x^{2}+1}\right] .
        \]

  \item $\dfrac{1+y'}{2\sqrt{x+y}}=1-y'$. Multiply by $2\sqrt{x+y}$:
        $1+y'=2\sqrt{x+y}(1-y')\Rightarrow y'(1+2\sqrt{x+y})
        =2\sqrt{x+y}-1\Rightarrow y'=\dfrac{2\sqrt{x+y}-1}{1+2\sqrt{x+y}}$.

  \item $y+xy'+2yy'=0\Rightarrow y'(x+2y)=-y\Rightarrow
        y'=-\dfrac{y}{x+2y}$. Differentiate again using the quotient
        rule:
        \[
          y''=-\frac{y'(x+2y)-y(1+2y')}{(x+2y)^{2}}
             =-\frac{y'(x+2y)-y-2yy'}{(x+2y)^{2}}
             =-\frac{xy'-y}{(x+2y)^{2}} .
        \]
        Substitute $y'=-y/(x+2y)$:
        \[
          xy'-y=-\frac{xy}{x+2y}-y=-\frac{xy+y(x+2y)}{x+2y}
               =-\frac{2xy+2y^{2}}{x+2y}=-\frac{2y(x+y)}{x+2y} .
        \]
        Hence
        \[
          y''=\frac{2y(x+y)}{(x+2y)^{3}} .
        \]

  \item $2x+2yy'=4y'\Rightarrow y'(2y-4)=-2x\Rightarrow
        y'=\dfrac{x}{2-y}$. At $(\sqrt{3},1)$: $y'=\dfrac{\sqrt{3}}{1}
        =\sqrt{3}$. Tangent: $y-1=\sqrt{3}(x-\sqrt{3})$, i.e.\
        $y=\sqrt{3}x-2$.

  \item $\sec^{2}(xy)\cdot(y+xy')=1\Rightarrow y+xy'=\cos^{2}(xy)
        \Rightarrow y'=\dfrac{\cos^{2}(xy)-y}{x}$.

  \item $\ln y=\dfrac{\ln x}{x}\Rightarrow\dfrac{y'}{y}
        =\dfrac{(1/x)\cdot x-\ln x\cdot 1}{x^{2}}=\dfrac{1-\ln x}{x^{2}}$.
        So $y'=x^{1/x}\cdot\dfrac{1-\ln x}{x^{2}}$.

  \item $\dfrac{y'}{y}+y+xy'=0\Rightarrow y'\left(\dfrac1y+x\right)=-y
        \Rightarrow y'=-\dfrac{y^{2}}{1+xy}$.

  \item $3y^{2}y'+y'=2x\Rightarrow y'=\dfrac{2x}{3y^{2}+1}$.

  \item Differentiate $x^{2/3}+y^{2/3}=2$: $\tfrac23x^{-1/3}
        +\tfrac23y^{-1/3}y'=0\Rightarrow y'=-\dfrac{y^{1/3}}{x^{1/3}}
        =-\sqrt[3]{y/x}$. At $(1,1)$: $y'=-1$. Tangent: $y-1=-(x-1)$,
        i.e.\ $y=-x+2$.

  \item $\cos(x+y)(1+y')=y'\Rightarrow\cos(x+y)+\cos(x+y)y'=y'
        \Rightarrow\cos(x+y)=y'[1-\cos(x+y)]\Rightarrow
        y'=\dfrac{\cos(x+y)}{1-\cos(x+y)}$.

  \item $\ln y=2x+3\ln(x+1)-\tfrac12\ln(x^{2}+5)$. Differentiate:
        \[
          \frac{y'}{y}=2+\frac{3}{x+1}-\frac{x}{x^{2}+5} .
        \]
        Hence
        \[
          y'=\frac{\mathrm{e}^{2x}(x+1)^{3}}{\sqrt{x^{2}+5}}
             \left[2+\frac{3}{x+1}-\frac{x}{x^{2}+5}\right] .
        \]

  \item $2x+2yy'=0\Rightarrow y'=-x/y$. Differentiate again:
        \[
          y''=-\frac{y\cdot 1-x\cdot y'}{y^{2}}
             =-\frac{y-x(-x/y)}{y^{2}}
             =-\frac{y+x^{2}/y}{y^{2}}
             =-\frac{y^{2}+x^{2}}{y^{3}}=-\frac{25}{y^{3}},
        \]
        using $x^{2}+y^{2}=25$.

  \item Differentiate $x^{1/2}y^{1/2}=\bar U$:
        $\tfrac12x^{-1/2}y^{1/2}+x^{1/2}\cdot\tfrac12y^{-1/2}y'=0$.
        Multiply by $2x^{1/2}y^{1/2}$: $y+xy'=0\Rightarrow y'=-y/x$. So
        $\dd{y}{x}=-\dfrac{y}{x}$ (negative slope along the indifference
        curve, standard MRS).

  \item $2x+2yy'=2+4y'\Rightarrow y'(2y-4)=2-2x\Rightarrow
        y'=\dfrac{1-x}{y-2}$. Horizontal tangent: $y'=0$ requires $x=1$
        (and $y\neq 2$). Plug into the curve: $1+y^{2}=2+4y\Rightarrow
        y^{2}-4y-1=0\Rightarrow y=2\pm\sqrt{5}$. So horizontal tangents at
        $(1,2+\sqrt{5})$ and $(1,2-\sqrt{5})$.

  \item Differentiate $\arctan(y/x)=x$:
        $\dfrac{1}{1+(y/x)^{2}}\cdot\dfrac{y'x-y}{x^{2}}=1$, i.e.\
        $\dfrac{y'x-y}{x^{2}+y^{2}}=1$. Hence $y'x-y=x^{2}+y^{2}$, so
        $y'=\dfrac{x^{2}+y^{2}+y}{x}$.

  \item $\ln y=x\ln x\Rightarrow\dfrac{y'}{y}=\ln x+1\Rightarrow
        y'=x^{x}(\ln x+1)$.

  \item Check $(2,4)$ on $x^{3}+y^{3}=9xy$: $8+64=72$, $9\cdot 8=72$.
        Good. Differentiate: $3x^{2}+3y^{2}y'=9y+9xy'\Rightarrow
        y'(3y^{2}-9x)=9y-3x^{2}\Rightarrow
        y'=\dfrac{9y-3x^{2}}{3y^{2}-9x}=\dfrac{3y-x^{2}}{y^{2}-3x}$.
        At $(2,4)$: $y'=\dfrac{12-4}{16-6}=\dfrac{8}{10}=\dfrac{4}{5}$.
        Tangent: $y-4=\tfrac45(x-2)$, i.e.\ $4x-5y+12=0$.

  \item $2yy'=4\Rightarrow y'=2/y$. Differentiate again:
        $y''=-\dfrac{2y'}{y^{2}}=-\dfrac{2(2/y)}{y^{2}}=-\dfrac{4}{y^{3}}$.

  \item $\mathrm{e}^{y}+x\mathrm{e}^{y}y'+y'=1\Rightarrow
        y'(x\mathrm{e}^{y}+1)=1-\mathrm{e}^{y}\Rightarrow
        y'=\dfrac{1-\mathrm{e}^{y}}{1+x\mathrm{e}^{y}}$.

  \item $\dfrac{2x}{9}+\dfrac{2y}{4}y'=0\Rightarrow y'=-\dfrac{4x}{9y}$.
        Set $-\dfrac{4x}{9y}=-1\Rightarrow 4x=9y\Rightarrow y=\dfrac{4x}{9}$.
        Plug into the ellipse: $\dfrac{x^{2}}{9}+\dfrac{(4x/9)^{2}}{4}=1
        \Rightarrow\dfrac{x^{2}}{9}+\dfrac{16x^{2}}{4\cdot 81}=1
        \Rightarrow\dfrac{x^{2}}{9}+\dfrac{4x^{2}}{81}=1
        \Rightarrow\dfrac{9x^{2}+4x^{2}}{81}=1\Rightarrow
        13x^{2}=81\Rightarrow x=\pm 9/\sqrt{13}$. Corresponding
        $y=\pm 4/\sqrt{13}$ (same sign as $x$). Points:
        $\left(\dfrac{9}{\sqrt{13}},\dfrac{4}{\sqrt{13}}\right)$ and
        $\left(-\dfrac{9}{\sqrt{13}},-\dfrac{4}{\sqrt{13}}\right)$.

  \item $\ln y=\cos(x)\ln(\sin x)\Rightarrow\dfrac{y'}{y}
        =-\sin(x)\ln(\sin x)+\cos(x)\cdot\dfrac{\cos x}{\sin x}
        =-\sin(x)\ln(\sin x)+\dfrac{\cos^{2}x}{\sin x}$. So
        \[
          y'=(\sin x)^{\cos x}\left[-\sin(x)\ln(\sin x)
             +\frac{\cos^{2}x}{\sin x}\right] .
        \]

  \item $2xy+x^{2}y'-y^{2}-x\cdot 2yy'+2x+2yy'=0$. Collect $y'$:
        $y'(x^{2}-2xy+2y)=y^{2}-2xy-2x$. So
        $y'=\dfrac{y^{2}-2xy-2x}{x^{2}-2xy+2y}$.

  \item $\tfrac13K^{-2/3}L^{2/3}\dd{K}{L}\cdot L^{2/3}\ldots$
        carefully: differentiate $K^{1/3}L^{2/3}=\bar Q$ with respect to
        $L$, treating $K=K(L)$:
        \[
          \tfrac13K^{-2/3}K'(L)\cdot L^{2/3}+K^{1/3}\cdot\tfrac23L^{-1/3}=0 .
        \]
        Multiply by $3K^{2/3}L^{1/3}$: $LK'(L)+2K=0\Rightarrow
        K'(L)=-\dfrac{2K}{L}$. So $\dd{K}{L}=-\dfrac{2K}{L}$ (the MRTS).

  \item $\cos(y)y'=1\Rightarrow y'=\dfrac{1}{\cos y}=\sec(y)$.
        Differentiate again: $y''=\sec(y)\tan(y)\cdot y'
        =\sec^{2}(y)\tan(y)$.

  \item At $x=2$, $y=2$. Differentiate $xy=4$: $y+xy'=0\Rightarrow
        y'=-y/x$. At $(2,2)$: $y'=-1$. Tangent: $y-2=-(x-2)$, i.e.\
        $y=-x+4$.

  \item $\ln y=\tfrac12\ln(x+1)+\tfrac13\ln(x-2)-4\ln(x+5)$.
        Differentiate:
        \[
          \frac{y'}{y}=\frac{1}{2(x+1)}+\frac{1}{3(x-2)}-\frac{4}{x+5} .
        \]
        Hence
        \[
          y'=\frac{\sqrt{x+1}\sqrt[3]{x-2}}{(x+5)^{4}}
             \left[\frac{1}{2(x+1)}+\frac{1}{3(x-2)}-\frac{4}{x+5}\right] .
        \]

  \item $\mathrm{e}^{y}y'\ln(x)+\dfrac{\mathrm{e}^{y}}{x}+2x=y'\Rightarrow
        y'[\mathrm{e}^{y}\ln(x)-1]=-\dfrac{\mathrm{e}^{y}}{x}-2x\Rightarrow
        y'=\dfrac{\mathrm{e}^{y}/x+2x}{1-\mathrm{e}^{y}\ln(x)}$.

  \item Differentiate $x^{2}+y^{2}=r^{2}$: $2x+2yy'=0\Rightarrow
        y'=-x/y$. The radius from origin to $(x_{0},y_{0})$ has slope
        $y_{0}/x_{0}$. Product of slopes:
        $(-x_{0}/y_{0})\cdot(y_{0}/x_{0})=-1$, confirming
        perpendicularity.

  \item $\ln y=x^{2}\ln(7)\Rightarrow\dfrac{y'}{y}=2x\ln(7)\Rightarrow
        y'=2x\ln(7)\cdot 7^{x^{2}}$.

  \item Differentiate $x^{2}+xy+y^{2}=3$: $2x+y+xy'+2yy'=0\Rightarrow
        y'(x+2y)=-(2x+y)\Rightarrow y'=-\dfrac{2x+y}{x+2y}$.
        Differentiate again using the quotient rule:
        \[
          y''=-\frac{(2+y')(x+2y)-(2x+y)(1+2y')}{(x+2y)^{2}} .
        \]
        Expand the numerator: $(2+y')(x+2y)=2x+4y+xy'+2yy'$, and
        $(2x+y)(1+2y')=2x+y+4xy'+2yy'$. Subtract:
        $N:=(2+y')(x+2y)-(2x+y)(1+2y')=(2x+4y+xy'+2yy')
        -(2x+y+4xy'+2yy')=3y-3xy'$. So
        $y''=-\dfrac{3y-3xy'}{(x+2y)^{2}}=\dfrac{3(xy'-y)}{(x+2y)^{2}}$.
        Substitute $y'=-\dfrac{2x+y}{x+2y}$:
        \[
          xy'-y=-\frac{x(2x+y)}{x+2y}-y
               =-\frac{x(2x+y)+y(x+2y)}{x+2y}
        \]
        \[
          \phantom{xy'-y}=-\frac{2x^{2}+xy+xy+2y^{2}}{x+2y}
               =-\frac{2(x^{2}+xy+y^{2})}{x+2y},
        \]
        using $x^{2}+xy+y^{2}=3$. Hence
        \[
          y''=\frac{3}{(x+2y)^{2}}\cdot\left(-\frac{6}{x+2y}\right)
             =-\frac{18}{(x+2y)^{3}} .
        \]

  \item Differentiate $px+qy=m$ w.r.t.\ $x$: $p+qy'=0\Rightarrow
        y'=-p/q$, as required (the budget line has constant slope).

  \item Differentiate $x^{2}-xy+y^{2}=7$: $2x-y-xy'+2yy'=0\Rightarrow
        y'(2y-x)=y-2x\Rightarrow y'=\dfrac{y-2x}{2y-x}$. At $(3,2)$:
        $y'=\dfrac{2-6}{4-3}=-4$. Tangent: $y-2=-4(x-3)$, i.e.\
        $y=-4x+14$.

  \item Differentiate $\cos(xy)=x+y$: $-\sin(xy)\cdot(y+xy')=1+y'
        \Rightarrow -y\sin(xy)-x\sin(xy)y'=1+y'\Rightarrow
        -y'(x\sin(xy)+1)=1+y\sin(xy)\Rightarrow
        y'=-\dfrac{1+y\sin(xy)}{1+x\sin(xy)}$.

  \item First verify $(3,3)$: $27+27=54$ and $6\cdot 9=54$. Good.
        Differentiate $x^{3}+y^{3}=6xy$: $3x^{2}+3y^{2}y'=6y+6xy'
        \Rightarrow y'(3y^{2}-6x)=6y-3x^{2}\Rightarrow
        y'=\dfrac{6y-3x^{2}}{3y^{2}-6x}=\dfrac{2y-x^{2}}{y^{2}-2x}$.
        At $(3,3)$: $y'=\dfrac{6-9}{9-6}=-1$. For $y''$, differentiate
        $y'=\dfrac{2y-x^{2}}{y^{2}-2x}$:
        \[
          y''=\frac{(2y'-2x)(y^{2}-2x)-(2y-x^{2})(2yy'-2)}{(y^{2}-2x)^{2}} .
        \]
        At $(3,3)$ with $y'=-1$: numerator pieces:
        \[
          (2y'-2x)(y^{2}-2x)=(-2-6)(9-6)=(-8)(3)=-24,
        \]
        \[
          (2y-x^{2})(2yy'-2)=(6-9)(6\cdot(-1)-2)=(-3)(-8)=24 .
        \]
        So numerator $=-24-24=-48$ and denominator $=(9-6)^{2}=9$.
        Hence $y''(3)=-\dfrac{48}{9}=-\dfrac{16}{3}$.
\end{solutions}

\section{Additional Problems}

\begin{problems}
  \item Compute $\dd{y}{x}$ implicitly for $x^{2}+9y^{2}=36$.
  \item Find the equation of the tangent line to $x^{2}+y^{2}=25$ at
        $(3,4)$.
  \item Use logarithmic differentiation to find $y'$ for $y=x^{\cos x}$
        ($x>0$).
  \item Compute $\dd{y}{x}$ implicitly for $x^{2}y+xy^{2}=6$.
  \item Compute $\dd{^{2}y}{x^{2}}$ implicitly for $x^{2}-y^{2}=4$.
  \item Find the equation of the tangent line to $x^{3}+xy+y^{3}=11$ at
        $(1,2)$.
  \item Use logarithmic differentiation to find $y'$ for
        \[
          y=\frac{(x+1)^{2}\sqrt{x-3}}{(2x+5)^{4}}\qquad (x>3).
        \]
  \item Compute $\dd{y}{x}$ implicitly for $\mathrm{e}^{xy}=x+y$.
  \item For an indifference curve $U(x,y)=x^{1/3}y^{2/3}=\bar U$
        (constant), compute $\dd{y}{x}$, the marginal rate of
        substitution.
  \item Compute $\dd{y}{x}$ implicitly for $\sin(x+y)=xy$.
  \item Find the equation of the tangent line to $x^{2}-2xy+3y^{2}=9$ at
        $(1,2)$.
  \item Compute $\dd{^{2}y}{x^{2}}$ implicitly for $x^{2}+y^{2}=9$.
  \item Find all points on the curve $x^{2}+y^{2}=4x-2y$ where the tangent
        line is vertical.
  \item For the production isoquant $K^{2/5}L^{3/5}=\bar Q$ (constant),
        compute $\dd{K}{L}$, the marginal rate of technical
        substitution.
  \item Use logarithmic differentiation to find $y'$ for
        $y=(x^{2}+1)^{\sin x}$ ($x\in\R$).

  \item The curve $k$ below is the set of points satisfying
        \[
          x^{3}-3x^{2}+y^{2}-2y=-1 .
        \]
        Verify that $(0,1)$, $(3,1)$, $(-1,3)$ and $(-1,-1)$ all lie on
        $k$.

\begin{gallery}
\begin{tikzpicture}
  \begin{axis}[geoaxis, xmin=-1.6, xmax=3.6, ymin=-1.9, ymax=3.9,
      xtick={-1,1,2,3}, ytick={-1,1,2,3}, title={the curve $k$}]
    \addplot[domain=-1.2:3, samples=241]{1+abs(x)*sqrt(3-x)};
    \addplot[domain=-1.2:3, samples=241]{1-abs(x)*sqrt(3-x)};
  \end{axis}
\end{tikzpicture}\hfill
\begin{tikzpicture}
  \begin{axis}[geoaxis, xmin=-1.4, xmax=1.4, ymin=-0.9, ymax=0.9,
      xtick={-1,1}, ytick={-0.5,0.5}, title={the curve $\ell$}]
    \addplot[domain=-1:1, samples=241]{x*sqrt(1-x^2)};
    \addplot[domain=-1:1, samples=241]{-x*sqrt(1-x^2)};
  \end{axis}
\end{tikzpicture}
\end{gallery}

  \item Name every point of $k$ at which the hypotheses of the implicit
        function theorem fail for writing $y$ as a function of $x$, and say
        which graphical feature identifies each one.
  \item Confirm the previous answer by computation: solve $F_{y}=0$
        together with $F=0$.
  \item Compute the slope of the tangent to $k$ at $(-1,3)$.
  \item Compute the slope of the tangent to $k$ at $(-1,-1)$, and explain
        why it is the negative of the previous answer.
  \item Find every point of $k$ at which the tangent is horizontal.
  \item Near $(3,1)$ no function $y(x)$ describes $k$. Show that
        $x=x(y)$ does exist there, and compute $\dd{x}{y}$ at that point.
  \item The curve $\ell$ above is $y^{2}=x^{2}\bigl(1-x^{2}\bigr)$. Write it
        as $F(x,y)=0$ and give $F_{x}$ and $F_{y}$.
  \item Name every point of $\ell$ at which the implicit function theorem
        fails, and the graphical feature of each.
  \item Compute the slope of the tangent to $\ell$ at
        $\bigl(\tfrac12,\tfrac{\sqrt3}{4}\bigr)$.
  \item Compute the slope of the tangent to $\ell$ at
        $\bigl(\tfrac12,-\tfrac{\sqrt3}{4}\bigr)$.
  \item Find the four points of $\ell$ at which the tangent is horizontal.
  \item At the origin two branches of $\ell$ cross. Separate them by solving
        for $y$, and compute the slope of each branch there.
  \item Explain in one or two sentences why $-F_{x}/F_{y}$ cannot deliver
        those two slopes, and what has to be done instead.
\end{problems}

\section{Additional Solutions}

\begin{solutions}
  \item $2x+18yy'=0\Rightarrow y'=-\dfrac{x}{9y}$.

  \item Check $(3,4)$: $9+16=25$. From $2x+2yy'=0$: $y'=-x/y$. At
        $(3,4)$: $y'=-3/4$. Tangent: $y-4=-\tfrac34(x-3)$, i.e.\
        $3x+4y=25$.

  \item $\ln y=\cos(x)\ln(x)\Rightarrow\dfrac{y'}{y}
        =-\sin(x)\ln(x)+\dfrac{\cos x}{x}$. So
        \[
          y'=x^{\cos x}\left[-\sin(x)\ln(x)+\frac{\cos x}{x}\right] .
        \]

  \item Differentiate $x^{2}y+xy^{2}=6$: $2xy+x^{2}y'+y^{2}+2xyy'=0
        \Rightarrow y'(x^{2}+2xy)=-(2xy+y^{2})
        \Rightarrow y'=-\dfrac{y(2x+y)}{x(x+2y)}$.

  \item $2x-2yy'=0\Rightarrow y'=x/y$. Differentiate again:
        \[
          y''=\frac{y-xy'}{y^{2}}=\frac{y-x(x/y)}{y^{2}}
             =\frac{y^{2}-x^{2}}{y^{3}}=\frac{-4}{y^{3}},
        \]
        using $x^{2}-y^{2}=4\Rightarrow y^{2}-x^{2}=-4$.

  \item Check $(1,2)$: $1+2+8=11$. Differentiate:
        $3x^{2}+y+xy'+3y^{2}y'=0\Rightarrow y'(x+3y^{2})=-(3x^{2}+y)
        \Rightarrow y'=-\dfrac{3x^{2}+y}{x+3y^{2}}$. At $(1,2)$:
        $y'=-\dfrac{3+2}{1+12}=-\dfrac{5}{13}$. Tangent:
        $y-2=-\tfrac{5}{13}(x-1)$, i.e.\ $5x+13y=31$.

  \item $\ln y=2\ln(x+1)+\tfrac12\ln(x-3)-4\ln(2x+5)$. Then
        \[
          \frac{y'}{y}=\frac{2}{x+1}+\frac{1}{2(x-3)}-\frac{8}{2x+5} .
        \]
        Hence
        \[
          y'=\frac{(x+1)^{2}\sqrt{x-3}}{(2x+5)^{4}}
             \left[\frac{2}{x+1}+\frac{1}{2(x-3)}-\frac{8}{2x+5}\right] .
        \]

  \item Differentiate $\mathrm{e}^{xy}=x+y$:
        $\mathrm{e}^{xy}(y+xy')=1+y'\Rightarrow
        y'(x\mathrm{e}^{xy}-1)=1-y\mathrm{e}^{xy}\Rightarrow
        y'=\dfrac{1-y\mathrm{e}^{xy}}{x\mathrm{e}^{xy}-1}$.

  \item Differentiate $x^{1/3}y^{2/3}=\bar U$:
        $\tfrac13x^{-2/3}y^{2/3}+x^{1/3}\cdot\tfrac23y^{-1/3}y'=0$.
        Multiply by $3x^{2/3}y^{1/3}$: $y+2xy'=0\Rightarrow y'=-y/(2x)$.
        So $\dd{y}{x}=-\dfrac{y}{2x}$ (the MRS).

  \item Differentiate $\sin(x+y)=xy$: $\cos(x+y)(1+y')=y+xy'
        \Rightarrow y'(\cos(x+y)-x)=y-\cos(x+y)\Rightarrow
        y'=\dfrac{y-\cos(x+y)}{\cos(x+y)-x}$.

  \item Check $(1,2)$: $1-4+12=9$. Good. Differentiate:
        $2x-2y-2xy'+6yy'=0\Rightarrow y'(6y-2x)=2y-2x\Rightarrow
        y'=\dfrac{y-x}{3y-x}$. At $(1,2)$: $y'=\dfrac{1}{5}$. Tangent:
        $y-2=\tfrac15(x-1)$, i.e.\ $x-5y+9=0$.

  \item $2x+2yy'=0\Rightarrow y'=-x/y$. Differentiate again:
        \[
          y''=-\frac{y-xy'}{y^{2}}=-\frac{y-x(-x/y)}{y^{2}}
             =-\frac{y+x^{2}/y}{y^{2}}
             =-\frac{y^{2}+x^{2}}{y^{3}}=-\frac{9}{y^{3}},
        \]
        using $x^{2}+y^{2}=9$.

  \item $2x+2yy'=4-2y'\Rightarrow y'(2y+2)=4-2x\Rightarrow
        y'=\dfrac{2-x}{y+1}$. Vertical tangent: $y'$ undefined when
        $y+1=0\Rightarrow y=-1$ (and $x\neq 2$). Plug into the curve:
        $x^{2}+1=4x+2\Rightarrow x^{2}-4x-1=0\Rightarrow
        x=2\pm\sqrt{5}$. Vertical tangents at $(2+\sqrt{5},-1)$ and
        $(2-\sqrt{5},-1)$.

  \item Differentiate $K^{2/5}L^{3/5}=\bar Q$ w.r.t.\ $L$, treating
        $K=K(L)$:
        \[
          \tfrac25K^{-3/5}K'(L)\cdot L^{3/5}+K^{2/5}\cdot\tfrac35L^{-2/5}=0 .
        \]
        Multiply by $5K^{3/5}L^{2/5}$: $2LK'(L)+3K=0\Rightarrow
        K'(L)=-\dfrac{3K}{2L}$. So $\dd{K}{L}=-\dfrac{3K}{2L}$ (the
        MRTS).

  \item $\ln y=\sin(x)\ln(x^{2}+1)\Rightarrow\dfrac{y'}{y}
        =\cos(x)\ln(x^{2}+1)+\sin(x)\cdot\dfrac{2x}{x^{2}+1}$. So
        \[
          y'=(x^{2}+1)^{\sin x}\left[\cos(x)\ln(x^{2}+1)
             +\frac{2x\sin x}{x^{2}+1}\right] .
        \]

  \item With $F(x,y)=x^{3}-3x^{2}+y^{2}-2y+1$:
        $F(0,1)=0-0+1-2+1=0$; \;$F(3,1)=27-27+1-2+1=0$;
        \;$F(-1,3)=-1-3+9-6+1=0$; \;$F(-1,-1)=-1-3+1+2+1=0$.
        All four are on $k$.
  \item Graphically: at $(0,1)$ the curve \textbf{crosses itself}, and at
        $(3,1)$ it has a \textbf{vertical tangent}. At every other point of
        $k$ a local $y(x)$ exists.
  \item $F_{y}=2y-2=0\iff y=1$. Substituting $y=1$ into $F=0$ gives
        $x^{3}-3x^{2}+1-2+1=x^{2}(x-3)=0$, so $x=0$ (a double root) or
        $x=3$: the points $(0,1)$ and $(3,1)$. At $(3,1)$,
        $F_{x}=3x^{2}-6x=9\neq0$, so this is the vertical-tangent case; at
        $(0,1)$, $F_{x}=0$ as well, the singular case.
  \item $F_{x}=3x^{2}-6x$, $F_{y}=2y-2$, so
        $y'=-\dfrac{3x^{2}-6x}{2y-2}$. At $(-1,3)$:
        $y'=-\dfrac{3+6}{6-2}=-\dfrac94$.
  \item At $(-1,-1)$: $y'=-\dfrac{3+6}{-2-2}=\dfrac94$. The equation
        contains $y$ only through $(y-1)^{2}$, so $k$ is symmetric about the
        line $y=1$; reflecting a curve in a horizontal line negates every
        slope, and $(-1,3)$ and $(-1,-1)$ are mirror images.
  \item Horizontal tangent needs $F_{x}=0$ and $F_{y}\neq0$:
        $3x^{2}-6x=3x(x-2)=0$ gives $x=0$ or $x=2$. The value $x=0$ forces
        $y=1$, where $F_{y}=0$ too --- that is the singular point, not a
        horizontal tangent. For $x=2$: $(y-1)^{2}=x^{2}(3-x)=4$, so $y=3$ or
        $y=-1$. The tangent is horizontal at $(2,3)$ and $(2,-1)$.
  \item At $(3,1)$ we have $F_{x}=9\neq0$, so the theorem applies with the
        roles of the variables exchanged and gives a function $x=x(y)$.
        Differentiating $F\bigl(x(y),y\bigr)=0$ with respect to $y$:
        $F_{x}\,x'(y)+F_{y}=0$, hence
        $\dd{x}{y}=-\dfrac{F_{y}}{F_{x}}=-\dfrac{0}{9}=0$ --- the curve is
        momentarily vertical, which is exactly a horizontal tangent for
        $x(y)$.
  \item $F(x,y)=y^{2}-x^{2}+x^{4}$, with $F_{x}=-2x+4x^{3}=2x(2x^{2}-1)$
        and $F_{y}=2y$.
  \item $F_{y}=2y=0\iff y=0$, and then $x^{2}(1-x^{2})=0$ gives
        $x=0,\pm1$. At $(\pm1,0)$: $F_{x}=\pm2\neq0$, \textbf{vertical
        tangents} --- the two tips of the figure eight. At $(0,0)$:
        $F_{x}=0$ too, the \textbf{self-intersection} in the middle.
  \item $y'=-\dfrac{2x(2x^{2}-1)}{2y}$. At
        $\bigl(\tfrac12,\tfrac{\sqrt3}{4}\bigr)$ the numerator is
        $2\cdot\tfrac12\bigl(\tfrac12-1\bigr)=-\tfrac12$ and the denominator
        is $\tfrac{\sqrt3}{2}$, so
        $y'=\dfrac{1/2}{\sqrt3/2}=\dfrac{1}{\sqrt3}=\dfrac{\sqrt3}{3}$.
  \item Only the sign of $y$ changes, so
        $y'=-\dfrac{\sqrt3}{3}$.
  \item $F_{x}=2x(2x^{2}-1)=0$ with $F_{y}\neq0$ gives $x=\pm
        \tfrac{1}{\sqrt2}$ (the root $x=0$ is the singular point). Then
        $y^{2}=x^{2}(1-x^{2})=\tfrac12\cdot\tfrac12=\tfrac14$, so
        $y=\pm\tfrac12$: the four points
        $\bigl(\pm\tfrac{1}{\sqrt2},\pm\tfrac12\bigr)$.
  \item Solving gives $y=x\sqrt{1-x^{2}}$ and $y=-x\sqrt{1-x^{2}}$. For the
        first,
        $y'=\sqrt{1-x^{2}}-\dfrac{x^{2}}{\sqrt{1-x^{2}}}$, which at $x=0$
        equals $1$; the second branch gives $-1$. The two branches cross the
        origin at $+45^{\circ}$ and $-45^{\circ}$.
  \item At the origin both partials vanish, so $-F_{x}/F_{y}$ is $0/0$: it
        carries no information, and a single number could not describe two
        different slopes anyway. The branches have to be separated first ---
        here by solving for $y$ --- and each differentiated on its own.
\end{solutions}

\chapter{Applications: Tangent, Linearization, Marginals, Elasticity, Optimization}

The derivative is the bridge between calculus and economics. Geometrically it
gives the slope of the tangent (and hence of the normal); analytically it
produces the linearisation that turns nonlinear functions into
easy-to-evaluate affine ones; in economics it is the marginal quantity
(cost, revenue, utility) and, after rescaling, the elasticity. The first-
and second-order conditions $f'=0$, $f''<0$ (or $>0$) drive one-variable
optimisation.

\section{Core Rules}

\begin{rulebox}{Tangent line}
At $(a,f(a))$: $y = f(a) + f'(a)(x-a)$.
\end{rulebox}

\begin{rulebox}{Normal line}
At $(a,f(a))$: perpendicular to the tangent, slope $-1/f'(a)$ when
$f'(a)\neq 0$:
\[
  y = f(a) - \frac{1}{f'(a)}(x-a).
\]
\end{rulebox}

\begin{rulebox}{Linear approximation (linearisation)}
For $x$ near $a$,
\[
  f(x) \approx L(x) := f(a) + f'(a)(x-a), \qquad \Delta f \approx f'(a)\,\Delta x .
\]
\end{rulebox}

\begin{rulebox}{Marginal quantities}
If $C(Q)$, $R(Q)$, $\pi(Q)=R(Q)-C(Q)$ are total cost, revenue and profit,
and $U(x)$ is utility, then
\[
  MC(Q)=C'(Q), \quad MR(Q)=R'(Q), \quad M\pi(Q)=\pi'(Q), \quad MU(x)=U'(x).
\]
Profit-maximising output satisfies $MR(Q^*)=MC(Q^*)$ (FOC) and
$\pi''(Q^*)<0$ (SOC).
\end{rulebox}

\begin{rulebox}{Elasticity}
The (point) elasticity of $f$ with respect to $x$ is
\[
  \varepsilon = \frac{x}{f(x)}f'(x) = \frac{\mathrm{d}f/f}{\mathrm{d}x/x}
    = \frac{\mathrm{d}\ln f}{\mathrm{d}\ln x} .
\]
For demand $q=q(p)$: $\displaystyle \varepsilon_d = \frac{p}{q}q'(p)$
(often negative). Demand is \emph{elastic} when $|\varepsilon_d|>1$,
\emph{inelastic} when $|\varepsilon_d|<1$, \emph{unit elastic} when
$|\varepsilon_d|=1$.
\end{rulebox}

\begin{rulebox}{1D unconstrained optimisation}
\begin{enumerate}[leftmargin=1.6em,itemsep=0.3ex]
  \item Solve $f'(x)=0$ on the relevant domain to find critical points.
  \item Classify with the second-derivative test: $f''(x^*)<0 \Rightarrow$
        local max, $f''(x^*)>0 \Rightarrow$ local min.
  \item Compare with boundary values and endpoint behaviour to identify
        global extrema.
\end{enumerate}
\end{rulebox}

\section{Worked Example}

\begin{worked}{Demand is $q=100-2p$. Find the price elasticity at $p=20$ and at the revenue-maximising price, and approximate $\sqrt{101}$ by linearisation}
\wstep{Step 1. Elasticity is a derivative times a ratio --- not a derivative.}
\[
  \varepsilon=\frac{\mathrm{d}q}{\mathrm{d}p}\cdot\frac{p}{q}
   =\frac{\%\ \text{change in }q}{\%\ \text{change in }p} .
\]
Here $\mathrm{d}q/\mathrm{d}p=-2$ everywhere, but the elasticity is not
constant, because $p/q$ moves along the curve.
\wtrap{$\varepsilon\neq\mathrm{d}q/\mathrm{d}p$. The slope has units
(quantity per euro) and changes if you re-denominate the price; the
elasticity is a pure number and does not. That invariance is the main
reason economists use it.}

\wstep{Step 2. Evaluate at $p=20$.}
$q=100-40=60$, so
\[
  \varepsilon=-2\cdot\frac{20}{60}=-\tfrac23\approx-0.67 .
\]
$|\varepsilon|<1$: demand is \textbf{inelastic} here. A $1\%$ price rise
costs only $0.67\%$ of quantity, so revenue \emph{rises} --- the firm is
pricing below the revenue maximum.
\wtrap{The sign is negative for any downward-sloping demand, and dropping it
is a real error even though the classification uses $|\varepsilon|$. Report
$-0.67$ and then say ``inelastic''; do not report $0.67$.}

\wstep{Step 3. Maximise revenue directly.}
\[
  R(p)=pq=p(100-2p)=100p-2p^{2},
  \qquad
  R'(p)=100-4p=0
  \quad\Longrightarrow\quad
  p^{*}=25 .
\]
$R''=-4<0$, so this is a maximum. Then $q^{*}=50$ and $R^{*}=1250$.

\wstep{Step 4. Check the elasticity there --- and find the general rule.}
\[
  \varepsilon(25)=-2\cdot\frac{25}{50}=-1 .
\]
Unit elastic, exactly at the revenue peak. That is no coincidence:
\[
  R=pq
  \;\Longrightarrow\;
  \frac{\mathrm{d}R}{\mathrm{d}p}=q+p\frac{\mathrm{d}q}{\mathrm{d}p}
  =q\Bigl(1+\varepsilon\Bigr),
\]
which vanishes exactly when $\varepsilon=-1$. Revenue rises where demand is
inelastic, falls where it is elastic, and peaks in between.

\wstep{Step 5. Marginal revenue, and the same fact in quantity terms.}
Inverting demand, $p=50-q/2$, so $R(q)=50q-q^{2}/2$ and
\[
  MR(q)=50-q .
\]
$MR=0$ at $q=50$ --- the same point. Note $MR$ falls twice as fast as the
demand curve $p=50-q/2$: the standard ``$MR$ has twice the slope'' result,
here visible as $-1$ against $-\tfrac12$.
\wtrap{Marginal revenue is $\mathrm{d}R/\mathrm{d}q$, not $R/q$ (that is
\emph{average} revenue, which equals the price). Confusing marginal with
average is the most consequential error in applied differentiation, and it
recurs for cost: $MC=C'(q)$, while $AC=C(q)/q$.}

\wstep{Step 6. Linearisation: the tangent line as a calculator.}
For $f(x)=\sqrt{x}$ near $a=100$, where the value is known exactly:
\[
  f(x)\approx f(a)+f'(a)(x-a),
  \qquad
  f'(x)=\frac{1}{2\sqrt x},\quad f'(100)=\frac{1}{20} .
\]
\[
  \sqrt{101}\approx10+\frac{1}{20}(1)=10.05 .
\]
The true value is $10.0499$, so the error is $10^{-4}$.
\wtrap{Expand about the nearest point whose value you actually know. Using
$a=101$ makes the formula exact and useless; using $a=64$ would be legal but
far away, and the error grows with $(x-a)^{2}$. Choose $a$ first, and choose
it well.}

\wstep{Step 7. The same tangent, three names.}
\[
  \underbrace{y=f(a)+f'(a)(x-a)}_{\text{tangent line}}
  \;=\;
  \underbrace{f(a)+f'(a)\,\Delta x}_{\text{linearisation}}
  \;=\;
  \underbrace{f(a)+\text{marginal}\times\text{one more unit}}_{\text{economics}} .
\]
``Marginal cost is the cost of one more unit'' is a linear approximation:
strictly $C'(q)$ is a rate, and $C(q+1)-C(q)\approx C'(q)$ with an error of
order $C''/2$. The approximation is good precisely when the curvature is
small.
\end{worked}

\section{Practice Problems}

\begin{problems}
  \item Find the equation of the tangent to $f(x)=x^3-2x+5$ at $x=1$.
  \item Use the linearisation of $f(x)=\sqrt{x}$ at $a=25$ to estimate
        $\sqrt{25.4}$.
  \item For $C(Q)=0.2Q^2+5Q+80$, find $MC(Q)$ and $MC(15)$.
  \item For demand $q(p)=200-5p$, compute $\varepsilon_d(p)$ and evaluate at
        $p=10$.
  \item Maximise profit $\pi(Q)=R(Q)-C(Q)$ where $R(Q)=60Q-Q^2$ and
        $C(Q)=Q^2+10Q+25$. Verify the SOC.
  \item Find the tangent to $f(x)=1/x$ at $x=2$.
  \item Use the linearisation of $f(x)=x^{1/4}$ at $a=16$ to estimate
        $16.3^{1/4}$.
  \item Inverse demand $p(Q)=80-0.4Q$. Find $R(Q)$, $MR(Q)$, and $MR(50)$.
  \item For the Cobb--Douglas production $Y=AK^\alpha L^\beta$ with $L$
        fixed, compute the output elasticity $\varepsilon_{Y,K}$.
  \item Minimise the average cost $AC(Q)=C(Q)/Q$ for
        $C(Q)=2Q^2+8Q+200$, over $Q>0$.
  \item Find the tangent to $f(x)=e^{2x}$ at $x=0$.
  \item Estimate $\ln(0.98)$ using the linearisation of $f(x)=\ln(x)$ at
        $a=1$.
  \item Using $C$ from problem 3 and $R$ from problem 8, find the
        profit-maximising output $Q^*$ and verify SOC.
  \item For demand $q(p)=80e^{-0.05p}$, compute $\varepsilon_d(p)$ and at
        $p=20$.
  \item A monopolist faces $p(Q)=150-Q$ and $C(Q)=Q^2+30Q+100$. Find $Q^*$,
        $p^*$, and $\pi^*$.
  \item Find the normal line to $f(x)=\sqrt{x}$ at $x=9$.
  \item Estimate $e^{-0.2}$ using the linearisation of $f(x)=e^x$ at $a=0$.
  \item A monopolist faces $p(Q)=200-3Q$ and $C(Q)=2Q^2+10Q+50$. Find
        $Q^*$, $p^*$, and $\pi^*$.
  \item For $q(p)=200-5p$, find the price at which demand is unit elastic.
  \item A firm sells at the fixed price $p=24$ per unit; cost is
        $C(Q)=Q^2+8Q+10$. Find the profit-maximising output and verify SOC.
  \item Find the normal line to $f(x)=\ln(x)$ at $x=e$.
  \item A firm's cost is $C(Q)=0.5Q^2+8Q+200$. Use linear approximation at
        $Q=20$ to estimate $C(21.5)$.
  \item Marginal cost is $MC(Q)=0.06Q^2-2.4Q+30$. Find the production level
        that minimises $MC$ and verify the SOC.
  \item Demand $q(p)=\dfrac{300}{p^2+4}$. Compute $\varepsilon_d(p)$ and
        evaluate at $p=2$.
  \item Maximise utility $U(x)=4\ln(x)-0.5x$ for $x>0$.
  \item Find the tangent to $f(x)=\cos(x)$ at $x=\pi/3$.
  \item Use linearisation of $f(x)=(1+x)^{10}$ at $a=0$ to estimate
        $(1.02)^{10}$.
  \item For $C(Q)=Q^3-9Q^2+30Q+40$, compute $MC(Q)$ and locate the $Q$
        where $MC$ attains its minimum.
  \item For demand $q(p)=Ap^{-\alpha}$ with $A,\alpha>0$, show that the
        (point) elasticity is the constant $-\alpha$.
  \item A producer faces inverse demand $p(Q)=100-2Q$ and constant
        marginal cost $MC=20$ (so $C(Q)=20Q+F$ with fixed cost $F$). Find
        $Q^*$, $p^*$, and verify SOC.
  \item Find the normal line to $f(x)=x^2$ at $x=1$.
  \item Use linearisation of $f(x)=\tan(x)$ at $a=0$ to estimate
        $\tan(0.05)$.
  \item For the cubic cost $C(Q)=Q^3-6Q^2+15Q+20$, find marginal cost and
        average cost, and the $Q>0$ where $AC$ is minimised. Verify SOC.
  \item For demand $q(p)=100-2p^2$ (defined on $0\le p\le\sqrt{50}$),
        compute $\varepsilon_d(p)$ at $p=4$ and classify.
  \item Maximise revenue $R(Q)=Q\cdot p(Q)$ when $p(Q)=120-2Q$ and verify
        SOC.
  \item Find the tangent and the normal to $f(x)=x^2-3x+2$ at $x=2$.
  \item For a firm with $C(Q)=4Q^2+20Q+100$, what is the change in cost
        (approximately) when output increases from 10 to 10.5?
  \item For $C(Q)=Q^2+10Q+9$ and $R(Q)=30Q$, find the profit-maximising
        output and the break-even points (where $\pi=0$).
  \item For demand $q(p)=400-p^2$, find the price at which
        $|\varepsilon_d|=1$.
  \item Find positive $x$ that maximises $f(x)=x^2e^{-x}$ on
        $(0,\infty)$. Verify SOC.
  \item Find the tangent to $f(x)=\arctan(x)$ at $x=1$.
  \item Estimate $\sqrt[3]{27.9}$ using the linearisation of
        $f(x)=x^{1/3}$ at $a=27$.
  \item For utility $U(x)=\sqrt{x}$ and consumption $x=100$, compute the
        marginal utility $MU(x)$ and its elasticity $\varepsilon_{U,x}$.
        Comment briefly on signs.
  \item A monopolist faces demand $q(p)=120-3p$ and constant unit cost
        $c=10$. Reformulate via the inverse demand to find $Q^*$, $p^*$,
        $\pi^*$. Verify SOC.
  \item Minimise $f(x)=x+\dfrac{16}{x}$ on $(0,\infty)$. Verify SOC.
  \item Find both the tangent and the normal to $f(x)=\sin(x)$ at
        $x=\pi/2$.
  \item For revenue $R(Q)=p(Q)\cdot Q$ with $p(Q)=100-Q$, express $MR$ in
        terms of $p$ and the (price-)elasticity of demand: verify the
        identity $MR = p\left(1+\dfrac{1}{\varepsilon_d}\right)$ at
        $Q=20$.
  \item Use linearisation of $f(x)=\ln(1+x)$ at $a=0$ to estimate
        $\ln(1.05)$, and compute the linearisation error.
  \item A firm has revenue $R(Q)=20Q-Q^2$ and cost
        $C(Q)=2Q+\frac12Q^2+10$, with the constraint that $Q\in[0,8]$.
        Find the profit-maximising $Q$ on this closed interval, comparing
        the interior critical point with the endpoints.
  \item For demand $q(p)=50/p$, compute $\varepsilon_d$ for all $p>0$ and
        explain what it implies for total revenue $R(p)=p\cdot q(p)$.
  \item Maximise $\pi(Q)=12Q-Q^2-\dfrac{Q^3}{3}$ for $Q\ge0$. (Take the
        unique non-negative critical point and verify SOC.)
  \item For inverse demand $p(Q)=50-Q$ and cost $C(Q)=10Q+Q^2$, compute
        $\pi(Q)$, $Q^*$, $p^*$, $\pi^*$ and also $\varepsilon_d$ at $Q^*$.
  \item A firm produces according to $Y=\sqrt{L}$ (a strictly concave
        one-input production function). The wage is $w=2$ and the output
        price is $p=8$. Choose $L\ge0$ to maximise profit
        $\pi(L)=p\cdot Y(L)-wL$. Verify SOC.
\end{problems}

\section{Solutions}

\begin{solutions}
  \item $f'(x)=3x^2-2$. $f(1)=4$, $f'(1)=1$. Tangent:
        $y=4+1\cdot(x-1)=x+3$.

  \item $f(25)=5$, $f'(x)=1/(2\sqrt{x})$, $f'(25)=1/10=0.1$. So
        $\sqrt{25.4}\approx 5+0.1\cdot0.4=5.04$. (True: $5.0399\ldots$)

  \item $MC(Q)=C'(Q)=0.4Q+5$. $MC(15)=6+5=11$.

  \item $q'(p)=-5$.
        $\varepsilon_d(p)=\dfrac{p}{200-5p}\cdot(-5)=\dfrac{-5p}{200-5p}
        =\dfrac{-p}{40-p}$. At $p=10$: $\varepsilon_d=-10/30=-1/3$, so
        $|\varepsilon_d|<1$, inelastic.

  \item $\pi(Q)=60Q-Q^2-Q^2-10Q-25=50Q-2Q^2-25$. FOC:
        $\pi'(Q)=50-4Q=0\Rightarrow Q^*=12.5$. SOC: $\pi''(Q)=-4<0$, so a
        maximum. $\pi(12.5)=625-312.5-25=287.5$.

  \item $f'(x)=-1/x^2$. $f(2)=1/2$, $f'(2)=-1/4$. Tangent:
        $y=\frac12-\frac14(x-2)=-\frac{x}{4}+1$.

  \item $f'(x)=\frac14x^{-3/4}$. $f(16)=2$,
        $f'(16)=\frac14\cdot16^{-3/4}=\frac14\cdot\frac18=\frac1{32}$. So
        $16.3^{1/4}\approx2+\frac{0.3}{32}=2+0.009375=2.009375$.

  \item $R(Q)=p(Q)Q=80Q-0.4Q^2$. $MR(Q)=80-0.8Q$. $MR(50)=80-40=40$.

  \item $\dfrac{\partial Y}{\partial K}=\alpha AK^{\alpha-1}L^\beta$. So
        \[
          \varepsilon_{Y,K}=\frac{K}{Y}\cdot\alpha AK^{\alpha-1}L^\beta
            =\frac{\alpha AK^\alpha L^\beta}{AK^\alpha L^\beta}=\alpha .
        \]
        The Cobb--Douglas exponent equals the elasticity.

  \item $AC(Q)=2Q+8+\dfrac{200}{Q}$.
        $AC'(Q)=2-\dfrac{200}{Q^2}=0\Rightarrow Q^2=100\Rightarrow Q^*=10$
        (positive root). $AC''(Q)=\dfrac{400}{Q^3}>0$ at $Q=10$, so a
        minimum. $AC(10)=20+8+20=48$.

  \item $f'(x)=2e^{2x}$. $f(0)=1$, $f'(0)=2$. Tangent: $y=1+2x$.

  \item $f(1)=0$, $f'(1)=1$. $\ln(0.98)\approx0+1\cdot(-0.02)=-0.02$.
        (True: $-0.02020\ldots$)

  \item Set $MR(Q)=MC(Q)$: $80-0.8Q=0.4Q+5\Rightarrow75=1.2Q\Rightarrow
        Q^*=62.5$. SOC: $\pi''(Q)=R''(Q)-C''(Q)=-0.8-0.4=-1.2<0$, so this
        is a maximum.

  \item $q'(p)=80\cdot(-0.05)e^{-0.05p}=-4e^{-0.05p}$. So
        \[
          \varepsilon_d(p)=\frac{p}{80e^{-0.05p}}\cdot(-4e^{-0.05p})
            =-0.05p .
        \]
        At $p=20$: $\varepsilon_d=-1$, exactly unit elastic.

  \item $R(Q)=p(Q)Q=150Q-Q^2$.
        $\pi(Q)=150Q-Q^2-Q^2-30Q-100=120Q-2Q^2-100$. FOC:
        $\pi'(Q)=120-4Q=0\Rightarrow Q^*=30$. SOC: $\pi''(Q)=-4<0$.
        $p^*=150-30=120$. $\pi^*=120\cdot30-2\cdot900-100
        =3600-1800-100=1700$.

  \item $f'(x)=\dfrac{1}{2\sqrt{x}}$. $f(9)=3$, $f'(9)=1/6$. Normal slope
        $=-1/f'(9)=-6$. Normal: $y=3-6(x-9)=-6x+57$.

  \item $f(0)=1$, $f'(0)=1$. $e^{-0.2}\approx1+1\cdot(-0.2)=0.8$. (True:
        $0.8187\ldots$)

  \item $R(Q)=p(Q)Q=200Q-3Q^2$.
        $\pi(Q)=200Q-3Q^2-2Q^2-10Q-50=190Q-5Q^2-50$. FOC:
        $\pi'(Q)=190-10Q=0\Rightarrow Q^*=19$. SOC: $\pi''(Q)=-10<0$, max.
        $p^*=200-57=143$. $\pi^*=190\cdot19-5\cdot361-50
        =3610-1805-50=1755$.

  \item Using problem 4: $|\varepsilon_d|=1 \Leftrightarrow
        \dfrac{p}{40-p}=1\Rightarrow p=40-p\Rightarrow p=20$. (At $p=20$,
        demand $q=200-100=100$.)

  \item $\pi(Q)=24Q-Q^2-8Q-10=16Q-Q^2-10$. FOC:
        $\pi'(Q)=16-2Q=0\Rightarrow Q^*=8$. SOC: $\pi''(Q)=-2<0$.
        $\pi(8)=128-64-10=54$.

  \item $f'(x)=1/x$. $f(e)=1$, $f'(e)=1/e$. Normal slope $=-e$. Normal:
        $y=1-e(x-e)=1-ex+e^2$.

  \item $C'(Q)=Q+8$. $C(20)=200+160+200=560$, $C'(20)=28$. So
        $C(21.5)\approx560+28\cdot1.5=602$. (Exact:
        $0.5\cdot462.25+172+200=603.125$; error $\approx1.1$.)

  \item $MC'(Q)=0.12Q-2.4=0\Rightarrow Q=20$. $MC''(Q)=0.12>0$, so $Q=20$
        minimises $MC$. $MC(20)=0.06\cdot400-48+30=24-48+30=6$.

  \item $q(p)=300(p^2+4)^{-1}$,
        $q'(p)=-\dfrac{600p}{(p^2+4)^2}$. Then
        \[
          \varepsilon_d(p)=\frac{p}{q}\cdot q'(p)
            =\frac{p(p^2+4)}{300}\cdot\left(-\frac{600p}{(p^2+4)^2}\right)
            =\frac{-2p^2}{p^2+4} .
        \]
        At $p=2$: $\varepsilon_d=-8/8=-1$ (unit elastic).

  \item $U'(x)=4/x-0.5=0\Rightarrow x^*=8$. $U''(x)=-4/x^2<0$, max.
        $U(8)=4\ln(8)-4=12\ln(2)-4\approx4.318$.

  \item $f'(x)=-\sin(x)$. $f(\pi/3)=1/2$, $f'(\pi/3)=-\sqrt3/2$. Tangent:
        $y=\frac12-\frac{\sqrt3}{2}\left(x-\frac{\pi}{3}\right)$.

  \item $f(0)=1$, $f'(x)=10(1+x)^9$, $f'(0)=10$. So
        $(1.02)^{10}\approx1+10\cdot0.02=1.2$. (True: $1.21899\ldots$,
        error $\approx0.019$.)

  \item $MC(Q)=C'(Q)=3Q^2-18Q+30$. $MC'(Q)=6Q-18=0\Rightarrow Q=3$.
        $MC''(Q)=6>0$, min. $MC(3)=27-54+30=3$.

  \item $q'(p)=A\cdot(-\alpha)p^{-\alpha-1}=-\alpha\dfrac{q}{p}$. Hence
        $\varepsilon_d=\dfrac{p}{q}q'(p)=\dfrac{p}{q}\cdot
        \left(-\alpha\dfrac{q}{p}\right)=-\alpha$. Constant elasticity
        demand has elasticity equal to the negative exponent.

  \item $R(Q)=p(Q)Q=100Q-2Q^2$. $\pi(Q)=100Q-2Q^2-20Q-F=80Q-2Q^2-F$. FOC:
        $\pi'(Q)=80-4Q=0\Rightarrow Q^*=20$. SOC: $\pi''(Q)=-4<0$, max.
        $p^*=100-40=60$. Cross-check via $MR=MC$: $100-4Q=20\Rightarrow
        Q=20$. Consistent. (Profit $\pi^*=80\cdot20-2\cdot400-F
        =1600-800-F=800-F$.)

  \item $f'(x)=2x$. $f(1)=1$, $f'(1)=2$. Normal slope $=-1/2$. Normal:
        $y=1-\frac12(x-1)=-\frac{x}{2}+\frac32$.

  \item $f(0)=0$, $f'(x)=\sec^2(x)$, $f'(0)=1$. So
        $\tan(0.05)\approx0+1\cdot0.05=0.05$. (True: $0.050042\ldots$)

  \item $MC(Q)=C'(Q)=3Q^2-12Q+15$. $AC(Q)=Q^2-6Q+15+20/Q$.
        $AC'(Q)=2Q-6-20/Q^2$. Set $2Q-6-20/Q^2=0$, multiply by $Q^2$:
        $2Q^3-6Q^2-20=0\Rightarrow Q^3-3Q^2-10=0$. Try $Q=$ roots: at
        $Q=3$, $27-27-10=-10$; at $Q=4$, $64-48-10=6$; root between 3 and
        4. Newton step from $Q_0=3.7$: $Q_0^3-3Q_0^2-10
        =50.653-41.07-10=-0.417$, derivative $3Q^2-6Q=3\cdot13.69-22.2
        =18.87$, $Q_1=3.7+0.417/18.87\approx3.722$. So $Q^*\approx3.72$.
        SOC: $AC''(Q)=2+40/Q^3>0$ for $Q>0$, so it is a minimum. (Note: at
        the optimum $AC(Q^*)=MC(Q^*)$, a standard identity.)

  \item $q'(p)=-4p$. At $p=4$: $q(4)=100-32=68$, $q'(4)=-16$.
        $\varepsilon_d(4)=\dfrac{4}{68}\cdot(-16)=-\dfrac{64}{68}
        =-\dfrac{16}{17}\approx-0.94$. Since $|\varepsilon_d|<1$, demand
        is \emph{inelastic} at $p=4$.

  \item $R(Q)=120Q-2Q^2$. $R'(Q)=120-4Q=0\Rightarrow Q^*=30$.
        $R''(Q)=-4<0$, max. $R(30)=3600-1800=1800$.

  \item $f'(x)=2x-3$. $f(2)=4-6+2=0$, $f'(2)=1$. Tangent:
        $y=0+1\cdot(x-2)=x-2$. Normal slope $=-1$. Normal:
        $y=-(x-2)=-x+2$.

  \item $C'(Q)=8Q+20$. $\Delta C\approx C'(10)\cdot0.5
        =(80+20)\cdot0.5=100\cdot0.5=50$.

  \item $\pi(Q)=30Q-Q^2-10Q-9=20Q-Q^2-9$. FOC:
        $\pi'(Q)=20-2Q=0\Rightarrow Q^*=10$. SOC: $\pi''(Q)=-2<0$, max.
        $\pi(10)=200-100-9=91$. Break-even: $20Q-Q^2-9=0\Rightarrow
        Q^2-20Q+9=0\Rightarrow Q=\dfrac{20\pm\sqrt{400-36}}{2}
        =10\pm\sqrt{91}$. Both positive: $Q\approx10-9.54=0.46$ and
        $Q\approx19.54$.

  \item $q(p)=400-p^2$, $q'(p)=-2p$.
        $\varepsilon_d=\dfrac{p}{400-p^2}\cdot(-2p)=\dfrac{-2p^2}{400-p^2}$.
        Set $\dfrac{2p^2}{400-p^2}=1\Rightarrow2p^2=400-p^2
        \Rightarrow3p^2=400\Rightarrow p=\dfrac{20}{\sqrt3}\approx11.55$.

  \item $f'(x)=2xe^{-x}-x^2e^{-x}=x(2-x)e^{-x}$. On $(0,\infty)$:
        $f'>0$ for $0<x<2$, $f'<0$ for $x>2$. Critical point at $x=2$.
        $f''(x)=(2-2x)e^{-x}-(2x-x^2)e^{-x}=(x^2-4x+2)e^{-x}$. At $x=2$:
        $f''(2)=(4-8+2)e^{-2}=-2e^{-2}<0$, confirming local (and global on
        $(0,\infty)$) maximum. $f(2)=4e^{-2}\approx0.541$.

  \item $f'(x)=\dfrac{1}{1+x^2}$. $f(1)=\pi/4$, $f'(1)=1/2$. Tangent:
        $y=\dfrac{\pi}{4}+\dfrac12(x-1)$.

  \item $f(27)=3$, $f'(x)=\frac13x^{-2/3}$,
        $f'(27)=\frac13\cdot27^{-2/3}=\frac13\cdot\frac19=\frac1{27}$. So
        $\sqrt[3]{27.9}\approx3+\frac1{27}\cdot0.9=3+1/30\approx3.0333$.
        (True: $3.03296\ldots$)

  \item $U'(x)=\dfrac{1}{2\sqrt{x}}>0$ for $x>0$ (monotone preferences);
        at $x=100$: $MU=\dfrac{1}{20}=0.05>0$. Elasticity:
        $\varepsilon_{U,x}=\dfrac{x}{U(x)}U'(x)=\dfrac{x}{\sqrt{x}}
        \cdot\dfrac{1}{2\sqrt{x}}=\dfrac12$. Constant elasticity $1/2$,
        positive (utility rises with $x$). Marginal utility itself is
        positive but its elasticity is positive and less than 1.

  \item Invert demand: $q=120-3p\Rightarrow p(Q)=40-Q/3$.
        $R(Q)=Qp=40Q-Q^2/3$. $\pi(Q)=40Q-Q^2/3-10Q=30Q-Q^2/3$. FOC:
        $\pi'(Q)=30-2Q/3=0\Rightarrow Q^*=45$. SOC: $\pi''(Q)=-2/3<0$,
        max. $p^*=40-15=25$, $\pi^*=30\cdot45-45^2/3=1350-675=675$.

  \item $f'(x)=1-16/x^2=0\Rightarrow x^2=16\Rightarrow x=4$ on
        $(0,\infty)$. $f''(x)=32/x^3>0$ for $x>0$, so $x=4$ is a minimum.
        $f(4)=4+4=8$.

  \item $f'(x)=\cos(x)$. $f(\pi/2)=1$, $f'(\pi/2)=0$. The tangent is
        horizontal: $y=1$. The normal is the vertical line $x=\pi/2$
        (slope $-1/0$ undefined; normal is vertical).

  \item At $Q=20$: $p=80$. $R(Q)=100Q-Q^2$, $MR(Q)=100-2Q$, $MR(20)=60$.
        From $q=100-p$, $q'(p)=-1$,
        $\varepsilon_d(p)=\dfrac{p}{q}\cdot(-1)=-\dfrac{p}{100-p}$. At
        $p=80$: $\varepsilon_d=-80/20=-4$. Then
        $p\left(1+\dfrac{1}{\varepsilon_d}\right)=80\left(1-\dfrac14\right)
        =80\cdot\dfrac34=60=MR$. Verified.

  \item $f(0)=0$, $f'(x)=1/(1+x)$, $f'(0)=1$. So
        $\ln(1.05)\approx0+1\cdot0.05=0.05$. True value:
        $\ln(1.05)=0.048790\ldots$ Linearisation error
        $\approx0.05-0.04879=0.00121$.

  \item $\pi(Q)=R(Q)-C(Q)=(20Q-Q^2)-(2Q+\frac12Q^2+10)
        =18Q-\frac32Q^2-10$. FOC: $\pi'(Q)=18-3Q=0\Rightarrow Q=6$.
        $Q=6\in[0,8]$. SOC: $\pi''(Q)=-3<0$. $\pi(6)=108-54-10=44$. Check
        endpoints: $\pi(0)=-10$, $\pi(8)=144-96-10=38$. So $Q^*=6$
        maximises profit; $\pi^*=44$.

  \item $q(p)=50/p$, $q'(p)=-50/p^2$.
        $\varepsilon_d=\dfrac{p}{q}q'(p)=\dfrac{p}{50/p}\cdot(-50/p^2)
        =\dfrac{p^2}{50}\cdot(-50/p^2)=-1$ for all $p>0$. The elasticity
        is identically $-1$, unit-elastic demand. Total revenue
        $R(p)=p\cdot50/p=50$, constant, independent of $p$. This is the
        hallmark of unit-elastic demand.

  \item $\pi'(Q)=12-2Q-Q^2$. Solve $Q^2+2Q-12=0\Rightarrow
        Q=\dfrac{-2+\sqrt{4+48}}{2}=\dfrac{-2+\sqrt{52}}{2}
        =-1+\sqrt{13}$. So $Q^*=\sqrt{13}-1\approx2.606$ (the positive
        root). SOC: $\pi''(Q)=-2-2Q$. At $Q^*$:
        $\pi''(Q^*)=-2-2(\sqrt{13}-1)=-2\sqrt{13}<0$, confirming maximum.

  \item $R(Q)=(50-Q)Q=50Q-Q^2$. $\pi(Q)=50Q-Q^2-10Q-Q^2=40Q-2Q^2$. FOC:
        $\pi'(Q)=40-4Q=0\Rightarrow Q^*=10$. SOC: $\pi''(Q)=-4<0$, max.
        $p^*=50-10=40$, $\pi^*=40\cdot10-2\cdot100=400-200=200$.
        Elasticity: $q(p)=50-p$, $q'(p)=-1$,
        $\varepsilon_d(p)=\dfrac{p}{50-p}\cdot(-1)$. At $p^*=40$:
        $\varepsilon_d=-40/10=-4$ (elastic, as expected at a monopolist's
        optimum where $|\varepsilon_d|>1$).

  \item $\pi(L)=8\sqrt{L}-2L$.
        $\pi'(L)=\dfrac{4}{\sqrt{L}}-2=0\Rightarrow\sqrt{L}=2\Rightarrow
        L^*=4$. SOC: $\pi''(L)=-\dfrac{2}{L^{3/2}}<0$ for $L>0$, so
        $L^*=4$ maximises profit. $Y(4)=2$,
        $\pi^*=8\cdot2-2\cdot4=16-8=8$.
\end{solutions}

\section{Additional Problems}

\begin{problems}
  \item Find the equation of the tangent to $f(x)=x^2+3x-1$ at $x=2$.
  \item Use the linearisation of $f(x)=\sqrt[3]{x}$ at $a=8$ to estimate
        $\sqrt[3]{8.6}$.
  \item For $C(Q)=0.05Q^2+8Q+150$, find $MC(Q)$ and $MC(20)$.
  \item For demand $q(p)=300-4p$, compute $\varepsilon_d(p)$ at $p=25$ and
        classify.
  \item Maximise profit $\pi(Q)=R(Q)-C(Q)$ where $R(Q)=90Q-1.5Q^2$ and
        $C(Q)=Q^2+20Q+30$.
  \item Find the tangent and normal to $f(x)=x^3-4x$ at $x=1$.
  \item Use the linearisation of $f(x)=\mathrm{e}^{x}$ at $a=0$ to estimate
        $\mathrm{e}^{0.1}$.
  \item Inverse demand $p(Q)=100-0.5Q$. Find $R(Q)$, $MR(Q)$, and $MR(40)$.
  \item For the Cobb--Douglas production $Y=AK^{\alpha}L^{1-\alpha}$ with
        $K$ fixed, compute the output elasticity $\varepsilon_{Y,L}$.
  \item Minimise the average cost $AC(Q)=C(Q)/Q$ for
        $C(Q)=3Q^2+12Q+108$, over $Q>0$.
  \item A monopolist faces $p(Q)=80-2Q$ and $C(Q)=Q^2+8Q+40$. Find $Q^*$,
        $p^*$, and $\pi^*$.
  \item For demand $q(p)=120\mathrm{e}^{-0.04p}$, compute
        $\varepsilon_d(p)$ and at $p=10$.
  \item For utility $U(x)=2\sqrt{x}-0.1x$ for $x>0$, find the maximiser
        and verify SOC.
  \item Estimate $\ln(1.04)$ using the linearisation of $f(x)=\ln(1+x)$
        at $a=0$.
  \item Minimise $f(x)=2x+\dfrac{50}{x}$ on $(0,\infty)$ and verify SOC.
\end{problems}

\section{Additional Solutions}

\begin{solutions}
  \item $f'(x)=2x+3$. $f(2)=9$, $f'(2)=7$. Tangent: $y=9+7(x-2)=7x-5$.

  \item $f(8)=2$, $f'(x)=\tfrac{1}{3}x^{-2/3}$,
        $f'(8)=\tfrac{1}{3}\cdot 8^{-2/3}=\tfrac{1}{3}\cdot\tfrac{1}{4}
        =\tfrac{1}{12}$. So
        $\sqrt[3]{8.6}\approx 2+\tfrac{0.6}{12}=2+0.05=2.05$. (True:
        $2.0494\ldots$)

  \item $MC(Q)=C'(Q)=0.1Q+8$. $MC(20)=2+8=10$.

  \item $q'(p)=-4$.
        $\varepsilon_d(p)=\dfrac{p}{300-4p}\cdot(-4)=\dfrac{-4p}{300-4p}$.
        At $p=25$: $\varepsilon_d=\dfrac{-100}{200}=-\dfrac{1}{2}$.
        $|\varepsilon_d|<1$: \emph{inelastic}.

  \item $\pi(Q)=90Q-1.5Q^2-Q^2-20Q-30=70Q-2.5Q^2-30$. FOC:
        $\pi'(Q)=70-5Q=0\Rightarrow Q^*=14$. SOC: $\pi''(Q)=-5<0$, max.
        $\pi(14)=980-490-30=460$.

  \item $f'(x)=3x^2-4$. $f(1)=-3$, $f'(1)=-1$. Tangent:
        $y=-3-(x-1)=-x-2$. Normal slope $=-1/(-1)=1$. Normal:
        $y=-3+1\cdot(x-1)=x-4$.

  \item $f(0)=1$, $f'(0)=1$. $\mathrm{e}^{0.1}\approx 1+0.1=1.1$. (True:
        $1.10517\ldots$)

  \item $R(Q)=p(Q)Q=100Q-0.5Q^2$. $MR(Q)=100-Q$. $MR(40)=60$.

  \item $\dfrac{\partial Y}{\partial L}=(1-\alpha)AK^{\alpha}L^{-\alpha}$.
        So
        \[
          \varepsilon_{Y,L}=\frac{L}{Y}\cdot(1-\alpha)AK^{\alpha}L^{-\alpha}
            =\frac{(1-\alpha)AK^{\alpha}L^{1-\alpha}}{AK^{\alpha}L^{1-\alpha}}
            =1-\alpha .
        \]
        Standard Cobb--Douglas: elasticity equals the exponent.

  \item $AC(Q)=3Q+12+\dfrac{108}{Q}$.
        $AC'(Q)=3-\dfrac{108}{Q^2}=0\Rightarrow Q^2=36\Rightarrow Q^*=6$.
        $AC''(Q)=\dfrac{216}{Q^3}>0$ at $Q=6$: minimum.
        $AC(6)=18+12+18=48$.

  \item $R(Q)=p(Q)Q=80Q-2Q^2$.
        $\pi(Q)=80Q-2Q^2-Q^2-8Q-40=72Q-3Q^2-40$. FOC:
        $\pi'(Q)=72-6Q=0\Rightarrow Q^*=12$. SOC: $\pi''(Q)=-6<0$.
        $p^*=80-24=56$. $\pi^*=72\cdot12-3\cdot144-40=864-432-40=392$.

  \item $q'(p)=120\cdot(-0.04)\mathrm{e}^{-0.04p}
        =-4.8\mathrm{e}^{-0.04p}$. So
        \[
          \varepsilon_d(p)=\frac{p}{120\mathrm{e}^{-0.04p}}
            \cdot(-4.8\mathrm{e}^{-0.04p})=-0.04p .
        \]
        At $p=10$: $\varepsilon_d=-0.4$. $|\varepsilon_d|<1$:
        \emph{inelastic}.

  \item $U'(x)=\tfrac{1}{\sqrt{x}}-0.1=0\Rightarrow\sqrt{x}=10
        \Rightarrow x^*=100$. $U''(x)=-\tfrac{1}{2}x^{-3/2}<0$ for $x>0$,
        max. $U(100)=20-10=10$.

  \item $f(0)=0$, $f'(0)=1$. $\ln(1.04)\approx 0+1\cdot0.04=0.04$.
        (True: $0.03922\ldots$)

  \item $f'(x)=2-\dfrac{50}{x^2}=0\Rightarrow x^2=25\Rightarrow x=5$ on
        $(0,\infty)$. $f''(x)=\dfrac{100}{x^3}>0$ for $x>0$, so $x=5$ is
        a minimum. $f(5)=10+10=20$.
\end{solutions}

\chapter{First-Order Partial Derivatives and the Gradient}

When a function depends on more than one input, we measure how it responds
to a change in \emph{one} variable at a time, holding the others fixed,
the ceteris paribus principle. This is exactly what a partial derivative
computes. From these one-variable slopes we assemble the gradient, which
packages all first-order information into a single vector and unlocks the
geometry of multivariate optimisation.

\section{Core Rules}

\begin{rulebox}{Topological vocabulary of domains}
\begin{itemize}[leftmargin=1.6em,itemsep=0.3ex]
  \item The \emph{open ball} (or \emph{neighbourhood}) of radius $r>0$ around $\mathbf{x}_0\in\R^n$ is
        $B_r(\mathbf{x}_0)=\{\mathbf{x}\in\R^n :
        \|\mathbf{x}-\mathbf{x}_0\|<r\}$.
  \item $\mathbf{x}_0\in M$ is an \emph{interior point} of $M$ if some open ball around
        $\mathbf{x}_0$ lies entirely in $M$.
  \item $M$ is \emph{open} if \emph{every} point of $M$ is
        an interior point.
  \item $\mathbf{x}_0$ is a \emph{boundary point} of
        $M$ if every ball around $\mathbf{x}_0$ contains points of $M$
        \emph{and} points of its complement.
  \item $M$ is \emph{closed} if it contains all
        its boundary points; equivalently, the complement
        $\R^n\setminus M$ is open.
  \item $M$ is \emph{bounded} if
        $M\subseteq B_r(\mathbf{0})$ for some radius $r$.
  \item $M\subseteq\R^n$ is \emph{compact} exactly when
        $M$ is closed \emph{and} bounded.
\end{itemize}
Economic examples: the open positive orthant $\{(x,y): x>0,\ y>0\}$
(the natural Cobb--Douglas domain) is open and unbounded; a budget set
$\{\mathbf{x}\in\R^n : \mathbf{p}\cdot\mathbf{x}\le m,\
\mathbf{x}\ge\mathbf{0}\}$ is closed and bounded, hence compact.
\end{rulebox}

\begin{rulebox}{Why it matters}
Differentiability (and hence every partial derivative below) is
defined at \emph{interior} points of the domain, where $f$ can be
probed in all directions. And by the \emph{Weierstrass theorem}, a
continuous function on a compact set attains its maximum and minimum:
this is why budget-constrained optimisation problems are
guaranteed to have solutions.
\end{rulebox}

\begin{rulebox}{Definition}
Let $f:\R^n\to\R$. The partial derivative of $f$ with respect to $x_i$ at
$\mathbf{x}=(x_1,\ldots,x_n)$ is
\[
  \pd{f}{x_i}(\mathbf{x}) \;=\; \lim_{h\to0}
    \frac{f(x_1,\ldots,x_i+h,\ldots,x_n)-f(x_1,\ldots,x_n)}{h} .
\]
\end{rulebox}

\begin{rulebox}{Practical computation}
To compute $\partial f/\partial x_i$, differentiate $f$ with respect to
$x_i$ and treat every other variable as a constant. All single-variable
rules (power, product, quotient, chain) apply unchanged.
\end{rulebox}

\begin{rulebox}{Notation}
The following all refer to the same object:
\[
  \pd{f}{x} = f_x = \partial_x f = D_x f .
\]
\end{rulebox}

\begin{rulebox}{Gradient}
The vector of all first-order partial derivatives is
\[
  \grad f(\mathbf{x}) \;=\;
  \begin{pmatrix} f_{x_1}(\mathbf{x}) \\ f_{x_2}(\mathbf{x}) \\ \vdots \\
    f_{x_n}(\mathbf{x}) \end{pmatrix} .
\]
\end{rulebox}

\begin{rulebox}{Geometric meaning}
\begin{itemize}[leftmargin=1.6em,itemsep=0.3ex]
  \item $\grad f(\mathbf{x}_0)$ points in the direction of \emph{steepest
        ascent} at $\mathbf{x}_0$.
  \item $\|\grad f(\mathbf{x}_0)\|$ is the maximal rate of increase at that
        point.
  \item $\grad f(\mathbf{x}_0)$ is \emph{perpendicular} to the level set
        $\{\mathbf{x} : f(\mathbf{x})=f(\mathbf{x}_0)\}$ at $\mathbf{x}_0$.
\end{itemize}
\end{rulebox}

\begin{rulebox}{Economic interpretation}
\begin{itemize}[leftmargin=1.6em,itemsep=0.3ex]
  \item For a production function $Q=F(K,L)$: $F_K=\mathrm{MP}_K$
        (marginal product of capital), $F_L=\mathrm{MP}_L$ (marginal
        product of labour).
  \item For a utility function $U(x,y)$: $U_x, U_y$ are the marginal
        utilities; the marginal rate of substitution is
        $\mathrm{MRS}_{xy}=U_x/U_y$.
  \item For production: the marginal rate of technical substitution is
        $\mathrm{MRTS}_{LK}=F_L/F_K$.
\end{itemize}
\end{rulebox}

\begin{rulebox}{Directional derivative}
Let $\mathbf{u}\in\R^n$ be a \emph{unit} vector ($\|\mathbf{u}\|=1$). The
directional derivative of $f$ at $\mathbf{x}_0$ in the direction
$\mathbf{u}$ is
\[
  D_{\mathbf{u}}f(\mathbf{x}_0) \;=\; \grad f(\mathbf{x}_0)\cdot\mathbf{u} .
\]
\end{rulebox}

\begin{rulebox}{Caution}
If you are given a non-unit direction $\mathbf{v}$, you \emph{must}
normalise first: $\mathbf{u}=\mathbf{v}/\|\mathbf{v}\|$. Forgetting this is
the single most common error in directional-derivative problems.
\end{rulebox}

\section{Practice Problems}

\begin{problems}
  \item Compute $f_x$ and $f_y$ for $f(x,y)=5x^4-2x^2y^3+y^5$.
  \item For $F(K,L)=12\,K^{1/4}L^{3/4}$, find $F_K$ and $F_L$ at
        $(K,L)=(16,81)$.
  \item Compute $f_x$ and $f_y$ for $f(x,y)=x^2e^{3y}$.
  \item Compute $\nabla f$ for $f(x,y)=4x^3y-xy^4+7$.
  \item Compute the directional derivative of $f(x,y)=x^2-4xy+y^3$ at
        $(2,1)$ in the direction $\mathbf{v}=(5,-12)$.
  \item For $F(K,L)=6K^{2/3}L^{1/3}$, compute $\mathrm{MP}_K$ and
        $\mathrm{MP}_L$ at $(K,L)=(27,8)$.
  \item Compute $f_x$ and $f_y$ for $f(x,y)=(3x-2y)^4$.
  \item Compute $f_x$ and $f_y$ for $f(x,y)=e^{xy^2}$.
  \item For $F(K,L)=25K^{0.6}L^{0.5}$, find $F_K$ and $F_L$.
  \item Compute $\nabla f$ for $f(x,y,z)=x^2+3yz-z^2e^x$.
  \item Compute the directional derivative of $f(x,y)=e^x\ln(1+y)$ at
        $(0,0)$ in the direction $\mathbf{v}=(1,2)$.
  \item For $U(x,y)=2\ln(x)+5\ln(y)$, derive $\mathrm{MRS}_{xy}=U_x/U_y$ as
        a function of $(x,y)$ and evaluate at $(3,5)$.
  \item Simplify and compute $f_x, f_y$ for
        $f(x,y)=\dfrac{x^2-y^2}{x+y}$.
  \item Compute $f_x$ and $f_y$ for $f(x,y)=\ln(x^2y+1)$.
  \item For $U(x_1,x_2,x_3)=x_1^{1/2}x_2^{1/3}x_3^{1/6}$, compute the
        three partials.
  \item Compute the directional derivative of $f(x,y,z)=xy-2yz+z^2$ at
        $(1,2,3)$ in the direction $\mathbf{v}=(1,2,2)$.
  \item For $F(K,L)=K^{0.25}L^{0.75}$, derive $\mathrm{MRTS}_{LK}=F_L/F_K$
        and evaluate at $(K,L)=(16,4)$.
  \item Expand and then compute $f_x, f_y$ for
        $f(x,y)=(x^2+y)(x-3y^2)$.
  \item Compute $f_x$ and $f_y$ for $f(x,y)=xe^{-y/x}$ with $x\neq0$.
  \item For $F(K,L)=(3K^{1/2}+2L^{1/2})^2$, find $F_K$ and $F_L$.
  \item Compute $\nabla f$ for $f(x,y)=\arctan(y/x)$ with $x>0$.
  \item For $f(x,y)=x^{1/2}y^{1/2}$ at $(K,L)=(4,9)$, find the unit
        direction of steepest ascent and the maximum rate of increase.
  \item For the CES function $F(K,L)=(K^{1/2}+L^{1/2})^2$, verify Euler's
        theorem $KF_K+LF_L=rF$ for the appropriate degree of homogeneity
        $r$.
  \item Compute $f_x$ and $f_y$ for $f(x,y)=\dfrac{2x+3y}{x^2+y^2}$.
  \item Compute $f_x$ and $f_y$ for $f(x,y)=\dfrac{e^x}{1+e^y}$.
  \item For $F(K,L)=A(\alpha K^\rho+(1-\alpha)L^\rho)^{\nu/\rho}$ (CES),
        find $F_L$.
  \item Compute the directional derivative of $f(x,y)=\ln(x^2+y^2)$ at
        $(3,4)$ in the direction $\mathbf{v}=(-1,1)$.
  \item Compute $f_x, f_y, f_z$ for $f(x,y,z)=x^2yz^3-4xy^2z+5z$.
  \item Compute $f_x$ and $f_y$ for $f(x,y)=\sqrt{4x^2+9y^2}$.
  \item For $F(K,L)=(K^{1/2}+L^{1/2})^4$, compute $F_K$ at
        $(K,L)=(4,9)$.
  \item Compute $\nabla f$ for $f(x,y,z)=\ln(x^2+y^2+z^2)$.
  \item Compute $f_x$ and $f_y$ for
        $f(x,y)=\ln\left(\sqrt{x^2+y^2}\right)$.
  \item Cost minimisation. A firm uses inputs $(x,y)$ with cost
        $C(x,y)=4x+9y$ and output $Q(x,y)=x^{1/2}y^{1/2}$. Compute
        $\nabla C$ and $\nabla Q$, and verify that at the cost-minimising
        input mix $\nabla C=\lambda\nabla Q$ for some $\lambda$.
  \item Compute $f_x$ and $f_y$ for $f(x,y)=(1+xy)^5$.
  \item Compute the directional derivative of $f(x,y)=x^3-2xy+y^2$ at
        $(1,2)$ in the direction $\mathbf{v}=(3,4)$.
  \item For $F(K,L)=10K^{0.3}L^{0.7}$ (Cobb--Douglas), compute $F_K$ and
        $F_L$ at $(K,L)=(100,100)$ and verify $KF_K+LF_L=F$.
  \item Compute $f_x$ and $f_y$ for $f(x,y)=\sin(xy)+\cos(x-y)$.
  \item Compute $\nabla f$ for $f(x,y)=x^2y-xy^2+3x-2y$ at the point
        $(1,-1)$.
  \item For $U(x,y)=x^\alpha y^{1-\alpha}$ (Cobb--Douglas,
        $0<\alpha<1$), derive $\mathrm{MRS}_{xy}=U_x/U_y$ and show it
        equals $\dfrac{\alpha}{1-\alpha}\cdot\dfrac{y}{x}$.
  \item Compute the directional derivative of $f(x,y)=e^{xy}$ at $(1,0)$
        in the direction $\mathbf{v}=(-3,4)$.
  \item Compute $f_x$ and $f_y$ for $f(x,y)=x^y$ (with $x>0$).
        \emph{Hint:} write $x^y=e^{y\ln x}$.
  \item For $F(K,L)=\min\{2K,3L\}$ (Leontief, fixed-proportions, on the
        smooth region away from the kink), compute $F_K$ and $F_L$ in
        each of the two regions $\{2K<3L\}$ and $\{2K>3L\}$.
  \item Compute $\nabla f$ for $f(x,y,z)=xyz+x^2-z^2$ at the point
        $(2,-1,3)$.
  \item Compute the unit vector of steepest \emph{descent} for
        $f(x,y)=x^2+4y^2$ at $(3,1)$, and the corresponding rate of
        decrease.
  \item For $f(x,y)=xe^y-ye^x$, compute $f_x, f_y$.
  \item Compute $f_x$ and $f_y$ for $f(x,y)=\dfrac{x}{x^2+y^2}$.
  \item For $U(x,y)=x+2\sqrt{y}$ (quasi-linear), derive
        $\mathrm{MRS}_{xy}$.
  \item Compute the directional derivative of $f(x,y,z)=x^2+y^2-z^2$ at
        $(2,-1,2)$ in the direction $\mathbf{v}=(2,2,1)$.
  \item Compute $f_x$ and $f_y$ for $f(x,y)=\tan(x+y)+\tan(x-y)$.
  \item For $F(K,L)=(\alpha K+\beta L)^r$ (linear-aggregate,
        $\alpha,\beta,r>0$), compute $F_K, F_L$, and verify
        $KF_K+LF_L=rF$.
  \item Compute $\nabla f$ for $f(x,y)=e^{-x}\sin(y)+e^x\cos(y)$.
  \item Compute the directional derivative of $f(x,y)=\sqrt{x^2+y^2}$ at
        $(3,4)$ in the direction $\mathbf{v}=(1,0)$.
  \item For $U(x,y)=\min\{2x,3y\}$, compute the marginal utilities on
        each smooth region $\{2x<3y\}$ and $\{2x>3y\}$.
  \item Compute $f_x, f_y$ for $f(x,y)=\dfrac{\ln(x+y)}{x-y}$ (assume
        $x>y>0$).
  \item For $F(K,L)=AK^\alpha L^\beta$ general Cobb--Douglas, derive
        $\mathrm{MRTS}_{LK}$ and show it depends only on the ratio $K/L$
        (not on scale).
  \item Compute the unit direction of steepest ascent for
        $f(x,y)=x^2-y^2$ at $(2,-1)$, and the maximum rate.
  \item Compute $f_x$ and $f_y$ for $f(x,y)=\cosh(x+y)+\sinh(x-y)$.
  \item For $f(x,y,z)=e^{x+y+z}$, compute $\nabla f$ at $(0,0,0)$ and the
        directional derivative in $\mathbf{v}=(1,1,1)$.
  \item Compute $f_x$ and $f_y$ for $f(x,y)=(\ln x)(\ln y)$ with
        $x,y>1$.
  \item For $F(K,L)=K+2K^{1/2}L^{1/2}+L$ (which equals
        $(K^{1/2}+L^{1/2})^2$), compute $F_K, F_L$ directly and confirm
        agreement with the CES form.
  \item Compute the directional derivative of $f(x,y)=x^2y+y^3$ at
        $(-1,2)$ in the direction $\mathbf{v}=(4,-3)$.
  \item For $U(x,y)=\alpha\ln x+\beta\ln y$ (log-Cobb--Douglas), derive
        $\mathrm{MRS}_{xy}$ and evaluate at $(x,y)=(\alpha,\beta)$.
  \item Compute $\nabla f$ for $f(x,y)=(x^2+y^2)^{-1/2}$ (Newtonian
        potential in 2D).
  \item Compute $f_x$ and $f_y$ for $f(x,y)=\arctan(xy)$.
  \item For $F(K,L)=(K^\rho+L^\rho)^{1/\rho}$ with $\rho<1$, $\rho\neq0$
        (CES with constant returns), compute $F_K$ and show
        $F_K=(K/F)^{\rho-1}$.
  \item Compute the directional derivative of $f(x,y,z)=xyz$ at
        $(2,3,-1)$ in the direction $\mathbf{v}=(2,-1,2)$.
\end{problems}

\section{Solutions}

\begin{solutions}
  \item $f_x=20x^3-4xy^3$, $f_y=-6x^2y^2+5y^4$.

  \item $F_K=12\cdot\frac14K^{-3/4}L^{3/4}=3(L/K)^{3/4}$;
        $F_L=12\cdot\frac34K^{1/4}L^{-1/4}=9(K/L)^{1/4}$. At $(16,81)$:
        $(L/K)^{3/4}=(81/16)^{3/4}=27/8$ so $F_K=81/8$;
        $(K/L)^{1/4}=(16/81)^{1/4}=2/3$ so $F_L=6$.

  \item $f_x=2xe^{3y}$, $f_y=3x^2e^{3y}$.

  \item $\nabla f=(12x^2y-y^4,\ 4x^3-4xy^3)^T$.

  \item $f_x=2x-4y$, $f_y=-4x+3y^2$. At $(2,1)$: $\nabla f=(0,-5)^T$.
        Normalise: $\|\mathbf{v}\|=\sqrt{25+144}=13$,
        $\mathbf{u}=(5/13,-12/13)^T$. Then
        \[
          D_{\mathbf{u}}f=0\cdot\frac{5}{13}+(-5)\cdot
            \left(-\frac{12}{13}\right)=\frac{60}{13} .
        \]

  \item $F_K=6\cdot\frac23K^{-1/3}L^{1/3}=4(L/K)^{1/3}$,
        $F_L=2(K/L)^{2/3}$. At $(27,8)$: $(L/K)^{1/3}=2/3$ so
        $\mathrm{MP}_K=8/3$; $(K/L)^{2/3}=(27/8)^{2/3}=9/4$ so
        $\mathrm{MP}_L=9/2$.

  \item Chain rule with $u=3x-2y$: $f_x=4(3x-2y)^3\cdot3=12(3x-2y)^3$,
        $f_y=-8(3x-2y)^3$.

  \item $f_x=y^2e^{xy^2}$, $f_y=2xye^{xy^2}$.

  \item $F_K=15K^{-0.4}L^{0.5}$, $F_L=12.5K^{0.6}L^{-0.5}$.

  \item $f_x=2x-z^2e^x$, $f_y=3z$, $f_z=3y-2ze^x$. So
        $\nabla f=(2x-z^2e^x,\ 3z,\ 3y-2ze^x)^T$.

  \item $f_x=e^x\ln(1+y)$, $f_y=e^x/(1+y)$. At $(0,0)$:
        $\nabla f=(0,1)^T$. Normalise: $\|\mathbf{v}\|=\sqrt5$,
        $\mathbf{u}=(1/\sqrt5,2/\sqrt5)^T$. So
        \[
          D_{\mathbf{u}}f=0\cdot\frac{1}{\sqrt5}+1\cdot\frac{2}{\sqrt5}
            =\frac{2\sqrt5}{5} .
        \]

  \item $U_x=2/x$, $U_y=5/y$. $\mathrm{MRS}_{xy}=U_x/U_y
        =(2/x)/(5/y)=2y/(5x)$. At $(3,5)$: $\mathrm{MRS}=10/15=2/3$.

  \item For $x+y\neq0$: $f(x,y)=x-y$, so $f_x=1$, $f_y=-1$.

  \item With $u=x^2y+1$: $f_x=2xy/(x^2y+1)$, $f_y=x^2/(x^2y+1)$.

  \item $U_{x_1}=\frac12x_1^{-1/2}x_2^{1/3}x_3^{1/6}$,
        $U_{x_2}=\frac13x_1^{1/2}x_2^{-2/3}x_3^{1/6}$,
        $U_{x_3}=\frac16x_1^{1/2}x_2^{1/3}x_3^{-5/6}$. Equivalently
        $U_{x_i}=(\alpha_i/x_i)U$ with $\alpha=(1/2,1/3,1/6)$.

  \item $f_x=y$, $f_y=x-2z$, $f_z=-2y+2z$. At $(1,2,3)$:
        $\nabla f=(2,-5,2)^T$. Normalise: $\|\mathbf{v}\|=\sqrt{1+4+4}=3$,
        $\mathbf{u}=(1/3,2/3,2/3)^T$. Then
        $D_{\mathbf{u}}f=(2-10+4)/3=-4/3$.

  \item $F_K=0.25K^{-0.75}L^{0.75}$, $F_L=0.75K^{0.25}L^{-0.25}$. So
        $\mathrm{MRTS}_{LK}=F_L/F_K=3K/L$. At $(16,4)$:
        $\mathrm{MRTS}=48/4=12$.

  \item Expand: $f=x^3-3x^2y^2+xy-3y^3$. Hence $f_x=3x^2-6xy^2+y$,
        $f_y=-6x^2y+x-9y^2$.

  \item With $u=-y/x$: $\partial u/\partial x=y/x^2$,
        $\partial u/\partial y=-1/x$. Product rule:
        \[
          f_x=e^{-y/x}+x\cdot e^{-y/x}\cdot\frac{y}{x^2}
            =e^{-y/x}\left(1+\frac{y}{x}\right), \qquad
          f_y=x\cdot e^{-y/x}\cdot\left(-\frac1x\right)=-e^{-y/x} .
        \]

  \item Let $g=3K^{1/2}+2L^{1/2}$, so $F=g^2$. Then
        \[
          F_K=2g\cdot\frac32K^{-1/2}=\frac{3g}{K^{1/2}}
            =\frac{3(3K^{1/2}+2L^{1/2})}{K^{1/2}}, \qquad
          F_L=2g\cdot L^{-1/2}=\frac{2(3K^{1/2}+2L^{1/2})}{L^{1/2}} .
        \]

  \item With $u=y/x$, $f=\arctan u$, $f'=1/(1+u^2)$:
        \[
          f_x=\frac{1}{1+y^2/x^2}\cdot\left(-\frac{y}{x^2}\right)
            =\frac{-y}{x^2+y^2}, \qquad
          f_y=\frac{x}{x^2+y^2} .
        \]
        So $\nabla f=(-y/(x^2+y^2),\ x/(x^2+y^2))^T$.

  \item $f_K=\frac12K^{-1/2}L^{1/2}$, $f_L=\frac12K^{1/2}L^{-1/2}$. At
        $(4,9)$: $\nabla f=(\frac12\cdot\frac32,\ \frac12\cdot\frac23)^T
        =(3/4,1/3)^T$. Then $\|\nabla f\|^2=9/16+1/9=97/144$, so the
        maximum rate is $\sqrt{97}/12\approx0.821$ in direction
        $\mathbf{u}^*=(9,4)^T/\sqrt{97}$.

  \item $F=(K^{1/2}+L^{1/2})^2$. Replacing $(K,L)\to(\lambda K,\lambda L)$
        gives $\lambda F$, so $r=1$. With $g=K^{1/2}+L^{1/2}$:
        $F_K=2g\cdot\frac12K^{-1/2}=g/K^{1/2}$, $F_L=g/L^{1/2}$. Hence
        \[
          KF_K+LF_L=K^{1/2}g+L^{1/2}g=g(K^{1/2}+L^{1/2})=g^2=F. \checkmark
        \]

  \item Quotient rule, $D=x^2+y^2$:
        \[
          f_x=\frac{2D-(2x+3y)(2x)}{D^2}=\frac{-2x^2-6xy+2y^2}{(x^2+y^2)^2},
          \qquad f_y=\frac{3x^2-4xy-3y^2}{(x^2+y^2)^2} .
        \]

  \item Treat denominator as constant for $f_x$: $f_x=e^x/(1+e^y)$. For
        $f_y$:
        \[
          f_y=e^x\cdot\frac{-e^y}{(1+e^y)^2}=-\frac{e^{x+y}}{(1+e^y)^2} .
        \]

  \item Let $S=\alpha K^\rho+(1-\alpha)L^\rho$, $F=AS^{\nu/\rho}$.
        Treating $K$ constant, $\partial S/\partial L=\rho(1-\alpha)L^{\rho-1}$:
        \[
          F_L=A\cdot\frac{\nu}{\rho}S^{\nu/\rho-1}\cdot\rho(1-\alpha)L^{\rho-1}
            =\nu(1-\alpha)AL^{\rho-1}S^{\nu/\rho-1} .
        \]

  \item $f=\ln(x^2+y^2)$, $f_x=2x/(x^2+y^2)$, $f_y=2y/(x^2+y^2)$. At
        $(3,4)$: $x^2+y^2=25$, $\nabla f=(6/25,8/25)^T$. Normalise:
        $\|\mathbf{v}\|=\sqrt2$, $\mathbf{u}=(-1/\sqrt2,1/\sqrt2)^T$. Then
        \[
          D_{\mathbf{u}}f=\frac{6}{25}\cdot\left(-\frac{1}{\sqrt2}\right)
            +\frac{8}{25}\cdot\frac{1}{\sqrt2}=\frac{2}{25\sqrt2}
            =\frac{\sqrt2}{25} .
        \]

  \item $f_x=2xyz^3-4y^2z$, $f_y=x^2z^3-8xyz$,
        $f_z=3x^2yz^2-4xy^2+5$.

  \item Chain rule, $u=4x^2+9y^2$:
        \[
          f_x=\frac{8x}{2\sqrt u}=\frac{4x}{\sqrt{4x^2+9y^2}}, \qquad
          f_y=\frac{9y}{\sqrt{4x^2+9y^2}} .
        \]

  \item With $g=K^{1/2}+L^{1/2}$, $F=g^4$,
        $F_K=4g^3\cdot\frac12K^{-1/2}=2g^3/K^{1/2}$. At $(4,9)$: $g=5$,
        $K^{1/2}=2$, so $F_K=2\cdot125/2=125$.

  \item Let $r^2=x^2+y^2+z^2$. Then $f=\ln(r^2)$ and $f_{x_i}=2x_i/r^2$.
        So $\nabla f=(2/(x^2+y^2+z^2))(x,y,z)^T$.

  \item Use $\ln\sqrt{u}=\frac12\ln u$: $f=\frac12\ln(x^2+y^2)$. So
        $f_x=x/(x^2+y^2)$, $f_y=y/(x^2+y^2)$.

  \item $\nabla C=(4,9)^T$. With $Q=x^{1/2}y^{1/2}$: $Q_x=\frac12(y/x)^{1/2}$,
        $Q_y=\frac12(x/y)^{1/2}$, so
        $\nabla Q=\frac12\left((y/x)^{1/2},(x/y)^{1/2}\right)^T$. The
        condition $\nabla C=\lambda\nabla Q$ requires
        $4/9=Q_x/Q_y=y/x$, i.e.\ $y=4x/9$. Along this locus
        $\nabla Q\propto(2/3,3/2)^T\propto(4,9)^T=\nabla C$. \checkmark

  \item Chain rule: $f_x=5(1+xy)^4\cdot y=5y(1+xy)^4$,
        $f_y=5x(1+xy)^4$.

  \item $f_x=3x^2-2y$, $f_y=-2x+2y$. At $(1,2)$:
        $\nabla f=(3-4,-2+4)^T=(-1,2)^T$. Normalise: $\|\mathbf{v}\|=5$,
        $\mathbf{u}=(3/5,4/5)^T$. Then
        \[
          D_{\mathbf{u}}f=-1\cdot\frac35+2\cdot\frac45
            =\frac{-3+8}{5}=1 .
        \]

  \item $F_K=10\cdot0.3K^{-0.7}L^{0.7}=3(L/K)^{0.7}$;
        $F_L=7(K/L)^{0.3}$. At $K=L=100$: $F_K=3$, $F_L=7$. And
        $F=10\cdot100^{0.3}\cdot100^{0.7}=10\cdot100=1000$. Then
        $KF_K+LF_L=100\cdot3+100\cdot7=1000=F$. \checkmark (homogeneity
        degree $r=0.3+0.7=1$).

  \item $f_x=y\cos(xy)-\sin(x-y)$, $f_y=x\cos(xy)+\sin(x-y)$.

  \item $f_x=2xy-y^2+3$, $f_y=x^2-2xy-2$. At $(1,-1)$: $f_x=-2-1+3=0$,
        $f_y=1+2-2=1$. So $\nabla f(1,-1)=(0,1)^T$.

  \item $U_x=\alpha x^{\alpha-1}y^{1-\alpha}$,
        $U_y=(1-\alpha)x^\alpha y^{-\alpha}$. So
        \[
          \mathrm{MRS}_{xy}=\frac{U_x}{U_y}
            =\frac{\alpha x^{\alpha-1}y^{1-\alpha}}{(1-\alpha)x^\alpha y^{-\alpha}}
            =\frac{\alpha}{1-\alpha}\cdot\frac{y}{x} . \checkmark
        \]

  \item $f_x=ye^{xy}$, $f_y=xe^{xy}$. At $(1,0)$: $\nabla f=(0,1)^T$.
        Normalise: $\|\mathbf{v}\|=5$, $\mathbf{u}=(-3/5,4/5)^T$. Then
        $D_{\mathbf{u}}f=0\cdot(-3/5)+1\cdot(4/5)=4/5$.

  \item Write $f=e^{y\ln x}$. Then $f_x=e^{y\ln x}\cdot y/x=yx^{y-1}$ and
        $f_y=e^{y\ln x}\cdot\ln x=x^y\ln x$.

  \item On $\{2K<3L\}$ the binding factor is capital: $F=2K$, so
        $F_K=2$, $F_L=0$. On $\{2K>3L\}$ labour binds: $F=3L$, so
        $F_K=0$, $F_L=3$. (Non-differentiable on $\{2K=3L\}$.)

  \item $f_x=yz+2x$, $f_y=xz$, $f_z=xy-2z$. At $(2,-1,3)$:
        $f_x=-3+4=1$, $f_y=6$, $f_z=-2-6=-8$. So
        $\nabla f=(1,6,-8)^T$.

  \item $\nabla f=(2x,8y)^T$. At $(3,1)$: $\nabla f=(6,8)^T$. Steepest
        descent is $-\nabla f/\|\nabla f\|$. $\|\nabla f\|=10$, so unit
        descent direction is $(-3/5,-4/5)^T$, with rate of decrease
        $-\|\nabla f\|=-10$ (i.e.\ $f$ decreases at rate 10).

  \item $f_x=e^y-ye^x$, $f_y=xe^y-e^x$.

  \item Quotient rule with $D=x^2+y^2$:
        \[
          f_x=\frac{D-x\cdot2x}{D^2}=\frac{y^2-x^2}{(x^2+y^2)^2}, \qquad
          f_y=\frac{-x\cdot2y}{D^2}=\frac{-2xy}{(x^2+y^2)^2} .
        \]

  \item $U_x=1$, $U_y=1/\sqrt{y}$. So $\mathrm{MRS}_{xy}=U_x/U_y
        =\sqrt{y}$ (depends only on $y$, quasi-linear utility).

  \item $f_x=2x$, $f_y=2y$, $f_z=-2z$. At $(2,-1,2)$:
        $\nabla f=(4,-2,-4)^T$. Normalise: $\|\mathbf{v}\|=\sqrt{4+4+1}=3$,
        $\mathbf{u}=(2/3,2/3,1/3)^T$. Then
        \[
          D_{\mathbf{u}}f=4\cdot\frac23+(-2)\cdot\frac23+(-4)\cdot\frac13
            =\frac{8-4-4}{3}=0 .
        \]
        (We are moving tangent to the level surface.)

  \item $f_x=\sec^2(x+y)+\sec^2(x-y)$, $f_y=\sec^2(x+y)-\sec^2(x-y)$.

  \item Let $S=\alpha K+\beta L$, $F=S^r$. Then $F_K=rS^{r-1}\alpha$,
        $F_L=rS^{r-1}\beta$. Hence
        \[
          KF_K+LF_L=rS^{r-1}(\alpha K+\beta L)=rS^{r-1}\cdot S=rS^r=rF.
            \checkmark
        \]

  \item $\nabla f=(-e^{-x}\sin y+e^x\cos y,\ e^{-x}\cos y-e^x\sin y)^T$.

  \item $f_x=x/\sqrt{x^2+y^2}$, $f_y=y/\sqrt{x^2+y^2}$. At $(3,4)$:
        $\sqrt{x^2+y^2}=5$, $\nabla f=(3/5,4/5)^T$. The direction
        $\mathbf{v}=(1,0)$ is already unit. So
        $D_{\mathbf{u}}f=3/5\cdot1+4/5\cdot0=3/5$.

  \item On $\{2x<3y\}$: $U=2x$, so $U_x=2$, $U_y=0$. On $\{2x>3y\}$:
        $U=3y$, so $U_x=0$, $U_y=3$. (Non-differentiable on the ray
        $2x=3y$.)

  \item Let $L=\ln(x+y)$, $D=x-y$. By the quotient rule:
        \[
          f_x=\frac{(1/(x+y))(x-y)-\ln(x+y)\cdot1}{(x-y)^2}
            =\frac{1/(x+y)\cdot(x-y)-\ln(x+y)}{(x-y)^2} .
        \]
        Simplifying
        $f_x=\dfrac{(x-y)/(x+y)-\ln(x+y)}{(x-y)^2}$. Similarly,
        \[
          f_y=\frac{(1/(x+y))(x-y)-\ln(x+y)\cdot(-1)}{(x-y)^2}
            =\frac{(x-y)/(x+y)+\ln(x+y)}{(x-y)^2} .
        \]

  \item $F_K=A\alpha K^{\alpha-1}L^\beta$, $F_L=A\beta K^\alpha L^{\beta-1}$.
        So
        \[
          \mathrm{MRTS}_{LK}=\frac{F_L}{F_K}=\frac{\beta}{\alpha}\cdot
            \frac{K}{L} .
        \]
        This depends only on the ratio $K/L$ and the exponent ratio
        $\beta/\alpha$: it is scale-invariant.

  \item $\nabla f=(2x,-2y)^T$. At $(2,-1)$: $\nabla f=(4,2)^T$. Direction
        of steepest ascent: $\mathbf{u}^*=(4,2)/\sqrt{20}=(2,1)/\sqrt5$.
        Maximum rate $=\|\nabla f\|=\sqrt{20}=2\sqrt5$.

  \item $f_x=\sinh(x+y)+\cosh(x-y)$, $f_y=\sinh(x+y)-\cosh(x-y)$.

  \item $\nabla f=(e^{x+y+z},e^{x+y+z},e^{x+y+z})^T$. At $(0,0,0)$:
        $\nabla f=(1,1,1)^T$. Normalise $\mathbf{v}=(1,1,1)$:
        $\|\mathbf{v}\|=\sqrt3$, $\mathbf{u}=(1,1,1)/\sqrt3$. Then
        $D_{\mathbf{u}}f=(1+1+1)/\sqrt3=\sqrt3$.

  \item $f_x=(\ln y)/x$, $f_y=(\ln x)/y$.

  \item $F_K=1+L^{1/2}/K^{1/2}$, $F_L=1+K^{1/2}/L^{1/2}$. From
        $F=(K^{1/2}+L^{1/2})^2$ with $g=K^{1/2}+L^{1/2}$:
        $F_K=g/K^{1/2}=(K^{1/2}+L^{1/2})/K^{1/2}=1+L^{1/2}/K^{1/2}$.
        \checkmark Same for $F_L$.

  \item $f_x=2xy$, $f_y=x^2+3y^2$. At $(-1,2)$: $\nabla f=(-4,1+12)^T
        =(-4,13)^T$. Normalise: $\|\mathbf{v}\|=5$,
        $\mathbf{u}=(4/5,-3/5)^T$. Then
        \[
          D_{\mathbf{u}}f=-4\cdot\frac45+13\cdot\left(-\frac35\right)
            =-\frac{16+39}{5}=-11 .
        \]

  \item $U_x=\alpha/x$, $U_y=\beta/y$.
        $\mathrm{MRS}_{xy}=\dfrac{\alpha/x}{\beta/y}=\dfrac{\alpha y}{\beta x}$.
        At $(\alpha,\beta)$: $\mathrm{MRS}=(\alpha\cdot\beta)/(\beta\cdot\alpha)=1$.

  \item $f=(x^2+y^2)^{-1/2}$. Chain rule:
        \[
          f_x=-\frac12(x^2+y^2)^{-3/2}\cdot2x=\frac{-x}{(x^2+y^2)^{3/2}},
          \qquad f_y=\frac{-y}{(x^2+y^2)^{3/2}} .
        \]
        So $\nabla f=-(x,y)^T/(x^2+y^2)^{3/2}$.

  \item Chain rule with $u=xy$:
        \[
          f_x=\frac{y}{1+x^2y^2}, \qquad f_y=\frac{x}{1+x^2y^2} .
        \]

  \item Let $S=K^\rho+L^\rho$, $F=S^{1/\rho}$. Then
        $F_K=\frac1\rho S^{1/\rho-1}\cdot\rho K^{\rho-1}
        =K^{\rho-1}S^{1/\rho-1}$. Since $S^{1/\rho}=F$, we have $S=F^\rho$
        and so $S^{1/\rho-1}=F^{1-\rho}\cdot F^{-1}\cdot F=F^{1-\rho}$.
        Thus $F_K=K^{\rho-1}F^{1-\rho}=(F/K)^{1-\rho}=(K/F)^{\rho-1}$.
        \checkmark

  \item $f_x=yz$, $f_y=xz$, $f_z=xy$. At $(2,3,-1)$:
        $\nabla f=(3\cdot(-1),2\cdot(-1),2\cdot3)^T=(-3,-2,6)^T$.
        Normalise $\mathbf{v}=(2,-1,2)$: $\|\mathbf{v}\|=3$,
        $\mathbf{u}=(2/3,-1/3,2/3)^T$. Then
        \[
          D_{\mathbf{u}}f=-3\cdot\frac23+(-2)\cdot\left(-\frac13\right)
            +6\cdot\frac23=-2+\frac23+4=2+\frac23=\frac83 .
        \]
\end{solutions}

\section{Additional Problems}

\begin{problems}
  \item Compute $f_x$ and $f_y$ for $f(x,y)=3x^5-2x^3y^2+4y^4$.
  \item For $F(K,L)=20\,K^{1/3}L^{2/3}$, find $F_K$ and $F_L$ at
        $(K,L)=(8,27)$.
  \item Compute $f_x$ and $f_y$ for $f(x,y)=e^{xy}$.
  \item Compute $\nabla f$ for $f(x,y)=2x^2y-3xy^3+5$.
  \item Compute the directional derivative of $f(x,y)=x^2+2xy-y^2$ at
        $(1,2)$ in the direction $\mathbf{v}=(3,4)$.
  \item For $F(K,L)=8K^{3/4}L^{1/4}$, compute $\mathrm{MP}_K$ and
        $\mathrm{MP}_L$ at $(K,L)=(16,16)$.
  \item Compute $f_x$ and $f_y$ for $f(x,y)=\ln(2x+3y)$.
  \item For $U(x,y)=3\ln(x)+4\ln(y)$, derive $\mathrm{MRS}_{xy}$ and
        evaluate at $(2,3)$.
  \item Compute $f_x$ and $f_y$ for $f(x,y)=(x+y^2)^4$.
  \item Compute $\nabla f$ for $f(x,y,z)=xy+yz+zx$.
  \item For a CES production $F(K,L)=5(K^{1/3}+L^{1/3})^3$, find $F_K$.
  \item Compute the directional derivative of $f(x,y)=x^2y$ at $(2,1)$
        in the direction $\mathbf{v}=(1,-1)$.
  \item For $U(x,y)=x^\alpha y^\beta$ (general Cobb--Douglas,
        $\alpha,\beta>0$), derive $\mathrm{MRS}_{xy}$ and show it depends
        only on the ratio $y/x$.
  \item Compute $f_x$ and $f_y$ for $f(x,y)=\dfrac{x-y}{x+y}$.
  \item For $f(K,L)=K^{1/2}L^{1/2}$ at $(K,L)=(16,9)$, find the unit
        direction of steepest ascent and the maximum rate of increase.
  \item Classify the open positive orthant
        $Q=\{(x,y)\in\R^2 : x>0,\ y>0\}$: is it open, closed, or
        neither? Is it bounded? Is it compact?
  \item Classify the budget set
        $B=\{(x,y)\in\R^2 : 2x+3y\le12,\ x\ge0,\ y\ge0\}$: open,
        closed, or neither? Bounded? Compact?
  \item Classify $M=\{(x,y)\in\R^2 : xy\ge1,\ x>0\}$: open, closed, or
        neither? Bounded? Compact? \emph{Hint:} check whether a
        sequence in $M$ can converge to a point with $x=0$.
  \item Classify the open disc
        $D=\{(x,y)\in\R^2 : x^2+y^2<4\}$: open, closed, or neither?
        Bounded? Compact?
  \item Classify (a) the half-open strip
        $S=\{(x,y)\in\R^2 : 0\le x<1\}$ and (b) the finite set
        $E=\{(0,0),\ (1,2),\ (3,1)\}$: open, closed, or neither?
        Bounded? Compact?
  \item Does the Weierstrass theorem guarantee a maximiser for the
        continuous function $u(x,y)=x^{1/2}y^{1/2}$ (a) on the budget
        set $B=\{(x,y): 2x+3y\le12,\ x\ge0,\ y\ge0\}$, (b) on the open
        positive orthant $Q=\{(x,y): x>0,\ y>0\}$? Justify each
        answer.
\end{problems}

\section{Additional Solutions}

\begin{solutions}
  \item $f_x=15x^4-6x^2y^2$, $f_y=-4x^3y+16y^3$.

  \item $F_K=20\cdot\frac13K^{-2/3}L^{2/3}=\frac{20}{3}(L/K)^{2/3}$;
        $F_L=20\cdot\frac23K^{1/3}L^{-1/3}=\frac{40}{3}(K/L)^{1/3}$. At
        $(8,27)$: $(L/K)^{2/3}=(27/8)^{2/3}=9/4$, so $F_K=15$;
        $(K/L)^{1/3}=(8/27)^{1/3}=2/3$, so
        $F_L=\frac{40}{3}\cdot\frac23=\frac{80}{9}$.

  \item $f_x=ye^{xy}$, $f_y=xe^{xy}$.

  \item $\nabla f=(4xy-3y^3,\ 2x^2-9xy^2)^T$.

  \item $f_x=2x+2y$, $f_y=2x-2y$. At $(1,2)$: $\nabla f=(6,-2)^T$.
        Normalise: $\|\mathbf{v}\|=5$, $\mathbf{u}=(3/5,4/5)^T$. Then
        \[
          D_{\mathbf{u}}f=6\cdot\frac35+(-2)\cdot\frac45
            =\frac{18-8}{5}=2 .
        \]

  \item $F_K=8\cdot\frac34K^{-1/4}L^{1/4}=6(L/K)^{1/4}$,
        $F_L=2(K/L)^{3/4}$. At $(16,16)$: $\mathrm{MP}_K=6$,
        $\mathrm{MP}_L=2$.

  \item $f_x=\dfrac{2}{2x+3y}$, $f_y=\dfrac{3}{2x+3y}$.

  \item $U_x=3/x$, $U_y=4/y$.
        $\mathrm{MRS}_{xy}=U_x/U_y=\dfrac{3y}{4x}$. At $(2,3)$:
        $\mathrm{MRS}=\dfrac98$.

  \item Chain: $f_x=4(x+y^2)^3$,
        $f_y=4(x+y^2)^3\cdot2y=8y(x+y^2)^3$.

  \item $f_x=y+z$, $f_y=x+z$, $f_z=x+y$. So
        $\nabla f=(y+z,\ x+z,\ x+y)^T$.

  \item Let $g=K^{1/3}+L^{1/3}$, so $F=5g^3$. Then
        $F_K=5\cdot3g^2\cdot\frac13K^{-2/3}=5g^2K^{-2/3}
        =\dfrac{5(K^{1/3}+L^{1/3})^2}{K^{2/3}}$.

  \item $f_x=2xy$, $f_y=x^2$. At $(2,1)$: $\nabla f=(4,4)^T$. Normalise:
        $\|\mathbf{v}\|=\sqrt2$, $\mathbf{u}=(1/\sqrt2,-1/\sqrt2)^T$.
        Then
        \[
          D_{\mathbf{u}}f=4\cdot\frac{1}{\sqrt2}
            +4\cdot\left(-\frac{1}{\sqrt2}\right)=0 .
        \]
        (Tangent to the level curve.)

  \item $U_x=\alpha x^{\alpha-1}y^\beta$,
        $U_y=\beta x^\alpha y^{\beta-1}$. So
        \[
          \mathrm{MRS}_{xy}=\frac{U_x}{U_y}
            =\frac{\alpha}{\beta}\cdot\frac{y}{x} .
        \]
        Depends only on $y/x$ (and the exponent ratio), not on the
        scale. \checkmark

  \item Quotient rule with $D=x+y$:
        \[
          f_x=\frac{D-(x-y)}{D^2}=\frac{2y}{(x+y)^2}, \qquad
          f_y=\frac{-D-(x-y)}{D^2}=\frac{-2x}{(x+y)^2} .
        \]

  \item $f_K=\frac12K^{-1/2}L^{1/2}$, $f_L=\frac12K^{1/2}L^{-1/2}$. At
        $(16,9)$: $f_K=\frac12\cdot\frac34=3/8$;
        $f_L=\frac12\cdot\frac43=2/3$. So $\nabla f=(3/8,2/3)^T$.
        Magnitude:
        $\|\nabla f\|=\sqrt{9/64+4/9}=\sqrt{81/576+256/576}
        =\sqrt{337/576}=\sqrt{337}/24\approx0.765$. Unit direction of
        steepest ascent:
        $\mathbf{u}^*=(3/8,2/3)^T/(\sqrt{337}/24)=(9,16)^T/\sqrt{337}$.

  \item \textbf{Open.} For any $(x,y)\in Q$ the ball of radius
        $r=\min\{x,y\}>0$ around $(x,y)$ stays inside $Q$, so every
        point is interior. \textbf{Not closed:} the boundary points
        $(x,0)$ and $(0,y)$ with $x,y\ge0$ (the coordinate half-axes,
        including the origin) do not belong to $Q$. \textbf{Not
        bounded:} $Q$ contains $(t,t)$ for every $t>0$, and
        $\|(t,t)\|=t\sqrt2\to\infty$. Hence $Q$ is open, unbounded and
        \textbf{not compact}.

  \item \textbf{Closed.} $B$ is the intersection of the three sets
        $\{2x+3y\le12\}$, $\{x\ge0\}$ and $\{y\ge0\}$, each defined by
        a non-strict inequality on a continuous function; each is
        closed, and an intersection of closed sets is closed (every
        boundary point of $B$ satisfies all three inequalities, so it
        lies in $B$). \textbf{Not open:} $(0,0)\in B$, but every ball
        around the origin contains points with negative coordinates,
        which are outside $B$; so $(0,0)$ is not interior.
        \textbf{Bounded:} from the constraints, $0\le x\le6$ and
        $0\le y\le4$, so $B\subseteq B_8(\mathbf{0})$. Closed and
        bounded in $\R^2$: $B$ is \textbf{compact}.

  \item First note the constraint $x>0$ together with $xy\ge1$ forces
        $y\ge1/x>0$. \textbf{Closed:} suppose
        $(x_n,y_n)\in M$ with $(x_n,y_n)\to(a,b)$. If $a=0$, then
        $y_n\to b$ finite would give $x_ny_n\to0<1$, contradicting
        $x_ny_n\ge1$; so $a>0$, and passing to the limit in
        $x_ny_n\ge1$ gives $ab\ge1$, hence $(a,b)\in M$. $M$ contains
        all its limit (boundary) points, so it is closed: even
        though it is written with a strict inequality, that inequality
        is never active. \textbf{Not open:} $(1,1)\in M$ lies on the
        hyperbola $xy=1$, and every ball around it contains points
        with $xy<1$. \textbf{Not bounded:} $(t,t)\in M$ for all
        $t\ge1$. Hence $M$ is closed, unbounded and \textbf{not
        compact}.

  \item \textbf{Open.} For $(x,y)\in D$ put
        $r=2-\sqrt{x^2+y^2}>0$; the ball $B_r(x,y)$ lies in $D$ by the
        triangle inequality. \textbf{Not closed:} the boundary circle
        $\{x^2+y^2=4\}$ is not contained in $D$ (e.g.\ $(2,0)\notin
        D$). \textbf{Bounded:} $D\subseteq B_3(\mathbf{0})$. Bounded
        but not closed: $D$ is \textbf{not compact}.

  \item (a) \textbf{Neither open nor closed.} The points $(0,y)$
        belong to $S$ but are boundary points (every ball around them
        contains points with $x<0$), so $S$ is not open; the boundary
        points $(1,y)$ do not belong to $S$, so $S$ is not closed.
        \textbf{Unbounded:} $(0,t)\in S$ for all $t$. Hence
        \textbf{not compact}. (b) \textbf{Closed:} around any point
        not in $E$ there is a ball avoiding all three points (a
        finite set always contains all its boundary points:
        indeed $E$ is its own boundary). \textbf{Not open:} no ball
        around $(0,0)$ lies inside $E$. \textbf{Bounded:}
        $E\subseteq B_4(\mathbf{0})$. Closed and bounded: $E$ is
        \textbf{compact}.

  \item (a) \textbf{Yes.} $B$ is compact (closed and bounded, see the
        budget-set problem above) and $u$ is continuous on $B$, so by
        the Weierstrass theorem $u$ attains both a maximum and a
        minimum on $B$; a maximiser exists. (b) \textbf{No.} $Q$ is
        not compact (it is neither closed nor bounded), so Weierstrass
        does not apply: and indeed no maximiser exists:
        $u(t,t)=t\to\infty$ as $t\to\infty$, so $u$ is unbounded above
        on $Q$.
\end{solutions}

\chapter{Curve Sketching: The Complete Discussion}

Everything needed to draw a curve has now been assembled in pieces:
domains and asymptotes in the functions volume, limits in the analysis
volume, the first derivative for monotonicity, the second for curvature.
What is missing is the habit of running all of it in one fixed order and
finishing with a picture. That is the \de{Kurvendiskussion}, and it is worth
drilling as a single routine for two reasons.

First, examinations ask for it directly. Second, and more usefully, the
routine is a self-checking mechanism. A sign table that says the function
falls on $(0,2)$ and a computed local maximum at $x=2$ cannot both be right;
the contradiction surfaces before the picture is drawn. A curve sketch is
the cheapest way there is to catch an algebra slip.

The chapter runs the protocol on rational, exponential, logarithmic and root
functions, and closes with the discussion an economist actually performs:
locating the inflection of a total cost curve, the minimum of average cost,
and the output that maximises profit --- all of them the same routine
applied to a curve with a name.

\section{Core Rules}

\begin{rulebox}{The protocol}
Always in this order, because each step feeds the next:
\begin{enumerate}[leftmargin=1.8em,itemsep=0.35ex]
  \item \textbf{Domain.} Denominators, even roots, logarithms. Note which
        points are excluded --- they are the candidates for poles.
  \item \textbf{Symmetry.} $f(-x)=f(x)$ even, $f(-x)=-f(x)$ odd. Either one
        halves the work.
  \item \textbf{Intercepts and roots.} $f(0)$, and the solutions of
        $f(x)=0$ with their multiplicities.
  \item \textbf{Behaviour at the edges.} Limits at $\pm\infty$ and at every
        excluded point, one side at a time. This produces the
        \textbf{asymptotes}: vertical at poles, horizontal or oblique at
        infinity.
  \item \textbf{First derivative.} Solve $f'(x)=0$ for the critical points,
        build a \textbf{sign table} for $f'$, and read off intervals of
        increase and decrease and the type of each extremum.
  \item \textbf{Second derivative.} Solve $f''(x)=0$, build a sign table
        for $f''$, and read off convex and concave intervals and the
        \textbf{inflection points}.
  \item \textbf{Value table} for any remaining doubt, then \textbf{draw}:
        asymptotes first as dashed guides, then the special points, then
        join them respecting monotonicity and curvature.
\end{enumerate}
\end{rulebox}

\begin{rulebox}{Sign tables}
A sign table is the bookkeeping device that prevents guesswork. List every
point where the expression can change sign --- its roots and its poles ---
in increasing order; they cut the domain into intervals. On each interval
the sign is constant, so \emph{one} test value settles it. For
$f'(x)=3x(x-2)$:
\[
  \begin{array}{c|ccccc}
    x & (-\infty,0) & 0 & (0,2) & 2 & (2,\infty)\\\hline
    f' & + & 0 & - & 0 & +\\
    f & \nearrow & \text{max} & \searrow & \text{min} & \nearrow
  \end{array}
\]
A sign change $+\to-$ is a local maximum, $-\to+$ a local minimum, and
\emph{no} change means no extremum, however much the derivative vanishes
(as at $x=0$ for $f=x^{3}$).
\end{rulebox}

\begin{rulebox}{First- and second-derivative tests}
Let $f'(x_{0})=0$.
\begin{itemize}[leftmargin=1.6em,itemsep=0.3ex]
  \item \textbf{First-derivative test:} the sign change of $f'$ across
        $x_{0}$ decides the type. Always applicable, and the only reliable
        test when $f''(x_{0})=0$.
  \item \textbf{Second-derivative test:} $f''(x_{0})>0$ gives a local
        minimum, $f''(x_{0})<0$ a local maximum. If $f''(x_{0})=0$ the test
        is \textbf{silent} --- fall back on the sign table. Compare
        $x^{4}$ (minimum), $-x^{4}$ (maximum) and $x^{3}$ (neither), all
        three with $f'=f''=0$ at the origin.
  \item \textbf{Inflection} at $x_{0}$ requires $f''$ to \emph{change
        sign} there. $f''(x_{0})=0$ alone is not enough: $f=x^{4}$ has
        $f''(0)=0$ and no inflection.
\end{itemize}
\end{rulebox}

\begin{rulebox}{Asymptotes, recalled}
For $f=p/q$ with $p,q$ polynomials, after cancelling common factors:
vertical asymptote at each zero of $q$; horizontal $y=0$ if
$\deg p<\deg q$; horizontal $y=$ ratio of leading coefficients if
$\deg p=\deg q$; and if $\deg p=\deg q+1$, polynomial division gives
\[
  f(x)=\underbrace{mx+c}_{\text{oblique asymptote}}+\frac{\text{remainder}}{q(x)} .
\]
For non-rational functions the same information comes from limits directly:
$\lim_{x\to\pm\infty}f(x)$, and for a possible oblique asymptote
$m=\lim_{x\to\infty}f(x)/x$ followed by $c=\lim_{x\to\infty}(f(x)-mx)$.
\end{rulebox}

\begin{gallery}
\hspace*{\fill}%
\begin{tikzpicture}
  \begin{axis}[fnaxis, width=8.4cm, height=7.4cm,
      xmin=-4.4, xmax=5.4, ymin=-6.4, ymax=10.4,
      xtick={-4,-2,2,4}, ytick={-4,4,8}]
    \addplot[gray!55, dashed, thin, samples=2, forget plot]
      coordinates {(1,-6.2)(1,10.2)};
    \addplot[gray!55, dashed, thin, samples=2, domain=-4.2:5.2, forget plot]{x+1};
    \addplot[domain=-4.2:0.86, samples=201]{x^2/(x-1)};
    \addplot[domain=1.14:5.2, samples=201]{x^2/(x-1)};
    \addplot[only marks, mark=*, mark size=1.3pt, black, forget plot]
      coordinates {(0,0) (2,4)};
    \node[font=\tiny, anchor=north east] at (axis cs:-0.1,-0.15) {$(0,0)$};
    \node[font=\tiny, anchor=west] at (axis cs:2.15,4) {$(2,4)$};
  \end{axis}
\end{tikzpicture}%
\hspace*{\fill}
\gcap{The finished discussion of $f(x)=x^{2}/(x-1)$: pole $x=1$, oblique
asymptote $y=x+1$, local maximum $(0,0)$ below the local minimum $(2,4)$.}
\end{gallery}

\begin{rulebox}{The economist's curve discussion}
Three facts fall straight out of the protocol applied to a total cost
function $C(q)$, and they are worth knowing as results:
\begin{itemize}[leftmargin=1.6em,itemsep=0.4ex]
  \item Marginal cost is $MC=C'$. The \textbf{inflection point of $C$} is
        exactly where $MC$ is smallest, since $MC'=C''$ changes sign there.
        The S-shaped cost curve turns from concave to convex at its most
        efficient scale.
  \item Average cost is $AC=C(q)/q$. Differentiating,
        \[
          AC'=\frac{C'(q)q-C(q)}{q^{2}}=\frac{MC-AC}{q},
        \]
        so $AC'=0$ exactly when $MC=AC$: \textbf{marginal cost cuts average
        cost at the minimum of average cost}. That is a theorem, not a
        coincidence of a particular example.
  \item Profit at price $p$ is $\pi=pq-C(q)$, so $\pi'=p-MC$ and the
        optimum has $p=MC$, confirmed as a maximum by $\pi''=-C''<0$,
        i.e.\ by cost being convex there.
\end{itemize}
\end{rulebox}

\section{Worked Example}

\begin{worked}{The complete discussion of $f(x)=\dfrac{x^{2}}{x-1}$}
\wstep{Step 1. Domain, and the reason to find it first.}
The denominator vanishes at $x=1$, so
$D_{f}=\R\setminus\{1\}$. Every later step --- sign tables, monotonicity,
convexity --- must be stated \emph{per interval}, and $x=1$ is the wall
between them.

\wstep{Step 2. Symmetry and roots.}
$f(-x)=\frac{x^{2}}{-x-1}$, which is neither $f(x)$ nor $-f(x)$: no
symmetry. The numerator vanishes at $x=0$ with multiplicity two, so the
curve \textbf{touches} the axis at the origin rather than crossing it, and
$f(0)=0$.

\wstep{Step 3. Behaviour at the pole, from both sides.}
Near $x=1$ the numerator is close to $1>0$, so the sign of $f$ is the sign
of $x-1$:
\[
  x\to1^{-}:\ \frac{1^{+}}{0^{-}}\to-\infty,
  \qquad
  x\to1^{+}:\ \frac{1^{+}}{0^{+}}\to+\infty .
\]
\wtrap{Take one-sided limits at every pole. Writing ``$f\to\infty$ at
$x=1$'' loses the sign information and makes the sketch impossible to draw
correctly.}

\wstep{Step 4. Behaviour at infinity: divide, do not guess.}
Numerator degree $2$ exceeds denominator degree $1$ by exactly one, so there
is an \textbf{oblique} asymptote. Polynomial division gives it:
\[
  f(x)=x+1+\frac{1}{x-1} .
\]
So $y=x+1$ is the asymptote, and the remainder $\frac{1}{x-1}$ says how the
curve sits relative to it: \emph{above} for $x>1$, \emph{below} for $x<1$.
\wtrap{Horizontal asymptote when the degrees are equal or the numerator's is
smaller; oblique when the numerator's degree is exactly one larger; neither
when it is larger by two or more. And the division is the whole method ---
guessing $y=x$ from the leading behaviour misses the $+1$.}

\wstep{Step 5. First derivative and critical points.}
Differentiating the divided form is easiest:
\[
  f'(x)=1-\frac{1}{(x-1)^{2}}=\frac{(x-1)^{2}-1}{(x-1)^{2}}=\frac{x(x-2)}{(x-1)^{2}} .
\]
Zeros at $x=0$ and $x=2$; the denominator is positive wherever it is
defined, so the sign of $f'$ is the sign of $x(x-2)$.
\begin{center}\footnotesize
\begin{tabular}{lcccc}
  \hline
  interval & $x<0$ & $0<x<1$ & $1<x<2$ & $x>2$\\
  \hline
  $f'$ & $+$ & $-$ & $-$ & $+$\\
  $f$ & rising & falling & falling & rising\\
  \hline
\end{tabular}
\end{center}

\wstep{Step 6. Classify, and notice something strange.}
$f''(x)=\dfrac{2}{(x-1)^{3}}$, so $f''(0)=-2<0$ and $f''(2)=2>0$:
\[
  \text{local maximum at }(0,0),
  \qquad
  \text{local minimum at }(2,4).
\]
The local maximum lies \textbf{below} the local minimum. That is impossible
for a continuous function --- and possible here only because the pole at
$x=1$ separates them. The two branches never see each other.
\wtrap{Local extrema are local. A local maximum of $0$ coexisting with a
local minimum of $4$ is not an arithmetic error, and neither point is a
global extremum: $f$ is unbounded in both directions.}

\wstep{Step 7. Convexity, and the absent inflection.}
$f''=\frac{2}{(x-1)^{3}}$ is never zero, so there is \textbf{no inflection
point}. It is negative for $x<1$ and positive for $x>1$:
concave on $(-\infty,1)$, convex on $(1,\infty)$. The change of curvature
happens at the pole, where the function does not exist --- so it is not an
inflection.
\wtrap{A change of convexity across a point of the domain is an inflection;
across a gap in the domain it is nothing at all. Check that the candidate
point is in the domain before naming it.}

\wstep{Step 8. Draw it.}
\begin{center}
\begin{tikzpicture}
\begin{axis}[fnaxis, width=8.4cm, height=6.0cm,
             xmin=-4.2, xmax=6.2, ymin=-9, ymax=13,
             xtick={-3,-1,1,2,4}, ytick={4,8},
             restrict y to domain=-40:40]
  \addplot[domain=-4:0.88, samples=200] {x^2/(x-1)};
  \addplot[domain=1.12:6, samples=200] {x^2/(x-1)};
  \addplot[accent!40, thick, dashed, samples=2, domain=-4:6, forget plot]
    {x+1};
  \addplot[gray!55, thin, dashed, samples=2, forget plot]
    coordinates {(1,-9) (1,13)};
  \addplot[only marks, mark=*, mark size=1.2pt, accent!70, forget plot]
    coordinates {(0,0) (2,4)};
\end{axis}
\end{tikzpicture}
\end{center}
Everything checks against the figure: the touch at the origin, the two
branches, the pole at $x=1$, the asymptote $y=x+1$ approached from below on
the left and from above on the right, and the maximum sitting lower than the
minimum.

\wstep{Step 9. The protocol, in the order that saves work.}
\begin{center}\footnotesize
\begin{tabular}{ll}
  \hline
  1. domain & 6. sign table for $f'$\\
  2. symmetry & 7. extrema, classified by $f''$\\
  3. roots and intercepts & 8. $f''$, convexity, inflections\\
  4. poles, one-sided limits & 9. sketch\\
  5. asymptotes (divide!) & 10. check the sketch against 1--8\\
  \hline
\end{tabular}
\end{center}
Step~10 is not optional. A sketch that contradicts the sign table means one
of them is wrong, and finding out which takes a minute --- far less than
losing the marks.
\end{worked}

\section{Practice Problems}

In problems 1--8 give the intervals of increase and decrease and all local
extrema, using a sign table.

\begin{problems}
  \item $f(x)=x^{3}-12x$
  \item $f(x)=x^{4}-8x^{2}$
  \item $f(x)=x\mathrm{e}^{-x}$
  \item $f(x)=\dfrac{x}{x^{2}+1}$
  \item $f(x)=x-\ln x$ on $x>0$
  \item $f(x)=x^{3}$ --- and explain why $f'(0)=0$ yields no extremum.
  \item $f(x)=\dfrac{x^{2}}{x-1}$
  \item $f(x)=\sqrt{x}-x$ on $x\ge0$
\end{problems}

In problems 9--14 give the convex and concave intervals and all inflection
points.

\begin{problems}[resume]
  \item $f(x)=x^{3}-3x^{2}$
  \item $f(x)=x^{4}-4x^{2}$
  \item $f(x)=\mathrm{e}^{-x^{2}}$
  \item $f(x)=x^{4}$ --- and explain why $f''(0)=0$ gives no inflection.
  \item $f(x)=\dfrac{\ln x}{x}$ on $x>0$
  \item $f(x)=x^{2}\mathrm{e}^{-x}$
\end{problems}

In problems 15--20 give every asymptote.

\begin{problems}[resume]
  \item $f(x)=\dfrac{x-1}{x+2}$
  \item $f(x)=\dfrac{1}{x^{2}-4}$
  \item $f(x)=\dfrac{x^{2}}{x-1}$
  \item $f(x)=\dfrac{x^{3}}{x^{2}+1}$
  \item $f(x)=\dfrac{\ln x}{x}$
  \item $f(x)=x+\mathrm{e}^{-x}$
\end{problems}

Problems 21--32 ask for a \emph{complete} discussion: domain, symmetry,
intercepts, asymptotes and edge behaviour, monotonicity and extrema,
curvature and inflections, and a sketch.

\begin{problems}[resume]
  \item $f(x)=x^{3}-3x^{2}$
  \item $f(x)=x^{4}-4x^{2}$
  \item $f(x)=\dfrac{x}{x^{2}+1}$
  \item $f(x)=\dfrac{x^{2}}{x-1}$
  \item $f(x)=\mathrm{e}^{-x^{2}}$
  \item $f(x)=x\mathrm{e}^{-x}$
  \item $f(x)=\dfrac{\ln x}{x}$
  \item $f(x)=\dfrac{x-1}{x+2}$
  \item $f(x)=x^{2}\mathrm{e}^{-x}$
  \item $f(x)=\dfrac{1}{x^{2}-4}$
  \item $f(x)=\sqrt{x}-x$
  \item $f(x)=x^{3}-3x+2$
\end{problems}

\begin{problems}[resume]
  \item A function satisfies $f'>0$ and $f''<0$ throughout an interval.
        Describe the shape of its graph there in one sentence.
  \item The graph of $f'$ is a parabola opening upwards with roots at
        $x=-1$ and $x=3$. What does $f$ do, and where are its extrema?
  \item The graph of $f'$ is a straight line through the origin with
        positive slope. What is $f$, up to a constant, and what is its
        shape?
  \item $f$ is differentiable with $f(1)=2$, $f'(1)=0$ and $f''(1)=-3$.
        Classify $x=1$ and give the tangent line there.
  \item Sketch a single curve with all of: domain $\R\setminus\{2\}$, a
        pole at $x=2$, horizontal asymptote $y=1$, a local minimum at
        $(0,0)$, and convex on $(2,\infty)$.
  \item Why can a function with a vertical asymptote at $x=a$ have a local
        maximum lying below a local minimum? Illustrate with problem 24.
\end{problems}

\begin{problems}[resume]
  \item For $C(q)=q^{3}-6q^{2}+15q$ on $q\ge0$: show $C$ is increasing
        throughout, find its inflection point, and show that marginal cost
        is minimal exactly there.
  \item For the same $C$: find the output minimising average cost and
        verify that $MC=AC$ at that output.
  \item With the same $C$ and price $p=15$, find the profit-maximising
        output and the maximal profit.
  \item For $R(q)=q(200-4q)$: discuss the curve on its economic domain and
        find the revenue-maximising quantity.
  \item For $C(q)=0.5q^{2}+2q+8$: find the minimum of $AC$ and confirm
        $MC=AC$ there.
  \item The adoption path is $y(t)=\dfrac{100}{1+\mathrm{e}^{-0.5t}}$.
        Find its asymptotes, show it is increasing, and locate its
        inflection point.
  \item Explain, using $AC'=(MC-AC)/q$, why average cost must be falling
        wherever marginal cost lies below it.
\end{problems}

\section{Solutions}

\begin{solutions}
  \item $f'=3x^{2}-12=3(x-2)(x+2)$: $+$ on $(-\infty,-2)$, $-$ on
        $(-2,2)$, $+$ on $(2,\infty)$. Local maximum $(-2,16)$, local
        minimum $(2,-16)$.
  \item $f'=4x^{3}-16x=4x(x-2)(x+2)$: $-$, $+$, $-$, $+$ on the four
        intervals cut by $-2,0,2$. Local minima $(\pm2,-16)$, local maximum
        $(0,0)$.
  \item $f'=(1-x)\mathrm{e}^{-x}$: positive for $x<1$, negative for $x>1$.
        Single local (and global) maximum $(1,1/\mathrm{e})$.
  \item $f'=\dfrac{1-x^{2}}{(x^{2}+1)^{2}}$: negative, positive, negative
        on $(-\infty,-1)$, $(-1,1)$, $(1,\infty)$. Local minimum
        $(-1,-\tfrac12)$, local maximum $(1,\tfrac12)$.
  \item $f'=1-\dfrac1x=\dfrac{x-1}{x}$: negative on $(0,1)$, positive on
        $(1,\infty)$. Local (global) minimum $(1,1)$.
  \item $f'=3x^{2}\ge0$ with equality only at $0$: the sign of $f'$ never
        changes, so $f$ increases throughout and $x=0$ is a
        \emph{saddle-like} stationary point, not an extremum. A vanishing
        derivative is necessary, never sufficient.
  \item $f'=\dfrac{x(x-2)}{(x-1)^{2}}$. The denominator is always positive,
        so the sign is that of $x(x-2)$: $+$, $-$, $-$, $+$ across
        $0,1,2$ (with $x=1$ excluded). Local maximum $(0,0)$, local minimum
        $(2,4)$; the function decreases on $(0,1)$ and on $(1,2)$.
  \item $f'=\dfrac{1}{2\sqrt{x}}-1$, zero when $\sqrt{x}=\tfrac12$, i.e.
        $x=\tfrac14$: positive before, negative after. Global maximum
        $\bigl(\tfrac14,\tfrac14\bigr)$.
  \item $f''=6x-6$: negative for $x<1$ (concave), positive for $x>1$
        (convex). Inflection $(1,-2)$.
  \item $f''=12x^{2}-8$: zero at $x=\pm\sqrt{2/3}$; convex outside,
        concave between. Inflections at
        $\bigl(\pm\sqrt{2/3},-\tfrac{20}{9}\bigr)$.
  \item $f'=-2x\mathrm{e}^{-x^{2}}$ and
        $f''=(4x^{2}-2)\mathrm{e}^{-x^{2}}$: zero at
        $x=\pm1/\sqrt2$, convex outside, concave between. Inflections at
        $\bigl(\pm1/\sqrt2,\mathrm{e}^{-1/2}\bigr)$.
  \item $f''=12x^{2}\ge0$, vanishing only at $x=0$ but \emph{not changing
        sign} there: the curve is convex on both sides, so there is no
        inflection. $f''=0$ is necessary, not sufficient.
  \item $f'=\dfrac{1-\ln x}{x^{2}}$, $f''=\dfrac{2\ln x-3}{x^{3}}$: zero at
        $\ln x=\tfrac32$, i.e. $x=\mathrm{e}^{3/2}$; concave before, convex
        after. Inflection at
        $\bigl(\mathrm{e}^{3/2},\tfrac32\mathrm{e}^{-3/2}\bigr)$.
  \item $f'=(2x-x^{2})\mathrm{e}^{-x}$,
        $f''=(x^{2}-4x+2)\mathrm{e}^{-x}$: zero at $x=2\pm\sqrt2$, convex
        outside those, concave between.
  \item Pole $x=-2$; equal degrees, so horizontal asymptote $y=1$.
  \item Poles $x=\pm2$; numerator degree smaller, so $y=0$.
  \item Pole $x=1$; degrees differ by one, and
        $f=x+1+\dfrac{1}{x-1}$, so the oblique asymptote is $y=x+1$.
  \item No real pole ($x^{2}+1>0$); degrees differ by one, and
        $x^{3}=x(x^{2}+1)-x$, so $f=x-\dfrac{x}{x^{2}+1}$ and the oblique
        asymptote is $y=x$.
  \item Vertical asymptote $x=0$ (as $x\to0^{+}$, $f\to-\infty$) and
        horizontal $y=0$ as $x\to\infty$, since $\ln x\ll x$.
  \item As $x\to+\infty$, $\mathrm{e}^{-x}\to0$, so $y=x$ is an oblique
        asymptote. As $x\to-\infty$ the exponential dominates and
        $f\to+\infty$ with no asymptote.
  \item Domain $\R$; no symmetry; $f=x^{2}(x-3)$ so roots $0$ (double,
        touching) and $3$; $f(0)=0$; ends $-\infty$ and $+\infty$. From
        problem 9's table: increasing on $(-\infty,0)$, decreasing on
        $(0,2)$, increasing on $(2,\infty)$, local maximum $(0,0)$, local
        minimum $(2,-4)$; concave then convex with inflection $(1,-2)$. See
        the figure.
  \item Domain $\R$; \emph{even}; $f=x^{2}(x^{2}-4)$ so roots $0$ (double)
        and $\pm2$; both ends $+\infty$; local maximum $(0,0)$, local
        minima $(\pm\sqrt2,-4)$; inflections
        $\bigl(\pm\sqrt{2/3},-\tfrac{20}{9}\bigr)$. A symmetric W.
  \item Domain $\R$; \emph{odd}; single root $0$; horizontal asymptote
        $y=0$ at both ends; local minimum $(-1,-\tfrac12)$, local maximum
        $(1,\tfrac12)$;
        $f''=\dfrac{2x(x^{2}-3)}{(x^{2}+1)^{3}}$ gives inflections at
        $(0,0)$ and $\bigl(\pm\sqrt3,\pm\tfrac{\sqrt3}{4}\bigr)$. See the
        figure.
  \item The worked example of the Core Rules: domain $\R\setminus\{1\}$,
        double root at $0$, pole $x=1$, oblique asymptote $y=x+1$, local
        maximum $(0,0)$, local minimum $(2,4)$, no inflection, concave left
        of the pole and convex right of it.
  \item Domain $\R$; \emph{even}; no roots ($f>0$ always); horizontal
        asymptote $y=0$ at both ends; global maximum $(0,1)$; inflections
        at $\bigl(\pm1/\sqrt2,\mathrm{e}^{-1/2}\bigr)$; convex outside
        them, concave between. The bell curve.
  \item Domain $\R$; no symmetry; single root $0$; as $x\to+\infty$,
        $f\to0^{+}$ (asymptote $y=0$), as $x\to-\infty$, $f\to-\infty$;
        global maximum $(1,1/\mathrm{e})$; inflection
        $(2,2/\mathrm{e}^{2})$, concave before, convex after. See the
        figure.
  \item Domain $(0,\infty)$; root $x=1$; vertical asymptote $x=0$ with
        $f\to-\infty$, horizontal asymptote $y=0$ as $x\to\infty$; global
        maximum $(\mathrm{e},1/\mathrm{e})$; inflection at
        $x=\mathrm{e}^{3/2}$. See the figure.
  \item Domain $\R\setminus\{-2\}$; root $x=1$; $f(0)=-\tfrac12$; pole
        $x=-2$ ($-\infty$ from the left, $+\infty$ from the right);
        horizontal asymptote $y=1$; $f'=\dfrac{3}{(x+2)^{2}}>0$ so strictly
        increasing on each branch with \emph{no} extrema;
        $f''=\dfrac{-6}{(x+2)^{3}}$, convex left of the pole, concave right
        of it.
  \item Domain $\R$; no symmetry; double root $0$; as $x\to+\infty$,
        $f\to0^{+}$; as $x\to-\infty$, $f\to+\infty$; local minimum
        $(0,0)$, local maximum $(2,4/\mathrm{e}^{2})$; inflections at
        $x=2\pm\sqrt2$.
  \item Domain $\R\setminus\{\pm2\}$; \emph{even}; no roots; poles
        $x=\pm2$; horizontal asymptote $y=0$;
        $f'=\dfrac{-2x}{(x^{2}-4)^{2}}$, so increasing for $x<0$ and
        decreasing for $x>0$, giving a local maximum
        $\bigl(0,-\tfrac14\bigr)$ --- a maximum with a negative value,
        since the middle branch lies entirely below the axis.
  \item Domain $[0,\infty)$; roots $0$ and $1$; $f\to-\infty$ as
        $x\to\infty$; global maximum $\bigl(\tfrac14,\tfrac14\bigr)$;
        $f''=-\tfrac14x^{-3/2}<0$, so concave throughout with no
        inflection; the tangent at $0$ is vertical.
  \item Domain $\R$; $f=(x-1)^{2}(x+2)$, so a double root at $1$ (touching)
        and a simple root at $-2$; $f(0)=2$; ends $\mp\infty$; local
        maximum $(-1,4)$, local minimum $(1,0)$; inflection $(0,2)$.
  \item It rises, but ever more slowly: increasing and concave, bending
        away from its own tangent downwards --- the shape of a production
        function with diminishing returns.
  \item $f'<0$ between the roots and $>0$ outside them, so $f$ increases on
        $(-\infty,-1)$, decreases on $(-1,3)$, increases on $(3,\infty)$: a
        local maximum at $x=-1$ and a local minimum at $x=3$.
  \item $f'(x)=kx$ with $k>0$ gives $f(x)=\tfrac{k}{2}x^{2}+c$: an upward
        parabola with vertex at $x=0$, convex everywhere.
  \item $f'(1)=0$ with $f''(1)<0$: a local \textbf{maximum} at $(1,2)$ by
        the second-derivative test. The tangent is horizontal, $y=2$.
  \item One valid answer: $f(x)=1-\dfrac{4}{(x-2)^{2}}$ has the pole and the
        asymptote $y=1$ but a maximum, not a minimum, so adjust ---
        $f(x)=\dfrac{x^{2}}{(x-2)^{2}}$ works: $f(0)=0$ is a local minimum,
        the pole is at $x=2$, and $f\to1$ at both ends. Any sketch with a
        left branch dipping to $(0,0)$ and rising to $+\infty$ at $x=2^{-}$,
        and a right branch descending convexly from $+\infty$ towards
        $y=1$, is correct.
  \item Because a local extremum is only a statement about a
        \emph{neighbourhood}. The pole cuts the domain into two intervals
        with no path between them, so the values on one side need not
        compare with those on the other. In problem 24 the maximum $(0,0)$
        lies on the left branch, the minimum $(2,4)$ on the right, and
        between them the function runs off to $\mp\infty$.
  \item $C'=3q^{2}-12q+15=3(q^{2}-4q+5)=3\bigl((q-2)^{2}+1\bigr)\ge3>0$, so
        $C$ increases for every $q$. $C''=6q-12$ vanishes at $q=2$ and
        changes sign there: inflection at $(2,14)$. Since $MC=C'$ and
        $MC'=C''$, marginal cost has its minimum exactly at $q=2$, where
        $MC=3$.
  \item $AC=q^{2}-6q+15=(q-3)^{2}+6$, minimal at $q=3$ with $AC=6$. There
        $MC=3\cdot9-12\cdot3+15=6=AC$, as the identity
        $AC'=(MC-AC)/q$ requires.
  \item $\pi=15q-(q^{3}-6q^{2}+15q)=-q^{3}+6q^{2}$. Then
        $\pi'=-3q^{2}+12q=-3q(q-4)$, zero at $q=4$ (and $q=0$), and
        $\pi''=-6q+12=-12<0$ at $q=4$: a maximum, with
        $\pi(4)=-64+96=32$.
  \item $R=200q-4q^{2}=-4(q-25)^{2}+2500$ on the economic domain
        $[0,50]$ (where the price $200-4q$ stays non-negative): a downward
        parabola, increasing on $[0,25)$, decreasing on $(25,50]$, concave
        throughout, with maximum revenue $2500$ at $q=25$ and zeros at both
        ends.
  \item $AC=0.5q+2+\dfrac{8}{q}$, so $AC'=0.5-\dfrac{8}{q^{2}}=0$ gives
        $q^{2}=16$, $q=4$, with $AC(4)=2+2+2=6$. And $MC=q+2=6$ at $q=4$:
        equal, as required.
  \item As $t\to\infty$, $y\to100$; as $t\to-\infty$, $y\to0$: two
        horizontal asymptotes. $y'=\dfrac{50\mathrm{e}^{-0.5t}}
        {(1+\mathrm{e}^{-0.5t})^{2}}>0$ always, so $y$ increases
        throughout. The inflection is where the numerator's growth turns,
        at $\mathrm{e}^{-0.5t}=1$, i.e. $t=0$, with $y=50$ --- exactly half
        the ceiling, which is the signature of the logistic curve.
  \item $q>0$, so the sign of $AC'$ is the sign of $MC-AC$. If
        $MC<AC$ then $AC'<0$ and average cost falls: producing one more
        unit costs less than the current average, which pulls the average
        down. Once $MC$ rises above $AC$ the average is pulled up, and the
        crossing point is therefore the minimum.
\end{solutions}

\begin{gallery}
\begin{tikzpicture}
  \begin{axis}[fnaxis, xmin=-1.6, xmax=3.9, ymin=-5.4, ymax=4.4,
      xtick={-1,1,2,3}, ytick={-4,-2,2}, title={$x^{3}-3x^{2}$}]
    \addplot[domain=-1.2:3.6, samples=201]{x^3-3*x^2};
    \addplot[only marks, mark=*, mark size=1.2pt, black, forget plot]
      coordinates {(0,0) (2,-4) (1,-2)};
  \end{axis}
\end{tikzpicture}\hfill
\begin{tikzpicture}
  \begin{axis}[fnaxis, xmin=-5.4, xmax=5.4, ymin=-0.75, ymax=0.75,
      xtick={-4,-2,2,4}, ytick={-0.5,0.5}, title={$x/(x^{2}+1)$}]
    \addplot[domain=-5.2:5.2, samples=201]{x/(x^2+1)};
    \addplot[only marks, mark=*, mark size=1.2pt, black, forget plot]
      coordinates {(1,0.5) (-1,-0.5) (0,0)};
  \end{axis}
\end{tikzpicture}
\end{gallery}

\begin{gallery}
\begin{tikzpicture}
  \begin{axis}[fnaxis, xmin=-1.4, xmax=5.4, ymin=-1.4, ymax=0.75,
      xtick={-1,1,2,4}, ytick={-1,0.5}, title={$x\mathrm{e}^{-x}$}]
    \addplot[domain=-1.2:5.2, samples=201]{x*exp(-x)};
    \addplot[only marks, mark=*, mark size=1.2pt, black, forget plot]
      coordinates {(1,0.3679) (2,0.2707)};
  \end{axis}
\end{tikzpicture}\hfill
\begin{tikzpicture}
  \begin{axis}[fnaxis, xmin=-0.5, xmax=9.4, ymin=-1.9, ymax=0.75,
      xtick={2,4,6,8}, ytick={-1,0.5}, title={$\ln x/x$}]
    \addplot[domain=0.14:9.2, samples=241]{ln(x)/x};
    \addplot[only marks, mark=*, mark size=1.2pt, black, forget plot]
      coordinates {(2.7183,0.3679) (4.4817,0.3347)};
  \end{axis}
\end{tikzpicture}
\gcap{Four of the completed discussions: problems 21, 23, 26 and 27, with
their extrema and inflection points marked.}
\end{gallery}

\section{Additional Problems}

\begin{problems}
  \item Monotonicity and extrema of $f(x)=2x^{3}-9x^{2}+12x$.
  \item Monotonicity and extrema of $f(x)=\dfrac{x^{2}}{x^{2}+3}$.
  \item Monotonicity and extrema of $f(x)=x\ln x$ on $x>0$.
  \item Convexity and inflections of $f(x)=x^{3}-6x^{2}+9x$.
  \item Convexity and inflections of $f(x)=\dfrac{1}{1+x^{2}}$.
  \item Convexity of $f(x)=\mathrm{e}^{x}+\mathrm{e}^{-x}$, and its
        minimum.
  \item All asymptotes of $f(x)=\dfrac{2x^{2}-3}{x+1}$.
  \item All asymptotes of $f(x)=\dfrac{x^{2}-1}{x^{2}+1}$.
  \item All asymptotes of $f(x)=x\mathrm{e}^{1/x}$.
  \item Complete discussion of $f(x)=x^{3}-6x^{2}+9x$.
  \item Complete discussion of $f(x)=\dfrac{x^{2}-4}{x^{2}}$.
  \item Complete discussion of $f(x)=(x-2)\mathrm{e}^{-x}$.
  \item Complete discussion of $f(x)=x^{2}\ln x$ on $x>0$.
  \item Complete discussion of $f(x)=\dfrac{4x}{x^{2}+4}$.
  \item Complete discussion of $f(x)=\sqrt{4-x^{2}}$.
  \item A cubic has a local maximum at $(0,4)$ and a local minimum at
        $(2,0)$. Determine it.
  \item A function has $f'(x)>0$ everywhere and $f''(x)<0$ everywhere, with
        horizontal asymptote $y=5$. Sketch it and give an example formula.
  \item Where does $f(x)=x^{4}-4x^{3}$ have inflections, and why does the
        second-derivative test fail at one of its critical points?
  \item Show that $f(x)=x+1/x$ has a local minimum at $x=1$ and a local
        maximum at $x=-1$, with the maximum value below the minimum value,
        and explain.
  \item For which $a$ does $f(x)=x^{3}+ax$ have two local extrema?
  \item For $C(q)=q^{3}-9q^{2}+30q$: find the inflection, the minimum of
        average cost, and verify $MC=AC$ there.
  \item For $C(q)=2q^{2}+50$ and price $p=40$: find the profit-maximising
        output, the profit, and the break-even outputs.
  \item A monopolist faces $p(q)=100-2q$ with $C(q)=q^{2}+20q$. Discuss the
        profit curve and find the optimum.
  \item Show that if $AC$ has a minimum at $q^{*}$ then the $AC$ curve is
        tangent to no ray from the origin except the one through
        $(q^{*},AC(q^{*}))$, and relate this to $MC=AC$.
  \item The production function $Q(L)=12L^{2}-L^{3}$ describes output from
        labour $L\in[0,12]$. Find where average product is maximal and show
        that marginal product equals average product there.
\end{problems}

\section{Additional Solutions}

\begin{solutions}
  \item $f'=6x^{2}-18x+12=6(x-1)(x-2)$: $+,-,+$. Local maximum $(1,5)$,
        local minimum $(2,4)$.
  \item $f'=\dfrac{6x}{(x^{2}+3)^{2}}$: negative for $x<0$, positive for
        $x>0$. Global minimum $(0,0)$; $f\to1$ at both ends, so no
        maximum.
  \item $f'=\ln x+1$, zero at $x=\mathrm{e}^{-1}$: negative before,
        positive after. Global minimum
        $\bigl(\mathrm{e}^{-1},-\mathrm{e}^{-1}\bigr)$.
  \item $f''=6x-12$: concave for $x<2$, convex for $x>2$; inflection
        $(2,2)$.
  \item $f'=\dfrac{-2x}{(1+x^{2})^{2}}$,
        $f''=\dfrac{6x^{2}-2}{(1+x^{2})^{3}}$: zero at
        $x=\pm1/\sqrt3$, convex outside, concave between. Inflections at
        $\bigl(\pm1/\sqrt3,\tfrac34\bigr)$.
  \item $f''=\mathrm{e}^{x}+\mathrm{e}^{-x}>0$ everywhere: convex, no
        inflection. $f'=\mathrm{e}^{x}-\mathrm{e}^{-x}=0$ at $x=0$, giving
        the global minimum $(0,2)$.
  \item Pole $x=-1$; degrees differ by one, and
        $2x^{2}-3=(x+1)(2x-2)-1$, so $f=2x-2-\dfrac{1}{x+1}$: oblique
        asymptote $y=2x-2$.
  \item No poles; equal degrees, so horizontal asymptote $y=1$ at both
        ends.
  \item As $x\to\pm\infty$, $\mathrm{e}^{1/x}\to1$ and more precisely
        $x\mathrm{e}^{1/x}=x+1+\tfrac{1}{2x}+\cdots$, so $y=x+1$ is an
        oblique asymptote. As $x\to0^{+}$, $f\to+\infty$ (vertical
        asymptote $x=0$); as $x\to0^{-}$, $f\to0$.
  \item $f=x(x-3)^{2}$: roots $0$ and $3$ (double); ends $\mp\infty$;
        $f'=3x^{2}-12x+9=3(x-1)(x-3)$ gives local maximum $(1,4)$ and local
        minimum $(3,0)$; $f''=6x-12$ gives inflection $(2,2)$; no
        symmetry, no asymptotes.
  \item $f=1-\dfrac{4}{x^{2}}$: domain $\R\setminus\{0\}$, \emph{even},
        roots $\pm2$, pole $x=0$ with $f\to-\infty$ from both sides,
        horizontal asymptote $y=1$; $f'=\dfrac{8}{x^{3}}$, so decreasing
        for $x<0$ and increasing for $x>0$ with no extrema;
        $f''=-\dfrac{24}{x^{4}}<0$: concave on both branches.
  \item Domain $\R$; root $x=2$; $f\to0^{+}$ as $x\to+\infty$
        (asymptote $y=0$) and $f\to-\infty$ as $x\to-\infty$;
        $f'=(3-x)\mathrm{e}^{-x}$ gives a maximum
        $(3,\mathrm{e}^{-3})$; $f''=(x-4)\mathrm{e}^{-x}$ gives an
        inflection $(4,2\mathrm{e}^{-4})$.
  \item Domain $(0,\infty)$; root $x=1$; $f\to0$ as $x\to0^{+}$ and
        $f\to\infty$ as $x\to\infty$; $f'=2x\ln x+x=x(2\ln x+1)$, zero at
        $x=\mathrm{e}^{-1/2}$, a minimum with value
        $-\tfrac{1}{2\mathrm{e}}$; $f''=2\ln x+3$, zero at
        $x=\mathrm{e}^{-3/2}$, an inflection.
  \item Domain $\R$; \emph{odd}; root $0$; horizontal asymptote $y=0$;
        $f'=\dfrac{4(4-x^{2})}{(x^{2}+4)^{2}}$ gives local minimum
        $(-2,-1)$ and local maximum $(2,1)$; inflections at $0$ and
        $\pm2\sqrt3$.
  \item Domain $[-2,2]$; \emph{even}; roots $\pm2$; range $[0,2]$; it is
        the upper semicircle of radius $2$, so concave throughout, maximum
        $(0,2)$, and the tangents at $x=\pm2$ are vertical. No asymptotes.
  \item Write $f=ax^{3}+bx^{2}+cx+d$. From the maximum at $(0,4)$:
        $d=4$ and $f'(0)=c=0$. From $f'(2)=12a+4b=0$ and $f(2)=8a+4b+4=0$:
        $b=-3a$ and $8a-12a+4=0$, so $a=1$, $b=-3$. Hence
        $f(x)=x^{3}-3x^{2}+4$.
  \item Increasing and concave with a ceiling: it rises steeply at first
        and flattens onto $y=5$ without reaching it. Example:
        $f(x)=5-\mathrm{e}^{-x}$.
  \item $f'=4x^{3}-12x^{2}=4x^{2}(x-3)$, critical at $x=0$ and $x=3$;
        $f''=12x^{2}-24x=12x(x-2)$, zero at $x=0$ and $x=2$ and changing
        sign at both, so inflections at $(0,0)$ and $(2,-16)$. At $x=0$ the
        second-derivative test fails because $f''(0)=0$; the sign table for
        $f'$ shows $4x^{2}(x-3)<0$ on both sides of $0$, so $x=0$ is no
        extremum --- it is the inflection.
  \item $f'=1-\dfrac{1}{x^{2}}$, zero at $x=\pm1$;
        $f''=\dfrac{2}{x^{3}}$, positive at $x=1$ (minimum, value $2$) and
        negative at $x=-1$ (maximum, value $-2$). The maximum value $-2$
        lies below the minimum value $2$ because the pole at $x=0$
        separates the two branches, so neither extremum is global.
  \item $f'=3x^{2}+a$ has two real roots exactly when $a<0$; then
        $x=\pm\sqrt{-a/3}$ and the sign of $f'$ changes at both, giving a
        maximum and a minimum. For $a\ge0$ the function is monotone
        increasing.
  \item $C''=6q-18=0$ at $q=3$: inflection, and $MC=3q^{2}-18q+30$ is
        minimal there with $MC=3$. $AC=q^{2}-9q+30=(q-4.5)^{2}+9.75$,
        minimal at $q=4.5$ with $AC=9.75$; and
        $MC(4.5)=60.75-81+30=9.75=AC$.
  \item $\pi=40q-2q^{2}-50$, so $\pi'=40-4q=0$ at $q=10$ with
        $\pi''=-4<0$: maximum profit $\pi(10)=400-200-50=150$. Break-even
        needs $\pi=0$: $2q^{2}-40q+50=0$, i.e. $q^{2}-20q+25=0$, so
        $q=10\pm\sqrt{75}$, about $1.34$ and $18.66$.
  \item $\pi=(100-2q)q-(q^{2}+20q)=80q-3q^{2}$, a downward parabola with
        zeros at $q=0$ and $q=80/3$, concave throughout. $\pi'=80-6q=0$ at
        $q=\tfrac{40}{3}\approx13.33$, with
        $\pi=\tfrac{1600}{3}\approx533.3$ and price
        $p=100-\tfrac{80}{3}\approx73.3$.
  \item The slope of the ray from the origin through $(q,C(q))$ is exactly
        $C(q)/q=AC(q)$. Minimising $AC$ therefore means finding the
        \emph{flattest} such ray, and the flattest ray that still meets the
        cost curve must touch it --- be tangent to it. At a tangency the
        ray's slope equals the curve's slope, i.e. $AC=MC$. That is the
        geometric content of $AC'=(MC-AC)/q$.
  \item $AP=Q/L=12L-L^{2}$, maximal where $12-2L=0$, i.e. $L=6$, with
        $AP=36$. And $MP=Q'=24L-3L^{2}$, so $MP(6)=144-108=36=AP$. The
        same tangency argument as for average cost applies, with the ray
        from the origin now touching the production curve.
\end{solutions}

\chapter{Functions of Several Variables: Surfaces and Level Curves}

A function of one variable is understood by looking at its graph. For a
function of two variables there are two pictures, and both are needed. The
\textbf{graph} is a surface in space: mountains, valleys, passes. The
\textbf{contour map} is what a cartographer draws instead --- the level
curves of that landscape, projected flat onto the plane. Economics almost
never draws the surface. It draws the contour map and gives the curves
names: isoquants, indifference curves, iso-profit lines.

The reason to spend a chapter on pictures before computing anything is that
most of multivariable calculus is easier to see than to calculate. Whether a
critical point is a maximum, a minimum or a saddle is visible at a glance in
a contour map. So is whether a function is convex, whether its upper contour
sets are convex, and where it is steep. And the geometry is not decoration:
the Lagrange method is the statement that two families of level curves touch
without crossing, and Kuhn--Tucker is a statement about which boundary of a
shaded region binds.

This chapter therefore does two things. It sets up domains, graphs, level
curves and the sub- and superlevel sets, with the vocabulary in both the
German and the English forms; and it drills reading properties off a picture
in both directions --- picture to property, and formula to picture.

\section{Core Rules}

\begin{rulebox}{Functions of two variables: domain, graph, range}
A function $f\colon D\to\R$ with $D\subseteq\R^{2}$ assigns a number
$f(x,y)$ to each point of a \textbf{region} $D$ in the plane.
\begin{itemize}[leftmargin=1.6em,itemsep=0.3ex]
  \item The \textbf{domain} $D$ is now a subset of the plane, not of the
        line, and the same three obstructions decide it: no zero
        denominator, no negative radicand under an even root, no
        non-positive argument in a logarithm. So
        $f(x,y)=\sqrt{9-x^{2}-y^{2}}$ has $D=\{(x,y):x^{2}+y^{2}\le9\}$, a
        closed disc, and $f(x,y)=\ln(x+y)$ has the open half-plane
        $x+y>0$.
  \item The \textbf{graph} is the surface
        $\{(x,y,f(x,y)) : (x,y)\in D\}$ in $\R^{3}$: above each point of
        the domain, a single height.
  \item The \textbf{range} is the set of heights attained, exactly as
        before.
\end{itemize}
Economic domains are usually the non-negative quadrant $K,L\ge0$, and often
its interior.
\end{rulebox}

\begin{rulebox}{Level curves and the contour map}
For $c\in\R$ the \textbf{level set} \de{(Niveaumenge, Niveaulinie)} is
\[
  N_{c}=\{(x,y)\in D : f(x,y)=c\} ,
\]
the set of points at height $c$. Drawing $N_{c}$ for several equally spaced
values of $c$ gives the \textbf{contour map}. To compute one, set
$f(x,y)=c$ and solve for $y$ --- the level set is a curve in the plane, and
often a familiar one.

Reading a contour map:
\begin{itemize}[leftmargin=1.6em,itemsep=0.3ex]
  \item \textbf{Contours cannot cross}, since a point has one height ---
        with one exception: at a \textbf{saddle} the level curve through
        the saddle point crosses itself.
  \item \textbf{Spacing is steepness.} Contours crowded together mean the
        surface rises fast; widely spaced contours mean gentle ground.
        Steepest ascent runs perpendicular to the contours.
  \item \textbf{Closed nested loops} shrinking to a point mark a
        \textbf{peak} or a \textbf{basin}; which one is decided by whether
        the labels increase or decrease inwards.
  \item A contour that \textbf{pinches and crosses} marks a saddle: two
        directions of ascent and two of descent alternate around it.
\end{itemize}
\end{rulebox}

\begin{rulebox}{Sub- and superlevel sets}
Two \emph{solid} sets accompany each level curve:
\[
  W(c)=\{(x,y): f(x,y)\le c\},
  \qquad
  B(c)=\{(x,y): f(x,y)\ge c\} .
\]
$W(c)$ is the \textbf{sublevel set} \de{(Subniveaumenge, untere
Konturmenge)}, $B(c)$ the \textbf{superlevel set} \de{(Superniveaumenge,
obere Konturmenge)}. Written through a point $P$ rather than a level,
$W(P)=W(f(P))$ and $B(P)=B(f(P))$: everything at most as high as $P$, and
everything at least as high. The same sets appear in the optimization
volume as $L_{c}$ and $U_{c}$ --- lower and upper contour set --- and in
consumer theory $B(P)$ is the \emph{better set} of the bundle $P$.

The link to curvature, drilled later and worth seeing now:
\begin{itemize}[leftmargin=1.6em,itemsep=0.3ex]
  \item $f$ convex $\Longrightarrow$ every $W(c)$ is convex;
  \item $f$ concave $\Longrightarrow$ every $B(c)$ is convex;
  \item every $B(c)$ convex is the \emph{definition} of
        \textbf{quasiconcave}, which is weaker: $f(x,y)=xy$ on the positive
        quadrant is quasiconcave but not concave.
\end{itemize}
\end{rulebox}

\begin{gallery}
\begin{tikzpicture}
  \begin{axis}[geoaxis, xmin=-3.2, xmax=3.2, ymin=-3.2, ymax=3.2,
      xtick={-2,2}, ytick={-2,2}, title={$W(4)$ for $f=x^{2}+y^{2}$}]
    \addplot[draw=none, name path=up, domain=-2:2, samples=201]{sqrt(4-x^2)};
    \addplot[draw=none, name path=lo, domain=-2:2, samples=201]{-sqrt(4-x^2)};
    \region{up}{lo}{}
    \addplot[domain=0:360, samples=181]({2*cos(x)},{2*sin(x)});
  \end{axis}
\end{tikzpicture}\hfill
\begin{tikzpicture}
  \begin{axis}[geoaxis, xmin=-0.4, xmax=8.4, ymin=-0.4, ymax=8.4,
      xtick={2,4,6,8}, ytick={2,4,6,8}, title={$B(4)$ for $f=xy$, $x,y>0$}]
    \addplot[draw=none, name path=hyp, domain=0.5:8, samples=201]{4/x};
    \addplot[draw=none, name path=top, domain=0.5:8, samples=2]{8.4};
    \region{hyp}{top}{}
    \addplot[domain=0.5:8, samples=201]{4/x};
  \end{axis}
\end{tikzpicture}
\gcap{A sublevel set of a convex function and a superlevel set of a
quasiconcave one. Both are convex sets; the second is the \emph{better set}
of a Cobb--Douglas bundle.}
\end{gallery}

\begin{rulebox}{The standard shapes}
Six surfaces cover most of what appears in exercises, and each has a
signature contour map.
\begin{itemize}[leftmargin=1.6em,itemsep=0.3ex]
  \item \textbf{Plane} $f=ax+by+c$: contours are \emph{parallel straight
        lines}, equally spaced. A budget line is one of them.
  \item \textbf{Paraboloid} $f=x^{2}+y^{2}$: a bowl; contours are
        concentric circles of radius $\sqrt{c}$, so they get
        \emph{closer together} outwards --- the bowl steepens.
  \item \textbf{Cone} $f=\sqrt{x^{2}+y^{2}}$: contours are concentric
        circles of radius $c$, \emph{equally spaced} --- constant slope.
  \item \textbf{Saddle} $f=x^{2}-y^{2}$: contours are hyperbolas; the level
        $c=0$ degenerates into the crossing pair $y=\pm x$.
  \item \textbf{Rotated saddle} $f=xy$: hyperbolas again, but with the
        \emph{axes} as asymptotes. Restricted to the positive quadrant this
        is the Cobb--Douglas picture.
  \item \textbf{Parabolic cylinder} $f=x^{2}$: the surface is a trough,
        constant in $y$; contours are \emph{pairs of vertical lines}
        $x=\pm\sqrt{c}$.
\end{itemize}
Circles alone do not identify the surface --- bowl, cone and dome all have
circular contours. The \emph{spacing} and the direction in which the labels
grow settle it.
\end{rulebox}

\begin{rulebox}{Reading curvature off the picture}
\begin{itemize}[leftmargin=1.6em,itemsep=0.3ex]
  \item \textbf{Convex} (bowl-shaped): every chord drawn between two points
        of the surface lies \emph{above} it; no part of the surface bulges
        over a chord. All sublevel sets are convex.
  \item \textbf{Concave} (dome-shaped): every chord lies \emph{below} the
        surface. All superlevel sets are convex.
  \item \textbf{Neither}: a saddle. Some chords lie above, some below.
  \item \textbf{Quasiconcave without being concave}: the superlevel sets
        are still convex --- nested ``better sets'' bulging towards the
        origin --- while the surface itself is not a dome. The Cobb--Douglas
        hill is the standard example.
  \item A \textbf{linear} function is both convex and concave, and both
        quasiconvex and quasiconcave.
\end{itemize}
\end{rulebox}

\begin{rulebox}{Economic readings}
\begin{itemize}[leftmargin=1.6em,itemsep=0.3ex]
  \item An \textbf{isoquant} is a level curve of a production function:
        all input combinations yielding the same output. An
        \textbf{indifference curve} is a level curve of utility.
  \item The slope of a level curve is the \textbf{marginal rate of
        substitution}, computed implicitly as $-f_{x}/f_{y}$
        (Chapter~5). Convex-to-the-origin level curves mean a diminishing
        rate of substitution.
  \item A \textbf{budget set} $\{p_{x}x+p_{y}y\le M,\ x,y\ge0\}$ is the
        sublevel set of a linear function intersected with the quadrant ---
        a triangle, and convex.
  \item The consumer's problem is to reach the highest indifference curve
        that still touches the budget triangle: the contour map plus the
        shaded set is the whole picture.
\end{itemize}
\end{rulebox}

\section{Worked Example}

\begin{worked}{Read the topography of $f(x,y)=x^{2}+4y^{2}$ and of $g(x,y)=x^{2}-y^{2}$: level curves, sub- and superlevel sets, and the surfaces}
\wstep{Step 1. A level curve is what you get by \emph{fixing the height}.}
Set $f(x,y)=c$ and solve for the curve in the plane:
\[
  x^{2}+4y^{2}=c .
\]
For $c>0$ divide by $c$ to reach the standard form of an ellipse:
\[
  \frac{x^{2}}{c}+\frac{y^{2}}{c/4}=1,
  \qquad\text{semi-axes } \sqrt c \text{ and } \tfrac{\sqrt c}{2} .
\]
So the contour map is a family of concentric ellipses, twice as wide as they
are tall. For $c=0$ the level ``curve'' collapses to the single point
$(0,0)$; for $c<0$ it is empty, since $f\ge0$ always.
\wtrap{Check which values of $c$ occur before drawing. Level sets can be
empty, a single point, or several disconnected pieces --- $x^{2}+4y^{2}=-1$
has no solutions, and a level set of $x^{2}-y^{2}$ has two branches.}

\wstep{Step 2. Spacing carries the steepness.}
The contours for $c=1,2,3,4$ have semi-axes $1,\sqrt2,\sqrt3,2$ in $x$. The
gaps shrink as you move out --- $1,\ 0.41,\ 0.32,\ 0.27$ --- so the surface
gets \emph{steeper} away from the origin. Contours that crowd together mean
a steep slope; contours far apart mean a plateau. That is the one skill a
contour map is for.

\wstep{Step 3. Sub- and superlevel sets.}
\[
  \{f\le c\}=\text{the filled ellipse: convex},
  \qquad
  \{f\ge c\}=\text{everything outside it: not convex.}
\]
$f$ is a convex function, so its \textbf{sub}level sets are convex. For a
concave function such as a utility $U$, it is the \textbf{super}level sets
$\{U\ge\bar u\}$ --- the ``at least as good as'' sets --- that are convex,
and that is precisely the convexity of preferences.
\wtrap{Convex function and convex set are different words. Match them
carefully: convex function $\Rightarrow$ sublevel sets convex; concave
function $\Rightarrow$ superlevel sets convex. The direction is chosen by
the inequality, not by the word ``convex''.}

\wstep{Step 4. The surface itself.}
$z=x^{2}+4y^{2}$ is an \textbf{elliptic paraboloid}: a bowl, with its
minimum at the origin, rising in every direction, steeper along $y$ than
along $x$ (the coefficient $4$ against $1$).

\wstep{Step 5. Now $g(x,y)=x^{2}-y^{2}$, where the sign flips.}
\[
  g=c:\quad x^{2}-y^{2}=c ,
\]
a family of \textbf{hyperbolas} --- opening left--right for $c>0$, up--down
for $c<0$. The level set $c=0$ degenerates into the \emph{pair of lines}
$y=\pm x$, which is where the two families change over.

\wstep{Step 6. Two pictures, side by side.}
\begin{center}
\begin{tikzpicture}
\begin{axis}[contouraxis, width=5.2cm, height=5.2cm, xmin=-2.6, xmax=2.6,
             ymin=-2.6, ymax=2.6, title={$x^{2}+4y^{2}=c$}]
  \foreach \c in {0.5,1,2,3,4} {
    \addplot[domain=0:360, samples=161, variable=\t]
      ({sqrt(\c)*cos(\t)}, {sqrt(\c)/2*sin(\t)});
  }
\end{axis}
\end{tikzpicture}
\hfill
\begin{tikzpicture}
\begin{axis}[contouraxis, width=5.2cm, height=5.2cm, xmin=-2.6, xmax=2.6,
             ymin=-2.6, ymax=2.6, title={$x^{2}-y^{2}=c$}]
  \foreach \c in {0.5,2} {
    \addplot[domain=-1.6:1.6, samples=121] ({sqrt(\c)*cosh(x)}, {sqrt(\c)*sinh(x)});
    \addplot[domain=-1.6:1.6, samples=121] ({-sqrt(\c)*cosh(x)}, {sqrt(\c)*sinh(x)});
    \addplot[domain=-1.6:1.6, samples=121] ({sqrt(\c)*sinh(x)}, {sqrt(\c)*cosh(x)});
    \addplot[domain=-1.6:1.6, samples=121] ({sqrt(\c)*sinh(x)}, {-sqrt(\c)*cosh(x)});
  }
  \addplot[gray!60, thin, samples=2, domain=-2.6:2.6, forget plot] {x};
  \addplot[gray!60, thin, samples=2, domain=-2.6:2.6, forget plot] {-x};
\end{axis}
\end{tikzpicture}
\end{center}
\begin{center}
\begin{tikzpicture}
\begin{axis}[surfaxis, title={$z=x^{2}+4y^{2}$}]
  \addplot3[surf, domain=-2:2, y domain=-1:1] {x^2+4*y^2};
\end{axis}
\end{tikzpicture}
\hfill
\begin{tikzpicture}
\begin{axis}[surfaxis, title={$z=x^{2}-y^{2}$}]
  \addplot3[surf, domain=-2:2, y domain=-2:2] {x^2-y^2};
\end{axis}
\end{tikzpicture}
\end{center}

\wstep{Step 7. Read the origin off each map.}
For $f$ the contours \emph{close around} the origin, so it is a minimum: go
any direction and you climb. For $g$ the contours \emph{cross} at the
origin, and the crossing is the signature of a \textbf{saddle}: along the
$x$-axis $g=x^{2}$ rises, along the $y$-axis $g=-y^{2}$ falls. The point is
a minimum in one direction and a maximum in another, so it is neither.
\wtrap{Closed loops shrinking to a point mean an extremum; crossing curves
mean a saddle. Checking only one direction --- for instance only along the
axes of $g$ where the answers disagree --- is exactly how a saddle gets
misreported as an optimum in Volume~7.}

\wstep{Step 8. The economic vocabulary.}
Level curves of a utility function are \textbf{indifference curves}; of a
production function, \textbf{isoquants}; of a cost function,
\textbf{isocost lines}. Their slope is the marginal rate of substitution,
computed in Chapter~5 by implicit differentiation, and the standard
``convex to the origin'' shape is exactly the convexity of the superlevel
sets in Step~3.
\end{worked}

\section{Practice Problems}

State the domain of each function as a region in the plane, in problems
1--8, and describe it in words.

\begin{problems}
  \item $f(x,y)=\sqrt{16-x^{2}-y^{2}}$
  \item $f(x,y)=\ln(x-y)$
  \item $f(x,y)=\dfrac{1}{x^{2}+y^{2}-1}$
  \item $f(x,y)=\dfrac{xy}{x-y}$
  \item $f(x,y)=\sqrt{x}+\sqrt{y}$
  \item $f(x,y)=\ln\bigl(xy\bigr)$
  \item $f(x,y)=\sqrt{x^{2}+y^{2}-4}$
  \item $f(K,L)=K^{0.4}L^{0.6}$, and state the economic domain too.
\end{problems}

In problems 9--20, give the level set $N_{c}$ for the stated level, solved
for $y$ where possible, and name the curve.

\begin{problems}[resume]
  \item $f(x,y)=x+2y$, \;$c=4$.
  \item $f(x,y)=x^{2}+y^{2}$, \;$c=9$.
  \item $f(x,y)=x^{2}+y^{2}$, \;$c=0$.
  \item $f(x,y)=x^{2}+4y^{2}$, \;$c=4$.
  \item $f(x,y)=xy$, \;$c=6$.
  \item $f(x,y)=x^{2}-y^{2}$, \;$c=0$.
  \item $f(x,y)=x^{2}-y^{2}$, \;$c=4$.
  \item $f(x,y)=y-x^{2}$, \;$c=2$.
  \item $f(x,y)=\sqrt{x^{2}+y^{2}}$, \;$c=3$.
  \item $f(x,y)=\mathrm{e}^{x+y}$, \;$c=1$.
  \item $f(K,L)=\sqrt{KL}$, \;$c=10$.
  \item $f(x,y)=2x+3y$, \;$c=12$; and say what this curve is called when
        $f$ is a cost.
\end{problems}

Problems 21--28 refer to the contour map below, which shows the level
curves of $f(x,y)=x^{3}-3x+y^{2}$ at the levels $c=-1,0,2,4$, together with
six marked points.

\begin{gallery}
\hspace*{\fill}%
\begin{tikzpicture}
  \begin{axis}[contouraxis, width=8.2cm, height=8.2cm,
      xmin=-2.6, xmax=2.6, ymin=-3.2, ymax=3.2,
      xtick={-2,-1,1,2}, ytick={-2,-1,1,2}]
    % c = 4
    \addplot[domain=-2.2:2.196, samples=201]{ sqrt(4-x^3+3*x)};
    \addplot[domain=-2.2:2.196, samples=201]{-sqrt(4-x^3+3*x)};
    % c = 2 (through the saddle at (-1,0))
    \addplot[domain=-2.2:2, samples=201]{ sqrt(2-x^3+3*x)};
    \addplot[domain=-2.2:2, samples=201]{-sqrt(2-x^3+3*x)};
    % c = 0, two pieces
    \addplot[domain=-2.2:-1.7321, samples=101]{ sqrt(-x^3+3*x)};
    \addplot[domain=-2.2:-1.7321, samples=101]{-sqrt(-x^3+3*x)};
    \addplot[domain=0:1.7321, samples=101]{ sqrt(-x^3+3*x)};
    \addplot[domain=0:1.7321, samples=101]{-sqrt(-x^3+3*x)};
    % c = -1, two pieces
    \addplot[domain=-2.2:-1.8794, samples=61]{ sqrt(-1-x^3+3*x)};
    \addplot[domain=-2.2:-1.8794, samples=61]{-sqrt(-1-x^3+3*x)};
    \addplot[domain=0.3473:1.5321, samples=101]{ sqrt(-1-x^3+3*x)};
    \addplot[domain=0.3473:1.5321, samples=101]{-sqrt(-1-x^3+3*x)};
    % marked points
    \addplot[only marks, mark=*, mark size=1.4pt, black, forget plot]
      coordinates {(0,0) (1,0) (-1,0) (2,1) (0,2) (-2,1)};
    \node[font=\tiny, anchor=north east] at (axis cs:-0.06,-0.12) {$P_{1}$};
    \node[font=\tiny, anchor=north west] at (axis cs:1.08,-0.12) {$P_{2}$};
    \node[font=\tiny, anchor=south east] at (axis cs:-1.12,0.12) {$P_{3}$};
    \node[font=\tiny, anchor=south west] at (axis cs:2.05,1.05) {$P_{4}$};
    \node[font=\tiny, anchor=south west] at (axis cs:0.05,2.05) {$P_{5}$};
    \node[font=\tiny, anchor=south east] at (axis cs:-2.05,1.05) {$P_{6}$};
  \end{axis}
\end{tikzpicture}%
\hspace*{\fill}
\gcap{Level curves of $f(x,y)=x^{3}-3x+y^{2}$ at $c=-1,0,2,4$.}
\end{gallery}

\begin{problems}[resume]
  \item Give $f(P_{i})$ for all six marked points
        $P_{1}=(0,0)$, $P_{2}=(1,0)$, $P_{3}=(-1,0)$, $P_{4}=(2,1)$,
        $P_{5}=(0,2)$, $P_{6}=(-2,1)$.
  \item Which of the six points lie on one and the same level curve?
  \item One of the drawn contours crosses itself. At which point, and what
        does that tell you about the surface there?
  \item One of the drawn contours has shrunk to a single point. Which
        level, which point, and what feature of the surface is it?
  \item From the picture alone, name the two critical points of $f$ and
        classify each as a local maximum, a local minimum, or a saddle.
  \item Confirm the previous answer by computation: solve
        $f_{x}=f_{y}=0$.
  \item Is the sublevel set $W(-1)$ connected? Answer from the picture and
        explain what the two pieces mean.
  \item Moving away from $P_{2}$, in which direction does $f$ increase
        fastest? Answer qualitatively from the contours, then check with
        $\grad f$ at a nearby point, say $(1.5,0.5)$.
\end{problems}

\begin{problems}[resume]
  \item Match each surface (a)--(f) below to its contour map (1)--(6).
  \item For each of the six functions in the previous problem, say whether
        it is convex, concave, both, or neither on all of $\R^{2}$.
\end{problems}

\begin{gallery}
\begin{tikzpicture}
  \begin{axis}[surfaxis, title={(a)}]
    \addplot3[surf, domain=-2:2, y domain=-2:2]{x+y};
  \end{axis}
\end{tikzpicture}\hfill
\begin{tikzpicture}
  \begin{axis}[surfaxis, title={(b)}, z post scale=1.2]
    \addplot3[surf, domain=-2:2, y domain=-2:2]{x^2+y^2};
  \end{axis}
\end{tikzpicture}\hfill
\begin{tikzpicture}
  \begin{axis}[surfaxis, title={(c)}, z post scale=1.2]
    \addplot3[surf, domain=-2:2, y domain=-2:2]{sqrt(x^2+y^2)};
  \end{axis}
\end{tikzpicture}
\end{gallery}

\begin{gallery}
\begin{tikzpicture}
  \begin{axis}[surfaxis, title={(d)}, z post scale=1.2]
    \addplot3[surf, domain=-2:2, y domain=-2:2]{x^2-y^2};
  \end{axis}
\end{tikzpicture}\hfill
\begin{tikzpicture}
  \begin{axis}[surfaxis, title={(e)}, z post scale=1.2, view={105}{32}]
    \addplot3[surf, domain=-2:2, y domain=-2:2]{x*y};
  \end{axis}
\end{tikzpicture}\hfill
\begin{tikzpicture}
  \begin{axis}[surfaxis, title={(f)}, z post scale=1.2]
    \addplot3[surf, domain=-2:2, y domain=-2:2]{x^2};
  \end{axis}
\end{tikzpicture}
\end{gallery}

\begin{gallery}
\begin{tikzpicture}
  \begin{axis}[contouraxis, xmin=-2.4, xmax=2.4, ymin=-2.4, ymax=2.4,
      xtick={-2,2}, ytick={-2,2}, title={(1)}]
    \foreach \c in {-2,-1,0,1,2} { \addplot[domain=-2.3:2.3]{\c-x}; }
  \end{axis}
\end{tikzpicture}\hfill
\begin{tikzpicture}
  \begin{axis}[contouraxis, xmin=-2.4, xmax=2.4, ymin=-2.4, ymax=2.4,
      xtick={-2,2}, ytick={-2,2}, title={(2)}]
    \foreach \r in {0.25,0.5,0.75,1,1.25,1.5,1.75,2} {
      \addplot[domain=0:360, samples=121]({\r*cos(x)},{\r*sin(x)}); }
  \end{axis}
\end{tikzpicture}\hfill
\begin{tikzpicture}
  \begin{axis}[contouraxis, xmin=-2.4, xmax=2.4, ymin=-2.4, ymax=2.4,
      xtick={-2,2}, ytick={-2,2}, title={(3)}]
    \foreach \r in {0.707,1,1.225,1.414,1.581,1.732,1.871,2} {
      \addplot[domain=0:360, samples=121]({\r*cos(x)},{\r*sin(x)}); }
  \end{axis}
\end{tikzpicture}
\end{gallery}

\begin{gallery}
\begin{tikzpicture}
  \begin{axis}[contouraxis, xmin=-2.4, xmax=2.4, ymin=-2.4, ymax=2.4,
      xtick={-2,2}, ytick={-2,2}, title={(4)}]
    \addplot[domain=-2.3:2.3]{x};
    \addplot[domain=-2.3:2.3]{-x};
    \foreach \s in {1,-1} {
      \addplot[domain=-1.4:1.4, samples=81]({\s*cosh(x)},{sinh(x)});
      \addplot[domain=-1.4:1.4, samples=81]({sinh(x)},{\s*cosh(x)}); }
  \end{axis}
\end{tikzpicture}\hfill
\begin{tikzpicture}
  \begin{axis}[contouraxis, xmin=-2.4, xmax=2.4, ymin=-2.4, ymax=2.4,
      xtick={-2,2}, ytick={-2,2}, title={(5)}]
    \foreach \c in {1,2} {
      \addplot[domain=0.42:2.3, samples=101]{\c/x};
      \addplot[domain=-2.3:-0.42, samples=101]{\c/x};
      \addplot[domain=0.42:2.3, samples=101]{-\c/x};
      \addplot[domain=-2.3:-0.42, samples=101]{-\c/x}; }
  \end{axis}
\end{tikzpicture}\hfill
\begin{tikzpicture}
  \begin{axis}[contouraxis, xmin=-2.4, xmax=2.4, ymin=-2.4, ymax=2.4,
      xtick={-2,2}, ytick={-2,2}, title={(6)}]
    \foreach \v in {-2,-1.414,-1,1,1.414,2} {
      \addplot[samples=2] coordinates {(\v,-2.3) (\v,2.3)}; }
  \end{axis}
\end{tikzpicture}
\end{gallery}

In problems 31--38, describe the stated sub- or superlevel set as a region,
and say whether it is convex.

\begin{problems}[resume]
  \item $W(4)$ for $f(x,y)=x^{2}+y^{2}$.
  \item $B(4)$ for $f(x,y)=x^{2}+y^{2}$.
  \item $W(6)$ for $f(x,y)=2x+3y$.
  \item $B(4)$ for $f(x,y)=xy$ restricted to $x,y>0$.
  \item $W(0)$ for $f(x,y)=x^{2}-y^{2}$.
  \item $B(1)$ for $f(x,y)=\sqrt{x^{2}+y^{2}}$.
  \item $W(P)$ for $f(x,y)=x^{2}+y^{2}$ with $P=(3,4)$.
  \item $B(P)$ for $f(K,L)=\sqrt{KL}$ with $P=(4,9)$, on $K,L>0$.
\end{problems}

\begin{problems}[resume]
  \item Sketch the level sets $W(2)$, $B(2)$, $W(-1)$ and $B(-1)$ of the
        one-variable function $g(x)=x^{3}-3x$ as subsets of the $x$-axis,
        and say which of the four are convex sets.
  \item Explain why for a function of one variable a level set
        $W(c)$ is convex exactly when it is an interval.
  \item A firm's isoquants are the curves $KL=\bar Q$. Sketch three of them
        for $\bar Q=1,4,9$ on $K,L>0$, and say whether the isoquants get
        closer together or further apart as output rises along the ray
        $K=L$.
  \item For the budget set $\{2x+5y\le100,\ x,y\ge0\}$: give the two
        intercepts, sketch the set, and say why it is convex.
  \item The indifference curves of $U(x,y)=xy$ and the budget line
        $2x+5y=100$ are drawn on the same axes. Describe, in words and
        without computing, the geometric condition that locates the optimal
        bundle.
  \item Two contour maps show circles as level curves. In the first the
        circles are equally spaced; in the second they crowd together
        outwards. Which is the cone and which the paraboloid, and why?
  \item A contour map has no closed loops and no crossings anywhere, only
        parallel straight lines. What can be said about the function?
\end{problems}

\section{Solutions}

\begin{solutions}
  \item $16-x^{2}-y^{2}\ge0\iff x^{2}+y^{2}\le16$: the closed disc of
        radius $4$ about the origin.
  \item $x-y>0$: the open half-plane strictly below the line $y=x$.
  \item $x^{2}+y^{2}\neq1$: the whole plane except the unit circle.
  \item $x\neq y$: the plane without the line $y=x$.
  \item $x\ge0$ and $y\ge0$: the closed non-negative quadrant.
  \item $xy>0$: the open first and third quadrants (the axes excluded).
  \item $x^{2}+y^{2}\ge4$: everything outside or on the circle of radius
        $2$ --- the complement of an open disc.
  \item The exponents are not integers, so $K\ge0$ and $L\ge0$: the closed
        non-negative quadrant. Economically the same set, usually taken
        open ($K,L>0$) so that marginal products exist.
  \item $y=\dfrac{4-x}{2}=2-\tfrac12x$: a straight line of slope
        $-\tfrac12$. All level sets are parallel lines.
  \item $x^{2}+y^{2}=9$, i.e. $y=\pm\sqrt{9-x^{2}}$: the circle of radius
        $3$.
  \item $x^{2}+y^{2}=0$ forces $x=y=0$: the level set is the single point
        $(0,0)$, the bottom of the bowl.
  \item $\dfrac{x^{2}}{4}+y^{2}=1$, so $y=\pm\sqrt{1-x^{2}/4}$: an ellipse
        with semi-axes $2$ and $1$.
  \item $y=6/x$: a hyperbola in the first and third quadrants, with both
        axes as asymptotes.
  \item $x^{2}=y^{2}\iff y=\pm x$: the \emph{two crossing lines} through
        the origin. This degenerate level is the signature of a saddle.
  \item $x^{2}-y^{2}=4$, so $y=\pm\sqrt{x^{2}-4}$, defined for
        $|x|\ge2$: a hyperbola opening left and right.
  \item $y=x^{2}+2$: a parabola, the standard parabola shifted up by $2$.
        All level sets are vertical translates of it.
  \item $\sqrt{x^{2}+y^{2}}=3\iff x^{2}+y^{2}=9$: the circle of radius
        $3$. Note that for the cone the radius equals $c$, not
        $\sqrt{c}$.
  \item $\mathrm{e}^{x+y}=1\iff x+y=0$, so $y=-x$: a straight line. Every
        level set of $\mathrm{e}^{x+y}$ is a line $x+y=\ln c$, defined only
        for $c>0$.
  \item $\sqrt{KL}=10\iff KL=100\iff L=100/K$: a hyperbola --- the
        isoquant for output $10$.
  \item $y=\dfrac{12-2x}{3}=4-\tfrac23x$: a straight line. For a cost
        function this is an \textbf{isocost line}; the analogous level
        curve of a budget constraint is the budget line.
  \item $f(x,y)=x^{3}-3x+y^{2}$ gives
        $f(P_{1})=0$, \;$f(P_{2})=1-3+0=-2$, \;$f(P_{3})=-1+3+0=2$,
        \;$f(P_{4})=8-6+1=3$, \;$f(P_{5})=0-0+4=4$,
        \;$f(P_{6})=-8+6+1=-1$.
  \item None of them share a level: the six values
        $0,-2,2,3,4,-1$ are distinct. $P_{3}$ lies on the drawn contour
        $c=2$, $P_{5}$ on $c=4$ and $P_{6}$ on $c=-1$; $P_{1}$, $P_{2}$ and
        $P_{4}$ lie between drawn contours.
  \item The contour $c=2$ crosses itself at $P_{3}=(-1,0)$. A
        self-crossing level curve marks a \textbf{saddle point}: from
        $(-1,0)$ the surface rises in two opposite directions and falls in
        the two others.
  \item The level $c=-2$ is a single point, $(1,0)$. Nested contours
        shrinking to a point mark a \textbf{local extremum}; since the
        levels decrease inwards, it is a local \emph{minimum}, with value
        $-2$.
  \item From the picture: a local minimum at $(1,0)$ (contours closing in
        on it) and a saddle at $(-1,0)$ (the crossing contour). There is no
        local maximum --- no contours close in on a point with the labels
        increasing inwards.
  \item $f_{x}=3x^{2}-3=0\Rightarrow x=\pm1$; $f_{y}=2y=0\Rightarrow y=0$.
        Critical points $(1,0)$ and $(-1,0)$. With $f_{xx}=6x$,
        $f_{yy}=2$, $f_{xy}=0$: at $(1,0)$ the determinant is $12>0$ with
        $f_{xx}>0$, a local minimum; at $(-1,0)$ it is $-12<0$, a saddle.
  \item No: $W(-1)$ has two pieces --- a small oval around the minimum
        $(1,0)$ and an unbounded piece on the far left. A sublevel set of a
        non-convex function may fall apart into components, which is
        exactly what convexity of $f$ would forbid.
  \item Away from $P_{2}=(1,0)$ the contours are closed ovals around it, so
        $f$ increases in \emph{every} direction, fastest where the ovals
        are most tightly packed --- to the right. Checking,
        $\grad f=(3x^{2}-3,\,2y)$, so at $(1.5,0.5)$,
        $\grad f=(3.75,\,1)$: pointing right and slightly up, and
        perpendicular to the contour there.
  \item (a)$\to$(1) plane, parallel lines. \;(b)$\to$(3) paraboloid:
        circles of radius $\sqrt{c}$, crowding outwards.
        \;(c)$\to$(2) cone: circles of radius $c$, equally spaced.
        \;(d)$\to$(4) saddle $x^{2}-y^{2}$: hyperbolas with the diagonals
        $y=\pm x$ as the degenerate level. \;(e)$\to$(5) $xy$: hyperbolas
        with the \emph{axes} as asymptotes. \;(f)$\to$(6) parabolic
        cylinder: pairs of vertical lines.
  \item (a) linear, hence \textbf{both} convex and concave. (b) convex (a
        bowl). (c) convex (a cone opening upwards; not differentiable at
        the origin, but still convex). (d) neither --- a saddle.
        (e) neither on $\R^{2}$; on the positive quadrant it is
        quasiconcave. (f) convex: $x^{2}$ is convex and adding no
        $y$-dependence preserves it.
  \item The closed disc $x^{2}+y^{2}\le4$ of radius $2$. Convex.
  \item The complement of the open disc, $x^{2}+y^{2}\ge4$: everything
        outside the circle. \textbf{Not} convex --- the segment joining
        $(-2,0)$ to $(2,0)$ passes through the origin, which is not in the
        set.
  \item The half-plane $2x+3y\le6$, bounded by the line through $(3,0)$ and
        $(0,2)$. Convex: every half-plane is.
  \item $\{(x,y): x,y>0,\ xy\ge4\}$, the region on the far side of the
        hyperbola $y=4/x$. Convex --- this is precisely the statement that
        $xy$ is quasiconcave on the quadrant.
  \item $x^{2}-y^{2}\le0\iff|y|\ge|x|$: the two opposite ``vertical'' sectors
        between the lines $y=\pm x$. Not convex: $(0,1)$ and $(0,-1)$ lie
        in it, but the midpoint $(0,0)$ is on the boundary and points just
        off the segment leave the set --- more simply, the set is not even
        connected in the way a convex set must be, since joining
        $(0,1)$ to $(1,2)$ stays inside but joining $(0,1)$ to $(0,-1)$
        passes through the pinch at the origin only.
  \item $\sqrt{x^{2}+y^{2}}\ge1$: everything outside or on the unit circle.
        Not convex, for the same reason as problem 32.
  \item $f(P)=9+16=25$, so $W(P)=\{x^{2}+y^{2}\le25\}$, the closed disc of
        radius $5$ through $P$. Convex.
  \item $f(P)=\sqrt{36}=6$, so
        $B(P)=\{(K,L):K,L>0,\ \sqrt{KL}\ge6\}=\{KL\ge36\}$: the region
        beyond the isoquant $L=36/K$. Convex.
  \item $g(x)=x^{3}-3x$ has $g(2)=2$, $g(-1)=2$, $g(1)=-2$, $g(-2)=-2$.
        $W(2)=\{x\le2\}$ minus nothing --- solving $x^{3}-3x\le2$ gives
        $(x+1)^{2}(x-2)\le0$, i.e. $x\le2$: the ray $(-\infty,2]$,
        \textbf{convex}. $B(2)=[2,\infty)\cup\{-1\}$: the ray from $2$
        \emph{plus the isolated point} $x=-1$, \textbf{not convex}.
        $W(-1)$: solving $x^{3}-3x\le-1$ gives
        $(-\infty,-1.879]\cup[0.347,1.532]$, two intervals, \textbf{not
        convex}. $B(-1)=[-1.879,0.347]\cup[1.532,\infty)$, also two
        pieces, \textbf{not convex}.
  \item In $\R^{1}$ a set is convex exactly when it contains the whole
        segment between any two of its points, and a subset of the line
        with that property is precisely an interval (possibly unbounded or
        empty). So ``convex'' and ``interval'' are the same statement one
        dimension down --- which is why the one-variable picture is the
        right place to first meet the definition.
  \item $L=\bar Q/K$: three hyperbola branches, $L=1/K$, $L=4/K$ and
        $L=9/K$. Along the ray $K=L$ they cut it at
        $K=1,2,3$: the isoquants for equally spaced \emph{square roots} of
        output, so for equally spaced output levels they move
        \emph{further apart}. Rising output requires ever larger input
        increases --- diminishing returns along the ray.
  \item Intercepts $(50,0)$ and $(0,20)$. The set is the triangle with
        vertices $(0,0)$, $(50,0)$, $(0,20)$. It is convex because it is an
        intersection of three half-planes ($2x+5y\le100$, $x\ge0$,
        $y\ge0$), and any intersection of convex sets is convex.
  \item The optimum is where the budget line is \textbf{tangent} to an
        indifference curve: the highest curve of the family that still
        touches the triangle touches it at exactly one point, and there the
        two curves have the same slope. Crossing would mean a higher curve
        is still reachable; missing the set entirely means the bundle is
        unaffordable.
  \item Equally spaced circles are the \textbf{cone}
        $\sqrt{x^{2}+y^{2}}$, whose level $c$ is a circle of radius $c$ ---
        constant slope, so equal height steps need equal horizontal steps.
        Circles crowding outwards are the \textbf{paraboloid}
        $x^{2}+y^{2}$, level $c$ having radius $\sqrt{c}$: the surface gets
        steeper, so equal height steps need ever smaller horizontal steps.
  \item Parallel straight contours, equally spaced or not, mean the
        function depends on the two variables only through one linear
        combination: $f(x,y)=h(ax+by)$ for some one-variable $h$. If the
        spacing is also constant, $h$ is linear and the surface is a plane.
        There are no critical points of extremum type anywhere.
\end{solutions}

\section{Additional Problems}

\begin{problems}
  \item Domain of $f(x,y)=\dfrac{\sqrt{y-x^{2}}}{x}$, as a region.
  \item Domain of $f(x,y)=\ln(4-x^{2}-y^{2})$, as a region.
  \item Domain and range of $f(x,y)=\mathrm{e}^{-(x^{2}+y^{2})}$.
  \item Level set of $f(x,y)=x^{2}+y^{2}-4x$ at $c=0$; name and locate the
        curve.
  \item Level set of $f(x,y)=x^{2}+y^{2}-4x$ at $c=-4$.
  \item Level set of $f(x,y)=y\,\mathrm{e}^{x}$ at $c=1$.
  \item Level set of $f(x,y)=\dfrac{y}{x}$ at $c=2$, and why the origin is
        excluded.
  \item Level sets of $f(x,y)=\max(x,y)$ at $c=1$: sketch and describe.
  \item For $f(x,y)=x^{2}+4y^{2}$, sketch the contours $c=4,16,36$ and say
        whether they crowd inwards or outwards.
  \item For $f(x,y)=x-y^{2}$, give the contours at $c=-1,0,1$ and describe
        the family.
  \item Give a function whose level curves are concentric circles that
        crowd \emph{inwards} towards the origin.
  \item Give a function of two variables whose level sets are all straight
        lines through the origin.
  \item $W(1)$ and $B(1)$ for $f(x,y)=|x|+|y|$: describe both and say which
        is convex.
  \item $W(0)$ for $f(x,y)=xy$ on all of $\R^{2}$: describe it and decide
        convexity.
  \item $B(2)$ for $f(x,y)=x+y$: describe it and decide convexity.
  \item $W(4)$ for $f(x,y)=x^{2}$: describe it and decide convexity.
  \item For $f(x,y)=x^{2}+y^{2}$ and $P=(1,2)$, give $W(P)$ and $B(P)$ and
        the level curve through $P$.
  \item A function has convex sublevel sets \emph{and} convex superlevel
        sets at every level. What kind of function can it be? Give an
        example.
  \item Sketch the surface $f(x,y)=4-x^{2}-y^{2}$ in words, give its
        contours, and classify it as convex or concave.
  \item Two functions have exactly the same family of level curves but
        different surfaces. Give such a pair and explain how the contour
        maps can nevertheless be told apart.
  \item Classify $f(x,y)=x^{2}+y^{2}$, $g(x,y)=-x^{2}-y^{2}$,
        $h(x,y)=x^{2}-y^{2}$ and $k(x,y)=3x-2y$ as convex, concave, both or
        neither.
  \item Which of the four functions in the previous problem are
        quasiconcave?
  \item The isoquants of a production function are straight lines with
        slope $-1$. What does that say about substitutability of the two
        inputs, and give such a production function.
  \item An indifference map consists of circles centred at $(4,4)$, with
        utility \emph{increasing} inwards. Interpret the point $(4,4)$
        economically and say why such preferences are unusual.
  \item On the contour map of problem 21, a walker starts at $P_{6}=(-2,1)$
        and always walks perpendicular to the contours, downhill. Describe
        where they end up and why the answer is not the minimum at
        $(1,0)$.
\end{problems}

\section{Additional Solutions}

\begin{solutions}
  \item $y-x^{2}\ge0$ and $x\neq0$: the region on or above the parabola
        $y=x^{2}$, with the $y$-axis removed.
  \item $4-x^{2}-y^{2}>0\iff x^{2}+y^{2}<4$: the \emph{open} disc of
        radius $2$ (boundary excluded, because the logarithm needs a
        strictly positive argument).
  \item $D=\R^{2}$. The exponent $-(x^{2}+y^{2})$ runs over $(-\infty,0]$,
        so the range is $(0,1]$ --- a bell-shaped hill with peak $1$ at the
        origin.
  \item $x^{2}-4x+y^{2}=0\iff(x-2)^{2}+y^{2}=4$: the circle of radius $2$
        centred at $(2,0)$, passing through the origin.
  \item $(x-2)^{2}+y^{2}=0$: the single point $(2,0)$ --- the minimum of
        the shifted bowl, with value $-4$.
  \item $y\,\mathrm{e}^{x}=1\iff y=\mathrm{e}^{-x}$: a decaying
        exponential curve.
  \item $y/x=2\iff y=2x$ with $x\neq0$: the line of slope $2$ through the
        origin, \emph{minus} the origin itself, since $f$ is undefined
        where $x=0$.
  \item $\max(x,y)=1$ is the union of the ray $\{x=1,\ y\le1\}$ and the ray
        $\{y=1,\ x\le1\}$: an L-shaped corner opening down-left, with its
        vertex at $(1,1)$.
  \item Ellipses $\dfrac{x^{2}}{c}+\dfrac{4y^{2}}{c}=1$ with semi-axes
        $\sqrt{c}$ and $\sqrt{c}/2$: for $c=4,16,36$ the semi-axes are
        $2,4,6$ and $1,2,3$. Equally spaced, in fact --- but for
        \emph{equally spaced} levels the contours would crowd outwards, as
        for any quadratic bowl.
  \item $x=y^{2}+c$: parabolas opening to the right, the whole family being
        horizontal translates of $x=y^{2}$ by $c$. For $c=-1,0,1$ the
        vertices sit at $(-1,0)$, $(0,0)$, $(1,0)$.
  \item Any function whose height grows ever more slowly with the radius,
        for instance $f(x,y)=\sqrt[4]{x^{2}+y^{2}}$ or
        $f(x,y)=\ln(x^{2}+y^{2})$: the level $c$ has radius growing faster
        than linearly in $c$, so equal height steps take ever larger
        radial steps and the contours thin out --- equivalently they crowd
        towards the origin.
  \item $f(x,y)=y/x$ (defined for $x\neq0$): the level $c$ is the punctured
        line $y=cx$. Any function of the ratio $y/x$ alone has this
        property; such functions are exactly the homogeneous functions of
        degree $0$.
  \item $W(1)=\{|x|+|y|\le1\}$ is the closed diamond (square rotated
        $45^{\circ}$) with vertices $(\pm1,0)$, $(0,\pm1)$ --- convex.
        $B(1)$ is everything outside it --- not convex.
  \item $xy\le0$ is the union of the second and fourth quadrants. Not
        convex: $(-1,1)$ and $(1,-1)$ both lie in it, but their midpoint
        $(0,0)$ does lie in it too --- take instead $(-1,2)$ and $(2,-1)$,
        whose midpoint $(0.5,0.5)$ has $xy>0$ and is outside.
  \item $x+y\ge2$: a closed half-plane, convex.
  \item $x^{2}\le4\iff-2\le x\le2$: the vertical strip between the lines
        $x=-2$ and $x=2$, convex (an intersection of two half-planes).
  \item $f(P)=1+4=5$: the level curve is the circle of radius $\sqrt5$,
        $W(P)$ the closed disc of radius $\sqrt5$, and $B(P)$ everything
        outside or on that circle.
  \item If every sublevel \emph{and} every superlevel set is convex, the
        function is both quasiconvex and quasiconcave. The level sets must
        then be ``flat'': any function of the form $h(ax+by)$ with $h$
        monotone qualifies, the simplest being a linear function
        $f(x,y)=ax+by$.
  \item A dome: the paraboloid turned upside down, with its peak of height
        $4$ above the origin and falling away in every direction. Contours
        are circles $x^{2}+y^{2}=4-c$ of radius $\sqrt{4-c}$, shrinking as
        $c$ rises to $4$. The function is \textbf{concave}.
  \item $f(x,y)=x^{2}+y^{2}$ and $g(x,y)=\sqrt{x^{2}+y^{2}}$ both have
        concentric circles as level sets. The maps differ in the
        \emph{labels}: for equally spaced levels $f$'s contours crowd
        outwards while $g$'s stay equally spaced. A contour map without its
        level labels does not determine the function.
  \item $f$ convex; $g$ concave; $h$ neither (a saddle); $k$ linear, hence
        both.
  \item $f$: quasiconvex, not quasiconcave (its superlevel sets are
        complements of discs). $g$: concave, hence quasiconcave. $h$:
        neither. $k$: linear, hence both quasiconvex and quasiconcave.
  \item Straight isoquants of slope $-1$ mean the marginal rate of
        technical substitution is constant at $1$: the inputs are
        \textbf{perfect substitutes}, one unit of either always replacing
        one unit of the other. The production function is
        $Q(K,L)=K+L$ (or any increasing function of $K+L$).
  \item $(4,4)$ is a \textbf{bliss point}: a most-preferred bundle, beyond
        which more of either good makes the consumer worse off. Standard
        preferences are monotone --- more is better --- so utility should
        increase without bound in the direction of both goods; circular
        indifference curves with utility rising inwards describe satiation
        instead.
  \item Walking perpendicular to the contours downhill from $P_{6}=(-2,1)$
        means following $-\grad f$. At $P_{6}$,
        $\grad f=(3\cdot4-3,\,2)=(9,2)$, so the walker moves left and
        slightly down --- into the region $x<-2$, where the contours of
        \emph{lower} level continue without bound. The function
        $x^{3}-3x+y^{2}$ tends to $-\infty$ as $x\to-\infty$, so the
        descent never ends: the walker never reaches the local minimum at
        $(1,0)$, because gradient descent finds a \emph{local} basin only
        if it starts in one.
\end{solutions}

\chapter{First-Order Partial Derivatives and the Gradient}

A function of one variable has a graph in the plane and a derivative that is
the slope of a tangent line. A function of two variables has a graph that is
a \emph{surface} in space, and there is no single slope any more: standing
at a point of a hillside, the ground rises steeply in one direction, falls
in another, and is level in a third. The first task of multivariable
calculus is to make that observation quantitative.

The device is to look one direction at a time. Freeze $y$ at a value and let
only $x$ move: the surface is cut by a vertical plane, the cut is an
ordinary curve, and the slope of \emph{that} curve is the partial derivative
$\partial f/\partial x$. This is the \emph{ceteris paribus} principle of
economics, and it is why partial derivatives are the natural language for
marginal products and marginal utilities.

Collecting the partial derivatives into one vector gives the gradient. It
turns out to carry all first-order information at once: it points in the
direction of steepest ascent, its length is the steepest slope available,
and it stands perpendicular to the level curve through the point. Those
three facts are proved below rather than asserted --- the third one in
particular is what the Lagrange method rests on several chapters later.

\section{Core Rules}

\begin{rulebox}{Topological vocabulary of domains}
\begin{itemize}[leftmargin=1.6em,itemsep=0.3ex]
  \item The \emph{open ball} (or \emph{neighbourhood}) of radius $r>0$ around $\mathbf{x}_0\in\R^n$ is
        $B_r(\mathbf{x}_0)=\{\mathbf{x}\in\R^n :
        \|\mathbf{x}-\mathbf{x}_0\|<r\}$.
  \item $\mathbf{x}_0\in M$ is an \emph{interior point} of $M$ if some open ball around
        $\mathbf{x}_0$ lies entirely in $M$.
  \item $M$ is \emph{open} if \emph{every} point of $M$ is
        an interior point.
  \item $\mathbf{x}_0$ is a \emph{boundary point} of
        $M$ if every ball around $\mathbf{x}_0$ contains points of $M$
        \emph{and} points of its complement.
  \item $M$ is \emph{closed} if it contains all
        its boundary points; equivalently, the complement
        $\R^n\setminus M$ is open.
  \item $M$ is \emph{bounded} if
        $M\subseteq B_r(\mathbf{0})$ for some radius $r$.
  \item $M\subseteq\R^n$ is \emph{compact} exactly when
        $M$ is closed \emph{and} bounded.
\end{itemize}
Economic examples: the open positive orthant $\{(x,y): x>0,\ y>0\}$
(the natural Cobb--Douglas domain) is open and unbounded; a budget set
$\{\mathbf{x}\in\R^n : \mathbf{p}\cdot\mathbf{x}\le m,\
\mathbf{x}\ge\mathbf{0}\}$ is closed and bounded, hence compact.
\end{rulebox}

\begin{rulebox}{Why it matters}
Differentiability (and hence every partial derivative below) is
defined at \emph{interior} points of the domain, where $f$ can be
probed in all directions. And by the \emph{Weierstrass theorem}, a
continuous function on a compact set attains its maximum and minimum:
this is why budget-constrained optimisation problems are
guaranteed to have solutions.
\end{rulebox}

\begin{rulebox}{Definition}
Let $f:\R^n\to\R$. The partial derivative of $f$ with respect to $x_i$ at
$\mathbf{x}=(x_1,\ldots,x_n)$ is
\[
  \pd{f}{x_i}(\mathbf{x}) \;=\; \lim_{h\to0}
    \frac{f(x_1,\ldots,x_i+h,\ldots,x_n)-f(x_1,\ldots,x_n)}{h} .
\]
Only the $i$-th slot changes; all other coordinates sit still. It is the
ordinary one-variable derivative of the function
$t\mapsto f(x_1,\ldots,t,\ldots,x_n)$, evaluated at $t=x_i$.
\end{rulebox}

\begin{rulebox}{What it measures: a slice of the surface}
Fix $y=y_{0}$. The vertical plane $y=y_{0}$ cuts the surface $z=f(x,y)$
along a curve, and along that curve only $x$ varies, so the curve is the
graph of the one-variable function $g(x)=f(x,y_{0})$. Then
\[
  \pd{f}{x}(x_{0},y_{0})=g'(x_{0}),
\]
the ordinary slope of the cut curve at $x_{0}$. Symmetrically,
$\partial f/\partial y$ is the slope of the cut made by the plane
$x=x_{0}$. A partial derivative is therefore never a slope \emph{of the
surface}; it is the slope of the surface \emph{along one chosen direction},
and different directions give different numbers. What the surface does in
all the remaining directions is the subject of the gradient below and of the
directional derivative.
\end{rulebox}

\begin{rulebox}{Practical computation}
To compute $\partial f/\partial x_i$, differentiate $f$ with respect to
$x_i$ and treat every other variable as a constant. All single-variable
rules (power, product, quotient, chain) apply unchanged. Worked example, for
$f(x,y)=x^{3}y^{2}+\mathrm{e}^{xy}-\ln y$:
\[
  \pd{f}{x}=3x^{2}y^{2}+y\,\mathrm{e}^{xy},
  \qquad
  \pd{f}{y}=2x^{3}y+x\,\mathrm{e}^{xy}-\frac{1}{y} .
\]
In the first computation $y$ is a number: $y^{2}$ is a constant factor, and
the chain rule on $\mathrm{e}^{xy}$ produces the inner derivative $y$. In
the second, $x$ is the constant and $\ln y$ finally contributes.
\end{rulebox}

\begin{rulebox}{Notation}
The following all refer to the same object:
\[
  \pd{f}{x} = f_x = \partial_x f = D_x f .
\]
\end{rulebox}

\begin{rulebox}{Gradient}
The vector of all first-order partial derivatives is
\[
  \grad f(\mathbf{x}) \;=\;
  \begin{pmatrix} f_{x_1}(\mathbf{x}) \\ f_{x_2}(\mathbf{x}) \\ \vdots \\
    f_{x_n}(\mathbf{x}) \end{pmatrix} .
\]
\end{rulebox}

\begin{rulebox}{Gradient and derivative: column and row}
The gradient is a \textbf{column} vector. The \textbf{derivative} of
$f\colon\R^{n}\to\R$ at $\mathbf{x}$, written $Df(\mathbf{x})$, is the same
list of partials arranged as a $1\times n$ \textbf{row}:
\[
  Df(\mathbf{x})=\bigl(f_{x_1}(\mathbf{x})\ \ f_{x_2}(\mathbf{x})\ \ \cdots\ \
  f_{x_n}(\mathbf{x})\bigr)=\grad f(\mathbf{x})^{T} .
\]
Same numbers, different shape. The row is the form that multiplies a column
of changes $\mathbf{h}$:
$Df(\mathbf{x})\,\mathbf{h}=\grad f(\mathbf{x})\cdot\mathbf{h}
=\sum_{i}f_{x_i}h_{i}$, which is the total differential of Chapter~12 and
the linear part of the tangent plane below. In Chapter~13 these rows are
stacked into the Jacobian matrix.
\end{rulebox}

\begin{rulebox}{Functions of $n$ variables written with sums}
Many economic functions are sums over goods, periods or inputs. To
differentiate with respect to one fixed variable $x_{k}$, ask which terms of
the sum contain $x_{k}$; every other term is a constant and drops out.
\begin{itemize}[leftmargin=1.6em,itemsep=0.3ex]
  \item \textbf{One function of each variable.} For
        $f=\sum_{i=1}^{n}g(x_{i})$ only the term $i=k$ survives:
        \[
          \pd{f}{x_{k}}=g'(x_{k}) .
        \]
        So $\sum_{i}x_{i}^{3}$ has $\partial f/\partial x_{k}=3x_{k}^{2}$.
  \item \textbf{The index inside the term.} $\sum_{i=1}^{n}x_{i}^{\,i}
        =x_{1}+x_{2}^{2}+\cdots+x_{n}^{n}$: the term containing $x_{k}$ is
        $x_{k}^{k}$, so $\partial f/\partial x_{k}=k\,x_{k}^{k-1}$.
        Likewise $\sum_{j=1}^{n}j/x_{j}^{4}=\sum_{j}j\,x_{j}^{-4}$ gives
        $\partial f/\partial x_{k}=-4k\,x_{k}^{-5}=-4k/x_{k}^{5}$.
  \item \textbf{The sum inside a function.} For $f=\bigl(\sum_{i}x_{i}\bigr)^{2}$
        use the chain rule; the inner derivative is
        $\partial\bigl(\sum_{i}x_{i}\bigr)/\partial x_{k}=1$, so
        $\partial f/\partial x_{k}=2\sum_{i=1}^{n}x_{i}$, the same for
        every $k$.
  \item \textbf{Products.} $f=\prod_{i}x_{i}$ gives
        $\partial f/\partial x_{k}=\prod_{i\neq k}x_{i}$: the product with
        the factor $x_{k}$ left out.
\end{itemize}
The summation index is a dummy and the differentiation variable is fixed:
call them by different letters ($i$ and $k$). Writing
$\partial/\partial x_{i}$ of a sum over $i$ is the standard source of
confusion.
\end{rulebox}

\begin{rulebox}{Exponentials to another base}
For $a>0$, $a^{x}=\mathrm{e}^{x\ln a}$, so
\[
  \frac{\mathrm{d}}{\mathrm{d}x}\,a^{x}=a^{x}\ln a,
  \qquad
  \pd{}{x}\,a^{g(x,y)}=a^{g(x,y)}\ln a\cdot g_{x}(x,y) .
\]
For example $\partial(2^{xy})/\partial x=y\,2^{xy}\ln2$. Keep the three
shapes apart: $x^{a}$ (variable base, power rule), $a^{x}$ (variable
exponent, factor $\ln a$), and $x^{y}$ (both, written as
$\mathrm{e}^{y\ln x}$).
\end{rulebox}

\begin{gallery}
\begin{tikzpicture}
  \begin{axis}[surfaxis, title={$f(K,L)=\sqrt{KL}$, sliced at $L=2$},
      xlabel={$K$}, ylabel={$L$}, zlabel={}, view={135}{26},
      zmin=0, zmax=3.2, z post scale=1.7]
    \addplot3[surf, domain=0.15:3, y domain=0.15:3, opacity=0.85]{sqrt(x*y)};
    \addplot3[samples=80, samples y=0, domain=0.15:3, line width=1.1pt,
      draw=black]({x},{2},{sqrt(2*x)});
  \end{axis}
\end{tikzpicture}\hfill
\begin{tikzpicture}
  \begin{axis}[contouraxis, xmin=-0.3, xmax=3.4, ymin=-0.3, ymax=3.4,
      xtick={1,2,3}, ytick={1,2,3},
      title={isoquants of $\sqrt{KL}$ with $\grad f$},
      xlabel={$K$}, ylabel={$L$}]
    \addplot[accent!40, line width=0.5pt, domain=0.22:3.3]{0.64/x};
    \addplot[accent!40, line width=0.5pt, domain=0.48:3.3]{1.44/x};
    \addplot[accent!40, line width=0.5pt, domain=0.85:3.3]{2.56/x};
    \addplot[accent!40, line width=0.5pt, domain=1.34:3.3]{4/x};
    \draw[-{Stealth[length=6pt,width=4pt]}, accent, line width=1pt]
      (axis cs:1,1) -- (axis cs:1.62,1.62);
    \draw[-{Stealth[length=6pt,width=4pt]}, accent, line width=1pt]
      (axis cs:2,0.5) -- (axis cs:2.16,1.14);
    \draw[-{Stealth[length=6pt,width=4pt]}, accent, line width=1pt]
      (axis cs:0.5,2) -- (axis cs:1.14,2.16);
    \addplot[only marks, mark=*, mark size=1pt, black, forget plot]
      coordinates {(1,1) (2,0.5) (0.5,2)};
  \end{axis}
\end{tikzpicture}
\gcap{Left: the plane $L=2$ cuts the surface; the slope of the cut curve is
$\partial f/\partial K$. Right: the gradient at three points, each one
perpendicular to the isoquant through it.}
\end{gallery}

\begin{rulebox}{Geometric meaning --- and why}
All three facts follow from one identity: the rate of change of $f$ in the
direction of a unit vector $\mathbf{u}$ is $\grad f\cdot\mathbf{u}$ (the
directional derivative below).
\begin{itemize}[leftmargin=1.6em,itemsep=0.4ex]
  \item \textbf{Steepest ascent.} Since
        $\grad f\cdot\mathbf{u}=\|\grad f\|\,\|\mathbf{u}\|\cos\theta
        =\|\grad f\|\cos\theta$, the rate is largest when $\cos\theta=1$,
        i.e. when $\mathbf{u}$ points \emph{along} $\grad f$. So
        $\grad f(\mathbf{x}_0)$ is the direction of steepest ascent, and
        $-\grad f$ of steepest descent.
  \item \textbf{Its length is that steepest rate.} Putting $\theta=0$ above
        leaves exactly $\|\grad f(\mathbf{x}_0)\|$.
  \item \textbf{Perpendicular to the level set.} Let $\mathbf{c}(t)$ be any
        curve running inside the level set through $\mathbf{x}_{0}$, so that
        $f(\mathbf{c}(t))$ is \emph{constant}. Differentiating with the
        chain rule gives
        \[
          0=\frac{\mathrm{d}}{\mathrm{d}t}f(\mathbf{c}(t))
           =\grad f(\mathbf{c}(t))\cdot\mathbf{c}'(t) .
        \]
        At $t_{0}$ this says $\grad f(\mathbf{x}_{0})$ is orthogonal to
        $\mathbf{c}'(t_{0})$, the tangent to the level set. Since the curve
        was arbitrary, the gradient is perpendicular to the whole level set.
\end{itemize}
The third fact is the one that matters most later: an optimum on a
constraint is a point where the level set of the objective and the
constraint set are \emph{tangent}, which is the same as their gradients
being parallel --- and that is the Lagrange condition.
\end{rulebox}

\begin{rulebox}{Economic interpretation}
\begin{itemize}[leftmargin=1.6em,itemsep=0.3ex]
  \item For a production function $Q=F(K,L)$: $F_K=\mathrm{MP}_K$
        (marginal product of capital), $F_L=\mathrm{MP}_L$ (marginal
        product of labour).
  \item For a utility function $U(x,y)$: $U_x, U_y$ are the marginal
        utilities; the marginal rate of substitution is
        $\mathrm{MRS}_{xy}=U_x/U_y$.
  \item For production: the marginal rate of technical substitution is
        $\mathrm{MRTS}_{LK}=F_L/F_K$.
\end{itemize}
\end{rulebox}

\begin{rulebox}{Directional derivative}
Let $\mathbf{u}\in\R^n$ be a \emph{unit} vector ($\|\mathbf{u}\|=1$). The
directional derivative of $f$ at $\mathbf{x}_0$ in the direction
$\mathbf{u}$ is
\[
  D_{\mathbf{u}}f(\mathbf{x}_0) \;=\; \grad f(\mathbf{x}_0)\cdot\mathbf{u} .
\]
\end{rulebox}

\begin{rulebox}{Caution}
If you are given a non-unit direction $\mathbf{v}$, you \emph{must}
normalise first: $\mathbf{u}=\mathbf{v}/\|\mathbf{v}\|$. Forgetting this is
the single most common error in directional-derivative problems.
\end{rulebox}

\begin{rulebox}{Differentiability in several variables}
Let $D\subseteq\R^{n}$ be open and $\mathbf{x}_{0}\in D$. The function $f$
is \textbf{differentiable} \de{(differenzierbar)} at $\mathbf{x}_{0}$ if
there is a vector $\mathbf{a}$ with
\[
  f(\mathbf{x}_{0}+\mathbf{h})=f(\mathbf{x}_{0})+\mathbf{a}\cdot\mathbf{h}
  +r(\mathbf{h}),
  \qquad
  \frac{r(\mathbf{h})}{\|\mathbf{h}\|}\to0\ \text{ as }\ \mathbf{h}\to\mathbf{0}.
\]
Near $\mathbf{x}_{0}$, $f$ is a linear function plus an error that vanishes
faster than the distance. For $n=2$ the graph of the linear part is a plane,
so differentiability says that the surface has a \textbf{tangent plane} at
the point: zoom in far enough and the surface looks flat.

If $f$ is differentiable at $\mathbf{x}_{0}$, then
\begin{itemize}[leftmargin=1.6em,itemsep=0.3ex]
  \item $f$ is continuous at $\mathbf{x}_{0}$;
  \item every partial derivative exists there, and
        $\mathbf{a}=\grad f(\mathbf{x}_{0})$ (take
        $\mathbf{h}=t\mathbf{e}_{k}$ and divide by $t$);
  \item every directional derivative exists and equals
        $\grad f(\mathbf{x}_{0})\cdot\mathbf{u}$.
\end{itemize}
\textbf{Sufficient criterion.} If all partials exist on an open
neighbourhood of $\mathbf{x}_{0}$ and are continuous there ($f$ is
$C^{1}$, \emph{continuously differentiable}), then $f$ is differentiable at
$\mathbf{x}_{0}$. Polynomials, exponentials, logarithms of positive
expressions, roots of positive expressions, quotients with non-zero
denominator, and all their sums, products and compositions are $C^{1}$ on
the open set where their formulas apply.

\textbf{Partials alone are not enough.} Take
$f(x,y)=\dfrac{xy}{x^{2}+y^{2}}$ with $f(0,0)=0$. On both axes $f=0$, so
$f_{x}(0,0)=f_{y}(0,0)=0$. But on the diagonal $y=x$,
$f=x^{2}/(2x^{2})=\tfrac12$ for every $x\neq0$. So $f$ is not even
continuous at the origin, and therefore not differentiable: the partials
see only two directions out of infinitely many.
\end{rulebox}

\begin{rulebox}{Where differentiability fails, and finding the set}
Compute the partials by the rules. The formulas are valid on the open set
where the rules applied; the points where a formula breaks down (a zero
denominator, a zero radicand, the kink of an absolute value or a minimum)
are the suspects, and each suspect is settled by the difference quotient.
\begin{itemize}[leftmargin=1.6em,itemsep=0.3ex]
  \item \textbf{Cone apex.} $f=\sqrt{(x-a)^{2}+(y-b)^{2}}$ has
        $f_{x}=(x-a)/f$, undefined at $(a,b)$. Directly,
        $f(a+h,b)-f(a,b)=\sqrt{h^{2}}=|h|$, and $|h|/h=\pm1$ has no limit:
        $f_{x}(a,b)$ does not exist, nor does $f_{y}(a,b)$. The same
        happens for the elliptic cone $\sqrt{4(x-a)^{2}+y^{2}}$ at $(a,0)$,
        with quotient $2|h|/h$. For the root of a sum of squares of
        \emph{linear} expressions, such as these cones, $f$ is
        differentiable exactly where the radicand is non-zero, and the
        graph has a sharp tip (or crease) where it vanishes. With higher
        powers under the root this can change: $\sqrt{x^{4}+y^{4}}$ is
        differentiable even at the origin, where its radicand is $0$.
  \item \textbf{$|x|+|y|$.} $f_{x}=1$ for $x>0$ and $-1$ for $x<0$; at
        $x=0$ the quotient is $|h|/h$ again. So $f_{x}$ fails on the whole
        $y$-axis and $f_{y}$ on the whole $x$-axis, and $f$ is
        differentiable exactly off the axes. The graph is an upside-down
        square pyramid, creased along both axes.
  \item \textbf{$\min\{ax,by\}$} with $a,b>0$. Where $ax<by$, $f=ax$;
        where $ax>by$, $f=by$. On the kink line $ax=by$ (in the quadrant
        $x,y>0$ of economics, a ray from the origin), raising $x$ leaves
        $f$ at $by$ (slope $0$) while lowering $x$ lowers $f$ at rate $a$:
        the one-sided slopes differ, so $f_{x}$ does not exist there, and
        nor does $f_{y}$. The graph is a roof whose ridge lies above that
        line.
\end{itemize}
A breaking formula is a warning, not a verdict: the partial may still
exist at that point, and only the difference quotient decides. And partials
that exist at a point still do not guarantee differentiability; continuity
of the partials near the point does.
\end{rulebox}

\begin{rulebox}{Proving differentiability from the definition}
When the $C^{1}$ criterion is not available, or a proof is asked for,
work with the remainder directly. At a point $\mathbf{x}^{*}$:
\begin{enumerate}[leftmargin=1.6em,itemsep=0.3ex]
  \item \textbf{Candidate.} Compute $\grad f(\mathbf{x}^{*})$; it is the
        only possible linear part. If a partial does not exist, stop: $f$
        is not differentiable.
  \item \textbf{Remainder.}
        $r(\mathbf{h})=f(\mathbf{x}^{*}+\mathbf{h})-f(\mathbf{x}^{*})
        -\grad f(\mathbf{x}^{*})\cdot\mathbf{h}$, expanded and simplified.
        For a polynomial, every term left over contains at least two
        factors $h_{i}$.
  \item \textbf{Estimate.} Show $|r(\mathbf{h})|/\|\mathbf{h}\|\to0$ with
        the standard inequalities
        \[
          |h_{i}|\le\|\mathbf{h}\|,\qquad
          \frac{h_{i}^{2}}{\|\mathbf{h}\|}\le\|\mathbf{h}\|,\qquad
          \frac{|h_{1}h_{2}|}{\|\mathbf{h}\|}\le\tfrac12\|\mathbf{h}\| ,
        \]
        and, for a product of three factors $h_{i}$, a bound
        $\|\mathbf{h}\|^{2}$.
  \item \textbf{Transcendental factors.} Replace them by their linear part
        plus a remainder: $\mathrm{e}^{t}=1+t+R(t)$ with $|R(t)|\le2t^{2}$
        for $|t|\le1$, and $\ln(1+t)=t+R(t)$ with $|R(t)|\le2t^{2}$ for
        $|t|\le\tfrac12$ (both from the Lagrange remainder of Chapter~4).
        The quadratic bound on $R$ is what makes $R/\|\mathbf{h}\|\to0$.
\end{enumerate}
\textbf{Roots at the origin.} For $f=(x^{2}+y^{2})^{p}$ with $p>0$,
$f(\mathbf{h})=\|\mathbf{h}\|^{2p}$. The difference quotient on an axis is
$|h|^{2p}/h$: it tends to $0$ if $2p>1$, and has no limit if $2p\le1$. So
$f$ is differentiable at the origin exactly when $p>\tfrac12$; the cone
$p=\tfrac12$ is the borderline case, and $p=\tfrac13$ is worse than a cone.
\end{rulebox}

\begin{rulebox}{Tangent plane and linearisation}
If $f$ is differentiable at $P=(x_{0},y_{0})$, the \textbf{tangent plane}
\de{(Tangentialebene)} to its graph at $(x_{0},y_{0},f(P))$ is
\[
  z=f(P)+f_{x}(P)\,(x-x_{0})+f_{y}(P)\,(y-y_{0}) .
\]
In $n$ variables, with $P=\mathbf{p}=(p_{1},\ldots,p_{n})$, the
\textbf{tangent hyperplane} is
\[
  z=f(\mathbf{p})+\grad f(\mathbf{p})\cdot(\mathbf{x}-\mathbf{p})
   =f(\mathbf{p})+\sum_{i=1}^{n}f_{x_i}(\mathbf{p})\,(x_{i}-p_{i}) .
\]
\begin{itemize}[leftmargin=1.6em,itemsep=0.3ex]
  \item The right-hand side is the \textbf{linearisation} $L(\mathbf{x})$
        of $f$ at $\mathbf{p}$, the multivariable version of the tangent
        line $f(a)+f'(a)(x-a)$ of Chapter~6. With
        $\mathrm{d}x_{i}=x_{i}-p_{i}$ it reads $L=f(\mathbf{p})+\mathrm{d}f$:
        the old value plus the total differential.
  \item The partials are evaluated \emph{at the point} $\mathbf{p}$; they
        are numbers, which is what makes the result a plane.
  \item For a point $Q$ near $P$, $f(Q)\approx L(Q)$. The difference
        $f(Q)-L(Q)$ is the remainder $r(\mathbf{h})$ of the definition, so
        comparing the two numbers measures how good the approximation is.
        For $f$ with continuous second partials it is at most of the order
        $\|Q-P\|^{2}$: halve the step and the error falls to about a
        quarter, or even less.
  \item Equivalently, $f_{x}(P)(x-x_{0})+f_{y}(P)(y-y_{0})-(z-f(P))=0$: the
        vector $(f_{x}(P),\,f_{y}(P),\,-1)$ is normal to the plane.
\end{itemize}
Where $f$ is not differentiable there is no tangent plane: a cone has none
at its apex.
\end{rulebox}

\section{Worked Examples}

\begin{worked}{For $f(x,y)=x^{3}y^{2}+\mathrm{e}^{2x}\ln y$, compute both partials and the gradient at $(1,1)$; then read the marginal products of $Q=10K^{0.3}L^{0.7}$}
\wstep{Step 1. Differentiating in $x$ means freezing $y$ --- literally.}
Treat $y$ as if it were the number $7$. Then $x^{3}y^{2}$ is
``$y^{2}$ times $x^{3}$'', a constant times a power, and $\ln y$ is a
constant multiplying $\mathrm{e}^{2x}$:
\[
  f_{x}=3x^{2}y^{2}+2\mathrm{e}^{2x}\ln y .
\]
\wtrap{Freezing means the frozen variable never produces a derivative
factor. $\frac{\partial}{\partial x}\bigl[\ln y\bigr]=0$, not $1/y$. The
symbol $\partial$ instead of $\mathrm{d}$ is the reminder that something is
being held fixed.}

\wstep{Step 2. Now freeze $x$ instead.}
$x^{3}$ is a constant multiplying $y^{2}$, and $\mathrm{e}^{2x}$ is a
constant multiplying $\ln y$:
\[
  f_{y}=2x^{3}y+\frac{\mathrm{e}^{2x}}{y} .
\]

\wstep{Step 3. Evaluate, and assemble the gradient.}
At $(1,1)$, with $\ln1=0$ and $\mathrm{e}^{2}\approx7.389$:
\[
  f_{x}(1,1)=3\cdot1\cdot1+2\mathrm{e}^{2}\cdot0=3,
  \qquad
  f_{y}(1,1)=2\cdot1\cdot1+\frac{\mathrm{e}^{2}}{1}=2+\mathrm{e}^{2}\approx9.389 .
\]
\[
  \grad f(1,1)=\begin{pmatrix}3\\2+\mathrm{e}^{2}\end{pmatrix}\approx\begin{pmatrix}3\\9.389\end{pmatrix}
\]
\wtrap{The gradient is a \textbf{vector}, and it is a vector in the
\emph{input} space --- here a direction in the $xy$-plane, not on the
surface. Reporting a single number, or a three-component object including
the height, are the two standard confusions.}

\wstep{Step 4. What each partial measures.}
$f_{x}$ is the slope of the curve obtained by slicing the surface with the
plane $y=1$ --- an ordinary one-variable derivative of the function
$x\mapsto f(x,1)$. Indeed $f(x,1)=x^{3}$, whose derivative at $x=1$ is $3$.
\checkmark\ Two partials, two slices, two ordinary derivatives.

\wstep{Step 5. What the gradient does that the partials alone do not.}
\begin{itemize}[leftmargin=1.6em,itemsep=0.2ex]
  \item It points in the direction of \textbf{steepest ascent}. Here that
        is roughly $(0.30,\,0.95)$ after normalising: mostly ``north''.
  \item Its length $\|\grad f\|=\sqrt{9+(2+\mathrm{e}^{2})^{2}}\approx9.86$
        is the steepest rate of climb available at that point.
  \item It is \textbf{perpendicular} to the level curve through the point.
\end{itemize}
The last is why the gradient governs constrained optimisation: at an
optimum the objective's gradient must be parallel to the constraint's, which
is the Lagrange condition.

\wstep{Step 6. Directional derivative: normalise, then take an inner product.}
In the direction of $\mathbf v=(3,4)$:
\[
  \mathbf u=\frac{\mathbf v}{\|\mathbf v\|}=\Bigl(\tfrac35,\tfrac45\Bigr),
  \qquad
  D_{\mathbf u}f=\grad f\cdot\mathbf u
  =\tfrac35\cdot3+\tfrac45(2+\mathrm{e}^{2})\approx9.31 .
\]
\wtrap{Normalise first. Using $(3,4)$ unnormalised gives $46.6$, which is
five times too large and has no meaning --- ``rate of change per unit
distance'' requires a unit vector. And $9.31<9.86$, as it must be: no
direction beats the gradient.}

\wstep{Step 7. The economic reading, on a production function.}
For $Q=10K^{0.3}L^{0.7}$:
\[
  MP_{K}=\pd{Q}{K}=3K^{-0.7}L^{0.7}=3\Bigl(\frac{L}{K}\Bigr)^{0.7},
  \qquad
  MP_{L}=\pd{Q}{L}=7\Bigl(\frac{K}{L}\Bigr)^{0.3} .
\]
At $K=100$, $L=50$: $MP_{K}=3(0.5)^{0.7}\approx1.85$ and
$MP_{L}=7(2)^{0.3}\approx8.62$. So one more unit of labour adds about $8.6$
units of output, one more of capital about $1.85$.

Two structural facts fall out. The \textbf{marginal rate of technical
substitution} is
\[
  \mathrm{MRTS}=\frac{MP_{L}}{MP_{K}}=\frac{0.7}{0.3}\cdot\frac{K}{L}\approx4.67 ,
\]
and each marginal product depends only on the \emph{ratio} $K/L$ --- the
signature of constant returns to scale, which is why Cobb--Douglas with
exponents summing to $1$ behaves the way it does.
\wtrap{$MP_{K}$ falls as $K$ rises (the exponent $-0.7$ is negative):
diminishing marginal product. But that is diminishing \emph{marginal},
not diminishing total --- output still increases. Confusing the two is the
most common misreading of a partial derivative in economics.}
\end{worked}

\begin{worked}{Find the tangent plane of $f(x,y)=x\ln(x+2y)$ at $P=(1,0)$, and test it at the nearby point $Q=(1.1,\,0.05)$}
\wstep{Step 1. Make sure there is a plane to find.}
The domain is the open half-plane $x+2y>0$, and $P$ lies in it since
$1+0>0$. On that half-plane $f$ is a polynomial times the logarithm of a
positive linear expression, so its partials exist and are continuous: $f$
is $C^{1}$, hence differentiable at $P$, and the tangent plane exists.

\wstep{Step 2. The height at $P$.}
$f(1,0)=1\cdot\ln1=0$. The plane touches the surface at $(1,0,0)$.

\wstep{Step 3. The two partials, then their values at $P$.}
In $x$ the function is a product of $x$ and $\ln(x+2y)$:
\[
  f_{x}=\ln(x+2y)+x\cdot\frac{1}{x+2y} .
\]
In $y$ the factor $x$ is a constant, and the chain rule supplies the inner
derivative $2$:
\[
  f_{y}=\frac{2x}{x+2y} .
\]
At $P$: $f_{x}(1,0)=\ln1+\tfrac11=1$ and $f_{y}(1,0)=\tfrac21=2$.
\wtrap{$x$ appears twice in $f$, once as a factor and once inside the
logarithm, so $f_{x}$ needs the product rule. Writing $f_{x}=x/(x+2y)$
alone loses the term $\ln(x+2y)$; here it happens to vanish at $P$, which
hides the error until the next point.}

\wstep{Step 4. Assemble the plane.}
\[
  z=f(P)+f_{x}(P)(x-1)+f_{y}(P)(y-0)=0+1\cdot(x-1)+2y=x+2y-1 .
\]
\wtrap{The partials go in as \emph{numbers} evaluated at $P$. Leaving them
as functions of $x$ and $y$ gives a curved surface, not a plane.}

\wstep{Step 5. Compare at $Q$.}
\[
  f(Q)=1.1\cdot\ln(1.1+0.1)=1.1\ln1.2\approx0.200554,
  \qquad
  L(Q)=1.1+0.1-1=0.200000 .
\]
The plane misses by about $0.000554$, while $f$ itself changed by about
$0.2$: the error is under $0.3\%$ of the change. The step has length
$\|Q-P\|=\sqrt{0.1^{2}+0.05^{2}}\approx0.112$. The error is far below that
length, and even below its square $0.0125$, as the rule box predicts: the
remainder vanishes faster than the distance.

\wstep{Step 6. The same number as a total differential.}
With $\mathrm{d}x=0.1$ and $\mathrm{d}y=0.05$,
\[
  \mathrm{d}f=f_{x}(P)\,\mathrm{d}x+f_{y}(P)\,\mathrm{d}y=0.1+0.1=0.2,
\]
and $L(Q)=f(P)+\mathrm{d}f=0.2$. \checkmark\ The tangent plane and the
total differential are one idea written two ways.
\end{worked}

\begin{worked}{Prove from the definition that $f(x,y)=x^{2}y$ is differentiable at every point; then decide at the origin for $(x^{2}+y^{2})^{1/3}$ and $(x^{2}+y^{2})^{2/3}$}
\wstep{Step 1. The candidate for the linear part.}
Fix a point $(a,b)$. If $f$ is differentiable there, the linear part must
be $\grad f(a,b)\cdot\mathbf{h}$, so compute the partials:
$f_{x}=2xy$, $f_{y}=x^{2}$, hence $\grad f(a,b)=(2ab,\,a^{2})$.

\wstep{Step 2. Form the remainder and let everything cancel that can.}
With $\mathbf{h}=(h_{1},h_{2})$,
\[
  f(a+h_{1},b+h_{2})=(a^{2}+2ah_{1}+h_{1}^{2})(b+h_{2})
  =a^{2}b+a^{2}h_{2}+2abh_{1}+2ah_{1}h_{2}+bh_{1}^{2}+h_{1}^{2}h_{2} .
\]
Subtract $f(a,b)=a^{2}b$ and the linear part $2abh_{1}+a^{2}h_{2}$:
\[
  r(\mathbf{h})=2a\,h_{1}h_{2}+b\,h_{1}^{2}+h_{1}^{2}h_{2} .
\]
Every surviving term contains at least \emph{two} factors $h_{i}$. That is
the sign that the candidate was right.

\wstep{Step 3. Bound the quotient by something that tends to $0$.}
Use $|h_{i}|\le\|\mathbf{h}\|$ and $|h_{1}h_{2}|\le\tfrac12\|\mathbf{h}\|^{2}$
(from $(|h_{1}|-|h_{2}|)^{2}\ge0$):
\[
  \frac{|r(\mathbf{h})|}{\|\mathbf{h}\|}
  \le\frac{2|a|\cdot\tfrac12\|\mathbf{h}\|^{2}+|b|\,\|\mathbf{h}\|^{2}
          +\|\mathbf{h}\|^{3}}{\|\mathbf{h}\|}
  =(|a|+|b|)\,\|\mathbf{h}\|+\|\mathbf{h}\|^{2}\;\longrightarrow\;0 .
\]
So $f$ is differentiable at $(a,b)$, for every $(a,b)$, with derivative
$(2ab,\,a^{2})$.
\wtrap{Checking $r/\|\mathbf{h}\|\to0$ along the axes, or along every
line, is not a proof: the bound must hold for \emph{all} small $\mathbf{h}$
at once. The bound may depend on the fixed point $(a,b)$, but only through
constants.}

\wstep{Step 4. A root at the origin: $g=(x^{2}+y^{2})^{1/3}$.}
First the partials, from the difference quotient:
\[
  \frac{g(h,0)-g(0,0)}{h}=\frac{(h^{2})^{1/3}}{h}
  =\frac{|h|^{2/3}}{h}=\pm|h|^{-1/3},
\]
which is unbounded as $h\to0$. So $g_{x}(0,0)$ does not exist (nor, by
symmetry, $g_{y}(0,0)$), and $g$ is \textbf{not differentiable} at the
origin, although it is continuous there. Its graph has a tip even sharper
than a cone's.

\wstep{Step 5. A higher power: $k=(x^{2}+y^{2})^{2/3}$.}
Now $\dfrac{k(h,0)-k(0,0)}{h}=\dfrac{|h|^{4/3}}{h}=\pm|h|^{1/3}\to0$, so both
partials are $0$ and the candidate linear part is $0$. The remainder is
$r(\mathbf{h})=k(\mathbf{h})=\|\mathbf{h}\|^{4/3}$, and
\[
  \frac{|r(\mathbf{h})|}{\|\mathbf{h}\|}=\|\mathbf{h}\|^{1/3}\to0 .
\]
So $k$ \textbf{is} differentiable at the origin, with the tangent plane
$z=0$. In general $(x^{2}+y^{2})^{p}=\|\mathbf{h}\|^{2p}$ at the origin, and
the verdict depends only on whether $2p>1$.
\end{worked}

\section{Practice Problems}

\begin{problems}
  \item Compute $f_x$ and $f_y$ for $f(x,y)=5x^4-2x^2y^3+y^5$.
  \item For $F(K,L)=12\,K^{1/4}L^{3/4}$, find $F_K$ and $F_L$ at
        $(K,L)=(16,81)$.
  \item Compute $f_x$ and $f_y$ for $f(x,y)=x^2e^{3y}$.
  \item Compute $\nabla f$ for $f(x,y)=4x^3y-xy^4+7$.
  \item Compute the directional derivative of $f(x,y)=x^2-4xy+y^3$ at
        $(2,1)$ in the direction $\mathbf{v}=(5,-12)$.
  \item For $F(K,L)=6K^{2/3}L^{1/3}$, compute $\mathrm{MP}_K$ and
        $\mathrm{MP}_L$ at $(K,L)=(27,8)$.
  \item Compute $f_x$ and $f_y$ for $f(x,y)=(3x-2y)^4$.
  \item Compute $f_x$ and $f_y$ for $f(x,y)=e^{xy^2}$.
  \item For $F(K,L)=25K^{0.6}L^{0.5}$, find $F_K$ and $F_L$.
  \item Compute $\nabla f$ for $f(x,y,z)=x^2+3yz-z^2e^x$.
  \item Compute the directional derivative of $f(x,y)=e^x\ln(1+y)$ at
        $(0,0)$ in the direction $\mathbf{v}=(1,2)$.
  \item For $U(x,y)=2\ln(x)+5\ln(y)$, derive $\mathrm{MRS}_{xy}=U_x/U_y$ as
        a function of $(x,y)$ and evaluate at $(3,5)$.
  \item Simplify and compute $f_x, f_y$ for
        $f(x,y)=\dfrac{x^2-y^2}{x+y}$.
  \item Compute $f_x$ and $f_y$ for $f(x,y)=\ln(x^2y+1)$.
  \item For $U(x_1,x_2,x_3)=x_1^{1/2}x_2^{1/3}x_3^{1/6}$, compute the
        three partials.
  \item Compute the directional derivative of $f(x,y,z)=xy-2yz+z^2$ at
        $(1,2,3)$ in the direction $\mathbf{v}=(1,2,2)$.
  \item For $F(K,L)=K^{0.25}L^{0.75}$, derive $\mathrm{MRTS}_{LK}=F_L/F_K$
        and evaluate at $(K,L)=(16,4)$.
  \item Expand and then compute $f_x, f_y$ for
        $f(x,y)=(x^2+y)(x-3y^2)$.
  \item Compute $f_x$ and $f_y$ for $f(x,y)=xe^{-y/x}$ with $x\neq0$.
  \item For $F(K,L)=(3K^{1/2}+2L^{1/2})^2$, find $F_K$ and $F_L$.
  \item Compute $\nabla f$ for $f(x,y)=\arctan(y/x)$ with $x>0$.
  \item For $f(x,y)=x^{1/2}y^{1/2}$ at $(K,L)=(4,9)$, find the unit
        direction of steepest ascent and the maximum rate of increase.
  \item For the CES function $F(K,L)=(K^{1/2}+L^{1/2})^2$, verify Euler's
        theorem $KF_K+LF_L=rF$ for the appropriate degree of homogeneity
        $r$.
  \item Compute $f_x$ and $f_y$ for $f(x,y)=\dfrac{2x+3y}{x^2+y^2}$.
  \item Compute $f_x$ and $f_y$ for $f(x,y)=\dfrac{e^x}{1+e^y}$.
  \item For $F(K,L)=A(\alpha K^\rho+(1-\alpha)L^\rho)^{\nu/\rho}$ (CES),
        find $F_L$.
  \item Compute the directional derivative of $f(x,y)=\ln(x^2+y^2)$ at
        $(3,4)$ in the direction $\mathbf{v}=(-1,1)$.
  \item Compute $f_x, f_y, f_z$ for $f(x,y,z)=x^2yz^3-4xy^2z+5z$.
  \item Compute $f_x$ and $f_y$ for $f(x,y)=\sqrt{4x^2+9y^2}$.
  \item For $F(K,L)=(K^{1/2}+L^{1/2})^4$, compute $F_K$ at
        $(K,L)=(4,9)$.
  \item Compute $\nabla f$ for $f(x,y,z)=\ln(x^2+y^2+z^2)$.
  \item Compute $f_x$ and $f_y$ for
        $f(x,y)=\ln\left(\sqrt{x^2+y^2}\right)$.
  \item Cost minimisation. A firm uses inputs $(x,y)$ with cost
        $C(x,y)=4x+9y$ and output $Q(x,y)=x^{1/2}y^{1/2}$. Compute
        $\nabla C$ and $\nabla Q$, and verify that at the cost-minimising
        input mix $\nabla C=\lambda\nabla Q$ for some $\lambda$.
  \item Compute $f_x$ and $f_y$ for $f(x,y)=(1+xy)^5$.
  \item Compute the directional derivative of $f(x,y)=x^3-2xy+y^2$ at
        $(1,2)$ in the direction $\mathbf{v}=(3,4)$.
  \item For $F(K,L)=10K^{0.3}L^{0.7}$ (Cobb--Douglas), compute $F_K$ and
        $F_L$ at $(K,L)=(100,100)$ and verify $KF_K+LF_L=F$.
  \item Compute $f_x$ and $f_y$ for $f(x,y)=\sin(xy)+\cos(x-y)$.
  \item Compute $\nabla f$ for $f(x,y)=x^2y-xy^2+3x-2y$ at the point
        $(1,-1)$.
  \item For $U(x,y)=x^\alpha y^{1-\alpha}$ (Cobb--Douglas,
        $0<\alpha<1$), derive $\mathrm{MRS}_{xy}=U_x/U_y$ and show it
        equals $\dfrac{\alpha}{1-\alpha}\cdot\dfrac{y}{x}$.
  \item Compute the directional derivative of $f(x,y)=e^{xy}$ at $(1,0)$
        in the direction $\mathbf{v}=(-3,4)$.
  \item Compute $f_x$ and $f_y$ for $f(x,y)=x^y$ (with $x>0$).
        \emph{Hint:} write $x^y=e^{y\ln x}$.
  \item For $F(K,L)=\min\{2K,3L\}$ (Leontief, fixed-proportions, on the
        smooth region away from the kink), compute $F_K$ and $F_L$ in
        each of the two regions $\{2K<3L\}$ and $\{2K>3L\}$.
  \item Compute $\nabla f$ for $f(x,y,z)=xyz+x^2-z^2$ at the point
        $(2,-1,3)$.
  \item Compute the unit vector of steepest \emph{descent} for
        $f(x,y)=x^2+4y^2$ at $(3,1)$, and the corresponding rate of
        decrease.
  \item For $f(x,y)=xe^y-ye^x$, compute $f_x, f_y$.
  \item Compute $f_x$ and $f_y$ for $f(x,y)=\dfrac{x}{x^2+y^2}$.
  \item For $U(x,y)=x+2\sqrt{y}$ (quasi-linear), derive
        $\mathrm{MRS}_{xy}$.
  \item Compute the directional derivative of $f(x,y,z)=x^2+y^2-z^2$ at
        $(2,-1,2)$ in the direction $\mathbf{v}=(2,2,1)$.
  \item Compute $f_x$ and $f_y$ for $f(x,y)=\tan(x+y)+\tan(x-y)$.
  \item For $F(K,L)=(\alpha K+\beta L)^r$ (linear-aggregate,
        $\alpha,\beta,r>0$), compute $F_K, F_L$, and verify
        $KF_K+LF_L=rF$.
  \item Compute $\nabla f$ for $f(x,y)=e^{-x}\sin(y)+e^x\cos(y)$.
  \item Compute the directional derivative of $f(x,y)=\sqrt{x^2+y^2}$ at
        $(3,4)$ in the direction $\mathbf{v}=(1,0)$.
  \item For $U(x,y)=\min\{2x,3y\}$, compute the marginal utilities on
        each smooth region $\{2x<3y\}$ and $\{2x>3y\}$.
  \item Compute $f_x, f_y$ for $f(x,y)=\dfrac{\ln(x+y)}{x-y}$ (assume
        $x>y>0$).
  \item For $F(K,L)=AK^\alpha L^\beta$ general Cobb--Douglas, derive
        $\mathrm{MRTS}_{LK}$ and show it depends only on the ratio $K/L$
        (not on scale).
  \item Compute the unit direction of steepest ascent for
        $f(x,y)=x^2-y^2$ at $(2,-1)$, and the maximum rate.
  \item Compute $f_x$ and $f_y$ for $f(x,y)=\cosh(x+y)+\sinh(x-y)$.
  \item For $f(x,y,z)=e^{x+y+z}$, compute $\nabla f$ at $(0,0,0)$ and the
        directional derivative in $\mathbf{v}=(1,1,1)$.
  \item Compute $f_x$ and $f_y$ for $f(x,y)=(\ln x)(\ln y)$ with
        $x,y>1$.
  \item For $F(K,L)=K+2K^{1/2}L^{1/2}+L$ (which equals
        $(K^{1/2}+L^{1/2})^2$), compute $F_K, F_L$ directly and confirm
        agreement with the CES form.
  \item Compute the directional derivative of $f(x,y)=x^2y+y^3$ at
        $(-1,2)$ in the direction $\mathbf{v}=(4,-3)$.
  \item For $U(x,y)=\alpha\ln x+\beta\ln y$ (log-Cobb--Douglas), derive
        $\mathrm{MRS}_{xy}$ and evaluate at $(x,y)=(\alpha,\beta)$.
  \item Compute $\nabla f$ for $f(x,y)=(x^2+y^2)^{-1/2}$ (Newtonian
        potential in 2D).
  \item Compute $f_x$ and $f_y$ for $f(x,y)=\arctan(xy)$.
  \item For $F(K,L)=(K^\rho+L^\rho)^{1/\rho}$ with $\rho<1$, $\rho\neq0$
        (CES with constant returns), compute $F_K$ and show
        $F_K=(K/F)^{\rho-1}$.
  \item Compute the directional derivative of $f(x,y,z)=xyz$ at
        $(2,3,-1)$ in the direction $\mathbf{v}=(2,-1,2)$.
\end{problems}

\section{Solutions}

\begin{solutions}
  \item $f_x=20x^3-4xy^3$, $f_y=-6x^2y^2+5y^4$.

  \item $F_K=12\cdot\frac14K^{-3/4}L^{3/4}=3(L/K)^{3/4}$;
        $F_L=12\cdot\frac34K^{1/4}L^{-1/4}=9(K/L)^{1/4}$. At $(16,81)$:
        $(L/K)^{3/4}=(81/16)^{3/4}=27/8$ so $F_K=81/8$;
        $(K/L)^{1/4}=(16/81)^{1/4}=2/3$ so $F_L=6$.

  \item $f_x=2xe^{3y}$, $f_y=3x^2e^{3y}$.

  \item $\nabla f=(12x^2y-y^4,\ 4x^3-4xy^3)^T$.

  \item $f_x=2x-4y$, $f_y=-4x+3y^2$. At $(2,1)$: $\nabla f=(0,-5)^T$.
        Normalise: $\|\mathbf{v}\|=\sqrt{25+144}=13$,
        $\mathbf{u}=(5/13,-12/13)^T$. Then
        \[
          D_{\mathbf{u}}f=0\cdot\frac{5}{13}+(-5)\cdot
            \left(-\frac{12}{13}\right)=\frac{60}{13} .
        \]

  \item $F_K=6\cdot\frac23K^{-1/3}L^{1/3}=4(L/K)^{1/3}$,
        $F_L=2(K/L)^{2/3}$. At $(27,8)$: $(L/K)^{1/3}=2/3$ so
        $\mathrm{MP}_K=8/3$; $(K/L)^{2/3}=(27/8)^{2/3}=9/4$ so
        $\mathrm{MP}_L=9/2$.

  \item Chain rule with $u=3x-2y$: $f_x=4(3x-2y)^3\cdot3=12(3x-2y)^3$,
        $f_y=-8(3x-2y)^3$.

  \item $f_x=y^2e^{xy^2}$, $f_y=2xye^{xy^2}$.

  \item $F_K=15K^{-0.4}L^{0.5}$, $F_L=12.5K^{0.6}L^{-0.5}$.

  \item $f_x=2x-z^2e^x$, $f_y=3z$, $f_z=3y-2ze^x$. So
        $\nabla f=(2x-z^2e^x,\ 3z,\ 3y-2ze^x)^T$.

  \item $f_x=e^x\ln(1+y)$, $f_y=e^x/(1+y)$. At $(0,0)$:
        $\nabla f=(0,1)^T$. Normalise: $\|\mathbf{v}\|=\sqrt5$,
        $\mathbf{u}=(1/\sqrt5,2/\sqrt5)^T$. So
        \[
          D_{\mathbf{u}}f=0\cdot\frac{1}{\sqrt5}+1\cdot\frac{2}{\sqrt5}
            =\frac{2\sqrt5}{5} .
        \]

  \item $U_x=2/x$, $U_y=5/y$. $\mathrm{MRS}_{xy}=U_x/U_y
        =(2/x)/(5/y)=2y/(5x)$. At $(3,5)$: $\mathrm{MRS}=10/15=2/3$.

  \item For $x+y\neq0$: $f(x,y)=x-y$, so $f_x=1$, $f_y=-1$.

  \item With $u=x^2y+1$: $f_x=2xy/(x^2y+1)$, $f_y=x^2/(x^2y+1)$.

  \item $U_{x_1}=\frac12x_1^{-1/2}x_2^{1/3}x_3^{1/6}$,
        $U_{x_2}=\frac13x_1^{1/2}x_2^{-2/3}x_3^{1/6}$,
        $U_{x_3}=\frac16x_1^{1/2}x_2^{1/3}x_3^{-5/6}$. Equivalently
        $U_{x_i}=(\alpha_i/x_i)U$ with $\alpha=(1/2,1/3,1/6)$.

  \item $f_x=y$, $f_y=x-2z$, $f_z=-2y+2z$. At $(1,2,3)$:
        $\nabla f=(2,-5,2)^T$. Normalise: $\|\mathbf{v}\|=\sqrt{1+4+4}=3$,
        $\mathbf{u}=(1/3,2/3,2/3)^T$. Then
        $D_{\mathbf{u}}f=(2-10+4)/3=-4/3$.

  \item $F_K=0.25K^{-0.75}L^{0.75}$, $F_L=0.75K^{0.25}L^{-0.25}$. So
        $\mathrm{MRTS}_{LK}=F_L/F_K=3K/L$. At $(16,4)$:
        $\mathrm{MRTS}=48/4=12$.

  \item Expand: $f=x^3-3x^2y^2+xy-3y^3$. Hence $f_x=3x^2-6xy^2+y$,
        $f_y=-6x^2y+x-9y^2$.

  \item With $u=-y/x$: $\partial u/\partial x=y/x^2$,
        $\partial u/\partial y=-1/x$. Product rule:
        \[
          f_x=e^{-y/x}+x\cdot e^{-y/x}\cdot\frac{y}{x^2}
            =e^{-y/x}\left(1+\frac{y}{x}\right), \qquad
          f_y=x\cdot e^{-y/x}\cdot\left(-\frac1x\right)=-e^{-y/x} .
        \]

  \item Let $g=3K^{1/2}+2L^{1/2}$, so $F=g^2$. Then
        \[
          F_K=2g\cdot\frac32K^{-1/2}=\frac{3g}{K^{1/2}}
            =\frac{3(3K^{1/2}+2L^{1/2})}{K^{1/2}}, \qquad
          F_L=2g\cdot L^{-1/2}=\frac{2(3K^{1/2}+2L^{1/2})}{L^{1/2}} .
        \]

  \item With $u=y/x$, $f=\arctan u$, $f'=1/(1+u^2)$:
        \[
          f_x=\frac{1}{1+y^2/x^2}\cdot\left(-\frac{y}{x^2}\right)
            =\frac{-y}{x^2+y^2}, \qquad
          f_y=\frac{x}{x^2+y^2} .
        \]
        So $\nabla f=(-y/(x^2+y^2),\ x/(x^2+y^2))^T$.

  \item $f_K=\frac12K^{-1/2}L^{1/2}$, $f_L=\frac12K^{1/2}L^{-1/2}$. At
        $(4,9)$: $\nabla f=(\frac12\cdot\frac32,\ \frac12\cdot\frac23)^T
        =(3/4,1/3)^T$. Then $\|\nabla f\|^2=9/16+1/9=97/144$, so the
        maximum rate is $\sqrt{97}/12\approx0.821$ in direction
        $\mathbf{u}^*=(9,4)^T/\sqrt{97}$.

  \item $F=(K^{1/2}+L^{1/2})^2$. Replacing $(K,L)\to(\lambda K,\lambda L)$
        gives $\lambda F$, so $r=1$. With $g=K^{1/2}+L^{1/2}$:
        $F_K=2g\cdot\frac12K^{-1/2}=g/K^{1/2}$, $F_L=g/L^{1/2}$. Hence
        \[
          KF_K+LF_L=K^{1/2}g+L^{1/2}g=g(K^{1/2}+L^{1/2})=g^2=F. \checkmark
        \]

  \item Quotient rule, $D=x^2+y^2$:
        \[
          f_x=\frac{2D-(2x+3y)(2x)}{D^2}=\frac{-2x^2-6xy+2y^2}{(x^2+y^2)^2},
          \qquad f_y=\frac{3x^2-4xy-3y^2}{(x^2+y^2)^2} .
        \]

  \item Treat denominator as constant for $f_x$: $f_x=e^x/(1+e^y)$. For
        $f_y$:
        \[
          f_y=e^x\cdot\frac{-e^y}{(1+e^y)^2}=-\frac{e^{x+y}}{(1+e^y)^2} .
        \]

  \item Let $S=\alpha K^\rho+(1-\alpha)L^\rho$, $F=AS^{\nu/\rho}$.
        Treating $K$ constant, $\partial S/\partial L=\rho(1-\alpha)L^{\rho-1}$:
        \[
          F_L=A\cdot\frac{\nu}{\rho}S^{\nu/\rho-1}\cdot\rho(1-\alpha)L^{\rho-1}
            =\nu(1-\alpha)AL^{\rho-1}S^{\nu/\rho-1} .
        \]

  \item $f=\ln(x^2+y^2)$, $f_x=2x/(x^2+y^2)$, $f_y=2y/(x^2+y^2)$. At
        $(3,4)$: $x^2+y^2=25$, $\nabla f=(6/25,8/25)^T$. Normalise:
        $\|\mathbf{v}\|=\sqrt2$, $\mathbf{u}=(-1/\sqrt2,1/\sqrt2)^T$. Then
        \[
          D_{\mathbf{u}}f=\frac{6}{25}\cdot\left(-\frac{1}{\sqrt2}\right)
            +\frac{8}{25}\cdot\frac{1}{\sqrt2}=\frac{2}{25\sqrt2}
            =\frac{\sqrt2}{25} .
        \]

  \item $f_x=2xyz^3-4y^2z$, $f_y=x^2z^3-8xyz$,
        $f_z=3x^2yz^2-4xy^2+5$.

  \item Chain rule, $u=4x^2+9y^2$:
        \[
          f_x=\frac{8x}{2\sqrt u}=\frac{4x}{\sqrt{4x^2+9y^2}}, \qquad
          f_y=\frac{9y}{\sqrt{4x^2+9y^2}} .
        \]

  \item With $g=K^{1/2}+L^{1/2}$, $F=g^4$,
        $F_K=4g^3\cdot\frac12K^{-1/2}=2g^3/K^{1/2}$. At $(4,9)$: $g=5$,
        $K^{1/2}=2$, so $F_K=2\cdot125/2=125$.

  \item Let $r^2=x^2+y^2+z^2$. Then $f=\ln(r^2)$ and $f_{x_i}=2x_i/r^2$.
        So $\nabla f=(2/(x^2+y^2+z^2))(x,y,z)^T$.

  \item Use $\ln\sqrt{u}=\frac12\ln u$: $f=\frac12\ln(x^2+y^2)$. So
        $f_x=x/(x^2+y^2)$, $f_y=y/(x^2+y^2)$.

  \item $\nabla C=(4,9)^T$. With $Q=x^{1/2}y^{1/2}$: $Q_x=\frac12(y/x)^{1/2}$
        and $Q_y=\frac12(x/y)^{1/2}$, so
        \[
          \nabla Q=\tfrac12\left((y/x)^{1/2},\,(x/y)^{1/2}\right)^{T} .
        \]
        The condition $\nabla C=\lambda\nabla Q$ requires
        $4/9=Q_x/Q_y=y/x$, i.e.\ $y=4x/9$. Along this locus
        $\nabla Q\propto(2/3,3/2)^T\propto(4,9)^T=\nabla C$. \checkmark

  \item Chain rule: $f_x=5(1+xy)^4\cdot y=5y(1+xy)^4$,
        $f_y=5x(1+xy)^4$.

  \item $f_x=3x^2-2y$, $f_y=-2x+2y$. At $(1,2)$:
        $\nabla f=(3-4,-2+4)^T=(-1,2)^T$. Normalise: $\|\mathbf{v}\|=5$,
        $\mathbf{u}=(3/5,4/5)^T$. Then
        \[
          D_{\mathbf{u}}f=-1\cdot\frac35+2\cdot\frac45
            =\frac{-3+8}{5}=1 .
        \]

  \item $F_K=10\cdot0.3K^{-0.7}L^{0.7}=3(L/K)^{0.7}$;
        $F_L=7(K/L)^{0.3}$. At $K=L=100$: $F_K=3$, $F_L=7$. And
        $F=10\cdot100^{0.3}\cdot100^{0.7}=10\cdot100=1000$. Then
        $KF_K+LF_L=100\cdot3+100\cdot7=1000=F$. \checkmark (homogeneity
        degree $r=0.3+0.7=1$).

  \item $f_x=y\cos(xy)-\sin(x-y)$, $f_y=x\cos(xy)+\sin(x-y)$.

  \item $f_x=2xy-y^2+3$, $f_y=x^2-2xy-2$. At $(1,-1)$: $f_x=-2-1+3=0$,
        $f_y=1+2-2=1$. So $\nabla f(1,-1)=(0,1)^T$.

  \item $U_x=\alpha x^{\alpha-1}y^{1-\alpha}$,
        $U_y=(1-\alpha)x^\alpha y^{-\alpha}$. So
        \[
          \mathrm{MRS}_{xy}=\frac{U_x}{U_y}
            =\frac{\alpha x^{\alpha-1}y^{1-\alpha}}{(1-\alpha)x^\alpha y^{-\alpha}}
            =\frac{\alpha}{1-\alpha}\cdot\frac{y}{x} . \checkmark
        \]

  \item $f_x=ye^{xy}$, $f_y=xe^{xy}$. At $(1,0)$: $\nabla f=(0,1)^T$.
        Normalise: $\|\mathbf{v}\|=5$, $\mathbf{u}=(-3/5,4/5)^T$. Then
        $D_{\mathbf{u}}f=0\cdot(-3/5)+1\cdot(4/5)=4/5$.

  \item Write $f=e^{y\ln x}$. Then $f_x=e^{y\ln x}\cdot y/x=yx^{y-1}$ and
        $f_y=e^{y\ln x}\cdot\ln x=x^y\ln x$.

  \item On $\{2K<3L\}$ the binding factor is capital: $F=2K$, so
        $F_K=2$, $F_L=0$. On $\{2K>3L\}$ labour binds: $F=3L$, so
        $F_K=0$, $F_L=3$. (Non-differentiable on $\{2K=3L\}$.)

  \item $f_x=yz+2x$, $f_y=xz$, $f_z=xy-2z$. At $(2,-1,3)$:
        $f_x=-3+4=1$, $f_y=6$, $f_z=-2-6=-8$. So
        $\nabla f=(1,6,-8)^T$.

  \item $\nabla f=(2x,8y)^T$. At $(3,1)$: $\nabla f=(6,8)^T$. Steepest
        descent is $-\nabla f/\|\nabla f\|$. $\|\nabla f\|=10$, so unit
        descent direction is $(-3/5,-4/5)^T$, with rate of decrease
        $-\|\nabla f\|=-10$ (i.e.\ $f$ decreases at rate 10).

  \item $f_x=e^y-ye^x$, $f_y=xe^y-e^x$.

  \item Quotient rule with $D=x^2+y^2$:
        \[
          f_x=\frac{D-x\cdot2x}{D^2}=\frac{y^2-x^2}{(x^2+y^2)^2}, \qquad
          f_y=\frac{-x\cdot2y}{D^2}=\frac{-2xy}{(x^2+y^2)^2} .
        \]

  \item $U_x=1$, $U_y=1/\sqrt{y}$. So $\mathrm{MRS}_{xy}=U_x/U_y
        =\sqrt{y}$ (depends only on $y$, quasi-linear utility).

  \item $f_x=2x$, $f_y=2y$, $f_z=-2z$. At $(2,-1,2)$:
        $\nabla f=(4,-2,-4)^T$. Normalise: $\|\mathbf{v}\|=\sqrt{4+4+1}=3$,
        $\mathbf{u}=(2/3,2/3,1/3)^T$. Then
        \[
          D_{\mathbf{u}}f=4\cdot\frac23+(-2)\cdot\frac23+(-4)\cdot\frac13
            =\frac{8-4-4}{3}=0 .
        \]
        (We are moving tangent to the level surface.)

  \item $f_x=\sec^2(x+y)+\sec^2(x-y)$, $f_y=\sec^2(x+y)-\sec^2(x-y)$.

  \item Let $S=\alpha K+\beta L$, $F=S^r$. Then $F_K=rS^{r-1}\alpha$,
        $F_L=rS^{r-1}\beta$. Hence
        \[
          KF_K+LF_L=rS^{r-1}(\alpha K+\beta L)=rS^{r-1}\cdot S=rS^r=rF.
            \checkmark
        \]

  \item $\nabla f=(-e^{-x}\sin y+e^x\cos y,\ e^{-x}\cos y-e^x\sin y)^T$.

  \item $f_x=x/\sqrt{x^2+y^2}$, $f_y=y/\sqrt{x^2+y^2}$. At $(3,4)$:
        $\sqrt{x^2+y^2}=5$, $\nabla f=(3/5,4/5)^T$. The direction
        $\mathbf{v}=(1,0)$ is already unit. So
        $D_{\mathbf{u}}f=3/5\cdot1+4/5\cdot0=3/5$.

  \item On $\{2x<3y\}$: $U=2x$, so $U_x=2$, $U_y=0$. On $\{2x>3y\}$:
        $U=3y$, so $U_x=0$, $U_y=3$. (Non-differentiable on the ray
        $2x=3y$.)

  \item Let $L=\ln(x+y)$, $D=x-y$. By the quotient rule:
        \[
          f_x=\frac{(1/(x+y))(x-y)-\ln(x+y)\cdot1}{(x-y)^2}
            =\frac{1/(x+y)\cdot(x-y)-\ln(x+y)}{(x-y)^2} .
        \]
        Simplifying
        $f_x=\dfrac{(x-y)/(x+y)-\ln(x+y)}{(x-y)^2}$. Similarly,
        \[
          f_y=\frac{(1/(x+y))(x-y)-\ln(x+y)\cdot(-1)}{(x-y)^2}
            =\frac{(x-y)/(x+y)+\ln(x+y)}{(x-y)^2} .
        \]

  \item $F_K=A\alpha K^{\alpha-1}L^\beta$, $F_L=A\beta K^\alpha L^{\beta-1}$.
        So
        \[
          \mathrm{MRTS}_{LK}=\frac{F_L}{F_K}=\frac{\beta}{\alpha}\cdot
            \frac{K}{L} .
        \]
        This depends only on the ratio $K/L$ and the exponent ratio
        $\beta/\alpha$: it is scale-invariant.

  \item $\nabla f=(2x,-2y)^T$. At $(2,-1)$: $\nabla f=(4,2)^T$. Direction
        of steepest ascent: $\mathbf{u}^*=(4,2)/\sqrt{20}=(2,1)/\sqrt5$.
        Maximum rate $=\|\nabla f\|=\sqrt{20}=2\sqrt5$.

  \item $f_x=\sinh(x+y)+\cosh(x-y)$, $f_y=\sinh(x+y)-\cosh(x-y)$.

  \item $\nabla f=(e^{x+y+z},e^{x+y+z},e^{x+y+z})^T$. At $(0,0,0)$:
        $\nabla f=(1,1,1)^T$. Normalise $\mathbf{v}=(1,1,1)$:
        $\|\mathbf{v}\|=\sqrt3$, $\mathbf{u}=(1,1,1)/\sqrt3$. Then
        $D_{\mathbf{u}}f=(1+1+1)/\sqrt3=\sqrt3$.

  \item $f_x=(\ln y)/x$, $f_y=(\ln x)/y$.

  \item $F_K=1+L^{1/2}/K^{1/2}$, $F_L=1+K^{1/2}/L^{1/2}$. From
        $F=(K^{1/2}+L^{1/2})^2$ with $g=K^{1/2}+L^{1/2}$:
        $F_K=g/K^{1/2}=(K^{1/2}+L^{1/2})/K^{1/2}=1+L^{1/2}/K^{1/2}$.
        \checkmark Same for $F_L$.

  \item $f_x=2xy$, $f_y=x^2+3y^2$. At $(-1,2)$: $\nabla f=(-4,1+12)^T
        =(-4,13)^T$. Normalise: $\|\mathbf{v}\|=5$,
        $\mathbf{u}=(4/5,-3/5)^T$. Then
        \[
          D_{\mathbf{u}}f=-4\cdot\frac45+13\cdot\left(-\frac35\right)
            =-\frac{16+39}{5}=-11 .
        \]

  \item $U_x=\alpha/x$, $U_y=\beta/y$.
        $\mathrm{MRS}_{xy}=\dfrac{\alpha/x}{\beta/y}=\dfrac{\alpha y}{\beta x}$.
        At $(\alpha,\beta)$: $\mathrm{MRS}=(\alpha\cdot\beta)/(\beta\cdot\alpha)=1$.

  \item $f=(x^2+y^2)^{-1/2}$. Chain rule:
        \[
          f_x=-\frac12(x^2+y^2)^{-3/2}\cdot2x=\frac{-x}{(x^2+y^2)^{3/2}},
          \qquad f_y=\frac{-y}{(x^2+y^2)^{3/2}} .
        \]
        So $\nabla f=-(x,y)^T/(x^2+y^2)^{3/2}$.

  \item Chain rule with $u=xy$:
        \[
          f_x=\frac{y}{1+x^2y^2}, \qquad f_y=\frac{x}{1+x^2y^2} .
        \]

  \item Let $S=K^\rho+L^\rho$, $F=S^{1/\rho}$. Then
        $F_K=\frac1\rho S^{1/\rho-1}\cdot\rho K^{\rho-1}
        =K^{\rho-1}S^{1/\rho-1}$. Since $S^{1/\rho}=F$, we have $S=F^\rho$
        and so $S^{1/\rho-1}=F^{1-\rho}\cdot F^{-1}\cdot F=F^{1-\rho}$.
        Thus $F_K=K^{\rho-1}F^{1-\rho}=(F/K)^{1-\rho}=(K/F)^{\rho-1}$.
        \checkmark

  \item $f_x=yz$, $f_y=xz$, $f_z=xy$. At $(2,3,-1)$:
        $\nabla f=(3\cdot(-1),2\cdot(-1),2\cdot3)^T=(-3,-2,6)^T$.
        Normalise $\mathbf{v}=(2,-1,2)$: $\|\mathbf{v}\|=3$,
        $\mathbf{u}=(2/3,-1/3,2/3)^T$. Then
        \[
          D_{\mathbf{u}}f=-3\cdot\frac23+(-2)\cdot\left(-\frac13\right)
            +6\cdot\frac23=-2+\frac23+4=2+\frac23=\frac83 .
        \]
\end{solutions}

\section{Additional Problems}

\begin{problems}
  \item Compute $f_x$ and $f_y$ for $f(x,y)=3x^5-2x^3y^2+4y^4$.
  \item For $F(K,L)=20\,K^{1/3}L^{2/3}$, find $F_K$ and $F_L$ at
        $(K,L)=(8,27)$.
  \item Compute $f_x$ and $f_y$ for $f(x,y)=e^{xy}$.
  \item Compute $\nabla f$ for $f(x,y)=2x^2y-3xy^3+5$.
  \item Compute the directional derivative of $f(x,y)=x^2+2xy-y^2$ at
        $(1,2)$ in the direction $\mathbf{v}=(3,4)$.
  \item For $F(K,L)=8K^{3/4}L^{1/4}$, compute $\mathrm{MP}_K$ and
        $\mathrm{MP}_L$ at $(K,L)=(16,16)$.
  \item Compute $f_x$ and $f_y$ for $f(x,y)=\ln(2x+3y)$.
  \item For $U(x,y)=3\ln(x)+4\ln(y)$, derive $\mathrm{MRS}_{xy}$ and
        evaluate at $(2,3)$.
  \item Compute $f_x$ and $f_y$ for $f(x,y)=(x+y^2)^4$.
  \item Compute $\nabla f$ for $f(x,y,z)=xy+yz+zx$.
  \item For a CES production $F(K,L)=5(K^{1/3}+L^{1/3})^3$, find $F_K$.
  \item Compute the directional derivative of $f(x,y)=x^2y$ at $(2,1)$
        in the direction $\mathbf{v}=(1,-1)$.
  \item For $U(x,y)=x^\alpha y^\beta$ (general Cobb--Douglas,
        $\alpha,\beta>0$), derive $\mathrm{MRS}_{xy}$ and show it depends
        only on the ratio $y/x$.
  \item Compute $f_x$ and $f_y$ for $f(x,y)=\dfrac{x-y}{x+y}$.
  \item For $f(K,L)=K^{1/2}L^{1/2}$ at $(K,L)=(16,9)$, find the unit
        direction of steepest ascent and the maximum rate of increase.
  \item Classify the open positive orthant
        $Q=\{(x,y)\in\R^2 : x>0,\ y>0\}$: is it open, closed, or
        neither? Is it bounded? Is it compact?
  \item Classify the budget set
        $B=\{(x,y)\in\R^2 : 2x+3y\le12,\ x\ge0,\ y\ge0\}$: open,
        closed, or neither? Bounded? Compact?
  \item Classify $M=\{(x,y)\in\R^2 : xy\ge1,\ x>0\}$: open, closed, or
        neither? Bounded? Compact? \emph{Hint:} check whether a
        sequence in $M$ can converge to a point with $x=0$.
  \item Classify the open disc
        $D=\{(x,y)\in\R^2 : x^2+y^2<4\}$: open, closed, or neither?
        Bounded? Compact?
  \item Classify (a) the half-open strip
        $S=\{(x,y)\in\R^2 : 0\le x<1\}$ and (b) the finite set
        $E=\{(0,0),\ (1,2),\ (3,1)\}$: open, closed, or neither?
        Bounded? Compact?
  \item Does the Weierstrass theorem guarantee a maximiser for the
        continuous function $u(x,y)=x^{1/2}y^{1/2}$ (a) on the budget
        set $B=\{(x,y): 2x+3y\le12,\ x\ge0,\ y\ge0\}$, (b) on the open
        positive orthant $Q=\{(x,y): x>0,\ y>0\}$? Justify each
        answer.
\end{problems}

In problems 22--29, for each function: state the domain; compute the
partials and the set on which they exist; decide, with a justification,
where $f$ is differentiable; and describe the graph near the points where
differentiability fails.

\begin{problems}[resume]
  \item $f(x,y)=\sqrt{(x-1)^{2}+(y+2)^{2}}$.
  \item $f(x,y)=\sqrt{4x^{2}+(y-3)^{2}}$.
  \item $f(x,y)=|x-2|+|y|$.
  \item $f(x,y)=\min\{2x,\,5y\}$ on $x,y>0$.
  \item $f(x,y)=\dfrac{3xy}{2x^{2}+y^{2}}$ for $(x,y)\neq(0,0)$, with
        $f(0,0)=0$. Also decide whether $f$ is continuous at the origin.
  \item $f(x,y)=\sqrt{x^{4}+y^{4}}$. The formula for $f_{x}$ breaks down
        at the origin; decide what really happens there.
  \item $f(x,y)=\sqrt{9-x^{2}-y^{2}}$.
  \item $f(x,y)=|x-y|$.
\end{problems}

In problems 30--37, find the tangent plane of $f$ at $P$ (the tangent
hyperplane, for three variables), then evaluate both $f$ and the plane at
the nearby point $Q$ and compare. Give both values to six decimals.

\begin{problems}[resume]
  \item $f(x,y)=\ln(2x+y^{2})$, \;$P=(0,1)$, \;$Q=(0.05,\,0.98)$.
  \item $f(x,y)=\mathrm{e}^{x-2y}$, \;$P=(2,1)$, \;$Q=(2.1,\,1.02)$.
  \item $f(x,y)=xy^{2}\,\mathrm{e}^{x-1}$, \;$P=(1,2)$,
        \;$Q=(0.98,\,2.03)$.
  \item $f(x,y)=\dfrac{x+y}{x-y}$, \;$P=(3,1)$, \;$Q=(3.1,\,0.9)$.
  \item $F(K,L)=K^{1/4}L^{3/4}$, \;$P=(16,1)$, \;$Q=(16.5,\,1.1)$; read
        the plane as an output estimate.
  \item $f(x,y,z)=\ln(x+2y+3z)$, \;$P=(2,-1,\tfrac13)$,
        \;$Q=(2.02,\,-0.99,\,0.35)$.
  \item $f(x,y,z)=x\,\mathrm{e}^{yz}$, \;$P=(2,0,5)$,
        \;$Q=(2.05,\,0.01,\,4.9)$.
  \item $f(x,y,z)=\dfrac{xy}{z}$, \;$P=(2,3,6)$, \;$Q=(2.1,\,2.9,\,6.2)$.
\end{problems}

In problems 38--46, compute the partial derivatives, and give the gradient
(a column) and the derivative (a row) at the stated point.

\begin{problems}[resume]
  \item $f(x_{1},\ldots,x_{n})=\sum_{i=1}^{n}\ln(1+x_{i}^{2})$, at
        $(1,1,\ldots,1)$.
  \item $f(x_{1},\ldots,x_{n})=\Bigl(\sum_{i=1}^{n}x_{i}\Bigr)^{3}$; then
        for $n=4$ at $(1,2,3,4)$.
  \item $f(x_{1},\ldots,x_{n})=\sum_{i=1}^{n}x_{i}^{\,i+1}$, at
        $(1,1,\ldots,1)$; then for $n=3$ at $(2,1,-1)$.
  \item $f(x_{1},\ldots,x_{n})=\sum_{j=1}^{n}\dfrac{j^{2}}{x_{j}^{3}}$ with
        all $x_{j}>0$, at $(1,1,\ldots,1)$.
  \item $f(x_{1},\ldots,x_{n})=\sum_{i=1}^{n}x_{i}^{2}
        -\dfrac1n\Bigl(\sum_{i=1}^{n}x_{i}\Bigr)^{2}$; show that the
        gradient vanishes exactly when all $x_{i}$ are equal, then evaluate
        for $n=3$ at $(1,2,6)$.
  \item $f(x,y,z)=y^{2}\ln(x^{2}+3yz)$, at $P=(1,1,1)$.
  \item $f(x,y,z)=x^{3}\,\mathrm{e}^{y^{2}-2xz}$, at $P=(1,0,1)$.
  \item $f(x,y)=2^{x}3^{y}+5^{xy}$, at $(1,1)$.
  \item $f(x_{1},\ldots,x_{n})=\sum_{i=1}^{n}3^{\,i x_{i}}$, at
        $(0,0,\ldots,0)$.
\end{problems}

In problems 47--52, decide differentiability directly from the definition,
by forming the remainder and estimating it, as in the third worked example.

\begin{problems}[resume]
  \item $f(x,y)=3xy$, at every point $(a,b)$.
  \item $f(x,y)=x^{2}+y^{2}$, first at the origin, then at every point.
  \item $f(x,y)=xy^{2}$, at every point $(a,b)$.
  \item $f(x,y)=y\,\mathrm{e}^{3x}$, at every point $(a,b)$.
  \item At the origin: (a) $f(x,y)=(x^{2}+4y^{2})^{1/3}$;
        (b) $g(x,y)=(x^{2}+y^{2})^{3/4}$.
  \item $f(x,y)=|xy|$: at the origin, and at the point $(2,0)$.
\end{problems}

In problems 53--54, $f(0,0)=0$. Show with the difference quotient that both
partials exist at the origin and are $0$; compute the limit along every
straight line through the origin and along the stated parabola; and decide
continuity and differentiability at the origin.

\begin{problems}[resume]
  \item $f(x,y)=\dfrac{x^{2}y}{x^{4}+y^{2}}$ for $(x,y)\neq(0,0)$; parabola
        $y=x^{2}$.
  \item $f(x,y)=\dfrac{2xy^{2}}{x^{2}+y^{4}}$ for $(x,y)\neq(0,0)$;
        parabola $x=y^{2}$.
\end{problems}

Problems 55--56 refer to the three surfaces below, each drawn over a square
around its special point.

\begin{gallery}
\begin{tikzpicture}
  \begin{axis}[surfaxis, title={(a)}, z post scale=1.1]
    \addplot3[surf, domain=-1:3, y domain=-2:2]{(x-1)^2+2*y^2};
  \end{axis}
\end{tikzpicture}\hfill
\begin{tikzpicture}
  \begin{axis}[surfaxis, title={(b)}, z post scale=1.1]
    \addplot3[surf, domain=-1:3, y domain=-2:2]{sqrt((x-1)^2+4*y^2)};
  \end{axis}
\end{tikzpicture}\hfill
\begin{tikzpicture}
  \begin{axis}[surfaxis, title={(c)}, z post scale=1.1]
    \addplot3[surf, domain=-2:2, y domain=-2:2]{2-abs(x)-abs(y)};
  \end{axis}
\end{tikzpicture}
\end{gallery}

\begin{problems}[resume]
  \item For each surface (a)--(c), say from the picture alone where the
        function is differentiable and where it is not, and why.
  \item Surface (b) has its tip above the point $(1,0)$, and its contours
        are ellipses twice as long in $x$ as in $y$. Which of the formulas
        $(x-1)^{2}+4y^{2}$, $\sqrt{(x-1)^{2}+4y^{2}}$ and
        $\sqrt{4(x-1)^{2}+y^{2}}$ can it be the graph of? Confirm the
        verdict of the previous problem for that formula with the
        difference quotient at $(1,0)$.
\end{problems}

\section{Additional Solutions}

\begin{solutions}
  \item $f_x=15x^4-6x^2y^2$, $f_y=-4x^3y+16y^3$.

  \item $F_K=20\cdot\frac13K^{-2/3}L^{2/3}=\frac{20}{3}(L/K)^{2/3}$;
        $F_L=20\cdot\frac23K^{1/3}L^{-1/3}=\frac{40}{3}(K/L)^{1/3}$. At
        $(8,27)$: $(L/K)^{2/3}=(27/8)^{2/3}=9/4$, so $F_K=15$;
        $(K/L)^{1/3}=(8/27)^{1/3}=2/3$, so
        $F_L=\frac{40}{3}\cdot\frac23=\frac{80}{9}$.

  \item $f_x=ye^{xy}$, $f_y=xe^{xy}$.

  \item $\nabla f=(4xy-3y^3,\ 2x^2-9xy^2)^T$.

  \item $f_x=2x+2y$, $f_y=2x-2y$. At $(1,2)$: $\nabla f=(6,-2)^T$.
        Normalise: $\|\mathbf{v}\|=5$, $\mathbf{u}=(3/5,4/5)^T$. Then
        \[
          D_{\mathbf{u}}f=6\cdot\frac35+(-2)\cdot\frac45
            =\frac{18-8}{5}=2 .
        \]

  \item $F_K=8\cdot\frac34K^{-1/4}L^{1/4}=6(L/K)^{1/4}$,
        $F_L=2(K/L)^{3/4}$. At $(16,16)$: $\mathrm{MP}_K=6$,
        $\mathrm{MP}_L=2$.

  \item $f_x=\dfrac{2}{2x+3y}$, $f_y=\dfrac{3}{2x+3y}$.

  \item $U_x=3/x$, $U_y=4/y$.
        $\mathrm{MRS}_{xy}=U_x/U_y=\dfrac{3y}{4x}$. At $(2,3)$:
        $\mathrm{MRS}=\dfrac98$.

  \item Chain: $f_x=4(x+y^2)^3$,
        $f_y=4(x+y^2)^3\cdot2y=8y(x+y^2)^3$.

  \item $f_x=y+z$, $f_y=x+z$, $f_z=x+y$. So
        $\nabla f=(y+z,\ x+z,\ x+y)^T$.

  \item Let $g=K^{1/3}+L^{1/3}$, so $F=5g^3$. Then
        $F_K=5\cdot3g^2\cdot\frac13K^{-2/3}=5g^2K^{-2/3}
        =\dfrac{5(K^{1/3}+L^{1/3})^2}{K^{2/3}}$.

  \item $f_x=2xy$, $f_y=x^2$. At $(2,1)$: $\nabla f=(4,4)^T$. Normalise:
        $\|\mathbf{v}\|=\sqrt2$, $\mathbf{u}=(1/\sqrt2,-1/\sqrt2)^T$.
        Then
        \[
          D_{\mathbf{u}}f=4\cdot\frac{1}{\sqrt2}
            +4\cdot\left(-\frac{1}{\sqrt2}\right)=0 .
        \]
        (Tangent to the level curve.)

  \item $U_x=\alpha x^{\alpha-1}y^\beta$,
        $U_y=\beta x^\alpha y^{\beta-1}$. So
        \[
          \mathrm{MRS}_{xy}=\frac{U_x}{U_y}
            =\frac{\alpha}{\beta}\cdot\frac{y}{x} .
        \]
        Depends only on $y/x$ (and the exponent ratio), not on the
        scale. \checkmark

  \item Quotient rule with $D=x+y$:
        \[
          f_x=\frac{D-(x-y)}{D^2}=\frac{2y}{(x+y)^2}, \qquad
          f_y=\frac{-D-(x-y)}{D^2}=\frac{-2x}{(x+y)^2} .
        \]

  \item $f_K=\frac12K^{-1/2}L^{1/2}$, $f_L=\frac12K^{1/2}L^{-1/2}$. At
        $(16,9)$: $f_K=\frac12\cdot\frac34=3/8$;
        $f_L=\frac12\cdot\frac43=2/3$. So $\nabla f=(3/8,2/3)^T$.
        Magnitude:
        $\|\nabla f\|=\sqrt{9/64+4/9}=\sqrt{81/576+256/576}
        =\sqrt{337/576}=\sqrt{337}/24\approx0.765$. Unit direction of
        steepest ascent:
        $\mathbf{u}^*=(3/8,2/3)^T/(\sqrt{337}/24)=(9,16)^T/\sqrt{337}$.

  \item \textbf{Open.} For any $(x,y)\in Q$ the ball of radius
        $r=\min\{x,y\}>0$ around $(x,y)$ stays inside $Q$, so every
        point is interior. \textbf{Not closed:} the boundary points
        $(x,0)$ and $(0,y)$ with $x,y\ge0$ (the coordinate half-axes,
        including the origin) do not belong to $Q$. \textbf{Not
        bounded:} $Q$ contains $(t,t)$ for every $t>0$, and
        $\|(t,t)\|=t\sqrt2\to\infty$. Hence $Q$ is open, unbounded and
        \textbf{not compact}.

  \item \textbf{Closed.} $B$ is the intersection of the three sets
        $\{2x+3y\le12\}$, $\{x\ge0\}$ and $\{y\ge0\}$, each defined by
        a non-strict inequality on a continuous function; each is
        closed, and an intersection of closed sets is closed (every
        boundary point of $B$ satisfies all three inequalities, so it
        lies in $B$). \textbf{Not open:} $(0,0)\in B$, but every ball
        around the origin contains points with negative coordinates,
        which are outside $B$; so $(0,0)$ is not interior.
        \textbf{Bounded:} from the constraints, $0\le x\le6$ and
        $0\le y\le4$, so $B\subseteq B_8(\mathbf{0})$. Closed and
        bounded in $\R^2$: $B$ is \textbf{compact}.

  \item First note the constraint $x>0$ together with $xy\ge1$ forces
        $y\ge1/x>0$. \textbf{Closed:} suppose
        $(x_n,y_n)\in M$ with $(x_n,y_n)\to(a,b)$. If $a=0$, then
        $y_n\to b$ finite would give $x_ny_n\to0<1$, contradicting
        $x_ny_n\ge1$; so $a>0$, and passing to the limit in
        $x_ny_n\ge1$ gives $ab\ge1$, hence $(a,b)\in M$. $M$ contains
        all its limit (boundary) points, so it is closed: even
        though it is written with a strict inequality, that inequality
        is never active. \textbf{Not open:} $(1,1)\in M$ lies on the
        hyperbola $xy=1$, and every ball around it contains points
        with $xy<1$. \textbf{Not bounded:} $(t,t)\in M$ for all
        $t\ge1$. Hence $M$ is closed, unbounded and \textbf{not
        compact}.

  \item \textbf{Open.} For $(x,y)\in D$ put
        $r=2-\sqrt{x^2+y^2}>0$; the ball $B_r(x,y)$ lies in $D$ by the
        triangle inequality. \textbf{Not closed:} the boundary circle
        $\{x^2+y^2=4\}$ is not contained in $D$ (e.g.\ $(2,0)\notin
        D$). \textbf{Bounded:} $D\subseteq B_3(\mathbf{0})$. Bounded
        but not closed: $D$ is \textbf{not compact}.

  \item (a) \textbf{Neither open nor closed.} The points $(0,y)$
        belong to $S$ but are boundary points (every ball around them
        contains points with $x<0$), so $S$ is not open; the boundary
        points $(1,y)$ do not belong to $S$, so $S$ is not closed.
        \textbf{Unbounded:} $(0,t)\in S$ for all $t$. Hence
        \textbf{not compact}. (b) \textbf{Closed:} around any point
        not in $E$ there is a ball avoiding all three points (a
        finite set always contains all its boundary points:
        indeed $E$ is its own boundary). \textbf{Not open:} no ball
        around $(0,0)$ lies inside $E$. \textbf{Bounded:}
        $E\subseteq B_4(\mathbf{0})$. Closed and bounded: $E$ is
        \textbf{compact}.

  \item (a) \textbf{Yes.} $B$ is compact (closed and bounded, see the
        budget-set problem above) and $u$ is continuous on $B$, so by
        the Weierstrass theorem $u$ attains both a maximum and a
        minimum on $B$; a maximiser exists. (b) \textbf{No.} $Q$ is
        not compact (it is neither closed nor bounded), so Weierstrass
        does not apply: and indeed no maximiser exists:
        $u(t,t)=t\to\infty$ as $t\to\infty$, so $u$ is unbounded above
        on $Q$.

  \item \textbf{Domain:} the radicand is a sum of squares, never
        negative, so $D=\R^{2}$. \textbf{Partials:} where the radicand is
        positive, i.e.\ for $(x,y)\neq(1,-2)$, the chain rule gives
        \[
          f_{x}=\frac{x-1}{\sqrt{(x-1)^{2}+(y+2)^{2}}},\qquad
          f_{y}=\frac{y+2}{\sqrt{(x-1)^{2}+(y+2)^{2}}} .
        \]
        At $(1,-2)$ use the quotient:
        $\bigl(f(1+h,-2)-f(1,-2)\bigr)/h=\sqrt{h^{2}}/h=|h|/h$, which is $1$
        for $h>0$ and $-1$ for $h<0$, so it has no limit. $f_{x}(1,-2)$
        does not exist, and in the same way neither does $f_{y}(1,-2)$. The
        partials exist exactly on $\R^{2}\setminus\{(1,-2)\}$.
        \textbf{Differentiability:} on that open set the partials are
        continuous ($f$ is $C^{1}$), so $f$ is differentiable there. At
        $(1,-2)$ it is not, since differentiability would force the
        partials to exist. Note that $f$ \emph{is} continuous at $(1,-2)$:
        continuity is not enough. \textbf{Graph:} a circular cone with its
        vertex at $(1,-2,0)$, opening upwards with slope $1$ in every
        direction; the contours are circles of radius $c$ about $(1,-2)$.
        At the sharp tip there is no tangent plane.

  \item \textbf{Domain:} $\R^{2}$, the radicand being a sum of squares.
        \textbf{Partials:} for $(x,y)\neq(0,3)$,
        \[
          f_{x}=\frac{4x}{\sqrt{4x^{2}+(y-3)^{2}}},\qquad
          f_{y}=\frac{y-3}{\sqrt{4x^{2}+(y-3)^{2}}} .
        \]
        At $(0,3)$: $f(h,3)-f(0,3)=\sqrt{4h^{2}}=2|h|$, so the quotient
        $2|h|/h=\pm2$ has no limit; and $f(0,3+h)-f(0,3)=|h|$ gives
        $\pm1$. Neither partial exists at $(0,3)$; both exist everywhere
        else. \textbf{Differentiability:} $C^{1}$, hence differentiable, on
        $\R^{2}\setminus\{(0,3)\}$; not differentiable at $(0,3)$.
        \textbf{Graph:} an elliptic cone with apex $(0,3,0)$. The contours
        $4x^{2}+(y-3)^{2}=c^{2}$ are ellipses about $(0,3)$ with semi-axes
        $c/2$ in the $x$-direction and $c$ in the $y$-direction: slope $2$
        across $x$, slope $1$ along $y$, and a sharp point at the apex.

  \item \textbf{Domain:} $\R^{2}$. \textbf{Partials:} for $x\neq2$,
        $f_{x}=1$ if $x>2$ and $f_{x}=-1$ if $x<2$; at a point $(2,y_{0})$,
        $\bigl(f(2+h,y_{0})-f(2,y_{0})\bigr)/h=|h|/h=\pm1$, so $f_{x}$ does
        not exist anywhere on the line $x=2$. Likewise $f_{y}=\pm1$ (the
        sign of $y$) for $y\neq0$, and $f_{y}$ does not exist on the line
        $y=0$. Both partials exist exactly off the two lines $x=2$ and
        $y=0$. \textbf{Differentiability:} off the two lines $f$ is, near
        each point, one of the four linear functions $\pm(x-2)\pm y$, so it
        is differentiable there; on the lines it is not, since a partial
        is missing. \textbf{Graph:} an upside-down square pyramid with its
        vertex at $(2,0,0)$; the contours $|x-2|+|y|=c$ are diamonds with
        corners $(2\pm c,0)$ and $(2,\pm c)$, and the surface is creased
        above the lines $x=2$ and $y=0$.

  \item \textbf{Domain:} the open quadrant $x,y>0$, as given.
        \textbf{Partials:} where $2x<5y$ (above the ray
        $y=\tfrac25x$), $f=2x$, so $f_{x}=2$, $f_{y}=0$. Where $2x>5y$,
        $f=5y$, so $f_{x}=0$, $f_{y}=5$. On the ray, at $(x_{0},y_{0})$ with
        $2x_{0}=5y_{0}$:
        \[
          f(x_{0}+h,y_{0})-f(x_{0},y_{0})=\min\{2x_{0}+2h,\,5y_{0}\}-5y_{0}
          =\min\{2h,\,0\},
        \]
        so the quotient is $0$ for $h>0$ and $2$ for $h<0$: no limit, no
        $f_{x}$. Similarly $f(x_{0},y_{0}+k)-f(x_{0},y_{0})=\min\{0,5k\}$
        gives one-sided slopes $0$ and $5$: no $f_{y}$. The partials exist
        exactly off the ray $y=\tfrac25x$. \textbf{Differentiability:}
        differentiable off the ray (locally linear), not on it.
        \textbf{Graph:} two planes, $z=2x$ and $z=5y$, meeting in a ridge
        above the ray: a roof. The contours are L-shaped corners, the level
        $c$ having its corner at $(c/2,\,c/5)$ on the ray, which is the
        Leontief isoquant picture.

  \item \textbf{Domain:} the denominator $2x^{2}+y^{2}$ vanishes only at
        the origin, where $f$ is defined separately, so $D=\R^{2}$.
        \textbf{Partials off the origin}, quotient rule with
        $N=2x^{2}+y^{2}$:
        \[
          f_{x}=\frac{3y\,N-3xy\cdot4x}{N^{2}}
               =\frac{3y\,(y^{2}-2x^{2})}{(2x^{2}+y^{2})^{2}},\qquad
          f_{y}=\frac{3x\,N-3xy\cdot2y}{N^{2}}
               =\frac{3x\,(2x^{2}-y^{2})}{(2x^{2}+y^{2})^{2}} .
        \]
        \textbf{At the origin}, from the definition: $f(h,0)=0$ and
        $f(0,h)=0$ for all $h$, so both quotients are $0$ and
        $f_{x}(0,0)=f_{y}(0,0)=0$. The partials exist \emph{everywhere}.
        \textbf{Continuity:} on the line $y=mx$, $x\neq0$,
        \[
          f(x,mx)=\frac{3mx^{2}}{(2+m^{2})x^{2}}=\frac{3m}{2+m^{2}},
        \]
        a constant that depends on $m$; on $y=x$ it is $1\neq f(0,0)$. So
        $f$ is not continuous at the origin, and hence not differentiable
        there, even though both partials exist. Away from the origin it is
        $C^{1}$ and differentiable. \textbf{Graph:} each ray from the origin
        carries its own constant height $3m/(2+m^{2})$, ranging between
        $-\tfrac{3\sqrt2}{4}$ and $\tfrac{3\sqrt2}{4}$ (attained at
        $m=\mp\sqrt2$); all these heights meet above the origin, so the
        surface is a twisted sheet torn at that point.

  \item \textbf{Domain:} $\R^{2}$. \textbf{Partials} off the origin:
        $f_{x}=\dfrac{2x^{3}}{\sqrt{x^{4}+y^{4}}}$,
        $f_{y}=\dfrac{2y^{3}}{\sqrt{x^{4}+y^{4}}}$. At the origin the
        formula fails, so use the quotient:
        $f(h,0)=\sqrt{h^{4}}=h^{2}$, and $h^{2}/h=h\to0$. So
        $f_{x}(0,0)=0$, and likewise $f_{y}(0,0)=0$: the partials exist
        everywhere, despite the broken formula. \textbf{Differentiability
        at the origin:} with $\mathbf{a}=\mathbf{0}$ the remainder is
        $r(\mathbf{h})=\sqrt{h_{1}^{4}+h_{2}^{4}}$. Since
        $h_{1}^{4}+h_{2}^{4}\le(h_{1}^{2}+h_{2}^{2})^{2}$ (the missing
        cross term $2h_{1}^{2}h_{2}^{2}$ is non-negative),
        \[
          0\le\frac{r(\mathbf{h})}{\|\mathbf{h}\|}
          \le\frac{h_{1}^{2}+h_{2}^{2}}{\|\mathbf{h}\|}=\|\mathbf{h}\|\to0 .
        \]
        So $f$ is differentiable at the origin with tangent plane $z=0$,
        and, being $C^{1}$ elsewhere, differentiable on all of $\R^{2}$.
        \textbf{Graph:} a bowl with a flat bottom at the origin; its
        contours $x^{4}+y^{4}=c^{2}$ are rounded squares. Nothing fails.

  \item \textbf{Domain:} $9-x^{2}-y^{2}\ge0$, the closed disc of radius $3$.
        \textbf{Partials:} for $x^{2}+y^{2}<9$,
        \[
          f_{x}=\frac{-x}{\sqrt{9-x^{2}-y^{2}}},\qquad
          f_{y}=\frac{-y}{\sqrt{9-x^{2}-y^{2}}},
        \]
        continuous on the open disc. On the boundary circle the radicand
        is $0$, the formulas divide by $0$, and, more basically, the points
        are not interior: every neighbourhood of a boundary point contains
        points outside the disc, where $f$ is not defined, so the two-sided
        difference quotients needed for the partials do not exist. \textbf{Differentiability:}
        $f$ is differentiable exactly on the open disc $x^{2}+y^{2}<9$.
        \textbf{Graph:} the upper hemisphere of radius $3$. Its tangent
        plane is horizontal at the top $(0,0,3)$, tilts more and more
        steeply outwards, and becomes vertical at the rim, where
        $\|\grad f\|\to\infty$.

  \item \textbf{Domain:} $\R^{2}$. \textbf{Partials:} where $x>y$,
        $f=x-y$, so $f_{x}=1$, $f_{y}=-1$; where $x<y$, $f=y-x$, so
        $f_{x}=-1$, $f_{y}=1$. At a point $(x_{0},x_{0})$ of the diagonal,
        $f(x_{0}+h,x_{0})-f(x_{0},x_{0})=|h|$, and $|h|/h=\pm1$ has no limit;
        the same holds in $y$. The partials exist exactly off the line
        $y=x$. \textbf{Differentiability:} differentiable off $y=x$, not on
        it, although $f$ is continuous everywhere. \textbf{Graph:} a
        V-shaped valley whose crease, at height $0$, runs along the line
        $y=x$; the contours are pairs of parallel lines $y=x\pm c$, the
        tilted-trough picture of Chapter~9 with a sharp floor.

  \item $f(P)=\ln(0+1)=0$. Chain rule:
        $f_{x}=\dfrac{2}{2x+y^{2}}$, $f_{y}=\dfrac{2y}{2x+y^{2}}$, so at
        $P$ (denominator $1$): $f_{x}=2$, $f_{y}=2$. Tangent plane:
        \[
          z=0+2(x-0)+2(y-1)=2x+2y-2 .
        \]
        At $Q$: $f(Q)=\ln(0.1+0.9604)=\ln1.0604\approx0.058646$, and
        $L(Q)=0.1+1.96-2=0.060000$. The plane overestimates by about
        $0.001354$.

  \item $f(P)=\mathrm{e}^{0}=1$. $f_{x}=\mathrm{e}^{x-2y}$,
        $f_{y}=-2\mathrm{e}^{x-2y}$, so at $P$: $f_{x}=1$, $f_{y}=-2$.
        Tangent plane:
        \[
          z=1+(x-2)-2(y-1)=x-2y+1 .
        \]
        At $Q$: $f(Q)=\mathrm{e}^{2.1-2.04}=\mathrm{e}^{0.06}\approx1.061837$,
        and $L(Q)=2.1-2.04+1=1.060000$. The plane underestimates by about
        $0.001837$: the exponential is convex, so it lies above all its
        tangent planes.

  \item $f(P)=1\cdot4\cdot\mathrm{e}^{0}=4$. Product rule in $x$, where
        both $x$ and $\mathrm{e}^{x-1}$ depend on $x$:
        \[
          f_{x}=y^{2}\mathrm{e}^{x-1}+xy^{2}\mathrm{e}^{x-1}
               =y^{2}\mathrm{e}^{x-1}(1+x),\qquad
          f_{y}=2xy\,\mathrm{e}^{x-1} .
        \]
        At $P$: $f_{x}=4\cdot1\cdot2=8$, $f_{y}=2\cdot1\cdot2\cdot1=4$.
        Tangent plane:
        \[
          z=4+8(x-1)+4(y-2) .
        \]
        At $Q$: $f(Q)=0.98\cdot2.03^{2}\cdot\mathrm{e}^{-0.02}
        =0.98\cdot4.1209\cdot0.980199\approx3.958515$, and
        $L(Q)=4-0.16+0.12=3.960000$. The plane is off by about
        $0.001485$.

  \item $f(P)=\dfrac{4}{2}=2$. Quotient rule with denominator $(x-y)^{2}$:
        \[
          f_{x}=\frac{(x-y)-(x+y)}{(x-y)^{2}}=\frac{-2y}{(x-y)^{2}},\qquad
          f_{y}=\frac{(x-y)+(x+y)}{(x-y)^{2}}=\frac{2x}{(x-y)^{2}} .
        \]
        At $P$: $f_{x}=-\tfrac24=-\tfrac12$, $f_{y}=\tfrac64=\tfrac32$.
        Tangent plane:
        \[
          z=2-\tfrac12(x-3)+\tfrac32(y-1) .
        \]
        At $Q$: $f(Q)=\dfrac{4.0}{2.2}\approx1.818182$, and
        $L(Q)=2-0.05-0.15=1.800000$. The error, about $0.018182$, is
        several times larger than in the neighbouring problems. The step
        keeps $x+y=4$ but moves the denominator $x-y$ from $2$ to $2.2$,
        and along that direction $f=4/(x-y)$ is strongly curved. The plane
        is a local approximation, and its quality depends on the curvature
        as much as on the length of the step.

  \item $F(P)=16^{1/4}\cdot1=2$.
        $F_{K}=\tfrac14K^{-3/4}L^{3/4}$, $F_{L}=\tfrac34K^{1/4}L^{-1/4}$; at
        $P$: $F_{K}=\tfrac14\cdot\tfrac18=\tfrac1{32}$ and
        $F_{L}=\tfrac34\cdot2=\tfrac32$. Tangent plane:
        \[
          z=2+\tfrac1{32}(K-16)+\tfrac32(L-1) .
        \]
        Economically: output $2$ plus marginal product times input change,
        for each input. At $Q$: $L(Q)=2+\tfrac{0.5}{32}+\tfrac32\cdot0.1
        =2+0.015625+0.15=2.165625$, while
        $F(Q)=16.5^{1/4}\cdot1.1^{3/4}\approx2.164789$. The estimate is
        high by about $0.000836$: with exponents summing to $1$ this
        Cobb--Douglas function is concave, so its graph lies below every
        tangent plane.

  \item $f(P)=\ln(2-2+1)=\ln1=0$. With $s=x+2y+3z$:
        $f_{x}=\tfrac1s$, $f_{y}=\tfrac2s$, $f_{z}=\tfrac3s$, and $s=1$ at
        $P$, so $\grad f(P)=(1,2,3)^{T}$. Tangent hyperplane:
        \[
          w=0+(x-2)+2(y+1)+3\bigl(z-\tfrac13\bigr)=x+2y+3z-1 .
        \]
        At $Q$: $s=2.02-1.98+1.05=1.09$, so $f(Q)=\ln1.09\approx0.086178$,
        and $L(Q)=1.09-1=0.090000$. The hyperplane overestimates by about
        $0.003822$, since $\ln$ is concave.

  \item $f(P)=2\mathrm{e}^{0}=2$. $f_{x}=\mathrm{e}^{yz}$,
        $f_{y}=xz\,\mathrm{e}^{yz}$, $f_{z}=xy\,\mathrm{e}^{yz}$; at $P$:
        $f_{x}=1$, $f_{y}=2\cdot5=10$, $f_{z}=2\cdot0=0$. Tangent
        hyperplane:
        \[
          w=2+(x-2)+10(y-0)+0\cdot(z-5)=x+10y .
        \]
        The change in $z$ does not enter at all, because $f_{z}(P)=0$. At
        $Q$: $f(Q)=2.05\,\mathrm{e}^{0.049}\approx2.152952$, and
        $L(Q)=2.05+0.1=2.150000$; error about $0.002952$.

  \item $f(P)=\tfrac66=1$. $f_{x}=\tfrac yz$, $f_{y}=\tfrac xz$,
        $f_{z}=-\tfrac{xy}{z^{2}}$; at $P$: $f_{x}=\tfrac12$,
        $f_{y}=\tfrac13$, $f_{z}=-\tfrac6{36}=-\tfrac16$. Tangent
        hyperplane:
        \[
          w=1+\tfrac12(x-2)+\tfrac13(y-3)-\tfrac16(z-6) .
        \]
        At $Q$: $L(Q)=1+0.05-0.033333-0.033333=0.983333$, and
        $f(Q)=\dfrac{2.1\cdot2.9}{6.2}=\dfrac{6.09}{6.2}\approx0.982258$;
        error about $0.001075$.

  \item Only the term $i=k$ contains $x_{k}$, so
        \[
          \pd{f}{x_{k}}=\frac{2x_{k}}{1+x_{k}^{2}},\qquad k=1,\ldots,n .
        \]
        At $(1,\ldots,1)$ every partial is $\tfrac22=1$:
        $\grad f=(1,1,\ldots,1)^{T}$ and $Df=(1\ \ 1\ \ \cdots\ \ 1)$.

  \item Chain rule with inner function $s=\sum_{i}x_{i}$, whose partial in
        $x_{k}$ is $1$:
        \[
          \pd{f}{x_{k}}=3\Bigl(\sum_{i=1}^{n}x_{i}\Bigr)^{2}\cdot1 ,
        \]
        the same for every $k$. For $n=4$ at $(1,2,3,4)$: $s=10$, so each
        partial is $300$, $\grad f=(300,300,300,300)^{T}$ and
        $Df=(300\ \ 300\ \ 300\ \ 300)$.

  \item The term containing $x_{k}$ is $x_{k}^{k+1}$, so
        $\pd{f}{x_{k}}=(k+1)\,x_{k}^{k}$. At $(1,\ldots,1)$:
        $\grad f=(2,3,\ldots,n+1)^{T}$, $Df=(2\ \ 3\ \ \cdots\ \ n+1)$. For
        $n=3$, $f=x_{1}^{2}+x_{2}^{3}+x_{3}^{4}$ with partials
        $2x_{1},\,3x_{2}^{2},\,4x_{3}^{3}$; at $(2,1,-1)$:
        $\grad f=(4,3,-4)^{T}$ and $Df=(4\ \ 3\ \ {-4})$.

  \item Write $f=\sum_{j}j^{2}x_{j}^{-3}$. The term with $j=k$ gives
        \[
          \pd{f}{x_{k}}=k^{2}\cdot(-3)x_{k}^{-4}=-\frac{3k^{2}}{x_{k}^{4}} .
        \]
        At $(1,\ldots,1)$: $\grad f=-3\,(1,4,9,\ldots,n^{2})^{T}$ and
        $Df=(-3\ \ {-12}\ \ {-27}\ \ \cdots\ \ {-3n^{2}})$.

  \item The first sum contributes $2x_{k}$; the second term, by the chain
        rule as in problem 39, contributes
        $-\tfrac1n\cdot2\sum_{i}x_{i}$. With the mean
        $\bar x=\tfrac1n\sum_{i}x_{i}$:
        \[
          \pd{f}{x_{k}}=2x_{k}-\frac2n\sum_{i=1}^{n}x_{i}=2\,(x_{k}-\bar x) .
        \]
        The gradient is zero exactly when $x_{k}=\bar x$ for every $k$,
        i.e.\ when all $x_{i}$ are equal. For $n=3$ at $(1,2,6)$:
        $\bar x=3$, so $\grad f=(-4,-2,6)^{T}$ and $Df=(-4\ \ {-2}\ \ 6)$.
        (The components always add up to $0$.)

  \item With $s=x^{2}+3yz$: $f_{x}=y^{2}\cdot\dfrac{2x}{s}$, since
        $y^{2}$ is a constant factor; $f_{y}$ needs the product rule,
        because $y$ appears in both factors:
        \[
          f_{y}=2y\ln s+y^{2}\cdot\frac{3z}{s};
        \]
        and $f_{z}=y^{2}\cdot\dfrac{3y}{s}=\dfrac{3y^{3}}{s}$. At
        $P=(1,1,1)$, $s=4$:
        \[
          f_{x}=\tfrac24=\tfrac12,\qquad
          f_{y}=2\ln4+\tfrac34\approx3.522589,\qquad
          f_{z}=\tfrac34 .
        \]
        So $\grad f(P)=\bigl(\tfrac12,\ 2\ln4+\tfrac34,\ \tfrac34\bigr)^{T}$
        and $Df(P)=\bigl(\tfrac12\ \ \ 2\ln4+\tfrac34\ \ \ \tfrac34\bigr)$.

  \item Write $E=\mathrm{e}^{y^{2}-2xz}$. In $x$ both factors move, so
        product and chain rule:
        \[
          f_{x}=3x^{2}E+x^{3}E\cdot(-2z)=x^{2}(3-2xz)\,E .
        \]
        In $y$ and $z$ only the exponential moves:
        $f_{y}=x^{3}E\cdot2y=2x^{3}y\,E$ and
        $f_{z}=x^{3}E\cdot(-2x)=-2x^{4}E$. At $P=(1,0,1)$,
        $E=\mathrm{e}^{-2}$:
        $f_{x}=1\cdot(3-2)\mathrm{e}^{-2}=\mathrm{e}^{-2}$, $f_{y}=0$,
        $f_{z}=-2\mathrm{e}^{-2}$. So
        $\grad f(P)=\mathrm{e}^{-2}(1,0,-2)^{T}\approx(0.135335,\,0,\,-0.270671)^{T}$
        and $Df(P)=\mathrm{e}^{-2}(1\ \ 0\ \ {-2})$.

  \item Use $\partial a^{g}/\partial x=a^{g}\ln a\cdot g_{x}$. In the first
        term $3^{y}$ is a constant factor; in the second the exponent $xy$
        has $x$-derivative $y$:
        \[
          f_{x}=2^{x}3^{y}\ln2+y\,5^{xy}\ln5,\qquad
          f_{y}=2^{x}3^{y}\ln3+x\,5^{xy}\ln5 .
        \]
        At $(1,1)$, $2^{x}3^{y}=6$ and $5^{xy}=5$:
        \[
          \grad f(1,1)=\begin{pmatrix}6\ln2+5\ln5\\6\ln3+5\ln5\end{pmatrix}
          \approx\begin{pmatrix}12.206073\\14.638863\end{pmatrix},
        \]
        and $Df(1,1)$ is the same pair written as a row.

  \item Only the term $i=k$ contains $x_{k}$, and its exponent $kx_{k}$ has
        derivative $k$:
        \[
          \pd{f}{x_{k}}=3^{\,kx_{k}}\ln3\cdot k=k\ln3\cdot3^{\,kx_{k}} .
        \]
        At the origin every $3^{0}=1$, so
        $\grad f=\ln3\,(1,2,\ldots,n)^{T}$ and
        $Df=\ln3\,(1\ \ 2\ \ \cdots\ \ n)$.

  \item Candidate: $f_{x}=3y$, $f_{y}=3x$, so $\grad f(a,b)=(3b,3a)$.
        Remainder:
        \[
          r(\mathbf{h})=3(a+h_{1})(b+h_{2})-3ab-3bh_{1}-3ah_{2}=3h_{1}h_{2} .
        \]
        Estimate:
        $\dfrac{|r(\mathbf{h})|}{\|\mathbf{h}\|}\le\dfrac{3\cdot\frac12
        \|\mathbf{h}\|^{2}}{\|\mathbf{h}\|}=\tfrac32\|\mathbf{h}\|\to0$.
        So $f$ is differentiable at every point, with derivative
        $(3b,\,3a)$.
  \item At the origin $\grad f=(0,0)$ and
        $r(\mathbf{h})=h_{1}^{2}+h_{2}^{2}=\|\mathbf{h}\|^{2}$, so
        $|r|/\|\mathbf{h}\|=\|\mathbf{h}\|\to0$: differentiable, with tangent
        plane $z=0$. At $(a,b)$: $\grad f=(2a,2b)$ and
        \[
          r(\mathbf{h})=(a+h_{1})^{2}+(b+h_{2})^{2}-a^{2}-b^{2}-2ah_{1}-2bh_{2}
          =h_{1}^{2}+h_{2}^{2},
        \]
        the same remainder, so the same estimate works everywhere.
  \item Candidate: $f_{x}=y^{2}$, $f_{y}=2xy$, so $\grad f(a,b)=(b^{2},2ab)$.
        Expand $(a+h_{1})(b^{2}+2bh_{2}+h_{2}^{2})
        =ab^{2}+2abh_{2}+ah_{2}^{2}+b^{2}h_{1}+2bh_{1}h_{2}+h_{1}h_{2}^{2}$
        and subtract $ab^{2}+b^{2}h_{1}+2abh_{2}$:
        \[
          r(\mathbf{h})=a\,h_{2}^{2}+2b\,h_{1}h_{2}+h_{1}h_{2}^{2} .
        \]
        Estimate:
        $\dfrac{|r|}{\|\mathbf{h}\|}\le|a|\,\|\mathbf{h}\|
        +|b|\,\|\mathbf{h}\|+\|\mathbf{h}\|^{2}\to0$. Differentiable
        everywhere, derivative $(b^{2},\,2ab)$.
  \item Candidate: $f_{x}=3y\mathrm{e}^{3x}$, $f_{y}=\mathrm{e}^{3x}$, so
        $\grad f(a,b)=\mathrm{e}^{3a}(3b,\,1)$. Write
        $\mathrm{e}^{3(a+h_{1})}=\mathrm{e}^{3a}\bigl(1+3h_{1}+R(3h_{1})\bigr)$:
        \[
          f(a+h_{1},b+h_{2})=\mathrm{e}^{3a}(b+h_{2})\bigl(1+3h_{1}+R\bigr)
          =\mathrm{e}^{3a}\bigl(b+3bh_{1}+h_{2}+3h_{1}h_{2}+(b+h_{2})R\bigr).
        \]
        Subtracting $b\mathrm{e}^{3a}$ and $\mathrm{e}^{3a}(3bh_{1}+h_{2})$
        leaves
        \[
          r(\mathbf{h})=\mathrm{e}^{3a}\bigl(3h_{1}h_{2}+(b+h_{2})\,R(3h_{1})\bigr).
        \]
        For $\|\mathbf{h}\|\le\tfrac13$ we have $|3h_{1}|\le1$, so
        $|R(3h_{1})|\le2\cdot9h_{1}^{2}\le18\|\mathbf{h}\|^{2}$, and
        \[
          \frac{|r|}{\|\mathbf{h}\|}\le\mathrm{e}^{3a}\Bigl(\tfrac32
          \|\mathbf{h}\|+18\bigl(|b|+\|\mathbf{h}\|\bigr)\|\mathbf{h}\|\Bigr)
          \to0 .
        \]
        Differentiable at every point.
  \item (a) On the $x$-axis $f(h,0)=(h^{2})^{1/3}=|h|^{2/3}$, so the
        difference quotient $|h|^{2/3}/h=\pm|h|^{-1/3}$ is unbounded:
        $f_{x}(0,0)$ does not exist (likewise $f_{y}$, with quotient
        $\pm4^{1/3}|h|^{-1/3}$). Not differentiable.
        (b) $g(\mathbf{h})=\|\mathbf{h}\|^{3/2}$. On an axis the quotient is
        $|h|^{3/2}/h=\pm|h|^{1/2}\to0$, so both partials are $0$. Then
        $r(\mathbf{h})=\|\mathbf{h}\|^{3/2}$ and
        $|r|/\|\mathbf{h}\|=\|\mathbf{h}\|^{1/2}\to0$: differentiable, with
        tangent plane $z=0$. The exponents decide: $2p=\tfrac23<1$ fails,
        $2p=\tfrac32>1$ works.
  \item \emph{Origin.} $f=0$ on both axes, so both partials are $0$. The
        remainder is $r(\mathbf{h})=|h_{1}h_{2}|\le\tfrac12\|\mathbf{h}\|^{2}$,
        so $|r|/\|\mathbf{h}\|\le\tfrac12\|\mathbf{h}\|\to0$: differentiable,
        with tangent plane $z=0$. \emph{At $(2,0)$.}
        $f(2,h)-f(2,0)=|2h|=2|h|$, and the quotient $2|h|/h=\pm2$ has no
        limit: $f_{y}(2,0)$ does not exist, so $f$ is not differentiable
        there. The graph is creased along both axes, except at the origin,
        where the two creases flatten out.
  \item \emph{Partials.} $f(h,0)=0$ and $f(0,h)=0$ for $h\neq0$, so both
        difference quotients are $0/h=0$: $f_{x}(0,0)=f_{y}(0,0)=0$.
        \emph{Lines.} On the $y$-axis $f=0$. On $y=mx$, $x\neq0$:
        \[
          f(x,mx)=\frac{mx^{3}}{x^{4}+m^{2}x^{2}}=\frac{mx}{x^{2}+m^{2}}
          \;\longrightarrow\;0\quad(x\to0),
        \]
        for $m\neq0$, and $f=0$ for $m=0$. Every straight line gives $0$.
        \emph{Parabola.} On $y=x^{2}$: $f=\dfrac{x^{4}}{x^{4}+x^{4}}=\tfrac12$
        for every $x\neq0$. So $f$ has no limit at the origin and is
        \textbf{not continuous} there, hence \textbf{not differentiable},
        although both partials exist and every line agrees.
  \item \emph{Partials.} $f=0$ on both axes, so $f_{x}(0,0)=f_{y}(0,0)=0$.
        \emph{Lines.} On $y=mx$:
        $f=\dfrac{2m^{2}x^{3}}{x^{2}+m^{4}x^{4}}=\dfrac{2m^{2}x}{1+m^{4}x^{2}}
        \to0$; on $x=0$, $f=0$. \emph{Parabola.} On $x=y^{2}$:
        $f=\dfrac{2y^{4}}{y^{4}+y^{4}}=1$ for $y\neq0$. Not continuous at the
        origin, hence not differentiable.
  \item (a) A smooth bowl, with no edges, creases or tips anywhere: it has a
        tangent plane at every point (horizontal at the bottom), so the
        function is differentiable everywhere. (b) A cone with a sharp tip:
        differentiable everywhere except at the tip, where every vertical
        cut through the tip is a V and no plane fits. (c) A pyramid whose
        four flat faces meet along two creases running over the axes:
        differentiable on the faces, not on the two creases (including
        the top point).
  \item The bowl $(x-1)^{2}+4y^{2}$ has no tip, so it is out. The level
        $c>0$ of $\sqrt{(x-1)^{2}+4y^{2}}$ is
        $\dfrac{(x-1)^{2}}{c^{2}}+\dfrac{y^{2}}{(c/2)^{2}}=1$, semi-axes $c$
        in $x$ and $c/2$ in $y$: twice as long in $x$, as in the picture.
        The third formula has semi-axes $c/2$ in $x$ and $c$ in $y$, the
        wrong way round. So (b) is $\sqrt{(x-1)^{2}+4y^{2}}$. At $(1,0)$:
        $f(1+h,0)-f(1,0)=\sqrt{h^{2}}=|h|$, and $|h|/h=\pm1$ has no limit;
        likewise $f(1,h)-f(1,0)=2|h|$. Neither partial exists, so $f$ is
        not differentiable at the tip, as the picture said.
\end{solutions}

\chapter{Higher-Order Partials, Young's Theorem, the Hessian, and Definiteness}

First-order partials describe slopes; second-order partials describe
\emph{curvature}. In one variable we used $f''$ to distinguish maxima from
minima. In several variables we need a whole \emph{matrix} of second
derivatives (the Hessian) and curvature is captured by its
definiteness, which we test with the leading principal minors (Sylvester's
criterion).

\section{Core Rules}

\begin{rulebox}{Direct second-order partials}
For $f(x,y)$,
\[
  f_{xx} \;=\; \pd{}{x}\!\left(\pd{f}{x}\right) \;=\; \frac{\partial^{2} f}{\partial x^{2}},
  \qquad
  f_{yy} \;=\; \frac{\partial^{2} f}{\partial y^{2}} .
\]
\end{rulebox}

\begin{rulebox}{Cross (mixed) partials}
\[
  f_{xy} \;=\; \pd{}{y}\!\left(\pd{f}{x}\right) \;=\; \frac{\partial^{2} f}{\partial y\,\partial x},
  \qquad
  f_{yx} \;=\; \pd{}{x}\!\left(\pd{f}{y}\right) .
\]
\end{rulebox}

\begin{rulebox}{Young's Theorem (Schwarz / Clairaut)}
If $f$ is $C^{2}$ on an open set (i.e.\ all second partials exist and
are continuous) then the order of differentiation is irrelevant:
\[
  f_{xy} \;=\; f_{yx} .
\]
This is the symmetry of the Hessian. All standard polynomial, exponential,
logarithmic, Cobb--Douglas, and CES forms are $C^{2}$ on their natural
domains.
\end{rulebox}

\begin{rulebox}{Hessian matrix}
For $f:\R^{n}\to\R$,
\[
  \Hess_f(\mathbf{x}) \;=\; \bigl[f_{x_ix_j}(\mathbf{x})\bigr]_{i,j=1}^{n}
  \;=\;
  \begin{pmatrix}
    f_{x_1x_1} & f_{x_1x_2} & \cdots & f_{x_1x_n} \\
    f_{x_2x_1} & f_{x_2x_2} & \cdots & f_{x_2x_n} \\
    \vdots     & \vdots     & \ddots & \vdots     \\
    f_{x_nx_1} & f_{x_nx_2} & \cdots & f_{x_nx_n}
  \end{pmatrix} .
\]
For $C^{2}$ functions $\Hess_f$ is symmetric.
\end{rulebox}

\begin{rulebox}{Definiteness}
A symmetric matrix $H$ is
\begin{itemize}[leftmargin=1.6em,itemsep=0.3ex]
  \item \emph{positive definite (PD)} if $\mathbf{v}^{T}H\mathbf{v}>0$ for
        every $\mathbf{v}\neq\mathbf{0}$ (curvature ``up'' in every
        direction $\Rightarrow$ local minimum;
  \item \emph{negative definite (ND)} if $\mathbf{v}^{T}H\mathbf{v}<0$ for
        every $\mathbf{v}\neq\mathbf{0}$) curvature ``down'' in every
        direction $\Rightarrow$ local maximum;
  \item \emph{indefinite} if $\mathbf{v}^{T}H\mathbf{v}$ takes both
        positive and negative values $\Rightarrow$ saddle point.
\end{itemize}
\end{rulebox}

\begin{rulebox}{Sylvester's criterion ($2\times2$ case)}
Let $H=\begin{pmatrix}a&b\\b&c\end{pmatrix}$ with leading principal minors
$D_1=a$ and $D_2=\det H=ac-b^{2}$. Then
\[
  \begin{aligned}
    D_1>0 \text{ and } D_2>0 &\;\Rightarrow\; H \text{ is PD}, \\
    D_1<0 \text{ and } D_2>0 &\;\Rightarrow\; H \text{ is ND}, \\
    D_2<0 &\;\Rightarrow\; H \text{ is indefinite (saddle)}, \\
    D_2=0 &\;\Rightarrow\; \text{test is inconclusive (semi-definite case)}.
  \end{aligned}
\]
Note: ND requires $D_1<0$ \emph{and} $D_2>0$; the alternating-sign
condition is what one generalises to higher $n$ (signs $-,+,-,+,\dots$ for
ND).
\end{rulebox}

\begin{rulebox}{Classification of critical points}
A point $\mathbf{x}_0$ is a \emph{critical point} of $f$ if
$\grad f(\mathbf{x}_0)=\mathbf{0}$. Then evaluate $\Hess_f(\mathbf{x}_0)$:
\begin{itemize}[leftmargin=1.6em,itemsep=0.3ex]
  \item $H$ PD at $\mathbf{x}_0$ $\Rightarrow$ strict local minimum;
  \item $H$ ND at $\mathbf{x}_0$ $\Rightarrow$ strict local maximum;
  \item $H$ indefinite at $\mathbf{x}_0$ $\Rightarrow$ saddle point;
  \item $D_2=0$ $\Rightarrow$ inconclusive, higher-order analysis
        required.
\end{itemize}
\end{rulebox}

\section{Worked Example}

\begin{worked}{Find and classify every critical point of $f(x,y)=x^{3}-3xy+y^{3}$}
\wstep{Step 1. First-order conditions: both partials zero, simultaneously.}
\[
  f_{x}=3x^{2}-3y=0
  \quad\Longrightarrow\quad y=x^{2},
\]
\[
  f_{y}=-3x+3y^{2}=0
  \quad\Longrightarrow\quad x=y^{2} .
\]
\wtrap{A critical point needs \emph{every} partial to vanish at the
\emph{same} point. Solving each equation separately and pairing the answers
arbitrarily is not the same thing --- the two equations form a system.}

\wstep{Step 2. Solve the system by substitution.}
Substituting $y=x^{2}$ into $x=y^{2}$:
\[
  x=x^{4}
  \;\Longrightarrow\;
  x^{4}-x=0
  \;\Longrightarrow\;
  x\bigl(x^{3}-1\bigr)=0
  \;\Longrightarrow\;
  x=0\ \text{or}\ x=1 .
\]
The corresponding $y=x^{2}$ gives the two critical points
\[
  (0,0)\qquad\text{and}\qquad(1,1).
\]
\wtrap{Do not divide by $x$ to get $x^{3}=1$. That silently discards the
root $x=0$ --- and $(0,0)$ turns out to be the more interesting of the two
points. Factor; never divide by something that may vanish.}

\wstep{Step 3. Build the Hessian once, symbolically.}
\[
  f_{xx}=6x,\qquad f_{yy}=6y,\qquad f_{xy}=f_{yx}=-3 ,
\]
\[
  \Hess f(x,y)=\begin{pmatrix}6x&-3\\-3&6y\end{pmatrix}.
\]
The two mixed partials agree, as Young's theorem guarantees for a function
with continuous second derivatives --- differentiating in $x$ then $y$ gives
the same result as $y$ then $x$. Computing both is a free check, and here
$\frac{\partial}{\partial y}(3x^{2}-3y)=-3$ and
$\frac{\partial}{\partial x}(-3x+3y^{2})=-3$. \checkmark

\wstep{Step 4. Classify $(0,0)$.}
\[
  \Hess f(0,0)=\begin{pmatrix}0&-3\\-3&0\end{pmatrix},
  \qquad
  \det=0-9=-9<0 .
\]
A negative determinant means the two eigenvalues have opposite signs, so the
form is indefinite:
\[
  \boxed{(0,0)\ \text{is a saddle point.}}
\]
No further test is needed, and the value of $f_{xx}$ is irrelevant here.
\wtrap{$\det\Hess<0$ settles the matter immediately --- saddle, whatever
$f_{xx}$ does. It is only when $\det\Hess>0$ that $f_{xx}$ has to be
consulted. Checking $f_{xx}$ first wastes effort and invites the wrong
conclusion.}

\wstep{Step 5. Classify $(1,1)$.}
\[
  \Hess f(1,1)=\begin{pmatrix}6&-3\\-3&6\end{pmatrix},
  \qquad
  \det=36-9=27>0,
  \qquad
  f_{xx}=6>0 .
\]
Positive determinant with a positive $f_{xx}$: positive definite, so
\[
  \boxed{(1,1)\ \text{is a strict local minimum},\qquad f(1,1)=1-3+1=-1 .}
\]

\wstep{Step 6. Verify the saddle directly, because it is cheap.}
Along the line $y=x$: $f(x,x)=2x^{3}-3x^{2}$, whose value near $0$ is
$\approx-3x^{2}<0$. Along $y=-x$: $f(x,-x)=x^{3}+3x^{2}-x^{3}=3x^{2}>0$. So
$f$ takes both signs arbitrarily close to $(0,0)$, where $f=0$ --- a saddle,
confirmed without any matrix.

\wstep{Step 7. The decision rule, complete.}
\begin{center}\footnotesize
\begin{tabular}{lll}
  \hline
  $\det\Hess$ & $f_{xx}$ & conclusion\\
  \hline
  $>0$ & $>0$ & strict local minimum\\
  $>0$ & $<0$ & strict local maximum\\
  $<0$ & --- & saddle point\\
  $=0$ & --- & \textbf{test fails}; investigate directly\\
  \hline
\end{tabular}
\end{center}
\wtrap{$\det\Hess=0$ means the test is silent, not that the point is a
saddle. $f=x^{4}+y^{4}$ (minimum), $f=-x^{4}-y^{4}$ (maximum) and
$f=x^{4}-y^{4}$ (saddle) all have the zero Hessian at the origin. When the
determinant vanishes you must argue directly, as in Step~6.}

\wstep{Step 8. Local is not global.}
Neither point is a global extremum: $f(x,0)=x^{3}$ runs to $\pm\infty$. On
an unbounded domain, second-order conditions can never certify a global
optimum by themselves --- that needs convexity of $f$ on the whole domain,
or a compact feasible set and Weierstrass. Volume~7 makes this the
difference between finding a candidate and having an answer.
\end{worked}

\section{Practice Problems}

\begin{problems}
  \item Compute all second partials of $f(x,y) = 4x^{3} - 5x^{2}y + 7y^{4}$
        and verify Young's theorem.
  \item Construct the Hessian of
        $f(x,y) = 5x^{2} + 6xy + 2y^{2} - 3x + 4y$.
  \item Classify the definiteness of
        $\Hess = \begin{pmatrix}6&2\\2&3\end{pmatrix}$ using Sylvester's
        criterion.
  \item Find and classify the critical points of
        $f(x,y) = x^{2} + 4y^{2} - 6x + 8y + 11$.
  \item Compute all second partials of $f(x,y) = (x+2y)^{3}$ and verify
        $f_{xy} = f_{yx}$.
  \item Construct the Hessian of $f(x,y) = x^{3} + 2xy^{2} - 4y$.
  \item Classify $\Hess = \begin{pmatrix}-4&2\\2&-5\end{pmatrix}$.
  \item Find and classify the critical points of
        $f(x,y) = -x^{2} - 3y^{2} + 8x - 6y + 5$.
  \item Compute all second partials of $f(x,y) = x^{2}\ln(y)$ ($y>0$) and
        verify Young's theorem.
  \item Construct the Hessian of $f(x,y) = \mathrm{e}^{2x-y}$.
  \item Classify $\Hess = \begin{pmatrix}2&4\\4&2\end{pmatrix}$.
  \item Find and classify the critical points of
        $f(x,y) = x^{2} - y^{2} + 6x - 4y$.
  \item Compute all second partials of $f(x,y) = \mathrm{e}^{x^{2}+y}$ and
        verify Young's theorem.
  \item Construct the Hessian of the Cobb--Douglas function
        $f(K,L) = K^{1/3}L^{2/3}$, $K,L>0$.
  \item Classify $\Hess = \begin{pmatrix}7&-3\\-3&1\end{pmatrix}$.
  \item Find and classify the critical points of
        $f(x,y) = x^{3} + y^{3} - 12x - 27y$.
  \item Compute all second partials of $f(x,y) = 1/(x+y)$ and verify
        Young's theorem.
  \item Construct the Hessian of
        $f(x,y,z) = x^{2} + 2y^{2} + 3z^{2} - xy + 2yz$.
  \item Classify $\Hess = \begin{pmatrix}-3&1\\1&-3\end{pmatrix}$.
  \item Find and classify the critical points of
        $f(x,y) = x^{2} + 4xy + y^{2} + 2y$.
  \item Compute all second partials of $f(x,y) = x^{1/2}y^{1/2}$
        ($x,y>0$) and verify $f_{xy} = f_{yx}$.
  \item Construct the Hessian of $f(x,y) = \ln(x^{2}+y^{2})$ at a generic
        point $(x,y)\neq(0,0)$.
  \item Find and classify the critical points of
        $f(x,y) = x^{4} + y^{4} - 4xy$.
  \item Compute all second partials of $f(x,y) = x\mathrm{e}^{y} + y\mathrm{e}^{x}$
        and verify Young's theorem.
  \item Construct the Hessian of $f(x,y) = x^{3} - 3xy^{2}$ (a harmonic
        function).
  \item Classify $\Hess = \begin{pmatrix}5&0\\0&-2\end{pmatrix}$ and state
        what this implies if $H$ is the Hessian at a critical point.
  \item Find and classify the critical points of
        $f(x,y) = xy - x^{2} - y^{2} + 6x + 3y$.
  \item Compute all second partials of $f(x,y) = \ln(xy)$ ($x,y>0$).
  \item Construct the Hessian of $f(x,y) = (x^{2}+y^{2})^{2}$ at a generic
        point.
  \item Classify $\Hess = \begin{pmatrix}-1&2\\2&-1\end{pmatrix}$.
  \item Find and classify the critical points of
        $f(x,y) = x^{3} - 3xy + y^{3}$.
  \item Compute all second partials of $f(x,y) = \sin(x)\cos(y)$ and
        verify Young's theorem.
  \item Construct the Hessian of $f(x,y,z) = \mathrm{e}^{xyz}$ at the point
        $(1,1,1)$ (evaluate numerically).
  \item Classify $\Hess = \begin{pmatrix}4&-2\\-2&4\end{pmatrix}$.
  \item Find and classify the critical points of the profit function
        $\pi(L,K) = 10L^{1/2}K^{1/2} - 2L - 3K$ (assume $L,K>0$).
  \item Compute all second partials of $f(x,y) = x^{2}y^{3}$ and verify
        Young's theorem.
  \item Construct the Hessian of
        $f(x,y) = \mathrm{e}^{x}\sin(y) - \mathrm{e}^{x}\cos(y)$.
  \item Classify $\Hess = \begin{pmatrix}9&6\\6&4\end{pmatrix}$. (Note the
        boundary case.)
  \item Find and classify the critical points of
        $f(x,y) = x^{2} + xy + y^{2} - 3x$.
  \item Compute all second partials of $f(x,y) = \arctan(y/x)$ ($x>0$) and
        verify Young's theorem.
  \item Construct the Hessian of the CES function
        $f(K,L) = (K^{1/2}+L^{1/2})^{2}$ at a generic point.
  \item Find and classify the critical points of
        $f(x,y) = -x^{4} - y^{4} + 8xy$.
  \item Compute all second partials of $f(x,y) = x^{2}\mathrm{e}^{-y}$ and
        verify Young's theorem.
  \item Construct the Hessian of $f(x,y) = \cos(xy)$ at the origin
        $(0,0)$.
  \item Classify $\Hess = \begin{pmatrix}0&3\\3&0\end{pmatrix}$.
  \item Find and classify the critical points of
        $f(x,y) = x^{3} + 3xy^{2} - 15x - 12y$.
  \item Construct the Hessian of $f(x,y) = \ln(1+x^{2}+y^{2})$ at the
        origin.
  \item Find and classify the critical points of
        $f(x,y) = \mathrm{e}^{-(x^{2}+y^{2})}$.
\end{problems}

\section{Solutions}

\begin{solutions}
  \item $f_x = 12x^{2} - 10xy$, $f_y = -5x^{2} + 28y^{3}$. Then
        $f_{xx} = 24x - 10y$, $f_{yy} = 84y^{2}$, and
        \[
          f_{xy} = \pd{}{y}\bigl(12x^{2}-10xy\bigr) = -10x, \qquad
          f_{yx} = \pd{}{x}\bigl(-5x^{2}+28y^{3}\bigr) = -10x . \; \checkmark
        \]

  \item $f_x = 10x + 6y - 3$, $f_y = 6x + 4y + 4$. Hessian
        \[
          \Hess = \begin{pmatrix}10&6\\6&4\end{pmatrix} .
        \]

  \item $D_1 = 6 > 0$, $D_2 = 18 - 4 = 14 > 0$. \textbf{PD.}

  \item $\grad f = (2x-6,\,8y+8) = (0,0) \Rightarrow x=3,\,y=-1$. Hessian
        $\Hess = \begin{pmatrix}2&0\\0&8\end{pmatrix}$: $D_1=2>0$,
        $D_2=16>0$, \textbf{PD}, so $(3,-1)$ is a strict \textbf{local
        minimum} with $f(3,-1) = 9+4-18-8+11 = -2$.

  \item Chain rule: $f_x = 3(x+2y)^{2}$, $f_y = 6(x+2y)^{2}$. So
        $f_{xx} = 6(x+2y)$, $f_{yy} = 24(x+2y)$, and
        \[
          f_{xy} = 12(x+2y) = f_{yx} . \; \checkmark
        \]

  \item $f_x = 3x^{2} + 2y^{2}$, $f_y = 4xy - 4$. Hessian
        \[
          \Hess = \begin{pmatrix}6x&4y\\4y&4x\end{pmatrix} .
        \]

  \item $D_1 = -4 < 0$, $D_2 = 20 - 4 = 16 > 0$. \textbf{ND.}

  \item $\grad f = (-2x+8,\,-6y-6) = (0,0) \Rightarrow x=4,\,y=-1$. Hessian
        $\Hess = \begin{pmatrix}-2&0\\0&-6\end{pmatrix}$: $D_1=-2<0$,
        $D_2=12>0$, \textbf{ND}, so $(4,-1)$ is a strict \textbf{local
        maximum} with $f(4,-1) = -16-3+32+6+5 = 24$.

  \item $f_x = 2x\ln(y)$, $f_y = x^{2}/y$. So $f_{xx} = 2\ln(y)$,
        $f_{yy} = -x^{2}/y^{2}$, and
        \[
          f_{xy} = 2x/y = f_{yx} . \; \checkmark
        \]

  \item $f_x = 2\mathrm{e}^{2x-y}$, $f_y = -\mathrm{e}^{2x-y}$. So
        $f_{xx} = 4\mathrm{e}^{2x-y}$, $f_{yy} = \mathrm{e}^{2x-y}$,
        $f_{xy} = -2\mathrm{e}^{2x-y}$. Hessian
        \[
          \Hess = \mathrm{e}^{2x-y}\begin{pmatrix}4&-2\\-2&1\end{pmatrix} .
        \]

  \item $D_1 = 2 > 0$, $D_2 = 4 - 16 = -12 < 0$. \textbf{Indefinite
        (saddle).}

  \item $\grad f = (2x+6,\,-2y-4) = (0,0) \Rightarrow x=-3,\,y=-2$. Hessian
        $\Hess = \begin{pmatrix}2&0\\0&-2\end{pmatrix}$: $D_2 = -4 < 0$,
        \textbf{indefinite}, so $(-3,-2)$ is a \textbf{saddle point}.

  \item Let $E = \mathrm{e}^{x^{2}+y}$. Then $f_x = 2xE$, $f_y = E$. So
        $f_{xx} = 2E + 4x^{2}E = (2+4x^{2})E$, $f_{yy} = E$, and
        \[
          f_{xy} = \pd{}{y}(2xE) = 2xE, \qquad
          f_{yx} = \pd{}{x}(E) = 2xE . \; \checkmark
        \]

  \item $f_K = \tfrac13K^{-2/3}L^{2/3}$, $f_L = \tfrac23K^{1/3}L^{-1/3}$.
        So
        \[
          f_{KK} = -\tfrac29K^{-5/3}L^{2/3}, \quad
          f_{LL} = -\tfrac29K^{1/3}L^{-4/3}, \quad
          f_{KL} = \tfrac29K^{-2/3}L^{-1/3} = f_{LK} . \; \checkmark
        \]
        Hessian
        \[
          \Hess = \tfrac29
          \begin{pmatrix}
            -K^{-5/3}L^{2/3} & K^{-2/3}L^{-1/3} \\
            K^{-2/3}L^{-1/3} & -K^{1/3}L^{-4/3}
          \end{pmatrix} .
        \]

  \item $D_1 = 7 > 0$, $D_2 = 7 - 9 = -2 < 0$. \textbf{Indefinite.}

  \item $\grad f = (3x^{2}-12,\,3y^{2}-27) = (0,0) \Rightarrow x^{2}=4,\,
        y^{2}=9$, giving $(\pm2,\pm3)$. Hessian
        $\Hess = \begin{pmatrix}6x&0\\0&6y\end{pmatrix}$, so $D_1=6x$,
        $D_2=36xy$.
        \begin{itemize}
          \item $(2,3)$: $D_1=12>0$, $D_2=216>0$, \textbf{PD}: local
                minimum, $f = 8+27-24-81 = -70$.
          \item $(2,-3)$: $D_2=-216<0$, \textbf{saddle}.
          \item $(-2,3)$: $D_2=-216<0$, \textbf{saddle}.
          \item $(-2,-3)$: $D_1=-12<0$, $D_2=216>0$, \textbf{ND}: local
                maximum, $f = -8-27+24+81 = 70$.
        \end{itemize}

  \item Let $S=x+y$, $f=S^{-1}$, $f_x = -S^{-2} = f_y$. So
        $f_{xx} = 2S^{-3} = f_{yy}$, and
        \[
          f_{xy} = 2S^{-3} = f_{yx} = \frac{2}{(x+y)^{3}} . \; \checkmark
        \]

  \item $f_x = 2x-y$, $f_y = 4y-x+2z$, $f_z = 6z+2y$. Hessian
        \[
          \Hess = \begin{pmatrix}2&-1&0\\-1&4&2\\0&2&6\end{pmatrix} .
        \]

  \item $D_1 = -3 < 0$, $D_2 = 9 - 1 = 8 > 0$. \textbf{ND.}

  \item $\grad f = (2x+4y,\,4x+2y+2) = (0,0)$. From the first, $x=-2y$;
        substituting, $-8y+2y+2=0 \Rightarrow y=1/3$, $x=-2/3$. Hessian
        $\Hess = \begin{pmatrix}2&4\\4&2\end{pmatrix}$: $D_1=2>0$,
        $D_2=4-16=-12<0$, \textbf{indefinite: saddle point} at
        $(-2/3,1/3)$.

  \item $f_x = \tfrac12x^{-1/2}y^{1/2}$, $f_y = \tfrac12x^{1/2}y^{-1/2}$.
        So
        \[
          f_{xx} = -\tfrac14x^{-3/2}y^{1/2}, \quad
          f_{yy} = -\tfrac14x^{1/2}y^{-3/2}, \quad
          f_{xy} = \tfrac14x^{-1/2}y^{-1/2} = f_{yx} . \; \checkmark
        \]

  \item Let $S=x^{2}+y^{2}$, $f=\ln S$, $f_x = 2x/S$, $f_y = 2y/S$. By the
        quotient rule:
        \[
          f_{xx} = \frac{2S-2x\cdot2x}{S^{2}} = \frac{2(y^{2}-x^{2})}{(x^{2}+y^{2})^{2}}, \quad
          f_{yy} = \frac{2(x^{2}-y^{2})}{(x^{2}+y^{2})^{2}}, \quad
          f_{xy} = \frac{-4xy}{(x^{2}+y^{2})^{2}} .
        \]
        Hessian
        \[
          \Hess = \frac{2}{(x^{2}+y^{2})^{2}}
          \begin{pmatrix}y^{2}-x^{2} & -2xy \\ -2xy & x^{2}-y^{2}\end{pmatrix} .
        \]
        Note $f_{xx}+f_{yy}=0$: $\ln(x^{2}+y^{2})$ is harmonic.

  \item $\grad f = (4x^{3}-4y,\,4y^{3}-4x) = (0,0) \Rightarrow y=x^{3}$
        and $x=y^{3}=x^{9}$. So $x(x^{8}-1)=0$, giving $x\in\{-1,0,1\}$.
        Critical points: $(0,0),(1,1),(-1,-1)$. Hessian
        $\Hess = \begin{pmatrix}12x^{2}&-4\\-4&12y^{2}\end{pmatrix}$.
        \begin{itemize}
          \item $(0,0)$: $\Hess = \begin{pmatrix}0&-4\\-4&0\end{pmatrix}$,
                $D_2=-16<0$, \textbf{saddle}.
          \item $(1,1)$: $\Hess = \begin{pmatrix}12&-4\\-4&12\end{pmatrix}$,
                $D_1=12>0$, $D_2=144-16=128>0$, \textbf{PD}: local
                minimum, $f=1+1-4=-2$.
          \item $(-1,-1)$: same $H$, \textbf{PD}: local minimum, $f=-2$.
        \end{itemize}

  \item $f_x = \mathrm{e}^{y}+y\mathrm{e}^{x}$, $f_y = x\mathrm{e}^{y}+\mathrm{e}^{x}$. So
        $f_{xx} = y\mathrm{e}^{x}$, $f_{yy} = x\mathrm{e}^{y}$, and
        \[
          f_{xy} = \mathrm{e}^{y}+\mathrm{e}^{x} = f_{yx} . \; \checkmark
        \]

  \item $f_x = 3x^{2}-3y^{2}$, $f_y = -6xy$. So $f_{xx}=6x$,
        $f_{yy}=-6x$, $f_{xy}=-6y=f_{yx}$. Hessian
        \[
          \Hess = \begin{pmatrix}6x&-6y\\-6y&-6x\end{pmatrix} .
        \]
        Note $f_{xx}+f_{yy}=0$: $f$ is harmonic.

  \item $D_1 = 5 > 0$, $D_2 = -10 < 0$. \textbf{Indefinite.} At a critical
        point this would mean a \textbf{saddle}.

  \item $\grad f = (y-2x+6,\,x-2y+3) = (0,0)$. So $2x-y=6$ and
        $x-2y=-3$. From the first, $y=2x-6$; substituting:
        $x-2(2x-6)=-3 \Rightarrow -3x+12=-3 \Rightarrow x=5$, $y=4$.
        Hessian $\Hess = \begin{pmatrix}-2&1\\1&-2\end{pmatrix}$:
        $D_1=-2<0$, $D_2=4-1=3>0$, \textbf{ND}, so $(5,4)$ is a
        \textbf{local maximum} with $f(5,4) = 20-25-16+30+12 = 21$.

  \item $\ln(xy) = \ln x + \ln y$. So $f_x = 1/x$, $f_y = 1/y$,
        $f_{xx} = -1/x^{2}$, $f_{yy} = -1/y^{2}$, $f_{xy} = 0 = f_{yx}$.
        \; $\checkmark$

  \item $f_x = 2(x^{2}+y^{2})\cdot2x = 4x(x^{2}+y^{2})$,
        $f_y = 4y(x^{2}+y^{2})$. Then
        \[
          f_{xx} = 4(x^{2}+y^{2})+4x\cdot2x = 4(3x^{2}+y^{2}), \quad
          f_{yy} = 4(x^{2}+3y^{2}), \quad
          f_{xy} = 4x\cdot2y = 8xy = f_{yx} .
        \]
        Hessian
        \[
          \Hess = 4\begin{pmatrix}3x^{2}+y^{2}&2xy\\2xy&x^{2}+3y^{2}\end{pmatrix} .
        \]

  \item $D_1 = -1 < 0$, $D_2 = 1 - 4 = -3 < 0$. Since $D_2<0$,
        \textbf{indefinite.}

  \item $\grad f = (3x^{2}-3y,\,-3x+3y^{2}) = (0,0) \Rightarrow y=x^{2}$
        and $x=y^{2}=x^{4}$, so $x(x^{3}-1)=0$, giving $x\in\{0,1\}$.
        Critical points: $(0,0)$ and $(1,1)$. Hessian
        $\Hess = \begin{pmatrix}6x&-3\\-3&6y\end{pmatrix}$.
        \begin{itemize}
          \item $(0,0)$: $\Hess = \begin{pmatrix}0&-3\\-3&0\end{pmatrix}$,
                $D_2=-9<0$, \textbf{saddle}.
          \item $(1,1)$: $\Hess = \begin{pmatrix}6&-3\\-3&6\end{pmatrix}$,
                $D_1=6>0$, $D_2=36-9=27>0$, \textbf{PD}: local minimum,
                $f=1-3+1=-1$.
        \end{itemize}

  \item $f_x = \cos(x)\cos(y)$, $f_y = -\sin(x)\sin(y)$. So
        $f_{xx} = -\sin(x)\cos(y)$, $f_{yy} = -\sin(x)\cos(y)$, and
        \[
          f_{xy} = -\cos(x)\sin(y) = f_{yx} . \; \checkmark
        \]

  \item Let $u=xyz$, $f=\mathrm{e}^{u}$. Then $f_x = yz\mathrm{e}^{u}$,
        $f_y = xz\mathrm{e}^{u}$, $f_z = xy\mathrm{e}^{u}$. Direct
        partials: $f_{xx} = (yz)^{2}\mathrm{e}^{u}$,
        $f_{yy} = (xz)^{2}\mathrm{e}^{u}$, $f_{zz} = (xy)^{2}\mathrm{e}^{u}$.
        Cross-partials (using the product rule on
        $f_x = yz\,\mathrm{e}^{xyz}$):
        \[
          f_{xy} = z\mathrm{e}^{xyz} + yz\cdot xz\mathrm{e}^{xyz}
                 = z\mathrm{e}^{xyz}(1+xyz) .
        \]
        Similarly $f_{xz} = y\mathrm{e}^{xyz}(1+xyz)$ and
        $f_{yz} = x\mathrm{e}^{xyz}(1+xyz)$. (Symmetry $f_{xy}=f_{yx}$
        etc.\ confirms Young's theorem.) At $(1,1,1)$: $u=1$,
        $\mathrm{e}^{u}=\mathrm{e}$, $1+xyz=2$, and $yz=xz=xy=1$, so
        \[
          \Hess(1,1,1) = \mathrm{e}\begin{pmatrix}1&2&2\\2&1&2\\2&2&1\end{pmatrix} .
        \]

  \item $D_1 = 4 > 0$, $D_2 = 16 - 4 = 12 > 0$. \textbf{PD.}

  \item $\pi_L = 5L^{-1/2}K^{1/2}-2$, $\pi_K = 5L^{1/2}K^{-1/2}-3$.
        Setting both to zero gives $5K^{1/2}/L^{1/2}=2$ and
        $5L^{1/2}/K^{1/2}=3$. Multiplying these two equations: $25=6$,
        which is a contradiction. So $\pi$ has \textbf{no interior
        critical point}. Why? The production function
        $Q=10L^{1/2}K^{1/2}$ is homogeneous of degree 1 (constant returns
        to scale), so along any ray $(\lambda L,\lambda K)$ revenue and
        cost both scale linearly: $\pi(\lambda L,\lambda K)=\lambda\pi(L,K)$.
        Hence profit is either identically zero, or unbounded above
        (depending on sign of $\pi$ along some direction). With
        $\max_{L/K}(10\sqrt{LK}-2L-3K)/K$: setting $t=L/K$ gives
        $\pi/K = 10\sqrt{t}-2t-3$, maximised at $t=25/4$ giving
        $\pi/K = 25-25/2-3 = 19/2 > 0$. So $\pi\to+\infty$ along the ray
        $L=(25/4)K$. The Hessian test does not apply because there is no
        critical point to classify. (Lesson: an inconsistent
        first-order system is itself a meaningful answer about the
        geometry of the problem.)

  \item $f_x = 2xy^{3}$, $f_y = 3x^{2}y^{2}$. So $f_{xx} = 2y^{3}$,
        $f_{yy} = 6x^{2}y$, and
        \[
          f_{xy} = 6xy^{2} = f_{yx} . \; \checkmark
        \]

  \item Let $g=\sin y-\cos y$, so $f=\mathrm{e}^{x}g$. Then
        $f_x = \mathrm{e}^{x}g$, $f_y = \mathrm{e}^{x}(\cos y+\sin y)$. So
        $f_{xx} = \mathrm{e}^{x}g = \mathrm{e}^{x}(\sin y-\cos y)$,
        $f_{yy} = \mathrm{e}^{x}(-\sin y+\cos y) = -\mathrm{e}^{x}g$, and
        \[
          f_{xy} = \mathrm{e}^{x}(\cos y+\sin y) = f_{yx} . \; \checkmark
        \]
        Hessian
        \[
          \Hess = \mathrm{e}^{x}
          \begin{pmatrix}
            \sin y-\cos y & \cos y+\sin y \\
            \cos y+\sin y & \cos y-\sin y
          \end{pmatrix} .
        \]

  \item $D_1 = 9 > 0$, $D_2 = 36 - 36 = 0$. \textbf{Inconclusive
        (positive semi-definite boundary case)}: Sylvester gives no
        verdict; we cannot rule out PSD with a zero eigenvalue. (In fact
        $H = (3,2)^{T}(3,2)$, so it is PSD with one zero eigenvalue.)

  \item $\grad f = (2x+y-3,\,x+2y) = (0,0)$. From the second: $x=-2y$;
        substituting: $-4y+y-3=0 \Rightarrow y=-1$, $x=2$. Critical point
        $(2,-1)$. Hessian $\Hess = \begin{pmatrix}2&1\\1&2\end{pmatrix}$:
        $D_1=2>0$, $D_2=4-1=3>0$, \textbf{PD}, so $(2,-1)$ is a strict
        \textbf{local minimum} with $f(2,-1) = 4-2+1-6 = -3$.

  \item With $u=y/x$, $f=\arctan u$: $f_x = -y/(x^{2}+y^{2})$,
        $f_y = x/(x^{2}+y^{2})$. Let $S=x^{2}+y^{2}$. Then
        \[
          f_{xx} = \frac{2xy}{S^{2}}, \qquad
          f_{yy} = \frac{-2xy}{S^{2}}, \qquad
          f_{xy} = \frac{-S+2y^{2}}{S^{2}} = \frac{y^{2}-x^{2}}{(x^{2}+y^{2})^{2}}
                  = f_{yx} . \; \checkmark
        \]
        (And $f_{xx}+f_{yy}=0$: $\arctan(y/x)$ is harmonic.)

  \item Let $g=K^{1/2}+L^{1/2}$. From sec06 sol.\ 23, $f_K = g/K^{1/2}$,
        $f_L = g/L^{1/2}$. Then
        \[
          f_{KK} = \frac{\tfrac12K^{-1/2}\cdot K^{1/2}-g\cdot\tfrac12K^{-1/2}}{K}
                  = \frac{\tfrac12-g/(2K^{1/2})}{K}
                  = \frac{1}{2K}-\frac{g}{2K^{3/2}} .
        \]
        More cleanly: $f_K = 1+(L/K)^{1/2}$, so
        $f_{KK} = -\tfrac12L^{1/2}K^{-3/2}$, and by symmetry
        $f_{LL} = -\tfrac12K^{1/2}L^{-3/2}$. Cross-partial:
        $f_{KL} = \tfrac12K^{-1/2}L^{-1/2} = f_{LK}$. \; $\checkmark$
        Hessian
        \[
          \Hess = \tfrac12
          \begin{pmatrix}
            -L^{1/2}K^{-3/2} & K^{-1/2}L^{-1/2} \\
            K^{-1/2}L^{-1/2} & -K^{1/2}L^{-3/2}
          \end{pmatrix} .
        \]

  \item $\grad f = (-4x^{3}+8y,\,-4y^{3}+8x) = (0,0) \Rightarrow x^{3}=2y$
        and $y^{3}=2x$. Substituting: $y=x^{3}/2$, then
        $(x^{3}/2)^{3}=2x \Rightarrow x^{9}/8=2x \Rightarrow x(x^{8}-16)=0$,
        so $x\in\{0,\pm16^{1/8}\}=\{0,\pm\sqrt2\}$. Points: $(0,0)$,
        $(\sqrt2,\sqrt2)$ (since $y=x^{3}/2=2\sqrt2/2=\sqrt2$), and
        $(-\sqrt2,-\sqrt2)$. Hessian
        $\Hess = \begin{pmatrix}-12x^{2}&8\\8&-12y^{2}\end{pmatrix}$.
        \begin{itemize}
          \item $(0,0)$: $\Hess = \begin{pmatrix}0&8\\8&0\end{pmatrix}$,
                $D_2=-64<0$, \textbf{saddle}.
          \item $(\sqrt2,\sqrt2)$: $x^{2}=y^{2}=2$,
                $\Hess = \begin{pmatrix}-24&8\\8&-24\end{pmatrix}$,
                $D_1=-24<0$, $D_2=576-64=512>0$, \textbf{ND}: local
                maximum, $f=-4-4+16=8$.
          \item $(-\sqrt2,-\sqrt2)$: same $H$, $f=-4-4+16=8$. \textbf{Local
                maximum.}
        \end{itemize}

  \item $f_x = 2x\mathrm{e}^{-y}$, $f_y = -x^{2}\mathrm{e}^{-y}$. So
        $f_{xx} = 2\mathrm{e}^{-y}$, $f_{yy} = x^{2}\mathrm{e}^{-y}$, and
        \[
          f_{xy} = -2x\mathrm{e}^{-y} = f_{yx} . \; \checkmark
        \]

  \item $f_x = -y\sin(xy)$, $f_y = -x\sin(xy)$. Then
        $f_{xx} = -y^{2}\cos(xy)$, $f_{yy} = -x^{2}\cos(xy)$, and
        (product rule)
        \[
          f_{xy} = -\sin(xy) - y\cdot x\cos(xy) = -\sin(xy) - xy\cos(xy) .
        \]
        At $(0,0)$: $\sin 0=0$, $\cos 0=1$. So $f_{xx}(0,0)=0$,
        $f_{yy}(0,0)=0$, $f_{xy}(0,0)=0$. Hessian
        \[
          \Hess(0,0) = \begin{pmatrix}0&0\\0&0\end{pmatrix} .
        \]

  \item $D_1 = 0$. Sylvester's stated criterion is inconclusive because
        $D_1$ is not strictly positive or negative. Directly:
        $D_2 = 0-9 = -9 < 0$, which by the third clause of Sylvester
        (sufficient on its own) gives \textbf{indefinite.}

  \item $\grad f = (3x^{2}+3y^{2}-15,\,6xy-12) = (0,0)$. From the second:
        $xy=2$. From the first: $x^{2}+y^{2}=5$. Combined with $xy=2$:
        $(x+y)^{2}=5+4=9$ so $x+y=\pm3$; $(x-y)^{2}=5-4=1$ so
        $x-y=\pm1$. Solutions $(x,y)\in\{(2,1),(1,2),(-2,-1),(-1,-2)\}$.
        Hessian $\Hess = \begin{pmatrix}6x&6y\\6y&6x\end{pmatrix}$, so
        $D_1=6x$, $D_2=36(x^{2}-y^{2})$.
        \begin{itemize}
          \item $(2,1)$: $D_1=12>0$, $D_2=36(4-1)=108>0$, \textbf{PD}:
                local minimum, $f=8+6-30-12=-28$.
          \item $(1,2)$: $D_1=6>0$, $D_2=36(1-4)=-108<0$, \textbf{saddle}.
          \item $(-2,-1)$: $D_1=-12<0$, $D_2=108>0$, \textbf{ND}: local
                maximum, $f=-8-6+30+12=28$.
          \item $(-1,-2)$: $D_1=-6<0$, $D_2=-108<0$, $D_2<0$ gives
                \textbf{saddle}.
        \end{itemize}

  \item Let $S=1+x^{2}+y^{2}$, $f=\ln S$, $f_x = 2x/S$, $f_y = 2y/S$.
        Then (using quotient rule):
        \[
          f_{xx} = \frac{2S-2x\cdot2x}{S^{2}} = \frac{2(1+y^{2}-x^{2})}{S^{2}}, \quad
          f_{yy} = \frac{2(1+x^{2}-y^{2})}{S^{2}}, \quad
          f_{xy} = \frac{-4xy}{S^{2}} = f_{yx} . \; \checkmark
        \]
        At the origin $(0,0)$: $S=1$, so
        \[
          \Hess(0,0) = \begin{pmatrix}2&0\\0&2\end{pmatrix} .
        \]

  \item $\grad f = \mathrm{e}^{-(x^{2}+y^{2})}(-2x,-2y) = (0,0)
        \Rightarrow x=y=0$, since $\mathrm{e}^{-r^{2}}>0$. Critical point
        $(0,0)$. Compute the Hessian: $f_x = -2x\mathrm{e}^{-(x^{2}+y^{2})}$,
        and at $(0,0)$, $f_{xx} = -2\mathrm{e}^{0}+(-2x)(-2x)\mathrm{e}^{0}\big|_{0,0}=-2$.
        Similarly $f_{yy}(0,0)=-2$, and
        $f_{xy}(0,0) = (-2x)(-2y)\mathrm{e}^{0}\big|_{0,0}=0$. Hessian
        $\Hess(0,0) = \begin{pmatrix}-2&0\\0&-2\end{pmatrix}$: $D_1=-2<0$,
        $D_2=4>0$, \textbf{ND: strict local maximum} at $(0,0)$ with
        $f(0,0)=1$. (Indeed, this is the global maximum, since $f\to0$
        at infinity and $f>0$ everywhere.)
\end{solutions}

\section{Additional Problems}

\begin{problems}
  \item Compute all second partials of $f(x,y) = 3x^{3} - 4x^{2}y + 5y^{3}$
        and verify Young's theorem.
  \item Construct the Hessian of
        $f(x,y) = 3x^{2} - 4xy + 2y^{2} + x - 5y$.
  \item Classify the definiteness of
        $\Hess = \begin{pmatrix}5&-1\\-1&4\end{pmatrix}$ using Sylvester's
        criterion.
  \item Find and classify the critical points of
        $f(x,y) = 2x^{2} + 3y^{2} - 4x + 6y + 1$.
  \item Compute all second partials of $f(x,y) = (2x+y)^{4}$ and verify
        $f_{xy} = f_{yx}$.
  \item Construct the Hessian of $f(x,y) = \mathrm{e}^{x+2y}$.
  \item Classify $\Hess = \begin{pmatrix}-3&1\\1&-2\end{pmatrix}$.
  \item Find and classify the critical points of
        $f(x,y) = -2x^{2} - 4y^{2} + 8x + 8y - 3$.
  \item Compute all second partials of $f(x,y) = xy\ln y$ ($y>0$) and
        verify Young's theorem.
  \item Construct the Hessian of $f(x,y) = x^{3}y + xy^{3}$.
  \item Classify $\Hess = \begin{pmatrix}2&3\\3&4\end{pmatrix}$.
  \item Find and classify the critical points of
        $f(x,y) = x^{2} - 2xy + 3y^{2} + 2x - 6y$.
  \item Compute all second partials of the Cobb--Douglas
        $f(K,L) = K^{1/2}L^{1/2}$ ($K,L>0$) and verify Young's theorem.
  \item Find and classify the critical points of
        $f(x,y) = x^{3} + y^{3} - 3x - 3y$.
  \item Classify $\Hess = \begin{pmatrix}-5&0\\0&-3\end{pmatrix}$ and state
        what it implies if $H$ is the Hessian at a critical point.
\end{problems}

\section{Additional Solutions}

\begin{solutions}
  \item $f_x = 9x^{2} - 8xy$, $f_y = -4x^{2} + 15y^{2}$. Then
        $f_{xx} = 18x - 8y$, $f_{yy} = 30y$, and
        \[
          f_{xy} = -8x = f_{yx} . \; \checkmark
        \]

  \item $f_x = 6x - 4y + 1$, $f_y = -4x + 4y - 5$. Hessian
        \[
          \Hess = \begin{pmatrix}6&-4\\-4&4\end{pmatrix} .
        \]

  \item $D_1 = 5 > 0$, $D_2 = 20 - 1 = 19 > 0$. \textbf{PD.}

  \item $\grad f = (4x-4,\,6y+6) = (0,0) \Rightarrow x=1,\,y=-1$. Hessian
        $\Hess = \begin{pmatrix}4&0\\0&6\end{pmatrix}$: $D_1=4>0$,
        $D_2=24>0$, \textbf{PD}, so $(1,-1)$ is a strict \textbf{local
        minimum} with $f(1,-1) = 2+3-4-6+1 = -4$.

  \item Chain rule: $f_x = 4(2x+y)^{3}\cdot2 = 8(2x+y)^{3}$,
        $f_y = 4(2x+y)^{3}$. So $f_{xx} = 48(2x+y)^{2}$,
        $f_{yy} = 12(2x+y)^{2}$, and
        \[
          f_{xy} = 24(2x+y)^{2} = f_{yx} . \; \checkmark
        \]

  \item $f_x = \mathrm{e}^{x+2y}$, $f_y = 2\mathrm{e}^{x+2y}$. So
        $f_{xx} = \mathrm{e}^{x+2y}$, $f_{yy} = 4\mathrm{e}^{x+2y}$,
        $f_{xy} = 2\mathrm{e}^{x+2y}$. Hessian
        \[
          \Hess = \mathrm{e}^{x+2y}\begin{pmatrix}1&2\\2&4\end{pmatrix} .
        \]
        (Note $\det\Hess = 0$: this is rank-1 PSD, consistent with $f$
        being a function of the single combination $x+2y$.)

  \item $D_1 = -3 < 0$, $D_2 = 6 - 1 = 5 > 0$. \textbf{ND.}

  \item $\grad f = (-4x+8,\,-8y+8) = (0,0) \Rightarrow x=2,\,y=1$. Hessian
        $\Hess = \begin{pmatrix}-4&0\\0&-8\end{pmatrix}$: $D_1=-4<0$,
        $D_2=32>0$, \textbf{ND}, so $(2,1)$ is a strict \textbf{local
        maximum} with $f(2,1) = -8-4+16+8-3 = 9$.

  \item $f_x = y\ln y$, $f_y = x\ln y + x = x(\ln y + 1)$. So
        $f_{xx} = 0$, $f_{yy} = x/y$, and
        \[
          f_{xy} = \ln y + 1 = f_{yx} . \; \checkmark
        \]

  \item $f_x = 3x^{2}y + y^{3}$, $f_y = x^{3} + 3xy^{2}$. So
        $f_{xx} = 6xy$, $f_{yy} = 6xy$,
        $f_{xy} = 3x^{2} + 3y^{2} = f_{yx}$. Hessian
        \[
          \Hess = \begin{pmatrix}
            6xy & 3x^{2}+3y^{2} \\
            3x^{2}+3y^{2} & 6xy
          \end{pmatrix} .
        \]

  \item $D_1 = 2 > 0$, $D_2 = 8 - 9 = -1 < 0$. \textbf{Indefinite
        (saddle).}

  \item $\grad f = (2x-2y+2,\,-2x+6y-6) = (0,0)$. From the first:
        $x = y-1$. Substitute: $-2(y-1)+6y-6 = 4y-4 = 0 \Rightarrow y=1$,
        $x=0$. Hessian $\Hess = \begin{pmatrix}2&-2\\-2&6\end{pmatrix}$:
        $D_1=2>0$, $D_2=12-4=8>0$, \textbf{PD}, so $(0,1)$ is a strict
        \textbf{local minimum} with $f(0,1) = 0-0+3+0-6 = -3$.

  \item $f_K = \tfrac12K^{-1/2}L^{1/2}$, $f_L = \tfrac12K^{1/2}L^{-1/2}$.
        So
        \[
          f_{KK} = -\tfrac14K^{-3/2}L^{1/2}, \quad
          f_{LL} = -\tfrac14K^{1/2}L^{-3/2}, \quad
          f_{KL} = \tfrac14K^{-1/2}L^{-1/2} = f_{LK} . \; \checkmark
        \]

  \item $\grad f = (3x^{2}-3,\,3y^{2}-3) = (0,0) \Rightarrow x^{2}=1,\,
        y^{2}=1$, giving $(\pm1,\pm1)$. Hessian
        $\Hess = \begin{pmatrix}6x&0\\0&6y\end{pmatrix}$, so $D_1=6x$,
        $D_2=36xy$.
        \begin{itemize}
          \item $(1,1)$: $D_1=6>0$, $D_2=36>0$, \textbf{PD}: local
                minimum, $f = 1+1-3-3 = -4$.
          \item $(1,-1)$: $D_2=-36<0$, \textbf{saddle}.
          \item $(-1,1)$: $D_2=-36<0$, \textbf{saddle}.
          \item $(-1,-1)$: $D_1=-6<0$, $D_2=36>0$, \textbf{ND}: local
                maximum, $f = -1-1+3+3 = 4$.
        \end{itemize}

  \item $D_1 = -5 < 0$, $D_2 = 15 > 0$. \textbf{ND}: at a critical point
        this would mean a \textbf{strict local maximum.}
\end{solutions}

\chapter{Multivariable Chain Rule and Total Differential}

When several inputs move at once, the total differential bundles the partial
derivatives into one linear approximation, and the multivariable chain rule
pipes that approximation through any composed function. Together they
generate the marginal rate of substitution along an indifference curve, the
marginal rate of technical substitution along an isoquant, and (as a
preview) the envelope theorem along an optimal value function.

\section{Core Rules}

\begin{rulebox}{Total Differential}
If $z=f(x,y)$ is differentiable at $(x,y)$, the \textbf{total differential}
is
\[
  \mathrm{d}z \;=\; f_{x}(x,y)\,\mathrm{d}x \;+\; f_{y}(x,y)\,\mathrm{d}y .
\]

\textbf{Interpretation.} $\mathrm{d}z$ is the best linear estimate of the
actual change $\Delta z$ when $x$ and $y$ shift by $\mathrm{d}x$ and
$\mathrm{d}y$. In economics, $\mathrm{d}U \approx MU_{x}\,\mathrm{d}x +
MU_{y}\,\mathrm{d}y$ tells how utility responds to a small basket change;
$\mathrm{d}Q \approx MP_{K}\,\mathrm{d}K + MP_{L}\,\mathrm{d}L$ does the same
for output.

For $n$ inputs, $z=f(x_{1},\ldots,x_{n})$:
\[
  \mathrm{d}z \;=\; \sum_{i=1}^{n} \pd{f}{x_{i}}\,\mathrm{d}x_{i} .
\]
\end{rulebox}

\begin{rulebox}{Chain Rule, Single Driver $t$}
If $z=f(x,y)$ with $x=x(t)$ and $y=y(t)$,
\[
  \dd{z}{t} \;=\; f_{x}\,\dd{x}{t} \;+\; f_{y}\,\dd{y}{t} .
\]
\end{rulebox}

\begin{rulebox}{Chain Rule, Two Drivers $s,t$}
If $z=f(x,y)$ with $x=x(s,t)$, $y=y(s,t)$,
\[
  \pd{z}{s} \;=\; f_{x}\,\pd{x}{s} + f_{y}\,\pd{y}{s},\qquad
  \pd{z}{t} \;=\; f_{x}\,\pd{x}{t} + f_{y}\,\pd{y}{t} .
\]
\end{rulebox}

\begin{rulebox}{Level Sets, MRS, MRTS}
On a level curve $f(x,y)=c$ the total differential is zero:
\[
  f_{x}\,\mathrm{d}x + f_{y}\,\mathrm{d}y = 0
  \;\Longrightarrow\;
  \dd{y}{x}\bigg|_{f=c} = -\frac{f_{x}}{f_{y}} .
\]
For a utility $U(x,y)$, the (signed) slope $-U_{x}/U_{y}$ is minus the
\textbf{MRS}; for a production function $Q(K,L)$, $-Q_{K}/Q_{L}$ is minus
the \textbf{MRTS}. We will state explicitly which sign convention we use in
each problem.
\end{rulebox}

\section{Worked Example}

\begin{worked}{With $z=x^{2}y$, $x=t^{2}$, $y=\ln t$, find $\mathrm{d}z/\mathrm{d}t$ two ways; then use the total differential on $Q=K^{0.3}L^{0.7}$}
\wstep{Step 1. Draw the dependency before writing a formula.}
$t$ drives both $x$ and $y$; $x$ and $y$ together drive $z$. So there are
\textbf{two paths} from $t$ to $z$, and the chain rule adds one term per
path:
\[
  \dd{z}{t}=\pd{z}{x}\dd{x}{t}+\pd{z}{y}\dd{y}{t} .
\]
\wtrap{One term per path. Losing a term is the standard error, and it is
invisible in the answer --- the result still looks like a derivative.
Sketching the tree (one node $t$, two arrows to $x$ and $y$, two arrows into
$z$) takes five seconds and prevents it.}

\wstep{Step 2. Note which symbol goes where.}
$\partial z/\partial x$ is \emph{partial}, because $z$ depends on two
things. $\mathrm{d}x/\mathrm{d}t$ is \emph{total}, because $x$ depends on
$t$ alone. And $\mathrm{d}z/\mathrm{d}t$ on the left is total, because once
the substitutions are made $z$ is a function of $t$ only. The notation is
carrying real information here.

\wstep{Step 3. Compute the four pieces.}
\[
  \pd{z}{x}=2xy,\qquad
  \pd{z}{y}=x^{2},\qquad
  \dd{x}{t}=2t,\qquad
  \dd{y}{t}=\frac1t .
\]

\wstep{Step 4. Assemble, then substitute the paths.}
\[
  \dd{z}{t}=2xy\cdot2t+x^{2}\cdot\frac1t
  =2t^{2}\ln t\cdot2t+\frac{t^{4}}{t}
  =\boxed{4t^{3}\ln t+t^{3}}
\]
\wtrap{Substitute $x=t^{2}$ and $y=\ln t$ only \emph{after} assembling the
formula. Substituting first turns the exercise into a one-variable problem
and defeats its purpose --- and in the many cases where $x(t)$ is not known
explicitly, it is not even possible.}

\wstep{Step 5. Check by direct substitution, which here is available.}
$z=(t^{2})^{2}\ln t=t^{4}\ln t$, so
\[
  \dd{z}{t}=4t^{3}\ln t+t^{4}\cdot\frac1t=4t^{3}\ln t+t^{3} . \checkmark
\]
Two routes, one answer. When both are available, use the direct one as the
check and the chain rule as the method.

\wstep{Step 6. The total differential: the chain rule with the $t$ removed.}
\[
  \mathrm{d}z=\pd{z}{x}\mathrm{d}x+\pd{z}{y}\mathrm{d}y
\]
is the first-order approximation to the change in $z$ caused by small
changes in $x$ and $y$ --- the two-variable linearisation. For
$Q=K^{0.3}L^{0.7}$:
\[
  \mathrm{d}Q=0.3K^{-0.7}L^{0.7}\,\mathrm{d}K+0.7K^{0.3}L^{-0.3}\,\mathrm{d}L .
\]

\wstep{Step 7. Divide by $Q$ and watch it turn into percentages.}
\[
  \frac{\mathrm{d}Q}{Q}
  =0.3\,\frac{\mathrm{d}K}{K}+0.7\,\frac{\mathrm{d}L}{L} .
\]
Every power cancels. This is the growth-accounting equation: the percentage
change in output is the \emph{exponent-weighted average} of the percentage
changes in inputs, and the exponents are the elasticities. If capital rises
$2\%$ while labour falls $1\%$:
\[
  \frac{\mathrm{d}Q}{Q}=0.3(0.02)+0.7(-0.01)=0.006-0.007=-0.001=-0.1\% .
\]
\wtrap{This identity holds only for Cobb--Douglas, where the elasticities
are the constant exponents. For a general $F(K,L)$ the weights are
$\frac{K F_{K}}{F}$ and $\frac{L F_{L}}{F}$, which vary from point to
point. Applying the constant-weight version to a CES function is a real
error.}

\wstep{Step 8. Two drivers instead of one.}
If $z=f(x,y)$ with $x=x(s,t)$ and $y=y(s,t)$, the same path-counting gives
one equation per driver:
\[
  \pd{z}{t}=\pd{z}{x}\pd{x}{t}+\pd{z}{y}\pd{y}{t},
  \qquad
  \pd{z}{s}=\pd{z}{x}\pd{x}{s}+\pd{z}{y}\pd{y}{s} .
\]
All derivatives are now partial, because two independent variables remain.
In matrix form these two lines are $\grad_{(s,t)}z=J^{\!\top}\grad_{(x,y)}z$
with $J$ the Jacobian of \chref{derivatives}{13}: the same rule, written once.
\end{worked}

\section{Practice Problems}

\begin{problems}
  \item Compute $\mathrm{d}z$ for $z=4x^{2}-5xy+2y^{2}$.
  \item Find $\dd{z}{t}$ for $z=x^{2}-y^{2}$ with $x=2t+1$, $y=t^{2}$.
  \item For $U(x,y)=x^{2}+4y=c$, find $\dd{y}{x}$ and state the unsigned
        MRS.
  \item Compute $\pd{z}{s}$ and $\pd{z}{t}$ for $z=x^{2}-3y$ with
        $x=s+2t$, $y=2s-t$.
  \item Compute $\mathrm{d}z$ for $z=x^{4}y^{-2}$.
  \item Find $\dd{z}{t}$ for $z=\mathrm{e}^{xy}$ with $x=\sin t$,
        $y=\cos t$.
  \item For the isoquant $Q(K,L)=K^{0.4}L^{0.6}=Q_{0}$, find $\dd{L}{K}$
        and the unsigned MRTS.
  \item Compute $\mathrm{d}z$ for $z=\ln(3x+2y)$.
  \item Compute $\pd{z}{s}$ and $\pd{z}{t}$ for $z=\mathrm{e}^{x}\cos y$
        with $x=st$, $y=s+t$.
  \item Find $\dd{z}{t}$ for $z=x^{2}y$ with $x=\ln t$, $y=t^{3}$.
  \item For $U(x,y)=x^{2/3}y^{1/3}=c$, find $\dd{y}{x}$ and the unsigned
        MRS.
  \item Compute $\mathrm{d}z$ for $z=\mathrm{e}^{x^{2}+y}$.
  \item Find $\dd{z}{t}$ for $z=\sqrt{x+y^{2}}$ with $x=t^{2}$, $y=t$.
  \item Compute $\pd{z}{s}$ and $\pd{z}{t}$ for $z=\ln(x^{2}+y^{2})$ with
        $x=s/t$, $y=t/s$.
  \item For $U(x,y)=\sqrt{x}+\sqrt{y}=c$, find $\dd{y}{x}$ and the unsigned
        MRS.
  \item Compute $\mathrm{d}z$ for $z=\dfrac{x+y}{x-y}$.
  \item Find $\dd{z}{t}$ for $z=\arctan(y/x)$ with $x=\mathrm{e}^{t}$,
        $y=\mathrm{e}^{2t}$.
  \item For $Q(K,L)=(aK^{1/2}+bL^{1/2})^{2}=Q_{0}$ (with $a,b>0$), find
        $\dd{L}{K}$ and the unsigned MRTS.
  \item Compute $\pd{z}{s}$ and $\pd{z}{t}$ for $z=x^{1/2}y^{1/2}$ with
        $x=s^{2}+t^{2}$, $y=st$.
  \item Compute $\mathrm{d}z$ for $z=x\sin y+y\cos x$.
  \item For $U(x,y)=\ln(x+1)+\ln(y+1)=c$, find $\dd{y}{x}$ and the
        unsigned MRS.
  \item Find $\dd{z}{t}$ for $z=x\ln y$ with $x=t^{2}+1$, $y=\mathrm{e}^{t}$.
  \item Compute $\mathrm{d}z$ for $z=x^{3}\mathrm{e}^{2y}$.
  \item Find $\dd{z}{t}$ for $z=x^{2}+y^{2}$ with $x=\cos t$, $y=\sin t$,
        and verify the result.
  \item For $U(x,y)=x+y^{2}=c$, find $\dd{y}{x}$ and state where the
        result is well-defined.
  \item Compute $\pd{z}{s}$ and $\pd{z}{t}$ for $z=xy^{2}$ with $x=s+t$,
        $y=s-t$.
  \item Compute $\mathrm{d}z$ for $z=\arctan(xy)$.
  \item Find $\dd{z}{t}$ for $z=\ln(x^{2}+y^{2})$ with $x=t$, $y=t^{2}$.
  \item For the isoquant $Q(K,L)=AK^{\alpha}L^{\beta}=Q_{0}$ (general
        Cobb--Douglas with $\alpha,\beta>0$), derive the MRTS in terms of
        $\alpha,\beta,K,L$.
  \item Compute $\mathrm{d}z$ for $z=\dfrac{x^{2}}{y}+\dfrac{y^{2}}{x}$.
  \item For $U(x,y)=\mathrm{e}^{x}+\mathrm{e}^{y}=c$, find $\dd{y}{x}$ and
        the unsigned MRS.
  \item Find $\dd{z}{t}$ for $z=\mathrm{e}^{x}\sin y$ with $x=2t$,
        $y=t^{2}$.
  \item Compute $\pd{z}{s}$ and $\pd{z}{t}$ for $z=x^{2}y-y^{3}$ with
        $x=s/t$, $y=st$.
  \item Compute $\mathrm{d}z$ for $z=(x+y)\mathrm{e}^{x-y}$.
  \item For the indifference curve $U(x,y)=\alpha\ln x+\beta\ln y=c$ (with
        $\alpha,\beta>0$), find $\dd{y}{x}$ and the MRS.
  \item Find $\dd{z}{t}$ for $z=\sqrt{xy}$ with $x=t^{3}$, $y=t$.
  \item Compute $\mathrm{d}z$ for $z=\ln(x^{2}+y^{2}+1)$.
  \item Compute $\pd{z}{s}$ and $\pd{z}{t}$ for $z=\mathrm{e}^{xy}$ with
        $x=s+t$, $y=s-t$.
  \item For $Q(K,L)=K+2\sqrt{KL}+L=Q_{0}$ (a perfect-square form), find
        the MRTS.
  \item Find $\dd{z}{t}$ for $z=\cos(xy)$ with $x=t$, $y=1/t$ (for $t>0$).
        Verify by direct substitution.
  \item For the indifference curve $U(x,y)=xy+x+y=c$, find $\dd{y}{x}$
        and the unsigned MRS.
  \item Use the total differential of $V(P,Y)=Y/P$ (real income) to
        approximate $\mathrm{d}V$ when $\mathrm{d}P=0.02P$ and
        $\mathrm{d}Y=0.03Y$. Interpret.
  \item Compute $\pd{z}{s}$ and $\pd{z}{t}$ for $z=x^{2}+2xy$ with
        $x=s\cos t$, $y=s\sin t$.
  \item For $U(x_{1},x_{2},x_{3})=x_{1}+2x_{2}+3x_{3}=c$, find
        $\pd{x_{3}}{x_{1}}$ and $\pd{x_{3}}{x_{2}}$ along the level set,
        and interpret as marginal rates of substitution.
\end{problems}

\section{Solutions}

\begin{solutions}
  \item $f_{x}=8x-5y$, $f_{y}=-5x+4y$.
        \[
          \mathrm{d}z \;=\; (8x-5y)\,\mathrm{d}x+(-5x+4y)\,\mathrm{d}y .
        \]

  \item $f_{x}=2x$, $f_{y}=-2y$, $\dd{x}{t}=2$, $\dd{y}{t}=2t$.
        \[
          \dd{z}{t}=2x\cdot 2+(-2y)\cdot 2t=4(2t+1)-4t^{3}=8t+4-4t^{3} .
        \]
        \textbf{Check.} $z=(2t+1)^{2}-t^{4}$, $\dd{z}{t}=8t+4-4t^{3}$.
        \checkmark

  \item $U_{x}=2x$, $U_{y}=4$. Signed slope $\dd{y}{x}=-2x/4=-x/2$.
        \textbf{MRS} $=x/2$.

  \item $f_{x}=2x$, $f_{y}=-3$; $x_{s}=1$, $x_{t}=2$, $y_{s}=2$,
        $y_{t}=-1$.
        \[
          \pd{z}{s}=2x(1)+(-3)(2)=2(s+2t)-6,\qquad
          \pd{z}{t}=2x(2)+(-3)(-1)=4(s+2t)+3 .
        \]

  \item $f_{x}=4x^{3}y^{-2}$, $f_{y}=-2x^{4}y^{-3}$.
        \[
          \mathrm{d}z \;=\; \frac{4x^{3}}{y^{2}}\,\mathrm{d}x
            -\frac{2x^{4}}{y^{3}}\,\mathrm{d}y .
        \]

  \item $f_{x}=y\mathrm{e}^{xy}$, $f_{y}=x\mathrm{e}^{xy}$;
        $\dd{x}{t}=\cos t$, $\dd{y}{t}=-\sin t$.
        \[
          \dd{z}{t}=\mathrm{e}^{xy}(y\cos t-x\sin t)
            =\mathrm{e}^{\sin t\cos t}(\cos^{2}t-\sin^{2}t)
            =\mathrm{e}^{\sin t\cos t}\cos(2t) .
        \]

  \item $Q_{K}=0.4K^{-0.6}L^{0.6}$, $Q_{L}=0.6K^{0.4}L^{-0.4}$.
        \[
          \frac{Q_{K}}{Q_{L}}=\frac{0.4}{0.6}\cdot\frac{L}{K}
            =\frac{2}{3}\,\frac{L}{K} .
        \]
        $\dd{L}{K}=-\tfrac{2}{3}L/K$. \textbf{MRTS} $=\tfrac{2}{3}L/K$.

  \item $f_{x}=3/(3x+2y)$, $f_{y}=2/(3x+2y)$.
        \[
          \mathrm{d}z \;=\; \frac{3\,\mathrm{d}x+2\,\mathrm{d}y}{3x+2y} .
        \]

  \item $f_{x}=\mathrm{e}^{x}\cos y$, $f_{y}=-\mathrm{e}^{x}\sin y$;
        $x_{s}=t$, $x_{t}=s$, $y_{s}=y_{t}=1$.
        \[
          \pd{z}{s}=\mathrm{e}^{st}\bigl(t\cos(s+t)-\sin(s+t)\bigr),\qquad
          \pd{z}{t}=\mathrm{e}^{st}\bigl(s\cos(s+t)-\sin(s+t)\bigr) .
        \]

  \item $f_{x}=2xy$, $f_{y}=x^{2}$; $\dd{x}{t}=1/t$, $\dd{y}{t}=3t^{2}$.
        \[
          \dd{z}{t}=2xy\cdot\frac{1}{t}+x^{2}\cdot 3t^{2}
            =\frac{2(\ln t)t^{3}}{t}+3t^{2}(\ln t)^{2}
            =2t^{2}\ln t+3t^{2}(\ln t)^{2} .
        \]

  \item $U_{x}=\tfrac{2}{3}x^{-1/3}y^{1/3}$,
        $U_{y}=\tfrac{1}{3}x^{2/3}y^{-2/3}$.
        \[
          \frac{U_{x}}{U_{y}}=2\,\frac{y}{x} .
        \]
        $\dd{y}{x}=-2y/x$. \textbf{MRS} $=2y/x$.

  \item $f_{x}=2x\mathrm{e}^{x^{2}+y}$, $f_{y}=\mathrm{e}^{x^{2}+y}$.
        \[
          \mathrm{d}z \;=\; \mathrm{e}^{x^{2}+y}(2x\,\mathrm{d}x+\mathrm{d}y) .
        \]

  \item $f_{x}=1/(2\sqrt{x+y^{2}})$, $f_{y}=y/\sqrt{x+y^{2}}$;
        $\dd{x}{t}=2t$, $\dd{y}{t}=1$.
        \[
          \dd{z}{t}=\frac{2t}{2\sqrt{x+y^{2}}}+\frac{y}{\sqrt{x+y^{2}}}
            =\frac{t+y}{\sqrt{x+y^{2}}}=\frac{2t}{\sqrt{2t^{2}}}
            =\frac{\sqrt{2}\,t}{|t|} .
        \]
        For $t>0$: $\dd{z}{t}=\sqrt{2}$. (Indeed $z=\sqrt{2}\,|t|$.)

  \item $f_{x}=2x/(x^{2}+y^{2})$, $f_{y}=2y/(x^{2}+y^{2})$;
        $x_{s}=1/t$, $x_{t}=-s/t^{2}$, $y_{s}=-t/s^{2}$, $y_{t}=1/s$.
        \[
          \pd{z}{s}=\frac{2}{x^{2}+y^{2}}\Bigl(\frac{x}{t}-\frac{yt}{s^{2}}\Bigr)
            =\frac{2}{x^{2}+y^{2}}\Bigl(\frac{s}{t^{2}}-\frac{t^{2}}{s^{3}}\Bigr) .
        \]
        \[
          \pd{z}{t}=\frac{2}{x^{2}+y^{2}}\Bigl(\frac{t}{s^{2}}-\frac{s^{2}}{t^{3}}\Bigr) .
        \]

  \item $U_{x}=1/(2\sqrt{x})$, $U_{y}=1/(2\sqrt{y})$.
        \[
          \frac{U_{x}}{U_{y}}=\sqrt{\frac{y}{x}} .
        \]
        $\dd{y}{x}=-\sqrt{y/x}$. \textbf{MRS} $=\sqrt{y/x}$.

  \item Quotient rule.
        \[
          f_{x}=\frac{(x-y)-(x+y)}{(x-y)^{2}}=\frac{-2y}{(x-y)^{2}},\qquad
          f_{y}=\frac{(x-y)+(x+y)}{(x-y)^{2}}=\frac{2x}{(x-y)^{2}} .
        \]
        \[
          \mathrm{d}z \;=\; \frac{-2y\,\mathrm{d}x+2x\,\mathrm{d}y}{(x-y)^{2}} .
        \]

  \item Let $u=y/x$. $f_{x}=-y/(x^{2}+y^{2})$, $f_{y}=x/(x^{2}+y^{2})$;
        $\dd{x}{t}=\mathrm{e}^{t}$, $\dd{y}{t}=2\mathrm{e}^{2t}$.
        \[
          \dd{z}{t}=\frac{-y\mathrm{e}^{t}+2x\mathrm{e}^{2t}}{x^{2}+y^{2}} .
        \]
        Substitute $x=\mathrm{e}^{t}$, $y=\mathrm{e}^{2t}$,
        $x^{2}+y^{2}=\mathrm{e}^{2t}(1+\mathrm{e}^{2t})$:
        \[
          \dd{z}{t}=\frac{-\mathrm{e}^{3t}+2\mathrm{e}^{3t}}
            {\mathrm{e}^{2t}(1+\mathrm{e}^{2t})}
            =\frac{\mathrm{e}^{t}}{1+\mathrm{e}^{2t}} .
        \]

  \item Let $A=aK^{1/2}+bL^{1/2}$, $Q=A^{2}$. $Q_{K}=2A\cdot(a/2)K^{-1/2}
        =aAK^{-1/2}$, $Q_{L}=bAL^{-1/2}$.
        \[
          \frac{Q_{K}}{Q_{L}}=\frac{a}{b}\sqrt{\frac{L}{K}} .
        \]
        $\dd{L}{K}=-(a/b)\sqrt{L/K}$. \textbf{MRTS} $=(a/b)\sqrt{L/K}$.

  \item $f_{x}=\tfrac12\sqrt{y/x}$, $f_{y}=\tfrac12\sqrt{x/y}$;
        $x_{s}=2s$, $x_{t}=2t$, $y_{s}=t$, $y_{t}=s$.
        \[
          \pd{z}{s}=s\sqrt{\frac{y}{x}}+\frac{t}{2}\sqrt{\frac{x}{y}},\qquad
          \pd{z}{t}=t\sqrt{\frac{y}{x}}+\frac{s}{2}\sqrt{\frac{x}{y}} .
        \]

  \item $f_{x}=\sin y-y\sin x$, $f_{y}=x\cos y+\cos x$.
        \[
          \mathrm{d}z \;=\; (\sin y-y\sin x)\,\mathrm{d}x
            +(x\cos y+\cos x)\,\mathrm{d}y .
        \]

  \item $U_{x}=1/(x+1)$, $U_{y}=1/(y+1)$.
        \[
          \frac{U_{x}}{U_{y}}=\frac{y+1}{x+1} .
        \]
        $\dd{y}{x}=-(y+1)/(x+1)$. \textbf{MRS} $=(y+1)/(x+1)$.

  \item $f_{x}=\ln y$, $f_{y}=x/y$; $\dd{x}{t}=2t$, $\dd{y}{t}=\mathrm{e}^{t}$.
        \[
          \dd{z}{t}=(\ln y)(2t)+\frac{x}{y}\mathrm{e}^{t}
            =t(2t)+\frac{(t^{2}+1)\mathrm{e}^{t}}{\mathrm{e}^{t}}
            =3t^{2}+1 .
        \]
        \textbf{Check.} $z=(t^{2}+1)t=t^{3}+t$, $\dd{z}{t}=3t^{2}+1$.
        \checkmark

  \item $f_{x}=3x^{2}\mathrm{e}^{2y}$, $f_{y}=2x^{3}\mathrm{e}^{2y}$.
        \[
          \mathrm{d}z \;=\; x^{2}\mathrm{e}^{2y}(3\,\mathrm{d}x+2x\,\mathrm{d}y) .
        \]

  \item $f_{x}=2x$, $f_{y}=2y$; $\dd{x}{t}=-\sin t$, $\dd{y}{t}=\cos t$.
        \[
          \dd{z}{t}=2x(-\sin t)+2y(\cos t)=-2\cos t\sin t+2\sin t\cos t=0 .
        \]
        \textbf{Check.} $z=\cos^{2}t+\sin^{2}t=1$, so $\dd{z}{t}=0$.
        \checkmark

  \item $U_{x}=1$, $U_{y}=2y$. Signed slope: $\dd{y}{x}=-1/(2y)$,
        well-defined for $y\neq 0$. The MRS in the usual sense is
        $1/(2y)$.

  \item $f_{x}=y^{2}$, $f_{y}=2xy$; $x_{s}=1$, $x_{t}=1$, $y_{s}=1$,
        $y_{t}=-1$.
        \[
          \pd{z}{s}=y^{2}(1)+2xy(1)=y^{2}+2xy=(s-t)^{2}+2(s+t)(s-t) .
        \]
        Simplify: $(s-t)^{2}+2(s^{2}-t^{2})=s^{2}-2st+t^{2}+2s^{2}-2t^{2}
          =3s^{2}-2st-t^{2}$.
        \[
          \pd{z}{t}=y^{2}(1)+2xy(-1)=y^{2}-2xy=(s-t)^{2}-2(s+t)(s-t)
            =-s^{2}-2st+3t^{2} .
        \]

  \item $f_{x}=y/(1+x^{2}y^{2})$, $f_{y}=x/(1+x^{2}y^{2})$.
        \[
          \mathrm{d}z \;=\; \frac{y\,\mathrm{d}x+x\,\mathrm{d}y}{1+x^{2}y^{2}} .
        \]

  \item $f_{x}=2x/(x^{2}+y^{2})$, $f_{y}=2y/(x^{2}+y^{2})$;
        $\dd{x}{t}=1$, $\dd{y}{t}=2t$.
        \[
          \dd{z}{t}=\frac{2x+2y(2t)}{x^{2}+y^{2}}
            =\frac{2t+4t^{3}}{t^{2}+t^{4}}
            =\frac{2t(1+2t^{2})}{t^{2}(1+t^{2})}
            =\frac{2(1+2t^{2})}{t(1+t^{2})} .
        \]

  \item $Q_{K}=A\alpha K^{\alpha-1}L^{\beta}$,
        $Q_{L}=A\beta K^{\alpha}L^{\beta-1}$.
        \[
          \frac{Q_{K}}{Q_{L}}=\frac{\alpha}{\beta}\cdot\frac{L}{K} .
        \]
        $\dd{L}{K}=-(\alpha/\beta)(L/K)$.
        \textbf{MRTS} $=(\alpha/\beta)(L/K)$.

  \item $f_{x}=2x/y-y^{2}/x^{2}$, $f_{y}=-x^{2}/y^{2}+2y/x$.
        \[
          \mathrm{d}z \;=\; \Bigl(\frac{2x}{y}-\frac{y^{2}}{x^{2}}\Bigr)\mathrm{d}x
            +\Bigl(\frac{2y}{x}-\frac{x^{2}}{y^{2}}\Bigr)\mathrm{d}y .
        \]

  \item $U_{x}=\mathrm{e}^{x}$, $U_{y}=\mathrm{e}^{y}$.
        \[
          \frac{U_{x}}{U_{y}}=\mathrm{e}^{x-y} .
        \]
        $\dd{y}{x}=-\mathrm{e}^{x-y}$. \textbf{MRS} $=\mathrm{e}^{x-y}$.

  \item $f_{x}=\mathrm{e}^{x}\sin y$, $f_{y}=\mathrm{e}^{x}\cos y$;
        $\dd{x}{t}=2$, $\dd{y}{t}=2t$.
        \[
          \dd{z}{t}=2\mathrm{e}^{x}\sin y+2t\mathrm{e}^{x}\cos y
            =2\mathrm{e}^{2t}(\sin t^{2}+t\cos t^{2}) .
        \]

  \item $f_{x}=2xy$, $f_{y}=x^{2}-3y^{2}$; $x_{s}=1/t$, $x_{t}=-s/t^{2}$,
        $y_{s}=t$, $y_{t}=s$. Note $xy=(s/t)(st)=s^{2}$.
        \[
          \pd{z}{s}=2xy\cdot\frac{1}{t}+(x^{2}-3y^{2})\cdot t
            =\frac{2s^{2}}{t}+t\Bigl(\frac{s^{2}}{t^{2}}-3s^{2}t^{2}\Bigr)
            =\frac{2s^{2}}{t}+\frac{s^{2}}{t}-3s^{2}t^{3}
            =\frac{3s^{2}}{t}-3s^{2}t^{3} .
        \]
        \[
          \pd{z}{t}=2xy\cdot\Bigl(-\frac{s}{t^{2}}\Bigr)+(x^{2}-3y^{2})\cdot s
            =-\frac{2s^{3}}{t^{2}}+s\Bigl(\frac{s^{2}}{t^{2}}-3s^{2}t^{2}\Bigr)
            =-\frac{s^{3}}{t^{2}}-3s^{3}t^{2} .
        \]
        \textbf{Check.} $z=(s/t)^{2}(st)-(st)^{3}=s^{3}/t-s^{3}t^{3}$. Then
        $\pd{z}{s}=3s^{2}/t-3s^{2}t^{3}$ and
        $\pd{z}{t}=-s^{3}/t^{2}-3s^{3}t^{2}$. \checkmark

  \item Product rule: $f_{x}=\mathrm{e}^{x-y}+(x+y)\mathrm{e}^{x-y}
        =(1+x+y)\mathrm{e}^{x-y}$, and
        $f_{y}=\mathrm{e}^{x-y}-(x+y)\mathrm{e}^{x-y}=(1-x-y)\mathrm{e}^{x-y}$.
        \[
          \mathrm{d}z \;=\; \mathrm{e}^{x-y}\bigl((1+x+y)\,\mathrm{d}x
            +(1-x-y)\,\mathrm{d}y\bigr) .
        \]

  \item $U_{x}=\alpha/x$, $U_{y}=\beta/y$.
        \[
          \frac{U_{x}}{U_{y}}=\frac{\alpha}{\beta}\cdot\frac{y}{x} .
        \]
        $\dd{y}{x}=-(\alpha/\beta)(y/x)$.
        \textbf{MRS} $=(\alpha/\beta)(y/x)$. (Standard log-Cobb-Douglas.)

  \item $f_{x}=\tfrac12\sqrt{y/x}$, $f_{y}=\tfrac12\sqrt{x/y}$;
        $\dd{x}{t}=3t^{2}$, $\dd{y}{t}=1$.
        \[
          \dd{z}{t}=\frac{3t^{2}}{2}\sqrt{\frac{y}{x}}
            +\frac{1}{2}\sqrt{\frac{x}{y}} .
        \]
        Substitute $y/x=1/t^{2}$, so $\sqrt{y/x}=1/t$ (for $t>0$), and
        $\sqrt{x/y}=t$:
        \[
          \dd{z}{t}=\frac{3t^{2}}{2}\cdot\frac{1}{t}+\frac{t}{2}
            =\frac{3t}{2}+\frac{t}{2}=2t .
        \]
        \textbf{Check.} $z=\sqrt{t^{3}\cdot t}=\sqrt{t^{4}}=t^{2}$,
        $\dd{z}{t}=2t$. \checkmark

  \item $f_{x}=2x/(x^{2}+y^{2}+1)$, $f_{y}=2y/(x^{2}+y^{2}+1)$.
        \[
          \mathrm{d}z \;=\; \frac{2x\,\mathrm{d}x+2y\,\mathrm{d}y}{x^{2}+y^{2}+1} .
        \]

  \item $f_{x}=y\mathrm{e}^{xy}$, $f_{y}=x\mathrm{e}^{xy}$; $x_{s}=1$,
        $x_{t}=1$, $y_{s}=1$, $y_{t}=-1$.
        \[
          \pd{z}{s}=y\mathrm{e}^{xy}+x\mathrm{e}^{xy}=(x+y)\mathrm{e}^{xy}
            =2s\,\mathrm{e}^{s^{2}-t^{2}} .
        \]
        \[
          \pd{z}{t}=y\mathrm{e}^{xy}-x\mathrm{e}^{xy}=(y-x)\mathrm{e}^{xy}
            =-2t\,\mathrm{e}^{s^{2}-t^{2}} .
        \]

  \item Note $Q=(\sqrt{K}+\sqrt{L})^{2}$. So $Q_{K}=2(\sqrt{K}+\sqrt{L})
        \cdot\dfrac{1}{2\sqrt{K}}=(\sqrt{K}+\sqrt{L})/\sqrt{K}$, and
        $Q_{L}=(\sqrt{K}+\sqrt{L})/\sqrt{L}$.
        \[
          \frac{Q_{K}}{Q_{L}}=\frac{\sqrt{L}}{\sqrt{K}}=\sqrt{\frac{L}{K}} .
        \]
        $\dd{L}{K}=-\sqrt{L/K}$. \textbf{MRTS} $=\sqrt{L/K}$.

  \item $f_{x}=-y\sin(xy)$, $f_{y}=-x\sin(xy)$; $\dd{x}{t}=1$,
        $\dd{y}{t}=-1/t^{2}$.
        \[
          \dd{z}{t}=-y\sin(xy)+\bigl(-x\sin(xy)\bigr)\Bigl(-\frac{1}{t^{2}}\Bigr)
            =-\sin(xy)\Bigl(y-\frac{x}{t^{2}}\Bigr) .
        \]
        With $x=t$, $y=1/t$: $xy=1$, and $y-x/t^{2}=1/t-1/t=0$. Hence
        $\dd{z}{t}=0$. \textbf{Check.} $z=\cos(xy)=\cos 1$ is constant.
        \checkmark

  \item $U_{x}=y+1$, $U_{y}=x+1$.
        \[
          \frac{U_{x}}{U_{y}}=\frac{y+1}{x+1} .
        \]
        $\dd{y}{x}=-(y+1)/(x+1)$. \textbf{MRS} $=(y+1)/(x+1)$.

  \item $V_{P}=-Y/P^{2}$, $V_{Y}=1/P$.
        \[
          \mathrm{d}V \;=\; -\frac{Y}{P^{2}}\,\mathrm{d}P+\frac{1}{P}\,\mathrm{d}Y
            =\frac{V}{1}\Bigl(-\frac{\mathrm{d}P}{P}+\frac{\mathrm{d}Y}{Y}\Bigr)\cdot 1
        \]
        since $V=Y/P$, we have $-Y\,\mathrm{d}P/P^{2}=-V\,\mathrm{d}P/P$
        and $\mathrm{d}Y/P=V\,\mathrm{d}Y/Y$. So
        \[
          \frac{\mathrm{d}V}{V} \;=\; \frac{\mathrm{d}Y}{Y}-\frac{\mathrm{d}P}{P}
            \;=\; 0.03-0.02 \;=\; 0.01 .
        \]
        Real income rises by about $1\%$. Interpretation: nominal income
        growth less inflation gives real income growth (log-linearised).

  \item $f_{x}=2x+2y$, $f_{y}=2x$; $x_{s}=\cos t$, $x_{t}=-s\sin t$,
        $y_{s}=\sin t$, $y_{t}=s\cos t$.
        \[
          \pd{z}{s}=(2x+2y)\cos t+2x\sin t .
        \]
        Substitute $x=s\cos t$, $y=s\sin t$: $2x+2y=2s(\cos t+\sin t)$,
        so
        \[
          \pd{z}{s}=2s(\cos t+\sin t)\cos t+2s\cos t\sin t
            =2s\cos^{2}t+4s\sin t\cos t
            =2s(\cos^{2}t+2\sin t\cos t) .
        \]
        \[
          \pd{z}{t}=(2x+2y)(-s\sin t)+2x(s\cos t)
            =-2s^{2}(\cos t+\sin t)\sin t+2s^{2}\cos^{2}t .
        \]
        Expanding: $-2s^{2}(\cos t\sin t+\sin^{2}t)+2s^{2}\cos^{2}t
          =2s^{2}(\cos^{2}t-\sin^{2}t-\sin t\cos t)
          =2s^{2}\bigl(\cos 2t-\tfrac12\sin 2t\bigr)$.

  \item On $U=c$, $\mathrm{d}U=\mathrm{d}x_{1}+2\,\mathrm{d}x_{2}
        +3\,\mathrm{d}x_{3}=0$. Solve for $\mathrm{d}x_{3}$:
        \[
          \mathrm{d}x_{3}=-\tfrac13\,\mathrm{d}x_{1}-\tfrac23\,\mathrm{d}x_{2}
          \;\Longrightarrow\;
          \pd{x_{3}}{x_{1}}\bigg|_{U=c}=-\tfrac13,\qquad
          \pd{x_{3}}{x_{2}}\bigg|_{U=c}=-\tfrac23 .
        \]
        Unsigned MRSes: $1/3$ (between $x_{3}$ and $x_{1}$) and $2/3$
        (between $x_{3}$ and $x_{2}$).
\end{solutions}

\section{Additional Problems}

\begin{problems}
  \item Compute $\mathrm{d}z$ for $z=3x^{2}+2xy-5y^{2}$.
  \item Find $\dd{z}{t}$ for $z=x^{2}+2y$ with $x=3t+1$, $y=t^{2}-1$.
  \item For $U(x,y)=x^{2}+y^{2}=c$, find $\dd{y}{x}$ and the unsigned MRS.
  \item Compute $\pd{z}{s}$ and $\pd{z}{t}$ for $z=x^{2}+y^{2}$ with
        $x=s+t$, $y=s-t$.
  \item Compute $\mathrm{d}z$ for $z=x^{3}\sqrt{y}$.
  \item Find $\dd{z}{t}$ for $z=\ln(x^{2}+y^{2})$ with $x=\mathrm{e}^{t}$,
        $y=\mathrm{e}^{-t}$.
  \item For the isoquant $Q(K,L)=K^{0.3}L^{0.7}=Q_{0}$, find $\dd{L}{K}$
        and the unsigned MRTS.
  \item Compute $\mathrm{d}z$ for $z=\mathrm{e}^{xy}$.
  \item Compute $\pd{z}{s}$ and $\pd{z}{t}$ for $z=xy^{2}$ with
        $x=s^{2}+t$, $y=s-t^{2}$.
  \item For $U(x,y)=\sqrt{xy}=c$, find $\dd{y}{x}$ and the unsigned MRS.
  \item Find $\dd{z}{t}$ for $z=\mathrm{e}^{x+y}$ with $x=\sin t$,
        $y=\cos t$.
  \item Compute $\mathrm{d}z$ for $z=\dfrac{x}{x+y}$.
  \item For $U(x,y)=x+\ln y=c$, find $\dd{y}{x}$ and the unsigned MRS.
  \item Find $\dd{z}{t}$ for $z=x^{2}y-y^{3}$ with $x=t^{2}$, $y=2t$.
  \item Use the total differential of $W(P,Q)=PQ$ (total revenue) to
        approximate $\mathrm{d}W$ when $\mathrm{d}P=0.01P$ and
        $\mathrm{d}Q=-0.005Q$. Interpret.
\end{problems}

In problems 16--18, approximate the change $\Delta f$ by the total
differential $\mathrm{d}f$ at the point $P$, then compute the exact change
and compare.

\begin{problems}[resume]
  \item $Q(x,y,z)=4\sqrt{x^{2}+yz}$ at $P=(3,2,8)$, with
        $\mathrm{d}x=0.1$, $\mathrm{d}y=-0.2$, $\mathrm{d}z=0.3$.
  \item $Y(x,y,z)=2\sqrt{xy+z^{2}}$ at $P=(2,8,3)$, with
        $\mathrm{d}x=0.05$, $\mathrm{d}y=0.1$, $\mathrm{d}z=-0.05$.
  \item $f(x,y)=\ln(x^{2}y)-xy$ at $P=(1,2)$, with
        $\mathrm{d}x=\mathrm{d}y=0.1$; explain why the approximation is
        poorer here.
\end{problems}

\section{Additional Solutions}

\begin{solutions}
  \item $f_{x}=6x+2y$, $f_{y}=2x-10y$.
        \[
          \mathrm{d}z \;=\; (6x+2y)\,\mathrm{d}x+(2x-10y)\,\mathrm{d}y .
        \]

  \item $f_{x}=2x$, $f_{y}=2$; $\dd{x}{t}=3$, $\dd{y}{t}=2t$.
        \[
          \dd{z}{t} \;=\; 2x\cdot 3+2\cdot 2t \;=\; 6(3t+1)+4t
            \;=\; 22t+6 .
        \]
        \textbf{Check.} $z=(3t+1)^{2}+2(t^{2}-1)=9t^{2}+6t+1+2t^{2}-2
        =11t^{2}+6t-1$, $\dd{z}{t}=22t+6$. \checkmark

  \item $U_{x}=2x$, $U_{y}=2y$. Signed slope: $\dd{y}{x}=-x/y$.
        \textbf{MRS} $=x/y$ (unsigned).

  \item $f_{x}=2x$, $f_{y}=2y$; $x_{s}=1$, $x_{t}=1$, $y_{s}=1$,
        $y_{t}=-1$.
        \[
          \pd{z}{s}=2x+2y=2(s+t)+2(s-t)=4s .
        \]
        \[
          \pd{z}{t}=2x-2y=2(s+t)-2(s-t)=4t .
        \]
        \textbf{Check.} $z=(s+t)^{2}+(s-t)^{2}=2(s^{2}+t^{2})$,
        $\pd{z}{s}=4s$, $\pd{z}{t}=4t$. \checkmark

  \item $f_{x}=3x^{2}\sqrt{y}$, $f_{y}=\dfrac{x^{3}}{2\sqrt{y}}$.
        \[
          \mathrm{d}z \;=\; 3x^{2}\sqrt{y}\,\mathrm{d}x
            +\frac{x^{3}}{2\sqrt{y}}\,\mathrm{d}y .
        \]

  \item $f_{x}=2x/(x^{2}+y^{2})$, $f_{y}=2y/(x^{2}+y^{2})$;
        $\dd{x}{t}=\mathrm{e}^{t}$, $\dd{y}{t}=-\mathrm{e}^{-t}$. With
        $x=\mathrm{e}^{t}$, $y=\mathrm{e}^{-t}$,
        $x^{2}+y^{2}=\mathrm{e}^{2t}+\mathrm{e}^{-2t}$.
        \[
          \dd{z}{t} \;=\;
            \frac{2\mathrm{e}^{t}\cdot\mathrm{e}^{t}
              +2\mathrm{e}^{-t}\cdot(-\mathrm{e}^{-t})}
              {\mathrm{e}^{2t}+\mathrm{e}^{-2t}}
          \;=\; \frac{2\mathrm{e}^{2t}-2\mathrm{e}^{-2t}}
              {\mathrm{e}^{2t}+\mathrm{e}^{-2t}}
          \;=\; \frac{2(\mathrm{e}^{2t}-\mathrm{e}^{-2t})}
              {\mathrm{e}^{2t}+\mathrm{e}^{-2t}}
          =2\tanh(2t) .
        \]

  \item $Q_{K}=0.3K^{-0.7}L^{0.7}$, $Q_{L}=0.7K^{0.3}L^{-0.3}$.
        \[
          \frac{Q_{K}}{Q_{L}}=\frac{0.3}{0.7}\cdot\frac{L}{K}
            =\frac{3}{7}\cdot\frac{L}{K} .
        \]
        $\dd{L}{K}=-(3/7)(L/K)$. \textbf{MRTS} $=(3/7)(L/K)$.

  \item $f_{x}=y\mathrm{e}^{xy}$, $f_{y}=x\mathrm{e}^{xy}$.
        \[
          \mathrm{d}z \;=\; \mathrm{e}^{xy}(y\,\mathrm{d}x+x\,\mathrm{d}y) .
        \]

  \item $f_{x}=y^{2}$, $f_{y}=2xy$; $x_{s}=2s$, $x_{t}=1$, $y_{s}=1$,
        $y_{t}=-2t$.
        \[
          \pd{z}{s}=y^{2}\cdot 2s+2xy\cdot 1=2sy^{2}+2xy .
        \]
        \[
          \pd{z}{t}=y^{2}\cdot 1+2xy\cdot(-2t)=y^{2}-4txy .
        \]

  \item $U_{x}=\tfrac{1}{2}\sqrt{y/x}$, $U_{y}=\tfrac{1}{2}\sqrt{x/y}$.
        \[
          \frac{U_{x}}{U_{y}}=\sqrt{\frac{y/x}{x/y}}=\frac{y}{x} .
        \]
        $\dd{y}{x}=-y/x$. \textbf{MRS} $=y/x$.

  \item $f_{x}=f_{y}=\mathrm{e}^{x+y}$; $\dd{x}{t}=\cos t$,
        $\dd{y}{t}=-\sin t$.
        \[
          \dd{z}{t}=\mathrm{e}^{x+y}(\cos t-\sin t)
            =\mathrm{e}^{\sin t+\cos t}(\cos t-\sin t) .
        \]

  \item Quotient rule with $u=x$, $v=x+y$:
        \[
          f_{x}=\frac{(x+y)-x}{(x+y)^{2}}=\frac{y}{(x+y)^{2}},\qquad
          f_{y}=\frac{-x}{(x+y)^{2}} .
        \]
        \[
          \mathrm{d}z=\frac{y\,\mathrm{d}x-x\,\mathrm{d}y}{(x+y)^{2}} .
        \]

  \item $U_{x}=1$, $U_{y}=1/y$. $U_{x}/U_{y}=y$. Signed slope
        $\dd{y}{x}=-y$. \textbf{MRS} $=y$.

  \item $f_{x}=2xy$, $f_{y}=x^{2}-3y^{2}$; $\dd{x}{t}=2t$,
        $\dd{y}{t}=2$.
        \[
          \dd{z}{t}=2xy\cdot 2t+(x^{2}-3y^{2})\cdot 2
            =4t\cdot t^{2}\cdot 2t+2(t^{4}-3\cdot 4t^{2})
            =8t^{4}+2t^{4}-24t^{2}=10t^{4}-24t^{2} .
        \]
        \textbf{Check.} $z=t^{4}\cdot 2t-8t^{3}=2t^{5}-8t^{3}$,
        $\dd{z}{t}=10t^{4}-24t^{2}$. \checkmark

  \item $W_{P}=Q$, $W_{Q}=P$.
        \[
          \mathrm{d}W=Q\,\mathrm{d}P+P\,\mathrm{d}Q
            =W\!\left(\frac{\mathrm{d}P}{P}+\frac{\mathrm{d}Q}{Q}\right) .
        \]
        So $\mathrm{d}W/W=0.01+(-0.005)=0.005$, i.e.\ revenue rises by
        about $0.5\%$. (Log-linear: revenue percent change is the sum of
        price and quantity percent changes.)

  \item At $P$: $x^{2}+yz=9+16=25$, so $Q(P)=4\cdot5=20$. Partials, by
        the chain rule with $w=x^{2}+yz$:
        \[
          Q_{x}=\frac{4\cdot2x}{2\sqrt w}=\frac{4x}{\sqrt w}=\frac{12}{5},\quad
          Q_{y}=\frac{2z}{\sqrt w}=\frac{16}{5},\quad
          Q_{z}=\frac{2y}{\sqrt w}=\frac45 .
        \]
        \[
          \mathrm{d}Q=2.4(0.1)+3.2(-0.2)+0.8(0.3)=0.24-0.64+0.24=-0.16 .
        \]
        Exact: $Q(3.1,1.8,8.3)=4\sqrt{9.61+14.94}=4\sqrt{24.55}
        \approx19.819183$, so $\Delta Q\approx-0.180817$. The differential
        gets the sign and the size about right; the error $0.021$ comes
        from the curvature of the root.
  \item At $P$: $xy+z^{2}=16+9=25$, $Y(P)=10$. With $w=xy+z^{2}$:
        \[
          Y_{x}=\frac{y}{\sqrt w}=\frac85,\quad Y_{y}=\frac{x}{\sqrt w}=\frac25,
          \quad Y_{z}=\frac{2z}{\sqrt w}=\frac65 .
        \]
        \[
          \mathrm{d}Y=1.6(0.05)+0.4(0.1)+1.2(-0.05)=0.08+0.04-0.06=0.06 .
        \]
        Exact: $Y(2.05,8.1,2.95)=2\sqrt{16.605+8.7025}=2\sqrt{25.3075}
        \approx10.061312$, so $\Delta Y\approx0.061312$: the error is about
        $0.0013$, two percent of the change.
  \item $f(P)=\ln2-2\approx-1.306853$. $f=2\ln x+\ln y-xy$, so
        $f_{x}=\dfrac2x-y=0$ and $f_{y}=\dfrac1y-x=-\tfrac12$ at $P$.
        \[
          \mathrm{d}f=0\cdot0.1-\tfrac12\cdot0.1=-0.05 .
        \]
        Exact: $f(1.1,2.1)=\ln(1.21\cdot2.1)-2.31=\ln2.541-2.31\approx
        -1.377442$, so $\Delta f\approx-0.070589$. The differential misses
        by $0.02$, about $30\%$ of the change. At $P$ the slope in the
        $x$-direction is $0$, so the linear term sees nothing of the
        $x$-step, while the curvature, with $f_{xx}=-2/x^{2}=-2$,
        $f_{xy}=-1$ and $f_{yy}=-1/y^{2}=-\tfrac14$, contributes about
        $\tfrac12\bigl(-2(0.01)+2(-1)(0.01)-\tfrac14(0.01)\bigr)\approx-0.021$:
        indeed $-0.05-0.021\approx-0.071$.
\end{solutions}

\chapter{Directional Derivatives and Vector-Valued Functions}

The partial derivatives $f_{x}$ and $f_{y}$ measure the response of $f$
when one input moves and the others stand still. Real adjustments are
rarely so tidy: a firm re-balances capital and labour together, a consumer
slides along a budget line, a hiker walks north-east rather than due
north. The \textbf{directional derivative} measures the rate of change of
$f$ per unit distance travelled in any prescribed direction, and a single
dot product with the gradient computes it. The gradient thereby earns its
geometric role: it points where $f$ climbs fastest, its length is that
fastest rate, and it stands perpendicular to the level curves, to the
indifference curves and isoquants of the previous chapters.

The second half of the chapter lets the function itself be vector-valued.
A map $\mathbf{f}:\R^{n}\to\R^{m}$ (a production plan sending inputs to
(output, cost), a coordinate change, a demand system) is differentiated
by collecting the partial derivatives of all components into the
\textbf{Jacobian matrix}, the correct ``derivative as linear
approximation'' for vector-valued functions. In this language the chain
rule becomes matrix multiplication. Determinants, invertibility and the
Implicit Function Theorem built on top of this matrix are the subject of
the next chapter.

\section{Core Rules}

\begin{rulebox}{Directional Derivative}
Let $f:\R^{n}\to\R$ be differentiable at $\mathbf{x}_{0}$ and let
$\mathbf{u}\in\R^{n}$ be a \textbf{unit} vector, $|\mathbf{u}|=1$. The
directional derivative of $f$ at $\mathbf{x}_{0}$ in the direction
$\mathbf{u}$ is
\[
  D_{\mathbf{u}}f(\mathbf{x}_{0})
  \;=\; \lim_{h\to 0}\frac{f(\mathbf{x}_{0}+h\mathbf{u})
        -f(\mathbf{x}_{0})}{h}
  \;=\; \grad f(\mathbf{x}_{0})\cdot\mathbf{u} .
\]
Two variables: $D_{\mathbf{u}}f=f_{x}u_{1}+f_{y}u_{2}$. The partials are
the special cases $\mathbf{u}=(1,0)$ and $\mathbf{u}=(0,1)$. The unit
length is what makes $D_{\mathbf{u}}f$ a rate \emph{per unit distance}.
\end{rulebox}

\begin{rulebox}{Normalising the Direction}
If a direction is given by a nonzero vector $\mathbf{v}$ that is not of
unit length, normalise first:
\[
  \mathbf{u} \;=\; \frac{\mathbf{v}}{|\mathbf{v}|}
  \;=\; \frac{\mathbf{v}}{\sqrt{v_{1}^{2}+\cdots+v_{n}^{2}}} .
\]
``The direction from $P$ to $Q$'' means $\mathbf{v}=Q-P$, then normalise.
Forgetting this step is a frequent error in directional-derivative
problems.
\end{rulebox}

\begin{rulebox}{Steepest Ascent and Descent}
Since $D_{\mathbf{u}}f=\grad f\cdot\mathbf{u}=|\grad f|\cos\theta$, where
$\theta$ is the angle between $\mathbf{u}$ and $\grad f$, at any point
with $\grad f(\mathbf{x}_{0})\neq\mathbf{0}$:
\begin{itemize}[leftmargin=1.6em,itemsep=0.3ex]
  \item $D_{\mathbf{u}}f$ is \emph{maximal} for
        $\mathbf{u}=\grad f/|\grad f|$, with maximum value $|\grad f|$
        (steepest ascent);
  \item $D_{\mathbf{u}}f$ is \emph{minimal} for
        $\mathbf{u}=-\grad f/|\grad f|$, with value $-|\grad f|$
        (steepest descent);
  \item $D_{\mathbf{u}}f=0$ exactly when $\mathbf{u}\perp\grad f$.
\end{itemize}
For output $Q(K,L)$, the gradient $(Q_{K},Q_{L})$ is the input adjustment
that raises output fastest per unit of adjustment, and $|\grad Q|$ is that
maximal marginal gain.

\textbf{A fixed amount of adjustment.} If the change vector
$(\mathrm{d}K,\mathrm{d}L)$ must have a given length $I$, output rises
fastest when it points along the gradient:
\[
  \mathrm{d}K:\mathrm{d}L=Q_{K}:Q_{L},\qquad
  (\mathrm{d}K,\mathrm{d}L)=\frac{I}{|\grad Q|}\,(Q_{K},Q_{L}),\qquad
  \mathrm{d}Q\approx I\,|\grad Q| .
\]
A \emph{budget} $\mathrm{d}K+\mathrm{d}L=I$ with $\mathrm{d}K,\mathrm{d}L\ge0$ is
a different constraint: the linear estimate $Q_{K}\mathrm{d}K+Q_{L}\mathrm{d}L$ is then largest at a
corner, with everything spent on the input with the larger marginal
product.
\end{rulebox}

\begin{rulebox}{Gradient and Level Curves}
Along the level set $f=c$ the function does not change, so every tangent
direction $\mathbf{u}$ of the level set satisfies
$D_{\mathbf{u}}f=\grad f\cdot\mathbf{u}=0$:
\[
  \grad f \;\perp\; \text{level curve (surface) through the point.}
\]
In two variables (with $f_{y}\neq0$) the tangent direction is proportional
to $(f_{y},-f_{x})$, consistent with the implicit slope
$\dd{y}{x}=-f_{x}/f_{y}$ of the previous chapters. Economically: the
marginal-utility vector $(U_{x},U_{y})$ is normal to the indifference
curve, and $(Q_{K},Q_{L})$ is normal to the isoquant.
\end{rulebox}

\begin{rulebox}{Tangent lines: to a level curve, and to the graph}
\textbf{Tangent to the level curve} $f(x,y)=c$ through
$\mathbf{p}=(x_{0},y_{0})$, where $\grad f(\mathbf{p})\neq\mathbf{0}$. The
gradient is a normal vector, so the tangent line is
\[
  \grad f(\mathbf{p})\cdot(\mathbf{x}-\mathbf{p})=0,
  \qquad\text{i.e.}\qquad
  f_{x}(\mathbf{p})(x-x_{0})+f_{y}(\mathbf{p})(y-y_{0})=0
\]
(normal form), or $\mathbf{x}=\mathbf{p}+t\mathbf{h}$ with
$\mathbf{h}\perp\grad f(\mathbf{p})$, for instance
$\mathbf{h}=(f_{y},-f_{x})$ (parametric form). Both describe the same
line in the plane.

\textbf{Tangent to the graph in a direction.} The vertical plane through
$P=(x_{0},y_{0})$ along the unit vector $\mathbf{u}$ cuts the surface
$z=f(x,y)$ in the curve $t\mapsto f(P+t\mathbf{u})$, whose slope at $t=0$
is $D_{\mathbf{u}}f(P)$. That is the geometric meaning of the directional
derivative: the slope of the surface in the direction $\mathbf{u}$. The
tangent line to the surface in that direction is a line in $\R^{3}$:
\[
  (x,y,z)=\bigl(x_{0},\,y_{0},\,f(P)\bigr)
          +t\,\bigl(u_{1},\,u_{2},\,D_{\mathbf{u}}f(P)\bigr),\qquad t\in\R .
\]
\begin{itemize}[leftmargin=1.6em,itemsep=0.3ex]
  \item $x$-direction ($\mathbf{u}=\mathbf{e}_{1}$):
        $\bigl(x_{0}+t,\;y_{0},\;f(P)+t\,f_{x}(P)\bigr)$;
        $y$-direction ($\mathbf{u}=\mathbf{e}_{2}$):
        $\bigl(x_{0},\;y_{0}+t,\;f(P)+t\,f_{y}(P)\bigr)$.
  \item For a direction vector $\mathbf{h}$ that is not normalised,
        $\bigl(P+t\mathbf{h},\;f(P)+t\,\grad f(P)\cdot\mathbf{h}\bigr)$ is
        the \emph{same} line, run through at a different speed. Only the
        slope, rise per unit of horizontal distance, needs the unit vector.
  \item All these lines lie in the tangent plane of \chref{derivatives}{10}.
\end{itemize}
\end{rulebox}

\begin{rulebox}{Second Directional Derivative}
For a $C^{2}$ function $f$, differentiating once more in the same unit
direction $\mathbf{u}$ gives
\[
  D^{2}_{\mathbf{u}}f
  \;=\; \mathbf{u}^{T}\Hess_{f}\,\mathbf{u}
  \;=\; \sum_{i,j} f_{x_{i}x_{j}}\,u_{i}u_{j} ;
  \qquad\text{two variables:}\quad
  D^{2}_{\mathbf{u}}f
  \;=\; f_{xx}u_{1}^{2}+2f_{xy}u_{1}u_{2}+f_{yy}u_{2}^{2} .
\]
It is the second derivative at $t=0$ of the profile
$t\mapsto f(\mathbf{x}_{0}+t\mathbf{u})$: the curvature of $f$ along the
line through $\mathbf{x}_{0}$ in direction $\mathbf{u}$. Definiteness of
the Hessian (\chref{derivatives}{11}) is exactly the statement that
$D^{2}_{\mathbf{u}}f$ has one sign for \emph{every} direction
$\mathbf{u}$.
\end{rulebox}

\begin{rulebox}{Curvature in a direction: computing it, and what it proves}
\textbf{The profile method.} Without the Hessian: set
$\varphi(t)=f(\mathbf{x}_{0}+t\mathbf{u})$, a function of one variable,
and differentiate twice. Then $\varphi'(0)=D_{\mathbf{u}}f(\mathbf{x}_{0})$
and $\varphi''(0)=D^{2}_{\mathbf{u}}f(\mathbf{x}_{0})$. Doing it both ways
is a built-in check.

\textbf{Non-unit directions.} For any $\mathbf{h}\neq\mathbf{0}$, with
$\mathbf{u}=\mathbf{h}/|\mathbf{h}|$,
\[
  \mathbf{h}^{T}H\mathbf{h}=|\mathbf{h}|^{2}\,D^{2}_{\mathbf{u}}f,
  \qquad
  \grad f\cdot\mathbf{h}=|\mathbf{h}|\,D_{\mathbf{u}}f .
\]
For the \emph{sign} of the curvature, normalising is unnecessary; for the
value of $D_{\mathbf{u}}f$ or $D^{2}_{\mathbf{u}}f$, divide by $|\mathbf{h}|$
or $|\mathbf{h}|^{2}$.

\textbf{Directional verdicts.} At $\mathbf{x}_{0}$, $f$ is
\emph{convex in the direction $\mathbf{u}$} if $D^{2}_{\mathbf{u}}f>0$,
\emph{concave in that direction} if $D^{2}_{\mathbf{u}}f<0$; the profile
$\varphi$ bends up or down there. A value $0$ gives no verdict.

\textbf{Refuting global convexity.} A convex $f$ has a PSD Hessian at
\emph{every} point, so $D^{2}_{\mathbf{u}}f\ge0$ everywhere in every
direction; a concave $f$ has $D^{2}_{\mathbf{u}}f\le0$ everywhere. Hence one
point and direction with $D^{2}_{\mathbf{u}}f<0$ shows that $f$ is not
convex, and one with $D^{2}_{\mathbf{v}}f>0$ (for example at a strict local
minimum with PD Hessian) shows that $f$ is not concave. Both together prove
that $f$ is \textbf{neither}.

\textbf{Where this goes next.} The box ``Curvature at a point and in a
direction'' in \chref{optimization}{3} turns these directional verdicts into
statements about convexity near a point and on a whole domain, including
the case of a semidefinite, singular Hessian where no local conclusion is
possible.
\end{rulebox}

\begin{rulebox}{Vector-Valued Functions and the Jacobian Matrix}
For $\mathbf{f}:\R^{n}\to\R^{m}$ with components
$\mathbf{f}=(f_{1},\ldots,f_{m})$, the \textbf{Jacobian matrix} at
$\mathbf{x}$ is the $m\times n$ matrix
\[
  J_{\mathbf{f}}(\mathbf{x}) \;=\;
  \begin{pmatrix}
    \partial f_{1}/\partial x_{1} & \cdots & \partial f_{1}/\partial x_{n}\\
    \vdots & \ddots & \vdots\\
    \partial f_{m}/\partial x_{1} & \cdots & \partial f_{m}/\partial x_{n}
  \end{pmatrix}:
  \quad\text{row $i$ is the gradient of $f_{i}$.}
\]
It is the best linear approximation of $\mathbf{f}$ near $\mathbf{x}_{0}$:
\[
  \mathbf{f}(\mathbf{x}_{0}+\mathbf{h})
  \;\approx\; \mathbf{f}(\mathbf{x}_{0})
  +J_{\mathbf{f}}(\mathbf{x}_{0})\,\mathbf{h} ,
\]
the matrix version of $f(x+h)\approx f(x)+f'(x)h$. Special cases: $m=1$
gives the gradient as a $1\times n$ row; $n=1$ gives the velocity
(tangent) column vector $\mathbf{r}'(t)$ of a curve.
\end{rulebox}

\begin{rulebox}{Differentiability of vector-valued maps}
$\mathbf{f}\colon D\to\R^{m}$, $D\subseteq\R^{n}$ open, is
\textbf{differentiable} at $\mathbf{x}_{0}$ if there is an $m\times n$
matrix $J$ with
\[
  \mathbf{f}(\mathbf{x}_{0}+\mathbf{h})=\mathbf{f}(\mathbf{x}_{0})
  +J\mathbf{h}+\mathbf{r}(\mathbf{h}),
  \qquad
  \frac{\|\mathbf{r}(\mathbf{h})\|}{\|\mathbf{h}\|}\to0\ \text{ as }
  \mathbf{h}\to\mathbf{0} .
\]
The ``$\approx$'' of the Jacobian box is made exact by the remainder.
\begin{itemize}[leftmargin=1.6em,itemsep=0.3ex]
  \item The matrix $J$ is then unique and equals the Jacobian
        $J_{\mathbf{f}}(\mathbf{x}_{0})$ of partial derivatives. It is
        called the \textbf{derivative} (or total derivative,
        \de{Ableitung, Jacobi-Matrix}) of $\mathbf{f}$ at $\mathbf{x}_{0}$.
  \item $\mathbf{f}$ is differentiable exactly when every component
        $f_{i}$ is: the remainder vector tends to $0$ faster than
        $\|\mathbf{h}\|$ exactly when each of its entries does.
  \item Existence of all partials is \emph{necessary} but not sufficient
        (the function $xy/(x^{2}+y^{2})$ of \chref{derivatives}{10} is a single component that fails).
        Continuous partials on an open set are sufficient: a $C^{1}$ map is
        differentiable.
  \item Geometry: for $n=m=1$ the linear part is the tangent line; for
        $n=2$, $m=1$ it is the tangent plane; for a curve ($n=1$) it is the
        tangent vector $\mathbf{r}'(t)$.
\end{itemize}
\end{rulebox}

\begin{rulebox}{The shape of a Jacobian, and which compositions exist}
$J_{\mathbf{f}}$ has one \textbf{row per output} and one \textbf{column per
input}: (number of outputs)$\times$(number of inputs).
\begin{itemize}[leftmargin=1.6em,itemsep=0.3ex]
  \item A curve $\R\to\R^{m}$ gives an $m\times1$ column; a scalar function
        $\R^{n}\to\R$ a $1\times n$ row; $\R^{3}\to\R^{2}$ a $2\times3$
        (wide) matrix; $\R^{2}\to\R^{4}$ a $4\times2$ (tall) one.
  \item \textbf{Linear and affine maps} have constant Jacobians:
        $J(A\mathbf{x}+\mathbf{b})=A$ for an $m\times n$ matrix $A$;
        $J(c\,\mathbf{x}-d\,\mathbf{a})=c\,I_{n}$; the identity map has
        $J=I_{n}$; $J(\mathbf{a}^{T}\mathbf{x})=\mathbf{a}^{T}$.
  \item \textbf{Patterns in $n$ dimensions.} If component $k$ depends on
        $x_{k}$ alone, $f_{k}=g_{k}(x_{k})$, the Jacobian is diagonal,
        $\operatorname{diag}\bigl(g_{1}'(x_{1}),\ldots,g_{n}'(x_{n})\bigr)$.
\end{itemize}
\textbf{Which composition exists.} $\mathbf{g}\circ\mathbf{f}$ means
``first $\mathbf{f}$, then $\mathbf{g}$''; it is defined when the number of
outputs of $\mathbf{f}$ equals the number of inputs of $\mathbf{g}$. With
$\mathbf{f}\colon\R^{n}\to\R^{m}$ and $\mathbf{g}\colon\R^{m}\to\R^{p}$ the
chain rule multiplies $(p\times m)(m\times n)$; the inner sizes must
match, which is the same condition. Check it before computing anything:
often only one of $\mathbf{g}\circ\mathbf{f}$ and $\mathbf{f}\circ\mathbf{g}$
makes sense.
\end{rulebox}

\begin{rulebox}{Chain Rule in Matrix Form}
If $\mathbf{f}:\R^{n}\to\R^{m}$ and $\mathbf{g}:\R^{m}\to\R^{p}$ are
differentiable, then $\mathbf{h}=\mathbf{g}\circ\mathbf{f}$ is
differentiable and
\[
  J_{\mathbf{h}}(\mathbf{x})
  \;=\; J_{\mathbf{g}}\bigl(\mathbf{f}(\mathbf{x})\bigr)\,
        J_{\mathbf{f}}(\mathbf{x}),
  \qquad
  (p\times n) \;=\; (p\times m)(m\times n) .
\]
Evaluate the \emph{outer} Jacobian at the \emph{inner value}
$\mathbf{f}(\mathbf{x})$; the order of the factors is fixed by the matrix
shapes. Every chain rule of \chref{derivatives}{12} is the case $p=1$: for
$z=f(\mathbf{r}(t))$,
\[
  \dd{z}{t} \;=\; \grad f\bigl(\mathbf{r}(t)\bigr)\cdot\mathbf{r}'(t),
\]
a $1\times m$ row times an $m\times 1$ column.
\end{rulebox}

\section{Worked Example}

\begin{worked}{For $f(x,y)=x^{2}+3xy-y^{2}$ at $P=(1,2)$: the derivative in the direction $(3,4)$, the steepest ascent, a direction of no change; then the Jacobian of $F(x,y)=(x^{2}-y,\;xy)$}
\wstep{Step 1. The gradient first --- every question below is an inner product with it.}
\[
  f_{x}=2x+3y,\qquad f_{y}=3x-2y,
\]
\[
  \grad f(1,2)=\bigl(2+6,\;3-4\bigr)=(8,-1).
\]

\wstep{Step 2. Normalise the direction, which the formula requires.}
\[
  \|(3,4)\|=5,
  \qquad
  \mathbf u=\Bigl(\tfrac35,\tfrac45\Bigr).
\]
\[
  D_{\mathbf u}f(P)=\grad f\cdot\mathbf u=8\cdot\tfrac35+(-1)\cdot\tfrac45
  =\frac{24-4}{5}=\boxed{4}.
\]
\wtrap{Without normalising you would get $\grad f\cdot(3,4)=20$, five times
too large. The directional derivative is a rate \emph{per unit distance
travelled}, so doubling the length of the direction vector must not double
the answer. If the problem hands you a non-unit vector --- and it usually
does --- divide by its length first.}

\wstep{Step 3. Steepest ascent and descent.}
The steepest increase is along the gradient itself:
\[
  \text{direction }\frac{(8,-1)}{\sqrt{65}},
  \qquad
  \text{rate }\|\grad f\|=\sqrt{65}\approx8.06 .
\]
The steepest \emph{decrease} is the opposite direction, at rate
$-\sqrt{65}$. Every other direction lies strictly between, since
\[
  D_{\mathbf u}f=\grad f\cdot\mathbf u=\|\grad f\|\cos\theta ,
\]
and $\cos\theta\in[-1,1]$. Our $4$ sits comfortably inside
$[-8.06,\,8.06]$. \checkmark
\wtrap{The value $8.06$ is the \emph{rate}; the direction is the normalised
gradient $(0.992,-0.124)$, not $(8,-1)$. Reporting the unnormalised gradient
as ``the direction'' conflates a direction with a magnitude.}

\wstep{Step 4. A direction of no change --- and what it is tangent to.}
$D_{\mathbf u}f=0$ requires $\mathbf u\perp\grad f$. Rotating $(8,-1)$ by a
right angle gives $(1,8)$, so
\[
  \mathbf u=\frac{(1,8)}{\sqrt{65}}
  \quad\Longrightarrow\quad
  D_{\mathbf u}f=\frac{8\cdot1+(-1)\cdot8}{\sqrt{65}}=0 . \checkmark
\]
Moving in that direction keeps $f$ constant to first order, so $(1,8)$ is
tangent to the level curve of $f$ through $P$ --- which is the geometric
content of ``the gradient is perpendicular to the level curve''.

\wstep{Step 5. A vector-valued function needs a matrix of derivatives.}
$F(x,y)=(x^{2}-y,\ xy)$ has two components, each with two partials. Collect
them with \textbf{one row per component function}:
\[
  J_{F}(x,y)=
  \begin{pmatrix}
    \pd{F_{1}}{x} & \pd{F_{1}}{y}\\[0.6ex]
    \pd{F_{2}}{x} & \pd{F_{2}}{y}
  \end{pmatrix}
  =\begin{pmatrix}2x&-1\\ y&x\end{pmatrix}.
\]
\wtrap{Rows are components, columns are variables. Transposing this is the
standard error, and it matters as soon as the Jacobian is multiplied by
anything. A check: for a scalar function ($m=1$) the Jacobian is a single
\emph{row}, and the gradient is conventionally a column --- so
$J=\grad^{\!\top}$.}

\wstep{Step 6. Evaluate, and read the determinant.}
\[
  J_{F}(1,2)=\begin{pmatrix}2&-1\\2&1\end{pmatrix},
  \qquad
  \det J_{F}(1,2)=2+2=4\neq0 .
\]
A non-zero Jacobian determinant says the map is locally invertible near
$(1,2)$: the system
\[
  x^{2}-y=u,\qquad xy=v
\]
can be solved for $(x,y)$ in terms of $(u,v)$ near that point. That is the
Implicit Function Theorem of \chref{derivatives}{14} in its multivariate form, and
$\det J$ is the condition it turns on.

\wstep{Step 7. Why the Jacobian is the right generalisation.}
\begin{center}\footnotesize
\begin{tabular}{lll}
  \hline
  $f:\R\to\R$ & derivative $f'$ & a number\\
  $f:\R^{n}\to\R$ & gradient $\grad f$ & a vector\\
  $F:\R^{n}\to\R^{m}$ & Jacobian $J_{F}$ & an $m\times n$ matrix\\
  \hline
\end{tabular}
\end{center}
In each row the object is the best \emph{linear} approximation at a point,
and the chain rule is matrix multiplication:
$J_{G\circ F}=J_{G}\cdot J_{F}$, with the order fixed by the order of the
composition. The one-variable chain rule $g'(f(x))f'(x)$ is the
$1\times1$ case.
\end{worked}

\section{Practice Problems}

\begin{problems}
  \item Find $D_{\mathbf{u}}f(1,1)$ for $f(x,y)=x^{2}y$ in the direction
        $\mathbf{u}=(1/\sqrt{2},\,1/\sqrt{2})$.
  \item For $f(x,y)=\ln(x^{2}+y^{2})$ at $(1,1)$, find the directional
        derivative toward $\mathbf{v}=(2,1)$.
  \item For $f(x,y)=x^{3}-3xy^{2}$ at $(2,1)$, find the directional
        derivative in the direction of the vector from $(2,1)$ to
        $(5,5)$.
  \item For $f(x,y,z)=xyz$ at $(1,2,3)$, find $D_{\mathbf{u}}f$ where
        $\mathbf{u}=\tfrac{1}{3}(1,2,2)$.
  \item For $f(x,y)=\mathrm{e}^{x+y}$ at $(0,0)$, find $D_{\mathbf{u}}f$
        where $\mathbf{u}$ makes a $60^{\circ}$ angle with the positive
        $x$-axis.
  \item For $f(x,y)=\sqrt{x^{2}+y^{2}}$ at $(3,4)$, find the directional
        derivative in the direction $(1,0)$.
  \item A consumer with utility $U(x,y)=xy$ currently consumes
        $(x,y)=(2,3)$. Her basket shifts in the direction
        $\mathbf{v}=(3,-4)$, more $x$, less $y$. Find the rate of
        change of utility per unit distance moved and interpret the sign.
  \item For $f(x,y)=x^{2}+xy+y^{2}$ at $(1,2)$, find the direction of
        steepest ascent and the maximum rate of increase.
  \item For $f(x,y)=x^{2}+2y^{2}$ at $(1,1)$, find the direction in which
        $f$ decreases fastest and the corresponding rate.
  \item Find all unit vectors $\mathbf{u}$ along which
        $f(x,y)=x^{2}-y^{2}$ has zero rate of change at $(2,1)$.
  \item For $f(x,y,z)=x^{2}+y^{2}-z^{2}$ at $(1,2,2)$, determine the
        maximum rate of increase and the direction in which it is
        attained.
  \item A hiker stands at $(1,1)$ on the hill $z=4-x^{2}-y^{2}$. In which
        direction should she walk to descend most quickly, and how steep
        is that descent?
  \item Output is $Q(K,L)=K^{1/2}L^{1/2}$ with current input mix
        $(K,L)=(4,9)$. Find the rate of change of output in the balanced
        expansion direction $\mathbf{u}=(1,1)/\sqrt{2}$; then find the
        input direction that raises output fastest and the corresponding
        maximal rate.
  \item Show that at every point of a level curve $f(x,y)=c$ the gradient
        $\grad f$ is orthogonal to the tangent direction of the curve.
        Conclude that the marginal-utility vector $(U_{x},U_{y})$ is
        normal to the indifference curve.
  \item For $f(x,y)=x^{2}y$ at $(1,2)$ with $\mathbf{u}=(3/5,\,4/5)$,
        compute both $D_{\mathbf{u}}f$ and the second directional
        derivative
        $D^{2}_{\mathbf{u}}f=\mathbf{u}^{T}\Hess_{f}\,\mathbf{u}$.
  \item For $f(x,y)=x^{3}+y^{3}-3xy$: show that $(1,1)$ is a critical
        point, then compute $D^{2}_{\mathbf{u}}f(1,1)$ for
        $\mathbf{u}=(1,1)/\sqrt{2}$ and for
        $\mathbf{v}=(1,-1)/\sqrt{2}$. What do the signs suggest?
  \item For $f(x,y)=\mathrm{e}^{x}\cos y$ at $(0,0)$, find all unit
        vectors $\mathbf{u}$ with $D^{2}_{\mathbf{u}}f=0$.
  \item Compute the Jacobian of $\mathbf{f}(x,y)=(x+y,\;x-y)$ and its
        determinant. What is special about the Jacobian of a linear map?
  \item Compute the Jacobian of
        $\mathbf{f}(x,y)=(x^{2}y,\;x+y^{2},\;xy)$, state its shape, and
        evaluate it at $(1,2)$.
  \item Compute the Jacobian of
        $\mathbf{f}(x,y,z)=(x+y,\;y+z,\;x+z)$ and its determinant.
  \item For $\mathbf{f}(x,y)=(x^{2}-y,\;xy)$: compute $J_{\mathbf{f}}$,
        then use the linear approximation
        $\mathbf{f}(\mathbf{x}_{0}+\mathbf{h})\approx
        \mathbf{f}(\mathbf{x}_{0})+J_{\mathbf{f}}(\mathbf{x}_{0})\mathbf{h}$
        at $\mathbf{x}_{0}=(1,1)$ to estimate $\mathbf{f}(1.1,\,0.9)$,
        and compare with the exact value.
  \item Demand and supply are $D(p,q)=100-p+0.5q$ and
        $S(p,q)=-10+2p+q$, and the two equilibrium conditions are
        $G_{1}=D-S=0$ and $G_{2}=p+q-20=0$. Build the Jacobian matrix
        $\partial(G_{1},G_{2})/\partial(p,q)$ and compute its
        determinant.
  \item Let $\mathbf{g}(x,y)=(x+y,\;x-y)$ and
        $\mathbf{F}(u,v)=(u^{2},\;uv)$. Compute the Jacobian of
        $\mathbf{F}\circ\mathbf{g}$ by the matrix chain rule
        $J_{\mathbf{F}\circ\mathbf{g}}
        =J_{\mathbf{F}}(\mathbf{g}(x,y))\,J_{\mathbf{g}}(x,y)$, and
        verify by composing first.
  \item For $f(x,y)=xy^{2}$ and the curve
        $\mathbf{r}(t)=(\mathrm{e}^{t},\,t)$, compute $\dd{z}{t}$ for
        $z=f(\mathbf{r}(t))$ as the matrix product
        $J_{f}(\mathbf{r}(t))\,J_{\mathbf{r}}(t)$ (a $1\times2$ row
        times a $2\times1$ column) and verify directly.
  \item A firm converts inputs into (output, cost) via
        $\mathbf{g}(K,L)=\bigl(K^{1/2}L^{1/2},\;2K+L\bigr)$ and earns
        profit $\pi(Q,C)=10Q-C$. Use the matrix chain rule to find the
        marginal profits $(\pi_{K},\pi_{L})$ at $(K,L)=(1,1)$, and verify
        by substituting first.
\end{problems}

\section{Solutions}

\begin{solutions}
  \item $\grad f=(2xy,\,x^{2})$. At $(1,1)$: $\grad f=(2,1)$, and
        $\mathbf{u}$ is already a unit vector.
        \[
          D_{\mathbf{u}}f=(2)\frac{1}{\sqrt{2}}+(1)\frac{1}{\sqrt{2}}
            =\frac{3}{\sqrt{2}}=\frac{3\sqrt{2}}{2} .
        \]

  \item $\grad f=\Bigl(\dfrac{2x}{x^{2}+y^{2}},\,
        \dfrac{2y}{x^{2}+y^{2}}\Bigr)$. At $(1,1)$: $\grad f=(1,1)$.
        Normalise: $|\mathbf{v}|=\sqrt{4+1}=\sqrt{5}$, so
        $\mathbf{u}=(2,1)/\sqrt{5}$.
        \[
          D_{\mathbf{u}}f=\frac{1\cdot 2+1\cdot 1}{\sqrt{5}}
            =\frac{3}{\sqrt{5}}=\frac{3\sqrt{5}}{5} .
        \]

  \item Direction vector $\mathbf{v}=(5,5)-(2,1)=(3,4)$,
        $|\mathbf{v}|=5$, so $\mathbf{u}=(3/5,\,4/5)$.
        $\grad f=(3x^{2}-3y^{2},\,-6xy)$; at $(2,1)$:
        $\grad f=(12-3,\,-12)=(9,-12)$.
        \[
          D_{\mathbf{u}}f=9\cdot\frac{3}{5}+(-12)\cdot\frac{4}{5}
            =\frac{27}{5}-\frac{48}{5}=-\frac{21}{5} .
        \]

  \item $\grad f=(yz,\,xz,\,xy)$. At $(1,2,3)$: $\grad f=(6,3,2)$. The
        vector $\mathbf{u}=\tfrac13(1,2,2)$ is unit since
        $\sqrt{1+4+4}/3=1$.
        \[
          D_{\mathbf{u}}f=\frac{6\cdot 1+3\cdot 2+2\cdot 2}{3}
            =\frac{6+6+4}{3}=\frac{16}{3} .
        \]

  \item $\mathbf{u}=(\cos 60^{\circ},\,\sin 60^{\circ})
        =(1/2,\,\sqrt{3}/2)$, a unit vector.
        $\grad f=(\mathrm{e}^{x+y},\,\mathrm{e}^{x+y})$; at $(0,0)$:
        $\grad f=(1,1)$.
        \[
          D_{\mathbf{u}}f=\frac{1}{2}+\frac{\sqrt{3}}{2}
            =\frac{1+\sqrt{3}}{2} .
        \]

  \item $\grad f=\Bigl(\dfrac{x}{\sqrt{x^{2}+y^{2}}},\,
        \dfrac{y}{\sqrt{x^{2}+y^{2}}}\Bigr)$. At $(3,4)$:
        $\grad f=(3/5,\,4/5)$. With $\mathbf{u}=(1,0)$ (already unit):
        \[
          D_{\mathbf{u}}f=\frac{3}{5} .
        \]

  \item $\grad U=(y,\,x)$; at $(2,3)$: $\grad U=(3,2)$. Normalise:
        $|\mathbf{v}|=\sqrt{9+16}=5$, $\mathbf{u}=(3/5,\,-4/5)$.
        \[
          D_{\mathbf{u}}U=3\cdot\frac{3}{5}+2\cdot\Bigl(-\frac{4}{5}\Bigr)
            =\frac{9}{5}-\frac{8}{5}=\frac{1}{5} .
        \]
        The rate is positive: the marginal-utility gain from the extra
        $x$ slightly outweighs the loss from the reduced $y$, so utility
        rises, at $1/5$ per unit distance moved.

  \item $\grad f=(2x+y,\,x+2y)$. At $(1,2)$: $\grad f=(4,5)$,
        $|\grad f|=\sqrt{16+25}=\sqrt{41}$.
        \[
          \mathbf{u}=\frac{(4,5)}{\sqrt{41}},\qquad
          \text{maximum rate}=\sqrt{41} .
        \]

  \item $\grad f=(2x,\,4y)$. At $(1,1)$: $\grad f=(2,4)$,
        $|\grad f|=\sqrt{4+16}=2\sqrt{5}$. Steepest descent:
        \[
          \mathbf{u}=-\frac{(2,4)}{2\sqrt{5}}=-\frac{(1,2)}{\sqrt{5}},
          \qquad \text{rate}=-2\sqrt{5} .
        \]

  \item $\grad f=(2x,\,-2y)$. At $(2,1)$: $\grad f=(4,-2)$. We need
        $\mathbf{u}\perp(4,-2)$: $4u_{1}-2u_{2}=0$, i.e.
        $u_{2}=2u_{1}$. Normalising $(1,2)$:
        \[
          \mathbf{u}=\pm\frac{(1,2)}{\sqrt{5}} .
        \]
        \textbf{Check.} $(4,-2)\cdot(1,2)=4-4=0$. \checkmark

  \item $\grad f=(2x,\,2y,\,-2z)$. At $(1,2,2)$: $\grad f=(2,4,-4)$,
        $|\grad f|=\sqrt{4+16+16}=6$. Maximum rate $6$, attained in the
        direction
        \[
          \mathbf{u}=\frac{(2,4,-4)}{6}=\frac{(1,2,-2)}{3} .
        \]

  \item $\grad z=(-2x,\,-2y)$; at $(1,1)$: $\grad z=(-2,-2)$,
        $|\grad z|=2\sqrt{2}$. Steepest descent direction:
        \[
          \mathbf{u}=-\frac{(-2,-2)}{2\sqrt{2}}=\frac{(1,1)}{\sqrt{2}} ,
        \]
        i.e. north-east, directly away from the summit above the origin.
        The slope of that descent is $-|\grad z|=-2\sqrt{2}$.

  \item $Q_{K}=\tfrac12K^{-1/2}L^{1/2}$, $Q_{L}=\tfrac12K^{1/2}L^{-1/2}$.
        At $(4,9)$:
        \[
          Q_{K}=\frac{1}{2}\cdot\frac{1}{2}\cdot 3=\frac{3}{4},\qquad
          Q_{L}=\frac{1}{2}\cdot 2\cdot\frac{1}{3}=\frac{1}{3},
        \]
        so $\grad Q=(3/4,\,1/3)$. Balanced direction
        $\mathbf{u}=(1,1)/\sqrt{2}$:
        \[
          D_{\mathbf{u}}Q=\frac{3/4+1/3}{\sqrt{2}}
            =\frac{13/12}{\sqrt{2}}=\frac{13\sqrt{2}}{24} .
        \]
        Fastest increase along $\grad Q/|\grad Q|$ with
        \[
          |\grad Q|=\sqrt{\frac{9}{16}+\frac{1}{9}}
            =\sqrt{\frac{81+16}{144}}=\frac{\sqrt{97}}{12},\qquad
          \mathbf{u}^{*}=\frac{12}{\sqrt{97}}\Bigl(\frac34,\frac13\Bigr)
            =\frac{(9,4)}{\sqrt{97}} .
        \]
        The maximal rate is $\sqrt{97}/12$ (larger than $13\sqrt{2}/24$,
        as it must be).

  \item Let $\mathbf{r}(t)=(x(t),y(t))$ be a differentiable curve lying
        in the level set, $f(\mathbf{r}(t))=c$ for all $t$.
        Differentiate with the chain rule:
        \[
          0=\dd{}{t}\,c=\grad f(\mathbf{r}(t))\cdot\mathbf{r}'(t) .
        \]
        Hence the tangent vector $\mathbf{r}'(t)$ is orthogonal to
        $\grad f$ at every point of the curve. For $U(x,y)=c$ an
        indifference curve, this says exactly that
        $(U_{x},U_{y})$ is normal to it.

  \item First order: $\grad f=(2xy,\,x^{2})$; at $(1,2)$:
        $\grad f=(4,1)$.
        \[
          D_{\mathbf{u}}f=4\cdot\frac{3}{5}+1\cdot\frac{4}{5}
            =\frac{12}{5}+\frac{4}{5}=\frac{16}{5} .
        \]
        Second order: $f_{xx}=2y=4$, $f_{xy}=2x=2$, $f_{yy}=0$ at
        $(1,2)$.
        \[
          D^{2}_{\mathbf{u}}f
            =4\Bigl(\frac{3}{5}\Bigr)^{2}
             +2\cdot 2\cdot\frac{3}{5}\cdot\frac{4}{5}
             +0\cdot\Bigl(\frac{4}{5}\Bigr)^{2}
            =\frac{36}{25}+\frac{48}{25}=\frac{84}{25} .
        \]

  \item $\grad f=(3x^{2}-3y,\,3y^{2}-3x)$; at $(1,1)$:
        $\grad f=(0,0)$, so $(1,1)$ is a critical point. Second
        partials: $f_{xx}=6x=6$, $f_{yy}=6y=6$, $f_{xy}=-3$, so
        $\Hess_{f}(1,1)=\begin{pmatrix}6&-3\\-3&6\end{pmatrix}$. For
        $\mathbf{u}=(1,1)/\sqrt{2}$:
        \[
          D^{2}_{\mathbf{u}}f=6\cdot\frac12
            +2(-3)\cdot\frac12+6\cdot\frac12=3-3+3=3 .
        \]
        For $\mathbf{v}=(1,-1)/\sqrt{2}$:
        \[
          D^{2}_{\mathbf{v}}f=6\cdot\frac12
            +2(-3)\cdot\Bigl(-\frac12\Bigr)+6\cdot\frac12=3+3+3=9 .
        \]
        Both are positive, the surface curves upward in both
        directions. Indeed $D_{1}=6>0$ and
        $D_{2}=36-9=27>0$, so the Hessian is positive definite and
        $(1,1)$ is a strict local minimum (\chref{derivatives}{11}).

  \item $f_{xx}=\mathrm{e}^{x}\cos y$, $f_{xy}=-\mathrm{e}^{x}\sin y$,
        $f_{yy}=-\mathrm{e}^{x}\cos y$. At $(0,0)$: $f_{xx}=1$,
        $f_{xy}=0$, $f_{yy}=-1$. Writing
        $\mathbf{u}=(\cos\theta,\,\sin\theta)$:
        \[
          D^{2}_{\mathbf{u}}f=\cos^{2}\theta-\sin^{2}\theta=\cos 2\theta .
        \]
        This vanishes for $2\theta=\tfrac{\pi}{2}+k\pi$, i.e.
        $\theta=\tfrac{\pi}{4}+k\tfrac{\pi}{2}$:
        \[
          \mathbf{u}=\pm\frac{(1,1)}{\sqrt{2}}
          \quad\text{and}\quad
          \mathbf{u}=\pm\frac{(1,-1)}{\sqrt{2}} .
        \]

  \item
        \[
          J_{\mathbf{f}}=\begin{pmatrix}1&1\\1&-1\end{pmatrix},\qquad
          \det J_{\mathbf{f}}=1\cdot(-1)-1\cdot 1=-2 .
        \]
        The map is linear, $\mathbf{f}(\mathbf{x})=A\mathbf{x}$ with
        $A=\bigl(\begin{smallmatrix}1&1\\1&-1\end{smallmatrix}\bigr)$,
        and for a linear map $f_{i}=\sum_{j}A_{ij}x_{j}$ gives
        $\partial f_{i}/\partial x_{j}=A_{ij}$: the Jacobian is the
        constant matrix $A$ itself: a linear map is its own linear
        approximation.

  \item With $f_{1}=x^{2}y$, $f_{2}=x+y^{2}$, $f_{3}=xy$:
        \[
          J_{\mathbf{f}}=\begin{pmatrix}2xy&x^{2}\\1&2y\\y&x\end{pmatrix}
          \;\in\;\R^{3\times 2} .
        \]
        A $3\times 2$ matrix: one row per component, one column per
        variable; being non-square, it has no determinant. At $(1,2)$:
        \[
          J_{\mathbf{f}}(1,2)=\begin{pmatrix}4&1\\1&4\\2&1\end{pmatrix} .
        \]

  \item
        \[
          J_{\mathbf{f}}=\begin{pmatrix}1&1&0\\0&1&1\\1&0&1\end{pmatrix} .
        \]
        Expanding along the first row:
        \[
          \det J_{\mathbf{f}}
            =1\,(1\cdot 1-1\cdot 0)-1\,(0\cdot 1-1\cdot 1)+0
            =1+1=2 .
        \]

  \item $J_{\mathbf{f}}=\begin{pmatrix}2x&-1\\y&x\end{pmatrix}$. At
        $\mathbf{x}_{0}=(1,1)$:
        \[
          \mathbf{f}(1,1)=(0,\,1),\qquad
          J_{\mathbf{f}}(1,1)=\begin{pmatrix}2&-1\\1&1\end{pmatrix} .
        \]
        With $\mathbf{h}=(0.1,\,-0.1)$:
        \[
          J_{\mathbf{f}}(1,1)\,\mathbf{h}
            =\begin{pmatrix}2(0.1)+(-1)(-0.1)\\
              1(0.1)+1(-0.1)\end{pmatrix}
            =\begin{pmatrix}0.3\\0\end{pmatrix},
        \]
        so $\mathbf{f}(1.1,\,0.9)\approx(0,1)+(0.3,\,0)=(0.3,\,1)$.
        Exact value:
        \[
          \mathbf{f}(1.1,\,0.9)=(1.21-0.9,\;1.1\cdot 0.9)
            =(0.31,\;0.99) .
        \]
        The approximation is off by only $0.01$ in each component.

  \item Excess demand:
        \[
          G_{1}=D-S=(100-p+0.5q)-(-10+2p+q)=110-3p-0.5q ,
        \]
        and $G_{2}=p+q-20$. Hence
        \[
          \frac{\partial(G_{1},G_{2})}{\partial(p,q)}
            =\begin{pmatrix}-3&-0.5\\1&1\end{pmatrix},\qquad
          \det=-3\cdot 1-(-0.5)\cdot 1=-2.5 .
        \]
        The determinant is nonzero: exactly the condition the next
        chapter's Implicit Function Theorem needs to solve the system
        locally for $(p,q)$ when a parameter shifts.

  \item Inner and outer Jacobians:
        \[
          J_{\mathbf{g}}=\begin{pmatrix}1&1\\1&-1\end{pmatrix},\qquad
          J_{\mathbf{F}}(u,v)=\begin{pmatrix}2u&0\\v&u\end{pmatrix} .
        \]
        Evaluate the outer at the inner value
        $(u,v)=\mathbf{g}(x,y)=(x+y,\,x-y)$ and multiply:
        \[
          J_{\mathbf{F}\circ\mathbf{g}}
            =\begin{pmatrix}2(x+y)&0\\x-y&x+y\end{pmatrix}
             \begin{pmatrix}1&1\\1&-1\end{pmatrix}
            =\begin{pmatrix}2(x+y)&2(x+y)\\
              (x-y)+(x+y)&(x-y)-(x+y)\end{pmatrix}
        \]
        \[
          \phantom{J_{\mathbf{F}\circ\mathbf{g}}}
            =\begin{pmatrix}2(x+y)&2(x+y)\\2x&-2y\end{pmatrix} .
        \]
        \textbf{Check.} Composing first:
        $(\mathbf{F}\circ\mathbf{g})(x,y)
        =\bigl((x+y)^{2},\,(x+y)(x-y)\bigr)
        =\bigl((x+y)^{2},\,x^{2}-y^{2}\bigr)$, whose Jacobian is
        \[
          \begin{pmatrix}2(x+y)&2(x+y)\\2x&-2y\end{pmatrix} .
          \quad\checkmark
        \]

  \item $J_{f}=\bigl(y^{2}\;\;2xy\bigr)$ (a row) and
        $J_{\mathbf{r}}(t)=\begin{pmatrix}\mathrm{e}^{t}\\1\end{pmatrix}$
        (a column). At $\mathbf{r}(t)=(\mathrm{e}^{t},\,t)$ the row is
        $\bigl(t^{2}\;\;2t\mathrm{e}^{t}\bigr)$, so
        \[
          \dd{z}{t}
            =\bigl(t^{2}\;\;2t\mathrm{e}^{t}\bigr)
             \begin{pmatrix}\mathrm{e}^{t}\\1\end{pmatrix}
            =t^{2}\mathrm{e}^{t}+2t\mathrm{e}^{t}
            =\mathrm{e}^{t}(t^{2}+2t) .
        \]
        \textbf{Check.} $z=\mathrm{e}^{t}\,t^{2}$, so
        $\dd{z}{t}=\mathrm{e}^{t}t^{2}+2t\mathrm{e}^{t}$. \checkmark

  \item Inner Jacobian:
        \[
          J_{\mathbf{g}}(K,L)=\begin{pmatrix}
            \tfrac12K^{-1/2}L^{1/2}&\tfrac12K^{1/2}L^{-1/2}\\
            2&1\end{pmatrix},\qquad
          J_{\mathbf{g}}(1,1)=\begin{pmatrix}1/2&1/2\\2&1\end{pmatrix} .
        \]
        Outer Jacobian (a row, constant since $\pi$ is linear):
        $J_{\pi}=(10\;\;-1)$. Matrix chain rule:
        \[
          (\pi_{K},\,\pi_{L})
            =(10\;\;-1)\begin{pmatrix}1/2&1/2\\2&1\end{pmatrix}
            =\bigl(10\cdot\tfrac12-1\cdot 2,\;
              10\cdot\tfrac12-1\cdot 1\bigr)
            =(3,\;4) .
        \]
        \textbf{Check.} Substituting first,
        $\pi(K,L)=10K^{1/2}L^{1/2}-2K-L$, so
        $\pi_{K}=5K^{-1/2}L^{1/2}-2=5-2=3$ and
        $\pi_{L}=5K^{1/2}L^{-1/2}-1=5-1=4$ at $(1,1)$. \checkmark
        At the current input mix an extra unit of capital raises profit
        by about $3$, an extra unit of labour by about $4$.
\end{solutions}

\section{Additional Problems}

\begin{problems}
  \item For $f(x,y)=x^{2}-y^{2}+3xy$ at $(1,2)$, find the directional
        derivative toward $\mathbf{v}=(4,3)$.
  \item For $f(x,y,z)=x^{2}+2y^{2}+3z^{2}$ at $(1,1,1)$, find the
        directional derivative toward $\mathbf{v}=(1,-1,1)$.
  \item For $U(x,y)=\ln x+2\ln y$ at $(1,2)$, find the direction of
        steepest ascent and the maximum rate of increase.
  \item Find all unit vectors $\mathbf{u}$ with $D_{\mathbf{u}}f=0$ for
        $f(x,y)=xy$ at $(2,1)$.
  \item At a point $P$ the directional derivatives of $f$ toward $(1,0)$
        and toward $(0,1)$ are $3$ and $-4$. Find $\grad f(P)$, the
        maximum rate of increase at $P$, and $D_{\mathbf{u}}f$ toward
        $\mathbf{u}=(3/5,\,4/5)$.
  \item For $f(x,y)=\ln(x^{2}+y^{2})$ at $(1,0)$ with
        $\mathbf{u}=(0,1)$, compute $D_{\mathbf{u}}f$ and
        $D^{2}_{\mathbf{u}}f$, and interpret the two results.
  \item Compute the Jacobian of
        $\mathbf{f}(x,y)=(x^{2}-y^{2},\;3xy+y^{2})$ and evaluate it at
        $(1,1)$.
  \item For $f(x,y)=xy$ and $\mathbf{r}(t)=(t^{2},\,t^{3})$, compute
        $\dd{z}{t}$ for $z=f(\mathbf{r}(t))$ by the matrix chain rule
        and verify directly.
  \item Let $\mathbf{g}(s,t)=(st,\;s+t)$ and
        $\mathbf{F}(x,y)=(x+y,\;xy)$. Compute
        $J_{\mathbf{F}\circ\mathbf{g}}$ by the matrix chain rule and
        verify by composing first.
  \item Inputs produce two intermediate goods $x_{1}=2K+L$ and
        $x_{2}=K+3L$, and final output is $y=x_{1}x_{2}$. Use the matrix
        chain rule to find $\pd{y}{K}$ and $\pd{y}{L}$ at $(K,L)=(1,1)$,
        and verify by expanding $y$ as a polynomial in $K$ and $L$.
\end{problems}

Problems 11--14: tangent lines to the graph of $f$ in a given direction.

\begin{problems}[resume]
  \item For $f(x,y)=x^{2}y+y^{3}$ at $P=(1,2)$, give the slope of the graph
        and the tangent line to the graph in the $x$-direction, in the
        $y$-direction, and in the direction of $\mathbf{v}=(3,-4)$.
  \item For $f(x,y)=\ln(x^{2}+3y)$ at $P=(2,0)$, give the slope of the
        graph in the direction of steepest ascent and the tangent line to
        the graph in that direction.
  \item For $f(x,y)=(x+y)\,\mathrm{e}^{x-y}$ at $P=(1,1)$, give the slope in
        the direction from $P$ towards $Q=(3,2)$ and the tangent line in
        that direction, once with the unit vector and once with
        $\mathbf{h}=Q-P$. Show that both describe the same line.
  \item For $Q(K,L)=4K^{1/2}L^{1/4}$ at $(16,16)$, give the tangent line to
        the graph in the $L$-direction, read it as an output estimate for
        $L=17$, and compare with the exact value.
\end{problems}

In problems 15--17, for the point $P$: (a) find the level $c=f(P)$; (b)
solve the level curve $f=c$ explicitly; (c) compute $\grad f(P)$ and check
that it is perpendicular to the tangent of the level curve; (d) give the
tangent line to the level curve in normal and in parametric form; (e) give
$D_{\mathbf{u}}f(P)$ in the direction of the gradient.

\begin{problems}[resume]
  \item $f(x,y)=\sqrt{2x^{2}+3y+1}$, \;$P=(1,2)$.
  \item $f(x,y)=x^{1/3}y^{2/3}$ on $x,y>0$, \;$P=(1,8)$.
  \item $f(x,y)=x\,\mathrm{e}^{y}+y$, \;$P=(2,0)$; solve the level curve
        for $x$.
\end{problems}

Problems 18--20: differentiability of vector-valued maps.

\begin{problems}[resume]
  \item Explain why
        $\mathbf{F}(x,y)=\bigl(x^{2}y,\;x-3y^{2},\;xy\bigr)$ is
        differentiable at every point. Give $J_{\mathbf{F}}(1,1)$, and
        compute the remainder vector $\mathbf{r}(\mathbf{h})$ at $(1,1)$
        to show directly that $\|\mathbf{r}(\mathbf{h})\|/\|\mathbf{h}\|\to0$.
  \item For $\mathbf{F}(x,y)=\bigl(\ln(3x^{2}+y^{4}),\;xy^{2},\;|y-1|\bigr)$,
        give the domain, the set on which $\mathbf{F}$ is differentiable,
        and the Jacobian there.
  \item Answer in words: when is a map $\mathbf{f}\colon\R^{n}\to\R^{m}$
        differentiable at $\mathbf{x}_{0}$; what is its derivative called
        and what shape does it have; how is it related to the partial
        derivatives of the components; and what is the geometric meaning
        for $n=m=1$ and for $n=2$, $m=1$?
\end{problems}

Problems 21--26: Jacobians of all shapes.

\begin{problems}[resume]
  \item $\mathbf{r}(t)=\bigl(\ln(1+t^{2}),\;\mathrm{e}^{-2t},\;\cos3t,\;5\bigr)$:
        the Jacobian, its shape, and its value at $t=0$.
  \item $\mathbf{F}(x,y,z)=(xyz,\;x^{2}-z)$ at $(1,2,3)$, and
        $g(x_{1},x_{2},x_{3},x_{4})=x_{1}x_{2}+x_{3}\mathrm{e}^{x_{4}}$ at
        $(1,1,1,0)$: shapes and values of the Jacobians.
  \item $\mathbf{F}(s,t)=\bigl(s+t,\;st,\;s^{2},\;\mathrm{e}^{s-t}\bigr)$:
        the Jacobian, its shape, and its value at $(1,1)$.
  \item Jacobians in $n$ dimensions: (a) the identity map
        $\mathbf{F}(\mathbf{x})=\mathbf{x}$ on $\R^{n}$; (b) the map
        $\mathbf{F}\colon\R^{n}\to\R^{n}$ with components
        $F_{k}(\mathbf{x})=k\,\mathrm{e}^{x_{k}}$, that is
        $(\mathrm{e}^{x_{1}},2\mathrm{e}^{x_{2}},\ldots,n\,\mathrm{e}^{x_{n}})$;
        (c) the curve $\mathbf{r}\colon\R\to\R^{n}$ with components
        $r_{k}(t)=t^{k}/k$, that is $(t,t^{2}/2,\ldots,t^{n}/n)$; (d)
        $f(\mathbf{x})=\sum_{k=1}^{n}x_{k}^{k}/k$.
  \item Jacobians of linear and affine maps: (a) $\mathbf{F}(\mathbf{x})=
        A\mathbf{x}$ with $A=\begin{pmatrix}2&-1&0\\1&0&3\\0&4&-2
        \end{pmatrix}$; (b) $\mathbf{F}(\mathbf{x})=3\mathbf{x}-2\mathbf{a}$
        for a fixed $\mathbf{a}\in\R^{n}$; (c) $f(\mathbf{x})=
        \mathbf{a}^{T}\mathbf{x}$; (d) $\mathbf{G}(\mathbf{x})=B\mathbf{x}
        +\mathbf{c}$ with $B$ a $2\times3$ matrix.
  \item Let $\mathbf{F}\colon\R^{3}\to\R^{2}$, $\mathbf{G}\colon\R^{2}\to
        \R^{4}$, $H\colon\R^{4}\to\R$ and $\mathbf{K}\colon\R\to\R^{3}$.
        Which of $\mathbf{G}\circ\mathbf{F}$, $\mathbf{F}\circ\mathbf{G}$,
        $H\circ\mathbf{G}$, $\mathbf{G}\circ H$, $\mathbf{K}\circ H$,
        $H\circ\mathbf{K}$, $\mathbf{F}\circ\mathbf{K}$,
        $\mathbf{K}\circ\mathbf{F}$ and $H\circ\mathbf{G}\circ\mathbf{F}$
        are defined, and what is the shape of the Jacobian of each one that
        is?
\end{problems}

Problems 27--30: the matrix chain rule with rectangular Jacobians. Verify
each result by composing first.

\begin{problems}[resume]
  \item $\mathbf{f}(x,y,z)=(x+yz,\;xz)$ and
        $\mathbf{g}(u,v)=(u,\;v,\;uv,\;u-v)$: the Jacobian of
        $\mathbf{g}\circ\mathbf{f}$ at $(1,2,1)$.
  \item $g(x,y,z)=x+2y+z^{2}$ and $\mathbf{f}(s)=(s,\;s^{2},\;\mathrm{e}^{s})$:
        say why both $\mathbf{f}\circ g$ and $g\circ\mathbf{f}$ exist; give
        the Jacobian of $\mathbf{f}\circ g$ at $(1,0,0)$ and the derivative
        of $g\circ\mathbf{f}$.
  \item $\mathbf{f}(x,y)=(x,\;y,\;xy)$ and $\mathbf{g}(u,v,w)=(u+w,\;vw)$:
        the Jacobian of $\mathbf{g}\circ\mathbf{f}$ at $(1,2)$ and of
        $\mathbf{f}\circ\mathbf{g}$ at $(1,1,1)$.
  \item $\mathbf{r}(t)=(\cos t,\;\sin t,\;t)$ and $F(x,y,z)=x^{2}+y^{2}+z^{2}$:
        $\dd{}{t}F(\mathbf{r}(t))$ by the chain rule, simplified.
\end{problems}

In problems 31--33, compute $D_{\mathbf{u}}f(P)$ and $D^{2}_{\mathbf{u}}f(P)$
for the direction of $\mathbf{h}$, both with the gradient and Hessian and
with the profile $\varphi(t)=f(P+t\mathbf{h})$, and say whether $f$ is
convex or concave in that direction at $P$.

\begin{problems}[resume]
  \item $f(x,y,z)=x\ln y+yz^{2}$, \;$P=(2,1,1)$, \;$\mathbf{h}=(1,2,-2)$.
  \item $f(x,y,z)=x\,\mathrm{e}^{yz}$, \;$P=(1,0,2)$, \;$\mathbf{h}=(2,1,2)$.
  \item $f(x,y,z)=xyz+\ln(xz)$, \;$P=(1,1,1)$, \;$\mathbf{h}=(1,-1,1)$.
  \item Redo practice problem 15 ($f=x^{2}y$ at $(1,2)$) with the
        \emph{unnormalised} $\mathbf{h}=(3,4)$: compute
        $\grad f\cdot\mathbf{h}$, $\mathbf{h}^{T}H\mathbf{h}$, and the
        profile derivatives $\varphi'(0)$, $\varphi''(0)$. Explain the
        factors $5$ and $25$.
\end{problems}

In problems 35--37, give the directional verdicts asked for, then decide
whether $f$ is convex, concave or neither on its whole domain.

\begin{problems}[resume]
  \item $f(x,y)=x^{3}-3x+y^{2}$: at $Q=(-1,1)$ in the direction of
        $\mathbf{h}=(3,4)$, compute $D_{\mathbf{u}}f$ and
        $D^{2}_{\mathbf{u}}f$; then use the point $(1,0)$.
  \item $f(x,y)=x^{4}-4xy+y^{2}$: classify the critical points, then give
        $D^{2}_{\mathbf{u}}f(0,0)$ for $\mathbf{u}=(1,1)/\sqrt2$ and for
        $\mathbf{u}=(0,1)$.
  \item $f(x,y)=x\ln y$ on $y>0$, at $P=(1,1)$: $D_{\mathbf{u}}f$ and
        $D^{2}_{\mathbf{u}}f$ for $\mathbf{u}=(1,0)$, $(0,1)$ and
        $(1,1)/\sqrt2$.
\end{problems}

Problems 38--40: statements to check. Justify every answer.

\begin{problems}[resume]
  \item True or false? (a) Walking along the contour line through $P$,
        the rate of change of $f$ is $0$. (b) The direction of steepest
        descent is $-\grad f(P)$. (c) If $\grad f(P)=\mathbf{0}$, then $P$
        is a peak or a valley. (d) $|\grad f(P)|$ is the slope of the
        steepest path through $P$. (e) The gradient at $P$ is tangent to
        the level curve through $P$.
  \item A hiker stands on a hillside with neither map nor formula. How can
        she find the direction of the gradient, and how does the slope in
        other directions compare with the slope in that one?
  \item True or false? (a) For $f(x,y)$ of two variables: if
        $D_{\mathbf{u}}f(P)=0$ for two non-parallel unit vectors $\mathbf{u}$, then $\grad f(P)=\mathbf{0}$. (b) If
        $D_{\mathbf{u}}f(P)=3$, then $D_{-\mathbf{u}}f(P)=-3$. (c) A
        function can have $D_{\mathbf{u}}f(P)=5$ for some unit vector
        $\mathbf{u}$ while $|\grad f(P)|=4$.
\end{problems}

Problems 41--42: an input change vector of fixed length is to raise output
as fast as possible.

\begin{problems}[resume]
  \item $Y=10K^{0.4}L^{0.6}$ at $K=32768$, $L=7776$ (so that
        $K^{0.4}=64$, $L^{0.6}=216$). The change vector
        $(\mathrm{d}K,\mathrm{d}L)$ has length $100$. Find the best ratio
        $\mathrm{d}K:\mathrm{d}L$, the amounts, the estimated and the exact
        rise in output; then compare with spending a budget
        $\mathrm{d}K+\mathrm{d}L=100$.
  \item $Y=\sqrt{KL}$ at $K=400$, $L=900$, with a change vector of length
        $20$: the same questions.
\end{problems}

\section{Additional Solutions}

\begin{solutions}
  \item $\grad f=(2x+3y,\,-2y+3x)$; at $(1,2)$:
        $\grad f=(2+6,\,-4+3)=(8,-1)$. Normalise:
        $|\mathbf{v}|=\sqrt{16+9}=5$, $\mathbf{u}=(4/5,\,3/5)$.
        \[
          D_{\mathbf{u}}f=8\cdot\frac{4}{5}+(-1)\cdot\frac{3}{5}
            =\frac{32}{5}-\frac{3}{5}=\frac{29}{5} .
        \]

  \item $\grad f=(2x,\,4y,\,6z)$; at $(1,1,1)$: $\grad f=(2,4,6)$.
        Normalise: $|\mathbf{v}|=\sqrt{1+1+1}=\sqrt{3}$,
        $\mathbf{u}=(1,-1,1)/\sqrt{3}$.
        \[
          D_{\mathbf{u}}f=\frac{2-4+6}{\sqrt{3}}
            =\frac{4}{\sqrt{3}}=\frac{4\sqrt{3}}{3} .
        \]

  \item $\grad U=(1/x,\,2/y)$; at $(1,2)$: $\grad U=(1,1)$,
        $|\grad U|=\sqrt{2}$. Steepest ascent along
        \[
          \mathbf{u}=\frac{(1,1)}{\sqrt{2}},\qquad
          \text{maximum rate}=\sqrt{2} .
        \]

  \item $\grad f=(y,\,x)$; at $(2,1)$: $\grad f=(1,2)$. We need
        $\mathbf{u}\perp(1,2)$: $u_{1}+2u_{2}=0$, so
        $\mathbf{u}=\pm(2,-1)/\sqrt{5}$.
        \textbf{Check.} $(1,2)\cdot(2,-1)=2-2=0$. \checkmark

  \item The directional derivatives toward $(1,0)$ and $(0,1)$ are
        precisely $f_{x}(P)$ and $f_{y}(P)$, so $\grad f(P)=(3,-4)$.
        Maximum rate: $|\grad f(P)|=\sqrt{9+16}=5$, attained toward
        $(3,-4)/5$. Toward $\mathbf{u}=(3/5,\,4/5)$:
        \[
          D_{\mathbf{u}}f=3\cdot\frac{3}{5}+(-4)\cdot\frac{4}{5}
            =\frac{9}{5}-\frac{16}{5}=-\frac{7}{5} .
        \]

  \item $\grad f=\Bigl(\dfrac{2x}{x^{2}+y^{2}},\,
        \dfrac{2y}{x^{2}+y^{2}}\Bigr)$; at $(1,0)$: $\grad f=(2,0)$, so
        \[
          D_{\mathbf{u}}f=(2,0)\cdot(0,1)=0 :
        \]
        $\mathbf{u}$ is tangent to the level curve $x^{2}+y^{2}=1$
        through $(1,0)$. Second partials:
        \[
          f_{xx}=\frac{2y^{2}-2x^{2}}{(x^{2}+y^{2})^{2}},\qquad
          f_{yy}=\frac{2x^{2}-2y^{2}}{(x^{2}+y^{2})^{2}},\qquad
          f_{xy}=\frac{-4xy}{(x^{2}+y^{2})^{2}} ;
        \]
        at $(1,0)$: $f_{xx}=-2$, $f_{yy}=2$, $f_{xy}=0$. Hence with
        $\mathbf{u}=(0,1)$:
        \[
          D^{2}_{\mathbf{u}}f=f_{yy}=2>0 .
        \]
        Along $\mathbf{u}$ the function is flat to first order but
        curves upward: moving off $(1,0)$ vertically raises $f$ at
        second order. \textbf{Check.}
        $f(1,t)=\ln(1+t^{2})\approx t^{2}$ for small $t$: derivative
        $0$, second derivative $2$. \checkmark

  \item
        \[
          J_{\mathbf{f}}=\begin{pmatrix}2x&-2y\\3y&3x+2y\end{pmatrix},
          \qquad
          J_{\mathbf{f}}(1,1)=\begin{pmatrix}2&-2\\3&5\end{pmatrix} .
        \]

  \item $J_{f}=\bigl(y\;\;x\bigr)$, evaluated at
        $\mathbf{r}(t)=(t^{2},\,t^{3})$: $\bigl(t^{3}\;\;t^{2}\bigr)$;
        $J_{\mathbf{r}}(t)=\begin{pmatrix}2t\\3t^{2}\end{pmatrix}$.
        \[
          \dd{z}{t}=\bigl(t^{3}\;\;t^{2}\bigr)
            \begin{pmatrix}2t\\3t^{2}\end{pmatrix}
            =2t^{4}+3t^{4}=5t^{4} .
        \]
        \textbf{Check.} $z=t^{2}\cdot t^{3}=t^{5}$, so
        $\dd{z}{t}=5t^{4}$. \checkmark

  \item Inner and outer Jacobians:
        \[
          J_{\mathbf{g}}(s,t)=\begin{pmatrix}t&s\\1&1\end{pmatrix},\qquad
          J_{\mathbf{F}}(x,y)=\begin{pmatrix}1&1\\y&x\end{pmatrix} .
        \]
        Evaluate the outer at $(x,y)=\mathbf{g}(s,t)=(st,\,s+t)$ and
        multiply:
        \[
          J_{\mathbf{F}\circ\mathbf{g}}
            =\begin{pmatrix}1&1\\s+t&st\end{pmatrix}
             \begin{pmatrix}t&s\\1&1\end{pmatrix}
            =\begin{pmatrix}t+1&s+1\\
              (s+t)t+st&(s+t)s+st\end{pmatrix}
            =\begin{pmatrix}t+1&s+1\\
              t^{2}+2st&s^{2}+2st\end{pmatrix} .
        \]
        \textbf{Check.} Composing first:
        $(\mathbf{F}\circ\mathbf{g})(s,t)
        =\bigl(st+s+t,\;s^{2}t+st^{2}\bigr)$, with Jacobian
        \[
          \begin{pmatrix}t+1&s+1\\2st+t^{2}&s^{2}+2st\end{pmatrix} .
          \quad\checkmark
        \]

  \item Inner map $\mathbf{g}(K,L)=(2K+L,\;K+3L)$ with constant
        Jacobian
        $J_{\mathbf{g}}=\begin{pmatrix}2&1\\1&3\end{pmatrix}$; outer
        function $y=x_{1}x_{2}$ with row Jacobian
        $\grad y=(x_{2}\;\;x_{1})$. At $(K,L)=(1,1)$:
        $(x_{1},x_{2})=(3,4)$, so the row is $(4\;\;3)$ and
        \[
          \Bigl(\pd{y}{K},\,\pd{y}{L}\Bigr)
            =(4\;\;3)\begin{pmatrix}2&1\\1&3\end{pmatrix}
            =\bigl(4\cdot 2+3\cdot 1,\;4\cdot 1+3\cdot 3\bigr)
            =(11,\;13) .
        \]
        \textbf{Check.} Expanding,
        $y=(2K+L)(K+3L)=2K^{2}+7KL+3L^{2}$, so
        $\pd{y}{K}=4K+7L=11$ and $\pd{y}{L}=7K+6L=13$ at $(1,1)$.
        \checkmark

  \item $f(P)=2+8=10$ and $\grad f=(2xy,\;x^{2}+3y^{2})=(4,13)$ at $P$.
        \begin{itemize}[leftmargin=1.2em,itemsep=0.2ex]
          \item $x$-direction: slope $f_{x}(P)=4$, line
                $(1+t,\;2,\;10+4t)$.
          \item $y$-direction: slope $f_{y}(P)=13$, line
                $(1,\;2+t,\;10+13t)$.
          \item $|\mathbf{v}|=5$, $\mathbf{u}=(\tfrac35,-\tfrac45)$,
                $D_{\mathbf{u}}f=\tfrac{12}{5}-\tfrac{52}{5}=-8$: the
                surface falls with slope $8$ that way. Line
                $\bigl(1+\tfrac35t,\;2-\tfrac45t,\;10-8t\bigr)$.
        \end{itemize}
  \item $f(P)=\ln4$. With $w=x^{2}+3y$, $\grad f=\bigl(\tfrac{2x}{w},\,
        \tfrac{3}{w}\bigr)=(1,\tfrac34)$ at $P$, where $w=4$. Steepest
        ascent: $\mathbf{u}=\grad f/|\grad f|$ with $|\grad f|=\tfrac54$, so
        $\mathbf{u}=(\tfrac45,\tfrac35)$ and the slope is
        $D_{\mathbf{u}}f=|\grad f|=\tfrac54$. Tangent line:
        \[
          (x,y,z)=\bigl(2+\tfrac45t,\;\tfrac35t,\;\ln4+\tfrac54t\bigr) .
        \]
  \item $f(P)=2\mathrm{e}^{0}=2$. Product rule:
        $f_{x}=\mathrm{e}^{x-y}(x+y+1)=3$,
        $f_{y}=\mathrm{e}^{x-y}(1-x-y)=-1$ at $P$. $\mathbf{h}=Q-P=(2,1)$,
        $|\mathbf{h}|=\sqrt5$, $\mathbf{u}=(2,1)/\sqrt5$:
        \[
          D_{\mathbf{u}}f=\frac{3\cdot2-1\cdot1}{\sqrt5}=\sqrt5\approx2.236 .
        \]
        With the unit vector:
        $\bigl(1+\tfrac{2}{\sqrt5}t,\;1+\tfrac{1}{\sqrt5}t,\;2+\sqrt5\,t\bigr)$.
        With $\mathbf{h}$, $\grad f\cdot\mathbf{h}=5$:
        $(1+2s,\;1+s,\;2+5s)$. Putting $s=t/\sqrt5$ in the second gives the
        first, so it is the same line; only the parameter runs at a
        different speed.
  \item $Q(16,16)=4\cdot4\cdot2=32$,
        $Q_{L}=K^{1/2}L^{-3/4}=4\cdot\tfrac18=\tfrac12$. Tangent line in the
        $L$-direction: $(16,\;16+t,\;32+\tfrac12t)$. At $L=17$ ($t=1$) it
        gives $32.5$; exactly, $Q(16,17)=16\cdot17^{1/4}\approx32.4887$. The
        slope $\tfrac12$ is the marginal product of labour, and the line
        slightly overestimates because $Q$ is concave in $L$.
  \item (a) $c=\sqrt{2+6+1}=3$. (b) $2x^{2}+3y+1=9$, so
        $y=\dfrac{8-2x^{2}}{3}$, a downward parabola. Its slope at $x=1$ is
        $y'=-\tfrac{4x}{3}=-\tfrac43$, tangent direction $(3,-4)$.
        (c) $\grad f=\Bigl(\dfrac{2x}{f},\,\dfrac{3}{2f}\Bigr)=
        \bigl(\tfrac23,\tfrac12\bigr)$, and
        $\bigl(\tfrac23,\tfrac12\bigr)\cdot(3,-4)=2-2=0$. \checkmark\ (d)
        Normal form: $\tfrac23(x-1)+\tfrac12(y-2)=0$, i.e.\
        $4(x-1)+3(y-2)=0$, $y=2-\tfrac43(x-1)$. Parametric:
        $(x,y)=(1,2)+t(3,-4)$. (e) $D_{\mathbf{u}}f=|\grad f|=
        \sqrt{\tfrac49+\tfrac14}=\sqrt{\tfrac{25}{36}}=\tfrac56$.
  \item (a) $c=1\cdot8^{2/3}=4$. (b) $y^{2/3}=4x^{-1/3}$, so
        $y=8x^{-1/2}$; slope $y'=-4x^{-3/2}=-4$ at $x=1$, direction $(1,-4)$.
        (c) $f_{x}=\tfrac13\dfrac fx=\tfrac43$,
        $f_{y}=\tfrac23\dfrac fy=\tfrac23\cdot\tfrac48=\tfrac13$, and
        $\bigl(\tfrac43,\tfrac13\bigr)\cdot(1,-4)=0$. \checkmark\ (d)
        $\tfrac43(x-1)+\tfrac13(y-8)=0$, i.e.\ $y=12-4x$; parametric
        $(1,8)+t(1,-4)$. (e) $|\grad f|=\sqrt{\tfrac{16}{9}+\tfrac19}=
        \tfrac{\sqrt{17}}{3}\approx1.374$.
  \item (a) $c=2\cdot1+0=2$. (b) $x\mathrm{e}^{y}=2-y$, so
        $x=(2-y)\mathrm{e}^{-y}$. Then
        $\dd{x}{y}=-\mathrm{e}^{-y}-(2-y)\mathrm{e}^{-y}=(y-3)\mathrm{e}^{-y}
        =-3$ at $y=0$: tangent direction $(\mathrm{d}x,\mathrm{d}y)=(-3,1)$.
        (c) $\grad f=(\mathrm{e}^{y},\;x\mathrm{e}^{y}+1)=(1,3)$, and
        $(1,3)\cdot(-3,1)=0$. \checkmark\ (d) $(x-2)+3y=0$, i.e.\
        $x=2-3y$; parametric $(2,0)+t(-3,1)$. (e) $|\grad f|=\sqrt{10}
        \approx3.162$.
  \item Each component is a polynomial, so its partials exist and are
        continuous everywhere; a $C^{1}$ map is differentiable, and a map is
        differentiable when each component is.
        $J_{\mathbf{F}}=\begin{pmatrix}2xy&x^{2}\\1&-6y\\y&x\end{pmatrix}$,
        and $J_{\mathbf{F}}(1,1)=\begin{pmatrix}2&1\\1&-6\\1&1\end{pmatrix}$.
        Remainder at $(1,1)$, component by component, with
        $\mathbf{F}(1,1)=(1,-2,1)$:
        \[\begin{aligned}
          r_{1}&=(1+h_{1})^{2}(1+h_{2})-1-2h_{1}-h_{2}
               =2h_{1}h_{2}+h_{1}^{2}+h_{1}^{2}h_{2},\\
          r_{2}&=(1+h_{1})-3(1+h_{2})^{2}+2-h_{1}+6h_{2}=-3h_{2}^{2},\\
          r_{3}&=(1+h_{1})(1+h_{2})-1-h_{1}-h_{2}=h_{1}h_{2} .
        \end{aligned}\]
        Each is bounded by a constant times $\|\mathbf{h}\|^{2}$ (for
        $\|\mathbf{h}\|\le1$), so
        $\|\mathbf{r}\|\le C\|\mathbf{h}\|^{2}$ and
        $\|\mathbf{r}\|/\|\mathbf{h}\|\to0$.
  \item \emph{Domain.} $3x^{2}+y^{4}>0$ fails only at the origin: $D=\R^{2}
        \setminus\{(0,0)\}$. \emph{Differentiability.} The first two
        components are $C^{1}$ on $D$. The third, $|y-1|$, has
        $\partial/\partial y$ equal to $+1$ above and $-1$ below the line
        $y=1$, and on the line the quotient $|h|/h$ has no limit. So
        $\mathbf{F}$ is differentiable exactly on
        $\{(x,y)\neq(0,0):y\neq1\}$, with
        \[
          J_{\mathbf{F}}=\begin{pmatrix}
            \dfrac{6x}{3x^{2}+y^{4}}&\dfrac{4y^{3}}{3x^{2}+y^{4}}\\[1.2ex]
            y^{2}&2xy\\ 0&\operatorname{sign}(y-1)\end{pmatrix} .
        \]
  \item $\mathbf{f}$ is differentiable at $\mathbf{x}_{0}$ if
        $\mathbf{f}(\mathbf{x}_{0}+\mathbf{h})=\mathbf{f}(\mathbf{x}_{0})
        +J\mathbf{h}+\mathbf{r}(\mathbf{h})$ for a matrix $J$, with
        $\|\mathbf{r}(\mathbf{h})\|/\|\mathbf{h}\|\to0$: near
        $\mathbf{x}_{0}$ the map is a linear map plus an error that vanishes
        faster than the step. The matrix is the derivative (total
        derivative), the Jacobian matrix, of shape $m\times n$. Its entries
        are the partials $\partial f_{i}/\partial x_{j}$; their existence is
        necessary but not sufficient, while continuity of all of them is
        sufficient, and $\mathbf{f}$ is differentiable exactly when each
        component is. For $n=m=1$ the linear part is the tangent line of
        the graph, for $n=2$, $m=1$ the tangent plane.
  \item $J_{\mathbf{r}}(t)=\Bigl(\dfrac{2t}{1+t^{2}},\;-2\mathrm{e}^{-2t},
        \;-3\sin3t,\;0\Bigr)^{T}$, a $4\times1$ column (four outputs, one
        input); the constant component gives $0$. At $t=0$:
        $(0,-2,0,0)^{T}$.
  \item $J_{\mathbf{F}}=\begin{pmatrix}yz&xz&xy\\2x&0&-1\end{pmatrix}$,
        $2\times3$; at $(1,2,3)$: $\begin{pmatrix}6&3&2\\2&0&-1\end{pmatrix}$.
        $J_{g}=\bigl(x_{2}\ \ x_{1}\ \ \mathrm{e}^{x_{4}}\ \
        x_{3}\mathrm{e}^{x_{4}}\bigr)$, a $1\times4$ row; at $(1,1,1,0)$:
        $(1\ \ 1\ \ 1\ \ 1)$.
  \item $J_{\mathbf{F}}=\begin{pmatrix}1&1\\t&s\\2s&0\\
        \mathrm{e}^{s-t}&-\mathrm{e}^{s-t}\end{pmatrix}$, $4\times2$ (four
        outputs, two inputs); at $(1,1)$:
        $\begin{pmatrix}1&1\\1&1\\2&0\\1&-1\end{pmatrix}$.
  \item (a) $F_{k}=x_{k}$, so $\partial F_{k}/\partial x_{j}$ is $1$ for
        $j=k$ and $0$ otherwise: $J=I_{n}$. (b) $F_{k}=k\mathrm{e}^{x_{k}}$
        depends on $x_{k}$ alone: $J=\operatorname{diag}\bigl(\mathrm{e}^{x_{1}},
        2\mathrm{e}^{x_{2}},\ldots,n\,\mathrm{e}^{x_{n}}\bigr)$, $n\times n$.
        (c) $r_{k}=t^{k}/k$ has $r_{k}'=t^{k-1}$: the $n\times1$ column
        $(1,t,t^{2},\ldots,t^{n-1})^{T}$. (d) The term with $x_{k}$ is
        $x_{k}^{k}/k$, with derivative $x_{k}^{k-1}$: the $1\times n$ row
        $(1\ \ x_{2}\ \ x_{3}^{2}\ \ \cdots\ \ x_{n}^{n-1})$.
  \item (a) Component $i$ is $\sum_{j}a_{ij}x_{j}$, with partials $a_{ij}$:
        $J=A$ at every point. (b) Component $k$ is $3x_{k}-2a_{k}$:
        $J=3I_{n}$, and the constant shift drops out. (c) $\partial f/
        \partial x_{k}=a_{k}$: $J=\mathbf{a}^{T}$, a $1\times n$ row. (d)
        $J=B$, $2\times3$, whatever $\mathbf{c}$ is.
  \item Match outputs of the inner map with inputs of the outer map.
        \begin{itemize}[leftmargin=1.2em,itemsep=0.2ex]
          \item $\mathbf{G}\circ\mathbf{F}$: $2=2$, defined,
                $\R^{3}\to\R^{4}$, $4\times3$. $\mathbf{F}\circ\mathbf{G}$:
                $4\neq3$, undefined.
          \item $H\circ\mathbf{G}$: $4=4$, defined, $\R^{2}\to\R$,
                $1\times2$. $\mathbf{G}\circ H$: $1\neq2$, undefined.
          \item $\mathbf{K}\circ H$: $1=1$, defined, $\R^{4}\to\R^{3}$,
                $3\times4$. $H\circ\mathbf{K}$: $3\neq4$, undefined.
          \item $\mathbf{F}\circ\mathbf{K}$: $3=3$, defined, $\R\to\R^{2}$,
                $2\times1$. $\mathbf{K}\circ\mathbf{F}$: $2\neq1$,
                undefined.
          \item $H\circ\mathbf{G}\circ\mathbf{F}$: $\R^{3}\to\R^{2}\to\R^{4}
                \to\R$, defined, $1\times3=(1\times4)(4\times2)(2\times3)$.
        \end{itemize}
  \item $\mathbf{f}(1,2,1)=(1+2,\;1)=(3,1)$.
        $J_{\mathbf{f}}=\begin{pmatrix}1&z&y\\z&0&x\end{pmatrix}
        =\begin{pmatrix}1&1&2\\1&0&1\end{pmatrix}$ and
        $J_{\mathbf{g}}=\begin{pmatrix}1&0\\0&1\\v&u\\1&-1\end{pmatrix}
        =\begin{pmatrix}1&0\\0&1\\1&3\\1&-1\end{pmatrix}$ at $(u,v)=(3,1)$.
        \[
          J_{\mathbf{g}\circ\mathbf{f}}=
          \begin{pmatrix}1&0\\0&1\\1&3\\1&-1\end{pmatrix}
          \begin{pmatrix}1&1&2\\1&0&1\end{pmatrix}
          =\begin{pmatrix}1&1&2\\1&0&1\\4&1&5\\0&1&1\end{pmatrix},
        \]
        a $(4\times2)(2\times3)=4\times3$ product. \textbf{Check.}
        $\mathbf{g}\circ\mathbf{f}=\bigl(x+yz,\;xz,\;x^{2}z+xyz^{2},\;
        x+yz-xz\bigr)$; the third row is
        $(2xz+yz^{2},\;xz^{2},\;x^{2}+2xyz)=(4,1,5)$ and the fourth
        $(1-z,\;z,\;y-x)=(0,1,1)$ at $(1,2,1)$. \checkmark
  \item $g\colon\R^{3}\to\R$ and $\mathbf{f}\colon\R\to\R^{3}$: in
        $\mathbf{f}\circ g$ the one output of $g$ feeds the one input of
        $\mathbf{f}$, and in $g\circ\mathbf{f}$ the three outputs of
        $\mathbf{f}$ feed the three inputs of $g$. Both exist.
        \emph{$\mathbf{f}\circ g\colon\R^{3}\to\R^{3}$.} The chain rule
        gives a column times a row,
        $J=\mathbf{f}'(g(\mathbf{x}))\,\grad g(\mathbf{x})^{T}$, a
        rank-one $3\times3$ matrix. At $(1,0,0)$: $g=1$,
        $\mathbf{f}'(1)=(1,\;2,\;\mathrm{e})^{T}$, $\grad g=(1,2,2z)=(1,2,0)$:
        \[
          J=\begin{pmatrix}1\\2\\\mathrm{e}\end{pmatrix}(1\ \ 2\ \ 0)
           =\begin{pmatrix}1&2&0\\2&4&0\\\mathrm{e}&2\mathrm{e}&0\end{pmatrix} .
        \]
        Check: $\mathbf{f}(g)=(g,\;g^{2},\;\mathrm{e}^{g})$ has rows
        $\grad g$, $2g\grad g$, $\mathrm{e}^{g}\grad g$. \checkmark\
        \emph{$g\circ\mathbf{f}\colon\R\to\R$.}
        $\grad g(\mathbf{f}(s))\cdot\mathbf{f}'(s)=(1,\;2,\;2\mathrm{e}^{s})
        \cdot(1,\;2s,\;\mathrm{e}^{s})=1+4s+2\mathrm{e}^{2s}$. Check:
        $g(\mathbf{f}(s))=s+2s^{2}+\mathrm{e}^{2s}$. \checkmark
  \item \emph{$\mathbf{g}\circ\mathbf{f}\colon\R^{2}\to\R^{2}$} at $(1,2)$:
        $\mathbf{f}(1,2)=(1,2,2)$,
        $J_{\mathbf{f}}=\begin{pmatrix}1&0\\0&1\\y&x\end{pmatrix}
        =\begin{pmatrix}1&0\\0&1\\2&1\end{pmatrix}$,
        $J_{\mathbf{g}}=\begin{pmatrix}1&0&1\\0&w&v\end{pmatrix}
        =\begin{pmatrix}1&0&1\\0&2&2\end{pmatrix}$, so
        \[
          J_{\mathbf{g}\circ\mathbf{f}}=\begin{pmatrix}1&0&1\\0&2&2\end{pmatrix}
          \begin{pmatrix}1&0\\0&1\\2&1\end{pmatrix}
          =\begin{pmatrix}3&1\\4&4\end{pmatrix} .
        \]
        Check: $\mathbf{g}(\mathbf{f}(x,y))=(x+xy,\;xy^{2})$, Jacobian
        $\begin{pmatrix}1+y&x\\y^{2}&2xy\end{pmatrix}$, which is
        $\begin{pmatrix}3&1\\4&4\end{pmatrix}$ at $(1,2)$. \checkmark\
        \emph{$\mathbf{f}\circ\mathbf{g}\colon\R^{3}\to\R^{3}$} at
        $(1,1,1)$: $\mathbf{g}(1,1,1)=(2,1)$, so
        $J_{\mathbf{f}}(2,1)=\begin{pmatrix}1&0\\0&1\\1&2\end{pmatrix}$ and
        $J_{\mathbf{g}}(1,1,1)=\begin{pmatrix}1&0&1\\0&1&1\end{pmatrix}$:
        \[
          J_{\mathbf{f}\circ\mathbf{g}}=\begin{pmatrix}1&0\\0&1\\1&2\end{pmatrix}
          \begin{pmatrix}1&0&1\\0&1&1\end{pmatrix}
          =\begin{pmatrix}1&0&1\\0&1&1\\1&2&3\end{pmatrix} .
        \]
        Check: $\mathbf{f}(\mathbf{g})=\bigl(u+w,\;vw,\;(u+w)vw\bigr)$; the
        last row is $\bigl(vw,\;(u+w)w,\;vw+(u+w)v\bigr)=(1,2,3)$.
        \checkmark
  \item $\grad F=(2x,2y,2z)$ and $\mathbf{r}'(t)=(-\sin t,\cos t,1)$:
        \[
          \dd{}{t}F(\mathbf{r}(t))=2\cos t(-\sin t)+2\sin t\cos t+2t\cdot1
          =2t .
        \]
        Check: $F(\mathbf{r}(t))=\cos^{2}t+\sin^{2}t+t^{2}=1+t^{2}$.
        \checkmark
  \item $|\mathbf{h}|=3$. $\grad f=(\ln y,\;x/y+z^{2},\;2yz)=(0,3,2)$, so
        $\grad f\cdot\mathbf{h}=0+6-4=2$ and $D_{\mathbf{u}}f=\tfrac23$.
        $H=\begin{pmatrix}0&1/y&0\\1/y&-x/y^{2}&2z\\0&2z&2y\end{pmatrix}
        =\begin{pmatrix}0&1&0\\1&-2&2\\0&2&2\end{pmatrix}$ and
        \[
          \mathbf{h}^{T}H\mathbf{h}=2\cdot1\cdot(1)(2)-2\cdot4
          +2\cdot2\cdot(2)(-2)+2\cdot4=4-8-16+8=-12,
        \]
        so $D^{2}_{\mathbf{u}}f=-\tfrac{12}{9}=-\tfrac43$. \emph{Profile.}
        $\varphi(t)=(2+t)\ln(1+2t)+(1+2t)(1-2t)^{2}$, and
        $(1+2t)(1-2t)^{2}=1-2t-4t^{2}+8t^{3}$. The log part has derivative
        $\ln(1+2t)+\dfrac{2(2+t)}{1+2t}$, equal to $4$ at $0$, and second
        derivative $\dfrac{2}{1+2t}+\dfrac{2(1+2t)-4(2+t)}{(1+2t)^{2}}$,
        equal to $2-6=-4$. So $\varphi'(0)=4-2=2$ and
        $\varphi''(0)=-4-8=-12$. \checkmark\ Concave in this direction.
  \item $|\mathbf{h}|=3$. $\grad f=\mathrm{e}^{yz}(1,\;xz,\;xy)=(1,2,0)$;
        $D_{\mathbf{u}}f=\tfrac{2+2+0}{3}=\tfrac43$. Second partials at $P$
        ($\mathrm{e}^{yz}=1$): $f_{xx}=0$, $f_{xy}=z\mathrm{e}^{yz}=2$,
        $f_{xz}=y\mathrm{e}^{yz}=0$, $f_{yy}=xz^{2}\mathrm{e}^{yz}=4$,
        $f_{yz}=x(1+yz)\mathrm{e}^{yz}=1$, $f_{zz}=xy^{2}\mathrm{e}^{yz}=0$.
        \[
          \mathbf{h}^{T}H\mathbf{h}=2\cdot2\cdot(2)(1)+4\cdot1+2\cdot1\cdot(1)(2)
          =8+4+4=16,\qquad D^{2}_{\mathbf{u}}f=\tfrac{16}{9} .
        \]
        \emph{Profile.} $\varphi(t)=(1+2t)\,\mathrm{e}^{t(2+2t)}
        =(1+2t)\,\mathrm{e}^{2t+2t^{2}}$. With $E=\mathrm{e}^{2t+2t^{2}}$,
        $E'=(2+4t)E$:
        $\varphi'=\bigl(2+(1+2t)(2+4t)\bigr)E$, so $\varphi'(0)=4$, and
        $\varphi''=\bigl(4+8t+4+8t\bigr)E+\bigl(2+(1+2t)(2+4t)\bigr)(2+4t)E$,
        so $\varphi''(0)=8+4\cdot2=16$. \checkmark\ Convex in this direction.
  \item $|\mathbf{h}|=\sqrt3$. $\grad f=(yz+\tfrac1x,\;xz,\;xy+\tfrac1z)
        =(2,1,2)$, $\grad f\cdot\mathbf{h}=2-1+2=3$,
        $D_{\mathbf{u}}f=3/\sqrt3=\sqrt3$.
        $H=\begin{pmatrix}-1/x^{2}&z&y\\z&0&x\\y&x&-1/z^{2}\end{pmatrix}
        =\begin{pmatrix}-1&1&1\\1&0&1\\1&1&-1\end{pmatrix}$, and
        \[
          \mathbf{h}^{T}H\mathbf{h}=-1+0-1+2\bigl[(1)(-1)+(1)(1)+(-1)(1)\bigr]
          =-2-2=-4,\qquad D^{2}_{\mathbf{u}}f=-\tfrac43 .
        \]
        \emph{Profile.} $\varphi(t)=(1+t)(1-t)(1+t)+\ln\bigl((1+t)^{2}\bigr)
        =1+t-t^{2}-t^{3}+2\ln(1+t)$. $\varphi'(t)=1-2t-3t^{2}+\tfrac{2}{1+t}$,
        $\varphi'(0)=3$; $\varphi''(t)=-2-6t-\tfrac{2}{(1+t)^{2}}$,
        $\varphi''(0)=-4$. \checkmark\ Concave in this direction.
  \item At $(1,2)$: $\grad f=(2xy,x^{2})=(4,1)$ and
        $H=\begin{pmatrix}2y&2x\\2x&0\end{pmatrix}=\begin{pmatrix}4&2\\2&0
        \end{pmatrix}$. With $\mathbf{h}=(3,4)$:
        $\grad f\cdot\mathbf{h}=12+4=16$ and
        $\mathbf{h}^{T}H\mathbf{h}=4\cdot9+2\cdot2\cdot12+0=84$. Profile:
        $\varphi(t)=(1+3t)^{2}(2+4t)=2+16t+42t^{2}+36t^{3}$, so
        $\varphi'(0)=16$ and $\varphi''(0)=84$. \checkmark\ Practice
        problem 15 found $D_{\mathbf{u}}f=\tfrac{16}{5}$ and
        $D^{2}_{\mathbf{u}}f=\tfrac{84}{25}$ with $\mathbf{u}=\mathbf{h}/5$:
        moving along $\mathbf{h}$ covers $5$ units of distance per unit of
        $t$, so first derivatives scale by $|\mathbf{h}|=5$ and second
        derivatives by $|\mathbf{h}|^{2}=25$. The sign, and so the
        convexity verdict, is the same either way.
  \item $\grad f=(3x^{2}-3,\;2y)$, $H=\begin{pmatrix}6x&0\\0&2\end{pmatrix}$.
        At $Q=(-1,1)$: $\grad f=(0,2)$, $|\mathbf{h}|=5$,
        $\mathbf{u}=(\tfrac35,\tfrac45)$, so $D_{\mathbf{u}}f=\tfrac85$;
        $H=\begin{pmatrix}-6&0\\0&2\end{pmatrix}$ and
        $D^{2}_{\mathbf{u}}f=\dfrac{-6\cdot9+2\cdot16}{25}=-\dfrac{22}{25}<0$:
        rising, but concave in this direction at $Q$. At $(1,0)$ the
        gradient is $\mathbf{0}$ and $H=\begin{pmatrix}6&0\\0&2\end{pmatrix}$
        is PD: a strict local minimum, with positive curvature in every
        direction. Negative curvature somewhere rules out convexity,
        positive curvature somewhere rules out concavity: $f$ is
        \textbf{neither}.
  \item $f_{x}=4x^{3}-4y$, $f_{y}=-4x+2y$. From $f_{y}=0$: $y=2x$; then
        $4x^{3}=8x$, so $x=0$ or $x^{2}=2$. Critical points $(0,0)$ and
        $\pm(\sqrt2,2\sqrt2)$. $H=\begin{pmatrix}12x^{2}&-4\\-4&2\end{pmatrix}$.
        At $\pm(\sqrt2,2\sqrt2)$: $\begin{pmatrix}24&-4\\-4&2\end{pmatrix}$,
        $D_{1}=24$, $D_{2}=48-16=32$: PD, strict local minima, with value
        $4-16+8=-4$. At $(0,0)$: $\begin{pmatrix}0&-4\\-4&2\end{pmatrix}$,
        $D_{2}=-16$: a saddle. There, for $\mathbf{u}=(1,1)/\sqrt2$,
        $D^{2}_{\mathbf{u}}f=\tfrac{0-8+2}{2}=-3<0$ (concave that way), and
        for $\mathbf{u}=(0,1)$, $D^{2}_{\mathbf{u}}f=2>0$ (convex that way).
        Globally: negative curvature at the origin and positive definite
        Hessians at the minima, so $f$ is neither convex nor concave.
  \item $\grad f=(\ln y,\;x/y)=(0,1)$ and
        $H=\begin{pmatrix}0&1/y\\1/y&-x/y^{2}\end{pmatrix}=
        \begin{pmatrix}0&1\\1&-1\end{pmatrix}$ at $P$.
        \begin{itemize}[leftmargin=1.2em,itemsep=0.2ex]
          \item $\mathbf{u}=(1,0)$: $D_{\mathbf{u}}f=0$,
                $D^{2}_{\mathbf{u}}f=0$. In fact $f(x,1)=0$: flat, no
                verdict beyond ``linear''.
          \item $\mathbf{u}=(0,1)$: $D_{\mathbf{u}}f=1$,
                $D^{2}_{\mathbf{u}}f=-1$: concave in this direction.
          \item $\mathbf{u}=(1,1)/\sqrt2$: $D_{\mathbf{u}}f=\tfrac{1}{\sqrt2}$,
                $D^{2}_{\mathbf{u}}f=\tfrac12(0+2-1)=\tfrac12$: convex in
                this direction.
        \end{itemize}
        Both signs at one point already decide it: $f$ is neither convex
        nor concave. Indeed $\det H=-1/y^{2}<0$ everywhere.
  \item (a) True: along the contour $f$ is constant; equivalently the
        tangent $\mathbf{u}$ is perpendicular to $\grad f$, so
        $D_{\mathbf{u}}f=0$. (b) True when $\grad f(P)\neq\mathbf{0}$: the
        unit vector $-\grad f/|\grad f|$ gives the most negative rate,
        $-|\grad f|$. (c) False: $f=xy$ has $\grad f(0,0)=\mathbf{0}$ at a
        saddle. (d) True: the largest directional derivative is
        $|\grad f(P)|$. (e) False: it is perpendicular (normal) to the level
        curve.
  \item She walks a short way in the direction in which she neither climbs
        nor descends: she is following the contour line. The gradient
        points at right angles to that path, on the uphill side. Its
        direction is the steepest one; in a direction at angle $\theta$ to
        it the slope is $|\grad f|\cos\theta$: smaller, zero along the
        contour, and negative once she turns downhill.
  \item (a) True: $\grad f\cdot\mathbf{u}_{1}=\grad f\cdot\mathbf{u}_{2}=0$
        with $\mathbf{u}_{1},\mathbf{u}_{2}$ independent forces $\grad f$ to
        be perpendicular to the whole plane, hence $\mathbf{0}$. With three
        variables two directions are not enough: $f(x,y,z)=z$ has
        $D_{\mathbf{e}_{1}}f=D_{\mathbf{e}_{2}}f=0$ but $\grad f=(0,0,1)$.
        (b) True:
        $D_{-\mathbf{u}}f=\grad f\cdot(-\mathbf{u})=-D_{\mathbf{u}}f$.
        (c) False: by Cauchy--Schwarz,
        $D_{\mathbf{u}}f=\grad f\cdot\mathbf{u}\le|\grad f|\,|\mathbf{u}|=4$.
  \item $Y=10\cdot64\cdot216=138240$. In output-ratio form
        $Y_{K}=0.4\,Y/K=\tfrac{27}{16}=1.6875$ and
        $Y_{L}=0.6\,Y/L=\tfrac{32}{3}\approx10.667$. Best ratio
        $\mathrm{d}K:\mathrm{d}L=Y_{K}:Y_{L}=81:512$.
        $|\grad Y|=\sqrt{1.6875^{2}+10.667^{2}}\approx10.7993$, so
        \[
          (\mathrm{d}K,\mathrm{d}L)=\frac{100}{10.7993}(1.6875,\;10.667)
          \approx(15.63,\;98.77),
          \qquad \mathrm{d}Y\approx100\cdot10.7993\approx1079.9 .
        \]
        Exactly, $Y(32783.63,\;7874.77)-Y\approx1077.5$. A budget
        $\mathrm{d}K+\mathrm{d}L=100$ is maximised in the linear estimate by
        spending everything on labour, the input with the larger marginal
        product: $\mathrm{d}Y\approx1066.7$ (exactly $1063.9$). That is
        less, but it also uses less: the steepest-ascent vector has
        $\mathrm{d}K+\mathrm{d}L\approx114.4$.
  \item $Y=\sqrt{400\cdot900}=600$, $Y_{K}=\dfrac{Y}{2K}=0.75$,
        $Y_{L}=\dfrac{Y}{2L}=\tfrac13$. Ratio $\mathrm{d}K:\mathrm{d}L=
        0.75:\tfrac13=9:4$. $|\grad Y|=\sqrt{0.5625+0.1111}\approx0.82074$,
        so $(\mathrm{d}K,\mathrm{d}L)\approx\tfrac{20}{0.82074}(0.75,\;0.3333)
        \approx(18.28,\;8.12)$ and $\mathrm{d}Y\approx20\cdot0.82074\approx
        16.41$; exactly $16.32$. A budget $\mathrm{d}K+\mathrm{d}L=20$ would
        go entirely into capital: $\mathrm{d}Y\approx15$ (exactly $14.82$).
\end{solutions}

\chapter{Inverse and Implicit Function Theorems, Jacobians, Homogeneity, Euler}

When a relation between variables is given implicitly, the Implicit
Function Theorem (IFT) extracts derivatives without solving for one
variable in terms of the other, exactly the situation behind
indifference curves and isoquants. The Jacobian generalises the derivative
to vector-valued maps; its determinant is the local area / volume scaling
factor, and the Inverse Function Theorem turns that determinant into a test
for whether the map can be undone near a point --- and, when it can, gives
the derivative of the inverse without ever computing the inverse.
Homogeneous functions and Euler's theorem then give a clean
structural identity that ties marginal products to total output,
underpinning constant-returns-to-scale models.

\section{Core Rules}

\begin{rulebox}{IFT, One Equation in Two Unknowns}
Suppose $F:\R^{2}\to\R$ is $C^{1}$ and $F(x_{0},y_{0})=0$ with
$F_{y}(x_{0},y_{0})\neq 0$. Then near $(x_{0},y_{0})$ there is a $C^{1}$
function $y=g(x)$ satisfying $F(x,g(x))=0$, and
\[
  \dd{y}{x} \;=\; -\frac{F_{x}}{F_{y}} .
\]
\end{rulebox}

\begin{rulebox}{What the condition $F_{y}\neq0$ rules out}
The theorem is local and its hypothesis is sharp. Where $F_{y}=0$ the
conclusion genuinely fails, and the failure is visible in the picture of the
curve $F=0$ (Chapter~5):
\begin{itemize}[leftmargin=1.6em,itemsep=0.3ex]
  \item $F_{y}=0$, $F_{x}\neq0$: a \textbf{vertical tangent}. No $y(x)$
        exists near the point, but $x=x(y)$ does, with
        $\dd{x}{y}=-F_{y}/F_{x}=0$. The theorem is symmetric in the two
        variables: at least one of the two representations survives as long
        as $\grad F\neq\mathbf{0}$.
  \item $F_{x}=F_{y}=0$: a \textbf{singular point}, typically a
        self-intersection where two branches with different slopes meet.
        Neither representation exists, and $-F_{x}/F_{y}$ degenerates to
        $0/0$.
\end{itemize}
So the working rule is: the level set of $F$ is a well-behaved curve, with a
well-defined tangent, exactly where $\grad F\neq\mathbf{0}$. This is the
same gradient condition that reappears as the \emph{constraint
qualification} in the Lagrange and Kuhn--Tucker chapters, where it rules out
precisely these kinks and crossings in the constraint set.
\end{rulebox}

\begin{rulebox}{IFT, Multivariate}
If $F(x_{1},\ldots,x_{n},y)=0$ with $F_{y}\neq 0$, then
$y=g(x_{1},\ldots,x_{n})$ locally and
\[
  \pd{y}{x_{i}} \;=\; -\frac{F_{x_{i}}}{F_{y}} .
\]
\end{rulebox}

\begin{rulebox}{MRS and MRTS via IFT}
For $U(x,y)=c$ set $F=U(x,y)-c$. Then $F_{x}=U_{x}$, $F_{y}=U_{y}$, so
\[
  \dd{y}{x}\bigg|_{U=c} = -\frac{U_{x}}{U_{y}}\quad\text{(signed)},\qquad
  \mathrm{MRS} = \frac{U_{x}}{U_{y}}\quad\text{(unsigned)} .
\]
Likewise on $Q(K,L)=Q_{0}$: $\dd{L}{K}=-Q_{K}/Q_{L}$,
$\mathrm{MRTS}=Q_{K}/Q_{L}$.
\end{rulebox}

\begin{rulebox}{Jacobian Matrix}
For $\mathbf{f}:\R^{n}\to\R^{m}$,
$\mathbf{f}(x_{1},\ldots,x_{n})=(f_{1},\ldots,f_{m})$, the \textbf{Jacobian}
is
\[
  J_{\mathbf{f}} \;=\;
  \begin{pmatrix}
    \partial f_{1}/\partial x_{1} & \cdots & \partial f_{1}/\partial x_{n}\\
    \vdots & \ddots & \vdots\\
    \partial f_{m}/\partial x_{1} & \cdots & \partial f_{m}/\partial x_{n}
  \end{pmatrix} \;\in\; \R^{m\times n} .
\]
For $m=n$, $\det J_{\mathbf{f}}$ is the local volume scaling factor. Whether
it vanishes is exactly what decides local invertibility, which is the next
three boxes.
\end{rulebox}

\begin{rulebox}{Inverse Function Theorem, one variable}
Let $f$ be continuously differentiable near $x_{0}$ and set $y_{0}=f(x_{0})$.
If
\[
  f'(x_{0})\neq0 ,
\]
then $f$ is invertible on some open interval around $x_{0}$, the inverse is
differentiable at $y_{0}$, and
\[
  \boxed{\;\bigl(f^{-1}\bigr)'(y_{0})=\frac{1}{f'(x_{0})}
        =\frac{1}{f'\bigl(f^{-1}(y_{0})\bigr)}\;}
\]
The formula is the chain rule read backwards. Differentiating the identity
$f\bigl(f^{-1}(y)\bigr)=y$ gives
$f'\bigl(f^{-1}(y)\bigr)\cdot\bigl(f^{-1}\bigr)'(y)=1$, and dividing
finishes it --- which is legitimate precisely because $f'\neq0$.

Note which point goes where: the derivative of the \emph{inverse} is
evaluated at $y_{0}$, the derivative of $f$ at $x_{0}$. Geometrically the
graph of $f^{-1}$ is the graph of $f$ mirrored in $y=x$, and mirroring turns
a slope into its reciprocal --- a vertical tangent for $f$ becoming a
horizontal one for $f^{-1}$, which is why $f'=0$ has to be excluded.

The derivatives of $\arcsin$, $\arccos$ and $\arctan$ in Chapter~30 are this
rule applied to $\sin$, $\cos$ and $\tan$; nothing else is needed to get
them.
\end{rulebox}

\begin{rulebox}{Inverse Function Theorem, several variables}
Let $\mathbf{f}:\R^{n}\to\R^{n}$ be continuously differentiable near
$\mathbf{x}_{0}$ and set $\mathbf{y}_{0}=\mathbf{f}(\mathbf{x}_{0})$. If
\[
  \det J_{\mathbf{f}}(\mathbf{x}_{0})\neq0 ,
\]
then $\mathbf{f}$ maps some open neighbourhood of $\mathbf{x}_{0}$
one-to-one onto an open neighbourhood of $\mathbf{y}_{0}$, the local inverse
is continuously differentiable, and its Jacobian is the \textbf{matrix
inverse} of the original:
\[
  \boxed{\;J_{\mathbf{f}^{-1}}(\mathbf{y}_{0})
        =\bigl[J_{\mathbf{f}}(\mathbf{x}_{0})\bigr]^{-1}\;}
\]
For $n=1$ the matrix is $1\times1$ and this is the previous box.
The condition $\det J\neq0$ is the exact analogue of $f'\neq0$: it says the
linear approximation at the point is invertible, and the theorem promises
that the map itself inherits that property nearby.

Note that the two named theorems of this chapter are the same statement
twice. Applying the Implicit Function Theorem to
$\mathbf{F}(\mathbf{x},\mathbf{y})=\mathbf{f}(\mathbf{x})-\mathbf{y}$
--- solving for $\mathbf{x}$ in terms of $\mathbf{y}$ --- has the condition
$\det J_{\mathbf{f}}\neq0$ and delivers exactly the inverse.
\end{rulebox}

\begin{rulebox}{Local, and only local}
Both theorems are \textbf{local}, and both directions of the obvious
over-reading are false.

\emph{A non-vanishing derivative everywhere does not give a global
inverse.} The map
$\mathbf{f}(x,y)=(\mathrm{e}^{x}\cos y,\ \mathrm{e}^{x}\sin y)$ has
$\det J=\mathrm{e}^{2x}>0$ at every point, yet
$\mathbf{f}(x,y)=\mathbf{f}(x,y+2\pi)$: it is invertible near each point and
nowhere invertible globally.

\emph{A vanishing derivative does not rule out an inverse.} $f(x)=x^{3}$ has
$f'(0)=0$, and is nevertheless a bijection of $\R$. What fails is not
invertibility but \emph{differentiability of the inverse}: $f^{-1}(y)=
y^{1/3}$ has a vertical tangent at $0$, exactly as the reciprocal formula
predicts.

So the condition is sufficient for a local inverse, necessary for a
\emph{differentiable} one, and silent about global behaviour.
\end{rulebox}

\begin{rulebox}{Homogeneous Functions}
$f:\R^{n}\to\R$ is \textbf{homogeneous of degree $k$} if for all $t>0$ and
all $\mathbf{x}$ in the domain,
\[
  f(t\mathbf{x}) \;=\; t^{k}f(\mathbf{x}) .
\]
Constant returns to scale (CRS) means $k=1$.
\end{rulebox}

\begin{rulebox}{Euler's Theorem}
If $f$ is differentiable and homogeneous of degree $k$ on an open cone,
then
\[
  \sum_{i=1}^{n} x_{i}\pd{f}{x_{i}} \;=\; k\,f(\mathbf{x}) .
\]
Two-variable case: $xf_{x}+yf_{y}=kf$. \textbf{Verification rule.} Always
compute the LHS $xf_{x}+yf_{y}$ and the RHS $kf$ \emph{separately}, then
check that they agree.
\end{rulebox}

\section{Worked Examples}

\begin{worked}{$F(x,y)=x^{2}y+y^{3}-10=0$ at $(1,2)$: does it define $y(x)$, and what is $y'$? Then verify Euler's theorem on $f(K,L)=K^{0.3}L^{0.7}$}
\wstep{Step 1. Check the point lies on the curve.}
$1^{2}\cdot2+2^{3}-10=2+8-10=0$. \checkmark

\wstep{Step 2. Test the hypothesis of the Implicit Function Theorem.}
The theorem needs $F_{y}\neq0$ at the point:
\[
  F_{y}=x^{2}+3y^{2},
  \qquad
  F_{y}(1,2)=1+12=13\neq0 .
\]
So there is a neighbourhood of $x=1$ on which the equation defines a unique
differentiable function $y(x)$ with $y(1)=2$.
\wtrap{The condition is on $F_{y}$ --- the derivative with respect to the
variable you are \textbf{solving for}. If you want $x$ as a function of $y$
instead, the condition is $F_{x}\neq0$. Testing the wrong partial can
license a conclusion that is false: at a point with $F_{y}=0$ the curve is
vertical and no function $y(x)$ exists there.}

\wstep{Step 3. Apply the formula.}
\[
  F_{x}=2xy,
  \qquad
  \dd{y}{x}=-\frac{F_{x}}{F_{y}}=-\frac{2xy}{x^{2}+3y^{2}},
  \qquad
  y'(1)=-\frac{4}{13}\approx-0.308 .
\]
\wtrap{The minus sign is part of the formula, and it is dropped constantly.
It comes from the chain rule applied to the identity
$F\bigl(x,y(x)\bigr)=0$: $F_{x}+F_{y}y'=0$, so $y'=-F_{x}/F_{y}$. Derive it
once and the sign stops being something to remember.}

\wstep{Step 4. Check by differentiating implicitly, from scratch.}
\[
  \frac{\mathrm{d}}{\mathrm{d}x}\bigl[x^{2}y+y^{3}-10\bigr]=0
  \;\Longrightarrow\;
  2xy+x^{2}y'+3y^{2}y'=0
  \;\Longrightarrow\;
  y'=\frac{-2xy}{x^{2}+3y^{2}} . \checkmark
\]
At $(1,2)$: $4+y'+12y'=0$, so $y'=-4/13$. The theorem's formula is not a new
technique --- it is Chapter~5 done once in general.

\wstep{Step 5. A numerical confirmation.}
Take $x=1.01$ and solve $1.0201y+y^{3}=10$ numerically: $y\approx1.996908$.
The linear prediction is $2-0.30769\times0.01=1.996923$. \checkmark\ The two
agree to four decimals over a step of $0.01$; the residual $1.5\times10^{-5}$
is the second-order term, and it should shrink by a factor of $100$ if the
step is halved twice --- which is the signature of a first-order
approximation.

\wstep{Step 6. Where this is used: comparative statics without solving.}
An equilibrium condition $F(p,\alpha)=0$ relates a price $p$ to a parameter
$\alpha$, usually with no closed-form solution for $p$. The theorem still
gives the response:
\[
  \dd{p}{\alpha}=-\frac{F_{\alpha}}{F_{p}} ,
\]
and its \emph{sign} --- which is often all that is asked --- follows from
the signs of two partials. This is the engine of every comparative-statics
result in economics.

\wstep{Step 7. Homogeneity: check the degree by scaling.}
$f$ is homogeneous of degree $r$ if $f(tK,tL)=t^{r}f(K,L)$ for all $t>0$.
For $f=K^{0.3}L^{0.7}$:
\[
  f(tK,tL)=(tK)^{0.3}(tL)^{0.7}=t^{0.3+0.7}K^{0.3}L^{0.7}=t\,f(K,L),
\]
so $r=1$: constant returns to scale. For Cobb--Douglas the degree is simply
the \emph{sum of the exponents}.
\wtrap{Scale \emph{every} argument by the same $t$, and require the identity
for all $t>0$. $f(K,L)=K^{0.3}L^{0.7}+5$ is \textbf{not} homogeneous --- the
constant does not scale --- and Euler's theorem fails for it. Homogeneity is
a strong property, not a formality.}

\wstep{Step 8. Euler's theorem, verified.}
For $f$ homogeneous of degree $r$,
\[
  K f_{K}+L f_{L}=r\,f .
\]
Here $f_{K}=0.3K^{-0.7}L^{0.7}$ and $f_{L}=0.7K^{0.3}L^{-0.3}$, so
\[
  Kf_{K}+Lf_{L}=0.3K^{0.3}L^{0.7}+0.7K^{0.3}L^{0.7}
  =(0.3+0.7)f=1\cdot f . \checkmark
\]

\wstep{Step 9. Read the economics --- this is the product-exhaustion theorem.}
Under constant returns to scale ($r=1$), paying each factor its marginal
product exhausts output exactly:
\[
  \underbrace{K\cdot MP_{K}}_{\text{capital's share}}
  +\underbrace{L\cdot MP_{L}}_{\text{labour's share}}
  =\underbrace{Q}_{\text{total output}} .
\]
Nothing is left over and nothing is missing. With increasing returns
($r>1$) the factor payments would exceed output, which is why perfect
competition and increasing returns cannot coexist --- a theorem of this
chapter, not a modelling assumption.
\end{worked}

\begin{worked}{The inverse theorem both ways: $f(x)=x^{5}+2x+1$ at $y=4$, and $\mathbf{f}(x,y)=(x^{2}-y^{2},\,2xy)$ at $(1,1)$}
\wstep{Step 1. One variable: check the hypothesis before anything else.}
\[
  f(x)=x^{5}+2x+1,
  \qquad
  f'(x)=5x^{4}+2\ \ge 2>0\quad\text{for every }x .
\]
The derivative never vanishes, so $f$ is strictly increasing and invertible
on all of $\R$, and its inverse is differentiable everywhere.

\wstep{Step 2. Locate the point --- in the \emph{domain}.}
We are asked about $y=4$, so we need the $x$ with $f(x)=4$:
\[
  x^{5}+2x+1=4
  \quad\Longrightarrow\quad
  x=1\ \ (1+2+1=4). \checkmark
\]
\wtrap{The formula needs $x_{0}=f^{-1}(y_{0})$, not $y_{0}$. Substituting
$4$ into $f'$ gives $5\cdot4^{4}+2=1282$ and the wrong answer $1/1282$.
Find the pre-image first; in an examination it is always a value you can
guess and verify, because no one can invert a quintic either.}

\wstep{Step 3. Apply the formula.}
\[
  \bigl(f^{-1}\bigr)'(4)=\frac{1}{f'(1)}=\frac{1}{5+2}=\boxed{\tfrac17}
\]
Note what was \emph{not} done: $f^{-1}$ was never computed. It cannot be ---
a general quintic has no closed-form root --- and that is exactly the point
of the theorem. It hands you a derivative of a function you cannot write
down.

\wstep{Step 4. Sanity-check it numerically.}
$f(1.01)=1.01^{5}+2.02+1\approx4.0710$, so a step of $0.0710$ in $y$ answers
to a step of $0.01$ in $x$:
\[
  \frac{\Delta x}{\Delta y}\approx\frac{0.01}{0.0710}\approx0.141
  \approx\tfrac17=0.1429 . \checkmark
\]

\wstep{Step 5. Several variables: build the Jacobian.}
For $\mathbf{f}(x,y)=(x^{2}-y^{2},\ 2xy)$,
\[
  J_{\mathbf{f}}=\begin{pmatrix}2x&-2y\\ 2y&2x\end{pmatrix},
  \qquad
  \det J_{\mathbf{f}}=4x^{2}+4y^{2}=4(x^{2}+y^{2}).
\]
This vanishes only at the origin, so the theorem applies at every other
point.

\wstep{Step 6. Evaluate at $(1,1)$ and invert the matrix.}
$\mathbf{f}(1,1)=(1-1,\,2)=(0,2)$, and
\[
  J_{\mathbf{f}}(1,1)=\begin{pmatrix}2&-2\\2&2\end{pmatrix},
  \qquad
  \det=4+4=8\neq0 .
\]
So a local inverse exists near $(0,2)$, and its Jacobian is the matrix
inverse:
\[
  J_{\mathbf{f}^{-1}}(0,2)
  =\begin{pmatrix}2&-2\\2&2\end{pmatrix}^{-1}
  =\frac18\begin{pmatrix}2&2\\-2&2\end{pmatrix}
  =\boxed{\begin{pmatrix}1/4&1/4\\-1/4&1/4\end{pmatrix}}
\]
Check by multiplying: $\begin{pmatrix}2&-2\\2&2\end{pmatrix}
\begin{pmatrix}1/4&1/4\\-1/4&1/4\end{pmatrix}
=\begin{pmatrix}1&0\\0&1\end{pmatrix}$. \checkmark
\wtrap{Invert the \emph{matrix}; do not take reciprocals entry by entry.
$\bigl[J\bigr]^{-1}$ is a matrix inverse, and the $2\times2$ rule swaps the
diagonal and negates the off-diagonal before dividing by the determinant.
Entrywise reciprocals would give $\begin{pmatrix}1/2&-1/2\\1/2&1/2
\end{pmatrix}$, which is not the inverse of anything relevant.}

\wstep{Step 7. Now show the inverse is only local.}
$\det J\neq0$ holds at every point except the origin, which invites the
conclusion that $\mathbf{f}$ is invertible on $\R^{2}\setminus\{0\}$. It is
not:
\[
  \mathbf{f}(-1,-1)=\bigl((-1)^{2}-(-1)^{2},\ 2(-1)(-1)\bigr)=(0,2)
  =\mathbf{f}(1,1).
\]
Two different points share an image, so no global inverse exists. In general
$\mathbf{f}(-x,-y)=\mathbf{f}(x,y)$: the map is two-to-one everywhere.
\wtrap{``$\det J\neq0$ everywhere'' gives a local inverse \emph{at each
point} and says nothing about fitting them together into one function. This
is the chapter's sharpest trap, and the map above is the standard
counterexample --- it is squaring, in disguise, and squaring is two-to-one.}

\wstep{Step 8. What the two halves have in common.}
\begin{center}\footnotesize
\begin{tabular}{lll}
  \hline
   & one variable & several variables\\
  \hline
  hypothesis & $f'(x_{0})\neq0$ & $\det J(\mathbf{x}_{0})\neq0$\\
  conclusion & local inverse, differentiable & the same\\
  derivative & $1/f'(x_{0})$ & $\bigl[J(\mathbf{x}_{0})\bigr]^{-1}$\\
  evaluated at & $y_{0}=f(x_{0})$ & $\mathbf{y}_{0}=\mathbf{f}(\mathbf{x}_{0})$\\
  says nothing about & global behaviour & global behaviour\\
  \hline
\end{tabular}
\end{center}
The right-hand column is the left-hand one with ``non-zero number''
replaced by ``invertible matrix'' --- which is what $\det\neq0$ means, and
why Volume~3 had to come first.
\end{worked}

\section{Practice Problems}

\begin{problems}
  \item Use IFT to find $\dd{y}{x}$ for $x^{3}+y^{3}=6xy$.
  \item Verify homogeneity and find the degree of
        $f(x,y)=x^{4}+5x^{2}y^{2}+y^{4}$.
  \item Compute the Jacobian and its determinant for
        $\mathbf{f}(x,y)=(3x-2y,\;x+4y)$.
  \item For $U(x,y)=x+2\ln y$, find the signed slope and unsigned MRS
        using IFT.
  \item Verify Euler's theorem for $f(x,y)=x^{2}+5xy-2y^{2}$ (degree $2$).
  \item Use IFT to find $\dd{y}{x}$ for $x^{2}y+xy^{3}=10$.
  \item Verify homogeneity and find the degree of
        $f(x,y)=\dfrac{x^{3}+y^{3}}{x-y}$.
  \item Compute the Jacobian and its determinant for
        $\mathbf{f}(x,y)=(x^{2}+y^{2},\;x^{2}-y^{2})$.
  \item For $U(x,y)=(x+1)(y+2)$, find the signed slope and unsigned MRS.
  \item Verify Euler's theorem for $Q(K,L)=4K^{1/4}L^{3/4}$ (degree $1$).
  \item Use IFT to find $\dd{y}{x}$ for $\mathrm{e}^{x+y}=x+y+1$.
  \item Verify homogeneity and find the degree of $f(x,y)=\sqrt{xy}$.
  \item Compute the Jacobian and its determinant for
        $\mathbf{f}(x,y)=(x\mathrm{e}^{y},\;y\mathrm{e}^{x})$.
  \item For $U(x,y)=x^{a}y^{b}$ ($a,b>0$), find the signed slope and MRS
        via IFT.
  \item Verify Euler's theorem for $Q(K,L)=K^{1/2}L^{1/2}+K+L$
        (degree $1$).
  \item Use IFT to find $\dd{y}{x}$ for $\ln(x^{2}+y)=x-y$.
  \item Verify homogeneity for
        $f(x,y,z)=\dfrac{xy+yz+zx}{x+y+z}$ and find the degree.
  \item Compute the Jacobian and its determinant for the cylindrical map
        \[
          \mathbf{f}(\rho,\theta,z)=(\rho\cos\theta,\;\rho\sin\theta,\;z).
        \]
  \item For the Leontief $U(x,y)=\min\{2x,3y\}$, discuss MRS on each
        smooth piece and at the kink.
  \item Verify Euler's theorem for the harmonic mean
        $f(x,y)=(x^{-1}+y^{-1})^{-1}$ (degree $1$).
  \item Use IFT to find $\dd{y}{x}$ for $x\cos y+y\cos x=1$.
  \item Verify homogeneity and find the degree of
        $f(x,y)=x\ln(y/x)+y\ln(x/y)$.
  \item Compute the Jacobian of
        $\mathbf{f}(x,y)=(x+y,\;xy,\;x^{2}+y^{2})$ and identify where the
        rank drops.
  \item For $Q(K,L)=5K+3L$, find the MRTS.
  \item Verify Euler's theorem for
        $Q(K,L)=\bigl(\alpha K^{\rho}+(1-\alpha)L^{\rho}\bigr)^{1/\rho}$
        (CES, degree $1$).
  \item Use IFT to find $\dd{y}{x}$ for $x^{y}=y^{x}$ with $x,y>0$.
  \item Verify whether $f(x,y)=x^{2}+xy+y$ is homogeneous; if so, find the
        degree.
  \item Compute the Jacobian and determinant of
        $\mathbf{f}(x,y)=(\mathrm{e}^{x}\cos y,\;\mathrm{e}^{x}\sin y)$.
        Interpret.
  \item For $Q(K,L)=K^{2/3}L^{1/3}$, find the MRTS via IFT.
  \item Verify Euler's theorem for $f(x,y,z)=x^{3}+y^{3}+z^{3}+3xyz$
        (degree $3$).
  \item Use IFT to find $\dd{y}{x}$ for $\sqrt{x}+\sqrt{y}+\sqrt{xy}=5$.
  \item Verify homogeneity and find the degree of
        $f(x,y)=\dfrac{x^{2}+y^{2}}{xy}$.
  % VERIFY: right-hand entries of the spherical map obscured by a printer's
  % page-turn mark in the scan; reconstructed from the standard spherical
  % parametrization, and confirmed by the worked solution below.
  \item Compute the Jacobian and determinant of the spherical map
        \[
          \mathbf{f}(\rho,\theta,\phi)=(\rho\sin\phi\cos\theta,\;
          \rho\sin\phi\sin\theta,\;\rho\cos\phi).
        \]
  \item For $U(x,y)=ax+by$ ($a,b>0$, perfect substitutes), find the MRS.
  \item Verify Euler's theorem for $f(x,y)=\dfrac{x^{3}+y^{3}}{x+y}$
        (degree $2$).
  \item Use IFT to find $\pd{z}{x}$ and $\pd{z}{y}$ for
        $F(x,y,z)=x^{2}+y^{2}+z^{2}-R^{2}=0$.
  \item Verify homogeneity of $f(x,y)=x\mathrm{e}^{y/x}$ and find the
        degree.
  \item Compute the Jacobian and determinant of
        $\mathbf{f}(u,v)=(u^{2}-v^{2},\;2uv)$.
  \item For $U(x,y)=\alpha\ln x+\beta\ln y$, find MRS via IFT.
  \item Verify Euler's theorem for $f(x,y)=\sqrt{x^{2}+y^{2}}$
        (degree $1$).
  \item Use IFT to find $\dd{y}{x}$ for $\sin(xy)+\cos(x+y)=1/2$.
  \item Decide whether $f(x,y)=x^{2}+y^{3}$ is homogeneous, with
        justification.
  \item Compute the Jacobian determinant of the polar map
        $\mathbf{f}(r,\theta)=(r\cos\theta,\;r\sin\theta)$ and connect to
        the area element.
  \item For $Q(K,L)=(K^{\rho}+L^{\rho})^{1/\rho}$ ($\rho\neq 0$), find the
        MRTS.
  \item Verify Euler's theorem for the Leontief-like
        $f(x,y)=\min\{x,y\}$ in the interior of each smooth region.
  \item Use IFT to find $\pd{y}{x_{1}}$ and $\pd{y}{x_{2}}$ at any point
        where $F(x_{1},x_{2},y)=x_{1}^{2}+x_{2}y+y^{3}-7=0$, $F_{y}\neq 0$.
  \item Verify homogeneity and find the degree of
        $f(x,y,z)=\dfrac{x^{2}y+y^{2}z+z^{2}x}{xyz}$.
  \item Compute the Jacobian determinant of
        $\mathbf{f}(x,y)=(x+\sin y,\;y+\sin x)$ and find where it
        vanishes.
  \item For $U(x,y)=(x^{\rho}+y^{\rho})^{1/\rho}$ (CES utility), find the
        MRS.
  \item Verify Euler's theorem (3-variable form) for
        $f(x,y,z)=xy+yz+zx$ (degree $2$).
  \item Use IFT to verify that on the unit circle $x^{2}+y^{2}=1$ the
        slope at $(\sqrt{3}/2,1/2)$ equals $-\sqrt{3}$.
  \item Show that $f(x,y)=(x^{1/2}+y^{1/2})^{2}$ is homogeneous of
        degree $1$, then verify Euler's theorem (LHS and RHS separately).
  \item Compute the Jacobian determinant of
        $\mathbf{f}(x,y)=(x/y,\;y/x)$.
  \item For the translog $\ln Q = \alpha_{0}+\alpha_{K}\ln K+\alpha_{L}\ln L$,
        find the MRTS in $(K,L)$ space.
  \item Verify Euler's theorem for the production function
        $Q(K,L,M)=K^{1/3}L^{1/3}M^{1/3}$ (three-factor CRS Cobb--Douglas).
  \item Compute the Jacobian and Jacobian determinant of
        $\mathbf{f}(x,y,z)=(x+y+z,\;xy+yz+zx,\;xyz)$, and explain the link
        to elementary symmetric polynomials.
  \item Show that $f(x)=x^{5}+2x+1$ is invertible on $\R$, and compute
        $\bigl(f^{-1}\bigr)'(4)$ without inverting $f$.
  \item For $f(x)=x^{3}+3x-2$, compute $\bigl(f^{-1}\bigr)'(2)$.
  \item For $f(x)=\mathrm{e}^{x}+x$, compute $\bigl(f^{-1}\bigr)'(1)$.
  \item For $f(x)=x+\ln x$ on $x>0$, compute $\bigl(f^{-1}\bigr)'(1)$.
  \item Derive $\bigl(f^{-1}\bigr)'(y)=1/f'\bigl(f^{-1}(y)\bigr)$ by
        differentiating the identity $f\bigl(f^{-1}(y)\bigr)=y$.
  \item $f(x)=x^{3}$ is a bijection of $\R$, yet $f'(0)=0$. Which
        conclusion of the Inverse Function Theorem fails at $0$, and which
        survives?
  \item Where does the Inverse Function Theorem apply to $f(x)=x^{2}$, and
        what is the largest interval containing $x=2$ on which $f$ is
        invertible?
  \item Give an interval containing $\pi/3$ on which $\sin$ is invertible,
        and compute $\bigl(\sin^{-1}\bigr)'\bigl(\tfrac{\sqrt3}{2}\bigr)$
        two ways: by the theorem, and from the known derivative of
        $\arcsin$.
  \item For $\mathbf{f}(x,y)=(x+y,\;x-y)$: compute $J$ and $\det J$, invert
        the map explicitly, and check that $J_{\mathbf{f}^{-1}}=J^{-1}$.
  \item For $\mathbf{f}(x,y)=(x^{2}-y^{2},\;2xy)$, compute $\det J$ at
        $(1,2)$ and say whether a local inverse exists there.
  \item Show that the map of the previous problem is \emph{not} globally
        injective, despite $\det J\neq0$ off the origin.
  \item For $\mathbf{f}(x,y)=(\mathrm{e}^{x}\cos y,\;\mathrm{e}^{x}\sin y)$,
        compute $\det J$ and reconcile ``non-zero everywhere'' with the fact
        that $\mathbf{f}$ is not injective.
  \item For the polar map
        $\mathbf{f}(r,\theta)=(r\cos\theta,\;r\sin\theta)$, give $\det J$
        and the set of points where the theorem gives nothing.
  \item For $\mathbf{f}(x,y)=(x+y^{3},\;y+x^{3})$, verify the hypothesis at
        the origin and give $J_{\mathbf{f}^{-1}}$ there.
  \item For $\mathbf{f}(x,y,z)=(x,\;x+y,\;x+y+z)$: compute $\det J$, write
        the inverse map explicitly, and confirm
        $J_{\mathbf{f}^{-1}}=J^{-1}$.
  \item A demand system $\mathbf{q}=\mathbf{f}(\mathbf{p})$ has
        $\det J_{\mathbf{f}}(\mathbf{p}_{0})\neq0$. What does the Inverse
        Function Theorem permit an economist to say, and what does it not?
\end{problems}

\section{Solutions}

\begin{solutions}
  \item $F=x^{3}+y^{3}-6xy$. $F_{x}=3x^{2}-6y$, $F_{y}=3y^{2}-6x$.
        \[
          \dd{y}{x} = -\frac{3x^{2}-6y}{3y^{2}-6x} = \frac{2y-x^{2}}{y^{2}-2x} .
        \]

  \item $f(tx,ty)=(tx)^{4}+5(tx)^{2}(ty)^{2}+(ty)^{4}
        =t^{4}(x^{4}+5x^{2}y^{2}+y^{4})=t^{4}f(x,y)$. \textbf{Degree}
        $k=4$.

  \item
        \[
          J=\begin{pmatrix} 3 & -2\\ 1 & 4\end{pmatrix},\qquad
          \det J=3\cdot 4-(-2)\cdot 1=14 .
        \]
        Globally invertible; areas scale by $14$.

  \item $U_{x}=1$, $U_{y}=2/y$. By IFT $\dd{y}{x}=-U_{x}/U_{y}=-y/2$.
        \textbf{MRS} $=y/2$.

  \item $f_{x}=2x+5y$, $f_{y}=5x-4y$.
        \[
          \mathrm{LHS}=x(2x+5y)+y(5x-4y)=2x^{2}+10xy-4y^{2} .
        \]
        \[
          \mathrm{RHS}=2f=2x^{2}+10xy-4y^{2} .
        \]
        $\mathrm{LHS}=\mathrm{RHS}$. \checkmark

  \item $F=x^{2}y+xy^{3}-10$. $F_{x}=2xy+y^{3}$, $F_{y}=x^{2}+3xy^{2}$.
        \[
          \dd{y}{x} = -\frac{2xy+y^{3}}{x^{2}+3xy^{2}}
            = -\frac{y(2x+y^{2})}{x(x+3y^{2})} .
        \]

  \item
        \[
          f(tx,ty)=\frac{(tx)^{3}+(ty)^{3}}{tx-ty}
            =\frac{t^{3}(x^{3}+y^{3})}{t(x-y)}=t^{2}f(x,y) .
        \]
        \textbf{Degree} $k=2$.

  \item
        \[
          J=\begin{pmatrix} 2x & 2y\\ 2x & -2y\end{pmatrix},\qquad
          \det J=(2x)(-2y)-(2y)(2x)=-8xy .
        \]
        Locally invertible whenever $x\neq 0$ and $y\neq 0$.

  \item $U_{x}=y+2$, $U_{y}=x+1$. $\dd{y}{x}=-(y+2)/(x+1)$.
        \textbf{MRS} $=(y+2)/(x+1)$.

  \item $Q_{K}=K^{-3/4}L^{3/4}$, $Q_{L}=3K^{1/4}L^{-1/4}$.
        \[
          \mathrm{LHS}=K\cdot K^{-3/4}L^{3/4}+L\cdot 3K^{1/4}L^{-1/4}
            =K^{1/4}L^{3/4}+3K^{1/4}L^{3/4}=4K^{1/4}L^{3/4} .
        \]
        \[
          \mathrm{RHS}=1\cdot Q=4K^{1/4}L^{3/4} .
        \]
        $\mathrm{LHS}=\mathrm{RHS}$. \checkmark

  \item $F=\mathrm{e}^{x+y}-x-y-1$. $F_{x}=\mathrm{e}^{x+y}-1=F_{y}$.
        \[
          \dd{y}{x} = -\frac{F_{x}}{F_{y}} = -1 .
        \]
        Valid away from $\mathrm{e}^{x+y}=1$, i.e., $x+y\neq 0$. (The
        implicit relation forces $x+y$ to be locally constant on any
        branch.)

  \item $f(tx,ty)=\sqrt{(tx)(ty)}=t\sqrt{xy}=t f(x,y)$ for $t>0$.
        \textbf{Degree} $k=1$.

  \item
        \[
          J=\begin{pmatrix} \mathrm{e}^{y} & x\mathrm{e}^{y}\\
                             y\mathrm{e}^{x} & \mathrm{e}^{x}\end{pmatrix},
          \qquad \det J=\mathrm{e}^{x+y}(1-xy) .
        \]
        Singular along $xy=1$.

  \item $U_{x}=ax^{a-1}y^{b}$, $U_{y}=bx^{a}y^{b-1}$.
        $U_{x}/U_{y}=(a/b)(y/x)$. By IFT $\dd{y}{x}=-(a/b)(y/x)$.
        \textbf{MRS} $=(a/b)(y/x)$.

  \item $Q_{K}=\tfrac12K^{-1/2}L^{1/2}+1$, $Q_{L}=\tfrac12K^{1/2}L^{-1/2}+1$.
        \[
          \mathrm{LHS}=K\bigl(\tfrac12K^{-1/2}L^{1/2}+1\bigr)
            +L\bigl(\tfrac12K^{1/2}L^{-1/2}+1\bigr)=K^{1/2}L^{1/2}+K+L .
        \]
        \[
          \mathrm{RHS}=1\cdot Q=K^{1/2}L^{1/2}+K+L .
        \]
        $\mathrm{LHS}=\mathrm{RHS}$. \checkmark

  \item $F=\ln(x^{2}+y)-x+y$. $F_{x}=\dfrac{2x}{x^{2}+y}-1$,
        $F_{y}=\dfrac{1}{x^{2}+y}+1$.
        \[
          \dd{y}{x} = -\frac{2x/(x^{2}+y)-1}{1/(x^{2}+y)+1}
            = -\frac{2x-(x^{2}+y)}{1+x^{2}+y}
            = \frac{x^{2}+y-2x}{1+x^{2}+y} .
        \]

  \item
        \[
          f(tx,ty,tz)=\frac{t^{2}(xy+yz+zx)}{t(x+y+z)}=t f(x,y,z) .
        \]
        \textbf{Degree} $k=1$.

  \item
        \[
          J=\begin{pmatrix}
            \cos\theta & -\rho\sin\theta & 0\\
            \sin\theta & \rho\cos\theta & 0\\
            0 & 0 & 1
          \end{pmatrix} .
        \]
        Expanding along the third row:
        $\det J = 1\cdot(\rho\cos^{2}\theta+\rho\sin^{2}\theta)=\rho$.
        Volume element $\mathrm{d}V=\rho\,\mathrm{d}\rho\,\mathrm{d}\theta\,\mathrm{d}z$.

  \item Leontief, non-differentiable on $\{2x=3y\}$. On $\{2x<3y\}$:
        $U=2x$, so $U_{x}=2$, $U_{y}=0$; signed slope $\dd{y}{x}$ is
        undefined (level curve is vertical), MRS is infinite: $x$ is
        the binding factor. On $\{2x>3y\}$: $U=3y$, $U_{x}=0$, $U_{y}=3$;
        signed slope $=0$, MRS $=0$: $y$ is binding. At the kink the
        two-sided derivative does not exist.

  \item $f=(x^{-1}+y^{-1})^{-1}=xy/(x+y)$. $f_{x}=y^{2}/(x+y)^{2}$,
        $f_{y}=x^{2}/(x+y)^{2}$.
        \[
          \mathrm{LHS}=x\cdot\frac{y^{2}}{(x+y)^{2}}
            +y\cdot\frac{x^{2}}{(x+y)^{2}}
            =\frac{xy(y+x)}{(x+y)^{2}}=\frac{xy}{x+y} .
        \]
        \[
          \mathrm{RHS}=1\cdot f=\frac{xy}{x+y} .
        \]
        $\mathrm{LHS}=\mathrm{RHS}$. \checkmark

  \item $F=x\cos y+y\cos x-1$. $F_{x}=\cos y-y\sin x$,
        $F_{y}=-x\sin y+\cos x$.
        \[
          \dd{y}{x} = -\frac{\cos y-y\sin x}{\cos x-x\sin y}
            = \frac{y\sin x-\cos y}{\cos x-x\sin y} .
        \]

  \item
        \[
          f(tx,ty)=(tx)\ln(ty/tx)+(ty)\ln(tx/ty)
            =tx\ln(y/x)+ty\ln(x/y)=t f(x,y)
        \]
        (the $t$ inside the log cancels). \textbf{Degree} $k=1$.

  \item
        \[
          J=\begin{pmatrix} 1 & 1\\ y & x\\ 2x & 2y\end{pmatrix}
            \in\R^{3\times 2} .
        \]
        Not square, so there is no determinant. The top $2\times 2$
        minor is $x-y$, which vanishes precisely on $x=y$; the rank
        drops to $1$ along that diagonal (the other minors also vanish
        there).

  \item Linear: $Q_{K}=5$, $Q_{L}=3$. $\dd{L}{K}=-5/3$. \textbf{MRTS}
        $=5/3$ (perfect substitutes with constant exchange rate).

  \item Let $S=\alpha K^{\rho}+(1-\alpha)L^{\rho}$, so $Q=S^{1/\rho}$.
        \[
          Q_{K}=\frac{1}{\rho}S^{1/\rho-1}\cdot\alpha\rho K^{\rho-1}
            =\alpha S^{1/\rho-1}K^{\rho-1},\qquad
          Q_{L}=(1-\alpha)S^{1/\rho-1}L^{\rho-1} .
        \]
        \[
          \mathrm{LHS}=K\alpha S^{1/\rho-1}K^{\rho-1}
            +L(1-\alpha)S^{1/\rho-1}L^{\rho-1}
            =S^{1/\rho-1}\bigl(\alpha K^{\rho}+(1-\alpha)L^{\rho}\bigr)
            =S^{1/\rho} .
        \]
        \[
          \mathrm{RHS}=1\cdot Q=S^{1/\rho} .
        \]
        $\mathrm{LHS}=\mathrm{RHS}$. \checkmark

  \item Take logs: $y\ln x-x\ln y=0$. Set $F=y\ln x-x\ln y$.
        \[
          F_{x}=\frac{y}{x}-\ln y,\qquad F_{y}=\ln x-\frac{x}{y} .
        \]
        \[
          \dd{y}{x} = -\frac{y/x-\ln y}{\ln x-x/y}
            = \frac{\ln y-y/x}{\ln x-x/y}
            = \frac{y(x\ln y-y)}{x(y\ln x-x)} .
        \]

  \item Test: $f(tx,ty)=(tx)^{2}+(tx)(ty)+(ty)=t^{2}x^{2}+t^{2}xy+ty$.
        This cannot be written as $t^{k}(x^{2}+xy+y)$ for a single $k$:
        the first two terms scale by $t^{2}$, the last by $t$.
        \textbf{Not homogeneous.}

  \item $f_{1}=\mathrm{e}^{x}\cos y$, $f_{2}=\mathrm{e}^{x}\sin y$.
        \[
          J=\begin{pmatrix} \mathrm{e}^{x}\cos y & -\mathrm{e}^{x}\sin y\\
                             \mathrm{e}^{x}\sin y & \mathrm{e}^{x}\cos y
            \end{pmatrix},\qquad
          \det J=\mathrm{e}^{2x}(\cos^{2}y+\sin^{2}y)=\mathrm{e}^{2x} .
        \]
        Always strictly positive: the map is locally invertible
        everywhere. (It is the complex exponential $w=\mathrm{e}^{x+iy}$;
        conformal, orientation-preserving.)

  \item $Q_{K}=\tfrac23K^{-1/3}L^{1/3}$, $Q_{L}=\tfrac13K^{2/3}L^{-2/3}$.
        \[
          \frac{Q_{K}}{Q_{L}}=2\,\frac{L}{K} .
        \]
        $\dd{L}{K}=-2L/K$. \textbf{MRTS} $=2L/K$.

  \item $f_{x}=3x^{2}+3yz$, $f_{y}=3y^{2}+3xz$, $f_{z}=3z^{2}+3xy$.
        \[
          \mathrm{LHS}=x(3x^{2}+3yz)+y(3y^{2}+3xz)+z(3z^{2}+3xy)
            =3(x^{3}+y^{3}+z^{3})+9xyz=3f .
        \]
        \[
          \mathrm{RHS}=3f .
        \]
        $\mathrm{LHS}=\mathrm{RHS}$. \checkmark

  \item $F=\sqrt{x}+\sqrt{y}+\sqrt{xy}-5$. Compute:
        \[
          F_{x}=\frac{1}{2\sqrt{x}}+\frac{y}{2\sqrt{xy}}
            =\frac{1}{2\sqrt{x}}+\frac{1}{2}\sqrt{\frac{y}{x}}
            =\frac{1+\sqrt{y}}{2\sqrt{x}} .
        \]
        By symmetry, $F_{y}=\dfrac{1+\sqrt{x}}{2\sqrt{y}}$.
        \[
          \dd{y}{x} = -\frac{F_{x}}{F_{y}}
            = -\frac{\sqrt{y}(1+\sqrt{y})}{\sqrt{x}(1+\sqrt{x})} .
        \]

  \item
        \[
          f(tx,ty)=\frac{t^{2}x^{2}+t^{2}y^{2}}{(tx)(ty)}
            =\frac{t^{2}(x^{2}+y^{2})}{t^{2}xy}=f(x,y) .
        \]
        \textbf{Degree} $k=0$.

  \item
        \[
          J=\begin{pmatrix}
            \sin\phi\cos\theta & -\rho\sin\phi\sin\theta & \rho\cos\phi\cos\theta\\
            \sin\phi\sin\theta & \rho\sin\phi\cos\theta & \rho\cos\phi\sin\theta\\
            \cos\phi & 0 & -\rho\sin\phi
          \end{pmatrix} .
        \]
        Expand along the third row:
        \[
          \det J = \cos\phi\Bigl[(-\rho\sin\phi\sin\theta)(\rho\cos\phi\sin\theta)
            -(\rho\cos\phi\cos\theta)(\rho\sin\phi\cos\theta)\Bigr]
        \]
        \[
          \qquad{}+(-\rho\sin\phi)\Bigl[(\sin\phi\cos\theta)(\rho\sin\phi\cos\theta)
            -(-\rho\sin\phi\sin\theta)(\sin\phi\sin\theta)\Bigr] .
        \]
        First bracket:
        \[
          -\rho^{2}\sin\phi\cos\phi(\sin^{2}\theta+\cos^{2}\theta)
            =-\rho^{2}\sin\phi\cos\phi .
        \]
        Multiplied by $\cos\phi$ this is
        $-\rho^{2}\sin\phi\cos^{2}\phi$. Second bracket:
        $\rho\sin^{2}\phi(\cos^{2}\theta+\sin^{2}\theta)=\rho\sin^{2}\phi$.
        Multiplied by $-\rho\sin\phi$: $-\rho^{2}\sin^{3}\phi$. Sum:
        $-\rho^{2}\sin\phi(\cos^{2}\phi+\sin^{2}\phi)=-\rho^{2}\sin\phi$.
        Thus $|\det J|=\rho^{2}\sin\phi$, the spherical volume element.

  \item $U_{x}=a$, $U_{y}=b$. $\dd{y}{x}=-a/b$. \textbf{MRS} $=a/b$
        (constant, perfect substitutes).

  \item Simplify: $f=\dfrac{(x+y)(x^{2}-xy+y^{2})}{x+y}=x^{2}-xy+y^{2}$.
        Then $f_{x}=2x-y$, $f_{y}=-x+2y$.
        \[
          \mathrm{LHS}=x(2x-y)+y(-x+2y)=2x^{2}-2xy+2y^{2} .
        \]
        \[
          \mathrm{RHS}=2f=2(x^{2}-xy+y^{2})=2x^{2}-2xy+2y^{2} .
        \]
        $\mathrm{LHS}=\mathrm{RHS}$. \checkmark

  \item $F=x^{2}+y^{2}+z^{2}-R^{2}$. $F_{x}=2x$, $F_{y}=2y$, $F_{z}=2z$.
        Provided $z\neq 0$:
        \[
          \pd{z}{x}=-\frac{F_{x}}{F_{z}}=-\frac{x}{z},\qquad
          \pd{z}{y}=-\frac{F_{y}}{F_{z}}=-\frac{y}{z} .
        \]
        (Geometrically, the outward normal on the sphere has direction
        $(x,y,z)/R$.)

  \item $f(tx,ty)=(tx)\mathrm{e}^{(ty)/(tx)}=tx\,\mathrm{e}^{y/x}
        =t\,f(x,y)$. \textbf{Degree} $k=1$.

  \item
        \[
          J=\begin{pmatrix} 2u & -2v\\ 2v & 2u\end{pmatrix},\qquad
          \det J=4u^{2}+4v^{2}=4(u^{2}+v^{2}) .
        \]
        Vanishes only at the origin. (This is the complex squaring
        $w=(u+iv)^{2}$.)

  \item $U_{x}=\alpha/x$, $U_{y}=\beta/y$. $U_{x}/U_{y}=(\alpha/\beta)(y/x)$.
        $\dd{y}{x}=-(\alpha/\beta)(y/x)$. \textbf{MRS}
        $=(\alpha/\beta)(y/x)$.

  \item $f=(x^{2}+y^{2})^{1/2}$. $f_{x}=x/\sqrt{x^{2}+y^{2}}$,
        $f_{y}=y/\sqrt{x^{2}+y^{2}}$.
        \[
          \mathrm{LHS}=\frac{x^{2}}{\sqrt{x^{2}+y^{2}}}
            +\frac{y^{2}}{\sqrt{x^{2}+y^{2}}}
            =\frac{x^{2}+y^{2}}{\sqrt{x^{2}+y^{2}}}=\sqrt{x^{2}+y^{2}} .
        \]
        \[
          \mathrm{RHS}=1\cdot f=\sqrt{x^{2}+y^{2}} .
        \]
        $\mathrm{LHS}=\mathrm{RHS}$. \checkmark

  \item $F=\sin(xy)+\cos(x+y)-1/2$. $F_{x}=y\cos(xy)-\sin(x+y)$,
        $F_{y}=x\cos(xy)-\sin(x+y)$.
        \[
          \dd{y}{x} = -\frac{y\cos(xy)-\sin(x+y)}{x\cos(xy)-\sin(x+y)}
            = \frac{\sin(x+y)-y\cos(xy)}{x\cos(xy)-\sin(x+y)} .
        \]

  \item $f(tx,ty)=t^{2}x^{2}+t^{3}y^{3}$. The two terms scale with
        different powers of $t$, so they cannot be factored as
        $t^{k}(x^{2}+y^{3})$. \textbf{Not homogeneous.}

  \item
        \[
          J=\begin{pmatrix} \cos\theta & -r\sin\theta\\
                             \sin\theta & r\cos\theta\end{pmatrix},\qquad
          \det J=r\cos^{2}\theta+r\sin^{2}\theta=r .
        \]
        Hence $|\det J|=r$, and the area element is
        $\mathrm{d}A=r\,\mathrm{d}r\,\mathrm{d}\theta$, the standard
        polar substitution.

  \item Let $S=K^{\rho}+L^{\rho}$. Then $Q=S^{1/\rho}$,
        $Q_{K}=S^{1/\rho-1}K^{\rho-1}$, $Q_{L}=S^{1/\rho-1}L^{\rho-1}$.
        \[
          \frac{Q_{K}}{Q_{L}}=\Bigl(\frac{K}{L}\Bigr)^{\rho-1}
            =\Bigl(\frac{L}{K}\Bigr)^{1-\rho} .
        \]
        $\dd{L}{K}=-(L/K)^{1-\rho}$. \textbf{MRTS} $=(L/K)^{1-\rho}$.

  \item On the smooth piece $\{x<y\}$: $f=x$, $f_{x}=1$, $f_{y}=0$.
        \[
          \mathrm{LHS}=x\cdot 1+y\cdot 0=x=f,\qquad \mathrm{RHS}=1\cdot f=x .
        \]
        $\mathrm{LHS}=\mathrm{RHS}$. \checkmark On $\{x>y\}$ by symmetry
        $\mathrm{LHS}=y=f=\mathrm{RHS}$. \checkmark Euler's identity
        holds wherever the function is differentiable. (At the ridge
        $x=y$ the function is not differentiable.)

  \item By IFT, $\pd{y}{x_{i}}=-F_{x_{i}}/F_{y}$ with
        $F_{y}=x_{2}+3y^{2}$:
        \[
          \pd{y}{x_{1}}=-\frac{2x_{1}}{x_{2}+3y^{2}},\qquad
          \pd{y}{x_{2}}=-\frac{y}{x_{2}+3y^{2}} .
        \]

  \item Numerator has degree $3$, denominator degree $3$. Concretely
        \[
          f(tx,ty,tz)=\frac{t^{3}(x^{2}y+y^{2}z+z^{2}x)}{t^{3}xyz}=f(x,y,z) .
        \]
        \textbf{Degree} $k=0$.

  \item $J=\begin{pmatrix} 1 & \cos y\\ \cos x & 1\end{pmatrix}$,
        $\det J=1-\cos x\cos y$. Vanishes when $\cos x\cos y=1$, i.e.,
        $\cos x=\cos y=1$ or $\cos x=\cos y=-1$, equivalently
        $(x,y)\in 2\pi\mathbb{Z}\times 2\pi\mathbb{Z}$ or
        $(2\mathbb{Z}+1)\pi\times(2\mathbb{Z}+1)\pi$.

  \item Let $S=x^{\rho}+y^{\rho}$, $U=S^{1/\rho}$.
        $U_{x}=S^{1/\rho-1}x^{\rho-1}$, $U_{y}=S^{1/\rho-1}y^{\rho-1}$.
        \[
          \frac{U_{x}}{U_{y}}=\Bigl(\frac{x}{y}\Bigr)^{\rho-1}
            =\Bigl(\frac{y}{x}\Bigr)^{1-\rho} .
        \]
        $\dd{y}{x}=-(y/x)^{1-\rho}$. \textbf{MRS} $=(y/x)^{1-\rho}$.

  \item $f_{x}=y+z$, $f_{y}=x+z$, $f_{z}=x+y$.
        \[
          \mathrm{LHS}=x(y+z)+y(x+z)+z(x+y)=2(xy+yz+zx)=2f .
        \]
        \[
          \mathrm{RHS}=2f .
        \]
        $\mathrm{LHS}=\mathrm{RHS}$. \checkmark

  \item $F=x^{2}+y^{2}-1$. $F_{x}=2x$, $F_{y}=2y$. $\dd{y}{x}=-x/y$.
        At $(\sqrt{3}/2,1/2)$: $\dd{y}{x}=-(\sqrt{3}/2)/(1/2)=-\sqrt{3}$.
        \checkmark

  \item $f(tx,ty)=\bigl((tx)^{1/2}+(ty)^{1/2}\bigr)^{2}
        =(t^{1/2})^{2}(\sqrt{x}+\sqrt{y})^{2}=t f(x,y)$. \textbf{Degree}
        $k=1$. For Euler's theorem, write $A=\sqrt{x}+\sqrt{y}$,
        $f=A^{2}$. Then $f_{x}=2A\cdot\dfrac{1}{2\sqrt{x}}=A/\sqrt{x}$,
        $f_{y}=A/\sqrt{y}$.
        \[
          \mathrm{LHS}=x\cdot\frac{A}{\sqrt{x}}+y\cdot\frac{A}{\sqrt{y}}
            =A(\sqrt{x}+\sqrt{y})=A^{2}=f .
        \]
        \[
          \mathrm{RHS}=1\cdot f=A^{2} .
        \]
        $\mathrm{LHS}=\mathrm{RHS}$. \checkmark

  \item
        \[
          J=\begin{pmatrix} 1/y & -x/y^{2}\\ -y/x^{2} & 1/x\end{pmatrix},
          \qquad \det J=\frac{1}{xy}-\frac{xy}{x^{2}y^{2}}
            =\frac{1}{xy}-\frac{1}{xy}=0 .
        \]
        The Jacobian is identically singular: because
        $(x/y)\cdot(y/x)=1$, the image lies on the curve $f_{1}f_{2}=1$
        (one-dimensional), so the map is not locally invertible
        anywhere.

  \item From $\ln Q=\alpha_{0}+\alpha_{K}\ln K+\alpha_{L}\ln L$, we get
        $Q=AK^{\alpha_{K}}L^{\alpha_{L}}$ with $A=\mathrm{e}^{\alpha_{0}}$
        (Cobb--Douglas). Then $Q_{K}=A\alpha_{K}K^{\alpha_{K}-1}L^{\alpha_{L}}$,
        $Q_{L}=A\alpha_{L}K^{\alpha_{K}}L^{\alpha_{L}-1}$.
        \[
          \frac{Q_{K}}{Q_{L}}=\frac{\alpha_{K}}{\alpha_{L}}\cdot\frac{L}{K} .
        \]
        $\dd{L}{K}=-(\alpha_{K}/\alpha_{L})(L/K)$. \textbf{MRTS}
        $=(\alpha_{K}/\alpha_{L})(L/K)$.

  \item $Q=K^{1/3}L^{1/3}M^{1/3}$ is homogeneous of degree
        $1/3+1/3+1/3=1$. $Q_{K}=\tfrac13K^{-2/3}L^{1/3}M^{1/3}=Q/(3K)$,
        and similarly $Q_{L}=Q/(3L)$, $Q_{M}=Q/(3M)$.
        \[
          \mathrm{LHS}=K\cdot\frac{Q}{3K}+L\cdot\frac{Q}{3L}
            +M\cdot\frac{Q}{3M}=\frac{Q}{3}+\frac{Q}{3}+\frac{Q}{3}=Q .
        \]
        \[
          \mathrm{RHS}=1\cdot Q=Q .
        \]
        $\mathrm{LHS}=\mathrm{RHS}$. \checkmark

  \item Write $f_{1}=x+y+z$, $f_{2}=xy+yz+zx$, $f_{3}=xyz$ (the
        elementary symmetric polynomials $e_{1},e_{2},e_{3}$). Then
        \[
          J=\begin{pmatrix}
            1 & 1 & 1\\
            y+z & x+z & x+y\\
            yz & xz & xy
          \end{pmatrix} .
        \]
        For the determinant, subtract column $1$ from columns $2$ and
        $3$ to get
        \[
          \det J = \det\begin{pmatrix}
            1 & 0 & 0\\
            y+z & x-y & x-z\\
            yz & z(x-y) & y(x-z)
          \end{pmatrix}
          = 1\cdot\det\begin{pmatrix}
            x-y & x-z\\ z(x-y) & y(x-z)
          \end{pmatrix}
        \]
        \[
          = (x-y)(x-z)\det\begin{pmatrix} 1 & 1\\ z & y\end{pmatrix}
          = (x-y)(x-z)(y-z) .
        \]
        Thus $\det J=-(x-y)(y-z)(z-x)$ (the Vandermonde discriminant, up
        to sign), which vanishes exactly when two of $x,y,z$ coincide:
        the symmetric-polynomial map fails to be locally invertible
        at points with repeated coordinates, consistent with the fact
        that an unordered triple $\{x,y,z\}$ has a unique sorted
        representative.
  \item $f'(x)=5x^{4}+2\ge2>0$ for all $x$, so $f$ is strictly increasing,
        hence injective, and the theorem applies everywhere. Solve
        $x^{5}+2x+1=4$: $x=1$ works ($1+2+1=4$), and monotonicity makes it
        the only root. So
        $\bigl(f^{-1}\bigr)'(4)=1/f'(1)=1/7$.
  \item $f(1)=1+3-2=2$, and $f'(x)=3x^{2}+3$, so $f'(1)=6$ and
        $\bigl(f^{-1}\bigr)'(2)=\tfrac16$. ($f'>0$ everywhere, so the
        inverse is global.)
  \item $f(0)=1+0=1$, and $f'(x)=\mathrm{e}^{x}+1$, so $f'(0)=2$ and
        $\bigl(f^{-1}\bigr)'(1)=\tfrac12$.
  \item $f(1)=1+\ln1=1$, and $f'(x)=1+1/x$, so $f'(1)=2$ and
        $\bigl(f^{-1}\bigr)'(1)=\tfrac12$.
  \item Differentiate both sides with respect to $y$. The left side is a
        composite, so the chain rule gives
        $f'\bigl(f^{-1}(y)\bigr)\cdot\bigl(f^{-1}\bigr)'(y)$; the right side
        gives $1$. Hence
        $\bigl(f^{-1}\bigr)'(y)=1/f'\bigl(f^{-1}(y)\bigr)$, the division
        being legitimate exactly when $f'\neq0$ there.
  \item Invertibility survives: $x^{3}$ is strictly increasing and onto, so
        $f^{-1}(y)=y^{1/3}$ exists globally. What fails is
        \emph{differentiability of the inverse at $0$}:
        $\bigl(y^{1/3}\bigr)'=\tfrac13y^{-2/3}\to\infty$ as $y\to0$. The
        hypothesis $f'\neq0$ is sufficient for a differentiable local
        inverse, not necessary for an inverse.
  \item $f'(x)=2x\neq0$ for $x\neq0$, so the theorem applies at every point
        except the origin. Around $x=2$ the largest interval of
        injectivity is $(0,\infty)$: on it $f$ is strictly increasing, while
        any interval containing $0$ in its interior contains a pair
        $\pm a$ with equal images.
  \item On $(-\pi/2,\pi/2)$ the sine is strictly increasing, and
        $\pi/3$ lies inside. By the theorem, with
        $\sin(\pi/3)=\sqrt3/2$ and $\cos(\pi/3)=\tfrac12$,
        \[
          \bigl(\sin^{-1}\bigr)'\Bigl(\tfrac{\sqrt3}{2}\Bigr)
          =\frac{1}{\cos(\pi/3)}=2 .
        \]
        Directly: $\arcsin'(y)=1/\sqrt{1-y^{2}}
        =1/\sqrt{1-\tfrac34}=1/\tfrac12=2$. \checkmark
  \item $J=\begin{pmatrix}1&1\\1&-1\end{pmatrix}$, $\det J=-2\neq0$
        everywhere. Solving $u=x+y$, $v=x-y$ gives
        $x=\tfrac{u+v}{2}$, $y=\tfrac{u-v}{2}$, so
        $J_{\mathbf{f}^{-1}}
        =\begin{pmatrix}1/2&1/2\\1/2&-1/2\end{pmatrix}$, which is indeed
        $J^{-1}=-\tfrac12\begin{pmatrix}-1&-1\\-1&1\end{pmatrix}$.
        \checkmark\ The map is linear with $\det\neq0$, so here local and
        global invertibility coincide.
  \item $J=\begin{pmatrix}2x&-2y\\2y&2x\end{pmatrix}$ and
        $\det J=4(x^{2}+y^{2})$, so at $(1,2)$ it is $4\cdot5=20\neq0$: a
        differentiable local inverse exists near
        $\mathbf{f}(1,2)=(-3,4)$.
  \item $\mathbf{f}(-x,-y)=\bigl((-x)^{2}-(-y)^{2},\,2(-x)(-y)\bigr)
        =(x^{2}-y^{2},\,2xy)=\mathbf{f}(x,y)$, so every point and its
        negative share an image --- for instance
        $\mathbf{f}(1,1)=\mathbf{f}(-1,-1)=(0,2)$. The map is two-to-one,
        and no global inverse exists. Local invertibility at every point
        does not assemble into global invertibility.
  \item $J=\mathrm{e}^{x}\begin{pmatrix}\cos y&-\sin y\\
        \sin y&\cos y\end{pmatrix}$, so
        $\det J=\mathrm{e}^{2x}(\cos^{2}y+\sin^{2}y)=\mathrm{e}^{2x}>0$
        everywhere. Yet $\mathbf{f}(x,y+2\pi)=\mathbf{f}(x,y)$. There is no
        contradiction: the theorem promises an inverse on \emph{some
        neighbourhood} of each point, and those neighbourhoods (width less
        than $2\pi$ in $y$) never join up into one.
  \item $J=\begin{pmatrix}\cos\theta&-r\sin\theta\\
        \sin\theta&r\cos\theta\end{pmatrix}$, so $\det J=r$. The theorem
        says nothing on the line $r=0$ --- the origin, where every $\theta$
        gives the same point and the map is infinitely many-to-one.
  \item $J=\begin{pmatrix}1&3y^{2}\\3x^{2}&1\end{pmatrix}$, which at
        $(0,0)$ is the identity, with $\det=1\neq0$. Hence a local inverse
        exists near $\mathbf{f}(0,0)=(0,0)$ and
        $J_{\mathbf{f}^{-1}}(0,0)=I^{-1}=I$.
  \item $J=\begin{pmatrix}1&0&0\\1&1&0\\1&1&1\end{pmatrix}$, lower
        triangular, so $\det J=1\neq0$. From $u=x$, $v=x+y$, $w=x+y+z$:
        $x=u$, $y=v-u$, $z=w-v$, giving
        $J_{\mathbf{f}^{-1}}=\begin{pmatrix}1&0&0\\-1&1&0\\
        0&-1&1\end{pmatrix}$. Multiplying the two matrices returns $I$.
        \checkmark
  \item It permits: near $\mathbf{p}_{0}$ the system can in principle be
        solved for prices in terms of quantities, the inverse demand
        $\mathbf{p}(\mathbf{q})$ is differentiable, and its derivative
        matrix is $\bigl[J_{\mathbf{f}}\bigr]^{-1}$ --- so the price
        responses can be computed from the quantity responses without ever
        solving the system. It does not permit: any claim that the same
        quantities cannot arise from a different price vector elsewhere.
        Global identification is a separate question, and $\det J\neq0$ is
        no evidence for it.
\end{solutions}

\section{Additional Problems}

\begin{problems}
  \item Use IFT to find $\dd{y}{x}$ for $x^{2}+2xy+y^{2}=9$.
  \item Verify homogeneity and find the degree of $f(x,y)=x^{3}y+xy^{3}$.
  \item Compute the Jacobian and its determinant for
        $\mathbf{f}(x,y)=(2x+y,\;x-y)$.
  \item For $U(x,y)=x^{2}+y$, find the signed slope and unsigned MRS
        using IFT.
  \item Verify Euler's theorem for $f(x,y)=x^{3}+2x^{2}y-y^{3}$
        (degree $3$).
  \item Use IFT to find $\dd{y}{x}$ for $x^{3}+y^{3}+xy=5$.
  \item Verify homogeneity and find the degree of
        $f(x,y)=\dfrac{x^{4}+y^{4}}{x^{2}+y^{2}}$.
  \item Compute the Jacobian and its determinant for
        $\mathbf{f}(x,y)=(x^{2}y,\;xy^{2})$.
  \item For $U(x,y)=x^{1/3}y^{2/3}$, find the signed slope and MRS via
        IFT.
  \item Verify Euler's theorem for $Q(K,L)=K^{1/2}L^{1/2}$ (degree $1$).
  \item Use IFT to find $\dd{y}{x}$ for $\ln(xy)+x+y=4$.
  \item Decide whether $f(x,y)=x+y^{2}$ is homogeneous, with
        justification.
  \item Compute the Jacobian determinant of
        $\mathbf{f}(x,y,z)=(x+y,\;y+z,\;x+z)$.
  \item For $Q(K,L)=aK+bL$ ($a,b>0$), find the MRTS.
  \item Verify Euler's theorem for $f(x,y,z)=xy+yz+xz$ (degree $2$,
        three-variable).
  \item For $f(x)=x^{7}+x^{3}+x$, compute $\bigl(f^{-1}\bigr)'(3)$.
  \item For $f(x)=2x+\sin x$, compute $\bigl(f^{-1}\bigr)'(0)$.
  \item For $f(x)=\ln x+x^{2}$ on $x>0$, compute $\bigl(f^{-1}\bigr)'(1)$.
  \item For $f(x)=\tan x$ on $(-\pi/2,\pi/2)$, compute
        $\bigl(f^{-1}\bigr)'(1)$ and check it against $\arctan'$.
  \item For $\mathbf{f}(x,y)=(2x+3y,\;x+2y)$, give $\det J$ and
        $J_{\mathbf{f}^{-1}}$, and say why the inverse is global here.
  \item For $\mathbf{f}(x,y)=(x^{3},\;y^{3})$: where does $\det J$ vanish,
        and is the map invertible there anyway?
  \item For $\mathbf{f}(x,y)=(x+\mathrm{e}^{y},\;y+\mathrm{e}^{x})$,
        evaluate $\det J$ at the origin and say what the theorem concludes.
  \item Let $\mathbf{f}(\mathbf{x})=A\mathbf{x}$ with $A$ an $n\times n$
        matrix. Show that $\det J$ is constant, and that local and global
        invertibility coincide.
  \item For $\mathbf{f}(x,y)=(x^{2}+y^{2},\;x^{2}-y^{2})$ on $x,y>0$,
        compute $\det J$ and decide local invertibility.
  \item State in one sentence each what $f'(x_{0})\neq0$ and
        $\det J(\mathbf{x}_{0})\neq0$ do and do not guarantee.
\end{problems}

\section{Additional Solutions}

\begin{solutions}
  \item $F=x^{2}+2xy+y^{2}-9=(x+y)^{2}-9$. $F_{x}=2(x+y)$,
        $F_{y}=2(x+y)$.
        \[
          \dd{y}{x}=-\frac{F_{x}}{F_{y}}=-1 .
        \]
        (The curve is two lines $x+y=\pm 3$, each with slope $-1$.)

  \item $f(tx,ty)=(tx)^{3}(ty)+(tx)(ty)^{3}=t^{4}(x^{3}y+xy^{3})
        =t^{4}f(x,y)$. \textbf{Degree $k=4$.}

  \item
        \[
          J=\begin{pmatrix} 2 & 1\\ 1 & -1\end{pmatrix},\qquad
          \det J=-2-1=-3 .
        \]
        Globally invertible; areas scale by $|{-3}|=3$ (with orientation
        reversal).

  \item $U_{x}=2x$, $U_{y}=1$. By IFT $\dd{y}{x}=-2x$. \textbf{MRS}
        $=2x$ (unsigned).

  \item $f_{x}=3x^{2}+4xy$, $f_{y}=2x^{2}-3y^{2}$.
        \[
          \mathrm{LHS}=x(3x^{2}+4xy)+y(2x^{2}-3y^{2})
          =3x^{3}+4x^{2}y+2x^{2}y-3y^{3}=3x^{3}+6x^{2}y-3y^{3} .
        \]
        \[
          \mathrm{RHS}=3f=3(x^{3}+2x^{2}y-y^{3})
          =3x^{3}+6x^{2}y-3y^{3} .
        \]
        $\mathrm{LHS}=\mathrm{RHS}$. \checkmark

  \item $F=x^{3}+y^{3}+xy-5$. $F_{x}=3x^{2}+y$, $F_{y}=3y^{2}+x$.
        \[
          \dd{y}{x}=-\frac{3x^{2}+y}{3y^{2}+x} .
        \]

  \item $f(tx,ty)=\dfrac{t^{4}(x^{4}+y^{4})}{t^{2}(x^{2}+y^{2})}
        =t^{2}f(x,y)$. \textbf{Degree $k=2$.}

  \item
        \[
          J=\begin{pmatrix} 2xy & x^{2}\\ y^{2} & 2xy\end{pmatrix},\qquad
          \det J=4x^{2}y^{2}-x^{2}y^{2}=3x^{2}y^{2} .
        \]
        Singular when $x=0$ or $y=0$.

  \item $U_{x}=\tfrac13 x^{-2/3}y^{2/3}$,
        $U_{y}=\tfrac23 x^{1/3}y^{-1/3}$.
        $U_{x}/U_{y}=\tfrac12\cdot y/x$. By IFT $\dd{y}{x}=-y/(2x)$.
        \textbf{MRS} $=y/(2x)$.

  \item $Q_{K}=\tfrac12 K^{-1/2}L^{1/2}$,
        $Q_{L}=\tfrac12 K^{1/2}L^{-1/2}$.
        \[
          \mathrm{LHS}=K\cdot\tfrac12 K^{-1/2}L^{1/2}
          +L\cdot\tfrac12 K^{1/2}L^{-1/2}
          =\tfrac12 K^{1/2}L^{1/2}+\tfrac12 K^{1/2}L^{1/2}
          =K^{1/2}L^{1/2} .
        \]
        \[
          \mathrm{RHS}=1\cdot Q=K^{1/2}L^{1/2} .
        \]
        $\mathrm{LHS}=\mathrm{RHS}$. \checkmark

  \item $F=\ln(xy)+x+y-4=\ln x+\ln y+x+y-4$. $F_{x}=1/x+1$,
        $F_{y}=1/y+1$.
        \[
          \dd{y}{x}=-\frac{1/x+1}{1/y+1}=-\frac{y(1+x)}{x(1+y)} .
        \]

  \item $f(tx,ty)=tx+t^{2}y^{2}$. The two terms scale with different
        powers of $t$, so they cannot share a common $t^{k}$ factor.
        \textbf{Not homogeneous.}

  \item
        \[
          J=\begin{pmatrix}
            1 & 1 & 0\\
            0 & 1 & 1\\
            1 & 0 & 1
          \end{pmatrix} .
        \]
        Expanding along the first row:
        $\det J=1\cdot(1\cdot 1-1\cdot 0)-1\cdot(0\cdot 1-1\cdot 1)+0
        =1-(-1)+0=2$.

  \item Linear: $Q_{K}=a$, $Q_{L}=b$. $\dd{L}{K}=-a/b$. \textbf{MRTS}
        $=a/b$ (constant, perfect substitutes).

  \item $f_{x}=y+z$, $f_{y}=x+z$, $f_{z}=x+y$.
        \[
          \mathrm{LHS}=x(y+z)+y(x+z)+z(x+y)=2(xy+yz+xz)=2f .
        \]
        \[
          \mathrm{RHS}=2f .
        \]
        $\mathrm{LHS}=\mathrm{RHS}$. \checkmark
  \item $f(1)=1+1+1=3$ and $f'(x)=7x^{6}+3x^{2}+1$, so $f'(1)=11$ and
        $\bigl(f^{-1}\bigr)'(3)=\tfrac{1}{11}$. ($f'>0$ everywhere, so the
        pre-image is unique.)
  \item $f(0)=0$ and $f'(x)=2+\cos x\ge1>0$, so $f'(0)=3$ and
        $\bigl(f^{-1}\bigr)'(0)=\tfrac13$.
  \item $f(1)=\ln1+1=1$ and $f'(x)=1/x+2x$, so $f'(1)=3$ and
        $\bigl(f^{-1}\bigr)'(1)=\tfrac13$.
  \item $\tan(\pi/4)=1$ and $f'(x)=1/\cos^{2}x$, so
        $f'(\pi/4)=1/(\sqrt2/2)^{2}=2$ and
        $\bigl(f^{-1}\bigr)'(1)=\tfrac12$. Directly:
        $\arctan'(1)=1/(1+1^{2})=\tfrac12$. \checkmark
  \item $J=\begin{pmatrix}2&3\\1&2\end{pmatrix}$ is constant with
        $\det J=4-3=1\neq0$, and
        $J_{\mathbf{f}^{-1}}=J^{-1}
        =\begin{pmatrix}2&-3\\-1&2\end{pmatrix}$. The map is linear, so the
        local inverse supplied by the theorem is the global one:
        $\mathbf{x}=A^{-1}\mathbf{y}$ is defined on all of $\R^{2}$.
  \item $J=\begin{pmatrix}3x^{2}&0\\0&3y^{2}\end{pmatrix}$ and
        $\det J=9x^{2}y^{2}$, which vanishes on both axes. The map is
        nevertheless a bijection of $\R^{2}$, with inverse
        $(u,v)\mapsto(u^{1/3},v^{1/3})$ --- but that inverse is not
        differentiable where $\det J=0$. Same lesson as $x^{3}$ in one
        variable.
  \item $J=\begin{pmatrix}1&\mathrm{e}^{y}\\ \mathrm{e}^{x}&1\end{pmatrix}$,
        so at $(0,0)$ it is
        $\begin{pmatrix}1&1\\1&1\end{pmatrix}$ with $\det J=0$. The theorem
        concludes \textbf{nothing}: its hypothesis fails, and that is not
        the same as showing no local inverse exists. A separate argument
        would be needed either way.
  \item Each component of $A\mathbf{x}$ is linear, so
        $\partial(A\mathbf{x})_{i}/\partial x_{j}=a_{ij}$ and $J=A$ at every
        point --- constant, with $\det J=\det A$. If $\det A\neq0$ then
        $\mathbf{x}=A^{-1}\mathbf{y}$ inverts the map on all of $\R^{n}$; if
        $\det A=0$ then $A$ has a non-trivial null space and the map is not
        injective anywhere. So for linear maps the local condition is
        already global.
  \item $J=\begin{pmatrix}2x&2y\\2x&-2y\end{pmatrix}$ and
        $\det J=-4xy-4xy=-8xy$, which is non-zero throughout $x,y>0$. So a
        differentiable local inverse exists at every point of the open
        quadrant.
  \item $f'(x_{0})\neq0$ guarantees an interval around $x_{0}$ on which $f$
        is invertible with a differentiable inverse of derivative
        $1/f'(x_{0})$; it does not guarantee that $f$ is injective anywhere
        beyond that interval. $\det J(\mathbf{x}_{0})\neq0$ guarantees the
        same locally in $n$ dimensions, with the inverse's Jacobian equal to
        $\bigl[J(\mathbf{x}_{0})\bigr]^{-1}$; it likewise says nothing about
        global injectivity.
\end{solutions}


\end{document}
"
\ifdefined\KindleEdition
  \documentclass[10pt,oneside]{book}
\else
  \documentclass[11pt,a4paper,twoside,openright]{book}
\fi
% Shared preamble for the Derivatives training sheet.
\usepackage[T1]{fontenc}
\usepackage[utf8]{inputenc}
\usepackage{lmodern}
\ifdefined\KindleEdition
  % 6-inch e-ink page (Kindle Paperwhite): the page *is* the screen, so the
  % margins are only what the eye needs, not what a printer needs.
  \usepackage[paperwidth=90mm,paperheight=122mm,
              top=4mm,bottom=6mm,left=4.5mm,right=4.5mm]{geometry}
\else
  \usepackage[a4paper,margin=2.4cm]{geometry}
\fi
\usepackage{amsmath,amssymb,amsthm}
\usepackage{microtype}
\usepackage{enumitem}
\usepackage{xcolor}
\usepackage[most]{tcolorbox}
\usepackage{fancyhdr}
\usepackage[hidelinks]{hyperref}

\definecolor{accent}{HTML}{1F4E79}
\definecolor{soft}{HTML}{F4F7FA}

% --- function graphs ------------------------------------------------------
\usepackage{pgfplots}
\pgfplotsset{compat=1.18}
\usetikzlibrary{arrows.meta}

% One consistent look for every function graph in these books.
\pgfplotsset{
  fnaxis/.style={
    width=5.9cm, height=4.6cm,
    axis lines=middle,
    axis line style={-{Stealth[length=4pt]}, gray!65!black, line width=0.3pt},
    tick style={gray!65!black},
    tick label style={font=\tiny, gray!50!black},
    label style={font=\scriptsize},
    xlabel={$x$}, ylabel={$y$},
    every axis x label/.style={at={(ticklabel* cs:1.0)}, anchor=west},
    every axis y label/.style={at={(ticklabel* cs:1.0)}, anchor=south east,
      xshift=-1pt},
    every axis plot/.append style={accent, very thick, smooth, samples=201},
    clip=true, enlargelimits=false,
    title style={font=\scriptsize, yshift=4pt},
  },
  fnaxis/.append style={/pgfplots/legend style={font=\tiny, draw=none,
    fill=none, at={(0.98,0.03)}, anchor=south east}},
}

% A row of small graphs, evenly spaced with \hfill between them, plus an
% optional caption line: \begin{gallery} <pic>\hfill<pic> \gcap{...} \end{gallery}
\newenvironment{gallery}
  {\par\medskip\noindent\ignorespaces}
  {\par\medskip\noindent}
\newcommand{\gcap}[1]{\par\smallskip{\centering\footnotesize #1\par}}

% --- plane regions, implicit curves, surfaces -----------------------------
\usepgfplotslibrary{fillbetween}
\definecolor{regionfill}{HTML}{9DB4D8}
\pgfplotsset{colormap={house}{color=(soft) color=(accent!45) color=(accent)}}

\pgfplotsset{
  % Squared, gridded axis for plane regions and implicitly given curves.
  geoaxis/.style={
    fnaxis,
    width=5.9cm, height=5.9cm,
    grid=both,
    grid style={gray!22, line width=0.2pt},
    minor tick num=1,
  },
  % Level-curve maps: many thin curves, no single curve emphasised.
  contouraxis/.style={
    geoaxis,
    every axis plot/.append style={accent!75, line width=0.5pt, smooth,
      samples=201},
  },
  % Surfaces: \addplot3[surf, domain=..., y domain=...] {expr};
  surfaxis/.style={
    width=6.0cm, height=5.0cm,
    view={130}{22},
    colormap name=house,
    ticks=none,
    xlabel={$x$}, ylabel={$y$}, zlabel={$f$},
    label style={font=\tiny, gray!50!black},
    title style={font=\scriptsize},
    every axis plot/.append style={samples=25, shader=faceted interp,
      draw=accent!35, line width=0.1pt},
  },
}

% Shorthand for a shaded region between two named paths.
% \region{<path A>}{<path B>}{<soft clip options>}
\newcommand{\region}[3]{%
  \addplot[regionfill, opacity=0.55, forget plot, sharp plot]
    fill between[of=#1 and #2, #3];}

% --- rule boxes -----------------------------------------------------------
\newtcolorbox{rulebox}[1]{%
  breakable, enhanced, colback=soft, colframe=accent,
  boxrule=0.4pt, left=7pt, right=7pt, top=6pt, bottom=6pt,
  fonttitle=\bfseries, coltitle=white, colbacktitle=accent,
  title={#1}}

% --- worked examples ------------------------------------------------------
% One fully narrated solution per chapter, carrying the chapter's central
% technique and the mistakes that technique invites:
%   \begin{worked}{Task} ... \wstep{Step 1.} ... \wtrap{...} ... \end{worked}
\newtcolorbox{worked}[1]{%
  breakable, enhanced, colback=white, colframe=accent!65,
  boxrule=0.4pt, left=7pt, right=7pt, top=6pt, bottom=6pt,
  fonttitle=\bfseries, coltitle=accent, colbacktitle=soft,
  title={#1}}

% Run-in step heading inside a worked example.
\newcommand{\wstep}[1]{\par\smallskip\noindent\textbf{#1}\enspace\ignorespaces}

% A pitfall, set off from the running solution so it can be found again.
\newcommand{\wtrap}[1]{%
  \par\smallskip
  {\leftskip=1.3em\noindent
   \llap{\textcolor{accent}{\rule[0.22ex]{0.9em}{0.6pt}}\hspace{0.4em}}%
   \textbf{Trap.}\enspace\itshape #1\par}%
  \smallskip}

% --- problem / solution lists --------------------------------------------
\newlist{problems}{enumerate}{1}
\setlist[problems]{label=\textbf{\arabic*.}, leftmargin=2.1em,
  itemsep=0.55ex, parsep=0pt, topsep=0.8ex}
\newlist{solutions}{enumerate}{1}
\setlist[solutions]{label=\textbf{\arabic*.}, leftmargin=2.1em,
  itemsep=1.0ex, parsep=0pt, topsep=0.8ex}

% --- running heads --------------------------------------------------------
\pagestyle{fancy}
\fancyhf{}
\fancyhead[LE,RO]{\thepage}
\fancyhead[RE]{\small\itshape\nouppercase{\leftmark}}
\fancyhead[LO]{\small\itshape\nouppercase{\rightmark}}
\setlength{\headheight}{14pt}
\renewcommand{\headrulewidth}{0.3pt}
\renewcommand{\chaptermark}[1]{\markboth{\chaptername\ \thechapter.\ #1}{}}
\renewcommand{\sectionmark}[1]{\markright{\thesection\ #1}}

% --- shorthands used throughout -------------------------------------------
\newcommand{\dd}[2]{\frac{\mathrm{d}#1}{\mathrm{d}#2}}
\newcommand{\pd}[2]{\frac{\partial #1}{\partial #2}}
\newcommand{\pdd}[3]{\frac{\partial^{2} #1}{\partial #2\,\partial #3}}
\newcommand{\R}{\mathbb{R}}
\newcommand{\grad}{\nabla}
\DeclareMathOperator{\sgn}{sgn}
\DeclareMathOperator{\Hess}{H}

% German gloss, used on first mention of a technical term
\newcommand{\de}[1]{\emph{#1}}

% --- cross-references to chapters ------------------------------------------
% Chapter files are shared by the standalone volumes and the compendium, so a
% hard-coded "Chapter 6" is wrong in one of them. Write \chref{derivatives}{6}
% instead: it prints "Chapter 6" inside the same volume, "Volume 5, Chapter 6"
% from another volume, and the running number ("Chapter 34") in the
% compendium. \chrefs{vol}{n} is the abbreviated form (Ch. / Vol., Ch.),
% \chnum{vol}{n} the bare number in the current build, and \volref{vol} names
% the volume ("Volume 5", or "Part V" in the compendium). Each volume's
% main.tex sets \thisvolume; the compendium's sets \incompendiumtrue.
\newif\ifincompendium
\newcommand*{\thisvolume}{}
\makeatletter
\@namedef{chref@vol@algebra}{1}\@namedef{chref@off@algebra}{0}
\@namedef{chref@vol@functions}{2}\@namedef{chref@off@functions}{9}
\@namedef{chref@vol@matrices}{3}\@namedef{chref@off@matrices}{14}
\@namedef{chref@vol@foundations}{4}\@namedef{chref@off@foundations}{24}
\@namedef{chref@vol@derivatives}{5}\@namedef{chref@off@derivatives}{28}
\@namedef{chref@vol@integrals}{6}\@namedef{chref@off@integrals}{42}
\@namedef{chref@vol@optimization}{7}\@namedef{chref@off@optimization}{50}
\@namedef{chref@vol@linprog}{8}\@namedef{chref@off@linprog}{60}
\@namedef{chref@vol@mixed}{9}\@namedef{chref@off@mixed}{68}
\newcommand*{\chref@check}[1]{%
  \@ifundefined{chref@vol@#1}{\PackageError{chref}{Unknown volume `#1'}{}}{}}
\newcommand*{\chnum}[2]{\chref@check{#1}%
  \ifincompendium\the\numexpr\csname chref@off@#1\endcsname+#2\relax
  \else#2\fi}
% \chref@same{vol}{same volume}{other volume}
\newcommand*{\chref@same}[3]{\def\chref@tmp{#1}%
  \ifx\chref@tmp\thisvolume\expandafter\@firstoftwo
  \else\expandafter\@secondoftwo\fi{#2}{#3}}
\newcommand*{\chref}[2]{\chref@check{#1}%
  \ifincompendium Chapter~\chnum{#1}{#2}%
  \else\chref@same{#1}{Chapter~#2}%
    {Volume~\csname chref@vol@#1\endcsname, Chapter~#2}\fi}
\newcommand*{\chrefs}[2]{\chref@check{#1}%
  \ifincompendium Ch.~\chnum{#1}{#2}%
  \else\chref@same{#1}{Ch.~#2}%
    {Vol.~\csname chref@vol@#1\endcsname, Ch.~#2}\fi}
\newcommand*{\volref}[1]{\chref@check{#1}%
  \ifincompendium Part~\@Roman{\csname chref@vol@#1\endcsname}%
  \else Volume~\csname chref@vol@#1\endcsname\fi}
\makeatother

% --- Kindle edition: compact layout for a 6-inch screen -------------------
% Everything below applies only when \KindleEdition is defined on the
% command line; the A4 edition is untouched.
\ifdefined\KindleEdition
  \usepackage{etoolbox}
  \usepackage{adjustbox}
  \usepackage{titlesec}

  % The 8cm measure leaves few break points; let TeX loosen a tight line
  % rather than push it into the margin.
  \setlength{\emergencystretch}{1.5em}

  % No running heads: on a 90mm page the header costs more than it gives.
  % A centred page number in the foot is kept so the contents list means
  % something.
  \pagestyle{plain}
  \fancypagestyle{plain}{\fancyhf{}\fancyfoot[C]{\small\thepage}%
    \renewcommand{\headrulewidth}{0pt}}

  % Compact headings: the book class leaves a third of a page above a
  % chapter title, which is unaffordable here.
  \titleformat{\chapter}[display]
    {\normalfont\Large\bfseries\color{accent}}{\normalsize\chaptertitlename\ \thechapter}{0.4em}{\Large}
  \titlespacing*{\chapter}{0pt}{0pt}{1.2em}
  \titlespacing*{\section}{0pt}{1.4ex plus .3ex}{0.7ex}
  \titlespacing*{\subsection}{0pt}{1.2ex plus .2ex}{0.5ex}

  % Boxes: trim the padding, which was set for a 16cm measure.
  \tcbset{left=3.5pt,right=3.5pt,top=3pt,bottom=3pt,
          fonttitle=\bfseries\small,
          halign title=flush left}

  % Problem and solution lists: reclaim the wide left margin.
  \setlist[problems]{label=\textbf{\arabic*.}, leftmargin=1.5em,
    itemsep=0.45ex, parsep=0pt, topsep=0.6ex}
  \setlist[solutions]{label=\textbf{\arabic*.}, leftmargin=1.5em,
    itemsep=0.8ex, parsep=0pt, topsep=0.6ex}

  % Figures and wide tables shrink to the measure instead of running off it.
  % adjustbox only scales down: anything already narrow is left alone.
  \BeforeBeginEnvironment{tikzpicture}{\begin{adjustbox}{max width=\linewidth}}
  \AfterEndEnvironment{tikzpicture}{\end{adjustbox}}
  \BeforeBeginEnvironment{tabular}{\begin{adjustbox}{max width=\linewidth}}
  \AfterEndEnvironment{tabular}{\end{adjustbox}}

  % Display maths was written for a 16cm measure. On 8cm the long chained
  % equations would run off the screen, so any display wider than the text
  % block is scaled down to fit; narrower ones are left at full size. The
  % whole corpus uses \[ ... \] and no other display environment, so this
  % one redefinition covers every case.
  \newsavebox{\kndbox}
  \renewcommand{\[}{\par\addvspace{0.7ex}%
    \begin{lrbox}{\kndbox}$\displaystyle}
  \renewcommand{\]}{$\end{lrbox}%
    {\centering
      \ifdim\wd\kndbox>\linewidth
        \resizebox{\linewidth}{!}{\usebox{\kndbox}}%
      \else
        \usebox{\kndbox}%
      \fi
     \par}%
    \addvspace{0.7ex}}

  % A narrow measure makes perfect justification impossible; prefer a loose
  % line to a line that runs into the margin.
  \sloppy
  \raggedbottom
\fi

\renewcommand*{\thisvolume}{derivatives}

\title{\Huge Derivatives \& Partial Derivatives\\[0.4em]
  \large A Training Sheet with Exercises and Worked Solutions}
\author{}
\date{}

\begin{document}
\frontmatter
\maketitle

\chapter*{About this sheet}
\addcontentsline{toc}{chapter}{About this sheet}

This sheet drills the computation of derivatives in one and
several variables, with the applications most used in economics: marginal quantities, elasticities, tangents and linearisations, rates of
substitution, and the second-order conditions behind optimisation.

The structure is the same in every chapter: the \textbf{Core Rules} section
states the rules compactly and says when each one applies; \textbf{Practice
Problems} drills them; \textbf{Solutions} works every problem through in
full, showing the intermediate step rather than only the answer.

Chapters 1--8 cover the one-variable theory: the differentiation rules and
the standard functions, higher derivatives and convexity, Taylor
approximation, implicit differentiation, the classical applications, the
indeterminate forms with L'H\^opital's rule, and a chapter that assembles all
of it into the complete curve discussion. Chapters 9--14 are the
multivariable half, opening with the geometry --- surfaces, level curves and
the sub- and superlevel sets --- before partial derivatives and the gradient
(with the topological vocabulary of domains), the Hessian and definiteness,
the multivariable chain rule and total differential, directional derivatives
and vector-valued functions, and finally the Implicit Function Theorem with
Jacobians and Euler's theorem. Most chapters close with a block of
\emph{additional problems} for a second pass through the material.

Work the problems before reading the solutions. A derivative you have
recognised is worth more than ten you have read.

\tableofcontents

\mainmatter
\chapter{Basic Differentiation Rules}

The derivative measures the instantaneous rate of change of a function. In
economics this is the language of \emph{marginal} quantities: marginal cost,
marginal product, marginal revenue. Before specialising to particular
functional forms, master the structural rules below; almost every derivative
computed later is built from a handful of them.

\section{Core Rules}

\begin{rulebox}{Definition (limit form)}
For a function $f:\R\to\R$ the derivative at $x$ is
\[
  f'(x) \;=\; \lim_{h\to 0}\frac{f(x+h)-f(x)}{h},
\]
provided this limit exists. Equivalent notations include $\dd{f}{x}$,
$\dfrac{\mathrm{d}}{\mathrm{d}x}f(x)$, and $Df(x)$.
\end{rulebox}

\begin{rulebox}{Linearity}
For constants $c,\lambda\in\R$ and differentiable $u,v$:
\begin{itemize}[leftmargin=1.6em,itemsep=0.3ex]
  \item \textbf{Constant rule:} $\dfrac{\mathrm{d}}{\mathrm{d}x}(c)=0$
  \item \textbf{Scalar multiple:} $(\lambda u)'=\lambda u'$
  \item \textbf{Sum rule:} $(u+v)'=u'+v'$
\end{itemize}
\end{rulebox}

\begin{rulebox}{Power rule}
For any real exponent $n$ (where the expression is defined):
\[
  \frac{\mathrm{d}}{\mathrm{d}x}x^{n} \;=\; n\,x^{n-1} .
\]
Special cases worth memorising: $(\sqrt{x})'=\dfrac{1}{2\sqrt{x}}$,
$\Bigl(\dfrac{1}{x}\Bigr)'=-\dfrac{1}{x^{2}}$, and $(x)'=1$.
\end{rulebox}

\begin{rulebox}{Product rule}
If $u,v$ are differentiable,
\[
  (uv)' \;=\; u'v+uv' .
\]
\end{rulebox}

\begin{rulebox}{Quotient rule}
If $u,v$ are differentiable and $v\neq 0$,
\[
  \left(\frac{u}{v}\right)' \;=\; \frac{u'v-uv'}{v^{2}} .
\]
\end{rulebox}

\begin{rulebox}{Chain rule}
If $f$ is differentiable at $g(x)$ and $g$ is differentiable at $x$,
\[
  (f\circ g)'(x) \;=\; f'\bigl(g(x)\bigr)\cdot g'(x) .
\]
In Leibniz notation, with $y=f(u)$ and $u=g(x)$, this reads
$\dd{y}{x}=\dd{y}{u}\cdot\dd{u}{x}$.
\end{rulebox}

\begin{rulebox}{Higher-order derivatives (preview)}
The second derivative is $f''(x)=(f'(x))'$, written also
$\dfrac{\mathrm{d}^{2}f}{\mathrm{d}x^{2}}$. In economics $f'>0$ and $f''<0$
characterise an increasing, concave production or utility function. Higher
orders $f^{(n)}$ appear in Taylor expansions and in concavity tests.
\end{rulebox}

\section{Practice Problems}

Differentiate each function with respect to $x$ (or the indicated variable).
Simplify where natural.

\begin{problems}
  \item $f(x)=x^{7}$
  \item $f(x)=(3x+5)^{6}$
  \item $f(x)=\dfrac{x}{x+1}$
  \item $f(x)=(x^{2}+1)(x^{3}-2x)$
  \item $f(x)=x^{5}-4x^{3}+9x-2$
  \item $f(x)=\sqrt{2x^{2}+1}$
  \item $C(q)=100+4q+0.05q^{2}$ \quad (find $C'(q)$, the marginal cost)
  \item $f(x)=\dfrac{2x-3}{x^{2}+1}$
  \item $f(x)=(3x-1)(x^{2}+4x+5)$
  \item $f(x)=\sqrt{x}$
  \item $f(x)=(x^{2}-4x+1)^{4}$
  \item $f(x)=\dfrac{x^{2}+1}{x^{2}-1}$
  \item $f(x)=x\sqrt{x+1}$
  \item $f(x)=6x^{4}+2x^{2}-x$
  \item $f(x)=\bigl((x+1)^{2}+1\bigr)^{3}$ \quad (double chain)
  \item $f(x)=\dfrac{\sqrt{x}}{x+1}$
  \item $f(x)=(x^{2}+2)(\sqrt{x}-3)$
  \item $f(x)=x^{2/3}$
  \item $R(q)=20q-0.1q^{2}$ \quad (find $R'(q)$, the marginal revenue)
  \item $f(x)=\dfrac{1}{(x^{2}+3)^{5}}$
  \item $f(x)=(x^{3}-x)(2x^{2}+7)$
  \item $f(x)=\dfrac{x^{3}-2x}{x+4}$
  \item $f(x)=\dfrac{1}{x^{4}}$
  \item $f(x)=\bigl(\sqrt{x}+2x\bigr)^{4}$
  \item $f(x)=5x^{4}-\tfrac{1}{2}x^{2}+7$
  \item $\Pi(L)=L\cdot(10-L)$ \quad (a profit-like product)
  \item $f(x)=\sqrt[3]{(x^{2}+x+1)^{2}}$
  \item Average cost: $\displaystyle AC(q)=\frac{100+2q^{2}}{q}$
  \item $f(x)=3\sqrt[3]{x}+\dfrac{2}{x^{2}}$
  \item $f(x)=(5x^{2}-1)\Bigl(\dfrac{1}{x}+2\Bigr)$
  \item $f(x)=\dfrac{1-x^{2}}{1+x^{2}}$
  \item $f(x)=\pi x^{3}-\sqrt{2}\,x+e$
  \item $f(x)=\bigl(\sqrt{x}+1\bigr)^{2}$ \quad \emph{(expand first)}
  \item $f(x)=x^{1/2}+x^{-1/2}$
  \item $f(x)=\dfrac{x^{3}-5x+1}{x}$
  \item $f(x)=\tfrac{1}{3}x^{6}+\tfrac{2}{5}x^{5}-4$
  \item $f(x)=(x^{2}-9)^{5}$
  \item $f(x)=\dfrac{x+2}{x-2}$
  \item $f(x)=(4x+3)(x^{2}-x)$
  \item $f(x)=2x^{8}-3x^{5}+x$
  \item $f(x)=\sqrt{4-x^{2}}$
  \item Profit: $\Pi(q)=(20-0.1q)q-(100+4q+0.05q^{2})$ \quad (find
        $\Pi'(q)$)
  \item $f(x)=\dfrac{1}{2x+7}$
  \item $f(x)=(x^{2}+1)^{3}(x-1)$
  \item $f(x)=\dfrac{x^{2}}{1+x}$
  \item $f(x)=7\sqrt{x}-\dfrac{3}{x}$
  \item $f(x)=\bigl(1+x^{3}\bigr)^{10}$
  \item $f(x)=(x+1)^{2}(x-1)^{3}$
  \item Marginal product: $Q(L)=5L^{1/2}$ \quad (find $Q'(L)$)
  \item $f(x)=\dfrac{3x+1}{(x+2)^{2}}$
  \item $f(x)=x^{3}-5x^{2}+7x-4$
  \item $f(x)=\sqrt{x^{2}+9}$
  \item $f(x)=(2x^{3}+1)(x^{2}+x)$
  \item $f(x)=\dfrac{x^{2}-1}{x^{2}+1}$
  \item $f(x)=\Bigl(\dfrac{1}{x}+x\Bigr)^{4}$
  \item $f(x)=-3x^{5}+2x^{-3}$
  \item $f(x)=\dfrac{1}{\sqrt{x^{2}+1}}$
  \item $f(x)=(x-1)(x-2)(x-3)$
  \item Cost: $C(q)=50+\sqrt{q^{2}+4}$ \quad (find $C'(q)$)
  \item $f(x)=\dfrac{x+1}{x^{2}+2}$
  \item $f(x)=(5-2x)^{7}$
  \item $f(x)=x^{2}\sqrt{1+x}$
  \item $f(x)=\dfrac{1}{x^{3}+1}$
  \item $f(x)=(x^{4}+3x)^{6}$
  \item $f(x)=\tfrac{2}{3}x^{3/2}-4x^{-1/2}$
  \item Revenue per unit time: $R(t)=(50+2t)^{3}$ \quad (find $R'(t)$)
  \item $f(x)=\dfrac{\sqrt{x}+1}{\sqrt{x}-1}$
  \item $f(x)=(x^{2}+1)(x^{2}-1)(x+3)$
  \item $f(x)=\sqrt{(x+1)(x+2)}$
  \item $f(x)=\dfrac{x^{2}+2x}{(x+1)^{2}}$
  \item $f(x)=\bigl(3x^{2}-2\sqrt{x}\bigr)^{5}$
\end{problems}

\section{Solutions}

\begin{solutions}
  \item $f'(x)=7x^{6}$ (power rule).

  \item Outer $u^{6}$, inner $u=3x+5$. $f'(x)=6(3x+5)^{5}\cdot 3
        =18(3x+5)^{5}$.

  \item Quotient rule with $u=x$, $v=x+1$, $u'=1$, $v'=1$:
        \[
          f'(x)=\frac{1\cdot(x+1)-x\cdot 1}{(x+1)^{2}}=\frac{1}{(x+1)^{2}} .
        \]

  \item Let $u=x^{2}+1$, $v=x^{3}-2x$. Then $u'=2x$, $v'=3x^{2}-2$.
        \[
          f'(x)=2x(x^{3}-2x)+(x^{2}+1)(3x^{2}-2)=5x^{4}-3x^{2}-2 .
        \]

  \item $f'(x)=5x^{4}-12x^{2}+9$.

  \item $f(x)=(2x^{2}+1)^{1/2}$, inner $u=2x^{2}+1$, $u'=4x$.
        \[
          f'(x)=\tfrac12(2x^{2}+1)^{-1/2}\cdot 4x=\frac{2x}{\sqrt{2x^{2}+1}} .
        \]

  \item $C'(q)=4+0.1q$. Marginal cost is linear in output.

  \item $u=2x-3$, $v=x^{2}+1$, $u'=2$, $v'=2x$.
        \[
          f'(x)=\frac{2(x^{2}+1)-(2x-3)(2x)}{(x^{2}+1)^{2}}
               =\frac{-2x^{2}+6x+2}{(x^{2}+1)^{2}} .
        \]

  \item $u=3x-1$, $v=x^{2}+4x+5$, $u'=3$, $v'=2x+4$.
        \[
          f'(x)=3(x^{2}+4x+5)+(3x-1)(2x+4)=9x^{2}+22x+11 .
        \]

  \item Write $\sqrt{x}=x^{1/2}$, so $f'(x)=\tfrac12x^{-1/2}
        =\dfrac{1}{2\sqrt{x}}$.

  \item Inner $u=x^{2}-4x+1$, $u'=2x-4$. $f'(x)=4(x^{2}-4x+1)^{3}(2x-4)
        =8(x-2)(x^{2}-4x+1)^{3}$.

  \item $u=x^{2}+1$, $v=x^{2}-1$, $u'=2x$, $v'=2x$.
        \[
          f'(x)=\frac{2x(x^{2}-1)-(x^{2}+1)(2x)}{(x^{2}-1)^{2}}
               =\frac{-4x}{(x^{2}-1)^{2}} .
        \]

  \item $u=x$, $v=(x+1)^{1/2}$, $u'=1$, $v'=\tfrac12(x+1)^{-1/2}$.
        \[
          f'(x)=\sqrt{x+1}+\frac{x}{2\sqrt{x+1}}
               =\frac{2(x+1)+x}{2\sqrt{x+1}}=\frac{3x+2}{2\sqrt{x+1}} .
        \]

  \item $f'(x)=24x^{3}+4x-1$.

  \item Two layers. Let $g(x)=(x+1)^{2}+1$, $g'(x)=2(x+1)$. Then
        \[
          f'(x)=3\,g(x)^{2}\cdot g'(x)=6(x+1)\bigl((x+1)^{2}+1\bigr)^{2} .
        \]

  \item $u=x^{1/2}$, $v=x+1$, $u'=\tfrac12x^{-1/2}$, $v'=1$.
        \[
          f'(x)=\frac{\frac{1}{2\sqrt{x}}(x+1)-\sqrt{x}}{(x+1)^{2}}
               =\frac{(x+1)-2x}{2\sqrt{x}(x+1)^{2}}
               =\frac{1-x}{2\sqrt{x}(x+1)^{2}} .
        \]

  \item $u=x^{2}+2$, $v=x^{1/2}-3$, $u'=2x$, $v'=\tfrac12x^{-1/2}$.
        \[
          f'(x)=2x(\sqrt{x}-3)+(x^{2}+2)\cdot\frac{1}{2\sqrt{x}}
               =2x\sqrt{x}-6x+\frac{x^{2}+2}{2\sqrt{x}} .
        \]

  \item $f'(x)=\tfrac23x^{-1/3}=\dfrac{2}{3\sqrt[3]{x}}$.

  \item $R'(q)=20-0.2q$. Marginal revenue falls with $q$.

  \item $f(x)=(x^{2}+3)^{-5}$, inner $u'=2x$.
        \[
          f'(x)=-5(x^{2}+3)^{-6}\cdot 2x=-\frac{10x}{(x^{2}+3)^{6}} .
        \]

  \item $u=x^{3}-x$, $v=2x^{2}+7$, $u'=3x^{2}-1$, $v'=4x$.
        \[
          f'(x)=(3x^{2}-1)(2x^{2}+7)+(x^{3}-x)(4x)=10x^{4}+15x^{2}-7 .
        \]

  \item $u=x^{3}-2x$, $v=x+4$, $u'=3x^{2}-2$, $v'=1$.
        \[
          f'(x)=\frac{(3x^{2}-2)(x+4)-(x^{3}-2x)}{(x+4)^{2}}
               =\frac{2x^{3}+12x^{2}-8}{(x+4)^{2}} .
        \]

  \item Write $f(x)=x^{-4}$, so $f'(x)=-4x^{-5}=-\dfrac{4}{x^{5}}$.

  \item Inner $u=\sqrt{x}+2x$, $u'=\dfrac{1}{2\sqrt{x}}+2$.
        \[
          f'(x)=4\bigl(\sqrt{x}+2x\bigr)^{3}\left(\frac{1}{2\sqrt{x}}+2\right) .
        \]

  \item $f'(x)=20x^{3}-x$.

  \item $\Pi'(L)=1\cdot(10-L)+L\cdot(-1)=10-2L$.

  \item $f(x)=(x^{2}+x+1)^{2/3}$, inner $u'=2x+1$.
        \[
          f'(x)=\tfrac23(x^{2}+x+1)^{-1/3}(2x+1)
               =\frac{2(2x+1)}{3\sqrt[3]{x^{2}+x+1}} .
        \]

  \item Easier: $AC(q)=100q^{-1}+2q$, so $AC'(q)=-100q^{-2}+2
        =2-\dfrac{100}{q^{2}}$.

  \item $f(x)=3x^{1/3}+2x^{-2}$, hence $f'(x)=x^{-2/3}-4x^{-3}
        =\dfrac{1}{x^{2/3}}-\dfrac{4}{x^{3}}$.

  \item $u=5x^{2}-1$, $v=x^{-1}+2$, $u'=10x$, $v'=-x^{-2}$.
        \[
          f'(x)=10x(x^{-1}+2)+(5x^{2}-1)(-x^{-2})
               =10+20x-5+x^{-2}=5+20x+\frac{1}{x^{2}} .
        \]

  \item $u=1-x^{2}$, $v=1+x^{2}$, $u'=-2x$, $v'=2x$.
        \[
          f'(x)=\frac{-2x(1+x^{2})-(1-x^{2})(2x)}{(1+x^{2})^{2}}
               =\frac{-4x}{(1+x^{2})^{2}} .
        \]

  \item $f'(x)=3\pi x^{2}-\sqrt{2}$. ($e$ is a constant.)

  \item Expand: $f(x)=x+2\sqrt{x}+1$, so $f'(x)=1+\dfrac{1}{\sqrt{x}}$.

  \item $f'(x)=\tfrac12x^{-1/2}-\tfrac12x^{-3/2}
        =\dfrac{1}{2\sqrt{x}}-\dfrac{1}{2x^{3/2}}$.

  \item Simplify first: $f(x)=x^{2}-5+x^{-1}$, so $f'(x)=2x-x^{-2}
        =2x-\dfrac{1}{x^{2}}$.

  \item $f'(x)=2x^{5}+2x^{4}$.

  \item Chain rule with $u=x^{2}-9$, $u'=2x$:
        \[
          f'(x)=5(x^{2}-9)^{4}\cdot 2x=10x(x^{2}-9)^{4} .
        \]

  \item $u=x+2$, $v=x-2$, $u'=1$, $v'=1$.
        \[
          f'(x)=\frac{1\cdot(x-2)-(x+2)\cdot 1}{(x-2)^{2}}
               =\frac{-4}{(x-2)^{2}} .
        \]

  \item Product rule with $u=4x+3$, $v=x^{2}-x$, $u'=4$, $v'=2x-1$.
        \[
          f'(x)=4(x^{2}-x)+(4x+3)(2x-1)=12x^{2}-2x-3 .
        \]

  \item $f'(x)=16x^{7}-15x^{4}+1$.

  \item $f(x)=(4-x^{2})^{1/2}$, inner $u'=-2x$.
        \[
          f'(x)=\tfrac12(4-x^{2})^{-1/2}\cdot(-2x)=-\frac{x}{\sqrt{4-x^{2}}} .
        \]

  \item Expand: $\Pi(q)=20q-0.1q^{2}-100-4q-0.05q^{2}
        =-100+16q-0.15q^{2}$.
        \[
          \Pi'(q)=16-0.3q .
        \]
        Marginal profit is zero at $q^{*}=16/0.3\approx 53.33$.

  \item Write $f(x)=(2x+7)^{-1}$. Chain rule: $f'(x)=-(2x+7)^{-2}\cdot 2
        =-\dfrac{2}{(2x+7)^{2}}$.

  \item Product rule with $u=(x^{2}+1)^{3}$, $v=x-1$. $u'=3(x^{2}+1)^{2}
        \cdot 2x=6x(x^{2}+1)^{2}$, $v'=1$.
        \[
          f'(x)=6x(x^{2}+1)^{2}(x-1)+(x^{2}+1)^{3}
               =(x^{2}+1)^{2}\bigl[6x(x-1)+(x^{2}+1)\bigr]
               =(x^{2}+1)^{2}(7x^{2}-6x+1) .
        \]

  \item Quotient rule with $u=x^{2}$, $v=1+x$, $u'=2x$, $v'=1$.
        \[
          f'(x)=\frac{2x(1+x)-x^{2}}{(1+x)^{2}}=\frac{x^{2}+2x}{(1+x)^{2}}
               =\frac{x(x+2)}{(1+x)^{2}} .
        \]

  \item $f(x)=7x^{1/2}-3x^{-1}$, so $f'(x)=\tfrac72x^{-1/2}+3x^{-2}
        =\dfrac{7}{2\sqrt{x}}+\dfrac{3}{x^{2}}$.

  \item Inner $u=1+x^{3}$, $u'=3x^{2}$.
        \[
          f'(x)=10(1+x^{3})^{9}\cdot 3x^{2}=30x^{2}(1+x^{3})^{9} .
        \]

  \item Product rule with $u=(x+1)^{2}$, $v=(x-1)^{3}$, $u'=2(x+1)$,
        $v'=3(x-1)^{2}$.
        \[
          f'(x)=2(x+1)(x-1)^{3}+3(x+1)^{2}(x-1)^{2}
               =(x+1)(x-1)^{2}\bigl[2(x-1)+3(x+1)\bigr]
               =(x+1)(x-1)^{2}(5x+1) .
        \]

  \item $Q(L)=5L^{1/2}$, so $Q'(L)=\tfrac52L^{-1/2}=\dfrac{5}{2\sqrt{L}}$.
        Marginal product is positive but diminishing.

  \item Quotient rule with $u=3x+1$, $v=(x+2)^{2}$, $u'=3$,
        $v'=2(x+2)$.
        \[
          f'(x)=\frac{3(x+2)^{2}-(3x+1)\cdot 2(x+2)}{(x+2)^{4}}
               =\frac{3(x+2)-2(3x+1)}{(x+2)^{3}}=\frac{-3x+4}{(x+2)^{3}} .
        \]

  \item $f'(x)=3x^{2}-10x+7$.

  \item $f(x)=(x^{2}+9)^{1/2}$, $u'=2x$.
        \[
          f'(x)=\tfrac12(x^{2}+9)^{-1/2}\cdot 2x=\frac{x}{\sqrt{x^{2}+9}} .
        \]

  \item Product rule with $u=2x^{3}+1$, $v=x^{2}+x$, $u'=6x^{2}$,
        $v'=2x+1$.
        \[
          f'(x)=6x^{2}(x^{2}+x)+(2x^{3}+1)(2x+1)=10x^{4}+8x^{3}+2x+1 .
        \]

  \item Quotient rule with $u=x^{2}-1$, $v=x^{2}+1$, $u'=2x$, $v'=2x$.
        \[
          f'(x)=\frac{2x(x^{2}+1)-(x^{2}-1)(2x)}{(x^{2}+1)^{2}}
               =\frac{4x}{(x^{2}+1)^{2}} .
        \]

  \item Inner $u=x^{-1}+x$, $u'=-x^{-2}+1$.
        \[
          f'(x)=4\Bigl(\frac{1}{x}+x\Bigr)^{3}\left(1-\frac{1}{x^{2}}\right) .
        \]

  \item $f'(x)=-15x^{4}-6x^{-4}=-15x^{4}-\dfrac{6}{x^{4}}$.

  \item $f(x)=(x^{2}+1)^{-1/2}$, inner $u'=2x$.
        \[
          f'(x)=-\tfrac12(x^{2}+1)^{-3/2}\cdot 2x=-\frac{x}{(x^{2}+1)^{3/2}} .
        \]

  \item Expand or use the product rule iteratively. Quickest:
        $(x-1)(x-2)(x-3)=x^{3}-6x^{2}+11x-6$, so $f'(x)=3x^{2}-12x+11$.

  \item $C(q)=50+(q^{2}+4)^{1/2}$, $u'=2q$.
        \[
          C'(q)=\tfrac12(q^{2}+4)^{-1/2}\cdot 2q=\frac{q}{\sqrt{q^{2}+4}} .
        \]

  \item Quotient rule with $u=x+1$, $v=x^{2}+2$, $u'=1$, $v'=2x$.
        \[
          f'(x)=\frac{1\cdot(x^{2}+2)-(x+1)(2x)}{(x^{2}+2)^{2}}
               =\frac{-x^{2}-2x+2}{(x^{2}+2)^{2}} .
        \]

  \item Inner $u=5-2x$, $u'=-2$.
        \[
          f'(x)=7(5-2x)^{6}\cdot(-2)=-14(5-2x)^{6} .
        \]

  \item Product rule with $u=x^{2}$, $v=(1+x)^{1/2}$, $u'=2x$,
        $v'=\tfrac12(1+x)^{-1/2}$.
        \[
          f'(x)=2x\sqrt{1+x}+\frac{x^{2}}{2\sqrt{1+x}}
               =\frac{4x(1+x)+x^{2}}{2\sqrt{1+x}}=\frac{5x^{2}+4x}{2\sqrt{1+x}} .
        \]

  \item $f(x)=(x^{3}+1)^{-1}$, inner $u'=3x^{2}$.
        \[
          f'(x)=-(x^{3}+1)^{-2}\cdot 3x^{2}=-\frac{3x^{2}}{(x^{3}+1)^{2}} .
        \]

  \item Inner $u=x^{4}+3x$, $u'=4x^{3}+3$.
        \[
          f'(x)=6(x^{4}+3x)^{5}(4x^{3}+3) .
        \]

  \item $f'(x)=\tfrac23\cdot\tfrac32x^{1/2}-4\cdot(-\tfrac12)x^{-3/2}
        =x^{1/2}+2x^{-3/2}=\sqrt{x}+\dfrac{2}{x^{3/2}}$.

  \item Inner $u=50+2t$, $u'=2$.
        \[
          R'(t)=3(50+2t)^{2}\cdot 2=6(50+2t)^{2} .
        \]

  \item Quotient rule with $u=\sqrt{x}+1$, $v=\sqrt{x}-1$,
        $u'=v'=\dfrac{1}{2\sqrt{x}}$.
        \[
          f'(x)=\frac{\frac{1}{2\sqrt{x}}(\sqrt{x}-1)
               -(\sqrt{x}+1)\cdot\frac{1}{2\sqrt{x}}}{(\sqrt{x}-1)^{2}}
               =\frac{-1/\sqrt{x}}{(\sqrt{x}-1)^{2}}
               =-\frac{1}{\sqrt{x}(\sqrt{x}-1)^{2}} .
        \]

  \item Group: note $(x^{2}+1)(x^{2}-1)=x^{4}-1$, so
        $f(x)=(x^{4}-1)(x+3)=x^{5}+3x^{4}-x-3$. Hence
        $f'(x)=5x^{4}+12x^{3}-1$.

  \item $f(x)=\bigl((x+1)(x+2)\bigr)^{1/2}=(x^{2}+3x+2)^{1/2}$, inner
        $u'=2x+3$.
        \[
          f'(x)=\tfrac12(x^{2}+3x+2)^{-1/2}(2x+3)
               =\frac{2x+3}{2\sqrt{(x+1)(x+2)}} .
        \]

  \item Simplify first: $x^{2}+2x=x(x+2)$ and $(x+1)^{2}=x^{2}+2x+1$.
        Quotient rule with $u=x^{2}+2x$, $v=(x+1)^{2}$, $u'=2x+2=2(x+1)$,
        $v'=2(x+1)$.
        \[
          f'(x)=\frac{2(x+1)(x+1)^{2}-(x^{2}+2x)\cdot 2(x+1)}{(x+1)^{4}}
               =\frac{2(x+1)^{2}-2(x^{2}+2x)}{(x+1)^{3}}=\frac{2}{(x+1)^{3}} .
        \]

  \item Inner $u=3x^{2}-2x^{1/2}$, $u'=6x-x^{-1/2}$.
        \[
          f'(x)=5\bigl(3x^{2}-2\sqrt{x}\bigr)^{4}\left(6x-\frac{1}{\sqrt{x}}\right) .
        \]
\end{solutions}

\section{Additional Problems}

\begin{problems}
  \item $f(x)=4x^{5}-7x^{3}+2x-9$
  \item $f(x)=(2x+3)^{8}$
  \item $f(x)=\dfrac{x-1}{x+1}$
  \item $f(x)=(x^{3}-2)(x^{2}+5)$
  \item $f(x)=\sqrt{3x^{2}+4}$
  \item $f(x)=\dfrac{x^{2}}{x^{3}+1}$
  \item Cost: $C(q)=200+6q+0.02q^{3}$ \quad (find $C'(q)$)
  \item $f(x)=(x^{2}+x)^{5}$
  \item $f(x)=\dfrac{1}{(2x-5)^{3}}$
  \item $f(x)=x\sqrt{x^{2}+1}$
  \item $f(x)=\bigl((x^{2}+1)^{2}-3\bigr)^{4}$ \quad (double chain)
  \item Marginal product: $Q(L)=8L^{2/3}$ \quad (find $Q'(L)$)
  \item $f(x)=\dfrac{x^{2}+3x}{(x-1)^{2}}$
  \item $f(x)=\sqrt[4]{x^{3}}$
  \item Profit: $\Pi(q)=(30-0.2q)\,q-(50+2q+0.1q^{2})$ \quad (find
        $\Pi'(q)$)
\end{problems}

\section{Additional Solutions}

\begin{solutions}
  \item $f'(x)=20x^{4}-21x^{2}+2$.

  \item Chain rule: $f'(x)=8(2x+3)^{7}\cdot 2=16(2x+3)^{7}$.

  \item Quotient rule with $u=x-1$, $v=x+1$, $u'=1$, $v'=1$:
        \[
          f'(x)=\frac{(x+1)-(x-1)}{(x+1)^{2}}=\frac{2}{(x+1)^{2}} .
        \]

  \item Product rule: $u=x^{3}-2$, $v=x^{2}+5$, $u'=3x^{2}$, $v'=2x$.
        \[
          f'(x)=3x^{2}(x^{2}+5)+(x^{3}-2)(2x)=5x^{4}+15x^{2}-4x .
        \]

  \item $f(x)=(3x^{2}+4)^{1/2}$, inner $u'=6x$.
        \[
          f'(x)=\tfrac12(3x^{2}+4)^{-1/2}\cdot 6x=\frac{3x}{\sqrt{3x^{2}+4}} .
        \]

  \item Quotient rule with $u=x^{2}$, $v=x^{3}+1$, $u'=2x$, $v'=3x^{2}$.
        \[
          f'(x)=\frac{2x(x^{3}+1)-x^{2}\cdot 3x^{2}}{(x^{3}+1)^{2}}
               =\frac{-x^{4}+2x}{(x^{3}+1)^{2}}
               =\frac{x(2-x^{3})}{(x^{3}+1)^{2}} .
        \]

  \item $C'(q)=6+0.06q^{2}$. Marginal cost grows quadratically with output.

  \item Chain rule: inner $u=x^{2}+x$, $u'=2x+1$.
        $f'(x)=5(x^{2}+x)^{4}(2x+1)$.

  \item $f(x)=(2x-5)^{-3}$. Chain rule: $f'(x)=-3(2x-5)^{-4}\cdot 2
        =-\dfrac{6}{(2x-5)^{4}}$.

  \item Product rule with $u=x$, $v=(x^{2}+1)^{1/2}$, $u'=1$,
        $v'=x/(x^{2}+1)^{1/2}$.
        \[
          f'(x)=\sqrt{x^{2}+1}+\frac{x^{2}}{\sqrt{x^{2}+1}}
               =\frac{2x^{2}+1}{\sqrt{x^{2}+1}} .
        \]

  \item Two layers. Let $g(x)=(x^{2}+1)^{2}-3$,
        $g'(x)=2(x^{2}+1)\cdot 2x=4x(x^{2}+1)$.
        \[
          f'(x)=4\,g(x)^{3}\cdot g'(x)
               =16x(x^{2}+1)\bigl((x^{2}+1)^{2}-3\bigr)^{3} .
        \]

  \item $Q'(L)=8\cdot\tfrac23L^{-1/3}=\dfrac{16}{3L^{1/3}}$. Marginal
        product is positive but diminishing.

  \item Quotient rule with $u=x^{2}+3x$, $v=(x-1)^{2}$, $u'=2x+3$,
        $v'=2(x-1)$.
        \[
          f'(x)=\frac{(2x+3)(x-1)^{2}-(x^{2}+3x)\cdot 2(x-1)}{(x-1)^{4}}
               =\frac{(2x+3)(x-1)-2(x^{2}+3x)}{(x-1)^{3}}
               =\frac{-5x-3}{(x-1)^{3}} .
        \]

  \item $f(x)=x^{3/4}$, so $f'(x)=\tfrac34x^{-1/4}=\dfrac{3}{4\sqrt[4]{x}}$.

  \item Expand: $\Pi(q)=30q-0.2q^{2}-50-2q-0.1q^{2}
        =-50+28q-0.3q^{2}$.
        \[
          \Pi'(q)=28-0.6q .
        \]
        Marginal profit is zero at $q^{*}=28/0.6\approx 46.67$.
\end{solutions}

\chapter{Derivatives of Standard Functions}

The basic rules of Section 1 reduce every derivative to a finite list of
\emph{building blocks}: exponentials, logarithms, trigonometric, inverse
trigonometric and hyperbolic functions. Memorise the table below. In
economics, exponentials govern continuous growth and discounting, while
logarithms encode percentage changes and elasticities.

\section{Core Rules}

\begin{rulebox}{Exponential functions}
\begin{itemize}[leftmargin=1.6em,itemsep=0.3ex]
  \item $\dfrac{\mathrm{d}}{\mathrm{d}x}\mathrm{e}^{x}=\mathrm{e}^{x}$
  \item $\dfrac{\mathrm{d}}{\mathrm{d}x}\mathrm{e}^{ax}=a\,\mathrm{e}^{ax}$
        \quad ($a$ constant; chain rule)
  \item $\dfrac{\mathrm{d}}{\mathrm{d}x}a^{x}=a^{x}\ln a$ \quad
        ($a>0$, $a\neq 1$)
\end{itemize}
\end{rulebox}

\begin{rulebox}{Logarithmic functions}
\begin{itemize}[leftmargin=1.6em,itemsep=0.3ex]
  \item $\dfrac{\mathrm{d}}{\mathrm{d}x}\ln x=\dfrac{1}{x}$ \quad ($x>0$)
  \item $\dfrac{\mathrm{d}}{\mathrm{d}x}\ln|x|=\dfrac{1}{x}$ \quad
        ($x\neq 0$)
  \item $\dfrac{\mathrm{d}}{\mathrm{d}x}\log_{a}x=\dfrac{1}{x\ln a}$ \quad
        ($a>0$, $a\neq 1$)
\end{itemize}
\end{rulebox}

\begin{rulebox}{Trigonometric functions}
Angles in radians.
\begin{itemize}[leftmargin=1.6em,itemsep=0.3ex]
  \item $(\sin x)'=\cos x$
  \item $(\cos x)'=-\sin x$
  \item $(\tan x)'=\sec^{2}x=\dfrac{1}{\cos^{2}x}$
  \item $(\cot x)'=-\csc^{2}x=-\dfrac{1}{\sin^{2}x}$
  \item $(\sec x)'=\sec x\tan x$
  \item $(\csc x)'=-\csc x\cot x$
\end{itemize}
\end{rulebox}

\begin{rulebox}{Inverse trigonometric functions}
\begin{itemize}[leftmargin=1.6em,itemsep=0.3ex]
  \item $(\arcsin x)'=\dfrac{1}{\sqrt{1-x^{2}}}$, \quad $|x|<1$
  \item $(\arccos x)'=-\dfrac{1}{\sqrt{1-x^{2}}}$, \quad $|x|<1$
  \item $(\arctan x)'=\dfrac{1}{1+x^{2}}$
\end{itemize}
\end{rulebox}

\begin{rulebox}{Hyperbolic functions}
Recall $\sinh x=\tfrac12(\mathrm{e}^{x}-\mathrm{e}^{-x})$,
$\cosh x=\tfrac12(\mathrm{e}^{x}+\mathrm{e}^{-x})$.
\begin{itemize}[leftmargin=1.6em,itemsep=0.3ex]
  \item $(\sinh x)'=\cosh x$
  \item $(\cosh x)'=\sinh x$ \quad (note: \emph{no} sign change)
  \item $(\tanh x)'=\operatorname{sech}^{2}x=1-\tanh^{2}x$
\end{itemize}
\end{rulebox}

\begin{rulebox}{Summary table}
\[
\begin{array}{lll}
\hline
f(x) & f'(x) & \text{Domain note} \\
\hline
\mathrm{e}^{x} & \mathrm{e}^{x} & \R \\
a^{x} & a^{x}\ln a & a>0,\ a\neq 1 \\
\ln x & 1/x & x>0 \\
\log_{a}x & 1/(x\ln a) & x>0 \\
\sin x & \cos x & \R \\
\cos x & -\sin x & \R \\
\tan x & \sec^{2}x & x\neq \tfrac{\pi}{2}+k\pi \\
\cot x & -\csc^{2}x & x\neq k\pi \\
\sec x & \sec x\tan x & x\neq \tfrac{\pi}{2}+k\pi \\
\csc x & -\csc x\cot x & x\neq k\pi \\
\arcsin x & 1/\sqrt{1-x^{2}} & |x|<1 \\
\arccos x & -1/\sqrt{1-x^{2}} & |x|<1 \\
\arctan x & 1/(1+x^{2}) & \R \\
\sinh x & \cosh x & \R \\
\cosh x & \sinh x & \R \\
\tanh x & \operatorname{sech}^{2}x & \R \\
\hline
\end{array}
\]
Combine each entry with the chain rule whenever the argument is a
non-trivial function of $x$.
\end{rulebox}

\section{Practice Problems}

Differentiate each of the following with respect to $x$ (or the indicated
variable).

\begin{problems}
  \item $f(x)=\mathrm{e}^{3x}$
  \item $f(x)=\ln(5x)$
  \item $f(x)=\sin(3x)$
  \item $f(x)=\arcsin(2x)$
  \item $f(x)=\sinh(2x)$
  \item $f(x)=\mathrm{e}^{-x^{2}}$
  \item $f(x)=\cos(x^{2})$
  \item $f(x)=\ln(x^{2}+1)$
  \item $f(x)=x\,\mathrm{e}^{x}$
  \item $f(x)=\arctan(x^{2})$
  \item $f(x)=\tan(2x+1)$
  \item Continuous compounding: $V(t)=V_{0}\,\mathrm{e}^{rt}$ \quad
        (constants $V_{0},r$)
  \item $f(x)=x\ln x$
  \item $f(x)=\cosh(x^{2})$
  \item $f(x)=\mathrm{e}^{\sqrt{x}}$
  \item $f(x)=\arccos(\sqrt{x})$
  \item $f(x)=x\sin x$
  \item $f(x)=\ln(\sqrt{x})$
  \item $f(x)=\sin^{2}x$ \quad (i.e.\ $(\sin x)^{2}$)
  \item $f(x)=2^{x}$
  \item $f(x)=\mathrm{e}^{x^{2}+1}$
  \item $f(x)=(\ln x)^{2}$
  \item $f(x)=\dfrac{\sin x}{x}$
  \item $f(x)=x\arctan x$
  \item $f(x)=\tanh(3x)$
  \item $f(x)=\log_{10}(x)$
  \item $f(x)=\sin(\cos x)$
  \item $f(x)=\mathrm{e}^{\sin x}$
  \item $f(x)=\arctan\Bigl(\dfrac{1}{x}\Bigr)$, \quad $x>0$
  \item $f(x)=\ln\Bigl(\dfrac{x}{x+1}\Bigr)$
  \item $f(x)=\sec(3x)$
  \item Log-utility: $U(c)=\ln(c)$ \quad (find marginal utility)
  \item $f(x)=\ln(\cos x)$
  \item $f(x)=\tan^{3}x$
  \item Logistic: $P(t)=\dfrac{1}{1+\mathrm{e}^{-t}}$
  \item $f(x)=\dfrac{\mathrm{e}^{x}}{1+\mathrm{e}^{x}}$
  \item $f(x)=\mathrm{e}^{x}\sin(2x)$
  \item Cobb-Douglas log: $f(K)=\ln(K^{\alpha})$ \quad ($\alpha$ constant)
  \item $f(x)=x^{x}$ \quad (use logarithmic differentiation)
  \item $f(x)=\mathrm{e}^{2x}\cos(3x)$
  \item $f(x)=\ln\bigl(x+\sqrt{x^{2}+1}\bigr)$
  \item $f(x)=\sin(x)\cos(x)$
  \item Discount factor: $D(t)=\mathrm{e}^{-rt}$ \quad (find $D'(t)$ and
        interpret)
  \item $f(x)=\arctan(\mathrm{e}^{x})$
  \item $f(x)=\tanh(x^{2}+1)$
  \item $f(x)=\dfrac{\ln x}{x}$
  \item $f(x)=\mathrm{e}^{x}\ln x$
  \item $f(x)=\sqrt{\sin x}$
  \item Exponential growth: $N(t)=1000\cdot\mathrm{e}^{0.04t}$ \quad (find
        $N'(t)$ and the growth rate $N'/N$)
  \item $f(x)=5^{x^{2}}$
  \item $f(x)=\arcsin(x^{2})$
  \item $f(x)=\cot(x)+\csc(x)$
  \item $f(x)=\sinh(x)\cosh(x)$
  \item Log of a product: $f(x)=\ln\bigl(x\sqrt{x^{2}+1}\bigr)$
  \item $f(x)=\cos^{4}(x)$
  \item $f(x)=\mathrm{e}^{x^{3}}$
  \item Iso-elastic utility: $\displaystyle U(c)=\frac{c^{1-\sigma}-1}{1-\sigma}$,
        $\sigma\neq 1$ \quad (find $U'(c)$)
  \item $f(x)=\arctan(\sqrt{x})$
  \item $f(x)=\tan(x)\ln(x)$
  \item $f(x)=\dfrac{1}{\ln x}$
  \item $f(x)=\mathrm{e}^{ax}\sin(bx)$ \quad (constants $a,b$)
  \item $f(x)=\log_{2}(x^{2}+1)$
  \item $f(x)=\cos(2x)\mathrm{e}^{-x}$
  \item $f(x)=\operatorname{arccot}(x)\cdot\arctan(x)$ \quad [use
        $\operatorname{arccot}(x)=\tfrac{\pi}{2}-\arctan(x)$]
  \item $f(x)=\ln(\ln(x))$, \quad $x>1$
  \item $f(x)=x^{2}\mathrm{e}^{-x}$
  \item Per-capita growth: $g(t)=\ln(N(t))$ where $N(t)=N_{0}\mathrm{e}^{rt}$
  \item $f(x)=\sin(\mathrm{e}^{x})$
  \item $f(x)=(1+\mathrm{e}^{x})^{3}$
  \item $f(x)=x^{\sin x}$ \quad (use logarithmic differentiation)
  \item $f(x)=\arctan(\tan x)$ \quad (restrict $x\in(-\tfrac{\pi}{2},
        \tfrac{\pi}{2})$)
\end{problems}

\section{Solutions}

\begin{solutions}
  \item Inner $u=3x$, $u'=3$. $f'(x)=\mathrm{e}^{3x}\cdot 3=3\mathrm{e}^{3x}$.

  \item Either expand $\ln(5x)=\ln 5+\ln x$, or chain:
        $f'(x)=\dfrac{1}{5x}\cdot 5=\dfrac{1}{x}$.

  \item Inner $u=3x$, $u'=3$. $f'(x)=\cos(3x)\cdot 3=3\cos(3x)$.

  \item Inner $u=2x$, $u'=2$.
        \[
          f'(x)=\frac{1}{\sqrt{1-(2x)^{2}}}\cdot 2=\frac{2}{\sqrt{1-4x^{2}}} .
        \]

  \item Inner $u=2x$, $u'=2$. $f'(x)=\cosh(2x)\cdot 2=2\cosh(2x)$.

  \item Inner $u=-x^{2}$, $u'=-2x$.
        $f'(x)=\mathrm{e}^{-x^{2}}\cdot(-2x)=-2x\,\mathrm{e}^{-x^{2}}$.

  \item Inner $u=x^{2}$, $u'=2x$.
        $f'(x)=-\sin(x^{2})\cdot 2x=-2x\sin(x^{2})$.

  \item Inner $u=x^{2}+1$, $u'=2x$. $f'(x)=\dfrac{2x}{x^{2}+1}$.

  \item Product rule: $f'(x)=1\cdot\mathrm{e}^{x}+x\cdot\mathrm{e}^{x}
        =(1+x)\,\mathrm{e}^{x}$.

  \item Inner $u=x^{2}$, $u'=2x$. $f'(x)=\dfrac{2x}{1+x^{4}}$.

  \item Inner $u=2x+1$, $u'=2$. $f'(x)=2\sec^{2}(2x+1)$.

  \item $V'(t)=V_{0}\cdot r\,\mathrm{e}^{rt}=r\,V(t)$. The growth rate is
        constant: $V'/V=r$.

  \item Product rule: $f'(x)=1\cdot\ln x+x\cdot\dfrac{1}{x}=\ln x+1$.

  \item Inner $u=x^{2}$, $u'=2x$. $f'(x)=\sinh(x^{2})\cdot 2x
        =2x\sinh(x^{2})$.

  \item Inner $u=\sqrt{x}$, $u'=\dfrac{1}{2\sqrt{x}}$.
        \[
          f'(x)=\mathrm{e}^{\sqrt{x}}\cdot\frac{1}{2\sqrt{x}}
               =\frac{\mathrm{e}^{\sqrt{x}}}{2\sqrt{x}} .
        \]

  \item Inner $u=\sqrt{x}=x^{1/2}$, $u'=\dfrac{1}{2\sqrt{x}}$.
        \[
          f'(x)=-\frac{1}{\sqrt{1-x}}\cdot\frac{1}{2\sqrt{x}}
               =-\frac{1}{2\sqrt{x-x^{2}}} .
        \]

  \item Product rule: $f'(x)=\sin x+x\cos x$.

  \item Simplify: $\ln(\sqrt{x})=\tfrac12\ln x$, so $f'(x)=\dfrac{1}{2x}$.

  \item Chain: $f'(x)=2\sin x\cdot\cos x=\sin(2x)$.

  \item $f'(x)=2^{x}\ln 2$ (general exponential rule).

  \item Inner $u=x^{2}+1$, $u'=2x$. $f'(x)=2x\,\mathrm{e}^{x^{2}+1}$.

  \item Chain on outer square: $f'(x)=2\ln x\cdot\dfrac{1}{x}
        =\dfrac{2\ln x}{x}$.

  \item Quotient rule:
        \[
          f'(x)=\frac{\cos x\cdot x-\sin x\cdot 1}{x^{2}}
               =\frac{x\cos x-\sin x}{x^{2}} .
        \]

  \item Product rule: $f'(x)=\arctan x+x\cdot\dfrac{1}{1+x^{2}}
        =\arctan x+\dfrac{x}{1+x^{2}}$.

  \item Inner $u=3x$, $u'=3$. $f'(x)=\operatorname{sech}^{2}(3x)\cdot 3
        =3\operatorname{sech}^{2}(3x)$.

  \item $f'(x)=\dfrac{1}{x\ln 10}$.

  \item Outer $\sin$, inner $\cos x$: $f'(x)=\cos(\cos x)\cdot(-\sin x)
        =-\sin x\cos(\cos x)$.

  \item Chain: $f'(x)=\mathrm{e}^{\sin x}\cdot\cos x=\cos x\,\mathrm{e}^{\sin x}$.

  \item Inner $u=1/x$, $u'=-1/x^{2}$.
        \[
          f'(x)=\frac{1}{1+(1/x)^{2}}\cdot\left(-\frac{1}{x^{2}}\right)
               =-\frac{1}{x^{2}+1} .
        \]
        (Aside: this confirms $\arctan x+\arctan(1/x)=\tfrac{\pi}{2}$ for
        $x>0$.)

  \item Expand $\ln\dfrac{x}{x+1}=\ln x-\ln(x+1)$, so
        $f'(x)=\dfrac{1}{x}-\dfrac{1}{x+1}=\dfrac{1}{x(x+1)}$.

  \item Inner $u=3x$, $u'=3$. $f'(x)=\sec(3x)\tan(3x)\cdot 3
        =3\sec(3x)\tan(3x)$.

  \item $U'(c)=\dfrac{1}{c}$. Marginal utility of consumption falls as $c$
        grows, diminishing marginal utility.

  \item Outer $\ln$, inner $\cos x$. $f'(x)=\dfrac{1}{\cos x}\cdot(-\sin x)
        =-\tan x$.

  \item Outer $u^{3}$, inner $\tan x$: $f'(x)=3\tan^{2}x\cdot\sec^{2}x$.

  \item Quotient rule with $u=1$, $v=1+\mathrm{e}^{-t}$. $u'=0$,
        $v'=-\mathrm{e}^{-t}$.
        \[
          P'(t)=\frac{\mathrm{e}^{-t}}{(1+\mathrm{e}^{-t})^{2}} .
        \]
        Equivalently, $P'(t)=P(t)\bigl(1-P(t)\bigr)$.

  \item Quotient rule. $u=\mathrm{e}^{x}$, $v=1+\mathrm{e}^{x}$,
        $u'=\mathrm{e}^{x}$, $v'=\mathrm{e}^{x}$.
        \[
          f'(x)=\frac{\mathrm{e}^{x}(1+\mathrm{e}^{x})-\mathrm{e}^{x}\cdot\mathrm{e}^{x}}
               {(1+\mathrm{e}^{x})^{2}}=\frac{\mathrm{e}^{x}}{(1+\mathrm{e}^{x})^{2}} .
        \]

  \item Product rule: $u=\mathrm{e}^{x}$, $v=\sin(2x)$, $u'=\mathrm{e}^{x}$,
        $v'=2\cos(2x)$.
        \[
          f'(x)=\mathrm{e}^{x}\bigl(\sin(2x)+2\cos(2x)\bigr) .
        \]

  \item Use the log power rule first: $f(K)=\alpha\ln K$, hence
        $f'(K)=\dfrac{\alpha}{K}$. (This is the elasticity-style structure
        that makes Cobb-Douglas convenient.)

  \item Logarithmic differentiation. Let $y=x^{x}$. Then $\ln y=x\ln x$.
        Differentiate:
        \[
          \frac{y'}{y}=\ln x+1,\qquad y'=x^{x}(\ln x+1) .
        \]

  \item Product rule with $u=\mathrm{e}^{2x}$, $v=\cos(3x)$,
        $u'=2\mathrm{e}^{2x}$, $v'=-3\sin(3x)$.
        \[
          f'(x)=2\mathrm{e}^{2x}\cos(3x)-3\mathrm{e}^{2x}\sin(3x)
               =\mathrm{e}^{2x}\bigl(2\cos(3x)-3\sin(3x)\bigr) .
        \]

  \item Inner $u=x+\sqrt{x^{2}+1}$. Note $u'=1+\dfrac{x}{\sqrt{x^{2}+1}}
        =\dfrac{\sqrt{x^{2}+1}+x}{\sqrt{x^{2}+1}}=\dfrac{u}{\sqrt{x^{2}+1}}$.
        Hence
        \[
          f'(x)=\frac{1}{u}\cdot\frac{u}{\sqrt{x^{2}+1}}=\frac{1}{\sqrt{x^{2}+1}} .
        \]
        (This is $\operatorname{arsinh}(x)$, so the result matches the
        derivative of inverse hyperbolic sine.)

  \item Product rule: $f'(x)=\cos x\cdot\cos x+\sin x\cdot(-\sin x)
        =\cos^{2}x-\sin^{2}x=\cos(2x)$.

  \item $D'(t)=-r\,\mathrm{e}^{-rt}=-r\,D(t)$. Discount factor decays at
        rate $r$; instantaneous proportional decay $D'/D=-r$.

  \item Inner $u=\mathrm{e}^{x}$, $u'=\mathrm{e}^{x}$.
        \[
          f'(x)=\frac{1}{1+\mathrm{e}^{2x}}\cdot\mathrm{e}^{x}
               =\frac{\mathrm{e}^{x}}{1+\mathrm{e}^{2x}} .
        \]

  \item Inner $u=x^{2}+1$, $u'=2x$. $f'(x)=\operatorname{sech}^{2}(x^{2}+1)
        \cdot 2x=2x\operatorname{sech}^{2}(x^{2}+1)$.

  \item Quotient rule with $u=\ln x$, $v=x$, $u'=1/x$, $v'=1$.
        \[
          f'(x)=\frac{(1/x)\cdot x-\ln x}{x^{2}}=\frac{1-\ln x}{x^{2}} .
        \]

  \item Product rule: $f'(x)=\mathrm{e}^{x}\ln x+\mathrm{e}^{x}\cdot\dfrac{1}{x}
        =\mathrm{e}^{x}\Bigl(\ln x+\dfrac{1}{x}\Bigr)$.

  \item $f(x)=(\sin x)^{1/2}$, inner $u'=\cos x$.
        \[
          f'(x)=\tfrac12(\sin x)^{-1/2}\cos x=\frac{\cos x}{2\sqrt{\sin x}} .
        \]

  \item $N'(t)=1000\cdot 0.04\,\mathrm{e}^{0.04t}=0.04\,N(t)$, so growth
        rate $N'/N=0.04=4\%$ per unit time.

  \item Rewrite $5^{x^{2}}=\mathrm{e}^{x^{2}\ln 5}$, so
        $f'(x)=\mathrm{e}^{x^{2}\ln 5}\cdot 2x\ln 5=2x\ln 5\cdot 5^{x^{2}}$.

  \item Inner $u=x^{2}$, $u'=2x$.
        \[
          f'(x)=\frac{2x}{\sqrt{1-x^{4}}} .
        \]

  \item $f'(x)=-\csc^{2}x-\csc x\cot x=-\csc x(\csc x+\cot x)$.

  \item Product rule: $f'(x)=\cosh^{2}x+\sinh^{2}x=\cosh(2x)$. (Use the
        hyperbolic double-angle identity.)

  \item Simplify first: $f(x)=\ln x+\tfrac12\ln(x^{2}+1)$, so
        \[
          f'(x)=\frac{1}{x}+\frac12\cdot\frac{2x}{x^{2}+1}
               =\frac{1}{x}+\frac{x}{x^{2}+1} .
        \]

  \item Outer $u^{4}$, inner $\cos x$, $u'=-\sin x$: $f'(x)=4\cos^{3}x\cdot
        (-\sin x)=-4\cos^{3}x\sin x$.

  \item Inner $u=x^{3}$, $u'=3x^{2}$. $f'(x)=3x^{2}\mathrm{e}^{x^{3}}$.

  \item $U'(c)=\dfrac{(1-\sigma)c^{-\sigma}}{1-\sigma}=c^{-\sigma}$.
        Marginal utility is positive and decreasing for $\sigma>0$.

  \item Inner $u=\sqrt{x}$, $u'=\dfrac{1}{2\sqrt{x}}$.
        \[
          f'(x)=\frac{1}{1+x}\cdot\frac{1}{2\sqrt{x}}
               =\frac{1}{2\sqrt{x}(1+x)} .
        \]

  \item Product rule: $f'(x)=\sec^{2}x\cdot\ln x+\tan x\cdot\dfrac{1}{x}$.

  \item $f(x)=(\ln x)^{-1}$, so $f'(x)=-(\ln x)^{-2}\cdot\dfrac{1}{x}
        =-\dfrac{1}{x(\ln x)^{2}}$.

  \item Product rule with $u=\mathrm{e}^{ax}$, $v=\sin(bx)$,
        $u'=a\mathrm{e}^{ax}$, $v'=b\cos(bx)$.
        \[
          f'(x)=\mathrm{e}^{ax}\bigl(a\sin(bx)+b\cos(bx)\bigr) .
        \]

  \item Use $\log_{2}(u)=\dfrac{\ln u}{\ln 2}$ with $u=x^{2}+1$:
        \[
          f'(x)=\frac{1}{\ln 2}\cdot\frac{2x}{x^{2}+1}
               =\frac{2x}{(x^{2}+1)\ln 2} .
        \]

  \item Product rule with $u=\cos(2x)$, $v=\mathrm{e}^{-x}$,
        $u'=-2\sin(2x)$, $v'=-\mathrm{e}^{-x}$.
        \[
          f'(x)=-2\sin(2x)\mathrm{e}^{-x}-\cos(2x)\mathrm{e}^{-x}
               =-\mathrm{e}^{-x}\bigl(2\sin(2x)+\cos(2x)\bigr) .
        \]

  \item Using $\operatorname{arccot}(x)=\tfrac{\pi}{2}-\arctan(x)$:
        \[
          f(x)=\Bigl(\frac{\pi}{2}-\arctan x\Bigr)\arctan x
              =\frac{\pi}{2}\arctan x-(\arctan x)^{2} .
        \]
        \[
          f'(x)=\frac{\pi}{2}\cdot\frac{1}{1+x^{2}}
               -2\arctan x\cdot\frac{1}{1+x^{2}}
               =\frac{\frac{\pi}{2}-2\arctan x}{1+x^{2}} .
        \]

  \item Outer $\ln$, inner $\ln x$. $f'(x)=\dfrac{1}{\ln x}\cdot\dfrac{1}{x}
        =\dfrac{1}{x\ln x}$.

  \item Product rule: $f'(x)=2x\,\mathrm{e}^{-x}+x^{2}(-\mathrm{e}^{-x})
        =x\,\mathrm{e}^{-x}(2-x)$.

  \item $g(t)=\ln(N_{0}\mathrm{e}^{rt})=\ln N_{0}+rt$, so $g'(t)=r$. Log
        of an exponential is linear: continuous growth rate is the slope
        of $\ln N$.

  \item Inner $u=\mathrm{e}^{x}$, $u'=\mathrm{e}^{x}$.
        $f'(x)=\cos(\mathrm{e}^{x})\cdot\mathrm{e}^{x}
        =\mathrm{e}^{x}\cos(\mathrm{e}^{x})$.

  \item Inner $u=1+\mathrm{e}^{x}$, $u'=\mathrm{e}^{x}$.
        $f'(x)=3(1+\mathrm{e}^{x})^{2}\cdot\mathrm{e}^{x}
        =3\mathrm{e}^{x}(1+\mathrm{e}^{x})^{2}$.

  \item Logarithmic differentiation. $\ln y=\sin x\cdot\ln x$.
        Differentiate:
        \[
          \frac{y'}{y}=\cos x\cdot\ln x+\sin x\cdot\frac{1}{x},
        \]
        so $y'=x^{\sin x}\Bigl(\cos x\ln x+\dfrac{\sin x}{x}\Bigr)$.

  \item For $x\in(-\tfrac{\pi}{2},\tfrac{\pi}{2})$ we have
        $\arctan(\tan x)=x$, so $f'(x)=1$. (Equivalently, chain rule:
        $\dfrac{1}{1+\tan^{2}x}\cdot\sec^{2}x=\dfrac{\sec^{2}x}{\sec^{2}x}=1$.)
\end{solutions}

\section{Additional Problems}

\begin{problems}
  \item $f(x)=\mathrm{e}^{5x}$
  \item $f(x)=\ln(7x+3)$
  \item $f(x)=\cos(4x)$
  \item $f(x)=\arctan(3x)$
  \item $f(x)=\mathrm{e}^{x^{3}}$
  \item $f(x)=\sin^{3}x$
  \item $f(x)=\ln(x^{3}+2)$
  \item $f(x)=x^{2}\,\mathrm{e}^{-x}$
  \item Continuous discount: $V(t)=1000\,\mathrm{e}^{-0.05t}$ \quad (find
        $V'(t)$)
  \item $f(x)=\dfrac{\ln x}{x^{2}}$
  \item $f(x)=\mathrm{e}^{x}\cos x$
  \item $f(x)=\arcsin(3x)$
  \item $f(x)=\tan(x^{2})$
  \item $f(x)=\ln(\sin x)$
  \item Log-utility variant: $U(x)=\ln(1+x)$ \quad (find $U'(x)$)
  \item $f(x)=3^{x^{2}+1}$
  \item $f(x)=(\ln x)^{3}$
  \item $f(x)=\mathrm{e}^{x^{2}}\sin x$
  \item $f(x)=\dfrac{\mathrm{e}^{x}-1}{\mathrm{e}^{x}+1}$
  \item $f(x)=(\cos x)^{x}$ \quad (use logarithmic differentiation,
        $\cos x>0$)
  \item $f(x)=\sinh(x)-\cosh(x)$
  \item $f(x)=\ln\Bigl(\dfrac{x^{2}}{x+1}\Bigr)$
  \item Per-capita capital: $k(t)=\ln(K(t)/L(t))$ with
        $K(t)=K_{0}\mathrm{e}^{at}$, $L(t)=L_{0}\mathrm{e}^{bt}$ \quad
        (find $k'(t)$)
  \item $f(x)=\arctan(\ln x)$
  \item $f(x)=\mathrm{e}^{\cos x}$
\end{problems}

\section{Additional Solutions}

\begin{solutions}
  \item Chain: $f'(x)=5\mathrm{e}^{5x}$.

  \item Chain: $f'(x)=\dfrac{7}{7x+3}$.

  \item Chain: $f'(x)=-4\sin(4x)$.

  \item Chain: $f'(x)=\dfrac{3}{1+9x^{2}}$.

  \item Chain: inner $u=x^{3}$, $u'=3x^{2}$. $f'(x)=3x^{2}\mathrm{e}^{x^{3}}$.

  \item Chain: $f'(x)=3\sin^{2}(x)\cos(x)$.

  \item Chain: $f'(x)=\dfrac{3x^{2}}{x^{3}+2}$.

  \item Product rule: $f'(x)=2x\,\mathrm{e}^{-x}-x^{2}\mathrm{e}^{-x}
        =x(2-x)\,\mathrm{e}^{-x}$.

  \item $V'(t)=1000\cdot(-0.05)\,\mathrm{e}^{-0.05t}=-50\,\mathrm{e}^{-0.05t}
        =-0.05\,V(t)$. Decays at rate $5\%$.

  \item Quotient rule with $u=\ln x$, $v=x^{2}$, $u'=1/x$, $v'=2x$.
        \[
          f'(x)=\frac{(1/x)x^{2}-\ln x\cdot 2x}{x^{4}}
               =\frac{x-2x\ln x}{x^{4}}=\frac{1-2\ln x}{x^{3}} .
        \]

  \item Product rule: $f'(x)=\mathrm{e}^{x}\cos x-\mathrm{e}^{x}\sin x
        =\mathrm{e}^{x}(\cos x-\sin x)$.

  \item Chain: $f'(x)=\dfrac{3}{\sqrt{1-9x^{2}}}$.

  \item Chain: inner $u=x^{2}$, $u'=2x$. $f'(x)=2x\sec^{2}(x^{2})$.

  \item Chain: $f'(x)=\dfrac{\cos x}{\sin x}=\cot x$.

  \item $U'(x)=\dfrac{1}{1+x}$. Marginal utility positive but diminishing
        as $x$ grows.

  \item Write $3^{x^{2}+1}=\mathrm{e}^{(x^{2}+1)\ln 3}$.
        $f'(x)=\mathrm{e}^{(x^{2}+1)\ln 3}\cdot 2x\ln 3
        =2x\ln 3\cdot 3^{x^{2}+1}$.

  \item Chain: $f'(x)=3(\ln x)^{2}\cdot\dfrac{1}{x}=\dfrac{3(\ln x)^{2}}{x}$.

  \item Product rule with $u=\mathrm{e}^{x^{2}}$, $v=\sin x$,
        $u'=2x\,\mathrm{e}^{x^{2}}$, $v'=\cos x$.
        \[
          f'(x)=2x\,\mathrm{e}^{x^{2}}\sin x+\mathrm{e}^{x^{2}}\cos x
               =\mathrm{e}^{x^{2}}(2x\sin x+\cos x) .
        \]

  \item Quotient rule with $u=\mathrm{e}^{x}-1$, $v=\mathrm{e}^{x}+1$,
        $u'=\mathrm{e}^{x}$, $v'=\mathrm{e}^{x}$.
        \[
          f'(x)=\frac{\mathrm{e}^{x}(\mathrm{e}^{x}+1)
                -(\mathrm{e}^{x}-1)\mathrm{e}^{x}}{(\mathrm{e}^{x}+1)^{2}}
               =\frac{2\mathrm{e}^{x}}{(\mathrm{e}^{x}+1)^{2}} .
        \]

  \item Logarithmic differentiation. $\ln y=x\ln(\cos x)$. Differentiate:
        \[
          \frac{y'}{y}=\ln(\cos x)+x\cdot\frac{-\sin x}{\cos x}
                      =\ln(\cos x)-x\tan x .
        \]
        So $y'=(\cos x)^{x}\bigl[\ln(\cos x)-x\tan x\bigr]$.

  \item $f'(x)=\cosh x-\sinh x$. (Equivalently, $f(x)=\mathrm{e}^{-x}$ and
        $f'(x)=-\mathrm{e}^{-x}=\cosh x-\sinh x$.)

  \item Simplify: $f(x)=2\ln x-\ln(x+1)$, so
        $f'(x)=\dfrac{2}{x}-\dfrac{1}{x+1}=\dfrac{x+2}{x(x+1)}$.

  \item $k(t)=\ln K_{0}+at-\ln L_{0}-bt=(a-b)t+\text{const}$, so
        $k'(t)=a-b$ (constant; difference of growth rates).

  \item Chain: outer $\arctan$, inner $\ln x$.
        $f'(x)=\dfrac{1}{1+(\ln x)^{2}}\cdot\dfrac{1}{x}
        =\dfrac{1}{x(1+(\ln x)^{2})}$.

  \item Chain: $f'(x)=\mathrm{e}^{\cos x}\cdot(-\sin x)
        =-\sin x\,\mathrm{e}^{\cos x}$.
\end{solutions}

\chapter{Higher-Order Derivatives, Convexity, Inflection}

Once $f'$ exists and is itself differentiable, we can differentiate again to
obtain $f''$, and iterate to higher orders. The second derivative records
how the slope changes: it controls convexity/concavity and powers the
second-order condition for optimisation. Higher derivatives of standard
functions follow simple closed-form patterns, and the Leibniz rule extends
the product rule to all orders.

\section{Core Rules}

\begin{rulebox}{Definition}
For $f:I\subseteq\R\to\R$ smooth enough, write
\[
  f''(x)=\frac{\mathrm{d}^{2}f}{\mathrm{d}x^{2}},\qquad
  f^{(n)}(x)=\frac{\mathrm{d}^{n}f}{\mathrm{d}x^{n}},\qquad
  f^{(0)}:=f .
\]
\end{rulebox}

\begin{rulebox}{Polynomials}
\[
  \frac{\mathrm{d}^{n}}{\mathrm{d}x^{n}}x^{k}
    =k(k-1)\cdots(k-n+1)\,x^{k-n} .
\]
In particular $(x^{n})^{(n)}=n!$ and $(x^{n})^{(n+1)}=0$.
\end{rulebox}

\begin{rulebox}{Exponentials}
\[
  \frac{\mathrm{d}^{n}}{\mathrm{d}x^{n}}\mathrm{e}^{ax}=a^{n}\mathrm{e}^{ax} .
\]
\end{rulebox}

\begin{rulebox}{Logarithm}
\[
  \frac{\mathrm{d}^{n}}{\mathrm{d}x^{n}}\ln(x)
    =\frac{(-1)^{n-1}(n-1)!}{x^{n}}\qquad\text{for } n\geq 1 .
\]
\end{rulebox}

\begin{rulebox}{Trig}
\[
  \frac{\mathrm{d}^{n}}{\mathrm{d}x^{n}}\sin(x)=\sin\Bigl(x+\frac{n\pi}{2}\Bigr),
  \qquad
  \frac{\mathrm{d}^{n}}{\mathrm{d}x^{n}}\cos(x)=\cos\Bigl(x+\frac{n\pi}{2}\Bigr) .
\]
\end{rulebox}

\begin{rulebox}{Convexity / concavity}
$f''\geq 0$ on $I$ implies $f$ is convex on $I$; $f''\leq 0$ implies
concave. Strict inequalities give strict (con)cavity.
\end{rulebox}

\begin{rulebox}{Inflection point}
$x_{0}$ is an inflection point if $f''$ \emph{changes sign} at $x_{0}$. The
condition $f''(x_{0})=0$ alone is insufficient (e.g.\ $x^{4}$ at $x_{0}=0$).
\end{rulebox}

\begin{rulebox}{Second-derivative test}
If $f'(x_{0})=0$ and $f''(x_{0})>0$, $x_{0}$ is a local minimum;
$f''(x_{0})<0$ gives a local maximum.
\end{rulebox}

\begin{rulebox}{Leibniz rule (higher product)}
\[
  (fg)^{(n)}=\sum_{k=0}^{n}\binom{n}{k}f^{(k)}g^{(n-k)} .
\]
\end{rulebox}

\begin{rulebox}{Taylor polynomial}
The Taylor polynomial of order $n$ at $a$:
\[
  T_{n}(x)=\sum_{k=0}^{n}\frac{f^{(k)}(a)}{k!}(x-a)^{k},
\]
with remainder
\[
  R_{n}(x)=\frac{f^{(n+1)}(\xi)}{(n+1)!}(x-a)^{n+1}
\]
for some $\xi$ between $a$ and $x$.
\end{rulebox}

\section{Practice Problems}

\begin{problems}
  \item Compute $f''(x)$ for $f(x)=2x^{4}-5x^{3}+x-7$.
  \item Determine the intervals on which $f(x)=x^{3}-6x^{2}+5$ is convex and
        on which it is concave.
  \item Find a closed form for $f^{(n)}(x)$ when $f(x)=\mathrm{e}^{-x}$.
  \item Find all inflection points of $f(x)=x^{6}$.
  \item Compute $f''(x)$ for $f(x)=\mathrm{e}^{-2x}$.
  \item For a cost function $C(Q)=Q^{3}-12Q^{2}+50Q+100$, determine the
        output level where marginal cost transitions from decreasing to
        increasing.
  \item Find a closed form for $f^{(n)}(x)$ when $f(x)=\ln(1+x)$.
  \item Compute $f''(x)$ for $f(x)=x^{2}\mathrm{e}^{x}$.
  \item Find all inflection points of $f(x)=x^{3}-9x^{2}+24x-3$.
  \item Determine the convexity/concavity intervals of $f(x)=x\mathrm{e}^{-x}$.
  \item Compute $f''(x)$ for $f(x)=\ln(2x+5)$.
  \item Find a closed form for $f^{(n)}(x)$ when $f(x)=\dfrac{1}{x+a}$
        (constant $a$).
  \item Show that $f(x)=x^{4}$ has $f''(0)=0$ but $x=0$ is \emph{not} an
        inflection point.
  \item Compute $f''(x)$ for $f(x)=3x^{2/3}-x^{-2}$.
  \item Determine the convexity intervals of $f(x)=\ln(1+\mathrm{e}^{x})$
        (softplus).
  \item Find all inflection points of $f(x)=x^{4}-12x^{2}+5$.
  \item Compute $f''(x)$ for $f(x)=x^{2}\ln(x)$.
  \item Find a closed form for $f^{(n)}(x)$ when $f(x)=\cos(x)$.
  \item For utility $U(x)=\ln(x)$, $x>0$, verify that $U$ is strictly
        concave (the standard ``risk-aversion'' condition).
  \item Compute $f''(x)$ for $f(x)=\tan(x)$.
  \item Find all inflection points of $f(x)=\dfrac{x^{2}}{x^{2}+3}$.
  \item Determine convexity/concavity intervals of $f(x)=\mathrm{e}^{x}-x$.
  \item Compute $f''(x)$ for $f(x)=\sin(x)\cos(x)$.
  \item Find a closed form for $f^{(n)}(x)$ when $f(x)=\dfrac{1}{(1-x)^{2}}$.
  \item Find all inflection points of $f(x)=x^{2}\mathrm{e}^{-x}$.
  \item Compute $f''(x)$ for $f(x)=\dfrac{x}{x-3}$.
  \item Determine convexity/concavity of $f(x)=x^{4}-8x^{3}+1$ and the
        corresponding inflection points.
  \item Find a closed form for $f^{(n)}(x)$ when $f(x)=\mathrm{e}^{ax}$ for
        a constant $a$.
  \item Compute $f''(x)$ for $f(x)=(x^{2}+4)^{3/2}$.
  \item For $f(x)=x^{5}$, determine whether $x=0$ is an inflection point.
  \item Compute $f''(x)$ for $f(x)=\sqrt{1+x^{2}}$.
  \item Determine the inflection points of $f(x)=\arctan(x)$.
  \item Find $f^{(n)}(x)$ for $f(x)=\sin(2x)$.
  \item Compute $f''(x)$ for $f(x)=\dfrac{1}{1+x^{2}}$.
  \item Determine convexity/concavity of $f(x)=x\ln(x)$, $x>0$.
  \item For the logistic function $f(x)=\dfrac{1}{1+\mathrm{e}^{-x}}$, find
        the inflection point.
  \item Compute $f''(x)$ for $f(x)=x^{3}\ln(x)$.
  \item For a profit function $\pi(Q)=-Q^{3}+15Q^{2}-27Q-10$, where is
        $\pi$ concave?
  \item Find $f^{(n)}(0)$ for $f(x)=\ln(1+x)$ at $x=0$ (i.e.\ the
        coefficients in its Taylor series at zero).
  \item Compute $f''(x)$ for $f(x)=\mathrm{e}^{x^{2}}$.
  \item Determine all inflection points of $f(x)=x^{5}-5x^{4}$.
  \item For $f(x)=x|x|$ (i.e.\ $x^{2}$ for $x\geq 0$, $-x^{2}$ for $x<0$),
        determine whether $x=0$ is an inflection point and whether
        $f''(0)$ exists.
  \item Compute $f''(x)$ for $f(x)=\arctan(x)$.
  \item Apply the Leibniz rule to compute $(x^{2}\mathrm{e}^{x})^{(n)}$ in
        closed form.
  \item Find the inflection points of $f(x)=(x-1)^{4}$.
  \item For $f(x)=\sin(x)$, find the Taylor polynomial $T_{5}(x)$ at $a=0$.
  \item Compute $f''(x)$ for $f(x)=\ln(\cos(x))$ on
        $\bigl(-\tfrac{\pi}{2},\tfrac{\pi}{2}\bigr)$.
  \item Find a closed form for $f^{(n)}(x)$ when $f(x)=x\mathrm{e}^{x}$,
        using the Leibniz rule.
  \item Determine the convexity intervals of $f(x)=\dfrac{1}{1+x^{2}}$ (the
        standard Cauchy-type bell curve).
  \item For a Cobb--Douglas production function $f(L)=AL^{\alpha}$ with
        $0<\alpha<1$ and $A>0$, verify that $f$ is strictly concave in
        $L>0$ (the law of diminishing returns).
\end{problems}

\section{Solutions}

\begin{solutions}
  \item $f'(x)=8x^{3}-15x^{2}+1$, $f''(x)=24x^{2}-30x$.

  \item $f'(x)=3x^{2}-12x$, $f''(x)=6x-12=6(x-2)$. So $f''(x)\geq 0$ iff
        $x\geq 2$: convex on $[2,\infty)$, concave on $(-\infty,2]$.

  \item Each differentiation introduces a factor of $-1$:
        $f^{(n)}(x)=(-1)^{n}\mathrm{e}^{-x}$.

  \item $f'(x)=6x^{5}$, $f''(x)=30x^{4}\geq 0$ for all $x$, with equality
        only at $x=0$. Since $f''$ does not change sign, $x=0$ is
        \emph{not} an inflection point. $f$ has no inflection points.

  \item $f'(x)=-2\mathrm{e}^{-2x}$, $f''(x)=4\mathrm{e}^{-2x}$.

  \item $MC(Q)=C'(Q)=3Q^{2}-24Q+50$. $MC'(Q)=6Q-24=0\Rightarrow Q=4$.
        $MC''(Q)=6>0$ so $Q=4$ minimises $MC$. At $Q=4$, $C''(Q)$ changes
        sign there: $C''(Q)=6Q-24$, negative for $Q<4$, positive for
        $Q>4$, so $Q=4$ is the inflection point of $C$ and the transition
        output for marginal cost.

  \item With $u=1+x$: $f'(x)=u^{-1}$, $f''(x)=-u^{-2}$,
        $f'''(x)=2u^{-3}$, in general
        \[
          f^{(n)}(x)=\frac{(-1)^{n-1}(n-1)!}{(1+x)^{n}}\qquad (n\geq 1) .
        \]

  \item $f'(x)=2x\mathrm{e}^{x}+x^{2}\mathrm{e}^{x}=(x^{2}+2x)\mathrm{e}^{x}$,
        then
        \[
          f''(x)=(2x+2)\mathrm{e}^{x}+(x^{2}+2x)\mathrm{e}^{x}
                =(x^{2}+4x+2)\mathrm{e}^{x} .
        \]

  \item $f'(x)=3x^{2}-18x+24$, $f''(x)=6x-18=6(x-3)$. Sign change at
        $x=3$. $f(3)=27-81+72-3=15$. Inflection at $(3,15)$.

  \item $f'(x)=\mathrm{e}^{-x}-x\mathrm{e}^{-x}=(1-x)\mathrm{e}^{-x}$.
        $f''(x)=-\mathrm{e}^{-x}-(1-x)\mathrm{e}^{-x}=(x-2)\mathrm{e}^{-x}$.
        Since $\mathrm{e}^{-x}>0$, the sign of $f''$ is that of $x-2$.
        Concave on $(-\infty,2]$, convex on $[2,\infty)$. Inflection at
        $x=2$.

  \item $f'(x)=\dfrac{2}{2x+5}$, $f''(x)=-\dfrac{4}{(2x+5)^{2}}$.

  \item $f(x)=(x+a)^{-1}$, so
        $f^{(n)}(x)=(-1)^{n}n!\,(x+a)^{-(n+1)}=\dfrac{(-1)^{n}n!}{(x+a)^{n+1}}$.

  \item $f'(x)=4x^{3}$, $f''(x)=12x^{2}$. Then $f''(0)=0$. But
        $f''(x)=12x^{2}\geq 0$ everywhere, no sign change. Hence $x=0$
        is not an inflection point. Indeed $f$ is convex on all of $\R$.

  \item $f'(x)=2x^{-1/3}+2x^{-3}$,
        $f''(x)=-\tfrac{2}{3}x^{-4/3}-6x^{-4}$.

  \item $f'(x)=\dfrac{\mathrm{e}^{x}}{1+\mathrm{e}^{x}}$. By the quotient
        rule
        \[
          f''(x)=\frac{\mathrm{e}^{x}(1+\mathrm{e}^{x})-\mathrm{e}^{x}\cdot\mathrm{e}^{x}}
                       {(1+\mathrm{e}^{x})^{2}}
                =\frac{\mathrm{e}^{x}}{(1+\mathrm{e}^{x})^{2}}>0
          \text{ for all } x .
        \]
        $f$ is strictly convex on all of $\R$.

  \item $f'(x)=4x^{3}-24x$, $f''(x)=12x^{2}-24=12(x^{2}-2)$. Sign changes
        at $x=\pm\sqrt{2}$. $f(\pm\sqrt{2})=4-24+5=-15$. Inflections at
        $(\pm\sqrt{2},-15)$.

  \item $f'(x)=2x\ln(x)+x$. Then $f''(x)=2\ln(x)+2+1=2\ln(x)+3$.

  \item $f^{(n)}(x)=\cos\bigl(x+\tfrac{n\pi}{2}\bigr)$.

  \item $U'(x)=1/x$, $U''(x)=-1/x^{2}<0$ for all $x>0$. So $U$ is strictly
        concave on $(0,\infty)$.

  \item $f'(x)=\sec^{2}(x)$. Then
        $f''(x)=2\sec(x)\cdot\sec(x)\tan(x)=2\sec^{2}(x)\tan(x)$.

  \item Write $f(x)=1-\dfrac{3}{x^{2}+3}$. Then
        $f'(x)=\dfrac{6x}{(x^{2}+3)^{2}}$ and
        \[
          f''(x)=\frac{6(x^{2}+3)^{2}-6x\cdot 2(x^{2}+3)\cdot 2x}
                       {(x^{2}+3)^{4}}
                =\frac{6(x^{2}+3)-24x^{2}}{(x^{2}+3)^{3}}
                =\frac{18-18x^{2}}{(x^{2}+3)^{3}} .
        \]
        Sign changes at $x=\pm 1$. $f(\pm 1)=1/4$. Inflections at
        $(\pm 1,\tfrac14)$.

  \item $f'(x)=\mathrm{e}^{x}-1$, $f''(x)=\mathrm{e}^{x}>0$ for all $x$.
        Strictly convex on $\R$. No concave region.

  \item $f(x)=\tfrac12\sin(2x)$, $f'(x)=\cos(2x)$, $f''(x)=-2\sin(2x)$.

  \item Note $\dfrac{1}{(1-x)^{2}}=\dfrac{\mathrm{d}}{\mathrm{d}x}
        \dfrac{1}{1-x}$. Since
        $\dfrac{\mathrm{d}^{k}}{\mathrm{d}x^{k}}\dfrac{1}{1-x}
        =\dfrac{k!}{(1-x)^{k+1}}$, we get
        \[
          f^{(n)}(x)=\frac{(n+1)!}{(1-x)^{n+2}} .
        \]

  \item $f'(x)=2x\mathrm{e}^{-x}-x^{2}\mathrm{e}^{-x}=(2x-x^{2})\mathrm{e}^{-x}$.
        $f''(x)=(2-2x)\mathrm{e}^{-x}-(2x-x^{2})\mathrm{e}^{-x}
               =(x^{2}-4x+2)\mathrm{e}^{-x}$. Roots: $x=2\pm\sqrt{2}$.
        Inflections at $x=2\pm\sqrt{2}$; $y$-values
        $(2\pm\sqrt{2})^{2}\mathrm{e}^{-(2\pm\sqrt{2})}
        =(6\pm4\sqrt{2})\mathrm{e}^{-(2\pm\sqrt{2})}$.

  \item Write $f(x)=1+\dfrac{3}{x-3}$. Then $f'(x)=-\dfrac{3}{(x-3)^{2}}$,
        $f''(x)=\dfrac{6}{(x-3)^{3}}$.

  \item $f'(x)=4x^{3}-24x^{2}$, $f''(x)=12x^{2}-48x=12x(x-4)$. Sign
        changes at $x=0$ and $x=4$. Convex on $(-\infty,0]\cup[4,\infty)$,
        concave on $[0,4]$. Inflections at $(0,1)$ and
        $(4,4^{4}-8\cdot 64+1)=(4,256-512+1)=(4,-255)$.

  \item Each derivative pulls out a factor $a$: $f^{(n)}(x)=a^{n}\mathrm{e}^{ax}$.

  \item $f'(x)=\tfrac32(x^{2}+4)^{1/2}\cdot 2x=3x\sqrt{x^{2}+4}$. Then
        \[
          f''(x)=3\sqrt{x^{2}+4}+3x\cdot\frac{x}{\sqrt{x^{2}+4}}
                =\frac{3(x^{2}+4)+3x^{2}}{\sqrt{x^{2}+4}}
                =\frac{6x^{2}+12}{\sqrt{x^{2}+4}}
                =\frac{6(x^{2}+2)}{\sqrt{x^{2}+4}} .
        \]

  \item $f'(x)=5x^{4}$, $f''(x)=20x^{3}$. $f''(x)<0$ for $x<0$,
        $f''(x)>0$ for $x>0$: sign change at $x=0$. Hence $x=0$ is an
        inflection point (with $f(0)=0$).

  \item $f'(x)=\dfrac{x}{\sqrt{1+x^{2}}}$. Then
        \[
          f''(x)=\frac{\sqrt{1+x^{2}}-x\cdot\frac{x}{\sqrt{1+x^{2}}}}
                       {1+x^{2}}
                =\frac{(1+x^{2})-x^{2}}{(1+x^{2})^{3/2}}
                =\frac{1}{(1+x^{2})^{3/2}} .
        \]

  \item $f'(x)=\dfrac{1}{1+x^{2}}$. Then
        $f''(x)=-\dfrac{2x}{(1+x^{2})^{2}}$. Sign change at $x=0$, with
        $f''(x)>0$ for $x<0$ and $f''(x)<0$ for $x>0$. Inflection at
        $(0,0)$.

  \item Write $\sin(2x)$. Each differentiation multiplies the
        argument-shift trick by $2$:
        \[
          f^{(n)}(x)=2^{n}\sin\Bigl(2x+\frac{n\pi}{2}\Bigr) .
        \]

  \item $f(x)=(1+x^{2})^{-1}$. $f'(x)=-2x(1+x^{2})^{-2}$. Then
        \[
          f''(x)=-2(1+x^{2})^{-2}+(-2x)\cdot(-2)(1+x^{2})^{-3}\cdot 2x
                =-\frac{2}{(1+x^{2})^{2}}+\frac{8x^{2}}{(1+x^{2})^{3}}
                =\frac{-2(1+x^{2})+8x^{2}}{(1+x^{2})^{3}}
                =\frac{6x^{2}-2}{(1+x^{2})^{3}} .
        \]

  \item $f'(x)=\ln(x)+1$, $f''(x)=1/x>0$ for all $x>0$. Strictly convex
        on $(0,\infty)$.

  \item $f(x)=\dfrac{1}{1+\mathrm{e}^{-x}}=\dfrac{\mathrm{e}^{x}}{1+\mathrm{e}^{x}}$.
        Then $f'(x)=\dfrac{\mathrm{e}^{x}}{(1+\mathrm{e}^{x})^{2}}$.
        Differentiate again using the quotient rule (with $u=\mathrm{e}^{x}$,
        $v=(1+\mathrm{e}^{x})^{2}$, $v'=2(1+\mathrm{e}^{x})\mathrm{e}^{x}$):
        \[
          f''(x)=\frac{\mathrm{e}^{x}(1+\mathrm{e}^{x})^{2}
                       -\mathrm{e}^{x}\cdot 2(1+\mathrm{e}^{x})\mathrm{e}^{x}}
                       {(1+\mathrm{e}^{x})^{4}}
                =\frac{\mathrm{e}^{x}(1+\mathrm{e}^{x})-2\mathrm{e}^{2x}}
                       {(1+\mathrm{e}^{x})^{3}}
                =\frac{\mathrm{e}^{x}(1-\mathrm{e}^{x})}{(1+\mathrm{e}^{x})^{3}} .
        \]
        Sign of $f''$ is the sign of $1-\mathrm{e}^{x}$: positive for
        $x<0$, negative for $x>0$. Inflection at $x=0$, where
        $f(0)=1/2$. Inflection point $(0,\tfrac12)$.

  \item $f'(x)=3x^{2}\ln(x)+x^{2}=x^{2}(3\ln(x)+1)$. Then
        \[
          f''(x)=2x(3\ln(x)+1)+x^{2}\cdot\frac{3}{x}
                =2x(3\ln(x)+1)+3x
                =6x\ln(x)+5x=x(6\ln(x)+5) .
        \]

  \item $\pi'(Q)=-3Q^{2}+30Q-27$, $\pi''(Q)=-6Q+30=6(5-Q)$. $\pi$ is
        concave on $\{Q:\pi''(Q)\leq 0\}=[5,\infty)$ and convex on
        $(-\infty,5]$. (Economic note: the profit function changes from
        convex to concave at the inflection $Q=5$.)

  \item $f^{(n)}(x)=\dfrac{(-1)^{n-1}(n-1)!}{(1+x)^{n}}$ for $n\geq 1$. At
        $x=0$: $f^{(n)}(0)=(-1)^{n-1}(n-1)!$ for $n\geq 1$, and $f(0)=0$.
        So the Taylor series is
        $\displaystyle\sum_{n\geq 1}\frac{(-1)^{n-1}}{n}x^{n}
        =x-\frac{x^{2}}{2}+\frac{x^{3}}{3}-\cdots$.

  \item $f'(x)=2x\mathrm{e}^{x^{2}}$. Then
        $f''(x)=2\mathrm{e}^{x^{2}}+2x\cdot 2x\mathrm{e}^{x^{2}}
               =(2+4x^{2})\mathrm{e}^{x^{2}}$.

  \item $f'(x)=5x^{4}-20x^{3}$, $f''(x)=20x^{3}-60x^{2}=20x^{2}(x-3)$.
        Now $20x^{2}\geq 0$, so the sign of $f''$ is determined by $x-3$
        except at $x=0$ where $f''=0$. Check signs: for $x<0$, $x^{2}>0$
        and $x-3<0$, so $f''<0$. For $0<x<3$, $x^{2}>0$ and $x-3<0$, so
        $f''<0$. No sign change at $x=0$. For $x>3$, $f''>0$. Sign change
        only at $x=3$. $f(3)=243-405=-162$. Inflection at $(3,-162)$.
        ($x=0$ is \emph{not} an inflection.)

  \item For $x>0$: $f(x)=x^{2}$, $f''(x)=2$. For $x<0$: $f(x)=-x^{2}$,
        $f''(x)=-2$. So $f''$ has different one-sided limits at $0$ ($2$
        from the right, $-2$ from the left). Hence $f''(0)$ does not
        exist. But $f'(x)=2|x|$ exists and is continuous everywhere, and
        $f''$ changes sign at $0$ (in the sense of one-sided values). The
        graph changes from concave (left) to convex (right) at $x=0$, so
        $x=0$ is an inflection point in the geometric sense even though
        $f''(0)$ itself is undefined.

  \item $f'(x)=\dfrac{1}{1+x^{2}}$. Then
        $f''(x)=-\dfrac{2x}{(1+x^{2})^{2}}$.

  \item Take $u=x^{2}$ and $v=\mathrm{e}^{x}$. Then $u'=2x$, $u''=2$,
        $u^{(k)}=0$ for $k\geq 3$, and $v^{(k)}=\mathrm{e}^{x}$. By
        Leibniz:
        \[
          (x^{2}\mathrm{e}^{x})^{(n)}
            =\binom{n}{0}x^{2}\mathrm{e}^{x}
             +\binom{n}{1}\cdot 2x\cdot\mathrm{e}^{x}
             +\binom{n}{2}\cdot 2\cdot\mathrm{e}^{x}
            =\mathrm{e}^{x}\bigl[x^{2}+2nx+n(n-1)\bigr] .
        \]

  \item $f'(x)=4(x-1)^{3}$, $f''(x)=12(x-1)^{2}\geq 0$ everywhere, with
        equality only at $x=1$. No sign change. Hence $x=1$ is \emph{not}
        an inflection point; $f$ is convex on all of $\R$.

  \item Using $f^{(k)}(0)=\sin(k\pi/2)$: values are $0,1,0,-1,0,1$ for
        $k=0,1,2,3,4,5$. Hence
        \[
          T_{5}(x)=x-\frac{x^{3}}{3!}+\frac{x^{5}}{5!}
                  =x-\frac{x^{3}}{6}+\frac{x^{5}}{120} .
        \]

  \item $f'(x)=-\tan(x)$. Then $f''(x)=-\sec^{2}(x)<0$ on
        $\bigl(-\tfrac{\pi}{2},\tfrac{\pi}{2}\bigr)$. So $f$ is strictly
        concave on that interval.

  \item Take $u=x$, $v=\mathrm{e}^{x}$. Then $u'=1$, $u^{(k)}=0$ for
        $k\geq 2$, and $v^{(k)}=\mathrm{e}^{x}$. By Leibniz:
        \[
          (x\mathrm{e}^{x})^{(n)}
            =\binom{n}{0}x\mathrm{e}^{x}+\binom{n}{1}\cdot 1\cdot\mathrm{e}^{x}
            =(x+n)\mathrm{e}^{x} .
        \]

  \item From problem 34: $f''(x)=\dfrac{6x^{2}-2}{(1+x^{2})^{3}}$. Sign of
        $f''$ equals sign of $6x^{2}-2$, i.e.\ positive iff
        $|x|>1/\sqrt{3}$. So $f$ is convex on
        $(-\infty,-1/\sqrt{3}]\cup[1/\sqrt{3},\infty)$ and concave on
        $[-1/\sqrt{3},1/\sqrt{3}]$. Inflection points at $x=\pm 1/\sqrt{3}$.

  \item $f'(L)=A\alpha L^{\alpha-1}$, $f''(L)=A\alpha(\alpha-1)L^{\alpha-2}$.
        With $A>0$, $0<\alpha<1$: $\alpha(\alpha-1)<0$ and
        $L^{\alpha-2}>0$, so $f''(L)<0$ for all $L>0$. Strictly concave on
        $(0,\infty)$, diminishing marginal product of labour.
\end{solutions}

\section{Additional Problems}

\begin{problems}
  \item Compute $f''(x)$ for $f(x)=3x^{5}-2x^{4}+x^{2}-7$.
  \item Determine the convexity/concavity intervals of
        $f(x)=x^{3}-3x^{2}-9x+5$.
  \item Find a closed form for $f^{(n)}(x)$ when $f(x)=\mathrm{e}^{3x}$.
  \item Find all inflection points of $f(x)=x^{4}-6x^{2}+8$.
  \item Compute $f''(x)$ for $f(x)=x\mathrm{e}^{-2x}$.
  \item For a cost function $C(Q)=Q^{3}-15Q^{2}+72Q+50$, determine the
        output level where marginal cost transitions from decreasing to
        increasing.
  \item Find a closed form for $f^{(n)}(x)$ when $f(x)=\ln(2-x)$, on
        $x<2$.
  \item Compute $f''(x)$ for $f(x)=\sqrt{x^{2}+4}$.
  \item Find all inflection points of $f(x)=x^{3}-12x^{2}+45x-7$.
  \item Determine the convexity intervals of $f(x)=\ln(x^{2}+1)$.
  \item Compute $f''(x)$ for $f(x)=\sin(x)\,\mathrm{e}^{x}$.
  \item Find a closed form for $f^{(n)}(x)$ when
        $f(x)=\dfrac{1}{(1-2x)}$.
  \item For utility $U(x)=\sqrt{x}$, $x>0$, verify that $U$ is strictly
        concave (diminishing marginal utility).
  \item Compute $f''(x)$ for $f(x)=x^{3}\mathrm{e}^{-x}$.
  \item Apply the Leibniz rule to compute $(x^{3}\mathrm{e}^{x})^{(n)}$ in
        closed form.
\end{problems}

\section{Additional Solutions}

\begin{solutions}
  \item $f'(x)=15x^{4}-8x^{3}+2x$, $f''(x)=60x^{3}-24x^{2}+2$.

  \item $f'(x)=3x^{2}-6x-9$, $f''(x)=6x-6=6(x-1)$. So $f''\geq 0$ iff
        $x\geq 1$: convex on $[1,\infty)$, concave on $(-\infty,1]$.
        Inflection at $x=1$, $f(1)=-6$.

  \item Each differentiation introduces a factor of $3$:
        $f^{(n)}(x)=3^{n}\mathrm{e}^{3x}$.

  \item $f'(x)=4x^{3}-12x$, $f''(x)=12x^{2}-12=12(x^{2}-1)$. Sign changes
        at $x=\pm 1$. $f(\pm 1)=1-6+8=3$. Inflections at $(\pm 1,3)$.

  \item Product rule:
        $f'(x)=\mathrm{e}^{-2x}-2x\mathrm{e}^{-2x}=(1-2x)\mathrm{e}^{-2x}$.
        Then
        \[
          f''(x)=-2\mathrm{e}^{-2x}-2(1-2x)\mathrm{e}^{-2x}
                =(4x-4)\mathrm{e}^{-2x}=4(x-1)\mathrm{e}^{-2x} .
        \]

  \item $MC(Q)=C'(Q)=3Q^{2}-30Q+72$. $MC'(Q)=6Q-30=0\Rightarrow Q=5$.
        $MC''(Q)=6>0$, so $Q=5$ minimises $MC$. Equivalently,
        $C''(Q)=6Q-30$ changes sign at $Q=5$: the cost function
        transitions from concave to convex at the inflection $Q=5$.

  \item With $u=2-x$: $f'(x)=-1/(2-x)$, $f''(x)=-1/(2-x)^{2}$,
        $f'''(x)=-2/(2-x)^{3}$. By induction:
        \[
          f^{(n)}(x)=-\frac{(n-1)!}{(2-x)^{n}}\qquad (n\geq 1) .
        \]

  \item $f(x)=(x^{2}+4)^{1/2}$. $f'(x)=\dfrac{x}{\sqrt{x^{2}+4}}$. Then
        \[
          f''(x)=\frac{\sqrt{x^{2}+4}-x\cdot\frac{x}{\sqrt{x^{2}+4}}}
                       {x^{2}+4}
                =\frac{(x^{2}+4)-x^{2}}{(x^{2}+4)^{3/2}}
                =\frac{4}{(x^{2}+4)^{3/2}} .
        \]

  \item $f'(x)=3x^{2}-24x+45$, $f''(x)=6x-24=6(x-4)$. Sign change at
        $x=4$. $f(4)=64-192+180-7=45$. Inflection at $(4,45)$.

  \item $f'(x)=\dfrac{2x}{x^{2}+1}$. By the quotient rule
        \[
          f''(x)=\frac{2(x^{2}+1)-2x\cdot 2x}{(x^{2}+1)^{2}}
                =\frac{2-2x^{2}}{(x^{2}+1)^{2}}
                =\frac{2(1-x^{2})}{(x^{2}+1)^{2}} .
        \]
        Sign of $f''$ matches sign of $1-x^{2}$. Convex on $[-1,1]$,
        concave on $(-\infty,-1]\cup[1,\infty)$. Inflections at
        $x=\pm 1$.

  \item $f'(x)=\cos(x)\mathrm{e}^{x}+\sin(x)\mathrm{e}^{x}
        =(\cos x+\sin x)\mathrm{e}^{x}$. Then
        \[
          f''(x)=(-\sin x+\cos x)\mathrm{e}^{x}
                 +(\cos x+\sin x)\mathrm{e}^{x}
                =2\cos(x)\,\mathrm{e}^{x} .
        \]

  \item $f(x)=(1-2x)^{-1}$. By induction:
        $f^{(n)}(x)=\dfrac{2^{n}\,n!}{(1-2x)^{n+1}}$.

  \item $U'(x)=\tfrac12 x^{-1/2}$, $U''(x)=-\tfrac14 x^{-3/2}<0$ for all
        $x>0$. So $U$ is strictly concave on $(0,\infty)$.

  \item $f'(x)=3x^{2}\mathrm{e}^{-x}-x^{3}\mathrm{e}^{-x}
        =(3x^{2}-x^{3})\mathrm{e}^{-x}$. Then
        \[
          f''(x)=(6x-3x^{2})\mathrm{e}^{-x}-(3x^{2}-x^{3})\mathrm{e}^{-x}
                =(x^{3}-6x^{2}+6x)\mathrm{e}^{-x}
                =x(x^{2}-6x+6)\mathrm{e}^{-x} .
        \]

  \item Take $u=x^{3}$ and $v=\mathrm{e}^{x}$. Then $u'=3x^{2}$,
        $u''=6x$, $u'''=6$, $u^{(k)}=0$ for $k\geq 4$, and
        $v^{(k)}=\mathrm{e}^{x}$. By Leibniz:
        \[
          (x^{3}\mathrm{e}^{x})^{(n)}
            =\binom{n}{0}x^{3}\mathrm{e}^{x}
             +\binom{n}{1}\cdot 3x^{2}\cdot\mathrm{e}^{x}
             +\binom{n}{2}\cdot 6x\cdot\mathrm{e}^{x}
             +\binom{n}{3}\cdot 6\cdot\mathrm{e}^{x}
            =\mathrm{e}^{x}\bigl[x^{3}+3nx^{2}+3n(n-1)x+n(n-1)(n-2)\bigr] .
        \]
\end{solutions}

\chapter{Taylor Approximation}

The previous chapter closed with a formula for the Taylor polynomial;
this chapter turns it into a working tool. Most functions that appear in
economics (utility, production, growth) are not polynomials, yet
polynomials are the easiest objects to evaluate and to reason about.
Taylor's theorem says that near a chosen base point every sufficiently
smooth function \emph{is} a polynomial, up to a controllable error. This
single idea underlies the linearisation of models around a steady state,
the log-linearisation trick $\ln(1+g)\approx g$ for growth rates, the
second-order conditions of optimisation, and quick numerical estimates
such as $\sqrt{4.1}$ or $\mathrm{e}^{0.1}$ without a calculator.

\section{Core Rules}

\begin{rulebox}{Taylor polynomial of order $n$}
Let $f$ be $n$ times differentiable at $a$. The \emph{Taylor polynomial}
of order $n$ of $f$ at $a$ is
\[
  T_{n}(x)=\sum_{k=0}^{n}\frac{f^{(k)}(a)}{k!}(x-a)^{k}
          =f(a)+f'(a)(x-a)+\frac{f''(a)}{2!}(x-a)^{2}+\cdots
           +\frac{f^{(n)}(a)}{n!}(x-a)^{n}.
\]
It is the \emph{unique} polynomial of degree $\leq n$ whose value and
first $n$ derivatives at $a$ all agree with those of $f$:
$T_{n}^{(k)}(a)=f^{(k)}(a)$ for $k=0,1,\ldots,n$.
\end{rulebox}

\begin{rulebox}{Linear and quadratic approximation ($n=1,2$)}
\[
  T_{1}(x)=f(a)+f'(a)(x-a),
  \qquad
  T_{2}(x)=f(a)+f'(a)(x-a)+\tfrac12 f''(a)(x-a)^{2}.
\]
$T_{1}$ is the tangent line at $a$ (the \emph{linearisation} of $f$);
$T_{2}$ is the best-fitting parabola, matching value, slope and
curvature at $a$. Each extra order typically buys one further correct
decimal place near $a$.
\end{rulebox}

\begin{rulebox}{Lagrange remainder and error bound}
If $f$ is $(n+1)$ times differentiable on an interval containing $a$ and
$x$, then
\[
  f(x)=T_{n}(x)+R_{n}(x),
  \qquad
  R_{n}(x)=\frac{f^{(n+1)}(\xi)}{(n+1)!}(x-a)^{n+1}
\]
for some $\xi$ strictly between $a$ and $x$. Since $\xi$ is unknown,
bound it: if $|f^{(n+1)}(t)|\leq M$ for all $t$ between $a$ and $x$,
then
\[
  \bigl|f(x)-T_{n}(x)\bigr|\leq\frac{M}{(n+1)!}\,|x-a|^{n+1}.
\]
The error shrinks like $|x-a|^{n+1}$: halving the distance to $a$
divides the error of $T_{n}$ by about $2^{n+1}$.
\end{rulebox}

\begin{rulebox}{Maclaurin case ($a=0$): standard expansions}
The Taylor polynomial at $a=0$ is called the \emph{Maclaurin}
polynomial. The classical expansions are
\[
\begin{aligned}
  \mathrm{e}^{x}&=1+x+\frac{x^{2}}{2!}+\frac{x^{3}}{3!}+\cdots
                  +\frac{x^{n}}{n!}+\cdots
                  &&\text{(all }x\text{)},\\
  \ln(1+x)&=x-\frac{x^{2}}{2}+\frac{x^{3}}{3}-\frac{x^{4}}{4}+\cdots
                  &&(-1<x\leq 1),\\
  (1+x)^{k}&=1+kx+\frac{k(k-1)}{2!}x^{2}
                  +\frac{k(k-1)(k-2)}{3!}x^{3}+\cdots
                  &&(|x|<1),\\
  \sin(x)&=x-\frac{x^{3}}{3!}+\frac{x^{5}}{5!}-\cdots
                  &&\text{(all }x\text{)},\\
  \cos(x)&=1-\frac{x^{2}}{2!}+\frac{x^{4}}{4!}-\cdots
                  &&\text{(all }x\text{)}.
\end{aligned}
\]
The first-order case $\ln(1+x)\approx x$ is the workhorse of
log-linearisation in economics: a small growth rate $g$ and the log
change $\ln(1+g)$ are interchangeable to first order.
\end{rulebox}

\begin{rulebox}{Quadratic approximation in economics}
Near a stationary point $a$ (where $f'(a)=0$) the linear term vanishes
and
\[
  f(x)\approx f(a)+\tfrac12 f''(a)(x-a)^{2}:
\]
locally, $f$ \emph{is} a parabola sitting on top of $f(a)$. If
$f''(a)<0$ the parabola opens downward and $a$ is a local maximum; if
$f''(a)>0$ it opens upward and $a$ is a local minimum. This is exactly
where the second-order condition comes from. In several variables the
role of $f''(a)$ is played by the Hessian, and the sign condition
becomes definiteness.
\end{rulebox}

\begin{rulebox}{Caution: validity is local}
A Taylor polynomial is a \emph{local} approximation: the error grows
like $|x-a|^{n+1}$, so accuracy is excellent near $a$ and can be
useless far away. For $\ln(1+x)$ and $(1+x)^{k}$ the expansions are
only valid on $(-1,1]$ resp.\ $(-1,1)$: outside, adding more terms
makes things \emph{worse}. Even inside: $T_{1}(x)=x$ for $\ln(1+x)$
gives $\ln(2)\approx 1$ against the true $0.693$, a $44\%$ error,
versus $0.5\%$ at $x=0.1$. Rule of thumb: expand around a base point
close to where you evaluate, and use the remainder bound to check.
\end{rulebox}

\section{Worked Example}

\begin{worked}{Expand $\ln(1+x)$ to third order about $0$, bound the error at $x=0.5$, and read the economics off the second-order term}
\wstep{Step 1. Collect the derivatives at the expansion point.}
\[
  f(x)=\ln(1+x),\quad
  f'=\frac{1}{1+x},\quad
  f''=\frac{-1}{(1+x)^{2}},\quad
  f'''=\frac{2}{(1+x)^{3}},\quad
  f^{(4)}=\frac{-6}{(1+x)^{4}} .
\]
At $a=0$: $f(0)=0$, $f'(0)=1$, $f''(0)=-1$, $f'''(0)=2$, and
$f^{(4)}(0)=-6$.
\wtrap{Evaluate \emph{after} differentiating, never before. Substituting
$x=0$ into $f'$ and then differentiating the resulting constant gives zeros
all the way down --- a surprisingly frequent way to produce the polynomial
$P_{3}(x)=x$.}

\wstep{Step 2. Assemble the Taylor polynomial.}
\[
  P_{n}(x)=\sum_{k=0}^{n}\frac{f^{(k)}(a)}{k!}(x-a)^{k}
  \quad\Longrightarrow\quad
  P_{3}(x)=0+x-\frac{1}{2!}x^{2}+\frac{2}{3!}x^{3}
  =\boxed{x-\frac{x^{2}}{2}+\frac{x^{3}}{3}}
\]
The factorials are what turn the raw derivatives $1,-1,2$ into the
coefficients $1,-\tfrac12,\tfrac13$.
\wtrap{Do not omit $k!$, and do not omit $(x-a)^{k}$ when $a\neq0$. Expanding
about $a=2$ and writing powers of $x$ rather than of $x-2$ is the standard
error, and it silently produces a polynomial that is wrong everywhere.}

\wstep{Step 3. Evaluate at $x=0.5$ and compare.}
\[
  P_{3}(0.5)=0.5-0.125+0.041\overline{6}=0.4166\overline{6},
  \qquad
  \ln(1.5)=0.405465 ,
\]
so the actual error is $0.011202$.

\wstep{Step 4. Bound the error \emph{before} knowing the true value.}
The Lagrange remainder is
\[
  R_{3}(x)=\frac{f^{(4)}(\xi)}{4!}x^{4}
  \qquad\text{for some }\xi\text{ between }0\text{ and }x .
\]
$\xi$ is unknown, so take the worst case. On $[0,0.5]$,
\[
  \bigl|f^{(4)}(\xi)\bigr|=\frac{6}{(1+\xi)^{4}}\le\frac{6}{1^{4}}=6
  \qquad(\text{largest at }\xi=0),
\]
\[
  |R_{3}(0.5)|\le\frac{6}{24}(0.5)^{4}=0.25\cdot0.0625=0.015625 .
\]
And indeed $0.011201\le0.015625$. \checkmark
\wtrap{Maximise $|f^{(n+1)}|$ over the whole interval between $a$ and $x$ ---
do not substitute $\xi=x$ or $\xi=a$ and call it a bound. Here $1/(1+\xi)^{4}$
is decreasing, so the worst case is at the \emph{left} end; for an
increasing derivative it would be at the right. Check which end before
committing.}

\wstep{Step 5. Note how fast the accuracy degrades with distance.}
At $x=0.1$ the same bound gives $\tfrac14(0.1)^{4}=0.000025$; at $x=0.5$ it
is $625$ times larger. The error carries the factor $x^{4}$, so halving the
distance to the expansion point cuts the error sixteenfold. At $x=1$ the
series still converges, but at $x=1.1$ it diverges outright --- the
expansion of $\ln(1+x)$ is valid only for $|x|<1$.
\wtrap{A Taylor polynomial is a \textbf{local} statement. It is excellent
near $a$, useless far from it, and there is no warning in the formula --- 
$P_{3}(5)=5-12.5+41.7=34.2$, while $\ln6=1.79$. Always ask how far from $a$
you are being asked to work.}

\wstep{Step 6. The first two terms are the economics.}
Keeping $n=1$ gives the \textbf{linear approximation} $\ln(1+x)\approx x$:
the familiar ``a log difference is a percentage change'', accurate to about
half a percentage point for $|x|\le0.1$. Keeping $n=2$ gives
\[
  f(x)\approx f(a)+f'(a)(x-a)+\tfrac12 f''(a)(x-a)^{2},
\]
whose quadratic term carries the \textbf{curvature}. In expected-utility
theory that term is exactly what produces the risk premium: a concave $u$
has $u''<0$, so the quadratic term is negative, and its size relative to
$u'$ is the Arrow--Pratt coefficient $-u''/u'$. Risk aversion \emph{is} the
second-order term of a Taylor expansion.

\wstep{Step 7. The expansions worth knowing by heart.}
\[
  \mathrm{e}^{x}=1+x+\tfrac{x^{2}}{2}+\tfrac{x^{3}}{6}+\cdots,
  \qquad
  \ln(1+x)=x-\tfrac{x^{2}}{2}+\tfrac{x^{3}}{3}-\cdots,
\]
\[
  (1+x)^{\alpha}=1+\alpha x+\tfrac{\alpha(\alpha-1)}{2}x^{2}+\cdots,
  \qquad
  \frac{1}{1-x}=1+x+x^{2}+\cdots
\]
The last is the geometric series of Volume~4, which is why a perpetuity
formula and a Taylor expansion are the same computation twice.
\end{worked}

\section{Practice Problems}

\begin{problems}
  \item Find $T_{1}$ of $f(x)=\ln(x)$ at $a=1$.
  \item Find $T_{1}$ of $f(x)=\sqrt{x}$ at $a=4$.
  \item Find $T_{1}$ of $f(x)=\dfrac{1}{x}$ at $a=1$.
  \item Find $T_{1}$ of $f(x)=x^{3}-2x$ at $a=2$.
  \item Find $T_{1}$ of $f(x)=\mathrm{e}^{2x}$ at $a=\ln 2$.
  \item Find $T_{2}$ of $f(x)=\mathrm{e}^{x}$ at $a=0$ and use it to
        approximate $\mathrm{e}^{0.1}$. Compare with the true value
        $\mathrm{e}^{0.1}=1.105171$.
  \item Find $T_{2}$ of $f(x)=\ln(x)$ at $a=1$ and use it to
        approximate $\ln(1.05)$. Compare with the true value
        $\ln(1.05)=0.048790$.
  \item Find $T_{2}$ of $f(x)=\sqrt{x}$ at $a=4$ and use it to
        approximate $\sqrt{4.1}$. Compare with the true value
        $\sqrt{4.1}=2.024846$.
  \item Find $T_{2}$ of $f(x)=\dfrac{1}{1-x}$ at $a=0$.
  \item Find $T_{2}$ of $f(x)=x\ln(x)$ at $a=1$.
  \item Find $T_{2}$ of $f(x)=\mathrm{e}^{-x^{2}}$ at $a=0$.
  \item Find $T_{3}$ of $f(x)=\ln(1+x)$ at $a=0$ and use it to
        approximate $\ln(0.9)$. Compare with the true value
        $\ln(0.9)=-0.105361$.
  \item Use the binomial expansion to find $T_{2}$ of
        $f(x)=(1+x)^{1/2}$ at $a=0$, and use it to approximate
        $\sqrt{1.1}$. Compare with the true value
        $\sqrt{1.1}=1.048809$.
  \item Find the Maclaurin polynomial $T_{4}$ of $f(x)=\cos(x)$ and use
        it to approximate $\cos(0.2)$. Compare with the true value
        $\cos(0.2)=0.980067$.
  \item A firm produces $Q(L)=10\sqrt{L}$ units with $L$ hours of
        labour. Find $T_{2}$ of $Q$ at $a=100$ and use it to
        approximate output at $L=104$. Compare with the true value
        $Q(104)=101.980390$.
  \item (Log-linearisation.) Output grows from $Y_{0}$ to
        $Y_{1}=Y_{0}(1+g)$. Show, using $T_{1}$ of $\ln(1+x)$ at $0$,
        that $\ln Y_{1}-\ln Y_{0}\approx g$, and quantify the
        approximation for $g=0.03$ given $\ln(1.03)=0.029559$.
  \item Use the Lagrange remainder to bound the error of the $T_{1}$
        approximation of $\sqrt{4.1}$ at $a=4$ (Problem 2), and compare
        the bound with the actual error.
  \item Use the Lagrange remainder to bound the error of the $T_{2}$
        approximation of $\mathrm{e}^{0.1}$ at $a=0$ (Problem 6), and
        compare the bound with the actual error.
  \item Find the smallest $n$ such that the Maclaurin polynomial
        $T_{n}$ of $\mathrm{e}^{x}$ approximates $\mathrm{e}^{0.5}$
        with error at most $10^{-4}$. (Use $\mathrm{e}^{0.5}<2$ in the
        remainder bound.)
  \item Prove, using $T_{1}$ with Lagrange remainder: if $f'(a)=0$ and
        $f''$ is continuous with $f''(a)<0$, then $a$ is a strict local
        maximum of $f$.
\end{problems}

\section{Solutions}

\begin{solutions}
  \item $f(1)=\ln(1)=0$; $f'(x)=1/x$, so $f'(1)=1$. Hence
        \[
          T_{1}(x)=0+1\cdot(x-1)=x-1 .
        \]

  \item $f(4)=2$; $f'(x)=\dfrac{1}{2\sqrt{x}}$, so
        $f'(4)=\dfrac{1}{2\cdot 2}=\dfrac14$. Hence
        \[
          T_{1}(x)=2+\tfrac14(x-4) .
        \]

  \item $f(1)=1$; $f'(x)=-1/x^{2}$, so $f'(1)=-1$. Hence
        \[
          T_{1}(x)=1-(x-1)=2-x .
        \]

  \item $f(2)=8-4=4$; $f'(x)=3x^{2}-2$, so $f'(2)=12-2=10$. Hence
        \[
          T_{1}(x)=4+10(x-2) .
        \]

  \item $f(\ln 2)=\mathrm{e}^{2\ln 2}=(\mathrm{e}^{\ln 2})^{2}=4$;
        $f'(x)=2\mathrm{e}^{2x}$, so $f'(\ln 2)=2\cdot 4=8$. Hence
        \[
          T_{1}(x)=4+8(x-\ln 2) .
        \]

  \item $f^{(k)}(x)=\mathrm{e}^{x}$ for every $k$, so
        $f(0)=f'(0)=f''(0)=1$ and
        \[
          T_{2}(x)=1+x+\frac{x^{2}}{2} .
        \]
        At $x=0.1$: $T_{2}(0.1)=1+0.1+\tfrac{0.01}{2}=1.105$. True
        value $1.105171$; the error is about $1.7\cdot 10^{-4}$.

  \item $f(1)=0$; $f'(x)=1/x$, $f'(1)=1$; $f''(x)=-1/x^{2}$,
        $f''(1)=-1$. Hence
        \[
          T_{2}(x)=(x-1)-\tfrac12(x-1)^{2} .
        \]
        At $x=1.05$:
        $T_{2}(1.05)=0.05-\tfrac12(0.0025)=0.05-0.00125=0.04875$. True
        value $0.048790$; the error is about $4.0\cdot 10^{-5}$.

  \item From Problem 2, $f(4)=2$ and $f'(4)=\tfrac14$. Further
        $f''(x)=-\tfrac14 x^{-3/2}$, so
        $f''(4)=-\tfrac14\cdot\tfrac18=-\tfrac{1}{32}$ and the
        quadratic coefficient is
        $\tfrac12 f''(4)=-\tfrac{1}{64}$:
        \[
          T_{2}(x)=2+\tfrac14(x-4)-\tfrac{1}{64}(x-4)^{2} .
        \]
        At $x=4.1$:
        \[
          T_{2}(4.1)=2+0.025-\tfrac{0.01}{64}
                    =2.025-0.00015625=2.02484375 .
        \]
        True value $2.024846$; the error is about $1.9\cdot 10^{-6}$
        (the linear estimate $2.025$ alone is off by
        $1.5\cdot 10^{-4}$).

  \item $f(0)=1$; $f'(x)=\dfrac{1}{(1-x)^{2}}$, $f'(0)=1$;
        $f''(x)=\dfrac{2}{(1-x)^{3}}$, $f''(0)=2$, so
        $\tfrac12 f''(0)=1$. Hence
        \[
          T_{2}(x)=1+x+x^{2} ,
        \]
        the start of the geometric series $1+x+x^{2}+\cdots$.

  \item $f(1)=1\cdot\ln(1)=0$; $f'(x)=\ln(x)+1$, $f'(1)=1$;
        $f''(x)=1/x$, $f''(1)=1$. Hence
        \[
          T_{2}(x)=(x-1)+\tfrac12(x-1)^{2} .
        \]

  \item $f(0)=1$; $f'(x)=-2x\mathrm{e}^{-x^{2}}$, $f'(0)=0$;
        $f''(x)=(-2+4x^{2})\mathrm{e}^{-x^{2}}$, $f''(0)=-2$. Hence
        \[
          T_{2}(x)=1-x^{2} .
        \]
        (Equivalently: substitute $u=-x^{2}$ into
        $\mathrm{e}^{u}\approx 1+u$.)

  \item $f(0)=0$; $f'(x)=\dfrac{1}{1+x}$, $f'(0)=1$;
        $f''(x)=-\dfrac{1}{(1+x)^{2}}$, $f''(0)=-1$;
        $f'''(x)=\dfrac{2}{(1+x)^{3}}$, $f'''(0)=2$. The coefficients
        are $1$, $-\tfrac{1}{2!}=-\tfrac12$ and
        $\tfrac{2}{3!}=\tfrac13$:
        \[
          T_{3}(x)=x-\frac{x^{2}}{2}+\frac{x^{3}}{3} .
        \]
        At $x=-0.1$:
        \[
          T_{3}(-0.1)=-0.1-\frac{0.01}{2}+\frac{-0.001}{3}
                     =-0.1-0.005-0.000333=-0.105333 .
        \]
        True value $-0.105361$; the error is about $2.7\cdot 10^{-5}$.

  \item With $k=\tfrac12$ in the binomial expansion: the linear
        coefficient is $k=\tfrac12$ and the quadratic coefficient is
        \[
          \frac{k(k-1)}{2!}
            =\frac{\tfrac12\cdot\bigl(-\tfrac12\bigr)}{2}
            =-\frac18 ,
        \]
        so
        \[
          T_{2}(x)=1+\frac{x}{2}-\frac{x^{2}}{8} .
        \]
        At $x=0.1$: $T_{2}(0.1)=1+0.05-0.00125=1.04875$. True value
        $1.048809$; the error is about $5.9\cdot 10^{-5}$.

  \item Derivatives of $\cos$ at $0$ cycle through
        $1,\,0,\,-1,\,0,\,1$ for $k=0,\ldots,4$, so
        \[
          T_{4}(x)=1-\frac{x^{2}}{2!}+\frac{x^{4}}{4!}
                  =1-\frac{x^{2}}{2}+\frac{x^{4}}{24} .
        \]
        At $x=0.2$: $(0.2)^{2}=0.04$ and $(0.2)^{4}=0.0016$, so
        \[
          T_{4}(0.2)=1-0.02+\frac{0.0016}{24}
                    =1-0.02+0.000067=0.980067 ,
        \]
        agreeing with the true value $0.980067$ to all six decimals
        shown (the next term is of size $\tfrac{(0.2)^{6}}{720}\approx
        9\cdot 10^{-8}$).

  \item $Q(100)=10\cdot 10=100$; $Q'(L)=5L^{-1/2}$,
        $Q'(100)=\tfrac{5}{10}=0.5$; $Q''(L)=-\tfrac52 L^{-3/2}$,
        $Q''(100)=-\tfrac{2.5}{1000}=-0.0025$, so
        $\tfrac12 Q''(100)=-0.00125$. Hence
        \[
          T_{2}(L)=100+0.5\,(L-100)-0.00125\,(L-100)^{2} .
        \]
        At $L=104$ (so $L-100=4$):
        \[
          T_{2}(104)=100+0.5\cdot 4-0.00125\cdot 16=102-0.02=101.98 .
        \]
        True value $101.980390$; the error is about
        $3.9\cdot 10^{-4}$, four extra labour hours raise output by
        just under $2$ units, the shortfall against the linear estimate
        $102$ reflecting diminishing returns ($Q''<0$).

  \item Taking logs, $\ln Y_{1}=\ln Y_{0}+\ln(1+g)$. For
        $f(x)=\ln(1+x)$ at $a=0$: $f(0)=0$ and $f'(0)=1$, so
        $T_{1}(x)=x$ and
        \[
          \ln Y_{1}-\ln Y_{0}=\ln(1+g)\approx g .
        \]
        For $g=0.03$ the log change is $\ln(1.03)=0.029559$, while the
        growth rate is $0.03$: the approximation is off by about
        $4.4\cdot 10^{-4}$, i.e.\ by roughly $1.5\%$ of the quantity
        itself. (The second-order term $-\tfrac12 g^{2}=-0.00045$
        accounts for almost all of the gap.)

  \item Here $n=1$, $a=4$, $x=4.1$ and $f''(t)=-\tfrac14 t^{-3/2}$.
        For $t\in[4,4.1]$, $t^{-3/2}$ is decreasing, so
        \[
          |f''(t)|\leq\tfrac14\cdot 4^{-3/2}=\tfrac14\cdot\tfrac18
                 =\tfrac{1}{32}=:M .
        \]
        The remainder bound gives
        \[
          |R_{1}(4.1)|\leq\frac{M}{2!}\,(0.1)^{2}
            =\frac{1}{32}\cdot\frac{0.01}{2}
            =\frac{0.01}{64}\approx 1.56\cdot 10^{-4} .
        \]
        Actual error: $T_{1}(4.1)=2.025$ versus $\sqrt{4.1}=2.024846$,
        i.e.\ $1.54\cdot 10^{-4}$, just inside the bound.

  \item Here $n=2$, $a=0$, $x=0.1$ and
        $f'''(t)=\mathrm{e}^{t}$. For $t\in(0,0.1)$,
        $\mathrm{e}^{t}<\mathrm{e}^{0.1}<1.2=:M$, so
        \[
          |R_{2}(0.1)|\leq\frac{M}{3!}\,(0.1)^{3}
            =\frac{1.2\cdot 0.001}{6}=2\cdot 10^{-4} .
        \]
        Actual error: $T_{2}(0.1)=1.105$ versus
        $\mathrm{e}^{0.1}=1.105171$, i.e.\ $1.7\cdot 10^{-4}$,
        consistent with the bound.

  \item With $f(x)=\mathrm{e}^{x}$, $a=0$, $x=0.5$:
        $f^{(n+1)}(\xi)=\mathrm{e}^{\xi}\leq\mathrm{e}^{0.5}<2$ for
        $\xi\in(0,0.5)$, so
        \[
          |R_{n}(0.5)|\leq\frac{2\,(0.5)^{n+1}}{(n+1)!} .
        \]
        Tabulate:
        \[
          n=3:\ \frac{2\cdot 0.0625}{24}\approx 5.2\cdot 10^{-3},
          \qquad
          n=4:\ \frac{2\cdot 0.03125}{120}\approx 5.2\cdot 10^{-4},
        \]
        \[
          n=5:\ \frac{2\cdot 0.015625}{720}\approx 4.3\cdot 10^{-5} .
        \]
        The first value below $10^{-4}$ occurs at $n=5$. Check with the
        true value $\mathrm{e}^{0.5}=1.648721$: $T_{4}(0.5)=1.648438$
        (error $2.8\cdot 10^{-4}$, indeed too large) while
        $T_{5}(0.5)=1.648698$ (error $2.3\cdot 10^{-5}$). So $n=5$.

  \item By Taylor's theorem with $n=1$: for each $x$ near $a$ there is
        $\xi$ between $a$ and $x$ with
        \[
          f(x)=f(a)+f'(a)(x-a)+\tfrac12 f''(\xi)(x-a)^{2}
              =f(a)+\tfrac12 f''(\xi)(x-a)^{2},
        \]
        using $f'(a)=0$. Since $f''$ is continuous and $f''(a)<0$,
        there is a neighbourhood $U$ of $a$ on which $f''<0$. For
        $x\in U$, $x\neq a$, the point $\xi$ also lies in $U$, so
        $f''(\xi)<0$ while $(x-a)^{2}>0$, giving
        \[
          f(x)-f(a)=\tfrac12 f''(\xi)(x-a)^{2}<0 .
        \]
        Hence $f(x)<f(a)$ for all $x\in U\setminus\{a\}$: $a$ is a
        strict local maximum. (This is the second-order condition in
        its cleanest form: near a stationary point the function
        \emph{is} its osculating parabola, up to higher-order terms.)
\end{solutions}

\section{Additional Problems}

\begin{problems}
  \item Find $T_{1}$ of $f(x)=(1+x)^{5}$ at $a=0$.
  \item Find $T_{1}$ of $f(x)=\sin(x)$ at $a=0$ and use it to
        approximate $\sin(0.1)$. Compare with the true value
        $\sin(0.1)=0.099833$.
  \item Find $T_{2}$ of $f(x)=x^{4}$ at $a=1$.
  \item Find $T_{2}$ of $f(x)=\dfrac{1}{x}$ at $a=2$.
  \item Find $T_{3}$ of $f(x)=\mathrm{e}^{x}$ at $a=0$ and use it to
        approximate $\mathrm{e}^{0.2}$. Compare with the true value
        $\mathrm{e}^{0.2}=1.221403$.
  \item Find $T_{2}$ of $f(x)=\ln(1+2x)$ at $a=0$.
  \item (Cobb--Douglas.) For $f(K)=K^{0.3}$, find $T_{2}$ at $a=1$ and
        use it to approximate $1.1^{0.3}$. Compare with the true value
        $1.1^{0.3}=1.029006$.
  \item Use the Lagrange remainder to bound the error of approximating
        $\ln(1.05)$ by $T_{1}(x)=x-1$ of $\ln(x)$ at $a=1$, and compare
        the bound with the actual error (true value
        $\ln(1.05)=0.048790$).
  \item Find the smallest $n$ such that the Maclaurin polynomial
        $T_{n}$ of $\ln(1+x)$ approximates $\ln(1.1)$ with error at
        most $10^{-6}$ (true value $\ln(1.1)=0.095310$).
  \item Suppose $|f'''(t)|\leq M$ for all $t$ between $a$ and $x$.
        Derive a bound for $|f(x)-T_{2}(x)|$ in terms of $M$ and
        $|x-a|$.
\end{problems}

\section{Additional Solutions}

\begin{solutions}
  \item $f(0)=1$; $f'(x)=5(1+x)^{4}$, $f'(0)=5$. Hence
        \[
          T_{1}(x)=1+5x .
        \]

  \item $f(0)=0$; $f'(x)=\cos(x)$, $f'(0)=1$. Hence $T_{1}(x)=x$ (the
        small-angle approximation) and $\sin(0.1)\approx 0.1$. True
        value $0.099833$; the error is about $1.7\cdot 10^{-4}$
        (matching the size of the first omitted term,
        $\tfrac{x^{3}}{3!}=\tfrac{0.001}{6}$).

  \item $f(1)=1$; $f'(x)=4x^{3}$, $f'(1)=4$; $f''(x)=12x^{2}$,
        $f''(1)=12$, so $\tfrac12 f''(1)=6$. Hence
        \[
          T_{2}(x)=1+4(x-1)+6(x-1)^{2} .
        \]

  \item $f(2)=\tfrac12$; $f'(x)=-1/x^{2}$, $f'(2)=-\tfrac14$;
        $f''(x)=2/x^{3}$, $f''(2)=\tfrac28=\tfrac14$, so
        $\tfrac12 f''(2)=\tfrac18$. Hence
        \[
          T_{2}(x)=\frac12-\frac14(x-2)+\frac18(x-2)^{2} .
        \]

  \item All derivatives equal $1$ at $0$, so
        \[
          T_{3}(x)=1+x+\frac{x^{2}}{2}+\frac{x^{3}}{6} .
        \]
        At $x=0.2$:
        \[
          T_{3}(0.2)=1+0.2+\frac{0.04}{2}+\frac{0.008}{6}
                    =1+0.2+0.02+0.001333=1.221333 .
        \]
        True value $1.221403$; the error is about $6.9\cdot 10^{-5}$.

  \item $f(0)=0$; $f'(x)=\dfrac{2}{1+2x}$, $f'(0)=2$;
        $f''(x)=-\dfrac{4}{(1+2x)^{2}}$, $f''(0)=-4$, so
        $\tfrac12 f''(0)=-2$. Hence
        \[
          T_{2}(x)=2x-2x^{2} .
        \]
        (Equivalently: substitute $u=2x$ into
        $\ln(1+u)\approx u-\tfrac12 u^{2}$.)

  \item $f(1)=1$; $f'(K)=0.3\,K^{-0.7}$, $f'(1)=0.3$;
        $f''(K)=0.3\cdot(-0.7)\,K^{-1.7}=-0.21\,K^{-1.7}$,
        $f''(1)=-0.21$, so $\tfrac12 f''(1)=-0.105$. Hence
        \[
          T_{2}(K)=1+0.3\,(K-1)-0.105\,(K-1)^{2} .
        \]
        At $K=1.1$:
        \[
          T_{2}(1.1)=1+0.03-0.105\cdot 0.01=1.03-0.00105=1.02895 .
        \]
        True value $1.029006$; the error is about $5.6\cdot 10^{-5}$.

  \item Here $n=1$, $a=1$, $x=1.05$ and $f''(t)=-1/t^{2}$. On
        $[1,1.05]$, $|f''(t)|=1/t^{2}\leq 1=:M$. The remainder bound
        gives
        \[
          |R_{1}(1.05)|\leq\frac{M}{2!}\,(0.05)^{2}
            =\frac{0.0025}{2}=1.25\cdot 10^{-3} .
        \]
        Actual error: $T_{1}(1.05)=0.05$ versus $0.048790$, i.e.\
        $1.21\cdot 10^{-3}$, close to the bound, because
        $1/t^{2}\approx 1$ throughout the interval.

  \item For $f(x)=\ln(1+x)$,
        $f^{(n+1)}(t)=\dfrac{(-1)^{n}\,n!}{(1+t)^{n+1}}$, so on
        $[0,0.1]$ we have $|f^{(n+1)}(t)|\leq n!$ and
        \[
          |R_{n}(0.1)|\leq\frac{n!}{(n+1)!}\,(0.1)^{n+1}
            =\frac{(0.1)^{n+1}}{n+1} .
        \]
        Tabulate:
        \[
          n=4:\ \frac{(0.1)^{5}}{5}=2\cdot 10^{-6},
          \qquad
          n=5:\ \frac{(0.1)^{6}}{6}\approx 1.7\cdot 10^{-7} .
        \]
        The first value below $10^{-6}$ occurs at $n=5$. Check against
        the true value $0.095310$: $T_{4}(0.1)=0.1-0.005+0.000333
        -0.000025=0.095308$ (error $1.9\cdot 10^{-6}$, indeed too
        large), while $T_{5}(0.1)=0.095308+0.000002=0.095310$ (error
        about $1.5\cdot 10^{-7}$). So $n=5$.

  \item Apply the Lagrange form with $n=2$: there is $\xi$ between $a$
        and $x$ with
        \[
          f(x)-T_{2}(x)=\frac{f'''(\xi)}{3!}(x-a)^{3},
        \]
        hence
        \[
          |f(x)-T_{2}(x)|\leq\frac{|f'''(\xi)|}{6}\,|x-a|^{3}
            \leq\frac{M}{6}\,|x-a|^{3} .
        \]
\end{solutions}

\chapter{Implicit Differentiation}

Most curves in the plane are not graphs. The circle $x^{2}+y^{2}=25$ fails
the vertical line test, the equation $\mathrm{e}^{y}+xy=1$ cannot be solved
for $y$ in closed form at all, and an indifference curve is handed to us as
a level set rather than as a formula. Yet each of them has a tangent at
almost every point, and economics asks for its slope constantly: that slope
\emph{is} the marginal rate of substitution.

Implicit differentiation gets it without solving for $y$. The move rests on
one idea, and the whole chapter is that idea applied repeatedly: pretend the
equation has already been solved near the point of interest, so that $y$ is
a function $y(x)$; substituting it makes the equation an identity in $x$;
identities may be differentiated. Because every $y$ in sight is secretly the
composite $y(x)$, the chain rule leaves a factor $y'$ behind wherever a $y$
was, and the resulting linear equation is solved for $y'$.

Two questions have to be kept apart, and the chapter answers both. \emph{Is}
there such a function $y(x)$ near the point --- and if so, what is its
slope? The first question is the Implicit Function Theorem, and the honest
answer is sometimes no: at a point where the curve crosses itself or stands
vertically, no function $y(x)$ describes it, and the slope formula returns a
division by zero to say so.

\section{Core Rules}

\begin{rulebox}{Why a factor $y'$ appears}
Suppose that near a point of the curve the equation really does pin down one
$y$ for each $x$, and call that function $y(x)$. Substituting it turns the
equation into an \textbf{identity in $x$}, true for every $x$ near $x_{0}$:
\[
  x^{2}+\bigl(y(x)\bigr)^{2}=25 .
\]
An identity may be differentiated on both sides. Every $y$ in the expression
is the composite $y(x)$, so it is differentiated by the \textbf{chain rule}
and leaves a factor $y'$ behind:
\[
  \frac{\mathrm{d}}{\mathrm{d}x}\bigl[y^{2}\bigr]
  =\frac{\mathrm{d}}{\mathrm{d}x}\bigl[(y(x))^{2}\bigr]
  =2\,y(x)\,y'(x)=2yy' .
\]
Hence $2x+2yy'=0$ and $y'=-x/y$. Nothing was solved for $y$, and the answer
is \emph{allowed} to contain $y$: it is the slope at a point of the curve,
and the two points with the same $x$ --- here $(3,4)$ and $(3,-4)$ --- have
slopes $-3/4$ and $+3/4$.
\end{rulebox}

\begin{rulebox}{Setup}
Suppose $F(x,y)=0$ defines $y=y(x)$ implicitly. Differentiate every term
with respect to $x$. Whenever an expression involves $y$, the chain rule
produces an extra factor $\dd{y}{x}$:
\[
  \frac{\mathrm{d}}{\mathrm{d}x}\bigl[g(y)\bigr]=g'(y)\cdot\dd{y}{x} .
\]
So $\frac{\mathrm{d}}{\mathrm{d}x}[\sin y]=\cos(y)\,y'$,
$\frac{\mathrm{d}}{\mathrm{d}x}[\mathrm{e}^{y}]=\mathrm{e}^{y}y'$, and
$\frac{\mathrm{d}}{\mathrm{d}x}[\ln y]=y'/y$.
\end{rulebox}

\begin{rulebox}{Mixed terms}
For a product such as $x^{3}y^{2}$ apply the product rule, treating $y$ as
a function of $x$:
\[
  \frac{\mathrm{d}}{\mathrm{d}x}\bigl[x^{3}y^{2}\bigr]=3x^{2}y^{2}+x^{3}\cdot 2yy' .
\]
\end{rulebox}

\begin{rulebox}{Solving for $y'$}
Collect all terms containing $y'$ on one side and factor:
\[
  A(x,y)\,y'+B(x,y)=0\;\Longrightarrow\;y'=-\frac{B(x,y)}{A(x,y)} .
\]
The equation is always \emph{linear} in $y'$, however complicated it looks,
because $y'$ only ever enters through one chain-rule step.
\end{rulebox}

\begin{rulebox}{The curve is a level set --- and that gives the formula at once}
Write the equation as $F(x,y)=0$. Then the curve is the \textbf{level set of
the surface $z=F(x,y)$ at height $0$}: the shoreline where the surface cuts
the horizontal plane $z=0$. Walking along the curve, $F$ never changes.
Writing that down with the chain rule for $F\bigl(x,y(x)\bigr)=0$,
\[
  \pd{F}{x}+\pd{F}{y}\cdot\dd{y}{x}=0
  \qquad\Longrightarrow\qquad
  \boxed{\;\dd{y}{x}=-\frac{F_{x}(x,y)}{F_{y}(x,y)}\;}
  \qquad\text{provided } F_{y}(x,y)\neq0 .
\]
This is not a new method, it is the same computation done once and for all:
for $F=x^{2}+y^{2}-25$ it returns $-2x/2y=-x/y$, as above. In economics the
curve is an indifference curve $U(x,y)=\bar u$ and the formula reads
$\dd{y}{x}=-U_{x}/U_{y}$: the marginal rate of substitution is the ratio of
marginal utilities.
\end{rulebox}

\begin{gallery}
\begin{tikzpicture}
  \begin{axis}[surfaxis, title={$z=F(x,y)=x^{3}-3x^{2}+y^{2}-2y+1$}]
    \addplot3[surf, domain=-1.2:3.2, y domain=-1.6:3.6]{x^3-3*x^2+y^2-2*y+1};
  \end{axis}
\end{tikzpicture}\hfill
\begin{tikzpicture}
  \begin{axis}[geoaxis, xmin=-1.6, xmax=3.6, ymin=-1.9, ymax=3.9,
      xtick={-1,1,2,3}, ytick={-1,1,2,3}, title={the level set $F=0$}]
    \addplot[domain=-1.2:3, samples=241]{1+abs(x)*sqrt(3-x)};
    \addplot[domain=-1.2:3, samples=241]{1-abs(x)*sqrt(3-x)};
    \addplot[only marks, mark=*, mark size=1.3pt, accent, forget plot]
      coordinates {(0,1) (3,1)};
  \end{axis}
\end{tikzpicture}
\gcap{The surface and its zero level set $k$. At $(0,1)$ the curve crosses
itself; at $(3,1)$ it stands vertical. Everywhere else a local $y(x)$
exists.}
\end{gallery}

\begin{rulebox}{When there is no $y(x)$: the Implicit Function Theorem}
\textbf{Theorem.} If $F$ is $C^{1}$ near $(x_{0},y_{0})$, if
$F(x_{0},y_{0})=0$, and if
\[
  F_{y}(x_{0},y_{0})\neq0 ,
\]
then there are neighbourhoods on which the equation $F(x,y)=0$ defines
exactly one differentiable function $y=y(x)$ with $y(x_{0})=y_{0}$, and its
derivative is $-F_{x}/F_{y}$.

The condition is exactly what the formula needs, and its failure is
visible in the picture. Where $F_{y}=0$:
\begin{itemize}[leftmargin=1.6em,itemsep=0.3ex]
  \item if $F_{x}\neq0$, the curve has a \textbf{vertical tangent}. No
        function $y(x)$ can describe it there --- but the roles may be
        swapped: $x=x(y)$ exists, with
        $\dd{x}{y}=-F_{y}/F_{x}$, and that derivative is $0$.
  \item if $F_{x}=0$ as well, the point is \textbf{singular}: typically the
        curve crosses itself, and \emph{two} branches with two different
        slopes pass through. No single function describes either the curve
        or its slope, and $-F_{x}/F_{y}$ degenerates to $0/0$. Separate the
        branches first, then differentiate each one.
\end{itemize}
Reading it off a picture: look for vertical tangents and self-intersections.
Those are the only places the theorem can fail. A failed condition is a
warning, not a proof that no $y(x)$ exists: $y^{3}-x=0$ has $F_{y}=0$ at the
origin, yet $y=x^{1/3}$ solves it for every $x$ (with a vertical tangent, so
not differentiably). Chapter~14 collects such cases.
\end{rulebox}

\begin{rulebox}{Logarithmic differentiation}
If $y=f(x)>0$ is a complicated product, quotient or power, take $\ln$
first:
\[
  \ln y=\ln f(x)
  \;\Longrightarrow\;
  \frac{y'}{y}=\frac{\mathrm{d}}{\mathrm{d}x}\ln f(x)
  \;\Longrightarrow\;
  y'=y\cdot\frac{\mathrm{d}}{\mathrm{d}x}\ln f(x) .
\]
Particularly useful for $y=u(x)^{v(x)}$: $\ln y=v(x)\ln u(x)$. This is
implicit differentiation applied to $\ln y-\ln f(x)=0$, which is why the
factor $y'/y$ appears.
\end{rulebox}

\begin{rulebox}{Higher-order implicit}
To find $y''$, differentiate $y'$ a second time with respect to $x$, again
remembering that any $y$ that appears contributes a factor $y'$. Then
substitute the expression already found for $y'$ to eliminate it.
\end{rulebox}

\begin{rulebox}{Tangent line via implicit slope}
The tangent at $(x_{0},y_{0})$ on $F(x,y)=0$ is
$y-y_{0}=y'(x_{0})\cdot(x-x_{0})$, where $y'(x_{0})$ is computed
implicitly. Both coordinates are substituted into $-F_{x}/F_{y}$ before the
line is written down.
\end{rulebox}

\section{Worked Example}

\begin{worked}{The folium $x^{3}+y^{3}=6xy$: find $y'$, the tangent at $(3,3)$, and every point where the tangent is vertical}
\wstep{Step 1. Confirm the point is on the curve.}
$27+27=54$ and $6\cdot3\cdot3=54$. \checkmark\ Do this first: an implicit
slope formula evaluated at a point that is not on the curve returns a
perfectly plausible number that means nothing.

\wstep{Step 2. Differentiate every term with respect to $x$, treating $y$ as $y(x)$.}
Left side: $\frac{\mathrm{d}}{\mathrm{d}x}\bigl[x^{3}+y^{3}\bigr]
=3x^{2}+3y^{2}y'$ --- the chain rule leaves a factor $y'$ behind wherever a
$y$ was.

Right side: $6xy$ is a \textbf{product} of two functions of $x$, so the
product rule applies:
\[
  \frac{\mathrm{d}}{\mathrm{d}x}\bigl[6xy\bigr]=6y+6xy' .
\]
\wtrap{The right side is where this problem is lost. Writing $6y$ alone
(treating $y$ as a constant) or $6xy'$ alone (forgetting the product rule)
are both common. Every occurrence of $y$ needs the chain rule, \emph{and}
every product needs the product rule --- they compound rather than replace
each other.}

\wstep{Step 3. Collect the $y'$ terms; the equation is always linear in $y'$.}
\[
  3x^{2}+3y^{2}y'=6y+6xy'
  \;\Longrightarrow\;
  3y^{2}y'-6xy'=6y-3x^{2}
  \;\Longrightarrow\;
  y'\bigl(3y^{2}-6x\bigr)=6y-3x^{2} .
\]
\[
  \boxed{\;y'=\frac{6y-3x^{2}}{3y^{2}-6x}=\frac{2y-x^{2}}{y^{2}-2x}\;}
\]
The answer contains $y$, and it must: two points of the curve can share an
$x$ and have different slopes, so $x$ alone cannot determine $y'$.

\wstep{Step 4. Evaluate at the point.}
\[
  y'(3,3)=\frac{2\cdot3-9}{9-6}=\frac{-3}{3}=-1 ,
\]
so the tangent is
\[
  y-3=-1(x-3)\quad\Longrightarrow\quad y=6-x .
\]
By the symmetry of the equation in $x$ and $y$, the curve is symmetric about
the line $y=x$, and $(3,3)$ lies on that line --- so a slope of $-1$,
perpendicular to it, is exactly what the geometry predicts. \checkmark

\wstep{Step 5. Cross-check with the level-set formula.}
Writing $F(x,y)=x^{3}+y^{3}-6xy$,
\[
  F_{x}=3x^{2}-6y,\qquad F_{y}=3y^{2}-6x,
  \qquad
  y'=-\frac{F_{x}}{F_{y}}=-\frac{3x^{2}-6y}{3y^{2}-6x}=\frac{6y-3x^{2}}{3y^{2}-6x},
\]
the same expression. The two routes --- differentiate term by term, or
apply $-F_{x}/F_{y}$ --- are the same computation; use whichever you find
faster, but know that they must agree.

\wstep{Step 6. Vertical tangents: where the denominator vanishes.}
$y'$ blows up when $F_{y}=3y^{2}-6x=0$, i.e.\ $x=y^{2}/2$. Substituting into
the curve:
\[
  \frac{y^{6}}{8}+y^{3}=6\cdot\frac{y^{2}}{2}\cdot y=3y^{3}
  \;\Longrightarrow\;
  \frac{y^{6}}{8}=2y^{3}
  \;\Longrightarrow\;
  y^{3}\bigl(y^{3}-16\bigr)=0 .
\]
So $y=0$ (giving $x=0$) or $y=16^{1/3}=2\sqrt[3]{2}\approx2.520$, giving
$x=y^{2}/2\approx3.175$. At the latter the tangent is genuinely vertical.
\wtrap{At the origin the numerator $2y-x^{2}$ also vanishes, so $y'$ is
$0/0$ there --- undefined for a different reason. The folium
\emph{crosses itself} at $(0,0)$, with two distinct tangents, so no single
function $y(x)$ describes the curve near it. That is precisely the case the
Implicit Function Theorem excludes, and a $0/0$ in the slope formula is how
it announces itself.}

\wstep{Step 7. The economic version of the same calculation.}
An indifference curve is a level set $U(x,y)=\bar u$, and the identical
manipulation gives
\[
  \frac{\mathrm{d}y}{\mathrm{d}x}=-\frac{U_{x}}{U_{y}}=-\mathrm{MRS} .
\]
For $U=x^{1/2}y^{1/2}=4$: $U_{x}=\tfrac12x^{-1/2}y^{1/2}$,
$U_{y}=\tfrac12x^{1/2}y^{-1/2}$, so $y'=-y/x$. At $(4,4)$ the slope is
$-1$; at $(8,2)$ it is $-\tfrac14$ --- the indifference curve flattens as
$x$ becomes abundant, which is diminishing marginal rate of substitution
stated as a derivative.
\end{worked}

\section{Practice Problems}

\begin{problems}
  \item Compute $\dd{y}{x}$ implicitly for $x^{2}+4y^{2}=16$.
  \item Find the equation of the tangent line to $x^{2}+4y^{2}=20$ at
        $(2,2)$.
  \item Use logarithmic differentiation to find $y'$ for $y=x^{\sin x}$
        ($x>0$, $\sin x>0$).
  \item Compute $\dd{^{2}y}{x^{2}}$ implicitly for $x^{2}+4y^{2}=4$.
  \item Compute $\dd{y}{x}$ implicitly for $x^{3}-y^{3}=3xy$.
  \item Compute $\dd{y}{x}$ implicitly for $x^{2}y^{3}=8$.
  \item Find the equation of the tangent line to $x^{2}y+y^{2}x=6$ at
        $(1,2)$.
  \item Use logarithmic differentiation to find $y'$ for $y=(\ln x)^{x}$
        ($x>1$).
  \item Compute $\dd{^{2}y}{x^{2}}$ implicitly for $x^{2}y=1$.
  \item Compute $\dd{y}{x}$ implicitly for $\cos(y)+x=y$.
  \item Find the equation of the tangent line to $y^{2}-x^{2}=3$ at
        $(1,2)$.
  \item Use logarithmic differentiation to find $y'$ for
        \[
          y=\frac{(2x+1)^{5}\sqrt{x^{2}+3}}{(x-4)^{3}}\qquad (x>4).
        \]
  \item Compute $\dd{y}{x}$ implicitly for $\mathrm{e}^{x+y}=xy$.
  \item Compute $\dd{^{2}y}{x^{2}}$ implicitly for $y^{3}=x^{2}$.
  \item Compute $\dd{y}{x}$ implicitly for $x^{2}+2xy+3y^{2}=12$.
  \item Find the equation of the tangent line to $\mathrm{e}^{y}+xy=\mathrm{e}+1$
        at $(1,1)$.
  \item Use logarithmic differentiation to find $y'$ for $y=(x^{2}+1)^{x}$
        ($x>0$).
  \item Compute $\dd{y}{x}$ implicitly for $\sqrt{x+y}=x-y$.
  \item Compute $\dd{^{2}y}{x^{2}}$ implicitly for $xy+y^{2}=1$.
  \item Find the equation of the tangent line to $x^{2}+y^{2}=4y$ at
        $(\sqrt{3},1)$.
  \item Compute $\dd{y}{x}$ implicitly for $\tan(xy)=x$.
  \item Use logarithmic differentiation to find $y'$ for $y=x^{1/x}$
        ($x>0$).
  \item Compute $\dd{y}{x}$ implicitly for $\ln(y)+xy=2$.
  \item Compute $\dd{y}{x}$ implicitly for $y^{3}+y=x^{2}$.
  \item Find the tangent line to the astroid $x^{2/3}+y^{2/3}=2$ at
        $(1,1)$.
  \item Compute $\dd{y}{x}$ implicitly for $\sin(x+y)=y$.
  \item Use logarithmic differentiation to find $y'$ for
        $y=\dfrac{\mathrm{e}^{2x}(x+1)^{3}}{\sqrt{x^{2}+5}}$ (assume $x>-1$).
  \item Compute $\dd{^{2}y}{x^{2}}$ implicitly for $x^{2}+y^{2}=25$.
  \item For an indifference curve given by $U(x,y)=x^{1/2}y^{1/2}=\bar U$
        (constant), compute $\dd{y}{x}$, the marginal rate of
        substitution.
  \item Find all points on the curve $x^{2}+y^{2}=2x+4y$ where the tangent
        line is horizontal.
  \item Compute $\dd{y}{x}$ implicitly for $\arctan(y/x)=x$.
  \item Use logarithmic differentiation to find $y'$ for $y=x^{x}$
        ($x>0$).
  \item Find the equation of the tangent line to $x^{3}+y^{3}=9xy$ at
        $(2,4)$ (point on the folium of Descartes).
  \item Compute $\dd{^{2}y}{x^{2}}$ implicitly for $y^{2}=4x$.
  \item Compute $\dd{y}{x}$ implicitly for $x\mathrm{e}^{y}+y=x$.
  \item Find all points on the ellipse $\dfrac{x^{2}}{9}+\dfrac{y^{2}}{4}=1$
        where the tangent line has slope $-1$.
  \item Use logarithmic differentiation to find $y'$ for
        $y=(\sin x)^{\cos x}$ on $(0,\pi/2)$.
  \item Compute $\dd{y}{x}$ implicitly for $x^{2}y-xy^{2}+x^{2}+y^{2}=0$.
  \item For the production isoquant $K^{1/3}L^{2/3}=\bar Q$ (constant),
        compute $\dd{K}{L}$, the marginal rate of technical
        substitution.
  \item Compute $\dd{^{2}y}{x^{2}}$ implicitly for $\sin(y)=x$.
  \item Find the tangent line to $xy=4$ at $x=2$.
  \item Use logarithmic differentiation to find $y'$ for
        $y=\dfrac{\sqrt{x+1}\sqrt[3]{x-2}}{(x+5)^{4}}$ (assume $x>2$).
  \item Compute $\dd{y}{x}$ implicitly for $\mathrm{e}^{y}\ln(x)+x^{2}=y$.
  \item Show that on the circle $x^{2}+y^{2}=r^{2}$ the tangent at
        $(x_{0},y_{0})$ is perpendicular to the radius. (I.e.\ derive the
        tangent slope and compare with the radius slope.)
  \item Use logarithmic differentiation to find $y'$ for $y=7^{x^{2}}$
        ($y>0$).
  \item Compute $\dd{^{2}y}{x^{2}}$ implicitly for $x^{2}+xy+y^{2}=3$.
  \item For the budget surface $px+qy=m$ (a line, $p,q,m>0$ constants),
        differentiate implicitly to recover $\dd{y}{x}=-p/q$.
  \item Find the tangent line to $x^{2}-xy+y^{2}=7$ at $(3,2)$.
  \item Compute $\dd{y}{x}$ implicitly for $\cos(xy)=x+y$.
  \item Compute $\dd{^{2}y}{x^{2}}$ implicitly for $x^{3}+y^{3}=6xy$ at the
        point $(3,3)$ (folium of Descartes). Express the answer as a
        number.
\end{problems}

\section{Solutions}

\begin{solutions}
  \item $2x+8yy'=0\Rightarrow y'=-\dfrac{x}{4y}$.

  \item Check $(2,2)$: $4+16=20$. From $2x+8yy'=0$: $y'=-x/(4y)$. At
        $(2,2)$: $y'=-2/8=-1/4$. Tangent: $y-2=-\tfrac14(x-2)$, i.e.\
        $x+4y=10$.

  \item $\ln y=\sin(x)\ln(x)\Rightarrow\dfrac{y'}{y}=\cos(x)\ln(x)+\dfrac{\sin x}{x}$.
        So
        \[
          y'=x^{\sin x}\left[\cos(x)\ln(x)+\frac{\sin x}{x}\right] .
        \]

  \item $2x+8yy'=0\Rightarrow y'=-\dfrac{x}{4y}$. Differentiate again:
        \[
          y''=-\frac{4y\cdot 1-x\cdot 4y'}{(4y)^{2}}=-\frac{y-xy'}{4y^{2}} .
        \]
        Substitute $y'=-x/(4y)$:
        $y-xy'=y+\dfrac{x^{2}}{4y}=\dfrac{4y^{2}+x^{2}}{4y}=\dfrac{4}{4y}
        =\dfrac{1}{y}$ (using $x^{2}+4y^{2}=4$). Hence
        $y''=-\dfrac{1}{4y^{3}}$.

  \item $3x^{2}-3y^{2}y'=3y+3xy'\Rightarrow -3(y^{2}+x)y'=3(y-x^{2})
        \Rightarrow y'=\dfrac{x^{2}-y}{y^{2}+x}$.

  \item $2xy^{3}+3x^{2}y^{2}y'=0\Rightarrow
        y'=-\dfrac{2xy^{3}}{3x^{2}y^{2}}=-\dfrac{2y}{3x}$.

  \item Differentiate $x^{2}y+y^{2}x=6$: $2xy+x^{2}y'+2yxy'+y^{2}=0
        \Rightarrow y'(x^{2}+2xy)=-(2xy+y^{2})
        \Rightarrow y'=-\dfrac{y(2x+y)}{x(x+2y)}$. At $(1,2)$:
        $y'=-\dfrac{2\cdot 4}{1\cdot 5}=-\dfrac{8}{5}$. Tangent:
        $y-2=-\tfrac85(x-1)$, i.e.\ $8x+5y=18$.

  \item $\ln y=x\ln(\ln x)\Rightarrow\dfrac{y'}{y}=\ln(\ln x)+x\cdot
        \dfrac{1/x}{\ln x}=\ln(\ln x)+\dfrac{1}{\ln x}$. So
        $y'=(\ln x)^{x}\left[\ln(\ln x)+\dfrac{1}{\ln x}\right]$.

  \item Explicit: $y=x^{-2}$, $y'=-2x^{-3}$, $y''=6x^{-4}$. Implicit
        check: $2xy+x^{2}y'=0\Rightarrow y'=-2y/x$. Differentiate:
        \[
          y''=-\frac{2y'x-2y}{x^{2}}=-\frac{2(-2y/x)x-2y}{x^{2}}
             =-\frac{-4y-2y}{x^{2}}=\frac{6y}{x^{2}}=\frac{6}{x^{4}} .
        \]

  \item $-\sin(y)y'+1=y'\Rightarrow y'(1+\sin y)=1\Rightarrow
        y'=\dfrac{1}{1+\sin y}$.

  \item $2yy'-2x=0\Rightarrow y'=x/y$. At $(1,2)$: $y'=1/2$. Tangent:
        $y-2=\tfrac12(x-1)$, i.e.\ $x-2y+3=0$.

  \item $\ln y=5\ln(2x+1)+\tfrac12\ln(x^{2}+3)-3\ln(x-4)$. Then
        \[
          \frac{y'}{y}=\frac{10}{2x+1}+\frac{x}{x^{2}+3}-\frac{3}{x-4} .
        \]
        Hence
        \[
          y'=\frac{(2x+1)^{5}\sqrt{x^{2}+3}}{(x-4)^{3}}
             \left[\frac{10}{2x+1}+\frac{x}{x^{2}+3}-\frac{3}{x-4}\right] .
        \]

  \item $\mathrm{e}^{x+y}(1+y')=y+xy'\Rightarrow
        y'(\mathrm{e}^{x+y}-x)=y-\mathrm{e}^{x+y}\Rightarrow
        y'=\dfrac{y-\mathrm{e}^{x+y}}{\mathrm{e}^{x+y}-x}$.

  \item $3y^{2}y'=2x\Rightarrow y'=\dfrac{2x}{3y^{2}}$. Differentiate
        again:
        \[
          y''=\frac{2\cdot 3y^{2}-2x\cdot 6yy'}{9y^{4}}
             =\frac{6y^{2}-12xy\cdot y'}{9y^{4}}
             =\frac{2}{3y^{2}}-\frac{4xy'}{3y^{3}} .
        \]
        Substitute $y'=2x/(3y^{2})$:
        \[
          y''=\frac{2}{3y^{2}}-\frac{4x}{3y^{3}}\cdot\frac{2x}{3y^{2}}
             =\frac{2}{3y^{2}}-\frac{8x^{2}}{9y^{5}}
             =\frac{6y^{3}-8x^{2}}{9y^{5}}
             =\frac{6x^{2}-8x^{2}}{9y^{5}}=-\frac{2x^{2}}{9y^{5}},
        \]
        using $y^{3}=x^{2}$. Equivalently, $y''=-\dfrac{2y^{3}}{9y^{5}}
        =-\dfrac{2}{9y^{2}}$.

  \item $2x+2y+2xy'+6yy'=0\Rightarrow y'(2x+6y)=-(2x+2y)\Rightarrow
        y'=-\dfrac{x+y}{x+3y}$.

  \item Differentiate $\mathrm{e}^{y}+xy=\mathrm{e}+1$:
        $\mathrm{e}^{y}y'+y+xy'=0\Rightarrow y'(\mathrm{e}^{y}+x)=-y
        \Rightarrow y'=-\dfrac{y}{\mathrm{e}^{y}+x}$. At $(1,1)$:
        $y'=-\dfrac{1}{\mathrm{e}+1}$. Tangent:
        $y-1=-\dfrac{1}{\mathrm{e}+1}(x-1)$.

  \item $\ln y=x\ln(x^{2}+1)\Rightarrow\dfrac{y'}{y}=\ln(x^{2}+1)
        +x\cdot\dfrac{2x}{x^{2}+1}=\ln(x^{2}+1)+\dfrac{2x^{2}}{x^{2}+1}$.
        So
        \[
          y'=(x^{2}+1)^{x}\left[\ln(x^{2}+1)+\frac{2x^{2}}{x^{2}+1}\right] .
        \]

  \item $\dfrac{1+y'}{2\sqrt{x+y}}=1-y'$. Multiply by $2\sqrt{x+y}$:
        $1+y'=2\sqrt{x+y}(1-y')\Rightarrow y'(1+2\sqrt{x+y})
        =2\sqrt{x+y}-1\Rightarrow y'=\dfrac{2\sqrt{x+y}-1}{1+2\sqrt{x+y}}$.

  \item $y+xy'+2yy'=0\Rightarrow y'(x+2y)=-y\Rightarrow
        y'=-\dfrac{y}{x+2y}$. Differentiate again using the quotient
        rule:
        \[
          y''=-\frac{y'(x+2y)-y(1+2y')}{(x+2y)^{2}}
             =-\frac{y'(x+2y)-y-2yy'}{(x+2y)^{2}}
             =-\frac{xy'-y}{(x+2y)^{2}} .
        \]
        Substitute $y'=-y/(x+2y)$:
        \[
          xy'-y=-\frac{xy}{x+2y}-y=-\frac{xy+y(x+2y)}{x+2y}
               =-\frac{2xy+2y^{2}}{x+2y}=-\frac{2y(x+y)}{x+2y} .
        \]
        Hence
        \[
          y''=\frac{2y(x+y)}{(x+2y)^{3}} .
        \]

  \item $2x+2yy'=4y'\Rightarrow y'(2y-4)=-2x\Rightarrow
        y'=\dfrac{x}{2-y}$. At $(\sqrt{3},1)$: $y'=\dfrac{\sqrt{3}}{1}
        =\sqrt{3}$. Tangent: $y-1=\sqrt{3}(x-\sqrt{3})$, i.e.\
        $y=\sqrt{3}x-2$.

  \item $\sec^{2}(xy)\cdot(y+xy')=1\Rightarrow y+xy'=\cos^{2}(xy)
        \Rightarrow y'=\dfrac{\cos^{2}(xy)-y}{x}$.

  \item $\ln y=\dfrac{\ln x}{x}\Rightarrow\dfrac{y'}{y}
        =\dfrac{(1/x)\cdot x-\ln x\cdot 1}{x^{2}}=\dfrac{1-\ln x}{x^{2}}$.
        So $y'=x^{1/x}\cdot\dfrac{1-\ln x}{x^{2}}$.

  \item $\dfrac{y'}{y}+y+xy'=0\Rightarrow y'\left(\dfrac1y+x\right)=-y
        \Rightarrow y'=-\dfrac{y^{2}}{1+xy}$.

  \item $3y^{2}y'+y'=2x\Rightarrow y'=\dfrac{2x}{3y^{2}+1}$.

  \item Differentiate $x^{2/3}+y^{2/3}=2$: $\tfrac23x^{-1/3}
        +\tfrac23y^{-1/3}y'=0\Rightarrow y'=-\dfrac{y^{1/3}}{x^{1/3}}
        =-\sqrt[3]{y/x}$. At $(1,1)$: $y'=-1$. Tangent: $y-1=-(x-1)$,
        i.e.\ $y=-x+2$.

  \item $\cos(x+y)(1+y')=y'\Rightarrow\cos(x+y)+\cos(x+y)y'=y'
        \Rightarrow\cos(x+y)=y'[1-\cos(x+y)]\Rightarrow
        y'=\dfrac{\cos(x+y)}{1-\cos(x+y)}$.

  \item $\ln y=2x+3\ln(x+1)-\tfrac12\ln(x^{2}+5)$. Differentiate:
        \[
          \frac{y'}{y}=2+\frac{3}{x+1}-\frac{x}{x^{2}+5} .
        \]
        Hence
        \[
          y'=\frac{\mathrm{e}^{2x}(x+1)^{3}}{\sqrt{x^{2}+5}}
             \left[2+\frac{3}{x+1}-\frac{x}{x^{2}+5}\right] .
        \]

  \item $2x+2yy'=0\Rightarrow y'=-x/y$. Differentiate again:
        \[
          y''=-\frac{y\cdot 1-x\cdot y'}{y^{2}}
             =-\frac{y-x(-x/y)}{y^{2}}
             =-\frac{y+x^{2}/y}{y^{2}}
             =-\frac{y^{2}+x^{2}}{y^{3}}=-\frac{25}{y^{3}},
        \]
        using $x^{2}+y^{2}=25$.

  \item Differentiate $x^{1/2}y^{1/2}=\bar U$:
        $\tfrac12x^{-1/2}y^{1/2}+x^{1/2}\cdot\tfrac12y^{-1/2}y'=0$.
        Multiply by $2x^{1/2}y^{1/2}$: $y+xy'=0\Rightarrow y'=-y/x$. So
        $\dd{y}{x}=-\dfrac{y}{x}$ (negative slope along the indifference
        curve, standard MRS).

  \item $2x+2yy'=2+4y'\Rightarrow y'(2y-4)=2-2x\Rightarrow
        y'=\dfrac{1-x}{y-2}$. Horizontal tangent: $y'=0$ requires $x=1$
        (and $y\neq 2$). Plug into the curve: $1+y^{2}=2+4y\Rightarrow
        y^{2}-4y-1=0\Rightarrow y=2\pm\sqrt{5}$. So horizontal tangents at
        $(1,2+\sqrt{5})$ and $(1,2-\sqrt{5})$.

  \item Differentiate $\arctan(y/x)=x$:
        $\dfrac{1}{1+(y/x)^{2}}\cdot\dfrac{y'x-y}{x^{2}}=1$, i.e.\
        $\dfrac{y'x-y}{x^{2}+y^{2}}=1$. Hence $y'x-y=x^{2}+y^{2}$, so
        $y'=\dfrac{x^{2}+y^{2}+y}{x}$.

  \item $\ln y=x\ln x\Rightarrow\dfrac{y'}{y}=\ln x+1\Rightarrow
        y'=x^{x}(\ln x+1)$.

  \item Check $(2,4)$ on $x^{3}+y^{3}=9xy$: $8+64=72$, $9\cdot 8=72$.
        Good. Differentiate: $3x^{2}+3y^{2}y'=9y+9xy'\Rightarrow
        y'(3y^{2}-9x)=9y-3x^{2}\Rightarrow
        y'=\dfrac{9y-3x^{2}}{3y^{2}-9x}=\dfrac{3y-x^{2}}{y^{2}-3x}$.
        At $(2,4)$: $y'=\dfrac{12-4}{16-6}=\dfrac{8}{10}=\dfrac{4}{5}$.
        Tangent: $y-4=\tfrac45(x-2)$, i.e.\ $4x-5y+12=0$.

  \item $2yy'=4\Rightarrow y'=2/y$. Differentiate again:
        $y''=-\dfrac{2y'}{y^{2}}=-\dfrac{2(2/y)}{y^{2}}=-\dfrac{4}{y^{3}}$.

  \item $\mathrm{e}^{y}+x\mathrm{e}^{y}y'+y'=1\Rightarrow
        y'(x\mathrm{e}^{y}+1)=1-\mathrm{e}^{y}\Rightarrow
        y'=\dfrac{1-\mathrm{e}^{y}}{1+x\mathrm{e}^{y}}$.

  \item $\dfrac{2x}{9}+\dfrac{2y}{4}y'=0\Rightarrow y'=-\dfrac{4x}{9y}$.
        Set $-\dfrac{4x}{9y}=-1\Rightarrow 4x=9y\Rightarrow y=\dfrac{4x}{9}$.
        Plug into the ellipse: $\dfrac{x^{2}}{9}+\dfrac{(4x/9)^{2}}{4}=1
        \Rightarrow\dfrac{x^{2}}{9}+\dfrac{16x^{2}}{4\cdot 81}=1
        \Rightarrow\dfrac{x^{2}}{9}+\dfrac{4x^{2}}{81}=1
        \Rightarrow\dfrac{9x^{2}+4x^{2}}{81}=1\Rightarrow
        13x^{2}=81\Rightarrow x=\pm 9/\sqrt{13}$. Corresponding
        $y=\pm 4/\sqrt{13}$ (same sign as $x$). Points:
        $\left(\dfrac{9}{\sqrt{13}},\dfrac{4}{\sqrt{13}}\right)$ and
        $\left(-\dfrac{9}{\sqrt{13}},-\dfrac{4}{\sqrt{13}}\right)$.

  \item $\ln y=\cos(x)\ln(\sin x)\Rightarrow\dfrac{y'}{y}
        =-\sin(x)\ln(\sin x)+\cos(x)\cdot\dfrac{\cos x}{\sin x}
        =-\sin(x)\ln(\sin x)+\dfrac{\cos^{2}x}{\sin x}$. So
        \[
          y'=(\sin x)^{\cos x}\left[-\sin(x)\ln(\sin x)
             +\frac{\cos^{2}x}{\sin x}\right] .
        \]

  \item $2xy+x^{2}y'-y^{2}-x\cdot 2yy'+2x+2yy'=0$. Collect $y'$:
        $y'(x^{2}-2xy+2y)=y^{2}-2xy-2x$. So
        $y'=\dfrac{y^{2}-2xy-2x}{x^{2}-2xy+2y}$.

  \item $\tfrac13K^{-2/3}L^{2/3}\dd{K}{L}\cdot L^{2/3}\ldots$
        carefully: differentiate $K^{1/3}L^{2/3}=\bar Q$ with respect to
        $L$, treating $K=K(L)$:
        \[
          \tfrac13K^{-2/3}K'(L)\cdot L^{2/3}+K^{1/3}\cdot\tfrac23L^{-1/3}=0 .
        \]
        Multiply by $3K^{2/3}L^{1/3}$: $LK'(L)+2K=0\Rightarrow
        K'(L)=-\dfrac{2K}{L}$. So $\dd{K}{L}=-\dfrac{2K}{L}$ (the MRTS).

  \item $\cos(y)y'=1\Rightarrow y'=\dfrac{1}{\cos y}=\sec(y)$.
        Differentiate again: $y''=\sec(y)\tan(y)\cdot y'
        =\sec^{2}(y)\tan(y)$.

  \item At $x=2$, $y=2$. Differentiate $xy=4$: $y+xy'=0\Rightarrow
        y'=-y/x$. At $(2,2)$: $y'=-1$. Tangent: $y-2=-(x-2)$, i.e.\
        $y=-x+4$.

  \item $\ln y=\tfrac12\ln(x+1)+\tfrac13\ln(x-2)-4\ln(x+5)$.
        Differentiate:
        \[
          \frac{y'}{y}=\frac{1}{2(x+1)}+\frac{1}{3(x-2)}-\frac{4}{x+5} .
        \]
        Hence
        \[
          y'=\frac{\sqrt{x+1}\sqrt[3]{x-2}}{(x+5)^{4}}
             \left[\frac{1}{2(x+1)}+\frac{1}{3(x-2)}-\frac{4}{x+5}\right] .
        \]

  \item $\mathrm{e}^{y}y'\ln(x)+\dfrac{\mathrm{e}^{y}}{x}+2x=y'\Rightarrow
        y'[\mathrm{e}^{y}\ln(x)-1]=-\dfrac{\mathrm{e}^{y}}{x}-2x\Rightarrow
        y'=\dfrac{\mathrm{e}^{y}/x+2x}{1-\mathrm{e}^{y}\ln(x)}$.

  \item Differentiate $x^{2}+y^{2}=r^{2}$: $2x+2yy'=0\Rightarrow
        y'=-x/y$. The radius from origin to $(x_{0},y_{0})$ has slope
        $y_{0}/x_{0}$. Product of slopes:
        $(-x_{0}/y_{0})\cdot(y_{0}/x_{0})=-1$, confirming
        perpendicularity.

  \item $\ln y=x^{2}\ln(7)\Rightarrow\dfrac{y'}{y}=2x\ln(7)\Rightarrow
        y'=2x\ln(7)\cdot 7^{x^{2}}$.

  \item Differentiate $x^{2}+xy+y^{2}=3$: $2x+y+xy'+2yy'=0\Rightarrow
        y'(x+2y)=-(2x+y)\Rightarrow y'=-\dfrac{2x+y}{x+2y}$.
        Differentiate again using the quotient rule:
        \[
          y''=-\frac{(2+y')(x+2y)-(2x+y)(1+2y')}{(x+2y)^{2}} .
        \]
        Expand the numerator: $(2+y')(x+2y)=2x+4y+xy'+2yy'$, and
        $(2x+y)(1+2y')=2x+y+4xy'+2yy'$. Subtract:
        $N:=(2+y')(x+2y)-(2x+y)(1+2y')=(2x+4y+xy'+2yy')
        -(2x+y+4xy'+2yy')=3y-3xy'$. So
        $y''=-\dfrac{3y-3xy'}{(x+2y)^{2}}=\dfrac{3(xy'-y)}{(x+2y)^{2}}$.
        Substitute $y'=-\dfrac{2x+y}{x+2y}$:
        \[
          xy'-y=-\frac{x(2x+y)}{x+2y}-y
               =-\frac{x(2x+y)+y(x+2y)}{x+2y}
        \]
        \[
          \phantom{xy'-y}=-\frac{2x^{2}+xy+xy+2y^{2}}{x+2y}
               =-\frac{2(x^{2}+xy+y^{2})}{x+2y},
        \]
        using $x^{2}+xy+y^{2}=3$. Hence
        \[
          y''=\frac{3}{(x+2y)^{2}}\cdot\left(-\frac{6}{x+2y}\right)
             =-\frac{18}{(x+2y)^{3}} .
        \]

  \item Differentiate $px+qy=m$ w.r.t.\ $x$: $p+qy'=0\Rightarrow
        y'=-p/q$, as required (the budget line has constant slope).

  \item Differentiate $x^{2}-xy+y^{2}=7$: $2x-y-xy'+2yy'=0\Rightarrow
        y'(2y-x)=y-2x\Rightarrow y'=\dfrac{y-2x}{2y-x}$. At $(3,2)$:
        $y'=\dfrac{2-6}{4-3}=-4$. Tangent: $y-2=-4(x-3)$, i.e.\
        $y=-4x+14$.

  \item Differentiate $\cos(xy)=x+y$: $-\sin(xy)\cdot(y+xy')=1+y'
        \Rightarrow -y\sin(xy)-x\sin(xy)y'=1+y'\Rightarrow
        -y'(x\sin(xy)+1)=1+y\sin(xy)\Rightarrow
        y'=-\dfrac{1+y\sin(xy)}{1+x\sin(xy)}$.

  \item First verify $(3,3)$: $27+27=54$ and $6\cdot 9=54$. Good.
        Differentiate $x^{3}+y^{3}=6xy$: $3x^{2}+3y^{2}y'=6y+6xy'
        \Rightarrow y'(3y^{2}-6x)=6y-3x^{2}\Rightarrow
        y'=\dfrac{6y-3x^{2}}{3y^{2}-6x}=\dfrac{2y-x^{2}}{y^{2}-2x}$.
        At $(3,3)$: $y'=\dfrac{6-9}{9-6}=-1$. For $y''$, differentiate
        $y'=\dfrac{2y-x^{2}}{y^{2}-2x}$:
        \[
          y''=\frac{(2y'-2x)(y^{2}-2x)-(2y-x^{2})(2yy'-2)}{(y^{2}-2x)^{2}} .
        \]
        At $(3,3)$ with $y'=-1$: numerator pieces:
        \[
          (2y'-2x)(y^{2}-2x)=(-2-6)(9-6)=(-8)(3)=-24,
        \]
        \[
          (2y-x^{2})(2yy'-2)=(6-9)(6\cdot(-1)-2)=(-3)(-8)=24 .
        \]
        So numerator $=-24-24=-48$ and denominator $=(9-6)^{2}=9$.
        Hence $y''(3)=-\dfrac{48}{9}=-\dfrac{16}{3}$.
\end{solutions}

\section{Additional Problems}

\begin{problems}
  \item Compute $\dd{y}{x}$ implicitly for $x^{2}+9y^{2}=36$.
  \item Find the equation of the tangent line to $x^{2}+y^{2}=25$ at
        $(3,4)$.
  \item Use logarithmic differentiation to find $y'$ for $y=x^{\cos x}$
        ($x>0$).
  \item Compute $\dd{y}{x}$ implicitly for $x^{2}y+xy^{2}=6$.
  \item Compute $\dd{^{2}y}{x^{2}}$ implicitly for $x^{2}-y^{2}=4$.
  \item Find the equation of the tangent line to $x^{3}+xy+y^{3}=11$ at
        $(1,2)$.
  \item Use logarithmic differentiation to find $y'$ for
        \[
          y=\frac{(x+1)^{2}\sqrt{x-3}}{(2x+5)^{4}}\qquad (x>3).
        \]
  \item Compute $\dd{y}{x}$ implicitly for $\mathrm{e}^{xy}=x+y$.
  \item For an indifference curve $U(x,y)=x^{1/3}y^{2/3}=\bar U$
        (constant), compute $\dd{y}{x}$, the marginal rate of
        substitution.
  \item Compute $\dd{y}{x}$ implicitly for $\sin(x+y)=xy$.
  \item Find the equation of the tangent line to $x^{2}-2xy+3y^{2}=9$ at
        $(1,2)$.
  \item Compute $\dd{^{2}y}{x^{2}}$ implicitly for $x^{2}+y^{2}=9$.
  \item Find all points on the curve $x^{2}+y^{2}=4x-2y$ where the tangent
        line is vertical.
  \item For the production isoquant $K^{2/5}L^{3/5}=\bar Q$ (constant),
        compute $\dd{K}{L}$, the marginal rate of technical
        substitution.
  \item Use logarithmic differentiation to find $y'$ for
        $y=(x^{2}+1)^{\sin x}$ ($x\in\R$).

  \item The curve $k$ below is the set of points satisfying
        \[
          x^{3}-3x^{2}+y^{2}-2y=-1 .
        \]
        Verify that $(0,1)$, $(3,1)$, $(-1,3)$ and $(-1,-1)$ all lie on
        $k$.

\begin{gallery}
\begin{tikzpicture}
  \begin{axis}[geoaxis, xmin=-1.6, xmax=3.6, ymin=-1.9, ymax=3.9,
      xtick={-1,1,2,3}, ytick={-1,1,2,3}, title={the curve $k$}]
    \addplot[domain=-1.2:3, samples=241]{1+abs(x)*sqrt(3-x)};
    \addplot[domain=-1.2:3, samples=241]{1-abs(x)*sqrt(3-x)};
  \end{axis}
\end{tikzpicture}\hfill
\begin{tikzpicture}
  \begin{axis}[geoaxis, xmin=-1.4, xmax=1.4, ymin=-0.9, ymax=0.9,
      xtick={-1,1}, ytick={-0.5,0.5}, title={the curve $\ell$}]
    \addplot[domain=-1:1, samples=241]{x*sqrt(1-x^2)};
    \addplot[domain=-1:1, samples=241]{-x*sqrt(1-x^2)};
  \end{axis}
\end{tikzpicture}
\end{gallery}

  \item Name every point of $k$ at which the hypotheses of the implicit
        function theorem fail for writing $y$ as a function of $x$, and say
        which graphical feature identifies each one.
  \item Confirm the previous answer by computation: solve $F_{y}=0$
        together with $F=0$.
  \item Compute the slope of the tangent to $k$ at $(-1,3)$.
  \item Compute the slope of the tangent to $k$ at $(-1,-1)$, and explain
        why it is the negative of the previous answer.
  \item Find every point of $k$ at which the tangent is horizontal.
  \item Near $(3,1)$ no function $y(x)$ describes $k$. Show that
        $x=x(y)$ does exist there, and compute $\dd{x}{y}$ at that point.
  \item The curve $\ell$ above is $y^{2}=x^{2}\bigl(1-x^{2}\bigr)$. Write it
        as $F(x,y)=0$ and give $F_{x}$ and $F_{y}$.
  \item Name every point of $\ell$ at which the implicit function theorem
        fails, and the graphical feature of each.
  \item Compute the slope of the tangent to $\ell$ at
        $\bigl(\tfrac12,\tfrac{\sqrt3}{4}\bigr)$.
  \item Compute the slope of the tangent to $\ell$ at
        $\bigl(\tfrac12,-\tfrac{\sqrt3}{4}\bigr)$.
  \item Find the four points of $\ell$ at which the tangent is horizontal.
  \item At the origin two branches of $\ell$ cross. Separate them by solving
        for $y$, and compute the slope of each branch there.
  \item Explain in one or two sentences why $-F_{x}/F_{y}$ cannot deliver
        those two slopes, and what has to be done instead.
\end{problems}

\section{Additional Solutions}

\begin{solutions}
  \item $2x+18yy'=0\Rightarrow y'=-\dfrac{x}{9y}$.

  \item Check $(3,4)$: $9+16=25$. From $2x+2yy'=0$: $y'=-x/y$. At
        $(3,4)$: $y'=-3/4$. Tangent: $y-4=-\tfrac34(x-3)$, i.e.\
        $3x+4y=25$.

  \item $\ln y=\cos(x)\ln(x)\Rightarrow\dfrac{y'}{y}
        =-\sin(x)\ln(x)+\dfrac{\cos x}{x}$. So
        \[
          y'=x^{\cos x}\left[-\sin(x)\ln(x)+\frac{\cos x}{x}\right] .
        \]

  \item Differentiate $x^{2}y+xy^{2}=6$: $2xy+x^{2}y'+y^{2}+2xyy'=0
        \Rightarrow y'(x^{2}+2xy)=-(2xy+y^{2})
        \Rightarrow y'=-\dfrac{y(2x+y)}{x(x+2y)}$.

  \item $2x-2yy'=0\Rightarrow y'=x/y$. Differentiate again:
        \[
          y''=\frac{y-xy'}{y^{2}}=\frac{y-x(x/y)}{y^{2}}
             =\frac{y^{2}-x^{2}}{y^{3}}=\frac{-4}{y^{3}},
        \]
        using $x^{2}-y^{2}=4\Rightarrow y^{2}-x^{2}=-4$.

  \item Check $(1,2)$: $1+2+8=11$. Differentiate:
        $3x^{2}+y+xy'+3y^{2}y'=0\Rightarrow y'(x+3y^{2})=-(3x^{2}+y)
        \Rightarrow y'=-\dfrac{3x^{2}+y}{x+3y^{2}}$. At $(1,2)$:
        $y'=-\dfrac{3+2}{1+12}=-\dfrac{5}{13}$. Tangent:
        $y-2=-\tfrac{5}{13}(x-1)$, i.e.\ $5x+13y=31$.

  \item $\ln y=2\ln(x+1)+\tfrac12\ln(x-3)-4\ln(2x+5)$. Then
        \[
          \frac{y'}{y}=\frac{2}{x+1}+\frac{1}{2(x-3)}-\frac{8}{2x+5} .
        \]
        Hence
        \[
          y'=\frac{(x+1)^{2}\sqrt{x-3}}{(2x+5)^{4}}
             \left[\frac{2}{x+1}+\frac{1}{2(x-3)}-\frac{8}{2x+5}\right] .
        \]

  \item Differentiate $\mathrm{e}^{xy}=x+y$:
        $\mathrm{e}^{xy}(y+xy')=1+y'\Rightarrow
        y'(x\mathrm{e}^{xy}-1)=1-y\mathrm{e}^{xy}\Rightarrow
        y'=\dfrac{1-y\mathrm{e}^{xy}}{x\mathrm{e}^{xy}-1}$.

  \item Differentiate $x^{1/3}y^{2/3}=\bar U$:
        $\tfrac13x^{-2/3}y^{2/3}+x^{1/3}\cdot\tfrac23y^{-1/3}y'=0$.
        Multiply by $3x^{2/3}y^{1/3}$: $y+2xy'=0\Rightarrow y'=-y/(2x)$.
        So $\dd{y}{x}=-\dfrac{y}{2x}$ (the MRS).

  \item Differentiate $\sin(x+y)=xy$: $\cos(x+y)(1+y')=y+xy'
        \Rightarrow y'(\cos(x+y)-x)=y-\cos(x+y)\Rightarrow
        y'=\dfrac{y-\cos(x+y)}{\cos(x+y)-x}$.

  \item Check $(1,2)$: $1-4+12=9$. Good. Differentiate:
        $2x-2y-2xy'+6yy'=0\Rightarrow y'(6y-2x)=2y-2x\Rightarrow
        y'=\dfrac{y-x}{3y-x}$. At $(1,2)$: $y'=\dfrac{1}{5}$. Tangent:
        $y-2=\tfrac15(x-1)$, i.e.\ $x-5y+9=0$.

  \item $2x+2yy'=0\Rightarrow y'=-x/y$. Differentiate again:
        \[
          y''=-\frac{y-xy'}{y^{2}}=-\frac{y-x(-x/y)}{y^{2}}
             =-\frac{y+x^{2}/y}{y^{2}}
             =-\frac{y^{2}+x^{2}}{y^{3}}=-\frac{9}{y^{3}},
        \]
        using $x^{2}+y^{2}=9$.

  \item $2x+2yy'=4-2y'\Rightarrow y'(2y+2)=4-2x\Rightarrow
        y'=\dfrac{2-x}{y+1}$. Vertical tangent: $y'$ undefined when
        $y+1=0\Rightarrow y=-1$ (and $x\neq 2$). Plug into the curve:
        $x^{2}+1=4x+2\Rightarrow x^{2}-4x-1=0\Rightarrow
        x=2\pm\sqrt{5}$. Vertical tangents at $(2+\sqrt{5},-1)$ and
        $(2-\sqrt{5},-1)$.

  \item Differentiate $K^{2/5}L^{3/5}=\bar Q$ w.r.t.\ $L$, treating
        $K=K(L)$:
        \[
          \tfrac25K^{-3/5}K'(L)\cdot L^{3/5}+K^{2/5}\cdot\tfrac35L^{-2/5}=0 .
        \]
        Multiply by $5K^{3/5}L^{2/5}$: $2LK'(L)+3K=0\Rightarrow
        K'(L)=-\dfrac{3K}{2L}$. So $\dd{K}{L}=-\dfrac{3K}{2L}$ (the
        MRTS).

  \item $\ln y=\sin(x)\ln(x^{2}+1)\Rightarrow\dfrac{y'}{y}
        =\cos(x)\ln(x^{2}+1)+\sin(x)\cdot\dfrac{2x}{x^{2}+1}$. So
        \[
          y'=(x^{2}+1)^{\sin x}\left[\cos(x)\ln(x^{2}+1)
             +\frac{2x\sin x}{x^{2}+1}\right] .
        \]

  \item With $F(x,y)=x^{3}-3x^{2}+y^{2}-2y+1$:
        $F(0,1)=0-0+1-2+1=0$; \;$F(3,1)=27-27+1-2+1=0$;
        \;$F(-1,3)=-1-3+9-6+1=0$; \;$F(-1,-1)=-1-3+1+2+1=0$.
        All four are on $k$.
  \item Graphically: at $(0,1)$ the curve \textbf{crosses itself}, and at
        $(3,1)$ it has a \textbf{vertical tangent}. At every other point of
        $k$ a local $y(x)$ exists.
  \item $F_{y}=2y-2=0\iff y=1$. Substituting $y=1$ into $F=0$ gives
        $x^{3}-3x^{2}+1-2+1=x^{2}(x-3)=0$, so $x=0$ (a double root) or
        $x=3$: the points $(0,1)$ and $(3,1)$. At $(3,1)$,
        $F_{x}=3x^{2}-6x=9\neq0$, so this is the vertical-tangent case; at
        $(0,1)$, $F_{x}=0$ as well, the singular case.
  \item $F_{x}=3x^{2}-6x$, $F_{y}=2y-2$, so
        $y'=-\dfrac{3x^{2}-6x}{2y-2}$. At $(-1,3)$:
        $y'=-\dfrac{3+6}{6-2}=-\dfrac94$.
  \item At $(-1,-1)$: $y'=-\dfrac{3+6}{-2-2}=\dfrac94$. The equation
        contains $y$ only through $(y-1)^{2}$, so $k$ is symmetric about the
        line $y=1$; reflecting a curve in a horizontal line negates every
        slope, and $(-1,3)$ and $(-1,-1)$ are mirror images.
  \item Horizontal tangent needs $F_{x}=0$ and $F_{y}\neq0$:
        $3x^{2}-6x=3x(x-2)=0$ gives $x=0$ or $x=2$. The value $x=0$ forces
        $y=1$, where $F_{y}=0$ too --- that is the singular point, not a
        horizontal tangent. For $x=2$: $(y-1)^{2}=x^{2}(3-x)=4$, so $y=3$ or
        $y=-1$. The tangent is horizontal at $(2,3)$ and $(2,-1)$.
  \item At $(3,1)$ we have $F_{x}=9\neq0$, so the theorem applies with the
        roles of the variables exchanged and gives a function $x=x(y)$.
        Differentiating $F\bigl(x(y),y\bigr)=0$ with respect to $y$:
        $F_{x}\,x'(y)+F_{y}=0$, hence
        $\dd{x}{y}=-\dfrac{F_{y}}{F_{x}}=-\dfrac{0}{9}=0$ --- the curve is
        momentarily vertical, which is exactly a horizontal tangent for
        $x(y)$.
  \item $F(x,y)=y^{2}-x^{2}+x^{4}$, with $F_{x}=-2x+4x^{3}=2x(2x^{2}-1)$
        and $F_{y}=2y$.
  \item $F_{y}=2y=0\iff y=0$, and then $x^{2}(1-x^{2})=0$ gives
        $x=0,\pm1$. At $(\pm1,0)$: $F_{x}=\pm2\neq0$, \textbf{vertical
        tangents} --- the two tips of the figure eight. At $(0,0)$:
        $F_{x}=0$ too, the \textbf{self-intersection} in the middle.
  \item $y'=-\dfrac{2x(2x^{2}-1)}{2y}$. At
        $\bigl(\tfrac12,\tfrac{\sqrt3}{4}\bigr)$ the numerator is
        $2\cdot\tfrac12\bigl(\tfrac12-1\bigr)=-\tfrac12$ and the denominator
        is $\tfrac{\sqrt3}{2}$, so
        $y'=\dfrac{1/2}{\sqrt3/2}=\dfrac{1}{\sqrt3}=\dfrac{\sqrt3}{3}$.
  \item Only the sign of $y$ changes, so
        $y'=-\dfrac{\sqrt3}{3}$.
  \item $F_{x}=2x(2x^{2}-1)=0$ with $F_{y}\neq0$ gives $x=\pm
        \tfrac{1}{\sqrt2}$ (the root $x=0$ is the singular point). Then
        $y^{2}=x^{2}(1-x^{2})=\tfrac12\cdot\tfrac12=\tfrac14$, so
        $y=\pm\tfrac12$: the four points
        $\bigl(\pm\tfrac{1}{\sqrt2},\pm\tfrac12\bigr)$.
  \item Solving gives $y=x\sqrt{1-x^{2}}$ and $y=-x\sqrt{1-x^{2}}$. For the
        first,
        $y'=\sqrt{1-x^{2}}-\dfrac{x^{2}}{\sqrt{1-x^{2}}}$, which at $x=0$
        equals $1$; the second branch gives $-1$. The two branches cross the
        origin at $+45^{\circ}$ and $-45^{\circ}$.
  \item At the origin both partials vanish, so $-F_{x}/F_{y}$ is $0/0$: it
        carries no information, and a single number could not describe two
        different slopes anyway. The branches have to be separated first ---
        here by solving for $y$ --- and each differentiated on its own.
\end{solutions}

\chapter{Applications: Tangent, Linearization, Marginals, Elasticity, Optimization}

The derivative is the bridge between calculus and economics. Geometrically it
gives the slope of the tangent (and hence of the normal); analytically it
produces the linearisation that turns nonlinear functions into
easy-to-evaluate affine ones; in economics it is the marginal quantity
(cost, revenue, utility) and, after rescaling, the elasticity. The first-
and second-order conditions $f'=0$, $f''<0$ (or $>0$) drive one-variable
optimisation.

\section{Core Rules}

\begin{rulebox}{Tangent line}
At $(a,f(a))$: $y = f(a) + f'(a)(x-a)$.
\end{rulebox}

\begin{rulebox}{Normal line}
At $(a,f(a))$: perpendicular to the tangent, slope $-1/f'(a)$ when
$f'(a)\neq 0$:
\[
  y = f(a) - \frac{1}{f'(a)}(x-a).
\]
\end{rulebox}

\begin{rulebox}{Linear approximation (linearisation)}
For $x$ near $a$,
\[
  f(x) \approx L(x) := f(a) + f'(a)(x-a), \qquad \Delta f \approx f'(a)\,\Delta x .
\]
\end{rulebox}

\begin{rulebox}{Marginal quantities}
If $C(Q)$, $R(Q)$, $\pi(Q)=R(Q)-C(Q)$ are total cost, revenue and profit,
and $U(x)$ is utility, then
\[
  MC(Q)=C'(Q), \quad MR(Q)=R'(Q), \quad M\pi(Q)=\pi'(Q), \quad MU(x)=U'(x).
\]
Profit-maximising output satisfies $MR(Q^*)=MC(Q^*)$ (FOC) and
$\pi''(Q^*)<0$ (SOC).
\end{rulebox}

\begin{rulebox}{Elasticity}
The (point) elasticity of $f$ with respect to $x$ is
\[
  \varepsilon = \frac{x}{f(x)}f'(x) = \frac{\mathrm{d}f/f}{\mathrm{d}x/x}
    = \frac{\mathrm{d}\ln f}{\mathrm{d}\ln x} .
\]
For demand $q=q(p)$: $\displaystyle \varepsilon_d = \frac{p}{q}q'(p)$
(often negative). Demand is \emph{elastic} when $|\varepsilon_d|>1$,
\emph{inelastic} when $|\varepsilon_d|<1$, \emph{unit elastic} when
$|\varepsilon_d|=1$.
\end{rulebox}

\begin{rulebox}{1D unconstrained optimisation}
\begin{enumerate}[leftmargin=1.6em,itemsep=0.3ex]
  \item Solve $f'(x)=0$ on the relevant domain to find critical points.
  \item Classify with the second-derivative test: $f''(x^*)<0 \Rightarrow$
        local max, $f''(x^*)>0 \Rightarrow$ local min.
  \item Compare with boundary values and endpoint behaviour to identify
        global extrema.
\end{enumerate}
\end{rulebox}

\section{Worked Example}

\begin{worked}{Demand is $q=100-2p$. Find the price elasticity at $p=20$ and at the revenue-maximising price, and approximate $\sqrt{101}$ by linearisation}
\wstep{Step 1. Elasticity is a derivative times a ratio --- not a derivative.}
\[
  \varepsilon=\frac{\mathrm{d}q}{\mathrm{d}p}\cdot\frac{p}{q}
   =\frac{\%\ \text{change in }q}{\%\ \text{change in }p} .
\]
Here $\mathrm{d}q/\mathrm{d}p=-2$ everywhere, but the elasticity is not
constant, because $p/q$ moves along the curve.
\wtrap{$\varepsilon\neq\mathrm{d}q/\mathrm{d}p$. The slope has units
(quantity per euro) and changes if you re-denominate the price; the
elasticity is a pure number and does not. That invariance is the main
reason economists use it.}

\wstep{Step 2. Evaluate at $p=20$.}
$q=100-40=60$, so
\[
  \varepsilon=-2\cdot\frac{20}{60}=-\tfrac23\approx-0.67 .
\]
$|\varepsilon|<1$: demand is \textbf{inelastic} here. A $1\%$ price rise
costs only $0.67\%$ of quantity, so revenue \emph{rises} --- the firm is
pricing below the revenue maximum.
\wtrap{The sign is negative for any downward-sloping demand, and dropping it
is a real error even though the classification uses $|\varepsilon|$. Report
$-0.67$ and then say ``inelastic''; do not report $0.67$.}

\wstep{Step 3. Maximise revenue directly.}
\[
  R(p)=pq=p(100-2p)=100p-2p^{2},
  \qquad
  R'(p)=100-4p=0
  \quad\Longrightarrow\quad
  p^{*}=25 .
\]
$R''=-4<0$, so this is a maximum. Then $q^{*}=50$ and $R^{*}=1250$.

\wstep{Step 4. Check the elasticity there --- and find the general rule.}
\[
  \varepsilon(25)=-2\cdot\frac{25}{50}=-1 .
\]
Unit elastic, exactly at the revenue peak. That is no coincidence:
\[
  R=pq
  \;\Longrightarrow\;
  \frac{\mathrm{d}R}{\mathrm{d}p}=q+p\frac{\mathrm{d}q}{\mathrm{d}p}
  =q\Bigl(1+\varepsilon\Bigr),
\]
which vanishes exactly when $\varepsilon=-1$. Revenue rises where demand is
inelastic, falls where it is elastic, and peaks in between.

\wstep{Step 5. Marginal revenue, and the same fact in quantity terms.}
Inverting demand, $p=50-q/2$, so $R(q)=50q-q^{2}/2$ and
\[
  MR(q)=50-q .
\]
$MR=0$ at $q=50$ --- the same point. Note $MR$ falls twice as fast as the
demand curve $p=50-q/2$: the standard ``$MR$ has twice the slope'' result,
here visible as $-1$ against $-\tfrac12$.
\wtrap{Marginal revenue is $\mathrm{d}R/\mathrm{d}q$, not $R/q$ (that is
\emph{average} revenue, which equals the price). Confusing marginal with
average is the most consequential error in applied differentiation, and it
recurs for cost: $MC=C'(q)$, while $AC=C(q)/q$.}

\wstep{Step 6. Linearisation: the tangent line as a calculator.}
For $f(x)=\sqrt{x}$ near $a=100$, where the value is known exactly:
\[
  f(x)\approx f(a)+f'(a)(x-a),
  \qquad
  f'(x)=\frac{1}{2\sqrt x},\quad f'(100)=\frac{1}{20} .
\]
\[
  \sqrt{101}\approx10+\frac{1}{20}(1)=10.05 .
\]
The true value is $10.0499$, so the error is $10^{-4}$.
\wtrap{Expand about the nearest point whose value you actually know. Using
$a=101$ makes the formula exact and useless; using $a=64$ would be legal but
far away, and the error grows with $(x-a)^{2}$. Choose $a$ first, and choose
it well.}

\wstep{Step 7. The same tangent, three names.}
\[
  \underbrace{y=f(a)+f'(a)(x-a)}_{\text{tangent line}}
  \;=\;
  \underbrace{f(a)+f'(a)\,\Delta x}_{\text{linearisation}}
  \;=\;
  \underbrace{f(a)+\text{marginal}\times\text{one more unit}}_{\text{economics}} .
\]
``Marginal cost is the cost of one more unit'' is a linear approximation:
strictly $C'(q)$ is a rate, and $C(q+1)-C(q)\approx C'(q)$ with an error of
order $C''/2$. The approximation is good precisely when the curvature is
small.
\end{worked}

\section{Practice Problems}

\begin{problems}
  \item Find the equation of the tangent to $f(x)=x^3-2x+5$ at $x=1$.
  \item Use the linearisation of $f(x)=\sqrt{x}$ at $a=25$ to estimate
        $\sqrt{25.4}$.
  \item For $C(Q)=0.2Q^2+5Q+80$, find $MC(Q)$ and $MC(15)$.
  \item For demand $q(p)=200-5p$, compute $\varepsilon_d(p)$ and evaluate at
        $p=10$.
  \item Maximise profit $\pi(Q)=R(Q)-C(Q)$ where $R(Q)=60Q-Q^2$ and
        $C(Q)=Q^2+10Q+25$. Verify the SOC.
  \item Find the tangent to $f(x)=1/x$ at $x=2$.
  \item Use the linearisation of $f(x)=x^{1/4}$ at $a=16$ to estimate
        $16.3^{1/4}$.
  \item Inverse demand $p(Q)=80-0.4Q$. Find $R(Q)$, $MR(Q)$, and $MR(50)$.
  \item For the Cobb--Douglas production $Y=AK^\alpha L^\beta$ with $L$
        fixed, compute the output elasticity $\varepsilon_{Y,K}$.
  \item Minimise the average cost $AC(Q)=C(Q)/Q$ for
        $C(Q)=2Q^2+8Q+200$, over $Q>0$.
  \item Find the tangent to $f(x)=e^{2x}$ at $x=0$.
  \item Estimate $\ln(0.98)$ using the linearisation of $f(x)=\ln(x)$ at
        $a=1$.
  \item Using $C$ from problem 3 and $R$ from problem 8, find the
        profit-maximising output $Q^*$ and verify SOC.
  \item For demand $q(p)=80e^{-0.05p}$, compute $\varepsilon_d(p)$ and at
        $p=20$.
  \item A monopolist faces $p(Q)=150-Q$ and $C(Q)=Q^2+30Q+100$. Find $Q^*$,
        $p^*$, and $\pi^*$.
  \item Find the normal line to $f(x)=\sqrt{x}$ at $x=9$.
  \item Estimate $e^{-0.2}$ using the linearisation of $f(x)=e^x$ at $a=0$.
  \item A monopolist faces $p(Q)=200-3Q$ and $C(Q)=2Q^2+10Q+50$. Find
        $Q^*$, $p^*$, and $\pi^*$.
  \item For $q(p)=200-5p$, find the price at which demand is unit elastic.
  \item A firm sells at the fixed price $p=24$ per unit; cost is
        $C(Q)=Q^2+8Q+10$. Find the profit-maximising output and verify SOC.
  \item Find the normal line to $f(x)=\ln(x)$ at $x=e$.
  \item A firm's cost is $C(Q)=0.5Q^2+8Q+200$. Use linear approximation at
        $Q=20$ to estimate $C(21.5)$.
  \item Marginal cost is $MC(Q)=0.06Q^2-2.4Q+30$. Find the production level
        that minimises $MC$ and verify the SOC.
  \item Demand $q(p)=\dfrac{300}{p^2+4}$. Compute $\varepsilon_d(p)$ and
        evaluate at $p=2$.
  \item Maximise utility $U(x)=4\ln(x)-0.5x$ for $x>0$.
  \item Find the tangent to $f(x)=\cos(x)$ at $x=\pi/3$.
  \item Use linearisation of $f(x)=(1+x)^{10}$ at $a=0$ to estimate
        $(1.02)^{10}$.
  \item For $C(Q)=Q^3-9Q^2+30Q+40$, compute $MC(Q)$ and locate the $Q$
        where $MC$ attains its minimum.
  \item For demand $q(p)=Ap^{-\alpha}$ with $A,\alpha>0$, show that the
        (point) elasticity is the constant $-\alpha$.
  \item A producer faces inverse demand $p(Q)=100-2Q$ and constant
        marginal cost $MC=20$ (so $C(Q)=20Q+F$ with fixed cost $F$). Find
        $Q^*$, $p^*$, and verify SOC.
  \item Find the normal line to $f(x)=x^2$ at $x=1$.
  \item Use linearisation of $f(x)=\tan(x)$ at $a=0$ to estimate
        $\tan(0.05)$.
  \item For the cubic cost $C(Q)=Q^3-6Q^2+15Q+20$, find marginal cost and
        average cost, and the $Q>0$ where $AC$ is minimised. Verify SOC.
  \item For demand $q(p)=100-2p^2$ (defined on $0\le p\le\sqrt{50}$),
        compute $\varepsilon_d(p)$ at $p=4$ and classify.
  \item Maximise revenue $R(Q)=Q\cdot p(Q)$ when $p(Q)=120-2Q$ and verify
        SOC.
  \item Find the tangent and the normal to $f(x)=x^2-3x+2$ at $x=2$.
  \item For a firm with $C(Q)=4Q^2+20Q+100$, what is the change in cost
        (approximately) when output increases from 10 to 10.5?
  \item For $C(Q)=Q^2+10Q+9$ and $R(Q)=30Q$, find the profit-maximising
        output and the break-even points (where $\pi=0$).
  \item For demand $q(p)=400-p^2$, find the price at which
        $|\varepsilon_d|=1$.
  \item Find positive $x$ that maximises $f(x)=x^2e^{-x}$ on
        $(0,\infty)$. Verify SOC.
  \item Find the tangent to $f(x)=\arctan(x)$ at $x=1$.
  \item Estimate $\sqrt[3]{27.9}$ using the linearisation of
        $f(x)=x^{1/3}$ at $a=27$.
  \item For utility $U(x)=\sqrt{x}$ and consumption $x=100$, compute the
        marginal utility $MU(x)$ and its elasticity $\varepsilon_{U,x}$.
        Comment briefly on signs.
  \item A monopolist faces demand $q(p)=120-3p$ and constant unit cost
        $c=10$. Reformulate via the inverse demand to find $Q^*$, $p^*$,
        $\pi^*$. Verify SOC.
  \item Minimise $f(x)=x+\dfrac{16}{x}$ on $(0,\infty)$. Verify SOC.
  \item Find both the tangent and the normal to $f(x)=\sin(x)$ at
        $x=\pi/2$.
  \item For revenue $R(Q)=p(Q)\cdot Q$ with $p(Q)=100-Q$, express $MR$ in
        terms of $p$ and the (price-)elasticity of demand: verify the
        identity $MR = p\left(1+\dfrac{1}{\varepsilon_d}\right)$ at
        $Q=20$.
  \item Use linearisation of $f(x)=\ln(1+x)$ at $a=0$ to estimate
        $\ln(1.05)$, and compute the linearisation error.
  \item A firm has revenue $R(Q)=20Q-Q^2$ and cost
        $C(Q)=2Q+\frac12Q^2+10$, with the constraint that $Q\in[0,8]$.
        Find the profit-maximising $Q$ on this closed interval, comparing
        the interior critical point with the endpoints.
  \item For demand $q(p)=50/p$, compute $\varepsilon_d$ for all $p>0$ and
        explain what it implies for total revenue $R(p)=p\cdot q(p)$.
  \item Maximise $\pi(Q)=12Q-Q^2-\dfrac{Q^3}{3}$ for $Q\ge0$. (Take the
        unique non-negative critical point and verify SOC.)
  \item For inverse demand $p(Q)=50-Q$ and cost $C(Q)=10Q+Q^2$, compute
        $\pi(Q)$, $Q^*$, $p^*$, $\pi^*$ and also $\varepsilon_d$ at $Q^*$.
  \item A firm produces according to $Y=\sqrt{L}$ (a strictly concave
        one-input production function). The wage is $w=2$ and the output
        price is $p=8$. Choose $L\ge0$ to maximise profit
        $\pi(L)=p\cdot Y(L)-wL$. Verify SOC.
\end{problems}

\section{Solutions}

\begin{solutions}
  \item $f'(x)=3x^2-2$. $f(1)=4$, $f'(1)=1$. Tangent:
        $y=4+1\cdot(x-1)=x+3$.

  \item $f(25)=5$, $f'(x)=1/(2\sqrt{x})$, $f'(25)=1/10=0.1$. So
        $\sqrt{25.4}\approx 5+0.1\cdot0.4=5.04$. (True: $5.0399\ldots$)

  \item $MC(Q)=C'(Q)=0.4Q+5$. $MC(15)=6+5=11$.

  \item $q'(p)=-5$.
        $\varepsilon_d(p)=\dfrac{p}{200-5p}\cdot(-5)=\dfrac{-5p}{200-5p}
        =\dfrac{-p}{40-p}$. At $p=10$: $\varepsilon_d=-10/30=-1/3$, so
        $|\varepsilon_d|<1$, inelastic.

  \item $\pi(Q)=60Q-Q^2-Q^2-10Q-25=50Q-2Q^2-25$. FOC:
        $\pi'(Q)=50-4Q=0\Rightarrow Q^*=12.5$. SOC: $\pi''(Q)=-4<0$, so a
        maximum. $\pi(12.5)=625-312.5-25=287.5$.

  \item $f'(x)=-1/x^2$. $f(2)=1/2$, $f'(2)=-1/4$. Tangent:
        $y=\frac12-\frac14(x-2)=-\frac{x}{4}+1$.

  \item $f'(x)=\frac14x^{-3/4}$. $f(16)=2$,
        $f'(16)=\frac14\cdot16^{-3/4}=\frac14\cdot\frac18=\frac1{32}$. So
        $16.3^{1/4}\approx2+\frac{0.3}{32}=2+0.009375=2.009375$.

  \item $R(Q)=p(Q)Q=80Q-0.4Q^2$. $MR(Q)=80-0.8Q$. $MR(50)=80-40=40$.

  \item $\dfrac{\partial Y}{\partial K}=\alpha AK^{\alpha-1}L^\beta$. So
        \[
          \varepsilon_{Y,K}=\frac{K}{Y}\cdot\alpha AK^{\alpha-1}L^\beta
            =\frac{\alpha AK^\alpha L^\beta}{AK^\alpha L^\beta}=\alpha .
        \]
        The Cobb--Douglas exponent equals the elasticity.

  \item $AC(Q)=2Q+8+\dfrac{200}{Q}$.
        $AC'(Q)=2-\dfrac{200}{Q^2}=0\Rightarrow Q^2=100\Rightarrow Q^*=10$
        (positive root). $AC''(Q)=\dfrac{400}{Q^3}>0$ at $Q=10$, so a
        minimum. $AC(10)=20+8+20=48$.

  \item $f'(x)=2e^{2x}$. $f(0)=1$, $f'(0)=2$. Tangent: $y=1+2x$.

  \item $f(1)=0$, $f'(1)=1$. $\ln(0.98)\approx0+1\cdot(-0.02)=-0.02$.
        (True: $-0.02020\ldots$)

  \item Set $MR(Q)=MC(Q)$: $80-0.8Q=0.4Q+5\Rightarrow75=1.2Q\Rightarrow
        Q^*=62.5$. SOC: $\pi''(Q)=R''(Q)-C''(Q)=-0.8-0.4=-1.2<0$, so this
        is a maximum.

  \item $q'(p)=80\cdot(-0.05)e^{-0.05p}=-4e^{-0.05p}$. So
        \[
          \varepsilon_d(p)=\frac{p}{80e^{-0.05p}}\cdot(-4e^{-0.05p})
            =-0.05p .
        \]
        At $p=20$: $\varepsilon_d=-1$, exactly unit elastic.

  \item $R(Q)=p(Q)Q=150Q-Q^2$.
        $\pi(Q)=150Q-Q^2-Q^2-30Q-100=120Q-2Q^2-100$. FOC:
        $\pi'(Q)=120-4Q=0\Rightarrow Q^*=30$. SOC: $\pi''(Q)=-4<0$.
        $p^*=150-30=120$. $\pi^*=120\cdot30-2\cdot900-100
        =3600-1800-100=1700$.

  \item $f'(x)=\dfrac{1}{2\sqrt{x}}$. $f(9)=3$, $f'(9)=1/6$. Normal slope
        $=-1/f'(9)=-6$. Normal: $y=3-6(x-9)=-6x+57$.

  \item $f(0)=1$, $f'(0)=1$. $e^{-0.2}\approx1+1\cdot(-0.2)=0.8$. (True:
        $0.8187\ldots$)

  \item $R(Q)=p(Q)Q=200Q-3Q^2$.
        $\pi(Q)=200Q-3Q^2-2Q^2-10Q-50=190Q-5Q^2-50$. FOC:
        $\pi'(Q)=190-10Q=0\Rightarrow Q^*=19$. SOC: $\pi''(Q)=-10<0$, max.
        $p^*=200-57=143$. $\pi^*=190\cdot19-5\cdot361-50
        =3610-1805-50=1755$.

  \item Using problem 4: $|\varepsilon_d|=1 \Leftrightarrow
        \dfrac{p}{40-p}=1\Rightarrow p=40-p\Rightarrow p=20$. (At $p=20$,
        demand $q=200-100=100$.)

  \item $\pi(Q)=24Q-Q^2-8Q-10=16Q-Q^2-10$. FOC:
        $\pi'(Q)=16-2Q=0\Rightarrow Q^*=8$. SOC: $\pi''(Q)=-2<0$.
        $\pi(8)=128-64-10=54$.

  \item $f'(x)=1/x$. $f(e)=1$, $f'(e)=1/e$. Normal slope $=-e$. Normal:
        $y=1-e(x-e)=1-ex+e^2$.

  \item $C'(Q)=Q+8$. $C(20)=200+160+200=560$, $C'(20)=28$. So
        $C(21.5)\approx560+28\cdot1.5=602$. (Exact:
        $0.5\cdot462.25+172+200=603.125$; error $\approx1.1$.)

  \item $MC'(Q)=0.12Q-2.4=0\Rightarrow Q=20$. $MC''(Q)=0.12>0$, so $Q=20$
        minimises $MC$. $MC(20)=0.06\cdot400-48+30=24-48+30=6$.

  \item $q(p)=300(p^2+4)^{-1}$,
        $q'(p)=-\dfrac{600p}{(p^2+4)^2}$. Then
        \[
          \varepsilon_d(p)=\frac{p}{q}\cdot q'(p)
            =\frac{p(p^2+4)}{300}\cdot\left(-\frac{600p}{(p^2+4)^2}\right)
            =\frac{-2p^2}{p^2+4} .
        \]
        At $p=2$: $\varepsilon_d=-8/8=-1$ (unit elastic).

  \item $U'(x)=4/x-0.5=0\Rightarrow x^*=8$. $U''(x)=-4/x^2<0$, max.
        $U(8)=4\ln(8)-4=12\ln(2)-4\approx4.318$.

  \item $f'(x)=-\sin(x)$. $f(\pi/3)=1/2$, $f'(\pi/3)=-\sqrt3/2$. Tangent:
        $y=\frac12-\frac{\sqrt3}{2}\left(x-\frac{\pi}{3}\right)$.

  \item $f(0)=1$, $f'(x)=10(1+x)^9$, $f'(0)=10$. So
        $(1.02)^{10}\approx1+10\cdot0.02=1.2$. (True: $1.21899\ldots$,
        error $\approx0.019$.)

  \item $MC(Q)=C'(Q)=3Q^2-18Q+30$. $MC'(Q)=6Q-18=0\Rightarrow Q=3$.
        $MC''(Q)=6>0$, min. $MC(3)=27-54+30=3$.

  \item $q'(p)=A\cdot(-\alpha)p^{-\alpha-1}=-\alpha\dfrac{q}{p}$. Hence
        $\varepsilon_d=\dfrac{p}{q}q'(p)=\dfrac{p}{q}\cdot
        \left(-\alpha\dfrac{q}{p}\right)=-\alpha$. Constant elasticity
        demand has elasticity equal to the negative exponent.

  \item $R(Q)=p(Q)Q=100Q-2Q^2$. $\pi(Q)=100Q-2Q^2-20Q-F=80Q-2Q^2-F$. FOC:
        $\pi'(Q)=80-4Q=0\Rightarrow Q^*=20$. SOC: $\pi''(Q)=-4<0$, max.
        $p^*=100-40=60$. Cross-check via $MR=MC$: $100-4Q=20\Rightarrow
        Q=20$. Consistent. (Profit $\pi^*=80\cdot20-2\cdot400-F
        =1600-800-F=800-F$.)

  \item $f'(x)=2x$. $f(1)=1$, $f'(1)=2$. Normal slope $=-1/2$. Normal:
        $y=1-\frac12(x-1)=-\frac{x}{2}+\frac32$.

  \item $f(0)=0$, $f'(x)=\sec^2(x)$, $f'(0)=1$. So
        $\tan(0.05)\approx0+1\cdot0.05=0.05$. (True: $0.050042\ldots$)

  \item $MC(Q)=C'(Q)=3Q^2-12Q+15$. $AC(Q)=Q^2-6Q+15+20/Q$.
        $AC'(Q)=2Q-6-20/Q^2$. Set $2Q-6-20/Q^2=0$, multiply by $Q^2$:
        $2Q^3-6Q^2-20=0\Rightarrow Q^3-3Q^2-10=0$. Try $Q=$ roots: at
        $Q=3$, $27-27-10=-10$; at $Q=4$, $64-48-10=6$; root between 3 and
        4. Newton step from $Q_0=3.7$: $Q_0^3-3Q_0^2-10
        =50.653-41.07-10=-0.417$, derivative $3Q^2-6Q=3\cdot13.69-22.2
        =18.87$, $Q_1=3.7+0.417/18.87\approx3.722$. So $Q^*\approx3.72$.
        SOC: $AC''(Q)=2+40/Q^3>0$ for $Q>0$, so it is a minimum. (Note: at
        the optimum $AC(Q^*)=MC(Q^*)$, a standard identity.)

  \item $q'(p)=-4p$. At $p=4$: $q(4)=100-32=68$, $q'(4)=-16$.
        $\varepsilon_d(4)=\dfrac{4}{68}\cdot(-16)=-\dfrac{64}{68}
        =-\dfrac{16}{17}\approx-0.94$. Since $|\varepsilon_d|<1$, demand
        is \emph{inelastic} at $p=4$.

  \item $R(Q)=120Q-2Q^2$. $R'(Q)=120-4Q=0\Rightarrow Q^*=30$.
        $R''(Q)=-4<0$, max. $R(30)=3600-1800=1800$.

  \item $f'(x)=2x-3$. $f(2)=4-6+2=0$, $f'(2)=1$. Tangent:
        $y=0+1\cdot(x-2)=x-2$. Normal slope $=-1$. Normal:
        $y=-(x-2)=-x+2$.

  \item $C'(Q)=8Q+20$. $\Delta C\approx C'(10)\cdot0.5
        =(80+20)\cdot0.5=100\cdot0.5=50$.

  \item $\pi(Q)=30Q-Q^2-10Q-9=20Q-Q^2-9$. FOC:
        $\pi'(Q)=20-2Q=0\Rightarrow Q^*=10$. SOC: $\pi''(Q)=-2<0$, max.
        $\pi(10)=200-100-9=91$. Break-even: $20Q-Q^2-9=0\Rightarrow
        Q^2-20Q+9=0\Rightarrow Q=\dfrac{20\pm\sqrt{400-36}}{2}
        =10\pm\sqrt{91}$. Both positive: $Q\approx10-9.54=0.46$ and
        $Q\approx19.54$.

  \item $q(p)=400-p^2$, $q'(p)=-2p$.
        $\varepsilon_d=\dfrac{p}{400-p^2}\cdot(-2p)=\dfrac{-2p^2}{400-p^2}$.
        Set $\dfrac{2p^2}{400-p^2}=1\Rightarrow2p^2=400-p^2
        \Rightarrow3p^2=400\Rightarrow p=\dfrac{20}{\sqrt3}\approx11.55$.

  \item $f'(x)=2xe^{-x}-x^2e^{-x}=x(2-x)e^{-x}$. On $(0,\infty)$:
        $f'>0$ for $0<x<2$, $f'<0$ for $x>2$. Critical point at $x=2$.
        $f''(x)=(2-2x)e^{-x}-(2x-x^2)e^{-x}=(x^2-4x+2)e^{-x}$. At $x=2$:
        $f''(2)=(4-8+2)e^{-2}=-2e^{-2}<0$, confirming local (and global on
        $(0,\infty)$) maximum. $f(2)=4e^{-2}\approx0.541$.

  \item $f'(x)=\dfrac{1}{1+x^2}$. $f(1)=\pi/4$, $f'(1)=1/2$. Tangent:
        $y=\dfrac{\pi}{4}+\dfrac12(x-1)$.

  \item $f(27)=3$, $f'(x)=\frac13x^{-2/3}$,
        $f'(27)=\frac13\cdot27^{-2/3}=\frac13\cdot\frac19=\frac1{27}$. So
        $\sqrt[3]{27.9}\approx3+\frac1{27}\cdot0.9=3+1/30\approx3.0333$.
        (True: $3.03296\ldots$)

  \item $U'(x)=\dfrac{1}{2\sqrt{x}}>0$ for $x>0$ (monotone preferences);
        at $x=100$: $MU=\dfrac{1}{20}=0.05>0$. Elasticity:
        $\varepsilon_{U,x}=\dfrac{x}{U(x)}U'(x)=\dfrac{x}{\sqrt{x}}
        \cdot\dfrac{1}{2\sqrt{x}}=\dfrac12$. Constant elasticity $1/2$,
        positive (utility rises with $x$). Marginal utility itself is
        positive but its elasticity is positive and less than 1.

  \item Invert demand: $q=120-3p\Rightarrow p(Q)=40-Q/3$.
        $R(Q)=Qp=40Q-Q^2/3$. $\pi(Q)=40Q-Q^2/3-10Q=30Q-Q^2/3$. FOC:
        $\pi'(Q)=30-2Q/3=0\Rightarrow Q^*=45$. SOC: $\pi''(Q)=-2/3<0$,
        max. $p^*=40-15=25$, $\pi^*=30\cdot45-45^2/3=1350-675=675$.

  \item $f'(x)=1-16/x^2=0\Rightarrow x^2=16\Rightarrow x=4$ on
        $(0,\infty)$. $f''(x)=32/x^3>0$ for $x>0$, so $x=4$ is a minimum.
        $f(4)=4+4=8$.

  \item $f'(x)=\cos(x)$. $f(\pi/2)=1$, $f'(\pi/2)=0$. The tangent is
        horizontal: $y=1$. The normal is the vertical line $x=\pi/2$
        (slope $-1/0$ undefined; normal is vertical).

  \item At $Q=20$: $p=80$. $R(Q)=100Q-Q^2$, $MR(Q)=100-2Q$, $MR(20)=60$.
        From $q=100-p$, $q'(p)=-1$,
        $\varepsilon_d(p)=\dfrac{p}{q}\cdot(-1)=-\dfrac{p}{100-p}$. At
        $p=80$: $\varepsilon_d=-80/20=-4$. Then
        $p\left(1+\dfrac{1}{\varepsilon_d}\right)=80\left(1-\dfrac14\right)
        =80\cdot\dfrac34=60=MR$. Verified.

  \item $f(0)=0$, $f'(x)=1/(1+x)$, $f'(0)=1$. So
        $\ln(1.05)\approx0+1\cdot0.05=0.05$. True value:
        $\ln(1.05)=0.048790\ldots$ Linearisation error
        $\approx0.05-0.04879=0.00121$.

  \item $\pi(Q)=R(Q)-C(Q)=(20Q-Q^2)-(2Q+\frac12Q^2+10)
        =18Q-\frac32Q^2-10$. FOC: $\pi'(Q)=18-3Q=0\Rightarrow Q=6$.
        $Q=6\in[0,8]$. SOC: $\pi''(Q)=-3<0$. $\pi(6)=108-54-10=44$. Check
        endpoints: $\pi(0)=-10$, $\pi(8)=144-96-10=38$. So $Q^*=6$
        maximises profit; $\pi^*=44$.

  \item $q(p)=50/p$, $q'(p)=-50/p^2$.
        $\varepsilon_d=\dfrac{p}{q}q'(p)=\dfrac{p}{50/p}\cdot(-50/p^2)
        =\dfrac{p^2}{50}\cdot(-50/p^2)=-1$ for all $p>0$. The elasticity
        is identically $-1$, unit-elastic demand. Total revenue
        $R(p)=p\cdot50/p=50$, constant, independent of $p$. This is the
        hallmark of unit-elastic demand.

  \item $\pi'(Q)=12-2Q-Q^2$. Solve $Q^2+2Q-12=0\Rightarrow
        Q=\dfrac{-2+\sqrt{4+48}}{2}=\dfrac{-2+\sqrt{52}}{2}
        =-1+\sqrt{13}$. So $Q^*=\sqrt{13}-1\approx2.606$ (the positive
        root). SOC: $\pi''(Q)=-2-2Q$. At $Q^*$:
        $\pi''(Q^*)=-2-2(\sqrt{13}-1)=-2\sqrt{13}<0$, confirming maximum.

  \item $R(Q)=(50-Q)Q=50Q-Q^2$. $\pi(Q)=50Q-Q^2-10Q-Q^2=40Q-2Q^2$. FOC:
        $\pi'(Q)=40-4Q=0\Rightarrow Q^*=10$. SOC: $\pi''(Q)=-4<0$, max.
        $p^*=50-10=40$, $\pi^*=40\cdot10-2\cdot100=400-200=200$.
        Elasticity: $q(p)=50-p$, $q'(p)=-1$,
        $\varepsilon_d(p)=\dfrac{p}{50-p}\cdot(-1)$. At $p^*=40$:
        $\varepsilon_d=-40/10=-4$ (elastic, as expected at a monopolist's
        optimum where $|\varepsilon_d|>1$).

  \item $\pi(L)=8\sqrt{L}-2L$.
        $\pi'(L)=\dfrac{4}{\sqrt{L}}-2=0\Rightarrow\sqrt{L}=2\Rightarrow
        L^*=4$. SOC: $\pi''(L)=-\dfrac{2}{L^{3/2}}<0$ for $L>0$, so
        $L^*=4$ maximises profit. $Y(4)=2$,
        $\pi^*=8\cdot2-2\cdot4=16-8=8$.
\end{solutions}

\section{Additional Problems}

\begin{problems}
  \item Find the equation of the tangent to $f(x)=x^2+3x-1$ at $x=2$.
  \item Use the linearisation of $f(x)=\sqrt[3]{x}$ at $a=8$ to estimate
        $\sqrt[3]{8.6}$.
  \item For $C(Q)=0.05Q^2+8Q+150$, find $MC(Q)$ and $MC(20)$.
  \item For demand $q(p)=300-4p$, compute $\varepsilon_d(p)$ at $p=25$ and
        classify.
  \item Maximise profit $\pi(Q)=R(Q)-C(Q)$ where $R(Q)=90Q-1.5Q^2$ and
        $C(Q)=Q^2+20Q+30$.
  \item Find the tangent and normal to $f(x)=x^3-4x$ at $x=1$.
  \item Use the linearisation of $f(x)=\mathrm{e}^{x}$ at $a=0$ to estimate
        $\mathrm{e}^{0.1}$.
  \item Inverse demand $p(Q)=100-0.5Q$. Find $R(Q)$, $MR(Q)$, and $MR(40)$.
  \item For the Cobb--Douglas production $Y=AK^{\alpha}L^{1-\alpha}$ with
        $K$ fixed, compute the output elasticity $\varepsilon_{Y,L}$.
  \item Minimise the average cost $AC(Q)=C(Q)/Q$ for
        $C(Q)=3Q^2+12Q+108$, over $Q>0$.
  \item A monopolist faces $p(Q)=80-2Q$ and $C(Q)=Q^2+8Q+40$. Find $Q^*$,
        $p^*$, and $\pi^*$.
  \item For demand $q(p)=120\mathrm{e}^{-0.04p}$, compute
        $\varepsilon_d(p)$ and at $p=10$.
  \item For utility $U(x)=2\sqrt{x}-0.1x$ for $x>0$, find the maximiser
        and verify SOC.
  \item Estimate $\ln(1.04)$ using the linearisation of $f(x)=\ln(1+x)$
        at $a=0$.
  \item Minimise $f(x)=2x+\dfrac{50}{x}$ on $(0,\infty)$ and verify SOC.
\end{problems}

\section{Additional Solutions}

\begin{solutions}
  \item $f'(x)=2x+3$. $f(2)=9$, $f'(2)=7$. Tangent: $y=9+7(x-2)=7x-5$.

  \item $f(8)=2$, $f'(x)=\tfrac{1}{3}x^{-2/3}$,
        $f'(8)=\tfrac{1}{3}\cdot 8^{-2/3}=\tfrac{1}{3}\cdot\tfrac{1}{4}
        =\tfrac{1}{12}$. So
        $\sqrt[3]{8.6}\approx 2+\tfrac{0.6}{12}=2+0.05=2.05$. (True:
        $2.0494\ldots$)

  \item $MC(Q)=C'(Q)=0.1Q+8$. $MC(20)=2+8=10$.

  \item $q'(p)=-4$.
        $\varepsilon_d(p)=\dfrac{p}{300-4p}\cdot(-4)=\dfrac{-4p}{300-4p}$.
        At $p=25$: $\varepsilon_d=\dfrac{-100}{200}=-\dfrac{1}{2}$.
        $|\varepsilon_d|<1$: \emph{inelastic}.

  \item $\pi(Q)=90Q-1.5Q^2-Q^2-20Q-30=70Q-2.5Q^2-30$. FOC:
        $\pi'(Q)=70-5Q=0\Rightarrow Q^*=14$. SOC: $\pi''(Q)=-5<0$, max.
        $\pi(14)=980-490-30=460$.

  \item $f'(x)=3x^2-4$. $f(1)=-3$, $f'(1)=-1$. Tangent:
        $y=-3-(x-1)=-x-2$. Normal slope $=-1/(-1)=1$. Normal:
        $y=-3+1\cdot(x-1)=x-4$.

  \item $f(0)=1$, $f'(0)=1$. $\mathrm{e}^{0.1}\approx 1+0.1=1.1$. (True:
        $1.10517\ldots$)

  \item $R(Q)=p(Q)Q=100Q-0.5Q^2$. $MR(Q)=100-Q$. $MR(40)=60$.

  \item $\dfrac{\partial Y}{\partial L}=(1-\alpha)AK^{\alpha}L^{-\alpha}$.
        So
        \[
          \varepsilon_{Y,L}=\frac{L}{Y}\cdot(1-\alpha)AK^{\alpha}L^{-\alpha}
            =\frac{(1-\alpha)AK^{\alpha}L^{1-\alpha}}{AK^{\alpha}L^{1-\alpha}}
            =1-\alpha .
        \]
        Standard Cobb--Douglas: elasticity equals the exponent.

  \item $AC(Q)=3Q+12+\dfrac{108}{Q}$.
        $AC'(Q)=3-\dfrac{108}{Q^2}=0\Rightarrow Q^2=36\Rightarrow Q^*=6$.
        $AC''(Q)=\dfrac{216}{Q^3}>0$ at $Q=6$: minimum.
        $AC(6)=18+12+18=48$.

  \item $R(Q)=p(Q)Q=80Q-2Q^2$.
        $\pi(Q)=80Q-2Q^2-Q^2-8Q-40=72Q-3Q^2-40$. FOC:
        $\pi'(Q)=72-6Q=0\Rightarrow Q^*=12$. SOC: $\pi''(Q)=-6<0$.
        $p^*=80-24=56$. $\pi^*=72\cdot12-3\cdot144-40=864-432-40=392$.

  \item $q'(p)=120\cdot(-0.04)\mathrm{e}^{-0.04p}
        =-4.8\mathrm{e}^{-0.04p}$. So
        \[
          \varepsilon_d(p)=\frac{p}{120\mathrm{e}^{-0.04p}}
            \cdot(-4.8\mathrm{e}^{-0.04p})=-0.04p .
        \]
        At $p=10$: $\varepsilon_d=-0.4$. $|\varepsilon_d|<1$:
        \emph{inelastic}.

  \item $U'(x)=\tfrac{1}{\sqrt{x}}-0.1=0\Rightarrow\sqrt{x}=10
        \Rightarrow x^*=100$. $U''(x)=-\tfrac{1}{2}x^{-3/2}<0$ for $x>0$,
        max. $U(100)=20-10=10$.

  \item $f(0)=0$, $f'(0)=1$. $\ln(1.04)\approx 0+1\cdot0.04=0.04$.
        (True: $0.03922\ldots$)

  \item $f'(x)=2-\dfrac{50}{x^2}=0\Rightarrow x^2=25\Rightarrow x=5$ on
        $(0,\infty)$. $f''(x)=\dfrac{100}{x^3}>0$ for $x>0$, so $x=5$ is
        a minimum. $f(5)=10+10=20$.
\end{solutions}

\chapter{First-Order Partial Derivatives and the Gradient}

When a function depends on more than one input, we measure how it responds
to a change in \emph{one} variable at a time, holding the others fixed,
the ceteris paribus principle. This is exactly what a partial derivative
computes. From these one-variable slopes we assemble the gradient, which
packages all first-order information into a single vector and unlocks the
geometry of multivariate optimisation.

\section{Core Rules}

\begin{rulebox}{Topological vocabulary of domains}
\begin{itemize}[leftmargin=1.6em,itemsep=0.3ex]
  \item The \emph{open ball} (or \emph{neighbourhood}) of radius $r>0$ around $\mathbf{x}_0\in\R^n$ is
        $B_r(\mathbf{x}_0)=\{\mathbf{x}\in\R^n :
        \|\mathbf{x}-\mathbf{x}_0\|<r\}$.
  \item $\mathbf{x}_0\in M$ is an \emph{interior point} of $M$ if some open ball around
        $\mathbf{x}_0$ lies entirely in $M$.
  \item $M$ is \emph{open} if \emph{every} point of $M$ is
        an interior point.
  \item $\mathbf{x}_0$ is a \emph{boundary point} of
        $M$ if every ball around $\mathbf{x}_0$ contains points of $M$
        \emph{and} points of its complement.
  \item $M$ is \emph{closed} if it contains all
        its boundary points; equivalently, the complement
        $\R^n\setminus M$ is open.
  \item $M$ is \emph{bounded} if
        $M\subseteq B_r(\mathbf{0})$ for some radius $r$.
  \item $M\subseteq\R^n$ is \emph{compact} exactly when
        $M$ is closed \emph{and} bounded.
\end{itemize}
Economic examples: the open positive orthant $\{(x,y): x>0,\ y>0\}$
(the natural Cobb--Douglas domain) is open and unbounded; a budget set
$\{\mathbf{x}\in\R^n : \mathbf{p}\cdot\mathbf{x}\le m,\
\mathbf{x}\ge\mathbf{0}\}$ is closed and bounded, hence compact.
\end{rulebox}

\begin{rulebox}{Why it matters}
Differentiability (and hence every partial derivative below) is
defined at \emph{interior} points of the domain, where $f$ can be
probed in all directions. And by the \emph{Weierstrass theorem}, a
continuous function on a compact set attains its maximum and minimum:
this is why budget-constrained optimisation problems are
guaranteed to have solutions.
\end{rulebox}

\begin{rulebox}{Definition}
Let $f:\R^n\to\R$. The partial derivative of $f$ with respect to $x_i$ at
$\mathbf{x}=(x_1,\ldots,x_n)$ is
\[
  \pd{f}{x_i}(\mathbf{x}) \;=\; \lim_{h\to0}
    \frac{f(x_1,\ldots,x_i+h,\ldots,x_n)-f(x_1,\ldots,x_n)}{h} .
\]
\end{rulebox}

\begin{rulebox}{Practical computation}
To compute $\partial f/\partial x_i$, differentiate $f$ with respect to
$x_i$ and treat every other variable as a constant. All single-variable
rules (power, product, quotient, chain) apply unchanged.
\end{rulebox}

\begin{rulebox}{Notation}
The following all refer to the same object:
\[
  \pd{f}{x} = f_x = \partial_x f = D_x f .
\]
\end{rulebox}

\begin{rulebox}{Gradient}
The vector of all first-order partial derivatives is
\[
  \grad f(\mathbf{x}) \;=\;
  \begin{pmatrix} f_{x_1}(\mathbf{x}) \\ f_{x_2}(\mathbf{x}) \\ \vdots \\
    f_{x_n}(\mathbf{x}) \end{pmatrix} .
\]
\end{rulebox}

\begin{rulebox}{Geometric meaning}
\begin{itemize}[leftmargin=1.6em,itemsep=0.3ex]
  \item $\grad f(\mathbf{x}_0)$ points in the direction of \emph{steepest
        ascent} at $\mathbf{x}_0$.
  \item $\|\grad f(\mathbf{x}_0)\|$ is the maximal rate of increase at that
        point.
  \item $\grad f(\mathbf{x}_0)$ is \emph{perpendicular} to the level set
        $\{\mathbf{x} : f(\mathbf{x})=f(\mathbf{x}_0)\}$ at $\mathbf{x}_0$.
\end{itemize}
\end{rulebox}

\begin{rulebox}{Economic interpretation}
\begin{itemize}[leftmargin=1.6em,itemsep=0.3ex]
  \item For a production function $Q=F(K,L)$: $F_K=\mathrm{MP}_K$
        (marginal product of capital), $F_L=\mathrm{MP}_L$ (marginal
        product of labour).
  \item For a utility function $U(x,y)$: $U_x, U_y$ are the marginal
        utilities; the marginal rate of substitution is
        $\mathrm{MRS}_{xy}=U_x/U_y$.
  \item For production: the marginal rate of technical substitution is
        $\mathrm{MRTS}_{LK}=F_L/F_K$.
\end{itemize}
\end{rulebox}

\begin{rulebox}{Directional derivative}
Let $\mathbf{u}\in\R^n$ be a \emph{unit} vector ($\|\mathbf{u}\|=1$). The
directional derivative of $f$ at $\mathbf{x}_0$ in the direction
$\mathbf{u}$ is
\[
  D_{\mathbf{u}}f(\mathbf{x}_0) \;=\; \grad f(\mathbf{x}_0)\cdot\mathbf{u} .
\]
\end{rulebox}

\begin{rulebox}{Caution}
If you are given a non-unit direction $\mathbf{v}$, you \emph{must}
normalise first: $\mathbf{u}=\mathbf{v}/\|\mathbf{v}\|$. Forgetting this is
the single most common error in directional-derivative problems.
\end{rulebox}

\section{Practice Problems}

\begin{problems}
  \item Compute $f_x$ and $f_y$ for $f(x,y)=5x^4-2x^2y^3+y^5$.
  \item For $F(K,L)=12\,K^{1/4}L^{3/4}$, find $F_K$ and $F_L$ at
        $(K,L)=(16,81)$.
  \item Compute $f_x$ and $f_y$ for $f(x,y)=x^2e^{3y}$.
  \item Compute $\nabla f$ for $f(x,y)=4x^3y-xy^4+7$.
  \item Compute the directional derivative of $f(x,y)=x^2-4xy+y^3$ at
        $(2,1)$ in the direction $\mathbf{v}=(5,-12)$.
  \item For $F(K,L)=6K^{2/3}L^{1/3}$, compute $\mathrm{MP}_K$ and
        $\mathrm{MP}_L$ at $(K,L)=(27,8)$.
  \item Compute $f_x$ and $f_y$ for $f(x,y)=(3x-2y)^4$.
  \item Compute $f_x$ and $f_y$ for $f(x,y)=e^{xy^2}$.
  \item For $F(K,L)=25K^{0.6}L^{0.5}$, find $F_K$ and $F_L$.
  \item Compute $\nabla f$ for $f(x,y,z)=x^2+3yz-z^2e^x$.
  \item Compute the directional derivative of $f(x,y)=e^x\ln(1+y)$ at
        $(0,0)$ in the direction $\mathbf{v}=(1,2)$.
  \item For $U(x,y)=2\ln(x)+5\ln(y)$, derive $\mathrm{MRS}_{xy}=U_x/U_y$ as
        a function of $(x,y)$ and evaluate at $(3,5)$.
  \item Simplify and compute $f_x, f_y$ for
        $f(x,y)=\dfrac{x^2-y^2}{x+y}$.
  \item Compute $f_x$ and $f_y$ for $f(x,y)=\ln(x^2y+1)$.
  \item For $U(x_1,x_2,x_3)=x_1^{1/2}x_2^{1/3}x_3^{1/6}$, compute the
        three partials.
  \item Compute the directional derivative of $f(x,y,z)=xy-2yz+z^2$ at
        $(1,2,3)$ in the direction $\mathbf{v}=(1,2,2)$.
  \item For $F(K,L)=K^{0.25}L^{0.75}$, derive $\mathrm{MRTS}_{LK}=F_L/F_K$
        and evaluate at $(K,L)=(16,4)$.
  \item Expand and then compute $f_x, f_y$ for
        $f(x,y)=(x^2+y)(x-3y^2)$.
  \item Compute $f_x$ and $f_y$ for $f(x,y)=xe^{-y/x}$ with $x\neq0$.
  \item For $F(K,L)=(3K^{1/2}+2L^{1/2})^2$, find $F_K$ and $F_L$.
  \item Compute $\nabla f$ for $f(x,y)=\arctan(y/x)$ with $x>0$.
  \item For $f(x,y)=x^{1/2}y^{1/2}$ at $(K,L)=(4,9)$, find the unit
        direction of steepest ascent and the maximum rate of increase.
  \item For the CES function $F(K,L)=(K^{1/2}+L^{1/2})^2$, verify Euler's
        theorem $KF_K+LF_L=rF$ for the appropriate degree of homogeneity
        $r$.
  \item Compute $f_x$ and $f_y$ for $f(x,y)=\dfrac{2x+3y}{x^2+y^2}$.
  \item Compute $f_x$ and $f_y$ for $f(x,y)=\dfrac{e^x}{1+e^y}$.
  \item For $F(K,L)=A(\alpha K^\rho+(1-\alpha)L^\rho)^{\nu/\rho}$ (CES),
        find $F_L$.
  \item Compute the directional derivative of $f(x,y)=\ln(x^2+y^2)$ at
        $(3,4)$ in the direction $\mathbf{v}=(-1,1)$.
  \item Compute $f_x, f_y, f_z$ for $f(x,y,z)=x^2yz^3-4xy^2z+5z$.
  \item Compute $f_x$ and $f_y$ for $f(x,y)=\sqrt{4x^2+9y^2}$.
  \item For $F(K,L)=(K^{1/2}+L^{1/2})^4$, compute $F_K$ at
        $(K,L)=(4,9)$.
  \item Compute $\nabla f$ for $f(x,y,z)=\ln(x^2+y^2+z^2)$.
  \item Compute $f_x$ and $f_y$ for
        $f(x,y)=\ln\left(\sqrt{x^2+y^2}\right)$.
  \item Cost minimisation. A firm uses inputs $(x,y)$ with cost
        $C(x,y)=4x+9y$ and output $Q(x,y)=x^{1/2}y^{1/2}$. Compute
        $\nabla C$ and $\nabla Q$, and verify that at the cost-minimising
        input mix $\nabla C=\lambda\nabla Q$ for some $\lambda$.
  \item Compute $f_x$ and $f_y$ for $f(x,y)=(1+xy)^5$.
  \item Compute the directional derivative of $f(x,y)=x^3-2xy+y^2$ at
        $(1,2)$ in the direction $\mathbf{v}=(3,4)$.
  \item For $F(K,L)=10K^{0.3}L^{0.7}$ (Cobb--Douglas), compute $F_K$ and
        $F_L$ at $(K,L)=(100,100)$ and verify $KF_K+LF_L=F$.
  \item Compute $f_x$ and $f_y$ for $f(x,y)=\sin(xy)+\cos(x-y)$.
  \item Compute $\nabla f$ for $f(x,y)=x^2y-xy^2+3x-2y$ at the point
        $(1,-1)$.
  \item For $U(x,y)=x^\alpha y^{1-\alpha}$ (Cobb--Douglas,
        $0<\alpha<1$), derive $\mathrm{MRS}_{xy}=U_x/U_y$ and show it
        equals $\dfrac{\alpha}{1-\alpha}\cdot\dfrac{y}{x}$.
  \item Compute the directional derivative of $f(x,y)=e^{xy}$ at $(1,0)$
        in the direction $\mathbf{v}=(-3,4)$.
  \item Compute $f_x$ and $f_y$ for $f(x,y)=x^y$ (with $x>0$).
        \emph{Hint:} write $x^y=e^{y\ln x}$.
  \item For $F(K,L)=\min\{2K,3L\}$ (Leontief, fixed-proportions, on the
        smooth region away from the kink), compute $F_K$ and $F_L$ in
        each of the two regions $\{2K<3L\}$ and $\{2K>3L\}$.
  \item Compute $\nabla f$ for $f(x,y,z)=xyz+x^2-z^2$ at the point
        $(2,-1,3)$.
  \item Compute the unit vector of steepest \emph{descent} for
        $f(x,y)=x^2+4y^2$ at $(3,1)$, and the corresponding rate of
        decrease.
  \item For $f(x,y)=xe^y-ye^x$, compute $f_x, f_y$.
  \item Compute $f_x$ and $f_y$ for $f(x,y)=\dfrac{x}{x^2+y^2}$.
  \item For $U(x,y)=x+2\sqrt{y}$ (quasi-linear), derive
        $\mathrm{MRS}_{xy}$.
  \item Compute the directional derivative of $f(x,y,z)=x^2+y^2-z^2$ at
        $(2,-1,2)$ in the direction $\mathbf{v}=(2,2,1)$.
  \item Compute $f_x$ and $f_y$ for $f(x,y)=\tan(x+y)+\tan(x-y)$.
  \item For $F(K,L)=(\alpha K+\beta L)^r$ (linear-aggregate,
        $\alpha,\beta,r>0$), compute $F_K, F_L$, and verify
        $KF_K+LF_L=rF$.
  \item Compute $\nabla f$ for $f(x,y)=e^{-x}\sin(y)+e^x\cos(y)$.
  \item Compute the directional derivative of $f(x,y)=\sqrt{x^2+y^2}$ at
        $(3,4)$ in the direction $\mathbf{v}=(1,0)$.
  \item For $U(x,y)=\min\{2x,3y\}$, compute the marginal utilities on
        each smooth region $\{2x<3y\}$ and $\{2x>3y\}$.
  \item Compute $f_x, f_y$ for $f(x,y)=\dfrac{\ln(x+y)}{x-y}$ (assume
        $x>y>0$).
  \item For $F(K,L)=AK^\alpha L^\beta$ general Cobb--Douglas, derive
        $\mathrm{MRTS}_{LK}$ and show it depends only on the ratio $K/L$
        (not on scale).
  \item Compute the unit direction of steepest ascent for
        $f(x,y)=x^2-y^2$ at $(2,-1)$, and the maximum rate.
  \item Compute $f_x$ and $f_y$ for $f(x,y)=\cosh(x+y)+\sinh(x-y)$.
  \item For $f(x,y,z)=e^{x+y+z}$, compute $\nabla f$ at $(0,0,0)$ and the
        directional derivative in $\mathbf{v}=(1,1,1)$.
  \item Compute $f_x$ and $f_y$ for $f(x,y)=(\ln x)(\ln y)$ with
        $x,y>1$.
  \item For $F(K,L)=K+2K^{1/2}L^{1/2}+L$ (which equals
        $(K^{1/2}+L^{1/2})^2$), compute $F_K, F_L$ directly and confirm
        agreement with the CES form.
  \item Compute the directional derivative of $f(x,y)=x^2y+y^3$ at
        $(-1,2)$ in the direction $\mathbf{v}=(4,-3)$.
  \item For $U(x,y)=\alpha\ln x+\beta\ln y$ (log-Cobb--Douglas), derive
        $\mathrm{MRS}_{xy}$ and evaluate at $(x,y)=(\alpha,\beta)$.
  \item Compute $\nabla f$ for $f(x,y)=(x^2+y^2)^{-1/2}$ (Newtonian
        potential in 2D).
  \item Compute $f_x$ and $f_y$ for $f(x,y)=\arctan(xy)$.
  \item For $F(K,L)=(K^\rho+L^\rho)^{1/\rho}$ with $\rho<1$, $\rho\neq0$
        (CES with constant returns), compute $F_K$ and show
        $F_K=(K/F)^{\rho-1}$.
  \item Compute the directional derivative of $f(x,y,z)=xyz$ at
        $(2,3,-1)$ in the direction $\mathbf{v}=(2,-1,2)$.
\end{problems}

\section{Solutions}

\begin{solutions}
  \item $f_x=20x^3-4xy^3$, $f_y=-6x^2y^2+5y^4$.

  \item $F_K=12\cdot\frac14K^{-3/4}L^{3/4}=3(L/K)^{3/4}$;
        $F_L=12\cdot\frac34K^{1/4}L^{-1/4}=9(K/L)^{1/4}$. At $(16,81)$:
        $(L/K)^{3/4}=(81/16)^{3/4}=27/8$ so $F_K=81/8$;
        $(K/L)^{1/4}=(16/81)^{1/4}=2/3$ so $F_L=6$.

  \item $f_x=2xe^{3y}$, $f_y=3x^2e^{3y}$.

  \item $\nabla f=(12x^2y-y^4,\ 4x^3-4xy^3)^T$.

  \item $f_x=2x-4y$, $f_y=-4x+3y^2$. At $(2,1)$: $\nabla f=(0,-5)^T$.
        Normalise: $\|\mathbf{v}\|=\sqrt{25+144}=13$,
        $\mathbf{u}=(5/13,-12/13)^T$. Then
        \[
          D_{\mathbf{u}}f=0\cdot\frac{5}{13}+(-5)\cdot
            \left(-\frac{12}{13}\right)=\frac{60}{13} .
        \]

  \item $F_K=6\cdot\frac23K^{-1/3}L^{1/3}=4(L/K)^{1/3}$,
        $F_L=2(K/L)^{2/3}$. At $(27,8)$: $(L/K)^{1/3}=2/3$ so
        $\mathrm{MP}_K=8/3$; $(K/L)^{2/3}=(27/8)^{2/3}=9/4$ so
        $\mathrm{MP}_L=9/2$.

  \item Chain rule with $u=3x-2y$: $f_x=4(3x-2y)^3\cdot3=12(3x-2y)^3$,
        $f_y=-8(3x-2y)^3$.

  \item $f_x=y^2e^{xy^2}$, $f_y=2xye^{xy^2}$.

  \item $F_K=15K^{-0.4}L^{0.5}$, $F_L=12.5K^{0.6}L^{-0.5}$.

  \item $f_x=2x-z^2e^x$, $f_y=3z$, $f_z=3y-2ze^x$. So
        $\nabla f=(2x-z^2e^x,\ 3z,\ 3y-2ze^x)^T$.

  \item $f_x=e^x\ln(1+y)$, $f_y=e^x/(1+y)$. At $(0,0)$:
        $\nabla f=(0,1)^T$. Normalise: $\|\mathbf{v}\|=\sqrt5$,
        $\mathbf{u}=(1/\sqrt5,2/\sqrt5)^T$. So
        \[
          D_{\mathbf{u}}f=0\cdot\frac{1}{\sqrt5}+1\cdot\frac{2}{\sqrt5}
            =\frac{2\sqrt5}{5} .
        \]

  \item $U_x=2/x$, $U_y=5/y$. $\mathrm{MRS}_{xy}=U_x/U_y
        =(2/x)/(5/y)=2y/(5x)$. At $(3,5)$: $\mathrm{MRS}=10/15=2/3$.

  \item For $x+y\neq0$: $f(x,y)=x-y$, so $f_x=1$, $f_y=-1$.

  \item With $u=x^2y+1$: $f_x=2xy/(x^2y+1)$, $f_y=x^2/(x^2y+1)$.

  \item $U_{x_1}=\frac12x_1^{-1/2}x_2^{1/3}x_3^{1/6}$,
        $U_{x_2}=\frac13x_1^{1/2}x_2^{-2/3}x_3^{1/6}$,
        $U_{x_3}=\frac16x_1^{1/2}x_2^{1/3}x_3^{-5/6}$. Equivalently
        $U_{x_i}=(\alpha_i/x_i)U$ with $\alpha=(1/2,1/3,1/6)$.

  \item $f_x=y$, $f_y=x-2z$, $f_z=-2y+2z$. At $(1,2,3)$:
        $\nabla f=(2,-5,2)^T$. Normalise: $\|\mathbf{v}\|=\sqrt{1+4+4}=3$,
        $\mathbf{u}=(1/3,2/3,2/3)^T$. Then
        $D_{\mathbf{u}}f=(2-10+4)/3=-4/3$.

  \item $F_K=0.25K^{-0.75}L^{0.75}$, $F_L=0.75K^{0.25}L^{-0.25}$. So
        $\mathrm{MRTS}_{LK}=F_L/F_K=3K/L$. At $(16,4)$:
        $\mathrm{MRTS}=48/4=12$.

  \item Expand: $f=x^3-3x^2y^2+xy-3y^3$. Hence $f_x=3x^2-6xy^2+y$,
        $f_y=-6x^2y+x-9y^2$.

  \item With $u=-y/x$: $\partial u/\partial x=y/x^2$,
        $\partial u/\partial y=-1/x$. Product rule:
        \[
          f_x=e^{-y/x}+x\cdot e^{-y/x}\cdot\frac{y}{x^2}
            =e^{-y/x}\left(1+\frac{y}{x}\right), \qquad
          f_y=x\cdot e^{-y/x}\cdot\left(-\frac1x\right)=-e^{-y/x} .
        \]

  \item Let $g=3K^{1/2}+2L^{1/2}$, so $F=g^2$. Then
        \[
          F_K=2g\cdot\frac32K^{-1/2}=\frac{3g}{K^{1/2}}
            =\frac{3(3K^{1/2}+2L^{1/2})}{K^{1/2}}, \qquad
          F_L=2g\cdot L^{-1/2}=\frac{2(3K^{1/2}+2L^{1/2})}{L^{1/2}} .
        \]

  \item With $u=y/x$, $f=\arctan u$, $f'=1/(1+u^2)$:
        \[
          f_x=\frac{1}{1+y^2/x^2}\cdot\left(-\frac{y}{x^2}\right)
            =\frac{-y}{x^2+y^2}, \qquad
          f_y=\frac{x}{x^2+y^2} .
        \]
        So $\nabla f=(-y/(x^2+y^2),\ x/(x^2+y^2))^T$.

  \item $f_K=\frac12K^{-1/2}L^{1/2}$, $f_L=\frac12K^{1/2}L^{-1/2}$. At
        $(4,9)$: $\nabla f=(\frac12\cdot\frac32,\ \frac12\cdot\frac23)^T
        =(3/4,1/3)^T$. Then $\|\nabla f\|^2=9/16+1/9=97/144$, so the
        maximum rate is $\sqrt{97}/12\approx0.821$ in direction
        $\mathbf{u}^*=(9,4)^T/\sqrt{97}$.

  \item $F=(K^{1/2}+L^{1/2})^2$. Replacing $(K,L)\to(\lambda K,\lambda L)$
        gives $\lambda F$, so $r=1$. With $g=K^{1/2}+L^{1/2}$:
        $F_K=2g\cdot\frac12K^{-1/2}=g/K^{1/2}$, $F_L=g/L^{1/2}$. Hence
        \[
          KF_K+LF_L=K^{1/2}g+L^{1/2}g=g(K^{1/2}+L^{1/2})=g^2=F. \checkmark
        \]

  \item Quotient rule, $D=x^2+y^2$:
        \[
          f_x=\frac{2D-(2x+3y)(2x)}{D^2}=\frac{-2x^2-6xy+2y^2}{(x^2+y^2)^2},
          \qquad f_y=\frac{3x^2-4xy-3y^2}{(x^2+y^2)^2} .
        \]

  \item Treat denominator as constant for $f_x$: $f_x=e^x/(1+e^y)$. For
        $f_y$:
        \[
          f_y=e^x\cdot\frac{-e^y}{(1+e^y)^2}=-\frac{e^{x+y}}{(1+e^y)^2} .
        \]

  \item Let $S=\alpha K^\rho+(1-\alpha)L^\rho$, $F=AS^{\nu/\rho}$.
        Treating $K$ constant, $\partial S/\partial L=\rho(1-\alpha)L^{\rho-1}$:
        \[
          F_L=A\cdot\frac{\nu}{\rho}S^{\nu/\rho-1}\cdot\rho(1-\alpha)L^{\rho-1}
            =\nu(1-\alpha)AL^{\rho-1}S^{\nu/\rho-1} .
        \]

  \item $f=\ln(x^2+y^2)$, $f_x=2x/(x^2+y^2)$, $f_y=2y/(x^2+y^2)$. At
        $(3,4)$: $x^2+y^2=25$, $\nabla f=(6/25,8/25)^T$. Normalise:
        $\|\mathbf{v}\|=\sqrt2$, $\mathbf{u}=(-1/\sqrt2,1/\sqrt2)^T$. Then
        \[
          D_{\mathbf{u}}f=\frac{6}{25}\cdot\left(-\frac{1}{\sqrt2}\right)
            +\frac{8}{25}\cdot\frac{1}{\sqrt2}=\frac{2}{25\sqrt2}
            =\frac{\sqrt2}{25} .
        \]

  \item $f_x=2xyz^3-4y^2z$, $f_y=x^2z^3-8xyz$,
        $f_z=3x^2yz^2-4xy^2+5$.

  \item Chain rule, $u=4x^2+9y^2$:
        \[
          f_x=\frac{8x}{2\sqrt u}=\frac{4x}{\sqrt{4x^2+9y^2}}, \qquad
          f_y=\frac{9y}{\sqrt{4x^2+9y^2}} .
        \]

  \item With $g=K^{1/2}+L^{1/2}$, $F=g^4$,
        $F_K=4g^3\cdot\frac12K^{-1/2}=2g^3/K^{1/2}$. At $(4,9)$: $g=5$,
        $K^{1/2}=2$, so $F_K=2\cdot125/2=125$.

  \item Let $r^2=x^2+y^2+z^2$. Then $f=\ln(r^2)$ and $f_{x_i}=2x_i/r^2$.
        So $\nabla f=(2/(x^2+y^2+z^2))(x,y,z)^T$.

  \item Use $\ln\sqrt{u}=\frac12\ln u$: $f=\frac12\ln(x^2+y^2)$. So
        $f_x=x/(x^2+y^2)$, $f_y=y/(x^2+y^2)$.

  \item $\nabla C=(4,9)^T$. With $Q=x^{1/2}y^{1/2}$: $Q_x=\frac12(y/x)^{1/2}$,
        $Q_y=\frac12(x/y)^{1/2}$, so
        $\nabla Q=\frac12\left((y/x)^{1/2},(x/y)^{1/2}\right)^T$. The
        condition $\nabla C=\lambda\nabla Q$ requires
        $4/9=Q_x/Q_y=y/x$, i.e.\ $y=4x/9$. Along this locus
        $\nabla Q\propto(2/3,3/2)^T\propto(4,9)^T=\nabla C$. \checkmark

  \item Chain rule: $f_x=5(1+xy)^4\cdot y=5y(1+xy)^4$,
        $f_y=5x(1+xy)^4$.

  \item $f_x=3x^2-2y$, $f_y=-2x+2y$. At $(1,2)$:
        $\nabla f=(3-4,-2+4)^T=(-1,2)^T$. Normalise: $\|\mathbf{v}\|=5$,
        $\mathbf{u}=(3/5,4/5)^T$. Then
        \[
          D_{\mathbf{u}}f=-1\cdot\frac35+2\cdot\frac45
            =\frac{-3+8}{5}=1 .
        \]

  \item $F_K=10\cdot0.3K^{-0.7}L^{0.7}=3(L/K)^{0.7}$;
        $F_L=7(K/L)^{0.3}$. At $K=L=100$: $F_K=3$, $F_L=7$. And
        $F=10\cdot100^{0.3}\cdot100^{0.7}=10\cdot100=1000$. Then
        $KF_K+LF_L=100\cdot3+100\cdot7=1000=F$. \checkmark (homogeneity
        degree $r=0.3+0.7=1$).

  \item $f_x=y\cos(xy)-\sin(x-y)$, $f_y=x\cos(xy)+\sin(x-y)$.

  \item $f_x=2xy-y^2+3$, $f_y=x^2-2xy-2$. At $(1,-1)$: $f_x=-2-1+3=0$,
        $f_y=1+2-2=1$. So $\nabla f(1,-1)=(0,1)^T$.

  \item $U_x=\alpha x^{\alpha-1}y^{1-\alpha}$,
        $U_y=(1-\alpha)x^\alpha y^{-\alpha}$. So
        \[
          \mathrm{MRS}_{xy}=\frac{U_x}{U_y}
            =\frac{\alpha x^{\alpha-1}y^{1-\alpha}}{(1-\alpha)x^\alpha y^{-\alpha}}
            =\frac{\alpha}{1-\alpha}\cdot\frac{y}{x} . \checkmark
        \]

  \item $f_x=ye^{xy}$, $f_y=xe^{xy}$. At $(1,0)$: $\nabla f=(0,1)^T$.
        Normalise: $\|\mathbf{v}\|=5$, $\mathbf{u}=(-3/5,4/5)^T$. Then
        $D_{\mathbf{u}}f=0\cdot(-3/5)+1\cdot(4/5)=4/5$.

  \item Write $f=e^{y\ln x}$. Then $f_x=e^{y\ln x}\cdot y/x=yx^{y-1}$ and
        $f_y=e^{y\ln x}\cdot\ln x=x^y\ln x$.

  \item On $\{2K<3L\}$ the binding factor is capital: $F=2K$, so
        $F_K=2$, $F_L=0$. On $\{2K>3L\}$ labour binds: $F=3L$, so
        $F_K=0$, $F_L=3$. (Non-differentiable on $\{2K=3L\}$.)

  \item $f_x=yz+2x$, $f_y=xz$, $f_z=xy-2z$. At $(2,-1,3)$:
        $f_x=-3+4=1$, $f_y=6$, $f_z=-2-6=-8$. So
        $\nabla f=(1,6,-8)^T$.

  \item $\nabla f=(2x,8y)^T$. At $(3,1)$: $\nabla f=(6,8)^T$. Steepest
        descent is $-\nabla f/\|\nabla f\|$. $\|\nabla f\|=10$, so unit
        descent direction is $(-3/5,-4/5)^T$, with rate of decrease
        $-\|\nabla f\|=-10$ (i.e.\ $f$ decreases at rate 10).

  \item $f_x=e^y-ye^x$, $f_y=xe^y-e^x$.

  \item Quotient rule with $D=x^2+y^2$:
        \[
          f_x=\frac{D-x\cdot2x}{D^2}=\frac{y^2-x^2}{(x^2+y^2)^2}, \qquad
          f_y=\frac{-x\cdot2y}{D^2}=\frac{-2xy}{(x^2+y^2)^2} .
        \]

  \item $U_x=1$, $U_y=1/\sqrt{y}$. So $\mathrm{MRS}_{xy}=U_x/U_y
        =\sqrt{y}$ (depends only on $y$, quasi-linear utility).

  \item $f_x=2x$, $f_y=2y$, $f_z=-2z$. At $(2,-1,2)$:
        $\nabla f=(4,-2,-4)^T$. Normalise: $\|\mathbf{v}\|=\sqrt{4+4+1}=3$,
        $\mathbf{u}=(2/3,2/3,1/3)^T$. Then
        \[
          D_{\mathbf{u}}f=4\cdot\frac23+(-2)\cdot\frac23+(-4)\cdot\frac13
            =\frac{8-4-4}{3}=0 .
        \]
        (We are moving tangent to the level surface.)

  \item $f_x=\sec^2(x+y)+\sec^2(x-y)$, $f_y=\sec^2(x+y)-\sec^2(x-y)$.

  \item Let $S=\alpha K+\beta L$, $F=S^r$. Then $F_K=rS^{r-1}\alpha$,
        $F_L=rS^{r-1}\beta$. Hence
        \[
          KF_K+LF_L=rS^{r-1}(\alpha K+\beta L)=rS^{r-1}\cdot S=rS^r=rF.
            \checkmark
        \]

  \item $\nabla f=(-e^{-x}\sin y+e^x\cos y,\ e^{-x}\cos y-e^x\sin y)^T$.

  \item $f_x=x/\sqrt{x^2+y^2}$, $f_y=y/\sqrt{x^2+y^2}$. At $(3,4)$:
        $\sqrt{x^2+y^2}=5$, $\nabla f=(3/5,4/5)^T$. The direction
        $\mathbf{v}=(1,0)$ is already unit. So
        $D_{\mathbf{u}}f=3/5\cdot1+4/5\cdot0=3/5$.

  \item On $\{2x<3y\}$: $U=2x$, so $U_x=2$, $U_y=0$. On $\{2x>3y\}$:
        $U=3y$, so $U_x=0$, $U_y=3$. (Non-differentiable on the ray
        $2x=3y$.)

  \item Let $L=\ln(x+y)$, $D=x-y$. By the quotient rule:
        \[
          f_x=\frac{(1/(x+y))(x-y)-\ln(x+y)\cdot1}{(x-y)^2}
            =\frac{1/(x+y)\cdot(x-y)-\ln(x+y)}{(x-y)^2} .
        \]
        Simplifying
        $f_x=\dfrac{(x-y)/(x+y)-\ln(x+y)}{(x-y)^2}$. Similarly,
        \[
          f_y=\frac{(1/(x+y))(x-y)-\ln(x+y)\cdot(-1)}{(x-y)^2}
            =\frac{(x-y)/(x+y)+\ln(x+y)}{(x-y)^2} .
        \]

  \item $F_K=A\alpha K^{\alpha-1}L^\beta$, $F_L=A\beta K^\alpha L^{\beta-1}$.
        So
        \[
          \mathrm{MRTS}_{LK}=\frac{F_L}{F_K}=\frac{\beta}{\alpha}\cdot
            \frac{K}{L} .
        \]
        This depends only on the ratio $K/L$ and the exponent ratio
        $\beta/\alpha$: it is scale-invariant.

  \item $\nabla f=(2x,-2y)^T$. At $(2,-1)$: $\nabla f=(4,2)^T$. Direction
        of steepest ascent: $\mathbf{u}^*=(4,2)/\sqrt{20}=(2,1)/\sqrt5$.
        Maximum rate $=\|\nabla f\|=\sqrt{20}=2\sqrt5$.

  \item $f_x=\sinh(x+y)+\cosh(x-y)$, $f_y=\sinh(x+y)-\cosh(x-y)$.

  \item $\nabla f=(e^{x+y+z},e^{x+y+z},e^{x+y+z})^T$. At $(0,0,0)$:
        $\nabla f=(1,1,1)^T$. Normalise $\mathbf{v}=(1,1,1)$:
        $\|\mathbf{v}\|=\sqrt3$, $\mathbf{u}=(1,1,1)/\sqrt3$. Then
        $D_{\mathbf{u}}f=(1+1+1)/\sqrt3=\sqrt3$.

  \item $f_x=(\ln y)/x$, $f_y=(\ln x)/y$.

  \item $F_K=1+L^{1/2}/K^{1/2}$, $F_L=1+K^{1/2}/L^{1/2}$. From
        $F=(K^{1/2}+L^{1/2})^2$ with $g=K^{1/2}+L^{1/2}$:
        $F_K=g/K^{1/2}=(K^{1/2}+L^{1/2})/K^{1/2}=1+L^{1/2}/K^{1/2}$.
        \checkmark Same for $F_L$.

  \item $f_x=2xy$, $f_y=x^2+3y^2$. At $(-1,2)$: $\nabla f=(-4,1+12)^T
        =(-4,13)^T$. Normalise: $\|\mathbf{v}\|=5$,
        $\mathbf{u}=(4/5,-3/5)^T$. Then
        \[
          D_{\mathbf{u}}f=-4\cdot\frac45+13\cdot\left(-\frac35\right)
            =-\frac{16+39}{5}=-11 .
        \]

  \item $U_x=\alpha/x$, $U_y=\beta/y$.
        $\mathrm{MRS}_{xy}=\dfrac{\alpha/x}{\beta/y}=\dfrac{\alpha y}{\beta x}$.
        At $(\alpha,\beta)$: $\mathrm{MRS}=(\alpha\cdot\beta)/(\beta\cdot\alpha)=1$.

  \item $f=(x^2+y^2)^{-1/2}$. Chain rule:
        \[
          f_x=-\frac12(x^2+y^2)^{-3/2}\cdot2x=\frac{-x}{(x^2+y^2)^{3/2}},
          \qquad f_y=\frac{-y}{(x^2+y^2)^{3/2}} .
        \]
        So $\nabla f=-(x,y)^T/(x^2+y^2)^{3/2}$.

  \item Chain rule with $u=xy$:
        \[
          f_x=\frac{y}{1+x^2y^2}, \qquad f_y=\frac{x}{1+x^2y^2} .
        \]

  \item Let $S=K^\rho+L^\rho$, $F=S^{1/\rho}$. Then
        $F_K=\frac1\rho S^{1/\rho-1}\cdot\rho K^{\rho-1}
        =K^{\rho-1}S^{1/\rho-1}$. Since $S^{1/\rho}=F$, we have $S=F^\rho$
        and so $S^{1/\rho-1}=F^{1-\rho}\cdot F^{-1}\cdot F=F^{1-\rho}$.
        Thus $F_K=K^{\rho-1}F^{1-\rho}=(F/K)^{1-\rho}=(K/F)^{\rho-1}$.
        \checkmark

  \item $f_x=yz$, $f_y=xz$, $f_z=xy$. At $(2,3,-1)$:
        $\nabla f=(3\cdot(-1),2\cdot(-1),2\cdot3)^T=(-3,-2,6)^T$.
        Normalise $\mathbf{v}=(2,-1,2)$: $\|\mathbf{v}\|=3$,
        $\mathbf{u}=(2/3,-1/3,2/3)^T$. Then
        \[
          D_{\mathbf{u}}f=-3\cdot\frac23+(-2)\cdot\left(-\frac13\right)
            +6\cdot\frac23=-2+\frac23+4=2+\frac23=\frac83 .
        \]
\end{solutions}

\section{Additional Problems}

\begin{problems}
  \item Compute $f_x$ and $f_y$ for $f(x,y)=3x^5-2x^3y^2+4y^4$.
  \item For $F(K,L)=20\,K^{1/3}L^{2/3}$, find $F_K$ and $F_L$ at
        $(K,L)=(8,27)$.
  \item Compute $f_x$ and $f_y$ for $f(x,y)=e^{xy}$.
  \item Compute $\nabla f$ for $f(x,y)=2x^2y-3xy^3+5$.
  \item Compute the directional derivative of $f(x,y)=x^2+2xy-y^2$ at
        $(1,2)$ in the direction $\mathbf{v}=(3,4)$.
  \item For $F(K,L)=8K^{3/4}L^{1/4}$, compute $\mathrm{MP}_K$ and
        $\mathrm{MP}_L$ at $(K,L)=(16,16)$.
  \item Compute $f_x$ and $f_y$ for $f(x,y)=\ln(2x+3y)$.
  \item For $U(x,y)=3\ln(x)+4\ln(y)$, derive $\mathrm{MRS}_{xy}$ and
        evaluate at $(2,3)$.
  \item Compute $f_x$ and $f_y$ for $f(x,y)=(x+y^2)^4$.
  \item Compute $\nabla f$ for $f(x,y,z)=xy+yz+zx$.
  \item For a CES production $F(K,L)=5(K^{1/3}+L^{1/3})^3$, find $F_K$.
  \item Compute the directional derivative of $f(x,y)=x^2y$ at $(2,1)$
        in the direction $\mathbf{v}=(1,-1)$.
  \item For $U(x,y)=x^\alpha y^\beta$ (general Cobb--Douglas,
        $\alpha,\beta>0$), derive $\mathrm{MRS}_{xy}$ and show it depends
        only on the ratio $y/x$.
  \item Compute $f_x$ and $f_y$ for $f(x,y)=\dfrac{x-y}{x+y}$.
  \item For $f(K,L)=K^{1/2}L^{1/2}$ at $(K,L)=(16,9)$, find the unit
        direction of steepest ascent and the maximum rate of increase.
  \item Classify the open positive orthant
        $Q=\{(x,y)\in\R^2 : x>0,\ y>0\}$: is it open, closed, or
        neither? Is it bounded? Is it compact?
  \item Classify the budget set
        $B=\{(x,y)\in\R^2 : 2x+3y\le12,\ x\ge0,\ y\ge0\}$: open,
        closed, or neither? Bounded? Compact?
  \item Classify $M=\{(x,y)\in\R^2 : xy\ge1,\ x>0\}$: open, closed, or
        neither? Bounded? Compact? \emph{Hint:} check whether a
        sequence in $M$ can converge to a point with $x=0$.
  \item Classify the open disc
        $D=\{(x,y)\in\R^2 : x^2+y^2<4\}$: open, closed, or neither?
        Bounded? Compact?
  \item Classify (a) the half-open strip
        $S=\{(x,y)\in\R^2 : 0\le x<1\}$ and (b) the finite set
        $E=\{(0,0),\ (1,2),\ (3,1)\}$: open, closed, or neither?
        Bounded? Compact?
  \item Does the Weierstrass theorem guarantee a maximiser for the
        continuous function $u(x,y)=x^{1/2}y^{1/2}$ (a) on the budget
        set $B=\{(x,y): 2x+3y\le12,\ x\ge0,\ y\ge0\}$, (b) on the open
        positive orthant $Q=\{(x,y): x>0,\ y>0\}$? Justify each
        answer.
\end{problems}

\section{Additional Solutions}

\begin{solutions}
  \item $f_x=15x^4-6x^2y^2$, $f_y=-4x^3y+16y^3$.

  \item $F_K=20\cdot\frac13K^{-2/3}L^{2/3}=\frac{20}{3}(L/K)^{2/3}$;
        $F_L=20\cdot\frac23K^{1/3}L^{-1/3}=\frac{40}{3}(K/L)^{1/3}$. At
        $(8,27)$: $(L/K)^{2/3}=(27/8)^{2/3}=9/4$, so $F_K=15$;
        $(K/L)^{1/3}=(8/27)^{1/3}=2/3$, so
        $F_L=\frac{40}{3}\cdot\frac23=\frac{80}{9}$.

  \item $f_x=ye^{xy}$, $f_y=xe^{xy}$.

  \item $\nabla f=(4xy-3y^3,\ 2x^2-9xy^2)^T$.

  \item $f_x=2x+2y$, $f_y=2x-2y$. At $(1,2)$: $\nabla f=(6,-2)^T$.
        Normalise: $\|\mathbf{v}\|=5$, $\mathbf{u}=(3/5,4/5)^T$. Then
        \[
          D_{\mathbf{u}}f=6\cdot\frac35+(-2)\cdot\frac45
            =\frac{18-8}{5}=2 .
        \]

  \item $F_K=8\cdot\frac34K^{-1/4}L^{1/4}=6(L/K)^{1/4}$,
        $F_L=2(K/L)^{3/4}$. At $(16,16)$: $\mathrm{MP}_K=6$,
        $\mathrm{MP}_L=2$.

  \item $f_x=\dfrac{2}{2x+3y}$, $f_y=\dfrac{3}{2x+3y}$.

  \item $U_x=3/x$, $U_y=4/y$.
        $\mathrm{MRS}_{xy}=U_x/U_y=\dfrac{3y}{4x}$. At $(2,3)$:
        $\mathrm{MRS}=\dfrac98$.

  \item Chain: $f_x=4(x+y^2)^3$,
        $f_y=4(x+y^2)^3\cdot2y=8y(x+y^2)^3$.

  \item $f_x=y+z$, $f_y=x+z$, $f_z=x+y$. So
        $\nabla f=(y+z,\ x+z,\ x+y)^T$.

  \item Let $g=K^{1/3}+L^{1/3}$, so $F=5g^3$. Then
        $F_K=5\cdot3g^2\cdot\frac13K^{-2/3}=5g^2K^{-2/3}
        =\dfrac{5(K^{1/3}+L^{1/3})^2}{K^{2/3}}$.

  \item $f_x=2xy$, $f_y=x^2$. At $(2,1)$: $\nabla f=(4,4)^T$. Normalise:
        $\|\mathbf{v}\|=\sqrt2$, $\mathbf{u}=(1/\sqrt2,-1/\sqrt2)^T$.
        Then
        \[
          D_{\mathbf{u}}f=4\cdot\frac{1}{\sqrt2}
            +4\cdot\left(-\frac{1}{\sqrt2}\right)=0 .
        \]
        (Tangent to the level curve.)

  \item $U_x=\alpha x^{\alpha-1}y^\beta$,
        $U_y=\beta x^\alpha y^{\beta-1}$. So
        \[
          \mathrm{MRS}_{xy}=\frac{U_x}{U_y}
            =\frac{\alpha}{\beta}\cdot\frac{y}{x} .
        \]
        Depends only on $y/x$ (and the exponent ratio), not on the
        scale. \checkmark

  \item Quotient rule with $D=x+y$:
        \[
          f_x=\frac{D-(x-y)}{D^2}=\frac{2y}{(x+y)^2}, \qquad
          f_y=\frac{-D-(x-y)}{D^2}=\frac{-2x}{(x+y)^2} .
        \]

  \item $f_K=\frac12K^{-1/2}L^{1/2}$, $f_L=\frac12K^{1/2}L^{-1/2}$. At
        $(16,9)$: $f_K=\frac12\cdot\frac34=3/8$;
        $f_L=\frac12\cdot\frac43=2/3$. So $\nabla f=(3/8,2/3)^T$.
        Magnitude:
        $\|\nabla f\|=\sqrt{9/64+4/9}=\sqrt{81/576+256/576}
        =\sqrt{337/576}=\sqrt{337}/24\approx0.765$. Unit direction of
        steepest ascent:
        $\mathbf{u}^*=(3/8,2/3)^T/(\sqrt{337}/24)=(9,16)^T/\sqrt{337}$.

  \item \textbf{Open.} For any $(x,y)\in Q$ the ball of radius
        $r=\min\{x,y\}>0$ around $(x,y)$ stays inside $Q$, so every
        point is interior. \textbf{Not closed:} the boundary points
        $(x,0)$ and $(0,y)$ with $x,y\ge0$ (the coordinate half-axes,
        including the origin) do not belong to $Q$. \textbf{Not
        bounded:} $Q$ contains $(t,t)$ for every $t>0$, and
        $\|(t,t)\|=t\sqrt2\to\infty$. Hence $Q$ is open, unbounded and
        \textbf{not compact}.

  \item \textbf{Closed.} $B$ is the intersection of the three sets
        $\{2x+3y\le12\}$, $\{x\ge0\}$ and $\{y\ge0\}$, each defined by
        a non-strict inequality on a continuous function; each is
        closed, and an intersection of closed sets is closed (every
        boundary point of $B$ satisfies all three inequalities, so it
        lies in $B$). \textbf{Not open:} $(0,0)\in B$, but every ball
        around the origin contains points with negative coordinates,
        which are outside $B$; so $(0,0)$ is not interior.
        \textbf{Bounded:} from the constraints, $0\le x\le6$ and
        $0\le y\le4$, so $B\subseteq B_8(\mathbf{0})$. Closed and
        bounded in $\R^2$: $B$ is \textbf{compact}.

  \item First note the constraint $x>0$ together with $xy\ge1$ forces
        $y\ge1/x>0$. \textbf{Closed:} suppose
        $(x_n,y_n)\in M$ with $(x_n,y_n)\to(a,b)$. If $a=0$, then
        $y_n\to b$ finite would give $x_ny_n\to0<1$, contradicting
        $x_ny_n\ge1$; so $a>0$, and passing to the limit in
        $x_ny_n\ge1$ gives $ab\ge1$, hence $(a,b)\in M$. $M$ contains
        all its limit (boundary) points, so it is closed: even
        though it is written with a strict inequality, that inequality
        is never active. \textbf{Not open:} $(1,1)\in M$ lies on the
        hyperbola $xy=1$, and every ball around it contains points
        with $xy<1$. \textbf{Not bounded:} $(t,t)\in M$ for all
        $t\ge1$. Hence $M$ is closed, unbounded and \textbf{not
        compact}.

  \item \textbf{Open.} For $(x,y)\in D$ put
        $r=2-\sqrt{x^2+y^2}>0$; the ball $B_r(x,y)$ lies in $D$ by the
        triangle inequality. \textbf{Not closed:} the boundary circle
        $\{x^2+y^2=4\}$ is not contained in $D$ (e.g.\ $(2,0)\notin
        D$). \textbf{Bounded:} $D\subseteq B_3(\mathbf{0})$. Bounded
        but not closed: $D$ is \textbf{not compact}.

  \item (a) \textbf{Neither open nor closed.} The points $(0,y)$
        belong to $S$ but are boundary points (every ball around them
        contains points with $x<0$), so $S$ is not open; the boundary
        points $(1,y)$ do not belong to $S$, so $S$ is not closed.
        \textbf{Unbounded:} $(0,t)\in S$ for all $t$. Hence
        \textbf{not compact}. (b) \textbf{Closed:} around any point
        not in $E$ there is a ball avoiding all three points (a
        finite set always contains all its boundary points:
        indeed $E$ is its own boundary). \textbf{Not open:} no ball
        around $(0,0)$ lies inside $E$. \textbf{Bounded:}
        $E\subseteq B_4(\mathbf{0})$. Closed and bounded: $E$ is
        \textbf{compact}.

  \item (a) \textbf{Yes.} $B$ is compact (closed and bounded, see the
        budget-set problem above) and $u$ is continuous on $B$, so by
        the Weierstrass theorem $u$ attains both a maximum and a
        minimum on $B$; a maximiser exists. (b) \textbf{No.} $Q$ is
        not compact (it is neither closed nor bounded), so Weierstrass
        does not apply: and indeed no maximiser exists:
        $u(t,t)=t\to\infty$ as $t\to\infty$, so $u$ is unbounded above
        on $Q$.
\end{solutions}

\chapter{Curve Sketching: The Complete Discussion}

Everything needed to draw a curve has now been assembled in pieces:
domains and asymptotes in the functions volume, limits in the analysis
volume, the first derivative for monotonicity, the second for curvature.
What is missing is the habit of running all of it in one fixed order and
finishing with a picture. That is the \de{Kurvendiskussion}, and it is worth
drilling as a single routine for two reasons.

First, examinations ask for it directly. Second, and more usefully, the
routine is a self-checking mechanism. A sign table that says the function
falls on $(0,2)$ and a computed local maximum at $x=2$ cannot both be right;
the contradiction surfaces before the picture is drawn. A curve sketch is
the cheapest way there is to catch an algebra slip.

The chapter runs the protocol on rational, exponential, logarithmic and root
functions, and closes with the discussion an economist actually performs:
locating the inflection of a total cost curve, the minimum of average cost,
and the output that maximises profit --- all of them the same routine
applied to a curve with a name.

\section{Core Rules}

\begin{rulebox}{The protocol}
Always in this order, because each step feeds the next:
\begin{enumerate}[leftmargin=1.8em,itemsep=0.35ex]
  \item \textbf{Domain.} Denominators, even roots, logarithms. Note which
        points are excluded --- they are the candidates for poles.
  \item \textbf{Symmetry.} $f(-x)=f(x)$ even, $f(-x)=-f(x)$ odd. Either one
        halves the work.
  \item \textbf{Intercepts and roots.} $f(0)$, and the solutions of
        $f(x)=0$ with their multiplicities.
  \item \textbf{Behaviour at the edges.} Limits at $\pm\infty$ and at every
        excluded point, one side at a time. This produces the
        \textbf{asymptotes}: vertical at poles, horizontal or oblique at
        infinity.
  \item \textbf{First derivative.} Solve $f'(x)=0$ for the critical points,
        build a \textbf{sign table} for $f'$, and read off intervals of
        increase and decrease and the type of each extremum.
  \item \textbf{Second derivative.} Solve $f''(x)=0$, build a sign table
        for $f''$, and read off convex and concave intervals and the
        \textbf{inflection points}.
  \item \textbf{Value table} for any remaining doubt, then \textbf{draw}:
        asymptotes first as dashed guides, then the special points, then
        join them respecting monotonicity and curvature.
\end{enumerate}
\end{rulebox}

\begin{rulebox}{Sign tables}
A sign table is the bookkeeping device that prevents guesswork. List every
point where the expression can change sign --- its roots and its poles ---
in increasing order; they cut the domain into intervals. On each interval
the sign is constant, so \emph{one} test value settles it. For
$f'(x)=3x(x-2)$:
\[
  \begin{array}{c|ccccc}
    x & (-\infty,0) & 0 & (0,2) & 2 & (2,\infty)\\\hline
    f' & + & 0 & - & 0 & +\\
    f & \nearrow & \text{max} & \searrow & \text{min} & \nearrow
  \end{array}
\]
A sign change $+\to-$ is a local maximum, $-\to+$ a local minimum, and
\emph{no} change means no extremum, however much the derivative vanishes
(as at $x=0$ for $f=x^{3}$).
\end{rulebox}

\begin{rulebox}{First- and second-derivative tests}
Let $f'(x_{0})=0$.
\begin{itemize}[leftmargin=1.6em,itemsep=0.3ex]
  \item \textbf{First-derivative test:} the sign change of $f'$ across
        $x_{0}$ decides the type. Always applicable, and the only reliable
        test when $f''(x_{0})=0$.
  \item \textbf{Second-derivative test:} $f''(x_{0})>0$ gives a local
        minimum, $f''(x_{0})<0$ a local maximum. If $f''(x_{0})=0$ the test
        is \textbf{silent} --- fall back on the sign table. Compare
        $x^{4}$ (minimum), $-x^{4}$ (maximum) and $x^{3}$ (neither), all
        three with $f'=f''=0$ at the origin.
  \item \textbf{Inflection} at $x_{0}$ requires $f''$ to \emph{change
        sign} there. $f''(x_{0})=0$ alone is not enough: $f=x^{4}$ has
        $f''(0)=0$ and no inflection.
\end{itemize}
\end{rulebox}

\begin{rulebox}{Asymptotes, recalled}
For $f=p/q$ with $p,q$ polynomials, after cancelling common factors:
vertical asymptote at each zero of $q$; horizontal $y=0$ if
$\deg p<\deg q$; horizontal $y=$ ratio of leading coefficients if
$\deg p=\deg q$; and if $\deg p=\deg q+1$, polynomial division gives
\[
  f(x)=\underbrace{mx+c}_{\text{oblique asymptote}}+\frac{\text{remainder}}{q(x)} .
\]
For non-rational functions the same information comes from limits directly:
$\lim_{x\to\pm\infty}f(x)$, and for a possible oblique asymptote
$m=\lim_{x\to\infty}f(x)/x$ followed by $c=\lim_{x\to\infty}(f(x)-mx)$.
\end{rulebox}

\begin{gallery}
\hspace*{\fill}%
\begin{tikzpicture}
  \begin{axis}[fnaxis, width=8.4cm, height=7.4cm,
      xmin=-4.4, xmax=5.4, ymin=-6.4, ymax=10.4,
      xtick={-4,-2,2,4}, ytick={-4,4,8}]
    \addplot[gray!55, dashed, thin, samples=2, forget plot]
      coordinates {(1,-6.2)(1,10.2)};
    \addplot[gray!55, dashed, thin, samples=2, domain=-4.2:5.2, forget plot]{x+1};
    \addplot[domain=-4.2:0.86, samples=201]{x^2/(x-1)};
    \addplot[domain=1.14:5.2, samples=201]{x^2/(x-1)};
    \addplot[only marks, mark=*, mark size=1.3pt, black, forget plot]
      coordinates {(0,0) (2,4)};
    \node[font=\tiny, anchor=north east] at (axis cs:-0.1,-0.15) {$(0,0)$};
    \node[font=\tiny, anchor=west] at (axis cs:2.15,4) {$(2,4)$};
  \end{axis}
\end{tikzpicture}%
\hspace*{\fill}
\gcap{The finished discussion of $f(x)=x^{2}/(x-1)$: pole $x=1$, oblique
asymptote $y=x+1$, local maximum $(0,0)$ below the local minimum $(2,4)$.}
\end{gallery}

\begin{rulebox}{The economist's curve discussion}
Three facts fall straight out of the protocol applied to a total cost
function $C(q)$, and they are worth knowing as results:
\begin{itemize}[leftmargin=1.6em,itemsep=0.4ex]
  \item Marginal cost is $MC=C'$. The \textbf{inflection point of $C$} is
        exactly where $MC$ is smallest, since $MC'=C''$ changes sign there.
        The S-shaped cost curve turns from concave to convex at its most
        efficient scale.
  \item Average cost is $AC=C(q)/q$. Differentiating,
        \[
          AC'=\frac{C'(q)q-C(q)}{q^{2}}=\frac{MC-AC}{q},
        \]
        so $AC'=0$ exactly when $MC=AC$: \textbf{marginal cost cuts average
        cost at the minimum of average cost}. That is a theorem, not a
        coincidence of a particular example.
  \item Profit at price $p$ is $\pi=pq-C(q)$, so $\pi'=p-MC$ and the
        optimum has $p=MC$, confirmed as a maximum by $\pi''=-C''<0$,
        i.e.\ by cost being convex there.
\end{itemize}
\end{rulebox}

\section{Worked Example}

\begin{worked}{The complete discussion of $f(x)=\dfrac{x^{2}}{x-1}$}
\wstep{Step 1. Domain, and the reason to find it first.}
The denominator vanishes at $x=1$, so
$D_{f}=\R\setminus\{1\}$. Every later step --- sign tables, monotonicity,
convexity --- must be stated \emph{per interval}, and $x=1$ is the wall
between them.

\wstep{Step 2. Symmetry and roots.}
$f(-x)=\frac{x^{2}}{-x-1}$, which is neither $f(x)$ nor $-f(x)$: no
symmetry. The numerator vanishes at $x=0$ with multiplicity two, so the
curve \textbf{touches} the axis at the origin rather than crossing it, and
$f(0)=0$.

\wstep{Step 3. Behaviour at the pole, from both sides.}
Near $x=1$ the numerator is close to $1>0$, so the sign of $f$ is the sign
of $x-1$:
\[
  x\to1^{-}:\ \frac{1^{+}}{0^{-}}\to-\infty,
  \qquad
  x\to1^{+}:\ \frac{1^{+}}{0^{+}}\to+\infty .
\]
\wtrap{Take one-sided limits at every pole. Writing ``$f\to\infty$ at
$x=1$'' loses the sign information and makes the sketch impossible to draw
correctly.}

\wstep{Step 4. Behaviour at infinity: divide, do not guess.}
Numerator degree $2$ exceeds denominator degree $1$ by exactly one, so there
is an \textbf{oblique} asymptote. Polynomial division gives it:
\[
  f(x)=x+1+\frac{1}{x-1} .
\]
So $y=x+1$ is the asymptote, and the remainder $\frac{1}{x-1}$ says how the
curve sits relative to it: \emph{above} for $x>1$, \emph{below} for $x<1$.
\wtrap{Horizontal asymptote when the degrees are equal or the numerator's is
smaller; oblique when the numerator's degree is exactly one larger; neither
when it is larger by two or more. And the division is the whole method ---
guessing $y=x$ from the leading behaviour misses the $+1$.}

\wstep{Step 5. First derivative and critical points.}
Differentiating the divided form is easiest:
\[
  f'(x)=1-\frac{1}{(x-1)^{2}}=\frac{(x-1)^{2}-1}{(x-1)^{2}}=\frac{x(x-2)}{(x-1)^{2}} .
\]
Zeros at $x=0$ and $x=2$; the denominator is positive wherever it is
defined, so the sign of $f'$ is the sign of $x(x-2)$.
\begin{center}\footnotesize
\begin{tabular}{lcccc}
  \hline
  interval & $x<0$ & $0<x<1$ & $1<x<2$ & $x>2$\\
  \hline
  $f'$ & $+$ & $-$ & $-$ & $+$\\
  $f$ & rising & falling & falling & rising\\
  \hline
\end{tabular}
\end{center}

\wstep{Step 6. Classify, and notice something strange.}
$f''(x)=\dfrac{2}{(x-1)^{3}}$, so $f''(0)=-2<0$ and $f''(2)=2>0$:
\[
  \text{local maximum at }(0,0),
  \qquad
  \text{local minimum at }(2,4).
\]
The local maximum lies \textbf{below} the local minimum. That is impossible
for a continuous function --- and possible here only because the pole at
$x=1$ separates them. The two branches never see each other.
\wtrap{Local extrema are local. A local maximum of $0$ coexisting with a
local minimum of $4$ is not an arithmetic error, and neither point is a
global extremum: $f$ is unbounded in both directions.}

\wstep{Step 7. Convexity, and the absent inflection.}
$f''=\frac{2}{(x-1)^{3}}$ is never zero, so there is \textbf{no inflection
point}. It is negative for $x<1$ and positive for $x>1$:
concave on $(-\infty,1)$, convex on $(1,\infty)$. The change of curvature
happens at the pole, where the function does not exist --- so it is not an
inflection.
\wtrap{A change of convexity across a point of the domain is an inflection;
across a gap in the domain it is nothing at all. Check that the candidate
point is in the domain before naming it.}

\wstep{Step 8. Draw it.}
\begin{center}
\begin{tikzpicture}
\begin{axis}[fnaxis, width=8.4cm, height=6.0cm,
             xmin=-4.2, xmax=6.2, ymin=-9, ymax=13,
             xtick={-3,-1,1,2,4}, ytick={4,8},
             restrict y to domain=-40:40]
  \addplot[domain=-4:0.88, samples=200] {x^2/(x-1)};
  \addplot[domain=1.12:6, samples=200] {x^2/(x-1)};
  \addplot[accent!40, thick, dashed, samples=2, domain=-4:6, forget plot]
    {x+1};
  \addplot[gray!55, thin, dashed, samples=2, forget plot]
    coordinates {(1,-9) (1,13)};
  \addplot[only marks, mark=*, mark size=1.2pt, accent!70, forget plot]
    coordinates {(0,0) (2,4)};
\end{axis}
\end{tikzpicture}
\end{center}
Everything checks against the figure: the touch at the origin, the two
branches, the pole at $x=1$, the asymptote $y=x+1$ approached from below on
the left and from above on the right, and the maximum sitting lower than the
minimum.

\wstep{Step 9. The protocol, in the order that saves work.}
\begin{center}\footnotesize
\begin{tabular}{ll}
  \hline
  1. domain & 6. sign table for $f'$\\
  2. symmetry & 7. extrema, classified by $f''$\\
  3. roots and intercepts & 8. $f''$, convexity, inflections\\
  4. poles, one-sided limits & 9. sketch\\
  5. asymptotes (divide!) & 10. check the sketch against 1--8\\
  \hline
\end{tabular}
\end{center}
Step~10 is not optional. A sketch that contradicts the sign table means one
of them is wrong, and finding out which takes a minute --- far less than
losing the marks.
\end{worked}

\section{Practice Problems}

In problems 1--8 give the intervals of increase and decrease and all local
extrema, using a sign table.

\begin{problems}
  \item $f(x)=x^{3}-12x$
  \item $f(x)=x^{4}-8x^{2}$
  \item $f(x)=x\mathrm{e}^{-x}$
  \item $f(x)=\dfrac{x}{x^{2}+1}$
  \item $f(x)=x-\ln x$ on $x>0$
  \item $f(x)=x^{3}$ --- and explain why $f'(0)=0$ yields no extremum.
  \item $f(x)=\dfrac{x^{2}}{x-1}$
  \item $f(x)=\sqrt{x}-x$ on $x\ge0$
\end{problems}

In problems 9--14 give the convex and concave intervals and all inflection
points.

\begin{problems}[resume]
  \item $f(x)=x^{3}-3x^{2}$
  \item $f(x)=x^{4}-4x^{2}$
  \item $f(x)=\mathrm{e}^{-x^{2}}$
  \item $f(x)=x^{4}$ --- and explain why $f''(0)=0$ gives no inflection.
  \item $f(x)=\dfrac{\ln x}{x}$ on $x>0$
  \item $f(x)=x^{2}\mathrm{e}^{-x}$
\end{problems}

In problems 15--20 give every asymptote.

\begin{problems}[resume]
  \item $f(x)=\dfrac{x-1}{x+2}$
  \item $f(x)=\dfrac{1}{x^{2}-4}$
  \item $f(x)=\dfrac{x^{2}}{x-1}$
  \item $f(x)=\dfrac{x^{3}}{x^{2}+1}$
  \item $f(x)=\dfrac{\ln x}{x}$
  \item $f(x)=x+\mathrm{e}^{-x}$
\end{problems}

Problems 21--32 ask for a \emph{complete} discussion: domain, symmetry,
intercepts, asymptotes and edge behaviour, monotonicity and extrema,
curvature and inflections, and a sketch.

\begin{problems}[resume]
  \item $f(x)=x^{3}-3x^{2}$
  \item $f(x)=x^{4}-4x^{2}$
  \item $f(x)=\dfrac{x}{x^{2}+1}$
  \item $f(x)=\dfrac{x^{2}}{x-1}$
  \item $f(x)=\mathrm{e}^{-x^{2}}$
  \item $f(x)=x\mathrm{e}^{-x}$
  \item $f(x)=\dfrac{\ln x}{x}$
  \item $f(x)=\dfrac{x-1}{x+2}$
  \item $f(x)=x^{2}\mathrm{e}^{-x}$
  \item $f(x)=\dfrac{1}{x^{2}-4}$
  \item $f(x)=\sqrt{x}-x$
  \item $f(x)=x^{3}-3x+2$
\end{problems}

\begin{problems}[resume]
  \item A function satisfies $f'>0$ and $f''<0$ throughout an interval.
        Describe the shape of its graph there in one sentence.
  \item The graph of $f'$ is a parabola opening upwards with roots at
        $x=-1$ and $x=3$. What does $f$ do, and where are its extrema?
  \item The graph of $f'$ is a straight line through the origin with
        positive slope. What is $f$, up to a constant, and what is its
        shape?
  \item $f$ is differentiable with $f(1)=2$, $f'(1)=0$ and $f''(1)=-3$.
        Classify $x=1$ and give the tangent line there.
  \item Sketch a single curve with all of: domain $\R\setminus\{2\}$, a
        pole at $x=2$, horizontal asymptote $y=1$, a local minimum at
        $(0,0)$, and convex on $(2,\infty)$.
  \item Why can a function with a vertical asymptote at $x=a$ have a local
        maximum lying below a local minimum? Illustrate with problem 24.
\end{problems}

\begin{problems}[resume]
  \item For $C(q)=q^{3}-6q^{2}+15q$ on $q\ge0$: show $C$ is increasing
        throughout, find its inflection point, and show that marginal cost
        is minimal exactly there.
  \item For the same $C$: find the output minimising average cost and
        verify that $MC=AC$ at that output.
  \item With the same $C$ and price $p=15$, find the profit-maximising
        output and the maximal profit.
  \item For $R(q)=q(200-4q)$: discuss the curve on its economic domain and
        find the revenue-maximising quantity.
  \item For $C(q)=0.5q^{2}+2q+8$: find the minimum of $AC$ and confirm
        $MC=AC$ there.
  \item The adoption path is $y(t)=\dfrac{100}{1+\mathrm{e}^{-0.5t}}$.
        Find its asymptotes, show it is increasing, and locate its
        inflection point.
  \item Explain, using $AC'=(MC-AC)/q$, why average cost must be falling
        wherever marginal cost lies below it.
\end{problems}

\section{Solutions}

\begin{solutions}
  \item $f'=3x^{2}-12=3(x-2)(x+2)$: $+$ on $(-\infty,-2)$, $-$ on
        $(-2,2)$, $+$ on $(2,\infty)$. Local maximum $(-2,16)$, local
        minimum $(2,-16)$.
  \item $f'=4x^{3}-16x=4x(x-2)(x+2)$: $-$, $+$, $-$, $+$ on the four
        intervals cut by $-2,0,2$. Local minima $(\pm2,-16)$, local maximum
        $(0,0)$.
  \item $f'=(1-x)\mathrm{e}^{-x}$: positive for $x<1$, negative for $x>1$.
        Single local (and global) maximum $(1,1/\mathrm{e})$.
  \item $f'=\dfrac{1-x^{2}}{(x^{2}+1)^{2}}$: negative, positive, negative
        on $(-\infty,-1)$, $(-1,1)$, $(1,\infty)$. Local minimum
        $(-1,-\tfrac12)$, local maximum $(1,\tfrac12)$.
  \item $f'=1-\dfrac1x=\dfrac{x-1}{x}$: negative on $(0,1)$, positive on
        $(1,\infty)$. Local (global) minimum $(1,1)$.
  \item $f'=3x^{2}\ge0$ with equality only at $0$: the sign of $f'$ never
        changes, so $f$ increases throughout and $x=0$ is a
        \emph{saddle-like} stationary point, not an extremum. A vanishing
        derivative is necessary, never sufficient.
  \item $f'=\dfrac{x(x-2)}{(x-1)^{2}}$. The denominator is always positive,
        so the sign is that of $x(x-2)$: $+$, $-$, $-$, $+$ across
        $0,1,2$ (with $x=1$ excluded). Local maximum $(0,0)$, local minimum
        $(2,4)$; the function decreases on $(0,1)$ and on $(1,2)$.
  \item $f'=\dfrac{1}{2\sqrt{x}}-1$, zero when $\sqrt{x}=\tfrac12$, i.e.
        $x=\tfrac14$: positive before, negative after. Global maximum
        $\bigl(\tfrac14,\tfrac14\bigr)$.
  \item $f''=6x-6$: negative for $x<1$ (concave), positive for $x>1$
        (convex). Inflection $(1,-2)$.
  \item $f''=12x^{2}-8$: zero at $x=\pm\sqrt{2/3}$; convex outside,
        concave between. Inflections at
        $\bigl(\pm\sqrt{2/3},-\tfrac{20}{9}\bigr)$.
  \item $f'=-2x\mathrm{e}^{-x^{2}}$ and
        $f''=(4x^{2}-2)\mathrm{e}^{-x^{2}}$: zero at
        $x=\pm1/\sqrt2$, convex outside, concave between. Inflections at
        $\bigl(\pm1/\sqrt2,\mathrm{e}^{-1/2}\bigr)$.
  \item $f''=12x^{2}\ge0$, vanishing only at $x=0$ but \emph{not changing
        sign} there: the curve is convex on both sides, so there is no
        inflection. $f''=0$ is necessary, not sufficient.
  \item $f'=\dfrac{1-\ln x}{x^{2}}$, $f''=\dfrac{2\ln x-3}{x^{3}}$: zero at
        $\ln x=\tfrac32$, i.e. $x=\mathrm{e}^{3/2}$; concave before, convex
        after. Inflection at
        $\bigl(\mathrm{e}^{3/2},\tfrac32\mathrm{e}^{-3/2}\bigr)$.
  \item $f'=(2x-x^{2})\mathrm{e}^{-x}$,
        $f''=(x^{2}-4x+2)\mathrm{e}^{-x}$: zero at $x=2\pm\sqrt2$, convex
        outside those, concave between.
  \item Pole $x=-2$; equal degrees, so horizontal asymptote $y=1$.
  \item Poles $x=\pm2$; numerator degree smaller, so $y=0$.
  \item Pole $x=1$; degrees differ by one, and
        $f=x+1+\dfrac{1}{x-1}$, so the oblique asymptote is $y=x+1$.
  \item No real pole ($x^{2}+1>0$); degrees differ by one, and
        $x^{3}=x(x^{2}+1)-x$, so $f=x-\dfrac{x}{x^{2}+1}$ and the oblique
        asymptote is $y=x$.
  \item Vertical asymptote $x=0$ (as $x\to0^{+}$, $f\to-\infty$) and
        horizontal $y=0$ as $x\to\infty$, since $\ln x\ll x$.
  \item As $x\to+\infty$, $\mathrm{e}^{-x}\to0$, so $y=x$ is an oblique
        asymptote. As $x\to-\infty$ the exponential dominates and
        $f\to+\infty$ with no asymptote.
  \item Domain $\R$; no symmetry; $f=x^{2}(x-3)$ so roots $0$ (double,
        touching) and $3$; $f(0)=0$; ends $-\infty$ and $+\infty$. From
        problem 9's table: increasing on $(-\infty,0)$, decreasing on
        $(0,2)$, increasing on $(2,\infty)$, local maximum $(0,0)$, local
        minimum $(2,-4)$; concave then convex with inflection $(1,-2)$. See
        the figure.
  \item Domain $\R$; \emph{even}; $f=x^{2}(x^{2}-4)$ so roots $0$ (double)
        and $\pm2$; both ends $+\infty$; local maximum $(0,0)$, local
        minima $(\pm\sqrt2,-4)$; inflections
        $\bigl(\pm\sqrt{2/3},-\tfrac{20}{9}\bigr)$. A symmetric W.
  \item Domain $\R$; \emph{odd}; single root $0$; horizontal asymptote
        $y=0$ at both ends; local minimum $(-1,-\tfrac12)$, local maximum
        $(1,\tfrac12)$;
        $f''=\dfrac{2x(x^{2}-3)}{(x^{2}+1)^{3}}$ gives inflections at
        $(0,0)$ and $\bigl(\pm\sqrt3,\pm\tfrac{\sqrt3}{4}\bigr)$. See the
        figure.
  \item The worked example of the Core Rules: domain $\R\setminus\{1\}$,
        double root at $0$, pole $x=1$, oblique asymptote $y=x+1$, local
        maximum $(0,0)$, local minimum $(2,4)$, no inflection, concave left
        of the pole and convex right of it.
  \item Domain $\R$; \emph{even}; no roots ($f>0$ always); horizontal
        asymptote $y=0$ at both ends; global maximum $(0,1)$; inflections
        at $\bigl(\pm1/\sqrt2,\mathrm{e}^{-1/2}\bigr)$; convex outside
        them, concave between. The bell curve.
  \item Domain $\R$; no symmetry; single root $0$; as $x\to+\infty$,
        $f\to0^{+}$ (asymptote $y=0$), as $x\to-\infty$, $f\to-\infty$;
        global maximum $(1,1/\mathrm{e})$; inflection
        $(2,2/\mathrm{e}^{2})$, concave before, convex after. See the
        figure.
  \item Domain $(0,\infty)$; root $x=1$; vertical asymptote $x=0$ with
        $f\to-\infty$, horizontal asymptote $y=0$ as $x\to\infty$; global
        maximum $(\mathrm{e},1/\mathrm{e})$; inflection at
        $x=\mathrm{e}^{3/2}$. See the figure.
  \item Domain $\R\setminus\{-2\}$; root $x=1$; $f(0)=-\tfrac12$; pole
        $x=-2$ ($-\infty$ from the left, $+\infty$ from the right);
        horizontal asymptote $y=1$; $f'=\dfrac{3}{(x+2)^{2}}>0$ so strictly
        increasing on each branch with \emph{no} extrema;
        $f''=\dfrac{-6}{(x+2)^{3}}$, convex left of the pole, concave right
        of it.
  \item Domain $\R$; no symmetry; double root $0$; as $x\to+\infty$,
        $f\to0^{+}$; as $x\to-\infty$, $f\to+\infty$; local minimum
        $(0,0)$, local maximum $(2,4/\mathrm{e}^{2})$; inflections at
        $x=2\pm\sqrt2$.
  \item Domain $\R\setminus\{\pm2\}$; \emph{even}; no roots; poles
        $x=\pm2$; horizontal asymptote $y=0$;
        $f'=\dfrac{-2x}{(x^{2}-4)^{2}}$, so increasing for $x<0$ and
        decreasing for $x>0$, giving a local maximum
        $\bigl(0,-\tfrac14\bigr)$ --- a maximum with a negative value,
        since the middle branch lies entirely below the axis.
  \item Domain $[0,\infty)$; roots $0$ and $1$; $f\to-\infty$ as
        $x\to\infty$; global maximum $\bigl(\tfrac14,\tfrac14\bigr)$;
        $f''=-\tfrac14x^{-3/2}<0$, so concave throughout with no
        inflection; the tangent at $0$ is vertical.
  \item Domain $\R$; $f=(x-1)^{2}(x+2)$, so a double root at $1$ (touching)
        and a simple root at $-2$; $f(0)=2$; ends $\mp\infty$; local
        maximum $(-1,4)$, local minimum $(1,0)$; inflection $(0,2)$.
  \item It rises, but ever more slowly: increasing and concave, bending
        away from its own tangent downwards --- the shape of a production
        function with diminishing returns.
  \item $f'<0$ between the roots and $>0$ outside them, so $f$ increases on
        $(-\infty,-1)$, decreases on $(-1,3)$, increases on $(3,\infty)$: a
        local maximum at $x=-1$ and a local minimum at $x=3$.
  \item $f'(x)=kx$ with $k>0$ gives $f(x)=\tfrac{k}{2}x^{2}+c$: an upward
        parabola with vertex at $x=0$, convex everywhere.
  \item $f'(1)=0$ with $f''(1)<0$: a local \textbf{maximum} at $(1,2)$ by
        the second-derivative test. The tangent is horizontal, $y=2$.
  \item One valid answer: $f(x)=1-\dfrac{4}{(x-2)^{2}}$ has the pole and the
        asymptote $y=1$ but a maximum, not a minimum, so adjust ---
        $f(x)=\dfrac{x^{2}}{(x-2)^{2}}$ works: $f(0)=0$ is a local minimum,
        the pole is at $x=2$, and $f\to1$ at both ends. Any sketch with a
        left branch dipping to $(0,0)$ and rising to $+\infty$ at $x=2^{-}$,
        and a right branch descending convexly from $+\infty$ towards
        $y=1$, is correct.
  \item Because a local extremum is only a statement about a
        \emph{neighbourhood}. The pole cuts the domain into two intervals
        with no path between them, so the values on one side need not
        compare with those on the other. In problem 24 the maximum $(0,0)$
        lies on the left branch, the minimum $(2,4)$ on the right, and
        between them the function runs off to $\mp\infty$.
  \item $C'=3q^{2}-12q+15=3(q^{2}-4q+5)=3\bigl((q-2)^{2}+1\bigr)\ge3>0$, so
        $C$ increases for every $q$. $C''=6q-12$ vanishes at $q=2$ and
        changes sign there: inflection at $(2,14)$. Since $MC=C'$ and
        $MC'=C''$, marginal cost has its minimum exactly at $q=2$, where
        $MC=3$.
  \item $AC=q^{2}-6q+15=(q-3)^{2}+6$, minimal at $q=3$ with $AC=6$. There
        $MC=3\cdot9-12\cdot3+15=6=AC$, as the identity
        $AC'=(MC-AC)/q$ requires.
  \item $\pi=15q-(q^{3}-6q^{2}+15q)=-q^{3}+6q^{2}$. Then
        $\pi'=-3q^{2}+12q=-3q(q-4)$, zero at $q=4$ (and $q=0$), and
        $\pi''=-6q+12=-12<0$ at $q=4$: a maximum, with
        $\pi(4)=-64+96=32$.
  \item $R=200q-4q^{2}=-4(q-25)^{2}+2500$ on the economic domain
        $[0,50]$ (where the price $200-4q$ stays non-negative): a downward
        parabola, increasing on $[0,25)$, decreasing on $(25,50]$, concave
        throughout, with maximum revenue $2500$ at $q=25$ and zeros at both
        ends.
  \item $AC=0.5q+2+\dfrac{8}{q}$, so $AC'=0.5-\dfrac{8}{q^{2}}=0$ gives
        $q^{2}=16$, $q=4$, with $AC(4)=2+2+2=6$. And $MC=q+2=6$ at $q=4$:
        equal, as required.
  \item As $t\to\infty$, $y\to100$; as $t\to-\infty$, $y\to0$: two
        horizontal asymptotes. $y'=\dfrac{50\mathrm{e}^{-0.5t}}
        {(1+\mathrm{e}^{-0.5t})^{2}}>0$ always, so $y$ increases
        throughout. The inflection is where the numerator's growth turns,
        at $\mathrm{e}^{-0.5t}=1$, i.e. $t=0$, with $y=50$ --- exactly half
        the ceiling, which is the signature of the logistic curve.
  \item $q>0$, so the sign of $AC'$ is the sign of $MC-AC$. If
        $MC<AC$ then $AC'<0$ and average cost falls: producing one more
        unit costs less than the current average, which pulls the average
        down. Once $MC$ rises above $AC$ the average is pulled up, and the
        crossing point is therefore the minimum.
\end{solutions}

\begin{gallery}
\begin{tikzpicture}
  \begin{axis}[fnaxis, xmin=-1.6, xmax=3.9, ymin=-5.4, ymax=4.4,
      xtick={-1,1,2,3}, ytick={-4,-2,2}, title={$x^{3}-3x^{2}$}]
    \addplot[domain=-1.2:3.6, samples=201]{x^3-3*x^2};
    \addplot[only marks, mark=*, mark size=1.2pt, black, forget plot]
      coordinates {(0,0) (2,-4) (1,-2)};
  \end{axis}
\end{tikzpicture}\hfill
\begin{tikzpicture}
  \begin{axis}[fnaxis, xmin=-5.4, xmax=5.4, ymin=-0.75, ymax=0.75,
      xtick={-4,-2,2,4}, ytick={-0.5,0.5}, title={$x/(x^{2}+1)$}]
    \addplot[domain=-5.2:5.2, samples=201]{x/(x^2+1)};
    \addplot[only marks, mark=*, mark size=1.2pt, black, forget plot]
      coordinates {(1,0.5) (-1,-0.5) (0,0)};
  \end{axis}
\end{tikzpicture}
\end{gallery}

\begin{gallery}
\begin{tikzpicture}
  \begin{axis}[fnaxis, xmin=-1.4, xmax=5.4, ymin=-1.4, ymax=0.75,
      xtick={-1,1,2,4}, ytick={-1,0.5}, title={$x\mathrm{e}^{-x}$}]
    \addplot[domain=-1.2:5.2, samples=201]{x*exp(-x)};
    \addplot[only marks, mark=*, mark size=1.2pt, black, forget plot]
      coordinates {(1,0.3679) (2,0.2707)};
  \end{axis}
\end{tikzpicture}\hfill
\begin{tikzpicture}
  \begin{axis}[fnaxis, xmin=-0.5, xmax=9.4, ymin=-1.9, ymax=0.75,
      xtick={2,4,6,8}, ytick={-1,0.5}, title={$\ln x/x$}]
    \addplot[domain=0.14:9.2, samples=241]{ln(x)/x};
    \addplot[only marks, mark=*, mark size=1.2pt, black, forget plot]
      coordinates {(2.7183,0.3679) (4.4817,0.3347)};
  \end{axis}
\end{tikzpicture}
\gcap{Four of the completed discussions: problems 21, 23, 26 and 27, with
their extrema and inflection points marked.}
\end{gallery}

\section{Additional Problems}

\begin{problems}
  \item Monotonicity and extrema of $f(x)=2x^{3}-9x^{2}+12x$.
  \item Monotonicity and extrema of $f(x)=\dfrac{x^{2}}{x^{2}+3}$.
  \item Monotonicity and extrema of $f(x)=x\ln x$ on $x>0$.
  \item Convexity and inflections of $f(x)=x^{3}-6x^{2}+9x$.
  \item Convexity and inflections of $f(x)=\dfrac{1}{1+x^{2}}$.
  \item Convexity of $f(x)=\mathrm{e}^{x}+\mathrm{e}^{-x}$, and its
        minimum.
  \item All asymptotes of $f(x)=\dfrac{2x^{2}-3}{x+1}$.
  \item All asymptotes of $f(x)=\dfrac{x^{2}-1}{x^{2}+1}$.
  \item All asymptotes of $f(x)=x\mathrm{e}^{1/x}$.
  \item Complete discussion of $f(x)=x^{3}-6x^{2}+9x$.
  \item Complete discussion of $f(x)=\dfrac{x^{2}-4}{x^{2}}$.
  \item Complete discussion of $f(x)=(x-2)\mathrm{e}^{-x}$.
  \item Complete discussion of $f(x)=x^{2}\ln x$ on $x>0$.
  \item Complete discussion of $f(x)=\dfrac{4x}{x^{2}+4}$.
  \item Complete discussion of $f(x)=\sqrt{4-x^{2}}$.
  \item A cubic has a local maximum at $(0,4)$ and a local minimum at
        $(2,0)$. Determine it.
  \item A function has $f'(x)>0$ everywhere and $f''(x)<0$ everywhere, with
        horizontal asymptote $y=5$. Sketch it and give an example formula.
  \item Where does $f(x)=x^{4}-4x^{3}$ have inflections, and why does the
        second-derivative test fail at one of its critical points?
  \item Show that $f(x)=x+1/x$ has a local minimum at $x=1$ and a local
        maximum at $x=-1$, with the maximum value below the minimum value,
        and explain.
  \item For which $a$ does $f(x)=x^{3}+ax$ have two local extrema?
  \item For $C(q)=q^{3}-9q^{2}+30q$: find the inflection, the minimum of
        average cost, and verify $MC=AC$ there.
  \item For $C(q)=2q^{2}+50$ and price $p=40$: find the profit-maximising
        output, the profit, and the break-even outputs.
  \item A monopolist faces $p(q)=100-2q$ with $C(q)=q^{2}+20q$. Discuss the
        profit curve and find the optimum.
  \item Show that if $AC$ has a minimum at $q^{*}$ then the $AC$ curve is
        tangent to no ray from the origin except the one through
        $(q^{*},AC(q^{*}))$, and relate this to $MC=AC$.
  \item The production function $Q(L)=12L^{2}-L^{3}$ describes output from
        labour $L\in[0,12]$. Find where average product is maximal and show
        that marginal product equals average product there.
\end{problems}

\section{Additional Solutions}

\begin{solutions}
  \item $f'=6x^{2}-18x+12=6(x-1)(x-2)$: $+,-,+$. Local maximum $(1,5)$,
        local minimum $(2,4)$.
  \item $f'=\dfrac{6x}{(x^{2}+3)^{2}}$: negative for $x<0$, positive for
        $x>0$. Global minimum $(0,0)$; $f\to1$ at both ends, so no
        maximum.
  \item $f'=\ln x+1$, zero at $x=\mathrm{e}^{-1}$: negative before,
        positive after. Global minimum
        $\bigl(\mathrm{e}^{-1},-\mathrm{e}^{-1}\bigr)$.
  \item $f''=6x-12$: concave for $x<2$, convex for $x>2$; inflection
        $(2,2)$.
  \item $f'=\dfrac{-2x}{(1+x^{2})^{2}}$,
        $f''=\dfrac{6x^{2}-2}{(1+x^{2})^{3}}$: zero at
        $x=\pm1/\sqrt3$, convex outside, concave between. Inflections at
        $\bigl(\pm1/\sqrt3,\tfrac34\bigr)$.
  \item $f''=\mathrm{e}^{x}+\mathrm{e}^{-x}>0$ everywhere: convex, no
        inflection. $f'=\mathrm{e}^{x}-\mathrm{e}^{-x}=0$ at $x=0$, giving
        the global minimum $(0,2)$.
  \item Pole $x=-1$; degrees differ by one, and
        $2x^{2}-3=(x+1)(2x-2)-1$, so $f=2x-2-\dfrac{1}{x+1}$: oblique
        asymptote $y=2x-2$.
  \item No poles; equal degrees, so horizontal asymptote $y=1$ at both
        ends.
  \item As $x\to\pm\infty$, $\mathrm{e}^{1/x}\to1$ and more precisely
        $x\mathrm{e}^{1/x}=x+1+\tfrac{1}{2x}+\cdots$, so $y=x+1$ is an
        oblique asymptote. As $x\to0^{+}$, $f\to+\infty$ (vertical
        asymptote $x=0$); as $x\to0^{-}$, $f\to0$.
  \item $f=x(x-3)^{2}$: roots $0$ and $3$ (double); ends $\mp\infty$;
        $f'=3x^{2}-12x+9=3(x-1)(x-3)$ gives local maximum $(1,4)$ and local
        minimum $(3,0)$; $f''=6x-12$ gives inflection $(2,2)$; no
        symmetry, no asymptotes.
  \item $f=1-\dfrac{4}{x^{2}}$: domain $\R\setminus\{0\}$, \emph{even},
        roots $\pm2$, pole $x=0$ with $f\to-\infty$ from both sides,
        horizontal asymptote $y=1$; $f'=\dfrac{8}{x^{3}}$, so decreasing
        for $x<0$ and increasing for $x>0$ with no extrema;
        $f''=-\dfrac{24}{x^{4}}<0$: concave on both branches.
  \item Domain $\R$; root $x=2$; $f\to0^{+}$ as $x\to+\infty$
        (asymptote $y=0$) and $f\to-\infty$ as $x\to-\infty$;
        $f'=(3-x)\mathrm{e}^{-x}$ gives a maximum
        $(3,\mathrm{e}^{-3})$; $f''=(x-4)\mathrm{e}^{-x}$ gives an
        inflection $(4,2\mathrm{e}^{-4})$.
  \item Domain $(0,\infty)$; root $x=1$; $f\to0$ as $x\to0^{+}$ and
        $f\to\infty$ as $x\to\infty$; $f'=2x\ln x+x=x(2\ln x+1)$, zero at
        $x=\mathrm{e}^{-1/2}$, a minimum with value
        $-\tfrac{1}{2\mathrm{e}}$; $f''=2\ln x+3$, zero at
        $x=\mathrm{e}^{-3/2}$, an inflection.
  \item Domain $\R$; \emph{odd}; root $0$; horizontal asymptote $y=0$;
        $f'=\dfrac{4(4-x^{2})}{(x^{2}+4)^{2}}$ gives local minimum
        $(-2,-1)$ and local maximum $(2,1)$; inflections at $0$ and
        $\pm2\sqrt3$.
  \item Domain $[-2,2]$; \emph{even}; roots $\pm2$; range $[0,2]$; it is
        the upper semicircle of radius $2$, so concave throughout, maximum
        $(0,2)$, and the tangents at $x=\pm2$ are vertical. No asymptotes.
  \item Write $f=ax^{3}+bx^{2}+cx+d$. From the maximum at $(0,4)$:
        $d=4$ and $f'(0)=c=0$. From $f'(2)=12a+4b=0$ and $f(2)=8a+4b+4=0$:
        $b=-3a$ and $8a-12a+4=0$, so $a=1$, $b=-3$. Hence
        $f(x)=x^{3}-3x^{2}+4$.
  \item Increasing and concave with a ceiling: it rises steeply at first
        and flattens onto $y=5$ without reaching it. Example:
        $f(x)=5-\mathrm{e}^{-x}$.
  \item $f'=4x^{3}-12x^{2}=4x^{2}(x-3)$, critical at $x=0$ and $x=3$;
        $f''=12x^{2}-24x=12x(x-2)$, zero at $x=0$ and $x=2$ and changing
        sign at both, so inflections at $(0,0)$ and $(2,-16)$. At $x=0$ the
        second-derivative test fails because $f''(0)=0$; the sign table for
        $f'$ shows $4x^{2}(x-3)<0$ on both sides of $0$, so $x=0$ is no
        extremum --- it is the inflection.
  \item $f'=1-\dfrac{1}{x^{2}}$, zero at $x=\pm1$;
        $f''=\dfrac{2}{x^{3}}$, positive at $x=1$ (minimum, value $2$) and
        negative at $x=-1$ (maximum, value $-2$). The maximum value $-2$
        lies below the minimum value $2$ because the pole at $x=0$
        separates the two branches, so neither extremum is global.
  \item $f'=3x^{2}+a$ has two real roots exactly when $a<0$; then
        $x=\pm\sqrt{-a/3}$ and the sign of $f'$ changes at both, giving a
        maximum and a minimum. For $a\ge0$ the function is monotone
        increasing.
  \item $C''=6q-18=0$ at $q=3$: inflection, and $MC=3q^{2}-18q+30$ is
        minimal there with $MC=3$. $AC=q^{2}-9q+30=(q-4.5)^{2}+9.75$,
        minimal at $q=4.5$ with $AC=9.75$; and
        $MC(4.5)=60.75-81+30=9.75=AC$.
  \item $\pi=40q-2q^{2}-50$, so $\pi'=40-4q=0$ at $q=10$ with
        $\pi''=-4<0$: maximum profit $\pi(10)=400-200-50=150$. Break-even
        needs $\pi=0$: $2q^{2}-40q+50=0$, i.e. $q^{2}-20q+25=0$, so
        $q=10\pm\sqrt{75}$, about $1.34$ and $18.66$.
  \item $\pi=(100-2q)q-(q^{2}+20q)=80q-3q^{2}$, a downward parabola with
        zeros at $q=0$ and $q=80/3$, concave throughout. $\pi'=80-6q=0$ at
        $q=\tfrac{40}{3}\approx13.33$, with
        $\pi=\tfrac{1600}{3}\approx533.3$ and price
        $p=100-\tfrac{80}{3}\approx73.3$.
  \item The slope of the ray from the origin through $(q,C(q))$ is exactly
        $C(q)/q=AC(q)$. Minimising $AC$ therefore means finding the
        \emph{flattest} such ray, and the flattest ray that still meets the
        cost curve must touch it --- be tangent to it. At a tangency the
        ray's slope equals the curve's slope, i.e. $AC=MC$. That is the
        geometric content of $AC'=(MC-AC)/q$.
  \item $AP=Q/L=12L-L^{2}$, maximal where $12-2L=0$, i.e. $L=6$, with
        $AP=36$. And $MP=Q'=24L-3L^{2}$, so $MP(6)=144-108=36=AP$. The
        same tangency argument as for average cost applies, with the ray
        from the origin now touching the production curve.
\end{solutions}

\chapter{Functions of Several Variables: Surfaces and Level Curves}

A function of one variable is understood by looking at its graph. For a
function of two variables there are two pictures, and both are needed. The
\textbf{graph} is a surface in space: mountains, valleys, passes. The
\textbf{contour map} is what a cartographer draws instead --- the level
curves of that landscape, projected flat onto the plane. Economics almost
never draws the surface. It draws the contour map and gives the curves
names: isoquants, indifference curves, iso-profit lines.

The reason to spend a chapter on pictures before computing anything is that
most of multivariable calculus is easier to see than to calculate. Whether a
critical point is a maximum, a minimum or a saddle is visible at a glance in
a contour map. So is whether a function is convex, whether its upper contour
sets are convex, and where it is steep. And the geometry is not decoration:
the Lagrange method is the statement that two families of level curves touch
without crossing, and Kuhn--Tucker is a statement about which boundary of a
shaded region binds.

This chapter therefore does two things. It sets up domains, graphs, level
curves and the sub- and superlevel sets, with the vocabulary in both the
German and the English forms; and it drills reading properties off a picture
in both directions --- picture to property, and formula to picture.

\section{Core Rules}

\begin{rulebox}{Functions of two variables: domain, graph, range}
A function $f\colon D\to\R$ with $D\subseteq\R^{2}$ assigns a number
$f(x,y)$ to each point of a \textbf{region} $D$ in the plane.
\begin{itemize}[leftmargin=1.6em,itemsep=0.3ex]
  \item The \textbf{domain} $D$ is now a subset of the plane, not of the
        line, and the same three obstructions decide it: no zero
        denominator, no negative radicand under an even root, no
        non-positive argument in a logarithm. So
        $f(x,y)=\sqrt{9-x^{2}-y^{2}}$ has $D=\{(x,y):x^{2}+y^{2}\le9\}$, a
        closed disc, and $f(x,y)=\ln(x+y)$ has the open half-plane
        $x+y>0$.
  \item The \textbf{graph} is the surface
        $\{(x,y,f(x,y)) : (x,y)\in D\}$ in $\R^{3}$: above each point of
        the domain, a single height.
  \item The \textbf{range} is the set of heights attained, exactly as
        before.
\end{itemize}
Economic domains are usually the non-negative quadrant $K,L\ge0$, and often
its interior.
\end{rulebox}

\begin{rulebox}{Level curves and the contour map}
For $c\in\R$ the \textbf{level set} \de{(Niveaumenge, Niveaulinie)} is
\[
  N_{c}=\{(x,y)\in D : f(x,y)=c\} ,
\]
the set of points at height $c$. Drawing $N_{c}$ for several equally spaced
values of $c$ gives the \textbf{contour map}. To compute one, set
$f(x,y)=c$ and solve for $y$ --- the level set is a curve in the plane, and
often a familiar one.

Reading a contour map:
\begin{itemize}[leftmargin=1.6em,itemsep=0.3ex]
  \item \textbf{Contours cannot cross}, since a point has one height ---
        with one exception: at a \textbf{saddle} the level curve through
        the saddle point crosses itself.
  \item \textbf{Spacing is steepness.} Contours crowded together mean the
        surface rises fast; widely spaced contours mean gentle ground.
        Steepest ascent runs perpendicular to the contours.
  \item \textbf{Closed nested loops} shrinking to a point mark a
        \textbf{peak} or a \textbf{basin}; which one is decided by whether
        the labels increase or decrease inwards.
  \item A contour that \textbf{pinches and crosses} marks a saddle: two
        directions of ascent and two of descent alternate around it.
\end{itemize}
\end{rulebox}

\begin{rulebox}{Sub- and superlevel sets}
Two \emph{solid} sets accompany each level curve:
\[
  W(c)=\{(x,y): f(x,y)\le c\},
  \qquad
  B(c)=\{(x,y): f(x,y)\ge c\} .
\]
$W(c)$ is the \textbf{sublevel set} \de{(Subniveaumenge, untere
Konturmenge)}, $B(c)$ the \textbf{superlevel set} \de{(Superniveaumenge,
obere Konturmenge)}. Written through a point $P$ rather than a level,
$W(P)=W(f(P))$ and $B(P)=B(f(P))$: everything at most as high as $P$, and
everything at least as high. The same sets appear in the optimization
volume as $L_{c}$ and $U_{c}$ --- lower and upper contour set --- and in
consumer theory $B(P)$ is the \emph{better set} of the bundle $P$.

The link to curvature, drilled later and worth seeing now:
\begin{itemize}[leftmargin=1.6em,itemsep=0.3ex]
  \item $f$ convex $\Longrightarrow$ every $W(c)$ is convex;
  \item $f$ concave $\Longrightarrow$ every $B(c)$ is convex;
  \item every $B(c)$ convex is the \emph{definition} of
        \textbf{quasiconcave}, which is weaker: $f(x,y)=xy$ on the positive
        quadrant is quasiconcave but not concave.
\end{itemize}
\end{rulebox}

\begin{gallery}
\begin{tikzpicture}
  \begin{axis}[geoaxis, xmin=-3.2, xmax=3.2, ymin=-3.2, ymax=3.2,
      xtick={-2,2}, ytick={-2,2}, title={$W(4)$ for $f=x^{2}+y^{2}$}]
    \addplot[draw=none, name path=up, domain=-2:2, samples=201]{sqrt(4-x^2)};
    \addplot[draw=none, name path=lo, domain=-2:2, samples=201]{-sqrt(4-x^2)};
    \region{up}{lo}{}
    \addplot[domain=0:360, samples=181]({2*cos(x)},{2*sin(x)});
  \end{axis}
\end{tikzpicture}\hfill
\begin{tikzpicture}
  \begin{axis}[geoaxis, xmin=-0.4, xmax=8.4, ymin=-0.4, ymax=8.4,
      xtick={2,4,6,8}, ytick={2,4,6,8}, title={$B(4)$ for $f=xy$, $x,y>0$}]
    \addplot[draw=none, name path=hyp, domain=0.5:8, samples=201]{4/x};
    \addplot[draw=none, name path=top, domain=0.5:8, samples=2]{8.4};
    \region{hyp}{top}{}
    \addplot[domain=0.5:8, samples=201]{4/x};
  \end{axis}
\end{tikzpicture}
\gcap{A sublevel set of a convex function and a superlevel set of a
quasiconcave one. Both are convex sets; the second is the \emph{better set}
of a Cobb--Douglas bundle.}
\end{gallery}

\begin{rulebox}{The standard shapes}
Six surfaces cover most of what appears in exercises, and each has a
signature contour map.
\begin{itemize}[leftmargin=1.6em,itemsep=0.3ex]
  \item \textbf{Plane} $f=ax+by+c$: contours are \emph{parallel straight
        lines}, equally spaced. A budget line is one of them.
  \item \textbf{Paraboloid} $f=x^{2}+y^{2}$: a bowl; contours are
        concentric circles of radius $\sqrt{c}$, so they get
        \emph{closer together} outwards --- the bowl steepens.
  \item \textbf{Cone} $f=\sqrt{x^{2}+y^{2}}$: contours are concentric
        circles of radius $c$, \emph{equally spaced} --- constant slope.
  \item \textbf{Saddle} $f=x^{2}-y^{2}$: contours are hyperbolas; the level
        $c=0$ degenerates into the crossing pair $y=\pm x$.
  \item \textbf{Rotated saddle} $f=xy$: hyperbolas again, but with the
        \emph{axes} as asymptotes. Restricted to the positive quadrant this
        is the Cobb--Douglas picture.
  \item \textbf{Parabolic cylinder} $f=x^{2}$: the surface is a trough,
        constant in $y$; contours are \emph{pairs of vertical lines}
        $x=\pm\sqrt{c}$.
\end{itemize}
Circles alone do not identify the surface --- bowl, cone and dome all have
circular contours. The \emph{spacing} and the direction in which the labels
grow settle it.
\end{rulebox}

\begin{rulebox}{Reading curvature off the picture}
\begin{itemize}[leftmargin=1.6em,itemsep=0.3ex]
  \item \textbf{Convex} (bowl-shaped): every chord drawn between two points
        of the surface lies \emph{above} it; no part of the surface bulges
        over a chord. All sublevel sets are convex.
  \item \textbf{Concave} (dome-shaped): every chord lies \emph{below} the
        surface. All superlevel sets are convex.
  \item \textbf{Neither}: a saddle. Some chords lie above, some below.
  \item \textbf{Quasiconcave without being concave}: the superlevel sets
        are still convex --- nested ``better sets'' bulging towards the
        origin --- while the surface itself is not a dome. The Cobb--Douglas
        hill is the standard example.
  \item A \textbf{linear} function is both convex and concave, and both
        quasiconvex and quasiconcave.
\end{itemize}
\end{rulebox}

\begin{rulebox}{Economic readings}
\begin{itemize}[leftmargin=1.6em,itemsep=0.3ex]
  \item An \textbf{isoquant} is a level curve of a production function:
        all input combinations yielding the same output. An
        \textbf{indifference curve} is a level curve of utility.
  \item The slope of a level curve is the \textbf{marginal rate of
        substitution}, computed implicitly as $-f_{x}/f_{y}$
        (Chapter~5). Convex-to-the-origin level curves mean a diminishing
        rate of substitution.
  \item A \textbf{budget set} $\{p_{x}x+p_{y}y\le M,\ x,y\ge0\}$ is the
        sublevel set of a linear function intersected with the quadrant ---
        a triangle, and convex.
  \item The consumer's problem is to reach the highest indifference curve
        that still touches the budget triangle: the contour map plus the
        shaded set is the whole picture.
\end{itemize}
\end{rulebox}

\section{Worked Example}

\begin{worked}{Read the topography of $f(x,y)=x^{2}+4y^{2}$ and of $g(x,y)=x^{2}-y^{2}$: level curves, sub- and superlevel sets, and the surfaces}
\wstep{Step 1. A level curve is what you get by \emph{fixing the height}.}
Set $f(x,y)=c$ and solve for the curve in the plane:
\[
  x^{2}+4y^{2}=c .
\]
For $c>0$ divide by $c$ to reach the standard form of an ellipse:
\[
  \frac{x^{2}}{c}+\frac{y^{2}}{c/4}=1,
  \qquad\text{semi-axes } \sqrt c \text{ and } \tfrac{\sqrt c}{2} .
\]
So the contour map is a family of concentric ellipses, twice as wide as they
are tall. For $c=0$ the level ``curve'' collapses to the single point
$(0,0)$; for $c<0$ it is empty, since $f\ge0$ always.
\wtrap{Check which values of $c$ occur before drawing. Level sets can be
empty, a single point, or several disconnected pieces --- $x^{2}+4y^{2}=-1$
has no solutions, and a level set of $x^{2}-y^{2}$ has two branches.}

\wstep{Step 2. Spacing carries the steepness.}
The contours for $c=1,2,3,4$ have semi-axes $1,\sqrt2,\sqrt3,2$ in $x$. The
gaps shrink as you move out --- $1,\ 0.41,\ 0.32,\ 0.27$ --- so the surface
gets \emph{steeper} away from the origin. Contours that crowd together mean
a steep slope; contours far apart mean a plateau. That is the one skill a
contour map is for.

\wstep{Step 3. Sub- and superlevel sets.}
\[
  \{f\le c\}=\text{the filled ellipse: convex},
  \qquad
  \{f\ge c\}=\text{everything outside it: not convex.}
\]
$f$ is a convex function, so its \textbf{sub}level sets are convex. For a
concave function such as a utility $U$, it is the \textbf{super}level sets
$\{U\ge\bar u\}$ --- the ``at least as good as'' sets --- that are convex,
and that is precisely the convexity of preferences.
\wtrap{Convex function and convex set are different words. Match them
carefully: convex function $\Rightarrow$ sublevel sets convex; concave
function $\Rightarrow$ superlevel sets convex. The direction is chosen by
the inequality, not by the word ``convex''.}

\wstep{Step 4. The surface itself.}
$z=x^{2}+4y^{2}$ is an \textbf{elliptic paraboloid}: a bowl, with its
minimum at the origin, rising in every direction, steeper along $y$ than
along $x$ (the coefficient $4$ against $1$).

\wstep{Step 5. Now $g(x,y)=x^{2}-y^{2}$, where the sign flips.}
\[
  g=c:\quad x^{2}-y^{2}=c ,
\]
a family of \textbf{hyperbolas} --- opening left--right for $c>0$, up--down
for $c<0$. The level set $c=0$ degenerates into the \emph{pair of lines}
$y=\pm x$, which is where the two families change over.

\wstep{Step 6. Two pictures, side by side.}
\begin{center}
\begin{tikzpicture}
\begin{axis}[contouraxis, width=5.2cm, height=5.2cm, xmin=-2.6, xmax=2.6,
             ymin=-2.6, ymax=2.6, title={$x^{2}+4y^{2}=c$}]
  \foreach \c in {0.5,1,2,3,4} {
    \addplot[domain=0:360, samples=161, variable=\t]
      ({sqrt(\c)*cos(\t)}, {sqrt(\c)/2*sin(\t)});
  }
\end{axis}
\end{tikzpicture}
\hfill
\begin{tikzpicture}
\begin{axis}[contouraxis, width=5.2cm, height=5.2cm, xmin=-2.6, xmax=2.6,
             ymin=-2.6, ymax=2.6, title={$x^{2}-y^{2}=c$}]
  \foreach \c in {0.5,2} {
    \addplot[domain=-1.6:1.6, samples=121] ({sqrt(\c)*cosh(x)}, {sqrt(\c)*sinh(x)});
    \addplot[domain=-1.6:1.6, samples=121] ({-sqrt(\c)*cosh(x)}, {sqrt(\c)*sinh(x)});
    \addplot[domain=-1.6:1.6, samples=121] ({sqrt(\c)*sinh(x)}, {sqrt(\c)*cosh(x)});
    \addplot[domain=-1.6:1.6, samples=121] ({sqrt(\c)*sinh(x)}, {-sqrt(\c)*cosh(x)});
  }
  \addplot[gray!60, thin, samples=2, domain=-2.6:2.6, forget plot] {x};
  \addplot[gray!60, thin, samples=2, domain=-2.6:2.6, forget plot] {-x};
\end{axis}
\end{tikzpicture}
\end{center}
\begin{center}
\begin{tikzpicture}
\begin{axis}[surfaxis, title={$z=x^{2}+4y^{2}$}]
  \addplot3[surf, domain=-2:2, y domain=-1:1] {x^2+4*y^2};
\end{axis}
\end{tikzpicture}
\hfill
\begin{tikzpicture}
\begin{axis}[surfaxis, title={$z=x^{2}-y^{2}$}]
  \addplot3[surf, domain=-2:2, y domain=-2:2] {x^2-y^2};
\end{axis}
\end{tikzpicture}
\end{center}

\wstep{Step 7. Read the origin off each map.}
For $f$ the contours \emph{close around} the origin, so it is a minimum: go
any direction and you climb. For $g$ the contours \emph{cross} at the
origin, and the crossing is the signature of a \textbf{saddle}: along the
$x$-axis $g=x^{2}$ rises, along the $y$-axis $g=-y^{2}$ falls. The point is
a minimum in one direction and a maximum in another, so it is neither.
\wtrap{Closed loops shrinking to a point mean an extremum; crossing curves
mean a saddle. Checking only one direction --- for instance only along the
axes of $g$ where the answers disagree --- is exactly how a saddle gets
misreported as an optimum in Volume~7.}

\wstep{Step 8. The economic vocabulary.}
Level curves of a utility function are \textbf{indifference curves}; of a
production function, \textbf{isoquants}; of a cost function,
\textbf{isocost lines}. Their slope is the marginal rate of substitution,
computed in Chapter~5 by implicit differentiation, and the standard
``convex to the origin'' shape is exactly the convexity of the superlevel
sets in Step~3.
\end{worked}

\section{Practice Problems}

State the domain of each function as a region in the plane, in problems
1--8, and describe it in words.

\begin{problems}
  \item $f(x,y)=\sqrt{16-x^{2}-y^{2}}$
  \item $f(x,y)=\ln(x-y)$
  \item $f(x,y)=\dfrac{1}{x^{2}+y^{2}-1}$
  \item $f(x,y)=\dfrac{xy}{x-y}$
  \item $f(x,y)=\sqrt{x}+\sqrt{y}$
  \item $f(x,y)=\ln\bigl(xy\bigr)$
  \item $f(x,y)=\sqrt{x^{2}+y^{2}-4}$
  \item $f(K,L)=K^{0.4}L^{0.6}$, and state the economic domain too.
\end{problems}

In problems 9--20, give the level set $N_{c}$ for the stated level, solved
for $y$ where possible, and name the curve.

\begin{problems}[resume]
  \item $f(x,y)=x+2y$, \;$c=4$.
  \item $f(x,y)=x^{2}+y^{2}$, \;$c=9$.
  \item $f(x,y)=x^{2}+y^{2}$, \;$c=0$.
  \item $f(x,y)=x^{2}+4y^{2}$, \;$c=4$.
  \item $f(x,y)=xy$, \;$c=6$.
  \item $f(x,y)=x^{2}-y^{2}$, \;$c=0$.
  \item $f(x,y)=x^{2}-y^{2}$, \;$c=4$.
  \item $f(x,y)=y-x^{2}$, \;$c=2$.
  \item $f(x,y)=\sqrt{x^{2}+y^{2}}$, \;$c=3$.
  \item $f(x,y)=\mathrm{e}^{x+y}$, \;$c=1$.
  \item $f(K,L)=\sqrt{KL}$, \;$c=10$.
  \item $f(x,y)=2x+3y$, \;$c=12$; and say what this curve is called when
        $f$ is a cost.
\end{problems}

Problems 21--28 refer to the contour map below, which shows the level
curves of $f(x,y)=x^{3}-3x+y^{2}$ at the levels $c=-1,0,2,4$, together with
six marked points.

\begin{gallery}
\hspace*{\fill}%
\begin{tikzpicture}
  \begin{axis}[contouraxis, width=8.2cm, height=8.2cm,
      xmin=-2.6, xmax=2.6, ymin=-3.2, ymax=3.2,
      xtick={-2,-1,1,2}, ytick={-2,-1,1,2}]
    % c = 4
    \addplot[domain=-2.2:2.196, samples=201]{ sqrt(4-x^3+3*x)};
    \addplot[domain=-2.2:2.196, samples=201]{-sqrt(4-x^3+3*x)};
    % c = 2 (through the saddle at (-1,0))
    \addplot[domain=-2.2:2, samples=201]{ sqrt(2-x^3+3*x)};
    \addplot[domain=-2.2:2, samples=201]{-sqrt(2-x^3+3*x)};
    % c = 0, two pieces
    \addplot[domain=-2.2:-1.7321, samples=101]{ sqrt(-x^3+3*x)};
    \addplot[domain=-2.2:-1.7321, samples=101]{-sqrt(-x^3+3*x)};
    \addplot[domain=0:1.7321, samples=101]{ sqrt(-x^3+3*x)};
    \addplot[domain=0:1.7321, samples=101]{-sqrt(-x^3+3*x)};
    % c = -1, two pieces
    \addplot[domain=-2.2:-1.8794, samples=61]{ sqrt(-1-x^3+3*x)};
    \addplot[domain=-2.2:-1.8794, samples=61]{-sqrt(-1-x^3+3*x)};
    \addplot[domain=0.3473:1.5321, samples=101]{ sqrt(-1-x^3+3*x)};
    \addplot[domain=0.3473:1.5321, samples=101]{-sqrt(-1-x^3+3*x)};
    % marked points
    \addplot[only marks, mark=*, mark size=1.4pt, black, forget plot]
      coordinates {(0,0) (1,0) (-1,0) (2,1) (0,2) (-2,1)};
    \node[font=\tiny, anchor=north east] at (axis cs:-0.06,-0.12) {$P_{1}$};
    \node[font=\tiny, anchor=north west] at (axis cs:1.08,-0.12) {$P_{2}$};
    \node[font=\tiny, anchor=south east] at (axis cs:-1.12,0.12) {$P_{3}$};
    \node[font=\tiny, anchor=south west] at (axis cs:2.05,1.05) {$P_{4}$};
    \node[font=\tiny, anchor=south west] at (axis cs:0.05,2.05) {$P_{5}$};
    \node[font=\tiny, anchor=south east] at (axis cs:-2.05,1.05) {$P_{6}$};
  \end{axis}
\end{tikzpicture}%
\hspace*{\fill}
\gcap{Level curves of $f(x,y)=x^{3}-3x+y^{2}$ at $c=-1,0,2,4$.}
\end{gallery}

\begin{problems}[resume]
  \item Give $f(P_{i})$ for all six marked points
        $P_{1}=(0,0)$, $P_{2}=(1,0)$, $P_{3}=(-1,0)$, $P_{4}=(2,1)$,
        $P_{5}=(0,2)$, $P_{6}=(-2,1)$.
  \item Which of the six points lie on one and the same level curve?
  \item One of the drawn contours crosses itself. At which point, and what
        does that tell you about the surface there?
  \item One of the drawn contours has shrunk to a single point. Which
        level, which point, and what feature of the surface is it?
  \item From the picture alone, name the two critical points of $f$ and
        classify each as a local maximum, a local minimum, or a saddle.
  \item Confirm the previous answer by computation: solve
        $f_{x}=f_{y}=0$.
  \item Is the sublevel set $W(-1)$ connected? Answer from the picture and
        explain what the two pieces mean.
  \item Moving away from $P_{2}$, in which direction does $f$ increase
        fastest? Answer qualitatively from the contours, then check with
        $\grad f$ at a nearby point, say $(1.5,0.5)$.
\end{problems}

\begin{problems}[resume]
  \item Match each surface (a)--(f) below to its contour map (1)--(6).
  \item For each of the six functions in the previous problem, say whether
        it is convex, concave, both, or neither on all of $\R^{2}$.
\end{problems}

\begin{gallery}
\begin{tikzpicture}
  \begin{axis}[surfaxis, title={(a)}]
    \addplot3[surf, domain=-2:2, y domain=-2:2]{x+y};
  \end{axis}
\end{tikzpicture}\hfill
\begin{tikzpicture}
  \begin{axis}[surfaxis, title={(b)}, z post scale=1.2]
    \addplot3[surf, domain=-2:2, y domain=-2:2]{x^2+y^2};
  \end{axis}
\end{tikzpicture}\hfill
\begin{tikzpicture}
  \begin{axis}[surfaxis, title={(c)}, z post scale=1.2]
    \addplot3[surf, domain=-2:2, y domain=-2:2]{sqrt(x^2+y^2)};
  \end{axis}
\end{tikzpicture}
\end{gallery}

\begin{gallery}
\begin{tikzpicture}
  \begin{axis}[surfaxis, title={(d)}, z post scale=1.2]
    \addplot3[surf, domain=-2:2, y domain=-2:2]{x^2-y^2};
  \end{axis}
\end{tikzpicture}\hfill
\begin{tikzpicture}
  \begin{axis}[surfaxis, title={(e)}, z post scale=1.2, view={105}{32}]
    \addplot3[surf, domain=-2:2, y domain=-2:2]{x*y};
  \end{axis}
\end{tikzpicture}\hfill
\begin{tikzpicture}
  \begin{axis}[surfaxis, title={(f)}, z post scale=1.2]
    \addplot3[surf, domain=-2:2, y domain=-2:2]{x^2};
  \end{axis}
\end{tikzpicture}
\end{gallery}

\begin{gallery}
\begin{tikzpicture}
  \begin{axis}[contouraxis, xmin=-2.4, xmax=2.4, ymin=-2.4, ymax=2.4,
      xtick={-2,2}, ytick={-2,2}, title={(1)}]
    \foreach \c in {-2,-1,0,1,2} { \addplot[domain=-2.3:2.3]{\c-x}; }
  \end{axis}
\end{tikzpicture}\hfill
\begin{tikzpicture}
  \begin{axis}[contouraxis, xmin=-2.4, xmax=2.4, ymin=-2.4, ymax=2.4,
      xtick={-2,2}, ytick={-2,2}, title={(2)}]
    \foreach \r in {0.25,0.5,0.75,1,1.25,1.5,1.75,2} {
      \addplot[domain=0:360, samples=121]({\r*cos(x)},{\r*sin(x)}); }
  \end{axis}
\end{tikzpicture}\hfill
\begin{tikzpicture}
  \begin{axis}[contouraxis, xmin=-2.4, xmax=2.4, ymin=-2.4, ymax=2.4,
      xtick={-2,2}, ytick={-2,2}, title={(3)}]
    \foreach \r in {0.707,1,1.225,1.414,1.581,1.732,1.871,2} {
      \addplot[domain=0:360, samples=121]({\r*cos(x)},{\r*sin(x)}); }
  \end{axis}
\end{tikzpicture}
\end{gallery}

\begin{gallery}
\begin{tikzpicture}
  \begin{axis}[contouraxis, xmin=-2.4, xmax=2.4, ymin=-2.4, ymax=2.4,
      xtick={-2,2}, ytick={-2,2}, title={(4)}]
    \addplot[domain=-2.3:2.3]{x};
    \addplot[domain=-2.3:2.3]{-x};
    \foreach \s in {1,-1} {
      \addplot[domain=-1.4:1.4, samples=81]({\s*cosh(x)},{sinh(x)});
      \addplot[domain=-1.4:1.4, samples=81]({sinh(x)},{\s*cosh(x)}); }
  \end{axis}
\end{tikzpicture}\hfill
\begin{tikzpicture}
  \begin{axis}[contouraxis, xmin=-2.4, xmax=2.4, ymin=-2.4, ymax=2.4,
      xtick={-2,2}, ytick={-2,2}, title={(5)}]
    \foreach \c in {1,2} {
      \addplot[domain=0.42:2.3, samples=101]{\c/x};
      \addplot[domain=-2.3:-0.42, samples=101]{\c/x};
      \addplot[domain=0.42:2.3, samples=101]{-\c/x};
      \addplot[domain=-2.3:-0.42, samples=101]{-\c/x}; }
  \end{axis}
\end{tikzpicture}\hfill
\begin{tikzpicture}
  \begin{axis}[contouraxis, xmin=-2.4, xmax=2.4, ymin=-2.4, ymax=2.4,
      xtick={-2,2}, ytick={-2,2}, title={(6)}]
    \foreach \v in {-2,-1.414,-1,1,1.414,2} {
      \addplot[samples=2] coordinates {(\v,-2.3) (\v,2.3)}; }
  \end{axis}
\end{tikzpicture}
\end{gallery}

In problems 31--38, describe the stated sub- or superlevel set as a region,
and say whether it is convex.

\begin{problems}[resume]
  \item $W(4)$ for $f(x,y)=x^{2}+y^{2}$.
  \item $B(4)$ for $f(x,y)=x^{2}+y^{2}$.
  \item $W(6)$ for $f(x,y)=2x+3y$.
  \item $B(4)$ for $f(x,y)=xy$ restricted to $x,y>0$.
  \item $W(0)$ for $f(x,y)=x^{2}-y^{2}$.
  \item $B(1)$ for $f(x,y)=\sqrt{x^{2}+y^{2}}$.
  \item $W(P)$ for $f(x,y)=x^{2}+y^{2}$ with $P=(3,4)$.
  \item $B(P)$ for $f(K,L)=\sqrt{KL}$ with $P=(4,9)$, on $K,L>0$.
\end{problems}

\begin{problems}[resume]
  \item Sketch the level sets $W(2)$, $B(2)$, $W(-1)$ and $B(-1)$ of the
        one-variable function $g(x)=x^{3}-3x$ as subsets of the $x$-axis,
        and say which of the four are convex sets.
  \item Explain why for a function of one variable a level set
        $W(c)$ is convex exactly when it is an interval.
  \item A firm's isoquants are the curves $KL=\bar Q$. Sketch three of them
        for $\bar Q=1,4,9$ on $K,L>0$, and say whether the isoquants get
        closer together or further apart as output rises along the ray
        $K=L$.
  \item For the budget set $\{2x+5y\le100,\ x,y\ge0\}$: give the two
        intercepts, sketch the set, and say why it is convex.
  \item The indifference curves of $U(x,y)=xy$ and the budget line
        $2x+5y=100$ are drawn on the same axes. Describe, in words and
        without computing, the geometric condition that locates the optimal
        bundle.
  \item Two contour maps show circles as level curves. In the first the
        circles are equally spaced; in the second they crowd together
        outwards. Which is the cone and which the paraboloid, and why?
  \item A contour map has no closed loops and no crossings anywhere, only
        parallel straight lines. What can be said about the function?
\end{problems}

\section{Solutions}

\begin{solutions}
  \item $16-x^{2}-y^{2}\ge0\iff x^{2}+y^{2}\le16$: the closed disc of
        radius $4$ about the origin.
  \item $x-y>0$: the open half-plane strictly below the line $y=x$.
  \item $x^{2}+y^{2}\neq1$: the whole plane except the unit circle.
  \item $x\neq y$: the plane without the line $y=x$.
  \item $x\ge0$ and $y\ge0$: the closed non-negative quadrant.
  \item $xy>0$: the open first and third quadrants (the axes excluded).
  \item $x^{2}+y^{2}\ge4$: everything outside or on the circle of radius
        $2$ --- the complement of an open disc.
  \item The exponents are not integers, so $K\ge0$ and $L\ge0$: the closed
        non-negative quadrant. Economically the same set, usually taken
        open ($K,L>0$) so that marginal products exist.
  \item $y=\dfrac{4-x}{2}=2-\tfrac12x$: a straight line of slope
        $-\tfrac12$. All level sets are parallel lines.
  \item $x^{2}+y^{2}=9$, i.e. $y=\pm\sqrt{9-x^{2}}$: the circle of radius
        $3$.
  \item $x^{2}+y^{2}=0$ forces $x=y=0$: the level set is the single point
        $(0,0)$, the bottom of the bowl.
  \item $\dfrac{x^{2}}{4}+y^{2}=1$, so $y=\pm\sqrt{1-x^{2}/4}$: an ellipse
        with semi-axes $2$ and $1$.
  \item $y=6/x$: a hyperbola in the first and third quadrants, with both
        axes as asymptotes.
  \item $x^{2}=y^{2}\iff y=\pm x$: the \emph{two crossing lines} through
        the origin. This degenerate level is the signature of a saddle.
  \item $x^{2}-y^{2}=4$, so $y=\pm\sqrt{x^{2}-4}$, defined for
        $|x|\ge2$: a hyperbola opening left and right.
  \item $y=x^{2}+2$: a parabola, the standard parabola shifted up by $2$.
        All level sets are vertical translates of it.
  \item $\sqrt{x^{2}+y^{2}}=3\iff x^{2}+y^{2}=9$: the circle of radius
        $3$. Note that for the cone the radius equals $c$, not
        $\sqrt{c}$.
  \item $\mathrm{e}^{x+y}=1\iff x+y=0$, so $y=-x$: a straight line. Every
        level set of $\mathrm{e}^{x+y}$ is a line $x+y=\ln c$, defined only
        for $c>0$.
  \item $\sqrt{KL}=10\iff KL=100\iff L=100/K$: a hyperbola --- the
        isoquant for output $10$.
  \item $y=\dfrac{12-2x}{3}=4-\tfrac23x$: a straight line. For a cost
        function this is an \textbf{isocost line}; the analogous level
        curve of a budget constraint is the budget line.
  \item $f(x,y)=x^{3}-3x+y^{2}$ gives
        $f(P_{1})=0$, \;$f(P_{2})=1-3+0=-2$, \;$f(P_{3})=-1+3+0=2$,
        \;$f(P_{4})=8-6+1=3$, \;$f(P_{5})=0-0+4=4$,
        \;$f(P_{6})=-8+6+1=-1$.
  \item None of them share a level: the six values
        $0,-2,2,3,4,-1$ are distinct. $P_{3}$ lies on the drawn contour
        $c=2$, $P_{5}$ on $c=4$ and $P_{6}$ on $c=-1$; $P_{1}$, $P_{2}$ and
        $P_{4}$ lie between drawn contours.
  \item The contour $c=2$ crosses itself at $P_{3}=(-1,0)$. A
        self-crossing level curve marks a \textbf{saddle point}: from
        $(-1,0)$ the surface rises in two opposite directions and falls in
        the two others.
  \item The level $c=-2$ is a single point, $(1,0)$. Nested contours
        shrinking to a point mark a \textbf{local extremum}; since the
        levels decrease inwards, it is a local \emph{minimum}, with value
        $-2$.
  \item From the picture: a local minimum at $(1,0)$ (contours closing in
        on it) and a saddle at $(-1,0)$ (the crossing contour). There is no
        local maximum --- no contours close in on a point with the labels
        increasing inwards.
  \item $f_{x}=3x^{2}-3=0\Rightarrow x=\pm1$; $f_{y}=2y=0\Rightarrow y=0$.
        Critical points $(1,0)$ and $(-1,0)$. With $f_{xx}=6x$,
        $f_{yy}=2$, $f_{xy}=0$: at $(1,0)$ the determinant is $12>0$ with
        $f_{xx}>0$, a local minimum; at $(-1,0)$ it is $-12<0$, a saddle.
  \item No: $W(-1)$ has two pieces --- a small oval around the minimum
        $(1,0)$ and an unbounded piece on the far left. A sublevel set of a
        non-convex function may fall apart into components, which is
        exactly what convexity of $f$ would forbid.
  \item Away from $P_{2}=(1,0)$ the contours are closed ovals around it, so
        $f$ increases in \emph{every} direction, fastest where the ovals
        are most tightly packed --- to the right. Checking,
        $\grad f=(3x^{2}-3,\,2y)$, so at $(1.5,0.5)$,
        $\grad f=(3.75,\,1)$: pointing right and slightly up, and
        perpendicular to the contour there.
  \item (a)$\to$(1) plane, parallel lines. \;(b)$\to$(3) paraboloid:
        circles of radius $\sqrt{c}$, crowding outwards.
        \;(c)$\to$(2) cone: circles of radius $c$, equally spaced.
        \;(d)$\to$(4) saddle $x^{2}-y^{2}$: hyperbolas with the diagonals
        $y=\pm x$ as the degenerate level. \;(e)$\to$(5) $xy$: hyperbolas
        with the \emph{axes} as asymptotes. \;(f)$\to$(6) parabolic
        cylinder: pairs of vertical lines.
  \item (a) linear, hence \textbf{both} convex and concave. (b) convex (a
        bowl). (c) convex (a cone opening upwards; not differentiable at
        the origin, but still convex). (d) neither --- a saddle.
        (e) neither on $\R^{2}$; on the positive quadrant it is
        quasiconcave. (f) convex: $x^{2}$ is convex and adding no
        $y$-dependence preserves it.
  \item The closed disc $x^{2}+y^{2}\le4$ of radius $2$. Convex.
  \item The complement of the open disc, $x^{2}+y^{2}\ge4$: everything
        outside the circle. \textbf{Not} convex --- the segment joining
        $(-2,0)$ to $(2,0)$ passes through the origin, which is not in the
        set.
  \item The half-plane $2x+3y\le6$, bounded by the line through $(3,0)$ and
        $(0,2)$. Convex: every half-plane is.
  \item $\{(x,y): x,y>0,\ xy\ge4\}$, the region on the far side of the
        hyperbola $y=4/x$. Convex --- this is precisely the statement that
        $xy$ is quasiconcave on the quadrant.
  \item $x^{2}-y^{2}\le0\iff|y|\ge|x|$: the two opposite ``vertical'' sectors
        between the lines $y=\pm x$. Not convex: $(0,1)$ and $(0,-1)$ lie
        in it, but the midpoint $(0,0)$ is on the boundary and points just
        off the segment leave the set --- more simply, the set is not even
        connected in the way a convex set must be, since joining
        $(0,1)$ to $(1,2)$ stays inside but joining $(0,1)$ to $(0,-1)$
        passes through the pinch at the origin only.
  \item $\sqrt{x^{2}+y^{2}}\ge1$: everything outside or on the unit circle.
        Not convex, for the same reason as problem 32.
  \item $f(P)=9+16=25$, so $W(P)=\{x^{2}+y^{2}\le25\}$, the closed disc of
        radius $5$ through $P$. Convex.
  \item $f(P)=\sqrt{36}=6$, so
        $B(P)=\{(K,L):K,L>0,\ \sqrt{KL}\ge6\}=\{KL\ge36\}$: the region
        beyond the isoquant $L=36/K$. Convex.
  \item $g(x)=x^{3}-3x$ has $g(2)=2$, $g(-1)=2$, $g(1)=-2$, $g(-2)=-2$.
        $W(2)=\{x\le2\}$ minus nothing --- solving $x^{3}-3x\le2$ gives
        $(x+1)^{2}(x-2)\le0$, i.e. $x\le2$: the ray $(-\infty,2]$,
        \textbf{convex}. $B(2)=[2,\infty)\cup\{-1\}$: the ray from $2$
        \emph{plus the isolated point} $x=-1$, \textbf{not convex}.
        $W(-1)$: solving $x^{3}-3x\le-1$ gives
        $(-\infty,-1.879]\cup[0.347,1.532]$, two intervals, \textbf{not
        convex}. $B(-1)=[-1.879,0.347]\cup[1.532,\infty)$, also two
        pieces, \textbf{not convex}.
  \item In $\R^{1}$ a set is convex exactly when it contains the whole
        segment between any two of its points, and a subset of the line
        with that property is precisely an interval (possibly unbounded or
        empty). So ``convex'' and ``interval'' are the same statement one
        dimension down --- which is why the one-variable picture is the
        right place to first meet the definition.
  \item $L=\bar Q/K$: three hyperbola branches, $L=1/K$, $L=4/K$ and
        $L=9/K$. Along the ray $K=L$ they cut it at
        $K=1,2,3$: the isoquants for equally spaced \emph{square roots} of
        output, so for equally spaced output levels they move
        \emph{further apart}. Rising output requires ever larger input
        increases --- diminishing returns along the ray.
  \item Intercepts $(50,0)$ and $(0,20)$. The set is the triangle with
        vertices $(0,0)$, $(50,0)$, $(0,20)$. It is convex because it is an
        intersection of three half-planes ($2x+5y\le100$, $x\ge0$,
        $y\ge0$), and any intersection of convex sets is convex.
  \item The optimum is where the budget line is \textbf{tangent} to an
        indifference curve: the highest curve of the family that still
        touches the triangle touches it at exactly one point, and there the
        two curves have the same slope. Crossing would mean a higher curve
        is still reachable; missing the set entirely means the bundle is
        unaffordable.
  \item Equally spaced circles are the \textbf{cone}
        $\sqrt{x^{2}+y^{2}}$, whose level $c$ is a circle of radius $c$ ---
        constant slope, so equal height steps need equal horizontal steps.
        Circles crowding outwards are the \textbf{paraboloid}
        $x^{2}+y^{2}$, level $c$ having radius $\sqrt{c}$: the surface gets
        steeper, so equal height steps need ever smaller horizontal steps.
  \item Parallel straight contours, equally spaced or not, mean the
        function depends on the two variables only through one linear
        combination: $f(x,y)=h(ax+by)$ for some one-variable $h$. If the
        spacing is also constant, $h$ is linear and the surface is a plane.
        There are no critical points of extremum type anywhere.
\end{solutions}

\section{Additional Problems}

\begin{problems}
  \item Domain of $f(x,y)=\dfrac{\sqrt{y-x^{2}}}{x}$, as a region.
  \item Domain of $f(x,y)=\ln(4-x^{2}-y^{2})$, as a region.
  \item Domain and range of $f(x,y)=\mathrm{e}^{-(x^{2}+y^{2})}$.
  \item Level set of $f(x,y)=x^{2}+y^{2}-4x$ at $c=0$; name and locate the
        curve.
  \item Level set of $f(x,y)=x^{2}+y^{2}-4x$ at $c=-4$.
  \item Level set of $f(x,y)=y\,\mathrm{e}^{x}$ at $c=1$.
  \item Level set of $f(x,y)=\dfrac{y}{x}$ at $c=2$, and why the origin is
        excluded.
  \item Level sets of $f(x,y)=\max(x,y)$ at $c=1$: sketch and describe.
  \item For $f(x,y)=x^{2}+4y^{2}$, sketch the contours $c=4,16,36$ and say
        whether they crowd inwards or outwards.
  \item For $f(x,y)=x-y^{2}$, give the contours at $c=-1,0,1$ and describe
        the family.
  \item Give a function whose level curves are concentric circles that
        crowd \emph{inwards} towards the origin.
  \item Give a function of two variables whose level sets are all straight
        lines through the origin.
  \item $W(1)$ and $B(1)$ for $f(x,y)=|x|+|y|$: describe both and say which
        is convex.
  \item $W(0)$ for $f(x,y)=xy$ on all of $\R^{2}$: describe it and decide
        convexity.
  \item $B(2)$ for $f(x,y)=x+y$: describe it and decide convexity.
  \item $W(4)$ for $f(x,y)=x^{2}$: describe it and decide convexity.
  \item For $f(x,y)=x^{2}+y^{2}$ and $P=(1,2)$, give $W(P)$ and $B(P)$ and
        the level curve through $P$.
  \item A function has convex sublevel sets \emph{and} convex superlevel
        sets at every level. What kind of function can it be? Give an
        example.
  \item Sketch the surface $f(x,y)=4-x^{2}-y^{2}$ in words, give its
        contours, and classify it as convex or concave.
  \item Two functions have exactly the same family of level curves but
        different surfaces. Give such a pair and explain how the contour
        maps can nevertheless be told apart.
  \item Classify $f(x,y)=x^{2}+y^{2}$, $g(x,y)=-x^{2}-y^{2}$,
        $h(x,y)=x^{2}-y^{2}$ and $k(x,y)=3x-2y$ as convex, concave, both or
        neither.
  \item Which of the four functions in the previous problem are
        quasiconcave?
  \item The isoquants of a production function are straight lines with
        slope $-1$. What does that say about substitutability of the two
        inputs, and give such a production function.
  \item An indifference map consists of circles centred at $(4,4)$, with
        utility \emph{increasing} inwards. Interpret the point $(4,4)$
        economically and say why such preferences are unusual.
  \item On the contour map of problem 21, a walker starts at $P_{6}=(-2,1)$
        and always walks perpendicular to the contours, downhill. Describe
        where they end up and why the answer is not the minimum at
        $(1,0)$.
\end{problems}

\section{Additional Solutions}

\begin{solutions}
  \item $y-x^{2}\ge0$ and $x\neq0$: the region on or above the parabola
        $y=x^{2}$, with the $y$-axis removed.
  \item $4-x^{2}-y^{2}>0\iff x^{2}+y^{2}<4$: the \emph{open} disc of
        radius $2$ (boundary excluded, because the logarithm needs a
        strictly positive argument).
  \item $D=\R^{2}$. The exponent $-(x^{2}+y^{2})$ runs over $(-\infty,0]$,
        so the range is $(0,1]$ --- a bell-shaped hill with peak $1$ at the
        origin.
  \item $x^{2}-4x+y^{2}=0\iff(x-2)^{2}+y^{2}=4$: the circle of radius $2$
        centred at $(2,0)$, passing through the origin.
  \item $(x-2)^{2}+y^{2}=0$: the single point $(2,0)$ --- the minimum of
        the shifted bowl, with value $-4$.
  \item $y\,\mathrm{e}^{x}=1\iff y=\mathrm{e}^{-x}$: a decaying
        exponential curve.
  \item $y/x=2\iff y=2x$ with $x\neq0$: the line of slope $2$ through the
        origin, \emph{minus} the origin itself, since $f$ is undefined
        where $x=0$.
  \item $\max(x,y)=1$ is the union of the ray $\{x=1,\ y\le1\}$ and the ray
        $\{y=1,\ x\le1\}$: an L-shaped corner opening down-left, with its
        vertex at $(1,1)$.
  \item Ellipses $\dfrac{x^{2}}{c}+\dfrac{4y^{2}}{c}=1$ with semi-axes
        $\sqrt{c}$ and $\sqrt{c}/2$: for $c=4,16,36$ the semi-axes are
        $2,4,6$ and $1,2,3$. Equally spaced, in fact --- but for
        \emph{equally spaced} levels the contours would crowd outwards, as
        for any quadratic bowl.
  \item $x=y^{2}+c$: parabolas opening to the right, the whole family being
        horizontal translates of $x=y^{2}$ by $c$. For $c=-1,0,1$ the
        vertices sit at $(-1,0)$, $(0,0)$, $(1,0)$.
  \item Any function whose height grows ever more slowly with the radius,
        for instance $f(x,y)=\sqrt[4]{x^{2}+y^{2}}$ or
        $f(x,y)=\ln(x^{2}+y^{2})$: the level $c$ has radius growing faster
        than linearly in $c$, so equal height steps take ever larger
        radial steps and the contours thin out --- equivalently they crowd
        towards the origin.
  \item $f(x,y)=y/x$ (defined for $x\neq0$): the level $c$ is the punctured
        line $y=cx$. Any function of the ratio $y/x$ alone has this
        property; such functions are exactly the homogeneous functions of
        degree $0$.
  \item $W(1)=\{|x|+|y|\le1\}$ is the closed diamond (square rotated
        $45^{\circ}$) with vertices $(\pm1,0)$, $(0,\pm1)$ --- convex.
        $B(1)$ is everything outside it --- not convex.
  \item $xy\le0$ is the union of the second and fourth quadrants. Not
        convex: $(-1,1)$ and $(1,-1)$ both lie in it, but their midpoint
        $(0,0)$ does lie in it too --- take instead $(-1,2)$ and $(2,-1)$,
        whose midpoint $(0.5,0.5)$ has $xy>0$ and is outside.
  \item $x+y\ge2$: a closed half-plane, convex.
  \item $x^{2}\le4\iff-2\le x\le2$: the vertical strip between the lines
        $x=-2$ and $x=2$, convex (an intersection of two half-planes).
  \item $f(P)=1+4=5$: the level curve is the circle of radius $\sqrt5$,
        $W(P)$ the closed disc of radius $\sqrt5$, and $B(P)$ everything
        outside or on that circle.
  \item If every sublevel \emph{and} every superlevel set is convex, the
        function is both quasiconvex and quasiconcave. The level sets must
        then be ``flat'': any function of the form $h(ax+by)$ with $h$
        monotone qualifies, the simplest being a linear function
        $f(x,y)=ax+by$.
  \item A dome: the paraboloid turned upside down, with its peak of height
        $4$ above the origin and falling away in every direction. Contours
        are circles $x^{2}+y^{2}=4-c$ of radius $\sqrt{4-c}$, shrinking as
        $c$ rises to $4$. The function is \textbf{concave}.
  \item $f(x,y)=x^{2}+y^{2}$ and $g(x,y)=\sqrt{x^{2}+y^{2}}$ both have
        concentric circles as level sets. The maps differ in the
        \emph{labels}: for equally spaced levels $f$'s contours crowd
        outwards while $g$'s stay equally spaced. A contour map without its
        level labels does not determine the function.
  \item $f$ convex; $g$ concave; $h$ neither (a saddle); $k$ linear, hence
        both.
  \item $f$: quasiconvex, not quasiconcave (its superlevel sets are
        complements of discs). $g$: concave, hence quasiconcave. $h$:
        neither. $k$: linear, hence both quasiconvex and quasiconcave.
  \item Straight isoquants of slope $-1$ mean the marginal rate of
        technical substitution is constant at $1$: the inputs are
        \textbf{perfect substitutes}, one unit of either always replacing
        one unit of the other. The production function is
        $Q(K,L)=K+L$ (or any increasing function of $K+L$).
  \item $(4,4)$ is a \textbf{bliss point}: a most-preferred bundle, beyond
        which more of either good makes the consumer worse off. Standard
        preferences are monotone --- more is better --- so utility should
        increase without bound in the direction of both goods; circular
        indifference curves with utility rising inwards describe satiation
        instead.
  \item Walking perpendicular to the contours downhill from $P_{6}=(-2,1)$
        means following $-\grad f$. At $P_{6}$,
        $\grad f=(3\cdot4-3,\,2)=(9,2)$, so the walker moves left and
        slightly down --- into the region $x<-2$, where the contours of
        \emph{lower} level continue without bound. The function
        $x^{3}-3x+y^{2}$ tends to $-\infty$ as $x\to-\infty$, so the
        descent never ends: the walker never reaches the local minimum at
        $(1,0)$, because gradient descent finds a \emph{local} basin only
        if it starts in one.
\end{solutions}

\chapter{First-Order Partial Derivatives and the Gradient}

A function of one variable has a graph in the plane and a derivative that is
the slope of a tangent line. A function of two variables has a graph that is
a \emph{surface} in space, and there is no single slope any more: standing
at a point of a hillside, the ground rises steeply in one direction, falls
in another, and is level in a third. The first task of multivariable
calculus is to make that observation quantitative.

The device is to look one direction at a time. Freeze $y$ at a value and let
only $x$ move: the surface is cut by a vertical plane, the cut is an
ordinary curve, and the slope of \emph{that} curve is the partial derivative
$\partial f/\partial x$. This is the \emph{ceteris paribus} principle of
economics, and it is why partial derivatives are the natural language for
marginal products and marginal utilities.

Collecting the partial derivatives into one vector gives the gradient. It
turns out to carry all first-order information at once: it points in the
direction of steepest ascent, its length is the steepest slope available,
and it stands perpendicular to the level curve through the point. Those
three facts are proved below rather than asserted --- the third one in
particular is what the Lagrange method rests on several chapters later.

\section{Core Rules}

\begin{rulebox}{Topological vocabulary of domains}
\begin{itemize}[leftmargin=1.6em,itemsep=0.3ex]
  \item The \emph{open ball} (or \emph{neighbourhood}) of radius $r>0$ around $\mathbf{x}_0\in\R^n$ is
        $B_r(\mathbf{x}_0)=\{\mathbf{x}\in\R^n :
        \|\mathbf{x}-\mathbf{x}_0\|<r\}$.
  \item $\mathbf{x}_0\in M$ is an \emph{interior point} of $M$ if some open ball around
        $\mathbf{x}_0$ lies entirely in $M$.
  \item $M$ is \emph{open} if \emph{every} point of $M$ is
        an interior point.
  \item $\mathbf{x}_0$ is a \emph{boundary point} of
        $M$ if every ball around $\mathbf{x}_0$ contains points of $M$
        \emph{and} points of its complement.
  \item $M$ is \emph{closed} if it contains all
        its boundary points; equivalently, the complement
        $\R^n\setminus M$ is open.
  \item $M$ is \emph{bounded} if
        $M\subseteq B_r(\mathbf{0})$ for some radius $r$.
  \item $M\subseteq\R^n$ is \emph{compact} exactly when
        $M$ is closed \emph{and} bounded.
\end{itemize}
Economic examples: the open positive orthant $\{(x,y): x>0,\ y>0\}$
(the natural Cobb--Douglas domain) is open and unbounded; a budget set
$\{\mathbf{x}\in\R^n : \mathbf{p}\cdot\mathbf{x}\le m,\
\mathbf{x}\ge\mathbf{0}\}$ is closed and bounded, hence compact.
\end{rulebox}

\begin{rulebox}{Why it matters}
Differentiability (and hence every partial derivative below) is
defined at \emph{interior} points of the domain, where $f$ can be
probed in all directions. And by the \emph{Weierstrass theorem}, a
continuous function on a compact set attains its maximum and minimum:
this is why budget-constrained optimisation problems are
guaranteed to have solutions.
\end{rulebox}

\begin{rulebox}{Definition}
Let $f:\R^n\to\R$. The partial derivative of $f$ with respect to $x_i$ at
$\mathbf{x}=(x_1,\ldots,x_n)$ is
\[
  \pd{f}{x_i}(\mathbf{x}) \;=\; \lim_{h\to0}
    \frac{f(x_1,\ldots,x_i+h,\ldots,x_n)-f(x_1,\ldots,x_n)}{h} .
\]
Only the $i$-th slot changes; all other coordinates sit still. It is the
ordinary one-variable derivative of the function
$t\mapsto f(x_1,\ldots,t,\ldots,x_n)$, evaluated at $t=x_i$.
\end{rulebox}

\begin{rulebox}{What it measures: a slice of the surface}
Fix $y=y_{0}$. The vertical plane $y=y_{0}$ cuts the surface $z=f(x,y)$
along a curve, and along that curve only $x$ varies, so the curve is the
graph of the one-variable function $g(x)=f(x,y_{0})$. Then
\[
  \pd{f}{x}(x_{0},y_{0})=g'(x_{0}),
\]
the ordinary slope of the cut curve at $x_{0}$. Symmetrically,
$\partial f/\partial y$ is the slope of the cut made by the plane
$x=x_{0}$. A partial derivative is therefore never a slope \emph{of the
surface}; it is the slope of the surface \emph{along one chosen direction},
and different directions give different numbers. What the surface does in
all the remaining directions is the subject of the gradient below and of the
directional derivative.
\end{rulebox}

\begin{rulebox}{Practical computation}
To compute $\partial f/\partial x_i$, differentiate $f$ with respect to
$x_i$ and treat every other variable as a constant. All single-variable
rules (power, product, quotient, chain) apply unchanged. Worked example, for
$f(x,y)=x^{3}y^{2}+\mathrm{e}^{xy}-\ln y$:
\[
  \pd{f}{x}=3x^{2}y^{2}+y\,\mathrm{e}^{xy},
  \qquad
  \pd{f}{y}=2x^{3}y+x\,\mathrm{e}^{xy}-\frac{1}{y} .
\]
In the first computation $y$ is a number: $y^{2}$ is a constant factor, and
the chain rule on $\mathrm{e}^{xy}$ produces the inner derivative $y$. In
the second, $x$ is the constant and $\ln y$ finally contributes.
\end{rulebox}

\begin{rulebox}{Notation}
The following all refer to the same object:
\[
  \pd{f}{x} = f_x = \partial_x f = D_x f .
\]
\end{rulebox}

\begin{rulebox}{Gradient}
The vector of all first-order partial derivatives is
\[
  \grad f(\mathbf{x}) \;=\;
  \begin{pmatrix} f_{x_1}(\mathbf{x}) \\ f_{x_2}(\mathbf{x}) \\ \vdots \\
    f_{x_n}(\mathbf{x}) \end{pmatrix} .
\]
\end{rulebox}

\begin{rulebox}{Gradient and derivative: column and row}
The gradient is a \textbf{column} vector. The \textbf{derivative} of
$f\colon\R^{n}\to\R$ at $\mathbf{x}$, written $Df(\mathbf{x})$, is the same
list of partials arranged as a $1\times n$ \textbf{row}:
\[
  Df(\mathbf{x})=\bigl(f_{x_1}(\mathbf{x})\ \ f_{x_2}(\mathbf{x})\ \ \cdots\ \
  f_{x_n}(\mathbf{x})\bigr)=\grad f(\mathbf{x})^{T} .
\]
Same numbers, different shape. The row is the form that multiplies a column
of changes $\mathbf{h}$:
$Df(\mathbf{x})\,\mathbf{h}=\grad f(\mathbf{x})\cdot\mathbf{h}
=\sum_{i}f_{x_i}h_{i}$, which is the total differential of Chapter~12 and
the linear part of the tangent plane below. In Chapter~13 these rows are
stacked into the Jacobian matrix.
\end{rulebox}

\begin{rulebox}{Functions of $n$ variables written with sums}
Many economic functions are sums over goods, periods or inputs. To
differentiate with respect to one fixed variable $x_{k}$, ask which terms of
the sum contain $x_{k}$; every other term is a constant and drops out.
\begin{itemize}[leftmargin=1.6em,itemsep=0.3ex]
  \item \textbf{One function of each variable.} For
        $f=\sum_{i=1}^{n}g(x_{i})$ only the term $i=k$ survives:
        \[
          \pd{f}{x_{k}}=g'(x_{k}) .
        \]
        So $\sum_{i}x_{i}^{3}$ has $\partial f/\partial x_{k}=3x_{k}^{2}$.
  \item \textbf{The index inside the term.} $\sum_{i=1}^{n}x_{i}^{\,i}
        =x_{1}+x_{2}^{2}+\cdots+x_{n}^{n}$: the term containing $x_{k}$ is
        $x_{k}^{k}$, so $\partial f/\partial x_{k}=k\,x_{k}^{k-1}$.
        Likewise $\sum_{j=1}^{n}j/x_{j}^{4}=\sum_{j}j\,x_{j}^{-4}$ gives
        $\partial f/\partial x_{k}=-4k\,x_{k}^{-5}=-4k/x_{k}^{5}$.
  \item \textbf{The sum inside a function.} For $f=\bigl(\sum_{i}x_{i}\bigr)^{2}$
        use the chain rule; the inner derivative is
        $\partial\bigl(\sum_{i}x_{i}\bigr)/\partial x_{k}=1$, so
        $\partial f/\partial x_{k}=2\sum_{i=1}^{n}x_{i}$, the same for
        every $k$.
  \item \textbf{Products.} $f=\prod_{i}x_{i}$ gives
        $\partial f/\partial x_{k}=\prod_{i\neq k}x_{i}$: the product with
        the factor $x_{k}$ left out.
\end{itemize}
The summation index is a dummy and the differentiation variable is fixed:
call them by different letters ($i$ and $k$). Writing
$\partial/\partial x_{i}$ of a sum over $i$ is the standard source of
confusion.
\end{rulebox}

\begin{rulebox}{Exponentials to another base}
For $a>0$, $a^{x}=\mathrm{e}^{x\ln a}$, so
\[
  \frac{\mathrm{d}}{\mathrm{d}x}\,a^{x}=a^{x}\ln a,
  \qquad
  \pd{}{x}\,a^{g(x,y)}=a^{g(x,y)}\ln a\cdot g_{x}(x,y) .
\]
For example $\partial(2^{xy})/\partial x=y\,2^{xy}\ln2$. Keep the three
shapes apart: $x^{a}$ (variable base, power rule), $a^{x}$ (variable
exponent, factor $\ln a$), and $x^{y}$ (both, written as
$\mathrm{e}^{y\ln x}$).
\end{rulebox}

\begin{gallery}
\begin{tikzpicture}
  \begin{axis}[surfaxis, title={$f(K,L)=\sqrt{KL}$, sliced at $L=2$},
      xlabel={$K$}, ylabel={$L$}, zlabel={}, view={135}{26},
      zmin=0, zmax=3.2, z post scale=1.7]
    \addplot3[surf, domain=0.15:3, y domain=0.15:3, opacity=0.85]{sqrt(x*y)};
    \addplot3[samples=80, samples y=0, domain=0.15:3, line width=1.1pt,
      draw=black]({x},{2},{sqrt(2*x)});
  \end{axis}
\end{tikzpicture}\hfill
\begin{tikzpicture}
  \begin{axis}[contouraxis, xmin=-0.3, xmax=3.4, ymin=-0.3, ymax=3.4,
      xtick={1,2,3}, ytick={1,2,3},
      title={isoquants of $\sqrt{KL}$ with $\grad f$},
      xlabel={$K$}, ylabel={$L$}]
    \addplot[accent!40, line width=0.5pt, domain=0.22:3.3]{0.64/x};
    \addplot[accent!40, line width=0.5pt, domain=0.48:3.3]{1.44/x};
    \addplot[accent!40, line width=0.5pt, domain=0.85:3.3]{2.56/x};
    \addplot[accent!40, line width=0.5pt, domain=1.34:3.3]{4/x};
    \draw[-{Stealth[length=6pt,width=4pt]}, accent, line width=1pt]
      (axis cs:1,1) -- (axis cs:1.62,1.62);
    \draw[-{Stealth[length=6pt,width=4pt]}, accent, line width=1pt]
      (axis cs:2,0.5) -- (axis cs:2.16,1.14);
    \draw[-{Stealth[length=6pt,width=4pt]}, accent, line width=1pt]
      (axis cs:0.5,2) -- (axis cs:1.14,2.16);
    \addplot[only marks, mark=*, mark size=1pt, black, forget plot]
      coordinates {(1,1) (2,0.5) (0.5,2)};
  \end{axis}
\end{tikzpicture}
\gcap{Left: the plane $L=2$ cuts the surface; the slope of the cut curve is
$\partial f/\partial K$. Right: the gradient at three points, each one
perpendicular to the isoquant through it.}
\end{gallery}

\begin{rulebox}{Geometric meaning --- and why}
All three facts follow from one identity: the rate of change of $f$ in the
direction of a unit vector $\mathbf{u}$ is $\grad f\cdot\mathbf{u}$ (the
directional derivative below).
\begin{itemize}[leftmargin=1.6em,itemsep=0.4ex]
  \item \textbf{Steepest ascent.} Since
        $\grad f\cdot\mathbf{u}=\|\grad f\|\,\|\mathbf{u}\|\cos\theta
        =\|\grad f\|\cos\theta$, the rate is largest when $\cos\theta=1$,
        i.e. when $\mathbf{u}$ points \emph{along} $\grad f$. So
        $\grad f(\mathbf{x}_0)$ is the direction of steepest ascent, and
        $-\grad f$ of steepest descent.
  \item \textbf{Its length is that steepest rate.} Putting $\theta=0$ above
        leaves exactly $\|\grad f(\mathbf{x}_0)\|$.
  \item \textbf{Perpendicular to the level set.} Let $\mathbf{c}(t)$ be any
        curve running inside the level set through $\mathbf{x}_{0}$, so that
        $f(\mathbf{c}(t))$ is \emph{constant}. Differentiating with the
        chain rule gives
        \[
          0=\frac{\mathrm{d}}{\mathrm{d}t}f(\mathbf{c}(t))
           =\grad f(\mathbf{c}(t))\cdot\mathbf{c}'(t) .
        \]
        At $t_{0}$ this says $\grad f(\mathbf{x}_{0})$ is orthogonal to
        $\mathbf{c}'(t_{0})$, the tangent to the level set. Since the curve
        was arbitrary, the gradient is perpendicular to the whole level set.
\end{itemize}
The third fact is the one that matters most later: an optimum on a
constraint is a point where the level set of the objective and the
constraint set are \emph{tangent}, which is the same as their gradients
being parallel --- and that is the Lagrange condition.
\end{rulebox}

\begin{rulebox}{Economic interpretation}
\begin{itemize}[leftmargin=1.6em,itemsep=0.3ex]
  \item For a production function $Q=F(K,L)$: $F_K=\mathrm{MP}_K$
        (marginal product of capital), $F_L=\mathrm{MP}_L$ (marginal
        product of labour).
  \item For a utility function $U(x,y)$: $U_x, U_y$ are the marginal
        utilities; the marginal rate of substitution is
        $\mathrm{MRS}_{xy}=U_x/U_y$.
  \item For production: the marginal rate of technical substitution is
        $\mathrm{MRTS}_{LK}=F_L/F_K$.
\end{itemize}
\end{rulebox}

\begin{rulebox}{Directional derivative}
Let $\mathbf{u}\in\R^n$ be a \emph{unit} vector ($\|\mathbf{u}\|=1$). The
directional derivative of $f$ at $\mathbf{x}_0$ in the direction
$\mathbf{u}$ is
\[
  D_{\mathbf{u}}f(\mathbf{x}_0) \;=\; \grad f(\mathbf{x}_0)\cdot\mathbf{u} .
\]
\end{rulebox}

\begin{rulebox}{Caution}
If you are given a non-unit direction $\mathbf{v}$, you \emph{must}
normalise first: $\mathbf{u}=\mathbf{v}/\|\mathbf{v}\|$. Forgetting this is
the single most common error in directional-derivative problems.
\end{rulebox}

\begin{rulebox}{Differentiability in several variables}
Let $D\subseteq\R^{n}$ be open and $\mathbf{x}_{0}\in D$. The function $f$
is \textbf{differentiable} \de{(differenzierbar)} at $\mathbf{x}_{0}$ if
there is a vector $\mathbf{a}$ with
\[
  f(\mathbf{x}_{0}+\mathbf{h})=f(\mathbf{x}_{0})+\mathbf{a}\cdot\mathbf{h}
  +r(\mathbf{h}),
  \qquad
  \frac{r(\mathbf{h})}{\|\mathbf{h}\|}\to0\ \text{ as }\ \mathbf{h}\to\mathbf{0}.
\]
Near $\mathbf{x}_{0}$, $f$ is a linear function plus an error that vanishes
faster than the distance. For $n=2$ the graph of the linear part is a plane,
so differentiability says that the surface has a \textbf{tangent plane} at
the point: zoom in far enough and the surface looks flat.

If $f$ is differentiable at $\mathbf{x}_{0}$, then
\begin{itemize}[leftmargin=1.6em,itemsep=0.3ex]
  \item $f$ is continuous at $\mathbf{x}_{0}$;
  \item every partial derivative exists there, and
        $\mathbf{a}=\grad f(\mathbf{x}_{0})$ (take
        $\mathbf{h}=t\mathbf{e}_{k}$ and divide by $t$);
  \item every directional derivative exists and equals
        $\grad f(\mathbf{x}_{0})\cdot\mathbf{u}$.
\end{itemize}
\textbf{Sufficient criterion.} If all partials exist on an open
neighbourhood of $\mathbf{x}_{0}$ and are continuous there ($f$ is
$C^{1}$, \emph{continuously differentiable}), then $f$ is differentiable at
$\mathbf{x}_{0}$. Polynomials, exponentials, logarithms of positive
expressions, roots of positive expressions, quotients with non-zero
denominator, and all their sums, products and compositions are $C^{1}$ on
the open set where their formulas apply.

\textbf{Partials alone are not enough.} Take
$f(x,y)=\dfrac{xy}{x^{2}+y^{2}}$ with $f(0,0)=0$. On both axes $f=0$, so
$f_{x}(0,0)=f_{y}(0,0)=0$. But on the diagonal $y=x$,
$f=x^{2}/(2x^{2})=\tfrac12$ for every $x\neq0$. So $f$ is not even
continuous at the origin, and therefore not differentiable: the partials
see only two directions out of infinitely many.
\end{rulebox}

\begin{rulebox}{Where differentiability fails, and finding the set}
Compute the partials by the rules. The formulas are valid on the open set
where the rules applied; the points where a formula breaks down (a zero
denominator, a zero radicand, the kink of an absolute value or a minimum)
are the suspects, and each suspect is settled by the difference quotient.
\begin{itemize}[leftmargin=1.6em,itemsep=0.3ex]
  \item \textbf{Cone apex.} $f=\sqrt{(x-a)^{2}+(y-b)^{2}}$ has
        $f_{x}=(x-a)/f$, undefined at $(a,b)$. Directly,
        $f(a+h,b)-f(a,b)=\sqrt{h^{2}}=|h|$, and $|h|/h=\pm1$ has no limit:
        $f_{x}(a,b)$ does not exist, nor does $f_{y}(a,b)$. The same
        happens for the elliptic cone $\sqrt{4(x-a)^{2}+y^{2}}$ at $(a,0)$,
        with quotient $2|h|/h$. For the root of a sum of squares of
        \emph{linear} expressions, such as these cones, $f$ is
        differentiable exactly where the radicand is non-zero, and the
        graph has a sharp tip (or crease) where it vanishes. With higher
        powers under the root this can change: $\sqrt{x^{4}+y^{4}}$ is
        differentiable even at the origin, where its radicand is $0$.
  \item \textbf{$|x|+|y|$.} $f_{x}=1$ for $x>0$ and $-1$ for $x<0$; at
        $x=0$ the quotient is $|h|/h$ again. So $f_{x}$ fails on the whole
        $y$-axis and $f_{y}$ on the whole $x$-axis, and $f$ is
        differentiable exactly off the axes. The graph is an upside-down
        square pyramid, creased along both axes.
  \item \textbf{$\min\{ax,by\}$} with $a,b>0$. Where $ax<by$, $f=ax$;
        where $ax>by$, $f=by$. On the kink line $ax=by$ (in the quadrant
        $x,y>0$ of economics, a ray from the origin), raising $x$ leaves
        $f$ at $by$ (slope $0$) while lowering $x$ lowers $f$ at rate $a$:
        the one-sided slopes differ, so $f_{x}$ does not exist there, and
        nor does $f_{y}$. The graph is a roof whose ridge lies above that
        line.
\end{itemize}
A breaking formula is a warning, not a verdict: the partial may still
exist at that point, and only the difference quotient decides. And partials
that exist at a point still do not guarantee differentiability; continuity
of the partials near the point does.
\end{rulebox}

\begin{rulebox}{Proving differentiability from the definition}
When the $C^{1}$ criterion is not available, or a proof is asked for,
work with the remainder directly. At a point $\mathbf{x}^{*}$:
\begin{enumerate}[leftmargin=1.6em,itemsep=0.3ex]
  \item \textbf{Candidate.} Compute $\grad f(\mathbf{x}^{*})$; it is the
        only possible linear part. If a partial does not exist, stop: $f$
        is not differentiable.
  \item \textbf{Remainder.}
        $r(\mathbf{h})=f(\mathbf{x}^{*}+\mathbf{h})-f(\mathbf{x}^{*})
        -\grad f(\mathbf{x}^{*})\cdot\mathbf{h}$, expanded and simplified.
        For a polynomial, every term left over contains at least two
        factors $h_{i}$.
  \item \textbf{Estimate.} Show $|r(\mathbf{h})|/\|\mathbf{h}\|\to0$ with
        the standard inequalities
        \[
          |h_{i}|\le\|\mathbf{h}\|,\qquad
          \frac{h_{i}^{2}}{\|\mathbf{h}\|}\le\|\mathbf{h}\|,\qquad
          \frac{|h_{1}h_{2}|}{\|\mathbf{h}\|}\le\tfrac12\|\mathbf{h}\| ,
        \]
        and, for a product of three factors $h_{i}$, a bound
        $\|\mathbf{h}\|^{2}$.
  \item \textbf{Transcendental factors.} Replace them by their linear part
        plus a remainder: $\mathrm{e}^{t}=1+t+R(t)$ with $|R(t)|\le2t^{2}$
        for $|t|\le1$, and $\ln(1+t)=t+R(t)$ with $|R(t)|\le2t^{2}$ for
        $|t|\le\tfrac12$ (both from the Lagrange remainder of Chapter~4).
        The quadratic bound on $R$ is what makes $R/\|\mathbf{h}\|\to0$.
\end{enumerate}
\textbf{Roots at the origin.} For $f=(x^{2}+y^{2})^{p}$ with $p>0$,
$f(\mathbf{h})=\|\mathbf{h}\|^{2p}$. The difference quotient on an axis is
$|h|^{2p}/h$: it tends to $0$ if $2p>1$, and has no limit if $2p\le1$. So
$f$ is differentiable at the origin exactly when $p>\tfrac12$; the cone
$p=\tfrac12$ is the borderline case, and $p=\tfrac13$ is worse than a cone.
\end{rulebox}

\begin{rulebox}{Tangent plane and linearisation}
If $f$ is differentiable at $P=(x_{0},y_{0})$, the \textbf{tangent plane}
\de{(Tangentialebene)} to its graph at $(x_{0},y_{0},f(P))$ is
\[
  z=f(P)+f_{x}(P)\,(x-x_{0})+f_{y}(P)\,(y-y_{0}) .
\]
In $n$ variables, with $P=\mathbf{p}=(p_{1},\ldots,p_{n})$, the
\textbf{tangent hyperplane} is
\[
  z=f(\mathbf{p})+\grad f(\mathbf{p})\cdot(\mathbf{x}-\mathbf{p})
   =f(\mathbf{p})+\sum_{i=1}^{n}f_{x_i}(\mathbf{p})\,(x_{i}-p_{i}) .
\]
\begin{itemize}[leftmargin=1.6em,itemsep=0.3ex]
  \item The right-hand side is the \textbf{linearisation} $L(\mathbf{x})$
        of $f$ at $\mathbf{p}$, the multivariable version of the tangent
        line $f(a)+f'(a)(x-a)$ of Chapter~6. With
        $\mathrm{d}x_{i}=x_{i}-p_{i}$ it reads $L=f(\mathbf{p})+\mathrm{d}f$:
        the old value plus the total differential.
  \item The partials are evaluated \emph{at the point} $\mathbf{p}$; they
        are numbers, which is what makes the result a plane.
  \item For a point $Q$ near $P$, $f(Q)\approx L(Q)$. The difference
        $f(Q)-L(Q)$ is the remainder $r(\mathbf{h})$ of the definition, so
        comparing the two numbers measures how good the approximation is.
        For $f$ with continuous second partials it is at most of the order
        $\|Q-P\|^{2}$: halve the step and the error falls to about a
        quarter, or even less.
  \item Equivalently, $f_{x}(P)(x-x_{0})+f_{y}(P)(y-y_{0})-(z-f(P))=0$: the
        vector $(f_{x}(P),\,f_{y}(P),\,-1)$ is normal to the plane.
\end{itemize}
Where $f$ is not differentiable there is no tangent plane: a cone has none
at its apex.
\end{rulebox}

\section{Worked Examples}

\begin{worked}{For $f(x,y)=x^{3}y^{2}+\mathrm{e}^{2x}\ln y$, compute both partials and the gradient at $(1,1)$; then read the marginal products of $Q=10K^{0.3}L^{0.7}$}
\wstep{Step 1. Differentiating in $x$ means freezing $y$ --- literally.}
Treat $y$ as if it were the number $7$. Then $x^{3}y^{2}$ is
``$y^{2}$ times $x^{3}$'', a constant times a power, and $\ln y$ is a
constant multiplying $\mathrm{e}^{2x}$:
\[
  f_{x}=3x^{2}y^{2}+2\mathrm{e}^{2x}\ln y .
\]
\wtrap{Freezing means the frozen variable never produces a derivative
factor. $\frac{\partial}{\partial x}\bigl[\ln y\bigr]=0$, not $1/y$. The
symbol $\partial$ instead of $\mathrm{d}$ is the reminder that something is
being held fixed.}

\wstep{Step 2. Now freeze $x$ instead.}
$x^{3}$ is a constant multiplying $y^{2}$, and $\mathrm{e}^{2x}$ is a
constant multiplying $\ln y$:
\[
  f_{y}=2x^{3}y+\frac{\mathrm{e}^{2x}}{y} .
\]

\wstep{Step 3. Evaluate, and assemble the gradient.}
At $(1,1)$, with $\ln1=0$ and $\mathrm{e}^{2}\approx7.389$:
\[
  f_{x}(1,1)=3\cdot1\cdot1+2\mathrm{e}^{2}\cdot0=3,
  \qquad
  f_{y}(1,1)=2\cdot1\cdot1+\frac{\mathrm{e}^{2}}{1}=2+\mathrm{e}^{2}\approx9.389 .
\]
\[
  \grad f(1,1)=\begin{pmatrix}3\\2+\mathrm{e}^{2}\end{pmatrix}\approx\begin{pmatrix}3\\9.389\end{pmatrix}
\]
\wtrap{The gradient is a \textbf{vector}, and it is a vector in the
\emph{input} space --- here a direction in the $xy$-plane, not on the
surface. Reporting a single number, or a three-component object including
the height, are the two standard confusions.}

\wstep{Step 4. What each partial measures.}
$f_{x}$ is the slope of the curve obtained by slicing the surface with the
plane $y=1$ --- an ordinary one-variable derivative of the function
$x\mapsto f(x,1)$. Indeed $f(x,1)=x^{3}$, whose derivative at $x=1$ is $3$.
\checkmark\ Two partials, two slices, two ordinary derivatives.

\wstep{Step 5. What the gradient does that the partials alone do not.}
\begin{itemize}[leftmargin=1.6em,itemsep=0.2ex]
  \item It points in the direction of \textbf{steepest ascent}. Here that
        is roughly $(0.30,\,0.95)$ after normalising: mostly ``north''.
  \item Its length $\|\grad f\|=\sqrt{9+(2+\mathrm{e}^{2})^{2}}\approx9.86$
        is the steepest rate of climb available at that point.
  \item It is \textbf{perpendicular} to the level curve through the point.
\end{itemize}
The last is why the gradient governs constrained optimisation: at an
optimum the objective's gradient must be parallel to the constraint's, which
is the Lagrange condition.

\wstep{Step 6. Directional derivative: normalise, then take an inner product.}
In the direction of $\mathbf v=(3,4)$:
\[
  \mathbf u=\frac{\mathbf v}{\|\mathbf v\|}=\Bigl(\tfrac35,\tfrac45\Bigr),
  \qquad
  D_{\mathbf u}f=\grad f\cdot\mathbf u
  =\tfrac35\cdot3+\tfrac45(2+\mathrm{e}^{2})\approx9.31 .
\]
\wtrap{Normalise first. Using $(3,4)$ unnormalised gives $46.6$, which is
five times too large and has no meaning --- ``rate of change per unit
distance'' requires a unit vector. And $9.31<9.86$, as it must be: no
direction beats the gradient.}

\wstep{Step 7. The economic reading, on a production function.}
For $Q=10K^{0.3}L^{0.7}$:
\[
  MP_{K}=\pd{Q}{K}=3K^{-0.7}L^{0.7}=3\Bigl(\frac{L}{K}\Bigr)^{0.7},
  \qquad
  MP_{L}=\pd{Q}{L}=7\Bigl(\frac{K}{L}\Bigr)^{0.3} .
\]
At $K=100$, $L=50$: $MP_{K}=3(0.5)^{0.7}\approx1.85$ and
$MP_{L}=7(2)^{0.3}\approx8.62$. So one more unit of labour adds about $8.6$
units of output, one more of capital about $1.85$.

Two structural facts fall out. The \textbf{marginal rate of technical
substitution} is
\[
  \mathrm{MRTS}=\frac{MP_{L}}{MP_{K}}=\frac{0.7}{0.3}\cdot\frac{K}{L}\approx4.67 ,
\]
and each marginal product depends only on the \emph{ratio} $K/L$ --- the
signature of constant returns to scale, which is why Cobb--Douglas with
exponents summing to $1$ behaves the way it does.
\wtrap{$MP_{K}$ falls as $K$ rises (the exponent $-0.7$ is negative):
diminishing marginal product. But that is diminishing \emph{marginal},
not diminishing total --- output still increases. Confusing the two is the
most common misreading of a partial derivative in economics.}
\end{worked}

\begin{worked}{Find the tangent plane of $f(x,y)=x\ln(x+2y)$ at $P=(1,0)$, and test it at the nearby point $Q=(1.1,\,0.05)$}
\wstep{Step 1. Make sure there is a plane to find.}
The domain is the open half-plane $x+2y>0$, and $P$ lies in it since
$1+0>0$. On that half-plane $f$ is a polynomial times the logarithm of a
positive linear expression, so its partials exist and are continuous: $f$
is $C^{1}$, hence differentiable at $P$, and the tangent plane exists.

\wstep{Step 2. The height at $P$.}
$f(1,0)=1\cdot\ln1=0$. The plane touches the surface at $(1,0,0)$.

\wstep{Step 3. The two partials, then their values at $P$.}
In $x$ the function is a product of $x$ and $\ln(x+2y)$:
\[
  f_{x}=\ln(x+2y)+x\cdot\frac{1}{x+2y} .
\]
In $y$ the factor $x$ is a constant, and the chain rule supplies the inner
derivative $2$:
\[
  f_{y}=\frac{2x}{x+2y} .
\]
At $P$: $f_{x}(1,0)=\ln1+\tfrac11=1$ and $f_{y}(1,0)=\tfrac21=2$.
\wtrap{$x$ appears twice in $f$, once as a factor and once inside the
logarithm, so $f_{x}$ needs the product rule. Writing $f_{x}=x/(x+2y)$
alone loses the term $\ln(x+2y)$; here it happens to vanish at $P$, which
hides the error until the next point.}

\wstep{Step 4. Assemble the plane.}
\[
  z=f(P)+f_{x}(P)(x-1)+f_{y}(P)(y-0)=0+1\cdot(x-1)+2y=x+2y-1 .
\]
\wtrap{The partials go in as \emph{numbers} evaluated at $P$. Leaving them
as functions of $x$ and $y$ gives a curved surface, not a plane.}

\wstep{Step 5. Compare at $Q$.}
\[
  f(Q)=1.1\cdot\ln(1.1+0.1)=1.1\ln1.2\approx0.200554,
  \qquad
  L(Q)=1.1+0.1-1=0.200000 .
\]
The plane misses by about $0.000554$, while $f$ itself changed by about
$0.2$: the error is under $0.3\%$ of the change. The step has length
$\|Q-P\|=\sqrt{0.1^{2}+0.05^{2}}\approx0.112$. The error is far below that
length, and even below its square $0.0125$, as the rule box predicts: the
remainder vanishes faster than the distance.

\wstep{Step 6. The same number as a total differential.}
With $\mathrm{d}x=0.1$ and $\mathrm{d}y=0.05$,
\[
  \mathrm{d}f=f_{x}(P)\,\mathrm{d}x+f_{y}(P)\,\mathrm{d}y=0.1+0.1=0.2,
\]
and $L(Q)=f(P)+\mathrm{d}f=0.2$. \checkmark\ The tangent plane and the
total differential are one idea written two ways.
\end{worked}

\begin{worked}{Prove from the definition that $f(x,y)=x^{2}y$ is differentiable at every point; then decide at the origin for $(x^{2}+y^{2})^{1/3}$ and $(x^{2}+y^{2})^{2/3}$}
\wstep{Step 1. The candidate for the linear part.}
Fix a point $(a,b)$. If $f$ is differentiable there, the linear part must
be $\grad f(a,b)\cdot\mathbf{h}$, so compute the partials:
$f_{x}=2xy$, $f_{y}=x^{2}$, hence $\grad f(a,b)=(2ab,\,a^{2})$.

\wstep{Step 2. Form the remainder and let everything cancel that can.}
With $\mathbf{h}=(h_{1},h_{2})$,
\[
  f(a+h_{1},b+h_{2})=(a^{2}+2ah_{1}+h_{1}^{2})(b+h_{2})
  =a^{2}b+a^{2}h_{2}+2abh_{1}+2ah_{1}h_{2}+bh_{1}^{2}+h_{1}^{2}h_{2} .
\]
Subtract $f(a,b)=a^{2}b$ and the linear part $2abh_{1}+a^{2}h_{2}$:
\[
  r(\mathbf{h})=2a\,h_{1}h_{2}+b\,h_{1}^{2}+h_{1}^{2}h_{2} .
\]
Every surviving term contains at least \emph{two} factors $h_{i}$. That is
the sign that the candidate was right.

\wstep{Step 3. Bound the quotient by something that tends to $0$.}
Use $|h_{i}|\le\|\mathbf{h}\|$ and $|h_{1}h_{2}|\le\tfrac12\|\mathbf{h}\|^{2}$
(from $(|h_{1}|-|h_{2}|)^{2}\ge0$):
\[
  \frac{|r(\mathbf{h})|}{\|\mathbf{h}\|}
  \le\frac{2|a|\cdot\tfrac12\|\mathbf{h}\|^{2}+|b|\,\|\mathbf{h}\|^{2}
          +\|\mathbf{h}\|^{3}}{\|\mathbf{h}\|}
  =(|a|+|b|)\,\|\mathbf{h}\|+\|\mathbf{h}\|^{2}\;\longrightarrow\;0 .
\]
So $f$ is differentiable at $(a,b)$, for every $(a,b)$, with derivative
$(2ab,\,a^{2})$.
\wtrap{Checking $r/\|\mathbf{h}\|\to0$ along the axes, or along every
line, is not a proof: the bound must hold for \emph{all} small $\mathbf{h}$
at once. The bound may depend on the fixed point $(a,b)$, but only through
constants.}

\wstep{Step 4. A root at the origin: $g=(x^{2}+y^{2})^{1/3}$.}
First the partials, from the difference quotient:
\[
  \frac{g(h,0)-g(0,0)}{h}=\frac{(h^{2})^{1/3}}{h}
  =\frac{|h|^{2/3}}{h}=\pm|h|^{-1/3},
\]
which is unbounded as $h\to0$. So $g_{x}(0,0)$ does not exist (nor, by
symmetry, $g_{y}(0,0)$), and $g$ is \textbf{not differentiable} at the
origin, although it is continuous there. Its graph has a tip even sharper
than a cone's.

\wstep{Step 5. A higher power: $k=(x^{2}+y^{2})^{2/3}$.}
Now $\dfrac{k(h,0)-k(0,0)}{h}=\dfrac{|h|^{4/3}}{h}=\pm|h|^{1/3}\to0$, so both
partials are $0$ and the candidate linear part is $0$. The remainder is
$r(\mathbf{h})=k(\mathbf{h})=\|\mathbf{h}\|^{4/3}$, and
\[
  \frac{|r(\mathbf{h})|}{\|\mathbf{h}\|}=\|\mathbf{h}\|^{1/3}\to0 .
\]
So $k$ \textbf{is} differentiable at the origin, with the tangent plane
$z=0$. In general $(x^{2}+y^{2})^{p}=\|\mathbf{h}\|^{2p}$ at the origin, and
the verdict depends only on whether $2p>1$.
\end{worked}

\section{Practice Problems}

\begin{problems}
  \item Compute $f_x$ and $f_y$ for $f(x,y)=5x^4-2x^2y^3+y^5$.
  \item For $F(K,L)=12\,K^{1/4}L^{3/4}$, find $F_K$ and $F_L$ at
        $(K,L)=(16,81)$.
  \item Compute $f_x$ and $f_y$ for $f(x,y)=x^2e^{3y}$.
  \item Compute $\nabla f$ for $f(x,y)=4x^3y-xy^4+7$.
  \item Compute the directional derivative of $f(x,y)=x^2-4xy+y^3$ at
        $(2,1)$ in the direction $\mathbf{v}=(5,-12)$.
  \item For $F(K,L)=6K^{2/3}L^{1/3}$, compute $\mathrm{MP}_K$ and
        $\mathrm{MP}_L$ at $(K,L)=(27,8)$.
  \item Compute $f_x$ and $f_y$ for $f(x,y)=(3x-2y)^4$.
  \item Compute $f_x$ and $f_y$ for $f(x,y)=e^{xy^2}$.
  \item For $F(K,L)=25K^{0.6}L^{0.5}$, find $F_K$ and $F_L$.
  \item Compute $\nabla f$ for $f(x,y,z)=x^2+3yz-z^2e^x$.
  \item Compute the directional derivative of $f(x,y)=e^x\ln(1+y)$ at
        $(0,0)$ in the direction $\mathbf{v}=(1,2)$.
  \item For $U(x,y)=2\ln(x)+5\ln(y)$, derive $\mathrm{MRS}_{xy}=U_x/U_y$ as
        a function of $(x,y)$ and evaluate at $(3,5)$.
  \item Simplify and compute $f_x, f_y$ for
        $f(x,y)=\dfrac{x^2-y^2}{x+y}$.
  \item Compute $f_x$ and $f_y$ for $f(x,y)=\ln(x^2y+1)$.
  \item For $U(x_1,x_2,x_3)=x_1^{1/2}x_2^{1/3}x_3^{1/6}$, compute the
        three partials.
  \item Compute the directional derivative of $f(x,y,z)=xy-2yz+z^2$ at
        $(1,2,3)$ in the direction $\mathbf{v}=(1,2,2)$.
  \item For $F(K,L)=K^{0.25}L^{0.75}$, derive $\mathrm{MRTS}_{LK}=F_L/F_K$
        and evaluate at $(K,L)=(16,4)$.
  \item Expand and then compute $f_x, f_y$ for
        $f(x,y)=(x^2+y)(x-3y^2)$.
  \item Compute $f_x$ and $f_y$ for $f(x,y)=xe^{-y/x}$ with $x\neq0$.
  \item For $F(K,L)=(3K^{1/2}+2L^{1/2})^2$, find $F_K$ and $F_L$.
  \item Compute $\nabla f$ for $f(x,y)=\arctan(y/x)$ with $x>0$.
  \item For $f(x,y)=x^{1/2}y^{1/2}$ at $(K,L)=(4,9)$, find the unit
        direction of steepest ascent and the maximum rate of increase.
  \item For the CES function $F(K,L)=(K^{1/2}+L^{1/2})^2$, verify Euler's
        theorem $KF_K+LF_L=rF$ for the appropriate degree of homogeneity
        $r$.
  \item Compute $f_x$ and $f_y$ for $f(x,y)=\dfrac{2x+3y}{x^2+y^2}$.
  \item Compute $f_x$ and $f_y$ for $f(x,y)=\dfrac{e^x}{1+e^y}$.
  \item For $F(K,L)=A(\alpha K^\rho+(1-\alpha)L^\rho)^{\nu/\rho}$ (CES),
        find $F_L$.
  \item Compute the directional derivative of $f(x,y)=\ln(x^2+y^2)$ at
        $(3,4)$ in the direction $\mathbf{v}=(-1,1)$.
  \item Compute $f_x, f_y, f_z$ for $f(x,y,z)=x^2yz^3-4xy^2z+5z$.
  \item Compute $f_x$ and $f_y$ for $f(x,y)=\sqrt{4x^2+9y^2}$.
  \item For $F(K,L)=(K^{1/2}+L^{1/2})^4$, compute $F_K$ at
        $(K,L)=(4,9)$.
  \item Compute $\nabla f$ for $f(x,y,z)=\ln(x^2+y^2+z^2)$.
  \item Compute $f_x$ and $f_y$ for
        $f(x,y)=\ln\left(\sqrt{x^2+y^2}\right)$.
  \item Cost minimisation. A firm uses inputs $(x,y)$ with cost
        $C(x,y)=4x+9y$ and output $Q(x,y)=x^{1/2}y^{1/2}$. Compute
        $\nabla C$ and $\nabla Q$, and verify that at the cost-minimising
        input mix $\nabla C=\lambda\nabla Q$ for some $\lambda$.
  \item Compute $f_x$ and $f_y$ for $f(x,y)=(1+xy)^5$.
  \item Compute the directional derivative of $f(x,y)=x^3-2xy+y^2$ at
        $(1,2)$ in the direction $\mathbf{v}=(3,4)$.
  \item For $F(K,L)=10K^{0.3}L^{0.7}$ (Cobb--Douglas), compute $F_K$ and
        $F_L$ at $(K,L)=(100,100)$ and verify $KF_K+LF_L=F$.
  \item Compute $f_x$ and $f_y$ for $f(x,y)=\sin(xy)+\cos(x-y)$.
  \item Compute $\nabla f$ for $f(x,y)=x^2y-xy^2+3x-2y$ at the point
        $(1,-1)$.
  \item For $U(x,y)=x^\alpha y^{1-\alpha}$ (Cobb--Douglas,
        $0<\alpha<1$), derive $\mathrm{MRS}_{xy}=U_x/U_y$ and show it
        equals $\dfrac{\alpha}{1-\alpha}\cdot\dfrac{y}{x}$.
  \item Compute the directional derivative of $f(x,y)=e^{xy}$ at $(1,0)$
        in the direction $\mathbf{v}=(-3,4)$.
  \item Compute $f_x$ and $f_y$ for $f(x,y)=x^y$ (with $x>0$).
        \emph{Hint:} write $x^y=e^{y\ln x}$.
  \item For $F(K,L)=\min\{2K,3L\}$ (Leontief, fixed-proportions, on the
        smooth region away from the kink), compute $F_K$ and $F_L$ in
        each of the two regions $\{2K<3L\}$ and $\{2K>3L\}$.
  \item Compute $\nabla f$ for $f(x,y,z)=xyz+x^2-z^2$ at the point
        $(2,-1,3)$.
  \item Compute the unit vector of steepest \emph{descent} for
        $f(x,y)=x^2+4y^2$ at $(3,1)$, and the corresponding rate of
        decrease.
  \item For $f(x,y)=xe^y-ye^x$, compute $f_x, f_y$.
  \item Compute $f_x$ and $f_y$ for $f(x,y)=\dfrac{x}{x^2+y^2}$.
  \item For $U(x,y)=x+2\sqrt{y}$ (quasi-linear), derive
        $\mathrm{MRS}_{xy}$.
  \item Compute the directional derivative of $f(x,y,z)=x^2+y^2-z^2$ at
        $(2,-1,2)$ in the direction $\mathbf{v}=(2,2,1)$.
  \item Compute $f_x$ and $f_y$ for $f(x,y)=\tan(x+y)+\tan(x-y)$.
  \item For $F(K,L)=(\alpha K+\beta L)^r$ (linear-aggregate,
        $\alpha,\beta,r>0$), compute $F_K, F_L$, and verify
        $KF_K+LF_L=rF$.
  \item Compute $\nabla f$ for $f(x,y)=e^{-x}\sin(y)+e^x\cos(y)$.
  \item Compute the directional derivative of $f(x,y)=\sqrt{x^2+y^2}$ at
        $(3,4)$ in the direction $\mathbf{v}=(1,0)$.
  \item For $U(x,y)=\min\{2x,3y\}$, compute the marginal utilities on
        each smooth region $\{2x<3y\}$ and $\{2x>3y\}$.
  \item Compute $f_x, f_y$ for $f(x,y)=\dfrac{\ln(x+y)}{x-y}$ (assume
        $x>y>0$).
  \item For $F(K,L)=AK^\alpha L^\beta$ general Cobb--Douglas, derive
        $\mathrm{MRTS}_{LK}$ and show it depends only on the ratio $K/L$
        (not on scale).
  \item Compute the unit direction of steepest ascent for
        $f(x,y)=x^2-y^2$ at $(2,-1)$, and the maximum rate.
  \item Compute $f_x$ and $f_y$ for $f(x,y)=\cosh(x+y)+\sinh(x-y)$.
  \item For $f(x,y,z)=e^{x+y+z}$, compute $\nabla f$ at $(0,0,0)$ and the
        directional derivative in $\mathbf{v}=(1,1,1)$.
  \item Compute $f_x$ and $f_y$ for $f(x,y)=(\ln x)(\ln y)$ with
        $x,y>1$.
  \item For $F(K,L)=K+2K^{1/2}L^{1/2}+L$ (which equals
        $(K^{1/2}+L^{1/2})^2$), compute $F_K, F_L$ directly and confirm
        agreement with the CES form.
  \item Compute the directional derivative of $f(x,y)=x^2y+y^3$ at
        $(-1,2)$ in the direction $\mathbf{v}=(4,-3)$.
  \item For $U(x,y)=\alpha\ln x+\beta\ln y$ (log-Cobb--Douglas), derive
        $\mathrm{MRS}_{xy}$ and evaluate at $(x,y)=(\alpha,\beta)$.
  \item Compute $\nabla f$ for $f(x,y)=(x^2+y^2)^{-1/2}$ (Newtonian
        potential in 2D).
  \item Compute $f_x$ and $f_y$ for $f(x,y)=\arctan(xy)$.
  \item For $F(K,L)=(K^\rho+L^\rho)^{1/\rho}$ with $\rho<1$, $\rho\neq0$
        (CES with constant returns), compute $F_K$ and show
        $F_K=(K/F)^{\rho-1}$.
  \item Compute the directional derivative of $f(x,y,z)=xyz$ at
        $(2,3,-1)$ in the direction $\mathbf{v}=(2,-1,2)$.
\end{problems}

\section{Solutions}

\begin{solutions}
  \item $f_x=20x^3-4xy^3$, $f_y=-6x^2y^2+5y^4$.

  \item $F_K=12\cdot\frac14K^{-3/4}L^{3/4}=3(L/K)^{3/4}$;
        $F_L=12\cdot\frac34K^{1/4}L^{-1/4}=9(K/L)^{1/4}$. At $(16,81)$:
        $(L/K)^{3/4}=(81/16)^{3/4}=27/8$ so $F_K=81/8$;
        $(K/L)^{1/4}=(16/81)^{1/4}=2/3$ so $F_L=6$.

  \item $f_x=2xe^{3y}$, $f_y=3x^2e^{3y}$.

  \item $\nabla f=(12x^2y-y^4,\ 4x^3-4xy^3)^T$.

  \item $f_x=2x-4y$, $f_y=-4x+3y^2$. At $(2,1)$: $\nabla f=(0,-5)^T$.
        Normalise: $\|\mathbf{v}\|=\sqrt{25+144}=13$,
        $\mathbf{u}=(5/13,-12/13)^T$. Then
        \[
          D_{\mathbf{u}}f=0\cdot\frac{5}{13}+(-5)\cdot
            \left(-\frac{12}{13}\right)=\frac{60}{13} .
        \]

  \item $F_K=6\cdot\frac23K^{-1/3}L^{1/3}=4(L/K)^{1/3}$,
        $F_L=2(K/L)^{2/3}$. At $(27,8)$: $(L/K)^{1/3}=2/3$ so
        $\mathrm{MP}_K=8/3$; $(K/L)^{2/3}=(27/8)^{2/3}=9/4$ so
        $\mathrm{MP}_L=9/2$.

  \item Chain rule with $u=3x-2y$: $f_x=4(3x-2y)^3\cdot3=12(3x-2y)^3$,
        $f_y=-8(3x-2y)^3$.

  \item $f_x=y^2e^{xy^2}$, $f_y=2xye^{xy^2}$.

  \item $F_K=15K^{-0.4}L^{0.5}$, $F_L=12.5K^{0.6}L^{-0.5}$.

  \item $f_x=2x-z^2e^x$, $f_y=3z$, $f_z=3y-2ze^x$. So
        $\nabla f=(2x-z^2e^x,\ 3z,\ 3y-2ze^x)^T$.

  \item $f_x=e^x\ln(1+y)$, $f_y=e^x/(1+y)$. At $(0,0)$:
        $\nabla f=(0,1)^T$. Normalise: $\|\mathbf{v}\|=\sqrt5$,
        $\mathbf{u}=(1/\sqrt5,2/\sqrt5)^T$. So
        \[
          D_{\mathbf{u}}f=0\cdot\frac{1}{\sqrt5}+1\cdot\frac{2}{\sqrt5}
            =\frac{2\sqrt5}{5} .
        \]

  \item $U_x=2/x$, $U_y=5/y$. $\mathrm{MRS}_{xy}=U_x/U_y
        =(2/x)/(5/y)=2y/(5x)$. At $(3,5)$: $\mathrm{MRS}=10/15=2/3$.

  \item For $x+y\neq0$: $f(x,y)=x-y$, so $f_x=1$, $f_y=-1$.

  \item With $u=x^2y+1$: $f_x=2xy/(x^2y+1)$, $f_y=x^2/(x^2y+1)$.

  \item $U_{x_1}=\frac12x_1^{-1/2}x_2^{1/3}x_3^{1/6}$,
        $U_{x_2}=\frac13x_1^{1/2}x_2^{-2/3}x_3^{1/6}$,
        $U_{x_3}=\frac16x_1^{1/2}x_2^{1/3}x_3^{-5/6}$. Equivalently
        $U_{x_i}=(\alpha_i/x_i)U$ with $\alpha=(1/2,1/3,1/6)$.

  \item $f_x=y$, $f_y=x-2z$, $f_z=-2y+2z$. At $(1,2,3)$:
        $\nabla f=(2,-5,2)^T$. Normalise: $\|\mathbf{v}\|=\sqrt{1+4+4}=3$,
        $\mathbf{u}=(1/3,2/3,2/3)^T$. Then
        $D_{\mathbf{u}}f=(2-10+4)/3=-4/3$.

  \item $F_K=0.25K^{-0.75}L^{0.75}$, $F_L=0.75K^{0.25}L^{-0.25}$. So
        $\mathrm{MRTS}_{LK}=F_L/F_K=3K/L$. At $(16,4)$:
        $\mathrm{MRTS}=48/4=12$.

  \item Expand: $f=x^3-3x^2y^2+xy-3y^3$. Hence $f_x=3x^2-6xy^2+y$,
        $f_y=-6x^2y+x-9y^2$.

  \item With $u=-y/x$: $\partial u/\partial x=y/x^2$,
        $\partial u/\partial y=-1/x$. Product rule:
        \[
          f_x=e^{-y/x}+x\cdot e^{-y/x}\cdot\frac{y}{x^2}
            =e^{-y/x}\left(1+\frac{y}{x}\right), \qquad
          f_y=x\cdot e^{-y/x}\cdot\left(-\frac1x\right)=-e^{-y/x} .
        \]

  \item Let $g=3K^{1/2}+2L^{1/2}$, so $F=g^2$. Then
        \[
          F_K=2g\cdot\frac32K^{-1/2}=\frac{3g}{K^{1/2}}
            =\frac{3(3K^{1/2}+2L^{1/2})}{K^{1/2}}, \qquad
          F_L=2g\cdot L^{-1/2}=\frac{2(3K^{1/2}+2L^{1/2})}{L^{1/2}} .
        \]

  \item With $u=y/x$, $f=\arctan u$, $f'=1/(1+u^2)$:
        \[
          f_x=\frac{1}{1+y^2/x^2}\cdot\left(-\frac{y}{x^2}\right)
            =\frac{-y}{x^2+y^2}, \qquad
          f_y=\frac{x}{x^2+y^2} .
        \]
        So $\nabla f=(-y/(x^2+y^2),\ x/(x^2+y^2))^T$.

  \item $f_K=\frac12K^{-1/2}L^{1/2}$, $f_L=\frac12K^{1/2}L^{-1/2}$. At
        $(4,9)$: $\nabla f=(\frac12\cdot\frac32,\ \frac12\cdot\frac23)^T
        =(3/4,1/3)^T$. Then $\|\nabla f\|^2=9/16+1/9=97/144$, so the
        maximum rate is $\sqrt{97}/12\approx0.821$ in direction
        $\mathbf{u}^*=(9,4)^T/\sqrt{97}$.

  \item $F=(K^{1/2}+L^{1/2})^2$. Replacing $(K,L)\to(\lambda K,\lambda L)$
        gives $\lambda F$, so $r=1$. With $g=K^{1/2}+L^{1/2}$:
        $F_K=2g\cdot\frac12K^{-1/2}=g/K^{1/2}$, $F_L=g/L^{1/2}$. Hence
        \[
          KF_K+LF_L=K^{1/2}g+L^{1/2}g=g(K^{1/2}+L^{1/2})=g^2=F. \checkmark
        \]

  \item Quotient rule, $D=x^2+y^2$:
        \[
          f_x=\frac{2D-(2x+3y)(2x)}{D^2}=\frac{-2x^2-6xy+2y^2}{(x^2+y^2)^2},
          \qquad f_y=\frac{3x^2-4xy-3y^2}{(x^2+y^2)^2} .
        \]

  \item Treat denominator as constant for $f_x$: $f_x=e^x/(1+e^y)$. For
        $f_y$:
        \[
          f_y=e^x\cdot\frac{-e^y}{(1+e^y)^2}=-\frac{e^{x+y}}{(1+e^y)^2} .
        \]

  \item Let $S=\alpha K^\rho+(1-\alpha)L^\rho$, $F=AS^{\nu/\rho}$.
        Treating $K$ constant, $\partial S/\partial L=\rho(1-\alpha)L^{\rho-1}$:
        \[
          F_L=A\cdot\frac{\nu}{\rho}S^{\nu/\rho-1}\cdot\rho(1-\alpha)L^{\rho-1}
            =\nu(1-\alpha)AL^{\rho-1}S^{\nu/\rho-1} .
        \]

  \item $f=\ln(x^2+y^2)$, $f_x=2x/(x^2+y^2)$, $f_y=2y/(x^2+y^2)$. At
        $(3,4)$: $x^2+y^2=25$, $\nabla f=(6/25,8/25)^T$. Normalise:
        $\|\mathbf{v}\|=\sqrt2$, $\mathbf{u}=(-1/\sqrt2,1/\sqrt2)^T$. Then
        \[
          D_{\mathbf{u}}f=\frac{6}{25}\cdot\left(-\frac{1}{\sqrt2}\right)
            +\frac{8}{25}\cdot\frac{1}{\sqrt2}=\frac{2}{25\sqrt2}
            =\frac{\sqrt2}{25} .
        \]

  \item $f_x=2xyz^3-4y^2z$, $f_y=x^2z^3-8xyz$,
        $f_z=3x^2yz^2-4xy^2+5$.

  \item Chain rule, $u=4x^2+9y^2$:
        \[
          f_x=\frac{8x}{2\sqrt u}=\frac{4x}{\sqrt{4x^2+9y^2}}, \qquad
          f_y=\frac{9y}{\sqrt{4x^2+9y^2}} .
        \]

  \item With $g=K^{1/2}+L^{1/2}$, $F=g^4$,
        $F_K=4g^3\cdot\frac12K^{-1/2}=2g^3/K^{1/2}$. At $(4,9)$: $g=5$,
        $K^{1/2}=2$, so $F_K=2\cdot125/2=125$.

  \item Let $r^2=x^2+y^2+z^2$. Then $f=\ln(r^2)$ and $f_{x_i}=2x_i/r^2$.
        So $\nabla f=(2/(x^2+y^2+z^2))(x,y,z)^T$.

  \item Use $\ln\sqrt{u}=\frac12\ln u$: $f=\frac12\ln(x^2+y^2)$. So
        $f_x=x/(x^2+y^2)$, $f_y=y/(x^2+y^2)$.

  \item $\nabla C=(4,9)^T$. With $Q=x^{1/2}y^{1/2}$: $Q_x=\frac12(y/x)^{1/2}$
        and $Q_y=\frac12(x/y)^{1/2}$, so
        \[
          \nabla Q=\tfrac12\left((y/x)^{1/2},\,(x/y)^{1/2}\right)^{T} .
        \]
        The condition $\nabla C=\lambda\nabla Q$ requires
        $4/9=Q_x/Q_y=y/x$, i.e.\ $y=4x/9$. Along this locus
        $\nabla Q\propto(2/3,3/2)^T\propto(4,9)^T=\nabla C$. \checkmark

  \item Chain rule: $f_x=5(1+xy)^4\cdot y=5y(1+xy)^4$,
        $f_y=5x(1+xy)^4$.

  \item $f_x=3x^2-2y$, $f_y=-2x+2y$. At $(1,2)$:
        $\nabla f=(3-4,-2+4)^T=(-1,2)^T$. Normalise: $\|\mathbf{v}\|=5$,
        $\mathbf{u}=(3/5,4/5)^T$. Then
        \[
          D_{\mathbf{u}}f=-1\cdot\frac35+2\cdot\frac45
            =\frac{-3+8}{5}=1 .
        \]

  \item $F_K=10\cdot0.3K^{-0.7}L^{0.7}=3(L/K)^{0.7}$;
        $F_L=7(K/L)^{0.3}$. At $K=L=100$: $F_K=3$, $F_L=7$. And
        $F=10\cdot100^{0.3}\cdot100^{0.7}=10\cdot100=1000$. Then
        $KF_K+LF_L=100\cdot3+100\cdot7=1000=F$. \checkmark (homogeneity
        degree $r=0.3+0.7=1$).

  \item $f_x=y\cos(xy)-\sin(x-y)$, $f_y=x\cos(xy)+\sin(x-y)$.

  \item $f_x=2xy-y^2+3$, $f_y=x^2-2xy-2$. At $(1,-1)$: $f_x=-2-1+3=0$,
        $f_y=1+2-2=1$. So $\nabla f(1,-1)=(0,1)^T$.

  \item $U_x=\alpha x^{\alpha-1}y^{1-\alpha}$,
        $U_y=(1-\alpha)x^\alpha y^{-\alpha}$. So
        \[
          \mathrm{MRS}_{xy}=\frac{U_x}{U_y}
            =\frac{\alpha x^{\alpha-1}y^{1-\alpha}}{(1-\alpha)x^\alpha y^{-\alpha}}
            =\frac{\alpha}{1-\alpha}\cdot\frac{y}{x} . \checkmark
        \]

  \item $f_x=ye^{xy}$, $f_y=xe^{xy}$. At $(1,0)$: $\nabla f=(0,1)^T$.
        Normalise: $\|\mathbf{v}\|=5$, $\mathbf{u}=(-3/5,4/5)^T$. Then
        $D_{\mathbf{u}}f=0\cdot(-3/5)+1\cdot(4/5)=4/5$.

  \item Write $f=e^{y\ln x}$. Then $f_x=e^{y\ln x}\cdot y/x=yx^{y-1}$ and
        $f_y=e^{y\ln x}\cdot\ln x=x^y\ln x$.

  \item On $\{2K<3L\}$ the binding factor is capital: $F=2K$, so
        $F_K=2$, $F_L=0$. On $\{2K>3L\}$ labour binds: $F=3L$, so
        $F_K=0$, $F_L=3$. (Non-differentiable on $\{2K=3L\}$.)

  \item $f_x=yz+2x$, $f_y=xz$, $f_z=xy-2z$. At $(2,-1,3)$:
        $f_x=-3+4=1$, $f_y=6$, $f_z=-2-6=-8$. So
        $\nabla f=(1,6,-8)^T$.

  \item $\nabla f=(2x,8y)^T$. At $(3,1)$: $\nabla f=(6,8)^T$. Steepest
        descent is $-\nabla f/\|\nabla f\|$. $\|\nabla f\|=10$, so unit
        descent direction is $(-3/5,-4/5)^T$, with rate of decrease
        $-\|\nabla f\|=-10$ (i.e.\ $f$ decreases at rate 10).

  \item $f_x=e^y-ye^x$, $f_y=xe^y-e^x$.

  \item Quotient rule with $D=x^2+y^2$:
        \[
          f_x=\frac{D-x\cdot2x}{D^2}=\frac{y^2-x^2}{(x^2+y^2)^2}, \qquad
          f_y=\frac{-x\cdot2y}{D^2}=\frac{-2xy}{(x^2+y^2)^2} .
        \]

  \item $U_x=1$, $U_y=1/\sqrt{y}$. So $\mathrm{MRS}_{xy}=U_x/U_y
        =\sqrt{y}$ (depends only on $y$, quasi-linear utility).

  \item $f_x=2x$, $f_y=2y$, $f_z=-2z$. At $(2,-1,2)$:
        $\nabla f=(4,-2,-4)^T$. Normalise: $\|\mathbf{v}\|=\sqrt{4+4+1}=3$,
        $\mathbf{u}=(2/3,2/3,1/3)^T$. Then
        \[
          D_{\mathbf{u}}f=4\cdot\frac23+(-2)\cdot\frac23+(-4)\cdot\frac13
            =\frac{8-4-4}{3}=0 .
        \]
        (We are moving tangent to the level surface.)

  \item $f_x=\sec^2(x+y)+\sec^2(x-y)$, $f_y=\sec^2(x+y)-\sec^2(x-y)$.

  \item Let $S=\alpha K+\beta L$, $F=S^r$. Then $F_K=rS^{r-1}\alpha$,
        $F_L=rS^{r-1}\beta$. Hence
        \[
          KF_K+LF_L=rS^{r-1}(\alpha K+\beta L)=rS^{r-1}\cdot S=rS^r=rF.
            \checkmark
        \]

  \item $\nabla f=(-e^{-x}\sin y+e^x\cos y,\ e^{-x}\cos y-e^x\sin y)^T$.

  \item $f_x=x/\sqrt{x^2+y^2}$, $f_y=y/\sqrt{x^2+y^2}$. At $(3,4)$:
        $\sqrt{x^2+y^2}=5$, $\nabla f=(3/5,4/5)^T$. The direction
        $\mathbf{v}=(1,0)$ is already unit. So
        $D_{\mathbf{u}}f=3/5\cdot1+4/5\cdot0=3/5$.

  \item On $\{2x<3y\}$: $U=2x$, so $U_x=2$, $U_y=0$. On $\{2x>3y\}$:
        $U=3y$, so $U_x=0$, $U_y=3$. (Non-differentiable on the ray
        $2x=3y$.)

  \item Let $L=\ln(x+y)$, $D=x-y$. By the quotient rule:
        \[
          f_x=\frac{(1/(x+y))(x-y)-\ln(x+y)\cdot1}{(x-y)^2}
            =\frac{1/(x+y)\cdot(x-y)-\ln(x+y)}{(x-y)^2} .
        \]
        Simplifying
        $f_x=\dfrac{(x-y)/(x+y)-\ln(x+y)}{(x-y)^2}$. Similarly,
        \[
          f_y=\frac{(1/(x+y))(x-y)-\ln(x+y)\cdot(-1)}{(x-y)^2}
            =\frac{(x-y)/(x+y)+\ln(x+y)}{(x-y)^2} .
        \]

  \item $F_K=A\alpha K^{\alpha-1}L^\beta$, $F_L=A\beta K^\alpha L^{\beta-1}$.
        So
        \[
          \mathrm{MRTS}_{LK}=\frac{F_L}{F_K}=\frac{\beta}{\alpha}\cdot
            \frac{K}{L} .
        \]
        This depends only on the ratio $K/L$ and the exponent ratio
        $\beta/\alpha$: it is scale-invariant.

  \item $\nabla f=(2x,-2y)^T$. At $(2,-1)$: $\nabla f=(4,2)^T$. Direction
        of steepest ascent: $\mathbf{u}^*=(4,2)/\sqrt{20}=(2,1)/\sqrt5$.
        Maximum rate $=\|\nabla f\|=\sqrt{20}=2\sqrt5$.

  \item $f_x=\sinh(x+y)+\cosh(x-y)$, $f_y=\sinh(x+y)-\cosh(x-y)$.

  \item $\nabla f=(e^{x+y+z},e^{x+y+z},e^{x+y+z})^T$. At $(0,0,0)$:
        $\nabla f=(1,1,1)^T$. Normalise $\mathbf{v}=(1,1,1)$:
        $\|\mathbf{v}\|=\sqrt3$, $\mathbf{u}=(1,1,1)/\sqrt3$. Then
        $D_{\mathbf{u}}f=(1+1+1)/\sqrt3=\sqrt3$.

  \item $f_x=(\ln y)/x$, $f_y=(\ln x)/y$.

  \item $F_K=1+L^{1/2}/K^{1/2}$, $F_L=1+K^{1/2}/L^{1/2}$. From
        $F=(K^{1/2}+L^{1/2})^2$ with $g=K^{1/2}+L^{1/2}$:
        $F_K=g/K^{1/2}=(K^{1/2}+L^{1/2})/K^{1/2}=1+L^{1/2}/K^{1/2}$.
        \checkmark Same for $F_L$.

  \item $f_x=2xy$, $f_y=x^2+3y^2$. At $(-1,2)$: $\nabla f=(-4,1+12)^T
        =(-4,13)^T$. Normalise: $\|\mathbf{v}\|=5$,
        $\mathbf{u}=(4/5,-3/5)^T$. Then
        \[
          D_{\mathbf{u}}f=-4\cdot\frac45+13\cdot\left(-\frac35\right)
            =-\frac{16+39}{5}=-11 .
        \]

  \item $U_x=\alpha/x$, $U_y=\beta/y$.
        $\mathrm{MRS}_{xy}=\dfrac{\alpha/x}{\beta/y}=\dfrac{\alpha y}{\beta x}$.
        At $(\alpha,\beta)$: $\mathrm{MRS}=(\alpha\cdot\beta)/(\beta\cdot\alpha)=1$.

  \item $f=(x^2+y^2)^{-1/2}$. Chain rule:
        \[
          f_x=-\frac12(x^2+y^2)^{-3/2}\cdot2x=\frac{-x}{(x^2+y^2)^{3/2}},
          \qquad f_y=\frac{-y}{(x^2+y^2)^{3/2}} .
        \]
        So $\nabla f=-(x,y)^T/(x^2+y^2)^{3/2}$.

  \item Chain rule with $u=xy$:
        \[
          f_x=\frac{y}{1+x^2y^2}, \qquad f_y=\frac{x}{1+x^2y^2} .
        \]

  \item Let $S=K^\rho+L^\rho$, $F=S^{1/\rho}$. Then
        $F_K=\frac1\rho S^{1/\rho-1}\cdot\rho K^{\rho-1}
        =K^{\rho-1}S^{1/\rho-1}$. Since $S^{1/\rho}=F$, we have $S=F^\rho$
        and so $S^{1/\rho-1}=F^{1-\rho}\cdot F^{-1}\cdot F=F^{1-\rho}$.
        Thus $F_K=K^{\rho-1}F^{1-\rho}=(F/K)^{1-\rho}=(K/F)^{\rho-1}$.
        \checkmark

  \item $f_x=yz$, $f_y=xz$, $f_z=xy$. At $(2,3,-1)$:
        $\nabla f=(3\cdot(-1),2\cdot(-1),2\cdot3)^T=(-3,-2,6)^T$.
        Normalise $\mathbf{v}=(2,-1,2)$: $\|\mathbf{v}\|=3$,
        $\mathbf{u}=(2/3,-1/3,2/3)^T$. Then
        \[
          D_{\mathbf{u}}f=-3\cdot\frac23+(-2)\cdot\left(-\frac13\right)
            +6\cdot\frac23=-2+\frac23+4=2+\frac23=\frac83 .
        \]
\end{solutions}

\section{Additional Problems}

\begin{problems}
  \item Compute $f_x$ and $f_y$ for $f(x,y)=3x^5-2x^3y^2+4y^4$.
  \item For $F(K,L)=20\,K^{1/3}L^{2/3}$, find $F_K$ and $F_L$ at
        $(K,L)=(8,27)$.
  \item Compute $f_x$ and $f_y$ for $f(x,y)=e^{xy}$.
  \item Compute $\nabla f$ for $f(x,y)=2x^2y-3xy^3+5$.
  \item Compute the directional derivative of $f(x,y)=x^2+2xy-y^2$ at
        $(1,2)$ in the direction $\mathbf{v}=(3,4)$.
  \item For $F(K,L)=8K^{3/4}L^{1/4}$, compute $\mathrm{MP}_K$ and
        $\mathrm{MP}_L$ at $(K,L)=(16,16)$.
  \item Compute $f_x$ and $f_y$ for $f(x,y)=\ln(2x+3y)$.
  \item For $U(x,y)=3\ln(x)+4\ln(y)$, derive $\mathrm{MRS}_{xy}$ and
        evaluate at $(2,3)$.
  \item Compute $f_x$ and $f_y$ for $f(x,y)=(x+y^2)^4$.
  \item Compute $\nabla f$ for $f(x,y,z)=xy+yz+zx$.
  \item For a CES production $F(K,L)=5(K^{1/3}+L^{1/3})^3$, find $F_K$.
  \item Compute the directional derivative of $f(x,y)=x^2y$ at $(2,1)$
        in the direction $\mathbf{v}=(1,-1)$.
  \item For $U(x,y)=x^\alpha y^\beta$ (general Cobb--Douglas,
        $\alpha,\beta>0$), derive $\mathrm{MRS}_{xy}$ and show it depends
        only on the ratio $y/x$.
  \item Compute $f_x$ and $f_y$ for $f(x,y)=\dfrac{x-y}{x+y}$.
  \item For $f(K,L)=K^{1/2}L^{1/2}$ at $(K,L)=(16,9)$, find the unit
        direction of steepest ascent and the maximum rate of increase.
  \item Classify the open positive orthant
        $Q=\{(x,y)\in\R^2 : x>0,\ y>0\}$: is it open, closed, or
        neither? Is it bounded? Is it compact?
  \item Classify the budget set
        $B=\{(x,y)\in\R^2 : 2x+3y\le12,\ x\ge0,\ y\ge0\}$: open,
        closed, or neither? Bounded? Compact?
  \item Classify $M=\{(x,y)\in\R^2 : xy\ge1,\ x>0\}$: open, closed, or
        neither? Bounded? Compact? \emph{Hint:} check whether a
        sequence in $M$ can converge to a point with $x=0$.
  \item Classify the open disc
        $D=\{(x,y)\in\R^2 : x^2+y^2<4\}$: open, closed, or neither?
        Bounded? Compact?
  \item Classify (a) the half-open strip
        $S=\{(x,y)\in\R^2 : 0\le x<1\}$ and (b) the finite set
        $E=\{(0,0),\ (1,2),\ (3,1)\}$: open, closed, or neither?
        Bounded? Compact?
  \item Does the Weierstrass theorem guarantee a maximiser for the
        continuous function $u(x,y)=x^{1/2}y^{1/2}$ (a) on the budget
        set $B=\{(x,y): 2x+3y\le12,\ x\ge0,\ y\ge0\}$, (b) on the open
        positive orthant $Q=\{(x,y): x>0,\ y>0\}$? Justify each
        answer.
\end{problems}

In problems 22--29, for each function: state the domain; compute the
partials and the set on which they exist; decide, with a justification,
where $f$ is differentiable; and describe the graph near the points where
differentiability fails.

\begin{problems}[resume]
  \item $f(x,y)=\sqrt{(x-1)^{2}+(y+2)^{2}}$.
  \item $f(x,y)=\sqrt{4x^{2}+(y-3)^{2}}$.
  \item $f(x,y)=|x-2|+|y|$.
  \item $f(x,y)=\min\{2x,\,5y\}$ on $x,y>0$.
  \item $f(x,y)=\dfrac{3xy}{2x^{2}+y^{2}}$ for $(x,y)\neq(0,0)$, with
        $f(0,0)=0$. Also decide whether $f$ is continuous at the origin.
  \item $f(x,y)=\sqrt{x^{4}+y^{4}}$. The formula for $f_{x}$ breaks down
        at the origin; decide what really happens there.
  \item $f(x,y)=\sqrt{9-x^{2}-y^{2}}$.
  \item $f(x,y)=|x-y|$.
\end{problems}

In problems 30--37, find the tangent plane of $f$ at $P$ (the tangent
hyperplane, for three variables), then evaluate both $f$ and the plane at
the nearby point $Q$ and compare. Give both values to six decimals.

\begin{problems}[resume]
  \item $f(x,y)=\ln(2x+y^{2})$, \;$P=(0,1)$, \;$Q=(0.05,\,0.98)$.
  \item $f(x,y)=\mathrm{e}^{x-2y}$, \;$P=(2,1)$, \;$Q=(2.1,\,1.02)$.
  \item $f(x,y)=xy^{2}\,\mathrm{e}^{x-1}$, \;$P=(1,2)$,
        \;$Q=(0.98,\,2.03)$.
  \item $f(x,y)=\dfrac{x+y}{x-y}$, \;$P=(3,1)$, \;$Q=(3.1,\,0.9)$.
  \item $F(K,L)=K^{1/4}L^{3/4}$, \;$P=(16,1)$, \;$Q=(16.5,\,1.1)$; read
        the plane as an output estimate.
  \item $f(x,y,z)=\ln(x+2y+3z)$, \;$P=(2,-1,\tfrac13)$,
        \;$Q=(2.02,\,-0.99,\,0.35)$.
  \item $f(x,y,z)=x\,\mathrm{e}^{yz}$, \;$P=(2,0,5)$,
        \;$Q=(2.05,\,0.01,\,4.9)$.
  \item $f(x,y,z)=\dfrac{xy}{z}$, \;$P=(2,3,6)$, \;$Q=(2.1,\,2.9,\,6.2)$.
\end{problems}

In problems 38--46, compute the partial derivatives, and give the gradient
(a column) and the derivative (a row) at the stated point.

\begin{problems}[resume]
  \item $f(x_{1},\ldots,x_{n})=\sum_{i=1}^{n}\ln(1+x_{i}^{2})$, at
        $(1,1,\ldots,1)$.
  \item $f(x_{1},\ldots,x_{n})=\Bigl(\sum_{i=1}^{n}x_{i}\Bigr)^{3}$; then
        for $n=4$ at $(1,2,3,4)$.
  \item $f(x_{1},\ldots,x_{n})=\sum_{i=1}^{n}x_{i}^{\,i+1}$, at
        $(1,1,\ldots,1)$; then for $n=3$ at $(2,1,-1)$.
  \item $f(x_{1},\ldots,x_{n})=\sum_{j=1}^{n}\dfrac{j^{2}}{x_{j}^{3}}$ with
        all $x_{j}>0$, at $(1,1,\ldots,1)$.
  \item $f(x_{1},\ldots,x_{n})=\sum_{i=1}^{n}x_{i}^{2}
        -\dfrac1n\Bigl(\sum_{i=1}^{n}x_{i}\Bigr)^{2}$; show that the
        gradient vanishes exactly when all $x_{i}$ are equal, then evaluate
        for $n=3$ at $(1,2,6)$.
  \item $f(x,y,z)=y^{2}\ln(x^{2}+3yz)$, at $P=(1,1,1)$.
  \item $f(x,y,z)=x^{3}\,\mathrm{e}^{y^{2}-2xz}$, at $P=(1,0,1)$.
  \item $f(x,y)=2^{x}3^{y}+5^{xy}$, at $(1,1)$.
  \item $f(x_{1},\ldots,x_{n})=\sum_{i=1}^{n}3^{\,i x_{i}}$, at
        $(0,0,\ldots,0)$.
\end{problems}

In problems 47--52, decide differentiability directly from the definition,
by forming the remainder and estimating it, as in the third worked example.

\begin{problems}[resume]
  \item $f(x,y)=3xy$, at every point $(a,b)$.
  \item $f(x,y)=x^{2}+y^{2}$, first at the origin, then at every point.
  \item $f(x,y)=xy^{2}$, at every point $(a,b)$.
  \item $f(x,y)=y\,\mathrm{e}^{3x}$, at every point $(a,b)$.
  \item At the origin: (a) $f(x,y)=(x^{2}+4y^{2})^{1/3}$;
        (b) $g(x,y)=(x^{2}+y^{2})^{3/4}$.
  \item $f(x,y)=|xy|$: at the origin, and at the point $(2,0)$.
\end{problems}

In problems 53--54, $f(0,0)=0$. Show with the difference quotient that both
partials exist at the origin and are $0$; compute the limit along every
straight line through the origin and along the stated parabola; and decide
continuity and differentiability at the origin.

\begin{problems}[resume]
  \item $f(x,y)=\dfrac{x^{2}y}{x^{4}+y^{2}}$ for $(x,y)\neq(0,0)$; parabola
        $y=x^{2}$.
  \item $f(x,y)=\dfrac{2xy^{2}}{x^{2}+y^{4}}$ for $(x,y)\neq(0,0)$;
        parabola $x=y^{2}$.
\end{problems}

Problems 55--56 refer to the three surfaces below, each drawn over a square
around its special point.

\begin{gallery}
\begin{tikzpicture}
  \begin{axis}[surfaxis, title={(a)}, z post scale=1.1]
    \addplot3[surf, domain=-1:3, y domain=-2:2]{(x-1)^2+2*y^2};
  \end{axis}
\end{tikzpicture}\hfill
\begin{tikzpicture}
  \begin{axis}[surfaxis, title={(b)}, z post scale=1.1]
    \addplot3[surf, domain=-1:3, y domain=-2:2]{sqrt((x-1)^2+4*y^2)};
  \end{axis}
\end{tikzpicture}\hfill
\begin{tikzpicture}
  \begin{axis}[surfaxis, title={(c)}, z post scale=1.1]
    \addplot3[surf, domain=-2:2, y domain=-2:2]{2-abs(x)-abs(y)};
  \end{axis}
\end{tikzpicture}
\end{gallery}

\begin{problems}[resume]
  \item For each surface (a)--(c), say from the picture alone where the
        function is differentiable and where it is not, and why.
  \item Surface (b) has its tip above the point $(1,0)$, and its contours
        are ellipses twice as long in $x$ as in $y$. Which of the formulas
        $(x-1)^{2}+4y^{2}$, $\sqrt{(x-1)^{2}+4y^{2}}$ and
        $\sqrt{4(x-1)^{2}+y^{2}}$ can it be the graph of? Confirm the
        verdict of the previous problem for that formula with the
        difference quotient at $(1,0)$.
\end{problems}

\section{Additional Solutions}

\begin{solutions}
  \item $f_x=15x^4-6x^2y^2$, $f_y=-4x^3y+16y^3$.

  \item $F_K=20\cdot\frac13K^{-2/3}L^{2/3}=\frac{20}{3}(L/K)^{2/3}$;
        $F_L=20\cdot\frac23K^{1/3}L^{-1/3}=\frac{40}{3}(K/L)^{1/3}$. At
        $(8,27)$: $(L/K)^{2/3}=(27/8)^{2/3}=9/4$, so $F_K=15$;
        $(K/L)^{1/3}=(8/27)^{1/3}=2/3$, so
        $F_L=\frac{40}{3}\cdot\frac23=\frac{80}{9}$.

  \item $f_x=ye^{xy}$, $f_y=xe^{xy}$.

  \item $\nabla f=(4xy-3y^3,\ 2x^2-9xy^2)^T$.

  \item $f_x=2x+2y$, $f_y=2x-2y$. At $(1,2)$: $\nabla f=(6,-2)^T$.
        Normalise: $\|\mathbf{v}\|=5$, $\mathbf{u}=(3/5,4/5)^T$. Then
        \[
          D_{\mathbf{u}}f=6\cdot\frac35+(-2)\cdot\frac45
            =\frac{18-8}{5}=2 .
        \]

  \item $F_K=8\cdot\frac34K^{-1/4}L^{1/4}=6(L/K)^{1/4}$,
        $F_L=2(K/L)^{3/4}$. At $(16,16)$: $\mathrm{MP}_K=6$,
        $\mathrm{MP}_L=2$.

  \item $f_x=\dfrac{2}{2x+3y}$, $f_y=\dfrac{3}{2x+3y}$.

  \item $U_x=3/x$, $U_y=4/y$.
        $\mathrm{MRS}_{xy}=U_x/U_y=\dfrac{3y}{4x}$. At $(2,3)$:
        $\mathrm{MRS}=\dfrac98$.

  \item Chain: $f_x=4(x+y^2)^3$,
        $f_y=4(x+y^2)^3\cdot2y=8y(x+y^2)^3$.

  \item $f_x=y+z$, $f_y=x+z$, $f_z=x+y$. So
        $\nabla f=(y+z,\ x+z,\ x+y)^T$.

  \item Let $g=K^{1/3}+L^{1/3}$, so $F=5g^3$. Then
        $F_K=5\cdot3g^2\cdot\frac13K^{-2/3}=5g^2K^{-2/3}
        =\dfrac{5(K^{1/3}+L^{1/3})^2}{K^{2/3}}$.

  \item $f_x=2xy$, $f_y=x^2$. At $(2,1)$: $\nabla f=(4,4)^T$. Normalise:
        $\|\mathbf{v}\|=\sqrt2$, $\mathbf{u}=(1/\sqrt2,-1/\sqrt2)^T$.
        Then
        \[
          D_{\mathbf{u}}f=4\cdot\frac{1}{\sqrt2}
            +4\cdot\left(-\frac{1}{\sqrt2}\right)=0 .
        \]
        (Tangent to the level curve.)

  \item $U_x=\alpha x^{\alpha-1}y^\beta$,
        $U_y=\beta x^\alpha y^{\beta-1}$. So
        \[
          \mathrm{MRS}_{xy}=\frac{U_x}{U_y}
            =\frac{\alpha}{\beta}\cdot\frac{y}{x} .
        \]
        Depends only on $y/x$ (and the exponent ratio), not on the
        scale. \checkmark

  \item Quotient rule with $D=x+y$:
        \[
          f_x=\frac{D-(x-y)}{D^2}=\frac{2y}{(x+y)^2}, \qquad
          f_y=\frac{-D-(x-y)}{D^2}=\frac{-2x}{(x+y)^2} .
        \]

  \item $f_K=\frac12K^{-1/2}L^{1/2}$, $f_L=\frac12K^{1/2}L^{-1/2}$. At
        $(16,9)$: $f_K=\frac12\cdot\frac34=3/8$;
        $f_L=\frac12\cdot\frac43=2/3$. So $\nabla f=(3/8,2/3)^T$.
        Magnitude:
        $\|\nabla f\|=\sqrt{9/64+4/9}=\sqrt{81/576+256/576}
        =\sqrt{337/576}=\sqrt{337}/24\approx0.765$. Unit direction of
        steepest ascent:
        $\mathbf{u}^*=(3/8,2/3)^T/(\sqrt{337}/24)=(9,16)^T/\sqrt{337}$.

  \item \textbf{Open.} For any $(x,y)\in Q$ the ball of radius
        $r=\min\{x,y\}>0$ around $(x,y)$ stays inside $Q$, so every
        point is interior. \textbf{Not closed:} the boundary points
        $(x,0)$ and $(0,y)$ with $x,y\ge0$ (the coordinate half-axes,
        including the origin) do not belong to $Q$. \textbf{Not
        bounded:} $Q$ contains $(t,t)$ for every $t>0$, and
        $\|(t,t)\|=t\sqrt2\to\infty$. Hence $Q$ is open, unbounded and
        \textbf{not compact}.

  \item \textbf{Closed.} $B$ is the intersection of the three sets
        $\{2x+3y\le12\}$, $\{x\ge0\}$ and $\{y\ge0\}$, each defined by
        a non-strict inequality on a continuous function; each is
        closed, and an intersection of closed sets is closed (every
        boundary point of $B$ satisfies all three inequalities, so it
        lies in $B$). \textbf{Not open:} $(0,0)\in B$, but every ball
        around the origin contains points with negative coordinates,
        which are outside $B$; so $(0,0)$ is not interior.
        \textbf{Bounded:} from the constraints, $0\le x\le6$ and
        $0\le y\le4$, so $B\subseteq B_8(\mathbf{0})$. Closed and
        bounded in $\R^2$: $B$ is \textbf{compact}.

  \item First note the constraint $x>0$ together with $xy\ge1$ forces
        $y\ge1/x>0$. \textbf{Closed:} suppose
        $(x_n,y_n)\in M$ with $(x_n,y_n)\to(a,b)$. If $a=0$, then
        $y_n\to b$ finite would give $x_ny_n\to0<1$, contradicting
        $x_ny_n\ge1$; so $a>0$, and passing to the limit in
        $x_ny_n\ge1$ gives $ab\ge1$, hence $(a,b)\in M$. $M$ contains
        all its limit (boundary) points, so it is closed: even
        though it is written with a strict inequality, that inequality
        is never active. \textbf{Not open:} $(1,1)\in M$ lies on the
        hyperbola $xy=1$, and every ball around it contains points
        with $xy<1$. \textbf{Not bounded:} $(t,t)\in M$ for all
        $t\ge1$. Hence $M$ is closed, unbounded and \textbf{not
        compact}.

  \item \textbf{Open.} For $(x,y)\in D$ put
        $r=2-\sqrt{x^2+y^2}>0$; the ball $B_r(x,y)$ lies in $D$ by the
        triangle inequality. \textbf{Not closed:} the boundary circle
        $\{x^2+y^2=4\}$ is not contained in $D$ (e.g.\ $(2,0)\notin
        D$). \textbf{Bounded:} $D\subseteq B_3(\mathbf{0})$. Bounded
        but not closed: $D$ is \textbf{not compact}.

  \item (a) \textbf{Neither open nor closed.} The points $(0,y)$
        belong to $S$ but are boundary points (every ball around them
        contains points with $x<0$), so $S$ is not open; the boundary
        points $(1,y)$ do not belong to $S$, so $S$ is not closed.
        \textbf{Unbounded:} $(0,t)\in S$ for all $t$. Hence
        \textbf{not compact}. (b) \textbf{Closed:} around any point
        not in $E$ there is a ball avoiding all three points (a
        finite set always contains all its boundary points:
        indeed $E$ is its own boundary). \textbf{Not open:} no ball
        around $(0,0)$ lies inside $E$. \textbf{Bounded:}
        $E\subseteq B_4(\mathbf{0})$. Closed and bounded: $E$ is
        \textbf{compact}.

  \item (a) \textbf{Yes.} $B$ is compact (closed and bounded, see the
        budget-set problem above) and $u$ is continuous on $B$, so by
        the Weierstrass theorem $u$ attains both a maximum and a
        minimum on $B$; a maximiser exists. (b) \textbf{No.} $Q$ is
        not compact (it is neither closed nor bounded), so Weierstrass
        does not apply: and indeed no maximiser exists:
        $u(t,t)=t\to\infty$ as $t\to\infty$, so $u$ is unbounded above
        on $Q$.

  \item \textbf{Domain:} the radicand is a sum of squares, never
        negative, so $D=\R^{2}$. \textbf{Partials:} where the radicand is
        positive, i.e.\ for $(x,y)\neq(1,-2)$, the chain rule gives
        \[
          f_{x}=\frac{x-1}{\sqrt{(x-1)^{2}+(y+2)^{2}}},\qquad
          f_{y}=\frac{y+2}{\sqrt{(x-1)^{2}+(y+2)^{2}}} .
        \]
        At $(1,-2)$ use the quotient:
        $\bigl(f(1+h,-2)-f(1,-2)\bigr)/h=\sqrt{h^{2}}/h=|h|/h$, which is $1$
        for $h>0$ and $-1$ for $h<0$, so it has no limit. $f_{x}(1,-2)$
        does not exist, and in the same way neither does $f_{y}(1,-2)$. The
        partials exist exactly on $\R^{2}\setminus\{(1,-2)\}$.
        \textbf{Differentiability:} on that open set the partials are
        continuous ($f$ is $C^{1}$), so $f$ is differentiable there. At
        $(1,-2)$ it is not, since differentiability would force the
        partials to exist. Note that $f$ \emph{is} continuous at $(1,-2)$:
        continuity is not enough. \textbf{Graph:} a circular cone with its
        vertex at $(1,-2,0)$, opening upwards with slope $1$ in every
        direction; the contours are circles of radius $c$ about $(1,-2)$.
        At the sharp tip there is no tangent plane.

  \item \textbf{Domain:} $\R^{2}$, the radicand being a sum of squares.
        \textbf{Partials:} for $(x,y)\neq(0,3)$,
        \[
          f_{x}=\frac{4x}{\sqrt{4x^{2}+(y-3)^{2}}},\qquad
          f_{y}=\frac{y-3}{\sqrt{4x^{2}+(y-3)^{2}}} .
        \]
        At $(0,3)$: $f(h,3)-f(0,3)=\sqrt{4h^{2}}=2|h|$, so the quotient
        $2|h|/h=\pm2$ has no limit; and $f(0,3+h)-f(0,3)=|h|$ gives
        $\pm1$. Neither partial exists at $(0,3)$; both exist everywhere
        else. \textbf{Differentiability:} $C^{1}$, hence differentiable, on
        $\R^{2}\setminus\{(0,3)\}$; not differentiable at $(0,3)$.
        \textbf{Graph:} an elliptic cone with apex $(0,3,0)$. The contours
        $4x^{2}+(y-3)^{2}=c^{2}$ are ellipses about $(0,3)$ with semi-axes
        $c/2$ in the $x$-direction and $c$ in the $y$-direction: slope $2$
        across $x$, slope $1$ along $y$, and a sharp point at the apex.

  \item \textbf{Domain:} $\R^{2}$. \textbf{Partials:} for $x\neq2$,
        $f_{x}=1$ if $x>2$ and $f_{x}=-1$ if $x<2$; at a point $(2,y_{0})$,
        $\bigl(f(2+h,y_{0})-f(2,y_{0})\bigr)/h=|h|/h=\pm1$, so $f_{x}$ does
        not exist anywhere on the line $x=2$. Likewise $f_{y}=\pm1$ (the
        sign of $y$) for $y\neq0$, and $f_{y}$ does not exist on the line
        $y=0$. Both partials exist exactly off the two lines $x=2$ and
        $y=0$. \textbf{Differentiability:} off the two lines $f$ is, near
        each point, one of the four linear functions $\pm(x-2)\pm y$, so it
        is differentiable there; on the lines it is not, since a partial
        is missing. \textbf{Graph:} an upside-down square pyramid with its
        vertex at $(2,0,0)$; the contours $|x-2|+|y|=c$ are diamonds with
        corners $(2\pm c,0)$ and $(2,\pm c)$, and the surface is creased
        above the lines $x=2$ and $y=0$.

  \item \textbf{Domain:} the open quadrant $x,y>0$, as given.
        \textbf{Partials:} where $2x<5y$ (above the ray
        $y=\tfrac25x$), $f=2x$, so $f_{x}=2$, $f_{y}=0$. Where $2x>5y$,
        $f=5y$, so $f_{x}=0$, $f_{y}=5$. On the ray, at $(x_{0},y_{0})$ with
        $2x_{0}=5y_{0}$:
        \[
          f(x_{0}+h,y_{0})-f(x_{0},y_{0})=\min\{2x_{0}+2h,\,5y_{0}\}-5y_{0}
          =\min\{2h,\,0\},
        \]
        so the quotient is $0$ for $h>0$ and $2$ for $h<0$: no limit, no
        $f_{x}$. Similarly $f(x_{0},y_{0}+k)-f(x_{0},y_{0})=\min\{0,5k\}$
        gives one-sided slopes $0$ and $5$: no $f_{y}$. The partials exist
        exactly off the ray $y=\tfrac25x$. \textbf{Differentiability:}
        differentiable off the ray (locally linear), not on it.
        \textbf{Graph:} two planes, $z=2x$ and $z=5y$, meeting in a ridge
        above the ray: a roof. The contours are L-shaped corners, the level
        $c$ having its corner at $(c/2,\,c/5)$ on the ray, which is the
        Leontief isoquant picture.

  \item \textbf{Domain:} the denominator $2x^{2}+y^{2}$ vanishes only at
        the origin, where $f$ is defined separately, so $D=\R^{2}$.
        \textbf{Partials off the origin}, quotient rule with
        $N=2x^{2}+y^{2}$:
        \[
          f_{x}=\frac{3y\,N-3xy\cdot4x}{N^{2}}
               =\frac{3y\,(y^{2}-2x^{2})}{(2x^{2}+y^{2})^{2}},\qquad
          f_{y}=\frac{3x\,N-3xy\cdot2y}{N^{2}}
               =\frac{3x\,(2x^{2}-y^{2})}{(2x^{2}+y^{2})^{2}} .
        \]
        \textbf{At the origin}, from the definition: $f(h,0)=0$ and
        $f(0,h)=0$ for all $h$, so both quotients are $0$ and
        $f_{x}(0,0)=f_{y}(0,0)=0$. The partials exist \emph{everywhere}.
        \textbf{Continuity:} on the line $y=mx$, $x\neq0$,
        \[
          f(x,mx)=\frac{3mx^{2}}{(2+m^{2})x^{2}}=\frac{3m}{2+m^{2}},
        \]
        a constant that depends on $m$; on $y=x$ it is $1\neq f(0,0)$. So
        $f$ is not continuous at the origin, and hence not differentiable
        there, even though both partials exist. Away from the origin it is
        $C^{1}$ and differentiable. \textbf{Graph:} each ray from the origin
        carries its own constant height $3m/(2+m^{2})$, ranging between
        $-\tfrac{3\sqrt2}{4}$ and $\tfrac{3\sqrt2}{4}$ (attained at
        $m=\mp\sqrt2$); all these heights meet above the origin, so the
        surface is a twisted sheet torn at that point.

  \item \textbf{Domain:} $\R^{2}$. \textbf{Partials} off the origin:
        $f_{x}=\dfrac{2x^{3}}{\sqrt{x^{4}+y^{4}}}$,
        $f_{y}=\dfrac{2y^{3}}{\sqrt{x^{4}+y^{4}}}$. At the origin the
        formula fails, so use the quotient:
        $f(h,0)=\sqrt{h^{4}}=h^{2}$, and $h^{2}/h=h\to0$. So
        $f_{x}(0,0)=0$, and likewise $f_{y}(0,0)=0$: the partials exist
        everywhere, despite the broken formula. \textbf{Differentiability
        at the origin:} with $\mathbf{a}=\mathbf{0}$ the remainder is
        $r(\mathbf{h})=\sqrt{h_{1}^{4}+h_{2}^{4}}$. Since
        $h_{1}^{4}+h_{2}^{4}\le(h_{1}^{2}+h_{2}^{2})^{2}$ (the missing
        cross term $2h_{1}^{2}h_{2}^{2}$ is non-negative),
        \[
          0\le\frac{r(\mathbf{h})}{\|\mathbf{h}\|}
          \le\frac{h_{1}^{2}+h_{2}^{2}}{\|\mathbf{h}\|}=\|\mathbf{h}\|\to0 .
        \]
        So $f$ is differentiable at the origin with tangent plane $z=0$,
        and, being $C^{1}$ elsewhere, differentiable on all of $\R^{2}$.
        \textbf{Graph:} a bowl with a flat bottom at the origin; its
        contours $x^{4}+y^{4}=c^{2}$ are rounded squares. Nothing fails.

  \item \textbf{Domain:} $9-x^{2}-y^{2}\ge0$, the closed disc of radius $3$.
        \textbf{Partials:} for $x^{2}+y^{2}<9$,
        \[
          f_{x}=\frac{-x}{\sqrt{9-x^{2}-y^{2}}},\qquad
          f_{y}=\frac{-y}{\sqrt{9-x^{2}-y^{2}}},
        \]
        continuous on the open disc. On the boundary circle the radicand
        is $0$, the formulas divide by $0$, and, more basically, the points
        are not interior: every neighbourhood of a boundary point contains
        points outside the disc, where $f$ is not defined, so the two-sided
        difference quotients needed for the partials do not exist. \textbf{Differentiability:}
        $f$ is differentiable exactly on the open disc $x^{2}+y^{2}<9$.
        \textbf{Graph:} the upper hemisphere of radius $3$. Its tangent
        plane is horizontal at the top $(0,0,3)$, tilts more and more
        steeply outwards, and becomes vertical at the rim, where
        $\|\grad f\|\to\infty$.

  \item \textbf{Domain:} $\R^{2}$. \textbf{Partials:} where $x>y$,
        $f=x-y$, so $f_{x}=1$, $f_{y}=-1$; where $x<y$, $f=y-x$, so
        $f_{x}=-1$, $f_{y}=1$. At a point $(x_{0},x_{0})$ of the diagonal,
        $f(x_{0}+h,x_{0})-f(x_{0},x_{0})=|h|$, and $|h|/h=\pm1$ has no limit;
        the same holds in $y$. The partials exist exactly off the line
        $y=x$. \textbf{Differentiability:} differentiable off $y=x$, not on
        it, although $f$ is continuous everywhere. \textbf{Graph:} a
        V-shaped valley whose crease, at height $0$, runs along the line
        $y=x$; the contours are pairs of parallel lines $y=x\pm c$, the
        tilted-trough picture of Chapter~9 with a sharp floor.

  \item $f(P)=\ln(0+1)=0$. Chain rule:
        $f_{x}=\dfrac{2}{2x+y^{2}}$, $f_{y}=\dfrac{2y}{2x+y^{2}}$, so at
        $P$ (denominator $1$): $f_{x}=2$, $f_{y}=2$. Tangent plane:
        \[
          z=0+2(x-0)+2(y-1)=2x+2y-2 .
        \]
        At $Q$: $f(Q)=\ln(0.1+0.9604)=\ln1.0604\approx0.058646$, and
        $L(Q)=0.1+1.96-2=0.060000$. The plane overestimates by about
        $0.001354$.

  \item $f(P)=\mathrm{e}^{0}=1$. $f_{x}=\mathrm{e}^{x-2y}$,
        $f_{y}=-2\mathrm{e}^{x-2y}$, so at $P$: $f_{x}=1$, $f_{y}=-2$.
        Tangent plane:
        \[
          z=1+(x-2)-2(y-1)=x-2y+1 .
        \]
        At $Q$: $f(Q)=\mathrm{e}^{2.1-2.04}=\mathrm{e}^{0.06}\approx1.061837$,
        and $L(Q)=2.1-2.04+1=1.060000$. The plane underestimates by about
        $0.001837$: the exponential is convex, so it lies above all its
        tangent planes.

  \item $f(P)=1\cdot4\cdot\mathrm{e}^{0}=4$. Product rule in $x$, where
        both $x$ and $\mathrm{e}^{x-1}$ depend on $x$:
        \[
          f_{x}=y^{2}\mathrm{e}^{x-1}+xy^{2}\mathrm{e}^{x-1}
               =y^{2}\mathrm{e}^{x-1}(1+x),\qquad
          f_{y}=2xy\,\mathrm{e}^{x-1} .
        \]
        At $P$: $f_{x}=4\cdot1\cdot2=8$, $f_{y}=2\cdot1\cdot2\cdot1=4$.
        Tangent plane:
        \[
          z=4+8(x-1)+4(y-2) .
        \]
        At $Q$: $f(Q)=0.98\cdot2.03^{2}\cdot\mathrm{e}^{-0.02}
        =0.98\cdot4.1209\cdot0.980199\approx3.958515$, and
        $L(Q)=4-0.16+0.12=3.960000$. The plane is off by about
        $0.001485$.

  \item $f(P)=\dfrac{4}{2}=2$. Quotient rule with denominator $(x-y)^{2}$:
        \[
          f_{x}=\frac{(x-y)-(x+y)}{(x-y)^{2}}=\frac{-2y}{(x-y)^{2}},\qquad
          f_{y}=\frac{(x-y)+(x+y)}{(x-y)^{2}}=\frac{2x}{(x-y)^{2}} .
        \]
        At $P$: $f_{x}=-\tfrac24=-\tfrac12$, $f_{y}=\tfrac64=\tfrac32$.
        Tangent plane:
        \[
          z=2-\tfrac12(x-3)+\tfrac32(y-1) .
        \]
        At $Q$: $f(Q)=\dfrac{4.0}{2.2}\approx1.818182$, and
        $L(Q)=2-0.05-0.15=1.800000$. The error, about $0.018182$, is
        several times larger than in the neighbouring problems. The step
        keeps $x+y=4$ but moves the denominator $x-y$ from $2$ to $2.2$,
        and along that direction $f=4/(x-y)$ is strongly curved. The plane
        is a local approximation, and its quality depends on the curvature
        as much as on the length of the step.

  \item $F(P)=16^{1/4}\cdot1=2$.
        $F_{K}=\tfrac14K^{-3/4}L^{3/4}$, $F_{L}=\tfrac34K^{1/4}L^{-1/4}$; at
        $P$: $F_{K}=\tfrac14\cdot\tfrac18=\tfrac1{32}$ and
        $F_{L}=\tfrac34\cdot2=\tfrac32$. Tangent plane:
        \[
          z=2+\tfrac1{32}(K-16)+\tfrac32(L-1) .
        \]
        Economically: output $2$ plus marginal product times input change,
        for each input. At $Q$: $L(Q)=2+\tfrac{0.5}{32}+\tfrac32\cdot0.1
        =2+0.015625+0.15=2.165625$, while
        $F(Q)=16.5^{1/4}\cdot1.1^{3/4}\approx2.164789$. The estimate is
        high by about $0.000836$: with exponents summing to $1$ this
        Cobb--Douglas function is concave, so its graph lies below every
        tangent plane.

  \item $f(P)=\ln(2-2+1)=\ln1=0$. With $s=x+2y+3z$:
        $f_{x}=\tfrac1s$, $f_{y}=\tfrac2s$, $f_{z}=\tfrac3s$, and $s=1$ at
        $P$, so $\grad f(P)=(1,2,3)^{T}$. Tangent hyperplane:
        \[
          w=0+(x-2)+2(y+1)+3\bigl(z-\tfrac13\bigr)=x+2y+3z-1 .
        \]
        At $Q$: $s=2.02-1.98+1.05=1.09$, so $f(Q)=\ln1.09\approx0.086178$,
        and $L(Q)=1.09-1=0.090000$. The hyperplane overestimates by about
        $0.003822$, since $\ln$ is concave.

  \item $f(P)=2\mathrm{e}^{0}=2$. $f_{x}=\mathrm{e}^{yz}$,
        $f_{y}=xz\,\mathrm{e}^{yz}$, $f_{z}=xy\,\mathrm{e}^{yz}$; at $P$:
        $f_{x}=1$, $f_{y}=2\cdot5=10$, $f_{z}=2\cdot0=0$. Tangent
        hyperplane:
        \[
          w=2+(x-2)+10(y-0)+0\cdot(z-5)=x+10y .
        \]
        The change in $z$ does not enter at all, because $f_{z}(P)=0$. At
        $Q$: $f(Q)=2.05\,\mathrm{e}^{0.049}\approx2.152952$, and
        $L(Q)=2.05+0.1=2.150000$; error about $0.002952$.

  \item $f(P)=\tfrac66=1$. $f_{x}=\tfrac yz$, $f_{y}=\tfrac xz$,
        $f_{z}=-\tfrac{xy}{z^{2}}$; at $P$: $f_{x}=\tfrac12$,
        $f_{y}=\tfrac13$, $f_{z}=-\tfrac6{36}=-\tfrac16$. Tangent
        hyperplane:
        \[
          w=1+\tfrac12(x-2)+\tfrac13(y-3)-\tfrac16(z-6) .
        \]
        At $Q$: $L(Q)=1+0.05-0.033333-0.033333=0.983333$, and
        $f(Q)=\dfrac{2.1\cdot2.9}{6.2}=\dfrac{6.09}{6.2}\approx0.982258$;
        error about $0.001075$.

  \item Only the term $i=k$ contains $x_{k}$, so
        \[
          \pd{f}{x_{k}}=\frac{2x_{k}}{1+x_{k}^{2}},\qquad k=1,\ldots,n .
        \]
        At $(1,\ldots,1)$ every partial is $\tfrac22=1$:
        $\grad f=(1,1,\ldots,1)^{T}$ and $Df=(1\ \ 1\ \ \cdots\ \ 1)$.

  \item Chain rule with inner function $s=\sum_{i}x_{i}$, whose partial in
        $x_{k}$ is $1$:
        \[
          \pd{f}{x_{k}}=3\Bigl(\sum_{i=1}^{n}x_{i}\Bigr)^{2}\cdot1 ,
        \]
        the same for every $k$. For $n=4$ at $(1,2,3,4)$: $s=10$, so each
        partial is $300$, $\grad f=(300,300,300,300)^{T}$ and
        $Df=(300\ \ 300\ \ 300\ \ 300)$.

  \item The term containing $x_{k}$ is $x_{k}^{k+1}$, so
        $\pd{f}{x_{k}}=(k+1)\,x_{k}^{k}$. At $(1,\ldots,1)$:
        $\grad f=(2,3,\ldots,n+1)^{T}$, $Df=(2\ \ 3\ \ \cdots\ \ n+1)$. For
        $n=3$, $f=x_{1}^{2}+x_{2}^{3}+x_{3}^{4}$ with partials
        $2x_{1},\,3x_{2}^{2},\,4x_{3}^{3}$; at $(2,1,-1)$:
        $\grad f=(4,3,-4)^{T}$ and $Df=(4\ \ 3\ \ {-4})$.

  \item Write $f=\sum_{j}j^{2}x_{j}^{-3}$. The term with $j=k$ gives
        \[
          \pd{f}{x_{k}}=k^{2}\cdot(-3)x_{k}^{-4}=-\frac{3k^{2}}{x_{k}^{4}} .
        \]
        At $(1,\ldots,1)$: $\grad f=-3\,(1,4,9,\ldots,n^{2})^{T}$ and
        $Df=(-3\ \ {-12}\ \ {-27}\ \ \cdots\ \ {-3n^{2}})$.

  \item The first sum contributes $2x_{k}$; the second term, by the chain
        rule as in problem 39, contributes
        $-\tfrac1n\cdot2\sum_{i}x_{i}$. With the mean
        $\bar x=\tfrac1n\sum_{i}x_{i}$:
        \[
          \pd{f}{x_{k}}=2x_{k}-\frac2n\sum_{i=1}^{n}x_{i}=2\,(x_{k}-\bar x) .
        \]
        The gradient is zero exactly when $x_{k}=\bar x$ for every $k$,
        i.e.\ when all $x_{i}$ are equal. For $n=3$ at $(1,2,6)$:
        $\bar x=3$, so $\grad f=(-4,-2,6)^{T}$ and $Df=(-4\ \ {-2}\ \ 6)$.
        (The components always add up to $0$.)

  \item With $s=x^{2}+3yz$: $f_{x}=y^{2}\cdot\dfrac{2x}{s}$, since
        $y^{2}$ is a constant factor; $f_{y}$ needs the product rule,
        because $y$ appears in both factors:
        \[
          f_{y}=2y\ln s+y^{2}\cdot\frac{3z}{s};
        \]
        and $f_{z}=y^{2}\cdot\dfrac{3y}{s}=\dfrac{3y^{3}}{s}$. At
        $P=(1,1,1)$, $s=4$:
        \[
          f_{x}=\tfrac24=\tfrac12,\qquad
          f_{y}=2\ln4+\tfrac34\approx3.522589,\qquad
          f_{z}=\tfrac34 .
        \]
        So $\grad f(P)=\bigl(\tfrac12,\ 2\ln4+\tfrac34,\ \tfrac34\bigr)^{T}$
        and $Df(P)=\bigl(\tfrac12\ \ \ 2\ln4+\tfrac34\ \ \ \tfrac34\bigr)$.

  \item Write $E=\mathrm{e}^{y^{2}-2xz}$. In $x$ both factors move, so
        product and chain rule:
        \[
          f_{x}=3x^{2}E+x^{3}E\cdot(-2z)=x^{2}(3-2xz)\,E .
        \]
        In $y$ and $z$ only the exponential moves:
        $f_{y}=x^{3}E\cdot2y=2x^{3}y\,E$ and
        $f_{z}=x^{3}E\cdot(-2x)=-2x^{4}E$. At $P=(1,0,1)$,
        $E=\mathrm{e}^{-2}$:
        $f_{x}=1\cdot(3-2)\mathrm{e}^{-2}=\mathrm{e}^{-2}$, $f_{y}=0$,
        $f_{z}=-2\mathrm{e}^{-2}$. So
        $\grad f(P)=\mathrm{e}^{-2}(1,0,-2)^{T}\approx(0.135335,\,0,\,-0.270671)^{T}$
        and $Df(P)=\mathrm{e}^{-2}(1\ \ 0\ \ {-2})$.

  \item Use $\partial a^{g}/\partial x=a^{g}\ln a\cdot g_{x}$. In the first
        term $3^{y}$ is a constant factor; in the second the exponent $xy$
        has $x$-derivative $y$:
        \[
          f_{x}=2^{x}3^{y}\ln2+y\,5^{xy}\ln5,\qquad
          f_{y}=2^{x}3^{y}\ln3+x\,5^{xy}\ln5 .
        \]
        At $(1,1)$, $2^{x}3^{y}=6$ and $5^{xy}=5$:
        \[
          \grad f(1,1)=\begin{pmatrix}6\ln2+5\ln5\\6\ln3+5\ln5\end{pmatrix}
          \approx\begin{pmatrix}12.206073\\14.638863\end{pmatrix},
        \]
        and $Df(1,1)$ is the same pair written as a row.

  \item Only the term $i=k$ contains $x_{k}$, and its exponent $kx_{k}$ has
        derivative $k$:
        \[
          \pd{f}{x_{k}}=3^{\,kx_{k}}\ln3\cdot k=k\ln3\cdot3^{\,kx_{k}} .
        \]
        At the origin every $3^{0}=1$, so
        $\grad f=\ln3\,(1,2,\ldots,n)^{T}$ and
        $Df=\ln3\,(1\ \ 2\ \ \cdots\ \ n)$.

  \item Candidate: $f_{x}=3y$, $f_{y}=3x$, so $\grad f(a,b)=(3b,3a)$.
        Remainder:
        \[
          r(\mathbf{h})=3(a+h_{1})(b+h_{2})-3ab-3bh_{1}-3ah_{2}=3h_{1}h_{2} .
        \]
        Estimate:
        $\dfrac{|r(\mathbf{h})|}{\|\mathbf{h}\|}\le\dfrac{3\cdot\frac12
        \|\mathbf{h}\|^{2}}{\|\mathbf{h}\|}=\tfrac32\|\mathbf{h}\|\to0$.
        So $f$ is differentiable at every point, with derivative
        $(3b,\,3a)$.
  \item At the origin $\grad f=(0,0)$ and
        $r(\mathbf{h})=h_{1}^{2}+h_{2}^{2}=\|\mathbf{h}\|^{2}$, so
        $|r|/\|\mathbf{h}\|=\|\mathbf{h}\|\to0$: differentiable, with tangent
        plane $z=0$. At $(a,b)$: $\grad f=(2a,2b)$ and
        \[
          r(\mathbf{h})=(a+h_{1})^{2}+(b+h_{2})^{2}-a^{2}-b^{2}-2ah_{1}-2bh_{2}
          =h_{1}^{2}+h_{2}^{2},
        \]
        the same remainder, so the same estimate works everywhere.
  \item Candidate: $f_{x}=y^{2}$, $f_{y}=2xy$, so $\grad f(a,b)=(b^{2},2ab)$.
        Expand $(a+h_{1})(b^{2}+2bh_{2}+h_{2}^{2})
        =ab^{2}+2abh_{2}+ah_{2}^{2}+b^{2}h_{1}+2bh_{1}h_{2}+h_{1}h_{2}^{2}$
        and subtract $ab^{2}+b^{2}h_{1}+2abh_{2}$:
        \[
          r(\mathbf{h})=a\,h_{2}^{2}+2b\,h_{1}h_{2}+h_{1}h_{2}^{2} .
        \]
        Estimate:
        $\dfrac{|r|}{\|\mathbf{h}\|}\le|a|\,\|\mathbf{h}\|
        +|b|\,\|\mathbf{h}\|+\|\mathbf{h}\|^{2}\to0$. Differentiable
        everywhere, derivative $(b^{2},\,2ab)$.
  \item Candidate: $f_{x}=3y\mathrm{e}^{3x}$, $f_{y}=\mathrm{e}^{3x}$, so
        $\grad f(a,b)=\mathrm{e}^{3a}(3b,\,1)$. Write
        $\mathrm{e}^{3(a+h_{1})}=\mathrm{e}^{3a}\bigl(1+3h_{1}+R(3h_{1})\bigr)$:
        \[
          f(a+h_{1},b+h_{2})=\mathrm{e}^{3a}(b+h_{2})\bigl(1+3h_{1}+R\bigr)
          =\mathrm{e}^{3a}\bigl(b+3bh_{1}+h_{2}+3h_{1}h_{2}+(b+h_{2})R\bigr).
        \]
        Subtracting $b\mathrm{e}^{3a}$ and $\mathrm{e}^{3a}(3bh_{1}+h_{2})$
        leaves
        \[
          r(\mathbf{h})=\mathrm{e}^{3a}\bigl(3h_{1}h_{2}+(b+h_{2})\,R(3h_{1})\bigr).
        \]
        For $\|\mathbf{h}\|\le\tfrac13$ we have $|3h_{1}|\le1$, so
        $|R(3h_{1})|\le2\cdot9h_{1}^{2}\le18\|\mathbf{h}\|^{2}$, and
        \[
          \frac{|r|}{\|\mathbf{h}\|}\le\mathrm{e}^{3a}\Bigl(\tfrac32
          \|\mathbf{h}\|+18\bigl(|b|+\|\mathbf{h}\|\bigr)\|\mathbf{h}\|\Bigr)
          \to0 .
        \]
        Differentiable at every point.
  \item (a) On the $x$-axis $f(h,0)=(h^{2})^{1/3}=|h|^{2/3}$, so the
        difference quotient $|h|^{2/3}/h=\pm|h|^{-1/3}$ is unbounded:
        $f_{x}(0,0)$ does not exist (likewise $f_{y}$, with quotient
        $\pm4^{1/3}|h|^{-1/3}$). Not differentiable.
        (b) $g(\mathbf{h})=\|\mathbf{h}\|^{3/2}$. On an axis the quotient is
        $|h|^{3/2}/h=\pm|h|^{1/2}\to0$, so both partials are $0$. Then
        $r(\mathbf{h})=\|\mathbf{h}\|^{3/2}$ and
        $|r|/\|\mathbf{h}\|=\|\mathbf{h}\|^{1/2}\to0$: differentiable, with
        tangent plane $z=0$. The exponents decide: $2p=\tfrac23<1$ fails,
        $2p=\tfrac32>1$ works.
  \item \emph{Origin.} $f=0$ on both axes, so both partials are $0$. The
        remainder is $r(\mathbf{h})=|h_{1}h_{2}|\le\tfrac12\|\mathbf{h}\|^{2}$,
        so $|r|/\|\mathbf{h}\|\le\tfrac12\|\mathbf{h}\|\to0$: differentiable,
        with tangent plane $z=0$. \emph{At $(2,0)$.}
        $f(2,h)-f(2,0)=|2h|=2|h|$, and the quotient $2|h|/h=\pm2$ has no
        limit: $f_{y}(2,0)$ does not exist, so $f$ is not differentiable
        there. The graph is creased along both axes, except at the origin,
        where the two creases flatten out.
  \item \emph{Partials.} $f(h,0)=0$ and $f(0,h)=0$ for $h\neq0$, so both
        difference quotients are $0/h=0$: $f_{x}(0,0)=f_{y}(0,0)=0$.
        \emph{Lines.} On the $y$-axis $f=0$. On $y=mx$, $x\neq0$:
        \[
          f(x,mx)=\frac{mx^{3}}{x^{4}+m^{2}x^{2}}=\frac{mx}{x^{2}+m^{2}}
          \;\longrightarrow\;0\quad(x\to0),
        \]
        for $m\neq0$, and $f=0$ for $m=0$. Every straight line gives $0$.
        \emph{Parabola.} On $y=x^{2}$: $f=\dfrac{x^{4}}{x^{4}+x^{4}}=\tfrac12$
        for every $x\neq0$. So $f$ has no limit at the origin and is
        \textbf{not continuous} there, hence \textbf{not differentiable},
        although both partials exist and every line agrees.
  \item \emph{Partials.} $f=0$ on both axes, so $f_{x}(0,0)=f_{y}(0,0)=0$.
        \emph{Lines.} On $y=mx$:
        $f=\dfrac{2m^{2}x^{3}}{x^{2}+m^{4}x^{4}}=\dfrac{2m^{2}x}{1+m^{4}x^{2}}
        \to0$; on $x=0$, $f=0$. \emph{Parabola.} On $x=y^{2}$:
        $f=\dfrac{2y^{4}}{y^{4}+y^{4}}=1$ for $y\neq0$. Not continuous at the
        origin, hence not differentiable.
  \item (a) A smooth bowl, with no edges, creases or tips anywhere: it has a
        tangent plane at every point (horizontal at the bottom), so the
        function is differentiable everywhere. (b) A cone with a sharp tip:
        differentiable everywhere except at the tip, where every vertical
        cut through the tip is a V and no plane fits. (c) A pyramid whose
        four flat faces meet along two creases running over the axes:
        differentiable on the faces, not on the two creases (including
        the top point).
  \item The bowl $(x-1)^{2}+4y^{2}$ has no tip, so it is out. The level
        $c>0$ of $\sqrt{(x-1)^{2}+4y^{2}}$ is
        $\dfrac{(x-1)^{2}}{c^{2}}+\dfrac{y^{2}}{(c/2)^{2}}=1$, semi-axes $c$
        in $x$ and $c/2$ in $y$: twice as long in $x$, as in the picture.
        The third formula has semi-axes $c/2$ in $x$ and $c$ in $y$, the
        wrong way round. So (b) is $\sqrt{(x-1)^{2}+4y^{2}}$. At $(1,0)$:
        $f(1+h,0)-f(1,0)=\sqrt{h^{2}}=|h|$, and $|h|/h=\pm1$ has no limit;
        likewise $f(1,h)-f(1,0)=2|h|$. Neither partial exists, so $f$ is
        not differentiable at the tip, as the picture said.
\end{solutions}

\chapter{Higher-Order Partials, Young's Theorem, the Hessian, and Definiteness}

First-order partials describe slopes; second-order partials describe
\emph{curvature}. In one variable we used $f''$ to distinguish maxima from
minima. In several variables we need a whole \emph{matrix} of second
derivatives (the Hessian) and curvature is captured by its
definiteness, which we test with the leading principal minors (Sylvester's
criterion).

\section{Core Rules}

\begin{rulebox}{Direct second-order partials}
For $f(x,y)$,
\[
  f_{xx} \;=\; \pd{}{x}\!\left(\pd{f}{x}\right) \;=\; \frac{\partial^{2} f}{\partial x^{2}},
  \qquad
  f_{yy} \;=\; \frac{\partial^{2} f}{\partial y^{2}} .
\]
\end{rulebox}

\begin{rulebox}{Cross (mixed) partials}
\[
  f_{xy} \;=\; \pd{}{y}\!\left(\pd{f}{x}\right) \;=\; \frac{\partial^{2} f}{\partial y\,\partial x},
  \qquad
  f_{yx} \;=\; \pd{}{x}\!\left(\pd{f}{y}\right) .
\]
\end{rulebox}

\begin{rulebox}{Young's Theorem (Schwarz / Clairaut)}
If $f$ is $C^{2}$ on an open set (i.e.\ all second partials exist and
are continuous) then the order of differentiation is irrelevant:
\[
  f_{xy} \;=\; f_{yx} .
\]
This is the symmetry of the Hessian. All standard polynomial, exponential,
logarithmic, Cobb--Douglas, and CES forms are $C^{2}$ on their natural
domains.
\end{rulebox}

\begin{rulebox}{Hessian matrix}
For $f:\R^{n}\to\R$,
\[
  \Hess_f(\mathbf{x}) \;=\; \bigl[f_{x_ix_j}(\mathbf{x})\bigr]_{i,j=1}^{n}
  \;=\;
  \begin{pmatrix}
    f_{x_1x_1} & f_{x_1x_2} & \cdots & f_{x_1x_n} \\
    f_{x_2x_1} & f_{x_2x_2} & \cdots & f_{x_2x_n} \\
    \vdots     & \vdots     & \ddots & \vdots     \\
    f_{x_nx_1} & f_{x_nx_2} & \cdots & f_{x_nx_n}
  \end{pmatrix} .
\]
For $C^{2}$ functions $\Hess_f$ is symmetric.
\end{rulebox}

\begin{rulebox}{Definiteness}
A symmetric matrix $H$ is
\begin{itemize}[leftmargin=1.6em,itemsep=0.3ex]
  \item \emph{positive definite (PD)} if $\mathbf{v}^{T}H\mathbf{v}>0$ for
        every $\mathbf{v}\neq\mathbf{0}$ (curvature ``up'' in every
        direction $\Rightarrow$ local minimum;
  \item \emph{negative definite (ND)} if $\mathbf{v}^{T}H\mathbf{v}<0$ for
        every $\mathbf{v}\neq\mathbf{0}$) curvature ``down'' in every
        direction $\Rightarrow$ local maximum;
  \item \emph{indefinite} if $\mathbf{v}^{T}H\mathbf{v}$ takes both
        positive and negative values $\Rightarrow$ saddle point.
\end{itemize}
\end{rulebox}

\begin{rulebox}{Sylvester's criterion ($2\times2$ case)}
Let $H=\begin{pmatrix}a&b\\b&c\end{pmatrix}$ with leading principal minors
$D_1=a$ and $D_2=\det H=ac-b^{2}$. Then
\[
  \begin{aligned}
    D_1>0 \text{ and } D_2>0 &\;\Rightarrow\; H \text{ is PD}, \\
    D_1<0 \text{ and } D_2>0 &\;\Rightarrow\; H \text{ is ND}, \\
    D_2<0 &\;\Rightarrow\; H \text{ is indefinite (saddle)}, \\
    D_2=0 &\;\Rightarrow\; \text{test is inconclusive (semi-definite case)}.
  \end{aligned}
\]
Note: ND requires $D_1<0$ \emph{and} $D_2>0$; the alternating-sign
condition is what one generalises to higher $n$ (signs $-,+,-,+,\dots$ for
ND).
\end{rulebox}

\begin{rulebox}{Classification of critical points}
A point $\mathbf{x}_0$ is a \emph{critical point} of $f$ if
$\grad f(\mathbf{x}_0)=\mathbf{0}$. Then evaluate $\Hess_f(\mathbf{x}_0)$:
\begin{itemize}[leftmargin=1.6em,itemsep=0.3ex]
  \item $H$ PD at $\mathbf{x}_0$ $\Rightarrow$ strict local minimum;
  \item $H$ ND at $\mathbf{x}_0$ $\Rightarrow$ strict local maximum;
  \item $H$ indefinite at $\mathbf{x}_0$ $\Rightarrow$ saddle point;
  \item $D_2=0$ $\Rightarrow$ inconclusive, higher-order analysis
        required.
\end{itemize}
\end{rulebox}

\section{Worked Example}

\begin{worked}{Find and classify every critical point of $f(x,y)=x^{3}-3xy+y^{3}$}
\wstep{Step 1. First-order conditions: both partials zero, simultaneously.}
\[
  f_{x}=3x^{2}-3y=0
  \quad\Longrightarrow\quad y=x^{2},
\]
\[
  f_{y}=-3x+3y^{2}=0
  \quad\Longrightarrow\quad x=y^{2} .
\]
\wtrap{A critical point needs \emph{every} partial to vanish at the
\emph{same} point. Solving each equation separately and pairing the answers
arbitrarily is not the same thing --- the two equations form a system.}

\wstep{Step 2. Solve the system by substitution.}
Substituting $y=x^{2}$ into $x=y^{2}$:
\[
  x=x^{4}
  \;\Longrightarrow\;
  x^{4}-x=0
  \;\Longrightarrow\;
  x\bigl(x^{3}-1\bigr)=0
  \;\Longrightarrow\;
  x=0\ \text{or}\ x=1 .
\]
The corresponding $y=x^{2}$ gives the two critical points
\[
  (0,0)\qquad\text{and}\qquad(1,1).
\]
\wtrap{Do not divide by $x$ to get $x^{3}=1$. That silently discards the
root $x=0$ --- and $(0,0)$ turns out to be the more interesting of the two
points. Factor; never divide by something that may vanish.}

\wstep{Step 3. Build the Hessian once, symbolically.}
\[
  f_{xx}=6x,\qquad f_{yy}=6y,\qquad f_{xy}=f_{yx}=-3 ,
\]
\[
  \Hess f(x,y)=\begin{pmatrix}6x&-3\\-3&6y\end{pmatrix}.
\]
The two mixed partials agree, as Young's theorem guarantees for a function
with continuous second derivatives --- differentiating in $x$ then $y$ gives
the same result as $y$ then $x$. Computing both is a free check, and here
$\frac{\partial}{\partial y}(3x^{2}-3y)=-3$ and
$\frac{\partial}{\partial x}(-3x+3y^{2})=-3$. \checkmark

\wstep{Step 4. Classify $(0,0)$.}
\[
  \Hess f(0,0)=\begin{pmatrix}0&-3\\-3&0\end{pmatrix},
  \qquad
  \det=0-9=-9<0 .
\]
A negative determinant means the two eigenvalues have opposite signs, so the
form is indefinite:
\[
  \boxed{(0,0)\ \text{is a saddle point.}}
\]
No further test is needed, and the value of $f_{xx}$ is irrelevant here.
\wtrap{$\det\Hess<0$ settles the matter immediately --- saddle, whatever
$f_{xx}$ does. It is only when $\det\Hess>0$ that $f_{xx}$ has to be
consulted. Checking $f_{xx}$ first wastes effort and invites the wrong
conclusion.}

\wstep{Step 5. Classify $(1,1)$.}
\[
  \Hess f(1,1)=\begin{pmatrix}6&-3\\-3&6\end{pmatrix},
  \qquad
  \det=36-9=27>0,
  \qquad
  f_{xx}=6>0 .
\]
Positive determinant with a positive $f_{xx}$: positive definite, so
\[
  \boxed{(1,1)\ \text{is a strict local minimum},\qquad f(1,1)=1-3+1=-1 .}
\]

\wstep{Step 6. Verify the saddle directly, because it is cheap.}
Along the line $y=x$: $f(x,x)=2x^{3}-3x^{2}$, whose value near $0$ is
$\approx-3x^{2}<0$. Along $y=-x$: $f(x,-x)=x^{3}+3x^{2}-x^{3}=3x^{2}>0$. So
$f$ takes both signs arbitrarily close to $(0,0)$, where $f=0$ --- a saddle,
confirmed without any matrix.

\wstep{Step 7. The decision rule, complete.}
\begin{center}\footnotesize
\begin{tabular}{lll}
  \hline
  $\det\Hess$ & $f_{xx}$ & conclusion\\
  \hline
  $>0$ & $>0$ & strict local minimum\\
  $>0$ & $<0$ & strict local maximum\\
  $<0$ & --- & saddle point\\
  $=0$ & --- & \textbf{test fails}; investigate directly\\
  \hline
\end{tabular}
\end{center}
\wtrap{$\det\Hess=0$ means the test is silent, not that the point is a
saddle. $f=x^{4}+y^{4}$ (minimum), $f=-x^{4}-y^{4}$ (maximum) and
$f=x^{4}-y^{4}$ (saddle) all have the zero Hessian at the origin. When the
determinant vanishes you must argue directly, as in Step~6.}

\wstep{Step 8. Local is not global.}
Neither point is a global extremum: $f(x,0)=x^{3}$ runs to $\pm\infty$. On
an unbounded domain, second-order conditions can never certify a global
optimum by themselves --- that needs convexity of $f$ on the whole domain,
or a compact feasible set and Weierstrass. Volume~7 makes this the
difference between finding a candidate and having an answer.
\end{worked}

\section{Practice Problems}

\begin{problems}
  \item Compute all second partials of $f(x,y) = 4x^{3} - 5x^{2}y + 7y^{4}$
        and verify Young's theorem.
  \item Construct the Hessian of
        $f(x,y) = 5x^{2} + 6xy + 2y^{2} - 3x + 4y$.
  \item Classify the definiteness of
        $\Hess = \begin{pmatrix}6&2\\2&3\end{pmatrix}$ using Sylvester's
        criterion.
  \item Find and classify the critical points of
        $f(x,y) = x^{2} + 4y^{2} - 6x + 8y + 11$.
  \item Compute all second partials of $f(x,y) = (x+2y)^{3}$ and verify
        $f_{xy} = f_{yx}$.
  \item Construct the Hessian of $f(x,y) = x^{3} + 2xy^{2} - 4y$.
  \item Classify $\Hess = \begin{pmatrix}-4&2\\2&-5\end{pmatrix}$.
  \item Find and classify the critical points of
        $f(x,y) = -x^{2} - 3y^{2} + 8x - 6y + 5$.
  \item Compute all second partials of $f(x,y) = x^{2}\ln(y)$ ($y>0$) and
        verify Young's theorem.
  \item Construct the Hessian of $f(x,y) = \mathrm{e}^{2x-y}$.
  \item Classify $\Hess = \begin{pmatrix}2&4\\4&2\end{pmatrix}$.
  \item Find and classify the critical points of
        $f(x,y) = x^{2} - y^{2} + 6x - 4y$.
  \item Compute all second partials of $f(x,y) = \mathrm{e}^{x^{2}+y}$ and
        verify Young's theorem.
  \item Construct the Hessian of the Cobb--Douglas function
        $f(K,L) = K^{1/3}L^{2/3}$, $K,L>0$.
  \item Classify $\Hess = \begin{pmatrix}7&-3\\-3&1\end{pmatrix}$.
  \item Find and classify the critical points of
        $f(x,y) = x^{3} + y^{3} - 12x - 27y$.
  \item Compute all second partials of $f(x,y) = 1/(x+y)$ and verify
        Young's theorem.
  \item Construct the Hessian of
        $f(x,y,z) = x^{2} + 2y^{2} + 3z^{2} - xy + 2yz$.
  \item Classify $\Hess = \begin{pmatrix}-3&1\\1&-3\end{pmatrix}$.
  \item Find and classify the critical points of
        $f(x,y) = x^{2} + 4xy + y^{2} + 2y$.
  \item Compute all second partials of $f(x,y) = x^{1/2}y^{1/2}$
        ($x,y>0$) and verify $f_{xy} = f_{yx}$.
  \item Construct the Hessian of $f(x,y) = \ln(x^{2}+y^{2})$ at a generic
        point $(x,y)\neq(0,0)$.
  \item Find and classify the critical points of
        $f(x,y) = x^{4} + y^{4} - 4xy$.
  \item Compute all second partials of $f(x,y) = x\mathrm{e}^{y} + y\mathrm{e}^{x}$
        and verify Young's theorem.
  \item Construct the Hessian of $f(x,y) = x^{3} - 3xy^{2}$ (a harmonic
        function).
  \item Classify $\Hess = \begin{pmatrix}5&0\\0&-2\end{pmatrix}$ and state
        what this implies if $H$ is the Hessian at a critical point.
  \item Find and classify the critical points of
        $f(x,y) = xy - x^{2} - y^{2} + 6x + 3y$.
  \item Compute all second partials of $f(x,y) = \ln(xy)$ ($x,y>0$).
  \item Construct the Hessian of $f(x,y) = (x^{2}+y^{2})^{2}$ at a generic
        point.
  \item Classify $\Hess = \begin{pmatrix}-1&2\\2&-1\end{pmatrix}$.
  \item Find and classify the critical points of
        $f(x,y) = x^{3} - 3xy + y^{3}$.
  \item Compute all second partials of $f(x,y) = \sin(x)\cos(y)$ and
        verify Young's theorem.
  \item Construct the Hessian of $f(x,y,z) = \mathrm{e}^{xyz}$ at the point
        $(1,1,1)$ (evaluate numerically).
  \item Classify $\Hess = \begin{pmatrix}4&-2\\-2&4\end{pmatrix}$.
  \item Find and classify the critical points of the profit function
        $\pi(L,K) = 10L^{1/2}K^{1/2} - 2L - 3K$ (assume $L,K>0$).
  \item Compute all second partials of $f(x,y) = x^{2}y^{3}$ and verify
        Young's theorem.
  \item Construct the Hessian of
        $f(x,y) = \mathrm{e}^{x}\sin(y) - \mathrm{e}^{x}\cos(y)$.
  \item Classify $\Hess = \begin{pmatrix}9&6\\6&4\end{pmatrix}$. (Note the
        boundary case.)
  \item Find and classify the critical points of
        $f(x,y) = x^{2} + xy + y^{2} - 3x$.
  \item Compute all second partials of $f(x,y) = \arctan(y/x)$ ($x>0$) and
        verify Young's theorem.
  \item Construct the Hessian of the CES function
        $f(K,L) = (K^{1/2}+L^{1/2})^{2}$ at a generic point.
  \item Find and classify the critical points of
        $f(x,y) = -x^{4} - y^{4} + 8xy$.
  \item Compute all second partials of $f(x,y) = x^{2}\mathrm{e}^{-y}$ and
        verify Young's theorem.
  \item Construct the Hessian of $f(x,y) = \cos(xy)$ at the origin
        $(0,0)$.
  \item Classify $\Hess = \begin{pmatrix}0&3\\3&0\end{pmatrix}$.
  \item Find and classify the critical points of
        $f(x,y) = x^{3} + 3xy^{2} - 15x - 12y$.
  \item Construct the Hessian of $f(x,y) = \ln(1+x^{2}+y^{2})$ at the
        origin.
  \item Find and classify the critical points of
        $f(x,y) = \mathrm{e}^{-(x^{2}+y^{2})}$.
\end{problems}

\section{Solutions}

\begin{solutions}
  \item $f_x = 12x^{2} - 10xy$, $f_y = -5x^{2} + 28y^{3}$. Then
        $f_{xx} = 24x - 10y$, $f_{yy} = 84y^{2}$, and
        \[
          f_{xy} = \pd{}{y}\bigl(12x^{2}-10xy\bigr) = -10x, \qquad
          f_{yx} = \pd{}{x}\bigl(-5x^{2}+28y^{3}\bigr) = -10x . \; \checkmark
        \]

  \item $f_x = 10x + 6y - 3$, $f_y = 6x + 4y + 4$. Hessian
        \[
          \Hess = \begin{pmatrix}10&6\\6&4\end{pmatrix} .
        \]

  \item $D_1 = 6 > 0$, $D_2 = 18 - 4 = 14 > 0$. \textbf{PD.}

  \item $\grad f = (2x-6,\,8y+8) = (0,0) \Rightarrow x=3,\,y=-1$. Hessian
        $\Hess = \begin{pmatrix}2&0\\0&8\end{pmatrix}$: $D_1=2>0$,
        $D_2=16>0$, \textbf{PD}, so $(3,-1)$ is a strict \textbf{local
        minimum} with $f(3,-1) = 9+4-18-8+11 = -2$.

  \item Chain rule: $f_x = 3(x+2y)^{2}$, $f_y = 6(x+2y)^{2}$. So
        $f_{xx} = 6(x+2y)$, $f_{yy} = 24(x+2y)$, and
        \[
          f_{xy} = 12(x+2y) = f_{yx} . \; \checkmark
        \]

  \item $f_x = 3x^{2} + 2y^{2}$, $f_y = 4xy - 4$. Hessian
        \[
          \Hess = \begin{pmatrix}6x&4y\\4y&4x\end{pmatrix} .
        \]

  \item $D_1 = -4 < 0$, $D_2 = 20 - 4 = 16 > 0$. \textbf{ND.}

  \item $\grad f = (-2x+8,\,-6y-6) = (0,0) \Rightarrow x=4,\,y=-1$. Hessian
        $\Hess = \begin{pmatrix}-2&0\\0&-6\end{pmatrix}$: $D_1=-2<0$,
        $D_2=12>0$, \textbf{ND}, so $(4,-1)$ is a strict \textbf{local
        maximum} with $f(4,-1) = -16-3+32+6+5 = 24$.

  \item $f_x = 2x\ln(y)$, $f_y = x^{2}/y$. So $f_{xx} = 2\ln(y)$,
        $f_{yy} = -x^{2}/y^{2}$, and
        \[
          f_{xy} = 2x/y = f_{yx} . \; \checkmark
        \]

  \item $f_x = 2\mathrm{e}^{2x-y}$, $f_y = -\mathrm{e}^{2x-y}$. So
        $f_{xx} = 4\mathrm{e}^{2x-y}$, $f_{yy} = \mathrm{e}^{2x-y}$,
        $f_{xy} = -2\mathrm{e}^{2x-y}$. Hessian
        \[
          \Hess = \mathrm{e}^{2x-y}\begin{pmatrix}4&-2\\-2&1\end{pmatrix} .
        \]

  \item $D_1 = 2 > 0$, $D_2 = 4 - 16 = -12 < 0$. \textbf{Indefinite
        (saddle).}

  \item $\grad f = (2x+6,\,-2y-4) = (0,0) \Rightarrow x=-3,\,y=-2$. Hessian
        $\Hess = \begin{pmatrix}2&0\\0&-2\end{pmatrix}$: $D_2 = -4 < 0$,
        \textbf{indefinite}, so $(-3,-2)$ is a \textbf{saddle point}.

  \item Let $E = \mathrm{e}^{x^{2}+y}$. Then $f_x = 2xE$, $f_y = E$. So
        $f_{xx} = 2E + 4x^{2}E = (2+4x^{2})E$, $f_{yy} = E$, and
        \[
          f_{xy} = \pd{}{y}(2xE) = 2xE, \qquad
          f_{yx} = \pd{}{x}(E) = 2xE . \; \checkmark
        \]

  \item $f_K = \tfrac13K^{-2/3}L^{2/3}$, $f_L = \tfrac23K^{1/3}L^{-1/3}$.
        So
        \[
          f_{KK} = -\tfrac29K^{-5/3}L^{2/3}, \quad
          f_{LL} = -\tfrac29K^{1/3}L^{-4/3}, \quad
          f_{KL} = \tfrac29K^{-2/3}L^{-1/3} = f_{LK} . \; \checkmark
        \]
        Hessian
        \[
          \Hess = \tfrac29
          \begin{pmatrix}
            -K^{-5/3}L^{2/3} & K^{-2/3}L^{-1/3} \\
            K^{-2/3}L^{-1/3} & -K^{1/3}L^{-4/3}
          \end{pmatrix} .
        \]

  \item $D_1 = 7 > 0$, $D_2 = 7 - 9 = -2 < 0$. \textbf{Indefinite.}

  \item $\grad f = (3x^{2}-12,\,3y^{2}-27) = (0,0) \Rightarrow x^{2}=4,\,
        y^{2}=9$, giving $(\pm2,\pm3)$. Hessian
        $\Hess = \begin{pmatrix}6x&0\\0&6y\end{pmatrix}$, so $D_1=6x$,
        $D_2=36xy$.
        \begin{itemize}
          \item $(2,3)$: $D_1=12>0$, $D_2=216>0$, \textbf{PD}: local
                minimum, $f = 8+27-24-81 = -70$.
          \item $(2,-3)$: $D_2=-216<0$, \textbf{saddle}.
          \item $(-2,3)$: $D_2=-216<0$, \textbf{saddle}.
          \item $(-2,-3)$: $D_1=-12<0$, $D_2=216>0$, \textbf{ND}: local
                maximum, $f = -8-27+24+81 = 70$.
        \end{itemize}

  \item Let $S=x+y$, $f=S^{-1}$, $f_x = -S^{-2} = f_y$. So
        $f_{xx} = 2S^{-3} = f_{yy}$, and
        \[
          f_{xy} = 2S^{-3} = f_{yx} = \frac{2}{(x+y)^{3}} . \; \checkmark
        \]

  \item $f_x = 2x-y$, $f_y = 4y-x+2z$, $f_z = 6z+2y$. Hessian
        \[
          \Hess = \begin{pmatrix}2&-1&0\\-1&4&2\\0&2&6\end{pmatrix} .
        \]

  \item $D_1 = -3 < 0$, $D_2 = 9 - 1 = 8 > 0$. \textbf{ND.}

  \item $\grad f = (2x+4y,\,4x+2y+2) = (0,0)$. From the first, $x=-2y$;
        substituting, $-8y+2y+2=0 \Rightarrow y=1/3$, $x=-2/3$. Hessian
        $\Hess = \begin{pmatrix}2&4\\4&2\end{pmatrix}$: $D_1=2>0$,
        $D_2=4-16=-12<0$, \textbf{indefinite: saddle point} at
        $(-2/3,1/3)$.

  \item $f_x = \tfrac12x^{-1/2}y^{1/2}$, $f_y = \tfrac12x^{1/2}y^{-1/2}$.
        So
        \[
          f_{xx} = -\tfrac14x^{-3/2}y^{1/2}, \quad
          f_{yy} = -\tfrac14x^{1/2}y^{-3/2}, \quad
          f_{xy} = \tfrac14x^{-1/2}y^{-1/2} = f_{yx} . \; \checkmark
        \]

  \item Let $S=x^{2}+y^{2}$, $f=\ln S$, $f_x = 2x/S$, $f_y = 2y/S$. By the
        quotient rule:
        \[
          f_{xx} = \frac{2S-2x\cdot2x}{S^{2}} = \frac{2(y^{2}-x^{2})}{(x^{2}+y^{2})^{2}}, \quad
          f_{yy} = \frac{2(x^{2}-y^{2})}{(x^{2}+y^{2})^{2}}, \quad
          f_{xy} = \frac{-4xy}{(x^{2}+y^{2})^{2}} .
        \]
        Hessian
        \[
          \Hess = \frac{2}{(x^{2}+y^{2})^{2}}
          \begin{pmatrix}y^{2}-x^{2} & -2xy \\ -2xy & x^{2}-y^{2}\end{pmatrix} .
        \]
        Note $f_{xx}+f_{yy}=0$: $\ln(x^{2}+y^{2})$ is harmonic.

  \item $\grad f = (4x^{3}-4y,\,4y^{3}-4x) = (0,0) \Rightarrow y=x^{3}$
        and $x=y^{3}=x^{9}$. So $x(x^{8}-1)=0$, giving $x\in\{-1,0,1\}$.
        Critical points: $(0,0),(1,1),(-1,-1)$. Hessian
        $\Hess = \begin{pmatrix}12x^{2}&-4\\-4&12y^{2}\end{pmatrix}$.
        \begin{itemize}
          \item $(0,0)$: $\Hess = \begin{pmatrix}0&-4\\-4&0\end{pmatrix}$,
                $D_2=-16<0$, \textbf{saddle}.
          \item $(1,1)$: $\Hess = \begin{pmatrix}12&-4\\-4&12\end{pmatrix}$,
                $D_1=12>0$, $D_2=144-16=128>0$, \textbf{PD}: local
                minimum, $f=1+1-4=-2$.
          \item $(-1,-1)$: same $H$, \textbf{PD}: local minimum, $f=-2$.
        \end{itemize}

  \item $f_x = \mathrm{e}^{y}+y\mathrm{e}^{x}$, $f_y = x\mathrm{e}^{y}+\mathrm{e}^{x}$. So
        $f_{xx} = y\mathrm{e}^{x}$, $f_{yy} = x\mathrm{e}^{y}$, and
        \[
          f_{xy} = \mathrm{e}^{y}+\mathrm{e}^{x} = f_{yx} . \; \checkmark
        \]

  \item $f_x = 3x^{2}-3y^{2}$, $f_y = -6xy$. So $f_{xx}=6x$,
        $f_{yy}=-6x$, $f_{xy}=-6y=f_{yx}$. Hessian
        \[
          \Hess = \begin{pmatrix}6x&-6y\\-6y&-6x\end{pmatrix} .
        \]
        Note $f_{xx}+f_{yy}=0$: $f$ is harmonic.

  \item $D_1 = 5 > 0$, $D_2 = -10 < 0$. \textbf{Indefinite.} At a critical
        point this would mean a \textbf{saddle}.

  \item $\grad f = (y-2x+6,\,x-2y+3) = (0,0)$. So $2x-y=6$ and
        $x-2y=-3$. From the first, $y=2x-6$; substituting:
        $x-2(2x-6)=-3 \Rightarrow -3x+12=-3 \Rightarrow x=5$, $y=4$.
        Hessian $\Hess = \begin{pmatrix}-2&1\\1&-2\end{pmatrix}$:
        $D_1=-2<0$, $D_2=4-1=3>0$, \textbf{ND}, so $(5,4)$ is a
        \textbf{local maximum} with $f(5,4) = 20-25-16+30+12 = 21$.

  \item $\ln(xy) = \ln x + \ln y$. So $f_x = 1/x$, $f_y = 1/y$,
        $f_{xx} = -1/x^{2}$, $f_{yy} = -1/y^{2}$, $f_{xy} = 0 = f_{yx}$.
        \; $\checkmark$

  \item $f_x = 2(x^{2}+y^{2})\cdot2x = 4x(x^{2}+y^{2})$,
        $f_y = 4y(x^{2}+y^{2})$. Then
        \[
          f_{xx} = 4(x^{2}+y^{2})+4x\cdot2x = 4(3x^{2}+y^{2}), \quad
          f_{yy} = 4(x^{2}+3y^{2}), \quad
          f_{xy} = 4x\cdot2y = 8xy = f_{yx} .
        \]
        Hessian
        \[
          \Hess = 4\begin{pmatrix}3x^{2}+y^{2}&2xy\\2xy&x^{2}+3y^{2}\end{pmatrix} .
        \]

  \item $D_1 = -1 < 0$, $D_2 = 1 - 4 = -3 < 0$. Since $D_2<0$,
        \textbf{indefinite.}

  \item $\grad f = (3x^{2}-3y,\,-3x+3y^{2}) = (0,0) \Rightarrow y=x^{2}$
        and $x=y^{2}=x^{4}$, so $x(x^{3}-1)=0$, giving $x\in\{0,1\}$.
        Critical points: $(0,0)$ and $(1,1)$. Hessian
        $\Hess = \begin{pmatrix}6x&-3\\-3&6y\end{pmatrix}$.
        \begin{itemize}
          \item $(0,0)$: $\Hess = \begin{pmatrix}0&-3\\-3&0\end{pmatrix}$,
                $D_2=-9<0$, \textbf{saddle}.
          \item $(1,1)$: $\Hess = \begin{pmatrix}6&-3\\-3&6\end{pmatrix}$,
                $D_1=6>0$, $D_2=36-9=27>0$, \textbf{PD}: local minimum,
                $f=1-3+1=-1$.
        \end{itemize}

  \item $f_x = \cos(x)\cos(y)$, $f_y = -\sin(x)\sin(y)$. So
        $f_{xx} = -\sin(x)\cos(y)$, $f_{yy} = -\sin(x)\cos(y)$, and
        \[
          f_{xy} = -\cos(x)\sin(y) = f_{yx} . \; \checkmark
        \]

  \item Let $u=xyz$, $f=\mathrm{e}^{u}$. Then $f_x = yz\mathrm{e}^{u}$,
        $f_y = xz\mathrm{e}^{u}$, $f_z = xy\mathrm{e}^{u}$. Direct
        partials: $f_{xx} = (yz)^{2}\mathrm{e}^{u}$,
        $f_{yy} = (xz)^{2}\mathrm{e}^{u}$, $f_{zz} = (xy)^{2}\mathrm{e}^{u}$.
        Cross-partials (using the product rule on
        $f_x = yz\,\mathrm{e}^{xyz}$):
        \[
          f_{xy} = z\mathrm{e}^{xyz} + yz\cdot xz\mathrm{e}^{xyz}
                 = z\mathrm{e}^{xyz}(1+xyz) .
        \]
        Similarly $f_{xz} = y\mathrm{e}^{xyz}(1+xyz)$ and
        $f_{yz} = x\mathrm{e}^{xyz}(1+xyz)$. (Symmetry $f_{xy}=f_{yx}$
        etc.\ confirms Young's theorem.) At $(1,1,1)$: $u=1$,
        $\mathrm{e}^{u}=\mathrm{e}$, $1+xyz=2$, and $yz=xz=xy=1$, so
        \[
          \Hess(1,1,1) = \mathrm{e}\begin{pmatrix}1&2&2\\2&1&2\\2&2&1\end{pmatrix} .
        \]

  \item $D_1 = 4 > 0$, $D_2 = 16 - 4 = 12 > 0$. \textbf{PD.}

  \item $\pi_L = 5L^{-1/2}K^{1/2}-2$, $\pi_K = 5L^{1/2}K^{-1/2}-3$.
        Setting both to zero gives $5K^{1/2}/L^{1/2}=2$ and
        $5L^{1/2}/K^{1/2}=3$. Multiplying these two equations: $25=6$,
        which is a contradiction. So $\pi$ has \textbf{no interior
        critical point}. Why? The production function
        $Q=10L^{1/2}K^{1/2}$ is homogeneous of degree 1 (constant returns
        to scale), so along any ray $(\lambda L,\lambda K)$ revenue and
        cost both scale linearly: $\pi(\lambda L,\lambda K)=\lambda\pi(L,K)$.
        Hence profit is either identically zero, or unbounded above
        (depending on sign of $\pi$ along some direction). With
        $\max_{L/K}(10\sqrt{LK}-2L-3K)/K$: setting $t=L/K$ gives
        $\pi/K = 10\sqrt{t}-2t-3$, maximised at $t=25/4$ giving
        $\pi/K = 25-25/2-3 = 19/2 > 0$. So $\pi\to+\infty$ along the ray
        $L=(25/4)K$. The Hessian test does not apply because there is no
        critical point to classify. (Lesson: an inconsistent
        first-order system is itself a meaningful answer about the
        geometry of the problem.)

  \item $f_x = 2xy^{3}$, $f_y = 3x^{2}y^{2}$. So $f_{xx} = 2y^{3}$,
        $f_{yy} = 6x^{2}y$, and
        \[
          f_{xy} = 6xy^{2} = f_{yx} . \; \checkmark
        \]

  \item Let $g=\sin y-\cos y$, so $f=\mathrm{e}^{x}g$. Then
        $f_x = \mathrm{e}^{x}g$, $f_y = \mathrm{e}^{x}(\cos y+\sin y)$. So
        $f_{xx} = \mathrm{e}^{x}g = \mathrm{e}^{x}(\sin y-\cos y)$,
        $f_{yy} = \mathrm{e}^{x}(-\sin y+\cos y) = -\mathrm{e}^{x}g$, and
        \[
          f_{xy} = \mathrm{e}^{x}(\cos y+\sin y) = f_{yx} . \; \checkmark
        \]
        Hessian
        \[
          \Hess = \mathrm{e}^{x}
          \begin{pmatrix}
            \sin y-\cos y & \cos y+\sin y \\
            \cos y+\sin y & \cos y-\sin y
          \end{pmatrix} .
        \]

  \item $D_1 = 9 > 0$, $D_2 = 36 - 36 = 0$. \textbf{Inconclusive
        (positive semi-definite boundary case)}: Sylvester gives no
        verdict; we cannot rule out PSD with a zero eigenvalue. (In fact
        $H = (3,2)^{T}(3,2)$, so it is PSD with one zero eigenvalue.)

  \item $\grad f = (2x+y-3,\,x+2y) = (0,0)$. From the second: $x=-2y$;
        substituting: $-4y+y-3=0 \Rightarrow y=-1$, $x=2$. Critical point
        $(2,-1)$. Hessian $\Hess = \begin{pmatrix}2&1\\1&2\end{pmatrix}$:
        $D_1=2>0$, $D_2=4-1=3>0$, \textbf{PD}, so $(2,-1)$ is a strict
        \textbf{local minimum} with $f(2,-1) = 4-2+1-6 = -3$.

  \item With $u=y/x$, $f=\arctan u$: $f_x = -y/(x^{2}+y^{2})$,
        $f_y = x/(x^{2}+y^{2})$. Let $S=x^{2}+y^{2}$. Then
        \[
          f_{xx} = \frac{2xy}{S^{2}}, \qquad
          f_{yy} = \frac{-2xy}{S^{2}}, \qquad
          f_{xy} = \frac{-S+2y^{2}}{S^{2}} = \frac{y^{2}-x^{2}}{(x^{2}+y^{2})^{2}}
                  = f_{yx} . \; \checkmark
        \]
        (And $f_{xx}+f_{yy}=0$: $\arctan(y/x)$ is harmonic.)

  \item Let $g=K^{1/2}+L^{1/2}$. From sec06 sol.\ 23, $f_K = g/K^{1/2}$,
        $f_L = g/L^{1/2}$. Then
        \[
          f_{KK} = \frac{\tfrac12K^{-1/2}\cdot K^{1/2}-g\cdot\tfrac12K^{-1/2}}{K}
                  = \frac{\tfrac12-g/(2K^{1/2})}{K}
                  = \frac{1}{2K}-\frac{g}{2K^{3/2}} .
        \]
        More cleanly: $f_K = 1+(L/K)^{1/2}$, so
        $f_{KK} = -\tfrac12L^{1/2}K^{-3/2}$, and by symmetry
        $f_{LL} = -\tfrac12K^{1/2}L^{-3/2}$. Cross-partial:
        $f_{KL} = \tfrac12K^{-1/2}L^{-1/2} = f_{LK}$. \; $\checkmark$
        Hessian
        \[
          \Hess = \tfrac12
          \begin{pmatrix}
            -L^{1/2}K^{-3/2} & K^{-1/2}L^{-1/2} \\
            K^{-1/2}L^{-1/2} & -K^{1/2}L^{-3/2}
          \end{pmatrix} .
        \]

  \item $\grad f = (-4x^{3}+8y,\,-4y^{3}+8x) = (0,0) \Rightarrow x^{3}=2y$
        and $y^{3}=2x$. Substituting: $y=x^{3}/2$, then
        $(x^{3}/2)^{3}=2x \Rightarrow x^{9}/8=2x \Rightarrow x(x^{8}-16)=0$,
        so $x\in\{0,\pm16^{1/8}\}=\{0,\pm\sqrt2\}$. Points: $(0,0)$,
        $(\sqrt2,\sqrt2)$ (since $y=x^{3}/2=2\sqrt2/2=\sqrt2$), and
        $(-\sqrt2,-\sqrt2)$. Hessian
        $\Hess = \begin{pmatrix}-12x^{2}&8\\8&-12y^{2}\end{pmatrix}$.
        \begin{itemize}
          \item $(0,0)$: $\Hess = \begin{pmatrix}0&8\\8&0\end{pmatrix}$,
                $D_2=-64<0$, \textbf{saddle}.
          \item $(\sqrt2,\sqrt2)$: $x^{2}=y^{2}=2$,
                $\Hess = \begin{pmatrix}-24&8\\8&-24\end{pmatrix}$,
                $D_1=-24<0$, $D_2=576-64=512>0$, \textbf{ND}: local
                maximum, $f=-4-4+16=8$.
          \item $(-\sqrt2,-\sqrt2)$: same $H$, $f=-4-4+16=8$. \textbf{Local
                maximum.}
        \end{itemize}

  \item $f_x = 2x\mathrm{e}^{-y}$, $f_y = -x^{2}\mathrm{e}^{-y}$. So
        $f_{xx} = 2\mathrm{e}^{-y}$, $f_{yy} = x^{2}\mathrm{e}^{-y}$, and
        \[
          f_{xy} = -2x\mathrm{e}^{-y} = f_{yx} . \; \checkmark
        \]

  \item $f_x = -y\sin(xy)$, $f_y = -x\sin(xy)$. Then
        $f_{xx} = -y^{2}\cos(xy)$, $f_{yy} = -x^{2}\cos(xy)$, and
        (product rule)
        \[
          f_{xy} = -\sin(xy) - y\cdot x\cos(xy) = -\sin(xy) - xy\cos(xy) .
        \]
        At $(0,0)$: $\sin 0=0$, $\cos 0=1$. So $f_{xx}(0,0)=0$,
        $f_{yy}(0,0)=0$, $f_{xy}(0,0)=0$. Hessian
        \[
          \Hess(0,0) = \begin{pmatrix}0&0\\0&0\end{pmatrix} .
        \]

  \item $D_1 = 0$. Sylvester's stated criterion is inconclusive because
        $D_1$ is not strictly positive or negative. Directly:
        $D_2 = 0-9 = -9 < 0$, which by the third clause of Sylvester
        (sufficient on its own) gives \textbf{indefinite.}

  \item $\grad f = (3x^{2}+3y^{2}-15,\,6xy-12) = (0,0)$. From the second:
        $xy=2$. From the first: $x^{2}+y^{2}=5$. Combined with $xy=2$:
        $(x+y)^{2}=5+4=9$ so $x+y=\pm3$; $(x-y)^{2}=5-4=1$ so
        $x-y=\pm1$. Solutions $(x,y)\in\{(2,1),(1,2),(-2,-1),(-1,-2)\}$.
        Hessian $\Hess = \begin{pmatrix}6x&6y\\6y&6x\end{pmatrix}$, so
        $D_1=6x$, $D_2=36(x^{2}-y^{2})$.
        \begin{itemize}
          \item $(2,1)$: $D_1=12>0$, $D_2=36(4-1)=108>0$, \textbf{PD}:
                local minimum, $f=8+6-30-12=-28$.
          \item $(1,2)$: $D_1=6>0$, $D_2=36(1-4)=-108<0$, \textbf{saddle}.
          \item $(-2,-1)$: $D_1=-12<0$, $D_2=108>0$, \textbf{ND}: local
                maximum, $f=-8-6+30+12=28$.
          \item $(-1,-2)$: $D_1=-6<0$, $D_2=-108<0$, $D_2<0$ gives
                \textbf{saddle}.
        \end{itemize}

  \item Let $S=1+x^{2}+y^{2}$, $f=\ln S$, $f_x = 2x/S$, $f_y = 2y/S$.
        Then (using quotient rule):
        \[
          f_{xx} = \frac{2S-2x\cdot2x}{S^{2}} = \frac{2(1+y^{2}-x^{2})}{S^{2}}, \quad
          f_{yy} = \frac{2(1+x^{2}-y^{2})}{S^{2}}, \quad
          f_{xy} = \frac{-4xy}{S^{2}} = f_{yx} . \; \checkmark
        \]
        At the origin $(0,0)$: $S=1$, so
        \[
          \Hess(0,0) = \begin{pmatrix}2&0\\0&2\end{pmatrix} .
        \]

  \item $\grad f = \mathrm{e}^{-(x^{2}+y^{2})}(-2x,-2y) = (0,0)
        \Rightarrow x=y=0$, since $\mathrm{e}^{-r^{2}}>0$. Critical point
        $(0,0)$. Compute the Hessian: $f_x = -2x\mathrm{e}^{-(x^{2}+y^{2})}$,
        and at $(0,0)$, $f_{xx} = -2\mathrm{e}^{0}+(-2x)(-2x)\mathrm{e}^{0}\big|_{0,0}=-2$.
        Similarly $f_{yy}(0,0)=-2$, and
        $f_{xy}(0,0) = (-2x)(-2y)\mathrm{e}^{0}\big|_{0,0}=0$. Hessian
        $\Hess(0,0) = \begin{pmatrix}-2&0\\0&-2\end{pmatrix}$: $D_1=-2<0$,
        $D_2=4>0$, \textbf{ND: strict local maximum} at $(0,0)$ with
        $f(0,0)=1$. (Indeed, this is the global maximum, since $f\to0$
        at infinity and $f>0$ everywhere.)
\end{solutions}

\section{Additional Problems}

\begin{problems}
  \item Compute all second partials of $f(x,y) = 3x^{3} - 4x^{2}y + 5y^{3}$
        and verify Young's theorem.
  \item Construct the Hessian of
        $f(x,y) = 3x^{2} - 4xy + 2y^{2} + x - 5y$.
  \item Classify the definiteness of
        $\Hess = \begin{pmatrix}5&-1\\-1&4\end{pmatrix}$ using Sylvester's
        criterion.
  \item Find and classify the critical points of
        $f(x,y) = 2x^{2} + 3y^{2} - 4x + 6y + 1$.
  \item Compute all second partials of $f(x,y) = (2x+y)^{4}$ and verify
        $f_{xy} = f_{yx}$.
  \item Construct the Hessian of $f(x,y) = \mathrm{e}^{x+2y}$.
  \item Classify $\Hess = \begin{pmatrix}-3&1\\1&-2\end{pmatrix}$.
  \item Find and classify the critical points of
        $f(x,y) = -2x^{2} - 4y^{2} + 8x + 8y - 3$.
  \item Compute all second partials of $f(x,y) = xy\ln y$ ($y>0$) and
        verify Young's theorem.
  \item Construct the Hessian of $f(x,y) = x^{3}y + xy^{3}$.
  \item Classify $\Hess = \begin{pmatrix}2&3\\3&4\end{pmatrix}$.
  \item Find and classify the critical points of
        $f(x,y) = x^{2} - 2xy + 3y^{2} + 2x - 6y$.
  \item Compute all second partials of the Cobb--Douglas
        $f(K,L) = K^{1/2}L^{1/2}$ ($K,L>0$) and verify Young's theorem.
  \item Find and classify the critical points of
        $f(x,y) = x^{3} + y^{3} - 3x - 3y$.
  \item Classify $\Hess = \begin{pmatrix}-5&0\\0&-3\end{pmatrix}$ and state
        what it implies if $H$ is the Hessian at a critical point.
\end{problems}

\section{Additional Solutions}

\begin{solutions}
  \item $f_x = 9x^{2} - 8xy$, $f_y = -4x^{2} + 15y^{2}$. Then
        $f_{xx} = 18x - 8y$, $f_{yy} = 30y$, and
        \[
          f_{xy} = -8x = f_{yx} . \; \checkmark
        \]

  \item $f_x = 6x - 4y + 1$, $f_y = -4x + 4y - 5$. Hessian
        \[
          \Hess = \begin{pmatrix}6&-4\\-4&4\end{pmatrix} .
        \]

  \item $D_1 = 5 > 0$, $D_2 = 20 - 1 = 19 > 0$. \textbf{PD.}

  \item $\grad f = (4x-4,\,6y+6) = (0,0) \Rightarrow x=1,\,y=-1$. Hessian
        $\Hess = \begin{pmatrix}4&0\\0&6\end{pmatrix}$: $D_1=4>0$,
        $D_2=24>0$, \textbf{PD}, so $(1,-1)$ is a strict \textbf{local
        minimum} with $f(1,-1) = 2+3-4-6+1 = -4$.

  \item Chain rule: $f_x = 4(2x+y)^{3}\cdot2 = 8(2x+y)^{3}$,
        $f_y = 4(2x+y)^{3}$. So $f_{xx} = 48(2x+y)^{2}$,
        $f_{yy} = 12(2x+y)^{2}$, and
        \[
          f_{xy} = 24(2x+y)^{2} = f_{yx} . \; \checkmark
        \]

  \item $f_x = \mathrm{e}^{x+2y}$, $f_y = 2\mathrm{e}^{x+2y}$. So
        $f_{xx} = \mathrm{e}^{x+2y}$, $f_{yy} = 4\mathrm{e}^{x+2y}$,
        $f_{xy} = 2\mathrm{e}^{x+2y}$. Hessian
        \[
          \Hess = \mathrm{e}^{x+2y}\begin{pmatrix}1&2\\2&4\end{pmatrix} .
        \]
        (Note $\det\Hess = 0$: this is rank-1 PSD, consistent with $f$
        being a function of the single combination $x+2y$.)

  \item $D_1 = -3 < 0$, $D_2 = 6 - 1 = 5 > 0$. \textbf{ND.}

  \item $\grad f = (-4x+8,\,-8y+8) = (0,0) \Rightarrow x=2,\,y=1$. Hessian
        $\Hess = \begin{pmatrix}-4&0\\0&-8\end{pmatrix}$: $D_1=-4<0$,
        $D_2=32>0$, \textbf{ND}, so $(2,1)$ is a strict \textbf{local
        maximum} with $f(2,1) = -8-4+16+8-3 = 9$.

  \item $f_x = y\ln y$, $f_y = x\ln y + x = x(\ln y + 1)$. So
        $f_{xx} = 0$, $f_{yy} = x/y$, and
        \[
          f_{xy} = \ln y + 1 = f_{yx} . \; \checkmark
        \]

  \item $f_x = 3x^{2}y + y^{3}$, $f_y = x^{3} + 3xy^{2}$. So
        $f_{xx} = 6xy$, $f_{yy} = 6xy$,
        $f_{xy} = 3x^{2} + 3y^{2} = f_{yx}$. Hessian
        \[
          \Hess = \begin{pmatrix}
            6xy & 3x^{2}+3y^{2} \\
            3x^{2}+3y^{2} & 6xy
          \end{pmatrix} .
        \]

  \item $D_1 = 2 > 0$, $D_2 = 8 - 9 = -1 < 0$. \textbf{Indefinite
        (saddle).}

  \item $\grad f = (2x-2y+2,\,-2x+6y-6) = (0,0)$. From the first:
        $x = y-1$. Substitute: $-2(y-1)+6y-6 = 4y-4 = 0 \Rightarrow y=1$,
        $x=0$. Hessian $\Hess = \begin{pmatrix}2&-2\\-2&6\end{pmatrix}$:
        $D_1=2>0$, $D_2=12-4=8>0$, \textbf{PD}, so $(0,1)$ is a strict
        \textbf{local minimum} with $f(0,1) = 0-0+3+0-6 = -3$.

  \item $f_K = \tfrac12K^{-1/2}L^{1/2}$, $f_L = \tfrac12K^{1/2}L^{-1/2}$.
        So
        \[
          f_{KK} = -\tfrac14K^{-3/2}L^{1/2}, \quad
          f_{LL} = -\tfrac14K^{1/2}L^{-3/2}, \quad
          f_{KL} = \tfrac14K^{-1/2}L^{-1/2} = f_{LK} . \; \checkmark
        \]

  \item $\grad f = (3x^{2}-3,\,3y^{2}-3) = (0,0) \Rightarrow x^{2}=1,\,
        y^{2}=1$, giving $(\pm1,\pm1)$. Hessian
        $\Hess = \begin{pmatrix}6x&0\\0&6y\end{pmatrix}$, so $D_1=6x$,
        $D_2=36xy$.
        \begin{itemize}
          \item $(1,1)$: $D_1=6>0$, $D_2=36>0$, \textbf{PD}: local
                minimum, $f = 1+1-3-3 = -4$.
          \item $(1,-1)$: $D_2=-36<0$, \textbf{saddle}.
          \item $(-1,1)$: $D_2=-36<0$, \textbf{saddle}.
          \item $(-1,-1)$: $D_1=-6<0$, $D_2=36>0$, \textbf{ND}: local
                maximum, $f = -1-1+3+3 = 4$.
        \end{itemize}

  \item $D_1 = -5 < 0$, $D_2 = 15 > 0$. \textbf{ND}: at a critical point
        this would mean a \textbf{strict local maximum.}
\end{solutions}

\chapter{Multivariable Chain Rule and Total Differential}

When several inputs move at once, the total differential bundles the partial
derivatives into one linear approximation, and the multivariable chain rule
pipes that approximation through any composed function. Together they
generate the marginal rate of substitution along an indifference curve, the
marginal rate of technical substitution along an isoquant, and (as a
preview) the envelope theorem along an optimal value function.

\section{Core Rules}

\begin{rulebox}{Total Differential}
If $z=f(x,y)$ is differentiable at $(x,y)$, the \textbf{total differential}
is
\[
  \mathrm{d}z \;=\; f_{x}(x,y)\,\mathrm{d}x \;+\; f_{y}(x,y)\,\mathrm{d}y .
\]

\textbf{Interpretation.} $\mathrm{d}z$ is the best linear estimate of the
actual change $\Delta z$ when $x$ and $y$ shift by $\mathrm{d}x$ and
$\mathrm{d}y$. In economics, $\mathrm{d}U \approx MU_{x}\,\mathrm{d}x +
MU_{y}\,\mathrm{d}y$ tells how utility responds to a small basket change;
$\mathrm{d}Q \approx MP_{K}\,\mathrm{d}K + MP_{L}\,\mathrm{d}L$ does the same
for output.

For $n$ inputs, $z=f(x_{1},\ldots,x_{n})$:
\[
  \mathrm{d}z \;=\; \sum_{i=1}^{n} \pd{f}{x_{i}}\,\mathrm{d}x_{i} .
\]
\end{rulebox}

\begin{rulebox}{Chain Rule, Single Driver $t$}
If $z=f(x,y)$ with $x=x(t)$ and $y=y(t)$,
\[
  \dd{z}{t} \;=\; f_{x}\,\dd{x}{t} \;+\; f_{y}\,\dd{y}{t} .
\]
\end{rulebox}

\begin{rulebox}{Chain Rule, Two Drivers $s,t$}
If $z=f(x,y)$ with $x=x(s,t)$, $y=y(s,t)$,
\[
  \pd{z}{s} \;=\; f_{x}\,\pd{x}{s} + f_{y}\,\pd{y}{s},\qquad
  \pd{z}{t} \;=\; f_{x}\,\pd{x}{t} + f_{y}\,\pd{y}{t} .
\]
\end{rulebox}

\begin{rulebox}{Level Sets, MRS, MRTS}
On a level curve $f(x,y)=c$ the total differential is zero:
\[
  f_{x}\,\mathrm{d}x + f_{y}\,\mathrm{d}y = 0
  \;\Longrightarrow\;
  \dd{y}{x}\bigg|_{f=c} = -\frac{f_{x}}{f_{y}} .
\]
For a utility $U(x,y)$, the (signed) slope $-U_{x}/U_{y}$ is minus the
\textbf{MRS}; for a production function $Q(K,L)$, $-Q_{K}/Q_{L}$ is minus
the \textbf{MRTS}. We will state explicitly which sign convention we use in
each problem.
\end{rulebox}

\section{Worked Example}

\begin{worked}{With $z=x^{2}y$, $x=t^{2}$, $y=\ln t$, find $\mathrm{d}z/\mathrm{d}t$ two ways; then use the total differential on $Q=K^{0.3}L^{0.7}$}
\wstep{Step 1. Draw the dependency before writing a formula.}
$t$ drives both $x$ and $y$; $x$ and $y$ together drive $z$. So there are
\textbf{two paths} from $t$ to $z$, and the chain rule adds one term per
path:
\[
  \dd{z}{t}=\pd{z}{x}\dd{x}{t}+\pd{z}{y}\dd{y}{t} .
\]
\wtrap{One term per path. Losing a term is the standard error, and it is
invisible in the answer --- the result still looks like a derivative.
Sketching the tree (one node $t$, two arrows to $x$ and $y$, two arrows into
$z$) takes five seconds and prevents it.}

\wstep{Step 2. Note which symbol goes where.}
$\partial z/\partial x$ is \emph{partial}, because $z$ depends on two
things. $\mathrm{d}x/\mathrm{d}t$ is \emph{total}, because $x$ depends on
$t$ alone. And $\mathrm{d}z/\mathrm{d}t$ on the left is total, because once
the substitutions are made $z$ is a function of $t$ only. The notation is
carrying real information here.

\wstep{Step 3. Compute the four pieces.}
\[
  \pd{z}{x}=2xy,\qquad
  \pd{z}{y}=x^{2},\qquad
  \dd{x}{t}=2t,\qquad
  \dd{y}{t}=\frac1t .
\]

\wstep{Step 4. Assemble, then substitute the paths.}
\[
  \dd{z}{t}=2xy\cdot2t+x^{2}\cdot\frac1t
  =2t^{2}\ln t\cdot2t+\frac{t^{4}}{t}
  =\boxed{4t^{3}\ln t+t^{3}}
\]
\wtrap{Substitute $x=t^{2}$ and $y=\ln t$ only \emph{after} assembling the
formula. Substituting first turns the exercise into a one-variable problem
and defeats its purpose --- and in the many cases where $x(t)$ is not known
explicitly, it is not even possible.}

\wstep{Step 5. Check by direct substitution, which here is available.}
$z=(t^{2})^{2}\ln t=t^{4}\ln t$, so
\[
  \dd{z}{t}=4t^{3}\ln t+t^{4}\cdot\frac1t=4t^{3}\ln t+t^{3} . \checkmark
\]
Two routes, one answer. When both are available, use the direct one as the
check and the chain rule as the method.

\wstep{Step 6. The total differential: the chain rule with the $t$ removed.}
\[
  \mathrm{d}z=\pd{z}{x}\mathrm{d}x+\pd{z}{y}\mathrm{d}y
\]
is the first-order approximation to the change in $z$ caused by small
changes in $x$ and $y$ --- the two-variable linearisation. For
$Q=K^{0.3}L^{0.7}$:
\[
  \mathrm{d}Q=0.3K^{-0.7}L^{0.7}\,\mathrm{d}K+0.7K^{0.3}L^{-0.3}\,\mathrm{d}L .
\]

\wstep{Step 7. Divide by $Q$ and watch it turn into percentages.}
\[
  \frac{\mathrm{d}Q}{Q}
  =0.3\,\frac{\mathrm{d}K}{K}+0.7\,\frac{\mathrm{d}L}{L} .
\]
Every power cancels. This is the growth-accounting equation: the percentage
change in output is the \emph{exponent-weighted average} of the percentage
changes in inputs, and the exponents are the elasticities. If capital rises
$2\%$ while labour falls $1\%$:
\[
  \frac{\mathrm{d}Q}{Q}=0.3(0.02)+0.7(-0.01)=0.006-0.007=-0.001=-0.1\% .
\]
\wtrap{This identity holds only for Cobb--Douglas, where the elasticities
are the constant exponents. For a general $F(K,L)$ the weights are
$\frac{K F_{K}}{F}$ and $\frac{L F_{L}}{F}$, which vary from point to
point. Applying the constant-weight version to a CES function is a real
error.}

\wstep{Step 8. Two drivers instead of one.}
If $z=f(x,y)$ with $x=x(s,t)$ and $y=y(s,t)$, the same path-counting gives
one equation per driver:
\[
  \pd{z}{t}=\pd{z}{x}\pd{x}{t}+\pd{z}{y}\pd{y}{t},
  \qquad
  \pd{z}{s}=\pd{z}{x}\pd{x}{s}+\pd{z}{y}\pd{y}{s} .
\]
All derivatives are now partial, because two independent variables remain.
In matrix form these two lines are $\grad_{(s,t)}z=J^{\!\top}\grad_{(x,y)}z$
with $J$ the Jacobian of \chref{derivatives}{13}: the same rule, written once.
\end{worked}

\section{Practice Problems}

\begin{problems}
  \item Compute $\mathrm{d}z$ for $z=4x^{2}-5xy+2y^{2}$.
  \item Find $\dd{z}{t}$ for $z=x^{2}-y^{2}$ with $x=2t+1$, $y=t^{2}$.
  \item For $U(x,y)=x^{2}+4y=c$, find $\dd{y}{x}$ and state the unsigned
        MRS.
  \item Compute $\pd{z}{s}$ and $\pd{z}{t}$ for $z=x^{2}-3y$ with
        $x=s+2t$, $y=2s-t$.
  \item Compute $\mathrm{d}z$ for $z=x^{4}y^{-2}$.
  \item Find $\dd{z}{t}$ for $z=\mathrm{e}^{xy}$ with $x=\sin t$,
        $y=\cos t$.
  \item For the isoquant $Q(K,L)=K^{0.4}L^{0.6}=Q_{0}$, find $\dd{L}{K}$
        and the unsigned MRTS.
  \item Compute $\mathrm{d}z$ for $z=\ln(3x+2y)$.
  \item Compute $\pd{z}{s}$ and $\pd{z}{t}$ for $z=\mathrm{e}^{x}\cos y$
        with $x=st$, $y=s+t$.
  \item Find $\dd{z}{t}$ for $z=x^{2}y$ with $x=\ln t$, $y=t^{3}$.
  \item For $U(x,y)=x^{2/3}y^{1/3}=c$, find $\dd{y}{x}$ and the unsigned
        MRS.
  \item Compute $\mathrm{d}z$ for $z=\mathrm{e}^{x^{2}+y}$.
  \item Find $\dd{z}{t}$ for $z=\sqrt{x+y^{2}}$ with $x=t^{2}$, $y=t$.
  \item Compute $\pd{z}{s}$ and $\pd{z}{t}$ for $z=\ln(x^{2}+y^{2})$ with
        $x=s/t$, $y=t/s$.
  \item For $U(x,y)=\sqrt{x}+\sqrt{y}=c$, find $\dd{y}{x}$ and the unsigned
        MRS.
  \item Compute $\mathrm{d}z$ for $z=\dfrac{x+y}{x-y}$.
  \item Find $\dd{z}{t}$ for $z=\arctan(y/x)$ with $x=\mathrm{e}^{t}$,
        $y=\mathrm{e}^{2t}$.
  \item For $Q(K,L)=(aK^{1/2}+bL^{1/2})^{2}=Q_{0}$ (with $a,b>0$), find
        $\dd{L}{K}$ and the unsigned MRTS.
  \item Compute $\pd{z}{s}$ and $\pd{z}{t}$ for $z=x^{1/2}y^{1/2}$ with
        $x=s^{2}+t^{2}$, $y=st$.
  \item Compute $\mathrm{d}z$ for $z=x\sin y+y\cos x$.
  \item For $U(x,y)=\ln(x+1)+\ln(y+1)=c$, find $\dd{y}{x}$ and the
        unsigned MRS.
  \item Find $\dd{z}{t}$ for $z=x\ln y$ with $x=t^{2}+1$, $y=\mathrm{e}^{t}$.
  \item Compute $\mathrm{d}z$ for $z=x^{3}\mathrm{e}^{2y}$.
  \item Find $\dd{z}{t}$ for $z=x^{2}+y^{2}$ with $x=\cos t$, $y=\sin t$,
        and verify the result.
  \item For $U(x,y)=x+y^{2}=c$, find $\dd{y}{x}$ and state where the
        result is well-defined.
  \item Compute $\pd{z}{s}$ and $\pd{z}{t}$ for $z=xy^{2}$ with $x=s+t$,
        $y=s-t$.
  \item Compute $\mathrm{d}z$ for $z=\arctan(xy)$.
  \item Find $\dd{z}{t}$ for $z=\ln(x^{2}+y^{2})$ with $x=t$, $y=t^{2}$.
  \item For the isoquant $Q(K,L)=AK^{\alpha}L^{\beta}=Q_{0}$ (general
        Cobb--Douglas with $\alpha,\beta>0$), derive the MRTS in terms of
        $\alpha,\beta,K,L$.
  \item Compute $\mathrm{d}z$ for $z=\dfrac{x^{2}}{y}+\dfrac{y^{2}}{x}$.
  \item For $U(x,y)=\mathrm{e}^{x}+\mathrm{e}^{y}=c$, find $\dd{y}{x}$ and
        the unsigned MRS.
  \item Find $\dd{z}{t}$ for $z=\mathrm{e}^{x}\sin y$ with $x=2t$,
        $y=t^{2}$.
  \item Compute $\pd{z}{s}$ and $\pd{z}{t}$ for $z=x^{2}y-y^{3}$ with
        $x=s/t$, $y=st$.
  \item Compute $\mathrm{d}z$ for $z=(x+y)\mathrm{e}^{x-y}$.
  \item For the indifference curve $U(x,y)=\alpha\ln x+\beta\ln y=c$ (with
        $\alpha,\beta>0$), find $\dd{y}{x}$ and the MRS.
  \item Find $\dd{z}{t}$ for $z=\sqrt{xy}$ with $x=t^{3}$, $y=t$.
  \item Compute $\mathrm{d}z$ for $z=\ln(x^{2}+y^{2}+1)$.
  \item Compute $\pd{z}{s}$ and $\pd{z}{t}$ for $z=\mathrm{e}^{xy}$ with
        $x=s+t$, $y=s-t$.
  \item For $Q(K,L)=K+2\sqrt{KL}+L=Q_{0}$ (a perfect-square form), find
        the MRTS.
  \item Find $\dd{z}{t}$ for $z=\cos(xy)$ with $x=t$, $y=1/t$ (for $t>0$).
        Verify by direct substitution.
  \item For the indifference curve $U(x,y)=xy+x+y=c$, find $\dd{y}{x}$
        and the unsigned MRS.
  \item Use the total differential of $V(P,Y)=Y/P$ (real income) to
        approximate $\mathrm{d}V$ when $\mathrm{d}P=0.02P$ and
        $\mathrm{d}Y=0.03Y$. Interpret.
  \item Compute $\pd{z}{s}$ and $\pd{z}{t}$ for $z=x^{2}+2xy$ with
        $x=s\cos t$, $y=s\sin t$.
  \item For $U(x_{1},x_{2},x_{3})=x_{1}+2x_{2}+3x_{3}=c$, find
        $\pd{x_{3}}{x_{1}}$ and $\pd{x_{3}}{x_{2}}$ along the level set,
        and interpret as marginal rates of substitution.
\end{problems}

\section{Solutions}

\begin{solutions}
  \item $f_{x}=8x-5y$, $f_{y}=-5x+4y$.
        \[
          \mathrm{d}z \;=\; (8x-5y)\,\mathrm{d}x+(-5x+4y)\,\mathrm{d}y .
        \]

  \item $f_{x}=2x$, $f_{y}=-2y$, $\dd{x}{t}=2$, $\dd{y}{t}=2t$.
        \[
          \dd{z}{t}=2x\cdot 2+(-2y)\cdot 2t=4(2t+1)-4t^{3}=8t+4-4t^{3} .
        \]
        \textbf{Check.} $z=(2t+1)^{2}-t^{4}$, $\dd{z}{t}=8t+4-4t^{3}$.
        \checkmark

  \item $U_{x}=2x$, $U_{y}=4$. Signed slope $\dd{y}{x}=-2x/4=-x/2$.
        \textbf{MRS} $=x/2$.

  \item $f_{x}=2x$, $f_{y}=-3$; $x_{s}=1$, $x_{t}=2$, $y_{s}=2$,
        $y_{t}=-1$.
        \[
          \pd{z}{s}=2x(1)+(-3)(2)=2(s+2t)-6,\qquad
          \pd{z}{t}=2x(2)+(-3)(-1)=4(s+2t)+3 .
        \]

  \item $f_{x}=4x^{3}y^{-2}$, $f_{y}=-2x^{4}y^{-3}$.
        \[
          \mathrm{d}z \;=\; \frac{4x^{3}}{y^{2}}\,\mathrm{d}x
            -\frac{2x^{4}}{y^{3}}\,\mathrm{d}y .
        \]

  \item $f_{x}=y\mathrm{e}^{xy}$, $f_{y}=x\mathrm{e}^{xy}$;
        $\dd{x}{t}=\cos t$, $\dd{y}{t}=-\sin t$.
        \[
          \dd{z}{t}=\mathrm{e}^{xy}(y\cos t-x\sin t)
            =\mathrm{e}^{\sin t\cos t}(\cos^{2}t-\sin^{2}t)
            =\mathrm{e}^{\sin t\cos t}\cos(2t) .
        \]

  \item $Q_{K}=0.4K^{-0.6}L^{0.6}$, $Q_{L}=0.6K^{0.4}L^{-0.4}$.
        \[
          \frac{Q_{K}}{Q_{L}}=\frac{0.4}{0.6}\cdot\frac{L}{K}
            =\frac{2}{3}\,\frac{L}{K} .
        \]
        $\dd{L}{K}=-\tfrac{2}{3}L/K$. \textbf{MRTS} $=\tfrac{2}{3}L/K$.

  \item $f_{x}=3/(3x+2y)$, $f_{y}=2/(3x+2y)$.
        \[
          \mathrm{d}z \;=\; \frac{3\,\mathrm{d}x+2\,\mathrm{d}y}{3x+2y} .
        \]

  \item $f_{x}=\mathrm{e}^{x}\cos y$, $f_{y}=-\mathrm{e}^{x}\sin y$;
        $x_{s}=t$, $x_{t}=s$, $y_{s}=y_{t}=1$.
        \[
          \pd{z}{s}=\mathrm{e}^{st}\bigl(t\cos(s+t)-\sin(s+t)\bigr),\qquad
          \pd{z}{t}=\mathrm{e}^{st}\bigl(s\cos(s+t)-\sin(s+t)\bigr) .
        \]

  \item $f_{x}=2xy$, $f_{y}=x^{2}$; $\dd{x}{t}=1/t$, $\dd{y}{t}=3t^{2}$.
        \[
          \dd{z}{t}=2xy\cdot\frac{1}{t}+x^{2}\cdot 3t^{2}
            =\frac{2(\ln t)t^{3}}{t}+3t^{2}(\ln t)^{2}
            =2t^{2}\ln t+3t^{2}(\ln t)^{2} .
        \]

  \item $U_{x}=\tfrac{2}{3}x^{-1/3}y^{1/3}$,
        $U_{y}=\tfrac{1}{3}x^{2/3}y^{-2/3}$.
        \[
          \frac{U_{x}}{U_{y}}=2\,\frac{y}{x} .
        \]
        $\dd{y}{x}=-2y/x$. \textbf{MRS} $=2y/x$.

  \item $f_{x}=2x\mathrm{e}^{x^{2}+y}$, $f_{y}=\mathrm{e}^{x^{2}+y}$.
        \[
          \mathrm{d}z \;=\; \mathrm{e}^{x^{2}+y}(2x\,\mathrm{d}x+\mathrm{d}y) .
        \]

  \item $f_{x}=1/(2\sqrt{x+y^{2}})$, $f_{y}=y/\sqrt{x+y^{2}}$;
        $\dd{x}{t}=2t$, $\dd{y}{t}=1$.
        \[
          \dd{z}{t}=\frac{2t}{2\sqrt{x+y^{2}}}+\frac{y}{\sqrt{x+y^{2}}}
            =\frac{t+y}{\sqrt{x+y^{2}}}=\frac{2t}{\sqrt{2t^{2}}}
            =\frac{\sqrt{2}\,t}{|t|} .
        \]
        For $t>0$: $\dd{z}{t}=\sqrt{2}$. (Indeed $z=\sqrt{2}\,|t|$.)

  \item $f_{x}=2x/(x^{2}+y^{2})$, $f_{y}=2y/(x^{2}+y^{2})$;
        $x_{s}=1/t$, $x_{t}=-s/t^{2}$, $y_{s}=-t/s^{2}$, $y_{t}=1/s$.
        \[
          \pd{z}{s}=\frac{2}{x^{2}+y^{2}}\Bigl(\frac{x}{t}-\frac{yt}{s^{2}}\Bigr)
            =\frac{2}{x^{2}+y^{2}}\Bigl(\frac{s}{t^{2}}-\frac{t^{2}}{s^{3}}\Bigr) .
        \]
        \[
          \pd{z}{t}=\frac{2}{x^{2}+y^{2}}\Bigl(\frac{t}{s^{2}}-\frac{s^{2}}{t^{3}}\Bigr) .
        \]

  \item $U_{x}=1/(2\sqrt{x})$, $U_{y}=1/(2\sqrt{y})$.
        \[
          \frac{U_{x}}{U_{y}}=\sqrt{\frac{y}{x}} .
        \]
        $\dd{y}{x}=-\sqrt{y/x}$. \textbf{MRS} $=\sqrt{y/x}$.

  \item Quotient rule.
        \[
          f_{x}=\frac{(x-y)-(x+y)}{(x-y)^{2}}=\frac{-2y}{(x-y)^{2}},\qquad
          f_{y}=\frac{(x-y)+(x+y)}{(x-y)^{2}}=\frac{2x}{(x-y)^{2}} .
        \]
        \[
          \mathrm{d}z \;=\; \frac{-2y\,\mathrm{d}x+2x\,\mathrm{d}y}{(x-y)^{2}} .
        \]

  \item Let $u=y/x$. $f_{x}=-y/(x^{2}+y^{2})$, $f_{y}=x/(x^{2}+y^{2})$;
        $\dd{x}{t}=\mathrm{e}^{t}$, $\dd{y}{t}=2\mathrm{e}^{2t}$.
        \[
          \dd{z}{t}=\frac{-y\mathrm{e}^{t}+2x\mathrm{e}^{2t}}{x^{2}+y^{2}} .
        \]
        Substitute $x=\mathrm{e}^{t}$, $y=\mathrm{e}^{2t}$,
        $x^{2}+y^{2}=\mathrm{e}^{2t}(1+\mathrm{e}^{2t})$:
        \[
          \dd{z}{t}=\frac{-\mathrm{e}^{3t}+2\mathrm{e}^{3t}}
            {\mathrm{e}^{2t}(1+\mathrm{e}^{2t})}
            =\frac{\mathrm{e}^{t}}{1+\mathrm{e}^{2t}} .
        \]

  \item Let $A=aK^{1/2}+bL^{1/2}$, $Q=A^{2}$. $Q_{K}=2A\cdot(a/2)K^{-1/2}
        =aAK^{-1/2}$, $Q_{L}=bAL^{-1/2}$.
        \[
          \frac{Q_{K}}{Q_{L}}=\frac{a}{b}\sqrt{\frac{L}{K}} .
        \]
        $\dd{L}{K}=-(a/b)\sqrt{L/K}$. \textbf{MRTS} $=(a/b)\sqrt{L/K}$.

  \item $f_{x}=\tfrac12\sqrt{y/x}$, $f_{y}=\tfrac12\sqrt{x/y}$;
        $x_{s}=2s$, $x_{t}=2t$, $y_{s}=t$, $y_{t}=s$.
        \[
          \pd{z}{s}=s\sqrt{\frac{y}{x}}+\frac{t}{2}\sqrt{\frac{x}{y}},\qquad
          \pd{z}{t}=t\sqrt{\frac{y}{x}}+\frac{s}{2}\sqrt{\frac{x}{y}} .
        \]

  \item $f_{x}=\sin y-y\sin x$, $f_{y}=x\cos y+\cos x$.
        \[
          \mathrm{d}z \;=\; (\sin y-y\sin x)\,\mathrm{d}x
            +(x\cos y+\cos x)\,\mathrm{d}y .
        \]

  \item $U_{x}=1/(x+1)$, $U_{y}=1/(y+1)$.
        \[
          \frac{U_{x}}{U_{y}}=\frac{y+1}{x+1} .
        \]
        $\dd{y}{x}=-(y+1)/(x+1)$. \textbf{MRS} $=(y+1)/(x+1)$.

  \item $f_{x}=\ln y$, $f_{y}=x/y$; $\dd{x}{t}=2t$, $\dd{y}{t}=\mathrm{e}^{t}$.
        \[
          \dd{z}{t}=(\ln y)(2t)+\frac{x}{y}\mathrm{e}^{t}
            =t(2t)+\frac{(t^{2}+1)\mathrm{e}^{t}}{\mathrm{e}^{t}}
            =3t^{2}+1 .
        \]
        \textbf{Check.} $z=(t^{2}+1)t=t^{3}+t$, $\dd{z}{t}=3t^{2}+1$.
        \checkmark

  \item $f_{x}=3x^{2}\mathrm{e}^{2y}$, $f_{y}=2x^{3}\mathrm{e}^{2y}$.
        \[
          \mathrm{d}z \;=\; x^{2}\mathrm{e}^{2y}(3\,\mathrm{d}x+2x\,\mathrm{d}y) .
        \]

  \item $f_{x}=2x$, $f_{y}=2y$; $\dd{x}{t}=-\sin t$, $\dd{y}{t}=\cos t$.
        \[
          \dd{z}{t}=2x(-\sin t)+2y(\cos t)=-2\cos t\sin t+2\sin t\cos t=0 .
        \]
        \textbf{Check.} $z=\cos^{2}t+\sin^{2}t=1$, so $\dd{z}{t}=0$.
        \checkmark

  \item $U_{x}=1$, $U_{y}=2y$. Signed slope: $\dd{y}{x}=-1/(2y)$,
        well-defined for $y\neq 0$. The MRS in the usual sense is
        $1/(2y)$.

  \item $f_{x}=y^{2}$, $f_{y}=2xy$; $x_{s}=1$, $x_{t}=1$, $y_{s}=1$,
        $y_{t}=-1$.
        \[
          \pd{z}{s}=y^{2}(1)+2xy(1)=y^{2}+2xy=(s-t)^{2}+2(s+t)(s-t) .
        \]
        Simplify: $(s-t)^{2}+2(s^{2}-t^{2})=s^{2}-2st+t^{2}+2s^{2}-2t^{2}
          =3s^{2}-2st-t^{2}$.
        \[
          \pd{z}{t}=y^{2}(1)+2xy(-1)=y^{2}-2xy=(s-t)^{2}-2(s+t)(s-t)
            =-s^{2}-2st+3t^{2} .
        \]

  \item $f_{x}=y/(1+x^{2}y^{2})$, $f_{y}=x/(1+x^{2}y^{2})$.
        \[
          \mathrm{d}z \;=\; \frac{y\,\mathrm{d}x+x\,\mathrm{d}y}{1+x^{2}y^{2}} .
        \]

  \item $f_{x}=2x/(x^{2}+y^{2})$, $f_{y}=2y/(x^{2}+y^{2})$;
        $\dd{x}{t}=1$, $\dd{y}{t}=2t$.
        \[
          \dd{z}{t}=\frac{2x+2y(2t)}{x^{2}+y^{2}}
            =\frac{2t+4t^{3}}{t^{2}+t^{4}}
            =\frac{2t(1+2t^{2})}{t^{2}(1+t^{2})}
            =\frac{2(1+2t^{2})}{t(1+t^{2})} .
        \]

  \item $Q_{K}=A\alpha K^{\alpha-1}L^{\beta}$,
        $Q_{L}=A\beta K^{\alpha}L^{\beta-1}$.
        \[
          \frac{Q_{K}}{Q_{L}}=\frac{\alpha}{\beta}\cdot\frac{L}{K} .
        \]
        $\dd{L}{K}=-(\alpha/\beta)(L/K)$.
        \textbf{MRTS} $=(\alpha/\beta)(L/K)$.

  \item $f_{x}=2x/y-y^{2}/x^{2}$, $f_{y}=-x^{2}/y^{2}+2y/x$.
        \[
          \mathrm{d}z \;=\; \Bigl(\frac{2x}{y}-\frac{y^{2}}{x^{2}}\Bigr)\mathrm{d}x
            +\Bigl(\frac{2y}{x}-\frac{x^{2}}{y^{2}}\Bigr)\mathrm{d}y .
        \]

  \item $U_{x}=\mathrm{e}^{x}$, $U_{y}=\mathrm{e}^{y}$.
        \[
          \frac{U_{x}}{U_{y}}=\mathrm{e}^{x-y} .
        \]
        $\dd{y}{x}=-\mathrm{e}^{x-y}$. \textbf{MRS} $=\mathrm{e}^{x-y}$.

  \item $f_{x}=\mathrm{e}^{x}\sin y$, $f_{y}=\mathrm{e}^{x}\cos y$;
        $\dd{x}{t}=2$, $\dd{y}{t}=2t$.
        \[
          \dd{z}{t}=2\mathrm{e}^{x}\sin y+2t\mathrm{e}^{x}\cos y
            =2\mathrm{e}^{2t}(\sin t^{2}+t\cos t^{2}) .
        \]

  \item $f_{x}=2xy$, $f_{y}=x^{2}-3y^{2}$; $x_{s}=1/t$, $x_{t}=-s/t^{2}$,
        $y_{s}=t$, $y_{t}=s$. Note $xy=(s/t)(st)=s^{2}$.
        \[
          \pd{z}{s}=2xy\cdot\frac{1}{t}+(x^{2}-3y^{2})\cdot t
            =\frac{2s^{2}}{t}+t\Bigl(\frac{s^{2}}{t^{2}}-3s^{2}t^{2}\Bigr)
            =\frac{2s^{2}}{t}+\frac{s^{2}}{t}-3s^{2}t^{3}
            =\frac{3s^{2}}{t}-3s^{2}t^{3} .
        \]
        \[
          \pd{z}{t}=2xy\cdot\Bigl(-\frac{s}{t^{2}}\Bigr)+(x^{2}-3y^{2})\cdot s
            =-\frac{2s^{3}}{t^{2}}+s\Bigl(\frac{s^{2}}{t^{2}}-3s^{2}t^{2}\Bigr)
            =-\frac{s^{3}}{t^{2}}-3s^{3}t^{2} .
        \]
        \textbf{Check.} $z=(s/t)^{2}(st)-(st)^{3}=s^{3}/t-s^{3}t^{3}$. Then
        $\pd{z}{s}=3s^{2}/t-3s^{2}t^{3}$ and
        $\pd{z}{t}=-s^{3}/t^{2}-3s^{3}t^{2}$. \checkmark

  \item Product rule: $f_{x}=\mathrm{e}^{x-y}+(x+y)\mathrm{e}^{x-y}
        =(1+x+y)\mathrm{e}^{x-y}$, and
        $f_{y}=\mathrm{e}^{x-y}-(x+y)\mathrm{e}^{x-y}=(1-x-y)\mathrm{e}^{x-y}$.
        \[
          \mathrm{d}z \;=\; \mathrm{e}^{x-y}\bigl((1+x+y)\,\mathrm{d}x
            +(1-x-y)\,\mathrm{d}y\bigr) .
        \]

  \item $U_{x}=\alpha/x$, $U_{y}=\beta/y$.
        \[
          \frac{U_{x}}{U_{y}}=\frac{\alpha}{\beta}\cdot\frac{y}{x} .
        \]
        $\dd{y}{x}=-(\alpha/\beta)(y/x)$.
        \textbf{MRS} $=(\alpha/\beta)(y/x)$. (Standard log-Cobb-Douglas.)

  \item $f_{x}=\tfrac12\sqrt{y/x}$, $f_{y}=\tfrac12\sqrt{x/y}$;
        $\dd{x}{t}=3t^{2}$, $\dd{y}{t}=1$.
        \[
          \dd{z}{t}=\frac{3t^{2}}{2}\sqrt{\frac{y}{x}}
            +\frac{1}{2}\sqrt{\frac{x}{y}} .
        \]
        Substitute $y/x=1/t^{2}$, so $\sqrt{y/x}=1/t$ (for $t>0$), and
        $\sqrt{x/y}=t$:
        \[
          \dd{z}{t}=\frac{3t^{2}}{2}\cdot\frac{1}{t}+\frac{t}{2}
            =\frac{3t}{2}+\frac{t}{2}=2t .
        \]
        \textbf{Check.} $z=\sqrt{t^{3}\cdot t}=\sqrt{t^{4}}=t^{2}$,
        $\dd{z}{t}=2t$. \checkmark

  \item $f_{x}=2x/(x^{2}+y^{2}+1)$, $f_{y}=2y/(x^{2}+y^{2}+1)$.
        \[
          \mathrm{d}z \;=\; \frac{2x\,\mathrm{d}x+2y\,\mathrm{d}y}{x^{2}+y^{2}+1} .
        \]

  \item $f_{x}=y\mathrm{e}^{xy}$, $f_{y}=x\mathrm{e}^{xy}$; $x_{s}=1$,
        $x_{t}=1$, $y_{s}=1$, $y_{t}=-1$.
        \[
          \pd{z}{s}=y\mathrm{e}^{xy}+x\mathrm{e}^{xy}=(x+y)\mathrm{e}^{xy}
            =2s\,\mathrm{e}^{s^{2}-t^{2}} .
        \]
        \[
          \pd{z}{t}=y\mathrm{e}^{xy}-x\mathrm{e}^{xy}=(y-x)\mathrm{e}^{xy}
            =-2t\,\mathrm{e}^{s^{2}-t^{2}} .
        \]

  \item Note $Q=(\sqrt{K}+\sqrt{L})^{2}$. So $Q_{K}=2(\sqrt{K}+\sqrt{L})
        \cdot\dfrac{1}{2\sqrt{K}}=(\sqrt{K}+\sqrt{L})/\sqrt{K}$, and
        $Q_{L}=(\sqrt{K}+\sqrt{L})/\sqrt{L}$.
        \[
          \frac{Q_{K}}{Q_{L}}=\frac{\sqrt{L}}{\sqrt{K}}=\sqrt{\frac{L}{K}} .
        \]
        $\dd{L}{K}=-\sqrt{L/K}$. \textbf{MRTS} $=\sqrt{L/K}$.

  \item $f_{x}=-y\sin(xy)$, $f_{y}=-x\sin(xy)$; $\dd{x}{t}=1$,
        $\dd{y}{t}=-1/t^{2}$.
        \[
          \dd{z}{t}=-y\sin(xy)+\bigl(-x\sin(xy)\bigr)\Bigl(-\frac{1}{t^{2}}\Bigr)
            =-\sin(xy)\Bigl(y-\frac{x}{t^{2}}\Bigr) .
        \]
        With $x=t$, $y=1/t$: $xy=1$, and $y-x/t^{2}=1/t-1/t=0$. Hence
        $\dd{z}{t}=0$. \textbf{Check.} $z=\cos(xy)=\cos 1$ is constant.
        \checkmark

  \item $U_{x}=y+1$, $U_{y}=x+1$.
        \[
          \frac{U_{x}}{U_{y}}=\frac{y+1}{x+1} .
        \]
        $\dd{y}{x}=-(y+1)/(x+1)$. \textbf{MRS} $=(y+1)/(x+1)$.

  \item $V_{P}=-Y/P^{2}$, $V_{Y}=1/P$.
        \[
          \mathrm{d}V \;=\; -\frac{Y}{P^{2}}\,\mathrm{d}P+\frac{1}{P}\,\mathrm{d}Y
            =\frac{V}{1}\Bigl(-\frac{\mathrm{d}P}{P}+\frac{\mathrm{d}Y}{Y}\Bigr)\cdot 1
        \]
        since $V=Y/P$, we have $-Y\,\mathrm{d}P/P^{2}=-V\,\mathrm{d}P/P$
        and $\mathrm{d}Y/P=V\,\mathrm{d}Y/Y$. So
        \[
          \frac{\mathrm{d}V}{V} \;=\; \frac{\mathrm{d}Y}{Y}-\frac{\mathrm{d}P}{P}
            \;=\; 0.03-0.02 \;=\; 0.01 .
        \]
        Real income rises by about $1\%$. Interpretation: nominal income
        growth less inflation gives real income growth (log-linearised).

  \item $f_{x}=2x+2y$, $f_{y}=2x$; $x_{s}=\cos t$, $x_{t}=-s\sin t$,
        $y_{s}=\sin t$, $y_{t}=s\cos t$.
        \[
          \pd{z}{s}=(2x+2y)\cos t+2x\sin t .
        \]
        Substitute $x=s\cos t$, $y=s\sin t$: $2x+2y=2s(\cos t+\sin t)$,
        so
        \[
          \pd{z}{s}=2s(\cos t+\sin t)\cos t+2s\cos t\sin t
            =2s\cos^{2}t+4s\sin t\cos t
            =2s(\cos^{2}t+2\sin t\cos t) .
        \]
        \[
          \pd{z}{t}=(2x+2y)(-s\sin t)+2x(s\cos t)
            =-2s^{2}(\cos t+\sin t)\sin t+2s^{2}\cos^{2}t .
        \]
        Expanding: $-2s^{2}(\cos t\sin t+\sin^{2}t)+2s^{2}\cos^{2}t
          =2s^{2}(\cos^{2}t-\sin^{2}t-\sin t\cos t)
          =2s^{2}\bigl(\cos 2t-\tfrac12\sin 2t\bigr)$.

  \item On $U=c$, $\mathrm{d}U=\mathrm{d}x_{1}+2\,\mathrm{d}x_{2}
        +3\,\mathrm{d}x_{3}=0$. Solve for $\mathrm{d}x_{3}$:
        \[
          \mathrm{d}x_{3}=-\tfrac13\,\mathrm{d}x_{1}-\tfrac23\,\mathrm{d}x_{2}
          \;\Longrightarrow\;
          \pd{x_{3}}{x_{1}}\bigg|_{U=c}=-\tfrac13,\qquad
          \pd{x_{3}}{x_{2}}\bigg|_{U=c}=-\tfrac23 .
        \]
        Unsigned MRSes: $1/3$ (between $x_{3}$ and $x_{1}$) and $2/3$
        (between $x_{3}$ and $x_{2}$).
\end{solutions}

\section{Additional Problems}

\begin{problems}
  \item Compute $\mathrm{d}z$ for $z=3x^{2}+2xy-5y^{2}$.
  \item Find $\dd{z}{t}$ for $z=x^{2}+2y$ with $x=3t+1$, $y=t^{2}-1$.
  \item For $U(x,y)=x^{2}+y^{2}=c$, find $\dd{y}{x}$ and the unsigned MRS.
  \item Compute $\pd{z}{s}$ and $\pd{z}{t}$ for $z=x^{2}+y^{2}$ with
        $x=s+t$, $y=s-t$.
  \item Compute $\mathrm{d}z$ for $z=x^{3}\sqrt{y}$.
  \item Find $\dd{z}{t}$ for $z=\ln(x^{2}+y^{2})$ with $x=\mathrm{e}^{t}$,
        $y=\mathrm{e}^{-t}$.
  \item For the isoquant $Q(K,L)=K^{0.3}L^{0.7}=Q_{0}$, find $\dd{L}{K}$
        and the unsigned MRTS.
  \item Compute $\mathrm{d}z$ for $z=\mathrm{e}^{xy}$.
  \item Compute $\pd{z}{s}$ and $\pd{z}{t}$ for $z=xy^{2}$ with
        $x=s^{2}+t$, $y=s-t^{2}$.
  \item For $U(x,y)=\sqrt{xy}=c$, find $\dd{y}{x}$ and the unsigned MRS.
  \item Find $\dd{z}{t}$ for $z=\mathrm{e}^{x+y}$ with $x=\sin t$,
        $y=\cos t$.
  \item Compute $\mathrm{d}z$ for $z=\dfrac{x}{x+y}$.
  \item For $U(x,y)=x+\ln y=c$, find $\dd{y}{x}$ and the unsigned MRS.
  \item Find $\dd{z}{t}$ for $z=x^{2}y-y^{3}$ with $x=t^{2}$, $y=2t$.
  \item Use the total differential of $W(P,Q)=PQ$ (total revenue) to
        approximate $\mathrm{d}W$ when $\mathrm{d}P=0.01P$ and
        $\mathrm{d}Q=-0.005Q$. Interpret.
\end{problems}

In problems 16--18, approximate the change $\Delta f$ by the total
differential $\mathrm{d}f$ at the point $P$, then compute the exact change
and compare.

\begin{problems}[resume]
  \item $Q(x,y,z)=4\sqrt{x^{2}+yz}$ at $P=(3,2,8)$, with
        $\mathrm{d}x=0.1$, $\mathrm{d}y=-0.2$, $\mathrm{d}z=0.3$.
  \item $Y(x,y,z)=2\sqrt{xy+z^{2}}$ at $P=(2,8,3)$, with
        $\mathrm{d}x=0.05$, $\mathrm{d}y=0.1$, $\mathrm{d}z=-0.05$.
  \item $f(x,y)=\ln(x^{2}y)-xy$ at $P=(1,2)$, with
        $\mathrm{d}x=\mathrm{d}y=0.1$; explain why the approximation is
        poorer here.
\end{problems}

\section{Additional Solutions}

\begin{solutions}
  \item $f_{x}=6x+2y$, $f_{y}=2x-10y$.
        \[
          \mathrm{d}z \;=\; (6x+2y)\,\mathrm{d}x+(2x-10y)\,\mathrm{d}y .
        \]

  \item $f_{x}=2x$, $f_{y}=2$; $\dd{x}{t}=3$, $\dd{y}{t}=2t$.
        \[
          \dd{z}{t} \;=\; 2x\cdot 3+2\cdot 2t \;=\; 6(3t+1)+4t
            \;=\; 22t+6 .
        \]
        \textbf{Check.} $z=(3t+1)^{2}+2(t^{2}-1)=9t^{2}+6t+1+2t^{2}-2
        =11t^{2}+6t-1$, $\dd{z}{t}=22t+6$. \checkmark

  \item $U_{x}=2x$, $U_{y}=2y$. Signed slope: $\dd{y}{x}=-x/y$.
        \textbf{MRS} $=x/y$ (unsigned).

  \item $f_{x}=2x$, $f_{y}=2y$; $x_{s}=1$, $x_{t}=1$, $y_{s}=1$,
        $y_{t}=-1$.
        \[
          \pd{z}{s}=2x+2y=2(s+t)+2(s-t)=4s .
        \]
        \[
          \pd{z}{t}=2x-2y=2(s+t)-2(s-t)=4t .
        \]
        \textbf{Check.} $z=(s+t)^{2}+(s-t)^{2}=2(s^{2}+t^{2})$,
        $\pd{z}{s}=4s$, $\pd{z}{t}=4t$. \checkmark

  \item $f_{x}=3x^{2}\sqrt{y}$, $f_{y}=\dfrac{x^{3}}{2\sqrt{y}}$.
        \[
          \mathrm{d}z \;=\; 3x^{2}\sqrt{y}\,\mathrm{d}x
            +\frac{x^{3}}{2\sqrt{y}}\,\mathrm{d}y .
        \]

  \item $f_{x}=2x/(x^{2}+y^{2})$, $f_{y}=2y/(x^{2}+y^{2})$;
        $\dd{x}{t}=\mathrm{e}^{t}$, $\dd{y}{t}=-\mathrm{e}^{-t}$. With
        $x=\mathrm{e}^{t}$, $y=\mathrm{e}^{-t}$,
        $x^{2}+y^{2}=\mathrm{e}^{2t}+\mathrm{e}^{-2t}$.
        \[
          \dd{z}{t} \;=\;
            \frac{2\mathrm{e}^{t}\cdot\mathrm{e}^{t}
              +2\mathrm{e}^{-t}\cdot(-\mathrm{e}^{-t})}
              {\mathrm{e}^{2t}+\mathrm{e}^{-2t}}
          \;=\; \frac{2\mathrm{e}^{2t}-2\mathrm{e}^{-2t}}
              {\mathrm{e}^{2t}+\mathrm{e}^{-2t}}
          \;=\; \frac{2(\mathrm{e}^{2t}-\mathrm{e}^{-2t})}
              {\mathrm{e}^{2t}+\mathrm{e}^{-2t}}
          =2\tanh(2t) .
        \]

  \item $Q_{K}=0.3K^{-0.7}L^{0.7}$, $Q_{L}=0.7K^{0.3}L^{-0.3}$.
        \[
          \frac{Q_{K}}{Q_{L}}=\frac{0.3}{0.7}\cdot\frac{L}{K}
            =\frac{3}{7}\cdot\frac{L}{K} .
        \]
        $\dd{L}{K}=-(3/7)(L/K)$. \textbf{MRTS} $=(3/7)(L/K)$.

  \item $f_{x}=y\mathrm{e}^{xy}$, $f_{y}=x\mathrm{e}^{xy}$.
        \[
          \mathrm{d}z \;=\; \mathrm{e}^{xy}(y\,\mathrm{d}x+x\,\mathrm{d}y) .
        \]

  \item $f_{x}=y^{2}$, $f_{y}=2xy$; $x_{s}=2s$, $x_{t}=1$, $y_{s}=1$,
        $y_{t}=-2t$.
        \[
          \pd{z}{s}=y^{2}\cdot 2s+2xy\cdot 1=2sy^{2}+2xy .
        \]
        \[
          \pd{z}{t}=y^{2}\cdot 1+2xy\cdot(-2t)=y^{2}-4txy .
        \]

  \item $U_{x}=\tfrac{1}{2}\sqrt{y/x}$, $U_{y}=\tfrac{1}{2}\sqrt{x/y}$.
        \[
          \frac{U_{x}}{U_{y}}=\sqrt{\frac{y/x}{x/y}}=\frac{y}{x} .
        \]
        $\dd{y}{x}=-y/x$. \textbf{MRS} $=y/x$.

  \item $f_{x}=f_{y}=\mathrm{e}^{x+y}$; $\dd{x}{t}=\cos t$,
        $\dd{y}{t}=-\sin t$.
        \[
          \dd{z}{t}=\mathrm{e}^{x+y}(\cos t-\sin t)
            =\mathrm{e}^{\sin t+\cos t}(\cos t-\sin t) .
        \]

  \item Quotient rule with $u=x$, $v=x+y$:
        \[
          f_{x}=\frac{(x+y)-x}{(x+y)^{2}}=\frac{y}{(x+y)^{2}},\qquad
          f_{y}=\frac{-x}{(x+y)^{2}} .
        \]
        \[
          \mathrm{d}z=\frac{y\,\mathrm{d}x-x\,\mathrm{d}y}{(x+y)^{2}} .
        \]

  \item $U_{x}=1$, $U_{y}=1/y$. $U_{x}/U_{y}=y$. Signed slope
        $\dd{y}{x}=-y$. \textbf{MRS} $=y$.

  \item $f_{x}=2xy$, $f_{y}=x^{2}-3y^{2}$; $\dd{x}{t}=2t$,
        $\dd{y}{t}=2$.
        \[
          \dd{z}{t}=2xy\cdot 2t+(x^{2}-3y^{2})\cdot 2
            =4t\cdot t^{2}\cdot 2t+2(t^{4}-3\cdot 4t^{2})
            =8t^{4}+2t^{4}-24t^{2}=10t^{4}-24t^{2} .
        \]
        \textbf{Check.} $z=t^{4}\cdot 2t-8t^{3}=2t^{5}-8t^{3}$,
        $\dd{z}{t}=10t^{4}-24t^{2}$. \checkmark

  \item $W_{P}=Q$, $W_{Q}=P$.
        \[
          \mathrm{d}W=Q\,\mathrm{d}P+P\,\mathrm{d}Q
            =W\!\left(\frac{\mathrm{d}P}{P}+\frac{\mathrm{d}Q}{Q}\right) .
        \]
        So $\mathrm{d}W/W=0.01+(-0.005)=0.005$, i.e.\ revenue rises by
        about $0.5\%$. (Log-linear: revenue percent change is the sum of
        price and quantity percent changes.)

  \item At $P$: $x^{2}+yz=9+16=25$, so $Q(P)=4\cdot5=20$. Partials, by
        the chain rule with $w=x^{2}+yz$:
        \[
          Q_{x}=\frac{4\cdot2x}{2\sqrt w}=\frac{4x}{\sqrt w}=\frac{12}{5},\quad
          Q_{y}=\frac{2z}{\sqrt w}=\frac{16}{5},\quad
          Q_{z}=\frac{2y}{\sqrt w}=\frac45 .
        \]
        \[
          \mathrm{d}Q=2.4(0.1)+3.2(-0.2)+0.8(0.3)=0.24-0.64+0.24=-0.16 .
        \]
        Exact: $Q(3.1,1.8,8.3)=4\sqrt{9.61+14.94}=4\sqrt{24.55}
        \approx19.819183$, so $\Delta Q\approx-0.180817$. The differential
        gets the sign and the size about right; the error $0.021$ comes
        from the curvature of the root.
  \item At $P$: $xy+z^{2}=16+9=25$, $Y(P)=10$. With $w=xy+z^{2}$:
        \[
          Y_{x}=\frac{y}{\sqrt w}=\frac85,\quad Y_{y}=\frac{x}{\sqrt w}=\frac25,
          \quad Y_{z}=\frac{2z}{\sqrt w}=\frac65 .
        \]
        \[
          \mathrm{d}Y=1.6(0.05)+0.4(0.1)+1.2(-0.05)=0.08+0.04-0.06=0.06 .
        \]
        Exact: $Y(2.05,8.1,2.95)=2\sqrt{16.605+8.7025}=2\sqrt{25.3075}
        \approx10.061312$, so $\Delta Y\approx0.061312$: the error is about
        $0.0013$, two percent of the change.
  \item $f(P)=\ln2-2\approx-1.306853$. $f=2\ln x+\ln y-xy$, so
        $f_{x}=\dfrac2x-y=0$ and $f_{y}=\dfrac1y-x=-\tfrac12$ at $P$.
        \[
          \mathrm{d}f=0\cdot0.1-\tfrac12\cdot0.1=-0.05 .
        \]
        Exact: $f(1.1,2.1)=\ln(1.21\cdot2.1)-2.31=\ln2.541-2.31\approx
        -1.377442$, so $\Delta f\approx-0.070589$. The differential misses
        by $0.02$, about $30\%$ of the change. At $P$ the slope in the
        $x$-direction is $0$, so the linear term sees nothing of the
        $x$-step, while the curvature, with $f_{xx}=-2/x^{2}=-2$,
        $f_{xy}=-1$ and $f_{yy}=-1/y^{2}=-\tfrac14$, contributes about
        $\tfrac12\bigl(-2(0.01)+2(-1)(0.01)-\tfrac14(0.01)\bigr)\approx-0.021$:
        indeed $-0.05-0.021\approx-0.071$.
\end{solutions}

\chapter{Directional Derivatives and Vector-Valued Functions}

The partial derivatives $f_{x}$ and $f_{y}$ measure the response of $f$
when one input moves and the others stand still. Real adjustments are
rarely so tidy: a firm re-balances capital and labour together, a consumer
slides along a budget line, a hiker walks north-east rather than due
north. The \textbf{directional derivative} measures the rate of change of
$f$ per unit distance travelled in any prescribed direction, and a single
dot product with the gradient computes it. The gradient thereby earns its
geometric role: it points where $f$ climbs fastest, its length is that
fastest rate, and it stands perpendicular to the level curves, to the
indifference curves and isoquants of the previous chapters.

The second half of the chapter lets the function itself be vector-valued.
A map $\mathbf{f}:\R^{n}\to\R^{m}$ (a production plan sending inputs to
(output, cost), a coordinate change, a demand system) is differentiated
by collecting the partial derivatives of all components into the
\textbf{Jacobian matrix}, the correct ``derivative as linear
approximation'' for vector-valued functions. In this language the chain
rule becomes matrix multiplication. Determinants, invertibility and the
Implicit Function Theorem built on top of this matrix are the subject of
the next chapter.

\section{Core Rules}

\begin{rulebox}{Directional Derivative}
Let $f:\R^{n}\to\R$ be differentiable at $\mathbf{x}_{0}$ and let
$\mathbf{u}\in\R^{n}$ be a \textbf{unit} vector, $|\mathbf{u}|=1$. The
directional derivative of $f$ at $\mathbf{x}_{0}$ in the direction
$\mathbf{u}$ is
\[
  D_{\mathbf{u}}f(\mathbf{x}_{0})
  \;=\; \lim_{h\to 0}\frac{f(\mathbf{x}_{0}+h\mathbf{u})
        -f(\mathbf{x}_{0})}{h}
  \;=\; \grad f(\mathbf{x}_{0})\cdot\mathbf{u} .
\]
Two variables: $D_{\mathbf{u}}f=f_{x}u_{1}+f_{y}u_{2}$. The partials are
the special cases $\mathbf{u}=(1,0)$ and $\mathbf{u}=(0,1)$. The unit
length is what makes $D_{\mathbf{u}}f$ a rate \emph{per unit distance}.
\end{rulebox}

\begin{rulebox}{Normalising the Direction}
If a direction is given by a nonzero vector $\mathbf{v}$ that is not of
unit length, normalise first:
\[
  \mathbf{u} \;=\; \frac{\mathbf{v}}{|\mathbf{v}|}
  \;=\; \frac{\mathbf{v}}{\sqrt{v_{1}^{2}+\cdots+v_{n}^{2}}} .
\]
``The direction from $P$ to $Q$'' means $\mathbf{v}=Q-P$, then normalise.
Forgetting this step is a frequent error in directional-derivative
problems.
\end{rulebox}

\begin{rulebox}{Steepest Ascent and Descent}
Since $D_{\mathbf{u}}f=\grad f\cdot\mathbf{u}=|\grad f|\cos\theta$, where
$\theta$ is the angle between $\mathbf{u}$ and $\grad f$, at any point
with $\grad f(\mathbf{x}_{0})\neq\mathbf{0}$:
\begin{itemize}[leftmargin=1.6em,itemsep=0.3ex]
  \item $D_{\mathbf{u}}f$ is \emph{maximal} for
        $\mathbf{u}=\grad f/|\grad f|$, with maximum value $|\grad f|$
        (steepest ascent);
  \item $D_{\mathbf{u}}f$ is \emph{minimal} for
        $\mathbf{u}=-\grad f/|\grad f|$, with value $-|\grad f|$
        (steepest descent);
  \item $D_{\mathbf{u}}f=0$ exactly when $\mathbf{u}\perp\grad f$.
\end{itemize}
For output $Q(K,L)$, the gradient $(Q_{K},Q_{L})$ is the input adjustment
that raises output fastest per unit of adjustment, and $|\grad Q|$ is that
maximal marginal gain.

\textbf{A fixed amount of adjustment.} If the change vector
$(\mathrm{d}K,\mathrm{d}L)$ must have a given length $I$, output rises
fastest when it points along the gradient:
\[
  \mathrm{d}K:\mathrm{d}L=Q_{K}:Q_{L},\qquad
  (\mathrm{d}K,\mathrm{d}L)=\frac{I}{|\grad Q|}\,(Q_{K},Q_{L}),\qquad
  \mathrm{d}Q\approx I\,|\grad Q| .
\]
A \emph{budget} $\mathrm{d}K+\mathrm{d}L=I$ with $\mathrm{d}K,\mathrm{d}L\ge0$ is
a different constraint: the linear estimate $Q_{K}\mathrm{d}K+Q_{L}\mathrm{d}L$ is then largest at a
corner, with everything spent on the input with the larger marginal
product.
\end{rulebox}

\begin{rulebox}{Gradient and Level Curves}
Along the level set $f=c$ the function does not change, so every tangent
direction $\mathbf{u}$ of the level set satisfies
$D_{\mathbf{u}}f=\grad f\cdot\mathbf{u}=0$:
\[
  \grad f \;\perp\; \text{level curve (surface) through the point.}
\]
In two variables (with $f_{y}\neq0$) the tangent direction is proportional
to $(f_{y},-f_{x})$, consistent with the implicit slope
$\dd{y}{x}=-f_{x}/f_{y}$ of the previous chapters. Economically: the
marginal-utility vector $(U_{x},U_{y})$ is normal to the indifference
curve, and $(Q_{K},Q_{L})$ is normal to the isoquant.
\end{rulebox}

\begin{rulebox}{Tangent lines: to a level curve, and to the graph}
\textbf{Tangent to the level curve} $f(x,y)=c$ through
$\mathbf{p}=(x_{0},y_{0})$, where $\grad f(\mathbf{p})\neq\mathbf{0}$. The
gradient is a normal vector, so the tangent line is
\[
  \grad f(\mathbf{p})\cdot(\mathbf{x}-\mathbf{p})=0,
  \qquad\text{i.e.}\qquad
  f_{x}(\mathbf{p})(x-x_{0})+f_{y}(\mathbf{p})(y-y_{0})=0
\]
(normal form), or $\mathbf{x}=\mathbf{p}+t\mathbf{h}$ with
$\mathbf{h}\perp\grad f(\mathbf{p})$, for instance
$\mathbf{h}=(f_{y},-f_{x})$ (parametric form). Both describe the same
line in the plane.

\textbf{Tangent to the graph in a direction.} The vertical plane through
$P=(x_{0},y_{0})$ along the unit vector $\mathbf{u}$ cuts the surface
$z=f(x,y)$ in the curve $t\mapsto f(P+t\mathbf{u})$, whose slope at $t=0$
is $D_{\mathbf{u}}f(P)$. That is the geometric meaning of the directional
derivative: the slope of the surface in the direction $\mathbf{u}$. The
tangent line to the surface in that direction is a line in $\R^{3}$:
\[
  (x,y,z)=\bigl(x_{0},\,y_{0},\,f(P)\bigr)
          +t\,\bigl(u_{1},\,u_{2},\,D_{\mathbf{u}}f(P)\bigr),\qquad t\in\R .
\]
\begin{itemize}[leftmargin=1.6em,itemsep=0.3ex]
  \item $x$-direction ($\mathbf{u}=\mathbf{e}_{1}$):
        $\bigl(x_{0}+t,\;y_{0},\;f(P)+t\,f_{x}(P)\bigr)$;
        $y$-direction ($\mathbf{u}=\mathbf{e}_{2}$):
        $\bigl(x_{0},\;y_{0}+t,\;f(P)+t\,f_{y}(P)\bigr)$.
  \item For a direction vector $\mathbf{h}$ that is not normalised,
        $\bigl(P+t\mathbf{h},\;f(P)+t\,\grad f(P)\cdot\mathbf{h}\bigr)$ is
        the \emph{same} line, run through at a different speed. Only the
        slope, rise per unit of horizontal distance, needs the unit vector.
  \item All these lines lie in the tangent plane of \chref{derivatives}{10}.
\end{itemize}
\end{rulebox}

\begin{rulebox}{Second Directional Derivative}
For a $C^{2}$ function $f$, differentiating once more in the same unit
direction $\mathbf{u}$ gives
\[
  D^{2}_{\mathbf{u}}f
  \;=\; \mathbf{u}^{T}\Hess_{f}\,\mathbf{u}
  \;=\; \sum_{i,j} f_{x_{i}x_{j}}\,u_{i}u_{j} ;
  \qquad\text{two variables:}\quad
  D^{2}_{\mathbf{u}}f
  \;=\; f_{xx}u_{1}^{2}+2f_{xy}u_{1}u_{2}+f_{yy}u_{2}^{2} .
\]
It is the second derivative at $t=0$ of the profile
$t\mapsto f(\mathbf{x}_{0}+t\mathbf{u})$: the curvature of $f$ along the
line through $\mathbf{x}_{0}$ in direction $\mathbf{u}$. Definiteness of
the Hessian (\chref{derivatives}{11}) is exactly the statement that
$D^{2}_{\mathbf{u}}f$ has one sign for \emph{every} direction
$\mathbf{u}$.
\end{rulebox}

\begin{rulebox}{Curvature in a direction: computing it, and what it proves}
\textbf{The profile method.} Without the Hessian: set
$\varphi(t)=f(\mathbf{x}_{0}+t\mathbf{u})$, a function of one variable,
and differentiate twice. Then $\varphi'(0)=D_{\mathbf{u}}f(\mathbf{x}_{0})$
and $\varphi''(0)=D^{2}_{\mathbf{u}}f(\mathbf{x}_{0})$. Doing it both ways
is a built-in check.

\textbf{Non-unit directions.} For any $\mathbf{h}\neq\mathbf{0}$, with
$\mathbf{u}=\mathbf{h}/|\mathbf{h}|$,
\[
  \mathbf{h}^{T}H\mathbf{h}=|\mathbf{h}|^{2}\,D^{2}_{\mathbf{u}}f,
  \qquad
  \grad f\cdot\mathbf{h}=|\mathbf{h}|\,D_{\mathbf{u}}f .
\]
For the \emph{sign} of the curvature, normalising is unnecessary; for the
value of $D_{\mathbf{u}}f$ or $D^{2}_{\mathbf{u}}f$, divide by $|\mathbf{h}|$
or $|\mathbf{h}|^{2}$.

\textbf{Directional verdicts.} At $\mathbf{x}_{0}$, $f$ is
\emph{convex in the direction $\mathbf{u}$} if $D^{2}_{\mathbf{u}}f>0$,
\emph{concave in that direction} if $D^{2}_{\mathbf{u}}f<0$; the profile
$\varphi$ bends up or down there. A value $0$ gives no verdict.

\textbf{Refuting global convexity.} A convex $f$ has a PSD Hessian at
\emph{every} point, so $D^{2}_{\mathbf{u}}f\ge0$ everywhere in every
direction; a concave $f$ has $D^{2}_{\mathbf{u}}f\le0$ everywhere. Hence one
point and direction with $D^{2}_{\mathbf{u}}f<0$ shows that $f$ is not
convex, and one with $D^{2}_{\mathbf{v}}f>0$ (for example at a strict local
minimum with PD Hessian) shows that $f$ is not concave. Both together prove
that $f$ is \textbf{neither}.

\textbf{Where this goes next.} The box ``Curvature at a point and in a
direction'' in \chref{optimization}{3} turns these directional verdicts into
statements about convexity near a point and on a whole domain, including
the case of a semidefinite, singular Hessian where no local conclusion is
possible.
\end{rulebox}

\begin{rulebox}{Vector-Valued Functions and the Jacobian Matrix}
For $\mathbf{f}:\R^{n}\to\R^{m}$ with components
$\mathbf{f}=(f_{1},\ldots,f_{m})$, the \textbf{Jacobian matrix} at
$\mathbf{x}$ is the $m\times n$ matrix
\[
  J_{\mathbf{f}}(\mathbf{x}) \;=\;
  \begin{pmatrix}
    \partial f_{1}/\partial x_{1} & \cdots & \partial f_{1}/\partial x_{n}\\
    \vdots & \ddots & \vdots\\
    \partial f_{m}/\partial x_{1} & \cdots & \partial f_{m}/\partial x_{n}
  \end{pmatrix}:
  \quad\text{row $i$ is the gradient of $f_{i}$.}
\]
It is the best linear approximation of $\mathbf{f}$ near $\mathbf{x}_{0}$:
\[
  \mathbf{f}(\mathbf{x}_{0}+\mathbf{h})
  \;\approx\; \mathbf{f}(\mathbf{x}_{0})
  +J_{\mathbf{f}}(\mathbf{x}_{0})\,\mathbf{h} ,
\]
the matrix version of $f(x+h)\approx f(x)+f'(x)h$. Special cases: $m=1$
gives the gradient as a $1\times n$ row; $n=1$ gives the velocity
(tangent) column vector $\mathbf{r}'(t)$ of a curve.
\end{rulebox}

\begin{rulebox}{Differentiability of vector-valued maps}
$\mathbf{f}\colon D\to\R^{m}$, $D\subseteq\R^{n}$ open, is
\textbf{differentiable} at $\mathbf{x}_{0}$ if there is an $m\times n$
matrix $J$ with
\[
  \mathbf{f}(\mathbf{x}_{0}+\mathbf{h})=\mathbf{f}(\mathbf{x}_{0})
  +J\mathbf{h}+\mathbf{r}(\mathbf{h}),
  \qquad
  \frac{\|\mathbf{r}(\mathbf{h})\|}{\|\mathbf{h}\|}\to0\ \text{ as }
  \mathbf{h}\to\mathbf{0} .
\]
The ``$\approx$'' of the Jacobian box is made exact by the remainder.
\begin{itemize}[leftmargin=1.6em,itemsep=0.3ex]
  \item The matrix $J$ is then unique and equals the Jacobian
        $J_{\mathbf{f}}(\mathbf{x}_{0})$ of partial derivatives. It is
        called the \textbf{derivative} (or total derivative,
        \de{Ableitung, Jacobi-Matrix}) of $\mathbf{f}$ at $\mathbf{x}_{0}$.
  \item $\mathbf{f}$ is differentiable exactly when every component
        $f_{i}$ is: the remainder vector tends to $0$ faster than
        $\|\mathbf{h}\|$ exactly when each of its entries does.
  \item Existence of all partials is \emph{necessary} but not sufficient
        (the function $xy/(x^{2}+y^{2})$ of \chref{derivatives}{10} is a single component that fails).
        Continuous partials on an open set are sufficient: a $C^{1}$ map is
        differentiable.
  \item Geometry: for $n=m=1$ the linear part is the tangent line; for
        $n=2$, $m=1$ it is the tangent plane; for a curve ($n=1$) it is the
        tangent vector $\mathbf{r}'(t)$.
\end{itemize}
\end{rulebox}

\begin{rulebox}{The shape of a Jacobian, and which compositions exist}
$J_{\mathbf{f}}$ has one \textbf{row per output} and one \textbf{column per
input}: (number of outputs)$\times$(number of inputs).
\begin{itemize}[leftmargin=1.6em,itemsep=0.3ex]
  \item A curve $\R\to\R^{m}$ gives an $m\times1$ column; a scalar function
        $\R^{n}\to\R$ a $1\times n$ row; $\R^{3}\to\R^{2}$ a $2\times3$
        (wide) matrix; $\R^{2}\to\R^{4}$ a $4\times2$ (tall) one.
  \item \textbf{Linear and affine maps} have constant Jacobians:
        $J(A\mathbf{x}+\mathbf{b})=A$ for an $m\times n$ matrix $A$;
        $J(c\,\mathbf{x}-d\,\mathbf{a})=c\,I_{n}$; the identity map has
        $J=I_{n}$; $J(\mathbf{a}^{T}\mathbf{x})=\mathbf{a}^{T}$.
  \item \textbf{Patterns in $n$ dimensions.} If component $k$ depends on
        $x_{k}$ alone, $f_{k}=g_{k}(x_{k})$, the Jacobian is diagonal,
        $\operatorname{diag}\bigl(g_{1}'(x_{1}),\ldots,g_{n}'(x_{n})\bigr)$.
\end{itemize}
\textbf{Which composition exists.} $\mathbf{g}\circ\mathbf{f}$ means
``first $\mathbf{f}$, then $\mathbf{g}$''; it is defined when the number of
outputs of $\mathbf{f}$ equals the number of inputs of $\mathbf{g}$. With
$\mathbf{f}\colon\R^{n}\to\R^{m}$ and $\mathbf{g}\colon\R^{m}\to\R^{p}$ the
chain rule multiplies $(p\times m)(m\times n)$; the inner sizes must
match, which is the same condition. Check it before computing anything:
often only one of $\mathbf{g}\circ\mathbf{f}$ and $\mathbf{f}\circ\mathbf{g}$
makes sense.
\end{rulebox}

\begin{rulebox}{Chain Rule in Matrix Form}
If $\mathbf{f}:\R^{n}\to\R^{m}$ and $\mathbf{g}:\R^{m}\to\R^{p}$ are
differentiable, then $\mathbf{h}=\mathbf{g}\circ\mathbf{f}$ is
differentiable and
\[
  J_{\mathbf{h}}(\mathbf{x})
  \;=\; J_{\mathbf{g}}\bigl(\mathbf{f}(\mathbf{x})\bigr)\,
        J_{\mathbf{f}}(\mathbf{x}),
  \qquad
  (p\times n) \;=\; (p\times m)(m\times n) .
\]
Evaluate the \emph{outer} Jacobian at the \emph{inner value}
$\mathbf{f}(\mathbf{x})$; the order of the factors is fixed by the matrix
shapes. Every chain rule of \chref{derivatives}{12} is the case $p=1$: for
$z=f(\mathbf{r}(t))$,
\[
  \dd{z}{t} \;=\; \grad f\bigl(\mathbf{r}(t)\bigr)\cdot\mathbf{r}'(t),
\]
a $1\times m$ row times an $m\times 1$ column.
\end{rulebox}

\section{Worked Example}

\begin{worked}{For $f(x,y)=x^{2}+3xy-y^{2}$ at $P=(1,2)$: the derivative in the direction $(3,4)$, the steepest ascent, a direction of no change; then the Jacobian of $F(x,y)=(x^{2}-y,\;xy)$}
\wstep{Step 1. The gradient first --- every question below is an inner product with it.}
\[
  f_{x}=2x+3y,\qquad f_{y}=3x-2y,
\]
\[
  \grad f(1,2)=\bigl(2+6,\;3-4\bigr)=(8,-1).
\]

\wstep{Step 2. Normalise the direction, which the formula requires.}
\[
  \|(3,4)\|=5,
  \qquad
  \mathbf u=\Bigl(\tfrac35,\tfrac45\Bigr).
\]
\[
  D_{\mathbf u}f(P)=\grad f\cdot\mathbf u=8\cdot\tfrac35+(-1)\cdot\tfrac45
  =\frac{24-4}{5}=\boxed{4}.
\]
\wtrap{Without normalising you would get $\grad f\cdot(3,4)=20$, five times
too large. The directional derivative is a rate \emph{per unit distance
travelled}, so doubling the length of the direction vector must not double
the answer. If the problem hands you a non-unit vector --- and it usually
does --- divide by its length first.}

\wstep{Step 3. Steepest ascent and descent.}
The steepest increase is along the gradient itself:
\[
  \text{direction }\frac{(8,-1)}{\sqrt{65}},
  \qquad
  \text{rate }\|\grad f\|=\sqrt{65}\approx8.06 .
\]
The steepest \emph{decrease} is the opposite direction, at rate
$-\sqrt{65}$. Every other direction lies strictly between, since
\[
  D_{\mathbf u}f=\grad f\cdot\mathbf u=\|\grad f\|\cos\theta ,
\]
and $\cos\theta\in[-1,1]$. Our $4$ sits comfortably inside
$[-8.06,\,8.06]$. \checkmark
\wtrap{The value $8.06$ is the \emph{rate}; the direction is the normalised
gradient $(0.992,-0.124)$, not $(8,-1)$. Reporting the unnormalised gradient
as ``the direction'' conflates a direction with a magnitude.}

\wstep{Step 4. A direction of no change --- and what it is tangent to.}
$D_{\mathbf u}f=0$ requires $\mathbf u\perp\grad f$. Rotating $(8,-1)$ by a
right angle gives $(1,8)$, so
\[
  \mathbf u=\frac{(1,8)}{\sqrt{65}}
  \quad\Longrightarrow\quad
  D_{\mathbf u}f=\frac{8\cdot1+(-1)\cdot8}{\sqrt{65}}=0 . \checkmark
\]
Moving in that direction keeps $f$ constant to first order, so $(1,8)$ is
tangent to the level curve of $f$ through $P$ --- which is the geometric
content of ``the gradient is perpendicular to the level curve''.

\wstep{Step 5. A vector-valued function needs a matrix of derivatives.}
$F(x,y)=(x^{2}-y,\ xy)$ has two components, each with two partials. Collect
them with \textbf{one row per component function}:
\[
  J_{F}(x,y)=
  \begin{pmatrix}
    \pd{F_{1}}{x} & \pd{F_{1}}{y}\\[0.6ex]
    \pd{F_{2}}{x} & \pd{F_{2}}{y}
  \end{pmatrix}
  =\begin{pmatrix}2x&-1\\ y&x\end{pmatrix}.
\]
\wtrap{Rows are components, columns are variables. Transposing this is the
standard error, and it matters as soon as the Jacobian is multiplied by
anything. A check: for a scalar function ($m=1$) the Jacobian is a single
\emph{row}, and the gradient is conventionally a column --- so
$J=\grad^{\!\top}$.}

\wstep{Step 6. Evaluate, and read the determinant.}
\[
  J_{F}(1,2)=\begin{pmatrix}2&-1\\2&1\end{pmatrix},
  \qquad
  \det J_{F}(1,2)=2+2=4\neq0 .
\]
A non-zero Jacobian determinant says the map is locally invertible near
$(1,2)$: the system
\[
  x^{2}-y=u,\qquad xy=v
\]
can be solved for $(x,y)$ in terms of $(u,v)$ near that point. That is the
Implicit Function Theorem of \chref{derivatives}{14} in its multivariate form, and
$\det J$ is the condition it turns on.

\wstep{Step 7. Why the Jacobian is the right generalisation.}
\begin{center}\footnotesize
\begin{tabular}{lll}
  \hline
  $f:\R\to\R$ & derivative $f'$ & a number\\
  $f:\R^{n}\to\R$ & gradient $\grad f$ & a vector\\
  $F:\R^{n}\to\R^{m}$ & Jacobian $J_{F}$ & an $m\times n$ matrix\\
  \hline
\end{tabular}
\end{center}
In each row the object is the best \emph{linear} approximation at a point,
and the chain rule is matrix multiplication:
$J_{G\circ F}=J_{G}\cdot J_{F}$, with the order fixed by the order of the
composition. The one-variable chain rule $g'(f(x))f'(x)$ is the
$1\times1$ case.
\end{worked}

\section{Practice Problems}

\begin{problems}
  \item Find $D_{\mathbf{u}}f(1,1)$ for $f(x,y)=x^{2}y$ in the direction
        $\mathbf{u}=(1/\sqrt{2},\,1/\sqrt{2})$.
  \item For $f(x,y)=\ln(x^{2}+y^{2})$ at $(1,1)$, find the directional
        derivative toward $\mathbf{v}=(2,1)$.
  \item For $f(x,y)=x^{3}-3xy^{2}$ at $(2,1)$, find the directional
        derivative in the direction of the vector from $(2,1)$ to
        $(5,5)$.
  \item For $f(x,y,z)=xyz$ at $(1,2,3)$, find $D_{\mathbf{u}}f$ where
        $\mathbf{u}=\tfrac{1}{3}(1,2,2)$.
  \item For $f(x,y)=\mathrm{e}^{x+y}$ at $(0,0)$, find $D_{\mathbf{u}}f$
        where $\mathbf{u}$ makes a $60^{\circ}$ angle with the positive
        $x$-axis.
  \item For $f(x,y)=\sqrt{x^{2}+y^{2}}$ at $(3,4)$, find the directional
        derivative in the direction $(1,0)$.
  \item A consumer with utility $U(x,y)=xy$ currently consumes
        $(x,y)=(2,3)$. Her basket shifts in the direction
        $\mathbf{v}=(3,-4)$, more $x$, less $y$. Find the rate of
        change of utility per unit distance moved and interpret the sign.
  \item For $f(x,y)=x^{2}+xy+y^{2}$ at $(1,2)$, find the direction of
        steepest ascent and the maximum rate of increase.
  \item For $f(x,y)=x^{2}+2y^{2}$ at $(1,1)$, find the direction in which
        $f$ decreases fastest and the corresponding rate.
  \item Find all unit vectors $\mathbf{u}$ along which
        $f(x,y)=x^{2}-y^{2}$ has zero rate of change at $(2,1)$.
  \item For $f(x,y,z)=x^{2}+y^{2}-z^{2}$ at $(1,2,2)$, determine the
        maximum rate of increase and the direction in which it is
        attained.
  \item A hiker stands at $(1,1)$ on the hill $z=4-x^{2}-y^{2}$. In which
        direction should she walk to descend most quickly, and how steep
        is that descent?
  \item Output is $Q(K,L)=K^{1/2}L^{1/2}$ with current input mix
        $(K,L)=(4,9)$. Find the rate of change of output in the balanced
        expansion direction $\mathbf{u}=(1,1)/\sqrt{2}$; then find the
        input direction that raises output fastest and the corresponding
        maximal rate.
  \item Show that at every point of a level curve $f(x,y)=c$ the gradient
        $\grad f$ is orthogonal to the tangent direction of the curve.
        Conclude that the marginal-utility vector $(U_{x},U_{y})$ is
        normal to the indifference curve.
  \item For $f(x,y)=x^{2}y$ at $(1,2)$ with $\mathbf{u}=(3/5,\,4/5)$,
        compute both $D_{\mathbf{u}}f$ and the second directional
        derivative
        $D^{2}_{\mathbf{u}}f=\mathbf{u}^{T}\Hess_{f}\,\mathbf{u}$.
  \item For $f(x,y)=x^{3}+y^{3}-3xy$: show that $(1,1)$ is a critical
        point, then compute $D^{2}_{\mathbf{u}}f(1,1)$ for
        $\mathbf{u}=(1,1)/\sqrt{2}$ and for
        $\mathbf{v}=(1,-1)/\sqrt{2}$. What do the signs suggest?
  \item For $f(x,y)=\mathrm{e}^{x}\cos y$ at $(0,0)$, find all unit
        vectors $\mathbf{u}$ with $D^{2}_{\mathbf{u}}f=0$.
  \item Compute the Jacobian of $\mathbf{f}(x,y)=(x+y,\;x-y)$ and its
        determinant. What is special about the Jacobian of a linear map?
  \item Compute the Jacobian of
        $\mathbf{f}(x,y)=(x^{2}y,\;x+y^{2},\;xy)$, state its shape, and
        evaluate it at $(1,2)$.
  \item Compute the Jacobian of
        $\mathbf{f}(x,y,z)=(x+y,\;y+z,\;x+z)$ and its determinant.
  \item For $\mathbf{f}(x,y)=(x^{2}-y,\;xy)$: compute $J_{\mathbf{f}}$,
        then use the linear approximation
        $\mathbf{f}(\mathbf{x}_{0}+\mathbf{h})\approx
        \mathbf{f}(\mathbf{x}_{0})+J_{\mathbf{f}}(\mathbf{x}_{0})\mathbf{h}$
        at $\mathbf{x}_{0}=(1,1)$ to estimate $\mathbf{f}(1.1,\,0.9)$,
        and compare with the exact value.
  \item Demand and supply are $D(p,q)=100-p+0.5q$ and
        $S(p,q)=-10+2p+q$, and the two equilibrium conditions are
        $G_{1}=D-S=0$ and $G_{2}=p+q-20=0$. Build the Jacobian matrix
        $\partial(G_{1},G_{2})/\partial(p,q)$ and compute its
        determinant.
  \item Let $\mathbf{g}(x,y)=(x+y,\;x-y)$ and
        $\mathbf{F}(u,v)=(u^{2},\;uv)$. Compute the Jacobian of
        $\mathbf{F}\circ\mathbf{g}$ by the matrix chain rule
        $J_{\mathbf{F}\circ\mathbf{g}}
        =J_{\mathbf{F}}(\mathbf{g}(x,y))\,J_{\mathbf{g}}(x,y)$, and
        verify by composing first.
  \item For $f(x,y)=xy^{2}$ and the curve
        $\mathbf{r}(t)=(\mathrm{e}^{t},\,t)$, compute $\dd{z}{t}$ for
        $z=f(\mathbf{r}(t))$ as the matrix product
        $J_{f}(\mathbf{r}(t))\,J_{\mathbf{r}}(t)$ (a $1\times2$ row
        times a $2\times1$ column) and verify directly.
  \item A firm converts inputs into (output, cost) via
        $\mathbf{g}(K,L)=\bigl(K^{1/2}L^{1/2},\;2K+L\bigr)$ and earns
        profit $\pi(Q,C)=10Q-C$. Use the matrix chain rule to find the
        marginal profits $(\pi_{K},\pi_{L})$ at $(K,L)=(1,1)$, and verify
        by substituting first.
\end{problems}

\section{Solutions}

\begin{solutions}
  \item $\grad f=(2xy,\,x^{2})$. At $(1,1)$: $\grad f=(2,1)$, and
        $\mathbf{u}$ is already a unit vector.
        \[
          D_{\mathbf{u}}f=(2)\frac{1}{\sqrt{2}}+(1)\frac{1}{\sqrt{2}}
            =\frac{3}{\sqrt{2}}=\frac{3\sqrt{2}}{2} .
        \]

  \item $\grad f=\Bigl(\dfrac{2x}{x^{2}+y^{2}},\,
        \dfrac{2y}{x^{2}+y^{2}}\Bigr)$. At $(1,1)$: $\grad f=(1,1)$.
        Normalise: $|\mathbf{v}|=\sqrt{4+1}=\sqrt{5}$, so
        $\mathbf{u}=(2,1)/\sqrt{5}$.
        \[
          D_{\mathbf{u}}f=\frac{1\cdot 2+1\cdot 1}{\sqrt{5}}
            =\frac{3}{\sqrt{5}}=\frac{3\sqrt{5}}{5} .
        \]

  \item Direction vector $\mathbf{v}=(5,5)-(2,1)=(3,4)$,
        $|\mathbf{v}|=5$, so $\mathbf{u}=(3/5,\,4/5)$.
        $\grad f=(3x^{2}-3y^{2},\,-6xy)$; at $(2,1)$:
        $\grad f=(12-3,\,-12)=(9,-12)$.
        \[
          D_{\mathbf{u}}f=9\cdot\frac{3}{5}+(-12)\cdot\frac{4}{5}
            =\frac{27}{5}-\frac{48}{5}=-\frac{21}{5} .
        \]

  \item $\grad f=(yz,\,xz,\,xy)$. At $(1,2,3)$: $\grad f=(6,3,2)$. The
        vector $\mathbf{u}=\tfrac13(1,2,2)$ is unit since
        $\sqrt{1+4+4}/3=1$.
        \[
          D_{\mathbf{u}}f=\frac{6\cdot 1+3\cdot 2+2\cdot 2}{3}
            =\frac{6+6+4}{3}=\frac{16}{3} .
        \]

  \item $\mathbf{u}=(\cos 60^{\circ},\,\sin 60^{\circ})
        =(1/2,\,\sqrt{3}/2)$, a unit vector.
        $\grad f=(\mathrm{e}^{x+y},\,\mathrm{e}^{x+y})$; at $(0,0)$:
        $\grad f=(1,1)$.
        \[
          D_{\mathbf{u}}f=\frac{1}{2}+\frac{\sqrt{3}}{2}
            =\frac{1+\sqrt{3}}{2} .
        \]

  \item $\grad f=\Bigl(\dfrac{x}{\sqrt{x^{2}+y^{2}}},\,
        \dfrac{y}{\sqrt{x^{2}+y^{2}}}\Bigr)$. At $(3,4)$:
        $\grad f=(3/5,\,4/5)$. With $\mathbf{u}=(1,0)$ (already unit):
        \[
          D_{\mathbf{u}}f=\frac{3}{5} .
        \]

  \item $\grad U=(y,\,x)$; at $(2,3)$: $\grad U=(3,2)$. Normalise:
        $|\mathbf{v}|=\sqrt{9+16}=5$, $\mathbf{u}=(3/5,\,-4/5)$.
        \[
          D_{\mathbf{u}}U=3\cdot\frac{3}{5}+2\cdot\Bigl(-\frac{4}{5}\Bigr)
            =\frac{9}{5}-\frac{8}{5}=\frac{1}{5} .
        \]
        The rate is positive: the marginal-utility gain from the extra
        $x$ slightly outweighs the loss from the reduced $y$, so utility
        rises, at $1/5$ per unit distance moved.

  \item $\grad f=(2x+y,\,x+2y)$. At $(1,2)$: $\grad f=(4,5)$,
        $|\grad f|=\sqrt{16+25}=\sqrt{41}$.
        \[
          \mathbf{u}=\frac{(4,5)}{\sqrt{41}},\qquad
          \text{maximum rate}=\sqrt{41} .
        \]

  \item $\grad f=(2x,\,4y)$. At $(1,1)$: $\grad f=(2,4)$,
        $|\grad f|=\sqrt{4+16}=2\sqrt{5}$. Steepest descent:
        \[
          \mathbf{u}=-\frac{(2,4)}{2\sqrt{5}}=-\frac{(1,2)}{\sqrt{5}},
          \qquad \text{rate}=-2\sqrt{5} .
        \]

  \item $\grad f=(2x,\,-2y)$. At $(2,1)$: $\grad f=(4,-2)$. We need
        $\mathbf{u}\perp(4,-2)$: $4u_{1}-2u_{2}=0$, i.e.
        $u_{2}=2u_{1}$. Normalising $(1,2)$:
        \[
          \mathbf{u}=\pm\frac{(1,2)}{\sqrt{5}} .
        \]
        \textbf{Check.} $(4,-2)\cdot(1,2)=4-4=0$. \checkmark

  \item $\grad f=(2x,\,2y,\,-2z)$. At $(1,2,2)$: $\grad f=(2,4,-4)$,
        $|\grad f|=\sqrt{4+16+16}=6$. Maximum rate $6$, attained in the
        direction
        \[
          \mathbf{u}=\frac{(2,4,-4)}{6}=\frac{(1,2,-2)}{3} .
        \]

  \item $\grad z=(-2x,\,-2y)$; at $(1,1)$: $\grad z=(-2,-2)$,
        $|\grad z|=2\sqrt{2}$. Steepest descent direction:
        \[
          \mathbf{u}=-\frac{(-2,-2)}{2\sqrt{2}}=\frac{(1,1)}{\sqrt{2}} ,
        \]
        i.e. north-east, directly away from the summit above the origin.
        The slope of that descent is $-|\grad z|=-2\sqrt{2}$.

  \item $Q_{K}=\tfrac12K^{-1/2}L^{1/2}$, $Q_{L}=\tfrac12K^{1/2}L^{-1/2}$.
        At $(4,9)$:
        \[
          Q_{K}=\frac{1}{2}\cdot\frac{1}{2}\cdot 3=\frac{3}{4},\qquad
          Q_{L}=\frac{1}{2}\cdot 2\cdot\frac{1}{3}=\frac{1}{3},
        \]
        so $\grad Q=(3/4,\,1/3)$. Balanced direction
        $\mathbf{u}=(1,1)/\sqrt{2}$:
        \[
          D_{\mathbf{u}}Q=\frac{3/4+1/3}{\sqrt{2}}
            =\frac{13/12}{\sqrt{2}}=\frac{13\sqrt{2}}{24} .
        \]
        Fastest increase along $\grad Q/|\grad Q|$ with
        \[
          |\grad Q|=\sqrt{\frac{9}{16}+\frac{1}{9}}
            =\sqrt{\frac{81+16}{144}}=\frac{\sqrt{97}}{12},\qquad
          \mathbf{u}^{*}=\frac{12}{\sqrt{97}}\Bigl(\frac34,\frac13\Bigr)
            =\frac{(9,4)}{\sqrt{97}} .
        \]
        The maximal rate is $\sqrt{97}/12$ (larger than $13\sqrt{2}/24$,
        as it must be).

  \item Let $\mathbf{r}(t)=(x(t),y(t))$ be a differentiable curve lying
        in the level set, $f(\mathbf{r}(t))=c$ for all $t$.
        Differentiate with the chain rule:
        \[
          0=\dd{}{t}\,c=\grad f(\mathbf{r}(t))\cdot\mathbf{r}'(t) .
        \]
        Hence the tangent vector $\mathbf{r}'(t)$ is orthogonal to
        $\grad f$ at every point of the curve. For $U(x,y)=c$ an
        indifference curve, this says exactly that
        $(U_{x},U_{y})$ is normal to it.

  \item First order: $\grad f=(2xy,\,x^{2})$; at $(1,2)$:
        $\grad f=(4,1)$.
        \[
          D_{\mathbf{u}}f=4\cdot\frac{3}{5}+1\cdot\frac{4}{5}
            =\frac{12}{5}+\frac{4}{5}=\frac{16}{5} .
        \]
        Second order: $f_{xx}=2y=4$, $f_{xy}=2x=2$, $f_{yy}=0$ at
        $(1,2)$.
        \[
          D^{2}_{\mathbf{u}}f
            =4\Bigl(\frac{3}{5}\Bigr)^{2}
             +2\cdot 2\cdot\frac{3}{5}\cdot\frac{4}{5}
             +0\cdot\Bigl(\frac{4}{5}\Bigr)^{2}
            =\frac{36}{25}+\frac{48}{25}=\frac{84}{25} .
        \]

  \item $\grad f=(3x^{2}-3y,\,3y^{2}-3x)$; at $(1,1)$:
        $\grad f=(0,0)$, so $(1,1)$ is a critical point. Second
        partials: $f_{xx}=6x=6$, $f_{yy}=6y=6$, $f_{xy}=-3$, so
        $\Hess_{f}(1,1)=\begin{pmatrix}6&-3\\-3&6\end{pmatrix}$. For
        $\mathbf{u}=(1,1)/\sqrt{2}$:
        \[
          D^{2}_{\mathbf{u}}f=6\cdot\frac12
            +2(-3)\cdot\frac12+6\cdot\frac12=3-3+3=3 .
        \]
        For $\mathbf{v}=(1,-1)/\sqrt{2}$:
        \[
          D^{2}_{\mathbf{v}}f=6\cdot\frac12
            +2(-3)\cdot\Bigl(-\frac12\Bigr)+6\cdot\frac12=3+3+3=9 .
        \]
        Both are positive, the surface curves upward in both
        directions. Indeed $D_{1}=6>0$ and
        $D_{2}=36-9=27>0$, so the Hessian is positive definite and
        $(1,1)$ is a strict local minimum (\chref{derivatives}{11}).

  \item $f_{xx}=\mathrm{e}^{x}\cos y$, $f_{xy}=-\mathrm{e}^{x}\sin y$,
        $f_{yy}=-\mathrm{e}^{x}\cos y$. At $(0,0)$: $f_{xx}=1$,
        $f_{xy}=0$, $f_{yy}=-1$. Writing
        $\mathbf{u}=(\cos\theta,\,\sin\theta)$:
        \[
          D^{2}_{\mathbf{u}}f=\cos^{2}\theta-\sin^{2}\theta=\cos 2\theta .
        \]
        This vanishes for $2\theta=\tfrac{\pi}{2}+k\pi$, i.e.
        $\theta=\tfrac{\pi}{4}+k\tfrac{\pi}{2}$:
        \[
          \mathbf{u}=\pm\frac{(1,1)}{\sqrt{2}}
          \quad\text{and}\quad
          \mathbf{u}=\pm\frac{(1,-1)}{\sqrt{2}} .
        \]

  \item
        \[
          J_{\mathbf{f}}=\begin{pmatrix}1&1\\1&-1\end{pmatrix},\qquad
          \det J_{\mathbf{f}}=1\cdot(-1)-1\cdot 1=-2 .
        \]
        The map is linear, $\mathbf{f}(\mathbf{x})=A\mathbf{x}$ with
        $A=\bigl(\begin{smallmatrix}1&1\\1&-1\end{smallmatrix}\bigr)$,
        and for a linear map $f_{i}=\sum_{j}A_{ij}x_{j}$ gives
        $\partial f_{i}/\partial x_{j}=A_{ij}$: the Jacobian is the
        constant matrix $A$ itself: a linear map is its own linear
        approximation.

  \item With $f_{1}=x^{2}y$, $f_{2}=x+y^{2}$, $f_{3}=xy$:
        \[
          J_{\mathbf{f}}=\begin{pmatrix}2xy&x^{2}\\1&2y\\y&x\end{pmatrix}
          \;\in\;\R^{3\times 2} .
        \]
        A $3\times 2$ matrix: one row per component, one column per
        variable; being non-square, it has no determinant. At $(1,2)$:
        \[
          J_{\mathbf{f}}(1,2)=\begin{pmatrix}4&1\\1&4\\2&1\end{pmatrix} .
        \]

  \item
        \[
          J_{\mathbf{f}}=\begin{pmatrix}1&1&0\\0&1&1\\1&0&1\end{pmatrix} .
        \]
        Expanding along the first row:
        \[
          \det J_{\mathbf{f}}
            =1\,(1\cdot 1-1\cdot 0)-1\,(0\cdot 1-1\cdot 1)+0
            =1+1=2 .
        \]

  \item $J_{\mathbf{f}}=\begin{pmatrix}2x&-1\\y&x\end{pmatrix}$. At
        $\mathbf{x}_{0}=(1,1)$:
        \[
          \mathbf{f}(1,1)=(0,\,1),\qquad
          J_{\mathbf{f}}(1,1)=\begin{pmatrix}2&-1\\1&1\end{pmatrix} .
        \]
        With $\mathbf{h}=(0.1,\,-0.1)$:
        \[
          J_{\mathbf{f}}(1,1)\,\mathbf{h}
            =\begin{pmatrix}2(0.1)+(-1)(-0.1)\\
              1(0.1)+1(-0.1)\end{pmatrix}
            =\begin{pmatrix}0.3\\0\end{pmatrix},
        \]
        so $\mathbf{f}(1.1,\,0.9)\approx(0,1)+(0.3,\,0)=(0.3,\,1)$.
        Exact value:
        \[
          \mathbf{f}(1.1,\,0.9)=(1.21-0.9,\;1.1\cdot 0.9)
            =(0.31,\;0.99) .
        \]
        The approximation is off by only $0.01$ in each component.

  \item Excess demand:
        \[
          G_{1}=D-S=(100-p+0.5q)-(-10+2p+q)=110-3p-0.5q ,
        \]
        and $G_{2}=p+q-20$. Hence
        \[
          \frac{\partial(G_{1},G_{2})}{\partial(p,q)}
            =\begin{pmatrix}-3&-0.5\\1&1\end{pmatrix},\qquad
          \det=-3\cdot 1-(-0.5)\cdot 1=-2.5 .
        \]
        The determinant is nonzero: exactly the condition the next
        chapter's Implicit Function Theorem needs to solve the system
        locally for $(p,q)$ when a parameter shifts.

  \item Inner and outer Jacobians:
        \[
          J_{\mathbf{g}}=\begin{pmatrix}1&1\\1&-1\end{pmatrix},\qquad
          J_{\mathbf{F}}(u,v)=\begin{pmatrix}2u&0\\v&u\end{pmatrix} .
        \]
        Evaluate the outer at the inner value
        $(u,v)=\mathbf{g}(x,y)=(x+y,\,x-y)$ and multiply:
        \[
          J_{\mathbf{F}\circ\mathbf{g}}
            =\begin{pmatrix}2(x+y)&0\\x-y&x+y\end{pmatrix}
             \begin{pmatrix}1&1\\1&-1\end{pmatrix}
            =\begin{pmatrix}2(x+y)&2(x+y)\\
              (x-y)+(x+y)&(x-y)-(x+y)\end{pmatrix}
        \]
        \[
          \phantom{J_{\mathbf{F}\circ\mathbf{g}}}
            =\begin{pmatrix}2(x+y)&2(x+y)\\2x&-2y\end{pmatrix} .
        \]
        \textbf{Check.} Composing first:
        $(\mathbf{F}\circ\mathbf{g})(x,y)
        =\bigl((x+y)^{2},\,(x+y)(x-y)\bigr)
        =\bigl((x+y)^{2},\,x^{2}-y^{2}\bigr)$, whose Jacobian is
        \[
          \begin{pmatrix}2(x+y)&2(x+y)\\2x&-2y\end{pmatrix} .
          \quad\checkmark
        \]

  \item $J_{f}=\bigl(y^{2}\;\;2xy\bigr)$ (a row) and
        $J_{\mathbf{r}}(t)=\begin{pmatrix}\mathrm{e}^{t}\\1\end{pmatrix}$
        (a column). At $\mathbf{r}(t)=(\mathrm{e}^{t},\,t)$ the row is
        $\bigl(t^{2}\;\;2t\mathrm{e}^{t}\bigr)$, so
        \[
          \dd{z}{t}
            =\bigl(t^{2}\;\;2t\mathrm{e}^{t}\bigr)
             \begin{pmatrix}\mathrm{e}^{t}\\1\end{pmatrix}
            =t^{2}\mathrm{e}^{t}+2t\mathrm{e}^{t}
            =\mathrm{e}^{t}(t^{2}+2t) .
        \]
        \textbf{Check.} $z=\mathrm{e}^{t}\,t^{2}$, so
        $\dd{z}{t}=\mathrm{e}^{t}t^{2}+2t\mathrm{e}^{t}$. \checkmark

  \item Inner Jacobian:
        \[
          J_{\mathbf{g}}(K,L)=\begin{pmatrix}
            \tfrac12K^{-1/2}L^{1/2}&\tfrac12K^{1/2}L^{-1/2}\\
            2&1\end{pmatrix},\qquad
          J_{\mathbf{g}}(1,1)=\begin{pmatrix}1/2&1/2\\2&1\end{pmatrix} .
        \]
        Outer Jacobian (a row, constant since $\pi$ is linear):
        $J_{\pi}=(10\;\;-1)$. Matrix chain rule:
        \[
          (\pi_{K},\,\pi_{L})
            =(10\;\;-1)\begin{pmatrix}1/2&1/2\\2&1\end{pmatrix}
            =\bigl(10\cdot\tfrac12-1\cdot 2,\;
              10\cdot\tfrac12-1\cdot 1\bigr)
            =(3,\;4) .
        \]
        \textbf{Check.} Substituting first,
        $\pi(K,L)=10K^{1/2}L^{1/2}-2K-L$, so
        $\pi_{K}=5K^{-1/2}L^{1/2}-2=5-2=3$ and
        $\pi_{L}=5K^{1/2}L^{-1/2}-1=5-1=4$ at $(1,1)$. \checkmark
        At the current input mix an extra unit of capital raises profit
        by about $3$, an extra unit of labour by about $4$.
\end{solutions}

\section{Additional Problems}

\begin{problems}
  \item For $f(x,y)=x^{2}-y^{2}+3xy$ at $(1,2)$, find the directional
        derivative toward $\mathbf{v}=(4,3)$.
  \item For $f(x,y,z)=x^{2}+2y^{2}+3z^{2}$ at $(1,1,1)$, find the
        directional derivative toward $\mathbf{v}=(1,-1,1)$.
  \item For $U(x,y)=\ln x+2\ln y$ at $(1,2)$, find the direction of
        steepest ascent and the maximum rate of increase.
  \item Find all unit vectors $\mathbf{u}$ with $D_{\mathbf{u}}f=0$ for
        $f(x,y)=xy$ at $(2,1)$.
  \item At a point $P$ the directional derivatives of $f$ toward $(1,0)$
        and toward $(0,1)$ are $3$ and $-4$. Find $\grad f(P)$, the
        maximum rate of increase at $P$, and $D_{\mathbf{u}}f$ toward
        $\mathbf{u}=(3/5,\,4/5)$.
  \item For $f(x,y)=\ln(x^{2}+y^{2})$ at $(1,0)$ with
        $\mathbf{u}=(0,1)$, compute $D_{\mathbf{u}}f$ and
        $D^{2}_{\mathbf{u}}f$, and interpret the two results.
  \item Compute the Jacobian of
        $\mathbf{f}(x,y)=(x^{2}-y^{2},\;3xy+y^{2})$ and evaluate it at
        $(1,1)$.
  \item For $f(x,y)=xy$ and $\mathbf{r}(t)=(t^{2},\,t^{3})$, compute
        $\dd{z}{t}$ for $z=f(\mathbf{r}(t))$ by the matrix chain rule
        and verify directly.
  \item Let $\mathbf{g}(s,t)=(st,\;s+t)$ and
        $\mathbf{F}(x,y)=(x+y,\;xy)$. Compute
        $J_{\mathbf{F}\circ\mathbf{g}}$ by the matrix chain rule and
        verify by composing first.
  \item Inputs produce two intermediate goods $x_{1}=2K+L$ and
        $x_{2}=K+3L$, and final output is $y=x_{1}x_{2}$. Use the matrix
        chain rule to find $\pd{y}{K}$ and $\pd{y}{L}$ at $(K,L)=(1,1)$,
        and verify by expanding $y$ as a polynomial in $K$ and $L$.
\end{problems}

Problems 11--14: tangent lines to the graph of $f$ in a given direction.

\begin{problems}[resume]
  \item For $f(x,y)=x^{2}y+y^{3}$ at $P=(1,2)$, give the slope of the graph
        and the tangent line to the graph in the $x$-direction, in the
        $y$-direction, and in the direction of $\mathbf{v}=(3,-4)$.
  \item For $f(x,y)=\ln(x^{2}+3y)$ at $P=(2,0)$, give the slope of the
        graph in the direction of steepest ascent and the tangent line to
        the graph in that direction.
  \item For $f(x,y)=(x+y)\,\mathrm{e}^{x-y}$ at $P=(1,1)$, give the slope in
        the direction from $P$ towards $Q=(3,2)$ and the tangent line in
        that direction, once with the unit vector and once with
        $\mathbf{h}=Q-P$. Show that both describe the same line.
  \item For $Q(K,L)=4K^{1/2}L^{1/4}$ at $(16,16)$, give the tangent line to
        the graph in the $L$-direction, read it as an output estimate for
        $L=17$, and compare with the exact value.
\end{problems}

In problems 15--17, for the point $P$: (a) find the level $c=f(P)$; (b)
solve the level curve $f=c$ explicitly; (c) compute $\grad f(P)$ and check
that it is perpendicular to the tangent of the level curve; (d) give the
tangent line to the level curve in normal and in parametric form; (e) give
$D_{\mathbf{u}}f(P)$ in the direction of the gradient.

\begin{problems}[resume]
  \item $f(x,y)=\sqrt{2x^{2}+3y+1}$, \;$P=(1,2)$.
  \item $f(x,y)=x^{1/3}y^{2/3}$ on $x,y>0$, \;$P=(1,8)$.
  \item $f(x,y)=x\,\mathrm{e}^{y}+y$, \;$P=(2,0)$; solve the level curve
        for $x$.
\end{problems}

Problems 18--20: differentiability of vector-valued maps.

\begin{problems}[resume]
  \item Explain why
        $\mathbf{F}(x,y)=\bigl(x^{2}y,\;x-3y^{2},\;xy\bigr)$ is
        differentiable at every point. Give $J_{\mathbf{F}}(1,1)$, and
        compute the remainder vector $\mathbf{r}(\mathbf{h})$ at $(1,1)$
        to show directly that $\|\mathbf{r}(\mathbf{h})\|/\|\mathbf{h}\|\to0$.
  \item For $\mathbf{F}(x,y)=\bigl(\ln(3x^{2}+y^{4}),\;xy^{2},\;|y-1|\bigr)$,
        give the domain, the set on which $\mathbf{F}$ is differentiable,
        and the Jacobian there.
  \item Answer in words: when is a map $\mathbf{f}\colon\R^{n}\to\R^{m}$
        differentiable at $\mathbf{x}_{0}$; what is its derivative called
        and what shape does it have; how is it related to the partial
        derivatives of the components; and what is the geometric meaning
        for $n=m=1$ and for $n=2$, $m=1$?
\end{problems}

Problems 21--26: Jacobians of all shapes.

\begin{problems}[resume]
  \item $\mathbf{r}(t)=\bigl(\ln(1+t^{2}),\;\mathrm{e}^{-2t},\;\cos3t,\;5\bigr)$:
        the Jacobian, its shape, and its value at $t=0$.
  \item $\mathbf{F}(x,y,z)=(xyz,\;x^{2}-z)$ at $(1,2,3)$, and
        $g(x_{1},x_{2},x_{3},x_{4})=x_{1}x_{2}+x_{3}\mathrm{e}^{x_{4}}$ at
        $(1,1,1,0)$: shapes and values of the Jacobians.
  \item $\mathbf{F}(s,t)=\bigl(s+t,\;st,\;s^{2},\;\mathrm{e}^{s-t}\bigr)$:
        the Jacobian, its shape, and its value at $(1,1)$.
  \item Jacobians in $n$ dimensions: (a) the identity map
        $\mathbf{F}(\mathbf{x})=\mathbf{x}$ on $\R^{n}$; (b) the map
        $\mathbf{F}\colon\R^{n}\to\R^{n}$ with components
        $F_{k}(\mathbf{x})=k\,\mathrm{e}^{x_{k}}$, that is
        $(\mathrm{e}^{x_{1}},2\mathrm{e}^{x_{2}},\ldots,n\,\mathrm{e}^{x_{n}})$;
        (c) the curve $\mathbf{r}\colon\R\to\R^{n}$ with components
        $r_{k}(t)=t^{k}/k$, that is $(t,t^{2}/2,\ldots,t^{n}/n)$; (d)
        $f(\mathbf{x})=\sum_{k=1}^{n}x_{k}^{k}/k$.
  \item Jacobians of linear and affine maps: (a) $\mathbf{F}(\mathbf{x})=
        A\mathbf{x}$ with $A=\begin{pmatrix}2&-1&0\\1&0&3\\0&4&-2
        \end{pmatrix}$; (b) $\mathbf{F}(\mathbf{x})=3\mathbf{x}-2\mathbf{a}$
        for a fixed $\mathbf{a}\in\R^{n}$; (c) $f(\mathbf{x})=
        \mathbf{a}^{T}\mathbf{x}$; (d) $\mathbf{G}(\mathbf{x})=B\mathbf{x}
        +\mathbf{c}$ with $B$ a $2\times3$ matrix.
  \item Let $\mathbf{F}\colon\R^{3}\to\R^{2}$, $\mathbf{G}\colon\R^{2}\to
        \R^{4}$, $H\colon\R^{4}\to\R$ and $\mathbf{K}\colon\R\to\R^{3}$.
        Which of $\mathbf{G}\circ\mathbf{F}$, $\mathbf{F}\circ\mathbf{G}$,
        $H\circ\mathbf{G}$, $\mathbf{G}\circ H$, $\mathbf{K}\circ H$,
        $H\circ\mathbf{K}$, $\mathbf{F}\circ\mathbf{K}$,
        $\mathbf{K}\circ\mathbf{F}$ and $H\circ\mathbf{G}\circ\mathbf{F}$
        are defined, and what is the shape of the Jacobian of each one that
        is?
\end{problems}

Problems 27--30: the matrix chain rule with rectangular Jacobians. Verify
each result by composing first.

\begin{problems}[resume]
  \item $\mathbf{f}(x,y,z)=(x+yz,\;xz)$ and
        $\mathbf{g}(u,v)=(u,\;v,\;uv,\;u-v)$: the Jacobian of
        $\mathbf{g}\circ\mathbf{f}$ at $(1,2,1)$.
  \item $g(x,y,z)=x+2y+z^{2}$ and $\mathbf{f}(s)=(s,\;s^{2},\;\mathrm{e}^{s})$:
        say why both $\mathbf{f}\circ g$ and $g\circ\mathbf{f}$ exist; give
        the Jacobian of $\mathbf{f}\circ g$ at $(1,0,0)$ and the derivative
        of $g\circ\mathbf{f}$.
  \item $\mathbf{f}(x,y)=(x,\;y,\;xy)$ and $\mathbf{g}(u,v,w)=(u+w,\;vw)$:
        the Jacobian of $\mathbf{g}\circ\mathbf{f}$ at $(1,2)$ and of
        $\mathbf{f}\circ\mathbf{g}$ at $(1,1,1)$.
  \item $\mathbf{r}(t)=(\cos t,\;\sin t,\;t)$ and $F(x,y,z)=x^{2}+y^{2}+z^{2}$:
        $\dd{}{t}F(\mathbf{r}(t))$ by the chain rule, simplified.
\end{problems}

In problems 31--33, compute $D_{\mathbf{u}}f(P)$ and $D^{2}_{\mathbf{u}}f(P)$
for the direction of $\mathbf{h}$, both with the gradient and Hessian and
with the profile $\varphi(t)=f(P+t\mathbf{h})$, and say whether $f$ is
convex or concave in that direction at $P$.

\begin{problems}[resume]
  \item $f(x,y,z)=x\ln y+yz^{2}$, \;$P=(2,1,1)$, \;$\mathbf{h}=(1,2,-2)$.
  \item $f(x,y,z)=x\,\mathrm{e}^{yz}$, \;$P=(1,0,2)$, \;$\mathbf{h}=(2,1,2)$.
  \item $f(x,y,z)=xyz+\ln(xz)$, \;$P=(1,1,1)$, \;$\mathbf{h}=(1,-1,1)$.
  \item Redo practice problem 15 ($f=x^{2}y$ at $(1,2)$) with the
        \emph{unnormalised} $\mathbf{h}=(3,4)$: compute
        $\grad f\cdot\mathbf{h}$, $\mathbf{h}^{T}H\mathbf{h}$, and the
        profile derivatives $\varphi'(0)$, $\varphi''(0)$. Explain the
        factors $5$ and $25$.
\end{problems}

In problems 35--37, give the directional verdicts asked for, then decide
whether $f$ is convex, concave or neither on its whole domain.

\begin{problems}[resume]
  \item $f(x,y)=x^{3}-3x+y^{2}$: at $Q=(-1,1)$ in the direction of
        $\mathbf{h}=(3,4)$, compute $D_{\mathbf{u}}f$ and
        $D^{2}_{\mathbf{u}}f$; then use the point $(1,0)$.
  \item $f(x,y)=x^{4}-4xy+y^{2}$: classify the critical points, then give
        $D^{2}_{\mathbf{u}}f(0,0)$ for $\mathbf{u}=(1,1)/\sqrt2$ and for
        $\mathbf{u}=(0,1)$.
  \item $f(x,y)=x\ln y$ on $y>0$, at $P=(1,1)$: $D_{\mathbf{u}}f$ and
        $D^{2}_{\mathbf{u}}f$ for $\mathbf{u}=(1,0)$, $(0,1)$ and
        $(1,1)/\sqrt2$.
\end{problems}

Problems 38--40: statements to check. Justify every answer.

\begin{problems}[resume]
  \item True or false? (a) Walking along the contour line through $P$,
        the rate of change of $f$ is $0$. (b) The direction of steepest
        descent is $-\grad f(P)$. (c) If $\grad f(P)=\mathbf{0}$, then $P$
        is a peak or a valley. (d) $|\grad f(P)|$ is the slope of the
        steepest path through $P$. (e) The gradient at $P$ is tangent to
        the level curve through $P$.
  \item A hiker stands on a hillside with neither map nor formula. How can
        she find the direction of the gradient, and how does the slope in
        other directions compare with the slope in that one?
  \item True or false? (a) For $f(x,y)$ of two variables: if
        $D_{\mathbf{u}}f(P)=0$ for two non-parallel unit vectors $\mathbf{u}$, then $\grad f(P)=\mathbf{0}$. (b) If
        $D_{\mathbf{u}}f(P)=3$, then $D_{-\mathbf{u}}f(P)=-3$. (c) A
        function can have $D_{\mathbf{u}}f(P)=5$ for some unit vector
        $\mathbf{u}$ while $|\grad f(P)|=4$.
\end{problems}

Problems 41--42: an input change vector of fixed length is to raise output
as fast as possible.

\begin{problems}[resume]
  \item $Y=10K^{0.4}L^{0.6}$ at $K=32768$, $L=7776$ (so that
        $K^{0.4}=64$, $L^{0.6}=216$). The change vector
        $(\mathrm{d}K,\mathrm{d}L)$ has length $100$. Find the best ratio
        $\mathrm{d}K:\mathrm{d}L$, the amounts, the estimated and the exact
        rise in output; then compare with spending a budget
        $\mathrm{d}K+\mathrm{d}L=100$.
  \item $Y=\sqrt{KL}$ at $K=400$, $L=900$, with a change vector of length
        $20$: the same questions.
\end{problems}

\section{Additional Solutions}

\begin{solutions}
  \item $\grad f=(2x+3y,\,-2y+3x)$; at $(1,2)$:
        $\grad f=(2+6,\,-4+3)=(8,-1)$. Normalise:
        $|\mathbf{v}|=\sqrt{16+9}=5$, $\mathbf{u}=(4/5,\,3/5)$.
        \[
          D_{\mathbf{u}}f=8\cdot\frac{4}{5}+(-1)\cdot\frac{3}{5}
            =\frac{32}{5}-\frac{3}{5}=\frac{29}{5} .
        \]

  \item $\grad f=(2x,\,4y,\,6z)$; at $(1,1,1)$: $\grad f=(2,4,6)$.
        Normalise: $|\mathbf{v}|=\sqrt{1+1+1}=\sqrt{3}$,
        $\mathbf{u}=(1,-1,1)/\sqrt{3}$.
        \[
          D_{\mathbf{u}}f=\frac{2-4+6}{\sqrt{3}}
            =\frac{4}{\sqrt{3}}=\frac{4\sqrt{3}}{3} .
        \]

  \item $\grad U=(1/x,\,2/y)$; at $(1,2)$: $\grad U=(1,1)$,
        $|\grad U|=\sqrt{2}$. Steepest ascent along
        \[
          \mathbf{u}=\frac{(1,1)}{\sqrt{2}},\qquad
          \text{maximum rate}=\sqrt{2} .
        \]

  \item $\grad f=(y,\,x)$; at $(2,1)$: $\grad f=(1,2)$. We need
        $\mathbf{u}\perp(1,2)$: $u_{1}+2u_{2}=0$, so
        $\mathbf{u}=\pm(2,-1)/\sqrt{5}$.
        \textbf{Check.} $(1,2)\cdot(2,-1)=2-2=0$. \checkmark

  \item The directional derivatives toward $(1,0)$ and $(0,1)$ are
        precisely $f_{x}(P)$ and $f_{y}(P)$, so $\grad f(P)=(3,-4)$.
        Maximum rate: $|\grad f(P)|=\sqrt{9+16}=5$, attained toward
        $(3,-4)/5$. Toward $\mathbf{u}=(3/5,\,4/5)$:
        \[
          D_{\mathbf{u}}f=3\cdot\frac{3}{5}+(-4)\cdot\frac{4}{5}
            =\frac{9}{5}-\frac{16}{5}=-\frac{7}{5} .
        \]

  \item $\grad f=\Bigl(\dfrac{2x}{x^{2}+y^{2}},\,
        \dfrac{2y}{x^{2}+y^{2}}\Bigr)$; at $(1,0)$: $\grad f=(2,0)$, so
        \[
          D_{\mathbf{u}}f=(2,0)\cdot(0,1)=0 :
        \]
        $\mathbf{u}$ is tangent to the level curve $x^{2}+y^{2}=1$
        through $(1,0)$. Second partials:
        \[
          f_{xx}=\frac{2y^{2}-2x^{2}}{(x^{2}+y^{2})^{2}},\qquad
          f_{yy}=\frac{2x^{2}-2y^{2}}{(x^{2}+y^{2})^{2}},\qquad
          f_{xy}=\frac{-4xy}{(x^{2}+y^{2})^{2}} ;
        \]
        at $(1,0)$: $f_{xx}=-2$, $f_{yy}=2$, $f_{xy}=0$. Hence with
        $\mathbf{u}=(0,1)$:
        \[
          D^{2}_{\mathbf{u}}f=f_{yy}=2>0 .
        \]
        Along $\mathbf{u}$ the function is flat to first order but
        curves upward: moving off $(1,0)$ vertically raises $f$ at
        second order. \textbf{Check.}
        $f(1,t)=\ln(1+t^{2})\approx t^{2}$ for small $t$: derivative
        $0$, second derivative $2$. \checkmark

  \item
        \[
          J_{\mathbf{f}}=\begin{pmatrix}2x&-2y\\3y&3x+2y\end{pmatrix},
          \qquad
          J_{\mathbf{f}}(1,1)=\begin{pmatrix}2&-2\\3&5\end{pmatrix} .
        \]

  \item $J_{f}=\bigl(y\;\;x\bigr)$, evaluated at
        $\mathbf{r}(t)=(t^{2},\,t^{3})$: $\bigl(t^{3}\;\;t^{2}\bigr)$;
        $J_{\mathbf{r}}(t)=\begin{pmatrix}2t\\3t^{2}\end{pmatrix}$.
        \[
          \dd{z}{t}=\bigl(t^{3}\;\;t^{2}\bigr)
            \begin{pmatrix}2t\\3t^{2}\end{pmatrix}
            =2t^{4}+3t^{4}=5t^{4} .
        \]
        \textbf{Check.} $z=t^{2}\cdot t^{3}=t^{5}$, so
        $\dd{z}{t}=5t^{4}$. \checkmark

  \item Inner and outer Jacobians:
        \[
          J_{\mathbf{g}}(s,t)=\begin{pmatrix}t&s\\1&1\end{pmatrix},\qquad
          J_{\mathbf{F}}(x,y)=\begin{pmatrix}1&1\\y&x\end{pmatrix} .
        \]
        Evaluate the outer at $(x,y)=\mathbf{g}(s,t)=(st,\,s+t)$ and
        multiply:
        \[
          J_{\mathbf{F}\circ\mathbf{g}}
            =\begin{pmatrix}1&1\\s+t&st\end{pmatrix}
             \begin{pmatrix}t&s\\1&1\end{pmatrix}
            =\begin{pmatrix}t+1&s+1\\
              (s+t)t+st&(s+t)s+st\end{pmatrix}
            =\begin{pmatrix}t+1&s+1\\
              t^{2}+2st&s^{2}+2st\end{pmatrix} .
        \]
        \textbf{Check.} Composing first:
        $(\mathbf{F}\circ\mathbf{g})(s,t)
        =\bigl(st+s+t,\;s^{2}t+st^{2}\bigr)$, with Jacobian
        \[
          \begin{pmatrix}t+1&s+1\\2st+t^{2}&s^{2}+2st\end{pmatrix} .
          \quad\checkmark
        \]

  \item Inner map $\mathbf{g}(K,L)=(2K+L,\;K+3L)$ with constant
        Jacobian
        $J_{\mathbf{g}}=\begin{pmatrix}2&1\\1&3\end{pmatrix}$; outer
        function $y=x_{1}x_{2}$ with row Jacobian
        $\grad y=(x_{2}\;\;x_{1})$. At $(K,L)=(1,1)$:
        $(x_{1},x_{2})=(3,4)$, so the row is $(4\;\;3)$ and
        \[
          \Bigl(\pd{y}{K},\,\pd{y}{L}\Bigr)
            =(4\;\;3)\begin{pmatrix}2&1\\1&3\end{pmatrix}
            =\bigl(4\cdot 2+3\cdot 1,\;4\cdot 1+3\cdot 3\bigr)
            =(11,\;13) .
        \]
        \textbf{Check.} Expanding,
        $y=(2K+L)(K+3L)=2K^{2}+7KL+3L^{2}$, so
        $\pd{y}{K}=4K+7L=11$ and $\pd{y}{L}=7K+6L=13$ at $(1,1)$.
        \checkmark

  \item $f(P)=2+8=10$ and $\grad f=(2xy,\;x^{2}+3y^{2})=(4,13)$ at $P$.
        \begin{itemize}[leftmargin=1.2em,itemsep=0.2ex]
          \item $x$-direction: slope $f_{x}(P)=4$, line
                $(1+t,\;2,\;10+4t)$.
          \item $y$-direction: slope $f_{y}(P)=13$, line
                $(1,\;2+t,\;10+13t)$.
          \item $|\mathbf{v}|=5$, $\mathbf{u}=(\tfrac35,-\tfrac45)$,
                $D_{\mathbf{u}}f=\tfrac{12}{5}-\tfrac{52}{5}=-8$: the
                surface falls with slope $8$ that way. Line
                $\bigl(1+\tfrac35t,\;2-\tfrac45t,\;10-8t\bigr)$.
        \end{itemize}
  \item $f(P)=\ln4$. With $w=x^{2}+3y$, $\grad f=\bigl(\tfrac{2x}{w},\,
        \tfrac{3}{w}\bigr)=(1,\tfrac34)$ at $P$, where $w=4$. Steepest
        ascent: $\mathbf{u}=\grad f/|\grad f|$ with $|\grad f|=\tfrac54$, so
        $\mathbf{u}=(\tfrac45,\tfrac35)$ and the slope is
        $D_{\mathbf{u}}f=|\grad f|=\tfrac54$. Tangent line:
        \[
          (x,y,z)=\bigl(2+\tfrac45t,\;\tfrac35t,\;\ln4+\tfrac54t\bigr) .
        \]
  \item $f(P)=2\mathrm{e}^{0}=2$. Product rule:
        $f_{x}=\mathrm{e}^{x-y}(x+y+1)=3$,
        $f_{y}=\mathrm{e}^{x-y}(1-x-y)=-1$ at $P$. $\mathbf{h}=Q-P=(2,1)$,
        $|\mathbf{h}|=\sqrt5$, $\mathbf{u}=(2,1)/\sqrt5$:
        \[
          D_{\mathbf{u}}f=\frac{3\cdot2-1\cdot1}{\sqrt5}=\sqrt5\approx2.236 .
        \]
        With the unit vector:
        $\bigl(1+\tfrac{2}{\sqrt5}t,\;1+\tfrac{1}{\sqrt5}t,\;2+\sqrt5\,t\bigr)$.
        With $\mathbf{h}$, $\grad f\cdot\mathbf{h}=5$:
        $(1+2s,\;1+s,\;2+5s)$. Putting $s=t/\sqrt5$ in the second gives the
        first, so it is the same line; only the parameter runs at a
        different speed.
  \item $Q(16,16)=4\cdot4\cdot2=32$,
        $Q_{L}=K^{1/2}L^{-3/4}=4\cdot\tfrac18=\tfrac12$. Tangent line in the
        $L$-direction: $(16,\;16+t,\;32+\tfrac12t)$. At $L=17$ ($t=1$) it
        gives $32.5$; exactly, $Q(16,17)=16\cdot17^{1/4}\approx32.4887$. The
        slope $\tfrac12$ is the marginal product of labour, and the line
        slightly overestimates because $Q$ is concave in $L$.
  \item (a) $c=\sqrt{2+6+1}=3$. (b) $2x^{2}+3y+1=9$, so
        $y=\dfrac{8-2x^{2}}{3}$, a downward parabola. Its slope at $x=1$ is
        $y'=-\tfrac{4x}{3}=-\tfrac43$, tangent direction $(3,-4)$.
        (c) $\grad f=\Bigl(\dfrac{2x}{f},\,\dfrac{3}{2f}\Bigr)=
        \bigl(\tfrac23,\tfrac12\bigr)$, and
        $\bigl(\tfrac23,\tfrac12\bigr)\cdot(3,-4)=2-2=0$. \checkmark\ (d)
        Normal form: $\tfrac23(x-1)+\tfrac12(y-2)=0$, i.e.\
        $4(x-1)+3(y-2)=0$, $y=2-\tfrac43(x-1)$. Parametric:
        $(x,y)=(1,2)+t(3,-4)$. (e) $D_{\mathbf{u}}f=|\grad f|=
        \sqrt{\tfrac49+\tfrac14}=\sqrt{\tfrac{25}{36}}=\tfrac56$.
  \item (a) $c=1\cdot8^{2/3}=4$. (b) $y^{2/3}=4x^{-1/3}$, so
        $y=8x^{-1/2}$; slope $y'=-4x^{-3/2}=-4$ at $x=1$, direction $(1,-4)$.
        (c) $f_{x}=\tfrac13\dfrac fx=\tfrac43$,
        $f_{y}=\tfrac23\dfrac fy=\tfrac23\cdot\tfrac48=\tfrac13$, and
        $\bigl(\tfrac43,\tfrac13\bigr)\cdot(1,-4)=0$. \checkmark\ (d)
        $\tfrac43(x-1)+\tfrac13(y-8)=0$, i.e.\ $y=12-4x$; parametric
        $(1,8)+t(1,-4)$. (e) $|\grad f|=\sqrt{\tfrac{16}{9}+\tfrac19}=
        \tfrac{\sqrt{17}}{3}\approx1.374$.
  \item (a) $c=2\cdot1+0=2$. (b) $x\mathrm{e}^{y}=2-y$, so
        $x=(2-y)\mathrm{e}^{-y}$. Then
        $\dd{x}{y}=-\mathrm{e}^{-y}-(2-y)\mathrm{e}^{-y}=(y-3)\mathrm{e}^{-y}
        =-3$ at $y=0$: tangent direction $(\mathrm{d}x,\mathrm{d}y)=(-3,1)$.
        (c) $\grad f=(\mathrm{e}^{y},\;x\mathrm{e}^{y}+1)=(1,3)$, and
        $(1,3)\cdot(-3,1)=0$. \checkmark\ (d) $(x-2)+3y=0$, i.e.\
        $x=2-3y$; parametric $(2,0)+t(-3,1)$. (e) $|\grad f|=\sqrt{10}
        \approx3.162$.
  \item Each component is a polynomial, so its partials exist and are
        continuous everywhere; a $C^{1}$ map is differentiable, and a map is
        differentiable when each component is.
        $J_{\mathbf{F}}=\begin{pmatrix}2xy&x^{2}\\1&-6y\\y&x\end{pmatrix}$,
        and $J_{\mathbf{F}}(1,1)=\begin{pmatrix}2&1\\1&-6\\1&1\end{pmatrix}$.
        Remainder at $(1,1)$, component by component, with
        $\mathbf{F}(1,1)=(1,-2,1)$:
        \[\begin{aligned}
          r_{1}&=(1+h_{1})^{2}(1+h_{2})-1-2h_{1}-h_{2}
               =2h_{1}h_{2}+h_{1}^{2}+h_{1}^{2}h_{2},\\
          r_{2}&=(1+h_{1})-3(1+h_{2})^{2}+2-h_{1}+6h_{2}=-3h_{2}^{2},\\
          r_{3}&=(1+h_{1})(1+h_{2})-1-h_{1}-h_{2}=h_{1}h_{2} .
        \end{aligned}\]
        Each is bounded by a constant times $\|\mathbf{h}\|^{2}$ (for
        $\|\mathbf{h}\|\le1$), so
        $\|\mathbf{r}\|\le C\|\mathbf{h}\|^{2}$ and
        $\|\mathbf{r}\|/\|\mathbf{h}\|\to0$.
  \item \emph{Domain.} $3x^{2}+y^{4}>0$ fails only at the origin: $D=\R^{2}
        \setminus\{(0,0)\}$. \emph{Differentiability.} The first two
        components are $C^{1}$ on $D$. The third, $|y-1|$, has
        $\partial/\partial y$ equal to $+1$ above and $-1$ below the line
        $y=1$, and on the line the quotient $|h|/h$ has no limit. So
        $\mathbf{F}$ is differentiable exactly on
        $\{(x,y)\neq(0,0):y\neq1\}$, with
        \[
          J_{\mathbf{F}}=\begin{pmatrix}
            \dfrac{6x}{3x^{2}+y^{4}}&\dfrac{4y^{3}}{3x^{2}+y^{4}}\\[1.2ex]
            y^{2}&2xy\\ 0&\operatorname{sign}(y-1)\end{pmatrix} .
        \]
  \item $\mathbf{f}$ is differentiable at $\mathbf{x}_{0}$ if
        $\mathbf{f}(\mathbf{x}_{0}+\mathbf{h})=\mathbf{f}(\mathbf{x}_{0})
        +J\mathbf{h}+\mathbf{r}(\mathbf{h})$ for a matrix $J$, with
        $\|\mathbf{r}(\mathbf{h})\|/\|\mathbf{h}\|\to0$: near
        $\mathbf{x}_{0}$ the map is a linear map plus an error that vanishes
        faster than the step. The matrix is the derivative (total
        derivative), the Jacobian matrix, of shape $m\times n$. Its entries
        are the partials $\partial f_{i}/\partial x_{j}$; their existence is
        necessary but not sufficient, while continuity of all of them is
        sufficient, and $\mathbf{f}$ is differentiable exactly when each
        component is. For $n=m=1$ the linear part is the tangent line of
        the graph, for $n=2$, $m=1$ the tangent plane.
  \item $J_{\mathbf{r}}(t)=\Bigl(\dfrac{2t}{1+t^{2}},\;-2\mathrm{e}^{-2t},
        \;-3\sin3t,\;0\Bigr)^{T}$, a $4\times1$ column (four outputs, one
        input); the constant component gives $0$. At $t=0$:
        $(0,-2,0,0)^{T}$.
  \item $J_{\mathbf{F}}=\begin{pmatrix}yz&xz&xy\\2x&0&-1\end{pmatrix}$,
        $2\times3$; at $(1,2,3)$: $\begin{pmatrix}6&3&2\\2&0&-1\end{pmatrix}$.
        $J_{g}=\bigl(x_{2}\ \ x_{1}\ \ \mathrm{e}^{x_{4}}\ \
        x_{3}\mathrm{e}^{x_{4}}\bigr)$, a $1\times4$ row; at $(1,1,1,0)$:
        $(1\ \ 1\ \ 1\ \ 1)$.
  \item $J_{\mathbf{F}}=\begin{pmatrix}1&1\\t&s\\2s&0\\
        \mathrm{e}^{s-t}&-\mathrm{e}^{s-t}\end{pmatrix}$, $4\times2$ (four
        outputs, two inputs); at $(1,1)$:
        $\begin{pmatrix}1&1\\1&1\\2&0\\1&-1\end{pmatrix}$.
  \item (a) $F_{k}=x_{k}$, so $\partial F_{k}/\partial x_{j}$ is $1$ for
        $j=k$ and $0$ otherwise: $J=I_{n}$. (b) $F_{k}=k\mathrm{e}^{x_{k}}$
        depends on $x_{k}$ alone: $J=\operatorname{diag}\bigl(\mathrm{e}^{x_{1}},
        2\mathrm{e}^{x_{2}},\ldots,n\,\mathrm{e}^{x_{n}}\bigr)$, $n\times n$.
        (c) $r_{k}=t^{k}/k$ has $r_{k}'=t^{k-1}$: the $n\times1$ column
        $(1,t,t^{2},\ldots,t^{n-1})^{T}$. (d) The term with $x_{k}$ is
        $x_{k}^{k}/k$, with derivative $x_{k}^{k-1}$: the $1\times n$ row
        $(1\ \ x_{2}\ \ x_{3}^{2}\ \ \cdots\ \ x_{n}^{n-1})$.
  \item (a) Component $i$ is $\sum_{j}a_{ij}x_{j}$, with partials $a_{ij}$:
        $J=A$ at every point. (b) Component $k$ is $3x_{k}-2a_{k}$:
        $J=3I_{n}$, and the constant shift drops out. (c) $\partial f/
        \partial x_{k}=a_{k}$: $J=\mathbf{a}^{T}$, a $1\times n$ row. (d)
        $J=B$, $2\times3$, whatever $\mathbf{c}$ is.
  \item Match outputs of the inner map with inputs of the outer map.
        \begin{itemize}[leftmargin=1.2em,itemsep=0.2ex]
          \item $\mathbf{G}\circ\mathbf{F}$: $2=2$, defined,
                $\R^{3}\to\R^{4}$, $4\times3$. $\mathbf{F}\circ\mathbf{G}$:
                $4\neq3$, undefined.
          \item $H\circ\mathbf{G}$: $4=4$, defined, $\R^{2}\to\R$,
                $1\times2$. $\mathbf{G}\circ H$: $1\neq2$, undefined.
          \item $\mathbf{K}\circ H$: $1=1$, defined, $\R^{4}\to\R^{3}$,
                $3\times4$. $H\circ\mathbf{K}$: $3\neq4$, undefined.
          \item $\mathbf{F}\circ\mathbf{K}$: $3=3$, defined, $\R\to\R^{2}$,
                $2\times1$. $\mathbf{K}\circ\mathbf{F}$: $2\neq1$,
                undefined.
          \item $H\circ\mathbf{G}\circ\mathbf{F}$: $\R^{3}\to\R^{2}\to\R^{4}
                \to\R$, defined, $1\times3=(1\times4)(4\times2)(2\times3)$.
        \end{itemize}
  \item $\mathbf{f}(1,2,1)=(1+2,\;1)=(3,1)$.
        $J_{\mathbf{f}}=\begin{pmatrix}1&z&y\\z&0&x\end{pmatrix}
        =\begin{pmatrix}1&1&2\\1&0&1\end{pmatrix}$ and
        $J_{\mathbf{g}}=\begin{pmatrix}1&0\\0&1\\v&u\\1&-1\end{pmatrix}
        =\begin{pmatrix}1&0\\0&1\\1&3\\1&-1\end{pmatrix}$ at $(u,v)=(3,1)$.
        \[
          J_{\mathbf{g}\circ\mathbf{f}}=
          \begin{pmatrix}1&0\\0&1\\1&3\\1&-1\end{pmatrix}
          \begin{pmatrix}1&1&2\\1&0&1\end{pmatrix}
          =\begin{pmatrix}1&1&2\\1&0&1\\4&1&5\\0&1&1\end{pmatrix},
        \]
        a $(4\times2)(2\times3)=4\times3$ product. \textbf{Check.}
        $\mathbf{g}\circ\mathbf{f}=\bigl(x+yz,\;xz,\;x^{2}z+xyz^{2},\;
        x+yz-xz\bigr)$; the third row is
        $(2xz+yz^{2},\;xz^{2},\;x^{2}+2xyz)=(4,1,5)$ and the fourth
        $(1-z,\;z,\;y-x)=(0,1,1)$ at $(1,2,1)$. \checkmark
  \item $g\colon\R^{3}\to\R$ and $\mathbf{f}\colon\R\to\R^{3}$: in
        $\mathbf{f}\circ g$ the one output of $g$ feeds the one input of
        $\mathbf{f}$, and in $g\circ\mathbf{f}$ the three outputs of
        $\mathbf{f}$ feed the three inputs of $g$. Both exist.
        \emph{$\mathbf{f}\circ g\colon\R^{3}\to\R^{3}$.} The chain rule
        gives a column times a row,
        $J=\mathbf{f}'(g(\mathbf{x}))\,\grad g(\mathbf{x})^{T}$, a
        rank-one $3\times3$ matrix. At $(1,0,0)$: $g=1$,
        $\mathbf{f}'(1)=(1,\;2,\;\mathrm{e})^{T}$, $\grad g=(1,2,2z)=(1,2,0)$:
        \[
          J=\begin{pmatrix}1\\2\\\mathrm{e}\end{pmatrix}(1\ \ 2\ \ 0)
           =\begin{pmatrix}1&2&0\\2&4&0\\\mathrm{e}&2\mathrm{e}&0\end{pmatrix} .
        \]
        Check: $\mathbf{f}(g)=(g,\;g^{2},\;\mathrm{e}^{g})$ has rows
        $\grad g$, $2g\grad g$, $\mathrm{e}^{g}\grad g$. \checkmark\
        \emph{$g\circ\mathbf{f}\colon\R\to\R$.}
        $\grad g(\mathbf{f}(s))\cdot\mathbf{f}'(s)=(1,\;2,\;2\mathrm{e}^{s})
        \cdot(1,\;2s,\;\mathrm{e}^{s})=1+4s+2\mathrm{e}^{2s}$. Check:
        $g(\mathbf{f}(s))=s+2s^{2}+\mathrm{e}^{2s}$. \checkmark
  \item \emph{$\mathbf{g}\circ\mathbf{f}\colon\R^{2}\to\R^{2}$} at $(1,2)$:
        $\mathbf{f}(1,2)=(1,2,2)$,
        $J_{\mathbf{f}}=\begin{pmatrix}1&0\\0&1\\y&x\end{pmatrix}
        =\begin{pmatrix}1&0\\0&1\\2&1\end{pmatrix}$,
        $J_{\mathbf{g}}=\begin{pmatrix}1&0&1\\0&w&v\end{pmatrix}
        =\begin{pmatrix}1&0&1\\0&2&2\end{pmatrix}$, so
        \[
          J_{\mathbf{g}\circ\mathbf{f}}=\begin{pmatrix}1&0&1\\0&2&2\end{pmatrix}
          \begin{pmatrix}1&0\\0&1\\2&1\end{pmatrix}
          =\begin{pmatrix}3&1\\4&4\end{pmatrix} .
        \]
        Check: $\mathbf{g}(\mathbf{f}(x,y))=(x+xy,\;xy^{2})$, Jacobian
        $\begin{pmatrix}1+y&x\\y^{2}&2xy\end{pmatrix}$, which is
        $\begin{pmatrix}3&1\\4&4\end{pmatrix}$ at $(1,2)$. \checkmark\
        \emph{$\mathbf{f}\circ\mathbf{g}\colon\R^{3}\to\R^{3}$} at
        $(1,1,1)$: $\mathbf{g}(1,1,1)=(2,1)$, so
        $J_{\mathbf{f}}(2,1)=\begin{pmatrix}1&0\\0&1\\1&2\end{pmatrix}$ and
        $J_{\mathbf{g}}(1,1,1)=\begin{pmatrix}1&0&1\\0&1&1\end{pmatrix}$:
        \[
          J_{\mathbf{f}\circ\mathbf{g}}=\begin{pmatrix}1&0\\0&1\\1&2\end{pmatrix}
          \begin{pmatrix}1&0&1\\0&1&1\end{pmatrix}
          =\begin{pmatrix}1&0&1\\0&1&1\\1&2&3\end{pmatrix} .
        \]
        Check: $\mathbf{f}(\mathbf{g})=\bigl(u+w,\;vw,\;(u+w)vw\bigr)$; the
        last row is $\bigl(vw,\;(u+w)w,\;vw+(u+w)v\bigr)=(1,2,3)$.
        \checkmark
  \item $\grad F=(2x,2y,2z)$ and $\mathbf{r}'(t)=(-\sin t,\cos t,1)$:
        \[
          \dd{}{t}F(\mathbf{r}(t))=2\cos t(-\sin t)+2\sin t\cos t+2t\cdot1
          =2t .
        \]
        Check: $F(\mathbf{r}(t))=\cos^{2}t+\sin^{2}t+t^{2}=1+t^{2}$.
        \checkmark
  \item $|\mathbf{h}|=3$. $\grad f=(\ln y,\;x/y+z^{2},\;2yz)=(0,3,2)$, so
        $\grad f\cdot\mathbf{h}=0+6-4=2$ and $D_{\mathbf{u}}f=\tfrac23$.
        $H=\begin{pmatrix}0&1/y&0\\1/y&-x/y^{2}&2z\\0&2z&2y\end{pmatrix}
        =\begin{pmatrix}0&1&0\\1&-2&2\\0&2&2\end{pmatrix}$ and
        \[
          \mathbf{h}^{T}H\mathbf{h}=2\cdot1\cdot(1)(2)-2\cdot4
          +2\cdot2\cdot(2)(-2)+2\cdot4=4-8-16+8=-12,
        \]
        so $D^{2}_{\mathbf{u}}f=-\tfrac{12}{9}=-\tfrac43$. \emph{Profile.}
        $\varphi(t)=(2+t)\ln(1+2t)+(1+2t)(1-2t)^{2}$, and
        $(1+2t)(1-2t)^{2}=1-2t-4t^{2}+8t^{3}$. The log part has derivative
        $\ln(1+2t)+\dfrac{2(2+t)}{1+2t}$, equal to $4$ at $0$, and second
        derivative $\dfrac{2}{1+2t}+\dfrac{2(1+2t)-4(2+t)}{(1+2t)^{2}}$,
        equal to $2-6=-4$. So $\varphi'(0)=4-2=2$ and
        $\varphi''(0)=-4-8=-12$. \checkmark\ Concave in this direction.
  \item $|\mathbf{h}|=3$. $\grad f=\mathrm{e}^{yz}(1,\;xz,\;xy)=(1,2,0)$;
        $D_{\mathbf{u}}f=\tfrac{2+2+0}{3}=\tfrac43$. Second partials at $P$
        ($\mathrm{e}^{yz}=1$): $f_{xx}=0$, $f_{xy}=z\mathrm{e}^{yz}=2$,
        $f_{xz}=y\mathrm{e}^{yz}=0$, $f_{yy}=xz^{2}\mathrm{e}^{yz}=4$,
        $f_{yz}=x(1+yz)\mathrm{e}^{yz}=1$, $f_{zz}=xy^{2}\mathrm{e}^{yz}=0$.
        \[
          \mathbf{h}^{T}H\mathbf{h}=2\cdot2\cdot(2)(1)+4\cdot1+2\cdot1\cdot(1)(2)
          =8+4+4=16,\qquad D^{2}_{\mathbf{u}}f=\tfrac{16}{9} .
        \]
        \emph{Profile.} $\varphi(t)=(1+2t)\,\mathrm{e}^{t(2+2t)}
        =(1+2t)\,\mathrm{e}^{2t+2t^{2}}$. With $E=\mathrm{e}^{2t+2t^{2}}$,
        $E'=(2+4t)E$:
        $\varphi'=\bigl(2+(1+2t)(2+4t)\bigr)E$, so $\varphi'(0)=4$, and
        $\varphi''=\bigl(4+8t+4+8t\bigr)E+\bigl(2+(1+2t)(2+4t)\bigr)(2+4t)E$,
        so $\varphi''(0)=8+4\cdot2=16$. \checkmark\ Convex in this direction.
  \item $|\mathbf{h}|=\sqrt3$. $\grad f=(yz+\tfrac1x,\;xz,\;xy+\tfrac1z)
        =(2,1,2)$, $\grad f\cdot\mathbf{h}=2-1+2=3$,
        $D_{\mathbf{u}}f=3/\sqrt3=\sqrt3$.
        $H=\begin{pmatrix}-1/x^{2}&z&y\\z&0&x\\y&x&-1/z^{2}\end{pmatrix}
        =\begin{pmatrix}-1&1&1\\1&0&1\\1&1&-1\end{pmatrix}$, and
        \[
          \mathbf{h}^{T}H\mathbf{h}=-1+0-1+2\bigl[(1)(-1)+(1)(1)+(-1)(1)\bigr]
          =-2-2=-4,\qquad D^{2}_{\mathbf{u}}f=-\tfrac43 .
        \]
        \emph{Profile.} $\varphi(t)=(1+t)(1-t)(1+t)+\ln\bigl((1+t)^{2}\bigr)
        =1+t-t^{2}-t^{3}+2\ln(1+t)$. $\varphi'(t)=1-2t-3t^{2}+\tfrac{2}{1+t}$,
        $\varphi'(0)=3$; $\varphi''(t)=-2-6t-\tfrac{2}{(1+t)^{2}}$,
        $\varphi''(0)=-4$. \checkmark\ Concave in this direction.
  \item At $(1,2)$: $\grad f=(2xy,x^{2})=(4,1)$ and
        $H=\begin{pmatrix}2y&2x\\2x&0\end{pmatrix}=\begin{pmatrix}4&2\\2&0
        \end{pmatrix}$. With $\mathbf{h}=(3,4)$:
        $\grad f\cdot\mathbf{h}=12+4=16$ and
        $\mathbf{h}^{T}H\mathbf{h}=4\cdot9+2\cdot2\cdot12+0=84$. Profile:
        $\varphi(t)=(1+3t)^{2}(2+4t)=2+16t+42t^{2}+36t^{3}$, so
        $\varphi'(0)=16$ and $\varphi''(0)=84$. \checkmark\ Practice
        problem 15 found $D_{\mathbf{u}}f=\tfrac{16}{5}$ and
        $D^{2}_{\mathbf{u}}f=\tfrac{84}{25}$ with $\mathbf{u}=\mathbf{h}/5$:
        moving along $\mathbf{h}$ covers $5$ units of distance per unit of
        $t$, so first derivatives scale by $|\mathbf{h}|=5$ and second
        derivatives by $|\mathbf{h}|^{2}=25$. The sign, and so the
        convexity verdict, is the same either way.
  \item $\grad f=(3x^{2}-3,\;2y)$, $H=\begin{pmatrix}6x&0\\0&2\end{pmatrix}$.
        At $Q=(-1,1)$: $\grad f=(0,2)$, $|\mathbf{h}|=5$,
        $\mathbf{u}=(\tfrac35,\tfrac45)$, so $D_{\mathbf{u}}f=\tfrac85$;
        $H=\begin{pmatrix}-6&0\\0&2\end{pmatrix}$ and
        $D^{2}_{\mathbf{u}}f=\dfrac{-6\cdot9+2\cdot16}{25}=-\dfrac{22}{25}<0$:
        rising, but concave in this direction at $Q$. At $(1,0)$ the
        gradient is $\mathbf{0}$ and $H=\begin{pmatrix}6&0\\0&2\end{pmatrix}$
        is PD: a strict local minimum, with positive curvature in every
        direction. Negative curvature somewhere rules out convexity,
        positive curvature somewhere rules out concavity: $f$ is
        \textbf{neither}.
  \item $f_{x}=4x^{3}-4y$, $f_{y}=-4x+2y$. From $f_{y}=0$: $y=2x$; then
        $4x^{3}=8x$, so $x=0$ or $x^{2}=2$. Critical points $(0,0)$ and
        $\pm(\sqrt2,2\sqrt2)$. $H=\begin{pmatrix}12x^{2}&-4\\-4&2\end{pmatrix}$.
        At $\pm(\sqrt2,2\sqrt2)$: $\begin{pmatrix}24&-4\\-4&2\end{pmatrix}$,
        $D_{1}=24$, $D_{2}=48-16=32$: PD, strict local minima, with value
        $4-16+8=-4$. At $(0,0)$: $\begin{pmatrix}0&-4\\-4&2\end{pmatrix}$,
        $D_{2}=-16$: a saddle. There, for $\mathbf{u}=(1,1)/\sqrt2$,
        $D^{2}_{\mathbf{u}}f=\tfrac{0-8+2}{2}=-3<0$ (concave that way), and
        for $\mathbf{u}=(0,1)$, $D^{2}_{\mathbf{u}}f=2>0$ (convex that way).
        Globally: negative curvature at the origin and positive definite
        Hessians at the minima, so $f$ is neither convex nor concave.
  \item $\grad f=(\ln y,\;x/y)=(0,1)$ and
        $H=\begin{pmatrix}0&1/y\\1/y&-x/y^{2}\end{pmatrix}=
        \begin{pmatrix}0&1\\1&-1\end{pmatrix}$ at $P$.
        \begin{itemize}[leftmargin=1.2em,itemsep=0.2ex]
          \item $\mathbf{u}=(1,0)$: $D_{\mathbf{u}}f=0$,
                $D^{2}_{\mathbf{u}}f=0$. In fact $f(x,1)=0$: flat, no
                verdict beyond ``linear''.
          \item $\mathbf{u}=(0,1)$: $D_{\mathbf{u}}f=1$,
                $D^{2}_{\mathbf{u}}f=-1$: concave in this direction.
          \item $\mathbf{u}=(1,1)/\sqrt2$: $D_{\mathbf{u}}f=\tfrac{1}{\sqrt2}$,
                $D^{2}_{\mathbf{u}}f=\tfrac12(0+2-1)=\tfrac12$: convex in
                this direction.
        \end{itemize}
        Both signs at one point already decide it: $f$ is neither convex
        nor concave. Indeed $\det H=-1/y^{2}<0$ everywhere.
  \item (a) True: along the contour $f$ is constant; equivalently the
        tangent $\mathbf{u}$ is perpendicular to $\grad f$, so
        $D_{\mathbf{u}}f=0$. (b) True when $\grad f(P)\neq\mathbf{0}$: the
        unit vector $-\grad f/|\grad f|$ gives the most negative rate,
        $-|\grad f|$. (c) False: $f=xy$ has $\grad f(0,0)=\mathbf{0}$ at a
        saddle. (d) True: the largest directional derivative is
        $|\grad f(P)|$. (e) False: it is perpendicular (normal) to the level
        curve.
  \item She walks a short way in the direction in which she neither climbs
        nor descends: she is following the contour line. The gradient
        points at right angles to that path, on the uphill side. Its
        direction is the steepest one; in a direction at angle $\theta$ to
        it the slope is $|\grad f|\cos\theta$: smaller, zero along the
        contour, and negative once she turns downhill.
  \item (a) True: $\grad f\cdot\mathbf{u}_{1}=\grad f\cdot\mathbf{u}_{2}=0$
        with $\mathbf{u}_{1},\mathbf{u}_{2}$ independent forces $\grad f$ to
        be perpendicular to the whole plane, hence $\mathbf{0}$. With three
        variables two directions are not enough: $f(x,y,z)=z$ has
        $D_{\mathbf{e}_{1}}f=D_{\mathbf{e}_{2}}f=0$ but $\grad f=(0,0,1)$.
        (b) True:
        $D_{-\mathbf{u}}f=\grad f\cdot(-\mathbf{u})=-D_{\mathbf{u}}f$.
        (c) False: by Cauchy--Schwarz,
        $D_{\mathbf{u}}f=\grad f\cdot\mathbf{u}\le|\grad f|\,|\mathbf{u}|=4$.
  \item $Y=10\cdot64\cdot216=138240$. In output-ratio form
        $Y_{K}=0.4\,Y/K=\tfrac{27}{16}=1.6875$ and
        $Y_{L}=0.6\,Y/L=\tfrac{32}{3}\approx10.667$. Best ratio
        $\mathrm{d}K:\mathrm{d}L=Y_{K}:Y_{L}=81:512$.
        $|\grad Y|=\sqrt{1.6875^{2}+10.667^{2}}\approx10.7993$, so
        \[
          (\mathrm{d}K,\mathrm{d}L)=\frac{100}{10.7993}(1.6875,\;10.667)
          \approx(15.63,\;98.77),
          \qquad \mathrm{d}Y\approx100\cdot10.7993\approx1079.9 .
        \]
        Exactly, $Y(32783.63,\;7874.77)-Y\approx1077.5$. A budget
        $\mathrm{d}K+\mathrm{d}L=100$ is maximised in the linear estimate by
        spending everything on labour, the input with the larger marginal
        product: $\mathrm{d}Y\approx1066.7$ (exactly $1063.9$). That is
        less, but it also uses less: the steepest-ascent vector has
        $\mathrm{d}K+\mathrm{d}L\approx114.4$.
  \item $Y=\sqrt{400\cdot900}=600$, $Y_{K}=\dfrac{Y}{2K}=0.75$,
        $Y_{L}=\dfrac{Y}{2L}=\tfrac13$. Ratio $\mathrm{d}K:\mathrm{d}L=
        0.75:\tfrac13=9:4$. $|\grad Y|=\sqrt{0.5625+0.1111}\approx0.82074$,
        so $(\mathrm{d}K,\mathrm{d}L)\approx\tfrac{20}{0.82074}(0.75,\;0.3333)
        \approx(18.28,\;8.12)$ and $\mathrm{d}Y\approx20\cdot0.82074\approx
        16.41$; exactly $16.32$. A budget $\mathrm{d}K+\mathrm{d}L=20$ would
        go entirely into capital: $\mathrm{d}Y\approx15$ (exactly $14.82$).
\end{solutions}

\chapter{Inverse and Implicit Function Theorems, Jacobians, Homogeneity, Euler}

When a relation between variables is given implicitly, the Implicit
Function Theorem (IFT) extracts derivatives without solving for one
variable in terms of the other, exactly the situation behind
indifference curves and isoquants. The Jacobian generalises the derivative
to vector-valued maps; its determinant is the local area / volume scaling
factor, and the Inverse Function Theorem turns that determinant into a test
for whether the map can be undone near a point --- and, when it can, gives
the derivative of the inverse without ever computing the inverse.
Homogeneous functions and Euler's theorem then give a clean
structural identity that ties marginal products to total output,
underpinning constant-returns-to-scale models.

\section{Core Rules}

\begin{rulebox}{IFT, One Equation in Two Unknowns}
Suppose $F:\R^{2}\to\R$ is $C^{1}$ and $F(x_{0},y_{0})=0$ with
$F_{y}(x_{0},y_{0})\neq 0$. Then near $(x_{0},y_{0})$ there is a $C^{1}$
function $y=g(x)$ satisfying $F(x,g(x))=0$, and
\[
  \dd{y}{x} \;=\; -\frac{F_{x}}{F_{y}} .
\]
\end{rulebox}

\begin{rulebox}{What the condition $F_{y}\neq0$ rules out}
The theorem is local and its hypothesis is sharp. Where $F_{y}=0$ the
conclusion genuinely fails, and the failure is visible in the picture of the
curve $F=0$ (Chapter~5):
\begin{itemize}[leftmargin=1.6em,itemsep=0.3ex]
  \item $F_{y}=0$, $F_{x}\neq0$: a \textbf{vertical tangent}. No $y(x)$
        exists near the point, but $x=x(y)$ does, with
        $\dd{x}{y}=-F_{y}/F_{x}=0$. The theorem is symmetric in the two
        variables: at least one of the two representations survives as long
        as $\grad F\neq\mathbf{0}$.
  \item $F_{x}=F_{y}=0$: a \textbf{singular point}, typically a
        self-intersection where two branches with different slopes meet.
        Neither representation exists, and $-F_{x}/F_{y}$ degenerates to
        $0/0$.
\end{itemize}
So the working rule is: the level set of $F$ is a well-behaved curve, with a
well-defined tangent, exactly where $\grad F\neq\mathbf{0}$. This is the
same gradient condition that reappears as the \emph{constraint
qualification} in the Lagrange and Kuhn--Tucker chapters, where it rules out
precisely these kinks and crossings in the constraint set.
\end{rulebox}

\begin{rulebox}{IFT, Multivariate}
If $F(x_{1},\ldots,x_{n},y)=0$ with $F_{y}\neq 0$, then
$y=g(x_{1},\ldots,x_{n})$ locally and
\[
  \pd{y}{x_{i}} \;=\; -\frac{F_{x_{i}}}{F_{y}} .
\]
\end{rulebox}

\begin{rulebox}{MRS and MRTS via IFT}
For $U(x,y)=c$ set $F=U(x,y)-c$. Then $F_{x}=U_{x}$, $F_{y}=U_{y}$, so
\[
  \dd{y}{x}\bigg|_{U=c} = -\frac{U_{x}}{U_{y}}\quad\text{(signed)},\qquad
  \mathrm{MRS} = \frac{U_{x}}{U_{y}}\quad\text{(unsigned)} .
\]
Likewise on $Q(K,L)=Q_{0}$: $\dd{L}{K}=-Q_{K}/Q_{L}$,
$\mathrm{MRTS}=Q_{K}/Q_{L}$.
\end{rulebox}

\begin{rulebox}{Jacobian Matrix}
For $\mathbf{f}:\R^{n}\to\R^{m}$,
$\mathbf{f}(x_{1},\ldots,x_{n})=(f_{1},\ldots,f_{m})$, the \textbf{Jacobian}
is
\[
  J_{\mathbf{f}} \;=\;
  \begin{pmatrix}
    \partial f_{1}/\partial x_{1} & \cdots & \partial f_{1}/\partial x_{n}\\
    \vdots & \ddots & \vdots\\
    \partial f_{m}/\partial x_{1} & \cdots & \partial f_{m}/\partial x_{n}
  \end{pmatrix} \;\in\; \R^{m\times n} .
\]
For $m=n$, $\det J_{\mathbf{f}}$ is the local volume scaling factor. Whether
it vanishes is exactly what decides local invertibility, which is the next
three boxes.
\end{rulebox}

\begin{rulebox}{Inverse Function Theorem, one variable}
Let $f$ be continuously differentiable near $x_{0}$ and set $y_{0}=f(x_{0})$.
If
\[
  f'(x_{0})\neq0 ,
\]
then $f$ is invertible on some open interval around $x_{0}$, the inverse is
differentiable at $y_{0}$, and
\[
  \boxed{\;\bigl(f^{-1}\bigr)'(y_{0})=\frac{1}{f'(x_{0})}
        =\frac{1}{f'\bigl(f^{-1}(y_{0})\bigr)}\;}
\]
The formula is the chain rule read backwards. Differentiating the identity
$f\bigl(f^{-1}(y)\bigr)=y$ gives
$f'\bigl(f^{-1}(y)\bigr)\cdot\bigl(f^{-1}\bigr)'(y)=1$, and dividing
finishes it --- which is legitimate precisely because $f'\neq0$.

Note which point goes where: the derivative of the \emph{inverse} is
evaluated at $y_{0}$, the derivative of $f$ at $x_{0}$. Geometrically the
graph of $f^{-1}$ is the graph of $f$ mirrored in $y=x$, and mirroring turns
a slope into its reciprocal --- a vertical tangent for $f$ becoming a
horizontal one for $f^{-1}$, which is why $f'=0$ has to be excluded.

The derivatives of $\arcsin$, $\arccos$ and $\arctan$ in Chapter~30 are this
rule applied to $\sin$, $\cos$ and $\tan$; nothing else is needed to get
them.
\end{rulebox}

\begin{rulebox}{Inverse Function Theorem, several variables}
Let $\mathbf{f}:\R^{n}\to\R^{n}$ be continuously differentiable near
$\mathbf{x}_{0}$ and set $\mathbf{y}_{0}=\mathbf{f}(\mathbf{x}_{0})$. If
\[
  \det J_{\mathbf{f}}(\mathbf{x}_{0})\neq0 ,
\]
then $\mathbf{f}$ maps some open neighbourhood of $\mathbf{x}_{0}$
one-to-one onto an open neighbourhood of $\mathbf{y}_{0}$, the local inverse
is continuously differentiable, and its Jacobian is the \textbf{matrix
inverse} of the original:
\[
  \boxed{\;J_{\mathbf{f}^{-1}}(\mathbf{y}_{0})
        =\bigl[J_{\mathbf{f}}(\mathbf{x}_{0})\bigr]^{-1}\;}
\]
For $n=1$ the matrix is $1\times1$ and this is the previous box.
The condition $\det J\neq0$ is the exact analogue of $f'\neq0$: it says the
linear approximation at the point is invertible, and the theorem promises
that the map itself inherits that property nearby.

Note that the two named theorems of this chapter are the same statement
twice. Applying the Implicit Function Theorem to
$\mathbf{F}(\mathbf{x},\mathbf{y})=\mathbf{f}(\mathbf{x})-\mathbf{y}$
--- solving for $\mathbf{x}$ in terms of $\mathbf{y}$ --- has the condition
$\det J_{\mathbf{f}}\neq0$ and delivers exactly the inverse.
\end{rulebox}

\begin{rulebox}{Local, and only local}
Both theorems are \textbf{local}, and both directions of the obvious
over-reading are false.

\emph{A non-vanishing derivative everywhere does not give a global
inverse.} The map
$\mathbf{f}(x,y)=(\mathrm{e}^{x}\cos y,\ \mathrm{e}^{x}\sin y)$ has
$\det J=\mathrm{e}^{2x}>0$ at every point, yet
$\mathbf{f}(x,y)=\mathbf{f}(x,y+2\pi)$: it is invertible near each point and
nowhere invertible globally.

\emph{A vanishing derivative does not rule out an inverse.} $f(x)=x^{3}$ has
$f'(0)=0$, and is nevertheless a bijection of $\R$. What fails is not
invertibility but \emph{differentiability of the inverse}: $f^{-1}(y)=
y^{1/3}$ has a vertical tangent at $0$, exactly as the reciprocal formula
predicts.

So the condition is sufficient for a local inverse, necessary for a
\emph{differentiable} one, and silent about global behaviour.
\end{rulebox}

\begin{rulebox}{Homogeneous Functions}
$f:\R^{n}\to\R$ is \textbf{homogeneous of degree $k$} if for all $t>0$ and
all $\mathbf{x}$ in the domain,
\[
  f(t\mathbf{x}) \;=\; t^{k}f(\mathbf{x}) .
\]
Constant returns to scale (CRS) means $k=1$.
\end{rulebox}

\begin{rulebox}{Euler's Theorem}
If $f$ is differentiable and homogeneous of degree $k$ on an open cone,
then
\[
  \sum_{i=1}^{n} x_{i}\pd{f}{x_{i}} \;=\; k\,f(\mathbf{x}) .
\]
Two-variable case: $xf_{x}+yf_{y}=kf$. \textbf{Verification rule.} Always
compute the LHS $xf_{x}+yf_{y}$ and the RHS $kf$ \emph{separately}, then
check that they agree.
\end{rulebox}

\section{Worked Examples}

\begin{worked}{$F(x,y)=x^{2}y+y^{3}-10=0$ at $(1,2)$: does it define $y(x)$, and what is $y'$? Then verify Euler's theorem on $f(K,L)=K^{0.3}L^{0.7}$}
\wstep{Step 1. Check the point lies on the curve.}
$1^{2}\cdot2+2^{3}-10=2+8-10=0$. \checkmark

\wstep{Step 2. Test the hypothesis of the Implicit Function Theorem.}
The theorem needs $F_{y}\neq0$ at the point:
\[
  F_{y}=x^{2}+3y^{2},
  \qquad
  F_{y}(1,2)=1+12=13\neq0 .
\]
So there is a neighbourhood of $x=1$ on which the equation defines a unique
differentiable function $y(x)$ with $y(1)=2$.
\wtrap{The condition is on $F_{y}$ --- the derivative with respect to the
variable you are \textbf{solving for}. If you want $x$ as a function of $y$
instead, the condition is $F_{x}\neq0$. Testing the wrong partial can
license a conclusion that is false: at a point with $F_{y}=0$ the curve is
vertical and no function $y(x)$ exists there.}

\wstep{Step 3. Apply the formula.}
\[
  F_{x}=2xy,
  \qquad
  \dd{y}{x}=-\frac{F_{x}}{F_{y}}=-\frac{2xy}{x^{2}+3y^{2}},
  \qquad
  y'(1)=-\frac{4}{13}\approx-0.308 .
\]
\wtrap{The minus sign is part of the formula, and it is dropped constantly.
It comes from the chain rule applied to the identity
$F\bigl(x,y(x)\bigr)=0$: $F_{x}+F_{y}y'=0$, so $y'=-F_{x}/F_{y}$. Derive it
once and the sign stops being something to remember.}

\wstep{Step 4. Check by differentiating implicitly, from scratch.}
\[
  \frac{\mathrm{d}}{\mathrm{d}x}\bigl[x^{2}y+y^{3}-10\bigr]=0
  \;\Longrightarrow\;
  2xy+x^{2}y'+3y^{2}y'=0
  \;\Longrightarrow\;
  y'=\frac{-2xy}{x^{2}+3y^{2}} . \checkmark
\]
At $(1,2)$: $4+y'+12y'=0$, so $y'=-4/13$. The theorem's formula is not a new
technique --- it is Chapter~5 done once in general.

\wstep{Step 5. A numerical confirmation.}
Take $x=1.01$ and solve $1.0201y+y^{3}=10$ numerically: $y\approx1.996908$.
The linear prediction is $2-0.30769\times0.01=1.996923$. \checkmark\ The two
agree to four decimals over a step of $0.01$; the residual $1.5\times10^{-5}$
is the second-order term, and it should shrink by a factor of $100$ if the
step is halved twice --- which is the signature of a first-order
approximation.

\wstep{Step 6. Where this is used: comparative statics without solving.}
An equilibrium condition $F(p,\alpha)=0$ relates a price $p$ to a parameter
$\alpha$, usually with no closed-form solution for $p$. The theorem still
gives the response:
\[
  \dd{p}{\alpha}=-\frac{F_{\alpha}}{F_{p}} ,
\]
and its \emph{sign} --- which is often all that is asked --- follows from
the signs of two partials. This is the engine of every comparative-statics
result in economics.

\wstep{Step 7. Homogeneity: check the degree by scaling.}
$f$ is homogeneous of degree $r$ if $f(tK,tL)=t^{r}f(K,L)$ for all $t>0$.
For $f=K^{0.3}L^{0.7}$:
\[
  f(tK,tL)=(tK)^{0.3}(tL)^{0.7}=t^{0.3+0.7}K^{0.3}L^{0.7}=t\,f(K,L),
\]
so $r=1$: constant returns to scale. For Cobb--Douglas the degree is simply
the \emph{sum of the exponents}.
\wtrap{Scale \emph{every} argument by the same $t$, and require the identity
for all $t>0$. $f(K,L)=K^{0.3}L^{0.7}+5$ is \textbf{not} homogeneous --- the
constant does not scale --- and Euler's theorem fails for it. Homogeneity is
a strong property, not a formality.}

\wstep{Step 8. Euler's theorem, verified.}
For $f$ homogeneous of degree $r$,
\[
  K f_{K}+L f_{L}=r\,f .
\]
Here $f_{K}=0.3K^{-0.7}L^{0.7}$ and $f_{L}=0.7K^{0.3}L^{-0.3}$, so
\[
  Kf_{K}+Lf_{L}=0.3K^{0.3}L^{0.7}+0.7K^{0.3}L^{0.7}
  =(0.3+0.7)f=1\cdot f . \checkmark
\]

\wstep{Step 9. Read the economics --- this is the product-exhaustion theorem.}
Under constant returns to scale ($r=1$), paying each factor its marginal
product exhausts output exactly:
\[
  \underbrace{K\cdot MP_{K}}_{\text{capital's share}}
  +\underbrace{L\cdot MP_{L}}_{\text{labour's share}}
  =\underbrace{Q}_{\text{total output}} .
\]
Nothing is left over and nothing is missing. With increasing returns
($r>1$) the factor payments would exceed output, which is why perfect
competition and increasing returns cannot coexist --- a theorem of this
chapter, not a modelling assumption.
\end{worked}

\begin{worked}{The inverse theorem both ways: $f(x)=x^{5}+2x+1$ at $y=4$, and $\mathbf{f}(x,y)=(x^{2}-y^{2},\,2xy)$ at $(1,1)$}
\wstep{Step 1. One variable: check the hypothesis before anything else.}
\[
  f(x)=x^{5}+2x+1,
  \qquad
  f'(x)=5x^{4}+2\ \ge 2>0\quad\text{for every }x .
\]
The derivative never vanishes, so $f$ is strictly increasing and invertible
on all of $\R$, and its inverse is differentiable everywhere.

\wstep{Step 2. Locate the point --- in the \emph{domain}.}
We are asked about $y=4$, so we need the $x$ with $f(x)=4$:
\[
  x^{5}+2x+1=4
  \quad\Longrightarrow\quad
  x=1\ \ (1+2+1=4). \checkmark
\]
\wtrap{The formula needs $x_{0}=f^{-1}(y_{0})$, not $y_{0}$. Substituting
$4$ into $f'$ gives $5\cdot4^{4}+2=1282$ and the wrong answer $1/1282$.
Find the pre-image first; in an examination it is always a value you can
guess and verify, because no one can invert a quintic either.}

\wstep{Step 3. Apply the formula.}
\[
  \bigl(f^{-1}\bigr)'(4)=\frac{1}{f'(1)}=\frac{1}{5+2}=\boxed{\tfrac17}
\]
Note what was \emph{not} done: $f^{-1}$ was never computed. It cannot be ---
a general quintic has no closed-form root --- and that is exactly the point
of the theorem. It hands you a derivative of a function you cannot write
down.

\wstep{Step 4. Sanity-check it numerically.}
$f(1.01)=1.01^{5}+2.02+1\approx4.0710$, so a step of $0.0710$ in $y$ answers
to a step of $0.01$ in $x$:
\[
  \frac{\Delta x}{\Delta y}\approx\frac{0.01}{0.0710}\approx0.141
  \approx\tfrac17=0.1429 . \checkmark
\]

\wstep{Step 5. Several variables: build the Jacobian.}
For $\mathbf{f}(x,y)=(x^{2}-y^{2},\ 2xy)$,
\[
  J_{\mathbf{f}}=\begin{pmatrix}2x&-2y\\ 2y&2x\end{pmatrix},
  \qquad
  \det J_{\mathbf{f}}=4x^{2}+4y^{2}=4(x^{2}+y^{2}).
\]
This vanishes only at the origin, so the theorem applies at every other
point.

\wstep{Step 6. Evaluate at $(1,1)$ and invert the matrix.}
$\mathbf{f}(1,1)=(1-1,\,2)=(0,2)$, and
\[
  J_{\mathbf{f}}(1,1)=\begin{pmatrix}2&-2\\2&2\end{pmatrix},
  \qquad
  \det=4+4=8\neq0 .
\]
So a local inverse exists near $(0,2)$, and its Jacobian is the matrix
inverse:
\[
  J_{\mathbf{f}^{-1}}(0,2)
  =\begin{pmatrix}2&-2\\2&2\end{pmatrix}^{-1}
  =\frac18\begin{pmatrix}2&2\\-2&2\end{pmatrix}
  =\boxed{\begin{pmatrix}1/4&1/4\\-1/4&1/4\end{pmatrix}}
\]
Check by multiplying: $\begin{pmatrix}2&-2\\2&2\end{pmatrix}
\begin{pmatrix}1/4&1/4\\-1/4&1/4\end{pmatrix}
=\begin{pmatrix}1&0\\0&1\end{pmatrix}$. \checkmark
\wtrap{Invert the \emph{matrix}; do not take reciprocals entry by entry.
$\bigl[J\bigr]^{-1}$ is a matrix inverse, and the $2\times2$ rule swaps the
diagonal and negates the off-diagonal before dividing by the determinant.
Entrywise reciprocals would give $\begin{pmatrix}1/2&-1/2\\1/2&1/2
\end{pmatrix}$, which is not the inverse of anything relevant.}

\wstep{Step 7. Now show the inverse is only local.}
$\det J\neq0$ holds at every point except the origin, which invites the
conclusion that $\mathbf{f}$ is invertible on $\R^{2}\setminus\{0\}$. It is
not:
\[
  \mathbf{f}(-1,-1)=\bigl((-1)^{2}-(-1)^{2},\ 2(-1)(-1)\bigr)=(0,2)
  =\mathbf{f}(1,1).
\]
Two different points share an image, so no global inverse exists. In general
$\mathbf{f}(-x,-y)=\mathbf{f}(x,y)$: the map is two-to-one everywhere.
\wtrap{``$\det J\neq0$ everywhere'' gives a local inverse \emph{at each
point} and says nothing about fitting them together into one function. This
is the chapter's sharpest trap, and the map above is the standard
counterexample --- it is squaring, in disguise, and squaring is two-to-one.}

\wstep{Step 8. What the two halves have in common.}
\begin{center}\footnotesize
\begin{tabular}{lll}
  \hline
   & one variable & several variables\\
  \hline
  hypothesis & $f'(x_{0})\neq0$ & $\det J(\mathbf{x}_{0})\neq0$\\
  conclusion & local inverse, differentiable & the same\\
  derivative & $1/f'(x_{0})$ & $\bigl[J(\mathbf{x}_{0})\bigr]^{-1}$\\
  evaluated at & $y_{0}=f(x_{0})$ & $\mathbf{y}_{0}=\mathbf{f}(\mathbf{x}_{0})$\\
  says nothing about & global behaviour & global behaviour\\
  \hline
\end{tabular}
\end{center}
The right-hand column is the left-hand one with ``non-zero number''
replaced by ``invertible matrix'' --- which is what $\det\neq0$ means, and
why Volume~3 had to come first.
\end{worked}

\section{Practice Problems}

\begin{problems}
  \item Use IFT to find $\dd{y}{x}$ for $x^{3}+y^{3}=6xy$.
  \item Verify homogeneity and find the degree of
        $f(x,y)=x^{4}+5x^{2}y^{2}+y^{4}$.
  \item Compute the Jacobian and its determinant for
        $\mathbf{f}(x,y)=(3x-2y,\;x+4y)$.
  \item For $U(x,y)=x+2\ln y$, find the signed slope and unsigned MRS
        using IFT.
  \item Verify Euler's theorem for $f(x,y)=x^{2}+5xy-2y^{2}$ (degree $2$).
  \item Use IFT to find $\dd{y}{x}$ for $x^{2}y+xy^{3}=10$.
  \item Verify homogeneity and find the degree of
        $f(x,y)=\dfrac{x^{3}+y^{3}}{x-y}$.
  \item Compute the Jacobian and its determinant for
        $\mathbf{f}(x,y)=(x^{2}+y^{2},\;x^{2}-y^{2})$.
  \item For $U(x,y)=(x+1)(y+2)$, find the signed slope and unsigned MRS.
  \item Verify Euler's theorem for $Q(K,L)=4K^{1/4}L^{3/4}$ (degree $1$).
  \item Use IFT to find $\dd{y}{x}$ for $\mathrm{e}^{x+y}=x+y+1$.
  \item Verify homogeneity and find the degree of $f(x,y)=\sqrt{xy}$.
  \item Compute the Jacobian and its determinant for
        $\mathbf{f}(x,y)=(x\mathrm{e}^{y},\;y\mathrm{e}^{x})$.
  \item For $U(x,y)=x^{a}y^{b}$ ($a,b>0$), find the signed slope and MRS
        via IFT.
  \item Verify Euler's theorem for $Q(K,L)=K^{1/2}L^{1/2}+K+L$
        (degree $1$).
  \item Use IFT to find $\dd{y}{x}$ for $\ln(x^{2}+y)=x-y$.
  \item Verify homogeneity for
        $f(x,y,z)=\dfrac{xy+yz+zx}{x+y+z}$ and find the degree.
  \item Compute the Jacobian and its determinant for the cylindrical map
        \[
          \mathbf{f}(\rho,\theta,z)=(\rho\cos\theta,\;\rho\sin\theta,\;z).
        \]
  \item For the Leontief $U(x,y)=\min\{2x,3y\}$, discuss MRS on each
        smooth piece and at the kink.
  \item Verify Euler's theorem for the harmonic mean
        $f(x,y)=(x^{-1}+y^{-1})^{-1}$ (degree $1$).
  \item Use IFT to find $\dd{y}{x}$ for $x\cos y+y\cos x=1$.
  \item Verify homogeneity and find the degree of
        $f(x,y)=x\ln(y/x)+y\ln(x/y)$.
  \item Compute the Jacobian of
        $\mathbf{f}(x,y)=(x+y,\;xy,\;x^{2}+y^{2})$ and identify where the
        rank drops.
  \item For $Q(K,L)=5K+3L$, find the MRTS.
  \item Verify Euler's theorem for
        $Q(K,L)=\bigl(\alpha K^{\rho}+(1-\alpha)L^{\rho}\bigr)^{1/\rho}$
        (CES, degree $1$).
  \item Use IFT to find $\dd{y}{x}$ for $x^{y}=y^{x}$ with $x,y>0$.
  \item Verify whether $f(x,y)=x^{2}+xy+y$ is homogeneous; if so, find the
        degree.
  \item Compute the Jacobian and determinant of
        $\mathbf{f}(x,y)=(\mathrm{e}^{x}\cos y,\;\mathrm{e}^{x}\sin y)$.
        Interpret.
  \item For $Q(K,L)=K^{2/3}L^{1/3}$, find the MRTS via IFT.
  \item Verify Euler's theorem for $f(x,y,z)=x^{3}+y^{3}+z^{3}+3xyz$
        (degree $3$).
  \item Use IFT to find $\dd{y}{x}$ for $\sqrt{x}+\sqrt{y}+\sqrt{xy}=5$.
  \item Verify homogeneity and find the degree of
        $f(x,y)=\dfrac{x^{2}+y^{2}}{xy}$.
  % VERIFY: right-hand entries of the spherical map obscured by a printer's
  % page-turn mark in the scan; reconstructed from the standard spherical
  % parametrization, and confirmed by the worked solution below.
  \item Compute the Jacobian and determinant of the spherical map
        \[
          \mathbf{f}(\rho,\theta,\phi)=(\rho\sin\phi\cos\theta,\;
          \rho\sin\phi\sin\theta,\;\rho\cos\phi).
        \]
  \item For $U(x,y)=ax+by$ ($a,b>0$, perfect substitutes), find the MRS.
  \item Verify Euler's theorem for $f(x,y)=\dfrac{x^{3}+y^{3}}{x+y}$
        (degree $2$).
  \item Use IFT to find $\pd{z}{x}$ and $\pd{z}{y}$ for
        $F(x,y,z)=x^{2}+y^{2}+z^{2}-R^{2}=0$.
  \item Verify homogeneity of $f(x,y)=x\mathrm{e}^{y/x}$ and find the
        degree.
  \item Compute the Jacobian and determinant of
        $\mathbf{f}(u,v)=(u^{2}-v^{2},\;2uv)$.
  \item For $U(x,y)=\alpha\ln x+\beta\ln y$, find MRS via IFT.
  \item Verify Euler's theorem for $f(x,y)=\sqrt{x^{2}+y^{2}}$
        (degree $1$).
  \item Use IFT to find $\dd{y}{x}$ for $\sin(xy)+\cos(x+y)=1/2$.
  \item Decide whether $f(x,y)=x^{2}+y^{3}$ is homogeneous, with
        justification.
  \item Compute the Jacobian determinant of the polar map
        $\mathbf{f}(r,\theta)=(r\cos\theta,\;r\sin\theta)$ and connect to
        the area element.
  \item For $Q(K,L)=(K^{\rho}+L^{\rho})^{1/\rho}$ ($\rho\neq 0$), find the
        MRTS.
  \item Verify Euler's theorem for the Leontief-like
        $f(x,y)=\min\{x,y\}$ in the interior of each smooth region.
  \item Use IFT to find $\pd{y}{x_{1}}$ and $\pd{y}{x_{2}}$ at any point
        where $F(x_{1},x_{2},y)=x_{1}^{2}+x_{2}y+y^{3}-7=0$, $F_{y}\neq 0$.
  \item Verify homogeneity and find the degree of
        $f(x,y,z)=\dfrac{x^{2}y+y^{2}z+z^{2}x}{xyz}$.
  \item Compute the Jacobian determinant of
        $\mathbf{f}(x,y)=(x+\sin y,\;y+\sin x)$ and find where it
        vanishes.
  \item For $U(x,y)=(x^{\rho}+y^{\rho})^{1/\rho}$ (CES utility), find the
        MRS.
  \item Verify Euler's theorem (3-variable form) for
        $f(x,y,z)=xy+yz+zx$ (degree $2$).
  \item Use IFT to verify that on the unit circle $x^{2}+y^{2}=1$ the
        slope at $(\sqrt{3}/2,1/2)$ equals $-\sqrt{3}$.
  \item Show that $f(x,y)=(x^{1/2}+y^{1/2})^{2}$ is homogeneous of
        degree $1$, then verify Euler's theorem (LHS and RHS separately).
  \item Compute the Jacobian determinant of
        $\mathbf{f}(x,y)=(x/y,\;y/x)$.
  \item For the translog $\ln Q = \alpha_{0}+\alpha_{K}\ln K+\alpha_{L}\ln L$,
        find the MRTS in $(K,L)$ space.
  \item Verify Euler's theorem for the production function
        $Q(K,L,M)=K^{1/3}L^{1/3}M^{1/3}$ (three-factor CRS Cobb--Douglas).
  \item Compute the Jacobian and Jacobian determinant of
        $\mathbf{f}(x,y,z)=(x+y+z,\;xy+yz+zx,\;xyz)$, and explain the link
        to elementary symmetric polynomials.
  \item Show that $f(x)=x^{5}+2x+1$ is invertible on $\R$, and compute
        $\bigl(f^{-1}\bigr)'(4)$ without inverting $f$.
  \item For $f(x)=x^{3}+3x-2$, compute $\bigl(f^{-1}\bigr)'(2)$.
  \item For $f(x)=\mathrm{e}^{x}+x$, compute $\bigl(f^{-1}\bigr)'(1)$.
  \item For $f(x)=x+\ln x$ on $x>0$, compute $\bigl(f^{-1}\bigr)'(1)$.
  \item Derive $\bigl(f^{-1}\bigr)'(y)=1/f'\bigl(f^{-1}(y)\bigr)$ by
        differentiating the identity $f\bigl(f^{-1}(y)\bigr)=y$.
  \item $f(x)=x^{3}$ is a bijection of $\R$, yet $f'(0)=0$. Which
        conclusion of the Inverse Function Theorem fails at $0$, and which
        survives?
  \item Where does the Inverse Function Theorem apply to $f(x)=x^{2}$, and
        what is the largest interval containing $x=2$ on which $f$ is
        invertible?
  \item Give an interval containing $\pi/3$ on which $\sin$ is invertible,
        and compute $\bigl(\sin^{-1}\bigr)'\bigl(\tfrac{\sqrt3}{2}\bigr)$
        two ways: by the theorem, and from the known derivative of
        $\arcsin$.
  \item For $\mathbf{f}(x,y)=(x+y,\;x-y)$: compute $J$ and $\det J$, invert
        the map explicitly, and check that $J_{\mathbf{f}^{-1}}=J^{-1}$.
  \item For $\mathbf{f}(x,y)=(x^{2}-y^{2},\;2xy)$, compute $\det J$ at
        $(1,2)$ and say whether a local inverse exists there.
  \item Show that the map of the previous problem is \emph{not} globally
        injective, despite $\det J\neq0$ off the origin.
  \item For $\mathbf{f}(x,y)=(\mathrm{e}^{x}\cos y,\;\mathrm{e}^{x}\sin y)$,
        compute $\det J$ and reconcile ``non-zero everywhere'' with the fact
        that $\mathbf{f}$ is not injective.
  \item For the polar map
        $\mathbf{f}(r,\theta)=(r\cos\theta,\;r\sin\theta)$, give $\det J$
        and the set of points where the theorem gives nothing.
  \item For $\mathbf{f}(x,y)=(x+y^{3},\;y+x^{3})$, verify the hypothesis at
        the origin and give $J_{\mathbf{f}^{-1}}$ there.
  \item For $\mathbf{f}(x,y,z)=(x,\;x+y,\;x+y+z)$: compute $\det J$, write
        the inverse map explicitly, and confirm
        $J_{\mathbf{f}^{-1}}=J^{-1}$.
  \item A demand system $\mathbf{q}=\mathbf{f}(\mathbf{p})$ has
        $\det J_{\mathbf{f}}(\mathbf{p}_{0})\neq0$. What does the Inverse
        Function Theorem permit an economist to say, and what does it not?
\end{problems}

\section{Solutions}

\begin{solutions}
  \item $F=x^{3}+y^{3}-6xy$. $F_{x}=3x^{2}-6y$, $F_{y}=3y^{2}-6x$.
        \[
          \dd{y}{x} = -\frac{3x^{2}-6y}{3y^{2}-6x} = \frac{2y-x^{2}}{y^{2}-2x} .
        \]

  \item $f(tx,ty)=(tx)^{4}+5(tx)^{2}(ty)^{2}+(ty)^{4}
        =t^{4}(x^{4}+5x^{2}y^{2}+y^{4})=t^{4}f(x,y)$. \textbf{Degree}
        $k=4$.

  \item
        \[
          J=\begin{pmatrix} 3 & -2\\ 1 & 4\end{pmatrix},\qquad
          \det J=3\cdot 4-(-2)\cdot 1=14 .
        \]
        Globally invertible; areas scale by $14$.

  \item $U_{x}=1$, $U_{y}=2/y$. By IFT $\dd{y}{x}=-U_{x}/U_{y}=-y/2$.
        \textbf{MRS} $=y/2$.

  \item $f_{x}=2x+5y$, $f_{y}=5x-4y$.
        \[
          \mathrm{LHS}=x(2x+5y)+y(5x-4y)=2x^{2}+10xy-4y^{2} .
        \]
        \[
          \mathrm{RHS}=2f=2x^{2}+10xy-4y^{2} .
        \]
        $\mathrm{LHS}=\mathrm{RHS}$. \checkmark

  \item $F=x^{2}y+xy^{3}-10$. $F_{x}=2xy+y^{3}$, $F_{y}=x^{2}+3xy^{2}$.
        \[
          \dd{y}{x} = -\frac{2xy+y^{3}}{x^{2}+3xy^{2}}
            = -\frac{y(2x+y^{2})}{x(x+3y^{2})} .
        \]

  \item
        \[
          f(tx,ty)=\frac{(tx)^{3}+(ty)^{3}}{tx-ty}
            =\frac{t^{3}(x^{3}+y^{3})}{t(x-y)}=t^{2}f(x,y) .
        \]
        \textbf{Degree} $k=2$.

  \item
        \[
          J=\begin{pmatrix} 2x & 2y\\ 2x & -2y\end{pmatrix},\qquad
          \det J=(2x)(-2y)-(2y)(2x)=-8xy .
        \]
        Locally invertible whenever $x\neq 0$ and $y\neq 0$.

  \item $U_{x}=y+2$, $U_{y}=x+1$. $\dd{y}{x}=-(y+2)/(x+1)$.
        \textbf{MRS} $=(y+2)/(x+1)$.

  \item $Q_{K}=K^{-3/4}L^{3/4}$, $Q_{L}=3K^{1/4}L^{-1/4}$.
        \[
          \mathrm{LHS}=K\cdot K^{-3/4}L^{3/4}+L\cdot 3K^{1/4}L^{-1/4}
            =K^{1/4}L^{3/4}+3K^{1/4}L^{3/4}=4K^{1/4}L^{3/4} .
        \]
        \[
          \mathrm{RHS}=1\cdot Q=4K^{1/4}L^{3/4} .
        \]
        $\mathrm{LHS}=\mathrm{RHS}$. \checkmark

  \item $F=\mathrm{e}^{x+y}-x-y-1$. $F_{x}=\mathrm{e}^{x+y}-1=F_{y}$.
        \[
          \dd{y}{x} = -\frac{F_{x}}{F_{y}} = -1 .
        \]
        Valid away from $\mathrm{e}^{x+y}=1$, i.e., $x+y\neq 0$. (The
        implicit relation forces $x+y$ to be locally constant on any
        branch.)

  \item $f(tx,ty)=\sqrt{(tx)(ty)}=t\sqrt{xy}=t f(x,y)$ for $t>0$.
        \textbf{Degree} $k=1$.

  \item
        \[
          J=\begin{pmatrix} \mathrm{e}^{y} & x\mathrm{e}^{y}\\
                             y\mathrm{e}^{x} & \mathrm{e}^{x}\end{pmatrix},
          \qquad \det J=\mathrm{e}^{x+y}(1-xy) .
        \]
        Singular along $xy=1$.

  \item $U_{x}=ax^{a-1}y^{b}$, $U_{y}=bx^{a}y^{b-1}$.
        $U_{x}/U_{y}=(a/b)(y/x)$. By IFT $\dd{y}{x}=-(a/b)(y/x)$.
        \textbf{MRS} $=(a/b)(y/x)$.

  \item $Q_{K}=\tfrac12K^{-1/2}L^{1/2}+1$, $Q_{L}=\tfrac12K^{1/2}L^{-1/2}+1$.
        \[
          \mathrm{LHS}=K\bigl(\tfrac12K^{-1/2}L^{1/2}+1\bigr)
            +L\bigl(\tfrac12K^{1/2}L^{-1/2}+1\bigr)=K^{1/2}L^{1/2}+K+L .
        \]
        \[
          \mathrm{RHS}=1\cdot Q=K^{1/2}L^{1/2}+K+L .
        \]
        $\mathrm{LHS}=\mathrm{RHS}$. \checkmark

  \item $F=\ln(x^{2}+y)-x+y$. $F_{x}=\dfrac{2x}{x^{2}+y}-1$,
        $F_{y}=\dfrac{1}{x^{2}+y}+1$.
        \[
          \dd{y}{x} = -\frac{2x/(x^{2}+y)-1}{1/(x^{2}+y)+1}
            = -\frac{2x-(x^{2}+y)}{1+x^{2}+y}
            = \frac{x^{2}+y-2x}{1+x^{2}+y} .
        \]

  \item
        \[
          f(tx,ty,tz)=\frac{t^{2}(xy+yz+zx)}{t(x+y+z)}=t f(x,y,z) .
        \]
        \textbf{Degree} $k=1$.

  \item
        \[
          J=\begin{pmatrix}
            \cos\theta & -\rho\sin\theta & 0\\
            \sin\theta & \rho\cos\theta & 0\\
            0 & 0 & 1
          \end{pmatrix} .
        \]
        Expanding along the third row:
        $\det J = 1\cdot(\rho\cos^{2}\theta+\rho\sin^{2}\theta)=\rho$.
        Volume element $\mathrm{d}V=\rho\,\mathrm{d}\rho\,\mathrm{d}\theta\,\mathrm{d}z$.

  \item Leontief, non-differentiable on $\{2x=3y\}$. On $\{2x<3y\}$:
        $U=2x$, so $U_{x}=2$, $U_{y}=0$; signed slope $\dd{y}{x}$ is
        undefined (level curve is vertical), MRS is infinite: $x$ is
        the binding factor. On $\{2x>3y\}$: $U=3y$, $U_{x}=0$, $U_{y}=3$;
        signed slope $=0$, MRS $=0$: $y$ is binding. At the kink the
        two-sided derivative does not exist.

  \item $f=(x^{-1}+y^{-1})^{-1}=xy/(x+y)$. $f_{x}=y^{2}/(x+y)^{2}$,
        $f_{y}=x^{2}/(x+y)^{2}$.
        \[
          \mathrm{LHS}=x\cdot\frac{y^{2}}{(x+y)^{2}}
            +y\cdot\frac{x^{2}}{(x+y)^{2}}
            =\frac{xy(y+x)}{(x+y)^{2}}=\frac{xy}{x+y} .
        \]
        \[
          \mathrm{RHS}=1\cdot f=\frac{xy}{x+y} .
        \]
        $\mathrm{LHS}=\mathrm{RHS}$. \checkmark

  \item $F=x\cos y+y\cos x-1$. $F_{x}=\cos y-y\sin x$,
        $F_{y}=-x\sin y+\cos x$.
        \[
          \dd{y}{x} = -\frac{\cos y-y\sin x}{\cos x-x\sin y}
            = \frac{y\sin x-\cos y}{\cos x-x\sin y} .
        \]

  \item
        \[
          f(tx,ty)=(tx)\ln(ty/tx)+(ty)\ln(tx/ty)
            =tx\ln(y/x)+ty\ln(x/y)=t f(x,y)
        \]
        (the $t$ inside the log cancels). \textbf{Degree} $k=1$.

  \item
        \[
          J=\begin{pmatrix} 1 & 1\\ y & x\\ 2x & 2y\end{pmatrix}
            \in\R^{3\times 2} .
        \]
        Not square, so there is no determinant. The top $2\times 2$
        minor is $x-y$, which vanishes precisely on $x=y$; the rank
        drops to $1$ along that diagonal (the other minors also vanish
        there).

  \item Linear: $Q_{K}=5$, $Q_{L}=3$. $\dd{L}{K}=-5/3$. \textbf{MRTS}
        $=5/3$ (perfect substitutes with constant exchange rate).

  \item Let $S=\alpha K^{\rho}+(1-\alpha)L^{\rho}$, so $Q=S^{1/\rho}$.
        \[
          Q_{K}=\frac{1}{\rho}S^{1/\rho-1}\cdot\alpha\rho K^{\rho-1}
            =\alpha S^{1/\rho-1}K^{\rho-1},\qquad
          Q_{L}=(1-\alpha)S^{1/\rho-1}L^{\rho-1} .
        \]
        \[
          \mathrm{LHS}=K\alpha S^{1/\rho-1}K^{\rho-1}
            +L(1-\alpha)S^{1/\rho-1}L^{\rho-1}
            =S^{1/\rho-1}\bigl(\alpha K^{\rho}+(1-\alpha)L^{\rho}\bigr)
            =S^{1/\rho} .
        \]
        \[
          \mathrm{RHS}=1\cdot Q=S^{1/\rho} .
        \]
        $\mathrm{LHS}=\mathrm{RHS}$. \checkmark

  \item Take logs: $y\ln x-x\ln y=0$. Set $F=y\ln x-x\ln y$.
        \[
          F_{x}=\frac{y}{x}-\ln y,\qquad F_{y}=\ln x-\frac{x}{y} .
        \]
        \[
          \dd{y}{x} = -\frac{y/x-\ln y}{\ln x-x/y}
            = \frac{\ln y-y/x}{\ln x-x/y}
            = \frac{y(x\ln y-y)}{x(y\ln x-x)} .
        \]

  \item Test: $f(tx,ty)=(tx)^{2}+(tx)(ty)+(ty)=t^{2}x^{2}+t^{2}xy+ty$.
        This cannot be written as $t^{k}(x^{2}+xy+y)$ for a single $k$:
        the first two terms scale by $t^{2}$, the last by $t$.
        \textbf{Not homogeneous.}

  \item $f_{1}=\mathrm{e}^{x}\cos y$, $f_{2}=\mathrm{e}^{x}\sin y$.
        \[
          J=\begin{pmatrix} \mathrm{e}^{x}\cos y & -\mathrm{e}^{x}\sin y\\
                             \mathrm{e}^{x}\sin y & \mathrm{e}^{x}\cos y
            \end{pmatrix},\qquad
          \det J=\mathrm{e}^{2x}(\cos^{2}y+\sin^{2}y)=\mathrm{e}^{2x} .
        \]
        Always strictly positive: the map is locally invertible
        everywhere. (It is the complex exponential $w=\mathrm{e}^{x+iy}$;
        conformal, orientation-preserving.)

  \item $Q_{K}=\tfrac23K^{-1/3}L^{1/3}$, $Q_{L}=\tfrac13K^{2/3}L^{-2/3}$.
        \[
          \frac{Q_{K}}{Q_{L}}=2\,\frac{L}{K} .
        \]
        $\dd{L}{K}=-2L/K$. \textbf{MRTS} $=2L/K$.

  \item $f_{x}=3x^{2}+3yz$, $f_{y}=3y^{2}+3xz$, $f_{z}=3z^{2}+3xy$.
        \[
          \mathrm{LHS}=x(3x^{2}+3yz)+y(3y^{2}+3xz)+z(3z^{2}+3xy)
            =3(x^{3}+y^{3}+z^{3})+9xyz=3f .
        \]
        \[
          \mathrm{RHS}=3f .
        \]
        $\mathrm{LHS}=\mathrm{RHS}$. \checkmark

  \item $F=\sqrt{x}+\sqrt{y}+\sqrt{xy}-5$. Compute:
        \[
          F_{x}=\frac{1}{2\sqrt{x}}+\frac{y}{2\sqrt{xy}}
            =\frac{1}{2\sqrt{x}}+\frac{1}{2}\sqrt{\frac{y}{x}}
            =\frac{1+\sqrt{y}}{2\sqrt{x}} .
        \]
        By symmetry, $F_{y}=\dfrac{1+\sqrt{x}}{2\sqrt{y}}$.
        \[
          \dd{y}{x} = -\frac{F_{x}}{F_{y}}
            = -\frac{\sqrt{y}(1+\sqrt{y})}{\sqrt{x}(1+\sqrt{x})} .
        \]

  \item
        \[
          f(tx,ty)=\frac{t^{2}x^{2}+t^{2}y^{2}}{(tx)(ty)}
            =\frac{t^{2}(x^{2}+y^{2})}{t^{2}xy}=f(x,y) .
        \]
        \textbf{Degree} $k=0$.

  \item
        \[
          J=\begin{pmatrix}
            \sin\phi\cos\theta & -\rho\sin\phi\sin\theta & \rho\cos\phi\cos\theta\\
            \sin\phi\sin\theta & \rho\sin\phi\cos\theta & \rho\cos\phi\sin\theta\\
            \cos\phi & 0 & -\rho\sin\phi
          \end{pmatrix} .
        \]
        Expand along the third row:
        \[
          \det J = \cos\phi\Bigl[(-\rho\sin\phi\sin\theta)(\rho\cos\phi\sin\theta)
            -(\rho\cos\phi\cos\theta)(\rho\sin\phi\cos\theta)\Bigr]
        \]
        \[
          \qquad{}+(-\rho\sin\phi)\Bigl[(\sin\phi\cos\theta)(\rho\sin\phi\cos\theta)
            -(-\rho\sin\phi\sin\theta)(\sin\phi\sin\theta)\Bigr] .
        \]
        First bracket:
        \[
          -\rho^{2}\sin\phi\cos\phi(\sin^{2}\theta+\cos^{2}\theta)
            =-\rho^{2}\sin\phi\cos\phi .
        \]
        Multiplied by $\cos\phi$ this is
        $-\rho^{2}\sin\phi\cos^{2}\phi$. Second bracket:
        $\rho\sin^{2}\phi(\cos^{2}\theta+\sin^{2}\theta)=\rho\sin^{2}\phi$.
        Multiplied by $-\rho\sin\phi$: $-\rho^{2}\sin^{3}\phi$. Sum:
        $-\rho^{2}\sin\phi(\cos^{2}\phi+\sin^{2}\phi)=-\rho^{2}\sin\phi$.
        Thus $|\det J|=\rho^{2}\sin\phi$, the spherical volume element.

  \item $U_{x}=a$, $U_{y}=b$. $\dd{y}{x}=-a/b$. \textbf{MRS} $=a/b$
        (constant, perfect substitutes).

  \item Simplify: $f=\dfrac{(x+y)(x^{2}-xy+y^{2})}{x+y}=x^{2}-xy+y^{2}$.
        Then $f_{x}=2x-y$, $f_{y}=-x+2y$.
        \[
          \mathrm{LHS}=x(2x-y)+y(-x+2y)=2x^{2}-2xy+2y^{2} .
        \]
        \[
          \mathrm{RHS}=2f=2(x^{2}-xy+y^{2})=2x^{2}-2xy+2y^{2} .
        \]
        $\mathrm{LHS}=\mathrm{RHS}$. \checkmark

  \item $F=x^{2}+y^{2}+z^{2}-R^{2}$. $F_{x}=2x$, $F_{y}=2y$, $F_{z}=2z$.
        Provided $z\neq 0$:
        \[
          \pd{z}{x}=-\frac{F_{x}}{F_{z}}=-\frac{x}{z},\qquad
          \pd{z}{y}=-\frac{F_{y}}{F_{z}}=-\frac{y}{z} .
        \]
        (Geometrically, the outward normal on the sphere has direction
        $(x,y,z)/R$.)

  \item $f(tx,ty)=(tx)\mathrm{e}^{(ty)/(tx)}=tx\,\mathrm{e}^{y/x}
        =t\,f(x,y)$. \textbf{Degree} $k=1$.

  \item
        \[
          J=\begin{pmatrix} 2u & -2v\\ 2v & 2u\end{pmatrix},\qquad
          \det J=4u^{2}+4v^{2}=4(u^{2}+v^{2}) .
        \]
        Vanishes only at the origin. (This is the complex squaring
        $w=(u+iv)^{2}$.)

  \item $U_{x}=\alpha/x$, $U_{y}=\beta/y$. $U_{x}/U_{y}=(\alpha/\beta)(y/x)$.
        $\dd{y}{x}=-(\alpha/\beta)(y/x)$. \textbf{MRS}
        $=(\alpha/\beta)(y/x)$.

  \item $f=(x^{2}+y^{2})^{1/2}$. $f_{x}=x/\sqrt{x^{2}+y^{2}}$,
        $f_{y}=y/\sqrt{x^{2}+y^{2}}$.
        \[
          \mathrm{LHS}=\frac{x^{2}}{\sqrt{x^{2}+y^{2}}}
            +\frac{y^{2}}{\sqrt{x^{2}+y^{2}}}
            =\frac{x^{2}+y^{2}}{\sqrt{x^{2}+y^{2}}}=\sqrt{x^{2}+y^{2}} .
        \]
        \[
          \mathrm{RHS}=1\cdot f=\sqrt{x^{2}+y^{2}} .
        \]
        $\mathrm{LHS}=\mathrm{RHS}$. \checkmark

  \item $F=\sin(xy)+\cos(x+y)-1/2$. $F_{x}=y\cos(xy)-\sin(x+y)$,
        $F_{y}=x\cos(xy)-\sin(x+y)$.
        \[
          \dd{y}{x} = -\frac{y\cos(xy)-\sin(x+y)}{x\cos(xy)-\sin(x+y)}
            = \frac{\sin(x+y)-y\cos(xy)}{x\cos(xy)-\sin(x+y)} .
        \]

  \item $f(tx,ty)=t^{2}x^{2}+t^{3}y^{3}$. The two terms scale with
        different powers of $t$, so they cannot be factored as
        $t^{k}(x^{2}+y^{3})$. \textbf{Not homogeneous.}

  \item
        \[
          J=\begin{pmatrix} \cos\theta & -r\sin\theta\\
                             \sin\theta & r\cos\theta\end{pmatrix},\qquad
          \det J=r\cos^{2}\theta+r\sin^{2}\theta=r .
        \]
        Hence $|\det J|=r$, and the area element is
        $\mathrm{d}A=r\,\mathrm{d}r\,\mathrm{d}\theta$, the standard
        polar substitution.

  \item Let $S=K^{\rho}+L^{\rho}$. Then $Q=S^{1/\rho}$,
        $Q_{K}=S^{1/\rho-1}K^{\rho-1}$, $Q_{L}=S^{1/\rho-1}L^{\rho-1}$.
        \[
          \frac{Q_{K}}{Q_{L}}=\Bigl(\frac{K}{L}\Bigr)^{\rho-1}
            =\Bigl(\frac{L}{K}\Bigr)^{1-\rho} .
        \]
        $\dd{L}{K}=-(L/K)^{1-\rho}$. \textbf{MRTS} $=(L/K)^{1-\rho}$.

  \item On the smooth piece $\{x<y\}$: $f=x$, $f_{x}=1$, $f_{y}=0$.
        \[
          \mathrm{LHS}=x\cdot 1+y\cdot 0=x=f,\qquad \mathrm{RHS}=1\cdot f=x .
        \]
        $\mathrm{LHS}=\mathrm{RHS}$. \checkmark On $\{x>y\}$ by symmetry
        $\mathrm{LHS}=y=f=\mathrm{RHS}$. \checkmark Euler's identity
        holds wherever the function is differentiable. (At the ridge
        $x=y$ the function is not differentiable.)

  \item By IFT, $\pd{y}{x_{i}}=-F_{x_{i}}/F_{y}$ with
        $F_{y}=x_{2}+3y^{2}$:
        \[
          \pd{y}{x_{1}}=-\frac{2x_{1}}{x_{2}+3y^{2}},\qquad
          \pd{y}{x_{2}}=-\frac{y}{x_{2}+3y^{2}} .
        \]

  \item Numerator has degree $3$, denominator degree $3$. Concretely
        \[
          f(tx,ty,tz)=\frac{t^{3}(x^{2}y+y^{2}z+z^{2}x)}{t^{3}xyz}=f(x,y,z) .
        \]
        \textbf{Degree} $k=0$.

  \item $J=\begin{pmatrix} 1 & \cos y\\ \cos x & 1\end{pmatrix}$,
        $\det J=1-\cos x\cos y$. Vanishes when $\cos x\cos y=1$, i.e.,
        $\cos x=\cos y=1$ or $\cos x=\cos y=-1$, equivalently
        $(x,y)\in 2\pi\mathbb{Z}\times 2\pi\mathbb{Z}$ or
        $(2\mathbb{Z}+1)\pi\times(2\mathbb{Z}+1)\pi$.

  \item Let $S=x^{\rho}+y^{\rho}$, $U=S^{1/\rho}$.
        $U_{x}=S^{1/\rho-1}x^{\rho-1}$, $U_{y}=S^{1/\rho-1}y^{\rho-1}$.
        \[
          \frac{U_{x}}{U_{y}}=\Bigl(\frac{x}{y}\Bigr)^{\rho-1}
            =\Bigl(\frac{y}{x}\Bigr)^{1-\rho} .
        \]
        $\dd{y}{x}=-(y/x)^{1-\rho}$. \textbf{MRS} $=(y/x)^{1-\rho}$.

  \item $f_{x}=y+z$, $f_{y}=x+z$, $f_{z}=x+y$.
        \[
          \mathrm{LHS}=x(y+z)+y(x+z)+z(x+y)=2(xy+yz+zx)=2f .
        \]
        \[
          \mathrm{RHS}=2f .
        \]
        $\mathrm{LHS}=\mathrm{RHS}$. \checkmark

  \item $F=x^{2}+y^{2}-1$. $F_{x}=2x$, $F_{y}=2y$. $\dd{y}{x}=-x/y$.
        At $(\sqrt{3}/2,1/2)$: $\dd{y}{x}=-(\sqrt{3}/2)/(1/2)=-\sqrt{3}$.
        \checkmark

  \item $f(tx,ty)=\bigl((tx)^{1/2}+(ty)^{1/2}\bigr)^{2}
        =(t^{1/2})^{2}(\sqrt{x}+\sqrt{y})^{2}=t f(x,y)$. \textbf{Degree}
        $k=1$. For Euler's theorem, write $A=\sqrt{x}+\sqrt{y}$,
        $f=A^{2}$. Then $f_{x}=2A\cdot\dfrac{1}{2\sqrt{x}}=A/\sqrt{x}$,
        $f_{y}=A/\sqrt{y}$.
        \[
          \mathrm{LHS}=x\cdot\frac{A}{\sqrt{x}}+y\cdot\frac{A}{\sqrt{y}}
            =A(\sqrt{x}+\sqrt{y})=A^{2}=f .
        \]
        \[
          \mathrm{RHS}=1\cdot f=A^{2} .
        \]
        $\mathrm{LHS}=\mathrm{RHS}$. \checkmark

  \item
        \[
          J=\begin{pmatrix} 1/y & -x/y^{2}\\ -y/x^{2} & 1/x\end{pmatrix},
          \qquad \det J=\frac{1}{xy}-\frac{xy}{x^{2}y^{2}}
            =\frac{1}{xy}-\frac{1}{xy}=0 .
        \]
        The Jacobian is identically singular: because
        $(x/y)\cdot(y/x)=1$, the image lies on the curve $f_{1}f_{2}=1$
        (one-dimensional), so the map is not locally invertible
        anywhere.

  \item From $\ln Q=\alpha_{0}+\alpha_{K}\ln K+\alpha_{L}\ln L$, we get
        $Q=AK^{\alpha_{K}}L^{\alpha_{L}}$ with $A=\mathrm{e}^{\alpha_{0}}$
        (Cobb--Douglas). Then $Q_{K}=A\alpha_{K}K^{\alpha_{K}-1}L^{\alpha_{L}}$,
        $Q_{L}=A\alpha_{L}K^{\alpha_{K}}L^{\alpha_{L}-1}$.
        \[
          \frac{Q_{K}}{Q_{L}}=\frac{\alpha_{K}}{\alpha_{L}}\cdot\frac{L}{K} .
        \]
        $\dd{L}{K}=-(\alpha_{K}/\alpha_{L})(L/K)$. \textbf{MRTS}
        $=(\alpha_{K}/\alpha_{L})(L/K)$.

  \item $Q=K^{1/3}L^{1/3}M^{1/3}$ is homogeneous of degree
        $1/3+1/3+1/3=1$. $Q_{K}=\tfrac13K^{-2/3}L^{1/3}M^{1/3}=Q/(3K)$,
        and similarly $Q_{L}=Q/(3L)$, $Q_{M}=Q/(3M)$.
        \[
          \mathrm{LHS}=K\cdot\frac{Q}{3K}+L\cdot\frac{Q}{3L}
            +M\cdot\frac{Q}{3M}=\frac{Q}{3}+\frac{Q}{3}+\frac{Q}{3}=Q .
        \]
        \[
          \mathrm{RHS}=1\cdot Q=Q .
        \]
        $\mathrm{LHS}=\mathrm{RHS}$. \checkmark

  \item Write $f_{1}=x+y+z$, $f_{2}=xy+yz+zx$, $f_{3}=xyz$ (the
        elementary symmetric polynomials $e_{1},e_{2},e_{3}$). Then
        \[
          J=\begin{pmatrix}
            1 & 1 & 1\\
            y+z & x+z & x+y\\
            yz & xz & xy
          \end{pmatrix} .
        \]
        For the determinant, subtract column $1$ from columns $2$ and
        $3$ to get
        \[
          \det J = \det\begin{pmatrix}
            1 & 0 & 0\\
            y+z & x-y & x-z\\
            yz & z(x-y) & y(x-z)
          \end{pmatrix}
          = 1\cdot\det\begin{pmatrix}
            x-y & x-z\\ z(x-y) & y(x-z)
          \end{pmatrix}
        \]
        \[
          = (x-y)(x-z)\det\begin{pmatrix} 1 & 1\\ z & y\end{pmatrix}
          = (x-y)(x-z)(y-z) .
        \]
        Thus $\det J=-(x-y)(y-z)(z-x)$ (the Vandermonde discriminant, up
        to sign), which vanishes exactly when two of $x,y,z$ coincide:
        the symmetric-polynomial map fails to be locally invertible
        at points with repeated coordinates, consistent with the fact
        that an unordered triple $\{x,y,z\}$ has a unique sorted
        representative.
  \item $f'(x)=5x^{4}+2\ge2>0$ for all $x$, so $f$ is strictly increasing,
        hence injective, and the theorem applies everywhere. Solve
        $x^{5}+2x+1=4$: $x=1$ works ($1+2+1=4$), and monotonicity makes it
        the only root. So
        $\bigl(f^{-1}\bigr)'(4)=1/f'(1)=1/7$.
  \item $f(1)=1+3-2=2$, and $f'(x)=3x^{2}+3$, so $f'(1)=6$ and
        $\bigl(f^{-1}\bigr)'(2)=\tfrac16$. ($f'>0$ everywhere, so the
        inverse is global.)
  \item $f(0)=1+0=1$, and $f'(x)=\mathrm{e}^{x}+1$, so $f'(0)=2$ and
        $\bigl(f^{-1}\bigr)'(1)=\tfrac12$.
  \item $f(1)=1+\ln1=1$, and $f'(x)=1+1/x$, so $f'(1)=2$ and
        $\bigl(f^{-1}\bigr)'(1)=\tfrac12$.
  \item Differentiate both sides with respect to $y$. The left side is a
        composite, so the chain rule gives
        $f'\bigl(f^{-1}(y)\bigr)\cdot\bigl(f^{-1}\bigr)'(y)$; the right side
        gives $1$. Hence
        $\bigl(f^{-1}\bigr)'(y)=1/f'\bigl(f^{-1}(y)\bigr)$, the division
        being legitimate exactly when $f'\neq0$ there.
  \item Invertibility survives: $x^{3}$ is strictly increasing and onto, so
        $f^{-1}(y)=y^{1/3}$ exists globally. What fails is
        \emph{differentiability of the inverse at $0$}:
        $\bigl(y^{1/3}\bigr)'=\tfrac13y^{-2/3}\to\infty$ as $y\to0$. The
        hypothesis $f'\neq0$ is sufficient for a differentiable local
        inverse, not necessary for an inverse.
  \item $f'(x)=2x\neq0$ for $x\neq0$, so the theorem applies at every point
        except the origin. Around $x=2$ the largest interval of
        injectivity is $(0,\infty)$: on it $f$ is strictly increasing, while
        any interval containing $0$ in its interior contains a pair
        $\pm a$ with equal images.
  \item On $(-\pi/2,\pi/2)$ the sine is strictly increasing, and
        $\pi/3$ lies inside. By the theorem, with
        $\sin(\pi/3)=\sqrt3/2$ and $\cos(\pi/3)=\tfrac12$,
        \[
          \bigl(\sin^{-1}\bigr)'\Bigl(\tfrac{\sqrt3}{2}\Bigr)
          =\frac{1}{\cos(\pi/3)}=2 .
        \]
        Directly: $\arcsin'(y)=1/\sqrt{1-y^{2}}
        =1/\sqrt{1-\tfrac34}=1/\tfrac12=2$. \checkmark
  \item $J=\begin{pmatrix}1&1\\1&-1\end{pmatrix}$, $\det J=-2\neq0$
        everywhere. Solving $u=x+y$, $v=x-y$ gives
        $x=\tfrac{u+v}{2}$, $y=\tfrac{u-v}{2}$, so
        $J_{\mathbf{f}^{-1}}
        =\begin{pmatrix}1/2&1/2\\1/2&-1/2\end{pmatrix}$, which is indeed
        $J^{-1}=-\tfrac12\begin{pmatrix}-1&-1\\-1&1\end{pmatrix}$.
        \checkmark\ The map is linear with $\det\neq0$, so here local and
        global invertibility coincide.
  \item $J=\begin{pmatrix}2x&-2y\\2y&2x\end{pmatrix}$ and
        $\det J=4(x^{2}+y^{2})$, so at $(1,2)$ it is $4\cdot5=20\neq0$: a
        differentiable local inverse exists near
        $\mathbf{f}(1,2)=(-3,4)$.
  \item $\mathbf{f}(-x,-y)=\bigl((-x)^{2}-(-y)^{2},\,2(-x)(-y)\bigr)
        =(x^{2}-y^{2},\,2xy)=\mathbf{f}(x,y)$, so every point and its
        negative share an image --- for instance
        $\mathbf{f}(1,1)=\mathbf{f}(-1,-1)=(0,2)$. The map is two-to-one,
        and no global inverse exists. Local invertibility at every point
        does not assemble into global invertibility.
  \item $J=\mathrm{e}^{x}\begin{pmatrix}\cos y&-\sin y\\
        \sin y&\cos y\end{pmatrix}$, so
        $\det J=\mathrm{e}^{2x}(\cos^{2}y+\sin^{2}y)=\mathrm{e}^{2x}>0$
        everywhere. Yet $\mathbf{f}(x,y+2\pi)=\mathbf{f}(x,y)$. There is no
        contradiction: the theorem promises an inverse on \emph{some
        neighbourhood} of each point, and those neighbourhoods (width less
        than $2\pi$ in $y$) never join up into one.
  \item $J=\begin{pmatrix}\cos\theta&-r\sin\theta\\
        \sin\theta&r\cos\theta\end{pmatrix}$, so $\det J=r$. The theorem
        says nothing on the line $r=0$ --- the origin, where every $\theta$
        gives the same point and the map is infinitely many-to-one.
  \item $J=\begin{pmatrix}1&3y^{2}\\3x^{2}&1\end{pmatrix}$, which at
        $(0,0)$ is the identity, with $\det=1\neq0$. Hence a local inverse
        exists near $\mathbf{f}(0,0)=(0,0)$ and
        $J_{\mathbf{f}^{-1}}(0,0)=I^{-1}=I$.
  \item $J=\begin{pmatrix}1&0&0\\1&1&0\\1&1&1\end{pmatrix}$, lower
        triangular, so $\det J=1\neq0$. From $u=x$, $v=x+y$, $w=x+y+z$:
        $x=u$, $y=v-u$, $z=w-v$, giving
        $J_{\mathbf{f}^{-1}}=\begin{pmatrix}1&0&0\\-1&1&0\\
        0&-1&1\end{pmatrix}$. Multiplying the two matrices returns $I$.
        \checkmark
  \item It permits: near $\mathbf{p}_{0}$ the system can in principle be
        solved for prices in terms of quantities, the inverse demand
        $\mathbf{p}(\mathbf{q})$ is differentiable, and its derivative
        matrix is $\bigl[J_{\mathbf{f}}\bigr]^{-1}$ --- so the price
        responses can be computed from the quantity responses without ever
        solving the system. It does not permit: any claim that the same
        quantities cannot arise from a different price vector elsewhere.
        Global identification is a separate question, and $\det J\neq0$ is
        no evidence for it.
\end{solutions}

\section{Additional Problems}

\begin{problems}
  \item Use IFT to find $\dd{y}{x}$ for $x^{2}+2xy+y^{2}=9$.
  \item Verify homogeneity and find the degree of $f(x,y)=x^{3}y+xy^{3}$.
  \item Compute the Jacobian and its determinant for
        $\mathbf{f}(x,y)=(2x+y,\;x-y)$.
  \item For $U(x,y)=x^{2}+y$, find the signed slope and unsigned MRS
        using IFT.
  \item Verify Euler's theorem for $f(x,y)=x^{3}+2x^{2}y-y^{3}$
        (degree $3$).
  \item Use IFT to find $\dd{y}{x}$ for $x^{3}+y^{3}+xy=5$.
  \item Verify homogeneity and find the degree of
        $f(x,y)=\dfrac{x^{4}+y^{4}}{x^{2}+y^{2}}$.
  \item Compute the Jacobian and its determinant for
        $\mathbf{f}(x,y)=(x^{2}y,\;xy^{2})$.
  \item For $U(x,y)=x^{1/3}y^{2/3}$, find the signed slope and MRS via
        IFT.
  \item Verify Euler's theorem for $Q(K,L)=K^{1/2}L^{1/2}$ (degree $1$).
  \item Use IFT to find $\dd{y}{x}$ for $\ln(xy)+x+y=4$.
  \item Decide whether $f(x,y)=x+y^{2}$ is homogeneous, with
        justification.
  \item Compute the Jacobian determinant of
        $\mathbf{f}(x,y,z)=(x+y,\;y+z,\;x+z)$.
  \item For $Q(K,L)=aK+bL$ ($a,b>0$), find the MRTS.
  \item Verify Euler's theorem for $f(x,y,z)=xy+yz+xz$ (degree $2$,
        three-variable).
  \item For $f(x)=x^{7}+x^{3}+x$, compute $\bigl(f^{-1}\bigr)'(3)$.
  \item For $f(x)=2x+\sin x$, compute $\bigl(f^{-1}\bigr)'(0)$.
  \item For $f(x)=\ln x+x^{2}$ on $x>0$, compute $\bigl(f^{-1}\bigr)'(1)$.
  \item For $f(x)=\tan x$ on $(-\pi/2,\pi/2)$, compute
        $\bigl(f^{-1}\bigr)'(1)$ and check it against $\arctan'$.
  \item For $\mathbf{f}(x,y)=(2x+3y,\;x+2y)$, give $\det J$ and
        $J_{\mathbf{f}^{-1}}$, and say why the inverse is global here.
  \item For $\mathbf{f}(x,y)=(x^{3},\;y^{3})$: where does $\det J$ vanish,
        and is the map invertible there anyway?
  \item For $\mathbf{f}(x,y)=(x+\mathrm{e}^{y},\;y+\mathrm{e}^{x})$,
        evaluate $\det J$ at the origin and say what the theorem concludes.
  \item Let $\mathbf{f}(\mathbf{x})=A\mathbf{x}$ with $A$ an $n\times n$
        matrix. Show that $\det J$ is constant, and that local and global
        invertibility coincide.
  \item For $\mathbf{f}(x,y)=(x^{2}+y^{2},\;x^{2}-y^{2})$ on $x,y>0$,
        compute $\det J$ and decide local invertibility.
  \item State in one sentence each what $f'(x_{0})\neq0$ and
        $\det J(\mathbf{x}_{0})\neq0$ do and do not guarantee.
\end{problems}

\section{Additional Solutions}

\begin{solutions}
  \item $F=x^{2}+2xy+y^{2}-9=(x+y)^{2}-9$. $F_{x}=2(x+y)$,
        $F_{y}=2(x+y)$.
        \[
          \dd{y}{x}=-\frac{F_{x}}{F_{y}}=-1 .
        \]
        (The curve is two lines $x+y=\pm 3$, each with slope $-1$.)

  \item $f(tx,ty)=(tx)^{3}(ty)+(tx)(ty)^{3}=t^{4}(x^{3}y+xy^{3})
        =t^{4}f(x,y)$. \textbf{Degree $k=4$.}

  \item
        \[
          J=\begin{pmatrix} 2 & 1\\ 1 & -1\end{pmatrix},\qquad
          \det J=-2-1=-3 .
        \]
        Globally invertible; areas scale by $|{-3}|=3$ (with orientation
        reversal).

  \item $U_{x}=2x$, $U_{y}=1$. By IFT $\dd{y}{x}=-2x$. \textbf{MRS}
        $=2x$ (unsigned).

  \item $f_{x}=3x^{2}+4xy$, $f_{y}=2x^{2}-3y^{2}$.
        \[
          \mathrm{LHS}=x(3x^{2}+4xy)+y(2x^{2}-3y^{2})
          =3x^{3}+4x^{2}y+2x^{2}y-3y^{3}=3x^{3}+6x^{2}y-3y^{3} .
        \]
        \[
          \mathrm{RHS}=3f=3(x^{3}+2x^{2}y-y^{3})
          =3x^{3}+6x^{2}y-3y^{3} .
        \]
        $\mathrm{LHS}=\mathrm{RHS}$. \checkmark

  \item $F=x^{3}+y^{3}+xy-5$. $F_{x}=3x^{2}+y$, $F_{y}=3y^{2}+x$.
        \[
          \dd{y}{x}=-\frac{3x^{2}+y}{3y^{2}+x} .
        \]

  \item $f(tx,ty)=\dfrac{t^{4}(x^{4}+y^{4})}{t^{2}(x^{2}+y^{2})}
        =t^{2}f(x,y)$. \textbf{Degree $k=2$.}

  \item
        \[
          J=\begin{pmatrix} 2xy & x^{2}\\ y^{2} & 2xy\end{pmatrix},\qquad
          \det J=4x^{2}y^{2}-x^{2}y^{2}=3x^{2}y^{2} .
        \]
        Singular when $x=0$ or $y=0$.

  \item $U_{x}=\tfrac13 x^{-2/3}y^{2/3}$,
        $U_{y}=\tfrac23 x^{1/3}y^{-1/3}$.
        $U_{x}/U_{y}=\tfrac12\cdot y/x$. By IFT $\dd{y}{x}=-y/(2x)$.
        \textbf{MRS} $=y/(2x)$.

  \item $Q_{K}=\tfrac12 K^{-1/2}L^{1/2}$,
        $Q_{L}=\tfrac12 K^{1/2}L^{-1/2}$.
        \[
          \mathrm{LHS}=K\cdot\tfrac12 K^{-1/2}L^{1/2}
          +L\cdot\tfrac12 K^{1/2}L^{-1/2}
          =\tfrac12 K^{1/2}L^{1/2}+\tfrac12 K^{1/2}L^{1/2}
          =K^{1/2}L^{1/2} .
        \]
        \[
          \mathrm{RHS}=1\cdot Q=K^{1/2}L^{1/2} .
        \]
        $\mathrm{LHS}=\mathrm{RHS}$. \checkmark

  \item $F=\ln(xy)+x+y-4=\ln x+\ln y+x+y-4$. $F_{x}=1/x+1$,
        $F_{y}=1/y+1$.
        \[
          \dd{y}{x}=-\frac{1/x+1}{1/y+1}=-\frac{y(1+x)}{x(1+y)} .
        \]

  \item $f(tx,ty)=tx+t^{2}y^{2}$. The two terms scale with different
        powers of $t$, so they cannot share a common $t^{k}$ factor.
        \textbf{Not homogeneous.}

  \item
        \[
          J=\begin{pmatrix}
            1 & 1 & 0\\
            0 & 1 & 1\\
            1 & 0 & 1
          \end{pmatrix} .
        \]
        Expanding along the first row:
        $\det J=1\cdot(1\cdot 1-1\cdot 0)-1\cdot(0\cdot 1-1\cdot 1)+0
        =1-(-1)+0=2$.

  \item Linear: $Q_{K}=a$, $Q_{L}=b$. $\dd{L}{K}=-a/b$. \textbf{MRTS}
        $=a/b$ (constant, perfect substitutes).

  \item $f_{x}=y+z$, $f_{y}=x+z$, $f_{z}=x+y$.
        \[
          \mathrm{LHS}=x(y+z)+y(x+z)+z(x+y)=2(xy+yz+xz)=2f .
        \]
        \[
          \mathrm{RHS}=2f .
        \]
        $\mathrm{LHS}=\mathrm{RHS}$. \checkmark
  \item $f(1)=1+1+1=3$ and $f'(x)=7x^{6}+3x^{2}+1$, so $f'(1)=11$ and
        $\bigl(f^{-1}\bigr)'(3)=\tfrac{1}{11}$. ($f'>0$ everywhere, so the
        pre-image is unique.)
  \item $f(0)=0$ and $f'(x)=2+\cos x\ge1>0$, so $f'(0)=3$ and
        $\bigl(f^{-1}\bigr)'(0)=\tfrac13$.
  \item $f(1)=\ln1+1=1$ and $f'(x)=1/x+2x$, so $f'(1)=3$ and
        $\bigl(f^{-1}\bigr)'(1)=\tfrac13$.
  \item $\tan(\pi/4)=1$ and $f'(x)=1/\cos^{2}x$, so
        $f'(\pi/4)=1/(\sqrt2/2)^{2}=2$ and
        $\bigl(f^{-1}\bigr)'(1)=\tfrac12$. Directly:
        $\arctan'(1)=1/(1+1^{2})=\tfrac12$. \checkmark
  \item $J=\begin{pmatrix}2&3\\1&2\end{pmatrix}$ is constant with
        $\det J=4-3=1\neq0$, and
        $J_{\mathbf{f}^{-1}}=J^{-1}
        =\begin{pmatrix}2&-3\\-1&2\end{pmatrix}$. The map is linear, so the
        local inverse supplied by the theorem is the global one:
        $\mathbf{x}=A^{-1}\mathbf{y}$ is defined on all of $\R^{2}$.
  \item $J=\begin{pmatrix}3x^{2}&0\\0&3y^{2}\end{pmatrix}$ and
        $\det J=9x^{2}y^{2}$, which vanishes on both axes. The map is
        nevertheless a bijection of $\R^{2}$, with inverse
        $(u,v)\mapsto(u^{1/3},v^{1/3})$ --- but that inverse is not
        differentiable where $\det J=0$. Same lesson as $x^{3}$ in one
        variable.
  \item $J=\begin{pmatrix}1&\mathrm{e}^{y}\\ \mathrm{e}^{x}&1\end{pmatrix}$,
        so at $(0,0)$ it is
        $\begin{pmatrix}1&1\\1&1\end{pmatrix}$ with $\det J=0$. The theorem
        concludes \textbf{nothing}: its hypothesis fails, and that is not
        the same as showing no local inverse exists. A separate argument
        would be needed either way.
  \item Each component of $A\mathbf{x}$ is linear, so
        $\partial(A\mathbf{x})_{i}/\partial x_{j}=a_{ij}$ and $J=A$ at every
        point --- constant, with $\det J=\det A$. If $\det A\neq0$ then
        $\mathbf{x}=A^{-1}\mathbf{y}$ inverts the map on all of $\R^{n}$; if
        $\det A=0$ then $A$ has a non-trivial null space and the map is not
        injective anywhere. So for linear maps the local condition is
        already global.
  \item $J=\begin{pmatrix}2x&2y\\2x&-2y\end{pmatrix}$ and
        $\det J=-4xy-4xy=-8xy$, which is non-zero throughout $x,y>0$. So a
        differentiable local inverse exists at every point of the open
        quadrant.
  \item $f'(x_{0})\neq0$ guarantees an interval around $x_{0}$ on which $f$
        is invertible with a differentiable inverse of derivative
        $1/f'(x_{0})$; it does not guarantee that $f$ is injective anywhere
        beyond that interval. $\det J(\mathbf{x}_{0})\neq0$ guarantees the
        same locally in $n$ dimensions, with the inverse's Jacobian equal to
        $\bigl[J(\mathbf{x}_{0})\bigr]^{-1}$; it likewise says nothing about
        global injectivity.
\end{solutions}


\end{document}
